% Generated mechanically from frozen input p01/ko/algebra.tex.
% Do not edit this generated file; edit the bound locale source instead.
% OK, start here.
%

\setcounter{chapter}{9}
\chapter{가환대수}



\phantomsection
\label{algebra-section-phantom}





\section{서론}
\label{algebra-section-introduction}

\noindent
이 문서에서는 가환대수의 기초를 설명한다. 참고문헌으로 \cite{MatCA}가 있다.






\section{규약}
\label{algebra-section-conventions}

\noindent
환은 $1$을 갖는 가환환이다. 영환도 환이며, 사실 소 이데알을 갖지 않는
유일한 환이다. Kronecker 기호 $\delta_{ij}$를 사용한다. $R \to S$가
환 준동형이고 $\mathfrak q$가 $S$의 소 이데알이면, $R$이 $S$의 부분환이
아니거나 $R \to S$가 단사가 아닌 경우에도 $R \to S$에 의한
$\mathfrak q$의 역상인 소 이데알을 나타내기 위해
``$\mathfrak p = R \cap \mathfrak q$''라는 표기를 사용한다.






\section{기본 개념}
\label{algebra-section-rings-basic}

\noindent
다음은 가환대수의 기본 개념을 열거한 것이다. 이 가운데 일부는 뒤의
본문에서 더 자세히 논의하고 일부는 이 목록에서 정의하지만, 나머지는
기초 개념으로 보고 정의하지 않는다. 아래의 이탤릭체 개념 대부분이
익숙하지 않다면, 계속 읽기 전에 대수학 입문서를 살펴보기를 권한다.

\begin{enumerate}
\item $R$은 {\it 환}이다,
\label{algebra-item-ring}
\item $x\in R$은 {\it 멱영원}이다,
\label{algebra-item-ring-element-nilpotent}
\item $x\in R$은 {\it 영인자}이다,
\label{algebra-item-ring-element-zerodivisor}
\item $x\in R$은 {\it 단원}이다,
\label{algebra-item-ring-element-unit}
\item $e \in R$은 {\it 멱등원}이다,
\label{algebra-item-ring-element-idempotent}
\item 멱등원 $e \in R$이 $e = 1$ 또는 $e = 0$이면 이를 {\it 자명하다}고 한다,
\label{algebra-item-idempotent-trivial}
\item $\varphi : R_1 \to R_2$는 {\it 환 준동형}이다,
\label{algebra-item-ring-homomorphism}
\item
\label{algebra-item-ring-homomorphism-finite-presentation}
$\varphi : R_1 \to R_2$는 {\it 유한 표시}이다, 또는
{\it $R_2$는 유한 표시 $R_1$-대수이다},
Definition \ref{algebra-definition-finite-type}을 보라,
\item
\label{algebra-item-ring-homomorphism-finite-type}
$\varphi : R_1 \to R_2$는 {\it 유한형}이다, 또는
{\it $R_2$는 유한형 $R_1$-대수이다},
Definition \ref{algebra-definition-finite-type}을 보라,
\item
\label{algebra-item-ring-homomorphism-finite}
$\varphi : R_1 \to R_2$는 {\it 유한}이다, 또는
{\it $R_2$는 유한 $R_1$-대수이다},
\item $R$은 {\it 정역}이다,
\label{algebra-item-ring-domain}
\item $R$은 {\it 축소환}이다,
\label{algebra-item-ring-reduced}
\item $R$은 {\it 뇌터 환}이다,
\label{algebra-item-ring-Noetherian}
\item $R$은 {\it 주 아이디얼 정역}, 즉 {\it PID}이다,
\label{algebra-item-ring-PID}
\item $R$은 {\it 유클리드 정역}이다,
\label{algebra-item-ring-Euclidean}
\item $R$은 {\it 유일분해정역}, 즉 {\it UFD}이다,
\label{algebra-item-ring-UFD}
\item $R$은 {\it 이산 값매김환}, 즉 {\it dvr}이다,
\label{algebra-item-ring-dvr}
\item $K$는 {\it 체}이다,
\label{algebra-item-field}
\item $L/K$는 {\it 체 확대}이다,
\label{algebra-item-field-extension}
\item $L/K$는 {\it 대수적 체 확대}이다,
\label{algebra-item-field-extension-algebraic}
\item $\{t_i\}_{i\in I}$는 $K$ 위의 $L$의 {\it 초월기저}이다,
\label{algebra-item-transcendence-basis}
\item $K$ 위의 $L$의 {\it 초월차수} $\text{trdeg}(L/K)$,
\label{algebra-item-transcendence-degree}
\item 체 $k$는 {\it 대수적으로 닫혀 있다},
\label{algebra-item-algebraically-closed}
\item
\label{algebra-item-extend-into-algebraically-closed}
$L/K$가 대수적이고 $\Omega/K$가 확대이며 $\Omega$가 대수적으로 닫혀 있으면,
$K$ 위의 사상을 연장하는 환 준동형 $L \to \Omega$가 존재한다,
\item $I \subset R$은 {\it 이데알}이다,
\label{algebra-item-ideal}
\item $I \subset R$은 {\it 근기 이데알}이다,
\label{algebra-item-ideal-radical}
\item $I$가 이데알이면 그 {\it 근기} $\sqrt{I}$가 있다,
\label{algebra-item-radical-ideal}
\item
\label{algebra-item-ideal-nilpotent}
$I \subset R$이 {\it 멱영 이데알}이라는 것은 어떤 $n \in \mathbf{N}$에
대하여 $I^n = 0$이라는 뜻이다,
\item
\label{algebra-item-ideal-locally-nilpotent}
$I \subset R$이 {\it 국소 멱영 이데알}이라는 것은 $I$의 모든 원소가
멱영원이라는 뜻이다,
\item $\mathfrak p \subset R$은 {\it 소 이데알}이다,
\label{algebra-item-prime-ideal}
\item
\label{algebra-item-prime-product-ideals}
$\mathfrak p \subset R$이 소 이데알이고 $I, J \subset R$이 이데알이며
$IJ\subset \mathfrak p$이면, $I \subset \mathfrak p$ 또는
$J \subset \mathfrak p$이다.
\item $\mathfrak m \subset R$은 {\it 극대 이데알}이다,
\label{algebra-item-maximal-ideal}
\item 영환이 아닌 모든 환은 극대 이데알을 갖는다,
\label{algebra-item-exists-maximal-ideal}
\item
\label{algebra-item-jacobson-radical}
$R$의 {\it 자코브슨 근기}는 $R$의 모든 극대 이데알의 교집합
$\text{rad}(R) = \bigcap_{\mathfrak m \subset R} \mathfrak m$이다,
\item 부분집합 $T \subset R$이 {\it 생성하는} 이데알 $(T)$,
\label{algebra-item-ideal-generated-by}
\item {\it 몫환} $R/I$,
\label{algebra-item-quotient-ring}
\item 환 $R$의 이데알 $I$가 소 이데알일 필요충분조건은 $R/I$가 정역인 것이다,
\label{algebra-item-characterize-prime-ideal}
\item
\label{algebra-item-characterize-maximal-ideal}
환 $R$의 이데알 $I$가 극대 이데알일 필요충분조건은 환 $R/I$가 체인 것이다,
\item
\label{algebra-item-inverse-image-ideal}
$\varphi : R_1 \to R_2$가 환 준동형이고 $I \subset R_2$가 이데알이면,
$\varphi^{-1}(I)$는 $R_1$의 이데알이다,
\item
\label{algebra-item-image-ideal}
$\varphi : R_1 \to R_2$가 환 준동형이고 $I \subset R_1$이 이데알이면,
$\varphi(I) \cdot R_2$ (때로는 $I \cdot R_2$ 또는 $IR_2$로 쓴다)는
$\varphi(I)$가 생성하는 $R_2$의 이데알이다,
\item
\label{algebra-item-inverse-image-prime}
$\varphi : R_1 \to R_2$가 환 준동형이고 $\mathfrak p \subset R_2$가
소 이데알이면, $\varphi^{-1}(\mathfrak p)$는 $R_1$의 소 이데알이다,
\item $M$은 {\it $R$-가군}이다,
\label{algebra-item-module}
\item
\label{algebra-item-annihilator}
$m \in M$에 대하여 $R$에서 $m$의 {\it 소멸자}는
$I = \{f \in R \mid fm = 0\}$이다,
\item $N \subset M$은 {\it $R$-부분가군}이다,
\label{algebra-item-submodule}
\item $M$은 {\it 뇌터 $R$-가군}이다,
\label{algebra-item-Noetherian-module}
\item $M$은 {\it 유한 $R$-가군}이다,
\label{algebra-item-finite-module}
\item $M$은 {\it 유한 생성 $R$-가군}이다,
\label{algebra-item-finitely-generated-module}
\item $M$은 {\it 유한 표시 $R$-가군}이다,
\label{algebra-item-finitely-presented-module}
\item $M$은 {\it 자유 $R$-가군}이다,
\label{algebra-item-free-module}
\item
\label{algebra-item-extension-free}
$0 \to K \to L \to M \to 0$이 $R$-가군들의 짧은 완전열이고 $K$, $M$이
자유가군이면, $L$도 자유가군이다,
\item $N \subset M \subset L$이 $R$-가군이면 $L/M = (L/N)/(M/N)$이다,
\label{algebra-item-isomorphism-theorem}
\item $S$는 {\it $R$의 곱셈적 부분집합}이다,
\label{algebra-item-multiplicative-subset}
\item $R$의 {\it 국소화} $R \to S^{-1}R$,
\label{algebra-item-localization-ring}
\item
\label{algebra-item-localization-zero}
$R$이 환이고 $S$가 $R$의 곱셈적 부분집합이면, $S^{-1}R$이 영환일
필요충분조건은 $S$가 $0$을 포함하는 것이다,
\item
\label{algebra-item-localize-nonzerodivisors}
$R$이 환이고 곱셈적 부분집합 $S$가 전부 영인자가 아닌 원소로 이루어져 있으면,
$R \to S^{-1}R$은 단사이다,
\item $\varphi : R_1 \to R_2$가 환 준동형이고 $S$가 $R_1$의 곱셈적
부분집합이면, $\varphi(S)$는 $R_2$의 곱셈적 부분집합이다,
\item
\label{algebra-item-products-multiplicative-subsets}
$S$, $S'$가 $R$의 곱셈적 부분집합이고 $SS'$가 곱들의 집합
$SS' = \{r \in R \mid \exists s\in S, \exists s' \in S', r = ss'\}$를
나타내면, $SS'$는 $R$의 곱셈적 부분집합이다,
\item
\label{algebra-item-localization-localization}
$S$, $S'$가 $R$의 곱셈적 부분집합이고 $\overline{S}$가 $(S')^{-1}R$에서
$S$의 상을 나타내면, $(SS')^{-1}R = \overline{S}^{-1}((S')^{-1}R)$이다,
\item $R$-가군 $M$의 {\it 국소화} $S^{-1}M$,
\label{algebra-item-localization-module}
\item
\label{algebra-item-localization-exact}
함자 $M \mapsto S^{-1}M$은 단사사상, 전사사상, 완전성을 보존한다,
\item
\label{algebra-item-localization-localization-module}
$S$, $S'$가 $R$의 곱셈적 부분집합이고 $M$이 $R$-가군이면,
$(SS')^{-1}M = S^{-1}((S')^{-1}M)$이다,
\item
\label{algebra-item-localize-ideal}
$R$이 환이고 $I$가 $R$의 이데알이며 $S$가 $R$의 곱셈적 부분집합이면,
$S^{-1}I$는 $S^{-1}R$의 이데알이고
$S^{-1}R/S^{-1}I = \overline{S}^{-1}(R/I)$이다. 여기서 $\overline{S}$는
$R/I$에서 $S$의 상이다,
\item
\label{algebra-item-ideal-in-localization}
$R$이 환이고 $S$가 $R$의 곱셈적 부분집합이면, $S^{-1}R$의 모든 이데알
$I'$은 $S^{-1}I$의 꼴이다. 여기서 $I$는 $R$에서 $I'$의 역상으로 택할 수 있다,
\item
\label{algebra-item-submodule-in-localization}
$R$이 환이고 $M$이 $R$-가군이며 $S$가 $R$의 곱셈적 부분집합이면,
$S^{-1}M$의 모든 부분가군 $N'$은 어떤 부분가군 $N \subset M$에 대하여
$S^{-1}N$의 꼴이다. 여기서 $N$은 $M$에서 $N'$의 역상으로 택할 수 있다,
\item $S = \{1, f, f^2, \ldots\}$이면 $R_f = S^{-1}R$이고 $M_f = S^{-1}M$이다,
\label{algebra-item-localize-f}
\item
\label{algebra-item-localize-p}
어떤 소 이데알 $\mathfrak p$에 대하여
$S = R \setminus \mathfrak p = \{x\in R \mid x\not\in \mathfrak p\}$이면,
$R_{\mathfrak p} = S^{-1}R$ 및 $M_{\mathfrak p} = S^{-1}M$로 쓰는 것이 관례이다,
\item {\it 국소환}은 극대 이데알을 정확히 하나 갖는 환이다,
\label{algebra-item-local-ring}
\item {\it 반국소환}은 극대 이데알을 유한 개 갖는 환이다,
\label{algebra-item-semi-local-ring}
\item
\label{algebra-item-localize-p-local-ring}
$\mathfrak p$가 $R$의 소 이데알이면, $R_{\mathfrak p}$는 극대 이데알
$\mathfrak p R_{\mathfrak p}$를 갖는 국소환이다,
\item
\label{algebra-item-residue-field}
환 $R$의 소 이데알 $\mathfrak p$의 {\it 잉여체} $\kappa(\mathfrak p)$는
정역 $R/\mathfrak p$의 분수체이다. 이는
$R_\mathfrak p/\mathfrak pR_\mathfrak p
= (R \setminus \mathfrak p)^{-1}R/\mathfrak p$와 같다,
\item $R$과 $M_1$, $M_2$가 주어졌을 때의 {\it 텐서곱} $M_1 \otimes_R M_2$,
\label{algebra-item-tensor-product}
\item
\label{algebra-item-cauchy-binet}
환 $R$ 위의 크기가 각각 $m \times n$ 및 $n \times m$인 행렬 $A$, $B$에
대하여 $R$에서 $\det(AB) = \sum \det(A_S)\det({}_SB)$가 성립한다. 여기서
합은 크기가 $m$인 부분집합 $S \subset \{1, \ldots, n\}$ 전체를 돌고,
$A_S$는 $S$에 대응하는 $A$의 열들로 이루어진 $m \times m$ 부분행렬이며,
${}_SB$는 $S$에 대응하는 $B$의 행들로 이루어진 $m \times m$ 부분행렬이다,
\item 기타.
\end{enumerate}




\section{뱀 보조정리}
\label{algebra-section-snake}

\noindent
뱀 보조정리와 그 변형들은 아벨 범주의 맥락에서
Homology, Section \ref{homology-section-abelian-categories}에서 논의한다.

\begin{lemma}
\label{algebra-lemma-snake}
\begin{reference}
\cite[III, Lemma 3.3]{Cartan-Eilenberg}
\end{reference}
행들이 완전한 아벨군의 가환도식
$$
\xymatrix{
& X \ar[r] \ar[d]^\alpha &
Y \ar[r] \ar[d]^\beta &
Z \ar[r] \ar[d]^\gamma &
0 \\
0 \ar[r] & U \ar[r] & V \ar[r] & W
}
$$
이 주어지면 표준적인 완전열
$$
\Ker(\alpha) \to \Ker(\beta) \to \Ker(\gamma)
\to
\Coker(\alpha) \to \Coker(\beta) \to \Coker(\gamma)
$$
이 존재한다. 또한 $X \to Y$가 단사이면 첫 번째 사상도 단사이고,
$V \to W$가 전사이면 마지막 사상도 전사이다.
\end{lemma}

\begin{proof}
사상 $\partial : \Ker(\gamma) \to \Coker(\alpha)$를 다음과 같이 정의한다.
$z \in \Ker(\gamma)$를 택하고 $z$로 가는 $y \in Y$를 택한다. 그러면
$\beta(y) \in V$는 $W$의 영원소로 가므로, $\beta(y)$는 어떤 $u \in U$의
상이다. $\partial z = \overline{u}$, 즉 $\alpha$의 여핵에서 $u$의 동치류로
정한다. 완전성의 증명은 생략한다.
\end{proof}

\section{유한 가군과 유한 표시 가군}
\label{algebra-section-module-finite-type}

\noindent
몇 가지 기본적인 표기와 보조정리를 제시한다.

\begin{definition}
\label{algebra-definition-module-finite-type}
$R$을 환, $M$을 $R$-가군이라 하자.
\begin{enumerate}
\item 어떤 $n \in \mathbf{N}$과 $x_1, \ldots, x_n \in M$이 존재하여
$M$의 모든 원소가 $x_i$들의 $R$-선형결합으로 표현되면, $M$을
{\it 유한 $R$-가군} 또는 {\it 유한 생성 $R$-가군}이라 한다.
동치로, 이는 어떤 $n \in \mathbf{N}$에 대하여 전사
$R^{\oplus n} \to M$이 존재한다는 뜻이다.
\item 어떤 정수 $n, m \in \mathbf{N}$과 완전열
$$
R^{\oplus m} \longrightarrow R^{\oplus n} \longrightarrow M \longrightarrow 0
$$
이 존재하면, $M$을 {\it 유한 표시 $R$-가군} 또는
{\it 유한 표시인 $R$-가군}이라 한다.
\end{enumerate}
\end{definition}

\noindent
비형식적으로 말해, $M$이 유한 표시 $R$-가군일 필요충분조건은
유한 생성이고 이 생성원들 사이의 관계식들로 이루어진 가군도
유한 생성인 것이다. 정의에 나오는 것과 같은 완전열의 선택을
$M$의 {\it 표시}라고 한다.

\begin{lemma}
\label{algebra-lemma-lift-map}
$R$을 환이라 하자. $\alpha : R^{\oplus n} \to M$과
$\beta : N \to M$을 가군 준동형이라 하자.
$\Im(\alpha) \subset \Im(\beta)$이면
$\alpha = \beta \circ \gamma$를 만족하는 $R$-가군 준동형
$\gamma : R^{\oplus n} \to N$이 존재한다.
\end{lemma}

\begin{proof}
$e_i = (0, \ldots, 0, 1, 0, \ldots, 0)$를 $R^{\oplus n}$의
제$i$ 기저벡터라 하자. 가정에 따라 $\alpha(e_i) = \beta(x_i)$를
만족하는 $x_i \in N$을 택한다.
$\gamma(a_1, \ldots, a_n) = \sum a_i x_i$로 놓는다.
구성상 $\alpha = \beta \circ \gamma$이다.
\end{proof}

\begin{lemma}
\label{algebra-lemma-extension}
$R$을 환이라 하자.
$$
0 \to M_1 \to M_2 \to M_3 \to 0
$$
을 $R$-가군들의 짧은 완전열이라 하자.
\begin{enumerate}
\item $M_1$과 $M_3$가 유한 $R$-가군이면, $M_2$도 유한
$R$-가군이다.
\item $M_1$과 $M_3$가 유한 표시 $R$-가군이면, $M_2$도 유한 표시
$R$-가군이다.
\item $M_2$가 유한 $R$-가군이면, $M_3$도 유한 $R$-가군이다.
\item $M_2$가 유한 표시 $R$-가군이고 $M_1$이 유한
$R$-가군이면, $M_3$는 유한 표시 $R$-가군이다.
\item $M_3$가 유한 표시 $R$-가군이고 $M_2$가 유한
$R$-가군이면, $M_1$은 유한 $R$-가군이다.
\end{enumerate}
\end{lemma}

\begin{proof}
(1)의 증명. $x_1, \ldots, x_n$이 $M_1$의 생성원이고
$y_1, \ldots, y_m \in M_2$의 $M_3$에서의 상들이 $M_3$를 생성하면,
$x_1, \ldots, x_n, y_1, \ldots, y_m$은 $M_2$를 생성한다.

\medskip\noindent
(3)은 정의에서 곧바로 따른다.

\medskip\noindent
(5)의 증명. $M_3$가 유한 표시이고 $M_2$가 유한이라고 가정하자.
표시
$$
R^{\oplus m} \to R^{\oplus n} \to M_3 \to 0
$$
를 택한다. Lemma \ref{algebra-lemma-lift-map}에 따라 다음 실선 도식이
가환하도록 하는 준동형 $R^{\oplus n} \to M_2$가 존재한다.
$$
\xymatrix{
& R^{\oplus m} \ar[r] \ar@{..>}[d] & R^{\oplus n} \ar[r] \ar[d] &
M_3 \ar[r] \ar[d]^{\text{id}} & 0 \\
0 \ar[r] & M_1 \ar[r] & M_2 \ar[r] & M_3 \ar[r] & 0
}
$$
이로써 점선 화살표가 생긴다. 뱀 보조정리
(Lemma \ref{algebra-lemma-snake})에 따라 동형
$$
\Coker(R^{\oplus m} \to M_1)
\cong
\Coker(R^{\oplus n} \to M_2)
$$
을 얻는다. 특히 $\Coker(R^{\oplus m} \to M_1)$은 유한
$R$-가군이다. $\Im(R^{\oplus m} \to M_1)$은 (3)에 따라 유한이므로,
(1)에 따라 $M_1$은 유한이다.

\medskip\noindent
(4)의 증명. $M_2$가 유한 표시이고 $M_1$이 유한이라고 가정하자.
표시 $R^{\oplus m} \to R^{\oplus n} \to M_2 \to 0$과 전사
$R^{\oplus k} \to M_1$을 택한다. Lemma \ref{algebra-lemma-lift-map}에 따라
합성 $R^{\oplus k} \to M_1 \to M_2$의 분해
$R^{\oplus k} \to R^{\oplus n} \to M_2$가 존재한다. 그러면
$R^{\oplus k + m} \to R^{\oplus n} \to M_3 \to 0$은 하나의 표시이다.

\medskip\noindent
(2)의 증명. $M_1$과 $M_3$가 유한 표시라고 가정하자.
(1)의 증명에 쓰인 논증으로 전사인 수직 화살표를 갖는 가환도식
$$
\xymatrix{
0 \ar[r] & R^{\oplus n} \ar[d] \ar[r] & R^{\oplus n + m} \ar[d] \ar[r] &
R^{\oplus m} \ar[d] \ar[r] & 0 \\
0 \ar[r] & M_1 \ar[r] & M_2 \ar[r] & M_3 \ar[r] & 0
}
$$
을 얻는다. 뱀 보조정리에 따라 짧은 완전열
$$
0 \to \Ker(R^{\oplus n} \to M_1) \to
\Ker(R^{\oplus n + m} \to M_2) \to
\Ker(R^{\oplus m} \to M_3) \to 0
$$
을 얻는다. (5)에 따라 양끝의 두 가군은 유한이고, 따라서 가운데
가군도 유한이다. (4)에 따라 $M_2$는 유한 표시이다.
\end{proof}

\begin{lemma}
\label{algebra-lemma-trivial-filter-finite-module}
\begin{slogan}
유한 가군에는 각 연이은 몫이 순환가군인 여과가 존재한다.
\end{slogan}
$R$을 환, $M$을 유한 $R$-가군이라 하자. 유한 $R$-부분가군들에 의한 여과
$$
0 = M_0 \subset M_1 \subset \ldots \subset M_n = M
$$
가 존재하여, 각 몫 $M_i/M_{i - 1}$은 $R$의 어떤 이데알 $I_i$에
대하여 $R/I_i$와 동형이다.
\end{lemma}

\begin{proof}
$M$의 생성원 수에 관하여 귀납한다. $x_1, \ldots, x_r \in M$을
생성원들이라 하자. $M' = Rx_1 \subset M$로 놓으면 $M/M'$에는
$r - 1$개의 생성원이 있으므로 귀납가설을 적용할 수 있다. 또한 분명히
$I_1 = \{f \in R \mid fx_1 = 0\}$에 대하여 $M' \cong R/I_1$이다.
\end{proof}

\begin{lemma}
\label{algebra-lemma-finite-over-subring}
$R \to S$를 환 준동형, $M$을 $S$-가군이라 하자.
$M$이 $R$-가군으로서 유한이면, $M$은 $S$-가군으로서도 유한이다.
\end{lemma}

\begin{proof}
실제로 $M$의 모든 $R$-생성집합은 $M$의 $S$-생성집합이기도 하다.
이는 $R$-가군 구조가 $S$ 안에서의 $R$의 상에 의해 유도되기 때문이다.
\end{proof}



\section{유한형 및 유한 표시 환 준동형}
\label{algebra-section-finite-type}

\begin{definition}
\label{algebra-definition-finite-type}
$R \to S$를 환 준동형이라 하자.
\begin{enumerate}
\item $n \in \mathbf{N}$과 $R$-대수의 전사
$R[x_1, \ldots, x_n] \to S$가 존재하면, $R \to S$가
{\it 유한형}이라고 하거나 {\it $S$가 유한형 $R$-대수}라고 한다.
% P01-E0263 / P01-E1272: duplicate records for the frozen English article error.
\item 어떤 정수 $n, m \in \mathbf{N}$과 다항식
$f_1, \ldots, f_m \in R[x_1, \ldots, x_n]$ 및 $R$-대수의 동형
$$
R[x_1, \ldots, x_n]/(f_1, \ldots, f_m) \cong S
$$
가 존재하면,
$R \to S$가 {\it 유한 표시}라고 한다.
\end{enumerate}
\end{definition}

\noindent
비형식적으로 말해, $R \to S$가 유한 표시일 필요충분조건은 $S$가
$R$-대수로서 유한 생성이고 생성원들 사이의 관계식들이 이루는
이데알이 유한 생성인 것이다. 정의와 같은 전사
$R[x_1, \ldots, x_n] \to S$의 선택을 때때로 $S$의 {\it 표시}라고 한다.

\begin{lemma}
\label{algebra-lemma-compose-finite-type}
유한형과 유한 표시라는 개념은 다음 보존 성질들을 갖는다.
\begin{enumerate}
\item 유한형 환 준동형들의 합성은 유한형이다.
\item 유한 표시 환 준동형들의 합성은 유한 표시이다.
\item $R \to S' \to S$가 주어지고 $R \to S$가 유한형이면,
$S' \to S$는 유한형이다.
\item $R \to S' \to S$가 주어지고 $R \to S$가 유한 표시이며
$R \to S'$가 유한형이면, $S' \to S$는 유한 표시이다.
\end{enumerate}
\end{lemma}

\begin{proof}
마지막 명제만 증명한다.
$S = R[x_1, \ldots, x_n]/(f_1, \ldots, f_m)$ 및
$S' = R[y_1, \ldots, y_a]/I$로 쓰자. $y_i$의 동치류
$\bar y_i$가 $S$의 $h_i \bmod (f_1, \ldots, f_m)$으로 간다고 하자.
그러면 분명히
$S = S'[x_1, \ldots, x_n]/(f_1, \ldots, f_m,
h_1 - \bar y_1, \ldots, h_a - \bar y_a)$이다.
\end{proof}

\begin{lemma}
\label{algebra-lemma-finite-presentation-independent}
$R \to S$를 유한 표시 환 준동형이라 하자. 임의의 전사
$\alpha : R[x_1, \ldots, x_n] \to S$에 대하여 $\alpha$의 핵은
$R[x_1, \ldots, x_n]$의 유한 생성 이데알이다.
\end{lemma}

\begin{proof}
$S = R[y_1, \ldots, y_m]/(f_1, \ldots, f_k)$로 쓰자.
$\alpha(x_i)$의 올림인 $g_i \in R[y_1, \ldots, y_m]$를 택하면,
$S = R[x_i, y_j]/(f_l, x_i - g_i)$임을 알 수 있다.
$\alpha(h_j)$가 $y_j \bmod (f_1, \ldots, f_k)$에 대응하도록
$h_j \in R[x_1, \ldots, x_n]$를 택한다. 준동형
$\psi : R[x_i, y_j] \to R[x_i]$, $x_i \mapsto x_i$,
$y_j \mapsto h_j$를 생각하자. 그러면 $\alpha$의 핵은
$(f_l, x_i - g_i)$의 $\psi$에 의한 상이므로 결론을 얻는다.
\end{proof}

\begin{lemma}
\label{algebra-lemma-finitely-presented-over-subring}
$R \to S$를 환 준동형, $M$을 $S$-가군이라 하자.
$R \to S$가 유한형이고 $M$이 $R$-가군으로서 유한 표시라고 가정하자.
그러면 $M$은 $S$-가군으로서 유한 표시이다.
\end{lemma}

\begin{proof}
이는 Lemma \ref{algebra-lemma-compose-finite-type}의 (4)의 증명과 비슷하다.
$S = R[x_1, \ldots, x_n]/J$라고 가정해도 좋다.
$M$을 $R$-가군으로서 생성하는 $y_1, \ldots, y_m \in M$을 택하고,
$R^{\oplus m} \to M$의 핵을 생성하는 관계식
$\sum a_{ij} y_j = 0$, $i = 1, \ldots, t$를 택한다.
각 $i = 1, \ldots, n$과 $j = 1, \ldots, m$에 대하여
어떤 $a_{ijk} \in R$을 사용하여
$$
x_i y_j = \sum a_{ijk} y_k
$$
로 쓰자. $y_1, \ldots, y_m$에 의해 생성되고 관계식
$\sum a_{ij} y_j = 0$, $i = 1, \ldots, t$ 및
$x_i y_j = \sum a_{ijk} y_k$, $i = 1, \ldots, n$,
$j = 1, \ldots, m$을 갖는 $S$-가군 $N$을 생각하자.
그러면 $N$은 표시
$$
S^{\oplus nm + t} \longrightarrow S^{\oplus m} \longrightarrow N
\longrightarrow 0
$$
를 갖는다. 구성상 전사 $\varphi : N \to M$이 있다.
증명을 끝내기 위해 $\varphi$가 단사임을 보이자.
어떤 $b_j \in S$에 대하여 $z = \sum b_j y_j \in N$이라고 하자.
$b_j$를 계수가 $R$에 있는 $x_1, \ldots, x_n$의 다항식으로
생각할 수 있다. $x_i y_j = \sum a_{ijk} y_k$ 꼴의 관계식을 적용하면
다항식들의 차수를 귀납적으로 낮출 수 있다. 따라서 어떤
$c_j \in R$에 대하여 $z = \sum c_j y_j$임을 알 수 있다.
그러므로 $\varphi(z) = 0$이면 벡터 $(c_1, \ldots, c_m)$은 벡터
$(a_{i1}, \ldots, a_{im})$들의 $R$-선형결합이고,
원하는 대로 $z = 0$이라는 결론을 얻는다.
\end{proof}





\section{유한 환 준동형}
\label{algebra-section-finite}

\noindent
정의는 다음과 같다.

\begin{definition}
\label{algebra-definition-finite-ring-map}
$\varphi : R \to S$를 환 준동형이라 하자. $S$가 $R$-가군으로서 유한이면
$\varphi : R \to S$가 {\it 유한}이라고 한다.
\end{definition}

\begin{lemma}
\label{algebra-lemma-finite-module-over-finite-extension}
$R \to S$를 유한 환 준동형, $M$을 $S$-가군이라 하자.
$M$이 $R$-가군으로서 유한일 필요충분조건은 $M$이 $S$-가군으로서 유한인 것이다.
\end{lemma}

\begin{proof}
한쪽 함의는 Lemma \ref{algebra-lemma-finite-over-subring}에서 따른다.
다른 쪽을 보기 위해 $M$이 $S$-가군으로서 유한이라고 가정하자.
$S$를 $R$-가군으로서 생성하는 $x_1, \ldots, x_n \in S$와
$M$을 $S$-가군으로서 생성하는 $y_1, \ldots, y_m \in M$을 택한다.
그러면 $x_i y_j$들이 $M$을 $R$-가군으로서 생성한다.
\end{proof}

\begin{lemma}
\label{algebra-lemma-finite-transitive}
$R \to S$와 $S \to T$가 유한 환 준동형이라고 하자.
그러면 $R \to T$는 유한이다.
\end{lemma}

\begin{proof}
$t_i$들이 $T$를 $S$-가군으로서 생성하고 $s_j$들이 $S$를
$R$-가군으로서 생성하면, $t_i s_j$들이 $T$를 $R$-가군으로서 생성한다.
(Lemma \ref{algebra-lemma-finite-module-over-finite-extension}에서도 따른다.)
\end{proof}

\begin{lemma}
\label{algebra-lemma-finite-finite-type}
$\varphi : R \to S$를 환 준동형이라 하자.
\begin{enumerate}
\item $\varphi$가 유한이면 $\varphi$는 유한형이다.
\item $S$가 $R$-가군으로서 유한 표시이면 $\varphi$는 유한 표시이다.
\end{enumerate}
\end{lemma}

\begin{proof}
(1)에 대하여, $x_1, \ldots, x_n \in S$가 $S$를 $R$-가군으로서
생성하면 $x_1, \ldots, x_n$은 $S$를 $R$-대수로서도 생성한다. (2)에 대하여,
$x_i$들을 $S$의 $R$-가군 생성원으로 볼 때 그들 사이의 관계식들의
생성집합을 $\sum r_j^ix_i = 0$, $j = 1, \ldots, m$이라 하자.
또한 어떤 $r_i \in R$과 $r_{ij}^k \in R$에 대하여
$1 = \sum r_ix_i$ 및 $x_ix_j = \sum r_{ij}^k x_k$로 쓰자. 그러면
$$
S = R[t_1, \ldots, t_n]/
(\sum r_j^it_i,\ 1 - \sum r_it_i,\ t_it_j - \sum r_{ij}^k t_k)
$$
가 $R$-대수로서 성립하므로 (2)가 증명된다.
\end{proof}

\noindent
유한 환 준동형에 관한 자세한 내용은
Section \ref{algebra-section-finite-ring-extensions}을 보라.

\section{여극한}
\label{algebra-section-colimits}

\noindent
본 절의 일부 내용은 범주론의 Sections \ref{categories-section-limits} --
\ref{categories-section-posets-limits}에서 다룬 여극한에 관한 일반 논의와
겹친다. 준순서 집합의 개념은 범주론의
Definition \ref{categories-definition-directed-set}에서 정의한다.
이는 부분순서 집합보다 약간 약한 개념이다.

\begin{definition}
\label{algebra-definition-directed-system}
$(I, \leq)$를 준순서 집합이라 하자.
{\it $I$ 위의 $R$-가군 계 $(M_i, \mu_{ij})$}는
$I$로 지표화된 $R$-가군의 족 $\{M_i\}_{i\in I}$와
$R$-가군 사상의 족
$\{\mu_{ij} : M_i \to M_j\}_{i \leq j}$로 이루어지며,
모든 $i \leq j \leq k$에 대하여
$$
\mu_{ii} = \text{id}_{M_i}\quad
\mu_{ik} = \mu_{jk}\circ \mu_{ij}
$$
를 만족한다. $I$가 유향 집합이면 $(M_i, \mu_{ij})$를
{\it 유향계}라고 한다.
\end{definition}

\noindent
이는 범주론의 Definition \ref{categories-definition-system-over-poset}과
Section \ref{categories-section-posets-limits}에서 정의한 개념과 같다.
임의의 범주에서 도표 또는 계의 여극한을 정의하는 방법은 범주론의
Definition \ref{categories-definition-colimit}을 보라.

\begin{lemma}
\label{algebra-lemma-colimit}
$(M_i, \mu_{ij})$를 준순서 집합 $I$ 위의 $R$-가군 계라 하자.
계 $(M_i, \mu_{ij})$의 여극한은 몫 $R$-가군
$(\bigoplus_{i\in I} M_i) /Q$이다. 여기서 $Q$는 다음 꼴의 모든 원소
$$
\iota_i(x_i) - \iota_j(\mu_{ij}(x_i))
$$
로 생성되는 $R$-부분가군이고,
$\iota_i : M_i \to \bigoplus_{i\in I} M_i$는 자연 포함사상이다.
여극한을 $M = \colim_i M_i$로 나타내고,
사영사상을 $\pi : \bigoplus_{i\in I} M_i \to M$로 나타내며,
$\phi_i = \pi \circ \iota_i : M_i \to M$로 둔다.
\end{lemma}

\begin{proof}
이 보조정리는 범주론의
Lemma \ref{categories-lemma-colimits-coproducts-coequalizers}의
특수한 경우이지만, 여기서는 이 경우를 직접 증명하기도 한다.
위의 구성에서 $\phi_i = \phi_j\circ \mu_{ij}$임을 주목하라.
쌍 $(M, \phi_i)$가 여극한임을 보이려면 다음 보편 성질을 만족함을
보여야 한다. 즉 $\psi_i : M_i \to Y$이고
$\psi_i = \psi_j\circ \mu_{ij}$인 임의의 다른 쌍
$(Y, \psi_i)$에 대하여, 다음 도표를 가환하게 하는 유일한
$R$-가군 준동형 $g : M \to Y$가 존재해야 한다.
$$
\xymatrix{
M_i \ar[rr]^{\mu_{ij}} \ar[dr]^{\phi_i} \ar[ddr]_{\psi_i} & &
M_j\ar[dl]_{\phi_j} \ar[ddl]^{\psi_j} \\
& M \ar[d]^{g}\\
& Y
}
$$
이는 직합 $\bigoplus M_i$의 직합항 $M_i$ 위에서 사상 $\psi_i$를
취하여 $g$를 정의할 수 있으므로 명백하다.
\end{proof}

\begin{lemma}
\label{algebra-lemma-directed-colimit}
$(M_i, \mu_{ij})$를 준순서 집합 $I$ 위의 $R$-가군 계라 하고
$I$가 유향이라고 가정하자. 계 $(M_i, \mu_{ij})$의 여극한은
다음과 같이 정의한 가군 $M$과 표준적으로 동형이다.
\begin{enumerate}
\item 집합으로서는
$$
M = \left(\coprod\nolimits_{i \in I} M_i\right)/\sim
$$
로 둔다. 여기서 $m \in M_i$와 $m' \in M_{i'}$에 대하여
$$
m \sim m' \Leftrightarrow
\mu_{ij}(m) = \mu_{i'j}(m')\text{ for some }j \geq i, i'
$$
로 정의한다.
\item 아벨군으로서는 $m \in M_i$와 $m' \in M_{i'}$에 대하여,
$M$에서 $m$과 $m'$의 동치류의 합을
$\mu_{ij}(m) + \mu_{i'j}(m')$의 동치류로 정의한다. 여기서
$j \in I$는 $i \leq j$와 $i' \leq j$를 만족하는 임의의 지표이다.
\item $R$-가군으로서는 $m \in M_i$와 $x \in R$에 대하여,
$M$에서 $x$와 $m$의 동치류의 곱을 $M$에서 $xm$의 동치류로 정의한다.
\end{enumerate}
표준 사상 $\phi_i : M_i \to M$은 표준 사상
$M_i \to \coprod_{i \in I} M_i$로부터 유도된다.
\end{lemma}

\begin{proof}
생략한다. 범주론의
Section \ref{categories-section-directed-colimits}과 비교하라.
\end{proof}

\begin{lemma}
\label{algebra-lemma-zero-directed-limit}
$(M_i, \mu_{ij})$를 유향계라 하자.
$M = \colim M_i$이고 $\mu_i : M_i \to M$라 하자.
그러면 $x_i \in M_i$에 대하여 $\mu_i(x_i) = 0$일 필요충분조건은
$\mu_{ij}(x_i) = 0$인 $j \geq i$가 존재하는 것이다.
\end{lemma}

\begin{proof}
이는 Lemma \ref{algebra-lemma-directed-colimit}에서 유향 여극한을 기술한
방식으로부터 명백하다.
\end{proof}

\begin{example}
\label{algebra-example-zero-colimit-different}
$a < b$와 $a < c$ 외에는 엄밀한 부등식이 없는 부분순서 집합
$I = \{a, b, c\}$를 생각하자.
$I$ 위의 계 $(M_a, M_b, M_c, \mu_{ab}, \mu_{ac})$는 세
$R$-가군 $M_a, M_b, M_c$와 두 $R$-가군 준동형
$\mu_{ab} : M_a \to M_b$ 및 $\mu_{ac} : M_a \to M_c$로 이루어진다.
이 계의 여극한은 단순히
$$
M := \colim_{i \in I} M_i = \Coker(M_a \to M_b \oplus M_c)
$$
이며, 여기서 사상은 $\mu_{ab} \oplus -\mu_{ac}$이다.
따라서 표준 사상 $M_a \to M$의 핵은
$\Ker(\mu_{ab}) + \Ker(\mu_{ac})$이다. 또한 표준 사상
$M_b \to M$의 핵은 $\Ker(\mu_{ac})$의 사상 $\mu_{ab}$에 의한 상이다.
그러므로 Lemma \ref{algebra-lemma-zero-directed-limit}의 결론은 일반적인
계에 대해서는 성립하지 않음이 명백하다.
\end{example}

\begin{definition}
\label{algebra-definition-homomorphism-directed-systems}
$(M_i, \mu_{ij})$와 $(N_i, \nu_{ij})$를 같은 준순서 집합 $I$ 위의
$R$-가군 계라 하자. $(M_i, \mu_{ij})$에서 $(N_i, \nu_{ij})$로 가는
{\it 계의 준동형} $\Phi$는 정의상 $R$-가군 준동형들의 족
$\phi_i : M_i \to N_i$로서, 모든 $i \leq j$에 대하여
$\phi_j \circ \mu_{ij} = \nu_{ij} \circ \phi_i$를 만족하는 것이다.
\end{definition}

\noindent
범주의 언어로 말하면 이는 대응하는 도표
$M : I \to \text{Mod}_R$와 $N : I \to \text{Mod}_R$ 사이의
함자 변환과 같은 개념이다. 다음 보조정리는 범주론의
Lemma \ref{categories-lemma-functorial-colimit}의 특수한 경우이다.

\begin{lemma}
\label{algebra-lemma-homomorphism-limit}
$(M_i, \mu_{ij})$와 $(N_i, \nu_{ij})$를 같은 준순서 집합 위의
$R$-가군 계라 하자. $(M_i, \mu_{ij})$에서 $(N_i, \nu_{ij})$로 가는
계의 사상 $\Phi = (\phi_i)$는 다음 유일한 준동형을 유도한다.
$$
\colim \phi_i : \colim M_i \longrightarrow \colim N_i
$$
이 준동형은 모든 $i \in I$에 대하여 다음 도표를 가환하게 한다.
$$
\xymatrix{
M_i \ar[r] \ar[d]_{\phi_i} & \colim M_i \ar[d]^{\colim \phi_i} \\
N_i \ar[r] & \colim N_i
}
$$
\end{lemma}

\begin{proof}
$M = \colim M_i$, $N = \colim N_i$, 그리고
$\phi = \colim \phi_i$라고 쓰자(마지막 것은 아직 구성해야 한다).
이하 별도의 언급 없이 Lemma \ref{algebra-lemma-colimit}에서 준 $M$과 $N$의
명시적 기술을 사용한다. 보조정리의 조건은 다음 도표가 가환한다는
조건과 동치이다.
$$
\xymatrix{
\bigoplus_{i\in I} M_i \ar[r] \ar[d]_{\bigoplus\phi_i} & M \ar[d]^\phi \\
\bigoplus_{i\in I} N_i \ar[r] & N
}
$$
따라서 $\phi$가 존재한다면 유일함이 명백하다.
$\phi$가 존재함을 보이려면 위쪽 수평 화살표의 핵이
$\bigoplus \phi_i$에 의해 아래쪽 수평 화살표의 핵으로 사상됨을
보이면 충분하다. 이를 위해 $j \leq k$와 $x_j \in M_j$를 택하자. 그러면
$$
(\bigoplus \phi_i)(x_j - \mu_{jk}(x_j))
=
\phi_j(x_j) - \phi_k(\mu_{jk}(x_j))
=
\phi_j(x_j) - \nu_{jk}(\phi_j(x_j))
$$
이고, 이는 요구한 대로 아래쪽 수평 화살표의 핵에 속한다.
\end{proof}

\begin{lemma}
\label{algebra-lemma-directed-colimit-exact}
\begin{slogan}
여과 여극한은 완전하다. 유향 여극한은 완전하다.
\end{slogan}
$I$를 유향 집합이라 하자.
$(L_i, \lambda_{ij})$, $(M_i, \mu_{ij})$, 그리고
$(N_i, \nu_{ij})$를 $I$ 위의 $R$-가군 계라 하자.
$\varphi_i : L_i \to M_i$와 $\psi_i : M_i \to N_i$를
$I$ 위의 계의 사상이라 하자. 모든 $i \in I$에 대하여
다음 $R$-가군 열이
$$
\xymatrix{
L_i \ar[r]^{\varphi_i} &
M_i \ar[r]^{\psi_i} &
N_i
}
$$
호몰로지가 $H_i$인 복합체라고 가정하자.
그러면 $R$-가군들 $H_i$는 $I$ 위의 계를 이루며,
다음 $R$-가군 열도 복합체이고,
$$
\xymatrix{
\colim_i L_i \ar[r]^\varphi &
\colim_i M_i \ar[r]^\psi &
\colim_i N_i
}
$$
그 호몰로지를 $H$로 나타내면
$$
H = \colim_i H_i.
$$
\end{lemma}

\begin{proof}
다음 열이 복합체임은 명백하다.
$
\xymatrix{
\colim_i L_i \ar[r]^\varphi &
\colim_i M_i \ar[r]^\psi &
\colim_i N_i
}
$
각 $i \in I$에 대하여 표준 $R$-가군 사상 $H_i \to H$가 있다.
이는 각 $[m] \in H_i = \Ker(\psi_i) / \Im(\varphi_i)$를
$m$의 $\colim_i M_i$에서의 상이
$H = \Ker(\psi) / \Im(\varphi)$에서 정하는 잉여류로 보낸다.
이 사상들은 사상 $\colim_i H_i \to H$를 유도한다.
이 사상이 전사이고 단사임을 보이는 일만 남는다.

\medskip\noindent
이하 별도의 언급 없이 Lemma \ref{algebra-lemma-directed-colimit}에서처럼
$I$ 위의 여극한을 기술한 방식을 반복하여 사용한다.
$h \in H$를 택하자.
$H = \Ker(\psi)/\Im(\varphi)$이므로, $h$는
$\Ker(\psi) \subset \colim_i M_i$의 원소 $[m]$이
법 $\Im(\varphi)$에 대하여 정하는 동치류이다.
$[m]$이 어떤 원소 $m \in M_i$로부터 오도록 $i$를 택한다.
$[m] \in \Ker(\psi)$이므로 가능한 방식으로,
$\nu_{ij}(\psi_i(m)) = 0$인 $j \geq i$를 택한다.
$i$를 $j$로, $m$을 $\mu_{ij}(m)$으로 바꾸면
$m \in \Ker(\psi_i)$라고 가정할 수 있다.
이는 사상 $\colim_i H_i \to H$가 전사임을 보인다.

\medskip\noindent
$h_i \in H_i$의 $H$에서의 상이 영이라고 가정하자.
$H_i = \Ker(\psi_i)/\Im(\varphi_i)$이므로 $h_i$를 원소
$m \in \Ker(\psi_i) \subset M_i$로 나타낼 수 있다.
$h_i$가 $H$에서 영이 된다는 가정은 $\colim_i M_i$에서 $m$의
동치류가 $\varphi$의 상에 속한다는 뜻이다.
따라서 $\varphi_j(l) = \mu_{ij}(m)$인 $j \geq i$와
$l \in L_j$가 존재한다. 이는 $H_j$에서 $h_i$의 상이 영임을
명백히 보인다. 이로써 $\colim_i H_i \to H$의 단사성이 증명된다.
\end{proof}

\begin{example}
\label{algebra-example-colimit-not-exact}
일반적으로 여극한을 취하는 것은 완전하지 않다.
Example \ref{algebra-example-zero-colimit-different}에서와 같이
$a < b$와 $a < c$ 외에는 엄밀한 부등식이 없는 부분순서 집합
$I = \{a, b, c\}$를 생각하자.
계의 사상
$(0, \mathbf{Z}, \mathbf{Z}, 0, 0) \to
(\mathbf{Z}, \mathbf{Z}, \mathbf{Z}, 1, 1)$을 생각하자.
Example \ref{algebra-example-zero-colimit-different}의 여극한에 대한 기술로부터,
이 계의 사상이 각 대상에서는 단사임에도 대응하는 여극한 사이의 사상은
단사가 아님을 알 수 있다. 따라서
Lemma \ref{algebra-lemma-directed-colimit-exact}의 결론은 일반적인
계에 대해서는 성립하지 않는다.
\end{example}

\begin{lemma}
\label{algebra-lemma-almost-directed-colimit-exact}
$\mathcal{I}$를 범주론의
Lemma \ref{categories-lemma-split-into-directed}의 가정을 만족하는
지표 범주라 하자. 그러면 $\mathcal{I}$ 위의 아벨군 도표의
여극한을 취하는 것은 완전하다(즉,
Lemma \ref{algebra-lemma-directed-colimit-exact}와 같은 결과가
이 상황에서도 성립한다).
\end{lemma}

\begin{proof}
범주론의 Lemma \ref{categories-lemma-split-into-directed}에 의해,
각 $\mathcal{I}_j$가 여과 범주이고 $J$는 공집합일 수도 있도록
$\mathcal{I} = \coprod_{j \in J} \mathcal{I}_j$로 쓸 수 있다.
범주론의 Lemma \ref{categories-lemma-directed-category-system}에 의해
지표 범주 $\mathcal{I}_j$ 위에서 여극한을 취하는 것은 어떤 유향 집합
위에서 여극한을 취하는 것과 같다. 따라서
Lemma \ref{algebra-lemma-directed-colimit-exact}를 이 여극한들에 적용할 수 있다.
이제 문제는 집합 $J$에 걸친 $R$-가군 범주의 쌍대곱이 완전함을
보이는 것으로 환원된다.
% P01-E0264 / P01-E1273: source wording candidates recorded in ERRATA_CANDIDATES.jsonl.
달리 말해, $R$-가군의 완전열 $L_j \to M_j \to N_j$에 대하여
다음 열이
$$
\bigoplus\nolimits_{j \in J} L_j
\longrightarrow
\bigoplus\nolimits_{j \in J} M_j
\longrightarrow
\bigoplus\nolimits_{j \in J} N_j
$$
완전함을 보여야 한다. 이는 직접 확인할 수 있으며 $J$가 공집합이어도
성립한다.
\end{proof}



\section{국소화}
\label{algebra-section-localization}

\begin{definition}
\label{algebra-definition-multiplicative-subset}
$R$을 환, $S$를 $R$의 부분집합이라 하자.
$1\in S$이고 $S$가 곱셈에 관하여 닫혀 있으면, 즉
$s, s' \in S \Rightarrow ss' \in S$이면
$S$를 {\it $R$의 곱셈적 부분집합}이라고 한다.
\end{definition}

\noindent
환 $A$와 곱셈적 부분집합 $S$가 주어졌다고 하자.
$A \times S$ 위의 관계를 다음과 같이 정의한다.
$$
(x, s) \sim (y, t)
\Leftrightarrow
\exists u \in S \text{ such that } (xt-ys)u = 0
$$
이것이 동치관계임은 쉽게 확인된다.
$x/s$ (또는 $\frac{x}{s}$)를 $(x, s)$의 동치류라 하고,
$S^{-1}A$를 모든 동치류의 집합이라 하자.
$S^{-1}A$에서 덧셈과 곱셈을 다음과 같이 정의한다.
$$
x/s + y/t = (xt + ys)/st, \quad
x/s \cdot y/t = xy/st
$$
이 연산들 아래에서 $S^{-1}A$가 환이 됨을 확인할 수 있다.

\begin{definition}
\label{algebra-definition-localization}
이 환을 {\it $S$에 관한 $A$의 국소화}라고 한다.
\end{definition}

\noindent
$A$에서 그 국소화 $S^{-1}A$로 가는 자연스러운 환 준동형
$$
A \longrightarrow S^{-1}A, \quad x \longmapsto x/1
$$
이 있으며, 이를 때때로 {\it 국소화 사상}이라고 한다.
일반적으로 국소화 사상은 단사가 아니며, $S$가 영인자를 전혀
포함하지 않으면 단사이다. 실제로 $x/1 = 0$이면 $A$에서
$xu = 0$인 $u\in S$가 존재하고, $S$에는 영인자가 없으므로
$x = 0$이다. 환의 국소화는 다음 보편 성질을 갖는다.

\begin{proposition}
\label{algebra-proposition-universal-property-localization}
$f : A \to B$를 $S$의 모든 원소를 $B$의 단원으로 보내는 환 준동형이라
하자. 그러면 다음 도표를 가환하게 하는 유일한 준동형
$g : S^{-1}A \to B$가 존재한다.
$$
\xymatrix{
A \ar[rr]^{f} \ar[dr] & & B \\
& S^{-1}A \ar[ur]_g
}
$$
\end{proposition}

\begin{proof}
존재성. 사상 $g$를 다음과 같이 정의한다.
$x/s\in S^{-1}A$에 대하여 $g(x/s) = f(x)f(s)^{-1}\in B$로 둔다.
정의로부터 이것이 잘 정의된 환 준동형임을 쉽게 확인할 수 있다.
또한 이 사상이 도표를 가환하게 함도 명백하다.

\medskip\noindent
유일성. 이제 $g' : S^{-1}A \to B$가 $g'(x/1) = f(x)$를 만족하면
$g = g'$임을 보이자. 도표의 가환성에 의해 $s \in S$에 대하여
$f(s) = g'(s/1)$이다. 그런데 $B$에서 $g'(1/s)f(s) = 1$이므로
$g'(1/s) = f(s)^{-1}$이고, 따라서
$g'(x/s) = g'(x/1)g'(1/s) = f(x)f(s)^{-1} = g(x/s)$이다.
\end{proof}

\begin{lemma}
\label{algebra-lemma-localization-zero}
국소화 $S^{-1}A$가 영환일 필요충분조건은 $0\in S$인 것이다.
\end{lemma}

\begin{proof}
$0\in S$이면 정의에 의해 임의의 쌍 $(a, s)\sim (0, 1)$이다.
$0\not \in S$이면 $S^{-1}A$에서 $1/1 \neq 0/1$임이 명백하다.
\end{proof}

\begin{lemma}
\label{algebra-lemma-localization-and-modules}
$R$을 환이라 하고 $S \subset R$을 곱셈적 부분집합이라 하자.
$S^{-1}R$-가군의 범주는 모든 $s \in S$가 $N$ 위에서 자기동형으로 작용하는
$R$-가군 $N$의 범주와 동치이다.
\end{lemma}

\begin{proof}
이 동치를 정의하는 함자는 $S^{-1}R$-가군 $M$에,
국소화 사상 $R \to S^{-1}R$을 통하여 $R$-가군으로 본 같은 가군을
대응시킨다. 반대로 $N$이 각 $s \in S$가 자기동형 $s_N$으로
작용하는 $R$-가군이면, $x/s$가 $x_N  \circ s_N^{-1}$로 작용하게 하여
$N$을 $S^{-1}R$-가군으로 볼 수 있다. 이 두 함자가 서로 준역임을
확인하는 일은 생략한다.
\end{proof}

\noindent
환의 국소화 개념은 가군의 국소화로 일반화할 수 있다.
$A$를 환, $S$를 $A$의 곱셈적 부분집합, $M$을 $A$-가군이라 하자.
$M \times S$ 위의 관계를 다음과 같이 정의한다.
$$
(m, s) \sim (n, t)
\Leftrightarrow
\exists u\in S \text{ such that } (mt-ns)u = 0
$$
이는 명백히 동치관계이다.
% P01-E0265: source naming-construction candidate recorded in ERRATA_CANDIDATES.jsonl.
$m/s$ (또는 $\frac{m}{s}$)를 $(m, s)$의 동치류라 하고,
$S^{-1}M$을 모든 동치류의 집합이라 하자.
덧셈과 스칼라곱을 다음과 같이 정의한다.
% P01-E0266 / P01-E1274: source scalar-multiplication candidate recorded in ERRATA_CANDIDATES.jsonl.
$$
m/s + n/t = (mt + ns)/st,\quad
m/s\cdot n/t = mn/st
$$
이로써 $S^{-1}M$은 $S^{-1}A$-가군이 됨이 명백하다.

\begin{definition}
\label{algebra-definition-localization-module}
$S^{-1}A$-가군 $S^{-1}M$을 $S$에 관한 $M$의 {\it 국소화}라고 한다.
\end{definition}

\noindent
$A$-가군 사상 $M \to S^{-1}M$, $m \mapsto m/1$이 있으며,
이를 때때로 {\it 국소화 사상}이라고 한다. 이 사상은 다음 보편 성질을
만족한다.

\begin{lemma}
\label{algebra-lemma-universal-property-localization-module}
% P01-E0267 / P01-E1275: source missing-verb candidate recorded in ERRATA_CANDIDATES.jsonl.
$R$을 환이라 하자. $S \subset R$을 곱셈적 부분집합이라 하자.
$M$, $N$을 $R$-가군이라 하자. $S$의 모든 원소가 $N$ 위에서
자기동형으로 작용한다고 가정하자. 그러면 국소화 사상으로 유도되는
표준 사상
$$
\Hom_R(S^{-1}M, N) \longrightarrow \Hom_R(M, N)
$$
은 동형이다.
\end{lemma}

\begin{proof}
이 사상이 잘 정의되고 $R$-선형임은 명백하다.
% P01-E0268: source math-mode candidate recorded in ERRATA_CANDIDATES.jsonl.
단사성: $\alpha \in \Hom_R(S^{-1}M, N)$이라 하고 임의의 원소
$m/s \in S^{-1}M$을 택하자. $s \cdot \alpha(m/s) = \alpha(m/1)$이므로
$ \alpha(m/s) =s^{-1}(\alpha (m/1))$이고, 따라서 $\alpha$는
$S^{-1}M$ 안의 $M$의 상 위에서 하는 작용에 의해 완전히 결정된다.
전사성: $\beta : M \rightarrow N$를 주어진 R-선형 사상이라 하자.
이를 $S^{-1}M$으로 “확장”할 수 있음을 보여야 한다.
집합의 사상
$$
M \times S \rightarrow N,\quad
(m,s) \mapsto s^{-1}\beta(m)
$$
을 정의한다. 이 사상은 위의 동치관계를 보존함이 명백하므로,
잘 정의된 사상 $\alpha : S^{-1}M \rightarrow N$으로 내려간다.
이 사상이 $R$-선형임을 보이는 일만 남는다. 이를 위해
$r, r' \in R$, $s, s' \in S$, 그리고 $m, m' \in M$을 택하자. 그러면
\begin{align*}
\alpha(r \cdot m/s + r' \cdot m' /s')
& =
\alpha((r \cdot s' \cdot m + r' \cdot s \cdot m') /(ss')) \\
& =
(ss')^{-1}\beta(r \cdot s' \cdot m + r' \cdot s \cdot m') \\
& =
(ss')^{-1} (r \cdot s' \beta (m) + r' \cdot s \beta (m')) \\
& =
r \alpha (m/s) + r' \alpha (m' /s')
\end{align*}
이고, 이로써 증명이 끝난다.
\end{proof}

\begin{example}
\label{algebra-example-localize-at-prime}
$A$를 환, $M$을 $A$-가군이라 하자.
다음은 국소화의 중요한 예들이다.
\begin{enumerate}
\item $A$의 소 이데알 $\mathfrak p$가 주어졌다고 하고
$S = A\setminus\mathfrak p$를 생각하자. $S$가 곱셈적 집합임은
즉시 확인된다. 이 경우 $A$와 $M$의 $S$에 관한 국소화를 각각
$A_\mathfrak p$와 $M_\mathfrak p$로 나타낸다.
이들을 {\it $\mathfrak p$에서의 $A$, 각각 $M$의 국소화}라고 한다.
\item $f\in A$라 하고 $S = \{1, f, f^2, \ldots\}$를 생각하자.
이는 명백히 $A$의 곱셈적 부분집합이다. 이 경우 국소화
$S^{-1}A$ (각각 $S^{-1}M$)를 $A_f$ (각각 $M_f$)로 나타낸다.
이를 {\it $f$에 관한 $A$, 각각 $M$의 국소화}라고 한다.
$A_f = 0$일 필요충분조건은 $A$에서 $f$가 멱영원인 것이다.
\item $S = \{f \in A \mid f \text{ is not a zerodivisor in }A\}$라 하자.
이는 $A$의 곱셈적 부분집합이다. 이 경우 환
$Q(A) = S^{-1}A$를 $A$의 {\it 전 몫환} 또는 {\it 전분수환}이라고 한다.
\item $A$가 정역이면 전 몫환 $Q(A)$는 $A$의 분수체이다.
체론의 Example \ref{fields-example-quotient-field}을 보라.
\end{enumerate}
\end{example}

\begin{lemma}
\label{algebra-lemma-localization-colimit}
$R$을 환이라 하자.
$S \subset R$을 곱셈적 부분집합이라 하자.
$M$을 $R$-가군이라 하자. 그러면
$$
S^{-1}M = \colim_{f \in S} M_f
$$
이다. 여기서 $S$ 위의 준순서는
$f \geq f' \Leftrightarrow f = f'f''$인 어떤 $f'' \in R$가 존재하는
것으로 주어지며, 이 경우 사상 $M_{f'} \to M_f$는
$m/(f')^e \mapsto m(f'')^e/f^e$로 주어진다.
\end{lemma}

\begin{proof}
생략한다. 힌트: Lemma \ref{algebra-lemma-universal-property-localization-module}의
보편 성질을 사용하라.
\end{proof}

\noindent
다음 문단에서 $A$는 환을 나타내고, $M, N$은 $A$ 위의 가군들을 나타낸다.

\medskip\noindent
$S$와 $S'$이 $A$의 곱셈적 집합이면
$$
SS' = \{ss' : s\in S, \ s'\in S'\}
$$
도 $A$의 곱셈적 집합임이 명백하다. 이때 다음이 성립한다.

\begin{proposition}
\label{algebra-proposition-localize-twice}
$\overline{S}$를 $S'^{-1}A$에서 $S$의 상이라 하자.
그러면 $(SS')^{-1}A$는 $\overline{S}^{-1}(S'^{-1}A)$와 동형이다.
\end{proposition}

\begin{proof}
$x\in A$를 $x/1\in (SS')^{-1}A$로 보내는 사상은 보편 성질에 의해
$x/s\in S'^{-1}A$를 $x/s \in (SS')^{-1}A$로 보내는 사상을 유도한다.
$\overline{S}$의 원소들의 상은 $(SS')^{-1}A$에서 가역이다.
보편 성질에 의해 $(x/s')/(s/1)$을 $x/ss'$로 보내는 사상
$f : \overline{S}^{-1}(S'^{-1}A) \to (SS')^{-1}A$를 얻는다.

\medskip\noindent
한편 $x\in A$를 $(x/1)/(1/1)$로 보내는
$A$에서 $\overline{S}^{-1}(S'^{-1}A)$로 가는 사상도 보편 성질에 의해
$x/ss'$를 $(x/s')/(s/1)$로 보내는 사상
$g : (SS')^{-1}A \to \overline{S}^{-1}(S'^{-1}A)$를 유도한다.
$f$와 $g$가 서로 역임은 즉시 확인되므로 둘 다 동형이다.
\end{proof}

\noindent
가군 $M$에 대해서는 다음이 성립한다.

\begin{proposition}
\label{algebra-proposition-localize-twice-module}
$S'^{-1}M$을 $A$-가군으로 보자. 그러면 $S^{-1}(S'^{-1}M)$은
$(SS')^{-1}M$과 동형이다.
\end{proposition}

\begin{proof}
% P01-E0269: source article/math-mode candidate recorded in ERRATA_CANDIDATES.jsonl.
주어진 $A$-가군 M에 대하여 $S^{-1}M$의 보편 성질을 아직 증명하지
않았음에 유의하라. 따라서 앞의 증명처럼 논증할 수 없고,
동형을 명시적으로 구성해야 한다.

\medskip\noindent
사상들을 다음과 같이 정의한다.
\begin{align*}
& f : S^{-1}(S'^{-1}M) \longrightarrow (SS')^{-1}M, \quad \frac{x/s'}{s}\mapsto
x/ss'\\
& g : (SS')^{-1}M \longrightarrow S^{-1}(S'^{-1}M), \quad x/t\mapsto
\frac{x/s'}{s}\ \text{for some }s\in S, s'\in S', \text{ and }
t = ss'
\end{align*}
% P01-E0270: source missing-preposition candidate recorded in ERRATA_CANDIDATES.jsonl.
이 준동형들이 잘 정의되는지, 즉 분수의 선택에 무관한지 확인해야 한다.
이는 쉽게 확인할 수 있으며, 이들이 서로 역임을 보이는 것도 간단하다.
\end{proof}

\noindent
% P01-E0271: source homomorphism-type candidate recorded in ERRATA_CANDIDATES.jsonl.
$u : M \to N$가 $A$-준동형이면 국소화는 실제로 $x/s$를 $u(x)/s$로
보내는 잘 정의된 $S^{-1}A$-준동형
$S^{-1}u : S^{-1}M \to S^{-1}N$을 유도한다.
이 구성이 함자적임은 즉시 확인되므로, $S^{-1}$은 실제로
$A$-가군의 범주에서 $S^{-1}A$-가군의 범주로 가는 함자이다.
더욱이 다음 명제에서 보이듯 이 함자는 완전하다.

\begin{proposition}
\label{algebra-proposition-localization-exact}
\begin{slogan}
국소화는 완전하다.
\end{slogan}
% P01-E0272: frozen base-ring switch recorded in ERRATA_CANDIDATES.jsonl.
$L\xrightarrow{u} M\xrightarrow{v} N$을 $R$-가군의 완전열이라 하자.
그러면 $S^{-1}L \to S^{-1}M \to S^{-1}N$도 완전하다.
\end{proposition}

\begin{proof}
먼저 국소화가 함자이므로
$S^{-1}L \to S^{-1}M \to S^{-1}N$가 복합체임은 명백하다.
다음으로 어떤 $x/s \in S^{-1}M$에 대하여 $x/s$가 $S^{-1}N$에서
영으로 사상된다고 가정하자. 그러면 정의에 의해
$N$에서 $v(xt) = v(x)t = 0$인 $t\in S$가 존재하므로
$xt \in \Ker(v)$이다. $L \to M \to N$의 완전성에 의해
어떤 $L$의 $y$에 대하여 $xt = u(y)$이다.
그러면 $x/s$는 $y/st$의 상이다. 이로써 완전성이 증명된다.
\end{proof}

\begin{lemma}
\label{algebra-lemma-localize-quotient-modules}
국소화는 몫을 보존한다. 즉 $N$이 $M$의 부분가군이면
$S^{-1}(M/N)\simeq (S^{-1}M)/(S^{-1}N)$이다.
\end{lemma}

\begin{proof}
완전열
$$
0 \longrightarrow N \longrightarrow M \longrightarrow M/N \longrightarrow 0
$$
로부터
$$
0 \longrightarrow S^{-1}N \longrightarrow S^{-1}M
\longrightarrow S^{-1}(M/N) \longrightarrow 0
$$
을 얻는다.
% P01-E0273 / P01-E1276: frozen result-type wording recorded in ERRATA_CANDIDATES.jsonl.
따라서 따름정리가 따른다.
\end{proof}

\noindent
% P01-E0274 / P01-E1276: frozen result-type wording recorded in ERRATA_CANDIDATES.jsonl.
앞의 따름정리에서 $A$의 이데알 $I$에 대하여 $N = I$와 $M = A$로
놓으면, $A$-가군으로서
$S^{-1}A/S^{-1}I \simeq S^{-1}(A/I)$임을 알 수 있다.
다음 명제는 이들이 환으로서 동형임을 보인다.

\begin{proposition}
\label{algebra-proposition-localize-quotient}
$I$를 $A$의 이데알, $S$를 $A$의 곱셈적 집합이라 하자.
그러면 $S^{-1}I$는 $S^{-1}A$의 이데알이고,
$\overline{S}$가 $A/I$에서 $S$의 상일 때
$\overline{S}^{-1}(A/I)$는 $S^{-1}A/S^{-1}I$와 동형이다.
\end{proposition}

\begin{proof}
$I$ 자체가 이데알이므로 $S^{-1}I$가 이데알이라는 사실은 명백하다.
다음을 정의하자.
$$
f : S^{-1}A\longrightarrow \overline{S}^{-1}(A/I), \quad x/s\mapsto
\overline{x}/\overline{s}
$$
여기서 $\overline{x}$와 $\overline{s}$는 $A/I$에서 각각
$x$와 $s$의 상이다. 이 증명에서 비슷한 표기를 계속 사용하겠다.
이 사상은 $S^{-1}A$의 보편 성질에 의해 잘 정의되고,
$S^{-1}I$는 그 핵에 포함된다. 따라서 사상
$$
\overline{f} : S^{-1}A/S^{-1}I \longrightarrow \overline{S}^{-1}(A/I), \quad
\overline{x/s}\mapsto \overline{x}/\overline{s}
$$
을 유도한다.

\medskip\noindent
한편 $x$를 $\overline{x/1}$로 보내는 사상
$A \to S^{-1}A/S^{-1}I$는 $\overline{x}$를 $\overline{x/1}$로 보내는
사상 $A/I \to S^{-1}A/S^{-1}I$를 유도한다.
$\overline{S}$의 상은 $S^{-1}A/S^{-1}I$에서 가역이므로,
보편 성질에 의해 사상
$$
g : \overline{S}^{-1}(A/I) \longrightarrow S^{-1}A/S^{-1}I, \quad
\frac{\overline{x}}{\overline{s}}\mapsto \overline{x/s}
$$
을 유도한다. 이제 $\overline{f}$와 $g$가 서로 역임이 명백하므로
둘 다 동형이다.
\end{proof}

\noindent
이제 부분가군이 국소화에서 어떻게 행동하는지 살펴보자.

\begin{lemma}
\label{algebra-lemma-submodule-localization}
$S^{-1}M$의 임의의 부분가군 $N'$은 어떤 $N\subset M$에 대하여
$S^{-1}N$ 꼴이다. 실제로 $N$을 $M$에서 $N'$의 역상으로 택할 수 있다.
\end{lemma}

\begin{proof}
$N$을 $M$에서 $N'$의 역상이라 하자.
% P01-E0275: frozen reversed inclusion and omitted reverse argument recorded in ERRATA_CANDIDATES.jsonl.
그러면 $S^{-1}N\supset N'$임을 알 수 있다.
이들이 같음을 보이기 위해 $S^{-1}N$의 $x/s$를 택하자.
여기서 $s\in S$이고 $x\in N$이다. 이로부터 $x/1\in N'$을 얻는다.
% P01-E0276: frozen base-ring switch recorded in ERRATA_CANDIDATES.jsonl.
$N'$은 $S^{-1}R$-부분가군이므로
$x/s = x/1\cdot 1/s\in N'$이다. 이로써 증명이 끝난다.
\end{proof}

\noindent
$M = A$와 $A$의 이데알 $N = I$를 택하면 다음 따름정리를 얻는다.
이는 Proposition \ref{algebra-proposition-localize-quotient}의 첫 번째 부분의
역으로 볼 수 있다.

\begin{lemma}
\label{algebra-lemma-ideal-in-localization}
\begin{slogan}
환의 국소화의 이데알은 이데알의 국소화이다.
\end{slogan}
$S^{-1}A$의 각 이데알 $I'$은 $S^{-1}I$ 꼴이며,
여기서 $I$를 $A$에서 $I'$의 역상으로 택할 수 있다.
\end{lemma}

\begin{proof}
Lemma \ref{algebra-lemma-submodule-localization}로부터 즉시 따른다.
\end{proof}

\section{내부 Hom}
\label{algebra-section-hom}

\noindent
$R$이 환이고 $M$, $N$이 $R$-가군이면
$$
\Hom_R(M, N) = \{ \varphi : M \to N\}
$$
은 $M$에서 $N$으로 가는 $R$-선형 사상들의 집합이다.
보통과 같이 $(\varphi + \psi)(m) = \varphi(m) + \psi(m)$로 놓으면
이 집합은 아벨군의 구조를 갖는다. 사실 $\Hom_R(M, N)$은 규칙
$(x \varphi)(m) = x \varphi(m) = \varphi(xm)$에 의해
$R$-가군이기도 하다.

\medskip\noindent
$R$-가군의 사상 $a : M \to M'$와 $b : N \to N'$가 주어지면
준동형에 $a$를 앞에서 합성하고 $b$를 뒤에서 합성할 수 있다.
이로부터 다음 가환도표를 얻는다.
$$
\xymatrix{
\Hom_R(M', N) \ar[d]_{- \circ a} \ar[r]_{b \circ -} &
\Hom_R(M', N') \ar[d]^{- \circ a} \\
\Hom_R(M, N) \ar[r]^{b \circ -} &
\Hom_R(M, N')
}
$$
사실 이 도표의 사상들은 $R$-가군 사상이다.
따라서 $\Hom_R$은 가법 함자
$$
\text{Mod}_R^{opp} \times \text{Mod}_R \longrightarrow \text{Mod}_R, \quad
(M, N) \longmapsto \Hom_R(M, N)
$$
를 정의한다.

\begin{lemma}
\label{algebra-lemma-hom-exact}
완전성과 $\Hom_R$. $R$을 환이라 하자. $M_1$, $M_2$, $M_3$를
$R$-가군이라 하고, $M_1 \to M_2$와 $M_2 \to M_3$를
$R$-가군 사상이라 하자.
\begin{enumerate}
\item $M_1 \to M_2 \to M_3 \to 0$가 완전일 필요충분조건은 모든
$R$-가군 $N$에 대하여
$0 \to \Hom_R(M_3, N) \to \Hom_R(M_2, N) \to \Hom_R(M_1, N)$이
완전인 것이다.
\item $0 \to M_1 \to M_2 \to M_3$가 완전일 필요충분조건은 모든
$R$-가군 $N$에 대하여
$0 \to \Hom_R(N, M_1) \to \Hom_R(N, M_2) \to \Hom_R(N, M_3)$가
완전인 것이다.
\end{enumerate}
\end{lemma}

\begin{proof}
 생략한다.
\end{proof}


\begin{lemma}
\label{algebra-lemma-hom-from-finitely-presented}
$R$을 환, $M$을 유한 표시 $R$-가군, $N$을 $R$-가군이라 하자.
\begin{enumerate}
\item $f \in R$에 대하여
$\Hom_R(M, N)_f = \Hom_{R_f}(M_f, N_f) = \Hom_R(M_f, N_f)$이고,
\item $R$의 곱셈적 부분집합 $S$에 대하여
$$
S^{-1}\Hom_R(M, N) = \Hom_{S^{-1}R}(S^{-1}M, S^{-1}N) =
\Hom_R(S^{-1}M, S^{-1}N).
$$
\end{enumerate}
\end{lemma}

\begin{proof}
(1)은 (2)의 특수한 경우이다.
(2)의 두 번째 등식은 Lemma \ref{algebra-lemma-localization-and-modules}로부터
따른다. 표시
$$
\bigoplus\nolimits_{j = 1, \ldots, m} R
\longrightarrow
\bigoplus\nolimits_{i = 1, \ldots, n} R
\to M \to 0.
$$
를 택하자. Lemma \ref{algebra-lemma-hom-exact}에 의해 이는 완전열
$$
0 \to
\Hom_R(M, N) \to
\bigoplus\nolimits_{i = 1, \ldots, n} N
\longrightarrow
\bigoplus\nolimits_{j = 1, \ldots, m} N.
$$
을 준다. $S$를 가역화하고
Proposition \ref{algebra-proposition-localization-exact}를 사용하면 완전열
$$
0 \to
S^{-1}\Hom_R(M, N) \to
\bigoplus\nolimits_{i = 1, \ldots, n} S^{-1}N
\longrightarrow
\bigoplus\nolimits_{j = 1, \ldots, m} S^{-1}N
$$
을 얻는다. 또한 $S^{-1}M$은 완전열
$$
\bigoplus\nolimits_{j = 1, \ldots, m} S^{-1}R
\longrightarrow
\bigoplus\nolimits_{i = 1, \ldots, n} S^{-1}R \to S^{-1}M \to 0
$$
에 속하고, 이것은 Lemma \ref{algebra-lemma-hom-exact}에 의해 완전열
$$
0 \to
\Hom_{S^{-1}R}(S^{-1}M, S^{-1}N) \to
\bigoplus\nolimits_{i = 1, \ldots, n} S^{-1}N
\longrightarrow
\bigoplus\nolimits_{j = 1, \ldots, m} S^{-1}N
$$
을 유도한다. 이는 위의 완전열과 같으므로 결과가 따른다.
\end{proof}

\section{유한 가군과 유한 표시 가군의 특징짓기}
\label{algebra-section-colim-and-hom}

\noindent
환 $R$ 위의 가군 $N$이 주어지면, $N$이 유한 가군인지 또는
유한 표시 가군인지를 함자 $\Hom_R(N, -)$로 특징지을 수 있다.

\begin{lemma}
\label{algebra-lemma-characterize-finite-module-hom}
$R$을 환, $N$을 $R$-가군이라 하자. 다음 조건들은 동치이다.
\begin{enumerate}
\item $N$은 유한 $R$-가군이다.
\item $R$-가군의 임의의 여과 여극한 $M = \colim M_i$에 대하여 사상
$\colim \Hom_R(N, M_i) \to \Hom_R(N, M)$은 단사이다.
\end{enumerate}
\end{lemma}

\begin{proof}
(1)을 가정하고 $N$의 생성원 $x_1, \ldots, x_m$을 택하자.
$N \to M_i$가 가군 사상이고 합성 $N \to M_i \to M$이 영이라 하자.
$M = \colim_{i' \geq i} M_{i'}$이므로 각
$j \in \{1, \ldots, m\}$에 대하여 $x_j$가 $M_{i'}$에서 영으로
사상되는 $i' \geq i$를 찾을 수 있다.
% P01-E0277 / P01-E1277: source article candidate recorded in ERRATA_CANDIDATES.jsonl.
$x_j$가 유한 개뿐이므로 이들 모두에 통하는 하나의 $i'$를 찾을 수 있다.
그러면 합성 $N \to M_i \to M_{i'}$가 영이므로 사상이 단사임,
즉 (2)가 성립함을 얻는다.

\medskip\noindent
(2)를 가정하자. 유한 부분집합 $E \subset N$에 대하여,
$N_E \subset N$을 $E$의 원소들로 생성되는 $R$-부분가군이라 하자.
그러면 $0 = \colim N/N_E$는 여과 여극한이다.
% P01-E0278: frozen maps-into wording recorded in ERRATA_CANDIDATES.jsonl.
따라서 어떤 $E$에 대하여 $\text{id} : N \to N$이 $N_E$ 안으로
사상됨을 알 수 있다. 즉 $N$은 유한 생성이다.
\end{proof}

\noindent
참조를 위하여, 가군의 원소들 사이에 관계식이 있다는 뜻을 정의한다.

\begin{definition}
\label{algebra-definition-relation}
$R$을 환, $M$을 $R$-가군이라 하자.
$n \geq 0$이고 $i = 1, \ldots, n$에 대하여 $x_i \in M$이라 하자.
$M$에서 $x_1, \ldots, x_n$ 사이의 {\it 관계식}은
$\sum_{i = 1, \ldots, n} f_i x_i = 0$을 만족하는 원소들의 열
$f_1, \ldots, f_n \in R$이다.
\end{definition}

\begin{lemma}
\label{algebra-lemma-module-colimit-fp}
$R$을 환, $M$을 $R$-가군이라 하자.
그러면 $M$은 $R$-가군들의 어떤 유향계 $(M_i, \mu_{ij})$의 여극한으로
쓸 수 있으며, 모든 $M_i$는 유한 표시 $R$-가군이다.
\end{lemma}

\begin{proof}
임의의 유한 부분집합 $S \subset M$과 $S$의 원소들 사이의
관계식들의 임의의 유한 모음 $E$를 생각하자.
따라서 각 $s \in S$에는 $x_s \in M$이 대응하고,
각 $e \in E$는 $\sum f_{e, s} x_s = 0$을 만족하는 원소들
$f_{e, s} \in R$의 벡터로 이루어진다.
$M_{S, E}$를 사상
$$
R^{\# E} \longrightarrow R^{\# S}, \quad
(g_e)_{e\in E} \longmapsto (\sum g_e f_{e, s})_{s\in S}.
$$
의 여핵이라 하자. 표준 사상 $M_{S, E} \to M$들이 있다.
$S \subset S'$이고 $E$의 원소들이 이 사상을 통해 $E'$의 관계식들에
대응하면, $M$으로 가는 사상들과 가환하는 명백한 사상
$M_{S, E} \to M_{S', E'}$가 있다.
$I$를 위와 같이 포함관계로 순서가 주어진 쌍 $(S, E)$의 집합이라 하자.
이 유향계의 여극한이 $M$임은 명백하다.
\end{proof}

\begin{lemma}
\label{algebra-lemma-characterize-finitely-presented-module-hom}
$R$을 환, $N$을 $R$-가군이라 하자. 다음 조건들은 동치이다.
\begin{enumerate}
\item $N$은 유한 표시 $R$-가군이다.
\item $R$-가군의 임의의 여과 여극한 $M = \colim M_i$에 대하여 사상
$\colim \Hom_R(N, M_i) \to \Hom_R(N, M)$은 전단사이다.
\end{enumerate}
\end{lemma}

\begin{proof}
(1)을 가정하고 $F_i$들이 유한 자유가군인 완전열
$F_{-1} \to F_0 \to N \to 0$을 택하자. 그러면 $R$-가군 $M$에 관하여
함자적인 완전열
$$
0 \to \Hom_R(N, M) \to \Hom_R(F_0, M) \to \Hom_R(F_{-1}, M)
$$
을 얻는다. 함자 $\Hom_R(F_i, M)$들은
$\Hom_R(R^{\oplus n}, M) = M^{\oplus n}$이므로 여과 여극한과 가환한다.
여과 여극한은 완전하므로
(Lemma \ref{algebra-lemma-directed-colimit-exact})
(2)가 성립함을 알 수 있다.

\medskip\noindent
(2)를 가정하자. Lemma \ref{algebra-lemma-module-colimit-fp}에 의해
각 $i$에 대하여 $N_i$가 유한 표시이도록 $N = \colim N_i$를
여과 여극한으로 쓸 수 있다.
따라서 어떤 $i$에 대하여 $\text{id}_N$은 $N_i$를 통해 분해된다.
이는 $N$이 유한 표시 $R$-가군(즉 $N_i$)의 직합항임을 뜻하므로
따라서 유한 표시이다.
\end{proof}

\section{텐서곱}
\label{algebra-section-tensor-product}

\begin{definition}
\label{algebra-definition-bilinear}
$R$을 환, $M, N, P$를 세 $R$-가군이라 하자.
사상 $f : M \times N \to P$에서 $M \times N$은 두 $R$-가군의
데카르트곱으로만 본다. 각 $x \in M$에 대하여 $N$에서 $P$로 가는
사상 $y\mapsto f(x, y)$가 $R$-선형이고, 각 $y\in N$에 대하여
사상 $x\mapsto f(x, y)$도 $R$-선형이면 이를 {\it $R$-쌍선형}이라고 한다.
\end{definition}

\begin{lemma}
\label{algebra-lemma-tensor-product}
$M, N$을 $R$-가군이라 하자. 다음 보편 성질을 갖는 쌍 $(T, g)$가
존재한다. 여기서 $T$는 $R$-가군이고
$g : M \times N \to T$는 $R$-쌍선형 사상이다.
임의의 $R$-가군 $P$와 임의의 $R$-쌍선형 사상
$f : M \times N \to P$에 대하여, $f = \tilde{f} \circ g$를 만족하는
유일한 $R$-선형 사상 $\tilde{f} : T \to P$가 존재한다.
달리 말하면 다음 도표가 가환한다.
$$
\xymatrix{
M \times N \ar[rr]^f \ar[dr]_g & & P\\
& T \ar[ur]_{\tilde f}
}
$$
더 나아가 $(T, g)$와 $(T', g')$가 이 성질을 갖는 두 쌍이면,
$j\circ g = g'$를 만족하는 유일한 동형
$j : T \to T'$가 존재한다.
\end{lemma}

\noindent
위의 보편 성질을 만족하는 $R$-가군 $T$를 $R$-가군 $M$과 $N$의
{\it 텐서곱}이라고 하며, $M \otimes_R N$으로 나타낸다.

\begin{proof}
먼저 그러한 $R$-가군 $T$의 존재를 증명한다.
$M, N$을 $R$-가군이라 하자.
$T$를 몫가군 $P/Q$라 하자. 여기서 $P$는 자유 $R$-가군
$R^{(M \times N)}$이고 $Q$는 다음 꼴의 모든 원소로 생성되는
$R$-가군이다($x\in M, y\in N$).
\begin{align*}
(x + x', y) - (x, y) - (x', y), \\
(x, y + y') - (x, y) - (x, y'), \\
(ax, y) - a(x, y), \\
(x, ay) - a(x, y)
\end{align*}
$\pi : M \times N \to T$를 자연 사상이라 하자.
쌍선형 조건을 확인할 때 위 관계식들로부터 알 수 있듯이 이 사상은
$R$-쌍선형이다.
% P01-E0280: source naming construction recorded in ERRATA_CANDIDATES.jsonl.
상 $\pi(x, y) = x \otimes y$로 나타내면 이 원소들은 $T$를 생성한다.
이제 $f : M \times N \to P$를 $R$-쌍선형 사상이라 하자.
사상 $f'(x \otimes y) = f(x, y)$를 확장하여
$f' : T \to P$를 정의할 수 있다. 명백히 $f = f'\circ \pi$이다.
더욱이 $f'$은 생성집합
$\{x \otimes y : x\in M, y\in N\}$에서의 값에 의해 유일하게 결정된다.
같은 성질을 만족하는 다른 쌍 $(T', g')$가 있다고 하자.
그러면 $g' = j\circ g$, $g = j'\circ g'$를 만족하는 유일한
$j : T \to T'$와 $j' : T' \to T$가 존재한다.
% P01-E0279 / P01-E1278: frozen reversed composites recorded in ERRATA_CANDIDATES.jsonl.
그런데 두 사상 $(j\circ j') \circ g$와 $g$는 보편 성질을 만족하므로
유일성에 의해 서로 같고, 따라서 $j'\circ j$는 $T$ 위의 항등사상이다.
마찬가지로 $(j'\circ j) \circ g' = g'$이고
$j\circ j'$는 $T'$ 위의 항등사상이다. 따라서 $j$는 동형이다.
\end{proof}

\begin{lemma}
\label{algebra-lemma-flip-tensor-product}
$M, N, P$를 $R$-가군이라 하자. 그러면 쌍선형 사상들
% P01-E1279: frozen direct-sum bilinear map recorded in ERRATA_CANDIDATES.jsonl.
\begin{align*}
(x, y) & \mapsto y \otimes x\\
(x + y, z) & \mapsto x \otimes z + y \otimes z\\
(r, x) & \mapsto rx
\end{align*}
은 다음 유일한 동형들을 유도한다.
\begin{align*}
M \otimes_R N & \to N \otimes_R M, \\
(M\oplus N)\otimes_R P & \to (M \otimes_R P)\oplus(N \otimes_R P),  \\
R \otimes_R M & \to M
\end{align*}
\end{lemma}

\begin{proof}
생략한다.
\end{proof}

\noindent
두 $R$-가군의 텐서곱을 유한 개의 $R$-가군으로 일반화하고,
다중 텐서곱과 다중선형 사상 사이의 대응을 세울 수 있다.
거의 같은 구성을 사용하면 다음을 증명할 수 있다.

\begin{lemma}
\label{algebra-lemma-multilinear}
$M_1, \ldots, M_r$를 $R$-가군이라 하자.
$R$-가군 T와 $R$-다중선형 사상
$g : M_1\times \ldots \times M_r \to T$로 이루어진 쌍 $(T, g)$가
다음 보편 성질을 만족하도록 존재한다.
임의의 $R$-다중선형 사상 $f : M_1\times \ldots \times M_r \to P$에
대하여 $f'\circ g = f$를 만족하는 유일한 $R$-가군 준동형
$f' : T \to P$가 존재한다. 그러한 가군 $T$는 유일한 동형을 제외하고
유일하다. 이를 $M_1\otimes_R \ldots \otimes_R M_r$로 나타내고,
보편 다중선형 사상을
$(m_1, \ldots, m_r) \mapsto m_1 \otimes \ldots \otimes m_r$로 나타낸다.
\end{lemma}

\begin{proof}
생략한다.
\end{proof}

\begin{lemma}
\label{algebra-lemma-transitive}
준동형들
$$
(M \otimes_R N)\otimes_R P \to
M \otimes_R N \otimes_R P \to
M \otimes_R (N \otimes_R P)
$$
에서
$f((x \otimes y)\otimes z) = x \otimes y \otimes z$이고
$g(x \otimes y \otimes z) = x \otimes (y \otimes z)$이며
$x\in M, y\in N, z\in P$라고 하자. 이들은 잘 정의되고 동형이다.
\end{lemma}

\begin{proof}
$f$가 잘 정의되고 동형임을 증명하겠다. 같은 증명이 $g$에도 적용된다.
임의의 $z\in P$를 고정하자. 그러면 $x\in M, y\in N$에 대하여 사상
$(x, y)\mapsto x \otimes y \otimes z$는 $x$와 $y$에 관해
$R$-쌍선형이므로, $f_z(x \otimes y) = x \otimes y \otimes z$로 보내는
준동형 $f_z : M \otimes N \to M \otimes N \otimes P$를 유도한다.
이제 $(w, z)\mapsto f_z(w)$로 주어지는
$(M \otimes N)\times P \to M \otimes N \otimes P$를 생각하자.
이 사상은 $R$-쌍선형이므로
$f : (M \otimes_R N)\otimes_R P \to M \otimes_R N \otimes_R P$를
유도하고 $f((x \otimes y)\otimes z) = x \otimes y \otimes z$이다.
역사상을 구성하기 위하여 사상
$\pi : M \times N \times P \to (M \otimes N)\otimes P$가
$R$-삼중선형임을 주목하자. 따라서 이는 보편 성질과 양립하는
$R$-선형 사상
$h : M \otimes N \otimes P \to (M \otimes N)\otimes P$를 유도한다.
여기서 $h(x \otimes y \otimes z) = (x \otimes y)\otimes z$임을 알 수 있다.
$f$와 $h$의 명시적 표현으로부터 $f\circ h$와 $h\circ f$는 각각
$M \otimes N \otimes P$와 $(M \otimes N)\otimes P$의 항등사상이다.
따라서 $f$는 원하는 동형이다.
\end{proof}

\noindent
귀납법을 사용하면 이것이 다중 텐서곱으로 확장됨을 알 수 있다.
Lemma \ref{algebra-lemma-flip-tensor-product}와 합치면
$R$-가군의 범주에서 텐서곱 연산은 결합적이고 가환적이며 분배적이다.

\begin{definition}
\label{algebra-definition-bimodule}
아벨군 $N$이 $A$-가군이면서 $B$-가군이고, 모든 $a \in A$와
$b \in B$에 대하여 $a$에 의한 곱셈과 $b$에 의한 곱셈이 가환하여
모든 $n \in N$에 대해 $b(an) = a(bn)$이면, 이 아벨군을
{\it $(A, B)$-쌍가군}이라고 한다. 이 상황에서는 보통 $B$의 작용을
오른쪽에 쓴다. 즉 $b \in B$와 $n \in N$에 대하여 $n$에 $b$를 곱한
결과를 $nb$로 나타낸다. 이 규약 아래에서 위의 양립성은 모든
$a\in A, b\in B, x\in N$에 대하여 $(ax)b = a(xb)$라는 것이다.
$(A, B)$-쌍가군 $N$을 나타내는 약어로 $_AN_B$를 사용한다.
\end{definition}

\begin{lemma}
\label{algebra-lemma-tensor-with-bimodule}
$A$-가군 $M$, $B$-가군 $P$, 그리고 $(A, B)$-쌍가군 $N$에 대하여,
가군 $(M \otimes_A N)\otimes_B P$와 $M \otimes_A(N \otimes_B P)$에는
모두 $(A, B)$-쌍가군 구조를 줄 수 있으며, 더욱이
$$
(M \otimes_A N)\otimes_B P \cong M \otimes_A(N \otimes_B P).
$$
\end{lemma}

\begin{proof}
선험적으로 $M \otimes_A N$은 $A$-가군이지만, 다음과 같이
$B$-가군 구조를 줄 수 있다.
$$
(x \otimes y)b = x \otimes yb, \quad x\in M, y\in N, b\in B
$$
따라서 $M \otimes_A N$은 $(A, B)$-쌍가군이 된다.
$N \otimes_B P$도 마찬가지이고, 따라서
$(M \otimes_A N)\otimes_B P$와 $M \otimes_A(N \otimes_B P)$도 그러하다.
Lemma \ref{algebra-lemma-transitive}에 의해 이 두 가군은 같은 사상을 통해
% P01-E0281: source duplicated wording recorded in ERRATA_CANDIDATES.jsonl.
$A$-가군으로서도 $B$-가군으로서도 동형이다.
\end{proof}

\begin{lemma}
\label{algebra-lemma-hom-from-tensor-product}
임의의 세 $R$-가군 $M, N, P$에 대하여
$$
\Hom_R(M \otimes_R N, P) \cong \Hom_R(M, \Hom_R(N, P))
$$
이다.
\end{lemma}

\begin{proof}
$R$-선형 사상 $\hat{f}\in \Hom_R(M \otimes_R N, P)$는
$R$-쌍선형 사상 $f : M \times N \to P$에 대응한다.
각 $x\in M$에 대하여 사상 $y\mapsto f(x, y)$는 보편 성질에 의해
$R$-선형이다. 따라서 $f$는 사상
$\phi_f : M \to \Hom_R(N, P)$에 대응한다. 이 사상은
$$
\phi_f(ax + y)(z) =
f(ax + y, z) = af(x, z)+f(y, z) =
(a\phi_f(x)+\phi_f(y))(z),
$$
이고, 이것이 모든 $a \in R$, $x \in M$, $y \in M$, $z \in N$에
대하여 성립하므로 $R$-선형이다. 반대로 임의의
$f \in \Hom_R(M, \Hom_R(N, P))$는
$(x, y)\mapsto f(x)(y)$로 주어지는 $R$-쌍선형 사상
$M \times N \to P$를 정의한다. 따라서 이는 두 가군
$\Hom_R(M \otimes_R N, P)$와 $\Hom_R(M, \Hom_R(N, P))$ 사이의
자연스러운 일대일 대응이다.
\end{proof}

\begin{lemma}[텐서곱은 여극한과 가환한다]
\label{algebra-lemma-tensor-products-commute-with-limits}
$(M_i, \mu_{ij})$를 준순서 집합 $I$ 위의 계라 하자.
$N$을 $R$-가군이라 하자. 그러면
$$
\colim (M_i \otimes N) \cong (\colim M_i)\otimes N.
$$
더욱이 이 동형은 준동형들
$\mu_i \otimes 1: M_i \otimes N \to M \otimes N$에 의해 유도된다.
여기서 $M = \colim_i M_i$이고 자연 사상은 $\mu_i : M_i \to M$이다.
\end{lemma}

\begin{proof}
첫 번째 증명. Lemma \ref{algebra-lemma-hom-from-tensor-product}에 의해 함자
$M' \mapsto M' \otimes_R N$은 함자 $N' \mapsto \Hom_R(N, N')$의
왼쪽 수반이다. 따라서 $M' \mapsto M' \otimes_R N$은 모든
여극한과 가환한다. 범주론의
Lemma \ref{categories-lemma-adjoint-exact}를 보라.

\medskip\noindent
두 번째 직접 증명. $P = \colim (M_i \otimes N)$라 하고 쌍대사영을
$\lambda_i : M_i \otimes N \to P$라 하자.
$M = \colim M_i$라 하고 쌍대사영을 $\mu_i : M_i \to M$라 하자.
그러면 모든 $i\leq j$에 대하여 다음 도표가 가환한다.
$$
\xymatrix{
M_i \otimes N \ar[r]_{\mu_i \otimes 1} \ar[d]_{\mu_{ij} \otimes 1} &
M \otimes N \ar[d]^{\text{id}} \\
M_j \otimes N \ar[r]^{\mu_j \otimes 1} &
M \otimes N
}
$$
Lemma \ref{algebra-lemma-homomorphism-limit}에 의해 이 사상들은
$\mu_i \otimes 1 = \psi \circ \lambda_i$를 만족하는 유일한 준동형
$\psi : P \to M \otimes N$을 유도한다.

\medskip\noindent
역사상을 구성하자. 각 $i\in I$에 대하여 표준 $R$-쌍선형 사상
$g_i : M_i \times N \to M_i \otimes N$이 있다.
이는 $\widehat{\phi} \circ (\mu_i \times 1) = \lambda_i \circ g_i$를
만족하는 유일한 사상 $\widehat{\phi} : M \times N \to P$를 유도한다.
이 사상은 $R$-쌍선형이다. 따라서 $R$-선형 사상
$\phi : M \otimes N \to P$를 유도한다.
다음 가환도표로부터
$$
\xymatrix{
M_i \times N \ar[r]^{g_i} \ar[d]^{\mu_i \times \text{id}} &
M_i \otimes N\ar[r]_{\text{id}} \ar[d]_{\lambda_i} &
M_i \otimes N \ar[d]_{\mu_i \otimes \text{id}} \ar[rd]^{\lambda_i} \\
M \times N \ar[r]^{\widehat{\phi}} &
P \ar[r]^{\psi} & M \otimes N \ar[r]^{\phi} & P
}
$$
$\psi\circ\widehat{\phi} = g$임을 알 수 있다. 여기서
$g : M \times N \to M \otimes N$은 표준 $R$-쌍선형 사상이다.
따라서 $\psi\circ\phi$는 $M \otimes N$ 위의 항등사상이다.
오른쪽 사각형과 삼각형으로부터 $\phi\circ\psi$도 $P$ 위의
항등사상이다.
\end{proof}

\begin{lemma}
\label{algebra-lemma-tensor-product-exact}
완전열
\begin{align*}
M_1\xrightarrow{f} M_2\xrightarrow{g} M_3 \to 0
\end{align*}
인 $R$-가군과 준동형의 열 및 임의의 $R$-가군 $N$이 주어졌다고 하자.
그러면 열
\begin{equation}
\label{algebra-equation-2ndex}
M_1\otimes N\xrightarrow{f \otimes 1} M_2\otimes N \xrightarrow{g \otimes 1}
M_3\otimes N \to 0
\end{equation}
은 완전하다. 달리 말해 함자 $- \otimes_R N$은
{\it 오른쪽 완전}이다. 이는 원래의 오른쪽 완전열의 각 항을
텐서하는 것이 완전성을 보존한다는 뜻이다.
\end{lemma}

\begin{proof}
모든 $R$-가군 $P$에 대하여 첫 번째 완전열에 함자
$\Hom(-, \Hom(N, P))$를 적용한다. 그러면
$$
0 \to
\Hom(M_3, \Hom(N, P)) \to
\Hom(M_2, \Hom(N, P)) \to
\Hom(M_1, \Hom(N, P))
$$
을 얻으며, 이는 Lemma \ref{algebra-lemma-hom-exact} (1)에 의해 완전하다.
Lemma \ref{algebra-lemma-hom-from-tensor-product}에 의해 이 열은
$$
0 \to \Hom(M_3 \otimes N, P) \to
\Hom(M_2 \otimes N, P) \to \Hom(M_1 \otimes N, P)
$$
이 되고, 따라서 이 열도 완전하다.
Lemma \ref{algebra-lemma-hom-exact} (1)을 다시 사용하면 원하는 완전열을 얻는다.
\end{proof}

\begin{remark}
\label{algebra-remark-tensor-product-not-exact}
그러나 일반적으로 텐서곱은 완전열을 보존하지 않는다.
달리 말해 $M_1 \to M_2 \to M_3$가 완전하더라도, 임의의
$R$-가군 $N$에 대하여
$M_1 \otimes N \to M_2 \otimes N \to M_3 \otimes N$이
완전일 필요는 없다.
\end{remark}











\begin{example}
\label{algebra-example-tensor-product-not-exact}
$\mathbf{Z}$-가군의 사상으로 본 단사 사상
$2 : \mathbf{Z}\to \mathbf{Z}$를 생각하자.
$N = \mathbf{Z}/2$라 하자. 그러면 유도된 사상
$\mathbf{Z} \otimes \mathbf{Z}/2 \to \mathbf{Z} \otimes \mathbf{Z}/2$는
단사가 아니다. 실제로
$x \otimes y\in \mathbf{Z} \otimes \mathbf{Z}/2$에 대하여
$$
(2 \otimes 1)(x \otimes y) = 2x \otimes y = x \otimes 2y = x \otimes 0 = 0
$$
이다. 따라서 유도된 사상은 영사상이지만
$\mathbf{Z} \otimes N\neq 0$이다.
\end{example}

\begin{remark}
\label{algebra-remark-flat-module}
$R$-가군 $N$에 대하여 함자 $-\otimes_R N$이 완전하다면,
즉 $N$과 텐서하는 것이 모든 완전열을 보존한다면,
% P01-E0282: source missing-article candidate recorded in ERRATA_CANDIDATES.jsonl.
$N$을 {\it 평탄} $R$-가군이라고 한다.
이는 나중에 Section \ref{algebra-section-flat}에서 논의한다.
\end{remark}

\begin{lemma}
\label{algebra-lemma-tensor-finiteness}
$R$을 환이라 하고 $M$과 $N$을 $R$-가군이라 하자.
\begin{enumerate}
\item $N$과 $M$이 유한이면 $M \otimes_R N$도 유한이다.
\item $N$과 $M$이 유한 표시이면 $M \otimes_R N$도 유한 표시이다.
\end{enumerate}
\end{lemma}

\begin{proof}
$M$이 유한이라고 가정하자. 그러면 표시
$0 \to K \to R^{\oplus n} \to M \to 0$을 택할 수 있다.
이는 Lemma \ref{algebra-lemma-tensor-product-exact}에 의해 완전열
$K \otimes_R N \to N^{\oplus n} \to M \otimes_R N \to 0$을 준다.
$N$도 유한이면 $M \otimes_R N$은 유한 가군의 몫이므로 유한이다.
Lemma \ref{algebra-lemma-extension}을 보라. 마찬가지로 $N$과 $M$이 모두
유한 표시이면 $K$가 유한이고, $M \otimes_R N$은 유한 표시 가군
$N^{\oplus n}$을 유한 가군, 즉 $K \otimes_R N$으로 나눈 몫임을
알 수 있다. 따라서 유한 표시이다. Lemma \ref{algebra-lemma-extension}을 보라.
\end{proof}

\begin{lemma}
\label{algebra-lemma-tensor-localization}
$M$을 $R$-가군이라 하자. 그러면 $S^{-1}R$-가군 $S^{-1}M$과
$S^{-1}R \otimes_R M$은 표준적으로 동형이며, 표준 동형
$f : S^{-1}R \otimes_R M \to S^{-1}M$은 다음과 같이 주어진다.
$$
f((a/s) \otimes m) = am/s, \forall a \in R, m \in M, s \in S
$$
\end{lemma}

\begin{proof}
사상 $f' : S^{-1}R \times M \to S^{-1}M$을
$f'(a/s, m) = am/s$로 주면 명백히 쌍선형이다.
따라서 보편 성질에 의해 이 사상은 보조정리의 명제에 나온 유일한
$S^{-1}R$-가군 준동형
$f : S^{-1}R \otimes_R M \to S^{-1}M$을 유도한다.
사실 $S^{-1}M$의 모든 원소는 $m/s$, $m\in M, s\in S$ 꼴이고,
$S^{-1}R \otimes_R M$의 모든 원소는 $1/s \otimes m$ 꼴이다.
후자를 보이려면 $S^{-1}R \otimes_R M$의 원소를 다음과 같이 쓰면 된다.
$$
\sum_k \frac{a_k}{s_k} \otimes m_k =
\sum_k \frac{a_k t_k}{s} \otimes m_k =
\frac{1}{s} \otimes \sum_k {a_k t_k}m_k = \frac{1}{s} \otimes m
$$
% P01-E0283: source standalone-fragment candidate recorded in ERRATA_CANDIDATES.jsonl.
여기서 $m = \sum_k {a_k t_k}m_k$이다. 그러면 $f$가 전사임은 명백하다.
또한 $f(\frac{1}{s} \otimes m) = m/s = 0$이면
$M$에서 $tm = 0$인 $t \in S$가 존재한다. 이때
$$
\frac{1}{s} \otimes m = \frac{1}{st} \otimes tm = \frac{1}{st} \otimes 0 = 0
$$
이므로 $f$는 단사이다.
\end{proof}

\begin{lemma}
\label{algebra-lemma-tensor-product-localization}
$M, N$을 $R$-가군이라 하자. 그러면 다음과 같이 주어지는 표준
$S^{-1}R$-가군 동형
$f : S^{-1}M \otimes_{S^{-1}R}S^{-1}N \to S^{-1}(M \otimes_R N)$이 있다.
$$
f((m/s)\otimes(n/t)) = (m \otimes n)/st
$$
\end{lemma}

\begin{proof}
Lemma \ref{algebra-lemma-tensor-with-bimodule}과
Lemma \ref{algebra-lemma-tensor-localization}을 반복하여 사용하면,
$S^{-1}R$이 $(R, S^{-1}R)$-쌍가군임에 유의할 때
다음 두 $S^{-1}R$-가군이 동형임을 알 수 있다.
\begin{align*}
S^{-1}(M \otimes_R N) & \cong S^{-1}R \otimes_R (M \otimes_R N)\\
 & \cong S^{-1}M \otimes_R N\\
 & \cong (S^{-1}M \otimes_{S^{-1}R}S^{-1}R)\otimes_R N\\
 & \cong S^{-1}M \otimes_{S^{-1}R}(S^{-1}R \otimes_R N)\\
 & \cong S^{-1}M \otimes_{S^{-1}R}S^{-1}N
\end{align*}
이 동형이 보조정리에서 말한 동형임은 쉽게 알 수 있다.
\end{proof}

\section{텐서대수}
\label{algebra-section-tensor-algebra}

\noindent
$R$을 환, $M$을 $R$-가군이라 하자.
$R$ 위의 $M$의 {\it 텐서대수}를 다음 비가환 $R$-대수로 정의한다.
$$
\text{T}(M) = \text{T}_R(M) =
\bigoplus\nolimits_{n \geq 0} \text{T}^n(M)
$$
여기서
$\text{T}^0(M) = R$,
$\text{T}^1(M) = M$,
$\text{T}^2(M) = M \otimes_R M$,
$\text{T}^3(M) = M \otimes_R M \otimes_R M$ 등이다.
곱셈은 순수 텐서 위에서 다음 규칙으로 정의한 뒤 선형성으로 확장한다.
$$
(x_1 \otimes x_2 \otimes \ldots \otimes x_n)
\cdot
(y_1 \otimes y_2 \otimes \ldots \otimes y_m)
=
x_1 \otimes x_2 \otimes \ldots \otimes x_n \otimes
y_1 \otimes y_2 \otimes \ldots \otimes y_m
$$

\medskip\noindent
$R$ 위의 $M$의 {\it 외대수 $\wedge(M)$}를 $\text{T}(M)$을
원소 $x \otimes x \in \text{T}^2(M)$들로 생성되는 양쪽 이데알로
나눈 몫으로 정의한다. $\wedge^n(M)$에서 순수 텐서
$x_1 \otimes \ldots \otimes x_n$의 상을
$x_1 \wedge \ldots \wedge x_n$으로 나타낸다.
이 원소들은 $\wedge^n(M)$을 생성하고, 각 $x_i$에 관해 $R$-선형이며,
두 $x_i$가 같으면 영이다. 즉 $x_1, x_2, \ldots, x_n$의 함수로서
교대적이다. $\wedge(M)$의 곱셈은 등급가환이다. 즉 모든
$x \in M$과 $y \in M$은 $x \wedge y = - y \wedge x$를 만족한다.

\medskip\noindent
그 예로 $M = Rx_1 \oplus \ldots \oplus Rx_n$이 유한 자유가군인 경우를
들 수 있다. 이 경우 $\wedge(M)$은 다음 원소들을 기저로 갖는
자유 $R$-가군이다.
$$
x_{i_1} \wedge \ldots \wedge x_{i_r}
$$
여기서 $0 \leq r \leq n$이고
$1 \leq i_1 < i_2 < \ldots < i_r \leq n$이다.

\medskip\noindent
$R$ 위의 $M$의 {\it 대칭대수 $\text{Sym}(M)$}를
$\text{T}(M)$을 원소
$x \otimes y - y \otimes x \in \text{T}^2(M)$들로 생성되는
양쪽 이데알로 나눈 몫으로 정의한다.
$\text{Sym}^n(M)$에서 순수 텐서 $x_1 \otimes \ldots \otimes x_n$의
상을 간단히 $x_1 \ldots x_n$으로 나타낸다.
이 원소들은 $\text{Sym}^n(M)$을 생성하고, 각 $x_i$에 관해
$R$-선형이며, 원소들의 열 $x_1, \ldots, x_n$이 열
$x_1', \ldots, x_n'$의 순열이면
$x_1 \ldots x_n = x_1' \ldots x_n'$이다.
따라서 $\text{Sym}(M)$이 가환임을 알 수 있다.

\medskip\noindent
그 예로 $M = Rx_1 \oplus \ldots \oplus Rx_n$이 유한 자유가군인
경우를 들 수 있다. 이 경우
$\text{Sym}(M) = R[x_1, \ldots, x_n]$은 다항식 대수이다.

\begin{lemma}
\label{algebra-lemma-free-tensor-algebra}
$R$을 환, $M$을 $R$-가군이라 하자.
$M$이 자유 $R$-가군이면 각 대칭승과 외승도 자유가군이다.
\end{lemma}

\begin{proof}
생략한다. 다만 위의 유한 자유가군인 경우를 보라.
\end{proof}

\begin{lemma}
\label{algebra-lemma-presentation-sym-exterior}
$R$을 환이라 하자.
$M_2 \to M_1 \to M \to 0$을 $R$-가군의 완전열이라 하자.
다음 완전열들이 있다.
$$
M_2 \otimes_R \text{Sym}^{n - 1}(M_1)
\to
\text{Sym}^n(M_1)
\to
\text{Sym}^n(M)
\to
0
$$
그리고 마찬가지로
$$
M_2 \otimes_R \wedge^{n - 1}(M_1)
\to
\wedge^n(M_1)
\to
\wedge^n(M)
\to
0
$$
이 있다.
\end{lemma}

\begin{proof}
생략한다.
\end{proof}

\begin{lemma}
\label{algebra-lemma-present-sym-wedge}
$R$을 환, $M$을 $R$-가군이라 하자.
$x_i$, $i \in I$를 $R$-가군으로서 $M$의 주어진 생성계라 하자.
$n \geq 2$라 하자. 다음 표준 완전열이 존재한다.
$$
\bigoplus_{1 \leq j_1 < j_2 \leq n}
\bigoplus_{i_1, i_2 \in I}
\text{T}^{n - 2}(M)
\oplus
\bigoplus_{1 \leq j_1 < j_2 \leq n}
\bigoplus_{i \in I}
\text{T}^{n - 2}(M)
\to
\text{T}^n(M)
\to
\wedge^n(M)
\to
0
$$
여기서 첫 번째 직합항의 순수 텐서
$m_1 \otimes \ldots \otimes m_{n - 2}$는 다음으로 사상된다.
\begin{align*}
\underbrace{
m_1 \otimes \ldots \otimes x_{i_1} \otimes \ldots
\otimes x_{i_2} \otimes \ldots \otimes m_{n - 2}
}_{\text{with } x_{i_1} \text{ and } x_{i_2}
\text{ occupying slots } j_1 \text{ and } j_2
\text{ in the tensor}} \\
+
\underbrace{
m_1 \otimes \ldots \otimes x_{i_2} \otimes \ldots
\otimes x_{i_1} \otimes \ldots \otimes m_{n - 2}
}_{\text{with } x_{i_2} \text{ and } x_{i_1}
\text{ occupying slots } j_1 \text{ and } j_2
\text{ in the tensor}}
\end{align*}
그리고 두 번째 직합항의
$m_1 \otimes \ldots \otimes m_{n - 2}$는 다음으로 사상된다.
$$
\underbrace{
m_1 \otimes \ldots \otimes x_i \otimes \ldots
\otimes x_i \otimes \ldots \otimes m_{n - 2}
}_{\text{with } x_{i} \text{ and } x_{i}
\text{ occupying slots } j_1 \text{ and } j_2
\text{ in the tensor}}
$$
또한 다음 표준 완전열이 있다.
$$
\bigoplus_{1 \leq j_1 < j_2 \leq n}
\bigoplus_{i_1, i_2 \in I}
\text{T}^{n - 2}(M)
\to
\text{T}^n(M)
\to
\text{Sym}^n(M)
\to
0
$$
여기서 순수 텐서 $m_1 \otimes \ldots \otimes m_{n - 2}$는
다음으로 사상된다.
\begin{align*}
\underbrace{
m_1 \otimes \ldots \otimes x_{i_1} \otimes \ldots
\otimes x_{i_2} \otimes \ldots \otimes m_{n - 2}
}_{\text{with } x_{i_1} \text{ and } x_{i_2}
\text{ occupying slots } j_1 \text{ and } j_2
\text{ in the tensor}} \\
-
\underbrace{
m_1 \otimes \ldots \otimes x_{i_2} \otimes \ldots
\otimes x_{i_1} \otimes \ldots \otimes m_{n - 2}
}_{\text{with } x_{i_2} \text{ and } x_{i_1}
\text{ occupying slots } j_1 \text{ and } j_2
\text{ in the tensor}}
\end{align*}
\end{lemma}

\begin{proof}
생략한다.
\end{proof}

\begin{lemma}
\label{algebra-lemma-present-wedge}
$A \to B$를 환 준동형, $M$을 $B$-가군이라 하자.
$n > 1$이라 하자. $A$-선형 사상
$M \otimes_A \ldots \otimes_A M \to \wedge^n_B(M)$의 핵은,
$i \not = j$, $m_1, \ldots, m_n \in M$일 때 $m_i = m_j$인 원소
$m_1 \otimes \ldots \otimes m_n$들과,
$i \not = j$, $m_1, \ldots, m_n \in M$, $b \in B$일 때의 원소
$m_1 \otimes \ldots \otimes bm_i \otimes \ldots \otimes m_n -
m_1 \otimes \ldots \otimes bm_j \otimes \ldots \otimes m_n$들에 의해
$A$-가군으로서 생성된다.
\end{lemma}

\begin{proof}
생략한다.
\end{proof}

\begin{lemma}
\label{algebra-lemma-colimit-tensor-algebra}
\begin{slogan}
텐서대수를 취하는 것은 여과 여극한과 가환한다.
\end{slogan}
$R$을 환, $M_i$를 $R$-가군의 유향계라 하자. 그러면
$\colim_i \text{T}(M_i) = \text{T}(\colim_i M_i)$이고,
대칭대수와 외대수에 대해서도 마찬가지이다.
\end{lemma}

\begin{proof}
생략한다. 힌트:
Lemma \ref{algebra-lemma-tensor-products-commute-with-limits}를 적용하라.
\end{proof}

\begin{lemma}
\label{algebra-lemma-tensor-algebra-localization}
% P01-E0284: frozen tensor-algebra symbol switch recorded in ERRATA_CANDIDATES.jsonl.
$R$을 환, $S \subset R$을 곱셈적 부분집합이라 하자.
그러면 임의의 $R$-가군 $M$에 대하여
$S^{-1}T_R(M) = T_{S^{-1}R}(S^{-1}M)$이다.
대칭대수와 외대수에 대해서도 마찬가지이다.
\end{lemma}

\begin{proof}
생략한다. 힌트: Lemma \ref{algebra-lemma-tensor-product-localization}을 적용하라.
\end{proof}

\section{기저변환}
\label{algebra-section-base-change}

\noindent
대수에서 기저변환을 다음과 같이 형식적으로 도입한다.

\begin{definition}
\label{algebra-definition-base-change}
$\varphi : R \to S$를 환 준동형, $M$을 $S$-가군이라 하자.
$R \to R'$을 임의의 환 준동형이라 하자.
$R \to R'$에 의한 $\varphi$의 {\it 기저변환}은 환 준동형
$R' \to S \otimes_R R'$이다. 이 상황에서 흔히
$S' = S \otimes_R R'$이라고 쓴다.
$S$-가군 $M$의 {\it 기저변환}은 $S'$-가군 $M \otimes_R R'$이다.
\end{definition}

\noindent
어떤 변수들의 모음 $x_i$, $i \in I$와 어떤 다항식들의 모음
$f_j \in R[x_i]$, $j \in J$에 대하여
$S = R[x_i]/(f_j)$라고 하자. 그러면
$S \otimes_R R' = R'[x_i]/(f'_j)$이다.
여기서 $f'_j \in R'[x_i]$는 $R \to R'$에 의해 유도된 사상
$R[x_i] \to R'[x_i]$ 아래에서 $f_j$의 상이다.
이 간단한 관찰이 기저변환을 이해하는 열쇠이다.

\begin{lemma}
\label{algebra-lemma-base-change-finiteness}
$R \to S$를 환 준동형, $M$을 $S$-가군이라 하자.
$R \to R'$을 환 준동형이라 하고,
$S' = S \otimes_R R'$와 $M' = M \otimes_R R'$을 그 기저변환들이라 하자.
\begin{enumerate}
\item $M$이 유한 $S$-가군이면 기저변환 $M'$은 유한 $S'$-가군이다.
\item $M$이 유한 표시 $S$-가군이면 기저변환 $M'$은 유한 표시
$S'$-가군이다.
\item $R \to S$가 유한형이면 기저변환 $R' \to S'$도 유한형이다.
\item $R \to S$가 유한 표시이면 기저변환 $R' \to S'$도 유한 표시이다.
\end{enumerate}
\end{lemma}

\begin{proof}
(1)의 증명. 전사인 $S$-선형 사상
$S^{\oplus n} \to M \to 0$을 택하자.
Lemmas \ref{algebra-lemma-flip-tensor-product}와
\ref{algebra-lemma-tensor-product-exact}에 의해 $R^\prime$과 텐서한 결과는
전사 $ {S^\prime}^{\oplus n} \to M^\prime \rightarrow 0$이다.
따라서 $M^\prime$은 유한 생성 $S^\prime$-가군이다.
(2)의 증명. 표시
$S^{\oplus m} \to S^{\oplus n} \to M \to 0$을 택하자.
Lemmas \ref{algebra-lemma-flip-tensor-product}와
\ref{algebra-lemma-tensor-product-exact}에 의해 $R^\prime$과 텐서한 결과는
$S^\prime$-가군 $M^\prime$의 유한 표시
$ {S^\prime}^{\oplus m} \to {S^\prime}^{\oplus n} \to M^\prime \to 0$을
준다. (3)의 증명. 가정에 의해 $I$를 유한하게 택할 수 있으므로
보조정리 앞의 관찰에서 따른다. (4)의 증명. 가정에 의해 $I$와 $J$를
유한하게 택할 수 있으므로 보조정리 앞의 관찰에서 따른다.
\end{proof}

\noindent
$\varphi : R \to S$를 환 준동형이라 하자. $S$-가군 $N$이 주어지면
규칙 $r \cdot n = \varphi(r)n$에 의해 $R$-가군 $N_R$을 얻는다.
이를 때때로 $N$의 $R$로의 {\it 제한}이라고 한다.

\begin{lemma}
\label{algebra-lemma-adjoint-tensor-restrict}
$R \to S$를 환 준동형이라 하자. 함자
$\text{Mod}_S \to \text{Mod}_R$, $N \mapsto N_R$ (제한)과
$\text{Mod}_R \to \text{Mod}_S$, $M \mapsto M \otimes_R S$
(기저변환)는 수반함자이다. 공식으로 쓰면
$$
\Hom_R(M, N_R) = \Hom_S(M \otimes_R S, N)
$$
이다.
\end{lemma}

\begin{proof}
$\alpha : M \to N_R$가 $R$-가군 사상이면
규칙 $\alpha'(m \otimes s) = s\alpha(m)$으로
$\alpha' : M \otimes_R S \to N$을 정의한다.
$\beta : M \otimes_R S \to N$가 $S$-가군 사상이면
규칙 $\beta'(m) = \beta(m \otimes 1)$로
$\beta' : M \to N_R$을 정의한다.
이 구성들이 서로 역임을 확인하는 일은 생략한다.
\end{proof}

\noindent
위 보조정리는 제한이 왼쪽 수반, 즉 기저변환을 가짐을 말한다.
제한은 오른쪽 수반도 갖는다.

\begin{lemma}
\label{algebra-lemma-adjoint-hom-restrict}
$R \to S$를 환 준동형이라 하자. 함자
$\text{Mod}_S \to \text{Mod}_R$, $N \mapsto N_R$ (제한)과
$\text{Mod}_R \to \text{Mod}_S$, $M \mapsto \Hom_R(S, M)$은
수반함자이다. 공식으로 쓰면
$$
\Hom_R(N_R, M) = \Hom_S(N, \Hom_R(S, M))
$$
이다.
\end{lemma}

\begin{proof}
$\alpha : N_R \to M$이 $R$-가군 사상이면
규칙 $\alpha'(n) = (s \mapsto \alpha(sn))$으로
$\alpha' : N \to \Hom_R(S, M)$을 정의한다.
$\beta : N \to \Hom_R(S, M)$가 $S$-가군 사상이면
규칙 $\beta'(n) = \beta(n)(1)$로
$\beta' : N_R \to M$을 정의한다.
이 구성들이 서로 역임을 확인하는 일은 생략한다.
\end{proof}

\begin{lemma}
\label{algebra-lemma-hom-from-tensor-product-variant}
$R \to S$를 환 준동형이라 하자. $S$-가군 $M, N$과
$R$-가군 $P$가 주어지면
$$
\Hom_R(M \otimes_S N, P) = \Hom_S(M, \Hom_R(N, P))
$$
이다.
\end{lemma}

\begin{proof}
이를 직접 증명할 수도 있지만,
Lemmas \ref{algebra-lemma-adjoint-hom-restrict}와
\ref{algebra-lemma-hom-from-tensor-product}의 귀결이기도 하다.
실제로
\begin{align*}
\Hom_R(M \otimes_S N, P)
& =
\Hom_S(M \otimes_S N, \Hom_R(S, P)) \\
& =
\Hom_S(M, \Hom_S(N, \Hom_R(S, P))) \\
& =
\Hom_S(M, \Hom_R(N, P))
\end{align*}
이며, 이는 원하는 결과이다.
\end{proof}

\section{잡사항}
\label{algebra-section-miscellany}

\noindent
본 절의 증명들은 기본 개념을 다룬 Section \ref{algebra-section-rings-basic}의
결과 외에는 어떤 결과도 참조하지 않아야 한다.

\begin{lemma}
\label{algebra-lemma-product-ideals-in-prime}
$R$을 환, $I$와 $J$를 두 이데알, $\mathfrak p$를 곱 $IJ$를 포함하는
소 이데알이라 하자. 그러면 $\mathfrak{p}$는 $I$ 또는 $J$를 포함한다.
\end{lemma}

\begin{proof}
반대로 가정하고 $x \in I \setminus \mathfrak p$와
$y \in J \setminus \mathfrak p$를 택하자. 그 곱은
$IJ \subset \mathfrak p$의 원소이며, 이는 $\mathfrak p$가
소 이데알이라는 가정에 모순이다.
\end{proof}

\begin{lemma}[소 이데알 회피]
\label{algebra-lemma-silly}
\begin{slogan}
1. 아핀 스킴에서 유한 개의 점이 열린 부분집합에 포함되면 이 점들은
더 작은 주 열린 부분집합에 포함된다.
2. 아핀 열린집합들은 아핀 스킴의 주어진 유한집합의 근방들 사이에서
공종이다.
\end{slogan}
$R$을 환이라 하자. $I_i \subset R$, $i = 1, \ldots, r$와
$J \subset R$을 이데알이라 하자. 다음을 가정한다.
\begin{enumerate}
\item 모든 $i = 1, \ldots, r$에 대하여 $J \not\subset I_i$이다.
\item $I_i$들 가운데 많아야 두 개를 제외하면 모두 소 이데알이다.
\end{enumerate}
그러면 모든 $i$에 대하여 $x\not\in I_i$인 $x \in J$가 존재한다.
\end{lemma}

\begin{proof}
$r = 1$일 때 결과는 참이다. $r = 2$이면
$x \not \in I_1$인 $x, y \in J$와 $y \not \in I_2$를 택하자.
$x \in I_2$이고 $y \in I_1$인 경우가 아니면 증명이 끝난다.
이 경우에도 원소 $x + y$는 $I_1$에 속할 수 없고
(속한다면 $x + y - y \in I_1$이기 때문이다),
$I_2$에도 속할 수 없다.

\medskip\noindent
$r \geq 3$일 때 결과가 $r - 1$에 대해 성립한다고 가정하자.
번호를 다시 매긴 뒤 $I_r$이 소 이데알이라고 가정할 수 있다.
$I_i$들 사이에는 포함관계가 없다고도 가정할 수 있다.
모든 $i = 1, \ldots, r - 1$에 대하여 $x \not \in I_i$인
$x \in J$를 택하자. $x \not\in I_r$이면 증명이 끝난다.
따라서 $x \in I_r$이라고 가정하자.
$J I_1 \ldots I_{r - 1} \subset I_r$이면
Lemma \ref{algebra-lemma-product-ideals-in-prime}에 의해 $J \subset I_r$이어서
모순이다. $y \not \in I_r$인
$y \in J I_1 \ldots I_{r - 1}$을 택하자. 그러면 $x + y$가 조건을 만족한다.
\end{proof}

\begin{lemma}
\label{algebra-lemma-silly-silly}
$R$을 환, $x \in R$, $I \subset R$을 이데알이라 하고,
$\mathfrak p_i$, $i = 1, \ldots, r$를 소 이데알들이라 하자.
모든 $i = 1, \ldots, r$에 대하여
$x + I \not \subset \mathfrak p_i$라고 가정하자.
그러면 모든 $i$에 대하여 $x + y \not \in \mathfrak p_i$인
$y \in I$가 존재한다.
\end{lemma}

\begin{proof}
$\mathfrak p_i$들 사이에는 포함관계가 없다고 가정할 수 있다.
순서를 바꾸어 $i < s$일 때 $x \not \in \mathfrak p_i$이고
$i \geq s$일 때 $x \in \mathfrak p_i$라고 가정할 수 있다.
$s = r + 1$이면 증명이 끝난다. 그렇지 않으면
$y \not \in \mathfrak p_s$인 $y \in I$를 찾을 수 있다.
$f \not \in \mathfrak p_s$인
$f \in \bigcap_{i < s} \mathfrak p_i$를 택하자.
% P01-E0285: frozen element-versus-subset wording recorded in ERRATA_CANDIDATES.jsonl.
그러면 $x + fy$는 $\mathfrak p_1, \ldots, \mathfrak p_s$에
포함되지 않는다. 따라서 $s$에 관한 귀납법으로 증명이 끝난다.
\end{proof}

\begin{lemma}[중국인의 나머지 정리]
\label{algebra-lemma-chinese-remainder}
$R$을 환이라 하자.
\begin{enumerate}
\item $a \not = b$일 때 $I_a + I_b = R$인 이데알
$I_1, \ldots, I_r$가 주어지면
$I_1 \cap \ldots \cap I_r = I_1I_2\ldots I_r$이고
$R/(I_1I_2\ldots I_r) \cong R/I_1 \times \ldots \times R/I_r$이다.
\item $\mathfrak m_1, \ldots, \mathfrak m_r$가 서로 다른 극대 이데알들이면
$a \not = b$일 때 $\mathfrak m_a + \mathfrak m_b = R$이고
위 결과를 적용할 수 있다.
\end{enumerate}
\end{lemma}

\begin{proof}
먼저 $I_1 \cap \ldots \cap I_r = I_1 \ldots I_r$임을 증명하자.
이는 유도된 환 준동형
$R/(I_1 \ldots I_r) \rightarrow R/I_1 \times \ldots \times R/I_r$이
단사임도 함의한다. 이데알은 임의의 환 원소와 곱하는 것에 관하여
닫혀 있으므로 포함관계
$I_1 \cap \ldots \cap I_r \supset I_1 \ldots I_r$는 항상 성립한다.
반대 포함관계를 증명하기 위하여 이데알들
$$
I_1 \ldots \hat I_i \ldots I_r,\quad i = 1, \ldots, r
$$
이 환 $R$을 생성한다고 주장한다. 이를 $r$에 관한 귀납법으로 증명한다.
$r = 2$일 때 성립한다. $r > 2$이면 $R$은 이데알들
$I_1 \ldots \hat I_i \ldots I_{r - 1}$,
$i = 1, \ldots, r - 1$의 합이다. 따라서 $I_r$은 이데알들
$I_1 \ldots \hat I_i \ldots I_r$, $i = 1, \ldots, r - 1$의 합이다.
이데알의 순서를 뒤집어 같은 논증을 적용하면 $I_1$은 이데알들
$I_1 \ldots \hat I_i \ldots I_r$, $i = 2, \ldots, r$의 합이다.
가정에 의해 $R = I_1 + I_r$이므로 $R$은 위에 표시한 이데알들의 합이다.
따라서 합이 일인 원소들
$a_i \in I_1 \ldots \hat I_i \ldots I_r$를 찾을 수 있다.
이 등식에 $I_1 \cap \ldots \cap I_r$의 원소를 곱하면
반대 포함관계를 얻는다. 이제 표준 사상
$R/(I_1 \ldots I_r) \rightarrow R/I_1 \times \ldots \times R/I_r$이
전사임을 보여야 한다. 이를 위해 위에서 얻은 등식
$1 = \sum_{i=1}^r a_i$ 위에서 이 사상이 어떻게 작용하는지 생각하자.
한편으로 환 준동형은 1을 1로 보내고, 다른 한편으로
임의의 $a_i$의 상은 $j \neq i$일 때 $R/I_j$에서 영이다.
따라서 $R/I_i$에서 $a_i$의 상은 항등원이다.
그러므로 임의의 원소
$(\bar{b_1}, \ldots, \bar{b_r}) \in R/I_1 \times \ldots \times R/I_r$가
주어지면 원소 $\sum_{i=1}^r a_i \cdot b_i$는 $R$에서의 역상이다.

\medskip\noindent
(2)를 보기 위하여, 서로 다른 극대 이데알의 정의 자체로부터
$a \neq b$일 때 $\mathfrak{m}_a + \mathfrak{m}_b = R$이므로
위 결과를 적용할 수 있다.
\end{proof}

\begin{lemma}
\label{algebra-lemma-matrix-left-inverse}
$R$을 환이라 하자. $n \geq m$이라 하자.
$A$를 $R$에 계수를 갖는 $n \times m$ 행렬이라 하자.
$J \subset R$을 $A$의 $m \times m$ 소행렬식들로 생성되는
이데알이라 하자.
\begin{enumerate}
\item 임의의 $f \in J$에 대하여
$BA = f 1_{m \times m}$인 $m \times n$ 행렬 $B$가 존재한다.
\item $f \in R$이고 어떤 $m \times n$ 행렬 $B$에 대하여
$BA = f 1_{m \times m}$이면 $f^m \in J$이다.
\end{enumerate}
\end{lemma}

\begin{proof}
$I \subset \{1, \ldots, n\}$이고 $|I| = m$일 때,
$E_I$를 사영
$$
R^{\oplus n} = \bigoplus\nolimits_{i \in \{1, \ldots, n\}} R
\longrightarrow \bigoplus\nolimits_{i \in I} R
$$
의 $m \times n$ 행렬이라 하고 $A_I = E_I A$로 둔다.
즉 $A_I$는 $A$에서 지표가 $I$에 속하는 행들을 취해 만든
$m \times m$ 행렬이다. $B_I$를 $A_I$의 수반행렬
(여인수 행렬의 전치)이라 하자. 즉
$A_I B_I = B_I A_I = \det(A_I) 1_{m \times m}$이다.
$A$의 $m \times m$ 소행렬식들은
$I \subset \{1, \ldots, n\}$이고 $|I| = m$인 모든 경우에 대한
행렬식 $\det A_I$이다. $f \in J$이면 어떤 $c_I \in R$에 대하여
$f = \sum c_I \det(A_I)$로 쓸 수 있다.
$B = \sum c_I B_I E_I$로 두면 (1)이 성립함을 알 수 있다.

\medskip\noindent
$f 1_{m \times m} = BA$이면 Cauchy--Binet 공식
(\ref{algebra-item-cauchy-binet})에 의해
$f^m = \sum b_I \det(A_I)$이다. 여기서 $b_I$는 $B$의 열들 가운데
지표가 $I$에 속하는 열들로 이루어진 $m \times m$ 행렬의 행렬식이다.
\end{proof}

\begin{lemma}
\label{algebra-lemma-matrix-right-inverse}
$R$을 환이라 하자. $n \geq m$이라 하자.
$A = (a_{ij})$를 $R$에 계수를 갖는 $n \times m$ 행렬이라 하고,
다음 블록 형태로 쓰자.
$$
A =
\left(
\begin{matrix}
A_1 \\
A_2
\end{matrix}
\right)
$$
여기서 $A_1$의 크기는 $m \times m$이다.
$B$를 $A_1$의 수반행렬(여인수 행렬의 전치)이라 하자. 그러면
$$
AB = 
\left(
\begin{matrix}
f 1_{m \times m} \\
C
\end{matrix}
\right)
$$
이다. 여기서 $f = \det(A_1)$이고, $c_{ij}$는 부호를 제외하면
$A$에서 행 $1, \ldots, \hat j, \ldots, m, i$에 대응하는
$m \times m$ 소행렬식의 행렬식이다.
\end{lemma}

\begin{proof}
% P01-E0286: frozen missing identity matrix recorded in ERRATA_CANDIDATES.jsonl.
수반행렬은 $A_1B = B A_1 = f$를 만족하므로 $AB$의 표현에서
첫 번째 블록은 옳다. 다음에 유의하라.
$$
c_{ij} = \sum\nolimits_k a_{ik}b_{kj} = \sum (-1)^{j + k}a_{ik} \det(A_1^{jk})
$$
% P01-E0287 / P01-E1280: frozen superscript switch recorded in ERRATA_CANDIDATES.jsonl.
여기서 $A_1^{ij}$는 $A_1$에서 $j$번째 행과 $k$번째 열을 제거한
행렬을 뜻한다. 마지막 표현은 보조정리의 명제에 나온 행렬의 행 전개이다.
\end{proof}

\begin{lemma}
\label{algebra-lemma-map-cannot-be-injective}
\begin{slogan}
유한 자유가군 사이의 사상은 정의역의 랭크가 공역의 랭크보다 크면
단사일 수 없다.
\end{slogan}
$R$을 영환이 아닌 환이라 하자. $n \geq 1$이라 하자.
$M$을 $< n$개의 원소로 생성되는 $R$-가군이라 하자.
그러면 임의의 $R$-가군 사상 $f : R^{\oplus n} \to M$은
영이 아닌 핵을 갖는다.
\end{lemma}

\begin{proof}
전사 $R^{\oplus n - 1} \to M$을 택하자.
사상 $f$를 사상 $f' : R^{\oplus n} \to R^{\oplus n - 1}$로
들어 올릴 수 있다(Lemma \ref{algebra-lemma-lift-map}).
$f'$이 영이 아닌 핵을 가짐을 보이면 충분하다.
사상 $f' : R^{\oplus n} \to R^{\oplus n - 1}$는 행렬
$A = (a_{ij})$로 주어진다. $a_{ij}$ 중 하나가 멱영원이 아니라고 하고,
이를 $a = a_{ij}$라 하자. 그러면 $R$을 국소화 $R_a$로 바꾸어
$a_{ij}$가 단원이라고 가정할 수 있다.
실제로 국소화는 완전하므로 국소화한 뒤 영이 아닌 핵을 찾으면
처음에도 영이 아닌 핵이 있었다.
Proposition \ref{algebra-proposition-localization-exact}를 보라.
% P01-E0288: frozen base-change terminology recorded in ERRATA_CANDIDATES.jsonl.
이 경우 $R^{\oplus n}$과 $R^{\oplus n - 1}$ 양쪽에서
기저변환을 하여 다음과 같은 경우로 환원할 수 있다.
$$
A =
\left(
\begin{matrix}
1 & 0 & 0 & \ldots \\
0 & a_{22} & a_{23} & \ldots \\
0 & a_{32} & \ldots \\
\ldots & \ldots
\end{matrix}
\right)
$$
따라서 이 경우에는 $n$에 관한 귀납법으로 증명이 끝난다.
그렇지 않으면 각 $a_{ij}$는 멱영원이다.
$I = (a_{ij}) \subset R$로 두자. 어떤 $m \geq 0$에 대하여
$I^{m + 1} = 0$임에 유의하라.
$I^m \not = 0$인 가장 큰 정수를 $m$이라 하자.
그러면 $(I^m)^{\oplus n}$이 사상의 핵에 포함됨을 알 수 있으므로
증명이 끝난다.
\end{proof}





\begin{lemma}
\label{algebra-lemma-rank}
\begin{slogan}
유한 자유가군의 랭크는 잘 정의된다.
\end{slogan}
$R$을 영환이 아닌 환이라 하자. $n, m \geq 0$을 정수라 하자.
$R^{\oplus n}$이 $R$-가군으로서 $R^{\oplus m}$과 동형이면
$n = m$이다.
\end{lemma}

\begin{proof}
Lemma \ref{algebra-lemma-map-cannot-be-injective}로부터 즉시 따른다.
\end{proof}

\section{케일리--해밀턴}
\label{algebra-section-cayley-hamilton}

\begin{lemma}
\label{algebra-lemma-charpoly}
$R$을 환이라 하자. $A = (a_{ij})$를 $R$에 계수를 갖는
$n \times n$ 행렬이라 하자. $P(x) \in R[x]$를 $A$의 특성다항식
($\det(x\text{id}_{n \times n} - A)$로 정의한다)이라 하자.
그러면 $\text{Mat}(n \times n, R)$에서 $P(A) = 0$이다.
\end{lemma}

\begin{proof}
다음 몇 단계를 거쳐 이 문제를 선형대수학의 잘 알려진
케일리--해밀턴 정리로 환원한다.
\begin{enumerate}
\item $\phi :S \rightarrow R$가 환 준동형이고 $b_{ij}$들이 이 사상 아래에서
$a_{ij}$들의 원상이라면, $\phi$가 환 준동형이므로 $S$와 $(b_{ij})$에 대해
명제를 보이면 충분하다.
\item $\psi :R \hookrightarrow S$가 단사 환 준동형이면, $a_{ij}$들을
$S$의 원소로 보아 $S$에 대해 결과를 보이면 충분함이 명백하다.
\item 따라서 먼저 다항식환의 경우인 $R = \mathbf{Z}[X_{ij}]$,
$a_{ij} = X_{ij}$로 환원한 다음, 케일리--해밀턴 정리를 적용할 수 있는
$R = \mathbf{Q}(X_{ij})$의 경우로 다시 환원할 수 있다.
\end{enumerate}
\end{proof}

\begin{lemma}
\label{algebra-lemma-charpoly-module}
$R$을 환이라 하자.
$M$을 유한 $R$-가군이라 하자.
$\varphi : M \to M$을 자기준동형이라 하자.
그러면 $M$의 자기준동형으로서 $P(\varphi) = 0$을 만족하는
일계수 다항식 $P \in R[T]$가 존재한다.
\end{lemma}

\begin{proof}
어떤 생성원들 $x_i \in M$에 대해
$(a_1, \ldots, a_n) \mapsto \sum a_ix_i$로 주어지는 전사
$R$-가군 사상 $R^{\oplus n} \to M$을 택한다.
$\varphi(x_i) = \sum a_{ij} x_j$를 만족하는
$(a_{i1}, \ldots, a_{in}) \in R^{\oplus n}$을 택한다.
달리 말해, $A = (a_{ij})$일 때 다음 도식은 가환한다.
$$
\xymatrix{
R^{\oplus n} \ar[d]_A \ar[r] & M \ar[d]^\varphi \\
R^{\oplus n} \ar[r] & M
}
$$
Lemma \ref{algebra-lemma-charpoly}에 의해 $P(A) = 0$을 만족하는
일계수 다항식 $P$가 존재한다. 따라서 $P(\varphi) = 0$이다.
\end{proof}

\begin{lemma}
\label{algebra-lemma-charpoly-module-ideal}
$R$을 환이라 하고 $I \subset R$을 이데알이라 하자.
$M$을 유한 $R$-가군이라 하자.
$\varphi(M) \subset IM$을 만족하는 자기준동형
$\varphi : M \to M$을 생각하자.
그러면 $a_j \in I^j$이고 $M$의 자기준동형으로서 $P(\varphi) = 0$을 만족하는
일계수 다항식
% P01-E0289: 동결 원문은 다항식의 소문자 t와 환 R[T]의 대문자 T를 혼용한다.
$P = t^n + a_1 t^{n - 1} + \ldots + a_n \in R[T]$
가 존재한다.
\end{lemma}

\begin{proof}
어떤 생성원들 $x_i \in M$에 대해
$(a_1, \ldots, a_n) \mapsto \sum a_ix_i$로 주어지는 전사
$R$-가군 사상 $R^{\oplus n} \to M$을 택한다.
$\varphi(x_i) = \sum a_{ij} x_j$를 만족하는
$(a_{i1}, \ldots, a_{in}) \in I^{\oplus n}$을 택한다.
달리 말해, $A = (a_{ij})$일 때 다음 도식은 가환한다.
$$
\xymatrix{
R^{\oplus n} \ar[d]_A \ar[r] & M \ar[d]^\varphi \\
I^{\oplus n} \ar[r] & M
}
$$
Lemma \ref{algebra-lemma-charpoly}에 의해 다항식
$P(t) = \det(t\text{id}_{n \times n} - A)$
은 요구된 성질을 모두 만족한다.
\end{proof}

\noindent
흥미로운 응용으로 다음의 놀라운 보조정리를 증명한다.

\begin{lemma}
\label{algebra-lemma-fun}
$R$을 환이라 하자.
$M$을 유한 $R$-가군이라 하자.
$\varphi : M \to M$을 전사 $R$-가군 사상이라 하자.
그러면 $\varphi$는 동형사상이다.
\end{lemma}

\begin{proof}[첫 번째 증명]
$R' = R[x]$라 쓰고, $x$가 $\varphi$로 작용하게 하여 $M$을 유한
$R'$-가군으로 생각하자. $I = (x) \subset R'$로 둔다. $\varphi$가
전사라는 가정에 의해 $IM = M$이다. 따라서 $M$을 $R'$-가군으로 보고,
이데알 $I$와 자기준동형 $\text{id}_M$에
Lemma \ref{algebra-lemma-charpoly-module-ideal}을 적용할 수 있다.
그 결과 $a_j \in I$일 때
$(1 + a_1 + \ldots + a_n)\text{id}_M = 0$을 얻는다.
어떤 $b_j(x) \in R[x]$에 대해 $a_j = b_j(x)x$라고 쓰자.
이를 다시 $\varphi$의 말로 옮기면
$\text{id}_M = -(\sum_{j = 1, \ldots, n} b_j(\varphi)) \varphi$이고,
따라서 $\varphi$는 가역이다.
\end{proof}

\begin{proof}[두 번째 증명]
$R$ 위에서 $M$의 생성원 수에 관해 귀납법을 쓴다. $M$이 한 원소로
생성되면 어떤 이데알 $I \subset R$에 대해 $M \cong R/I$이다.
이 경우 $R$을 $R/I$로 바꾸어 $M = R$이라 할 수 있다.
이때 $\varphi : R \to R$는 어떤 원소 $r \in R$에 의한 $M$ 위의
곱셈으로 주어진다. $\varphi$가 전사이면 $\varphi$가 $1$을 상으로 가져야 하므로
$r$은 가역이고, 따라서 $\varphi$는 가역이다.

\medskip\noindent
이제 $n - 1$개 원소로 생성되는 가군에 대해 보조정리를 증명했다고 하고,
$n$개 원소로 생성되는 가군 $M$을 살펴보자. $A$를 환 $R[t]$라 하고,
$t$가 $\varphi$로 작용하게 하여 $M$을 $A$-가군으로 보자. $M$은
$R$ 위에서 유한이므로 $R[t]$ 위에서도 유한하다. 또한 우리가
$\varphi$의 단사성을 증명하려는 것은 집합론적 성질이므로,
$A$ 위의 자기준동형 $t : M \to M$의 단사성을 증명해도 된다.
이로써 자기준동형이 바탕환의 한 원소에 의한 곱셈인 경우로 문제가
환원되었다. $M$의 생성원들 가운데 처음 $n - 1$개가 생성하는
부분-$A$-가군을 $M' \subset M$이라 하고, 다음 도식을 생각하자.
$$
\xymatrix{
0 \ar[r] & M' \ar[r]\ar[d]^{\varphi\mid_{M'}} & M\ar[d]^\varphi \ar[r] &
M/M' \ar[d]^{\varphi \bmod M'} \ar[r] & 0 \\
0 \ar[r] & M' \ar[r]                                  & M \ar[r]
           & M/M' \ar[r]                                  & 0,
}
$$
여기서 $\varphi$는 바탕의 한 원소에 의한 곱셈이고 $M'$과 $M/M'$도
각각 $A$-가군이므로, $\varphi$의 $M'$에 대한 제한과 $\varphi$가
몫 $M/M'$ 위에 유도하는 사상은 잘 정의된다. $n = 1$인 경우에 의해
사상 $M/M' \to M/M'$는 동형사상이다. 도식 추적을 하면
$\varphi|_{M'}$가 전사이고, 따라서 귀납가정에 의해 $\varphi|_{M'}$는
동형사상이다. 뱀 보조정리에 의해 가운데 세로 사상도 동형사상이다.
\end{proof}

\section{환의 스펙트럼}
\label{algebra-section-spectrum-ring}

\noindent
환의 스펙트럼을 위상공간으로 보는 내용은 대수학 장에 두고,
아핀 스킴은 스킴 장에 두기로 임의로 정한다.

\begin{definition}
\label{algebra-definition-spectrum-ring}
$R$을 환이라 하자.
\begin{enumerate}
\item $R$의 {\it 스펙트럼}은 $R$의 소 이데알들의 집합이다.
보통 $\Spec(R)$로 나타낸다.
\item 부분집합 $T \subset R$가 주어졌을 때, $V(T) \subset \Spec(R)$를
$T$를 포함하는 소 이데알들의 집합, 즉 $V(T) = \{ \mathfrak p \in
\Spec(R) \mid \forall f\in T, f\in \mathfrak p\}$로 둔다.
\item 원소 $f \in R$가 주어졌을 때, $D(f) \subset \Spec(R)$를
$f$를 포함하지 않는 소 이데알들의 집합으로 둔다.
\end{enumerate}
\end{definition}

\begin{lemma}
\label{algebra-lemma-Zariski-topology}
$R$을 환이라 하자.
\begin{enumerate}
\item 환 $R$의 스펙트럼이 공집합인 것과 $R$이 영환인 것은 동치이다.
\item 영환이 아닌 모든 환에는 극대 이데알이 존재한다.
\item 영환이 아닌 모든 환에는 극소 소 이데알이 존재한다.
\item 이데알 $I \subset R$와 소 이데알 $I \subset \mathfrak p$가
주어지면, $I$ 위에서 극소인 소 이데알 $\mathfrak q$, 즉
$I \subset \mathfrak q \subset \mathfrak p$가 존재한다.
\item $T \subset R$이고 $(T)$가 $R$에서 $T$가 생성하는 이데알이면
$V((T)) = V(T)$이다.
\item $I$가 이데알이고 $\sqrt{I}$가 그 근기이면
(\ref{algebra-item-radical-ideal})의 기본 개념에 의해 $V(I) = V(\sqrt{I})$이다.
\item $R$의 이데알 $I$가 주어지면
$\sqrt{I} = \bigcap_{I \subset \mathfrak p} \mathfrak p$이다.
\item 이데알 $I$에 대해 $V(I) = \emptyset$인 것과 $I$가 단위 이데알인
것은 동치이다.
\item $I$, $J$가 $R$의 이데알이면 $V(I) \cup V(J) =
V(I \cap J)$이다.
\item $(I_a)_{a\in A}$가 $R$의 이데알들의 집합이면
$\bigcap_{a\in A} V(I_a) = V(\bigcup_{a\in A} I_a)$이다.
\item $f \in R$이면 $D(f) \amalg V(f) = \Spec(R)$이다.
\item $f \in R$이면 $D(f) = \emptyset$인 것과 $f$가 멱영원인 것은
동치이다.
\item 어떤 단위원 $u \in R$에 대해 $f = u f'$이면
$D(f) = D(f')$이다.
\item $I \subset R$가 이데알이고 $\mathfrak p$가
$\mathfrak p \not\in V(I)$인 $R$의 소 이데알이면,
$\mathfrak p \in D(f)$이고 $D(f) \cap V(I) = \emptyset$인
$f \in R$가 존재한다.
\item $f, g \in R$이면 $D(fg) = D(f) \cap D(g)$이다.
\item $i \in I$에 대해 $f_i \in R$이면
$\bigcup_{i\in I} D(f_i)$는 $\Spec(R)$에서
$V(\{f_i \}_{i\in I})$의 여집합이다.
\item $f \in R$이고 $D(f) = \Spec(R)$이면 $f$는 단위원이다.
\end{enumerate}
\end{lemma}

\begin{proof}
아래의 대응하는 항목에서 각 부분을 다룬다.
\begin{enumerate}
\item (2) 또는 (3)의 직접적인 결과이다.
\item $R$의 모든 진 이데알들의 집합을 $\mathfrak{A}$라 하자.
$(0) \in \mathfrak{A}$가 진 이데알이므로 이 집합은 공집합이 아니며,
포함관계로 순서가 주어진다. $A$를 $\mathfrak A$의 전순서 부분집합이라
하자. 그러면 $\bigcup_{I \in A} I$는 실제로 이데알이다.
% P01-E1281: 동결 원문은 왼쪽 피연산자 1을 수학 모드 밖에 둔다.
모든 $I \in A$에 대해 1 $\notin I$이므로 그 합집합은 1을 포함하지 않고,
따라서 진 이데알이다. 그러므로 $\bigcup_{I \in A} I$는
$\mathfrak{A}$의 원소이며 집합 $A$의 상계이다. 이제 초른의 보조정리에
의해 $\mathfrak{A}$에는 극대 원소가 존재하고, 이것이 구하려던
극대 이데알이다.
\item $R$은 영환이 아니므로 소 이데알인 극대 이데알을 하나 포함한다.
따라서 $R$의 모든 소 이데알들의 집합 $\mathfrak{A}$는 공집합이 아니다.
$\mathfrak{A}$에는 역포함관계로 순서를 준다. $A$를 $\mathfrak{A}$의
전순서 부분집합이라 하자. $J = \bigcap_{I \in A} I$가 이데알임은
꽤 명백하다. 그러나 이것이 소 이데알임은 그만큼 명백하지 않다.
$xy \in J$라 하자. 그러면 모든 $I \in A$에 대해 $xy \in I$이다.
이제 $B = \{I \in A | y \in I\}$라 하고
$K = \bigcap_{I \in B} I$라 하자. $A$가 전순서 집합이므로
$K = J$이거나(이 경우 $y \in J$이므로 끝난다), 아니면 $K \supset J$이고
$K$에 진부분집합으로 포함되는 모든 $I \in A$에 대해 $y \notin I$이다.
하지만 이 이데알들은 소 이데알이므로 그 모든 $I, x \in I$이다.
따라서 $x \in J$이다. 어느 경우든 원하는 대로 $J$는 소 이데알이다.
그러므로 초른의 보조정리에 의해 극대 원소를 얻는데, 이 경우에는
극소 소 이데알이다.
\item $\mathfrak{p}$에 포함되고 $I$를 포함하는 소 이데알들만
고려하면 (3)과 정확히 같은 논증이다.
\item $(T)$는 $T$를 포함하는 가장 작은 이데알이다. 따라서 어떤
이데알 $I$에 대해 $T \subset I$이면 $(T) \subset I$이기도 하다.
그러므로 $I \in V(T)$이면 $I \in V((T))$이기도 하다.
반대쪽 포함관계는 명백하다.
\item $I \subset \sqrt{I}, V(\sqrt{I}) \subset V(I)$이다.
이제 $\mathfrak{p} \in V(I)$라 하고 $x \in \sqrt{I}$라 하자.
어떤 $n$에 대해 $x^n \in I$이다. 따라서 $x^n \in \mathfrak{p}$이다.
그런데 $\mathfrak{p}$가 소 이데알이므로 단순한 귀납논증으로
$x \in \mathfrak{p}$를 얻는다. 따라서
$\sqrt{I} \subset \mathfrak{p}$이고 $\mathfrak{p} \in V(\sqrt{I})$이다.
\item $f \in R \setminus \sqrt{I}$라 하자. 그러면 모든 $n$에 대해
$f^n \notin I$이다. 따라서 $S = \{1, f, f^2, \ldots\}$는 $0$을
포함하지 않는 곱셈적 부분집합이다. $S^{-1}I$를 포함하는 소 이데알
$\bar{\mathfrak{p}} \subset S^{-1}R$를 택한다. $\bar{\mathfrak{p}}$의
$R$에서의 원상을 $\mathfrak{p}$라 하면, 이는 $I$를 포함하고 $S$와
만나지 않는 소 이데알이다. 이로부터
$\bigcap_{I \subset \mathfrak p} \mathfrak p \subset \sqrt{I}$를 얻는다.
이제 $a \in \sqrt{I}$이면 어떤 $n$에 대해 $a^n \in I$이다.
따라서 $I \subset \mathfrak{p}$이면 $a^n \in \mathfrak{p}$이다.
$\mathfrak{p}$가 소 이데알이므로 $a \in \mathfrak{p}$이다.
이로써 등식이 증명되었다.
\item $I$가 단위 이데알이 아닌 것과 $I$가 어떤 극대 이데알에 포함되는
것은 동치이다(이를 보려면 환 $R/I$에 (2)를 적용하라). 그 극대
이데알은 소 이데알이다.
\item $\mathfrak{p} \in V(I) \cup V(J)$이면
$I \subset \mathfrak{p}$ 또는 $J \subset \mathfrak{p}$이고,
따라서 $I \cap J \subset \mathfrak{p}$이다. 이제
$I \cap J \subset \mathfrak{p}$이면 $IJ \subset \mathfrak{p}$이고,
% P01-E0290 / P01-E1282: 동결 원문은 or를 of로 잘못 쓰고 이데알 포함을 원소 관계로 서술한다.
$\mathfrak{p}$가 소 이데알이므로 $I$ 또는 $J$가 $\mathfrak{p}$ 안에 있다.
\item $\mathfrak{p} \in \bigcap_{a \in A} V(I_a) \Leftrightarrow
I_a \subset \mathfrak{p}, \forall a \in A \Leftrightarrow
\mathfrak{p} \in V(\bigcup_{a\in A} I_a)$
\item $\mathfrak{p}$가 소 이데알이고 $f \in R$이면
$f \in \mathfrak{p}$ 또는 $f \notin \mathfrak{p}$이며(엄밀히 둘 중
하나만 성립한다), 이것이 바로 서로소 합집합이라는 말의 뜻이다.
\item $a \in R$가 멱영원이면 어떤 $n$에 대해 $a^n = 0$이다.
따라서 임의의 소 이데알에 대해 $a^n \in \mathfrak{p}$이다.
귀납법으로 $a \in \mathfrak{p}$임을 보일 수 있으므로
$D(a) = \emptyset$이다. 한편 (7)에서 보았듯이 $a \in R$가
멱영원이 아니면 이를 포함하지 않는 소 이데알이 존재한다.
\item $u$가 가역이므로
$f \in \mathfrak{p} \Leftrightarrow uf \in \mathfrak{p}$이다.
\item $\mathfrak{p} \notin V(I)$이면
$\exists f \in I \setminus \mathfrak{p}$이다. 그러면
$f \notin \mathfrak{p}$이므로 $\mathfrak{p} \in D(f)$이다.
또한 $\mathfrak{q} \in D(f)$이면 $f \notin \mathfrak{q}$이므로
$I$는 $\mathfrak{q}$에 포함되지 않는다. 따라서
$D(f) \cap V(I) = \emptyset$이다.
\item $fg \in \mathfrak{p}$이면 $f \in \mathfrak{p}$ 또는
$g \in \mathfrak{p}$이다. 따라서 $f \notin \mathfrak{p}$이고
$g \notin \mathfrak{p}$이면 $fg \notin \mathfrak{p}$이다.
$\mathfrak{p}$가 이데알이므로 $fg \notin \mathfrak{p}$이면
$f \notin \mathfrak{p}$이고 $g \notin \mathfrak{p}$이다.
\item $\mathfrak{p} \in \bigcup_{i \in I} D(f_i) \Leftrightarrow \exists i \in
I, f_i \notin \mathfrak{p} \Leftrightarrow \mathfrak{p} \in \Spec(R)
\setminus V(\{f_i\}_{i \in I})$
\item $D(f) = \Spec(R)$이면 $V(f) = \emptyset$이고,
따라서 $fR = R$이므로 $f$는 단위원이다.
\end{enumerate}
\end{proof}

\noindent
이 보조정리에 의해 Definition \ref{algebra-definition-spectrum-ring}의 부분집합
$V(T)$들은 $\Spec(R)$ 위의 한 위상의 닫힌집합들을 이룬다.
또한 집합 $D(f)$들은 열린집합이며 이 위상의 기저를 이룬다.

\begin{definition}
\label{algebra-definition-Zariski-topology}
$R$을 환이라 하자. 닫힌집합들이 $V(T)$인 $\Spec(R)$ 위의 위상을
{\it 자리스키} 위상이라 한다. 열린 부분집합 $D(f)$들을
$\Spec(R)$의 {\it 표준 열린집합}이라 한다.
\end{definition}

\noindent
문맥상 $\Spec(R)$를 단지 집합으로 보는지 위상공간으로 보는지는
명확할 것이다.

\begin{lemma}
\label{algebra-lemma-spec-functorial}
\begin{slogan}
스펙트럼의 함자성
\end{slogan}
$\varphi : R \to R'$가 환 준동형이라고 하자. 유도된 사상
$$
\Spec(\varphi) : \Spec(R') \longrightarrow \Spec(R),
\quad
\mathfrak p' \longmapsto \varphi^{-1}(\mathfrak p')
$$
은 자리스키 위상에 대해 연속이다. 실제로 임의의 원소 $f \in R$에 대해
$\Spec(\varphi)^{-1}(D(f)) = D(\varphi(f))$이다.
\end{lemma}

\begin{proof}
(\ref{algebra-item-inverse-image-prime})의 기본 개념에 의해
$\mathfrak p := \varphi^{-1}(\mathfrak p')$는 실제로 $R$의 소 이데알이다.
보조정리의 마지막 명제는 정의에서 바로 따르며, 첫 번째 명제를 함의한다.
\end{proof}

\noindent
$\varphi' : R' \to R''$가 두 번째 환 준동형이면 합성
$$
\Spec(R'')
\longrightarrow
\Spec(R')
\longrightarrow
\Spec(R)
$$
은 $\Spec(\varphi' \circ \varphi)$와 같다. 달리 말해, $\Spec$은
환의 범주에서 위상공간의 범주로 가는 반변 함자이다.

\begin{lemma}
\label{algebra-lemma-spec-localization}
$R$을 환이라 하고 $S \subset R$을 곱셈적 부분집합이라 하자.
사상 $R \to S^{-1}R$는 $\Spec$의 함자성에 의해 위상동형
$$
\Spec(S^{-1}R)
\longrightarrow
\{\mathfrak p \in \Spec(R) \mid S \cap \mathfrak p = \emptyset \}
$$
을 유도한다. 여기서 오른쪽에는 $\Spec(R)$의 자리스키 위상에서 유도된
위상을 준다. 역함수는
$\mathfrak p \mapsto S^{-1}\mathfrak p = \mathfrak p(S^{-1}R)$로 주어진다.
\end{lemma}

\begin{proof}
보조정리의 화살표 오른쪽을 $D$라 하자.
소 이데알 $\mathfrak p' \subset S^{-1}R$를 택하고
% P01-E0291: 동결 원문에는 be 동사가 빠져 있다.
$\mathfrak p$를 $R$에서 $\mathfrak p'$의 원상이라 하자.
$\mathfrak p'$는 $1$을 포함하지 않으므로 $\mathfrak p$는 $S$의 어떤
원소도 포함하지 않는다. 따라서 $\mathfrak p \in D$이고 상은 $D$에
포함된다. 이제 $\mathfrak p \in D$라 하자. 가정에 의해 상
$\overline{S}$는 $0$을 포함하지 않는다.
(\ref{algebra-item-localization-zero})의 기본 개념에 의해
$\overline{S}^{-1}(R/\mathfrak p)$는 영환이 아니다.
(\ref{algebra-item-localize-ideal})의 기본 개념에 의해
$S^{-1}R / S^{-1}\mathfrak p = \overline{S}^{-1}(R/\mathfrak p)$는
정역이므로 $S^{-1}\mathfrak p$는 소 이데알이다. 이 환들의 등식은
$R$에서 $S^{-1}\mathfrak p$의 원상이 $\mathfrak p$와 같다는 것도
보여 준다. 실제로 (\ref{algebra-item-localize-nonzerodivisors})의 기본 개념에
의해 $R/\mathfrak p \to \overline{S}^{-1}(R/\mathfrak p)$가 단사이다.
따라서 사상 $\Spec(S^{-1}R) \to \Spec(R)$는 주어진 역함수를 가지며
$D$ 위로 전단사이다. Lemma \ref{algebra-lemma-spec-functorial}에 의해 연속이다.
마지막으로 $D(g) \subset \Spec(S^{-1}R)$를 표준 열린집합이라 하자.
어떤 $h\in R$와 $s\in S$에 대해 $g = h/s$라고 쓰자.
$g$와 $h/1$은 단위원만큼 차이나므로 $\Spec(S^{-1}R)$에서
$D(g) = D(h/1)$이다. 따라서 Lemma \ref{algebra-lemma-spec-functorial}과 위의
전단사성에 의해 $D(g) = D(h/1)$의 상은 $D \cap D(h)$이다.
이로부터 이 사상은 열린사상이기도 하다.
\end{proof}

\begin{lemma}
\label{algebra-lemma-standard-open}
$R$을 환이라 하고 $f \in R$이라 하자. 사상 $R \to R_f$는
$\Spec$의 함자성에 의해 위상동형
$$
\Spec(R_f) \longrightarrow D(f) \subset \Spec(R).
$$
을 유도한다. 역함수는 $\mathfrak p \mapsto \mathfrak p \cdot R_f$로
주어진다.
\end{lemma}

\begin{proof}
Lemma \ref{algebra-lemma-spec-localization}의 특수한 경우이다.
\end{proof}

\noindent
스펙트럼의 모든 ``아핀 열린집합''이 표준 열린집합인 것은 아니다.
Example \ref{algebra-example-affine-open-not-standard}을 보라.

\begin{lemma}
\label{algebra-lemma-spec-closed}
$R$을 환이라 하고 $I \subset R$을 이데알이라 하자. 사상
$R \to R/I$는 $\Spec$의 함자성에 의해 위상동형
$$
\Spec(R/I) \longrightarrow V(I) \subset \Spec(R).
$$
을 유도한다. 역함수는 $\mathfrak p \mapsto \mathfrak p / I$로 주어진다.
\end{lemma}

\begin{proof}
상이 $V(I)$에 포함됨은 즉시 알 수 있다. 반대로
$\mathfrak p \in V(I)$이면 $\mathfrak p \supset I$이므로 이데알
$\mathfrak p /I \subset R/I$를 생각할 수 있다.
(\ref{algebra-item-isomorphism-theorem})의 기본 개념을 쓰면
$(R/I)/(\mathfrak p/I) = R/\mathfrak p$가 정역이고, 따라서
$\mathfrak p/I$는 소 이데알이다. 이로부터 $D(f + I)$의 상이
$D(f) \cap V(I)$임을 즉시 알 수 있으며, 따라서 이 사상은
위상동형이다.
\end{proof}



\begin{lemma}
\label{algebra-lemma-quasi-compact}
\begin{slogan}
환의 스펙트럼은 준콤팩트이다
\end{slogan}
$R$을 환이라 하자. 공간 $\Spec(R)$는 준콤팩트이다.
\end{lemma}

\begin{proof}
% P01-E0292: 동결 원문은 유한 부분덮개가 아니라 유한 덮개에 의한 세분을 말한다.
$\Spec(R)$의 표준 열린집합들에 의한 임의의 덮개를 유한 덮개로
세분할 수 있음을 보이면 충분하다. 따라서 $R$의 원소들의 집합
$\{f_i\}_{i\in I}$에 대해 $\Spec(R) = \cup D(f_i)$라고 하자.
이는 $\cap V(f_i) = \emptyset$이라는 뜻이다.
Lemma \ref{algebra-lemma-Zariski-topology}에 따르면 이는
$V(\{f_i \}) = \emptyset$이라는 뜻이다. 같은 보조정리에 따르면 이는
$f_i$들이 생성하는 이데알이 $R$의 단위 이데알이라는 뜻이다.
따라서 $J \subset I$가 유한일 때 $1$을 다음과 같은 {\it 유한} 합으로
쓸 수 있다: $1 = \sum_{i \in J} r_i f_i$. 그러면
$\Spec(R) = \cup_{i \in J} D(f_i)$가 따른다.
\end{proof}

\begin{lemma}
\label{algebra-lemma-topology-spec}
$R$을 환이라 하자. $X = \Spec(R)$ 위의 위상은 다음 성질들을 갖는다.
\begin{enumerate}
\item $X$는 준콤팩트이다.
\item $X$의 위상에는 준콤팩트 열린집합들로 이루어진 기저가 있다.
\item 임의의 두 준콤팩트 열린집합의 교집합은 준콤팩트이다.
\end{enumerate}
\end{lemma}

\begin{proof}
환의 스펙트럼은 준콤팩트이다. Lemma \ref{algebra-lemma-quasi-compact}을 보라.
위상의 기저는 표준 열린집합들
$D(f) = \Spec(R_f)$
(Lemma \ref{algebra-lemma-standard-open})로 이루어지고, 첫 번째 언급에 의해
이들은 준콤팩트이다. 두 표준 열린집합의 교집합은
$D(f) \cap D(g) = D(fg)$이므로 준콤팩트이다. 임의의 두 준콤팩트
열린집합 $U, V \subset X$가 주어지면
$U = D(f_1) \cup \ldots \cup D(f_n)$ 및
$V = D(g_1) \cup \ldots \cup D(g_m)$로 쓸 수 있다. 그러면
% P01-E0293: 동결 원문의 마지막 합집합에는 지표 범위가 표시되지 않는다.
$U \cap V = \bigcup D(f_ig_j)$는 준콤팩트이다.
\end{proof}

\section{국소환}
\label{algebra-section-local-rings}

\noindent
국소환은 대수기하학의 핵심적인 도구이다.

\begin{definition}
\label{algebra-definition-local-ring}
{\it 국소환}은 극대 이데알을 정확히 하나 갖는 환이다.
$R$이 국소환이면 그 극대 이데알을 흔히 $\mathfrak m_R$로 나타내며,
체 $R/\mathfrak m_R$를 국소환 $R$의 {\it 잉여체}라 한다.
흔히 ``$(R, \mathfrak m)$을 국소환이라 하자'' 또는
``$(R, \mathfrak m, \kappa)$를 국소환이라 하자''라고 말하여
$R$이 국소환이고 $\mathfrak m$이 유일한 극대 이데알이며
$\kappa = R/\mathfrak m$이 잉여체임을 나타낸다.
{\it 국소환의 국소 준동형}은 $R$과 $S$가 국소환이고
$\varphi(\mathfrak m_R) \subset \mathfrak m_S$를 만족하는 환 사상
$\varphi : R \to S$이다. $R$과 $S$가 국소환임이 주어져 있다면
``{\it 국소환 사상 $\varphi : R \to S$}''이라는 말은 $\varphi$가
국소환의 국소 준동형이라는 뜻이다.
\end{definition}

\noindent
체는 국소환이다. 체 사이의 모든 환 사상은 국소환의 국소 준동형이다.

\medskip\noindent
환 $R$을 소 이데알 $\mathfrak p$에서 국소화한 $R_\mathfrak p$는
극대 이데알 $\mathfrak p R_\mathfrak p$를 갖는 국소환이다.
실제로 Lemma \ref{algebra-lemma-spec-localization}에 의해 $R_\mathfrak p$의 모든
소 이데알은 소 이데알 $\mathfrak p R_\mathfrak p$에 포함된다
(따라서 이것은 $R_\mathfrak p$의 유일한 극대 이데알이다).
$R_\mathfrak p$의 잉여체를 $\kappa(\mathfrak p)$로 나타내고,
이를 {$\mathfrak p$의 잉여체}라 한다. Proposition
\ref{algebra-proposition-localize-quotient}에 의해 $\kappa(\mathfrak p)$를
정역 $R/\mathfrak p$의 분수체와 동일시할 수 있다. 합성
$$
\Spec(\kappa(\mathfrak p)) \to \Spec(R_\mathfrak p) \to \Spec(R)
$$
을 통해 정의역의 유일한 점은 공역의 점 $\mathfrak p$로 간다.

\medskip\noindent
$\varphi : R \to S$를 환 사상이라 하자. $\mathfrak q \subset S$를
소 이데알이라 하고 $R$의 소 이데알
$\mathfrak p = \varphi^{-1}(\mathfrak q)$를 생각하자.
$\varphi(\mathfrak p) \subset \mathfrak q$이므로 유도된 환 사상
$$
R_\mathfrak p \to S_\mathfrak q,\quad r/g \mapsto \varphi(r)/\varphi(g)
$$
은 국소환 사상이고, 유도된 잉여체의 사상
$\kappa(\mathfrak p) \to \kappa(\mathfrak q)$를 얻는다.


\begin{example}
\label{algebra-example-not-local}
$R$이 국소환이고 $\mathfrak p \subset R$가 극대가 아닌 소 이데알이면
$R \to R_\mathfrak p$는 국소 준동형이 아니다.
\end{example}

\begin{lemma}
\label{algebra-lemma-characterize-local-ring}
$R$을 환이라 하자. 다음 조건들은 서로 동치이다.
\begin{enumerate}
\item $R$은 국소환이다.
\item $\Spec(R)$에는 닫힌점이 정확히 하나 있다.
\item $R$에는 극대 이데알 $\mathfrak m$이 있고,
$R \setminus \mathfrak m$의 모든 원소는 단위원이다.
\item $R$은 영환이 아니며, 모든 $x \in R$에 대해 $x$ 또는
$1 - x$ 중 적어도 하나는 가역이다.
\end{enumerate}
\end{lemma}

\begin{proof}
$R$을 환이라 하고 $\mathfrak m$을 극대 이데알이라 하자.
$x \in R \setminus \mathfrak m$이고 $x$가 단위원이 아니면
$x$를 포함하는 극대 이데알 $\mathfrak m'$이 존재한다. 따라서
$R$에는 적어도 두 극대 이데알이 있다. 역으로 $\mathfrak m'$이 또 다른
극대 이데알이면 $x \in \mathfrak m'$, $x \not \in \mathfrak m$인
그런 원소를 택하자. 명백히 $x$는 단위원이 아니다. 이로써 (1)과 (3)의
동치가 증명되었다. (1)과 (2)의 동치는 동어반복이다.
$R$이 국소환이면 $x$는 $\mathfrak m$에 속하거나 속하지 않으므로
(4)가 성립한다. (4)가 성립하고 $\mathfrak m$, $\mathfrak m'$이 서로
다른 극대 이데알이라 하자. 중국인의 나머지 정리
(Lemma \ref{algebra-lemma-chinese-remainder})에 의해
$x \bmod \mathfrak m' = 0$이고 $x \bmod \mathfrak m = 1$인
$x \in R$를 택할 수 있다. 이 원소 $x$는 가역이 아니며
$1 - x$도 가역이 아니므로 모순이다. 따라서 (4)와 (1)은 동치이다.
\end{proof}

\begin{lemma}
\label{algebra-lemma-characterize-local-ring-map}
$\varphi : R \to S$를 환 사상이라 하고 $R$과 $S$가 국소환이라고 하자.
다음 조건들은 서로 동치이다.
\begin{enumerate}
\item $\varphi$는 국소환 사상이다.
\item $\varphi(\mathfrak m_R) \subset \mathfrak m_S$이다.
\item $\varphi^{-1}(\mathfrak m_S) = \mathfrak m_R$이다.
\item 임의의 $x \in R$에 대해 $\varphi(x)$가 $S$에서 가역이면
$x$는 $R$에서 가역이다.
\end{enumerate}
\end{lemma}

\begin{proof}
조건 (1)과 (2)는 정의에 의해 동치이다. (3)이 성립하면 (2)가 성립한다.
역으로 (2)가 성립하면 $\varphi^{-1}(\mathfrak m_S)$는 극대 이데알
$\mathfrak m_R$을 포함하는 소 이데알이므로
$\varphi^{-1}(\mathfrak m_S) = \mathfrak m_R$이다. 마지막으로
Lemma \ref{algebra-lemma-characterize-local-ring}에 의해 (4)는 (2)의 대우이다.
\end{proof}

\begin{remark}
\label{algebra-remark-fundamental-diagram}
환 사상 $\varphi : R \to S$와 소 이데알 $\mathfrak p \subset R$에
따르는 기본적인 가환도식은 다음과 같다.
$$
\xymatrix{
\kappa(\mathfrak p) \otimes_R S =
S_{\mathfrak p}/{\mathfrak p}S_{\mathfrak p}
&
S_{\mathfrak p} \ar[l]
&
S \ar[r] \ar[l]
&
S/\mathfrak pS \ar[r]
&
(R \setminus \mathfrak p)^{-1}S/\mathfrak pS
\\
\kappa(\mathfrak p) =
R_{\mathfrak p}/{\mathfrak p}R_{\mathfrak p} \ar[u]
&
R_{\mathfrak p} \ar[u] \ar[l]
&
R \ar[u] \ar[r] \ar[l]
&
R/\mathfrak p \ar[u] \ar[r]
&
\kappa(\mathfrak p) \ar[u]
}
$$
이 도식에서 가장 왼쪽 세로열과 가장 오른쪽 세로열은 동일하다.
스펙트럼 위에서 수평 사상들은 각자의 상으로의 위상동형을 유도하고,
정사각형들은 위상공간의 섬유곱 정사각형을 유도한다
(Lemmas \ref{algebra-lemma-spec-localization} 및 \ref{algebra-lemma-spec-closed} 참조).
따라서 $\mathfrak p$가 Spec 위 사상의 상에 속하는 것과
$S \otimes_R \kappa(\mathfrak p)$가 영환이 아닌 것은 동치이다.
$\mathfrak p$ 위에 놓이는 소 이데알 $\mathfrak q \subset S$, 즉
$\mathfrak p = \varphi^{-1}(\mathfrak q)$인 것이 존재하면 이 도식을
다음 도식으로 확장할 수 있다.
$$
\xymatrix{
\kappa(\mathfrak q) = S_{\mathfrak q}/{\mathfrak q}S_{\mathfrak q}
&
S_{\mathfrak q} \ar[l]
&
S \ar[r] \ar[l]
&
S/\mathfrak q \ar[r]
&
\kappa(\mathfrak q)
\\
\kappa(\mathfrak p) \otimes_R S =
S_{\mathfrak p}/{\mathfrak p}S_{\mathfrak p} \ar[u]
&
S_{\mathfrak p} \ar[u] \ar[l]
&
S \ar[u] \ar[r] \ar[l]
&
S/\mathfrak pS \ar[u] \ar[r]
&
(R \setminus \mathfrak p)^{-1}S/\mathfrak pS \ar[u]
\\
\kappa(\mathfrak p) =
R_{\mathfrak p}/{\mathfrak p}R_{\mathfrak p} \ar[u]
&
R_{\mathfrak p} \ar[u] \ar[l]
&
R \ar[u] \ar[r] \ar[l]
&
R/\mathfrak p \ar[u] \ar[r]
&
\kappa(\mathfrak p) \ar[u]
}
$$
이 도식에서도 여전히 가장 왼쪽 세로열과 가장 오른쪽 세로열은 동일하고,
스펙트럼 위에서 수평 사상들은 각자의 상으로의 위상동형을 유도한다.
\end{remark}

\begin{lemma}
\label{algebra-lemma-in-image}
$\varphi : R \to S$를 환 사상이라 하고 $\mathfrak p$를 $R$의
소 이데알이라 하자. 다음 조건들은 서로 동치이다.
\begin{enumerate}
\item $\mathfrak p$는 $\Spec(S) \to \Spec(R)$의 상에 속한다.
\item $S \otimes_R \kappa(\mathfrak p) \not = 0$이다.
\item $S_{\mathfrak p}/\mathfrak p S_{\mathfrak p} \not = 0$이다.
\item $(S/\mathfrak pS)_{\mathfrak p} \not = 0$이다.
\item $\mathfrak p = \varphi^{-1}(\mathfrak pS)$이다.
\end{enumerate}
\end{lemma}

\begin{proof}
첫 두 조건의 동치는 이미 Remark \ref{algebra-remark-fundamental-diagram}에서
보았다. 나머지는 이를 다시 표현한 것뿐이다.
\end{proof}

\begin{remark}
\label{algebra-remark-local-ring-fibre}
$R \to S$를 환 사상이라 하자. $\mathfrak q$를 $R$의 소 이데알
$\mathfrak p$ 위에 놓이는 $S$의 소 이데알이라 하자.
Remark \ref{algebra-remark-fundamental-diagram}에 따르면 소 이데알
$\mathfrak q$는 섬유환 $F = S \otimes_R \kappa(\mathfrak p)$의 유일한
소 이데알 $\overline{\mathfrak q}$에 대응한다. 그러면
$$
F_{\overline{\mathfrak q}}
\cong S_\mathfrak q \otimes_{R_\mathfrak p} \kappa(\mathfrak p)
\cong S_\mathfrak q/\mathfrak p S_\mathfrak q
$$
이다. 실제로 명백한 환 사상
$F \to S_\mathfrak q \otimes_{R_\mathfrak p} \kappa(\mathfrak p)$가 있고,
이는 $F \to F_{\overline{\mathfrak q}}$와 동형임을 쉽게 알 수 있다.
두 번째 동형은 $\kappa(\mathfrak p)$가 $R_\mathfrak p$를
$\mathfrak pR_\mathfrak p$로 나눈 몫이라는 사실에서 따른다.
\end{remark}

\section{환의 야코브슨 근기}
\label{algebra-section-radical}

\noindent
환 $R$의 {\it 야코브슨 근기} $\text{rad}(R)$는 $R$의 모든 극대
이데알의 교집합임을 상기하자. $R$이 국소환이면 $\text{rad}(R)$는
$R$의 극대 이데알이다.

\begin{lemma}
\label{algebra-lemma-contained-in-radical}
$R$을 야코브슨 근기 $\text{rad}(R)$를 갖는 환이라 하자.
$I \subset R$를 이데알이라 하자. 다음 조건들은 서로 동치이다.
\begin{enumerate}
\item $I \subset \text{rad}(R)$이다.
\item $1 + I$의 모든 원소는 $R$의 단위원이다.
\end{enumerate}
이 경우 $R/I$의 단위원으로 가는 $R$의 모든 원소는 단위원이다.
\end{lemma}

\begin{proof}
$f \in \text{rad}(R)$이면 $R$의 모든 극대 이데알 $\mathfrak m$에 대해
$f \in \mathfrak m$이다. 따라서 $R$의 모든 극대 이데알
$\mathfrak m$에 대해 $1 + f \not \in \mathfrak m$이다.
그러므로 $\Spec(R)$의 닫힌 부분집합 $V(1 + f)$는 공집합이다.
Lemma \ref{algebra-lemma-Zariski-topology}에 의해 $1 + f$는 단위원이다.

\medskip\noindent
역으로 모든 $f \in I$에 대해 $1 + f$가 단위원이라고 하자.
$\mathfrak m$이 극대 이데알이고 $I \not \subset \mathfrak m$이면
$I + \mathfrak m = R$이다. 따라서 어떤 $g \in \mathfrak m$과
$f \in I$에 대해 $1 = f + g$이다. 그러면 $g = 1 + (-f)$는
단위원이 아니므로 모순이다.

\medskip\noindent
마지막 명제를 위해 $f \in R$가 $R/I$의 단위원으로 간다고 하자.
$f \bmod I$의 곱셈 역원으로 가는 $g \in R$를 찾을 수 있다.
그러면 $fg = 1 \bmod I$이다. 따라서 (2)에 의해 $fg$는 $R$의
단위원이고, 이로부터 $f$가 단위원임이 따른다.
\end{proof}

\begin{lemma}
\label{algebra-lemma-surjective-on-spec-units}
$\varphi : R \to S$를 유도된 사상 $\Spec(S) \to \Spec(R)$가
전사인 환 사상이라 하자. 원소 $x \in R$가 단위원인 것과
$\varphi(x) \in S$가 단위원인 것은 동치이다.
\end{lemma}

\begin{proof}
$x$가 단위원이면 $\varphi(x)$도 단위원이다. 역으로 $\varphi(x)$가
단위원이면 모든 $\mathfrak q \in \Spec(S)$에 대해
$\varphi(x) \not \in \mathfrak q$이다. 따라서 모든
$\mathfrak q \in \Spec(S)$에 대해
$x \not \in \varphi^{-1}(\mathfrak q) = \Spec(\varphi)(\mathfrak q)$이다.
$\Spec(\varphi)$가 전사이므로 Lemma \ref{algebra-lemma-Zariski-topology}의
(17)에 의해 $x$가 단위원이라고 결론 내린다.
\end{proof}

\section{나카야마의 보조정리}
\label{algebra-section-nakayama}

\noindent
\cite{MatCA}에서 인용한다. ``이 간단하지만 중요한 보조정리는
T.~Nakayama, G.~Azumaya, W.~Krull에 의한 것이다. 우선권은 분명하지
않으며, 보통 나카야마의 보조정리라고 부르지만 고 Nakayama 교수는
그 이름을 좋아하지 않았다.''

\begin{lemma}[나카야마의 보조정리]
\label{algebra-lemma-NAK}
\begin{reference}
\cite[1.M Lemma (NAK) page 11]{MatCA}
\end{reference}
\begin{history}
\cite{MatCA}에서 인용한다. ``이 간단하지만 중요한 보조정리는
T.~Nakayama, G.~Azumaya, W.~Krull에 의한 것이다. 우선권은 분명하지
않으며, 보통 나카야마의 보조정리라고 부르지만 고 Nakayama 교수는
그 이름을 좋아하지 않았다.''
\end{history}
$R$을 야코브슨 근기 $\text{rad}(R)$를 갖는 환이라 하자.
$M$을 $R$-가군이라 하고 $I \subset R$를 이데알이라 하자.
\begin{enumerate}
\item
\label{algebra-item-nakayama}
$IM = M$이고 $M$이 유한이면 $fM = 0$인 $f \in 1 + I$가 존재한다.
\item $IM = M$이고 $M$이 유한이며 $I \subset \text{rad}(R)$이면
$M = 0$이다.
\item $N, N' \subset M$이고 $M = N + IN'$이며 $N'$이 유한이면,
$fM \subset N$이고 $M_f = N_f$인 $f \in 1 + I$가 존재한다.
\item $N, N' \subset M$이고 $M = N + IN'$이며 $N'$이 유한이고
$I \subset \text{rad}(R)$이면 $M = N$이다.
\item $N \to M$이 가군 사상이고 $N/IN \to M/IM$이 전사이며
$M$이 유한이면, $N_f \to M_f$가 전사인 $f \in 1 + I$가 존재한다.
\item $N \to M$이 가군 사상이고 $N/IN \to M/IM$이 전사이며
$M$이 유한이고 $I \subset \text{rad}(R)$이면 $N \to M$은 전사이다.
\item $x_1, \ldots, x_n \in M$이 $M/IM$을 생성하고 $M$이 유한이면,
$x_1, \ldots, x_n$이 $R_f$ 위에서 $M_f$를 생성하는
$f \in 1 + I$가 존재한다.
\item $x_1, \ldots, x_n \in M$이 $M/IM$을 생성하고 $M$이 유한이며
$I \subset \text{rad}(R)$이면, $M$은 $x_1, \ldots, x_n$으로 생성된다.
\item $IM = M$이고 $I$가 멱영 이데알이면 $M = 0$이다.
\item $N, N' \subset M$이고 $M = N + IN'$이며 $I$가 멱영 이데알이면
$M = N$이다.
\item $N \to M$이 가군 사상이고 $I$가 멱영 이데알이며
$N/IN \to M/IM$이 전사이면 $N \to M$은 전사이다.
\item $\{x_\alpha\}_{\alpha \in A}$가 $M/IM$을 생성하는 $M$의
원소들의 집합이고 $I$가 멱영 이데알이면, $M$은 $x_\alpha$들로
생성된다.
\end{enumerate}
\end{lemma}

\begin{proof}
(\ref{algebra-item-nakayama})의 증명. $R$ 위에서 $M$의 생성원
$y_1, \ldots, y_m$을 택한다. 각 $i$에 대해 $M = IM$이므로
$z_{ij} \in I$인 $y_i = \sum z_{ij} y_j$로 쓸 수 있다.
달리 말해 $\sum_j (\delta_{ij} - z_{ij})y_j = 0$이다.
$m \times m$ 행렬 $A = (\delta_{ij} - z_{ij})$의 행렬식을 $f$라 하자.
$f \in 1 + I$임에 유의하라(행렬 $A$는 $I$을 법으로 성분별로
$m \times m$ 단위행렬과 합동이기 때문이다).
Lemma \ref{algebra-lemma-matrix-left-inverse} (1)에 의해
$BA = f 1_{m \times m}$인 $m \times m$ 행렬 $B$가 존재한다.
성분별로 쓰면 모든 $h$와 $j$에 대해
$\sum_{i} b_{hi} a_{ij} = f \delta_{hj}$이다. 따라서 모든 $h$에 대해
$\sum_{i, j} b_{hi} a_{ij} y_j
= \sum_{j} f \delta_{hj} y_j = f y_h$이다. 달리 말해 모든 $h$에 대해
$0 = f y_h$이다(각 $i$가 $\sum_j a_{ij} y_j = 0$을 만족하기 때문이다).
이는 $f$가 $M$을 영멸함을 뜻한다.

\medskip\noindent
Lemma \ref{algebra-lemma-contained-in-radical}에 의해
% P01-E0294: 동결 원문은 invertible element 앞의 관사를 빠뜨린다.
$1 + \text{rad}(R)$의 원소는 $R$의 가역원이다. 따라서
(\ref{algebra-item-nakayama})가 (2)를 함의한다. $N'$이 유한이므로 유한인
$M/N$에 (1)을 적용하여 (3)을 얻는다. $N'$이 유한이므로 유한인
$M/N$에 (2)를 적용하여 (4)를 얻는다. $M$과 부분가군
$\Im(N \to M)$ 및 $M$에 (3)을 적용하여 (5)를 얻는다. $M$과 부분가군
$\Im(N \to M)$ 및 $M$에 (4)를 적용하여 (6)을 얻는다.
사상 $R^{\oplus n} \to M$,
$(a_1, \ldots, a_n) \mapsto a_1x_1 + \ldots + a_nx_n$에 (5)를 적용하여
(7)을 얻는다. 사상 $R^{\oplus n} \to M$,
$(a_1, \ldots, a_n) \mapsto a_1x_1 + \ldots + a_nx_n$에 (6)을 적용하여
(8)을 얻는다.

\medskip\noindent
(9)가 성립하는 것은 $M = IM$이면 모든 $n \geq 0$에 대해
$M = I^nM$이고, $I$가 멱영 이데알이라는 것은 어떤 $n \gg 0$에 대해
$I^n = 0$이라는 뜻이기 때문이다. (10), (11), (12)는 위에서 쓴
논증으로 (9)에서 따른다.
\end{proof}

\begin{lemma}
\label{algebra-lemma-NAK-localization}
$R$을 환, $S \subset R$을 곱셈적 부분집합, $I \subset R$를 이데알,
$M$을 유한 $R$-가군이라 하자. $x_1, \ldots, x_r \in M$이
$S^{-1}(R/I)$-가군으로서 $S^{-1}(M/IM)$을 생성하면,
$x_1, \ldots, x_r$이 $R_f$-가군으로서 $M_f$를 생성하는
$f \in S + I$가 존재한다.\footnote{특수한 경우: (I) $I = 0$.
보조정리는 $x_1, \ldots, x_r$이 $S^{-1}M$을 생성하면 어떤
$f \in S$에 대해 $x_1, \ldots, x_r$이 $M_f$를 생성한다는 뜻이다.
(II) $I = \mathfrak p$가 소 이데알이고 $S = R \setminus \mathfrak p$이다.
보조정리는 $x_1, \ldots, x_r$이 $M \otimes_R \kappa(\mathfrak p)$를
생성하면 $f \in R$, $f \not \in \mathfrak p$인 어떤 원소에 대해
$x_1, \ldots, x_r$이 $M_f$를 생성한다는 뜻이다.}
\end{lemma}

\begin{proof}
특수한 경우 $I = 0$. $R$ 위에서 $M$의 생성원을
$y_1, \ldots, y_s$라 하자. $S^{-1}M$은 $x_1, \ldots, x_r$으로
생성되므로 각 $i$에 대해 어떤 $a_{ij} \in R$와 $s_{ij} \in S$를 써서
$S^{-1}M$에서 $y_i = \sum (a_{ij}/s_{ij})x_j$로 쓸 수 있다.
$s_{ij}$들의 곱 $s \in S$를 곱하면 어떤 $a'_{ij} \in R$에 대해
$S^{-1}M$에서 $sy_i = \sum a'_{ij}x_j$임을 알 수 있다.
이는 다시 $M$에서 $t_isy_i = \sum t_ia'_{ij}x_j$인
$t_i \in S$가 존재한다는 뜻이다. 따라서 $t \in S$가 $t_i$들의
곱이면, $M_{st}$에서 $y_i$는 $x_1, \ldots, x_r$이 생성하는
$R_{st}$-부분가군에 속한다. 그러므로
% P01-E0295: 동결 원문은 복수 주어에 단수 동사 generates를 쓴다.
$x_1, \ldots, x_r$은 $M_{st}$를 생성한다.

\medskip\noindent
일반적인 경우. 특수한 경우에 의해 $x_1, \ldots, x_r$이
$(R/I)_s$ 위에서 $(M/IM)_s$를 생성하는 $s \in S$를 찾을 수 있다.
Lemma \ref{algebra-lemma-NAK}에 의해 $x_1, \ldots, x_r$이 $(R_s)_g$ 위에서
$(M_s)_g$를 생성하는 $g \in 1 + I_s \subset R_s$를 찾을 수 있다.
$g = 1 + i/s'$라고 쓰면 $f = ss' + is$가 원하는 원소이다.
세부사항은 생략한다.
\end{proof}

\begin{lemma}
\label{algebra-lemma-when-surjective-local}
$A \to B$를 국소환들의 국소 준동형이라 하자. 다음을 가정하자.
\begin{enumerate}
\item $B$는 $A$-가군으로서 유한이다.
\item $\mathfrak m_B$는 유한 생성 이데알이다.
\item $A \to B$는 잉여체 위의 동형사상을 유도한다.
\item $\mathfrak m_A/\mathfrak m_A^2 \to \mathfrak m_B/\mathfrak m_B^2$는
전사이다.
\end{enumerate}
그러면 $A \to B$는 전사이다.
\end{lemma}

\begin{proof}
$A \to B$가 전사임을 보이기 위해 이를 $A$-가군의 사상으로 보고
Lemma \ref{algebra-lemma-NAK} (6)을 적용한다. 그러면
$A/\mathfrak m_A \to B/\mathfrak m_AB$가 전사임을 보이면 충분하다.
$A/\mathfrak m_A = B/\mathfrak m_B$이므로
$\mathfrak m_AB \to \mathfrak m_B$가 전사임을 보이면 충분하다.
$\mathfrak m_AB \to \mathfrak m_B$를 $B$-가군의 사상으로 보고
Lemma \ref{algebra-lemma-NAK} (6)을 적용한다. 그러면
$\mathfrak m_AB/\mathfrak m_A\mathfrak m_B \to \mathfrak m_B/\mathfrak m_B^2$가
전사임을 보이면 충분하다. 이는 가정 (4)에서 따른다.
\end{proof}

\section{스펙트럼의 열린 부분집합과 닫힌 부분집합}
\label{algebra-section-open-and-closed}

\noindent
스펙트럼에서 동시에 열린집합이고 닫힌집합인 부분집합은 환의
멱등원에 대응한다.

\begin{lemma}
\label{algebra-lemma-idempotent-spec}
$R$을 환이라 하고 $e \in R$를 멱등원이라 하자. 이 경우
$$
\Spec(R) = D(e) \amalg D(1-e).
$$
\end{lemma}

\begin{proof}
정역의 멱등원 $e$는 $1$ 또는 $0$임에 유의하라. 따라서
\begin{eqnarray*}
D(e)
& = &
\{ \mathfrak p \in \Spec(R)
\mid
e \not\in \mathfrak p \} \\
& = &
\{ \mathfrak p \in \Spec(R)
\mid
e \not = 0\text{ in }\kappa(\mathfrak p) \} \\
& = &
\{ \mathfrak p \in \Spec(R)
\mid
e = 1\text{ in }\kappa(\mathfrak p) \}
\end{eqnarray*}
이다. 마찬가지로
\begin{eqnarray*}
D(1-e)
& = &
\{ \mathfrak p \in \Spec(R)
\mid
1 - e \not\in \mathfrak p \} \\
& = &
\{ \mathfrak p \in \Spec(R)
\mid
e \not = 1\text{ in }\kappa(\mathfrak p) \} \\
& = &
\{ \mathfrak p \in \Spec(R)
\mid
e = 0\text{ in }\kappa(\mathfrak p) \}
\end{eqnarray*}
이다. 임의의 잉여체에서 $e$의 상은 $1$ 또는 $0$이므로
$D(e)$와 $D(1-e)$는 $\Spec(R)$ 전체를 덮는다.
\end{proof}

\begin{lemma}
\label{algebra-lemma-spec-product}
$R_1$과 $R_2$를 환이라 하고 $R = R_1 \times R_2$라 하자.
사상 $R \to R_1$, $(x, y) \mapsto x$와 $R \to R_2$,
$(x, y) \mapsto y$는 연속사상 $\Spec(R_1) \to \Spec(R)$와
$\Spec(R_2) \to \Spec(R)$를 유도한다. 유도된 사상
$$
\Spec(R_1) \amalg \Spec(R_2)
\longrightarrow
\Spec(R)
$$
은 위상동형이다. 달리 말해 $R = R_1\times R_2$의 스펙트럼은
$R_1$의 스펙트럼과 $R_2$의 스펙트럼의 서로소 합집합이다.
\end{lemma}

\begin{proof}
$e_1 = (1, 0)$이고 $e_2 = (0, 1)$일 때 $1 = e_1 + e_2$라고 쓰자.
$e_1$과 $e_2 = 1 - e_1$은 멱등원임에 유의하라.
$R_1 = R_{e_1}$이 $R$을 $e_1$에서 국소화한 것임을 보이는 일은
독자에게 맡긴다. $e_2$에 대해서도 마찬가지이다. 따라서 보조정리의
명제는 Lemma \ref{algebra-lemma-idempotent-spec}과
Lemma \ref{algebra-lemma-standard-open}에서 따른다.
\end{proof}

\noindent
함수의 붙이기 보조정리를 도입한 뒤에 다음 보조정리를 다시 증명한다.
Section \ref{algebra-section-tilde-module-sheaf}를 보라.

\begin{lemma}
\label{algebra-lemma-disjoint-decomposition}
$R$을 환이라 하자. 동시에 열린집합이고 닫힌집합인 각
$U \subset \Spec(R)$에 대해 $U = D(e)$를 만족하는 유일한 멱등원
$e \in R$가 존재한다. 이로부터 동시에 열린집합이고 닫힌집합인
부분집합 $U \subset \Spec(R)$와 멱등원 $e \in R$ 사이의 일대일
대응을 얻는다.
\end{lemma}

\begin{proof}
$U \subset \Spec(R)$가 열린집합이면서 닫힌집합이라고 하자.
$U$가 닫힌집합이므로 Lemma \ref{algebra-lemma-quasi-compact}에 의해
준콤팩트이고, 그 여집합도 마찬가지이다. 이를 표준 열린집합들의
유한 합집합 $U = \bigcup_{i = 1}^n D(f_i)$로 쓰자. 마찬가지로
$\Spec(R) \setminus U = \bigcup_{j = 1}^m D(g_j)$를 표준
열린집합들의 유한 합집합으로 쓰자.
$\emptyset = D(f_i) \cap D(g_j) = D(f_i g_j)$이므로
Lemma \ref{algebra-lemma-Zariski-topology}에 의해 $f_i g_j$는 멱영원이다.
$I = (f_1, \ldots, f_n) \subset R$ 및
$J = (g_1, \ldots, g_m) \subset R$로 두자. $V(J)$는 $U$와 같고,
$V(I)$는 $U$의 여집합과 같으므로
$\Spec(R) = V(I) \amalg V(J)$이다. 위에서 언급한 멱영성에 의해 충분히
큰 정수 $N$에 대해 $(IJ)^N = (0)$이다.
$\bigcup D(f_i) \cup \bigcup D(g_j) = \Spec(R)$이므로
Lemma \ref{algebra-lemma-Zariski-topology}에 의해 $I + J = R$이다.
이 등식을 $2N$제곱하면 $I^N + J^N = R$을 얻는다.
$x \in I^N$이고 $y \in J^N$일 때 $1 = x + y$라고 쓰자.
$I^N J^N = (0)$이므로 $0 = xy = x(1 - x)$이다. 따라서
$x = x^2$는 멱등원이고 $I^N \subset I$에 속한다. 멱등원
$y = 1 - x$는 $J^N \subset J$에 속한다. 이는 모든
$\mathfrak p \in V(J)$에 대해 멱등원 $x$가 잉여체
$\kappa(\mathfrak p)$에서 $1$로 가고, 모든 $\mathfrak p \in V(I)$에
대해 $\kappa(\mathfrak p)$에서 $x$가 $0$으로 감을 보여 준다.

\medskip\noindent
유일성을 보이기 위해 $e_1, e_2$가 $R$의 서로 다른 멱등원이라고 하자.
$e_1 \in \mathfrak p$이고 $e_2 \not \in \mathfrak p$이거나 그 반대인
소 이데알 $\mathfrak p$가 존재함을 보여야 한다.
$e_i' = 1 - e_i$라고 쓰자. $e_1 \not = e_2$이면
$0 \not = e_1 - e_2  = e_1(e_2 + e_2') - (e_1 + e_1')e_2
= e_1 e_2' - e_1' e_2$이다. 따라서 멱등원
$e_1 e_2' \not = 0$ 또는 $e_1' e_2 \not = 0$이다.
% P01-E0296: 동결 원문은 영이 아닌이라는 필수 조건 없이 멱등원은 멱영원이 아니라고 한다.
멱등원은 멱영원이 아니므로 Lemma \ref{algebra-lemma-Zariski-topology}에 의해
$e_1e_2' \not \in \mathfrak p$ 또는 $e_1'e_2 \not \in \mathfrak p$인
소 이데알 $\mathfrak p$를 찾을 수 있다. 이것이 원하는 소 이데알을
준다는 것은 쉽게 알 수 있다.
\end{proof}

\begin{lemma}
\label{algebra-lemma-characterize-spec-connected}
$R$을 영환이 아닌 환이라 하자. $\Spec(R)$가 연결인 것과 $R$에
자명하지 않은 멱등원이 없는 것은 동치이다.
\end{lemma}

\begin{proof}
Lemma \ref{algebra-lemma-disjoint-decomposition}과 연결 위상공간의 정의에서
명백히 따른다.
\end{proof}

\begin{lemma}
\label{algebra-lemma-ideal-is-squared-union-connected}
$R$의 유한 생성 이데알 $I \subset R$가 $I = I^2$을 만족한다고 하자.
그러면 다음이 성립한다.
\begin{enumerate}
\item $I = (e)$인 멱등원 $e \in R$가 존재한다.
\item 멱등원 $e' = 1 - e \in R$에 대해 $R/I \cong R_{e'}$이다.
\item $V(I)$는 $\Spec(R)$에서 열린집합이면서 닫힌집합이다.
\end{enumerate}
\end{lemma}

\begin{proof}
나카야마의 Lemma \ref{algebra-lemma-NAK}에 의해 $fI = 0$인 원소
$f = 1 + i$, $i \in I$가 존재한다. 그러면 $f^2 = f + fi = f$이므로
멱등원이다. 멱등원 $e = 1 - f = -i \in I$를 생각하자.
$j \in I$에 대해 $ej = j - fj = j$이므로 $I = (e)$이다.
이로써 (1)이 증명되었다.

\medskip\noindent
(2)와 (3)은 (1)에서 따른다. 실제로
$V(I) = V(e) = \Spec(R) \setminus D(e)$이고, 이는
Lemma \ref{algebra-lemma-idempotent-spec} 또는
Lemma \ref{algebra-lemma-disjoint-decomposition}에 의해 열린집합이면서
닫힌집합이다. 이로써 (3)이 증명되었다. (2)를 위해 사상
$R \to R_{e'}$가 전사임에 유의하라. 실제로 $R_{e'}$에서
$x/(e')^n = x/e' = xe'/(e')^2 = xe'/e' = x/1$이다.
사상 $R \to R_{e'}$의 핵은 $e'$의 양의 거듭제곱에 의해 영멸되는
$R$의 원소들의 집합이다. $e'$은 멱등원이므로 이는 $e'$에 의해
영멸되는 원소들의 이데알이고, $e + e' = 1$이 서로 직교하는
두 멱등원이므로 이 이데알은 $I = (e)$이다. 이로써 (2)가 증명되었다.
\end{proof}

\section{스펙트럼의 연결 성분}
\label{algebra-section-connected-components}

\noindent
스펙트럼의 연결 성분은 처음 생각하는 것만큼 이해하기 쉽지 않다.
우리는 국소 연결 공간의 위상에 익숙하지만, 일반적으로 환의 스펙트럼은
국소 연결이 아니기 때문이다.

\begin{lemma}
\label{algebra-lemma-closed-union-connected-components}
$R$을 환이라 하고 $T \subset \Spec(R)$를 스펙트럼의 부분집합이라 하자.
다음 조건들은 서로 동치이다.
\begin{enumerate}
\item $T$는 닫힌집합이며 $\Spec(R)$의 연결 성분들의 합집합이다.
\item $T$는 $\Spec(R)$에서 동시에 열린집합이고 닫힌집합인
부분집합들의 교집합이다.
\item $T = V(I)$이고, 여기서 $I \subset R$는 멱등원들로 생성되는
이데알이다.
\end{enumerate}
또한 (3)의 이데알이 존재하면 유일하다.
\end{lemma}

\begin{proof}
Lemma \ref{algebra-lemma-topology-spec}과
Topology, Lemma \ref{topology-lemma-closed-union-connected-components}에
의해 (1)과 (2)는 동치이다. (2)를 가정하고
$T = \bigcap U_\alpha$라고 쓰자. 여기서
$U_\alpha \subset \Spec(R)$는 열린집합이면서 닫힌집합이다.
Lemma \ref{algebra-lemma-disjoint-decomposition}에 의해 어떤 멱등원
$e_\alpha \in R$에 대해 $U_\alpha = D(e_\alpha)$이다.
이제 $I = (1 - e_\alpha)$로 두면 $T = V(I)$이므로 (3)이 성립한다.
마지막으로 (3)을 가정하자. 어떤 멱등원들의 모음 $e_\alpha$에 대해
$T = V(I)$ 및 $I = (e_\alpha)$라고 쓰자. 그러면
$T = \bigcap V(e_\alpha) = \bigcap D(1 - e_\alpha)$임이 명백하다.

\medskip\noindent
$I$가 멱등원들로 생성되는 이데알이라고 하자.
$V(I) \subset V(e)$인 멱등원 $e \in R$를 생각하자.
Lemma \ref{algebra-lemma-Zariski-topology}에 의해 어떤 $n \geq 1$에 대해
$e^n \in I$이다. $e$가 멱등원이므로 이는 $e \in I$라는 뜻이다.
따라서 $I$는 $T \subset V(e)$를 만족하는 멱등원 $e$들로 정확히
생성된다. 달리 말해 이데알 $I$는 닫힌 부분집합 $T$에 의해 완전히
결정되며, 이로써 유일성이 증명된다.
\end{proof}

\begin{lemma}
\label{algebra-lemma-connected-component}
$R$을 환이라 하자. $\Spec(R)$의 연결 성분은 $V(I)$ 꼴이다.
여기서 $I$는 $R$의 모든 멱등원이 $R/I$에서 $0$ 또는 $1$로 가도록
하는, 멱등원들로 생성되는 이데알이다.
\end{lemma}

\begin{proof}
$\mathfrak p$를 $R$의 소 이데알이라 하자.
% P01-E0297: 동결 원문은 have see라는 동사 오기를 포함한다.
Lemma \ref{algebra-lemma-topology-spec}에 의해 위상공간 $\Spec(R)$에 대해
Topology, Lemma \ref{topology-lemma-connected-component-intersection}의
가정들이 성립함을 알 수 있다. 따라서 $\Spec(R)$에서 $\mathfrak p$의
연결 성분은 $\mathfrak p$를 포함하면서 동시에 열린집합이고
닫힌집합인 $\Spec(R)$의 부분집합들의 교집합이다. 그러므로 이는
$V(I)$와 같다. 여기서 $I$는 멱등원 $e \in R$들 중
$\kappa(\mathfrak p)$에서 $e$가 $0$으로 가는 것들로 생성된다. Lemma
\ref{algebra-lemma-disjoint-decomposition}을 보라. 이 모음에 속하지 않는
임의의 멱등원 $e$는 $R/I$에서 명백히 $1$로 간다.
\end{proof}

\section{붙이기 성질}
\label{algebra-section-more-glueing}

\noindent
이 절에서는 다음 꼴의 몇 가지 표준 결과를 모은다. 어떤 성질이
표준 열린 덮개의 모든 원소에 대해 참이면 그 성질은 전체에서도
참이다. 실제로 다음 보조정리처럼 국소환의 수준에서 확인하는 것만으로
충분한 경우가 많다.

\begin{lemma}
\label{algebra-lemma-characterize-zero-local}
$R$을 환이라 하자.
\begin{enumerate}
\item $R$-가군 $M$의 원소 $x$에 대해 다음 조건들은 서로 동치이다.
\begin{enumerate}
\item $x = 0$이다.
\item 모든 $\mathfrak p \in \Spec(R)$에 대해 $x$는 $M_\mathfrak p$에서
영으로 간다.
\item $R$의 모든 극대 이데알 $\mathfrak m$에 대해 $x$는
$M_{\mathfrak m}$에서 영으로 간다.
\end{enumerate}
달리 말해 사상 $M \to \prod_{\mathfrak m} M_{\mathfrak m}$은 단사이다.
\item $R$-가군 $M$에 대해 다음 조건들은 서로 동치이다.
\begin{enumerate}
\item $M$은 영가군이다.
\item 모든 $\mathfrak p \in \Spec(R)$에 대해 $M_{\mathfrak p}$는
영가군이다.
\item $R$의 모든 극대 이데알 $\mathfrak m$에 대해
$M_{\mathfrak m}$는 영가군이다.
\end{enumerate}
\item $R$-가군들의 복합체 $M_1 \to M_2 \to M_3$에 대해
다음 조건들은 서로 동치이다.
\begin{enumerate}
\item $M_1 \to M_2 \to M_3$는 완전하다.
\item $R$의 모든 소 이데알 $\mathfrak p$에 대해 국소화
$M_{1, \mathfrak p} \to M_{2, \mathfrak p} \to M_{3, \mathfrak p}$는
완전하다.
\item $R$의 모든 극대 이데알 $\mathfrak m$에 대해 국소화
$M_{1, \mathfrak m} \to M_{2, \mathfrak m} \to M_{3, \mathfrak m}$는
완전하다.
\end{enumerate}
\item $R$-가군들의 사상 $f : M \to M'$에 대해 다음 조건들은
서로 동치이다.
\begin{enumerate}
\item $f$는 단사이다.
\item $R$의 모든 소 이데알 $\mathfrak p$에 대해
$f_{\mathfrak p} : M_\mathfrak p \to M'_\mathfrak p$는 단사이다.
\item $R$의 모든 극대 이데알 $\mathfrak m$에 대해
$f_{\mathfrak m} : M_\mathfrak m \to M'_\mathfrak m$는 단사이다.
\end{enumerate}
\item $R$-가군들의 사상 $f : M \to M'$에 대해 다음 조건들은
서로 동치이다.
\begin{enumerate}
\item $f$는 전사이다.
\item $R$의 모든 소 이데알 $\mathfrak p$에 대해
$f_{\mathfrak p} : M_\mathfrak p \to M'_\mathfrak p$는 전사이다.
\item $R$의 모든 극대 이데알 $\mathfrak m$에 대해
$f_{\mathfrak m} : M_\mathfrak m \to M'_\mathfrak m$는 전사이다.
\end{enumerate}
\item $R$-가군들의 사상 $f : M \to M'$에 대해 다음 조건들은
서로 동치이다.
\begin{enumerate}
\item $f$는 전단사이다.
\item $R$의 모든 소 이데알 $\mathfrak p$에 대해
$f_{\mathfrak p} : M_\mathfrak p \to M'_\mathfrak p$는 전단사이다.
\item $R$의 모든 극대 이데알 $\mathfrak m$에 대해
$f_{\mathfrak m} : M_\mathfrak m \to M'_\mathfrak m$는 전단사이다.
\end{enumerate}
\end{enumerate}
\end{lemma}

\begin{proof}
(1)에서와 같이 $x \in M$이라 하자.
$I = \{f \in R \mid fx = 0\}$로 두자. $I$가 이데알임은 쉽게 알 수
있다(이는 $x$의 영멸 이데알이다). 조건 (1)(c)는 모든 극대 이데알
$\mathfrak m$에 대해 $fx =0$인 $f \in R \setminus \mathfrak m$이
존재한다는 뜻이다. 달리 말해 $V(I)$는 닫힌점을 포함하지 않는다.
Lemma \ref{algebra-lemma-Zariski-topology}에 의해 $I$는 단위 이데알이다.
따라서 $x$는 영이고, 즉 (1)(a)가 성립한다. 이로써 (1)이 증명되었다.

\medskip\noindent
$M$의 모든 원소에 (1)을 동시에 적용하면 (2)가 따른다.

\medskip\noindent
(3)의 증명. 수열의 호몰로지를 $H$, 즉
$H = \Ker(M_2 \to M_3)/\Im(M_1 \to M_2)$라 하자.
Proposition \ref{algebra-proposition-localization-exact}에 의해
$H_\mathfrak p$는 수열
$M_{1, \mathfrak p} \to M_{2, \mathfrak p} \to M_{3, \mathfrak p}$의
호몰로지이다. 따라서 (3)은 (2)의 결과이다.

\medskip\noindent
(4)와 (5)는 (3)의 특수한 경우이다. (6)은 (4)와 (5)를 결합하면
형식적으로 따른다.
\end{proof}

\begin{lemma}
\label{algebra-lemma-cover}
\begin{slogan}
가군과 대수의 자리스키 국소적 성질
\end{slogan}
$R$을 환, $M$을 $R$-가군, $S$를 $R$-대수라 하자.
$f_1, \ldots, f_n$이 $\bigcup D(f_i) = \Spec(R)$를 만족하는
$R$의 원소들의 유한 목록, 달리 말해 $(f_1, \ldots, f_n) = R$이라고
하자.
\begin{enumerate}
\item 각 $M_{f_i} = 0$이면 $M = 0$이다.
\item 각 $M_{f_i}$가 유한 $R_{f_i}$-가군이면 $M$은 유한
$R$-가군이다.
\item 각 $M_{f_i}$가 유한 표시 $R_{f_i}$-가군이면 $M$은 유한 표시
$R$-가군이다.
\item $M \to N$을 $R$-가군들의 사상이라 하자. 각 $i$에 대해
$M_{f_i} \to N_{f_i}$가 동형사상이면 $M \to N$은 동형사상이다.
\item $0 \to M'' \to M \to M' \to 0$을 $R$-가군들의 복합체라 하자.
각 $i$에 대해 $0 \to M''_{f_i} \to M_{f_i} \to M'_{f_i} \to 0$이
완전하면 $0 \to M'' \to M \to M' \to 0$은 완전하다.
\item 각 $R_{f_i}$가 뇌터환이면 $R$은 뇌터환이다.
\item 각 $S_{f_i}$가 유한형 $R_{f_i}$-대수이면 $S$는 유한형
$R$-대수이다.
\item 각 $S_{f_i}$가 $R_{f_i}$ 위에서 유한 표시이면 $S$는 유한 표시
$R$-대수이다.
\end{enumerate}
\end{lemma}

\begin{proof}
각 부분을 차례로 증명한다.
\begin{enumerate}
\item Proposition \ref{algebra-proposition-localize-twice}에 의해 가정은 모든
$\mathfrak p \in \Spec(R)$에 대해 $M_\mathfrak p = 0$임을 함의하므로,
Lemma \ref{algebra-lemma-characterize-zero-local}에 의해 결론을 얻는다.
\item 각 $i$에 대해 $M_{f_i}$의 유한 생성집합 $X_i$를 택한다.
일반성을 잃지 않고 $X_i$의 원소들이 국소화 사상
$M \rightarrow M_{f_i}$의 상에 있다고 할 수 있으므로, $M$에서
$X_i$의 원소들의 원상들로 이루어진 유한 집합 $Y_i$를 택한다.
이 집합들의 합집합을 $Y$라 하자. 이 역시 유한 집합이다.
기저 원소 $e_y$를 $y$로 보내는 명백한 $R$-선형사상
$R^Y \rightarrow M$을 생각하자. 가정에 의해 이 사상은 $R$의 임의의
소 이데알 $\mathfrak p$에서 국소화한 뒤 전사이므로
Lemma \ref{algebra-lemma-characterize-zero-local}에 의해 전사이고,
$M$은 유한 생성이다.
\item (2)에 의해 짧은 완전열을 얻는다.
% P01-E0298: 동결 원문은 덮개의 크기 n을 자유가군의 랭크로도 재사용한다.
$$
0 \rightarrow K \rightarrow R^n \rightarrow M \rightarrow 0
$$
국소화는 완전 함자이고 $M_{f_i}$는 유한 표시이므로
Lemma \ref{algebra-lemma-extension}에 의해 모든 $1 \leq i \leq n$에 대해
$K_{f_i}$가 유한 생성임을 알 수 있다. (2)에 의해 $K$는 유한
$R$-가군이고, 따라서 $M$은 유한 표시이다.
\item Proposition \ref{algebra-proposition-localize-twice}에 의해 가정은 모든
소 이데알에서의 국소화에 유도된 준동형이 동형사상임을 함의하므로,
Lemma \ref{algebra-lemma-characterize-zero-local}에 의해 결론을 얻는다.
\item Proposition \ref{algebra-proposition-localize-twice}에 의해 가정은 모든
소 이데알에서 국소화한 수열이 짧은 완전열임을 함의하므로,
Lemma \ref{algebra-lemma-characterize-zero-local}에 의해 결론을 얻는다.
\item $R$의 모든 이데알에 유한 생성집합이 있음을 보이겠다.
이를 위해 $I \subset R$를 임의의 이데알이라 하자.
Proposition \ref{algebra-proposition-localization-exact}에 의해 각
$I_{f_i} \subset R_{f_i}$는 이데알이다. 가정에 의해 이들은 모두
유한 생성이므로 (2)에 의해 결론을 얻는다.
\item 각 $i$에 대해 $S_{f_i}$의 유한 생성집합 $X_i$를 택한다.
일반성을 잃지 않고 $X_i$의 원소들이 국소화 사상
$S \rightarrow S_{f_i}$의 상에 있다고 할 수 있으므로, $S$에서
$X_i$의 원소들의 원상들로 이루어진 유한 집합 $Y_i$를 택한다.
이 집합들의 합집합을 $Y$라 하자. 이 역시 유한 집합이다.
$Y$가 유도하는 대수 준동형 $R[X_y]_{y \in Y} \rightarrow S$를
생각하자. 이는 대수 준동형이므로 그 상 $T$는 $R$-가군 $S$의
$R$-부분가군이고, 따라서 몫가군 $S/T$를 생각할 수 있다.
가정에 의해 이는 $f_i$들에서 국소화하면 영가군이므로 (1)에 의해
영가군이다. 따라서 $S$는 유한형 $R$-대수이다.
\item 앞 항목에 의해 전사 $R$-대수 준동형
$R[X_1, \ldots, X_n] \rightarrow S$가 존재한다. 그 핵을 $K$라 하자.
이는 $R[X_1, \ldots, X_n]$의 이데알이고, 각 $f_i$에서의 국소화에서
유한 생성이다. $f_i$들은 $R$에서 단위 이데알을 생성하므로
$R[X_1, \ldots, X_n]$에서도 단위 이데알을 생성한다. 이제 (2)를
적용하면 증명이 끝난다.
\end{enumerate}
\end{proof}

\begin{lemma}
\label{algebra-lemma-cover-upstairs}
$R \to S$를 환 사상이라 하자. $g_1, \ldots, g_n$이
$\bigcup D(g_i) = \Spec(S)$를 만족하는 $S$의 원소들의 유한 목록,
달리 말해 $(g_1, \ldots, g_n) = S$라고 하자.
\begin{enumerate}
\item 각 $S_{g_i}$가 $R$ 위에서 유한형이면 $S$는 $R$ 위에서
유한형이다.
\item 각 $S_{g_i}$가 $R$ 위에서 유한 표시이면 $S$는 $R$ 위에서
유한 표시이다.
\end{enumerate}
\end{lemma}

\begin{proof}
$\sum h_i g_i = 1$인 $h_1, \ldots, h_n \in S$를 택한다.

\medskip\noindent
(1)의 증명. 각 $i$에 대해 $R$-대수로서 $S_{g_i}$를 생성하는
원소들의 유한 목록 $x_{i, j} \in S_{g_i}$,
$j = 1, \ldots, m_i$를 택한다. 어떤 $y_{i, j} \in S$와
$n_{i, j} \ge 0$에 대해
$x_{i, j} = y_{i, j}/g_i^{n_{i, j}}$라고 쓰자.
$g_1, \ldots, g_n$, $h_1, \ldots, h_n$ 및
$y_{i, j}$, $i = 1, \ldots, n$, $j = 1, \ldots, m_i$로 생성되는
$R$-부분대수 $S' \subset S$를 생각하자.
국소화는 완전하므로(Proposition \ref{algebra-proposition-localization-exact})
$S'_{g_i} \to S_{g_i}$는 단사이다. 한편 $y_{i, j}$를 택한 방식에
의해 전사이다. $h_1, \ldots, h_n \in S'$이므로 원소들
$g_1, \ldots, g_n$은 $S'$에서 단위 이데알을 생성한다. 따라서
$S' \to S$를 $S'$-가군 사상으로 보면 Lemma \ref{algebra-lemma-cover}에 의해
동형사상이다.

\medskip\noindent
(2)의 증명. 이미 $S$가 유한형임을 안다. 어떤 이데알 $J$에 대해
$S = R[x_1, \ldots, x_m]/J$라고 쓰자. 각 $i$에 대해 $g_i$의 원상
$g'_i \in R[x_1, \ldots, x_m]$를 택하고, $h_i$의 원상
$h'_i \in R[x_1, \ldots, x_m]$를 택한다. 그러면
$$
S_{g_i} = R[x_1, \ldots, x_m, y_i]/(J_i + (1 - y_ig'_i))
$$
이다. 여기서 $J_i$는 $J$가 생성하는
$R[x_1, \ldots, x_m, y_i]$의 이데알이다. 작은 세부사항은 생략한다.
Lemma \ref{algebra-lemma-finite-presentation-independent}에 의해 원소들의 유한
목록 $f_{i, j} \in J$, $j = 1, \ldots, m_i$를 택하여, $J_i$에서
$f_{i, j}$들의 상과 $1 - y_ig'_i$가 이데알
$J_i + (1 - y_ig'_i)$를 생성하게 할 수 있다. 다음과 같이 두자.
$$
S' = R[x_1, \ldots, x_m]/\left(\sum h'_ig'_i - 1, f_{i, j}; 
i = 1, \ldots, n, j = 1, \ldots, m_i\right)
$$
전사 $R$-대수 사상 $S' \to S$가 존재한다. $S'$에서 원소
$g'_1, \ldots, g'_n$들의 동치류는 단위 이데알을 생성하며, 구성에
의해 사상 $S'_{g'_i} \to S_{g_i}$들은 단사이다. 따라서 (1)에서처럼
결론을 얻는다.
\end{proof}

\section{함수 붙이기}
\label{algebra-section-tilde-module-sheaf}

\noindent
이 절에서는 표준 열린집합들에 의한 열린 덮개
$$
\Spec(R) = \bigcup\nolimits_{i = 1}^n D(f_i)
$$
가 주어지고, 각 $i$에 대해 $h_i \in R_{f_i}$가 주어져
$R_{f_i f_j}$의 원소로서 $h_i = h_j$를 만족하면, $h \in R$가 유일하게
존재하여 $h$의 $R_{f_i}$에서의 상이 $h_i$가 됨을 보인다. 이 결과는 두 가지로
해석할 수 있다.
\begin{enumerate}
\item 규칙 $D(f) \mapsto R_f$는 표준 열린집합들 위의 환의 층이다.
Sheaves, Section \ref{sheaves-section-bases}를 보라.
\item $R_f$의 원소를 $D(f)$ 위의 ``대수적'' 또는 ``정칙'' 함수로
생각하면, 이들은 위상다양체의 연속함수 또는 미분가능 다양체의
미분가능 함수와 같은 방식으로 붙여진다.
\end{enumerate}

\begin{lemma}
\label{algebra-lemma-cover-module}
$R$을 환이라 하자. $f_1, \ldots, f_n$을 단위 이데알을 생성하는
$R$의 원소들이라 하고 $M$을 $R$-가군이라 하자. 수열
$$
0 \to
M \xrightarrow{\alpha}
\bigoplus\nolimits_{i = 1}^n M_{f_i} \xrightarrow{\beta}
\bigoplus\nolimits_{i, j = 1}^n M_{f_i f_j}
$$
은 완전하다. 여기서
$\alpha(m) = (m/1, \ldots, m/1)$이고
$\beta(m_1/f_1^{e_1}, \ldots, m_n/f_n^{e_n})
= (m_i/f_i^{e_i} - m_j/f_j^{e_j})_{(i, j)}$이다.
\end{lemma}

\begin{proof}
임의의 극대 이데알 $\mathfrak m$에서 이 수열을 국소화한 것이
완전함을 보이면 충분하다. Lemma \ref{algebra-lemma-characterize-zero-local}을
보라. $f_1, \ldots, f_n$은 단위 이데알을 생성하므로
$f_i \not \in \mathfrak m$인 $i$가 존재한다. 지표를 다시 매겨
$i = 1$이라고 할 수 있다.
Proposition \ref{algebra-proposition-localize-twice-module}에 의해
$(M_{f_i})_\mathfrak m = (M_\mathfrak m)_{f_i}$이고
$(M_{f_if_j})_\mathfrak m = (M_\mathfrak m)_{f_if_j}$이다.
특히 $f_1$이 단위원이므로
$(M_{f_1})_\mathfrak m = M_\mathfrak m$이고
$(M_{f_1 f_i})_\mathfrak m = (M_\mathfrak m)_{f_i}$이다.
수열의 사상들은 Lemma \ref{algebra-lemma-universal-property-localization-module}과
$M$ 위의 항등사상에서 오는 표준 사상들임에 유의하라.
지금까지의 내용을 종합하여 $R$을 $R_\mathfrak m$으로, $M$을
$M_\mathfrak m$으로, $f_i$들을 $R_\mathfrak m$에서의 상으로,
$f_1$을 $1 \in R_\mathfrak m$으로 바꾸면 $f_1 = 1$인 경우로 환원된다.

\medskip\noindent
$f_1 = 1$이라고 하자. 이제 $\alpha$의 단사성은 자명하다.
$m = (m_i) \in \bigoplus_{i = 1}^n M_{f_i}$가 $\beta$의 핵에 속한다고
하자. 그러면 $m_1 \in M_{f_1} = M$이다. 또한 $\beta(m) = 0$이면
$m_1$과 $m_i$는 $M_{f_1f_i} = M_{f_i}$의 같은 원소로 간다.
따라서 $\alpha(m_1) = m$이고 증명이 끝난다.
\end{proof}

\begin{lemma}
\label{algebra-lemma-standard-covering}
$R$을 환이라 하고,
% P01-E0299: 동결 원문의 원소 목록에는 생략부호 뒤 쉼표가 빠져 있다.
$f_1, f_2, \ldots f_n\in R$이 $R$에서 단위 이데알을 생성한다고 하자.
그러면 다음 수열은 완전하다.
$$
0 \longrightarrow
R \longrightarrow
\bigoplus\nolimits_i R_{f_i} \longrightarrow
\bigoplus\nolimits_{i, j}R_{f_if_j}
$$
여기서 사상 $\alpha : R \longrightarrow \bigoplus_i R_{f_i}$와
$\beta : \bigoplus_i R_{f_i} \longrightarrow \bigoplus_{i, j} R_{f_if_j}$는
다음과 같이 정의된다.
$$
\alpha(x) = \left(\frac{x}{1}, \ldots, \frac{x}{1}\right)
\text{ and }
\beta\left(\frac{x_1}{f_1^{r_1}}, \ldots, \frac{x_n}{f_n^{r_n}}\right)
=
\left(\frac{x_i}{f_i^{r_i}}-\frac{x_j}{f_j^{r_j}}~\text{in}~R_{f_if_j}\right).
$$
\end{lemma}

\begin{proof}
Lemma \ref{algebra-lemma-cover-module}의 특수한 경우이다.
\end{proof}

\noindent
다음 내용은 위에서 이미 보았지만 편의를 위해 명시적으로 진술한다.

\begin{lemma}
\label{algebra-lemma-disjoint-implies-product}
$R$을 환이라 하자. $U$와 $V$가 모두 열린집합이고
$\Spec(R) = U \amalg V$이면, Lemma \ref{algebra-lemma-spec-product}의 사상들을
통해 $U \cong \Spec(R_1)$ 및 $V \cong \Spec(R_2)$이고
$R \cong R_1 \times R_2$이다. 또한 $R_1$과 $R_2$는 둘 다 환 $R$의
국소화인 동시에 몫이다.
\end{lemma}

\begin{proof}
Lemma \ref{algebra-lemma-disjoint-decomposition}에 의해 어떤 멱등원 $e$에 대해
$U = D(e)$이고 $V = D(1-e)$이다.
Lemma \ref{algebra-lemma-standard-covering}에 의해
$R \cong R_e \times R_{1 - e}$이다(명백히 $R_{e(1-e)} = 0$이므로
붙이기 조건은 자명하다. 물론 이 경우 곱 분해를 직접 증명하는 것도
자명하다). 이로써 보조정리가 따른다.
\end{proof}

\begin{lemma}
\label{algebra-lemma-when-injective-covering}
$R$을 환이라 하고 $f_1, \ldots, f_n \in R$이라 하자.
$M$을 $R$-가군이라 하자. $M \to \bigoplus M_{f_i}$가 단사인 것과
$$
M \longrightarrow \bigoplus\nolimits_{i = 1, \ldots, n} M, \quad
m \longmapsto (f_1m, \ldots, f_nm)
$$
가 단사인 것은 동치이다.
\end{lemma}

\begin{proof}
사상 $M \to \bigoplus M_{f_i}$가 단사인 것과 다음 조건은 동치이다.
$f_i^{e_i}m = 0$, $i = 1, \ldots, n$을 만족하는 모든
$m \in M$과 $e_1, \ldots, e_n \geq 1$에 대해 $m = 0$이다.
이 조건은 표시된 사상이 단사임을 명백히 함의한다. 역으로 표시된
사상이 단사이고, $f_i^{e_i}m = 0$, $i = 1, \ldots, n$을 만족하는
$m \in M$과 $e_1, \ldots, e_n \geq 1$이 주어졌다고 하자.
모든 $i$에 대해 $e_i = 1$이면 표시된 사상의 단사성에서 즉시
$m = 0$을 얻는다. 이제 $e = \sum e_i$에 관한 귀납법으로 이러한
모든 자료에 대해 성립함을 보이자. 기초 단계는 $e = n$이고 방금
다루었다. 어떤 $e_i > 1$이면 $m' = f_im$으로 두자. 귀납가정에 의해
$m' = 0$이다. 따라서 $f_i m = 0$이고, 즉 $e_i = 1$로 택하여 $e$를
줄일 수 있으므로 결론을 얻는다.
\end{proof}

\noindent
다음 보조정리는 평탄 내림이라는 더 일반적인 맥락에서 진술하고
증명하는 편이 낫다. 하지만 위의 내용과 잘 맞으므로 여기에 진술하는
것도 타당하다.

\begin{lemma}
\label{algebra-lemma-glue-modules}
$R$을 환이라 하고 $f_1, \ldots, f_n \in R$이라 하자. 다음 자료가
주어졌다고 하자.
\begin{enumerate}
\item 각 $i$에 대한 $R_{f_i}$-가군 $M_i$.
\item 각 쌍 $i, j$에 대한 $R_{f_if_j}$-가군 동형사상
$\psi_{ij} : (M_i)_{f_j} \to (M_j)_{f_i}$.
\end{enumerate}
이 자료가 모든 도식
$$
\xymatrix{
(M_i)_{f_jf_k}
\ar[rd]_{\psi_{ij}}
\ar[rr]^{\psi_{ik}}
& &
(M_k)_{f_if_j} \\
&
(M_j)_{f_if_k} \ar[ru]_{\psi_{jk}}
}
$$
이 가환한다는 ``코사이클 조건''을 만족한다고 하자(모든 삼중항
$i, j, k$에 대해). 이 자료로 다음을 정의한다.
$$
M = \Ker\left(
\bigoplus\nolimits_{1 \leq i \leq n} M_i
\longrightarrow
\bigoplus\nolimits_{1 \leq i, j \leq n} (M_i)_{f_j}
\right)
$$
여기서 $(m_1, \ldots, m_n)$은 $(i, j)$번째 성분이
$m_i/1 - \psi_{ji}(m_j/1)$인 원소로 간다. 그러면 자연스러운 사상
$M \to M_i$는 동형사상 $M_{f_i} \to M_i$를 유도한다. 또한 모든
$m \in M$에 대해 $\psi_{ij}(m/1) = m/1$이다(명백한 표기법을 썼다).
\end{lemma}

\begin{proof}
$M_{f_1} \to M_1$이 동형사상임을 보이기 위해 $R_{f_1}$의 모든
소 이데알 $\mathfrak p'$에서 국소화한 것이 동형사상임을 보이면
충분하다. Lemma \ref{algebra-lemma-characterize-zero-local}을 보라.
어떤 소 이데알 $\mathfrak p \subset R$에 대해
$\mathfrak p' = \mathfrak p R_{f_1}$이고 $f_1 \not \in \mathfrak p$라고
쓰자. Lemma \ref{algebra-lemma-standard-open}을 보라. 국소화는 완전하므로
(Proposition \ref{algebra-proposition-localization-exact}) 다음을 얻는다.
\begin{align*}
(M_{f_1})_{\mathfrak p'} & =
M_\mathfrak p \\
& =
\Ker\left(
\bigoplus\nolimits_{1 \leq i \leq n} M_{i, \mathfrak p}
\longrightarrow
\bigoplus\nolimits_{1 \leq i, j \leq n} ((M_i)_{f_j})_\mathfrak p
\right) \\
& =
\Ker\left(
\bigoplus\nolimits_{1 \leq i \leq n} M_{i, \mathfrak p}
\longrightarrow
\bigoplus\nolimits_{1 \leq i, j \leq n} (M_{i, \mathfrak p})_{f_j}
\right)
\end{align*}
여기서 Proposition \ref{algebra-proposition-localize-twice-module}도 사용했다.
$f_1$은 $R_\mathfrak p$의 단위원이므로, $R$을 $R_\mathfrak p$로,
$f_i$를 $R_\mathfrak p$에서의 $f_i$의 상으로, $M$을 $M_\mathfrak p$로,
$f_1$을 $1$로 바꾸어 $f_1 = 1$인 경우로 환원된다.

\medskip\noindent
$f_1 = 1$이라고 하자. 그러면 $j = 2, \ldots, n$에 대해
$\psi_{1j} : (M_1)_{f_j} \to M_j$는 동형사상이다. 이 동형사상들을
사용하여 $M_j = (M_1)_{f_j}$로 동일시하면,
$\psi_{ij} : (M_1)_{f_if_j} \to (M_1)_{f_if_j}$가 표준 동일시임을
알 수 있다. 따라서 복합체
$$
0 \to M_1 \to
\bigoplus\nolimits_{1 \leq i \leq n} (M_1)_{f_i}
\longrightarrow
\bigoplus\nolimits_{1 \leq i, j \leq n}
(M_1)_{f_if_j}
$$
는 Lemma \ref{algebra-lemma-cover-module}에 의해 완전하다. 따라서 이 경우
첫 번째 사상은 $M_1$을 $M$과 동일시하며, 나머지는 모두 명백하다.
\end{proof}

\section{영인자와 전분수환}
\label{algebra-section-total-quotient-ring}

\noindent
극소 소 이데알에서의 국소환은 다음 성질들을 갖는다.

\begin{lemma}
\label{algebra-lemma-minimal-prime-reduced-ring}
$\mathfrak p$를 환 $R$의 극소 소 이데알이라 하자.
$R_{\mathfrak p}$의 극대 이데알의 모든 원소는 멱영원이다.
$R$이 축소환이면 $R_{\mathfrak p}$는 체이다.
\end{lemma}

\begin{proof}
${\mathfrak p}R_{\mathfrak p}$의 어떤 원소 $x$가 멱영원이 아니면
$D(x) \not = \emptyset$이다. Lemma \ref{algebra-lemma-Zariski-topology}을 보라.
이는 $\mathfrak p$의 극소성에 모순이다. $R$이 축소환이면
${\mathfrak p}R_{\mathfrak p} = 0$이고, 따라서 체이다.
\end{proof}

\begin{lemma}
\label{algebra-lemma-reduced-ring-sub-product-fields}
$R$을 축소환이라 하자. 그러면 다음이 성립한다.
\begin{enumerate}
\item $R$은 체들의 곱의 부분환이다.
\item $R \to \prod_{\mathfrak p\text{ minimal}} R_{\mathfrak p}$는
체들의 곱으로의 매입이다.
\item $\bigcup_{\mathfrak p\text{ minimal}} \mathfrak p$는 $R$의
영인자들의 집합이다.
\end{enumerate}
\end{lemma}

\begin{proof}
Lemma \ref{algebra-lemma-minimal-prime-reduced-ring}에 의해 각 환
$R_\mathfrak p$는 체이다. 특히 환 사상 $R \to R_\mathfrak p$의 핵은
$\mathfrak p$이다. Lemma \ref{algebra-lemma-Zariski-topology}에 의해
$\bigcap_{\mathfrak p} \mathfrak p = (0)$이다. 따라서 (2)와 (1)이
성립한다. $x y = 0$이고 $y \not = 0$이면 어떤 극소 소 이데알
$\mathfrak p$에 대해 $y \not \in \mathfrak p$이다. 따라서
$x \in \mathfrak p$이다. 그러므로 $R$의 모든 영인자는
$\bigcup_{\mathfrak p\text{ minimal}} \mathfrak p$에 속한다.
역으로 어떤 극소 소 이데알 $\mathfrak p$에 대해
$x \in \mathfrak p$라고 하자. 그러면 $x$는 $R_\mathfrak p$에서
영으로 가므로, $y \not \in \mathfrak p$이고 $xy = 0$인
$y \in R$가 존재한다. 달리 말해 $x$는 영인자이다. 이로써 (3)과
보조정리의 증명이 끝난다.
\end{proof}

\noindent
환 $R$의 전분수환 $Q(R)$는 Example \ref{algebra-example-localize-at-prime}에서
도입했다.

\begin{lemma}
\label{algebra-lemma-total-ring-fractions}
$R$을 환이라 하고 $S \subset R$를 비영인자들로 이루어진 곱셈적
부분집합이라 하자. 그러면 $Q(R) \cong Q(S^{-1}R)$이다.
특히 $Q(R) \cong Q(Q(R))$이다.
\end{lemma}

\begin{proof}
$x \in S^{-1}R$가 비영인자이고 어떤 $r \in R$, $f \in S$에 대해
$x = r/f$이면, $r$은 $R$의 비영인자이다. 이로부터 보조정리가 따른다.
\end{proof}

\noindent
붙이기 결과들을 적용하여 Example \ref{algebra-example-localize-at-prime}에서
도입한 전분수환 $Q(R)$에 관한 내용을 증명할 수 있다.

\begin{lemma}
\label{algebra-lemma-total-ring-fractions-no-embedded-points}
$R$을 환이라 하자. $R$에 유한 개의 극소 소 이데알
$\mathfrak q_1, \ldots, \mathfrak q_t$가 있고,
$\mathfrak q_1 \cup \ldots \cup \mathfrak q_t$가 $R$의 영인자들의
집합이라고 하자. 그러면 전분수환 $Q(R)$는
$R_{\mathfrak q_1} \times \ldots \times R_{\mathfrak q_t}$와 같다.
\end{lemma}

\begin{proof}
임의의 비영인자는
% P01-E0300: 동결 원문은 집합의 원소 관계를 contained in으로 서술한다.
$R \setminus \mathfrak q_i$에 포함되므로 자연스러운 사상
$Q(R) \to R_{\mathfrak q_i}$들이 있다. 따라서 자연스러운 사상
$Q(R) \to R_{\mathfrak q_1} \times \ldots \times R_{\mathfrak q_t}$가
있다. 극소가 아닌 임의의 소 이데알 $\mathfrak p \subset R$에 대해
Lemma \ref{algebra-lemma-silly}로부터
$\mathfrak p \not \subset \mathfrak q_1 \cup \ldots \cup \mathfrak q_t$이다.
따라서 $\Spec(Q(R)) = \{\mathfrak q_1, \ldots, \mathfrak q_t\}$이다
($\Spec(R)$의 부분집합으로 보았다. Lemma
\ref{algebra-lemma-spec-localization}을 보라). 그러므로 $\Spec(Q(R))$는
유한 이산집합이고, Lemma \ref{algebra-lemma-disjoint-implies-product}에 의해
$\Spec(A_i) = \{\mathfrak{q}_i\}$인
$Q(R) = A_1 \times \ldots \times A_t$가 따른다. 또한 $A_i$는
$R$의 국소화인 국소환이다. 따라서 $A_i \cong R_{\mathfrak q_i}$이다.
\end{proof}

\section{스펙트럼의 기약 성분}
\label{algebra-section-irreducible}

\noindent
환의 스펙트럼의 기약 성분이 환의 극소 소 이데알에 대응함을 보인다.

\begin{lemma}
\label{algebra-lemma-irreducible}
$R$을 환이라 하자.
\begin{enumerate}
\item 소 이데알 $\mathfrak p \subset R$에 대해 자리스키 위상에서
$\{\mathfrak p\}$의 폐포는 $V(\mathfrak p)$이다. 식으로 쓰면
$\overline{\{\mathfrak p\}} = V(\mathfrak p)$이다.
\item $\Spec(R)$의 기약 닫힌 부분집합들은 정확히 소 이데알
$\mathfrak p \subset R$에 대한 부분집합 $V(\mathfrak p)$들이다.
\item $\Spec(R)$의 기약 성분들(Topology, Definition
\ref{topology-definition-irreducible-components} 참조)은 정확히
극소 소 이데알 $\mathfrak p \subset R$에 대한 부분집합
$V(\mathfrak p)$들이다.
\end{enumerate}
\end{lemma}

\begin{proof}
$ \mathfrak p \in V(I)$이면 $I \subset \mathfrak p$임에 유의하라.
따라서 $\overline{\{\mathfrak p\}} = V(\mathfrak p)$임이 명백하다.
특히 $V(\mathfrak p)$는 한원소집합의 폐포이므로 기약이다.
두 번째 명제가 세 번째 명제를 함의한다. 두 번째 명제를 보이기 위해
$I$가 근기 이데알인 $V(I) \subset \Spec(R)$를 생각하자.
$I$가 소 이데알이 아니면 $a, b\in R$, $a, b\not \in I$이고
$ab\in I$인 원소를 택한다. 이 경우
$V(I, a) \cup V(I, b) = V(I)$이지만,
Lemma \ref{algebra-lemma-Zariski-topology}에 의해 $V(I, b) = V(I)$도
$V(I, a) = V(I)$도 아니다. 따라서 $V(I)$는 기약이 아니다.
\end{proof}

\noindent
달리 말해 이 보조정리는 $\Spec(R)$의 모든 기약 닫힌 부분집합이 어떤
소 이데알 $\mathfrak p$에 대한 $V(\mathfrak p)$ 꼴임을 보여 준다.
$V(\mathfrak p) = \overline{\{\mathfrak p\}}$이므로 각 기약 닫힌
부분집합에는 유일한 일반점이 있다. Topology, Definition
\ref{topology-definition-generic-point}를 보라. 특히 $\Spec(R)$는
소버 위상공간이다. 이 사실을 다음 보조정리에 기록한다.

\begin{lemma}
\label{algebra-lemma-spec-spectral}
환의 스펙트럼은 스펙트럴 공간이다. Topology, Definition
\ref{topology-definition-spectral-space}를 보라.
\end{lemma}

\begin{proof}
형식적으로 이는 Lemma \ref{algebra-lemma-irreducible}과
Lemma \ref{algebra-lemma-topology-spec}에서 따른다. 위의 논의도 보라.
\end{proof}

\begin{lemma}
\label{algebra-lemma-irreducible-components-containing-x}
$R$을 환이라 하고 $\mathfrak p \subset R$를 소 이데알이라 하자.
\begin{enumerate}
\item $\mathfrak p$를 지나는 $\Spec(R)$의 기약 닫힌 부분집합들의
집합은 소 이데알 $\mathfrak q \subset R_{\mathfrak p}$들과 일대일로
대응한다.
\item $\mathfrak p$를 지나는 $\Spec(R)$의 기약 성분들의 집합은
극소 소 이데알 $\mathfrak q \subset R_{\mathfrak p}$들과 일대일로
대응한다.
\end{enumerate}
\end{lemma}

\begin{proof}
Lemma \ref{algebra-lemma-irreducible}과 $\Spec(R_\mathfrak p)$를 기술하는
Lemma \ref{algebra-lemma-spec-localization}에서 따른다. 그 기술에 따르면
$\Spec(R_\mathfrak p)$는 $\mathfrak q \subset \mathfrak p$인 $R$의
소 이데알 $\mathfrak q$들에 대응한다.
\end{proof}

\begin{lemma}
\label{algebra-lemma-standard-open-containing-maximal-point}
$R$을 환이라 하고 $\mathfrak p$를 $R$의 극소 소 이데알이라 하자.
$W \subset \Spec(R)$를 점 $\mathfrak p$를 포함하지 않는 준콤팩트
열린집합이라 하자. 그러면 $f \not \in \mathfrak p$이고
$D(f) \cap W = \emptyset$인 $f \in R$가 존재한다.
\end{lemma}

\begin{proof}
$W$가 준콤팩트이므로 이를 표준 아핀 열린집합 $D(g_i)$,
$i = 1, \ldots, n$의 유한 합집합으로 쓸 수 있다.
$\mathfrak p \not \in W$이므로 모든 $i$에 대해
$g_i \in \mathfrak p$이다. Lemma
\ref{algebra-lemma-minimal-prime-reduced-ring}에 의해 각 $g_i$는
$R_{\mathfrak p}$에서 멱영원이다. 따라서 모든 $i$에 대해 어떤
$n_i > 0$에 대해 $f g_i^{n_i} = 0$인 $f \in R$,
$f \not \in \mathfrak p$를 찾을 수 있다. 그러면 $D(f)$가 원하는
열린집합이다.
\end{proof}

\begin{lemma}
\label{algebra-lemma-ring-with-only-minimal-primes}
$R$을 환이라 하고 $X = \Spec(R)$를 위상공간으로 보자.
다음 조건들은 서로 동치이다.
\begin{enumerate}
\item $X$는 프로유한이다.
\item $X$는 하우스도르프 공간이다.
\item $X$는 완전 비연결이다.
\item $X$의 모든 준콤팩트 열린집합은 닫힌집합이다.
\item 소 이데알들 사이에 자명하지 않은 포함관계가 없다.
\item 모든 소 이데알은 극대 이데알이다.
\item 모든 소 이데알은 극소이다.
\item 모든 표준 열린집합 $D(f) \subset X$는 닫힌집합이다.
% P01-E0301: 동결 원문은 미완성 편집 자리표시자를 동치조건으로 남겨 두었다.
\item 여기에 더 추가한다.
\end{enumerate}
\end{lemma}

\begin{proof}
첫 번째 증명. (5), (6), (7)이 동치임은 명백하다. 모든 준콤팩트
열린집합은 표준 열린집합들의 유한 합집합이므로 (4)와 (8)이 동치임도
명백하다. (7) $\Rightarrow$ (4)는
Lemma \ref{algebra-lemma-standard-open-containing-maximal-point}에서 따른다.
(4)가 성립한다고 하자. $\mathfrak p, \mathfrak p'$를 $R$의 서로 다른
소 이데알이라 하자. $f \in \mathfrak p'$, $f \not \in \mathfrak p$인
원소를 택한다(필요하면 $\mathfrak p$와 $\mathfrak p'$을 바꾼다).
그러면 $\mathfrak p' \not \in D(f)$이고 $\mathfrak p \in D(f)$이다.
(4)에 의해 열린집합 $D(f)$는 닫힌집합이기도 하다. 따라서
$\mathfrak p$와 $\mathfrak p'$은 합집합이 $X$인 서로소 열린 근방들에
속한다. 그러므로 $X$는 하우스도르프 공간이고 완전 비연결이다.
따라서 (4) $\Rightarrow$ (2) 및 (3)이다. (3)이 성립하면
$\Spec(R)$의 점들 사이에 특수화가 존재할 수 없으므로 (5)가 성립한다.
$X$가 하우스도르프 공간이면 모든 점이 닫힌점이므로 (2)는 (6)을
함의한다. 따라서 (2), (3), (4), (5), (6), (7), (8)은 동치이다.
모든 프로유한 공간은 하우스도르프 공간이므로 (1)은 (2)를 함의한다.
$X$가 (2)와 (3)을 만족하면, $X$는 Lemma
\ref{algebra-lemma-quasi-compact}에 의해 준콤팩트이므로 Topology, Lemma
\ref{topology-lemma-profinite}에 의해 프로유한이다.

\medskip\noindent
두 번째 증명. (4)와 (8)의 동치를 제외하면 이는
Lemma \ref{algebra-lemma-spec-spectral}과 순수한 위상적 사실들에서 따른다.
Topology, Lemma \ref{topology-lemma-characterize-profinite-spectral}을
보라.
\end{proof}

\section{환의 스펙트럼의 예}
\label{algebra-section-examples-spectra}

\noindent
이 절에서는 스펙트럼의 몇 가지 예를 모은다.

\begin{example}
\label{algebra-example-spec-Zxmodx2minus4}
이 예에서는 $X = \Spec(\mathbf{Z}[x]/(x^2 - 4))$를 기술한다.
$\mathfrak{p}$를 $X$의 임의의 소 이데알이라 하자. 자연스러운 환 사상을
$\phi : \mathbf{Z} \to \mathbf{Z}[x]/(x^2 - 4)$라 하자. 그러면
$ \phi^{-1}(\mathfrak p)$는 $\mathbf{Z}$의 소 이데알이다.
$ \phi^{-1}(\mathfrak p) = (2)$이면 $\mathfrak p$가 $2$를 포함하므로,
사상 $ \mathbf{Z}[x]/(x^2 - 4) \to  \mathbf{Z}[x]/(x^2 - 4, 2)$를 통해
이는
$\mathbf{Z}[x]/(x^2 - 4, 2) \cong (\mathbf{Z}/2\mathbf{Z})[x]/(x^2)$의
소 이데알에 대응한다. $(\mathbf{Z}/2\mathbf{Z})[x]/(x^2)$의 임의의
소 이데알은 $(x^2)$을 포함하는 $(\mathbf{Z}/2\mathbf{Z})[x]$의
소 이데알에 대응한다. 따라서 그러한 소 이데알은 $x$를 포함한다.
$(\mathbf{Z}/2\mathbf{Z}) \cong (\mathbf{Z}/2\mathbf{Z})[x]/(x)$가
체이므로 $(x)$는 극대 이데알이다. 모든 소 이데알은 $(x)$를 포함하고
$(x)$는 극대이므로 이 환에는 소 이데알 $(x)$가 하나뿐이다.
따라서 이 경우 $\mathfrak p = (2, x)$이다. 이제 $q > 2$에 대해
$ \phi^{-1}(\mathfrak p) = (q)$이면, $\mathfrak p$가 $q$를 포함하므로
사상 $ \mathbf{Z}[x]/(x^2 - 4) \to  \mathbf{Z}[x]/(x^2 - 4, q)$를 통해
이는
$\mathbf{Z}[x]/(x^2 - 4, q) \cong (\mathbf{Z}/q\mathbf{Z})[x]/(x^2 - 4)$의
소 이데알에 대응한다. $(\mathbf{Z}/q\mathbf{Z})[x]/(x^2 - 4)$의
임의의 소 이데알은 $(x^2 - 4) = (x -2)(x + 2)$를 포함하는
$(\mathbf{Z}/q\mathbf{Z})[x]$의 소 이데알에 대응한다. 따라서 이 소
이데알들은 $x -2$ 또는 $x + 2$를 포함해야 한다.
$(\mathbf{Z}/q\mathbf{Z})[x]$가 PID이므로 영이 아닌 모든 소 이데알은
극대이고, 따라서 $(x-2)(x + 2)$를 포함하는
$(\mathbf{Z}/q\mathbf{Z})[x]$의 소 이데알은 정확히 두 개,
즉 $(x-2)$와 $(x + 2)$이다. 결론적으로
$2 \neq -2 \in \mathbf{Z}/(q)$이므로 두 소 이데알
$(q, x-2)$와 $(q, x + 2)$가 존재한다. 마지막으로
$\phi^{-1}(\mathfrak p) = (0)$인 경우를 다룬다. $\mathfrak p$는
$(x^2 - 4) = (x -2)(x + 2)$를 포함하는 $\mathbf{Z}[x]$의 소 이데알에
대응한다. 따라서 $\mathfrak p$는 $(x-2)$ 또는 $(x + 2)$를 포함한다.
그러므로 $\mathfrak p$는 가정에 의해 $\mathbf{Z}$와 오직 $0$에서만
만나는 $\mathbf{Z}[x]/(x - 2)$ 또는 $\mathbf{Z}[x]/(x + 2)$의
소 이데알에 대응한다. $\mathbf{Z}[x]/(x - 2) \cong \mathbf{Z}$이고
$\mathbf{Z}[x]/(x + 2) \cong \mathbf{Z}$이므로, 이는 $\mathfrak p$가
이 환들 중 하나에서 $0$에 대응해야 함을 뜻한다. 따라서 원래 환에서
$\mathfrak p = (x - 2)$ 또는 $\mathfrak p = (x + 2)$이다.
\end{example}

\begin{example}
\label{algebra-example-spec-Zx}
이 예에서는 $X = \Spec(\mathbf{Z}[x])$를 기술한다.
$\mathfrak p \in X$를 고정하자. $\phi : \mathbf{Z} \to \mathbf{Z}[x]$라
하면 $\phi^{-1}(\mathfrak p) \in \Spec(\mathbf{Z})$이다.
% P01-E0302 / P01-E0303: 동결 원문의 양의 표수 분류는 영 소 이데알을 빠뜨리고 정수 올림의 기약성 조건과 괄호를 잘못 처리한다.
$q$가 소수이고 $q > 0$일 때 $\phi^{-1}(\mathfrak p) = (q)$이면
$\mathfrak p$는 $(\mathbf{Z}/(q))[x]$의 소 이데알에 대응하며,
이는 $(\mathbf{Z}/(q))[x]$에서 기약인 다항식으로 생성되어야 한다.
이 다항식의 최소 차수 대표를 택하면 $\mathbf{Z}[x]$에서도 기약이다.
따라서 이 경우 $\mathfrak p = (q, f_q)$이고, 여기서 $f_q$는
$\mathbf{Z}[x]$의 기약다항식이며 $(\mathbf{Z}/(q) [x])$에서 보아도
기약이다. 이제 $\phi^{-1}(\mathfrak p) = (0)$이라고 하자.
% P01-E0304: 동결 원문의 일반 섬유 논의는 영 이데알을 빠뜨린다.
이 경우 $\mathfrak p$는 상수가 아닌 다항식들로 생성되어야 하며,
$\mathfrak p$가 소 이데알이므로 이들을 $\mathbf{Z}[x]$에서 기약이라고
할 수 있다. 가우스의 보조정리에 의해 이 다항식들은 $\mathbf{Q}[x]$에서도
기약이다. $\mathbf{Q}[x]$는 유클리드 정역이므로, $\mathfrak p$를
생성하는 서로 다른 기약원 $f, g$가 적어도 둘 있으면 어떤
$a, b \in \mathbf{Q}[x]$에 대해 $1 = af + bg$이다. 공통분모를 곱하면
$\bar{a}, \bar{b} \in \mathbf{Z}[x]$ 및 영이 아닌
$m \in \mathbf{Z}$에 대해 $m = \bar{a}f + \bar{b} g$이다.
이는 모순이다. 따라서 $\mathfrak p$는 $\mathbf{Z}[x]$의 기약다항식
하나로 생성된다.
\end{example}

\begin{example}
\label{algebra-example-spec-kxy}
이 예에서는 $k$가 임의의 체일 때 $X = \Spec(k[x, y])$를 기술한다.
명백히 $(0)$은 소 이데알이고, $k[x, y]$가 유일분해정역이므로
기약다항식으로 생성되는 임의의 주 이데알도 소 이데알이다. 이제
$\mathfrak p$가 주 이데알이 아닌 $X$의 원소라고 하자. $k[x, y]$는
뇌터 유일분해정역이므로 소 이데알 $\mathfrak p$는 유한 개의
기약다항식 $(f_1, \ldots, f_n)$으로 생성된다. 이제 서로 동반이 아닌
$k[x, y]$의 기약다항식 $f, g$에 대해
$(f, g) \cap k[x] \neq 0$임을 주장한다. 이를 보이기 위해서는
$k(x)[y]$에서 보았을 때 $f$와 $g$가 서로소임을 보이면 충분하다.
이 경우 $k(x)[y]$는 유클리드 정역이므로 유클리드 알고리즘을 적용하고
분모를 없애면
% P01-E0305: 동결 원문은 분모를 없앤 뒤의 계수 a,b를 k[x]에 둔다.
$p, a, b \in k[x]$에 대해 $p = af + bg$를 얻는다. 그렇지 않다고,
즉 어떤 비단위원 $h \in k(x)[y]$가 $f$와 $g$를 모두 나눈다고 하자.
그러면 가우스의 보조정리에 의해 어떤 $a, b \in k(x)$에 대해
% P01-E0306: 동결 원문은 y에 의존하는 ah,bh를 k[x]에 둔다.
$ah, bh \in k[x]$이고 $ah | f$ 및 $bh | g$이다. 기약성에 의해
$ah = f$이고 $bh = g$이다($h \notin k(x)$이기 때문이다).
따라서 $k(x)[y]$에서 $\frac{a}{b} g = f$이므로 $f, g $는 동반이다.
$k(x)$는 $k[x]$의 분수체이므로 공약수를 갖지 않는 원소
$r , s \in k[x]$에 대해 $g = \frac{r}{s} f $로 쓸 수 있다.
이는 $k[x, y]$에서 $sg = rf$임을 뜻하고, $k[x, y]$가
유일분해정역이므로 $s$는 $f$를 나누어야 한다. 따라서 $s = 1$ 또는
$s = f$이다. $s = f$이면 $r = g$이고, 따라서 $f, g \in k[x]$이므로
$k(x)$의 단위원이며 $k(x)[y]$에서 서로소여야 한다. 이는 가정에
모순이다. $s = 1$이면 $g = rf$이고 역시 모순이다. 따라서 유클리드
정역 $k(x)[y]$에서 $f, g$는 서로소여야 한다. 이로써 $\mathfrak p$가
어떤 기약다항식 $p \in k[x] \subset k[x, y]$를 포함하는 경우로
환원되었다. 위의 논의에 의해 $\mathfrak p$는 환
$k[x, y]/(p) = k(\alpha)[y]$의 소 이데알에 대응한다. 여기서
$\alpha$는 최소다항식이 $p$인 $k$ 위의 대수적 원소이다. 이 환은
PID이므로 임의의 소 이데알은 $(0)$ 또는 $k(\alpha)[y]$의
기약다항식에 대응한다. 따라서 $\mathfrak p$는 $(p)$ 또는 $(p, f)$
꼴이고, 여기서 $f$는 몫 $k[x, y]/(p)$에서 기약인 $k[x, y]$의
다항식이다.
\end{example}

\begin{example}
\label{algebra-example-affine-open-not-standard}
환
$$
R = \{ f \in \mathbf{Q}[z]\text{ with }f(0) = f(1) \}.
$$
을 생각하자. 사상
$$
\varphi : \mathbf{Q}[A, B] \to R
$$
을 $\varphi(A) = z^2-z$ 및 $\varphi(B) = z^3-z^2$로 정의한다.
$(A^3 - B^2 + AB) \subset \Ker(\varphi)$이고
$A^3 - B^2 + AB$가 기약임은 쉽게 확인할 수 있다. $\varphi$가
전사라고 하자. 그러면 $R$은 정역의 부분환이어서 정역이므로
$\Ker(\varphi)$는 $\mathbf{Q}[A, B]$의 소 이데알이어야 한다.
% P01-E0307: 동결 원문의 이 곡선식을 포함하는 극대 이데알 분류는 유효하지 않다.
$(A^3-B^2 + AB)$를 포함하는 소 이데알은 $(A^3-B^2 + AB)$ 자체와,
$f$가 $g$를 법으로 기약인 $f, g\in\mathbf{Q}[A, B]$에 대한 임의의
극대 이데알 $(f, g)$이다. 그러나 $R$은 체가 아니므로 핵은
$(A^3-B^2 + AB)$여야 한다. 따라서
% P01-E0308: 동결 원문은 유도된 동형사상의 화살표 방향을 거꾸로 쓴다.
$\varphi$는 동형사상 $R \to \mathbf{Q}[A, B]/(A^3-B^2 + AB)$를 준다.

\medskip\noindent
$\varphi$가 전사임을 보이려면 임의의 $f\in R$를 $A(z) = z^2-z$와
$B(z) = z^3-z^2$에 대한 $\mathbf{Q}$-계수 다항식으로 나타내야 한다.
관계 $zA(z) = B(z)$에 유의하라. $a = f(0) = f(1)$이라 하자.
그러면 $z(z-1)$이 $f(z)-a$를 나누어야 하므로
$f(z) = z(z-1)g(z)+a = A(z)g(z)+a$라고 쓸 수 있다.
$\deg(g) < 2$이면
% P01-E0309: 동결 원문은 선형인 g를 다음 경우에서 쓰일 h로 잘못 부른다.
$h(z) = c_1z + c_0$이고
$f(z) = A(z)(c_1z + c_0)+a = c_1B(z)+c_0A(z)+a$이므로 끝난다.
$\deg(g)\geq 2$이면 다항식 나눗셈 알고리즘에 의해
$g(z) = A(z)h(z)+b_1z + b_0$로 쓸 수 있다
($\deg(h)\leq\deg(g)-2$). 따라서
% P01-E0310: 동결 원문의 전개식은 앞 등식에 있던 상수항 a를 빠뜨린다.
$f(z) = A(z)^2h(z)+b_1B(z)+b_0A(z)$이다. $h(z)$에 나눗셈을 적용하고
반복하면 $f(z)$를 $A(z)$와 $B(z)$의 다항식으로 나타낼 수 있다.
따라서 $\varphi$는 전사이다.

\medskip\noindent
$a \in \mathbf{Q}$, $a \neq 0, \frac{1}{2}, 1$이라 하고
$$
R_a = \{ f \in \mathbf{Q}[z, \frac{1}{z-a}]\text{ with }f(0) = f(1)
\}.
$$
을 생각하자. 이 역시 유한 생성 $\mathbf{Q}$-대수이다. 실제로 함수
$z^2-z$, $z^3-z$, $\frac{a^2-a}{z-a}+z$가 $\mathbf{Q}$-대수로서
$R_a$를 생성함은 쉽게 확인할 수 있다. 다음 포함관계들이 있다.
$$
R\subset R_a\subset\mathbf{Q}[z, \frac{1}{z-a}], \quad
R\subset\mathbf{Q}[z]\subset\mathbf{Q}[z, \frac{1}{z-a}].
$$
환 T와 곱셈적 부분집합 $S\subset T$에 대해 환 사상
$T \to S^{-1}T$가 스펙트럼 위의 사상
$\Spec(S^{-1}T) \to \Spec(T)$를 유도하며, 이 사상은 부분집합
$$
\{\mathfrak p \in \Spec(T) \mid S \cap \mathfrak p = \emptyset\}
\subset \Spec(T).
$$
위로의 위상동형임을 상기하자(Lemma \ref{algebra-lemma-spec-localization}).
어떤 $f\in T$에 대해 $S = \{ 1, f, f^2, \ldots\}$이면 이는
% P01-E0311: 동결 원문은 표준 열린집합을 스펙트럼이 아니라 환 T의 부분집합으로 둔다.
열린집합 $D(f)\subset T$이다. 이제 환 사상 $R \to R_a$에 대해
이에 대응하는 성질을 확인한다. 포함관계 $R\subset R_a$가 유도하는 사상
$\theta : \Spec(R_a) \to \Spec(R)$가 $\Spec(R)$의 열린 부분집합
위로의 위상동형임을, $\theta$가 단사 국소 위상동형임을 확인하여
보이겠다. 이를 위해 이제 기술할 두 개의 특별 열린집합에 의한
$\Spec(R_a)$의 열린 덮개를 사용한다.
임의의 $r\in\mathbf{Q}$에 대해 $r$에서의 값매김으로 주어지는 준동형을
$\text{ev}_r : R \to \mathbf{Q}$라 하자. $r = 0$ 및 $r = 1-a$에 대해
이는 준동형 $\text{ev}_r' : R_a \to \mathbf{Q}$로 확장된다
(후자의 경우 $a\neq\frac{1}{2}$이므로 $\frac{1}{z-a}$가
$z = 1-a$에서 잘 정의되기 때문이다). 그러나 $\text{ev}_a$는
$R_a$로 확장되지 않는다. $\mathfrak{m}_r = \Ker(\text{ev}_r)$라고 쓰자.
그러면
$$
\mathfrak{m}_0 = (z^2-z, z^3-z),
$$
$$
\mathfrak{m}_a = ((z-1 + a)(z-a), (z^2-1 + a)(z-a)), \text{ and}
$$
$$
\mathfrak{m}_{1-a} = ((z-1 + a)(z-a), (z-1 + a)(z^2-a)).
$$
이다. 이를 확인하기 위해 오른쪽들이 왼쪽들에 명백히 포함됨에
유의하라. 그런 다음 생성원들을 $A$와 $B$로 나타내고
$R$을 $\mathbf{Q}[A, B]/(A^3-B^2 + AB)$로 보아 오른쪽들이
극대 이데알임을 확인하라. $\mathfrak{m}_a$는 $\theta$의 상에 속하지
않음에 유의하라. 실제로
$$
(z^2 - z)^2(z - a)\left(\frac{a^2 - a}{z - a} + z\right) =
(z^2 - z)^2(a^2 - a) + (z^2 - z)^2(z - a)z
$$
이다. 왼쪽은 $(z^2 - z)(z - a)$가 $\mathfrak m_a$에 속하고
$(z^2 - z)(\frac{a^2 - a}{z - a} + z)$가 $R_a$에 속하므로
$\mathfrak m_a R_a$에 속한다. 마찬가지로 원소
$(z^2 - z)^2(z - a)z$는 $(z^2 - z)$가 $R_a$에 속하고
$(z^2 - z)(z - a)$가 $\mathfrak m_a$에 속하므로
$\mathfrak m_a R_a$에 속한다. $a \not \in \{0, 1\}$이므로
$(z^2 - z)^2 \in \mathfrak m_a R_a$라고 결론 내린다. 따라서
$R_a$의 어떤 이데알 $I$도 $I \cap R = \mathfrak m_a$를 만족할 수
없다. 그런 $I$는 $(z^2 - z)^2$을 포함해야 하지만 이 원소는 $R$에는
속하고 $\mathfrak m_a$에는 속하지 않기 때문이다. 특별 열린집합
$D((z-1 + a)(z-a))\subset\Spec(R)$는 닫힌집합
$\{\mathfrak{m}_a, \mathfrak{m}_{1-a}\}$의 여집합과 같다.
이제 $R_{(z-1 + a)(z-a)} = (R_a)_{(z-1 + a)(z-a)}$임을 확인하라.
이 국소화된 환을 $R'$이라 부르면 사상 $R \to R'$은
$R \to R_a \to R'$로 분해된다. 따라서 Lemma
\ref{algebra-lemma-spec-localization}에 의해 이 사상들은
$\Spec(R') \subset \Spec(R_a)$와 $\Spec(R') \subset \Spec(R)$를
열린 부분집합으로 나타낸다. 그러므로
$D((z-1 + a)(z-a))$에 제한한
$\theta : \Spec(R_a) \to \Spec(R)$는 열린 부분집합 위로의
위상동형이다.
마찬가지로 $D((z^2 + z + 2a-2)(z-a)) \subset \Spec(R_a)$에 제한한
$\theta$는 열린 부분집합
$D((z^2 + z + 2a-2)(z-a)) \subset \Spec(R)$ 위로의 위상동형이다.
$z^2 + z + 2a-2$가 $\mathbf{Q}$ 위에서 기약인지에 따라, 앞의 특별
열린집합의 여집합은 닫힌점 $\mathfrak{m}_a$와 함께 한 개 또는 두 개의
닫힌점으로 이루어진다. 또한 $R_a$에서 원소
% P01-E0312: 동결 원문은 두 열린집합의 첫 생성원을 2a-2에서 2a-a로 바꾼다.
$(z^2 + z + 2a-a)(z-a)$와 $(z-1 + a)(z-a)$가 생성하는 이데알은
$R_a$ 전체이므로 이 두 특별 열린집합은 $\Spec(R_a)$를 덮는다.
따라서 $\theta$가 $\Spec(R)-\{\mathfrak{m}_a\}$ 위로의 위상동형임을
보이려면, 이 한 개 또는 두 개의 점이 결코 $\mathfrak{m}_{1-a}$와
같지 않음을 보이면 충분하다. 실제로 그러하다. $1-a$가
$z^2 + z + 2a-2$의 근인 것과 $a = 0$ 또는 $a = 1$인 것은 동치인데,
두 경우 모두 일어나지 않기 때문이다.

\medskip\noindent
이 위상동형은 $R$의 한 원소에서의 국소화와 같은 거동을 보이지만,
% P01-E0313: 동결 원문은 z-a를 역으로 만드는 국소화를 극대 이데알에서의 국소환이라고 잘못 부른다.
$\mathbf{Q}[z, \frac{1}{z-a}]$가 극대 이데알 $(z-a)$에서의
$\mathbf{Q}[z]$의 국소화인 반면 환 $R_a$는 $R$의 국소화가
{\it 아니다}. 실제로
% P01-E0314: 동결 원문은 모든 국소화가 새 단위원을 만든다고 과도하게 일반화한다.
임의의 국소화 $S^{-1}R$는 원래 환 $R$보다 더 많은 단위원을 만든다.
$R$의 단위원들은 $\mathbf{Q}$의 단위원들인 $\mathbf{Q}^\times$이다.
실제로 $R_a$의 단위원들이 $\mathbf{Q}^*$임은 쉽게 알 수 있다.
즉 $\mathbf{Q}[z, \frac{1}{z - a}]$의 단위원들은
$c \in \mathbf{Q}^*$ 및 $n \in \mathbf{Z}$에 대한 $c (z - a)^n$이고,
이들이 $R_a$에 속하는 것은 $n = 0$인 경우뿐임이 명백하다.
따라서 $R_a$에는 $R$보다 더 많은 단위원이 없으며, 그러므로
$R$의 국소화일 수 없다.

\medskip\noindent
$a\neq 0, 1$이라는 사실은 $\frac{1}{z-a}$가 $z = 0, 1$에서 의미를
갖도록 하기 위해 사용했다. $a\neq 1/2$라는 사실은 몇 군데에서
사용했다. (1) $R_a$ 위의 $\text{ev}_{1-a}$의 핵을 말할 수 있도록
하기 위해 사용했으며, 이는 $\mathfrak{m}_{1-a}$가 $R_a$의 점임을
보장한다(즉 $R_a$에는 $R$의 점이 정확히 하나 빠져 있다).
(2) 마지막에
% P01-E0315: 동결 원문은 앞에서 정의한 n 대신 새 지수 k와 ell을 설명 없이 도입한다.
$(z-a)^{k + \ell}$가 $R$에 속할 수 있는 것은 $k = \ell = 0$인 경우뿐임을
결론 내리기 위해 사용했다. 실제로 $a = 1/2$이면 $k + \ell$가 짝수인
한 이는 $R$에 속한다. 따라서 이 경우 $R_a$에는 $R$보다 더 많은
단위원이 실제로 존재하고, $R_a$는 $R$의 국소화일 수도 있다.
\end{example}

\section{소 이데알에 관한 메타 관찰}
\label{algebra-section-oka-families}

\noindent
이 절은 CRing 프로젝트에서 가져왔다. $R$을 환이라 하고
$S \subset R$를 곱셈적 부분집합이라 하자.
Lemma \ref{algebra-lemma-spec-localization}의 한 결과는 $S$와 만나지 않는다는
성질에 관해 극대인 이데알 $I \subset R$가 소 이데알이라는 것이다.
그 이유는 $S^{-1}R$의 어떤 극대 이데알 $\mathfrak m$에 대해
$I = R \cap \mathfrak m$이기 때문이다. 여러 이데알의 성질에서 극대인
것들이 소 이데알이라는 사실이 드러난다. 이를 보는 일반적인 방법은
\cite{Lam-Reyes}에서 개발되었다. 이 절에서는 잠시 본론을 벗어나 이
현상을 설명한다.

\medskip\noindent
$R$을 환이라 하자. $I$가 $R$의 이데알이고 $a \in R$이면 다음과 같이
정의한다.
$$
(I : a) = \left\{ x \in R \mid xa \in I\right\}.
$$
더 일반적으로 $J \subset R$가 이데알이면 다음과 같이 정의한다.
$$
(I : J) = \left\{ x \in R \mid xJ \subset I\right\}.
$$

\begin{lemma}
\label{algebra-lemma-colon}
$R$을 환이라 하자. 주 이데알 $J \subset R$와 임의의 이데알
$I \subset J$에 대해 $I = J (I : J)$이다.
\end{lemma}

\begin{proof}
$J = (a)$라고 하자. 그러면 $(I : J) = (I : a)$이다.
$I \subset J$이므로 임의의 $y \in I$는 어떤 $x \in (I : a)$에 대해
$y = xa$ 꼴이다. 따라서 $I \subset J (I : J)$이다. 역으로
$x \in (I : a)$이면 $xJ = (xa) \subset I$이고, 이는 반대쪽
포함관계를 증명한다.
\end{proof}

\noindent
$\mathcal{F}$를 $R$의 이데알들의 모음이라 하자. $\mathcal{F}$의
여집합에서 극대인 원소들이 소 이데알임을 보장하는 조건에 관심을 둔다.

\begin{definition}
\label{algebra-definition-oka-family}
$R$을 환이라 하고 $\mathcal{F}$를 $R$의 이데알들의 집합이라 하자.
$R \in \mathcal{F}$이고, 이데알 $I \subset R$와 어떤 $a \in R$에 대해
$(I : a), (I, a) \in \mathcal{F}$이면 언제나 $I \in \mathcal{F}$일 때
$\mathcal{F}$를 {\it 오카 족}이라 한다.
\end{definition}

\noindent
오카 족의 몇 가지 예를 들자. 첫 번째 예는 이 절의 서론에서 논의한
기본 예이다.

\begin{example}
\label{algebra-example-oka-family-not-meet-multiplicative-set}
$R$을 환이라 하고 $S$를 $R$의 곱셈적 부분집합이라 하자.
$\mathcal{F} = \{I \subset R \mid I \cap S \not = \emptyset\}$가 오카
족임을 주장한다. 실제로 어떤 $a \in R$에 대해
$(I : a), (I, a) \in \mathcal{F}$라고 하자.
$s \in (I, a) \cap S$와 $s' \in (I : a) \cap S$를 택하면
$ss' \in I \cap S$이므로 $I \in \mathcal{F}$이다. 따라서
$\mathcal{F}$는 오카 족이다.
\end{example}

\begin{example}
\label{algebra-example-oka-family-finitely-generated}
$R$을 환, $I \subset R$를 이데알, $a \in R$를 원소라 하자.
$(I : a)$가 $a_1, \ldots, a_n$으로 생성되고, $(I, a)$가
$b_1, \ldots, b_m \in I$일 때 $a, b_1, \ldots, b_m$으로 생성되면,
$I$는 $aa_1, \ldots, aa_n, b_1, \ldots, b_m$으로 생성된다.
이를 보기 위해 $x \in I$이면 $x \in (I, a)$이므로
$a, b_1, \ldots, b_m$의 선형결합이지만, $a$의 계수는
$(I:a)$에 속해야 함에 유의하라. 그 결과 유한 생성 이데알들의 족은
오카 족이다.
\end{example}

\begin{example}
\label{algebra-example-oka-family-principal}
환 $R$의 주 이데알들의 족이 오카 족임을 보이자. 실제로
$I \subset R$가 이데알이고 $a \in R$이며 $(I, a)$와 $(I : a)$가
주 이데알이라고 하자. $(I : a) = (I : (I, a))$임에 유의하라.
$J = (I, a)$로 두면 $J$와 $(I : J)$는 모두 주 이데알이다.
Lemma \ref{algebra-lemma-colon}에 의해 $I = J (I : J)$이다. 따라서 지금의
상황에서는 $J = (I, a)$와 $(I : J)$가 주 이데알이므로 $I$도
주 이데알이다.
\end{example}

\begin{example}
\label{algebra-example-oka-family-bound-cardinality}
$R$을 환이라 하고 $\kappa$를 무한 기수라 하자. 많아야 $\kappa$개의
원소로 생성될 수 있는 이데알들의 족은 오카 족이다. 논증은
Example \ref{algebra-example-oka-family-finitely-generated}의 논증과 비슷하므로
생략한다.
\end{example}

\begin{example}
\label{algebra-example-oka-family-property-modules}
$A$를 환, $I \subset A$를 이데알, $a \in A$를 원소라 하자.
첫 번째 화살표가 $a$에 의한 곱셈으로 주어지는 짧은 완전열
$0 \to A/(I : a) \to A/I \to A/(I, a) \to 0$이 있다. 따라서
$P$가 확대 아래에서 안정적이고 $0$에 대해 성립하는 $A$-가군의
성질이면, $A/I$가 $P$를 갖는 이데알 $I$들의 족은 오카 족이다.
\end{example}

\begin{proposition}
\label{algebra-proposition-oka}
$\mathcal{F}$가 이데알들의 오카 족이면 $\mathcal{F}$의 여집합의
임의의 극대 원소는 소 이데알이다.
\end{proposition}

\begin{proof}
$I \not \in \mathcal{F}$가 $\mathcal{F}$에 속하지 않는다는 성질에
관해 극대이지만 $I$가 소 이데알은 아니라고 하자.
$R \in \mathcal{F}$이므로 $I \not = R$임에 유의하라. $I$가 소
이데알이 아니므로 $ab \in I$인 $a, b \in R - I$를 찾을 수 있다.
그러면 $(I, a) \neq I$이고 $(I : a)$는 $b \not \in I$를 포함하므로
$(I : a) \neq I$이기도 하다. 따라서 $(I : a), (I, a)$는 둘 다
$I$를 진하게 포함하므로 $\mathcal{F}$에 속해야 한다. 오카 조건에
의해 $I \in \mathcal{F}$이고, 이는 모순이다.
\end{proof}

\noindent
이제 위의 예들 대부분을 환의 소 이데알에 관한 보조정리로 바꿀 수 있다.

\begin{lemma}
\label{algebra-lemma-simple}
$R$을 환이라 하고 $S$를 $R$의 곱셈적 부분집합이라 하자.
$I \cap S = \emptyset$이라는 성질에 관해 극대인 이데알
$I \subset R$는 소 이데알이다.
\end{lemma}

\begin{proof}
이는 이 절의 서론에서 논의한 예이다. 다른 증명으로는
Example \ref{algebra-example-oka-family-not-meet-multiplicative-set}과
Proposition \ref{algebra-proposition-oka}를 결합하라.
\end{proof}

\begin{lemma}
\label{algebra-lemma-cohen}
$R$을 환이라 하자.
\begin{enumerate}
\item 유한 생성이 아니라는 성질에 관해 극대인 이데알
$I \subset R$는 소 이데알이다.
\item $R$의 모든 소 이데알이 유한 생성이면 $R$의 모든 이데알이
유한 생성이다.\footnote{뒤에서 $R$이 뇌터환이라고 말할 것이다.}
\end{enumerate}
\end{lemma}

\begin{proof}
첫 번째 명제는 Example \ref{algebra-example-oka-family-finitely-generated}와
Proposition \ref{algebra-proposition-oka}의 즉각적인 결과이다. 두 번째를 위해
유한 생성이 아닌 이데알 $I \subset R$가 존재한다고 하자.
유한 생성이 아닌 이데알들의 전순서 사슬 $\left\{I_\alpha\right\}$의
합집합은 유한 생성이 아니다. 실제로 $I = \bigcup I_\alpha$가
$a_1, \ldots, a_n$으로 생성된다면 모든 생성원이 어떤
$I_\alpha $에 속하고, 따라서 그것을 생성하게 된다. 초른의 보조정리에
의해 유한 생성이 아니라는 성질에 관해 극대인 이데알이 존재한다.
첫 번째 부분에 의해 이 이데알은 소 이데알이다.
\end{proof}

\begin{lemma}
\label{algebra-lemma-primes-principal}
$R$을 환이라 하자.
\begin{enumerate}
\item 주 이데알이 아니라는 성질에 관해 극대인 이데알
$I \subset R$는 소 이데알이다.
\item $R$의 모든 소 이데알이 주 이데알이면 $R$의 모든 이데알이
주 이데알이다.
\end{enumerate}
\end{lemma}

\begin{proof}
첫 번째 부분은 Example \ref{algebra-example-oka-family-principal}과
Proposition \ref{algebra-proposition-oka}에서 따른다. 두 번째를 위해
주 이데알이 아닌 이데알 $I \subset R$가 존재한다고 하자.
% P01-E0316 / P01-E1283: 동결 원문의 관계절에는 are가 빠져 있다.
주 이데알이 아닌 이데알들의 전순서 사슬 $\left\{I_\alpha\right\}$의
합집합은 주 이데알이 아니다. 실제로 $I = \bigcup I_\alpha$가
$a$로 생성된다면 $a$는 어떤 $I_\alpha $에 속하고 $a$가 그것을
생성하게 된다. 초른의 보조정리에 의해 주 이데알이 아니라는 성질에
관해 극대인 이데알이 존재한다. 첫 번째 부분에 의해 이 이데알은
반드시 소 이데알이다.
\end{proof}

\begin{lemma}
\label{algebra-lemma-characterize-domain}
$R$을 환이라 하자.
\begin{enumerate}
\item 비영인자를 포함하지 않는 이데알들 가운데 극대인 이데알은
소 이데알이다.
\item $R$이 영환이 아니고 $R$의 영이 아닌 모든 소 이데알이
비영인자를 포함하면 $R$은 정역이다.
\end{enumerate}
\end{lemma}

\begin{proof}
비영인자들의 집합을 $S$라 하자. 이는 $R$의 곱셈적 부분집합이다.
따라서 $S$와 만나지 않는다는 성질에 관해 극대인 임의의 이데알은
소 이데알이다. Lemma \ref{algebra-lemma-simple}을 보라. 그러므로 영이 아닌
모든 소 이데알이 비영인자를 포함하면 $(0)$은 소 이데알이고, 즉
$R$은 정역이다.
\end{proof}

\begin{remark}
\label{algebra-remark-cohen-bound-cardinality}
$R$을 환이라 하고 $\kappa$를 무한 기수라 하자.
Example \ref{algebra-example-oka-family-bound-cardinality}와
Proposition \ref{algebra-proposition-oka}을 적용하면 $\kappa$개의 원소로
생성되지 않는다는 성질에 관해 극대인 임의의 이데알이 소 이데알임을
알 수 있다. 이 결과는 그다지 유용하지 않다. $R$의 모든 소 이데알은
$\aleph_0$개의 원소로 생성되지만 어떤 이데알은 그렇지 않은 환이
존재하기 때문이다. 구체적으로 $k$를 체라 하고, $T$를 기수가
$\aleph_0$보다 큰 집합이라 하며 다음과 같이 두자.
$$
R = k[\{x_n\}_{n \geq 1}, \{z_{t, n}\}_{t \in T, n \geq 0}]/
(x_n^2, z_{t, n}^2, x_n z_{t, n} - z_{t, n - 1})
$$
이는 유일한 소 이데알 $\mathfrak m = (x_n)$을 갖는 국소환이다.
그러나 이데알 $(z_{t, n})$은 가산 개의 원소로 생성될 수 없다.
\end{remark}

\begin{example}
\label{algebra-example-noetherian-topology-on-spec}
\begin{reference}
2020년 11월 12일 Lukas Heger의 논평.
\end{reference}
$R$을 환이라 하고 $X = \Spec(R)$라 하자. $X$의 닫힌 부분집합들은
% P01-E0317 / P01-E1284: 동결 원문의 radical ideas는 radical ideals의 오기이다.
$R$의 근기 이데알에 대응하므로(Lemma \ref{algebra-lemma-Zariski-topology}),
$X$가 뇌터 위상공간인 것과 근기 이데알들에 대해 오름사슬조건이
성립하는 것은 동치이다. 이는 모든 근기 이데알이 유한 생성 이데알의
근기인 것과 동치이다(세부사항은 생략한다). 다음과 같이 두자.
$$
\mathcal{F} = \{I \subset R \mid
\sqrt{I} = \sqrt{(f_1, \ldots, f_n)}\text{ for some }n
\text{ and }f_1, \ldots, f_n \in R\}.
$$
임의의 이데알 $I \subset R$와 임의의 원소 $a \in R$에 대해 성립하는
항등식
$$
\sqrt{I} = \sqrt{(I, a)(I : a)}
$$
을 사용하여 $\mathcal{F}$가 오카 족임을 보이는 일은 독자에게 맡긴다.
한편 어느 것도 $\mathcal{F}$에 속하지 않는 이데알들의 전순서 사슬
$\{I_\alpha\}$가 있으면 그 합집합 $I = \bigcup I_\alpha$도
$\mathcal{F}$에 속할 수 없다. 그렇지 않고
$\sqrt{I} = \sqrt{(f_1, \ldots, f_n)}$이면 어떤 $e$에 대해
$f_i^e \in I$이고, 이어서 $i$와 무관한 어떤 $\alpha$에 대해
$f_i^e \in I_\alpha$이며, 그러면
$\sqrt{I_\alpha} = \sqrt{(f_1, \ldots, f_n)}$이므로 모순이다.
따라서 $\mathcal{F}$에 속하지 않는 이데알들의 집합이 공집합이 아니면
극대 원소를 갖고, Lemma \ref{algebra-lemma-cohen}에서와 똑같이 $X$가 뇌터
위상공간인 것과 $R$의 모든 소 이데알이 어떤
$f_1, \ldots, f_n \in R$에 대한
$\sqrt{(f_1, \ldots, f_n)}$과 같은 것은 동치라고 결론 내린다.
이 결과가 필요해지면 여기에 주의 깊게 진술하고 증명할 것이다.
\end{example}

\section{유한 표시 환 사상의 상}
\label{algebra-section-images-finite-presentation}

\noindent
이 절에서는 환 준동형 $R \to S$가 유도하는 사상
$\Spec(S) \to \Spec(R)$의 위상에 관한 몇 가지 결과, 주로 슈발레 정리를
증명한다. 이를 위해 Topology, Section
\ref{topology-section-constructible}에서 정의한 구성가능 집합,
준콤팩트 집합, 레트로콤팩트 집합 등의 개념을 사용한다.

\begin{lemma}
\label{algebra-lemma-qc-open}
$U \subset \Spec(R)$를 열린집합이라 하자. 다음 조건들은 동치이다.
\begin{enumerate}
\item $U$는 $\Spec(R)$에서 레트로콤팩트이다.
\item $U$는 준콤팩트이다.
\item $U$는 표준 열린집합들의 유한 합집합이다.
\item 유한 생성 이데알 $I \subset R$가 존재하여
% P01-E0318: 동결 원문에서 X는 도입되지 않았으며 그대로 보존한다.
$X \setminus V(I) = U$이다.
\end{enumerate}
\end{lemma}

\begin{proof}
$\Spec(R)$가 준콤팩트이므로 (1) $\Rightarrow$ (2)이다. Lemma
\ref{algebra-lemma-quasi-compact}을 보라. 표준 열린집합들이 위상의 기저를
이루므로 (2) $\Rightarrow$ (3)이다. 이제 (3) $\Rightarrow$ (1)을
증명하자. $U = \bigcup_{i = 1\ldots n} D(f_i)$라 하자. $U$가
$\Spec(R)$에서 레트로콤팩트임을 보이려면 $\Spec(R)$의 임의의
준콤팩트 열린집합 $V$에 대하여 $U \cap V$가 준콤팩트임을 보이면
충분하다. (2) $\Rightarrow$ (3)에 따라 가능한 방식으로
$V = \bigcup_{j = 1\ldots m} D(g_j)$라고 쓰자. 각 표준 열린집합은
어떤 환의 스펙트럼과 위상동형이므로 준콤팩트이다. Lemma
\ref{algebra-lemma-standard-open}과 Lemma \ref{algebra-lemma-quasi-compact}을 보라.
따라서
$U \cap V =
(\bigcup_{i = 1\ldots n} D(f_i)) \cap (\bigcup_{j = 1\ldots m} D(g_j))
= \bigcup_{i, j} D(f_i g_j)$는 준콤팩트 열린집합들의 유한 합집합이므로
준콤팩트이다. 끝으로 (4)와 (3)의 동치는 Lemma
\ref{algebra-lemma-Zariski-topology}에서 따름에 유의하라.
\end{proof}

\begin{lemma}
\label{algebra-lemma-affine-map-quasi-compact}
$\varphi : R \to S$를 환 준동형이라 하자. 유도된 연속사상
$f : \Spec(S) \to \Spec(R)$는 준콤팩트이다. 임의의 구성가능 집합
$E \subset \Spec(R)$에 대하여 역상 $f^{-1}(E)$는 $\Spec(S)$에서
구성가능하다.
\end{lemma}

\begin{proof}
먼저 임의의 준콤팩트 열린집합 $U \subset \Spec(R)$의 역상이
준콤팩트임을 보이자. Lemma \ref{algebra-lemma-qc-open}에 따라 $U$를 표준
% P01-E0319: 동결 원문의 finite open은 finite union의 오기이다.
열린집합들의 유한 합집합으로 쓸 수 있다. 따라서 Lemma
\ref{algebra-lemma-spec-functorial}에 의해 $f^{-1}(U)$는 표준 열린집합들의
유한 합집합이다. 그러므로 Lemma \ref{algebra-lemma-qc-open}을 다시 적용하면
$f^{-1}(U)$는 준콤팩트이다. 두 번째 명제는 이제 Topology, Lemma
\ref{topology-lemma-inverse-images-constructibles}에서 따른다.
\end{proof}

\begin{lemma}
\label{algebra-lemma-constructible}
$R$을 환이라 하자. $\Spec(R)$의 부분집합이 구성가능일 필요충분조건은
$f, g_1, \ldots, g_m \in R$에 대하여
$D(f) \cap V(g_1, \ldots, g_m)$ 꼴의 부분집합들의 유한 합집합으로
쓸 수 있는 것이다.
\end{lemma}

\begin{proof}
Lemma \ref{algebra-lemma-qc-open}에 의해 부분집합 $D(f)$와
$V(g_1, \ldots, g_m)$의 여집합은 레트로콤팩트 열린집합이다. 따라서
$D(f) \cap V(g_1, \ldots, g_m)$는 구성가능 부분집합이고, 그러한
부분집합들의 임의의 유한 합집합도 구성가능하다. 역으로
$T \subset \Spec(R)$가 구성가능이라고 하자. Topology, Definition
\ref{topology-definition-constructible}에 따라 $U, V \subset \Spec(R)$가
레트로콤팩트 열린집합인 경우 $T = U \cap V^c$라고 가정해도 된다.
Lemma \ref{algebra-lemma-qc-open}에 의해
$U = \bigcup_{i = 1, \ldots, n} D(f_i)$ 및
$V = \bigcup_{j = 1, \ldots, m} D(g_j)$라고 쓸 수 있다. 그러면
$T = \bigcup_{i = 1, \ldots, n} \big(D(f_i) \cap V(g_1, \ldots, g_m)\big)$이다.
\end{proof}

\begin{lemma}
\label{algebra-lemma-constructible-is-image}
$R$을 환이라 하고 $T \subset \Spec(R)$를 구성가능 집합이라 하자.
그러면 유한 표시 환 준동형 $R \to S$가 존재하여 $T$는
$\Spec(S)$의 $\Spec(R)$에서의 상이다.
\end{lemma}

\begin{proof}
환들의 유한 곱의 스펙트럼은 각 스펙트럼의 서로소 합집합이다. Lemma
\ref{algebra-lemma-spec-product}를 보라. 따라서 $T = T_1 \cup T_2$이고
$T_1$ 및 $T_2$에 대해 결과가 성립하면 $T$에 대해서도 성립한다.
Lemma \ref{algebra-lemma-constructible}에 따라
$T = D(f) \cap V(g_1, \ldots, g_m)$이라고 가정해도 된다. 이 경우
$T$는 사상
$\Spec((R/(g_1, \ldots, g_m))_f) \to \Spec(R)$의 상이다. Lemma
\ref{algebra-lemma-standard-open}과 Lemma \ref{algebra-lemma-spec-closed}를 보라.
\end{proof}

\begin{lemma}
\label{algebra-lemma-open-fp}
$R$을 환이라 하자. $f$를 $R$의 원소라 하고 $S = R_f$로 두자.
그러면 $\Spec(S)$의 구성가능 부분집합의 상은 $\Spec(R)$에서
구성가능하다.
\end{lemma}

\begin{proof}
Lemma \ref{algebra-lemma-qc-open}을 별도의 언급 없이 거듭 사용한다.
$U, V$를 $\Spec(S)$의 준콤팩트 열린집합이라 하자.
$U \cap V^c$의 상이 구성가능임을 보이겠다. Lemma
\ref{algebra-lemma-standard-open}의 식별 $\Spec(S) = D(f)$ 아래에서
$U, V$는 $\Spec(R)$의 준콤팩트 열린집합 $U', V'$에 대응한다.
따라서 $U' \cap (V')^c$가 $\Spec(R)$에서 구성가능임을 보이면
충분하며, 이는 명백하다.
\end{proof}

\begin{lemma}
\label{algebra-lemma-closed-fp}
$R$을 환이라 하자. $I$를 $R$의 유한 생성 이데알이라 하고
$S = R/I$로 두자. 그러면 $\Spec(S)$의 구성가능 부분집합의 상은
$\Spec(R)$에서 구성가능하다.
\end{lemma}

\begin{proof}
$I = (f_1, \ldots, f_m)$이면 $V(I)$는 $\bigcup D(f_i)$의 여집합이다.
Lemma \ref{algebra-lemma-Zariski-topology}을 보라. 따라서 Lemma
\ref{algebra-lemma-qc-open}에 의해 이는 구성가능하다. 사상 $R \to S$를
$f \mapsto \overline{f}$로 나타내자. $\overline{U}, \overline{V}$가
$\Spec(S)$의 레트로콤팩트 열린집합일 때
$\overline{U} \cap \overline{V}^c$의 $\Spec(R)$에서의 상이
구성가능임을 보이면 된다. Lemma \ref{algebra-lemma-qc-open}에 의해
$\overline{U} = \bigcup D(\overline{g_i})$라고 쓸 수 있다.
$U = \bigcup D({g_i})$로 두면 $\overline{U}$의 상은
$U \cap V(I)$이고, 이는 $\Spec(R)$에서 구성가능하다. 마찬가지로
$\overline{V}$의 상은 $\Spec(R)$의 어떤 레트로콤팩트 열린집합
$V$에 대하여 $V \cap V(I)$이다. 따라서
$\overline{U} \cap \overline{V}^c$의 상은 원하는 대로
$U \cap V(I) \cap V^c$이다.
\end{proof}

\begin{lemma}
\label{algebra-lemma-affineline-open}
$R$을 환이라 하자. 사상 $\Spec(R[x]) \to \Spec(R)$는 열린사상이고,
임의의 표준 열린집합의 상은 준콤팩트 열린집합이다.
\end{lemma}

\begin{proof}
표준 열린집합 $D(f)$, $f\in R[x]$의 상이 준콤팩트 열린집합임을
보이면 충분하다. $D(f)$의 상은
$\Spec(R[x]_f) \to \Spec(R)$의 상이다. $\mathfrak p \subset R$를
소 이데알이라 하고, $\overline{f}$를 $\kappa(\mathfrak p)[x]$에서의
$f$의 상이라 하자. Lemma \ref{algebra-lemma-in-image}에서 보았듯이
$\mathfrak p$가 상에 속할 필요충분조건은
$R[x]_f \otimes_R \kappa(\mathfrak p) =
\kappa(\mathfrak p)[x]_{\overline{f}}$가 영환이 아닌 것이다. 이는
$f$가 $\kappa(\mathfrak p)[x]$에서 영으로 가지 않는다는 조건,
다시 말해 $f$의 어떤 계수가 $\mathfrak p$에 속하지 않는다는 조건과
정확히 같다. 따라서 $f = a_d x^d + \ldots + a_0$이면 $D(f)$의 상은
$D(a_d) \cup \ldots \cup D(a_0)$이다.
\end{proof}

\noindent
아래에서 사용할 특성다항식의 한 성질을 증명한다.

\begin{lemma}
\label{algebra-lemma-characteristic-polynomial-prime}
$R \to A$를 환 준동형이라 하자. $R$-가군으로서
$A \cong R^{\oplus n}$이라고 가정하고 $f \in A$라 하자. 곱셈사상
$m_f: A \to A$는 $R$-선형이므로 특성다항식
$P(T) = T^n + r_{n-1}T^{n-1} + \ldots + r_0 \in R[T]$를 갖는다.
임의의 소 이데알 $\mathfrak{p} \in \Spec(R)$에 대하여 $f$가
$A \otimes_R \kappa(\mathfrak{p})$ 위에서 멱영적으로 작용할
필요충분조건은 $\mathfrak p \in V(r_0, \ldots, r_{n-1})$인 것이다.
\end{lemma}

\begin{proof}
곱셈사상
$m_{\bar f}: A \otimes_R \kappa(\mathfrak p) \to
A \otimes_R \kappa(\mathfrak p)$의 특성다항식
$\bar P(T) \in \kappa(\mathfrak p)[T]$가 단지 $P(T)$의
$\kappa(\mathfrak p)[T]$에서의 상임을 보이면 결과는 쉽게 따른다.
여기서 이 사상은 $A \otimes_R \kappa(\mathfrak p)$의 원소에
$\bar f$를 곱하며,
% P01-E0320: 동결 원문은 f의 상을 kappa(p)의 원소로 잘못 기술한다.
그 원소는 $\kappa(\mathfrak p)$에서 본 $f$의 상이다.
$R$의 원소를 성분으로 갖는 $m_f$의 행렬을 $(a_{ij})$라 하되,
$R$-가군 $A$의 기저 $e_1, \ldots, e_n$을 사용하자. 그러면
$A \otimes_R \kappa(\mathfrak p) \cong (R \otimes_R
\kappa(\mathfrak p))^{\oplus n} = \kappa(\mathfrak p)^n$이고, 이는
$\kappa(\mathfrak p)$ 위의 $n$차원 벡터공간이며 기저는
$e_1 \otimes 1, \ldots, e_n \otimes 1$이다. 상은
$\bar f = f \otimes 1$이고, 따라서 곱셈사상 $m_{\bar f}$의 행렬은
$(a_{ij} \otimes 1)$이다. 그러므로 그 특성다항식은 정확히
$P(T)$의 상이다.

\medskip\noindent
선형대수학에서 $n$차원 벡터공간의 선형변환이 멱영적으로 작용할
필요충분조건은 그 특성다항식이 $T^n$인 것임을 안다(특성다항식은
최소다항식의 어떤 거듭제곱을 나누기 때문이다). 따라서 $f$가
$A \otimes_R \kappa(\mathfrak p)$ 위에서 멱영적으로 작용할
필요충분조건은 $\bar P(T) = T^n$인 것이다. 이는 모든
$0 \leq i \leq n - 1$에 대하여 $r_i \in \mathfrak p$일
필요충분조건이고, 즉 $\mathfrak p \in
V(r_0, \ldots, r_{n - 1}).$인 경우이다.
\end{proof}

\begin{lemma}
\label{algebra-lemma-affineline-special}
$R$을 환이라 하고 $f, g \in R[x]$를 다항식들이라 하자. $g$의
최고차항 계수가 $R$의 단위원이라고 가정하자.
% P01-E0321: 동결 원문의 There exists elements는 문법적 오류이다.
원소들 $r_i\in R$, $i = 1\ldots, n$이 존재하여
$D(f) \cap V(g)$의 $\Spec(R)$에서의 상은
$\bigcup_{i = 1, \ldots, n} D(r_i)$이다.
\end{lemma}

\begin{proof}
$g = ux^d + a_{d-1}x^{d-1} + \ldots + a_0$라고 쓰자. 여기서 $d$는
$g$의 차수이고, 따라서 $u \in R^*$이다. 환 $A = R[x]/(g)$를
생각하자. 이는 $R$-가군으로서 유한 자유이고, 그 기저는
$1, x, \ldots, x^{d-1}$의 상들이다. $A$ 위에서 $f$의 상에 의한
곱셈을 생각하자. 이는 $R$-가군 사상이다. 따라서 이 사상의
특성다항식을 $P(T) \in R[T]$라 둘 수 있다.
$P(T) = T^d + r_{d-1} T^{d-1} + \ldots + r_0$라고 쓰자.
$r_0, \ldots, r_{d-1}$이 원하는 성질을 가짐을 주장한다. 아래에서는
특성다항식의 다음 성질을 사용한다.
$$
\mathfrak p \in V(r_0, \ldots, r_{d-1})
\Leftrightarrow
\text{multiplication by }f\text{ is nilpotent on }
A \otimes_R \kappa(\mathfrak p).
$$
이는 Lemma \ref{algebra-lemma-characteristic-polynomial-prime}에서 증명하였다.

\medskip\noindent
$\mathfrak q\in D(f) \cap V(g)$라고 하고
$\mathfrak p = \mathfrak q \cap R$라 하자. 그러면 $f$에 의한 곱셈과
양립하는 영이 아닌 사상
$A \otimes_R \kappa(\mathfrak p) \to \kappa(\mathfrak q)$가 있다.
또한 $f$는 $\kappa(\mathfrak q)$ 위에서 단위원으로 작용한다. 따라서
$\mathfrak p \not \in  V(r_0, \ldots, r_{d-1})$라고 결론 내린다.

\medskip\noindent
다른 한편, $R$의 어떤 소 이데알 $\mathfrak p$와 어떤
$0 \leq i \leq d - 1$에 대하여 $r_i \not\in \mathfrak p$라고 하자.
그러면 $f$에 의한 곱셈은 대수
$A \otimes_R \kappa(\mathfrak p)$ 위에서 멱영적이지 않다. 따라서
$f$의 상을 포함하지 않는 소 이데알
$\overline{\mathfrak q} \subset A \otimes_R \kappa(\mathfrak p)$가
존재한다. $R[x]$에서 $\overline{\mathfrak q}$의 역상은
$\mathfrak p$로 가는 $D(f) \cap V(g)$의 원소이다.
\end{proof}

\begin{theorem}[슈발레 정리]
\label{algebra-theorem-chevalley}
$R \to S$가 유한 표시라고 가정하자. $\Spec(S)$의 구성가능 부분집합의
$\Spec(R)$에서의 상은 구성가능하다.
\end{theorem}

\begin{proof}
$S = R[x_1, \ldots, x_n]/(f_1, \ldots, f_m)$이라고 쓰자. 사상
$R \to S$는
$R \to R[x_1] \to R[x_1, x_2]
\to \ldots \to R[x_1, \ldots, x_{n-1}] \to S$로 분해할 수 있다.
따라서 $S = R[x]/(f_1, \ldots, f_m)$이라고 가정해도 된다. 이 경우
사상을 $R \to R[x] \to S$로 분해하고, Lemma
\ref{algebra-lemma-closed-fp}에 의해 $S = R[x]$인 경우로 귀착된다.
% P01-E0322: 동결 원문에는 it suffices의 it이 빠져 있다.
Lemma \ref{algebra-lemma-qc-open}에 따라
$T = (\bigcup_{i = 1\ldots n} D(f_i)) \cap V(g_1, \ldots, g_m)$이고
$f_i , g_j \in R[x]$일 때 그 상이 $\Spec(R)$에서 구성가능임을 보이면
충분하다. 구성가능 집합들의 유한 합집합은 구성가능이므로
$n = 1$인 경우, 즉 $T = D(f) \cap V(g_1, \ldots, g_m)$인 경우만
다루면 충분하다.

\medskip\noindent
$c \in R$이면
$$
\Spec(R) =
V(c) \amalg D(c) =
\Spec(R/(c)) \amalg \Spec(R_c),
$$
이고, 이에 대응하여 $\Spec(R[x]) =
V(c) \amalg D(c) = \Spec(R/(c)[x]) \amalg
\Spec(R_c[x])$이다. $T = D(f) \cap V(g_1, \ldots, g_m)$와 각 부분의
교집합은 여전히 같은 모양이며, 이때 $f$, $g_i$를 각각
$R/(c)[x]$, $R_c[x]$에서의 상으로 바꾼다. $T$의 $\Spec(R)$에서의
상은 $T \cap V(c)$의 상과 $T \cap D(c)$의 상의 합집합임에 유의하라.
Lemma \ref{algebra-lemma-open-fp}와 Lemma \ref{algebra-lemma-closed-fp}을 사용하면
두 부분의 상이 각각 $\Spec(R/(c))$와 $\Spec(R_c)$에서 구성가능임을
증명하는 것으로 충분하다.

\medskip\noindent
위와 같이 $T = D(f) \cap V(g_1, \ldots, g_m)$이고
$\deg(g_1) \leq \deg(g_2) \leq \ldots \leq \deg(g_m)$이라고
가정하자. $m$과 $g_i$들의 차수에 관해 귀납법을 사용한다.
$d_1 = \deg(g_1)$, 즉 $g_1 = c x^{d_1} + l.o.t$라 하자. 여기서
$c \in R$는 영이 아니다. $R$을 $R/(c)$와 $R_c$ 두 조각으로
나누면 $g_1$의 차수를 낮추거나(이는 귀납가정으로 다루어진다),
$c$가 가역인 경우로 귀착된다. $c$가 가역이고 $m > 1$이면
$g_2 = c' x^{d_2} + l.o.t$라고 쓰자. 이 경우
$g_2' = g_2 - (c'/c) x^{d_2 - d_1} g_1$을 생각하자. 이데알
$(g_1, g_2, \ldots, g_m)$과 $(g_1, g_2', g_3, \ldots, g_m)$이
같으므로
% P01-E0323: 동결 원문의 g_3 뒤에는 쉼표가 빠져 있다.
$T = D(f) \cap V(g_1, g_2', g_3\ldots, g_m)$임을 알 수 있다.
그런데 여기서 $g_2'$의 차수는 $g_2$의 차수보다 엄밀히 작으므로
이 경우는 귀납가정으로 다루어진다.

\medskip\noindent
% P01-E0324: 동결 원문의 bases case는 base cases의 오기이다.
위 귀납법의 기저 경우는 (a) $g$의 최고차항 계수가 가역인
$T = D(f) \cap V(g)$인 경우와 (b) $T = D(f)$인 경우이다. 이 두
경우는 Lemma \ref{algebra-lemma-affineline-special}과 Lemma
\ref{algebra-lemma-affineline-open}에서 다루었다.
\end{proof}

\section{상에 관한 추가 결과}
\label{algebra-section-more-images}

\noindent
이 절에서는 환 준동형에 대해 $\Spec$에서의 상과 관련된 몇 가지
보조정리를 더 모은다. 예를 들어 Section \ref{algebra-section-going-up}도 보라.

\begin{lemma}
\label{algebra-lemma-generic-finite-presentation}
$R \subset S$를 정역들의 포함이라 하자. $R \to S$가 유한형이라고
가정하자. 영이 아닌 $f \in R$와 영이 아닌 $g \in S$가 존재하여
$R_f \to S_{fg}$는 유한 표시이다.
\end{lemma}

\begin{proof}
$R$ 위에서 $S$의 생성원 수에 대한 귀납법을 쓴다. 증명 도중 영이 아닌
어떤 $f \in R$에 대하여 $R$을 $R_f$로, $S$를 $S_f$로 바꾸어도 된다.

\medskip\noindent
$S$가 $R$ 위에서 하나의 원소로 생성된다고 하자. 그러면 어떤 소 이데알
$\mathfrak q \subset R[x]$에 대하여 $S = R[x]/\mathfrak q$이다.
$\mathfrak q = (0)$이면 증명할 것이 없다. $\mathfrak q \not = (0)$이면
$x$에 관한 차수가 최소인 영이 아닌 원소 $h \in \mathfrak q$를 택하자.
$a_i \in R$ 및 $f \not = 0$에 대하여
$h = f x^d + a_{d - 1} x^{d - 1} + \ldots + a_0$라고 쓰자. $R$과
$S$에서 $f$를 가역화한 뒤에는 $h$가 일계수라고 가정해도 된다.
전사 $R$-대수 사상 $R[x]/(h) \to S$를 얻는다. $R$-가군으로서
$R[x]/(h) = R \oplus Rx \oplus \ldots \oplus Rx^{d - 1}$이고, $d$의
최소성에 의해 $R[x]/(h)$는 $S$로 단사적으로 사상된다. 따라서
$R[x]/(h) \cong S$는 $R$ 위에서 유한 표시이다.

\medskip\noindent
$S$가 $R$ 위에서 $n > 1$개의 원소로 생성된다고 하자.
$x_1, \ldots, x_n \in S$가 $S$를 생성한다고 하자. $S' \subset S$를
$x_1, \ldots, x_{n-1}$이 생성하는 $R$-부분대수라 하자.
% P01-E0325: 동결 원문에는 the induction hypothesis의 the가 빠져 있다.
귀납가정에 의해 영이 아닌 $f\in R$와 $g \in S'$가 존재하여
$R_f \to S'_{fg}$가 유한 표시임을 안다. 다음으로 귀납가정을
$S'_{fg} \to S_{fg}$에 적용하면 $f' \in S'_{fg}$와
$g' \in S_{fg}$가 존재하여 $S'_{fgf'} \to S_{fgf'g'}$가 유한
표시임을 알 수 있다. 결론을 내리는 일은 독자에게 맡긴다.
\end{proof}

\begin{lemma}
\label{algebra-lemma-characterize-image-finite-type}
$R \to S$를 유한형 환 준동형이라 하자. $X = \Spec(R)$ 및
$Y = \Spec(S)$로 나타내자.
% P01-E0326: 동결 원문의 Write ... the induced map은 문법적으로 불완전하다.
유도된 스펙트럼의 사상을 $f : Y \to X$로 쓰자.
$E \subset Y = \Spec(S)$를 구성가능 집합이라 하자. 점
$\xi \in X$가 $f(E)$에 속하면
$\overline{\{\xi\}} \cap f(E)$는 $\overline{\{\xi\}}$의 조밀한
열린 부분집합을 포함한다.
\end{lemma}

\begin{proof}
$\xi \in X$를 $f(E)$의 점이라 하자. $\xi$로 가는 점
$\eta \in E$를 택하자. $\mathfrak p \subset R$를 $\xi$에 대응하는
소 이데알이라 하고, $\mathfrak q \subset S$를 $\eta$에 대응하는
소 이데알이라 하자. 다음 도표를 생각하자.
$$
\xymatrix{
\eta \ar[r] \ar@{|->}[d] & E \cap Y' \ar[r] \ar[d] &
Y' = \Spec(S/\mathfrak q) \ar[r] \ar[d] &
Y \ar[d] \\
\xi \ar[r] & f(E) \cap X' \ar[r] &
X' = \Spec(R/\mathfrak p) \ar[r] &
X
}
$$
Lemma \ref{algebra-lemma-affine-map-quasi-compact}에 의해 집합 $E \cap Y'$는
$Y'$에서 구성가능하다. 따라서 $X$를 $X'$으로, $Y$를 $Y'$으로
바꾸어도 된다. 그러므로 $R \subset S$가 정역들의 포함이고, $\xi$는
$X$의 일반점이며, $\eta$는 $Y$의 일반점이라고 가정해도 된다.
Lemma \ref{algebra-lemma-generic-finite-presentation}와 슈발레 정리
(Theorem \ref{algebra-theorem-chevalley})를 함께 적용하면 조밀한 열린집합
$U \subset X$, $V \subset Y$가 존재하여 $f(V) \subset U$이고,
$f : V \to U$는 구성가능 집합을 구성가능 집합으로 보낸다.
$E \cap V$는 $V$에서 구성가능임에 유의하라. Topology, Lemma
\ref{topology-lemma-open-immersion-constructible-inverse-image}를 보라.
따라서 $f(E \cap V)$는 $U$에서 구성가능하고 $\xi$를 포함한다.
Topology, Lemma \ref{topology-lemma-generic-point-in-constructible}에
의해 $f(E \cap V)$는 조밀한 열린집합 $U' \subset U$를 포함한다.
\end{proof}

\noindent
이 절의 끝에서는 구성가능 집합과 무관한 스펙트럼 사상의 상에 관한
몇 가지 결과를 더 제시한다.

\begin{lemma}
\label{algebra-lemma-surjective-spec-radical-ideal}
$\varphi : R \to S$를 환 준동형이라 하자. 다음 조건들은 동치이다.
\begin{enumerate}
\item 사상 $\Spec(S) \to \Spec(R)$는 전사이다.
\item 임의의 이데알 $I \subset R$에 대하여 $\sqrt{IS}$의 $R$에서의
역상은 $\sqrt{I}$와 같다.
\item 임의의 근기 이데알 $I \subset R$에 대하여 $IS$의 $R$에서의
역상은 $I$와 같다.
\item $R$의 모든 소 이데알 $\mathfrak p$에 대하여
$\mathfrak p S$의 $R$에서의 역상은 $\mathfrak p$이다.
\end{enumerate}
이 경우 임의의 기저변환 뒤에도 같은 명제가 성립한다. 즉 환 준동형
$R \to R'$가 주어지면 환 준동형 $R' \to R' \otimes_R S$도 동치인
성질 (1), (2), (3)을 갖는다.
\end{lemma}

\begin{proof}
$J \subset S$가 이데알이면
$\sqrt{\varphi^{-1}(J)} = \varphi^{-1}(\sqrt{J})$이다. 따라서 (2)와
(3)은 동치이다. (3) $\Rightarrow$ (4)는 즉각적이다. $I \subset R$가
근기 이데알이면 Lemma \ref{algebra-lemma-Zariski-topology}에 의해
$I = \bigcap_{I \subset \mathfrak p} \mathfrak p$이다. 따라서
(4) $\Rightarrow$ (2)이다. Lemma \ref{algebra-lemma-in-image}에 의해
$\mathfrak p = \varphi^{-1}(\mathfrak p S)$일 필요충분조건은
$\mathfrak p$가 상에 속하는 것이다. 따라서 (1) $\Leftrightarrow$ (4)이다.
그러므로 (1), (2), (3), (4)는 동치이다.

\medskip\noindent
(1)이 성립한다고 가정하자. $R \to R'$를 환 준동형이라 하자.
$\mathfrak p' \subset R'$를 $R$의 소 이데알 $\mathfrak p$ 위에 놓인
소 이데알이라 하자. $\mathfrak p'$가
$\Spec(R' \otimes_R S) \to \Spec(R')$의 상에 속함을 보이려면
Lemma \ref{algebra-lemma-in-image}에 따라
$(R' \otimes_R S) \otimes_{R'} \kappa(\mathfrak p')$가 영이 아님을
보여야 한다. 그런데
$$
(R' \otimes_R S) \otimes_{R'} \kappa(\mathfrak p') =
S \otimes_R \kappa(\mathfrak p)
\otimes_{\kappa(\mathfrak p)} \kappa(\mathfrak p')
$$
이다. 가정에 의해 $S \otimes_R \kappa(\mathfrak p)$는 영이 아니고
$\kappa(\mathfrak p) \to \kappa(\mathfrak p')$는 체의 확대이므로
이 환도 영이 아니다.
\end{proof}

\begin{lemma}
\label{algebra-lemma-domain-image-dense-set-points-generic-point}
$R$을 정역이라 하고 $\varphi : R \to S$를 환 준동형이라 하자.
다음 조건들은 동치이다.
\begin{enumerate}
\item 환 준동형 $R \to S$는 단사이다.
\item $\Spec(S) \to \Spec(R)$의 상은 조밀한 점들의 집합을 포함한다.
\item $R$에서의 역상이 $(0)$인 소 이데알
$\mathfrak q \subset S$가 존재한다.
\end{enumerate}
\end{lemma}

\begin{proof}
$K$를 정역 $R$의 분수체라 하자. $R \to S$가 단사라고 가정하자.
국소화는 완전하므로 $K \to S \otimes_R K$도 단사이다. 따라서
Lemma \ref{algebra-lemma-in-image}에 의해 $(0)$으로 가는 소 이데알이 있다.

\medskip\noindent
$(0)$은 $\Spec(R)$에서 조밀하므로 마지막 조건은 두 번째 조건을
함의함에 유의하라.

\medskip\noindent
두 번째 조건이 성립한다고 하자. $f \in R$, $f \not = 0$라 하자.
$R$이 정역이므로
% P01-E0327: 동결 원문은 V(f)를 Spec(R)가 아니라 R의 부분집합이라 한다.
$V(f)$는 $R$의 진닫힌 부분집합이다. 가정에 의해
$\varphi(f) \not \in \mathfrak q$인 $S$의 소 이데알 $\mathfrak q$가
존재한다. 따라서 $\varphi(f) \not = 0$이다. 그러므로
$R \to S$는 단사이다.
\end{proof}

\begin{lemma}
\label{algebra-lemma-injective-minimal-primes-in-image}
$R \subset S$를 단사 환 준동형이라 하자. 그러면
$\Spec(S) \to \Spec(R)$는 모든 극소 소 이데알을 만난다.
\end{lemma}

\begin{proof}
$\mathfrak p \subset R$를 극소 소 이데알이라 하자. 이 경우
$R_{\mathfrak p}$는 유일한 소 이데알을 갖는다. 따라서
$S_{\mathfrak p}$가 영이 아님을 보이면 충분하다. 이는 국소화가
완전하다는 사실에서 따른다. Proposition
\ref{algebra-proposition-localization-exact}를 보라.
\end{proof}

\begin{lemma}
\label{algebra-lemma-image-dense-generic-points}
$R \to S$를 환 준동형이라 하자. 다음 조건들은 동치이다.
\begin{enumerate}
\item $R \to S$의 핵은 멱영원들로 이루어진다.
\item $R$의 극소 소 이데알들은 $\Spec(S) \to \Spec(R)$의 상에 속한다.
\item $\Spec(S) \to \Spec(R)$의 상은 $\Spec(R)$에서 조밀하다.
\end{enumerate}
\end{lemma}

\begin{proof}
$I = \Ker(R \to S)$라 하자. Lemma \ref{algebra-lemma-Zariski-topology}에 의해
$\sqrt{(0)} = \bigcap_{\mathfrak q \subset S} \mathfrak q$임에 유의하라.
따라서 $\sqrt{I} = \bigcap_{\mathfrak q \subset S} R \cap \mathfrak q$이다.
그러므로 $V(I) = V(\sqrt{I})$는 $\Spec(S) \to \Spec(R)$의 상의
폐포이다. 이로써 (1)과 (3)이 동치임을 알 수 있다. (2)가 (3)을
함의하는 것은 명백하다. 끝으로 (1)을 가정하자. 스펙트럼의 위상과
사상을 바꾸지 않고 $R$을 $R/I$로, $S$를 $S/IS$로 바꾸어도 된다.
따라서 (1) $\Rightarrow$ (2)는 Lemma
\ref{algebra-lemma-injective-minimal-primes-in-image}에서 따른다.
\end{proof}

\begin{lemma}
\label{algebra-lemma-minimal-prime-image-minimal-prime}
$R \to S$를 환 준동형이라 하자. 극소 소 이데알
$\mathfrak p \subset R$가 $\Spec(S) \to \Spec(R)$의 상에 속하면,
이는 어떤 극소 소 이데알의 상이다.
\end{lemma}

\begin{proof}
$\mathfrak p = \mathfrak q \cap R$라고 하자. Lemma
\ref{algebra-lemma-Zariski-topology}에 따라 $\mathfrak r \subset \mathfrak q$인
극소 소 이데알 $\mathfrak r \subset S$를 택하자. $\mathfrak p$의
최소성에 의해 $\mathfrak p = \mathfrak r \cap R$임을 안다.
\end{proof}

\begin{lemma}
\label{algebra-lemma-intersection-not-zero}
$A \subset B$를 정역들의 포함이라 하고, 이 포함이 분수체들의
대수적 확대를 유도한다고 하자. $J \subset B$가 영이 아닌 이데알이면
$A \cap J$도 영이 아니다. 따라서 $\Spec(B)$의 진닫힌 부분집합의 상은
$\Spec(A)$에서 조밀하지 않다.
\end{lemma}

\begin{proof}
$x \in J$를 영이 아닌 원소라 하자. $x$는 $A$의 분수체 위에서
대수적이므로, $a_0, a_d \not = 0$이며
$d \geq 1$ 및 $a_0, \ldots, a_d \in A$가 존재하여 $B$에서
$a_d x^d + a_{d - 1} x^{d - 1} + \ldots + a_0 = 0$이다. 그러면
$a_0 \in A \cap J$이다.
\end{proof}

\section{뇌터환}
\label{algebra-section-Noetherian}

\noindent
환 $R$의 모든 이데알이 유한 생성이면 $R$을 {\it 뇌터환}이라고 한다.
이는 $R$의 이데알들에 대한 상승 사슬 조건과 동치임이 명백하다.
Lemma \ref{algebra-lemma-cohen}에 의해 $R$의 모든 소 이데알이 유한 생성임을
확인하는 것으로 충분하다.

\begin{lemma}
\label{algebra-lemma-Noetherian-permanence}
\begin{slogan}
뇌터 성질은 유한형 확대와 국소화를 거쳐도 안정적이다.
\end{slogan}
뇌터환 위에서 유한 생성된 모든 환은 뇌터환이다. 뇌터환의 임의의
국소화도 뇌터환이다.
\end{lemma}

\begin{proof}
국소화에 관한 명제는 $J \subset S^{-1}R$인 임의의 이데알이
$I \cdot S^{-1}R$ 꼴이라는 사실에서 따른다. 뇌터환에서의 임의의
몫 $R/I$도 뇌터환이다. 실제로 환 $R$은 뇌터환이고 이 몫의 임의의 이데알
$\overline{J} \subset R/I$는 어떤 이데알 $I \subset J \subset R$에 대한
$J/I$ 꼴이다. 따라서 $R$이 뇌터환이면 $R[X]$도 뇌터환임을 보이면
충분하다. $R[X]$의 이데알들에 대한 상승 사슬
$J_1 \subset J_2 \subset \ldots$을 생각하자. $I_{i, d}$를 $J_i$의
차수 $d$ 다항식의 최고차항 계수로 나타나는 $R$의 원소들이 생성하는
이데알로 정의하자. $i \leq i'$ 및 $d \leq d'$이면
$I_{i, d} \subset I_{i', d'}$임은 명백하다. $R$의 상승 사슬 조건에
의해 모든 $I_{i, d}$ 가운데 서로 다른 이데알은 유한 개뿐이다.
(힌트: $\mathbf{N} \times \mathbf{N}$의 무한 부분집합은 증가하는
무한 수열을 포함한다.)
$i_0$를 충분히 크게 택하여 모든 $i \geq i_0$ 및 모든 $d$에 대하여
$I_{i, d} = I_{i_0, d}$가 되게 하자. 어떤 $i \geq i_0$에 대해
$f \in J_i$라고 하자. 차수 $d = \deg(f)$에 대한 귀납법으로
$f \in J_{i_0}$임을 보인다. 즉, $f$와 같은 최고차항 계수를 갖고
차수가 $d$인 $g\in J_{i_0}$가 존재한다. 귀납가정에 의해
$f - g \in J_{i_0}$이고, 증명이 끝난다.
\end{proof}

\begin{lemma}
\label{algebra-lemma-Noetherian-power-series}
$R$이 뇌터환이면 형식적 멱급수환
$R[[x_1, \ldots, x_n]]$도 뇌터환이다.
\end{lemma}

\begin{proof}
$R[[x_1, \ldots, x_{n + 1}]] \cong
R[[x_1, \ldots, x_n]][[x_{n + 1}]]$이므로, $R$이 뇌터환이면
$R[[x]]$가 뇌터환이라는 명제를 증명하면 충분하다. 이데알
$I \subset R[[x]]$를 택하자. $I$가 유한 생성 이데알임을 보여야 한다.
각 정수 d에 대하여
$I_d = \{a \in R \mid ax^d + \text{h.o.t.} \in I\}$라고 두자.
그러면 $R$이 뇌터환이므로 $I_0 \subset I_1 \subset \ldots$는 안정화된다.
$I_{d_0} = I_{d_0 + 1} = \ldots$가 되도록 $d_0$를 택하자.
$d \leq d_0$인 각 $d$에 대해
$f_{d, j} \in I \cap (x^d)$, $j = 1, \ldots, n_d$인 원소를 택하여
$f_{d, j} = a_{d, j}x^d + \text{h.o.t}$라고 쓸 때
$I_d = (a_{d, j})$가 되게 하자. 다음과 같이 두자.
$I' = (\{f_{d, j}\}_{d = 0, \ldots, d_0, j = 1, \ldots, n_d})$.
그러면 $I' \subset I$임은 명백하다. $f \in I$를 택하자. 먼저
$c_{d, i} \in R$를 택하여
$$
f - \sum c_{d, i} f_{d, i} \in (x^{d_0 + 1}) \cap I.
$$
다음으로 $c_{i, 1} \in R$, $i = 1, \ldots, n_{d_0}$를 택하여
$$
f - \sum c_{d, i} f_{d, i} - \sum c_{i, 1}xf_{d_0, i}
\in (x^{d_0 + 2}) \cap I.
$$
다음으로 $c_{i, 2} \in R$, $i = 1, \ldots, n_{d_0}$를 택하여
$$
f - \sum c_{d, i} f_{d, i} - \sum c_{i, 1}xf_{d_0, i}
- \sum c_{i, 2}x^2f_{d_0, i}
\in (x^{d_0 + 3}) \cap I.
$$
이 과정을 계속하면 결국
$$
f = \sum c_{d, i} f_{d, i} +
\sum\nolimits_i (\sum\nolimits_e c_{i, e} x^e)f_{d_0, i}
$$
를 얻는다. 이는 원하는 대로 $I'$에 속한다.
\end{proof}

\noindent
다음 보조정리는 쉽지만 ``절대 뇌터 환 감소''라고 불리는 기법에서
유한형 $\mathbf{Z}$-대수가 자주 등장하기 때문에 유용하다.

\begin{lemma}
\label{algebra-lemma-obvious-Noetherian}
체 위의 유한형 대수는 뇌터환이다.
유한형 $\mathbf{Z}$-대수는 뇌터환이다.
\end{lemma}

\begin{proof}
이는 Lemma \ref{algebra-lemma-Noetherian-permanence}와 체가 뇌터환이라는
사실, 그리고
% P01-E0329: 동결 원문은 a Noetherian ring에서 관사가 빠져 있다.
$\mathbf{Z}$가 주 이데알 정역이므로 뇌터환이라는 사실에서 즉시 따른다.
\end{proof}

\begin{lemma}
\label{algebra-lemma-Noetherian-finite-type-is-finite-presentation}
$R$을 뇌터환이라 하자.
\begin{enumerate}
\item 유한 $R$-가군은 유한 표시이다.
\item 유한 $R$-가군의 부분가군은 유한하다.
\item 유한형 $R$-대수는 $R$ 위에서 유한 표시이다.
\end{enumerate}
\end{lemma}

\begin{proof}
$M$을 유한 $R$-가군이라 하자. Lemma
\ref{algebra-lemma-trivial-filter-finite-module}에 의해 $M$의 유한 여과를
찾을 수 있으며, 그 연속적인 몫들은 $R/I$ 꼴이다. 모든 이데알이
유한 생성이므로 각 몫 $R/I$는 유한 표시이다. 따라서 Lemma
\ref{algebra-lemma-extension}에 의해 $M$은 유한 표시이다. 이것이 (1)을
증명한다.

\medskip\noindent
$N \subset M$을 부분가군이라 하자. $M$이 유한이므로 몫 $M/N$도
유한하다. 따라서 (1)에 의해 $M/N$은 유한 표시이다. 그러므로
Lemma \ref{algebra-lemma-extension}의 (5)에 의해 $N$은 유한임을 알 수 있다.
이것이 (2)를 증명한다.

\medskip\noindent
(3)을 보이기 위해 $R[x_1, \ldots, x_n]$의 임의의 이데알이 Lemma
\ref{algebra-lemma-Noetherian-permanence}에 의해 유한 생성임에 유의하라.
\end{proof}

\begin{lemma}
\label{algebra-lemma-Noetherian-topology}
$R$이 뇌터환이면 $\Spec(R)$는 뇌터 위상 공간이다. Topology,
Definition \ref{topology-definition-noetherian}을 참조하라.
\end{lemma}

\begin{proof}
이는 $\Spec(R)$의 모든 닫힌 부분집합이 근기 이데알 $I$에 대해
$V(I)$ 꼴로 유일하게 주어진다는 사실 때문이다. Lemma
\ref{algebra-lemma-Zariski-topology}를 참조하라. 그리고 이 대응은 포함을
뒤집는다. 따라서 정의로부터 결과가 따른다.
\end{proof}

\begin{lemma}
\label{algebra-lemma-Noetherian-irreducible-components}
\begin{slogan}
뇌터 아핀 스킴은 유한 개의 일반점을 갖는다.
\end{slogan}
$R$이 뇌터환이면 $\Spec(R)$는 유한 개의 기약 성분을 갖는다. 다시
말해, $R$은 유한 개의 극소 소 이데알을 갖는다.
\end{lemma}

\begin{proof}
Lemma \ref{algebra-lemma-Noetherian-topology}와 Topology, Lemma
\ref{topology-lemma-Noetherian}에 의해 기약 성분이 유한 개임을
알 수 있다. Lemma \ref{algebra-lemma-irreducible}에 의해 이 성분들은
$R$의 극소 소 이데알에 대응한다.
\end{proof}

\begin{lemma}
\label{algebra-lemma-Noetherian-base-change-finite-type}
$R \to S$를 환 사상이라 하자. $R \to R'$가 유한형이라고 하자.
$S$가 뇌터환이면 기저변환 $S' = R' \otimes_R S$도 뇌터환이다.
\end{lemma}

\begin{proof}
Lemma \ref{algebra-lemma-base-change-finiteness}에 의해 유한형 성질은
기저변환에 대해 안정적이다. 따라서 $S \to S'$는 유한형이다.
$S$가 뇌터환이므로 Lemma \ref{algebra-lemma-Noetherian-permanence}를
적용할 수 있다.
\end{proof}

\begin{lemma}
\label{algebra-lemma-Noetherian-field-extension}
$k$를 체라 하고 $R$을 뇌터 $k$-대수라 하자. $K/k$가 유한 생성
체 확대이면 $K \otimes_k R$는 뇌터환이다.
\end{lemma}

\begin{proof}
$K/k$가 유한 생성 체 확대이므로, 유한 생성 $k$-대수
$B \subset K$가 존재하여 $K$가 $B$의 분수체가 된다. 다시 말해
$K = S^{-1}B$이고 $S = B \setminus \{0\}$이다. 그러면
$K \otimes_k R = S^{-1}(B \otimes_k R)$이다. Lemma
\ref{algebra-lemma-Noetherian-base-change-finite-type}에 의해
$B \otimes_k R$는 뇌터환이다. 마지막으로 Lemma
\ref{algebra-lemma-Noetherian-permanence}에 의해
$K \otimes_k R = S^{-1}(B \otimes_k R)$는 뇌터환이다.
\end{proof}

\noindent
때때로 유용한 재미있는 보조정리들을 소개한다.

\begin{lemma}
\label{algebra-lemma-subring-of-local-ring}
$R$을 환이라 하고 $\mathfrak p \subset R$를 소 이데알이라 하자.
다음 각각의 경우에 $R_f \to R_\mathfrak p$가 단사이고
$f \in R$, $f \not \in \mathfrak p$인 원소가 존재한다.
\begin{enumerate}
\item $R$이 정역인 경우,
\item $R$이 뇌터환인 경우, 또는
\item $R$이 축약환이고 유한 개의 극소 소 이데알을 갖는 경우.
\end{enumerate}
\end{lemma}

\begin{proof}
$R$이 정역이면 $R \subset R_\mathfrak p$이므로 $f = 1$로 충분하다.
$R$이 뇌터환이면 $R \to R_\mathfrak p$의 핵 $I$는 유한 생성
이데알이고, $IR_f = 0$이 되도록 $f \in R$, $f \not \in \mathfrak p$를
찾을 수 있다. 이 $f$에 대해 $R_f \to R_\mathfrak p$는 단사이므로
$f$가 조건을 만족한다. $R$이 축약환이고 극소 소 이데알
$\mathfrak p_1, \ldots, \mathfrak p_n$을 유한 개 갖는다면
$f \in \bigcap_{\mathfrak p_i \not \subset \mathfrak p} \mathfrak p_i$,
$f \not \in \mathfrak p$인 것을 택할 수 있다. 실제로
$\mathfrak{p}_i\not\subset \mathfrak{p}$이면
$f_i \in \mathfrak{p}_i$, $f_i \not\in \mathfrak{p}$인 원소가
존재하며 $f = \prod f_i$가 조건을 만족한다. 이 $f$에 대해
$R_f \subset R_\mathfrak p$인 것은 $R_f$의 극소 소 이데알들이
$R_\mathfrak p$의 극소 소 이데알들에 대응하기 때문이다. 세부 사항은
생략하고 Lemma \ref{algebra-lemma-reduced-ring-sub-product-fields}를 적용한다.
\end{proof}

\begin{lemma}
\label{algebra-lemma-surjective-endo-noetherian-ring-is-iso}
뇌터환의 모든 전사 자기준동형은 동형이다.
\end{lemma}

\begin{proof}
$f : R \to R$가 그러한 자기준동형이지만 단사가 아니라고 하자. 그러면
$$
\Ker(f) \subset \Ker(f \circ f) \subset
\Ker(f \circ f \circ f) \subset \ldots
$$
는 이데알들의 엄격히 증가하는 사슬이 된다.
\end{proof}

\section{국소 멱영 이데알}
\label{algebra-section-locally-nilpotent}

\noindent
정의는 다음과 같다.

\begin{definition}
\label{algebra-definition-locally-nilpotent-ideal}
$R$을 환이라 하자. $I \subset R$를 이데알이라 하자.
$I$의 모든 $x \in I$에 대해 어떤 $n \in \mathbf{N}$이 존재하여
$x^n = 0$이면 $I$를 {\it 국소 멱영}이라고 한다. 어떤
$n \in \mathbf{N}$이 존재하여 $I^n = 0$이면 {\it 멱영}이라고
한다.
\end{definition}

\begin{example}
\label{algebra-example-locally-nilpotent-not-nilpotent}
$R = k[x_n | n \in \mathbf{N}]$를 체 $k$ 위에서 무한히 많은
변수를 갖는 다항식환이라 하자. $I$를 원소 $x_n^n$ ($n \in \mathbf{N}$)들이
생성하는 이데알이라 하고 $S = R/I$라 하자.
그러면 $J \subset S$는 $x_n$의 상들, $n \in \mathbf{N}$이 생성하는
이데알로서 국소 멱영이지만 멱영은 아니다. 실제로 멱영원들의
$S$-선형결합은 멱영원이므로, $J$가 국소 멱영임을 보이려면 그
생성원들이 모두 멱영원임을 관찰하는 것으로 충분하다 (이는 분명하다).
한편 각 $n \in \mathbf{N}$에 대해 $x_{n + 1}^n \not \in I$이므로
$J^n \not = 0$이다. 따라서 $J$는 멱영이 아니다.
\end{example}

\begin{lemma}
\label{algebra-lemma-locally-nilpotent}
$R \to R'$를 환 사상이라 하고 $I \subset R$를 국소 멱영 이데알이라
하자. 그러면 $IR'$는 $R'$의 국소 멱영 이데알이다.
\end{lemma}

\begin{proof}
멱영원 $x, y \in R'$에 대해서는 $x + y$도 멱영원이라는 사실에서
따른다. 즉 $x^n = 0$이고 $y^m = 0$이면
$(x + y)^{n + m - 1} = 0$이다.
\end{proof}

\begin{lemma}
\label{algebra-lemma-locally-nilpotent-unit}
$R$을 환이라 하고 $I \subset R$를 국소 멱영 이데알이라 하자.
$R$의 원소 $x$가 가역원일 필요충분조건은 $R/I$에서 $x$의 상이
가역원인 것이다.
\end{lemma}

\begin{proof}
$x$가 $R$에서 가역원이면 그 상은 분명히 $R/I$에서 가역원이다.
역을 증명하면 된다.

$R/I$에서 $y \in R$의 상이 $x$의 상의 역원이라고 하자. 그러면
어떤 $z \in I$에 대해 $xy = 1 - z$이다. 이는 $xR$을 법으로
$1\equiv z$라는 뜻이다. $z$는 국소 멱영 이데알 $I$에 속하므로
충분히 큰 $N$에 대해 $z^N = 0$이다. 따라서 $xR$을 법으로
$1 = 1^N \equiv z^N = 0$이다. 다시 말해 $x$는 $1$을 나누므로
가역원이다.
\end{proof}

\begin{lemma}
\label{algebra-lemma-Noetherian-power}
\begin{slogan}
뇌터환의 이데알은 그 이데알의 각 원소가 멱영원이면 멱영이다.
\end{slogan}
$R$을 뇌터환이라 하자. $I, J$를 $R$의 이데알이라 하자.
$J \subset \sqrt{I}$라고 가정하자. 그러면 어떤 $n$에 대해
$J^n \subset I$이다. 특히 뇌터환에서는 ``국소 멱영 이데알''과
``멱영 이데알''의 개념이 일치한다.
\end{lemma}

\begin{proof}
$J = (f_1, \ldots, f_s)$라 하자. 가정에 의해
$f_i^{d_i} \in I$이다. $n = d_1 + d_2 + \ldots + d_s + 1$로
두자.
\end{proof}

\begin{lemma}
\label{algebra-lemma-lift-idempotents}
$R$을 환이라 하고 $I \subset R$를 국소 멱영 이데알이라 하자.
그러면 $R \to R/I$는 멱등원들 사이에 전단사 대응을 유도한다.
\end{lemma}

\begin{proof}[보조정리 \ref{algebra-lemma-lift-idempotents}의 첫 번째 증명]
$I$가 국소 멱영이므로 모든 소 이데알에 포함된다. 따라서
$\Spec(R/I) = V(I) = \Spec(R)$이다. 그러므로 이 보조정리는
Lemma \ref{algebra-lemma-disjoint-decomposition}에서 따른다.
\end{proof}

\begin{proof}[보조정리 \ref{algebra-lemma-lift-idempotents}의 두 번째 증명]
$\overline{e} \in R/I$가 멱등원이라고 하자. $\overline{e}$를
$R$의 멱등원으로 들어 올려야 한다.

\medskip\noindent
먼저 $\overline{e}$의 임의의 올림 $f \in R$를 택하고
$x = f^2 - f$로 두자. 그러면 $x \in I$이므로 $x$는 멱영원이다
($I$가 국소 멱영이기 때문이다). 이제 $J$를 $x$가 생성하는 $R$의
이데알이라 하자. $J$는 단지 국소 멱영인 것뿐 아니라 멱영이다.
이는 멱영원 $x$가 생성하기 때문이다.

\medskip\noindent
이제 어떤 $k \geq 1$에 대해 $\overline{e}$의 올림
$e \in R$를 찾아 $e^2 - e \in J^k$가 되었다고 하자. 다음과 같이
두자.
$e' = e - (2e - 1)(e^2 - e) = 3e^2 - 2e^3$.
이는 $\overline{e}$의 또 다른 올림이다 (멱등원인
$\overline{e}$로부터 $e^2 - e \in I$이기 때문이다). 그러면
$$
(e')^2 - e' = (4e^2 - 4e - 3)(e^2 - e)^2 \in J^{2k}
$$
이다. 이는 간단한 계산으로 확인된다.

\medskip\noindent
따라서 $e^2 - e \in J^k$인 $\overline{e}$의 올림 $e$에서
시작하여, $(e')^2 - e' \in J^{2k}$인 $\overline{e}$의 올림 $e'$를
얻었다. 이 방법으로 근사를 계속 개선할 수 있다
($k = 1$에 대해 조건을 만족하는 $e = f$에서 시작한다).
결국 $J^k = 0$인 단계에 도달한다. 그 단계에서
$e^2 - e \in J^k = 0$인 $\overline{e}$의 올림 $e$를 얻으므로,
이 $e$는 멱등원이다.

\medskip\noindent
따라서 $\overline{e} \in R/I$인 임의의 멱등원 $\overline{e}$가
$R$의 멱등원으로 올라감을 보였다. 이제 이 올림의 유일성을 증명해야 한다. 실제로
두 올림 $e_1$, $e_2$를 택하자. $e_1 = e_2$임을 보여야 한다.

\medskip\noindent
$e_1$과 $e_2$의 정의에 의해 $e_1 \equiv e_2 \mod I$이고,
$e_1$과 $e_2$는 모두 멱등원이다. $e_1 \equiv e_2 \mod I$로부터
$e_1 - e_2 \in I$임을 알 수 있으므로, $e_1 - e_2$는 멱영원이다
($I$가 국소 멱영이기 때문이다). 간단한 계산
($e_1$과 $e_2$의 멱등성을 사용한다)으로
$(e_1 - e_2)^3 = e_1 - e_2$임을 얻는다. 이를 사용하고 귀납하면,
임의의 양의 홀수 $k$에 대해 $(e_1 - e_2)^k = e_1 - e_2$이다.
충분히 큰 모든 $k$는 $(e_1 - e_2)^k = 0$을 만족한다
($e_1 - e_2$가 멱영원이므로). 따라서 $e_1 - e_2 = 0$이고
$e_1 = e_2$이다. 증명이 끝난다.
\end{proof}

\begin{lemma}
\label{algebra-lemma-lift-idempotents-noncommutative}
$A$를 비가환일 수도 있는 대수라 하자. $x = e^2 - e$가 멱영원인
$e \in A$를 택하자. 다음 꼴의 멱등원이 존재한다.
$e' = e + x(\sum a_{i, j}e^ix^j) \in A$
여기서 $a_{i, j} \in \mathbf{Z}$이다.
\end{lemma}

\begin{proof}
$R_n = \mathbf{Z}[e]/((e^2 - e)^n)$이라는 환을 생각하자.
각 $R_n$에 대해 결과를 증명할 수 있으면 보조정리가 따른다는 것은
분명하다. $R_n$에서 이데알 $I = (e^2 - e)$를 택하고 Lemma
\ref{algebra-lemma-lift-idempotents}를 적용한다.
\end{proof}

\begin{lemma}
\label{algebra-lemma-lift-nth-roots}
$R$을 환이라 하고 $I \subset R$를 국소 멱영 이데알이라 하자.
$R/I$에서 가역인 정수 $n \geq 1$을 택하자.
\begin{enumerate}
\item $n$제곱 사상 $1 + I \to 1 + I$, $1 + x \mapsto (1 + x)^n$은
전단사이다.
\item $R$의 가역원이 $n$제곱이라는 것은 $R/I$에서 그 상이
$n$제곱이라는 것과 동치이다.
\end{enumerate}
\end{lemma}

\begin{proof}
$a \in R$를 가역원이라 하고, $R/I$에서 그 상이 $b^n$의 상과
같다고 하자 ($b \in R$). $b$는 가역원이다
(Lemma \ref{algebra-lemma-locally-nilpotent-unit}). 따라서 어떤 $x \in I$에
대해 $ab^{-n} = 1 + x$이고, (1)에 의해 $ab^{-n} = c^n$이다.
그러므로 (1)로부터 (2)가 따른다.

\medskip\noindent
(1)의 증명. 이는 $1 + x \mapsto (1 + x)^n$ 사상에 역사가
존재하기 때문이다. 구체적으로 $1 + x$를 다음으로 보내는 사상을
생각할 수 있다.
\begin{align*}
(1 + x)^{1/n}
& =
1 + {1/n \choose 1}x +
{1/n \choose 2}x^2 +
{1/n \choose 3}x^3 + \ldots \\
& =
1 + \frac{1}{n} x + \frac{1 - n}{2n^2}x^2 +
\frac{(1 - n)(1 - 2n)}{6n^3}x^3 + \ldots
\end{align*}
초등 미적분에서와 같다. $x^k = 0$인 $k \gg 0$이므로 급수는
유한하고, 각 계수 ${1/n \choose k} \in \mathbf{Z}[1/n]$이므로
이는 의미가 있다 (세부 사항은 생략한다. Lemma
\ref{algebra-lemma-locally-nilpotent-unit}에 의해 $n$이 $R$에서 가역임을
관찰하라).
\end{proof}

\section{호기심}
\label{algebra-section-curiosity}

\noindent
보조정리 \ref{algebra-lemma-disjoint-implies-product}는 어떤 이데알 $I \subset R$에 대해
$V(I)$가 열린 집합일 때 무슨 일이 일어나는지를 설명한다.
그렇다면 $\Spec(S^{-1}R)$가 $\Spec(R)$에서 닫힌 집합이면 어떠한가?
다음 두 보조정리가 부분적인 답을 준다. 자세한 내용은
절 \ref{algebra-section-pure-ideals}을(를) 참조하라.

\begin{lemma}
\label{algebra-lemma-invert-closed-quotient}
$R$을 환이라 하자. $S \subset R$를 곱셈 부분집합이라 하자.
사상 $\Spec(S^{-1}R) \to \Spec(R)$의 상이 닫혀 있다고 하자.
그러면 어떤 이데알 $I \subset R$에 대해 $S^{-1}R \cong R/I$이다.
\end{lemma}

\begin{proof}
$I = \Ker(R \to S^{-1}R)$라 두면 $V(I)$가 상을 포함한다.
어떤 이데알 $I' \subset R$에 대해 상이 닫힌 부분집합
$V(I') \subset \Spec(R)$라고 하자. 그러면 $V(I') \subset V(I)$이다.
$f \in I'$이면 $f/1 \in S^{-1}R$가 모든 소 아이디얼에 포함됨을 알 수 있다.
따라서 어떤 $n \geq 1$에 대해 $f^n$은 $S^{-1}R$에서 영원으로 사상된다
(보조정리 \ref{algebra-lemma-Zariski-topology}).
그러므로 $V(I') = V(I)$이다.
이로부터 모든 $g \in S$가 $I$를 법으로 가역임을 알 수 있다.
따라서 환 사상 $R/I \to S^{-1}R$ 및 $S^{-1}R \to R/I$를 얻는다.
첫 번째 사상은 $I$의 선택에 의해 단사이다.
두 번째 사상은 $S^{-1}R \to S^{-1}(R/I) = R/I$이고,
국소화가 완전하므로 그 핵은 $S^{-1}I$이다.
$S^{-1}I = 0$이므로 두 번째 사상도 단사임을 알 수 있다.
따라서 $S^{-1}R \cong R/I$이다.
\end{proof}

\begin{lemma}
\label{algebra-lemma-invert-closed-split}
$R$을 환이라 하자. $S \subset R$를 곱셈 부분집합이라 하자.
사상 $\Spec(S^{-1}R) \to \Spec(R)$의 상이 닫혀 있다고 하자.
$R$이 뇌터환이거나, $\Spec(R)$가 뇌터 위상 공간이거나,
$S$가 모노이드로서 유한 생성이면, 어떤 환 $R'$에 대해
$R \cong S^{-1}R \times R'$이다.
\end{lemma}

\begin{proof}
보조정리 \ref{algebra-lemma-invert-closed-quotient}에 의해 어떤 이데알 $I \subset R$에 대해
$S^{-1}R \cong R/I$이다. 보조정리 \ref{algebra-lemma-disjoint-implies-product}에 의해
$V(I)$가 열려 있음을 보이는 것으로 충분하다.
$R$이 뇌터환이면 $\Spec(R)$는 뇌터 위상 공간이다
(보조정리 \ref{algebra-lemma-Noetherian-topology}).
$\Spec(R)$가 뇌터 위상 공간이면 그 여집합
$\Spec(R) \setminus V(I)$는 준콤팩트이다
(위상론, 보조정리 \ref{topology-lemma-Noetherian-quasi-compact}).
따라서 어떤 유한한 $f_1, \ldots, f_n \in I$가 존재하여
$V(I) = V(f_1, \ldots, f_n)$이다.
각 $f_i$가 $S^{-1}R$에서 영으로 사상되므로, 어떤 $g \in S$가 존재하여
$i = 1, \ldots, n$에 대해 $gf_i = 0$이다. 따라서 원하는 대로
$D(g) = V(I)$이다.
$S$가 모노이드로서 유한 생성인 경우, $S$가
$g_1, \ldots, g_m$으로 생성된다고 하자. 그러면
$S^{-1}R \cong R_{g_1 \ldots g_m}$이고
$V(I) = D(g_1 \ldots g_m)$임을 얻는다.
\end{proof}
\section{힐베르트 영점정리}
\label{algebra-section-nullstellensatz}

\begin{theorem}[힐베르트 영점정리]
\label{algebra-theorem-nullstellensatz}
$k$를 체라 하자.
\begin{enumerate}
\item
\label{algebra-item-finite-kappa}
임의의 극대 이데알 $\mathfrak m \subset k[x_1, \ldots, x_n]$에 대해
체확대 $\kappa(\mathfrak m)/k$는 유한하다.
\item
\label{algebra-item-polynomial-ring-Jacobson}
임의의 근기 이데알 $I \subset k[x_1, \ldots, x_n]$은
이를 포함하는 극대 이데알들의 교집합이다.
\end{enumerate}
같은 결과가 유한형 $k$-대수에서도 성립한다.
\end{theorem}

\begin{proof}
이 정리의 부분 (\ref{algebra-item-finite-kappa})을 다항식 대수
$k[x_1, \ldots, x_n]$인 경우에 증명하면 충분하다. 임의의 유한 생성
$k$-대수는 이러한 다항식 대수의 몫이기 때문이다.
이를 $n$에 대한 귀납법으로 증명한다. $n = 0$인 경우는 자명하다.
$\mathfrak m$이 $k[x_1, \ldots, x_n]$의 극대 이데알이라고 하자.
$\mathfrak p \subset k[x_n]$를 $\mathfrak m$과 $k[x_n]$의 교집합이라 하자.

\medskip\noindent
$\mathfrak p \not = (0)$이면, $\mathfrak p$는 극대이고
기약인 단위원 다항식 $P$ 하나로 생성된다
($k[x_n]$에서 유클리드 알고리즘을 사용하기 때문이다). 그러면
$k' = k[x_n]/\mathfrak p$는 $k$의 유한 체확대이며
$\kappa(\mathfrak m)$에 포함된다. 이 경우 다음 전사 사상을 얻는다.
$$
k'[x_1, \ldots, x_{n-1}]
\to
k'[x_1, \ldots, x_n] =
k' \otimes_k k[x_1, \ldots, x_n]
\longrightarrow
\kappa(\mathfrak m)
$$
따라서 귀납가정에 의해 $\kappa(\mathfrak m)$은 $k'$의 유한
체확대이다. 그러므로 $\kappa(\mathfrak m)$은 $k$ 위에서도 유한하다.

\medskip\noindent
$\mathfrak p = (0)$이면 환 확대
$k[x_n] \subset k[x_1, \ldots, x_n]/\mathfrak m$를 생각한다.
이는 유한 생성 환 확대이므로, 보조정리
\ref{algebra-lemma-obvious-Noetherian} 및
\ref{algebra-lemma-Noetherian-finite-type-is-finite-presentation}에 의해
유한 표시이다. 따라서
$\Spec(k[x_1, \ldots, x_n]/\mathfrak m)$의
$\Spec(k[x_n])$에서의 상은 정리
\ref{algebra-theorem-chevalley}에 의해 구성가능하다. 상은
$(0)$을 포함하므로, 어떤 $f\in k[x_n]$가 0이 아니어서
표준 열린집합 $D(f)$를 포함한다. $D(f)$는 분명 무한집합이므로
$k[x_1, \ldots, x_n]/\mathfrak m$가 체라는 가정과 모순이다
(따라서 그 스펙트럼은 한 점으로 이루어진다).

\medskip\noindent
(\ref{algebra-item-polynomial-ring-Jacobson})의 증명. $R$을 $k$ 위에서
유한형인 환이라 하고, $I \subset R$를 근기 이데알이라 하자.
$f \in R$, $f \not \in I$를 택하자. $I \subset \mathfrak m$이고
$f \not \in \mathfrak m$인 극대 이데알 $\mathfrak m \subset R$을
찾아야 한다. $(R/I)_f$는 영이 아니다. 실제로 이 환에서 $1 = 0$이라면
$f^n \in I$를 뜻하고, $I$가 근기 이데알이므로 $f \in I$가 되어
$f$에 대한 가정과 모순이다.
그러므로 보조정리 \ref{algebra-lemma-Zariski-topology}에 의해
극대 이데알 $\mathfrak m'$을 $(R/I)_f$에서 택할 수 있다.
$R$에서 $\mathfrak m'$의 역상인 $\mathfrak m \subset R$을 택하자.
그러면 $I \subset \mathfrak m$이고 $f \not \in \mathfrak m$이다.
$\mathfrak m$이 $R$의 극대 이데알임을 보이면 끝난다. 분명히
$$
k \subset R/\mathfrak m \subset \kappa(\mathfrak m').
$$
부분 (\ref{algebra-item-finite-kappa})에 의해 체확대
$\kappa(\mathfrak m')/k$는 유한하다. 따라서 보조정리
\ref{fields-lemma-subalgebra-algebraic-extension-field}에 의해
$R/\mathfrak m$은 체이다. 그러므로 $\mathfrak m$은 극대이고
증명이 끝난다.
\end{proof}
\begin{lemma}
\label{algebra-lemma-field-finite-type-over-domain}
$R$을 환이라 하자. $K$를 체라 하자.
$R \subset K$이고 $K$가 $R$ 위에서 유한형이면,
어떤 $f \in R$가 존재하여 $R_f$는 체이고,
$K/R_f$는 유한 체확대이다.
\end{lemma}

\begin{proof}
보조정리 \ref{algebra-lemma-characterize-image-finite-type}에 의해
공집합이 아닌 열린집합 $U \subset \Spec(R)$가
$\Spec(K) \to \Spec(R)$의 상 $\{(0)\}$에 포함되어 존재한다.
$f \in R$, $f \not = 0$을 택하여 $D(f) \subset U$, 즉
$D(f) = \{(0)\}$가 되게 하자. 그러면 $R_f$는 스펙트럼이
정확히 한 점을 갖는 정역이고, $R_f$는 체이다. 이제 $K$는 체 $R_f$ 위에서
유한 생성 대수이고, 힐베르트 영점정리
(정리 \ref{algebra-theorem-nullstellensatz})에 의해 $R_f$의 유한 체확대이다.
\end{proof}
\section{야콥슨 환}
\label{algebra-section-ring-jacobson}

\noindent
환 $R$을 잡자. $\Spec(R)$의 닫힌 점들은 $R$의 극대 이데알이다.
대수기하학에서 자연스럽게 나타나는 환들은 극대 이데알을 많이 갖는
경우가 흔하다. 예를 들어 체 또는 $\mathbf{Z}$ 위의 유한형 대수들이
그렇다. 이러한 환들이 야콥슨 환의 예임을 보일 것이다.

\begin{definition}
\label{algebra-definition-ring-jacobson}
환 $R$을 환이라 하자. 모든 근기 이데알
$I$가 이를 포함하는 극대 이데알들의 교집합이면
$R$을 {\it 야콥슨 환}이라 한다.
\end{definition}

\begin{lemma}
\label{algebra-lemma-finite-type-field-Jacobson}
체 위에서 유한형인 임의의 대수는 야콥슨 환이다.
\end{lemma}

\begin{proof}
이는 정리 \ref{algebra-theorem-nullstellensatz} 및 정의
\ref{algebra-definition-ring-jacobson}에서 따른다.
\end{proof}

\begin{lemma}
\label{algebra-lemma-jacobson-prime}
환 $R$을 환이라 하자. $R$의 모든 소 이데알이
그것을 포함하는 극대 이데알들의 교집합이면,
$R$은 야콥슨 환이다.
\end{lemma}

\begin{proof}
이는 모든 근기 이데알 $I \subset R$이
이를 포함하는 소 이데알들의 교집합이라는 사실에서
즉시 따른다. 보조정리 \ref{algebra-lemma-Zariski-topology}를
참조하라.
\end{proof}

\begin{lemma}
\label{algebra-lemma-jacobson}
환 $R$이 야콥슨 환일 필요충분조건은 $\Spec(R)$가
야콥슨임이다. 위상론, 정의
\ref{topology-definition-space-jacobson}을 참조하라.
\end{lemma}

\begin{proof}
$R$이 야콥슨 환이라고 하자. $Z \subset \Spec(R)$를
닫힌 부분집합이라 하자. $Z$에서 닫힌 점들의 집합이
$Z$에서 조밀함을 보여야 한다. $U \subset \Spec(R)$를
$U \cap Z$가 공집합이 아닌 열린집합이라 하자.
$Z \cap U$가 $\Spec(R)$의 닫힌 점을 포함함을 보여야 한다.
표준 열린집합들이 $\Spec(R)$의 위상에 대한 기저를 이루므로
$U = D(f)$라고 가정해도 된다. 보조정리
\ref{algebra-lemma-Zariski-topology}에 따라 $Z = V(I)$이고
$I$가 근기 이데알이라고 가정해도 된다. 또한
$f \not \in I$임을 알 수 있다. 가정에 의해 $I \subset \mathfrak m$이고
$f \not\in \mathfrak m$인 극대 이데알
$\mathfrak m \subset R$이 존재한다. 따라서
원하는 대로 $\mathfrak m \in D(f) \cap V(I) = U \cap Z$이다.

\medskip\noindent
반대로 $\Spec(R)$가 야콥슨이라고 하자.
$I \subset R$를 근기 이데알이라 하자. $I$를 포함하는
극대 이데알들의 교집합을
$J = \cap_{I \subset \mathfrak m} \mathfrak m$이라고 하자.
분명히 $J$는 근기 이데알이고, $V(J) \subset V(I)$이며,
$V(J)$는 $V(I)$의 닫힌 점을 모두 포함하는 $V(I)$의
가장 작은 닫힌 부분집합이다. 가정에 의해 $V(J) = V(I)$이다.
그런데 보조정리 \ref{algebra-lemma-Zariski-topology}에 따르면
자리스키 닫힌 집합과 근기 이데알 사이에 전단사가 있으므로,
원하는 대로 $I = J$이다.
\end{proof}
\begin{lemma}
\label{algebra-lemma-characterize-jacobson}
환 $R$을 환이라 하자. $R$이 야콥슨 환이 아니면 다음 조건을
만족하는 소 이데알 $\mathfrak p \subset R$와 원소 $f \in R$가 존재한다.
\begin{enumerate}
\item $\mathfrak p$는 극대 이데알이 아니다.
\item $f \not \in \mathfrak p$이다.
\item $V(\mathfrak p) \cap D(f) = \{\mathfrak p\}$이다.
\item $(R/\mathfrak p)_f$는 체이다.
\end{enumerate}
반대로 $R$이 야콥슨 환이면, (1)과 (2)를 만족하는 임의의 쌍
$(\mathfrak p, f)$에 대해 집합 $V(\mathfrak p) \cap D(f)$는 무한하다.
\end{lemma}

\begin{proof}
$R$이 야콥슨 환이 아니라고 하자. 보조정리
\ref{algebra-lemma-jacobson}에 의해 이는 닫힌 부분집합
$T \subset \Spec(R)$가 존재하여 그 닫힌 점들의 집합
$T_0 \subset T$가 $T$에서 조밀하지 않다는 뜻이다.
$T_0 \subset V(f)$이지만 $T \not \subset V(f)$가 되도록
$f \in R$을 택하자. $T = V(I)$이면
$T \cap D(f)$는 $\Spec((R/I)_f)$와 위상동형임에 유의하라.
보조정리 \ref{algebra-lemma-spec-closed} 및
\ref{algebra-lemma-standard-open}을 참조하라. 모든 환은 극대 이데알을
가지므로 (보조정리 \ref{algebra-lemma-Zariski-topology})
공간 $T \cap D(f)$의 닫힌 점 $t$를 택할 수 있다.
그러면 $t$는 극대가 아닌 소 이데알
$\mathfrak p \subset R$에 대응한다($t \not \in T_0$이기 때문이다).
따라서 (1)이 성립한다. 구성에 의해 $f \not \in \mathfrak p$이므로
(2)도 성립한다. $t$가 $T \cap D(f)$의 닫힌 점이므로
$V(\mathfrak p) \cap D(f) = \{\mathfrak p\}$임을 알 수 있고,
즉 (3)이 성립한다. 따라서 $(R/\mathfrak p)_f$는 스펙트럼이
한 점인 정역이고, 따라서 (4)가 성립한다
(예를 들어 보조정리 \ref{algebra-lemma-characterize-local-ring} 및
\ref{algebra-lemma-minimal-prime-reduced-ring}을 함께 사용하라).

\medskip\noindent
반대로 $R$이 야콥슨 환이고 $(\mathfrak p, f)$가 (1)과 (2)를
만족한다고 하자. 만일
$V(\mathfrak p) \cap D(f) =
\{\mathfrak p, \mathfrak q_1, \ldots, \mathfrak q_t\}$라면
$\mathfrak p \not = \mathfrak q_i$이므로 어떤 원소
$g \in R$가 존재하여 $g \not \in \mathfrak p$이지만 모든 $i$에 대해
$g \in \mathfrak q_i$이다. 따라서
$V(\mathfrak p) \cap D(fg) = \{\mathfrak p\}$가 되는데,
$\Spec(R)$가 야콥슨 위상 공간이므로 $\Spec(R)$의 각 국소 폐쇄
부분집합이 적어도 하나의 닫힌 점을 포함한다는 사실에 모순이다.
\end{proof}

\begin{lemma}
\label{algebra-lemma-pid-jacobson}
환 $\mathbf{Z}$는 야콥슨 환이다. 더 일반적으로 다음 조건을
만족하는 환 $R$을 잡자.
\begin{enumerate}
\item $R$은 정역이다.
\item $R$은 뇌터환이다.
\item 모든 영이 아닌 소 이데알은 극대 이데알이다.
\item $R$은 무한히 많은 극대 이데알을 갖는다.
\end{enumerate}
그러면 $R$은 야콥슨 환이다.
\end{lemma}

\begin{proof}
$R$이 (1), (2), (3), (4)를 만족한다고 하자. 이 명제는
$(0) = \bigcap_{\mathfrak m \subset R} \mathfrak m$임을 뜻한다.
$R$이 무한히 많은 극대 이데알을 가지므로, 임의의 영이 아닌
$x \in R$이 유한 개 이하의 극대 이데알에 포함됨을 보이면 충분하다.
다시 말해 $V(x)$가 유한함을 보이면 된다.
보조정리 \ref{algebra-lemma-spec-closed}에 의해 $V(x)$는
$\Spec(R/xR)$와 위상동형이다. 가정 (3)에 의해 $R/xR$의
모든 소 이데알은 극소이고, 따라서 $\Spec(R/xR)$의 기약 성분에
대응한다(보조정리 \ref{algebra-lemma-irreducible}).
$R/xR$가 뇌터환이므로 위상 공간 $\Spec(R/xR)$는
뇌터이다(보조정리 \ref{algebra-lemma-Noetherian-topology}) 그리고
유한 개의 기약 성분을 갖는다(위상론, 보조정리
\ref{topology-lemma-Noetherian}). 따라서 원하는 대로 $V(x)$는
유한하다.
\end{proof}
\begin{example}
\label{algebra-example-infinite-product-fields-jacobson}
무한 집합 $A$를 잡자.
각 $\alpha \in A$에 대해 체 $k_\alpha$를 택하자.
$R = \prod_{\alpha\in A} k_\alpha$가 야콥슨 환임을 보이자.
먼저 임의의 원소 $f \in R$가
$f = ue$ 꼴이고, $u \in R$은 단위원이며 $e\in R$은
멱등원임에 유의하라(세부 사항은 독자에게 맡긴다).
따라서 $D(f) = D(e)$이고, $R_f = R_e = R/(1-e)$는
$R$의 몫이다.
사실 이러한 성질을 갖는 임의의 환은 야콥슨 환이다.
구체적으로 $\mathfrak p \subset R$가 소 이데알이고
$f \in R$, $f \not \in \mathfrak p$라고 하자. $\mathfrak p \subset \mathfrak m$이고
$f \not\in \mathfrak m$인 $R$의 극대 이데알
$\mathfrak m$을 찾아야 한다.
$R_f$가 $R$의 몫이므로 $R_f$의 임의의 극대 이데알은
$f$를 포함하지 않는 $R$의 극대 이데알에 대응한다.
따라서 $\mathfrak p R_f$를 포함하는 $R_f$의 극대 이데알을
택하면 결과가 따른다.
\end{example}

\begin{example}
\label{algebra-example-not-jacobson}
유한 개의 극대 이데알 $\mathfrak m_i$, $i = 1, \ldots, n$을
갖는 정역 $R$은 체인 경우를 제외하면 야콥슨 환이 아니다.
실제로 이 경우 $(0)$은 극대 이데알들의 교집합이 아니다.
$(0) \not =
\mathfrak m_1 \cap \mathfrak m_2 \cap \ldots \cap \mathfrak m_n
\supset \mathfrak m_1 \cdot \mathfrak m_2 \cdot \ldots
\cdot \mathfrak m_n \not = 0$이기 때문이다.
특히 이산 값매김환 또는 적어도 두 소 이데알을 갖는 국소환은
야콥슨 환이 아니다.
\end{example}

\begin{lemma}
\label{algebra-lemma-finite-residue-extension-closed}
환 사상 $R \to S$를 잡자.
$\mathfrak m \subset R$를 극대 이데알이라 하자.
$\mathfrak q \subset S$를 $\mathfrak m$ 위에 놓이는 소 이데알이라
하고, $\kappa(\mathfrak q)/\kappa(\mathfrak m)$가 대수적 체확대라고 하자.
그러면 $\mathfrak q$는 $S$의 극대 이데알이다.
\end{lemma}

\begin{proof}
다음 도식을 생각하자.
$$
\xymatrix{
S \ar[r] & S/\mathfrak q \ar[r] & \kappa(\mathfrak q) \\
R \ar[r] \ar[u] & R/\mathfrak m \ar[u]
}
$$
$\kappa(\mathfrak m) \subset S/\mathfrak q \subset
\kappa(\mathfrak q)$임을 알 수 있다. 체확대
$\kappa(\mathfrak m) \subset \kappa(\mathfrak q)$가
대수적이므로 $\kappa(\mathfrak m)$과 $\kappa(\mathfrak q)$
사이에 있는 임의의 환은 체이다
(체론, 보조정리 \ref{fields-lemma-subalgebra-algebraic-extension-field}).
따라서 $S/\mathfrak q$는 체이고, 결과적으로
$\kappa(\mathfrak q)$와 같다.
\end{proof}

\begin{lemma}
\label{algebra-lemma-dimension}
$k$가 체이고 $V$가 $k$ 위의 영이 아닌 벡터 공간이라고 하자.
$V$의 차원(이는 기수이다)이 $k$의 기수보다 작다고 가정하자.
그러면 임의의 선형 연산자 $T : V \to V$에 대해
$P(t) \in k[t]$인 최고차항 계수가 1인 다항식이
존재하여 $P(T)$는 가역이 아니다.
\end{lemma}

\begin{proof}
그렇지 않다면 $V$는 체 $k(t)$ 위의 벡터 공간 구조를 상속한다.
그런데 $k(t)$의 $k$ 위 차원은 $k$의 기수 이상이다.
이는 예를 들어 원소
$\frac{1}{t - \lambda}$들이 $k$-선형 독립이라는 사실에서 따른다.
\end{proof}

\noindent
힐베르트 영점정리의 또 다른 형태를 보자.

\begin{theorem}
\label{algebra-theorem-uncountable-nullstellensatz}
체 $k$를 잡자. $S$를 원소 $\{x_i\}_{i \in I}$에 의해
$k$ 위에서 생성되는 $k$-대수라 하자. $I$의 기수가
$k$의 기수보다 작다고 가정하자. 그러면
\begin{enumerate}
\item 모든 극대 이데알 $\mathfrak m \subset S$에 대해 체확대
$\kappa(\mathfrak m)/k$는 대수적이다.
\item $S$는 야콥슨 환이다.
\end{enumerate}
\end{theorem}

\begin{proof}
$I$가 유한하면 결과는 힐베르트 영점정리,
정리 \ref{algebra-theorem-nullstellensatz}에서 따른다. 이제 증명의
나머지 부분에서는 $I$가 무한하다고 가정한다. $S$가 이 환의
몫이므로, 변수 $x_i$들의 다항식 환에서 극대인
$\mathfrak m \subset k[\{x_i\}_{i \in I}]$에 대해 결과를
증명하면 충분하다.
$I$가 무한하므로 단항식들의 집합
$x_{i_1}^{e_1} \ldots x_{i_r}^{e_r}$, $i_1, \ldots, i_r \in I$
및 $e_1, \ldots, e_r \geq 0$의 기수는 $I$의 기수 이하이다.
$I \times \ldots \times I$의 기수가 $I$의 기수이고,
$\bigcup_{n \geq 0} I^n$의 기수도 같은 기수이기 때문이다.
($I$가 유한하면 이는 참이 아니며, 이 경우 이 증명은
$k$가 비가산일 때에만 작동한다.)
모순을 얻기 위해 $k$ 위에서 초월적인 $T \in \kappa(\mathfrak m)$를 택하자.
원소 $T$에 의한 곱셈으로 주어지는 $k$-선형 사상
$T : \kappa(\mathfrak m) \to \kappa(\mathfrak m)$는
모든 최고차항 계수가 1인 다항식 $P(t) \in k[t]$에 대해
$P(T)$가 가역이라는 성질을 갖는다.
또한 $\kappa(\mathfrak m)$의 $k$ 위 차원은 $I$의 기수 이하이다.
이는 $k$ 위의 벡터 공간
$k[\{x_i\}_{i \in I}]$의 몫이기 때문이다
(위에서 본 것처럼 그 차원은 $\# I$이다).
이는 보조정리 \ref{algebra-lemma-dimension}에 모순이다.

\medskip\noindent
$S$가 야콥슨 환임을 보이기 위해 다음과 같이 논증하자.
그렇지 않다고 하면 보조정리
\ref{algebra-lemma-characterize-jacobson}에 의해 소 이데알
$\mathfrak q \subset S$와 원소 $f \in S$, $f \not \in \mathfrak q$가
존재하여 $\mathfrak q$가 극대가 아니고
$(S/\mathfrak q)_f$가 체이다.
그런데 $(S/\mathfrak q)_f$는 많아야 $\# I + 1$개의 원소로
생성된다. 따라서 체확대 $(S/\mathfrak q)_f/k$는
대수적이다(증명의 첫 부분에 의해).
이는 $\kappa(\mathfrak q)$가 $k$의 대수적 체확대임을 뜻하므로,
\ref{algebra-lemma-finite-residue-extension-closed}에 의해
$\mathfrak q$는 극대 이데알이다. 이는 모순이며
증명을 끝낸다.
\end{proof}
\begin{lemma}
\label{algebra-lemma-base-change-Jacobson}
체 $k$를 잡자. $S$를 $k$-대수라 하자.
기수가 $S$의 기수보다 큰 임의의 체확대 $K/k$에 대해
다음이 성립한다.
\begin{enumerate}
\item $S_K$의 모든 극대 이데알 $\mathfrak m$에 대해
체 $\kappa(\mathfrak m)$은 $K$ 위에서 대수적이다.
\item $S_K$는 야콥슨 환이다.
\end{enumerate}
\end{lemma}

\begin{proof}
$K$의 기수가 $S$의 기수보다 크도록 $k \subset K$를 택하자.
$S$의 원소들이 $K$-대수 $S_K$를 생성하므로 정리
\ref{algebra-theorem-uncountable-nullstellensatz}를 적용할 수 있다.
\end{proof}

\begin{example}
\label{algebra-example-countable-trick-does-not-work}
앞의 정리 \ref{algebra-theorem-uncountable-nullstellensatz}의 증명에서
사용한 요령은 $k$가 가산체이고 $I$도 가산일 때에는
실제로 작동하지 않는다. $k$를 가산체라 하자. $x$를 변수라
하자. $k(x)$를 $x$의 유리함수체라 하자.
다항식 대수 $R = k[x, \{x_f\}_{f \in k[x]-\{0\}}]$를 생각하자.
$I = (\{fx_f - 1\}_{f\in k[x] - \{0\}})$라 하자. $I$는
$R$의 진부분 이데알임에 유의하라.
$I \subset \mathfrak m$인 극대 이데알을 택하자.
그러면 $k \subset R/\mathfrak m$은 $k(x)$와 동형이고,
$k$ 위에서 대수적이지 않다.
\end{example}

\begin{lemma}
\label{algebra-lemma-Jacobson-invert-element}
$R$을 야콥슨 환이라 하자. $f \in R$이라 하자. 환 $R_f$는
야콥슨 환이고, $R_f$의 극대 이데알들은 $f$를 포함하지 않는
$R$의 극대 이데알들에 대응한다.
\end{lemma}

\begin{proof}
위상론, 보조정리 \ref{topology-lemma-jacobson-inherited}에 의해
$D(f) = \Spec(R_f)$는 야콥슨이고, $D(f)$의 닫힌 점들은
$\Spec(R)$에서 $D(f)$에 속하는 닫힌 점들에 대응함을 알 수 있다.
따라서 보조정리 \ref{algebra-lemma-jacobson}에 의해 $R_f$는 야콥슨 환이다.
\end{proof}

\begin{example}
\label{algebra-example-localize-not-preserve-closed-points}
$R$이 야콥슨 환이 아닐 때 보조정리
\ref{algebra-lemma-Jacobson-invert-element}가 거짓임을 보이는 간단한
예를 보자. 환 $R = \mathbf{Z}_{(2)}$를 생각하자. 즉, 소 이데알
$(2)$에서 $\mathbf{Z}$를 국소화한 것이다. 원소 $2$에서 $R$을
국소화한 것은 $\mathbf{Q}$와 동형이다. 즉
$R_2 \cong \mathbf{Q}$이다. 분명히 사상
$R \to R_2$는 $\Spec(\mathbf{Q})$의 닫힌 점을
$\Spec(R)$의 일반점으로 보낸다.
\end{example}
\begin{example}
\label{algebra-example-infinite-localize-not-preserve-closed-points}
$R$이 야콥슨 환이지만 무한히 많은 원소에서 국소화하면
보조정리 \ref{algebra-lemma-Jacobson-invert-element}가 거짓임을 보이는
간단한 예를 보자.
즉, $R = \mathbf{Z}$라 하고 모든 영이 아닌 원소들의 집합에서
$R$의 국소화
$(R \setminus \{0\})^{-1}R \cong \mathbf{Q}$를 생각하자.
분명히 사상 $\mathbf{Z} \to \mathbf{Q}$는
$\Spec(\mathbf{Q})$의 닫힌 점을
$\Spec(\mathbf{Z})$의 일반점으로 보낸다.
\end{example}

\begin{lemma}
\label{algebra-lemma-Jacobson-mod-ideal}
$R$을 야콥슨 환이라 하자. $I \subset R$를 이데알이라 하자.
환 $R/I$는 야콥슨 환이고 $R/I$의 극대 이데알들은 $I$를
포함하는 $R$의 극대 이데알들에 대응한다.
\end{lemma}

\begin{proof}
증명은 보조정리 \ref{algebra-lemma-Jacobson-invert-element}의
증명과 같다.
\end{proof}

\begin{lemma}
\label{algebra-lemma-silly-jacobson}
$R$을 야콥슨 환이라 하자. $K$를 체라 하자. $R \subset K$이고
$K$가 $R$ 위에서 유한형이면, $R$은 체이고 $K/R$은
유한 체확대이다.
\end{lemma}

\begin{proof}
먼저 $R$이 정역임에 유의하라.
보조정리 \ref{algebra-lemma-field-finite-type-over-domain}에 의해
어떤 영이 아닌 $f \in R$에 대해 $R_f$는 체이고
$K/R_f$는 유한 체확대이다. 따라서 $(0)$은 $R_f$의
극대 이데알이고, 보조정리 \ref{algebra-lemma-Jacobson-invert-element}에 의해
$(0)$은 $R$의 극대 이데알임을 결론낼 수 있다.
\end{proof}

\begin{proposition}
\label{algebra-proposition-Jacobson-permanence}
$R$을 야콥슨 환이라 하자. $R \to S$를 유한형 환 사상이라 하자. 그러면
\begin{enumerate}
\item 환 $S$는 야콥슨 환이다.
\item 사상 $\Spec(S) \to \Spec(R)$는 닫힌 점을
닫힌 점으로 보낸다.
\item $\mathfrak m' \subset S$가 $\mathfrak m \subset R$ 위에 놓인 극대 이데알이면
체확대 $\kappa(\mathfrak m')/\kappa(\mathfrak m)$는
유한하다.
\end{enumerate}
\end{proposition}

\begin{proof}
$\mathfrak m' \subset S$를 극대 이데알이라 하고
$R \cap \mathfrak m' = \mathfrak m$이라 하자.
그러면 $R/\mathfrak m \to S/\mathfrak m'$는
보조정리 \ref{algebra-lemma-silly-jacobson}의 조건을 만족한다.
이는 보조정리 \ref{algebra-lemma-Jacobson-mod-ideal}에 따른 것이다.
따라서 $R/\mathfrak m$은 체이고
$\mathfrak m$은 극대 이데알이며, 유도된 잉여체 확대는 유한하다.
이로써 (2)와 (3)이 증명된다.  

\medskip\noindent
만일 $S$가 야콥슨 환이 아니라면, 보조정리 \ref{algebra-lemma-characterize-jacobson}에 의해
$S$의 극대가 아닌 소 이데알 $\mathfrak q$와
$g \in S$, $g \not\in \mathfrak q$가 존재하여
$(S/\mathfrak q)_g$가 체가 된다.
모순을 얻기 위해 $\mathfrak q$가 극대 이데알임을 보이자.
$\mathfrak p = \mathfrak q \cap R$라 하자. 그러면
$R/\mathfrak p \to (S/\mathfrak q)_g$는
보조정리 \ref{algebra-lemma-silly-jacobson}의 조건을 만족한다.
이는 보조정리 \ref{algebra-lemma-Jacobson-mod-ideal}에 따른 것이다.
따라서 $R/\mathfrak p$는 체이고 체확대
$\kappa(\mathfrak p) \to (S/\mathfrak q)_g = \kappa(\mathfrak q)$는
유한하므로 대수적이다. 그러면
보조정리 \ref{algebra-lemma-finite-residue-extension-closed}에 의해
$\mathfrak q$는 $S$의 극대 이데알이다. 이는 모순이다.
\end{proof}

\begin{lemma}
\label{algebra-lemma-corollary-jacobson}
$\mathbf{Z}$ 위의 모든 유한형 대수는 야콥슨 환이다.
\end{lemma}

\begin{proof}
보조정리 \ref{algebra-lemma-pid-jacobson}와
명제 \ref{algebra-proposition-Jacobson-permanence}를 결합하면 된다.
\end{proof}

\begin{lemma}
\label{algebra-lemma-image-finite-type-map-Jacobson-rings}
$R \to S$를 야콥슨 환들 사이의 유한형 환 사상이라 하자.
$X = \Spec(R)$와 $Y = \Spec(S)$로 나타내자.
유도된 스펙트럼 사상을 $f : Y \to X$로 쓰자.
$E \subset Y = \Spec(S)$를 구성가능 집합이라 하자.
위상 공간의 닫힌 점들의 집합을 아래첨자 ${}_0$로 나타내자.
\begin{enumerate}
\item $f(E)_0 = f(E_0) = X_0 \cap f(E)$이다.
\item 점 $\xi \in X$가 $f(E)$에 속할 필요충분조건은
$\overline{\{\xi\}} \cap f(E_0)$가 $\overline{\{\xi\}}$에서 조밀한 것이다.
\end{enumerate}
\end{lemma}

\begin{proof}
다음과 같은 연속사상들의 가환도식이 있다.
$$
\xymatrix{
E \ar[r] \ar[d] & Y \ar[d] \\
f(E) \ar[r] & X
}
$$
$x \in f(E)$가 $f(E)$에서 닫혀 있다고 하자. 그러면
$f^{-1}(\{x\})\cap E$는 공집합이 아니고 $E$에서 닫혀 있다.
두 포함
$$
f^{-1}(\{x\}) \cap E \subset E \subset Y
$$
에 위상론 보조정리 \ref{topology-lemma-jacobson-inherited}를 적용하면,
$Y$에서 닫힌 점 $y \in f^{-1}(\{x\}) \cap E$가 존재함을 알 수 있다.
달리 말하면, $x$로 가는 $y \in Y_0$와 $y \in E_0$가 존재한다.
따라서 $x \in f(E_0)$이다.
이로써 $f(E)_0 \subset f(E_0)$임을 증명했다.
명제 \ref{algebra-proposition-Jacobson-permanence}에 의해
$f(E_0) \subset X_0 \cap f(E)$이다.
포함 $X_0 \cap f(E) \subset f(E)_0$는 자명하다.
이로써 첫 번째 주장이 증명된다.

\medskip\noindent
$\xi \in f(E)$라 하자.
보조정리 \ref{algebra-lemma-characterize-image-finite-type}에 따르면
집합 $f(E) \cap \overline{\{\xi\}}$는
$\overline{\{\xi\}}$의 조밀한 열린부분집합을 포함한다.
$X$가 야콥슨 공간이므로
$f(E) \cap \overline{\{\xi\}}$는 닫힌 점들의 조밀한 집합을 포함한다.
위상론 보조정리 \ref{topology-lemma-jacobson-inherited}를 보라.
따라서 보조정리의 (1)로 결론을 얻는다.

\medskip\noindent
반대로, $\overline{\{\xi\}} \cap f(E_0)$가
$\overline{\{\xi\}}$에서 조밀하다고 하자.
보조정리 \ref{algebra-lemma-constructible-is-image}에 의해,
$E$가 $Y' := \Spec(S') \to Y$의 상이 되도록 하는
유한 표시 환 사상 $S \to S'$가 존재한다.
환 사상 $S \to S'$에 보조정리의 첫 번째 부분을 적용하면
$E_0$는 $Y'_0$의 상이다. 따라서 $S$를 $S'$로 바꾸어
$E = Y$라고 가정할 수 있다. $\xi$가
$\mathfrak p \subset R$에 대응한다고 하자. 도식
$$
\xymatrix{
S \ar[r] & S/\mathfrak p S \\
R \ar[r] \ar[u] & R/\mathfrak p \ar[u]
}
$$
을 생각하자. 이 도식과 $V(\mathfrak p)$에서
$f(Y_0) \cap V(\mathfrak p)$가 조밀하다는 사실은
사상 $R/\mathfrak p \to S/\mathfrak p S$가
보조정리 \ref{algebra-lemma-domain-image-dense-set-points-generic-point}의
조건 (2)를 만족함을 보인다.
따라서 $(0)$으로 가는 소 이데알
$\overline{\mathfrak q} \subset S/\mathfrak pS$가 존재한다.
달리 말하면, $S$에서 $\overline{\mathfrak q}$의 역상
$\mathfrak q$는 원하는 대로 $\mathfrak p$로 간다.
\end{proof}

\noindent
위 보조정리의 결론은 상의 닫힌 점들의 집합으로부터
$f$의 상을 알아낼 수 있다는 것이다.
사상이 유한 표시인 경우에는 구성가능 집합의 상이 구성가능함을
알기 때문에 조금 더 좋다. 이를 진술하기 전에 몇 가지 표기를 도입하자.
구성가능 집합들의 집합을 $\text{Constr}(X)$로 나타내자.
$R \to S$를 환 사상이라 하자.
$X = \Spec(R)$와 $Y = \Spec(S)$로 나타내자.
유도된 스펙트럼 사상을 $f : Y \to X$로 쓰자.
위상 공간의 닫힌 점들의 집합을 아래첨자 ${}_0$로 나타내자.


\begin{lemma}
\label{algebra-lemma-conclude-jacobson-Noetherian}
위와 같은 표기를 사용하자. $R$이 뇌터 야콥슨 환이라고 하자.
또한 $R \to S$가 유한형이라고 하자.
다음과 같은 가환도식이 있다.
$$
\xymatrix{
\text{Constr}(Y) \ar[r]^{E \mapsto E_0} \ar[d]^{E \mapsto f(E)} &
\text{Constr}(Y_0) \ar[d]^{E \mapsto f(E)} \\
\text{Constr}(X) \ar[r]^{E \mapsto E_0} &
\text{Constr}(X_0)
}
$$
여기서 가로 화살표들은 위상론
보조정리 \ref{topology-lemma-jacobson-equivalent-constructible}의 전단사이다.
\end{lemma}

\begin{proof}
$R \to S$가 유한형이므로 유한 표시이다.
보조정리 \ref{algebra-lemma-Noetherian-finite-type-is-finite-presentation}를 보라.
% P01-E1336
따라서 $X$ 안의 구성가능 집합의 상은 슈발레 정리
(정리 \ref{algebra-theorem-chevalley})에 의해 $Y$ 안에서 구성가능하다.
이를 보조정리 \ref{algebra-lemma-image-finite-type-map-Jacobson-rings}와
결합하면 보조정리가 따른다.
\end{proof}

\noindent
야콥슨 환의 사용을 보여 주기 위해 다음 두 예를 제시한다.

\begin{example}
\label{algebra-example-product-matrices-zero}
$k$를 체라 하자. 공간 $\Spec(k[x, y]/(xy))$는
두 개의 기약 성분, 즉 $x$축과 $y$축을 갖는다.
이를 일반화하여 다음과 같이 두자.
$$
R = k[x_{11}, x_{12}, x_{21}, x_{22}, y_{11}, y_{12}, y_{21}, y_{22}]/
\mathfrak a,
$$
여기서 $\mathfrak a$는
$k[x_{11}, x_{12}, x_{21}, x_{22}, y_{11}, y_{12}, y_{21}, y_{22}]$
안에서 다음 $2 \times 2$ 곱셈 행렬의 성분들이 생성하는 이데알이다.
$$
\left(
\begin{matrix}
x_{11} & x_{12}\\
x_{21} & x_{22}
\end{matrix}
\right)
 \left(
\begin{matrix}
y_{11} & y_{12}\\
y_{21} & y_{22}
\end{matrix}
\right).
$$
이 예에서 $\Spec(R)$를 기술할 것이다.

\medskip\noindent
$\Spec(k[x, y]/(xy))$에 관한 명제를 증명하기 위해 다음과 같이 논증한다.
% P01-E1337
$\mathfrak p \subset k[x, y]$가 $xy$를 포함하는 임의의 이데알이면
$x$ 또는 $y$가 $\mathfrak p$에 포함되어야 한다. 따라서 그러한
극소 소 이데알은 $(x)$와 $(y)$뿐이다. $k$가
대수적으로 닫혀 있는 경우, 이 성분들의 $\text{max-Spec}$은
$y$축과 $x$축의 점 집합으로 시각화할 수 있다.

\medskip\noindent
일반화의 경우에는 다음 다항식환의 스펙트럼의 닫힌 점들을
$k[x_{11}, x_{12}, x_{21}, x_{22}, y_{11}, y_{12}, y_{21}, y_{22}])$
행렬들의 공간과 동일시할 수 있음에 유의하자.
$$
\left\{ (X, Y) \in \text{Mat}(2, k)\times \text{Mat}(2, k) \mid
X = \left(
\begin{matrix}
x_{11} & x_{12}\\
x_{21} & x_{22}
\end{matrix}
\right),
Y= \left(
\begin{matrix}
y_{11} & y_{12}\\
y_{21} & y_{22}
\end{matrix}
\right)
\right\}
$$
적어도 $k$가 대수적으로 닫혀 있다면 그렇다.
이제 행렬들의 공간 $\{(X, Y)\}$ 위에 다음 군 작용을 정의하자.
$\text{GL}(2, k)\times \text{GL}(2, k)\times \text{GL}(2, k)$
의 작용을 다음과 같이 정의한다.
$$
(g_1, g_2, g_3) \times (X, Y) \mapsto ((g_1Xg_2^{-1}, g_2Yg_3^{-1})).
$$
여기서 다음 대수 집합도 관찰하자.
$$
\text{GL}(2, k)\times \text{GL}(2, k)\times \text{GL}(2, k) \subset
\text{Mat}(2, k)\times \text{Mat}(2, k) \times \text{Mat}(2, k)
$$
이는 다음 정역의 극대 스펙트럼이므로 기약이다.
$$
k[x_{11}, x_{12}, \ldots, z_{21}, z_{22}, (x_{11}x_{22}-x_{12}x_{21})^{-1}
, (y_{11}y_{22}-y_{12}y_{21})^{-1}, (z_{11}z_{22}-z_{12}z_{21})^{-1}].
$$
% P01-E1338
기약 대수 집합의 상은 여전히 기약이므로,
집합 $\{(X, Y)\in \text{Mat}(2, k)\times \text{Mat}(2, k)|XY = 0\}$의
궤도를 분류하고 그 폐포를 취하는 것으로 충분하다.
표준 선형대수에 의해 다음 세 경우로 줄어든다.
\begin{enumerate}
\item $g_1Xg_2^{-1} = I_{2\times 2}$가 되도록 하는
$\exists (g_1, g_2)$가 존재한다.
그러면 $Y$는 필연적으로 $0$이고, 대수 집합으로서 이는
군 작용에 대해 불변이다. 따라서 이 궤도는 소 이데알
$(y_{11}, y_{12}, y_{21}, y_{22})$가 정의하는
기약 대수 집합 안에 포함된다. 폐포를 취하면
$(y_{11}, y_{12}, y_{21}, y_{22})$가 실제로 한 성분임을 알 수 있다.
\item $\exists (g_1, g_2)$가 다음을 만족한다.
$$
g_1Xg_2^{-1} = \left(
\begin{matrix}
1 & 0 \\
0 & 0
\end{matrix}
\right).
$$
이 경우는 $X$가 계수 1인 행렬일 필요충분조건이며,
더 나아가 그러한 $X$에 의해 $Y$가 소거되는 것은 다음과 동치이다.
$$
x_{11}y_{11}+x_{12}y_{21} = 0; \quad x_{11}y_{12}+x_{12}y_{22} = 0;
$$
$$
x_{21}y_{11}+x_{22}y_{21} = 0; \quad x_{21}y_{12}+x_{22}y_{22} = 0.
$$
계수 1인 $X$를 하나 고정하자. 위 방정식들을 만족하는
0이 아닌 $Y$들은 다음 이유로 기약 대수 집합을 이룬다.
($Y = 0$인 경우는 앞의 경우에 포함된다.)
$0 = g_1Xg_2^{-1}g_2Y$이므로
$$
g_2Y = \left(
\begin{matrix}
0 & 0 \\
y_{21}' & y_{22}'
\end{matrix}
\right).
$$
오른쪽에서 $g_3$에 의한 추가 $\text{GL}(2, k)$ 작용을 취하면
$g_2Y$를 다음 형태로 만들 수 있다.
$$
g_2Yg_3^{-1} = \left(
\begin{matrix}
0 & 0 \\
0 & 1
\end{matrix}
\right),
$$
따라서 그러한 $Y$들은 이 작용에 의한
$\text{GL}(2, k)$의 상과 동형인 기약 대수 집합을 이룬다.
마지막으로 $X$에 대한 ‘계수 1’ 조건은
기약 대수 집합 $\det X = x_{11}x_{22} - x_{12}x_{21} = 0$의
열린 조밀 부분집합을 이룬다는 점에 주목하자.
그러므로 다음 다섯 방정식 모두가 0이 아닌 행렬쌍의 공간의
열린 부분집합에서 기약 성분을 정의한다.
$(x_{11}y_{11}+x_{12}y_{21}, x_{11}y_{12}+x_{12}y_{22}, x_{21}y_{11}
+x_{22}y_{21}, x_{21}y_{12}+x_{22}y_{22}, x_{11}x_{22}-x_{12}x_{21})$
위의 궤적이 열린 조밀 부분집합인 기약 성분으로
$\det X = 0$, $\det Y = 0$라는 두 방정식이
$\Spec(R)$를 자른다는 것을 보일 수 있다.
\item $\exists (g_1, g_2)$가 $g_1Xg_2^{-1} = 0$을 만족하거나,
동치로 $X = 0$이다. 그러면 $Y$는 임의일 수 있고 이 성분은
따라서 $(x_{11}, x_{12}, x_{21}, x_{22})$로 정의된다.
\end{enumerate}
\end{example}


\begin{example}
\label{algebra-example-idempotent-matrices}
또 다른 예로 다음을 생각하자.
$R = k[\{t_{ij}\}_{i, j = 1}^{n}]/\mathfrak a$라 하자. 여기서 $\mathfrak a$는
곱셈 행렬 $T^2-T$의 성분들이 생성하는 이데알이고,
$T = (t_{ij})$이다. 선형대수에 의해, 사상
$g, T \mapsto gTg^{-1}$로 정의되는 $GL(n, k)$ 작용 아래에서
$T$는 그 계수에 의해 분류되며 각 $T$는
$r$개의 1과 $n-r$개의 0을 갖는
$\text{diag}(1, \ldots, 1, 0, \ldots, 0)$와 켤레이다.
따라서 그러한 $\text{diag}(1, \ldots, 1, 0, \ldots, 0)$의 각 궤도는
기약 성분을 이루고 모든 멱등 행렬은 그러한 궤도 중 하나에 포함된다.
이제 서로 다른 두 궤도가 반드시 서로소임을 보이자.
이를 위해서는 서로 다른 궤도에서 서로 다른 값을 취하는
다항식 함수들을 만들기만 하면 된다. 표수가 0인 경우에는
그러한 함수로
$f(t_{ij}) = trace(T) = \sum_{i = 1}^nt_{ii}$를 택할 수 있다.
양의 표수인 경우에는 조금 더 까다롭다. $T \neq 0$이어도
$trace(T) = 0$일 수 있기 때문이다. 예를 들어 $char = 3$인 경우
$$
trace\left(
\begin{matrix}
1 & & \\
& 1 & \\
& & 1
\end{matrix}
\right) = 3 = 0
$$
이다. 어쨌든 이 성분들은 다른 함수들을 사용하여 분리할 수 있다.
예를 들어 표수 3인 경우 $tr(\wedge^3T)$는
$diag(1, 1, 1)$에 대응하는 성분에서 1의 값을 취하고
다른 성분에서는 0의 값을 취한다.
\end{example}

\section{유한 및 정수적 환 확대}
\label{algebra-section-finite-ring-extensions}

\noindent
유한 환 사상과 정수적 환 사상에 관한 자명한 보조정리들이다.
정의를 상기하자.

\begin{definition}
\label{algebra-definition-integral-ring-map}
$\varphi : R \to S$를 환 사상이라 하자.
\begin{enumerate}
\item 원소 $s \in S$가 $R$ 위에서 {\it 정수적}이라는 것은
일계수 다항식 $P(x) \in R[x]$가 존재하여
$P^\varphi(s) = 0$이고, 여기서 $P^\varphi(x) \in S[x]$는
$\varphi : R[x] \to S[x]$에 의한 $P$의 상인 것을 뜻한다.
\item 환 사상 $\varphi$가 {\it 정수적}이라는 것은
모든 $s \in S$가 $R$ 위에서 정수적인 것을 뜻한다.
\end{enumerate}
\end{definition}

\begin{lemma}
\label{algebra-lemma-characterize-integral-element}
$\varphi : R \to S$를 환 사상이라 하자. $y \in S$라 하자.
$1 \in M$이고 $yM \subset M$인 $S$의 유한 $R$-부분가군
$M$이 존재한다면, $y$는 $R$ 위에서 정수적이다.
\end{lemma}

\begin{proof}
사상 $\varphi : M \to M$, $x \mapsto y \cdot x$를 생각하자.
보조정리 \ref{algebra-lemma-charpoly-module}에 의해
$P(\varphi) = 0$인 일계수 다항식 $P \in R[T]$가 존재한다.
환 $S$에서 다음을 얻는다.
$P(y) = P(y) \cdot 1 = P(\varphi)(1) = 0$.
\end{proof}

\begin{lemma}
\label{algebra-lemma-finite-is-integral}
유한 환 사상은 정수적이다.
\end{lemma}

\begin{proof}
$R \to S$를 유한이라 하자. $y \in S$라 하자.
 보조정리 \ref{algebra-lemma-characterize-integral-element}을
$M = S$에 적용하면 $y$가 $R$ 위에서 정수적임을 알 수 있다.
\end{proof}

\begin{lemma}
\label{algebra-lemma-characterize-integral}
$\varphi : R \to S$를 환 사상이라 하자. $s_1, \ldots, s_n$을
$S$의 유한 원소 집합이라 하자.
이 경우 모든 $i = 1, \ldots, n$에 대해 $s_i$가
$R$ 위에서 정수적일 필요충분조건은
모든 $s_i$를 포함하고 $R$ 위에서 유한한
$R$-부분대수 $S' \subset S$가 존재하는 것이다.
\end{lemma}

\begin{proof}
각 $s_i$가 정수적이면 $\varphi(R)$와 $s_i$들이 생성하는
부분대수는 $R$ 위에서 유한하다. 실제로 $s_i$가 $R$ 위에서
차수 $d_i$인 일계수 방정식을 만족한다면, 이 부분대수는
$0 \leq e_i \leq d_i - 1$인 원소
$s_1^{e_1} \ldots s_n^{e_n}$들에 의해 $R$-가군으로 생성된다.
반대로 모든 $s_i$를 포함하는 유한 $R$-부분대수
$S'$가 주어졌다고 하자. 그러면 보조정리
\ref{algebra-lemma-finite-is-integral}에 의해 모든 $s_i$는 정수적이다.
\end{proof}

\begin{lemma}
\label{algebra-lemma-characterize-finite-in-terms-of-integral}
$R \to S$를 환 사상이라 하자. 다음은 서로 동치이다.
\begin{enumerate}
\item $R \to S$는 유한하다.
\item $R \to S$는 정수적이고 유한형이다.
\item $R$ 위의 대수로서 $S$를 생성하는
$x_1, \ldots, x_n \in S$가 존재하고 각 $x_i$는 $R$ 위에서 정수적이다.
\end{enumerate}
\end{lemma}

 \begin{proof}
보조정리 \ref{algebra-lemma-characterize-integral}에서 즉시 따른다.
\end{proof}

\begin{lemma}
\label{algebra-lemma-integral-transitive}
\begin{slogan}
정수적 환 사상들의 합성은 정수적이다.
\end{slogan}
$R \to S$와 $S \to T$가 정수적 환 사상이라고 하자.
그러면 $R \to T$는 정수적이다.
\end{lemma}

\begin{proof}
$t \in T$라 하자. $P(x) \in S[x]$를 $P(t) = 0$인
일계수 다항식이라 하자.
보조정리 \ref{algebra-lemma-characterize-integral}을 $P$의 계수들의
유한 집합에 적용하자. 그러면 $t$는 $R$ 위에서 유한한
부분대수 $S' \subset S$ 위에서 정수적이다. 다시 보조정리
\ref{algebra-lemma-characterize-integral}을 적용하여 $S'$ 위에서 유한하고
$t$를 포함하는 부분대수 $T' \subset T$를 찾는다.
$R \to S' \to T'$에 보조정리 \ref{algebra-lemma-finite-transitive}를
적용하면 $T'$가 $R$ 위에서 유한임을 알 수 있다.
이제 보조정리 \ref{algebra-lemma-finite-is-integral}에 의해
$t$가 $R$ 위에서 정수적이라는 결론을 얻는다.
\end{proof}

\begin{lemma}
\label{algebra-lemma-integral-closure-is-ring}
$R \to S$를 환 준동형이라 하자. 집합
$$
S' = \{s \in S \mid s\text{ is integral over }R\}
$$
은 $S$의 $R$-부분대수이다.
\end{lemma}

\begin{proof}
이는 보조정리 \ref{algebra-lemma-characterize-integral}과
\ref{algebra-lemma-finite-is-integral}에서 자명하다.
\end{proof}

\begin{lemma}
\label{algebra-lemma-finite-product-integral}
$R_i\to S_i$를 $i = 1, \ldots, n$에 대한 환 사상들이라 하자.
$R$과 $S$를 각각 $R_i$들과 $S_i$들의 곱이라 하자.
그러면 원소 $s = (s_1, \ldots, s_n) \in S$가 $R$ 위에서
정수적일 필요충분조건은 각 $s_i$가 $R_i$ 위에서 정수적인 것이다.
\end{lemma}

\begin{proof}
생략한다.
\end{proof}

\begin{definition}
\label{algebra-definition-integral-closure}
$R \to S$를 환 사상이라 하자.
$R$ 위에서 정수적인 원소들의 환 $S' \subset S$를
보조정리 \ref{algebra-lemma-integral-closure-is-ring}에 따라
$S$에서 $R$의 {\it 정수적 폐포}라고 부른다.
$R \subset S$이면 $R = S'$일 때 $R$가 $S$에서
{\it 정수적으로 닫혔다}고 한다.
\end{definition}

\noindent
특히 $R \to S$가 정수적일 필요충분조건은
$S$에서 $R$의 정수적 폐포가 $S$ 전체인 것이다.

\begin{lemma}
\label{algebra-lemma-finite-product-integral-closure}
$R_i\to S_i$를 $i = 1, \ldots, n$에 대한 환 사상들이라 하자.
$S_i$에서 $R_i$의 정수적 폐포를 $S'_i$로 나타내자.
또한 $R$과 $S$를 각각 $R_i$들과 $S_i$들의 곱이라 하자.
그러면 $S$에서 $R$의 정수적 폐포는 $S'_i$들의 곱이다.
특히 $R \to S$가 정수적으로 닫힐 필요충분조건은 각
$R_i \to S_i$가 정수적으로 닫힌 것이다.
\end{lemma}

\begin{proof}
이는 보조정리 \ref{algebra-lemma-finite-product-integral}에서 즉시 따른다.
\end{proof}

\begin{lemma}
\label{algebra-lemma-integral-closure-localize}
정수적 폐포는 국소화와 가환한다. $A \to B$가 환 사상이고
$S \subset A$가 곱셈 부분집합이면, $S^{-1}B$에서
$S^{-1}A$의 정수적 폐포는 $S^{-1}B'$이다. 여기서
$B' \subset B$는 $B$에서 $A$의 정수적 폐포이다.
\end{lemma}

\begin{proof}
국소화는 완전하므로 $S^{-1}B' \subset S^{-1}B$임을 알 수 있다.
$x \in B'$와 $f \in S$를 택하자. 그러면 어떤 $a_i \in A$에 대하여
$B$에서
$x^d + \sum_{i = 1, \ldots, d} a_i x^{d - i} = 0$
이다. 따라서
$$
(x/f)^d + \sum\nolimits_{i = 1, \ldots, d} a_i/f^i (x/f)^{d - i} = 0
$$
도 $S^{-1}B$에서 성립한다. 이로부터 $S^{-1}B'$가
$S^{-1}A$에서 $S^{-1}B$의 정수적 폐포에 포함됨을 알 수 있다.
반대로 $x/f \in S^{-1}B$가 $S^{-1}A$ 위에서 정수적이라고 하자.
그러면 어떤 $a_i \in A$와 $f_i \in S$에 대하여
$$
(x/f)^d + \sum\nolimits_{i = 1, \ldots, d} (a_i/f_i) (x/f)^{d - i} = 0
$$
가 $S^{-1}B$에서 성립한다. 이는
$$
(f'f_1 \ldots f_d x)^d +
\sum\nolimits_{i = 1, \ldots, d}
f^i(f')^if_1^i \ldots f_i^{i - 1} \ldots f_d^i a_i
(f'f_1 \ldots f_dx)^{d - i} = 0
$$
인 어떤 $f' \in S$가 존재한다는 뜻이다. 따라서
$f'f_1\ldots f_dx \in B'$이고, 그러므로 원하는 대로
$x/f \in S^{-1}B'$이다.
\end{proof}

\begin{lemma}
\label{algebra-lemma-integral-closure-stalks}
\begin{slogan}
환 위의 대수의 원소가 그 환 위에서 정수적일 필요충분조건은
환의 모든 소 이데알에서 국소적으로 정수적인 것이다.
\end{slogan}
$\varphi : R \to S$를 환 사상이라 하자.
$x \in S$라 하자. 다음은 서로 동치이다.
\begin{enumerate}
\item $x$는 $R$ 위에서 정수적이다.
\item 모든 소 이데알 $\mathfrak p \subset R$에 대하여
$x \in S_{\mathfrak p}$가 $R_{\mathfrak p}$ 위에서 정수적이다.
\end{enumerate}
\end{lemma}

\begin{proof}
(1)이 (2)를 함은 자명하다. (2)를 가정하자.
$\varphi(R)$와 $x$가 생성하는 $R$-대수 $S' \subset S$를 생각하자.
$\mathfrak p$를 $R$의 소 이데알이라 하자. 그러면 어떤
$a_i \in R_{\mathfrak p}$에 대하여 $S_{\mathfrak p}$에서
$x^d + \sum_{i = 1, \ldots, d} \varphi(a_i) x^{d - i} = 0$
임을 안다. 따라서 나타나는 분모들을 살펴보면,
어떤 $f \in R$, $f \not \in \mathfrak p$가 존재하여
$a_i \in R_f$이고
$x^d + \sum_{i = 1, \ldots, d} \varphi(a_i) x^{d - i} = 0$
가 $S_f$에서 성립한다. 이는 $S'_f$가 $R_f$ 위에서 유한임을 뜻한다.
$\mathfrak p$가 임의였고 $\Spec(R)$가 준콤팩트이므로
(보조정리 \ref{algebra-lemma-quasi-compact}) $S'_{f_i}$가 $R_{f_i}$ 위에서
유한하도록 $R$의 단위원을 생성하는 유한한 원소
$f_1, \ldots, f_n \in R$을 찾을 수 있다.
따라서 보조정리 \ref{algebra-lemma-cover}에서 $S'$가 $R$ 위에서
유한임을 얻는다. 그러므로 보조정리
\ref{algebra-lemma-characterize-integral}에 의해 $x$는 $R$ 위에서 정수적이다.
\end{proof}
\begin{lemma}
\label{algebra-lemma-base-change-integral}
\begin{slogan}
기저변환에서 정수성과 유한성은 보존된다.
\end{slogan}
환 사상 $R \to S$와 $R \to R'$를 잡자.
$S' = R' \otimes_R S$로 놓자.
\begin{enumerate}
\item $R \to S$가 정수적이면 $R' \to S'$도 정수적이다.
\item $R \to S$가 유한이면 $R' \to S'$도 유한이다.
\end{enumerate}
\end{lemma}

\begin{proof}
(1)을 증명하자.
$s_i \in S$를 $R$ 위에서 $S$를 생성하는 원소들이라 하자.
각 원소는 일계수 다항식 $P_i$를 만족한다.
$R$ 위에서. 따라서 원소 $1 \otimes s_i \in S'$가 생성하고
$S'$는 $R'$ 위에서 이에 대응하는 다항식을 만족한다.
$P_i'$는 $R'$ 위에서 $S'$를 생성하고 $R'$ 위에서 생성하므로
$S'$가 $R'$ 위에서 정수적임을 알 수 있다.
(2)의 증명은 생략한다.
\end{proof}

\begin{lemma}
\label{algebra-lemma-integral-local}
환 사상 $R \to S$를 잡자.
$f_1, \ldots, f_n \in R$가 단위원을 생성한다고 하자.
\begin{enumerate}
\item 각 $R_{f_i} \to S_{f_i}$가 정수적이면 $R \to S$도 정수적이다.
\item 각 $R_{f_i} \to S_{f_i}$가 유한이면 $R \to S$도 유한이다.
\end{enumerate}
\end{lemma}

\begin{proof}
(1)을 증명하자.
$s \in S$를 잡자. $I \subset R[x]$인 다항식들의 아이디얼을 생각하자.
다항식 $P$가 $P(s) = 0$을 만족하는 경우를 모은 것이다. $J \subset R$를
$I$의 원소들의 최고차항 계수들의 아이디얼(!)이라 하자.
가정과 분모 제거에 의해 $f_i^{n_i} \in J$임을 모든 $i$에 대해 알 수 있고,
어떤 $n_i \geq 0$이다. 따라서 $J$는 $1$을 포함한다.
$s$는 $R$ 위에서 정수적이다.
(2)의 증명은 생략한다.
\end{proof}

\begin{lemma}
\label{algebra-lemma-integral-permanence}
$A \to B \to C$를 환 사상이라 하자.
\begin{enumerate}
\item $A \to C$가 정수적이면 $B \to C$도 정수적이다.
\item $A \to C$가 유한이면 $B \to C$도 유한이다.
\end{enumerate}
\end{lemma}

\begin{proof}
생략한다.
\end{proof}

\begin{lemma}
\label{algebra-lemma-integral-closure-transitive}
$A \to B \to C$를 환 사상이라 하자.
$B'$를 $B$에서 $A$의 정수적 폐포라 하고,
$C'$를 $C$에서 $B'$의 정수적 폐포라 하자. 그러면
$C'$는 $C$에서 $A$의 정수적 폐포이다.
\end{lemma}

\begin{proof}
생략한다.
\end{proof}

\begin{lemma}
\label{algebra-lemma-integral-overring-surjective}
$R \to S$가 $R \subset S$를 만족하는 정수적 환 확장이라고 하자.
그러면 $\varphi : \Spec(S) \to \Spec(R)$는 전사이다.
\end{lemma}

\begin{proof}
$\mathfrak p \subset R$를 소 이데알이라 하자.
$\mathfrak pS_{\mathfrak p} \not = S_{\mathfrak p}$임을 보여야 한다.
이는 보조정리 \ref{algebra-lemma-in-image}를 보라.
국소화 $R_{\mathfrak p} \to S_{\mathfrak p}$는
국소화가 완전하므로 단사이고, 보조정리
\ref{algebra-lemma-integral-closure-localize} 또는
\ref{algebra-lemma-base-change-integral}에 의해 정수적이다.
따라서 $R$, $S$를 $R_{\mathfrak p}$, $S_{\mathfrak p}$로 바꾸어도 되며,
$R$이 극대 이데알 $\mathfrak m$을 갖는 국소환이라고 가정할 수 있다.
그러므로 $\mathfrak mS \not = S$임을 보이면 충분하다.
모순을 얻기 위해 $f_i \in \mathfrak m$ 및 $s_i \in S$에 대해
$1 = \sum f_i s_i$라고 하자.
$R \subset S' \subset S$이고 $R \to S'$가 유한이며 $s_i \in S'$가
되도록 잡자. 이는 보조정리
\ref{algebra-lemma-characterize-integral}에 의해 가능하다.
방정식 $1 = \sum f_i s_i$는 유한 $R$-가군 $S'$가
$S' = \mathfrak m S'$를 만족함을 뜻한다.
따라서 나카야마 보조정리 \ref{algebra-lemma-NAK}에 의해 $S' = 0$이다.
이는 모순이다.
\end{proof}

\begin{lemma}
\label{algebra-lemma-integral-under-field}
$R$를 환이라 하자. $K$를 체라 하자.
$R \subset K$이고 $K$가 $R$ 위에서 정수적이면,
$R$은 체이고 $K$는 대수적 확장이다.
$R \subset K$이고 $K$가 $R$ 위에서 유한이면,
$R$은 체이고 $K$는 유한 대수적 확장이다.
\end{lemma}

\begin{proof}
$R \subset K$가 정수적이라고 하자.
보조정리 \ref{algebra-lemma-integral-overring-surjective}에 의해
$\Spec(R)$가 $1$개의 점만 가짐을 알 수 있다.
분명히 $R$은 정역이므로
$R = R_{(0)}$은 체이다(보조정리 \ref{algebra-lemma-minimal-prime-reduced-ring}).
나머지 주장들은 이로부터 즉시 따른다.
\end{proof}

\begin{lemma}
\label{algebra-lemma-integral-over-field}
$k$를 체라 하자. $S$를 $k$ 위의 $k$-대수라 하자.
\begin{enumerate}
\item $S$가 정역이고 $k$ 위에서 유한 차원이면 $S$는 체이다.
\item $S$가 $k$ 위에서 정수적이고 정역이면 $S$는 체이다.
\item $S$가 $k$ 위에서 정수적이면 $S$의 모든 소 이데알은
극대 이데알이다(더 많은 귀결은 보조정리
\ref{algebra-lemma-ring-with-only-minimal-primes}를 보라).
\end{enumerate}
\end{lemma}

\begin{proof}
소 이데알에 대한 주장은
``정수적 $+$ 정역 $\Rightarrow$ 체''라는 주장에 따른다.
$S$가 $k$ 위에서 정수적이고 $S$가 정역이라고 하자.
$s \in S$를 잡자. 보조정리
\ref{algebra-lemma-characterize-integral}에 의해 s를 포함하는
유한 차원 $k$-부분대수 $k \subset S' \subset S$를 찾을 수 있다.
따라서 $s$를 포함하는 $S$가 체임을 얻는다.
첫 번째 주장을 증명할 수 있으면 된다. $S$가
$k$ 위에서 유한 차원인 정역이라고 하자. $s\in S$를 잡자.
$S$가 정역이므로 곱셈 사상 $s : S \to S$는
차원에 의해 전사이다. 따라서 원소 $s_1 \in S$가 존재하여
$ss_1 = 1$이다. 그러므로 $S$는 체이다.
\end{proof}
\begin{lemma}
\label{algebra-lemma-integral-no-inclusion}
$R \to S$가 정수적이라고 하자.
$\mathfrak q, \mathfrak q' \in \Spec(S)$가 서로 다른 소 이데알이고
$\Spec(R)$에서 같은 상을 갖는다고 하자.
그러면 $\mathfrak q \subset \mathfrak q'$도 아니고
$\mathfrak q' \subset \mathfrak q$도 아니다.
\end{lemma}

\begin{proof}
$\mathfrak p \subset R$를 그 상이라 하자.
비고 \ref{algebra-remark-fundamental-diagram}에 의해
소 이데알 $\mathfrak q, \mathfrak q'$는
$S \otimes_R \kappa(\mathfrak p)$의 아이디얼에 대응한다.
따라서 이 보조정리는 보조정리 \ref{algebra-lemma-integral-over-field}에서 따른다.
\end{proof}

\begin{lemma}
\label{algebra-lemma-finite-finite-fibres}
$R \to S$가 유한이라고 하자.
그러면 $\Spec(S) \to \Spec(R)$의 올들은 유한하다.
\end{lemma}

\begin{proof}
비고 \ref{algebra-remark-fundamental-diagram}의 논의에 의해
올들은 $S \otimes_R \kappa(\mathfrak p)$인 환들의 스펙트럼이다.
$R \to S$가 유한이므로 이 올 환들은
$\kappa(\mathfrak p)$ 위에서 유한하고, 따라서 보조정리
\ref{algebra-lemma-Noetherian-permanence}에 의해 뇌터 환이다.
보조정리
\ref{algebra-lemma-integral-no-inclusion}에 의해
$S \otimes_R \kappa(\mathfrak p)$의 모든 소 이데알은 극소
소 이데알이다. 따라서 보조정리
\ref{algebra-lemma-Noetherian-irreducible-components}에 의해
그러한 소 이데알은 많아야 유한 개이다.
\end{proof}

\begin{lemma}
\label{algebra-lemma-integral-going-up}
$R \to S$를 환 사상이라 하고 $S$가 $R$ 위에서 정수적이라고 하자.
$\mathfrak p \subset \mathfrak p' \subset R$를 소 이데알이라 하자.
$\mathfrak q$를 $\mathfrak p$로 사상되는 $S$의 소 이데알이라 하자.
그러면 $\mathfrak q \subset \mathfrak q'$이고
$\mathfrak p'$로 사상되는 소 이데알 $\mathfrak q'$가 존재한다.
\end{lemma}

\begin{proof}
$R$을 $R/\mathfrak p$로, $S$를 $S/\mathfrak q$로 바꾸어도 된다.
그러면 정역의 정수적 확장 $R \subset S$와
$\mathfrak p' \subset R$라는 상황으로 줄어든다.
보조정리 \ref{algebra-lemma-integral-overring-surjective}에 의해 결론을 얻는다.
\end{proof}
\noindent
위 보조정리에서 표현된 성질을 환 사상 $R \to S$의
``going up 성질''이라고 부른다. 정의 \ref{algebra-definition-going-up-down}을 보라.

\begin{lemma}
\label{algebra-lemma-finite-finitely-presented-extension}
$R \to S$를 유한이고 유한 표시된 환 사상이라 하자.
$M$을 $S$-가군이라 하자.
그러면 $M$이 $R$-가군으로 유한 표시되는 것은
$M$이 $S$-가군으로 유한 표시되는 것과 동치이다.
\end{lemma}

\begin{proof}
한 방향의 함의는 보조정리
\ref{algebra-lemma-finitely-presented-over-subring}에서 따른다.
다른 방향을 보이자. $M$이 $S$-가군으로 유한 표시된다고 하자.
다음 표시를 택하자.
$$
S^{\oplus m} \longrightarrow
S^{\oplus n} \longrightarrow
M \longrightarrow 0
$$
$S$가 $R$-가군으로 유한이므로 $S^{\oplus n} \to M$의 핵은
유한 $R$-가군이다. 따라서 보조정리 \ref{algebra-lemma-extension}에 의해
$S$가 $R$-가군으로 유한 표시됨을 보이면 충분하다.

\medskip\noindent
$y_1, \ldots, y_n \in S$를 잡아 $y_1, \ldots, y_n$이
$R$-가군으로 $S$를 생성한다고 하자.
보조정리 \ref{algebra-lemma-characterize-integral-element}에 의해
각 $y_i$는 $R$ 위에서 정수적이다. $P_i(x) \in R[x]$이고
$P_i(y_i) = 0$인 일계수 다항식을 택하자. 다음 환을 생각하자.
$$
S' = R[x_1, \ldots, x_n]/(P_1(x_1), \ldots, P_n(x_n))
$$
그러면 보조정리 \ref{algebra-lemma-compose-finite-type}에 의해
$S$가 $S'$-대수로 유한 표시됨을 알 수 있다.
$S' \to S$가 전사이므로, 보조정리
\ref{algebra-lemma-finite-presentation-independent}에 의해 핵
$J = \Ker(S' \to S)$는 아이디얼로 유한 생성된다.
따라서 $J$는 정의로부터 즉시 유한 $S'$-가군이다.
그러므로 $S = \Coker(J \to S')$는 보조정리 \ref{algebra-lemma-extension}에 의해
$S'$-가군으로 유한 표시된다.
따라서 첫 번째 문단과 같이 논하면 $S'$가 $R$-가군으로
유한 표시됨을 보이면 충분하다. 실제로 $S'$는
$0 \leq e_i < \deg(P_i)$인 단항식
$x_1^{e_1} \ldots x_n^{e_n}$들을 기저로 갖는
$R$-가군으로서 자유이다. 즉 $R \to S'$를 다음 합성으로 쓸 수 있다.
$$
R \to R[x_1]/(P_1(x_1)) \to R[x_1, x_2]/(P_1(x_1), P_2(x_2)) \to
\ldots \to S'
$$
이 수열에서 $i$번째 환은 $(i - 1)$번째 환 위의
가군으로서 $1, x_i, \ldots, x_i^{\deg(P_i) - 1}$을 기저로 갖는 자유 가군이다.
따라서 귀납법으로 결과가 쉽게 따른다. 세부 사항은 생략한다.
\end{proof}
\begin{lemma}
\label{algebra-lemma-silly-normal}
$R$를 환이라 하자. $x, y \in R$를 영인자가 아닌 원소라 하자.
$R[x/y] \subset R_{xy}$를 $x/y$로 생성되는 $R$-부분대수라 하고,
마찬가지로 $R[y/x]$ 및 $R[x/y, y/x]$를 정의하자.
$R$이 $R_x$ 또는 $R_y$에서 정수적으로 닫혀 있다고 하자. 그러면 수열
$$
0 \to R \xrightarrow{(-1, 1)} R[x/y] \oplus R[y/x] \xrightarrow{(1, 1)}
R[x/y, y/x] \to 0
$$
은 $R$-가군의 짧은 완전열이다.
\end{lemma}

\begin{proof}
$x/y \cdot y/x = 1$이므로
$R[x/y] \oplus R[y/x] \to R[x/y, y/x]$가 전사임은 분명하다.
$\alpha \in R[x/y] \cap R[y/x]$라 하자. 중간에서의 완전성을
보이려면 $\alpha \in R$임을 증명해야 한다. 가정에 의해
$$
\alpha = a_0 + a_1 x/y + \ldots + a_n (x/y)^n =
b_0 + b_1 y/x + \ldots + b_m(y/x)^m
$$
어떤 $n, m \geq 0$ 및 $a_i, b_j \in R$에 대해 쓸 수 있다.
$N > \max(n, m)$인 N을 택하자.
$R_{xy}$의 원소
$$
(x/y)^N, (x/y)^{N - 1}, \ldots, x/y, 1, y/x, \ldots, (y/x)^{N - 1}, (y/x)^N
$$
들이 생성하는 유한 $R$-부분가군 $M$을 생각하자.
$\alpha M \subset M$임을 보이자. 실제로
$0 \leq i \leq N$에 대해
$(x/y)^i (b_0 + b_1 y/x + \ldots + b_m(y/x)^m) \in M$이고,
또한 $0 \leq i \leq N$에 대해
$(y/x)^i (a_0 + a_1 x/y + \ldots + a_n(x/y)^n) \in M$임은 분명하다.
따라서 보조정리 \ref{algebra-lemma-characterize-integral-element}에 의해
$\alpha$는 $R$ 위에서 정수적이다. $\alpha \in R_x$임에 유의하자.
그러므로 $R$이 $R_x$에서 정수적으로 닫혀 있으면 원하는 대로
$\alpha \in R$이다.
\end{proof}
\section{정규 환}
\label{algebra-section-normal-rings}

\noindent
먼저 정역의 정규성 개념을 도입한 다음,
(매우 일반적인) 정규 환의 개념을 도입하자.

\begin{definition}
\label{algebra-definition-domain-normal}
정역 $R$이 그 분수체에서 정수적으로 닫혀 있으면
{\it normal}이라고 부른다.
\end{definition}

\begin{lemma}
\label{algebra-lemma-integral-closure-in-normal}
환 사상 $R \to S$를 잡자.
$S$가 정규 정역이면 $S$에서 $R$의 정수적 폐포는
정규 정역이다.
\end{lemma}

\begin{proof}
생략한다.
\end{proof}
\noindent
다음 개념은 정규성을 연구할 때 때때로 유용하다.

\begin{definition}
\label{algebra-definition-almost-integral}
$R$를 정역이라 하자.
\begin{enumerate}
\item R의 분수체의 원소 $g$가, 어떤 $r \in R$, $r\not = 0$가
존재하여 모든 $n \geq 0$에 대해 $rg^n \in R$이면
분수체의 $R$에 대해 {\it almost\ integral\ over\ $R$}이라고 부른다.
\item 정역 $R$이 그 분수체의 모든
분수체의 $R$의 거의 정수적 원소가
$R$에 포함될 때 {\it completely\ normal}이라고 부른다.
\end{enumerate}
\end{definition}

\noindent
다음 보조정리는 뇌터 정역이 정규인 것과 완전히 정규인 것이
동치임을 보여 준다.

\begin{lemma}
\label{algebra-lemma-almost-integral}
$R$를 분수체 $K$를 갖는 정역이라 하자.
$u, v \in K$가 $R$ 위에서 거의 정수적이면
$u + v$와 $uv$도 그러하다. $R$ 위에서 정수적인
$g \in K$의 원소는 $R$ 위에서 거의 정수적이다.
$R$가 뇌터 환이면 역도 성립한다.
\end{lemma}

\begin{proof}
모든 $n \geq 0$에 대해 $ru^n \in R$이고
모든 $n \geq 0$에 대해 $v^nr' \in R$이면,
모든 $n \geq 0$에 대해 $(uv)^nrr'$와 $(u + v)^nrr'$가
$R$에 속한다. 따라서 첫 번째 주장이 따른다.
$g \in K$가 $R$ 위에서 정수적이라고 하자.
이 경우 어떤 $d > 0$가 존재하여
환 $R[g]$가 $R$-가군으로 $1, g, \ldots, g^d$에 의해 생성된다.
$1, g, \ldots, g^d \in K$의 공통 분모 $r \in R$을 잡자.
그러면 $rR[g] \subset R$이고, 따라서 $g$는 $R$ 위에서 거의 정수적이다.

\medskip\noindent
$R$가 뇌터이고 $g \in K$가 $R$ 위에서 거의 정수적이라고 하자.
정의에 나타난 대로 $r \in R$, $r\not = 0$를 잡자.
그러면 $R[g] \subset \frac{1}{r}R$가 $R$-가군으로 성립한다.
$R$가 뇌터이므로 이는 $R[g]$가 $R$ 위에서 유한임을 뜻한다.
따라서 보조정리 \ref{algebra-lemma-finite-is-integral}에 의해
$g$는 $R$ 위에서 정수적이다.
\end{proof}
\begin{lemma}
\label{algebra-lemma-localize-normal-domain}
정규 정역의 임의의 국소화는 정규이다.
\end{lemma}

\begin{proof}
$R$을 정규 정역이라 하고 $S \subset R$를 곱셈 부분집합이라 하자.
$R$의 분수체의 원소 $g$가 $S^{-1}R$ 위에서 정수적이라고 하자.
$a_i \in S^{-1}R$이고 $P(g) = 0$을 만족하는 다항식
$P = x^d + \sum_{j < d} a_j x^j$를 택하자.
모든 $i$에 대해 $sa_i \in R$가 되도록 $s \in S$를 택하자.
그러면 $sg$는 계수 $s^{d-j}a_j$가 $R$에 속하는 일계수 다항식
$x^d + \sum_{j < d} s^{d-j}a_j x^j$을 만족한다.
따라서 $R$이 정규이므로 $sg \in R$이다.
그러므로 $g \in S^{-1}R$이다.
\end{proof}
\begin{lemma}
\label{algebra-lemma-PID-normal}
주 아이디얼 정역은 정규이다.
\end{lemma}

\begin{proof}
$R$을 주 아이디얼 정역이라 하자.
$g = a/b$를 $R$의 분수체의 원소이면서 $R$ 위에서 정수적인 원소라 하자.
$R$이 주 아이디얼 정역이므로 $a$와 $b$의 공통 인자를 나누어 내고
$(a, b) = R$이라고 가정할 수 있다. 이 경우
$r_i \in R$인 임의의 방정식
$(a/b)^n + r_{n-1} (a/b)^{n-1} + \ldots + r_0 = 0$
은 $a^n \in (b)$를 뜻한다. 이는 $(a, b) = R$과 모순이므로
$b$는 $R$의 단위원이어야 한다.
\end{proof}

\begin{lemma}
\label{algebra-lemma-prepare-polynomial-ring-normal}
$R$를 분수체 $K$를 갖는 정역이라 하자.
$f = \sum \alpha_i x^i$가 $K[x]$의 원소라고 하자.
\begin{enumerate}
\item $f$가 $R[x]$ 위에서 정수적이면 모든 $\alpha_i$는
$R$ 위에서 정수적이고,
\item $f$가 $R[x]$ 위에서 거의 정수적이면 모든 $\alpha_i$는
$R$ 위에서 거의 정수적이다.
\end{enumerate}
\end{lemma}

\begin{proof}
먼저 두 번째 명제를 증명하자.
$f = \alpha_0 + \alpha_1 x + \ldots + \alpha_r x^r$라고 쓰고
$\alpha_r \not = 0$이라고 하자.
가정에 의해 $h = b_0 + b_1 x + \ldots + b_s x^s \in R[x]$가 존재하여,
$b_s \not = 0$이고 모든 $n \geq 0$에 대해 $f^n h \in R[x]$이다.
이는 모든 $n \geq 0$에 대해 $b_s \alpha_r^n \in R$임을 뜻한다.
따라서 $\alpha_r$는 $R$ 위에서 거의 정수적이다.
거의 정수적 원소들의 집합은 부분환을 이루므로
(보조정리 \ref{algebra-lemma-almost-integral}), 다음을 얻는다.
$f - \alpha_r x^r = \alpha_0 + \alpha_1 x + \ldots + \alpha_{r - 1} x^{r - 1}$는
$R[x]$ 위에서 거의 정수적이다. $r$에 대한 귀납법으로 증명이 끝난다.

\medskip\noindent
첫 번째 명제를 증명하기 위해 절대 뇌터 환 감소를 사용한다.
즉, $\alpha_i = a_i / b_i$라고 쓰고
$P(t) = t^d + \sum_{j < d} f_j t^j$를
$f_j \in R[x]$인 계수를 갖고 $P(f) = 0$을 만족하는 다항식이라고 하자.
$f_j = \sum f_{ji}x^i$라고 쓰자. $R$의 원소들
$a_i, b_i, f_{ji}$의 유한 목록으로 생성되는 부분환
$R_0 \subset R$를 생각하자. 이는 정역이다. 그 분수체를
$K_0$라고 하자. $R_0$는 유한형
$\mathbf{Z}$-대수이므로 뇌터 환이다
(보조정리 \ref{algebra-lemma-obvious-Noetherian} 참조).
$R$에서 원소들 $a_i, b_i, f_{ji}$ 사이에 성립하는 모든 항등식이
$R_0$에서도 성립하므로, 여전히 $f \in K_0[x]$는 $R_0[x]$ 위에서
정수적이다. 보조정리 \ref{algebra-lemma-almost-integral}에 의해
$f$는 $R_0[x]$ 위에서 거의 정수적이다. 이 보조정리의 두 번째 명제에 의해
$\alpha_i$들은 $R_0$ 위에서 거의 정수적이다. 그리고 $R_0$가 뇌터 환이므로
이들은 $R_0$ 위에서 정수적이다
(보조정리 \ref{algebra-lemma-almost-integral} 참조).
그러므로 물론 이들은 $R$ 위에서도 정수적이다.
\end{proof}

\begin{lemma}
\label{algebra-lemma-polynomial-domain-normal}
$R$을 정규 정역이라 하자.
그러면 $R[x]$는 정규 정역이다.
\end{lemma}

\begin{proof}
$R$이 체 $K$인 경우에는 이 결과가 성립한다. 왜냐하면
$K[x]$는 유클리드 정역이고, 따라서 주 아이디얼 정역이며,
보조정리 \ref{algebra-lemma-PID-normal}에 의해 정규이기 때문이다.
$R[x]$의 분수체의 원소 $g$가 $R[x]$ 위에서 정수적이라고 하자.
$R$의 분수체인 $K$에 대해 $g$가 $K[x]$ 위에서도 정수적이므로
$g = \alpha_d x^d + \alpha_{d-1}x^{d-1} +
\ldots + \alpha_0$로 쓸 수 있고, $\alpha_i \in K$이다.
보조정리 \ref{algebra-lemma-prepare-polynomial-ring-normal}에 의해
$\alpha_i$들은 $R$ 위에서 정수적이고,
따라서 $R$의 원소이다.
\end{proof}

\begin{lemma}
\label{algebra-lemma-power-series-over-Noetherian-normal-domain}
$R$을 뇌터 정규 정역이라 하자. 그러면 $R[[x]]$는
뇌터 정규 정역이다.
\end{lemma}

\begin{proof}
멱급수환은 보조정리 \ref{algebra-lemma-Noetherian-power-series}에 의해
뇌터 환이다.
$f, g \in R[[x]]$가 영이 아닌 원소라고 하자. 이때
$w = f/g$가 $R[[x]]$ 위에서 정수적이라고 하자.
$R$의 분수체를 $K$라고 하자. 로랑 급수환
$K((x)) = K[[x]][1/x]$는 체이므로, 다음과 같이 쓸 수 있다.
어떤 $n \in \mathbf{Z}$, $a_i \in K$, $a_n \not = 0$에 대해
$w = a_n x^n + a_{n + 1} x^{n + 1} + \ldots$이다.
보조정리 \ref{algebra-lemma-almost-integral}에 의해
$R[[x]]$의 영이 아닌 원소
$h = b_m x^m + b_{m + 1} x^{m + 1} + \ldots$가 존재하고
$b_m \not = 0$이며 모든 $e \geq 1$에 대해
$w^e h \in R[[x]]$이다. 따라서 $n \geq 0$이고
모든 $e \geq 1$에 대해 $b_m a_n^e \in R$이다.
$R$가 뇌터 환이므로 같은 보조정리에 의해 $a_n \in R$이다.
이제 $a_n, a_{n + 1}, \ldots, a_{N - 1} \in R$라고 하자.
그러면 같은 논리를
$w - a_n x^n - \ldots - a_{N - 1} x^{N - 1} = a_N x^N + \ldots$에
적용할 수 있다. 이 과정을 통해 모든 $a_i \in R$임을 알 수 있으므로
보조정리가 증명된다.
\end{proof}

\begin{lemma}
\label{algebra-lemma-normality-is-local}
$R$를 정역이라 하자. 다음은 서로 동치이다.
\begin{enumerate}
\item 정역 $R$은 정규 정역이다.
\item 모든 소 아이디얼 $\mathfrak p \subset R$에 대해 국소환
$R_{\mathfrak p}$는 정규 정역이고,
\item 모든 극대 아이디얼 $\mathfrak m$에 대해 환 $R_{\mathfrak m}$은
정규 정역이다.
\end{enumerate}
\end{lemma}

\begin{proof}
(1) $\Rightarrow$ (2)는 보조정리 \ref{algebra-lemma-localize-normal-domain}에서 따른다.
함의 (2) $\Rightarrow$ (3)은 자명하다. 함의
(3) $\Rightarrow$ (1)은 임의의 정역 $R$에 대해 분수체 안에서
$$
R = \bigcap\nolimits_{\mathfrak m} R_{\mathfrak m}
$$
가 성립한다는 사실에서 따른다. 실제로 $R$의 분수체의 원소인 $g$에 대해
아이디얼 $I = \{x \in R \mid xg \in R\}$는 어떤 극대 아이디얼
$\mathfrak m$에도 포함되지 않으므로 $I = R$이다.
\end{proof}

\noindent
보조정리 \ref{algebra-lemma-normality-is-local}은 다음 정의가
정의 \ref{algebra-definition-domain-normal}과 양립함을 보여 준다.
(이는 EGA의 정의이다 -- \cite[IV, 5.13.5 and 0, 4.1.4]{EGA} 참조.)

\begin{definition}
\label{algebra-definition-ring-normal}
환 $R$이 모든 소 아이디얼
$\mathfrak p \subset R$에 대해 국소화 $R_{\mathfrak p}$가
정규 정역일 때 (정의 \ref{algebra-definition-domain-normal} 참조)
{\it 정규}라고 부른다.
\end{definition}

\noindent
정규 환은 모든 소 아이디얼에서의 국소화들의 곱의 부분환이므로
$R$은 축소 환임에 유의하자
(예를 들어 보조정리 \ref{algebra-lemma-characterize-zero-local} 참조).

\begin{lemma}
\label{algebra-lemma-normal-ring-integrally-closed}
정규 환은 그 전체 분수환에서 정수적으로 닫혀 있다.
\end{lemma}

\begin{proof}
$R$을 정규 환이라 하자. $x \in Q(R)$를 $R$의 전체 분수환의
원소이면서 $R$ 위에서 정수적인 것으로 하자.
$I = \{f \in R, fx \in R\}$로 놓자. $\mathfrak p \subset R$를
소 아이디얼이라 하자. $R \to R_{\mathfrak p}$가 평탄하므로
$R_{\mathfrak p} \subset Q(R) \otimes_R R_{\mathfrak p}$임을 안다.
$R_{\mathfrak p}$가 정규 정역이므로 $x \otimes 1$은
$R_{\mathfrak p}$의 원소이다. 따라서 $a, f \in R$, $f \not \in \mathfrak p$인
a, f를 잡아 $x \otimes 1 = a \otimes 1/f$로 쓸 수 있다.
이는 $fx - a$가 $Q(R) \otimes_R R_{\mathfrak p} = Q(R)_{\mathfrak p}$에서
0으로 사상됨을 뜻하고, 다시 말해 $f' \in R$, $f' \not \in \mathfrak p$인
원소가 존재하여 $f'fx = f'a$가 $R$에서 성립한다는 뜻이다.
다르게 말하면 $ff' \in I$이다. 따라서 $I$는 $R$의 소 아이디얼 중
어느 것에도 포함되지 않는 아이디얼이고, 즉 $I = R$이며 $x \in R$이다.
\end{proof}

\begin{lemma}
\label{algebra-lemma-localization-normal-ring}
정규 환의 국소화는 정규 환이다.
\end{lemma}

\begin{proof}
생략한다.
\end{proof}

\begin{lemma}
\label{algebra-lemma-polynomial-ring-normal}
$R$을 정규 환이라 하자. 그러면 $R[x]$는 정규 환이다.
\end{lemma}

\begin{proof}
$\mathfrak q$를 $R[x]$의 소 아이디얼이라 하자.
$\mathfrak p = R \cap \mathfrak q$로 놓자.
그러면 보조정리 \ref{algebra-lemma-polynomial-domain-normal}에 의해
$R_{\mathfrak p}[x]$는 정규 정역이다.
따라서 보조정리 \ref{algebra-lemma-localize-normal-domain}에 의해
$(R[x])_{\mathfrak q}$는 정규 정역이다.
\end{proof}

\begin{lemma}
\label{algebra-lemma-finite-product-normal}
정규 환들의 유한 곱은 정규이다.
\end{lemma}

\begin{proof}
두 정규 환, 즉 $R$과 $S$의 곱이 정규임을 보이면 충분하다.
보조정리 \ref{algebra-lemma-disjoint-decomposition}에 의해 $R\times S$의
소 아이디얼은 $\mathfrak{p}\times S$와 $R\times
\mathfrak{q}$, where $\mathfrak{p}$ and $\mathfrak{q}$ are primes of $R$
% P01-E1339
그리고 $S$ respectively. Localization yields
$(R\times S)_{\mathfrak{p}\times S}=R_{\mathfrak{p}}$이고,
이는 가정에 의해 정규 정역이다. $S$에 대해서도 마찬가지이다.
\end{proof}

\begin{lemma}
\label{algebra-lemma-characterize-reduced-ring-normal}
$R$를 환이라 하자. $R$이 축소되어 있고 극소 소 아이디얼이 유한 개뿐이라고 가정하자.
다음 명제들은 서로 동치이다.
\begin{enumerate}
\item $R$은 정규 환이다.
\item $R$은 그 전체 분수환에서 정수적으로 닫혀 있다.
\item $R$은 정규 정역들의 유한 곱이다.
\end{enumerate}
\end{lemma}

\begin{proof}
함의 (1) $\Rightarrow$ (2)와
(3) $\Rightarrow$ (1)은 일반적으로 성립한다.
보조정리 \ref{algebra-lemma-normal-ring-integrally-closed}와
\ref{algebra-lemma-finite-product-normal}을 참조하라.

\medskip\noindent
$\mathfrak p_1, \ldots, \mathfrak p_n$을 $R$의 극소 소 아이디얼들이라고 하자.
보조정리 \ref{algebra-lemma-reduced-ring-sub-product-fields}와
\ref{algebra-lemma-total-ring-fractions-no-embedded-points}에 의해
$Q(R) = R_{\mathfrak p_1} \times \ldots \times R_{\mathfrak p_n}$이고,
보조정리 \ref{algebra-lemma-minimal-prime-reduced-ring}에 의해 각 인자는 체이다.
$Q(R)$의 $i$번째 멱등원소
$e_i = (0, \ldots, 0, 1, 0, \ldots, 0)$라고 하자.

\medskip\noindent
$R$이 $Q(R)$에서 정수적으로 닫혀 있으면 특히 멱등원소 $e_i$들을
포함한다. 따라서 $R$은 $n$개 정역의 곱이다
(절 \ref{algebra-section-connected-components} 및
\ref{algebra-section-tilde-module-sheaf} 참조). 각 인자는
분수체가 $R_{\mathfrak p_i}$인 $R/\mathfrak p_i$의 형태이다.
보조정리 \ref{algebra-lemma-finite-product-integral-closure}에 의해 각 사상
$R/\mathfrak p_i \to R_{\mathfrak p_i}$는 정수적으로 닫혀 있다.
따라서 $R$은 정규 정역들의 유한 곱이다.
\end{proof}

\begin{lemma}
\label{algebra-lemma-colimit-normal-ring}
$(R_i, \varphi_{ii'})$를 환들의 유향계라고 하자
(범주론, 정의 \ref{algebra-definition-directed-system}).
각 $R_i$가 정규 환이면
$R = \colim_i R_i$도 정규 환이다.
\end{lemma}

\begin{proof}
$\mathfrak p \subset R$를 소 아이디얼이라 하자.
$\mathfrak p_i = R_i \cap \mathfrak p$로 놓자
(통상적인 표기상의 남용이다).
그러면
$R_{\mathfrak p} = \colim_i (R_i)_{\mathfrak p_i}$이다.
각 $(R_i)_{\mathfrak p_i}$가 정규 정역이므로
이제 정규 정역에 대한 보조정리의 명제를 증명하면 된다.
$a, b \in R$이고 $a/b$가 일계수 다항식
$P(T) \in R[T]$를 만족한다고 하자. 그러면 충분히 큰
$i \in I$를 택하여 $a, b$가 $R_i$ 위의 원소
$a_i, b_i$에서 오고, $P$가 일계수 다항식
$P_i\in R_i[T]$에서 오며, $P_i(a_i/b_i)=0$이 되게 할 수 있다.
$R_i$가 정규이므로 $a_i/b_i \in R_i$이고, 따라서 $a/b \in R$이다.
\end{proof}







\section{정규환 위 정수적 확대의 하강}
\label{algebra-section-going-down-integral-over-normal}

\noindent
먼저 아이디얼 위에서 정수적인 원소라는 개념을 조금 다룬 다음,
절 제목에서 언급한 정리를 증명한다.

\begin{definition}
\label{algebra-definition-integral-over-ideal}
$\varphi : R \to S$를 환 사상이라 하자.
$I \subset R$를 아이디얼이라 하자.
원소 $g \in S$에 대해, 계수 $a_j \in I^{d-j}$를 갖는
일계수 다항식
$P = x^d + \sum_{j < d} a_j x^j$가 존재하여
$P^\varphi(g) = 0$이 $S$에서 성립하면
이 원소가 {\it $I$ 위에서 정수적}이라고 한다.
\end{definition}

\noindent
이는 주로 $\varphi = \text{id}_R : R \to R$인 경우에 사용한다.
이 경우 $I$ 위에서 정수적인 원소들의 집합 $I'$를
{\it $I$의 정수적 폐포}라고 부른다. $I'$가
$R$의 아이디얼임을 보일 것이다(물론 $I \subset I'$이다).

\begin{lemma}
\label{algebra-lemma-characterize-integral-ideal}
$\varphi : R \to S$를 환 사상이라 하자.
$I \subset R$를 아이디얼이라 하자.
$A = \sum I^nt^n \subset R[t]$를
$R \oplus It \subset R[t]$로 생성되는 다항식환의 부분환이라 하자.
원소 $s \in S$가 $I$ 위에서 정수적인 것은
원소 $st \in S[t]$가 $A$ 위에서 정수적인 것과 동치이다.
\end{lemma}

\begin{proof}
$st$가 $A$ 위에서 정수적이라고 하자.
$P = x^d + \sum_{j < d} a_j x^j$를
$A$의 계수를 갖고 $P^\varphi(st) = 0$을 만족하는 일계수 다항식이라 하자.
$a_j' \in A$를 $a_j$의 차수 $d-j$ 부분이라 하자.
다르게 말하면, 어떤 $a_j'' \in I^{d-j}$에 대해
$a_j' = a_j'' t^{d-j}$이다.
차수를 고려하면 여전히
$(st)^d + \sum_{j < d} \varphi(a_j'') t^{d-j} (st)^j = 0$이다.
따라서 $s^d + \sum_{j < d} \varphi(a_j'') s^j = 0$이고,
$s$는 $I$ 위에서 정수적이다.

\medskip\noindent
$s$가 $I$ 위에서 정수적이라고 하자.
$a_j \in I^{d-j}$인
$P = x^d + \sum_{j < d} a_j x^j$가 있다고 하자.
그러면 $A$에 계수를 갖는 다항식
$Q = x^d + \sum_{j < d} (a_j t^{d-j}) x^j$를 곧바로 얻고,
이는 $st$가 $A$ 위에서 정수적임을 증명한다.
\end{proof}

\begin{lemma}
\label{algebra-lemma-integral-over-ideal-is-submodule}
$\varphi : R \to S$를 환 사상이라 하자.
$I \subset R$를 아이디얼이라 하자.
$I$ 위에서 정수적인 $S$의 원소들의 집합은
$S$의 $R$-부분가군을 이룬다.
더 나아가 $s \in S$가 $R$ 위에서 정수적이고
$s'$가 $I$ 위에서 정수적이면,
$ss'$도 $I$ 위에서 정수적이다.
\end{lemma}

\begin{proof}
더 언급하지 않고 보조정리 \ref{algebra-lemma-integral-closure-is-ring}을 사용한다.
덧셈에 대해 닫혀 있음은, 그 표기를 따를
보조정리 \ref{algebra-lemma-characterize-integral-ideal}의 특징짓기에서 명백하다.
$R$ 위에서 정수적인 임의의 원소 $s \in S$는
$S[t]$의 차수 $0$인 원소 $s$에 대응하며, 이는
$A$ 위에서 정수적이다($R \subset A$이기 때문이다).
% P01-E1340
따라서 $S[t]$에서 $s$를 곱하는 것은
$A$ 위에서 정수적이라는 성질을 보존하고,
\end{proof}

\begin{lemma}
\label{algebra-lemma-integral-integral-over-ideal}
$\varphi : R \to S$가 정수적이라고 하자.
$I \subset R$를 아이디얼이라 하자.
그러면 $IS$의 모든 원소는 $I$ 위에서 정수적이다.
\end{lemma}

\begin{proof}
보조정리 \ref{algebra-lemma-integral-over-ideal-is-submodule}에서 즉시 따른다.
\end{proof}

\begin{lemma}
\label{algebra-lemma-polynomials-divide}
$K$를 체라 하자. $n, m \in \mathbf{N}$이고
$a_0, \ldots, a_{n - 1}, b_0, \ldots, b_{m - 1} \in K$라 하자.
다항식 $x^n + a_{n - 1}x^{n - 1} + \ldots + a_0$가
$K[x]$에서 다항식 $x^m + b_{m - 1} x^{m - 1} + \ldots + b_0$를 나누면,
다음이 성립한다.
\begin{enumerate}
\item $a_0, \ldots, a_{n - 1}$은 원소
$b_0, \ldots, b_{m - 1}$을 포함하는 $K$의 임의의 부분환
$R_0$ 위에서 정수적이고,
\item 각 $a_i$는 원소
$a_0, \ldots, a_{n - 1}, b_0, \ldots, b_{m - 1}$을 포함하는
임의의 부분환 $R \subset K$에 대해 다음 아이디얼에 속한다.
$$
\sqrt{(b_0, \ldots, b_{m-1})R}
$$
\end{enumerate}
\end{lemma}

\begin{proof}
다음과 같이 쓸 수 있게 하는 체확대 $L/K$를 택하자.
$x^m + b_{m - 1} x^{m - 1} + \ldots + b_0 =
\prod_{i = 1}^m (x - \beta_i)$이고 $\beta_i \in L$이다.
체론, 절 \ref{fields-section-splitting-fieds}을 참조하라.
각 $\beta_i$는 $R_0$ 위에서 정수적이다.
각 $a_i$는 $\beta_1, \ldots, \beta_m$의 동차 다항식이므로,
$a_i$에 대해서도 같은 결론을 얻는다
(보조정리 \ref{algebra-lemma-integral-closure-is-ring}을 사용하라).
이는 (1)을 증명한다.

\medskip\noindent
$R$을 (2)에서와 같이 잡자. 다음이 성립하도록
$c_0, \ldots, c_{m - n - 1} \in K$를 택하자.
$$
\begin{matrix}
x^m + b_{m - 1} x^{m - 1} + \ldots + b_0 =  \\
(x^n + a_{n - 1}x^{n - 1} + \ldots + a_0)
(x^{m - n} + c_{m - n - 1}x^{m - n - 1}+ \ldots + c_0).
\end{matrix}
$$
이 등식은 다음을 뜻한다.
\begin{align*}
c_{m - n - 1} & = b_{m - 1} - a_{n - 1}, \\
c_{m - n - 2} & = b_{m - 2} - a_{n - 2} - a_{n - 1}c_{m - n - 1}, \\
\ldots
\end{align*}
따라서 모든 $j$에 대해 $c_j \in R$이다.
근기 $\sqrt{(b_0, \ldots, b_{m - 1})}$로 나누어
축소 환 $\overline{R}$를 얻는다.
상 $\overline{a}_i \in \overline{R}$가 0임을 보이면 된다.
그리고 $\overline{R}[x]$에서 다음 관계가 성립한다.
$$
\begin{matrix}
x^m = x^m + \overline{b}_{m - 1} x^{m - 1} + \ldots + \overline{b}_0 = \\
(x^n + \overline{a}_{n - 1}x^{n - 1} + \ldots + \overline{a}_0)
(x^{m - n} + \overline{c}_{m - n - 1}x^{m - n - 1}+ \ldots + \overline{c}_0).
\end{matrix}
$$
이로부터 모든 $i$에 대해 $\overline{a}_i = 0$임을 쉽게 알 수 있다.
실제로 보조정리 \ref{algebra-lemma-minimal-prime-reduced-ring}에 의해
극소 소 아이디얼 $\mathfrak{p}$에서 $\overline{R}$의 국소화는 체이고,
$\overline{R}_{\mathfrak p}[x]$는 유일분해정역이다.
따라서 $f = x^n + \sum \overline{a}_i x^i$는
$x^n$과 동반이며, $f$가 일계수이므로
$\overline{R}_{\mathfrak p}[x]$에서 $f = x^n$이다.
그러면 $s \in \overline{R}$, $s \not\in \mathfrak p$인 원소가 존재하여
$s(f - x^n) = 0$이다. 따라서 모든 $\overline{a}_i$는
$\mathfrak p$에 속하고, 보조정리
\ref{algebra-lemma-reduced-ring-sub-product-fields}에 의해 결론이 따른다.
\end{proof}

\begin{lemma}
\label{algebra-lemma-minimal-polynomial-normal-domain}
$R \subset S$를 정역들의 포함이라 하자.
$R$이 정규라고 가정하자. $g \in S$가
$R$ 위에서 정수적이라고 하자. 그러면 $g$의 최소다항식은
$R$에 계수를 갖는다.
\end{lemma}

\begin{proof}
$P = x^m + b_{m-1} x^{m-1} + \ldots + b_0$를
$R$에 계수를 갖고 $P(g) = 0$을 만족하는 다항식이라 하자.
$Q = x^n + a_{n-1}x^{n-1} + \ldots + a_0$를
$R$의 분수체 $K$ 위에서 $g$의 최소다항식이라 하자.
그러면 $Q$는 $K[x]$에서 $P$를 나눈다.
보조정리 \ref{algebra-lemma-polynomials-divide}에 의해
$a_i$들이 $R$ 위에서 정수적임을 안다.
$R$이 정규이므로 이들은 $R$에 속한다.
\end{proof}

\begin{proposition}
\label{algebra-proposition-going-down-normal-integral}
$R \subset S$를 정역들의 포함이라 하자.
$R$이 정규이고 $S$가 $R$ 위에서 정수적이라고 가정하자.
$\mathfrak p \subset \mathfrak p' \subset R$를 소 아이디얼들이라 하자.
$\mathfrak q'$를 $\mathfrak p' = R \cap \mathfrak q'$를 만족하는
$S$의 소 아이디얼이라 하자.
그러면 $\mathfrak q \subset \mathfrak q'$이고
$\mathfrak p = R \cap \mathfrak q$를 만족하는 소 아이디얼
$\mathfrak q$가 존재한다. 다르게 말하면
$R \to S$에 대해 하강 성질이 성립한다
(정의 \ref{algebra-definition-going-up-down} 참조).
\end{proposition}

\begin{proof}
$\mathfrak p$, $\mathfrak p'$와 $\mathfrak q'$를
명제에서와 같이 잡자. $\mathfrak q \subset \mathfrak q'$이고
$R \cap \mathfrak q = \mathfrak p$를 만족하는 소 아이디얼
$\mathfrak q$가 존재함을 보여야 한다. 이는
$\mathfrak p$로 사상되는 $S_{\mathfrak q'}$의 소 아이디얼을 찾는 것과 같다.
보조정리 \ref{algebra-lemma-in-image}에 따르면
$\mathfrak p S_{\mathfrak q'} \cap R
= \mathfrak p$임을 보이면 된다.
$z \in \mathfrak p S_{\mathfrak q'} \cap R$를 택하자.
$y \in \mathfrak pS$이고 $g \in S$, $g \not\in \mathfrak q'$인
원소를 택하여 $z = y/g$로 쓸 수 있다.
다르게 쓰면 $zg = y$이다.

\medskip\noindent
보조정리 \ref{algebra-lemma-integral-integral-over-ideal}에 의해
$b_i \in \mathfrak p$인 일계수 다항식
$P = x^m + b_{m-1} x^{m-1} + \ldots + b_0$가 존재하여
$P(y) = 0$이다.

\medskip\noindent
보조정리 \ref{algebra-lemma-minimal-polynomial-normal-domain}에 의해
$K$ 위에서 $g$의 최소다항식은 $R$에 계수를 갖는다.
이를 $Q = x^n + a_{n-1} x^{n-1} + \ldots
+ a_0$로 쓰자. 모든 $a_i$, $i = n-1, \ldots, 0$가
$\mathfrak p$에 속하는 것은 아니다. 그렇지 않으면
% P01-E1341
$g^n = \sum_{j < n} a_j g^j \in \mathfrak pS
\subset \mathfrak p'S
\subset \mathfrak q'$
이 되어 모순이다.

\medskip\noindent
$y = zg$이므로 위의 식에서 곧바로
$Q' = x^n + za_{n-1} x^{n-1} + \ldots + z^{n}a_0$가
$y$의 최소다항식임을 안다. 따라서
$Q'$는 $P$를 나누고, 보조정리 \ref{algebra-lemma-polynomials-divide}에 의해
$z^ja_{n - j} \in \sqrt{(b_0, \ldots, b_{m-1})}
\subset \mathfrak p$, $j =  1, \ldots, n$임을 안다.
모든 $a_i$, $i = n-1, \ldots, 0$가
$\mathfrak p$에 속하는 것은 아니므로, 원하는 대로
$z \in \mathfrak p$를 얻는다.
\end{proof}
\section{평탄 가군과 평탄 환 사상}
\label{algebra-section-flat}

\noindent
자주 사용되는 한 결과는 $M = \colim_{i\in \mathcal{I}} M_i$가
$R$-가군들의 귀납적 극한이고 $N$이 $R$-가군이면
$$
M \otimes N
=
\colim_{i\in \mathcal{I}} M_i \otimes_R N,
$$
이라는 것이다(보조정리 \ref{algebra-lemma-tensor-products-commute-with-limits} 참조).
이 성질은 보통 {\it $\otimes$가 귀납적 극한과 교환한다}고 표현한다.
또 다른 자주 사용되는 결과는
$0 \to N_1 \to N_2 \to N_3 \to 0$이 완전열이고
$M$이 임의의 $R$-가군이면
$$
M \otimes_R N_1
\to
M \otimes_R N_2
\to
M \otimes_R N_3
\to
0
$$
도 여전히 완전하다는 것이다
(보조정리 \ref{algebra-lemma-tensor-product-exact} 참조).
이 두 성질은 함자 $N \mapsto M \otimes_R N$이
{\it 오른쪽 완전}임을 말해 준다.
범주론, 절 \ref{categories-section-exact-functor} 및
호몰로지, 절 \ref{homology-section-functors}를 참조하라.
$R$-가군 $M$은 $N \mapsto N \otimes_R M$도 왼쪽 완전,
즉 완전할 때 평탄하다고 한다. 정확한 정의는 다음과 같다.

\begin{definition}
\label{algebra-definition-flat}
$R$을 환이라 하자.
\begin{enumerate}
\item $R$-가군 $M$은, $R$-가군의 임의의 완전열
$N_1 \to N_2 \to N_3$에 대해 열
$M \otimes_R N_1 \to M \otimes_R N_2 \to M \otimes_R N_3$도
완전하면 {\it 평탄}이라고 한다.
\item $R$-가군 $M$은, $R$-가군의 복합체
$N_1 \to N_2 \to N_3$가 완전인 것이
열 $M \otimes_R N_1 \to M \otimes_R N_2 \to M \otimes_R N_3$가
완전인 것과 동치이면 {\it 충실히 평탄}이라고 한다.
\item 환 사상 $R \to S$는
$S$가 $R$-가군으로서 평탄하면 {\it 평탄}이라고 한다.
\item 환 사상 $R \to S$는
$S$가 $R$-가군으로서 충실히 평탄하면 {\it 충실히 평탄}이라고 한다.
\end{enumerate}
\end{definition}

\noindent
평탄성 조건을 사용하는 한 예는 다음과 같다.

\begin{lemma}
\label{algebra-lemma-flat-intersect-ideals}
$R$을 환이라 하자. $I, J \subset R$를 아이디얼들이라 하고,
$M$을 평탄한 $R$-가군이라 하자. 그러면
$IM \cap JM = (I \cap J)M$이다.
\end{lemma}

\begin{proof}
완전열 $0 \to I \cap J \to R \to R/I \oplus R/J$를 생각하자.
평탄 가군 $M$과 텐서곱하여 완전열
$$
0 \to (I \cap J) \otimes_R M \to M \to M/IM \oplus M/JM
$$
을 얻는다. $M \to M/IM \oplus M/JM$의 핵은
$IM \cap JM$과 같으므로 결론이 따른다.
\end{proof}

\begin{lemma}
\label{algebra-lemma-colimit-flat}
$R$을 환이라 하자. $\{M_i, \varphi_{ii'}\}$를
평탄한 $R$-가군들의 유향계라고 하자. 그러면
$\colim_i M_i$는 평탄한 $R$-가군이다.
\end{lemma}

\begin{proof}
이는 $\otimes$가 귀납적 극한과 교환하고 유향 귀납적 극한이 완전하기
때문이다. 보조정리 \ref{algebra-lemma-directed-colimit-exact}을 참조하라.
\end{proof}

\begin{lemma}
\label{algebra-lemma-composition-flat}
(충실히) 평탄한 환 사상들의 합성은 (충실히) 평탄하다.
$R \to R'$가 (충실히) 평탄하고 $M'$가
(충실히) 평탄한 $R'$-가군이면, $M'$는
(충실히) 평탄한 $R$-가군이다.
\end{lemma}

\begin{proof}
첫 번째 명제는 두 번째 명제의 특수한 경우이므로,
후자를 증명하면 충분함이 명백하다.
$R \to R'$를 평탄한 환 사상이라 하고 $M'$를 평탄한
$R'$-가군이라 하자. $M'$가 평탄한 $R$-가군임을 보여야 한다.
$N_1 \to N_2 \to N_3$를 $R$-가군들의 완전 복합체라고 하자.
그러면 복합체
$R' \otimes_R N_1 \to
R' \otimes_R N_2 \to R' \otimes_R N_3$는 완전하다
($R'$가 평탄한 $R$-가군이기 때문이다). 따라서 복합체
$M' \otimes_{R'} \left(R' \otimes_R N_1\right)
\to M' \otimes_{R'} \left(R' \otimes_R N_2\right)
\to M' \otimes_{R'} \left(R' \otimes_R N_3\right)$도 완전하다
($M'$가 평탄한 $R'$-가군이기 때문이다). 임의의 $R$-가군 $N$에 대해
$M' \otimes_{R'} \left(R' \otimes_R N\right)
\cong \left(M' \otimes_{R'} R'\right) \otimes_R N
\cong M' \otimes_R N$이라는 자연스러운 동형이 있으므로
(보조정리 \ref{algebra-lemma-tensor-with-bimodule} 및
\ref{algebra-lemma-flip-tensor-product}에 의해), 이 복합체는
$M' \otimes_R N_1 \to M' \otimes_R N_2 \to M' \otimes_R N_3$와
동형이고, 따라서 후자도 완전하다.
이는 $M'$가 평탄한 $R$-가군임을 보인다.
이 논증을 거꾸로 추적하면, $R \to R'$가 충실히 평탄하고
$M'$가 충실히 평탄한 $R'$-가군일 때
$M'$가 충실히 평탄한 $R$-가군임을 보일 수 있다.
\end{proof}

\begin{lemma}
\label{algebra-lemma-flat}
$M$을 $R$-가군이라 하자. 다음은 서로 동치이다.
\begin{enumerate}
\item
\label{algebra-item-flat}
$M$은 $R$ 위에서 평탄하다.
\item
\label{algebra-item-injective}
$R$-가군의 모든 단사 사상 $N \subset N'$에 대해
사상 $N \otimes_R M \to N'\otimes_R M$은 단사이다.
\item
\label{algebra-item-f-ideal}
모든 아이디얼 $I \subset R$에 대해 사상
$I \otimes_R M \to R \otimes_R M = M$은 단사이다.
\item
\label{algebra-item-ffg-ideal}
모든 유한 생성 아이디얼 $I \subset R$에 대해 사상
$I \otimes_R M \to R \otimes_R M = M$은 단사이다.
\end{enumerate}
\end{lemma}

\begin{proof}
(\ref{algebra-item-flat})이면 (\ref{algebra-item-injective})이고, 다시
(\ref{algebra-item-f-ideal})이며, 다시 (\ref{algebra-item-ffg-ideal})인 것은
모두 자명하다. 따라서 (\ref{algebra-item-ffg-ideal})이면
(\ref{algebra-item-flat})임을 증명한다.
$N_1 \to N_2 \to N_3$가 완전하다고 하자.
$K = \Ker(N_2 \to N_3)$ 및 $Q = \Im(N_2 \to N_3)$로 놓자.
그러면 다음 사상들을 얻는다.
$$
N_1 \otimes_R M \to
K \otimes_R M \to
N_2 \otimes_R M \to
Q \otimes_R M \to
N_3 \otimes_R M
$$
첫 번째와 세 번째 화살표는 전사이다.
따라서 두 번째와 네 번째 화살표가 단사임을 보이면
증명이 끝난다\footnote{이를 더 자세히 설명하자.
두 번째와 네 번째 화살표가 단사임을 안다고 가정하자.
완전열 $K \to N_2 \to Q \to 0$에
보조정리 \ref{algebra-lemma-tensor-product-exact}을 적용하면
열 $K \otimes_R M \to N_2 \otimes_R M \to
Q \otimes_R M \to 0$이 완전함을 얻는다. 따라서
$\Ker \left(N_2 \otimes_R M \to Q \otimes_R M\right)
= \Im \left(K \otimes_R M \to N_2 \otimes_R M\right)$이다.
또한 $N_1 \otimes_R M \to K \otimes_R M$의 전사성 때문에
$\Im \left(K \otimes_R M \to N_2 \otimes_R M\right)
= \Im \left(N_1 \otimes_R M \to N_2 \otimes_R M\right)$이고,
$Q \otimes_R M \to N_3 \otimes_R M$의 단사성 때문에
$\Ker \left(N_2 \otimes_R M \to Q \otimes_R M\right)
= \Ker \left(N_2 \otimes_R M \to N_3 \otimes_R M\right)$이다.
따라서
$\Ker \left(N_2 \otimes_R M \to N_3 \otimes_R M\right)
= \Im \left(N_1 \otimes_R M \to N_2 \otimes_R M\right)$이고,
이는 함자 $- \otimes_R M$이 완전함을 보이므로
$M$은 평탄하다.}.
따라서 $- \otimes_R M$이 단사인 $R$-가군 사상을
단사인 $R$-가군 사상으로 보냄을 보이면 충분하다.

\medskip\noindent
$K \to N$을 단사인 $R$-가군 사상이라고 하고,
$x \in \Ker(K \otimes_R M \to N \otimes_R M)$라 하자.
$x$가 0임을 보여야 한다.
$R$-가군 $K$는 그 유한 $R$-부분가군들의 합집합이다.
따라서 $K \otimes_R M$은 $K_i$가 $K$의 모든 유한
$R$-부분가군을 달릴 때 $K_i \otimes_R M$ 꼴의
$R$-가군들의 귀납적 극한이다
(텐서곱은 귀납적 극한과 교환하기 때문이다).
따라서 어떤 $i$에 대해 우리의 $x$는 원소
$x_i \in K_i \otimes_R M$에서 온다.
그러므로 $K$가 유한 $R$-가군이라고 가정해도 된다.
이를 가정하자. 단사 사상 $K \to N$을 포함으로 보아
$K \subset N$이라고 하자.

\medskip\noindent
$R$-가군 $N$은 $K$를 포함하는 유한 $R$-부분가군들의 합집합이다.
따라서 $N \otimes_R M$은 $N_i$가 $K$를 포함하는
$N$의 모든 유한 $R$-부분가군을 달릴 때
$N_i \otimes_R M$ 꼴의 $R$-가군들의 귀납적 극한이다
(다시 텐서곱이 귀납적 극한과 교환하기 때문이다).
이는 유향계에 대한 귀납적 극한임에 유의하자
($N$의 두 유한 부분가군의 합은 다시 유한하기 때문이다).
따라서 보조정리 \ref{algebra-lemma-zero-directed-limit}에 의해
원소 $x \in K \otimes_R M$은 이 $R$-가군들 가운데 적어도 하나인
$N_i \otimes_R M$에서 0으로 사상된다
($x$가 $N \otimes_R M$에서 0으로 사상되기 때문이다).
따라서 $N$이 유한 $R$-가군이라고 가정해도 된다.

\medskip\noindent
$N$이 유한 $R$-가군이라고 가정하자.
어떤 $L \subset L' \subset R^{\oplus n}$에 대해
$N = R^{\oplus n}/L$ 및 $K = L'/L$로 쓰자.
임의의 $R$-부분가군 $G \subset R^{\oplus n}$에 대해,
합성
$G \otimes_R M \to R^n \otimes_R M = M^{\oplus n}$으로 얻는
표준 사상 $G \otimes_R M \to M^{\oplus n}$이 있다.
$L \otimes_R M \to M^{\oplus n}$ 및
$L' \otimes_R M \to M^{\oplus n}$이 단사임을 보이면 충분하다.
실제로 그렇다면
$K \otimes_R M = L' \otimes_R M/L \otimes_R M \to M^{\oplus n}/L \otimes_R M$
도 단사임을 알 수 있다\footnote{이는
$L' \otimes_R M$과 $L \otimes_R M$을 $M^{\oplus n}$의
부분가군으로 식별하면 명백해진다.
% P01-E1342
그렇게 해도 되는 것은 사상 $L \otimes_R M \to M^{\oplus n}$과
$L' \otimes_R M \to M^{\oplus n}$이 단사이고,
명백한 사상 $L' \otimes_R M \to L \otimes_R M$과 교환하기 때문이다.}.

\medskip\noindent
따라서 $L \subset R^{\oplus n}$이 $R$-부분가군일 때
$L \otimes_R M \to M^{\oplus n}$이 단사임을 보이면 충분하다.
$n$에 대한 귀납법으로 이를 증명한다. 기초 단계 $n = 1$은 아래에서 다룬다.
귀납 단계에서는 $n > 1$이라고 가정하고
$L' = L \cap R \oplus 0^{\oplus n - 1}$로 놓는다.
그러면 $L'' = L/L'$는 $R^{\oplus n - 1}$의 부분가군이다.
다음 도식을 얻는다.
$$
\xymatrix{
&
L' \otimes_R M \ar[r] \ar[d] &
L \otimes_R M \ar[r] \ar[d] &
L'' \otimes_R M \ar[r] \ar[d] &
0 \\
0 \ar[r] &
M \ar[r] &
M^{\oplus n} \ar[r] &
M^{\oplus n - 1} \ar[r] & 0
}
$$
귀납 가정과 기초 단계에 의해 왼쪽과 오른쪽 수직 화살표는 단사이다.
행들은 완전하다. 따라서 가운데 수직 화살표도 단사이다.

\medskip\noindent
위 귀납법의 기초 단계는 $L \subset R$이 아이디얼인 경우이다.
다르게 말하면 $R$의 임의의 아이디얼 $I$에 대해
$I \otimes_R M \to M$이 단사임을 보여야 한다.
$I$가 유한 생성이면 이것이 참임을 안다.
하지만 $I = \bigcup I_\alpha$는 그 안에 포함된 유한 생성 아이디얼
$I_\alpha$들의 합집합이다. 다르게 말하면
$I = \colim I_\alpha$이다. $\otimes$가 귀납적 극한과 교환하므로
$I \otimes_R M = \colim I_\alpha \otimes_R M$이고,
모든 사상 $I_\alpha \otimes_R M \to M$은 가정에 의해 단사이므로
$I \otimes_R M \to M$도 단사이다.
\end{proof}


\begin{lemma}
\label{algebra-lemma-colimit-rings-flat}
$\{R_i, \varphi_{ii'}\}$를 유향집합 $I$ 위의 환들의 계라고 하자.
$R = \colim_i R_i$로 놓자.
\begin{enumerate}
\item $M$이 $R$-가군이고 모든 $i$에 대해 $M$이
$R_i$-가군으로서 평탄하면, $M$은 평탄한 $R$-가군이다.
\item $i \in I$에 대해 $M_i$를 평탄한 $R_i$-가군이라 하고,
$i' \geq i$에 대해 $f_{ii'} : M_i \to M_{i'}$를
$f_{i' i''} \circ f_{i i'} = f_{i i''}$를 만족하는
$\varphi_{ii'}$-선형 사상이라 하자. 그러면
$M = \colim_{i \in I} M_i$는 평탄한 $R$-가군이다.
\end{enumerate}
\end{lemma}

\begin{proof}
(1)은 모든 $i$에 대해 $M_i = M$이고
$f_{i i'} = \text{id}_M$인 (2)의 특수한 경우이다. (2)를 증명하자.
$\mathfrak a \subset R$를 유한 생성 아이디얼이라 하자.
보조정리 \ref{algebra-lemma-flat}에 의해
$\mathfrak a \otimes_R M \to M$이 단사임을 보이면 충분하다.
어떤 $i \in I$와 유한 생성 아이디얼
$\mathfrak a' \subset R_i$를 찾아 $\mathfrak a = \mathfrak a'R$로 쓸 수 있다.
그러면 $\mathfrak a = \colim_{i' \geq i} \mathfrak a'R_{i'}$이다.
$\otimes$가 귀납적 극한과 교환하므로 사상
$\mathfrak a \otimes_R M \to M$은 다음 사상들의 귀납적 극한이다.
$$
\mathfrak a'R_{i'} \otimes_{R_{i'}} M_{i'} \longrightarrow M_{i'}
$$
가정에 의해 이 사상들은 모두 단사이다. $I$ 위의 귀납적 극한은
보조정리 \ref{algebra-lemma-directed-colimit-exact}에 의해 완전하므로 결론이 따른다.
\end{proof}

\begin{lemma}
\label{algebra-lemma-flat-base-change}
$M$이 $R$ 위에서 (충실히) 평탄하고 $R \to R'$가
환 사상이라고 하자. 그러면 $M \otimes_R R'$는
$R'$ 위에서 (충실히) 평탄하다.
\end{lemma}

\begin{proof}
임의의 $R'$-가군 $N$에 대해 표준 동형
$N \otimes_{R'} (R'\otimes_R M)
= N \otimes_R M$이 있다. 따라서 함자
$-\otimes_{R'}(R'\otimes_R M)$의 원하는 완전성 성질은
함자 $-\otimes_R M$의 대응하는 완전성 성질에서 따른다.
\end{proof}

\begin{lemma}
\label{algebra-lemma-flatness-descends}
$R \to R'$를 충실히 평탄한 환 사상이라 하자.
$M$을 $R$ 위의 가군이라 하고 $M' = R' \otimes_R M$로 놓자.
그러면 $M$이 $R$ 위에서 평탄인 것과
$M'$가 $R'$ 위에서 평탄인 것은 동치이다.
\end{lemma}

\begin{proof}
보조정리 \ref{algebra-lemma-flat-base-change}에 의해 $M$이 평탄하면
$M'$도 평탄하다. 역으로 $M'$가 평탄하다고 하자.
$N_1 \to N_2 \to N_3$를 $R$-가군들의 완전열이라 하자.
$N_1 \otimes_R M \to N_2 \otimes_R M \to N_3 \otimes_R M$이
완전함을 보이고자 한다. $R \to R'$가 평탄하므로
$N_1 \otimes_R R' \to N_2 \otimes_R R' \to N_3 \otimes_R R'$가
완전함을 안다. $M'$의 평탄성에 의해
$N_1 \otimes_R R' \otimes_{R'} M'
\to N_2 \otimes_R R' \otimes_{R'} M'
\to N_3 \otimes_R R' \otimes_{R'} M'$는 완전하다.
이를 다음과 같이 쓸 수 있다.
$N_1 \otimes_R M \otimes_R R'
\to N_2 \otimes_R M \otimes_R R'
\to N_3 \otimes_R M \otimes_R R'$.
마지막으로 충실한 평탄성에 의해
$N_1 \otimes_R M \to N_2 \otimes_R M \to N_3 \otimes_R M$은
완전하다.
\end{proof}

\begin{lemma}
\label{algebra-lemma-flatness-descends-more-general}
$R$을 환이라 하자. $S \to S'$를 $R$-대수들의 평탄한 사상이라 하자.
$M$을 $S$ 위의 가군이라 하고 $M' = S' \otimes_S M$로 놓자.
\begin{enumerate}
\item $M$이 $R$ 위에서 평탄하면 $M'$도 $R$ 위에서 평탄하다.
\item $S \to S'$가 충실히 평탄하면, $M$이 $R$ 위에서 평탄인 것과
$M'$가 $R$ 위에서 평탄인 것은 동치이다.
\end{enumerate}
\end{lemma}

\begin{proof}
$N \to N'$를 $R$-가군들의 단사 사상이라 하자.
$S \to S'$의 평탄성에 의해
$$
\Ker(N \otimes_R M \to N' \otimes_R M) \otimes_S S'
=
\Ker(N \otimes_R M' \to N' \otimes_R M')
$$
이다. $M$이 $R$ 위에서 평탄하면 좌변은 0이고,
보조정리 \ref{algebra-lemma-flat}의 평탄성에 대한 두 번째 특징짓기에 의해
$M'$가 $R$ 위에서 평탄함을 얻는다.
$M'$가 $R$ 위에서 평탄하면 우변이 0이다.
여기에 $S \to S'$가 충실히 평탄하다는 조건까지 있으면
$\Ker(N \otimes_R M \to N' \otimes_R M)$이 0임을 얻고,
이는 다시 $M$이 $R$ 위에서 평탄함을 보인다.
\end{proof}

\begin{lemma}
\label{algebra-lemma-flat-permanence}
$R \to S$를 환 사상이라 하고 $M$을 $S$-가군이라 하자.
$M$이 $R$-가군으로서 평탄하고 $S$-가군으로서 충실히 평탄하면,
$R \to S$는 평탄하다.
\end{lemma}

\begin{proof}
$N_1 \to N_2 \to N_3$를 $R$-가군들의 완전열이라 하자.
가정에 의해 $N_1 \otimes_R M \to N_2 \otimes_R M \to N_3 \otimes_R M$은
완전하다. 이를 다음과 같이 쓸 수 있다.
$$
N_1 \otimes_R S \otimes_S M
\to
N_2 \otimes_R S \otimes_S M
\to
N_3 \otimes_R S \otimes_S M.
$$
$M$이 $S$ 위에서 충실히 평탄하므로
$N_1 \otimes_R S \to N_2 \otimes_R S \to N_3 \otimes_R S$가
완전함을 얻는다. 따라서 $R \to S$는 평탄하다.
\end{proof}

\noindent
$R$을 환이라 하자.
$M$을 $R$-가군이라 하자.
$\sum f_i x_i = 0$을 $M$에서의 관계식이라 하자.
정수 $m \geq 0$, 원소들 $y_j \in M$, $j = 1, \ldots, m$, 그리고
원소들 $a_{ij} \in R$, $i = 1, \ldots, n$, $j = 1, \ldots, m$이 존재하여
다음을 만족하면 관계식 $\sum f_i x_i$를 {\it 자명하다}고 한다.
$$
x_i = \sum\nolimits_j a_{ij} y_j, \forall i,
\quad\text{and}\quad
0 = \sum\nolimits_i f_ia_{ij}, \forall j.
$$

\begin{lemma}[평탄성의 방정식 판정법]
\label{algebra-lemma-flat-eq}
$R$ 위의 가군 $M$이 평탄인 것과
$M$에서의 모든 관계식이 자명한 것은 동치이다.
\end{lemma}

\begin{proof}
$M$이 평탄하고 $\sum f_i x_i = 0$이 $M$에서의 관계식이라고 하자.
$I = (f_1, \ldots, f_n)$로 놓고,
$K = \Ker(R^n \to I, (a_1, \ldots, a_n) \mapsto \sum_i a_i f_i)$로 놓자.
그러면 짧은 완전열 $0 \to K \to R^n \to I \to 0$을 얻는다.
$\sum f_i \otimes x_i$는 $I \otimes_R M$의 원소이고
$R \otimes_R M = M$에서 0으로 사상된다. 평탄성에 의해
$\sum f_i \otimes x_i$는 $I \otimes_R M$에서 0이다.
따라서 $K \otimes_R M$의 어떤 원소가
$\sum e_i \otimes x_i \in R^n \otimes_R M$으로 사상된다.
여기서 $e_i$는 $R^n$의 $i$번째 기저 원소이다.
이 원소를 $\sum k_j \otimes y_j$로 쓰고,
$R^n$에서 $k_j$의 상을 $\sum a_{ij} e_i$로 쓰면 결과를 얻는다.

\medskip\noindent
모든 관계식이 자명하다고 가정하자. $I$를 유한 생성 아이디얼이라 하고,
$x = \sum f_i \otimes x_i$를 $I \otimes_R M$의 원소로서
$R \otimes_R M = M$에서 0으로 사상되는 것이라 하자.
이는 정확히 $\sum f_i x_i$가 $M$에서의 관계식이라는 뜻이다.
그 관계식이 자명하다는 사실로부터 $x$가 0임이 쉽게 따른다. 실제로
$$
x
=
\sum f_i \otimes x_i
=
\sum f_i \otimes \left(\sum a_{ij}y_j\right)
=
\sum \left(\sum f_i a_{ij}\right) \otimes y_j
=
0
$$
이다.
\end{proof}

\begin{lemma}
\label{algebra-lemma-flat-tor-zero}
$R$이 환이고, $0 \to M'' \to M' \to M \to 0$이 짧은 완전열이며,
$N$이 $R$-가군이라고 하자. $M$이 평탄하면
$N \otimes_R M'' \to N \otimes_R M'$는 단사이다. 즉, 열
$$
0 \to N \otimes_R M'' \to N \otimes_R M' \to N \otimes_R M \to 0
$$
은 짧은 완전열이다.
\end{lemma}

\begin{proof}
$R^{(I)} \to N$을 자유 가군에서 $N$으로 가는 전사 사상이라 하고,
그 핵을 $K$라 하자. 다음 도식에 뱀 보조정리를 적용하면 결과가 따른다.
$$
\begin{matrix}
 & & 0 & & 0 & & 0 & & \\
 & & \uparrow & & \uparrow & & \uparrow & & \\
 & & M''\otimes_R N & \to & M' \otimes_R N & \to & M \otimes_R N & \to & 0 \\
 & & \uparrow & & \uparrow & & \uparrow & & \\
0 & \to & (M'')^{(I)} & \to & (M')^{(I)} & \to & M^{(I)} & \to & 0 \\
 & & \uparrow & & \uparrow & & \uparrow & & \\
 & & M''\otimes_R K & \to & M' \otimes_R K & \to & M \otimes_R K & \to & 0 \\
 & & & & & & \uparrow & & \\
 & & & & & & 0 & &
\end{matrix}
$$
그 행과 열들은 완전하다. 자유 가군 $R^{(I)}$와 텐서곱하는 것은
완전하므로 가운데 행이 완전하다.
\end{proof}


\begin{lemma}
\label{algebra-lemma-flat-ses}
$0 \to M' \to M \to M'' \to 0$이 $R$-가군들의
짧은 완전열이라고 하자.
$M'$와 $M''$가 평탄하면 $M$도 평탄하다.
$M$과 $M''$가 평탄하면 $M'$도 평탄하다.
\end{lemma}

\begin{proof}
모든 아이디얼 $I \subset R$에 대해 사상
$N \otimes_R I \to N$이 단사이면 가군 $N$이 평탄하다는
판정법을 사용하겠다. 보조정리 \ref{algebra-lemma-flat}을 참조하라.
아이디얼 $I \subset R$를 생각하자.
다음 도식을 생각하자.
$$
\begin{matrix}
0 & \to & M' & \to & M & \to & M'' & \to & 0 \\
& & \uparrow & & \uparrow & & \uparrow & & \\
& & M'\otimes_R I & \to & M \otimes_R I & \to & M''\otimes_R I & \to & 0
\end{matrix}
$$
그 행들은 완전하다. 이로부터 첫 번째 주장이 곧바로 증명된다.
두 번째 주장은 $M''$가 평탄하면 보조정리
\ref{algebra-lemma-flat-tor-zero}에 의해 아래쪽 왼쪽 수평 화살표가
단사라는 사실에서 따른다.
\end{proof}

\begin{lemma}
\label{algebra-lemma-easy-ff}
$R$을 환이라 하자.
$M$을 $R$-가군이라 하자.
다음은 서로 동치이다.
\begin{enumerate}
\item $M$은 충실히 평탄하다.
\item $M$은 평탄하고, 모든 $R$-가군 준동형
$\alpha : N \to N'$에 대해 $\alpha = 0$인 것과
$\alpha \otimes \text{id}_M = 0$인 것은 동치이다.
\end{enumerate}
\end{lemma}

\begin{proof}
$M$이 충실히 평탄하면
$0 \to \Ker(\alpha) \to N \to N'$가 완전인 것과
$M$과 텐서곱한 뒤에도 같은 열이 완전인 것은 동치이다.
이는 (1)에서 (2)가 따름을 증명한다.
반대로 (2)를 가정하자. $N_1 \to N_2 \to N_3$를 복합체라 하고,
복합체 $N_1 \otimes_R M \to N_2 \otimes_R M \to N_3\otimes_R M$이
완전하다고 가정하자. $x \in \Ker(N_2 \to N_3)$를 택하고 사상
$\alpha : R \to N_2/\Im(N_1)$,
$r \mapsto rx + \Im(N_1)$를 생각하자.
복합체 $-\otimes_R M$의 완전성에 의해
$\alpha \otimes \text{id}_M$가 영 사상임을 안다.
가정에 의해 $\alpha$가 영 사상이다.
따라서 $x $는 $N_1 \to N_2$의 상에 속한다.
\end{proof}

\begin{lemma}
\label{algebra-lemma-ff}
\begin{slogan}
평탄 가군이 충실히 평탄인 것과 그 섬유들이 모두 영이 아닌 것은 동치이다.
\end{slogan}
$M$을 평탄한 $R$-가군이라 하자.
다음은 서로 동치이다.
\begin{enumerate}
\item $M$은 충실히 평탄하다.
% P01-E1298
\item 모든 영이 아닌 $R$-가군 $N$에 대해 텐서곱
$M \otimes_R N$은 영이 아니다.
\item 모든 $\mathfrak p \in \Spec(R)$에 대해 텐서곱
$M \otimes_R \kappa(\mathfrak p)$은 영이 아니다.
\item $R$의 모든 극대 아이디얼 $\mathfrak m$에 대해 텐서곱
$M \otimes_R \kappa(\mathfrak m) = M/{\mathfrak m}M$은 영이 아니다.
\end{enumerate}
\end{lemma}

\begin{proof}
$M$이 충실히 평탄하고 $N \not = 0$이라고 하자.
보조정리 \ref{algebra-lemma-easy-ff}에 의해 영이 아닌 사상
$1 : N \to N$은 영이 아닌 사상
$M \otimes_R N \to M \otimes_R N$을 유도하므로
$M \otimes_R N \not = 0$이다. 따라서 (1)은 (2)를 함의한다.
함의 (2) $\Rightarrow$ (3) $\Rightarrow$ (4)는 자명하다.

\medskip\noindent
(4)를 가정하자. $N_1 \to N_2 \to N_3$가 복합체이고,
$N_1 \otimes_R M \to N_2\otimes_R M \to
N_3\otimes_R M$이 완전하다고 가정하자.
$H$를 이 복합체의 코호몰로지라 하자. 즉,
$H = \Ker(N_2 \to N_3)/\Im(N_1 \to N_2)$이다.
증명을 끝내기 위해 $H = 0$임을 보이겠다.
평탄성에 의해 $H \otimes_R M = 0$임을 안다.
$x \in H$를 택하고
$I = \{f \in R \mid fx = 0 \}$를 그 소멸자라 하자.
$R/I \subset H$이므로 $M$의 평탄성에 의해
$M/IM \subset H \otimes_R M = 0$을 얻는다.
$I \not =  R$이면 극대 아이디얼
$I \subset \mathfrak m \subset R$를 택할 수 있다.
이는 곧바로 모순을 준다.
\end{proof}

\begin{lemma}
\label{algebra-lemma-ff-rings}
$R \to S$를 평탄한 환 사상이라 하자.
다음은 서로 동치이다.
\begin{enumerate}
\item $R \to S$는 충실히 평탄하다.
\item 유도된 $\Spec$ 사이의 사상은 전사이다.
\item 모든 닫힌점 $x \in \Spec(R)$는
사상 $\Spec(S) \to \Spec(R)$의 상에 속한다.
\end{enumerate}
\end{lemma}

\begin{proof}
이는 보조정리 \ref{algebra-lemma-ff}에서 곧바로 따른다.
실제로 주석 \ref{algebra-remark-fundamental-diagram}에서
$\mathfrak p$가 상에 속하는 것과 환
$S \otimes_R \kappa(\mathfrak p)$이 영이 아닌 것은
동치임을 보았다.
\end{proof}

\begin{lemma}
\label{algebra-lemma-local-flat-ff}
국소환들 사이의 평탄한 국소환 준동형은 충실히 평탄하다.
\end{lemma}

\begin{proof}
보조정리 \ref{algebra-lemma-ff-rings}에서 곧바로 따른다.
\end{proof}

\noindent
평탄성은 국소화와 잘 맞물린다.

\begin{lemma}
\label{algebra-lemma-flat-localization}
$R$을 환이라 하고 $S \subset R$를 곱셈적 부분집합이라 하자.
\begin{enumerate}
\item 국소화 $S^{-1}R$는 평탄한 $R$-대수이다.
\item $M$이 $S^{-1}R$-가군이면, $M$이 평탄한 $R$-가군인 것과
$M$이 평탄한 $S^{-1}R$-가군인 것은 동치이다.
\item $M$이 $R$-가군이라고 하자. $M$이 평탄한 $R$-가군인 것과
$R$의 모든 소 아이디얼 $\mathfrak p$에 대해 $M_{\mathfrak p}$가
평탄한 $R_{\mathfrak p}$-가군인 것은 동치이다.
\item $M$이 $R$-가군이라고 하자. $M$이 평탄한 $R$-가군인 것과
$R$의 모든 극대 아이디얼 $\mathfrak m$에 대해
$M_{\mathfrak m}$이 평탄한 $R_{\mathfrak m}$-가군인 것은 동치이다.
\item $R \to A$가 환 사상이고, $M$이 $A$-가군이며,
$g_1, \ldots, g_m \in A$가 $A$의 단위 아이디얼을 생성한다고 하자.
그러면 $M$이 $R$ 위에서 평탄인 것과 각 국소화
$M_{g_i}$가 $R$ 위에서 평탄인 것은 동치이다.
\item $R \to A$가 환 사상이고 $M$이 $A$-가군이라고 하자.
그러면 $M$이 평탄한 $R$-가군인 것과 $A$의 모든 소 아이디얼
$\mathfrak q$에 대해 국소화 $M_{\mathfrak q}$가
평탄한 $R_{\mathfrak p}$-가군인 것은 동치이다
(여기서 $\mathfrak p$는 $\mathfrak q$ 아래에 놓인 $R$의 소 아이디얼이다).
\item $R \to A$가 환 사상이고 $M$이 $A$-가군이라고 하자.
그러면 $M$이 평탄한 $R$-가군인 것과 $A$의 모든 극대 아이디얼
$\mathfrak m$에 대해 국소화 $M_{\mathfrak m}$이
평탄한 $R_{\mathfrak p}$-가군인 것은 동치이다
(여기서 $\mathfrak p = R \cap \mathfrak m$이다).
\end{enumerate}
\end{lemma}

\begin{proof}
보조정리의 마지막 명제를 증명하자.
증명에서 국소화가 완전하고 텐서곱과 교환한다는 사실을 거듭 사용한다.
절 \ref{algebra-section-localization} 및 \ref{algebra-section-tensor-product}를 참조하라.

\medskip\noindent
$R \to A$가 환 사상이고 $M$이 $A$-가군이라고 하자.
$A$의 모든 극대 아이디얼 $\mathfrak m$에 대해
$M_{\mathfrak m}$이 평탄한 $R_{\mathfrak p}$-가군이라고 가정하자
(여기서 $\mathfrak p = R \cap \mathfrak m$이다).
$I \subset R$를 아이디얼이라 하자.
사상 $I \otimes_R M \to M$이 단사임을 보여야 한다.
이를 $A$-가군들의 사상으로 생각할 수 있다.
가정에 의해 국소화
$(I \otimes_R M)_{\mathfrak m} \to M_{\mathfrak m}$은 단사이다.
왜냐하면
$(I \otimes_R M)_{\mathfrak m} =
I_{\mathfrak p} \otimes_{R_{\mathfrak p}} M_{\mathfrak m}$
이기 때문이다. 따라서 보조정리 \ref{algebra-lemma-characterize-zero-local}에 의해
$I \otimes_R M \to M$의 핵은 0이다.
따라서 $M$은 $R$ 위에서 평탄하다.

\medskip\noindent
역으로 $M$이 $R$ 위에서 평탄하다고 하자.
$R$의 소 아이디얼 $\mathfrak p$ 위에 놓인 $A$의 소 아이디얼
$\mathfrak q$를 택하자. $I \subset R_{\mathfrak p}$를
아이디얼이라 하자. 사상
$I \otimes_{R_{\mathfrak p}} M_{\mathfrak q} \to M_{\mathfrak q}$가
단사임을 보여야 한다. 어떤 아이디얼 $J \subset R$에 대해
$I = J_{\mathfrak p}$로 쓸 수 있다. 그러면 사상
$I \otimes_{R_{\mathfrak p}} M_{\mathfrak q} \to M_{\mathfrak q}$는
단사 사상 $J \otimes_R M \to M$을 $\mathfrak q$에서
국소화한 것에 불과하다. 국소화는 완전하므로
$M_{\mathfrak q}$가 평탄한 $R_{\mathfrak p}$-가군임을 알 수 있다.

\medskip\noindent
이는 (7)과 (6)을 증명한다. 나머지 명제들은 마지막 명제에서
곧바로 따른다(증명은 생략한다).
\end{proof}

\begin{lemma}
\label{algebra-lemma-flat-going-down}
$R \to S$가 평탄하다고 하자. $\mathfrak p \subset \mathfrak p'$를
$R$의 소 아이디얼들이라 하고, $\mathfrak q' \subset S$를
$\mathfrak p'$로 사상되는 $S$의 소 아이디얼이라 하자.
그러면 $\mathfrak p$로 사상되는 소 아이디얼
$\mathfrak q \subset \mathfrak q'$가 존재한다.
\end{lemma}

\begin{proof}
보조정리 \ref{algebra-lemma-flat-localization}에 의해 국소환 사상
$R_{\mathfrak p'} \to S_{\mathfrak q'}$는 평탄하다.
보조정리 \ref{algebra-lemma-local-flat-ff}에 의해 이 국소환 사상은
충실히 평탄하다. 보조정리 \ref{algebra-lemma-ff-rings}에 의해
$\mathfrak p R_{\mathfrak p'}$로 사상되는 소 아이디얼이 존재한다.
이 소 아이디얼을 $S$로 역상시키면 원하는 것을 얻는다.
\end{proof}

\noindent
이 보조정리에 기술된 $R \to S$의 성질을 “하강 성질”이라고 한다.
정의 \ref{algebra-definition-going-up-down}를 참조하라.

\begin{lemma}
\label{algebra-lemma-colimit-faithfully-flat}
$R$을 환이라 하고 $\{S_i, \varphi_{ii'}\}$를
충실히 평탄한 $R$-대수들의 유향계라 하자.
그러면 $S = \colim_i S_i$는 충실히 평탄한 $R$-대수이다.
\end{lemma}

\begin{proof}
보조정리 \ref{algebra-lemma-colimit-flat}에 의해 $S$가 평탄함을 안다.
$\mathfrak m \subset R$를 극대 아이디얼이라 하자.
보조정리 \ref{algebra-lemma-ff-rings}에 의해 환
$S_i/\mathfrak m S_i$들 가운데 어느 것도 영환이 아니다.
따라서 $1$이 0과 같지 않으므로
$S/\mathfrak mS = \colim S_i/\mathfrak mS_i$도 영환이 아니다.
그러므로 $\Spec(S) \to \Spec(R)$의 상은 $\mathfrak m$을 포함하고,
보조정리 \ref{algebra-lemma-ff-rings}에 의해 $R \to S$가
충실히 평탄함을 알 수 있다.
\end{proof}





\section{지지집합과 소멸자}
\label{algebra-section-supp-and-ann}

\noindent
몇 가지 매우 기본적인 정의와 보조정리를 제시한다.

\begin{definition}
\label{algebra-definition-support-module}
$R$을 환이라 하고 $M$을 $R$-가군이라 하자.
{\it $M$의 지지집합}은 다음 집합이다.
$$
\text{Supp}(M)
=
\{
\mathfrak p \in \Spec(R)
\mid
M_{\mathfrak p} \not = 0
\}
$$
\end{definition}

\begin{lemma}
\label{algebra-lemma-support-zero}
\begin{slogan}
환 위의 가군의 지지집합이 공집합인 것과 그 가군이 영가군인 것은 동치이다.
\end{slogan}
$R$을 환이라 하고 $M$을 $R$-가군이라 하자. 그러면
$$
M = (0) \Leftrightarrow \text{Supp}(M) = \emptyset.
$$
\end{lemma}

\begin{proof}
실제로 보조정리 \ref{algebra-lemma-characterize-zero-local}은
$M$이 영이 아니면 $\text{Supp}(M)$가 항상 극대 아이디얼 하나를
포함한다는 것까지 보여 준다.
\end{proof}

\begin{definition}
\label{algebra-definition-annihilator}
$R$을 환이라 하고 $M$을 $R$-가군이라 하자.
\begin{enumerate}
\item 원소 $m \in M$이 주어졌을 때 {\it $m$의 소멸자}는
다음 아이디얼이다.
$$
\text{Ann}_R(m) = \text{Ann}(m) = \{f \in R \mid fm = 0\}.
$$
\item {\it $M$의 소멸자}는 다음 아이디얼이다.
$$
\text{Ann}_R(M) = \text{Ann}(M) = \{f \in R \mid fm = 0\ \forall m \in M\}.
$$
\end{enumerate}
\end{definition}

\begin{lemma}
\label{algebra-lemma-annihilator-flat-base-change}
$R \to S$를 평탄한 환 사상이라 하고, $M$을 $R$-가군,
$m \in M$을 원소라 하자. 그러면
$\text{Ann}_R(m) S = \text{Ann}_S(m \otimes 1)$이다.
$M$이 유한 $R$-가군이면
$\text{Ann}_R(M) S = \text{Ann}_S(M \otimes_R S)$이다.
\end{lemma}

\begin{proof}
$I = \text{Ann}_R(m)$로 놓자. 정의에 의해 완전열
$0 \to I \to R \to M$이 있으며, 여기서 사상 $R \to M$은
$f$를 $fm$으로 보낸다. 평탄성을 사용하여 완전열
$0 \to I \otimes_R S \to S \to M \otimes_R S$를 얻고,
이는 첫 번째 주장을 증명한다. $m_1, \ldots, m_n$이 $M$의
생성원들의 집합이면
$\text{Ann}_R(M) = \bigcap \text{Ann}_R(m_i)$이다. 마찬가지로
$\text{Ann}_S(M \otimes_R S) = \bigcap \text{Ann}_S(m_i \otimes 1)$이다.
$I_i = \text{Ann}_R(m_i)$로 놓자. 이제
$\bigcap_{i = 1, \ldots, n} (I_i S) = (\bigcap_{i = 1, \ldots, n} I_i)S$
임을 보이면 충분하다. 이는 보조정리
\ref{algebra-lemma-flat-intersect-ideals}이다.
\end{proof}

\begin{lemma}
\label{algebra-lemma-support-closed}
$R$을 환이라 하고 $M$을 $R$-가군이라 하자.
$M$이 유한이면 $\text{Supp}(M)$는 닫혀 있다.
더 정확히, $I = \text{Ann}(M)$가 $M$의 소멸자이면
$V(I) = \text{Supp}(M)$이다.
\end{lemma}

\begin{proof}
$V(I) = \text{Supp}(M)$임을 보이겠다.

\medskip\noindent
$\mathfrak p \in \text{Supp}(M)$라고 하자. 그러면
$M_{\mathfrak p} \not = 0$이다. $M_\mathfrak p$에서의 상이
영이 아닌 원소 $m \in M$을 택하자. 국소화 $M_\mathfrak p$의
구성에 의해 $m$의 소멸자는 $\mathfrak p$에 포함된다.
따라서 특히 $I = \text{Ann}(M)$도 $\mathfrak p$에 포함되어야 한다.

\medskip\noindent
역으로 $\mathfrak p \not \in \text{Supp}(M)$라고 하자.
그러면 $M_{\mathfrak p} = 0$이다.
$x_1, \ldots, x_r \in M$을 생성원들이라 하자.
보조정리 \ref{algebra-lemma-localization-colimit}에 의해
$f \in R$, $f\not\in \mathfrak p$인 원소가 존재하여
$M_f$에서 $x_i/1 = 0$이다. 따라서 어떤 $n_i \geq 1$에 대해
$f^{n_i} x_i = 0$이다. 그러므로 $n = \max\{n_i\}$에 대해
$f^nM = 0$이고, 원하는 결론을 얻는다.
\end{proof}

\begin{lemma}
\label{algebra-lemma-support-base-change}
$R \to R'$를 환 사상이라 하고 $M$을 유한 $R$-가군이라 하자.
그러면 $\text{Supp}(M \otimes_R R')$는
$\text{Supp}(M)$의 역상이다.
\end{lemma}

\begin{proof}
$\mathfrak p \in \text{Supp}(M)$라 하자.
나카야마 보조정리(보조정리 \ref{algebra-lemma-NAK})에 의해
$$
M \otimes_R \kappa(\mathfrak p) = M_\mathfrak p/\mathfrak p M_\mathfrak p
$$
는 영이 아닌 $\kappa(\mathfrak p)$-벡터 공간이다.
따라서 $\mathfrak p$ 위에 놓인 모든 소 아이디얼
$\mathfrak p' \subset R'$에 대해
$$
(M \otimes_R R')_{\mathfrak p'}/\mathfrak p' (M \otimes_R R')_{\mathfrak p'} =
(M \otimes_R R') \otimes_{R'} \kappa(\mathfrak p') =
M \otimes_R \kappa(\mathfrak p) \otimes_{\kappa(\mathfrak p)}
\kappa(\mathfrak p')
$$
가 영이 아님을 안다. 이는
$\mathfrak p' \in \text{Supp}(M \otimes_R R')$임을 뜻한다.
역으로, $\mathfrak p' \subset R'$가 임의의 소 아이디얼
$\mathfrak p \subset R$ 위에 놓이면
$$
(M \otimes_R R')_{\mathfrak p'} =
M_\mathfrak p \otimes_{R_\mathfrak p} R'_{\mathfrak p'}.
$$
따라서 $\mathfrak p \subset R$ 위에 놓인
$\mathfrak p' \in \text{Supp}(M \otimes_R R')$가 있으면
$\mathfrak p \in \text{Supp}(M)$이다.
\end{proof}

\begin{lemma}
\label{algebra-lemma-support-element}
$R$을 환, $M$을 $R$-가군, $m \in M$을 원소라 하자.
그러면 $\mathfrak p \in V(\text{Ann}(m))$인 것과
$m$이 $M_\mathfrak p$에서 0으로 사상되지 않는 것은 동치이다.
\end{lemma}

\begin{proof}
$M$을 $Rm \subset M$로 바꾸어도 된다. 그러면
(1) $\text{Ann}(m) = \text{Ann}(M)$이고,
(2) $m$이 $M_\mathfrak p$에서 0으로 사상되지 않는 것과
$\mathfrak p \in \text{Supp}(M)$인 것은 동치이다.
이제 결과는 보조정리 \ref{algebra-lemma-support-closed}에서 따른다.
\end{proof}

\begin{lemma}
\label{algebra-lemma-support-finite-presentation-constructible}
$R$을 환이라 하고 $M$을 $R$-가군이라 하자.
$M$이 유한 표시 $R$-가군이면 $\text{Supp}(M)$는
$\Spec(R)$의 닫힌 부분집합이고 그 여집합은 준콤팩트이다.
\end{lemma}

\begin{proof}
다음 표시를 택하자.
$$
R^{\oplus m} \longrightarrow R^{\oplus n} \longrightarrow M \to 0
$$
$A \in \text{Mat}(n \times m, R)$를 첫 번째 사상의 행렬이라 하자.
나카야마 보조정리 \ref{algebra-lemma-NAK}에 의해
$$
M_{\mathfrak p} \not = 0 \Leftrightarrow
M \otimes \kappa(\mathfrak p) \not = 0 \Leftrightarrow
\text{rank}(A \bmod \mathfrak p) < n.
$$
임을 안다. 따라서 $I$가 $A$의 $n \times n$ 소행렬식들로 생성되는
$R$의 아이디얼이면 $\text{Supp}(M) = V(I)$이다.
$I$는 유한 생성이므로 $I = (f_1, \ldots, f_t)$로 쓸 수 있고,
$\Spec(R) \setminus V(I)$가 표준 열린집합들
$D(f_i)$의 유한 합집합이므로 준콤팩트임을 알 수 있다.
\end{proof}

\begin{lemma}
\label{algebra-lemma-support-quotient}
$R$을 환이라 하고 $M$을 $R$-가군이라 하자.
\begin{enumerate}
\item $M$이 유한이면 $M/IM$의 지지집합은
$\text{Supp}(M) \cap V(I)$이다.
\item $N \subset M$이면 $\text{Supp}(N) \subset
\text{Supp}(M)$이다.
\item $Q$가 $M$의 몫가군이면 $\text{Supp}(Q) \subset
\text{Supp}(M)$이다.
\item $0 \to N \to M \to Q \to 0$이 짧은 완전열이면
$\text{Supp}(M) = \text{Supp}(Q) \cup \text{Supp}(N)$이다.
\end{enumerate}
\end{lemma}

\begin{proof}
함자들 $M \mapsto M_{\mathfrak p}$는 완전하다.
이로부터 첫 번째를 제외한 모든 주장이 곧바로 따른다.
첫 번째 주장을 위해서는 $M_\mathfrak p \not = 0$이고
$I \subset \mathfrak p$이면
$(M/IM)_{\mathfrak p} = M_\mathfrak p/IM_\mathfrak p \not = 0$임을
보여야 한다. 이는 나카야마 보조정리 \ref{algebra-lemma-NAK}에서 따른다.
\end{proof}





\section{상승 성질과 하강 성질}
\label{algebra-section-going-up}

\noindent
$\mathfrak p$, $\mathfrak p'$를 환 $R$의 소 아이디얼들이라 하자.
$X = \Spec(R)$에 자리스키 위상을 주자.
$\mathfrak p$에 대응하는 점을 $x \in X$로,
$\mathfrak p'$에 대응하는 점을 $x' \in X$로 쓰자.
그러면 다음이 성립한다.
$$
x' \leadsto x \Leftrightarrow \mathfrak p' \subset \mathfrak p.
$$
말로 쓰면, $x$가 $x'$의 특수화인 것과
$\mathfrak p' \subset \mathfrak p$인 것은 동치이다.
용어와 표기법은 위상수학, 절
\ref{topology-section-specialization}을 참조하라.

\begin{definition}
\label{algebra-definition-going-up-down}
$\varphi : R \to S$를 환 사상이라 하자.
\begin{enumerate}
\item $R$의 소 아이디얼들 $\mathfrak p \subset \mathfrak p'$와
$\mathfrak p$ 위에 놓인 $S$의 소 아이디얼 $\mathfrak q$가 주어졌을 때,
(a) $\mathfrak q \subset \mathfrak q'$이고 (b) $\mathfrak q'$가
$\mathfrak p'$ 위에 놓이는 $S$의 소 아이디얼 $\mathfrak q'$가
존재하면 $\varphi : R \to S$가 {\it 상승 성질}을 만족한다고 한다.
\item $R$의 소 아이디얼들 $\mathfrak p \subset \mathfrak p'$와
$\mathfrak p'$ 위에 놓인 $S$의 소 아이디얼 $\mathfrak q'$가 주어졌을 때,
(a) $\mathfrak q \subset \mathfrak q'$이고 (b) $\mathfrak q$가
$\mathfrak p$ 위에 놓이는 $S$의 소 아이디얼 $\mathfrak q$가
존재하면 $\varphi : R \to S$가 {\it 하강 성질}을 만족한다고 한다.
\end{enumerate}
\end{definition}

\noindent
% P01-E1300
지금까지 다음 경우들을 보았다.
\begin{enumerate}
\item 정수적인 환 사상은 상승 성질을 만족한다.
보조정리 \ref{algebra-lemma-integral-going-up}을 참조하라.
\item 특수한 경우로 유한 환 사상은 상승 성질을 만족한다.
\item 특수한 경우로 몫사상 $R \to R/I$는 상승 성질을 만족한다.
\item 평탄한 환 사상은 하강 성질을 만족한다.
보조정리 \ref{algebra-lemma-flat-going-down}을 참조하라.
\item 특수한 경우로 모든 국소화는 하강 성질을 만족한다.
\item 정역들의 확대 $R \subset S$에서 $R$이 정규이고
$S$가 $R$ 위에서 정수적이면 하강 성질을 만족한다.
명제 \ref{algebra-proposition-going-down-normal-integral}을 참조하라.
\end{enumerate}
하강 성질이 성립하는 또 다른 경우는 다음과 같다.

\begin{lemma}
\label{algebra-lemma-open-going-down}
$R \to S$를 환 사상이라 하자. 유도된 사상
$\varphi : \Spec(S) \to \Spec(R)$가 열린 사상이면
$R \to S$는 하강 성질을 만족한다.
\end{lemma}

\begin{proof}
$\mathfrak p \subset \mathfrak p' \subset R$이고
$\mathfrak q' \subset S$가 $\mathfrak p'$ 위에 놓인다고 하자.
$\varphi$가 열린 사상이므로 모든 $g \in S$,
$g \not \in \mathfrak q'$에 대해 $\mathfrak p$가
$D(g) \subset \Spec(S)$의 상에 속함을 안다. 다르게 말하면
$S_g \otimes_R \kappa(\mathfrak p)$은 영이 아니다.
$S_{\mathfrak q'}$가 이 $S_g$들의 유향 귀납적 극한이므로
$S_{\mathfrak q'} \otimes_R \kappa(\mathfrak p)$도 영이 아니다.
보조정리 \ref{algebra-lemma-localization-colimit} 및
\ref{algebra-lemma-tensor-products-commute-with-limits}를 참조하라.
따라서 원하는 대로 $\mathfrak p$는
$\Spec(S_{\mathfrak q'}) \to \Spec(R)$의 상에 속한다.
\end{proof}

\begin{lemma}
\label{algebra-lemma-going-up-down-specialization}
$R \to S$를 환 사상이라 하자.
\begin{enumerate}
\item $R \to S$가 하강 성질을 만족하는 것과 사상
$\Spec(S) \to \Spec(R)$을 따라 일반화를 올릴 수 있는 것은 동치이다.
위상수학, 정의 \ref{topology-definition-lift-specializations}를 참조하라.
\item $R \to S$가 상승 성질을 만족하는 것과 사상
$\Spec(S) \to \Spec(R)$을 따라 특수화를 올릴 수 있는 것은 동치이다.
위상수학, 정의 \ref{topology-definition-lift-specializations}를 참조하라.
\end{enumerate}
\end{lemma}

\begin{proof}
생략한다.
\end{proof}

\begin{lemma}
\label{algebra-lemma-going-up-down-composition}
$R \to S$와 $S \to T$가 하강 성질을 만족하는 환 사상들이라고 하자.
그러면 $R \to T$도 하강 성질을 만족한다.
상승 성질에 대해서도 마찬가지이다.
\end{lemma}

\begin{proof}
보조정리 \ref{algebra-lemma-going-up-down-specialization}에 따르면 이는
위상수학, 보조정리
\ref{topology-lemma-lift-specialization-composition}에서 따른다
\end{proof}

\begin{lemma}
\label{algebra-lemma-image-stable-specialization-closed}
$R \to S$를 환 사상이라 하고, $T \subset \Spec(R)$를
$\Spec(S)$의 상이라 하자. $T$가 특수화에 대해 안정이면
$T$는 닫혀 있다.
\end{lemma}

\begin{proof}
두 가지 증명을 제시한다.

\medskip\noindent
첫 번째 증명. $\mathfrak p \subset R$를 소 아이디얼이라 하고,
이에 대응하는 $\Spec(R)$의 점이 $T$의 폐포에 속한다고 하자.
이는 모든 $f \in R$, $f \not \in \mathfrak p$에 대해
$D(f) \cap T \not = \emptyset$이라는 뜻이다.
$D(f) \cap T$는 $\Spec(S_f)$의 $\Spec(R)$ 안에서의 상임에 유의하자.
따라서 $S_f \not = 0$임을 얻는다. 다르게 말하면 환 $S_f$에서
$1 \not = 0$이다. $S_{\mathfrak p}$는 환들 $S_f$의
유향 귀납적 극한이므로 $S_{\mathfrak p}$에서도 $1 \not = 0$이다.
다르게 말하면 $S_{\mathfrak p} \not = 0$이고,
$\Spec(S_{\mathfrak p}) \to \Spec(S) \to \Spec(R)$의 상을 생각하면
$\mathfrak p' \subset \mathfrak p$를 만족하는
$\mathfrak p' \in T$가 존재함을 알 수 있다.
$T$가 특수화에 대해 닫혀 있다고 가정했으므로
원하는 대로 $\mathfrak p$는 $T$의 점이다.

\medskip\noindent
두 번째 증명. $I = \Ker(R \to S)$로 놓자.
$R$을 $R/I$로 바꾸어도 된다. 이 경우 환 사상
$R \to S$는 단사이다.
보조정리 \ref{algebra-lemma-injective-minimal-primes-in-image}에 의해
$R$의 모든 극소 소 아이디얼은 상 $T$에 포함된다.
따라서 $T$가 특수화에 대해 안정이면 모든 소 아이디얼을 포함한다.
\end{proof}

\begin{lemma}
\label{algebra-lemma-going-up-closed}
$R \to S$를 환 사상이라 하자. 다음은 서로 동치이다.
\begin{enumerate}
\item $R \to S$에 대해 상승 성질이 성립한다.
\item 사상 $\Spec(S) \to \Spec(R)$는 닫힌 사상이다.
\end{enumerate}
\end{lemma}

\begin{proof}
위상공간들의 닫힌 사상을 따라 특수화를 올릴 수 있다는 것은
일반적인 사실이다. 위상수학, 보조정리
\ref{topology-lemma-closed-open-map-specialization}를 참조하라.
따라서 두 번째 조건은 첫 번째 조건을 함의한다.

\medskip\noindent
$R \to S$에 대해 상승 성질이 성립한다고 하자.
$V(I) \subset \Spec(S)$를 닫힌집합이라 하자.
$V(I)$의 $\Spec(R)$ 안에서의 상이 닫혀 있음을 보이고자 한다.
환 사상 $S \to S/I$는 분명히 상승 성질을 만족한다.
따라서 보조정리 \ref{algebra-lemma-going-up-down-composition}에 의해
$R \to S \to S/I$도 상승 성질을 만족한다.
$S$를 $S/I$로 바꾸면 $\Spec(S)$의 $\Spec(R)$ 안에서의 상
$T$가 닫혀 있음을 보이면 충분하다.
위상수학, 보조정리
\ref{topology-lemma-open-closed-specialization} 및
\ref{topology-lemma-lift-specializations-images}에 의해
이 상은 특수화에 대해 안정이다. 따라서 결과는 보조정리
\ref{algebra-lemma-image-stable-specialization-closed}에서 따른다.
\end{proof}

\begin{lemma}
\label{algebra-lemma-constructible-stable-specialization-closed}
$R$을 환이라 하고 $E \subset \Spec(R)$를 구성가능 부분집합이라 하자.
\begin{enumerate}
\item $E$가 특수화에 대해 안정이면 $E$는 닫혀 있다.
\item $E$가 일반화에 대해 안정이면 $E$는 열려 있다.
\end{enumerate}
\end{lemma}

\begin{proof}
첫 번째 증명. 첫 번째 주장은 보조정리
\ref{algebra-lemma-image-stable-specialization-closed}와 보조정리
\ref{algebra-lemma-constructible-is-image}를 결합하면 따른다.
두 번째 주장은 구성가능 집합의 여집합이 구성가능하다는 사실
(위상수학, 보조정리 \ref{topology-lemma-constructible} 참조),
이 보조정리의 첫 번째 부분, 그리고 위상수학, 보조정리
\ref{topology-lemma-open-closed-specialization}에서 따른다.

\medskip\noindent
두 번째 증명. 보조정리 \ref{algebra-lemma-spec-spectral}에 의해
$\Spec(R)$는 스펙트럴 공간이므로, 이는 위상수학, 보조정리
\ref{topology-lemma-constructible-stable-specialization-closed}의
특수한 경우이다.
\end{proof}

\begin{proposition}
\label{algebra-proposition-fppf-open}
$R \to S$가 평탄하고 유한 표시라고 하자.
그러면 $\Spec(S) \to \Spec(R)$는 열린 사상이다.
더 일반적으로 이는 하강 성질을 만족하는 모든 유한 표시 환 사상
$R \to S$에 대해 성립한다.
\end{proposition}

\begin{proof}
$R \to S$가 평탄하면 보조정리 \ref{algebra-lemma-flat-going-down}에 의해
$R \to S$는 하강 성질을 만족한다. 따라서 이 보조정리를 증명하기 위해
$R \to S$가 유한 표시이고 하강 성질을 만족한다고 가정해도 된다.

\medskip\noindent
표준 열린집합들 $D(g) \subset \Spec(S)$, $g \in S$가
위상의 기저를 이루므로, $D(g)$의 상이 열려 있음을 보이면 충분하다.
$\Spec(S_g) \to \Spec(S)$는 $\Spec(S_g)$에서 $D(g)$ 위로의
위상동형임을 상기하자(보조정리 \ref{algebra-lemma-standard-open}).
$S \to S_g$가 하강 성질을 만족하므로(위 참조),
보조정리 \ref{algebra-lemma-going-up-down-composition}에 의해
$R \to S_g$도 하강 성질을 만족한다. 따라서 $S$를 $S_g$로
바꾸고 나면 그 상이 열려 있음을 보이는 것으로 충분하다.
슈발레 정리(정리 \ref{algebra-theorem-chevalley})에 의해 그 상은
구성가능 집합 $E$이다. 또한 $R \to S$가 하강 성질을 만족하므로
$E$는 일반화에 대해 안정이다. 위상수학, 보조정리
\ref{topology-lemma-open-closed-specialization} 및
\ref{topology-lemma-lift-specializations-images}를 참조하라.
따라서 보조정리 \ref{algebra-lemma-constructible-stable-specialization-closed}에
의해 $E$는 열려 있다.
\end{proof}

\begin{lemma}
\label{algebra-lemma-same-image}
$k$를 체라 하고 $R$, $S$를 $k$-대수들이라 하자.
$S' \subset S$를 $k$-부분대수라 하고 $f \in S' \otimes_k R$라 하자.
다음 가환 도식에서 대각선 화살표들의 상은 서로 같다.
$$
\xymatrix{
\Spec((S \otimes_k R)_f) \ar[rd] \ar[rr] & &
\Spec((S' \otimes_k R)_f) \ar[ld] \\
& \Spec(R) &
}
$$
\end{lemma}

\begin{proof}
$\mathfrak p \subset R$가 남서쪽 화살표의 상에 속한다고 하자.
이는 보조정리 \ref{algebra-lemma-in-image}에 의해
$$
(S' \otimes_k R)_f \otimes_R \kappa(\mathfrak p)
=
(S' \otimes_k \kappa(\mathfrak p))_f
$$
가 영환이 아니라는 뜻이다. 즉,
$S' \otimes_k \kappa(\mathfrak p)$은 영환이 아니고,
그 안에서 $f$의 상은 멱영이 아니다. 환 사상
$S' \otimes_k \kappa(\mathfrak p) \to S \otimes_k \kappa(\mathfrak p)$은
단사이다. 따라서 $S \otimes_k \kappa(\mathfrak p)$도 영환이 아니고,
그 안에서 $f$의 상도 멱영이 아니다.
그러므로 $(S \otimes_k R)_f \otimes_R \kappa(\mathfrak p)$은
영환이 아니다. 따라서 보조정리 \ref{algebra-lemma-in-image}에 의해
원하는 대로 $\mathfrak p$가 남동쪽 화살표의 상에 속함을 안다.
\end{proof}

\begin{lemma}
\label{algebra-lemma-map-into-tensor-algebra-open}
$k$를 체라 하자.
$R$과 $S$를 $k$-대수들이라 하자.
사상 $\Spec(S \otimes_k R) \to \Spec(R)$는 열린 사상이다.
\end{lemma}

\begin{proof}
$f \in S \otimes_k R$라 하자.
표준 열린집합 $D(f)$의 상이 열려 있음을 보이면 충분하다.
$S' \subset S$를 $f \in S' \otimes_k R$를 만족하는
유한형 $k$-부분대수라 하자. 사상 $R \to S' \otimes_k R$는
평탄하고 유한 표시이므로, 명제 \ref{algebra-proposition-fppf-open}에 의해
$\Spec((S' \otimes_k R)_f) \to \Spec(R)$의 상 $U$는 열려 있다.
보조정리 \ref{algebra-lemma-same-image}에 의해 이것은 $D(f)$의 상이기도 하므로
결론이 따른다.
\end{proof}

\noindent
때때로 유용한 까다로운 보조정리를 하나 제시한다.

\begin{lemma}
\label{algebra-lemma-unique-prime-over-localize-below}
$R \to S$를 환 사상이라 하자.
$\mathfrak p \subset R$를 소 아이디얼이라 하자.
다음을 가정하자.
\begin{enumerate}
\item $\mathfrak p$ 위에 놓인 유일한 소 아이디얼
$\mathfrak q \subset S$가 존재한다.
\item 다음 가운데 하나가 성립한다.
\begin{enumerate}
\item $R \to S$에 대해 상승 성질이 성립한다.
\item $R \to S$에 대해 하강 성질이 성립하고,
$R$의 각 소 아이디얼 위에 놓이는 $S$의 소 아이디얼은 많아야 하나이다.
\end{enumerate}
\end{enumerate}
그러면 $S_{\mathfrak p} = S_{\mathfrak q}$이다.
\end{lemma}

\begin{proof}
$\Spec(S_{\mathfrak p})$의 점에 대응하는 임의의 소 아이디얼
$\mathfrak q' \subset S$를 생각하자. 이는
$\mathfrak p' = R \cap \mathfrak q'$가 $\mathfrak p$에
포함된다는 뜻이다. 그림으로 나타내면 다음과 같다.
$$
\xymatrix{
\mathfrak q' \ar@{-}[d] \ar@{-}[r] & ? \ar@{-}[r] \ar@{-}[d] & S \ar@{-}[d] \\
\mathfrak p' \ar@{-}[r] & \mathfrak p \ar@{-}[r] & R
}
$$
(1)과 (2)(a)를 가정하자.
상승 성질에 의해 소 아이디얼 $\mathfrak q'' \subset S$가 존재하며,
$\mathfrak q' \subset \mathfrak q''$이고 $\mathfrak q''$는
$\mathfrak p$ 위에 놓인다.
$\mathfrak q$의 유일성에 의해
$\mathfrak q'' = \mathfrak q$임을 얻는다. 다르게 말하면
$\mathfrak q'$는 $\Spec(S_{\mathfrak q})$의 점을 정한다.

\medskip\noindent
(1)과 (2)(b)를 가정하자.
하강 성질에 의해 $\mathfrak p'$ 위에 놓이는 소 아이디얼
$\mathfrak q'' \subset \mathfrak q$가 존재한다.
$\mathfrak p'$ 위에 놓인 소 아이디얼의 유일성에 의해
$\mathfrak q' = \mathfrak q''$임을 안다. 다르게 말하면
$\mathfrak q'$는 $\Spec(S_{\mathfrak q})$의 점을 정한다.

\medskip\noindent
두 경우 모두 사상
$\Spec(S_{\mathfrak q}) \to \Spec(S_{\mathfrak p})$가
전단사임을 얻는다. 이는 분명 $S - \mathfrak q$의 모든 원소가
$S_{\mathfrak p}$에서 모두 가역이라는 뜻이다. 다르게 말하면
$S_{\mathfrak p} = S_{\mathfrak q}$이다.
\end{proof}

\noindent
다음 보조정리는 평탄한 환 사상에 대한 하강 성질을 일반화한다.

\begin{lemma}
\label{algebra-lemma-going-down-flat-module}
$R \to S$를 환 사상이라 하고, $N$을 $R$ 위에서 평탄한
유한 $S$-가군이라 하자.
$\text{Supp}(N) \subset \Spec(S)$에 유도위상을 주자.
그러면 $\text{Supp}(N) \to \Spec(R)$을 따라 일반화를 올릴 수 있다.
\end{lemma}

\begin{proof}
명제의 뜻은 다음과 같다.
$\mathfrak p \subset \mathfrak p' \subset R$를 소 아이디얼들이라 하자.
% P01-E1343
$\mathfrak q' \subset S$를
$\mathfrak q' \in \text{Supp}(N)$인 소 아이디얼이라 하자.
그러면 $\mathfrak p$ 위에 놓이고
$\mathfrak q \in \text{Supp}(N)$인 소 아이디얼
$\mathfrak q \subset \mathfrak q'$가 존재한다.
$N$이 $R$ 위에서 평탄하므로 보조정리
\ref{algebra-lemma-flat-localization}에 의해 $N_{\mathfrak q'}$는
$R_{\mathfrak p'}$ 위에서 평탄하다.
$N_{\mathfrak q'}$는 $S_{\mathfrak q'}$ 위에서 유한이고,
$\mathfrak q' \in \text{Supp}(N)$이므로 영이 아니다.
따라서 나카야마 보조정리 \ref{algebra-lemma-NAK}에 의해
$N_{\mathfrak q'} \otimes_{S_{\mathfrak q'}} \kappa(\mathfrak q')$는
영이 아니다. 그러므로
$N_{\mathfrak q'} \otimes_{R_{\mathfrak p'}} \kappa(\mathfrak p')$도
영이 아니다. 보조정리 \ref{algebra-lemma-ff}에 의해
$N_{\mathfrak q'} \otimes_{R_{\mathfrak p'}} \kappa(\mathfrak p)$도
영이 아님을 얻는다. 유한 영이 아닌 가군
$N_{\mathfrak q'} \otimes_{R_{\mathfrak p'}} \kappa(\mathfrak p)$의
소멸자를
$J \subset S_{\mathfrak q'} \otimes_{R_{\mathfrak p'}} \kappa(\mathfrak p)$라 하자.
$J$는 진 아이디얼이므로
$S_{\mathfrak q'} \otimes_{R_{\mathfrak p'}} \kappa(\mathfrak p)/J$의
소 아이디얼에 대응하는 소 아이디얼 $\mathfrak q \subset S$를
택할 수 있다. 이 소 아이디얼은 $N$의 지지집합에 속하고,
$\mathfrak p$ 위에 놓이며, 원하는 대로 $\mathfrak q'$에 포함된다.
\end{proof}



\section{분리 확대}
\label{algebra-section-separability}

\noindent
이 절에서는 비대수적 체확대에 대한 분리성을 다룬다.
이는 기하적으로 기약인 대수의 개념과 밀접하게 관련된다.
정의 \ref{algebra-definition-geometrically-reduced}를 참조하라.

\begin{definition}
\label{algebra-definition-separable-field-extension}
$K/k$를 체확대라 하자.
\begin{enumerate}
\item $K/k$의 초월기저 $\{x_i; i \in I\}$가 존재하여
$K/k(x_i; i \in I)$가 분리 가능 대수적 확대이면,
$K$가 {\it $k$ 위에서 분리적으로 생성된다}고 한다.
\item 모든 부분확대 $k \subset K' \subset K$에서 $K'$가
$k$ 위에서 유한 생성되었을 때마다 $K'/k$가 분리적으로 생성되면,
$K$가 {\it $k$ 위에서 분리 가능하다}고 한다.
\end{enumerate}
\end{definition}

\noindent
이 다소 어색한 정의만으로는 분리적으로 생성된 체확대 자체가
분리 가능한지 분명하지 않다. 실제로는 그렇다는 것이
보조정리 \ref{algebra-lemma-separably-generated-separable}에서 밝혀진다.

\begin{lemma}
\label{algebra-lemma-subextensions-are-separable}
$K/k$를 분리 가능 체확대라 하자.
임의의 부분확대 $K/K'/k$에 대해 체확대 $K'/k$는 분리 가능하다.
\end{lemma}

\begin{proof}
정의에서 바로 따른다.
\end{proof}

\begin{lemma}
\label{algebra-lemma-generating-finitely-generated-separable-field-extensions}
% P01-E1344
$K/k$를 분리적으로 생성되고 유한 생성된 체확대라 하자.
$r = \text{trdeg}_k(K)$라 놓자. 그러면 다음을 만족하는
$K$의 원소 $x_1, \ldots, x_{r + 1}$가 존재한다.
\begin{enumerate}
\item $x_1, \ldots, x_r$는 $k$ 위의 $K$의 초월기저이다.
\item $K = k(x_1, \ldots, x_{r + 1})$이다.
\item $x_{r + 1}$는 $k(x_1, \ldots, x_r)$ 위에서 분리 가능하다.
\end{enumerate}
\end{lemma}

\begin{proof}
정의와 체론, 보조정리 \ref{fields-lemma-primitive-element}를 결합하면 된다.
\end{proof}

\begin{lemma}
\label{algebra-lemma-make-separably-generated}
$K/k$를 유한 생성 체확대라 하자. 다음과 같은 도식이 존재한다.
$$
\xymatrix{
K \ar[r] & K' \\
k \ar[u] \ar[r] & k' \ar[u]
}
$$
여기서 $k'/k$, $K'/K$는 유한 순수 비분리 체확대이고,
$K'/k'$는 분리적으로 생성된 체확대이다.
\end{lemma}

\begin{proof}
이 보조정리는 $k$의 표수가 $p > 0$일 때에만 흥미롭다.
$x_1, \ldots, x_r$를 $k$ 위의 $K$의 초월기저로 택하자.
$K$는 $k$ 위에서 유한 생성되므로
$k(x_1, \ldots, x_r) \subset K$는 유한 확대이다.
체론, 보조정리 \ref{fields-lemma-separable-first}에서 얻는 부분확대를
$K/K_{sep}/k(x_1, \ldots, x_r)$라 하자.
$K = K_{sep}$이면 증명이 끝난다.
$d = [K : K_{sep}]$에 관해 귀납법을 사용하겠다.

\medskip\noindent
$d > 1$이라 하자. $\alpha = \beta^p \in K_{sep}$이고
$\beta \not \in K_{sep}$인 $\beta \in K$를 택하자.
$P = T^n + a_1T^{n - 1} + \ldots + a_n$를
$k(x_1, \ldots, x_r)$ 위의 $\alpha$의 최소다항식이라 하자.
각 $a_i$가 $k'(x_1^{1/p}, \ldots, x_r^{1/p})$에서
$p$제곱이 되도록 $p$제곱근들을 첨가하여 얻는
유한 순수 비분리 확대 $k'/k$를 택하자.
그러한 확대가 존재하며, 자세한 내용은 생략한다.
다음 도식에 들어맞는 체를 $L$이라 하자.
$$
\xymatrix{
K \ar[r] & L \\
k(x_1, \ldots, x_r) \ar[u] \ar[r] & k'(x_1^{1/p}, \ldots, x_r^{1/p}) \ar[u]
}
$$
$L$은 $K$와 $k'(x_1^{1/p}, \ldots, x_r^{1/p})$의
합성체라고 가정해도 되고, 실제로 그렇게 하자.
체론, 보조정리 \ref{fields-lemma-separable-first}에서 얻는 부분확대를
$L/L_{sep}/k'(x_1^{1/p}, \ldots, x_r^{1/p})$라 하자.
그러면 $L_{sep}$는 $K_{sep}$와
$k'(x_1^{1/p}, \ldots, x_r^{1/p})$의 합성체이다.
원소 $\alpha \in L_{sep}$는 계수가 모두
$k'(x_1^{1/p}, \ldots, x_r^{1/p})$에서 $p$제곱이고
근들이 서로 다른 다항식 $P$의 영점이다.
체론, 보조정리 \ref{fields-lemma-pth-root}에 의해
어떤 $\alpha' \in L_{sep}$에 대해
$\alpha = (\alpha')^p$임을 안다.
이는 분명 $\beta$가 $\alpha' \in L_{sep}$로 사상된다는 뜻이다.
다르게 말하면 다음과 같은 체의 탑을 얻는다.
$$
\xymatrix{
K \ar[r] & L \\
K_{sep}(\beta) \ar[r] \ar[u] & L_{sep} \ar[u] \\
K_{sep} \ar[r] \ar[u] & L_{sep} \ar@{=}[u] \\
k(x_1, \ldots, x_r) \ar[u] \ar[r] & k'(x_1^{1/p}, \ldots, x_r^{1/p}) \ar[u] \\
k \ar[r] \ar[u] & k' \ar[u]
}
$$
따라서 이 구성은
$[L : L_{sep}] < [K : K_{sep}]$인 새로운 상황을 준다.
귀납가정에 의해 확대 $L/k'$에 대해 보조정리와 같은
$k' \subset k''$ 및 $L \subset L'$를 찾을 수 있다.
그러면 확대 $k''/k$와 $L'/K$가 확대 $K/k$에 대해 원하는 조건을 만족한다.
이로써 보조정리가 증명되었다.
\end{proof}



\section{기하적으로 축소인 대수}
\label{algebra-section-geometrically-reduced}

\noindent
기하적으로 축소인 대수에 관한 주된 결과는
보조정리 \ref{algebra-lemma-geometrically-reduced-finite-purely-inseparable-extension}이다.
독자에게 정의를 읽은 뒤 그 보조정리로 건너뛸 것을 권한다.

\begin{definition}
\label{algebra-definition-geometrically-reduced}
$k$를 체라 하고 $S$를 $k$-대수라 하자.
모든 체확대 $K/k$에 대해 $K$-대수 $K \otimes_k S$가
축소이면, $S$가 {\it $k$ 위에서 기하적으로 축소}라고 한다.
\end{definition}

\noindent
$k$를 체라 하고 $S$를 축소 $k$-대수라 하자.
$S$가 기하적으로 축소인지 확인하려면
$\overline{k} \otimes_k S$가 축소인지 확인하는 것으로 충분하다.
여기서 $\overline{k}$는 $k$의 대수적 폐포를 뜻한다.
실제로는 유한 순수 비분리 체확대 $k'/k$에 대해서만 이를 확인해도 충분하다.
보조정리
\ref{algebra-lemma-geometrically-reduced-finite-purely-inseparable-extension}를 참조하라.

\begin{lemma}
\label{algebra-lemma-subalgebra-separable}
기하적으로 축소인 대수의 기본 성질들은 다음과 같다.
$k$를 체라 하고 $S$를 $k$-대수라 하자.
\begin{enumerate}
\item $S$가 $k$ 위에서 기하적으로 축소이면 모든
$k$-부분대수도 그러하다.
\item $S$의 모든 유한 생성 $k$-부분대수가 기하적으로 축소이면
$S$도 기하적으로 축소이다.
\item 기하적으로 축소인 $k$-대수들의 유향 귀납적 극한은
기하적으로 축소이다.
\item $S$가 $k$ 위에서 기하적으로 축소이면 $S$의 임의의 국소화도
$k$ 위에서 기하적으로 축소이다.
\end{enumerate}
\end{lemma}

\begin{proof}
% P01-E1345
생략한다. 두 번째와 세 번째 성질은 텐서곱이 귀납적 극한과
가환한다는 사실에서 따른다.
\end{proof}

\begin{lemma}
\label{algebra-lemma-geometrically-reduced-permanence}
$k$를 체라 하자.
$R$이 $k$ 위에서 기하적으로 축소이고 $S \subset R$가
곱셈적 부분집합이면, 국소화 $S^{-1}R$도 $k$ 위에서
기하적으로 축소이다.
$R$이 $k$ 위에서 기하적으로 축소이면 $R[x]$도 $k$ 위에서
기하적으로 축소이다.
\end{lemma}

\begin{proof}
생략한다. 힌트: 축소환의 국소화는 축소이고,
국소화는 텐서곱과 가환한다.
\end{proof}

\noindent
다음 보조정리들의 증명에서는 다음 관찰을 거듭 사용한다.
$R' \subset R$와 $S' \subset S$가 $k$-대수의 포함이라고 하자.
그러면 사상 $R' \otimes_k S' \to R \otimes_k S$는 단사이다.

\begin{lemma}
\label{algebra-lemma-limit-argument}
$k$를 체라 하고 $R$, $S$를 $k$-대수들이라 하자.
\begin{enumerate}
\item $R \otimes_k S$가 축소가 아니면,
$R' \otimes_k S'$가 축소가 아니게 하는 유한 생성 부분대수
$R' \subset R$, $S' \subset S$가 존재한다.
\item $R \otimes_k S$가 영이 아닌 영인자를 가지면,
$R' \otimes_k S'$가 영이 아닌 영인자를 가지게 하는 유한 생성 부분대수
$R' \subset R$, $S' \subset S$가 존재한다.
\item $R \otimes_k S$가 자명하지 않은 멱등원을 가지면,
$R' \otimes_k S'$가 자명하지 않은 멱등원을 가지게 하는 유한 생성 부분대수
$R' \subset R$, $S' \subset S$가 존재한다.
\end{enumerate}
\end{lemma}

\begin{proof}
$z \in R \otimes_k S$가 멱영이라고 하자. 이를
$z = \sum_{i = 1, \ldots, n} x_i \otimes y_i$로 쓸 수 있다.
따라서 $x_i$들이 생성하는 $k$-부분대수를 $R'$로,
$y_i$들이 생성하는 $k$-부분대수를 $S'$로 택할 수 있다.
두 번째와 세 번째 주장도 같은 방법으로 증명된다.
\end{proof}

\begin{lemma}
\label{algebra-lemma-geometrically-reduced-any-reduced-base-change}
$k$를 체라 하자.
$S$를 기하적으로 축소인 $k$-대수라 하고,
$R$을 임의의 축소 $k$-대수라 하자.
그러면 $R \otimes_k S$는 축소이다.
\end{lemma}

\begin{proof}
보조정리 \ref{algebra-lemma-limit-argument}에 의해
$R$이 $k$ 위에서 유한형이라고 가정해도 된다.
그러면 축소 뇌터 환인 $R$은 체들의 유한 곱에 매입된다.
보조정리 \ref{algebra-lemma-total-ring-fractions-no-embedded-points},
\ref{algebra-lemma-Noetherian-irreducible-components}, 그리고
\ref{algebra-lemma-minimal-prime-reduced-ring}을 참조하라.
따라서 $R$이 체들의 유한 곱이라고 가정해도 된다.
이 경우 $R \otimes_k S$가 축소라는 것은
정의 \ref{algebra-definition-geometrically-reduced}에서 따른다.
\end{proof}

\begin{lemma}
\label{algebra-lemma-separable-extension-preserves-reducedness}
$k$를 체라 하자.
$S$를 축소 $k$-대수라 하자.
$K/k$를 분리 가능 체확대 또는 분리적으로 생성된 체확대라 하자.
그러면 $K \otimes_k S$는 축소이다.
\end{lemma}

\begin{proof}
$k \subset K$가 분리 가능하다고 하자.
보조정리 \ref{algebra-lemma-limit-argument}에 의해
$S$가 $k$ 위에서 유한형이고 $K$가 $k$ 위에서
유한 생성된다고 가정해도 된다.
그러면 $S$는 자신의 전분수환인 체들의 유한 곱에 매입된다.
보조정리 \ref{algebra-lemma-minimal-prime-reduced-ring}과
\ref{algebra-lemma-total-ring-fractions-no-embedded-points}를 참조하라.
따라서 실제로 $S$가 정역이라고 가정해도 된다.
보조정리
\ref{algebra-lemma-generating-finitely-generated-separable-field-extensions}와 같이
$x_1, \ldots, x_{r + 1} \in K$를 택하자.
$P \in k(x_1, \ldots, x_r)[T]$를 $x_{r + 1}$의
최소다항식이라 하자. 이는 분리 가능 다항식이다.
$k[x_1, \ldots, x_r] \otimes_k S = S[x_1, \ldots, x_r]$가
정역임은 쉽게 알 수 있다. 따라서
$k(x_1, \ldots, x_r) \otimes_k S$도
$S[x_1, \ldots, x_r]$의 국소화이므로 정역이다.
환확대
$k(x_1, \ldots, x_r) \otimes_k S \subset K \otimes_k S$는
하나의 원소 $x_{r + 1}$와 하나의 방정식 $P$로 생성된다.
그러므로 $K \otimes_k S$는 $F[T]/(P)$에 매입된다.
여기서 $F$는 $k(x_1, \ldots, x_r) \otimes_k S$의 분수체이다.
$P$가 분리 가능하므로 이는 체들의 유한 곱이고, 결론이 따른다.

\medskip\noindent
이 시점에는 분리적으로 생성된 체확대가 분리 가능한지 아직 모르므로
이 경우에도 보조정리를 증명해야 한다.
$\{x_i\}_{i \in I}$를 $k$ 위의 $K$의 분리 초월기저라고 하자.
임의의 유한한 원소 집합 $\lambda_j \in K$에 대해
다음을 유한 분리 가능 확대가 되게 하는 유한 부분집합
$T \subset I$가 존재한다.
$k(\{x_i\}_{i\in T}) \subset k(\{x_i\}_{i \in T} \cup \{\lambda_j\})$
따라서 $K$는 $k$의 유한 생성되고 분리적으로 생성된 확대들의
유향 귀납적 극한이다. 그러므로 앞 문단의 논증이 이 경우에도 적용된다.
\end{proof}

\begin{lemma}
\label{algebra-lemma-generic-points-geometrically-reduced}
$k$를 체라 하고 $S$를 $k$-대수라 하자.
$S$가 축소이고, $S$의 모든 극소 소 아이디얼 $\mathfrak p$에 대해
$S_{\mathfrak p}$가 기하적으로 축소라고 가정하자.
그러면 $S$는 기하적으로 축소이다.
\end{lemma}

\begin{proof}
$S$가 축소이므로 사상
$S \to \prod_{\mathfrak p\text{ minimal}} S_{\mathfrak p}$는 단사이다.
보조정리 \ref{algebra-lemma-reduced-ring-sub-product-fields}를 참조하라.
$K/k$가 체확대이면 다음 사상들은 단사이다.
$$
S \otimes_k K \to (\prod S_\mathfrak p) \otimes_k K \to
\prod S_\mathfrak p \otimes_k K
$$
첫 번째 사상은 $k \to K$가 평탄하므로 단사이고,
두 번째 사상은 $K$가 자유 $k$-가군이므로 직접 확인하면 단사이다.
$S_\mathfrak p$가 기하적으로 축소이므로 오른쪽 환은 축소이다.
따라서 $S \otimes_k K$는 축소환의 부분환으로서 축소이다.
\end{proof}

\begin{lemma}
\label{algebra-lemma-separable-algebraic-diagonal}
$k'/k$를 분리 가능 대수적 확대라 하자.
곱셈 사상 $k' \otimes_k k' \to k'$가
$k' \otimes_k k' \to S^{-1}(k' \otimes_k k')$와 동일시되게 하는
곱셈적 부분집합 $S \subset k' \otimes_k k'$가 존재한다.
\end{lemma}

\begin{proof}
먼저 $k'/k$가 유한 분리 가능하다고 하자.
그러면 체론, 보조정리 \ref{fields-lemma-primitive-element}에 의해
$k' = k(\alpha)$이다.
$P \in k[x]$를 $k$ 위의 $\alpha$의 최소다항식이라 하자.
그러면 체론, 절 \ref{fields-section-separable-extensions}에 의해
$P$는 기약이고 분리 가능하며 일계수인 다항식이다.
% P01-E1346
그러면 $k'[x]/(P) \to k' \otimes_k k'$,
$\sum \alpha_i x^i \mapsto \alpha_i \otimes \alpha^i$는 동형이다.
$k'[x]$에서 $P = (x - \alpha) Q$로 인수분해할 수 있고,
$P$가 분리 가능하므로 $Q(\alpha) \not = 0$임을 안다.
그러면 $k'[x]/(P)$에서 $Q$가 생성하는 곱셈적 집합 $S'$가
원하는 성질을 만족한다. 즉,
$k' = (S')^{-1}(k'[x]/(P))$이다.
구조의 운반에 의해 $k' \otimes_k k'$ 안에서 $S'$의 상 $S$도
원하는 성질을 만족한다.

\medskip\noindent
일반적인 경우에는 주어진 체를 $k$ 위의 유한 부분체 확대들의 합집합
$k' = \bigcup k_i$로 쓴다. 각 $i$에 대해
$k_i = S_i^{-1}(k_i \otimes_k k_i)$가 되게 하는 곱셈적 부분집합
$S_i \subset k_i \otimes_k k_i$가 존재한다.
$S$를 $\bigcup S_i \subset k' \otimes_k k'$의 곱셈적 폐포라 하자.
% P01-E1347
곱셈 사상은 분명 $S$의 모든 원소를 $k'$의 가역원으로 보낸다.
$k' \otimes_k k'$가 환들 $k_i \otimes_k k_i$의 합집합이라는 사실을
사용하면, 곱셈 사상의 핵에 속하는 모든 원소가
$S^{-1}(k' \otimes_k k')$에서 영으로 사상됨을 알 수 있다.
국소화의 완전성을 사용하면 결과가 따른다.
\end{proof}

\begin{lemma}
\label{algebra-lemma-geometrically-reduced-over-separable-algebraic}
$k'/k$를 분리 가능 대수적 체확대라 하자.
$A$를 $k'$ 위의 대수라 하자.
그러면 $A$가 $k$ 위에서 기하적으로 축소일 필요충분조건은
$k'$ 위에서 기하적으로 축소인 것이다.
\end{lemma}

\begin{proof}
$A$가 $k'$ 위에서 기하적으로 축소라고 하자.
$K/k$를 체확대라 하자.
그러면 보조정리 \ref{algebra-lemma-separable-extension-preserves-reducedness}에 의해
$K \otimes_k k'$는 축소환이다.
따라서 보조정리
\ref{algebra-lemma-geometrically-reduced-any-reduced-base-change}에 의해
$K \otimes_k A = (K \otimes_k k') \otimes_{k'} A$가 축소임을 얻는다.

\medskip\noindent
$A$가 $k$ 위에서 기하적으로 축소라고 하자.
$K/k'$를 체확대라 하자. 그러면
$$
K \otimes_{k'} A = (K \otimes_k A) \otimes_{(k' \otimes_k k')} k'
$$
보조정리 \ref{algebra-lemma-separable-algebraic-diagonal}에 의해
$k' \otimes_k k' \to k'$는 국소화이므로,
$K \otimes_{k'} A$는 축소 대수의 국소화이고 따라서 축소이다.
\end{proof}




\section{분리 확대, 계속}
\label{algebra-section-separability-continued}

\noindent
이 절에서는 절 \ref{algebra-section-separability}에서 시작한 논의를 계속한다.

\begin{lemma}
\label{algebra-lemma-mini-separability}
$k$를 표수가 $p > 1$인 체라 하자. $K/k$를
$x_1, \ldots, x_{n + 1} \in K$로 생성되는 체확대라 하고
다음을 가정하자.
\begin{enumerate}
\item $\{x_1, \ldots, x_n\}$은 $K/k$의 초월기저이다.
\item $K$의 모든 $k$-선형독립 부분집합 $\{a_1, \ldots, a_m\}$에 대해
집합 $\{a^p_1, \ldots, a_m^p\}$은 $k$-선형독립이다.
\end{enumerate}
그러면 $\{ x_1, \ldots, \widehat{x}_j, \ldots, x_{n+1}\}$가
$K / k$의 분리 초월기저가 되게 하는 $1 \leq j \leq n+1$이 존재한다.
\end{lemma}

\begin{proof}
가정에 의해 $x_{n + 1}$은 $k(x_1, \ldots, x_n)$ 위에서 대수적이므로,
$F(x_1, \ldots, x_{n+1}) = 0$을 만족하는 영이 아닌 다항식
$F \in k[X_1, \ldots, X_{n + 1}]$가 존재한다.
전차수가 최소인 $F$를 택하자. 그러면 $F$는 기약이다.
그렇지 않으면 적어도 하나의 기약인 인수도 같은 성질을 가져야 하기 때문이다.

\medskip\noindent
% P01-E1348
어떤 $i$에 대해서는 $F$에 나타나는 $X_i$의 지수들이
모두 $p$의 배수인 것은 아니라고 주장한다.
모순을 얻기 위해 $F$에 나타나는 모든 지수가 $p$의 배수라고 하자.
그러면 $F$의 계수들을 $\lambda_\alpha$라 할 때 집합
$$
\{x_1^{\alpha_1} \ldots x^{\alpha_{n+1}}_{n+1} \mid \lambda_\alpha \neq 0\}
\subset
K
$$
은 $k$-선형종속이다. 가정 (2)에 의해 집합
$$
\{x_1^{\alpha_1 / p} \ldots x^{\alpha_{n+1} / p}_{n+1}
\mid \lambda_\alpha \neq 0 \}
$$
도 $k$-선형종속이고, 이는 $\deg(F)$의 최소성에 모순이다.

\medskip\noindent
$F$에 $X_i$의 $p$제곱이 아닌 거듭제곱이 나타나게 하는
$i$를 택하자. 그러면 $x_i$는
$$
L = k(x_1, \ldots, x_{i - 1}, x_{i + 1}, \ldots, x_{n+1})
$$
위에서 대수적이다. 체론, 보조정리
\ref{fields-lemma-transcendence-degree}에 의해
$$
x_1, \ldots, x_{i - 1}, x_{i + 1}, \ldots, x_{n+1}
$$
은
$K/k$의 초월기저이다. 따라서 $L$은
$$
x_1, \ldots, x_{i - 1}, x_{i + 1}, \ldots, x_{n + 1}
$$
을
부정원으로 하는 $k$ 위의 다항식환의 분수체이다.
가우스 보조정리에 의해
$$
P(T) =
F(x_1, \ldots, x_{i - 1}, T, x_{i + 1}, \ldots, x_{n + 1}) \in L[T]
$$
는 기약이다. 구성에 의해 $P(T)$는 $L[T^p]$에 속하지 않는다.
따라서 원하는 대로 $K/L$은 분리 가능하다.
\end{proof}

\noindent
$p$를 소수라 하고 $k$를 표수가 $p$인 체라 하자.
이 경우 $k$의 모든 원소의 $p$제곱근을 $k$에 첨가하여 얻는
$k$의 확대를 $k^{1/p}$로 쓴다.
다르게 말하면, 이는 $k$의 원소들의 $p$제곱근들로 생성되는
$k$의 대수적 폐포의 부분체이다.

\begin{lemma}
\label{algebra-lemma-characterize-separable-field-extensions}
$k$를 표수가 $p > 0$인 체라 하고 $K/k$를 체확대라 하자.
다음은 서로 동치이다.
\begin{enumerate}
\item $K$는 $k$ 위에서 분리 가능하다.
\item $K$의 모든 $k$-선형독립 부분집합 $\{a_1, \ldots, a_m\}$에 대해
집합 $\{a^p_1, \ldots, a_m^p\}$은 $k$-선형독립이다.
\item 환 $K \otimes_k k^{1/p}$는 축소이다.
\item $K$는 $k$ 위에서 기하적으로 축소이다.
\end{enumerate}
\end{lemma}

\begin{proof}
(1) $\Rightarrow$ (4)는 보조정리
\ref{algebra-lemma-separable-extension-preserves-reducedness}에서 따른다.
(4) $\Rightarrow$ (3)은 자명하다.

\medskip\noindent
(3)을 가정하자. 다음 환 준동형을 생각하자.
$m : K \otimes_k k^{1/p} \rightarrow K$
$$
\lambda \otimes \mu \rightarrow \lambda^p \mu^p
$$
모든 $x \in K \otimes_k k^{1/p}$에 대해
$x^p = m(x) \otimes 1$임에 유의하자.
$K \otimes_k k^{1/p}$가 축소이므로 $m$은 단사이다.
$\{a_1, \ldots, a_m\} \subset K$가 $k$-선형독립이면
$\{a_1 \otimes 1, \ldots, a_m \otimes 1\}$은
$k^{1/p}$-선형독립이다.
$m$의 단사성에 의해 $a_1^p, \ldots, a_m^p$의
자명하지 않은 어떤 $k$-선형결합도 영이 아님을 얻는다.
% P01-E1349
따라서 (3)은 (2)를 함의한다.

\medskip\noindent
(2)를 가정하자. (1)을 증명하기 위해 $K$가 $k$ 위에서
유한 생성된다고 가정해도 되고, $K$가 $k$ 위에서
분리적으로 생성됨을 증명하면 된다.
$\{x_1, \ldots, x_d\}$를 $K/k$의 초월기저라 하자.
체론, 보조정리 \ref{fields-lemma-algebraic-finitely-generated}에 의해
$K' = k(x_1, \ldots, x_d)$라 하면 $[K : K'] < \infty$이다.
비분리 차수 $[K : K']_i$가 최소가 되도록 초월기저를 택하자.
$K / K'$가 분리 가능하면 증명이 끝난다.
그렇지 않다고 가정하여 모순을 얻자.
그러면 $K'$ 위에서 분리 가능하지 않은
$x_{d + 1} \in K$가 존재하고, 특히
$[K'(x_{d+1}) : K']_i > 1$이다.
% P01-E1350
보조정리 \ref{algebra-lemma-mini-separability}에 의해
$K'' = k(x_1, \ldots, \widehat{x}_j, \ldots, x_{d+1})$가
$[K'(x_{d+1}) : K'']_i = 1$을 만족하게 하는
$1 \leq j \leq n + 1$이 존재한다.
곱셈성에 의해 $[K : K'']_i < [K : K']_i$이고,
원하는 모순을 얻는다.
\end{proof}

\begin{lemma}
\label{algebra-lemma-separably-generated-separable}
분리적으로 생성된 체확대는 분리 가능하다.
\end{lemma}

\begin{proof}
보조정리 \ref{algebra-lemma-separable-extension-preserves-reducedness}와
보조정리 \ref{algebra-lemma-characterize-separable-field-extensions}를 결합하면 된다.
\end{proof}

\noindent
다음 보조정리에서는 정의 \ref{algebra-definition-perfection}에서 정의하는
완전 폐포의 개념을 사용한다.

\begin{lemma}
\label{algebra-lemma-geometrically-reduced-finite-purely-inseparable-extension}
$k$를 체라 하고 $S$를 $k$-대수라 하자.
다음은 서로 동치이다.
\begin{enumerate}
\item $k$의 모든 유한 순수 비분리 확대 $k'$에 대해
$k' \otimes_k S$가 축소이다.
\item $k^{1/p} \otimes_k S$가 축소이다.
\item $k^{perf} \otimes_k S$가 축소이다.
여기서 $k^{perf}$는 $k$의 완전 폐포이다.
\item $\overline{k} \otimes_k S$가 축소이다.
여기서 $\overline{k}$는 $k$의 대수적 폐포이다.
\item $S$는 $k$ 위에서 기하적으로 축소이다.
\end{enumerate}
\end{lemma}

\begin{proof}
모든 유한 순수 비분리 확대 $k'/k$는 $k^{perf}$에 매입된다.
또한 $k^{1/p}$는 $k^{perf}$에 매입되고,
후자는 $\overline{k}$에 매입된다.
따라서 (5) $\Rightarrow$ (4) $\Rightarrow$ (3) $\Rightarrow$ (2)와
(3) $\Rightarrow$ (1)은 명백하다.

\medskip\noindent
(1) $\Rightarrow$ (5)를 증명하자.
$k$의 모든 유한 순수 비분리 확대 $k'$에 대해
$k' \otimes_k S$가 축소라고 가정하자.
$K/k$를 체확대라 하자. $K \otimes_k S$가 축소임을 보여야 한다.
보조정리 \ref{algebra-lemma-limit-argument}에 의해 $K/k$가 유한 생성
체확대인 경우로 환원된다. 보조정리
\ref{algebra-lemma-make-separably-generated}와 같은 도식
$$
\xymatrix{
K \ar[r] & K' \\
k \ar[u] \ar[r] & k' \ar[u]
}
$$
을 택하자. 가정에 의해 $k' \otimes_k S$는 축소이다.
보조정리 \ref{algebra-lemma-separable-extension-preserves-reducedness}에 의해
$K' \otimes_k S$가 축소임이 따른다.
따라서 원하는 대로 $K \otimes_k S$가 축소임을 얻는다.

\medskip\noindent
마지막으로 (2) $\Rightarrow$ (5)를 증명하자.
$k^{1/p} \otimes_k S$가 축소라고 가정하자. 그러면 $S$는 축소이다.
더 나아가 각 극소 소 아이디얼 $\mathfrak p$에서의 국소화
$S_{\mathfrak p}$에 대해 환
$k^{1/p}\otimes_k S_{\mathfrak p}$는
$k^{1/p} \otimes_k S$의 국소화이므로 축소이다.
그러나 보조정리 \ref{algebra-lemma-minimal-prime-reduced-ring}에 의해
$S_{\mathfrak p}$는 체이고,
보조정리 \ref{algebra-lemma-characterize-separable-field-extensions}에 의해
$S_{\mathfrak p}$는 기하적으로 축소이다.
보조정리 \ref{algebra-lemma-generic-points-geometrically-reduced}에 의해
$S$가 기하적으로 축소임이 따른다.
\end{proof}



\section{완전체}
\label{algebra-section-perfect-fields}

\noindent
정의는 다음과 같다.

\begin{definition}
\label{algebra-definition-perfect}
$k$를 체라 하자. $k$의 모든 체확대가 $k$ 위에서 분리 가능하면
$k$를 {\it 완전체}라고 한다.
\end{definition}

\begin{lemma}
\label{algebra-lemma-perfect}
체 $k$가 완전체일 필요충분조건은 그것이 표수가 $0$인 체이거나,
모든 원소가 $p$제곱근을 갖는 표수 $p > 0$인 체인 것이다.
\end{lemma}

\begin{proof}
표수가 영인 경우는 명백하다.
$k$의 표수가 $p > 0$이라고 하자.
$k$가 완전체이면 $k$의 한 원소의 $p$제곱근을 첨가한 모든 체확대가
자명해야 하므로, $k$의 모든 원소는 $p$제곱근을 갖는다.
반대로 모든 원소가 $p$제곱근을 가지면 $k = k^{1/p}$이고,
보조정리 \ref{algebra-lemma-characterize-separable-field-extensions}에 의해
$k$의 모든 체확대는 분리 가능하다.
\end{proof}

\begin{lemma}
\label{algebra-lemma-make-separable}
$K/k$를 유한 생성 체확대라 하자. 다음 도식이 존재한다.
$$
\xymatrix{
K \ar[r] & K' \\
k \ar[u] \ar[r] & k' \ar[u]
}
$$
여기서 $k'/k$, $K'/K$는 유한 순수 비분리 체확대이고,
$K'/k'$는 분리 가능 체확대이다.
이 상황에서는 $K' = k'K$가 합성체라고 가정할 수 있으며,
또한 $K' = (k' \otimes_k K)_{red}$라고 가정할 수 있다.
\end{lemma}

\begin{proof}
보조정리 \ref{algebra-lemma-make-separably-generated}에 의해
$K'/k'$가 분리적으로 생성되는 이러한 도식을 찾을 수 있다.
보조정리 \ref{algebra-lemma-separably-generated-separable}에 의해
이는 $K'$가 $k'$ 위에서 분리 가능함을 뜻한다.
합성체 $k'K$는 $K'/k'$의 부분확대이므로,
보조정리 \ref{algebra-lemma-subextensions-are-separable}에 의해
$k' \subset k'K$는 분리 가능하다.
어떤 $n \gg 0$에 대해 사상 $x \mapsto x^{p^n}$가 환
$(k' \otimes_k K)_{red}$를 $K$ 안으로 보내므로,
이 환은 정역이다. 따라서 보조정리
\ref{algebra-lemma-integral-over-field}에 의해 체이다.
그러므로 $(k' \otimes_k K)_{red} \to K'$는 이를
$k'K$ 위로 동형으로 보낸다.
\end{proof}

\begin{lemma}
\label{algebra-lemma-perfection}
\begin{slogan}
모든 체는 유일한 완전 폐포를 갖는다.
\end{slogan}
모든 체 $k$에 대해 $k'$가 완전체인 순수 비분리 확대
$k'/k$가 존재한다. 체확대 $k'/k$는 유일한 동형을 제외하고 유일하다.
\end{lemma}

\begin{proof}
$k$의 표수가 영이면 $k' = k$가 유일한 선택이다.
$k$의 표수가 $p > 0$이라고 하자.
모든 $n > 0$에 대해 다음을 만족하는 유일한 대수적 확대
$k \subset k^{1/p^n}$이 존재한다.
(a) 모든 원소 $\lambda \in k$는 $k^{1/p^n}$ 안에서
$p^n$제곱근을 갖고, (b) 모든 원소 $\mu \in k^{1/p^n}$에 대해
$\mu^{p^n} \in k$이다.
실제로 환 사상
$k \to k^{1/p^n} = k$, $x \mapsto x^{p^n}$를 생각하면 된다.
이는 단사이고 (a)와 (b)를 만족한다.
$y \mapsto y^p$를 통한 $k$의 확대들로서
$k^{1/p^n} \subset k^{1/p^{n + 1}}$임은 명백하다.
그러면 $k' = \bigcup k^{1/p^n}$로 둘 수 있다.
몇 가지 세부사항은 생략한다.
\end{proof}

\begin{definition}
\label{algebra-definition-perfection}
$k$를 체라 하자. 보조정리 \ref{algebra-lemma-perfection}의 체확대
$k'/k$를 $k$의 {\it 완전 폐포}라고 한다. 표기는 $k^{perf}/k$이다.
\end{definition}

\noindent
$k'/k$가 임의의 대수적 순수 비분리 확대이면,
$k'$는 $k^{perf}$의 부분확대이다. 즉, $k^{perf}/k'/k$이다.
실제로 보조정리 \ref{algebra-lemma-perfection}의 유일성에 의해
$(k')^{perf}$는 $k^{perf}$와 동형이다.

\begin{lemma}
\label{algebra-lemma-perfect-reduced}
$k$를 완전체라 하자.
% P01-E1351
임의의 축소 $k$ 대수는 $k$ 위에서 기하적으로 축소이다.
$R$, $S$를 $k$-대수들이라 하자.
$R$과 $S$가 모두 축소라고 가정하자.
그러면 $k$-대수 $R \otimes_k S$는 축소이다.
\end{lemma}

\begin{proof}
첫 번째 주장은 보조정리
\ref{algebra-lemma-geometrically-reduced-finite-purely-inseparable-extension}에서 따른다.
두 번째 주장에는 첫 번째 주장과 보조정리
\ref{algebra-lemma-geometrically-reduced-any-reduced-base-change}를 사용하라.
\end{proof}









\section{보편 위상동형사상}
\label{algebra-section-universal-homeomorphism}

\noindent
$k'/k$를 대수적 순수 비분리 체확대라 하자.
그러면 임의의 $k$-대수 $R$에 대해 환 사상
$R \to k' \otimes_k R$는 스펙트럼들의 위상동형사상을 유도한다.
그 이유는 아래의 조금 더 일반적인 보조정리
\ref{algebra-lemma-p-ring-map}이다.

\begin{lemma}
\label{algebra-lemma-surjective-locally-nilpotent-kernel}
$\varphi : R \to S$를 핵이 국소적으로 멱영인 전사 사상이라 하자.
그러면 $\varphi$는 스펙트럼들의 위상동형사상과
잉여체들의 동형들을 유도한다.
임의의 환 사상 $R \to R'$에 대해 환 사상
$R' \to R' \otimes_R S$도 핵이 국소적으로 멱영인 전사 사상이다.
\end{lemma}

\begin{proof}
보조정리 \ref{algebra-lemma-spec-closed}에 의해
$\Spec(S) \to \Spec(R)$는 닫힌부분집합
$V(\Ker(\varphi))$ 위로의 위상동형사상이다.
$R$의 모든 소 아이디얼은 $R$의 모든 멱영원을 포함하므로
물론 $V(\Ker(\varphi)) = \Spec(R)$이다.
이는 잉여체에 관한 주장도 함의한다.
텐서곱의 오른쪽 완전성에 의해
$\Ker(\varphi)R'$는 전사 사상 $R' \to R' \otimes_R S$의 핵이다.
따라서 마지막 주장은 보조정리 \ref{algebra-lemma-locally-nilpotent}에서 따른다.
\end{proof}

\begin{lemma}
\label{algebra-lemma-powers-field}
\begin{reference}
\cite[Lemma 3.1.6]{Alper-adequate}
\end{reference}
$k'/k$를 체확대라 하자. 다음은 서로 동치이다.
\begin{enumerate}
\item 각 $x \in k'$에 대해 $x^n \in k$가 되게 하는
$n > 0$이 존재한다.
\item $k' = k$이거나, $k$와 $k'$의 표수가 $p > 0$이고
$k'/k$가 순수 비분리 확대이거나 $k$와 $k'$가
$\mathbf{F}_p$의 대수적 확대이다.
\end{enumerate}
\end{lemma}

\begin{proof}
(2)에 열거된 각 가능성이 (1)을 만족함을 관찰하자.
따라서 $k'/k$가 (1)을 만족한다고 가정하고
(2)의 경우들 가운데 하나에 속함을 증명하자.
$k = k'$인 경우를 버리면 $k' \not = k$라고 가정해도 된다.
$k'/k$가 대수적인 것은 명백하다.
따라서 $k'/k$가 자명하지 않은 유한 확대라고 가정해도 된다.
체론, 보조정리 \ref{fields-lemma-separable-first}에서 얻는
분리 가능 부분확대를 $k'/k'_{sep}/k$라 하자.
% P01-E1352
$k = k'_{sep}$이거나 $k$가 $\mathbf{F}_p$ 위에서
대수적임을 보여야 한다. 따라서 $k'/k$가 자명하지 않은
유한 분리 가능 확대라고 가정하고,
$k$가 $\mathbf{F}_p$ 위에서 대수적임을 보이면 된다.

\medskip\noindent
$x \in k'$, $x \not \in k$를 택하자.
$x^n \in k$이고 $(x + 1)^m \in k$가 되게 하는
$n, m > 0$을 택하자. $\overline{k}$를 $k$의 대수적 폐포라 하자.
$\sigma(x) \not = \tau(x)$를 만족하는 매입
$\sigma, \tau : k' \to \overline{k}$를 택할 수 있다.
이는 체론, 절 \ref{fields-section-separable-extensions}의 논의에서 따른다.
더 정확히는 $k'$를 $x$로 생성되는 $k$-확대로 바꾼 뒤
체론, 보조정리 \ref{fields-lemma-count-embeddings}에서 따른다.
그러면 $\overline{k}$ 안의 어떤 $n$차 단위원근 $\zeta$에 대해
$\sigma(x) = \zeta \tau(x)$임을 안다.
마찬가지로 어떤 $m$차 단위원근 $\zeta' \in \overline{k}$에 대해
$\sigma(x + 1) = \zeta' \tau(x + 1)$임을 안다.
$\sigma(x + 1) \not = \tau(x + 1)$이므로
$\zeta' \not = 1$이다. 이제
$$
\zeta' (\tau(x) + 1) =
\zeta' \tau(x + 1) =
\sigma(x + 1) =
\sigma(x) + 1 =
\zeta \tau(x) + 1
$$
에서
$$
\tau(x) (\zeta' - \zeta) = 1 - \zeta'
$$
를 얻는다. 따라서 $\zeta' \not = \zeta$이고
$$
\tau(x) = (1 - \zeta')/(\zeta' - \zeta)
$$
이다. 그러므로 $k'$의 원소 가운데 $k$에 속하지 않는 것은 모두
소체 위에서 대수적이다. $k'$는 $k'$의 원소 가운데 $k$에
속하지 않는 것들로 소체 위에서 생성되므로,
$k'$와 이에 따라 $k$도 소체 위에서 대수적이다.

\medskip\noindent
마지막으로 $k$의 표수가 $0$이면 위의 논의는 다음과 같이
모순을 낳는다. 독자에게 직접 증명을 찾아볼 것을 권한다.
모든 유리수 $y$에 대해 마찬가지로
$\sigma(x + y) = \zeta_y\tau(x + y)$를 만족하는 단위원근
$\zeta_y$를 얻는다. 그러면
$$
\zeta \tau(x) + y = \zeta_y(\tau(x) + y)
$$
이고, 위의 $\tau(x)$ 공식에 의해
$\zeta_y \in \mathbf{Q}(\zeta, \zeta')$임을 얻는다.
수체 $\mathbf{Q}(\zeta, \zeta')$에는 단위원근이 유한 개뿐이므로,
$\zeta_y = \zeta_{y'}$인 서로 다른 두 유리수 $y, y'$를 찾는다.
그러면
$$
y - y' =
\sigma(x + y) - \sigma(x + y') =
\zeta_y(\tau(x + y)) - \zeta_{y'}\tau(x + y') = \zeta_y(y - y')
$$
이고, 이는 $\zeta_y = 1$을 함의하여 모순이다.
\end{proof}

\begin{lemma}
\label{algebra-lemma-powers}
$\varphi : R \to S$를 환 사상이라 하자. 다음을 가정하자.
\begin{enumerate}
\item 모든 $x \in S$에 대해 $x^n$이 $\varphi$의 상에 속하게 하는
$n > 0$이 존재한다.
\item $\Ker(\varphi)$는 국소적으로 멱영이다.
\end{enumerate}
그러면 $\varphi$는 스펙트럼들의 위상동형사상을 유도하고,
보조정리 \ref{algebra-lemma-powers-field}의 동치 조건들을 만족하는
잉여체 확대들을 유도한다.
\end{lemma}

\begin{proof}
(1)과 (2)를 가정하자.
$S$의 두 소 아이디얼 $\mathfrak q, \mathfrak q'$가
$R$의 같은 소 아이디얼 $\mathfrak p$ 위에 놓인다고 하자.
$x \in S$가 $x \in \mathfrak q$, $x \not \in \mathfrak q'$를
만족한다고 하자. 그러면 모든 $n > 0$에 대해
$x^n \in \mathfrak q$이고 $x^n \not \in \mathfrak q'$이다.
어떤 $n > 0$과 $y \in R$에 대해 $x^n = \varphi(y)$이면
$$
x^n \in \mathfrak q \Rightarrow y \in \mathfrak p \Rightarrow
x^n \in \mathfrak q'
$$
이고, 이는 모순이다. 따라서 위와 같은 $x$는 존재하지 않으며
$\mathfrak q = \mathfrak q'$임을 얻는다. 즉, 스펙트럼 위의 사상은
단사이다. 가정 (2)에 의해 핵 $I = \Ker(\varphi)$는
모든 소 아이디얼에 포함되므로, 위상공간으로서
$\Spec(R) = \Spec(R/I)$이다. 유도된 사상 $R/I \to S$는
% P01-E1353
가정 (1)에 의해 정수적이므로, 보조정리
\ref{algebra-lemma-integral-overring-surjective}에 의해
$\Spec(S) \to \Spec(R/I)$는 전사이다.
위의 결과들을 결합하면 $\Spec(S) \to \Spec(R)$가 전단사임을 안다.
임의의 $x \in S$를 택하고, 어떤 $n > 0$에 대해
$\varphi(y) = x^n$이 되게 하는 $y \in R$를 택하면,
위의 전단사를 통해 열린집합 $D(x) \subset \Spec(S)$는
열린집합 $D(y) \subset \Spec(R)$에 대응한다.
따라서 사상 $\Spec(S) \to \Spec(R)$는 위상동형사상이다.

\medskip\noindent
잉여체에 관한 주장을 보기 위해, 소 아이디얼
$\mathfrak q \subset S$가 소 아이디얼
$\mathfrak p \subset R$ 위에 놓인다고 하자.
$x \in \kappa(\mathfrak q)$라 하자.
$\kappa(\mathfrak q)$를 국소환 $S_\mathfrak q$의 잉여체로 생각하면,
$x$는 $y \in S$, $z \in S$, $z \not \in \mathfrak q$인 어떤
$y/z \in S_\mathfrak q$의 상이다.
$y^n, z^m$이 $\varphi$의 상에 속하게 하는 $n, m > 0$을 택하자.
그러면 $x^{nm}$은 $(y/z)^{nm} = (y^n)^m/(z^m)^n$의 잉여류이고,
이는 $R_\mathfrak p \to S_\mathfrak q$의 상에 속한다.
따라서 $x^{nm}$은
$\kappa(\mathfrak p) \to \kappa(\mathfrak q)$의 상에 속한다.
\end{proof}

\begin{lemma}
\label{algebra-lemma-2-3-ring-map}
$\varphi : R \to S$를 환 사상이라 하고 다음을 가정하자.
\begin{enumerate}
\item[(a)] $S$는 $x^2, x^3 \in \varphi(R)$을 만족하는 원소
$x$들로 $R$-대수로서 생성된다.
\item[(b)] $\Ker(\varphi)$는 국소적으로 멱영이다.
\end{enumerate}
그러면 $\varphi$는 잉여체들의 동형과
스펙트럼들의 위상동형사상을 유도한다.
임의의 환 사상 $R \to R'$에 대해 환 사상
$R' \to R' \otimes_R S$도 (a)와 (b)를 만족한다.
\end{lemma}

\begin{proof}
(a)와 (b)를 가정하자.
$S$는 $R$ 위에서 정수적이므로 스펙트럼 위의 사상은 닫힌 사상이다.
보조정리 \ref{algebra-lemma-going-up-closed}와
\ref{algebra-lemma-integral-going-up}을 참조하라.
보조정리 \ref{algebra-lemma-image-dense-generic-points}에 의해 그 상은 조밀하다.
따라서 $\Spec(S) \to \Spec(R)$는 전사이다.
$\mathfrak q \subset S$가 $\mathfrak p \subset R$ 위에 놓인
소 아이디얼이면, 체확대 $\kappa(\mathfrak q)/\kappa(\mathfrak p)$는
제곱과 세제곱이 $\kappa(\mathfrak p)$에 속하는 원소
$\alpha \in \kappa(\mathfrak q)$들로 생성된다.
그러므로 분명 $\alpha \in \kappa(\mathfrak p)$이고
$\kappa(\mathfrak q) = \kappa(\mathfrak p)$임을 얻는다.
$\mathfrak q, \mathfrak q'$가 $\mathfrak p$ 위에 놓인
서로 다른 두 소 아이디얼이면, (a)에서와 같은 $S$의 생성원
$x$들 가운데 적어도 하나는
$\kappa(\mathfrak q) = \kappa(\mathfrak p)$와
$\kappa(\mathfrak q') = \kappa(\mathfrak p)$ 안에서 서로 다른 상을 갖는다.
이는 $x^2$과 $x^3$이 둘 다 같은 상을 갖는다는 사실에 모순이다.
따라서 $\Spec(S) \to \Spec(R)$는 단사이고,
이미 보인 바에 의해 위상동형사상이다.

\medskip\noindent
$\varphi$는 스펙트럼들의 위상동형사상을 유도하므로 특히
스펙트럼 위에서 전사이고, 이 성질은 임의의 기저변환으로 보존된다.
보조정리 \ref{algebra-lemma-surjective-spec-radical-ideal}을 참조하라.
따라서 임의의 $R \to R'$에 대해 환 사상
$R' \to R' \otimes_R S$의 핵은 멱영원들로 이루어진다.
보조정리 \ref{algebra-lemma-image-dense-generic-points}를 참조하라.
다르게 말하면 $R' \to R' \otimes_R S$에 대해 (b)가 성립한다.
(a)가 기저변환으로 보존되는 것은 명백하다.
\end{proof}

\begin{lemma}
\label{algebra-lemma-help-with-powers}
$p$를 소수라 하고 $n, m > 0$을 두 정수라 하자.
다음을 만족하는 정수 $a$가 존재한다.
$(x + y)^{p^a}, p^a(x + y) \in \mathbf{Z}[x^{p^n}, p^nx, y^{p^m}, p^my]$.
\end{lemma}

\begin{proof}
$a \geq n, m$이면 $p^a(x + y)$에 관한 주장은 명백하다.
실제로 $a \gg n, m$을 택하자. 다음과 같이 쓰자.
$$
(x + y)^{p^a}  = \sum\nolimits_{i, j \geq 0, i + j = p^a}
{p^a \choose i, j} x^iy^j
$$
$i + j = p^a$인 모든 $i, j \geq 0$에 대해
$i = q p^n + r$, $r \in \{0, \ldots, p^n - 1\}$ 및
$j = q' p^m + r'$, $r' \in \{0, \ldots, p^m - 1\}$로 쓰자.
다음 조건이 성립하면
$(x + y)^{p^a} \in \mathbf{Z}[x^{p^n}, p^nx, y^{p^m}, p^my]$이다.
$$
p^{nr + mr'} \text{ divides } {p^a \choose i, j}
$$
$r = r' = 0$이면 이 나눗셈 성질이 성립한다.
$r \not = 0$이면 다음과 같이 쓴다.
$$
{p^a \choose i, j} = \frac{p^a}{i} {p^a - 1 \choose i - 1, j}
$$
$r \not = 0$이므로 유리수 $p^a/i$의 $p$-진 값매김은
적어도 $a - (n - 1)$이다. 이는 $i$가 $p^n$으로
나누어떨어지지 않기 때문이다.
따라서 이 경우 $ {p^a \choose i, j}$는 $p^{a - n + 1}$로 나누어떨어진다.
마찬가지로 $r' \not = 0$이면 $ {p^a \choose i, j}$가
$p^{a - m + 1}$로 나누어떨어짐을 안다.
$a = np^n + mp^m + n + m$으로 택하면 충분하다.
\end{proof}

\begin{lemma}
\label{algebra-lemma-p-ring-map-field}
$k'/k$를 체확대라 하고 $p$를 소수라 하자.
다음은 서로 동치이다.
\begin{enumerate}
\item $k'$는 $x^{p^n} \in k$와 $p^nx \in k$를 만족하는
어떤 $n > 0$이 존재하는 원소 $x$들로 $k$ 위의 체확대로서 생성된다.
\item $k = k'$이거나,
% P01-E1354
$k$와 $k'$의 표수가 $p$이고 $k'/k$가 순수 비분리이다.
\end{enumerate}
\end{lemma}

\begin{proof}
$x \in k'$라 하자. $x^{p^n} \in k$이고 $p^nx \in k$인
$n > 0$이 존재하며 표수가 $p$가 아니면 $x \in k$이다.
표수가 $p$이면 $x^{p^n} \in k$이고,
따라서 $x$는 $k$ 위에서 순수 비분리임을 얻는다.
\end{proof}

\begin{lemma}
\label{algebra-lemma-p-ring-map}
$\varphi : R \to S$를 환 사상이라 하고 $p$를 소수라 하자.
다음을 가정하자.
\begin{enumerate}
\item[(a)] $S$는 $x^{p^n} \in \varphi(R)$와
$p^nx \in \varphi(R)$를 만족하는 어떤 $n > 0$이 존재하는 원소
$x$들로 $R$-대수로서 생성된다.
\item[(b)] $\Ker(\varphi)$는 국소적으로 멱영이다.
\end{enumerate}
그러면 $\varphi$는 스펙트럼들의 위상동형사상을 유도하고,
보조정리 \ref{algebra-lemma-p-ring-map-field}의 동치 조건들을 만족하는
잉여체 확대들을 유도한다.
임의의 환 사상 $R \to R'$에 대해 환 사상
$R' \to R' \otimes_R S$도 (a)와 (b)를 만족한다.
\end{lemma}

\begin{proof}
(a)와 (b)를 가정하자. (b)는 보조정리
\ref{algebra-lemma-powers}의 조건 (2)와 동치임에 유의하자.
$T \subset S$를 $x^{p^n} , p^n x \in \varphi(R)$가 되게 하는
정수 $n > 0$이 존재하는 원소 $x \in S$들의 집합이라 하자.
$T = S$라고 주장한다. 이는 보조정리 \ref{algebra-lemma-powers}의
조건 (1)이 성립함을 증명하므로, $\varphi$가 스펙트럼들의
위상동형사상을 유도함이 따른다.
% P01-E1355
가정 (a)에 의해 $T \subset S$가 $R$-부분대수임을 보이면 충분하다.
$x \in T$이고 $y \in R$이면 $yx \in T$임은 명백하다.
$x, y \in T$이고
$x^{p^n}, y^{p^m}, p^n x, p^m y \in \varphi(R)$를 만족하는
$n, m > 0$이 주어졌다고 하자.
그러면 $(xy)^{p^{n + m}}, p^{n + m}xy \in \varphi(R)$이므로
$xy \in T$이다. 보조정리 \ref{algebra-lemma-help-with-powers}에 의해
$x + y \in T$이고, 주장이 증명되었다.

\medskip\noindent
$\varphi$는 스펙트럼들의 위상동형사상을 유도하므로 특히
스펙트럼 위에서 전사이고, 이 성질은 임의의 기저변환으로 보존된다.
보조정리 \ref{algebra-lemma-surjective-spec-radical-ideal}을 참조하라.
따라서 임의의 $R \to R'$에 대해 환 사상
$R' \to R' \otimes_R S$의 핵은 멱영원들로 이루어진다.
보조정리 \ref{algebra-lemma-image-dense-generic-points}를 참조하라.
다르게 말하면 $R' \to R' \otimes_R S$에 대해 (b)가 성립한다.
(a)가 기저변환으로 보존되는 것은 명백하다.
마지막으로 잉여체에 관한 조건은 (a)에서 따른다.
$R$-대수로서 $S$의 생성원들은 잉여체 확대의 생성원들로 사상되기 때문이다.
\end{proof}

\begin{lemma}
\label{algebra-lemma-radicial}
$\varphi : R \to S$를 환 사상이라 하고 다음을 가정하자.
\begin{enumerate}
\item $\varphi$는 스펙트럼의 단사 사상을 유도한다.
\item $\varphi$는 순수 비분리 잉여체 확대들을 유도한다.
\end{enumerate}
그러면 임의의 환 사상 $R \to R'$에 대해
$R' \to R' \otimes_R S$도 성질 (1)과 (2)를 만족한다.
\end{lemma}

\begin{proof}
$S' = R' \otimes_R S$라 놓으면 다음과 같은 환의 스펙트럼들의
연속사상 가환 도식을 얻는다.
$$
\xymatrix{
\Spec(S') \ar[r] \ar[d] & \Spec(S) \ar[d] \\
\Spec(R') \ar[r] & \Spec(R)
}
$$
$\mathfrak p' \subset R'$를 $\mathfrak p \subset R$ 위에 놓인
소 아이디얼이라 하자. $\mathfrak p$ 위에 놓인 $S$의
소 아이디얼이 없으면 $\mathfrak p'$ 위에 놓인 $S'$의
소 아이디얼도 없다. 그렇지 않으면 주석
\ref{algebra-remark-fundamental-diagram}에 의해
$F = S \otimes_R \kappa(\mathfrak p)$의 유일한 소 아이디얼
$\mathfrak r$가 존재하고, 그 잉여체는
$\kappa(\mathfrak p)$ 위에서 순수 비분리이다.
다음 환 사상들을 생각하자.
$$
\kappa(\mathfrak p) \to F \to \kappa(\mathfrak r)
$$
보조정리 \ref{algebra-lemma-minimal-prime-reduced-ring}에 의해
아이디얼 $\mathfrak r \subset F$는 국소적으로 멱영이므로,
환 사상 $F \to \kappa(\mathfrak r)$에 보조정리
\ref{algebra-lemma-surjective-locally-nilpotent-kernel}을 적용할 수 있다.
환 사상 $\kappa(\mathfrak p) \to \kappa(\mathfrak r)$에는
보조정리 \ref{algebra-lemma-p-ring-map}을 적용할 수 있다.
따라서 다음 사상들에서 합성과 두 번째 화살표는
$$
\kappa(\mathfrak p') \to
\kappa(\mathfrak p') \otimes_{\kappa(\mathfrak p)} F \to
\kappa(\mathfrak p') \otimes_{\kappa(\mathfrak p)} \kappa(\mathfrak r)
$$
스펙트럼 위의 전단사와 순수 비분리 잉여체 확대들을 유도한다.
이는 첫 번째 사상에 대해서도 같은 결과를 함의한다. 이제
$$
\kappa(\mathfrak p') \otimes_{\kappa(\mathfrak p)} F =
\kappa(\mathfrak p') \otimes_{\kappa(\mathfrak p)}
\kappa(\mathfrak p) \otimes_R S =
\kappa(\mathfrak p') \otimes_R S =
\kappa(\mathfrak p') \otimes_{R'} R' \otimes_R S
$$
이므로 주석 \ref{algebra-remark-fundamental-diagram}의 논의에 의해 결론이 따른다.
\end{proof}

\begin{lemma}
\label{algebra-lemma-radicial-integral}
$\varphi : R \to S$를 환 사상이라 하고 다음을 가정하자.
\begin{enumerate}
\item $\varphi$는 정수적이다.
\item $\varphi$는 스펙트럼의 단사 사상을 유도한다.
\item $\varphi$는 순수 비분리 잉여체 확대들을 유도한다.
\end{enumerate}
그러면 $\varphi$는 $\Spec(S)$에서 $\Spec(R)$의 닫힌부분집합
위로의 위상동형사상을 유도한다. 또한 임의의 환 사상
$R \to R'$에 대해 $R' \to R' \otimes_R S$도
성질 (1), (2), (3)을 만족한다.
\end{lemma}

\begin{proof}
보조정리 \ref{algebra-lemma-going-up-closed}와
\ref{algebra-lemma-integral-going-up}에 의해 스펙트럼 위의 사상은 닫힌 사상이다.
보조정리 \ref{algebra-lemma-radicial}과
\ref{algebra-lemma-base-change-integral}에 의해 이 성질들은 기저변환으로 보존된다.
\end{proof}

\begin{lemma}
\label{algebra-lemma-radicial-integral-bijective}
$\varphi : R \to S$를 환 사상이라 하고 다음을 가정하자.
\begin{enumerate}
\item $\varphi$는 정수적이다.
% P01-E1356
\item $\varphi$는 스펙트럼의 전단사 사상을 유도한다.
\item $\varphi$는 순수 비분리 잉여체 확대들을 유도한다.
\end{enumerate}
그러면 $\varphi$는 스펙트럼들의 위상동형사상을 유도한다.
또한 임의의 환 사상 $R \to R'$에 대해
$R' \to R' \otimes_R S$도 성질 (1), (2), (3)을 만족한다.
\end{lemma}

\begin{proof}
보조정리 \ref{algebra-lemma-radicial-integral}과
\ref{algebra-lemma-surjective-spec-radical-ideal}에서 따른다.
\end{proof}

\begin{lemma}
\label{algebra-lemma-universally-bijective}
$\varphi : R \to S$를 다음을 만족하는 환 사상이라 하자.
\begin{enumerate}
\item $\varphi$의 핵은 국소적으로 멱영이다.
\item $S$는 어떤 $n > 0$과 다항식 $P(T) \in R[T]$가 존재하여
$S[T]$에서 그 상이 $(T - x)^n$이 되게 하는 원소
$x$들로 $R$-대수로서 생성된다.
\end{enumerate}
그러면 $\Spec(S) \to \Spec(R)$는 위상동형사상이고,
$R \to S$는 순수 비분리 잉여체 확대들을 유도한다.
또한 조건 (1)과 (2)는 임의의 기저변환으로 보존된다.
\end{lemma}

\begin{proof}
$R$을 $R/\Ker(\varphi)$로 바꾸어도 된다.
보조정리 \ref{algebra-lemma-surjective-locally-nilpotent-kernel}을 참조하라.
가정 (2)에 의해 $S$는 $R$ 위에서 정수적인 원소들로
$R$ 위에서 생성된다. 따라서 $R \subset S$는 정수적이다.
보조정리 \ref{algebra-lemma-integral-closure-is-ring}을 참조하라.
특히 $\Spec(S) \to \Spec(R)$는 전사이고 닫힌 사상이다.
보조정리 \ref{algebra-lemma-integral-overring-surjective},
\ref{algebra-lemma-going-up-closed}, 그리고
\ref{algebra-lemma-integral-going-up}을 참조하라.

\medskip\noindent
(2)의 생성원 가운데 하나를 $x \in S$라 하자. 즉,
% P01-E1357
$(T - x)^n \in R[T]$가 되게 하는 $n > 0$이 존재한다.
$\mathfrak p \subset R$를 소 아이디얼이라 하자.
위의 결과와 보조정리 \ref{algebra-lemma-in-image}에 의해 환
$\kappa(\mathfrak p) \otimes_R S$는 영이 아니다.
$\kappa(\mathfrak p)$의 표수가 영이면
$nx \in R$에서 $1 \otimes x$가
$\kappa(\mathfrak p) \to \kappa(\mathfrak p) \otimes_R S$의
상에 속함을 알 수 있다. 따라서
$\kappa(\mathfrak p) \to \kappa(\mathfrak p) \otimes_R S$는 동형이다.
$\kappa(\mathfrak p)$의 표수가 $p > 0$이면
$m$이 $p$와 서로소가 되도록 $n = p^k m$으로 쓰자.
$\kappa(\mathfrak p) \otimes_R S[T]$에서
$$
(T - 1 \otimes x)^n = ((T - 1 \otimes x)^{p^k})^m =
(T^{p^k} - 1 \otimes x^{p^k})^m
$$
이다. 따라서 $mx^{p^k} \in R$임을 안다. 이는
$1 \otimes x^{p^k}$가
$\kappa(\mathfrak p) \to \kappa(\mathfrak p) \otimes_R S$의
상에 속함을 뜻한다. 따라서
$\kappa(\mathfrak p) \to \kappa(\mathfrak p) \otimes_R S$에
보조정리 \ref{algebra-lemma-p-ring-map}을 적용할 수 있다.
두 경우 모두 $\kappa(\mathfrak p) \otimes_R S$는 유일한
소 아이디얼을 가지며, 그 잉여체는
$\kappa(\mathfrak p)$ 위에서 순수 비분리임을 얻는다.
주석 \ref{algebra-remark-fundamental-diagram}에 의해
$\varphi$가 스펙트럼 위에서 전단사임을 얻는다.

\medskip\noindent
기저변환에 관한 주장은 자명하다.
\end{proof}











\section{기하적으로 기약인 대수}
\label{algebra-section-algebras-over-fields}

\noindent
체 $k$ 위의 대수 $S$가 {\it 기하적으로 기약}이라는 것은
모든 체확대 $k'/k$에 대해 대수 $S \otimes_k k'$가 유일한
극소 소 아이디얼을 갖는다는 뜻이다. 이 절에서는 이 개념에
관련된 이론을 조금 전개한다.

\begin{lemma}
\label{algebra-lemma-flat-fibres-irreducible}
$R \to S$를 환 사상이라 하자. 다음을 가정하자.
\begin{enumerate}
\item[(a)] $\Spec(R)$는 기약이다.
\item[(b)] $R \to S$는 평탄하다.
\item[(c)] $R \to S$는 유한 표시이다.
\item[(d)] $R$의 조밀한 소 아이디얼들의 모음에 속하는 모든
$\mathfrak p$에 대해 섬유환 $S \otimes_R \kappa(\mathfrak p)$의
스펙트럼이 기약이다.
\end{enumerate}
그러면 $\Spec(S)$는 기약이다.
더 일반적으로, (b) $+$ (c)를 합친 조건을
‘사상 $\Spec(S) \to \Spec(R)$가 열린 사상이다’로 바꾸어도 참이다.
\end{lemma}

\begin{proof}
가정 (b)와 (c)에 의해 스펙트럼 사이의 사상은 열린 사상이다.
명제 \ref{algebra-proposition-fppf-open}을 참조하라.
따라서 위상론, 보조정리
\ref{topology-lemma-irreducible-on-top}에서 이 보조정리가 따른다.
\end{proof}

\begin{lemma}
\label{algebra-lemma-separably-closed-irreducible}
$k$를 분리폐체라 하자.
$R$, $S$를 $k$-대수들이라 하자. $R$과 $S$가 각각 유일한
극소 소 아이디얼을 가지면 $R \otimes_k S$도 그러하다.
\end{lemma}

\begin{proof}
$k \subset \overline{k}$를 완전 폐포라 하자.
정의 \ref{algebra-definition-perfection}을 참조하라.
가정에 의해 $\overline{k}$는 대수적으로 닫혀 있다.
환 사상 $R \to R \otimes_k \overline{k}$,
$S \to S \otimes_k \overline{k}$, 그리고
$R \otimes_k S \to (R \otimes_k S) \otimes_k \overline{k}
= (R \otimes_k \overline{k}) \otimes_{\overline{k}} (S \otimes_k \overline{k})$
는 보조정리 \ref{algebra-lemma-p-ring-map}의 가정들을 만족한다.
따라서 $k$가 대수적으로 닫혀 있다고 가정해도 된다.

\medskip\noindent
$R$과 $S$를 각각 그 축소환으로 대체해도 된다.
따라서 $R$과 $S$가 정역이라고 가정해도 된다.
보조정리 \ref{algebra-lemma-perfect-reduced}에 의해 $R \otimes_k S$는
축소이다. 따라서 그 스펙트럼이 가약일 필요충분조건은
영이 아닌 영인자를 포함하는 것이다. 보조정리
\ref{algebra-lemma-limit-argument}에 의해 다음 경우로 환원된다.
% P01-E1358
기저체가 대수적으로 닫혀 있고 $R$과 $S$가 $k$ 위의 유한형 정역인 경우이다.

\medskip\noindent
환 사상 $R \to R \otimes_k S$는 유한 표시이고 평탄하다.
더 나아가 $R$의 모든 극대 아이디얼 $\mathfrak m$에 대해
힐베르트 영점 정리 \ref{algebra-theorem-nullstellensatz}에 의해
$k \cong R/\mathfrak m$이므로
$(R \otimes_k S) \otimes_R R/\mathfrak m \cong S$이다.
또한 보조정리 \ref{algebra-lemma-finite-type-field-Jacobson}에 의해
$\Spec(R)$는 야콥슨이므로 극대 아이디얼들의 집합은
$R$의 스펙트럼에서 조밀하다.
따라서 보조정리 \ref{algebra-lemma-flat-fibres-irreducible}를
환 사상 $R \to R \otimes_k S$에 적용할 수 있고,
원하는 대로 $R \otimes_k S$의 스펙트럼이 기약임을 얻는다.
\end{proof}

\begin{lemma}
\label{algebra-lemma-geometrically-irreducible}
$k$를 체라 하자.
$R$을 $k$-대수라 하자.
다음은 서로 동치이다.
\begin{enumerate}
\item 모든 체확대 $k'/k$에 대해
$R \otimes_k k'$의 스펙트럼이 기약이다.
\item 모든 유한 분리 가능 체확대 $k'/k$에 대해
$R \otimes_k k'$의 스펙트럼이 기약이다.
\item $\overline{k}$가 $k$의 분리 대수적 폐포일 때
$R \otimes_k \overline{k}$의 스펙트럼이 기약이다.
\item $\overline{k}$가 $k$의 대수적 폐포일 때
$R \otimes_k \overline{k}$의 스펙트럼이 기약이다.
\end{enumerate}
\end{lemma}

\begin{proof}
(1)이 (2)를 함의하는 것은 명백하다.

\medskip\noindent
% P01-E1359
(2)를 가정하고 $\overline{k}$를 $k$의 분리 대수적 폐포라 하자.
$\mathfrak q_i \subset R \otimes_k \overline{k}$, $i = 1, 2$가
두 극소 소 아이디얼이라고 하자. 모든 유한 부분확대
$\overline{k}/k'/k$에 대해 확대 $k'/k$는 분리 가능하고
환 사상 $R \otimes_k k' \to R \otimes_k \overline{k}$는 평탄하다.
따라서 평탄 환 사상에 대해 하강이 성립하므로
(보조정리 \ref{algebra-lemma-flat-going-down})
$\mathfrak p_i = (R \otimes_k k') \cap \mathfrak q_i$는
극소 소 아이디얼이다. 가정 (2)에 의해
$\mathfrak p_1 = \mathfrak p_2$임을 얻는다.
$\overline{k} = \bigcup k'$이므로
$\mathfrak q_1 = \mathfrak q_2$라고 결론 내릴 수 있다.
따라서 $\Spec(R \otimes_k \overline{k})$는 기약이다.

\medskip\noindent
(3)을 가정하고 $\overline{k}$를 $k$의 대수적 폐포라 하자.
$\overline{k}/\overline{k}'/k$를 이에 대응하는
$k$의 분리 대수적 폐포라 하자. 그러면
$\overline{k}/\overline{k}'$는 순수 비분리 확대이거나
(양의 표수에서) 자명한 확대이다.
따라서 $R \otimes_k \overline{k}' \to R \otimes_k \overline{k}$는,
예를 들어 보조정리 \ref{algebra-lemma-p-ring-map}에 의해,
스펙트럼 위의 위상동형사상을 유도한다. 이로써 (4)를 얻는다.

\medskip\noindent
(4)를 가정하자. $k'/k$를 임의의 체확대라 하고
$\overline{k}$를 $k$의 대수적 폐포라 하자.
$k'$와 $\overline{k}$가 모두 어떤 체 $F$의 부분체와 동형이 되도록
$F$를 택할 수 있다. 그러면
$$
R \otimes_k F = (R \otimes_k \overline{k}) \otimes_{\overline{k}} F
$$
이고, 보조정리 \ref{algebra-lemma-separably-closed-irreducible}에 의해
$R \otimes_k F$가 유일한 극소 소 아이디얼을 가짐을 알 수 있다.
마지막으로 환 사상 $R \otimes_k k' \to R \otimes_k F$는
평탄하고 단사이므로, $R \otimes_k k'$의 모든 극소 소 아이디얼은
$R \otimes_k F$의 한 극소 소 아이디얼의 상이다
(보조정리 \ref{algebra-lemma-injective-minimal-primes-in-image}와
하강을 사용하라). 따라서 그러한 극소 소 아이디얼은 하나뿐이고
증명이 끝난다.
\end{proof}

\begin{definition}
\label{algebra-definition-geometrically-irreducible}
$k$를 체라 하자.
$S$를 $k$-대수라 하자.
모든 체확대 $k'/k$에 대해
$S \otimes_k k'$의 스펙트럼이
기약이면\footnote{기약 공간은 공집합이 아니다.},
$S$가 {\it $k$ 위에서 기하적으로 기약}이라고 한다.
\end{definition}

\noindent
보조정리 \ref{algebra-lemma-geometrically-irreducible}에 의해,
유한 분리 가능 체확대 $k'/k$들에 대해서만 확인하거나
$k'$가 $k$의 분리 대수적 폐포인 경우만 확인해도 충분하다.

\begin{lemma}
\label{algebra-lemma-separably-closed-irreducible-implies-geometric}
$k$를 체라 하자.
$R$을 $k$-대수라 하자.
$k$가 분리폐체이면, $R$이 $k$ 위에서 기하적으로 기약일
필요충분조건은 $R$의 스펙트럼이 기약인 것이다.
\end{lemma}

\begin{proof}
정의 \ref{algebra-definition-geometrically-irreducible} 뒤의
주석에서 바로 따른다.
\end{proof}

\begin{lemma}
\label{algebra-lemma-subalgebra-geometrically-irreducible}
$k$를 체라 하자. $S$를 $k$-대수라 하자.
\begin{enumerate}
\item $S$가 $k$ 위에서 기하적으로 기약이면 모든
$k$-부분대수도 그러하다.
\item $S$의 모든 유한 생성 $k$-부분대수가 기하적으로 기약이면
$S$도 기하적으로 기약이다.
\item 기하적으로 기약인 $k$-대수들의 유향 귀납적 극한은
기하적으로 기약이다.
\end{enumerate}
\end{lemma}

\begin{proof}
$S' \subset S$를 부분대수라 하자. 그러면 임의의 확대 $k'/k$에 대해
환 사상 $S' \otimes_k k' \to S \otimes_k k'$도 단사이다.
따라서 보조정리 \ref{algebra-lemma-injective-minimal-primes-in-image}와
(연속사상에 의한 기약 공간의 상은 기약이라는 사실에 의해)
(1)이 따른다.
% P01-E1360
두 번째와 세 번째 성질은 텐서곱이 귀납적 극한과 가환한다는
사실에서 따른다.
\end{proof}

\begin{lemma}
\label{algebra-lemma-geometrically-irreducible-any-base-change}
$k$를 체라 하자.
$S$를 기하적으로 기약인 $k$-대수라 하자.
$R$을 임의의 $k$-대수라 하자.
사상
$$
\Spec(R \otimes_k S) \longrightarrow \Spec(R)
$$
은 기약 성분들 사이의 전단사를 유도한다.
\end{lemma}

\begin{proof}
기약 성분들은 극소 소 아이디얼들에 대응한다는 것을 상기하자
(보조정리 \ref{algebra-lemma-irreducible}).
$R \to R \otimes_k S$가 평탄하므로 하강
(보조정리 \ref{algebra-lemma-flat-going-down})에 의해
$R \otimes_k S$의 모든 극소 소 아이디얼은 $R$의 극소 소 아이디얼
위에 놓인다. 역으로 $\mathfrak p \subset R$가 (극소) 소 아이디얼이면
$$
R \otimes_k S/\mathfrak p(R \otimes_k S)
=
(R/\mathfrak p) \otimes_k S
\subset
\kappa(\mathfrak p) \otimes_k S
$$
이다. 이는 $R \to R \otimes_k S$가 평탄하기 때문이다.
가정에 의해 환 $\kappa(\mathfrak p) \otimes_k S$의
스펙트럼은 기약이다. 따라서
$R \otimes_k S/\mathfrak p(R \otimes_k S)$는
유일한 극소 소 아이디얼을 갖는다
(보조정리 \ref{algebra-lemma-injective-minimal-primes-in-image}).
다르게 말하면 기약 집합 $V(\mathfrak p)$의 역상은 기약이다.
따라서 보조정리가 따른다.
\end{proof}

\noindent
기하적으로 기약인 체확대라는 개념에 관해 몇 가지를 설명하자.

\begin{lemma}
\label{algebra-lemma-field-extension-geometrically-irreducible}
$K/k$를 체확대라 하자. $k$가 $K$ 안에서 대수적으로 닫혀 있으면,
$K$는 $k$ 위에서 기하적으로 기약이다.
\end{lemma}

\begin{proof}
$k$가 $K$ 안에서 대수적으로 닫혀 있다고 가정하자.
정의 \ref{algebra-definition-geometrically-irreducible}과
보조정리 \ref{algebra-lemma-geometrically-irreducible}에 의해,
모든 유한 분리 가능 확대 $k'/k$에 대해
$K \otimes_k k'$의 스펙트럼이 기약임을 보이면 충분하다.
$k'$가 어떤 $\alpha \in k'$에 의해 $k$ 위에서 생성된다고 하자.
체론, 보조정리 \ref{fields-lemma-primitive-element}을 참조하라.
$P = T^d + a_1 T^{d - 1} + \ldots + a_d \in k[T]$를
$\alpha$의 최소다항식이라 하자. 그러면
$K \otimes_k k' \cong K[T]/(P)$이다.
$K[T]/(P)$의 스펙트럼이 가약이 될 수 있는 유일한 경우는
$P$가 $K[T]$에서 가약인 경우이다. 모순을 얻기 위해
$P = P_1 P_2$가 $K[T]$에서 자명하지 않게 인수분해된다고 가정하자.
보조정리 \ref{algebra-lemma-polynomials-divide}에 의해
$P_1$과 $P_2$의 계수들은 $k$ 위에서 대수적이다.
가정에 의해 $P_1$과 $P_2$의 계수들은 $k$에 속하는데,
이는 $P$가 $k$ 위에서 기약이라는 사실과 모순이다.
\end{proof}

\begin{lemma}
\label{algebra-lemma-geometrically-irreducible-transitive}
$K/k$를 기하적으로 기약인 체확대라 하자.
$S$를 기하적으로 기약인 $K$-대수라 하자.
그러면 $S$는 $k$ 위에서 기하적으로 기약이다.
\end{lemma}

\begin{proof}
정의 \ref{algebra-definition-geometrically-irreducible}과
보조정리 \ref{algebra-lemma-geometrically-irreducible}에 의해,
모든 유한 분리 가능 확대 $k'/k$에 대해
$S \otimes_k k'$의 스펙트럼이 기약임을 보이면 충분하다.
$K$가 $k$ 위에서 기하적으로 기약이므로
$K' = K \otimes_k k'$는 $K$의 유한 분리 가능 체확대이다.
따라서 $S$가 $K$ 위에서 기하적으로 기약이라는 가정에 의해
$S \otimes_k k' = S \otimes_K K'$의 스펙트럼은 기약이다.
\end{proof}

\begin{lemma}
\label{algebra-lemma-geometrically-irreducible-base-change-transcendental}
$K/k$를 체확대라 하자. 다음은 서로 동치이다.
\begin{enumerate}
\item $K$는 $k$ 위에서 기하적으로 기약이다.
\item 순수 초월확대들 사이에 유도되는 확대 $K(t)/k(t)$는
기하적으로 기약이다.
\end{enumerate}
\end{lemma}

\begin{proof}
% P01-E1361
(1)을 가정하자. $\Omega$를 $k(t)$의 대수적 폐포라 하자.
정의 \ref{algebra-definition-geometrically-irreducible}에 의해
다음 환의 스펙트럼이 기약임을 얻는다.
$$
K \otimes_k \Omega = K \otimes_k k(t) \otimes_{k(t)} \Omega
$$
$K(t)$는 $K \otimes_k k(T)$의 국소화이므로
$K(t) \otimes_{k(t)} \Omega$의 스펙트럼이 기약이라고 결론 내릴 수 있다.
따라서 보조정리 \ref{algebra-lemma-geometrically-irreducible}에 의해
$K(t)/k(t)$는 기하적으로 기약이다.

\medskip\noindent
(2)를 가정하자. $k'/k$를 체확대라 하자.
$K \otimes_k k'$가 유일한 극소 소 아이디얼을 가짐을 보여야 한다.
다음 환의 스펙트럼이 기약임, 즉 유일한 극소 소 아이디얼을 가짐을 안다.
$$
K(t) \otimes_{k(t)} k'(t)
$$
단사 사상 $K \otimes_k k' \to K(t) \otimes_{k(t)} k'(t)$가
존재하므로(자세한 내용은 생략한다), 보조정리
\ref{algebra-lemma-injective-minimal-primes-in-image}와
\ref{algebra-lemma-minimal-prime-image-minimal-prime}에 의해 결론이 따른다.
\end{proof}

\begin{lemma}
\label{algebra-lemma-geometrically-irreducible-add-transcendental}
$K/L/M$을 체들의 탑이라 하고 $L/M$이 기하적으로 기약이라고 하자.
$x \in K$가 $L$ 위에서 초월적이라고 하자.
그러면 $L(x)/M(x)$는 기하적으로 기약이다.
\end{lemma}

\begin{proof}
체 $L(x)$와 $M(x)$가 각각 $L$과 $M$의 순수 초월확대이므로
보조정리
\ref{algebra-lemma-geometrically-irreducible-base-change-transcendental}에서 따른다.
\end{proof}

\begin{lemma}
\label{algebra-lemma-geometrically-irreducible-separable-elements}
$K/k$를 체확대라 하자. 다음은 서로 동치이다.
\begin{enumerate}
\item $K/k$는 기하적으로 기약이다.
\item $k$ 위에서 분리 대수적인 모든 원소 $\alpha \in K$는
$k$에 속한다.
\end{enumerate}
\end{lemma}

\begin{proof}
(1)을 가정하고 $\alpha \in K$가 $k$ 위에서 분리 대수적이라고 하자.
그러면 $k' = k(\alpha)$는 $K$에 포함된 $k$의 유한 분리 가능 확대이다.
보조정리 \ref{algebra-lemma-subalgebra-geometrically-irreducible}에 의해
확대 $k'/k$는 기하적으로 기약이다.
특히 $k' \otimes_k \overline{k}$의 스펙트럼은 기약이다.
(따라서 이 환이 체들의 곱이면 인자가 정확히 하나이다.)
체론, 보조정리
\ref{fields-lemma-finite-separable-tensor-alg-closed}에 의해
$\Hom_k(k', \overline{k})$가 하나의 원소를 가짐을 얻고, 이어서
체론, 보조정리 \ref{fields-lemma-separable-equality}에 의해
$k' = k$임을 얻는다. 따라서 (2)가 성립한다.

\medskip\noindent
(2)를 가정하자. $k' \subset K$를 $k$ 위에서 대수적인 원소들로
이루어진 부분체라 하자.
보조정리 \ref{algebra-lemma-field-extension-geometrically-irreducible}에 의해
확대 $K/k'$는 기하적으로 기약이다.
가정에 의해 $k'/k$는 순수 비분리 확대이다.
보조정리 \ref{algebra-lemma-p-ring-map}에 의해 확대
$k'/k$는 기하적으로 기약이다. 따라서 보조정리
\ref{algebra-lemma-geometrically-irreducible-transitive}에 의해
$K/k$가 기하적으로 기약임을 얻는다.
\end{proof}

\begin{lemma}
\label{algebra-lemma-make-geometrically-irreducible}
$K/k$를 체확대라 하자. $k$ 위에서 분리 대수적인 원소들로
이루어진 부분확대 $K/k'/k$를 생각하자.
그러면 $K$는 $k'$ 위에서 기하적으로 기약이다.
$K/k$가 유한 생성 체확대이면 $[k' : k] < \infty$이다.
\end{lemma}

\begin{proof}
첫 번째 주장은 보조정리
\ref{algebra-lemma-geometrically-irreducible-separable-elements}와,
분리 대수적 확대의 추이성에 의해 $k'$ 위에서 분리 대수적인 원소들이
$k'$에 속한다는 사실에서 바로 따른다. 체론, 보조정리
\ref{fields-lemma-separable-permanence}를 참조하라.
$K/k$가 유한 생성이면 체론, 보조정리
\ref{fields-lemma-algebraic-closure-in-finitely-generated}에 의해
$k'$는 $k$ 위에서 유한이다.
\end{proof}

\begin{lemma}
\label{algebra-lemma-Galois-orbit}
$K/k$를 체확대라 하자.
$\overline{k}/k$를 분리 대수적 폐포라 하자.
그러면 $\text{Gal}(\overline{k}/k)$는
$\overline{k} \otimes_k K$의 소 아이디얼들 위에 추이적으로 작용한다.
\end{lemma}

\begin{proof}
보조정리 \ref{algebra-lemma-make-geometrically-irreducible}에서 얻은
부분확대를 $K/k'/k$라 하자.
% P01-E1362
$k \subset \overline{k}$가 정수적이므로
$\overline{k} \otimes_k K$와 $\overline{k} \otimes_k k'$의
모든 소 아이디얼은 극대이다. 보조정리
\ref{algebra-lemma-integral-no-inclusion}을 참조하라.
보조정리 \ref{algebra-lemma-geometrically-irreducible-any-base-change}에 의해
사상
$$
\Spec(\overline{k} \otimes_k K) \to \Spec(\overline{k} \otimes_k k')
$$
은 전단사이다. 실제로 (1) 모든 소 아이디얼이 극소 소 아이디얼이고,
(2) $\overline{k} \otimes_k K = (\overline{k} \otimes_k k') \otimes_{k'} K$이며,
(3) $K$가 $k'$ 위에서 기하적으로 기약이기 때문이다.
따라서 $\overline{k} \otimes_k k'$의 소 아이디얼들 위의
$\text{Gal}(\overline{k}/k)$의 작용에 대해 보조정리를 증명하면 충분하다.

\medskip\noindent
$\overline{k} \otimes_k k'$의 모든 소 아이디얼이 극대이므로
그 잉여체들은 $\overline{k}$와 동형이다.
따라서 $\overline{k} \otimes_k k'$의 소 아이디얼들은
$\Hom_k(k', \overline{k})$의 원소들과 일대일로 대응한다.
여기서 $\sigma \in \Hom_k(k', \overline{k})$에 대응하는 것은
사상 $1 \otimes \sigma : \overline{k} \otimes_k k' \to \overline{k}$의
핵 $\mathfrak p_\sigma$이다.
특히 $\text{Gal}(\overline{k}/k)$는 원하는 대로
이 집합 위에 추이적으로 작용한다.
\end{proof}




\section{기하적으로 연결인 대수}
\label{algebra-section-geometrically-connected}

\begin{lemma}
\label{algebra-lemma-separably-closed-connected}
$k$를 분리폐체라 하자.
% P01-E1363
$R$, $S$를 $k$-대수들이라 하자. $\Spec(R)$와
$\Spec(S)$가 연결이면 $\Spec(R \otimes_k S)$도 연결이다.
\end{lemma}

\begin{proof}
$\Spec(R)$가 연결일 필요충분조건은 $R$이 자명하지 않은
멱등원을 갖지 않는 것임을 상기하자.
보조정리 \ref{algebra-lemma-characterize-spec-connected}를 참조하라.
따라서 보조정리 \ref{algebra-lemma-limit-argument}에 의해
$R$과 $S$가 $k$ 위에서 유한형이라고 가정해도 된다.
이 경우 $R$과 $S$는 뇌터 환이고 극소 소 아이디얼을
유한 개만 갖는다. 보조정리
\ref{algebra-lemma-Noetherian-irreducible-components}를 참조하라.
그러므로 $n$과 $m$이 각각 $\Spec(R)$과 $\Spec(S)$의
기약 성분의 개수일 때 $n + m$에 관한 귀납법으로 논증할 수 있다.
$n$이나 $m$ 가운데 하나가 영인 경우는 물론 자명하다.
$n = m = 1$, 즉 $\Spec(R)$과 $\Spec(S)$가 각각 하나의
기약 성분을 가지는 경우에는 보조정리
\ref{algebra-lemma-separably-closed-irreducible}에 의해 결과가 성립한다.
$n > 1$이라고 하자. $\mathfrak p \subset R$를
기약 닫힌 부분집합 $T \subset \Spec(R)$에 대응하는
극소 소 아이디얼이라 하자.
$T' \subset \Spec(R)$를 나머지 $n - 1$개 기약 성분들의 합집합이라 하자.
$T' = V(I) = \Spec(R/I)$가 되는 아이디얼 $I \subset R$를 택하자
(보조정리 \ref{algebra-lemma-spec-closed}).
극소 소 아이디얼을 주의 깊게 택하면 $T'$가 연결이 되도록 할 수도 있다.
위상론, 보조정리
\ref{topology-lemma-remove-irreducible-connected}를 참조하라.
그러면 $T \cup T' = \Spec(R)$이고,
$\Spec(R)$가 연결이라는 가정에 의해
$T \cap T' = V(\mathfrak p + I) = \Spec(R/(\mathfrak p + I))$는
공집합이 아니다.
$\Spec(R \otimes_k S)$에서 $T$의 역상은
$\Spec(R/\mathfrak p \otimes_k S)$이고,
% P01-E1364
$\Spec(R \otimes_k S)$에서 $T'$의 역상은
$\Spec(R/I \otimes_k S)$이다. 귀납가정에 의해 둘 다 연결이다.
$T \cap T'$의 역상은
$\Spec(R/(\mathfrak p + I) \otimes_k S)$이고, 이는 공집합이 아니다.
따라서 $\Spec(R \otimes_k S)$는 연결이다.
\end{proof}

\begin{lemma}
\label{algebra-lemma-geometrically-connected}
$k$를 체라 하자.
$R$을 $k$-대수라 하자.
다음은 서로 동치이다.
\begin{enumerate}
\item 모든 체확대 $k'/k$에 대해
$R \otimes_k k'$의 스펙트럼이 연결이다.
\item 모든 유한 분리 가능 체확대 $k'/k$에 대해
$R \otimes_k k'$의 스펙트럼이 연결이다.
\end{enumerate}
\end{lemma}

\begin{proof}
임의의 체확대 $k'/k$에 대해
$R \otimes_k k'$의 스펙트럼이 연결인 것은
$R \otimes_k k'$가 자명하지 않은 멱등원을 갖지 않는 것과 동치이다.
보조정리 \ref{algebra-lemma-characterize-spec-connected}를 참조하라.
(2)를 가정하자. $k \subset \overline{k}$를
$k$의 분리 대수적 폐포라 하자.
보조정리 \ref{algebra-lemma-limit-argument}을 사용하면,
(2)는 $R \otimes_k \overline{k}$가 자명하지 않은
멱등원을 갖지 않는 것과 동치임을 알 수 있다.
임의의 체확대 $k'/k$에 대해 $k' \subset \overline{k}'$이고
$\overline{k}'/\overline{k}$인 체확대를 찾을 수 있다.
보조정리 \ref{algebra-lemma-separably-closed-connected}에 의해
$R \otimes_k \overline{k}'$는 자명하지 않은 멱등원을 갖지 않는다.
$R \otimes_k k'$가 자명하지 않은 멱등원을 가지면
$R \otimes_k \overline{k}'$도 그러한 멱등원을 가지므로 모순이다.
\end{proof}

\begin{definition}
\label{algebra-definition-geometrically-connected}
$k$를 체라 하자.
$S$를 $k$-대수라 하자.
모든 체확대 $k'/k$에 대해
$S \otimes_k k'$의 스펙트럼이 연결이면
$S$가 {\it $k$ 위에서 기하적으로 연결}이라고 한다.
\end{definition}

\noindent
보조정리 \ref{algebra-lemma-geometrically-connected}에 의해,
유한 분리 가능 체확대 $k'/k$들에 대해서만 확인해도 충분하다.

\begin{lemma}
\label{algebra-lemma-separably-closed-connected-implies-geometric}
$k$를 체라 하자.
$R$을 $k$-대수라 하자.
$k$가 분리폐체이면, $R$이 $k$ 위에서 기하적으로 연결일
필요충분조건은 $R$의 스펙트럼이 연결인 것이다.
\end{lemma}

\begin{proof}
정의 \ref{algebra-definition-geometrically-connected} 뒤의
주석에서 바로 따른다.
\end{proof}

\begin{lemma}
\label{algebra-lemma-subalgebra-geometrically-connected}
$k$를 체라 하자. $S$를 $k$-대수라 하자.
\begin{enumerate}
\item $S$가 $k$ 위에서 기하적으로 연결이면 모든
$k$-부분대수도 그러하다.
\item $S$의 모든 유한 생성 $k$-부분대수가 기하적으로 연결이면
$S$도 기하적으로 연결이다.
\item 기하적으로 연결인 $k$-대수들의 유향 귀납적 극한은
기하적으로 연결이다.
\end{enumerate}
\end{lemma}

\begin{proof}
이는 자명하지 않은 멱등원이 존재하지 않는다는 말로 연결성을
특징짓는 것에서 따른다.
% P01-E1365
두 번째와 세 번째 성질은 텐서곱이 귀납적 극한과 가환한다는
사실에서 따른다.
\end{proof}

\noindent
다음 보조정리는 더 일반적인 대수다양체, 보조정리
\ref{varieties-lemma-bijection-connected-components}에 의해 대체된다.

\begin{lemma}
\label{algebra-lemma-geometrically-connected-any-base-change}
$k$를 체라 하자.
$S$를 기하적으로 연결인 $k$-대수라 하자.
$R$을 임의의 $k$-대수라 하자.
사상
$$
R \longrightarrow R \otimes_k S
$$
은 멱등원들의 집합 사이의 전단사를 유도하고, 사상
$$
\Spec(R \otimes_k S) \longrightarrow \Spec(R)
$$
은 연결 성분들의 집합 사이의 전단사를 유도한다.
\end{lemma}

\begin{proof}
두 번째 주장은 첫 번째 주장과 보조정리
\ref{algebra-lemma-connected-component}에서 따른다.
보조정리 \ref{algebra-lemma-subalgebra-geometrically-connected}와
\ref{algebra-lemma-limit-argument}에 의해 $R$과 $S$가
$k$ 위에서 유한형이라고 가정해도 된다.
그러면 $R \otimes_k S$도 $k$ 위에서 유한형이다.
이 경우 모든 환은 뇌터 환이므로 그 스펙트럼들은 연결 성분을
유한 개만 갖는다는 점에 유의하라
(기약 성분을 유한 개만 갖기 때문이다. 보조정리
\ref{algebra-lemma-Noetherian-irreducible-components}를 참조하라).
특히 여기서 모든 연결 성분은 열린집합이다!
따라서 보조정리 \ref{algebra-lemma-disjoint-implies-product}에 의해
이 경우 첫 번째 주장은 두 번째 주장과 동치이다.
이를 증명하자.
대수 $S$가 기하적으로 연결이고 영이 아니므로
$X = \Spec(R \otimes_k S) \to \Spec(R) = Y$의 모든 섬유는
연결이고 공집합이 아니다.
또한 $R \to R \otimes_k S$가 평탄하고 유한 표시이므로
사상 $X \to Y$는 열린 사상이다
(명제 \ref{algebra-proposition-fppf-open}).
% P01-E1366
위상론, 보조정리
\ref{topology-lemma-connected-fibres-connected-components}에 의해
$X \to Y$는 연결 성분들의 집합 사이의 전단사를 유도한다.
\end{proof}





\section{기하적으로 정역인 대수}
\label{algebra-section-geometrically-integral}

\noindent
정의는 다음과 같다.

\begin{definition}
\label{algebra-definition-geometrically-integral}
$k$를 체라 하자.
$S$를 $k$-대수라 하자.
% P01-E1367
모든 체확대 $k'/k$에 대해 $S \otimes_k k'$가 정역이면,
$S$가 {\it $k$ 위에서 기하적으로 정역}이라고 한다.
\end{definition}

\noindent
% P01-E1368
기하적으로 정역인 대수에 관한 모든 문제는
기하적으로 축소이고 기약인 대수에 관한 문제로 바꿀 수 있다.

\begin{lemma}
\label{algebra-lemma-geometrically-integral}
$k$를 체라 하자.
$S$를 $k$-대수라 하자.
이 경우 $S$가 $k$ 위에서 기하적으로 정역일 필요충분조건은
$S$가 $k$ 위에서 기하적으로 기약이고 기하적으로 축소인 것이다.
\end{lemma}

\begin{proof}
생략한다.
\end{proof}

\begin{lemma}
\label{algebra-lemma-characterize-geometrically-integral}
$k$를 체라 하자. $S$를 $k$-대수라 하자.
다음은 서로 동치이다.
\begin{enumerate}
\item $S$는 $k$ 위에서 기하적으로 정역이다.
\item 모든 유한 체확대 $k'/k$에 대해
$S \otimes_k k'$는 정역이다.
\item $\overline{k}$가 $k$의 대수적 폐포일 때
$S \otimes_k \overline{k}$는 정역이다.
\end{enumerate}
\end{lemma}

\begin{proof}
보조정리 \ref{algebra-lemma-geometrically-integral},
\ref{algebra-lemma-geometrically-reduced-finite-purely-inseparable-extension}, 그리고
\ref{algebra-lemma-geometrically-irreducible}에서 따른다.
\end{proof}

\begin{lemma}
\label{algebra-lemma-geometrically-integral-any-integral-base-change}
$k$를 체라 하고 $S$를 기하적으로 정역인 $k$-대수라 하자.
$R$을 정역인 $k$-대수라 하자. 그러면
$R \otimes_k S$는 정역이다.
\end{lemma}

\begin{proof}
보조정리 \ref{algebra-lemma-geometrically-reduced-any-reduced-base-change}에 의해
환 $R \otimes_k S$는 축소이고, 보조정리
\ref{algebra-lemma-geometrically-irreducible-any-base-change}에 의해
환 $R \otimes_k S$는 기약이다
(스펙트럼이 기약 성분을 하나만 갖는다).
따라서 $R \otimes_k S$는 정역이다.
\end{proof}




\section{값매김환}
\label{algebra-section-valuation-rings}

\noindent
몇 가지 정의부터 제시한다.

\begin{definition}
\label{algebra-definition-valuation-ring}
값매김환.
\begin{enumerate}
\item $K$를 체라 하자. $A$, $B$를 $K$에 포함된 국소환들이라 하자.
$A \subset B$이고 $\mathfrak m_A = A \cap \mathfrak m_B$이면
$B$가 $A$를 {\it 지배한다}고 한다.
\item $A$를 환이라 하자. $A$가 국소 정역이고,
$A$의 분수체에 포함된 국소환들 사이의 지배 관계에 대해
$A$가 극대이면 $A$를 {\it 값매김환}이라고 한다.
\item $A$를 분수체 $K$를 갖는 값매김환이라 하자.
$R \subset K$가 $K$의 부분환일 때 $R \subset A$이면
$A$가 $R$에 {\it 중심을 둔다}고 한다.
\end{enumerate}
\end{definition}

\noindent
이 정의에 따르면 체도 값매김환이다.

\begin{lemma}
\label{algebra-lemma-dominate}
$K$를 체라 하고 $A \subset K$를 국소 부분환이라 하자.
그러면 $A$를 지배하고 분수체가 $K$인 값매김환이 존재한다.
\end{lemma}

\begin{proof}
$K$의 국소 부분환들의 모음을 지배 관계로 순서화한
부분순서집합으로 생각하자.
$\{A_i\}_{i \in I}$가 $K$의 국소 부분환들의 전순서화된
모음이라고 하자. 그러면 $B = \bigcup A_i$는 모든 $A_i$를
지배하는 국소 부분환이다.
따라서 초른의 보조정리에 의해 다음을 보이면 충분하다.
$A \subset K$가 분수체가 $K$가 아닌 국소환이면,
$A$를 지배하는 국소환 $B \subset K$, $B \not = A$가 존재한다.

\medskip\noindent
$A$의 분수체에 속하지 않는 $t \in K$를 택하자.
$t$가 $A$ 위에서 초월적이면 $A[t] \subset K$이고,
$A[t]_{(t, \mathfrak m)} \subset K$는 $A$와 다르면서
$A$를 지배하는 국소환이다.
$t$가 $A$ 위에서 대수적이라고 하자.
그러면 어떤 영이 아닌 $a \in A$에 대해 $at$는 $A$ 위에서 정수적이다.
이 경우 $A$와 $ta$로 생성되는 부분환 $A' \subset K$는
$A$ 위에서 유한이다.
보조정리 \ref{algebra-lemma-integral-overring-surjective}에 의해
$\mathfrak m$ 위에 놓이는 소 아이디얼 $\mathfrak m' \subset A'$가 존재한다.
그러면 $A'_{\mathfrak m'}$는 $A$를 지배한다.
$A = A'_{\mathfrak m'}$이면 $t$가 $A$의 분수체에 속하는데,
이는 그렇지 않다는 가정과 모순이다.
따라서 원하는 대로 $A \not = A'_{\mathfrak m'}$이다.
\end{proof}

\begin{lemma}
\label{algebra-lemma-valuation-ring-normal}
$A$를 값매김환이라 하자.
그러면 $A$는 정규 정역이다.
\end{lemma}

\begin{proof}
$x$가 $A$의 분수체에 속하고 $A$ 위에서 정수적이라고 하자.
% P01-E1369
$A'$를 $A$와 $x$로 생성되는 $K$의 부분환이라 하자.
$A\subset A'$는 정수적 확대이므로 보조정리
\ref{algebra-lemma-integral-overring-surjective}에 의해
$\mathfrak m$ 위에 놓이는 소 아이디얼
$\mathfrak m' \subset A'$가 존재한다.
그러면 $A'_{\mathfrak m'}$는 $A$를 지배한다.
$A$가 값매김환이므로 $A=A'_{\mathfrak m'}$이고, 따라서 $x\in A$이다.
\end{proof}


\begin{lemma}
\label{algebra-lemma-valuation-ring-x-or-x-inverse}
$A$를 극대 아이디얼 $\mathfrak m$과 분수체 $K$를 갖는
값매김환이라 하자.
$x \in K$라 하자. 그러면 $x \in A$이거나 $x^{-1} \in A$이거나
둘 다이다.
\end{lemma}

\begin{proof}
$x$가 $A$에 속하지 않는다고 가정하자.
$A'$를 $A$와 $x$로 생성되는 $K$의 부분환이라 하자.
$A$가 값매김환이므로 $\mathfrak m$ 위에 놓이는
$A'$의 소 아이디얼은 없다.
$\mathfrak m$이 극대이므로 $V(\mathfrak m A') = \emptyset$이다.
그러면 보조정리 \ref{algebra-lemma-Zariski-topology}에 의해
$\mathfrak m A' = A'$이다.
따라서 $t_i \in \mathfrak m$인 원소들을 사용하여
$1 = \sum_{i = 0}^d t_i x^i$로 쓸 수 있다.
이는
$(1 - t_0) (x^{-1})^d - \sum t_i (x^{-1})^{d - i} = 0$을 함의한다.
특히 $x^{-1}$은 $A$ 위에서 정수적이고, 따라서 보조정리
\ref{algebra-lemma-valuation-ring-normal}에 의해 $x^{-1} \in A$이다.
\end{proof}

\begin{lemma}
\label{algebra-lemma-x-or-x-inverse-valuation-ring}
$A \subset K$를 체 $K$의 부분환이라 하자.
모든 $x \in K$에 대해 $x \in A$이거나 $x^{-1} \in A$이거나
둘 다라고 가정하자.
그러면 $A$는 분수체가 $K$인 값매김환이다.
\end{lemma}

\begin{proof}
$A$가 $K$가 아니면 $A$는 체가 아니므로 영이 아닌
극대 아이디얼 $\mathfrak m$이 존재한다.
$\mathfrak m'$이 두 번째 극대 아이디얼이면,
$x \in \mathfrak m$, $y \not \in \mathfrak m$,
$x \not \in \mathfrak m'$, $y \in \mathfrak m'$이 되도록
$x, y \in A$를 택하자.
그러면 $x/y \in A$도 $y/x \in A$도 아니므로
보조정리의 가정에 모순이다. 따라서 $A$는 국소환이다.
$A$를 지배하면서 $K$에 포함된 국소환을 $A'$라 하자.
$x \in A'$라 하자. $x \in A$임을 보여야 한다.
그렇지 않으면 $x^{-1} \in A$이고, 물론
$x^{-1} \in \mathfrak m_A$이다. 그러나 그러면
$x^{-1} \in \mathfrak m_{A'}$이고, 이는 $x \in A'$에 모순이다.
\end{proof}

\begin{lemma}
\label{algebra-lemma-colimit-valuation-rings}
\begin{slogan}
값매김환은 여과 직접극한 아래에서 보존된다.
\end{slogan}
$I$를 유향집합이라 하자. $(A_i, \varphi_{ij})$를
$I$ 위의 값매김환들의 계라 하자.
그러면 $A = \colim A_i$는 값매김환이다.
\end{lemma}

\begin{proof}
$A$가 정역인 것은 명백하다. $a, b \in A$라 하자.
보조정리 \ref{algebra-lemma-x-or-x-inverse-valuation-ring}에 의해
$A$에서 $a | b$이거나 $b | a$임을 보이면 충분하다.
$a_i, b_i \in A_i$가 존재하여 $a, b$로 사상되도록
$i$를 충분히 크게 택하자.
$A_i$ 안의 $a_i, b_i$에 보조정리
\ref{algebra-lemma-valuation-ring-x-or-x-inverse}를 적용하면
$A$ 안의 $a, b$에 대한 결과를 얻는다.
\end{proof}

\begin{lemma}
\label{algebra-lemma-valuation-ring-cap-field}
$L/K$를 체확대라 하자. $B \subset L$가 값매김환이면
$A = K \cap B$는 값매김환이다.
\end{lemma}

\begin{proof}
$L$을 $B$의 분수체 $F$로, $K$를 $K \cap F$로 대체할 수 있다.
그러면 보조정리 \ref{algebra-lemma-valuation-ring-x-or-x-inverse}와
\ref{algebra-lemma-x-or-x-inverse-valuation-ring}을 함께 적용하여
보조정리를 얻는다.
\end{proof}

\begin{lemma}
\label{algebra-lemma-valuation-ring-cap-field-finite}
$L/K$를 대수적 체확대라 하자. $B \subset L$가
분수체가 $L$인 값매김환이고 체는 아니라고 하자.
그러면 $A = K \cap B$는 체가 아닌 값매김환이다.
\end{lemma}

\begin{proof}
보조정리 \ref{algebra-lemma-valuation-ring-cap-field}에 의해
$A$는 값매김환이다. $A$가 체이면 $A = K$이다.
그러면 $A = K \subset B$는 정수적 확대이므로
$B$의 소 아이디얼들 사이에는 진포함이 없다
(보조정리 \ref{algebra-lemma-integral-no-inclusion}).
이는 $B$가 체가 아닌 국소 정역이라는 가정과 모순이다.
\end{proof}

\begin{lemma}
\label{algebra-lemma-make-valuation-rings}
$A$를 값매김환이라 하자. 모든 소 아이디얼
$\mathfrak p \subset A$에 대해 몫 $A/\mathfrak p$는 값매김환이다.
국소화 $A_\mathfrak p$도 그러하며, 사실 $A$의 임의의 국소화도 그러하다.
\end{lemma}

\begin{proof}
보조정리 \ref{algebra-lemma-x-or-x-inverse-valuation-ring}에서 주어진
값매김환의 특징짓기를 사용하라.
\end{proof}

\begin{lemma}
\label{algebra-lemma-stack-valuation-rings}
$A'$를 잉여체가 $K$인 값매김환이라 하자.
$A$를 분수체가 $K$인 값매김환이라 하자.
그러면
$C = \{\lambda \in A' \mid \lambda \bmod \mathfrak m_{A'} \in A\}$는
값매김환이다.
\end{lemma}

\begin{proof}
$\mathfrak m_{A'} \subset C$이고 $C/\mathfrak m_{A'} = A$임에 유의하자.
특히 $C$의 분수체는 $A'$의 분수체와 같다.
보조정리 \ref{algebra-lemma-x-or-x-inverse-valuation-ring}의 판정법을
사용하여 이 보조정리를 증명하겠다.
$x$를 $C$의 분수체의 원소라 하자.
그 보조정리에 의해 $x \in A'$라고 가정해도 된다.
$x \in \mathfrak m_{A'}$이면 $x \in C$임을 알 수 있다.
그렇지 않으면 $x$는 $A'$의 가역원이고 $x^{-1} \in A'$도 성립한다.
그러므로 같은 보조정리에 의해 $x$ 또는 $x^{-1}$ 가운데 하나가
$A$의 원소로 사상된다.
\end{proof}

\begin{lemma}
\label{algebra-lemma-find-valuation-rings}
$A$를 분수체가 $K$인 정규 정역이라 하자.
\begin{enumerate}
\item 모든 $x \in K$, $x \not \in A$에 대해,
분수체가 $K$이고 $x \not \in V$인 값매김환
$A \subset V \subset K$가 존재한다.
\item $A$가 국소환이면, 더 나아가 $A$를 지배하는 $V$를 택할 수 있다.
\end{enumerate}
다르게 말하면 $A$는 $K$ 안에서 $A$를 포함하는 모든 값매김환들의
교집합이고, $A$가 국소환이면 $A$는 $K$ 안에서 $A$를 지배하는
모든 값매김환들의 교집합이다.
\end{lemma}

\begin{proof}
$x \in K$, $x \not \in A$라고 하자. $B = A[x^{-1}]$를 생각하자.
그러면 $x \not \in B$이다. 실제로
$x = a_0 + a_1x^{-1} + \ldots + a_d x^{-d}$이면
$x^{d + 1} - a_0x^d - \ldots - a_d = 0$이므로 $x$는
$A$ 위에서 정수적이다. 이는 $A$가 정규라는 사실과 모순이다.
따라서 $x^{-1}$은 $B$의 가역원이 아니다.
그러므로 $V(x^{-1}) \subset \Spec(B)$는 공집합이 아니며
(보조정리 \ref{algebra-lemma-Zariski-topology}),
$x^{-1} \in \mathfrak p$인 소 아이디얼 $\mathfrak p \subset B$를
택할 수 있다. $B_\mathfrak p$를 지배하는 값매김환
$V \subset K$를 택하자(보조정리 \ref{algebra-lemma-dominate}).
$x^{-1} \in \mathfrak m_V$이므로 $x \not \in V$이다.

\medskip\noindent
$A$가 국소환이면
$x^{-1} B + \mathfrak m_A B \not = B$라고 주장한다.
실제로 $a_i \in A$와 $a'_i \in \mathfrak m_A$에 대해
$1 = (a_0 + a_1x^{-1} + \ldots + a_d x^{-d})x^{-1} +
a'_0 + \ldots + a'_d x^{-d}$이면 다음을 얻는다.
$$
(1 - a'_0) x^{d + 1} - (a_0 + a'_1) x^d - \ldots - a_d = 0
$$
$a'_0 \in \mathfrak m_A$이므로 $1 - a'_0$은 $A$의 가역원이다.
따라서 앞에서와 마찬가지로 $x$가 $A$ 위에서 정수적이라는
모순을 얻는다.
% P01-E1370
그런 다음 $\mathfrak p \supset x^{-1} B + \mathfrak m_A B$인
그러한 소 아이디얼을 택하면 $A$를 지배하는 $V$를 얻는다.
\end{proof}

\noindent
{\it 전순서 아벨군}은 아벨군 $\Gamma$와 그 위의 전순서
$\geq$로 이루어진 쌍 $(\Gamma, \geq)$이며,
모든 $\gamma, \gamma', \gamma'' \in \Gamma$에 대해
$\gamma \geq \gamma' \Rightarrow
\gamma + \gamma'' \geq \gamma' + \gamma''$를 만족한다.

\begin{lemma}
\label{algebra-lemma-valuation-group}
$A$를 분수체가 $K$인 값매김환이라 하자.
$\Gamma = K^*/A^*$라 두자(군법은 덧셈으로 쓴다).
$\gamma, \gamma' \in \Gamma$에 대해
$\gamma - \gamma'$가 $A - \{0\} \to \Gamma$의 상에 속할
필요충분조건이 $\gamma \geq \gamma'$인 것으로 정의하자.
그러면 $(\Gamma, \geq)$는 전순서 아벨군이다.
\end{lemma}

\begin{proof}
생략한다. 보조정리
\ref{algebra-lemma-valuation-ring-x-or-x-inverse}에서 쉽게 따른다.
$A = K$인 경우에는 유일한 전순서를 갖춘 영군
$\Gamma = \{0\}$을 얻는다는 점에 유의하라.
\end{proof}

\begin{definition}
\label{algebra-definition-value-group}
$A$를 값매김환이라 하자.
\begin{enumerate}
\item 보조정리 \ref{algebra-lemma-valuation-group}의 전순서 아벨군
$(\Gamma, \geq)$를 값매김환 $A$의 {\it 값군}이라고 한다.
\item 사상 $v : A - \{0\} \to \Gamma$와
$v : K^* \to \Gamma$를 $A$에 연관된 {\it 값매김}이라고 한다.
\item $\Gamma \cong \mathbf{Z}$이면 값매김환 $A$를
{\it 이산 값매김환}이라고 한다.
\end{enumerate}
\end{definition}

\noindent
$\Gamma \cong \mathbf{Z}$이면 $1 \geq 0$을 만족하는 그러한
동형사상은 유일하다. 이와 같이 동형사상을 택하면 순서는
정수의 통상적인 순서가 된다.

\begin{lemma}
\label{algebra-lemma-properties-valuation}
$A$를 값매김환이라 하자.
값매김 $v : A -\{0\} \to \Gamma_{\geq 0}$은 다음 성질을 갖는다.
\begin{enumerate}
\item $v(a) = 0 \Leftrightarrow a \in A^*$,
\item $v(ab) = v(a) + v(b)$,
\item $a + b \not = 0$이면 $v(a + b) \geq \min(v(a), v(b))$.
\end{enumerate}
\end{lemma}

\begin{proof}
생략한다.
\end{proof}

\begin{lemma}
\label{algebra-lemma-characterize-valuation-ring}
$A$를 환이라 하자. 다음은 서로 동치이다.
\begin{enumerate}
\item $A$는 값매김환이다.
\item $A$는 국소 정역이고 $A$의 모든 유한 생성 아이디얼은 주아이디얼이다.
\end{enumerate}
\end{lemma}

\begin{proof}
$A$가 값매김환이라고 하고 $f_1, \ldots, f_n \in A$라 하자.
$v(f_j)$들 가운데 $v(f_i)$가 최소가 되도록 $i$를 택하자.
그러면 $(f_i) = (f_1, \ldots, f_n)$이다.
역으로 $A$가 국소 정역이고 $A$의 모든 유한 생성 아이디얼이
주아이디얼이라고 가정하자.
$f, g \in A$를 택하고 $(f, g) = (h)$로 쓰자.
그러면 어떤 $a, b, c, d \in A$에 대해
$f = ah$, $g = bh$, $h = cf + dg$이다.
따라서 $ac + bd = 1$이고 $a$나 $b$ 가운데 하나는 가역원이다.
즉 $g/f$나 $f/g$ 가운데 하나는 $A$의 원소이다.
보조정리 \ref{algebra-lemma-x-or-x-inverse-valuation-ring}에 의해
$A$는 값매김환이다.
\end{proof}

\begin{lemma}
\label{algebra-lemma-valuation-valuation-ring}
$(\Gamma, \geq)$를 전순서 아벨군이라 하자.
$K$를 체라 하자. $v : K^* \to \Gamma$를 아벨군의 준동형이라 하고,
$a, b, a + b$가 영이 아닌 $a, b \in K$에 대해
$v(a + b) \geq \min(v(a), v(b))$라고 가정하자. 그러면
$$
A =
\{
x \in K \mid x = 0 \text{ or } v(x) \geq 0
\}
$$
는 값군이 $\Im(v) \subset \Gamma$이고 극대 아이디얼이
$$
\mathfrak m =
\{
x \in K \mid x = 0 \text{ or } v(x) > 0
\}
$$
이며 가역원군이
$$
A^* =
\{
x \in K^* \mid v(x) = 0
\}.
$$
인 값매김환이다.
\end{lemma}

\begin{proof}
생략한다.
\end{proof}

\noindent
$(\Gamma, \geq)$를 전순서 아벨군이라 하자.
{\it $\Gamma$의 아이디얼}은 다음 성질을 갖는 부분집합
$I \subset \Gamma$이다. $I$의 모든 원소는 $\geq 0$이고,
$\gamma \in I$, $\gamma' \geq \gamma$이면 $\gamma' \in I$이다.
그러한 아이디얼이 $0 \not \in I$이고
$\gamma + \gamma' \in I, \gamma, \gamma' \geq 0
\Rightarrow \gamma \in I \text{ or } \gamma' \in I$를 만족하면
{\it 소 아이디얼}이라고 한다.

\begin{lemma}
\label{algebra-lemma-ideals-valuation-ring}
$A$를 값매김환이라 하자.
$A$의 아이디얼들은 $\Gamma$의 아이디얼들과 $1 - 1$로 대응한다.
이 전단사는 포함관계를 보존하고 소 아이디얼을 소 아이디얼로 보낸다.
\end{lemma}

\begin{proof}
생략한다.
\end{proof}

\begin{lemma}
\label{algebra-lemma-valuation-ring-Noetherian-discrete}
값매김환이 뇌터 환일 필요충분조건은
이산 값매김환이거나 체인 것이다.
\end{lemma}

\begin{proof}
$A$가 이산 값매김환이고 값매김
$v : A \setminus \{0\} \to \mathbf{Z}$가
$\Im(v) = \mathbf{Z}_{\geq 0}$이 되도록 정규화되어 있다고 하자.
보조정리 \ref{algebra-lemma-ideals-valuation-ring}에 의해 $A$의 아이디얼들은
부분집합 $I_n = \{0\} \cup v^{-1}(\mathbf{Z}_{\geq n})$이다.
$v(x) = n$인 임의의 원소 $x \in A$가 $I_n$을 생성함은 명백하다.
따라서 $A$는 주아이디얼 정역이고, 특히 뇌터 환이다.

\medskip\noindent
$A$가 값군 $\Gamma$를 갖는 뇌터 값매김환이라고 하자.
보조정리 \ref{algebra-lemma-ideals-valuation-ring}에 의해
$\Gamma$의 아이디얼들에 대해 상승 사슬 조건이 성립한다.
$A$가 체가 아니라고 가정해도 된다. 즉,
$\gamma > 0$인 $\gamma \in \Gamma$가 존재한다.
$\gamma > 0$인 부분집합 $\gamma + \Gamma_{\geq 0}$들에
상승 사슬 조건을 적용하면 $0$보다 큰 최소 원소
$\gamma_0$가 존재함을 알 수 있다.
$\gamma \in \Gamma$가 $\gamma > 0$인 원소라고 하자.
$\gamma$, $\gamma - \gamma_0$, $\gamma - 2\gamma_0$ 등의
원소들로 이루어진 수열을 생각하자.
상승 사슬 조건에 의해 이들이 모두 $> 0$일 수는 없다.
$\gamma - n \gamma_0$를 마지막으로 $\geq 0$인 원소라 하자.
$\gamma_0$의 최소성에 의해 $0 = \gamma - n \gamma_0$이다.
따라서 원하는 대로 $\Gamma$는 순환군이다.
\end{proof}




\section{뇌터 환에 관한 추가 결과}
\label{algebra-section-Noetherian-again}


\begin{lemma}
\label{algebra-lemma-Noetherian-basic}
$R$을 뇌터 환이라 하자.
모든 유한 $R$-모듈은 유한 표시이다.
유한 $R$-모듈의 모든 부분모듈은 유한이다.
유한 $R$-모듈의 $R$-부분모듈들에 대해 상승 사슬 조건이 성립한다.
\end{lemma}

\begin{proof}
먼저 유한 $R$-모듈 $M$의 모든 부분모듈 $N$이 유한임을 보이자.
$M$의 생성원 개수에 관한 귀납법을 사용한다.
이 개수가 $1$이면 어떤 아이디얼 $I \subset J \subset R$에 대해
$N = J/I \subset M = R/I$이다.
따라서 뇌터성의 정의에서 결과가 따른다.
$M$의 생성원 개수가 $1$보다 크면, $M'$과 $M''$의 생성원 개수가
더 적은 짧은 완전열
$0 \to M' \to M \to M'' \to 0$을 찾을 수 있다.
$N' = M' \cap N$, $N'' = \Im(N \to M'')$로 두면
$N$에 대해서도 같은 형태의 짧은 완전열을 얻는다.
$N$의 생성원 개수는 $N'$의 생성원 개수와 $N''$의 생성원 개수의
합 이하이므로 귀납가정에서 결과가 따른다.

\medskip\noindent
$M$이 유한 표시임을 보이려면 표시 $R^n \to M$의 핵에
앞의 결과를 적용하면 된다.

\medskip\noindent
$M$의 $R$-부분모듈들에 대한 상승 사슬 조건이
$M$의 모든 부분모듈이 유한 $R$-모듈이라는 조건과 동치임은
잘 알려져 있고 쉽게 증명된다. 증명은 생략한다.
\end{proof}

\begin{lemma}[아르틴--리스]
\label{algebra-lemma-Artin-Rees}
% P01-E0362
$R$을 뇌터 환이라 하고 $I \subset R$를 아이디얼이라 하자.
$N \subset M$을 유한 $R$-모듈들이라 하자.
모든 $n \geq c$에 대해
$I^n M \cap N  =  I^{n-c}(I^cM \cap N)$이 되는 상수 $c > 0$이 존재한다.
\end{lemma}

\begin{proof}
환 $S = R \oplus I \oplus I^2 \oplus \ldots
= \bigoplus_{n \geq 0} I^n$을 생각하자. 관례로 $I^0 = R$이다.
곱셈은 $R$에서의 곱셈에 의해
$I^n \times I^m$을 $I^{n + m}$ 안으로 보낸다.
$I = (f_1, \ldots, f_t)$이면 $S$는 뇌터 환
$R[X_1, \ldots, X_t]$의 몫임에 유의하라.
이 사상은 단항식 $X_1^{e_1}\ldots X_t^{e_t}$를
$f_1^{e_1}\ldots f_t^{e_t}$로 보낼 뿐이다. 따라서 $S$는 뇌터 환이다.
마찬가지로 모듈 $M \oplus IM \oplus I^2M \oplus \ldots
= \bigoplus_{n \geq 0} I^nM$을 생각하자.
이는 유한 생성 $S$-모듈이다.
실제로 $x_1, \ldots, x_r$가 $R$ 위에서 $M$을 생성하면,
이들은 $S$ 위에서도 $\bigoplus_{n \geq 0} I^nM$을 생성한다.
이제 부분모듈 $\bigoplus_{n \geq 0} I^nM \cap N$을 생각하자.
쉽게 확인할 수 있듯이 이는 $S$-부분모듈이다.
보조정리 \ref{algebra-lemma-Noetherian-basic}에 의해 이는 유한 생성
$S$-모듈이다. 그 생성원들을
$\xi_j \in \bigoplus_{n \geq 0} I^nM \cap N$, $j = 1, \ldots, s$라 하자.
각 $\xi_j$를 동차 성분들로 분해하여, 어떤 $d_j$에 대해
각각 $\xi_j \in I^{d_j}M \cap N$이라고 가정해도 된다.
$c = \max\{d_j\}$로 두자. 그러면 모든 $n \geq c$에 대해
$I^nM \cap N$의 모든 원소는 $h_j \in I^{n - d_j}$인
$\sum h_j \xi_j$ 꼴이다.
이 사실과 자명한 포함관계
$I^{n-d_j}(I^{d_j}M \cap N) \subset I^{n-c}(I^cM \cap N)$에서
보조정리가 따른다.
\end{proof}

\begin{lemma}
\label{algebra-lemma-map-AR}
뇌터 환 $R$ 위의 유한 생성 모듈들의 완전열
$0 \to K \to M \xrightarrow{f} N$이 주어졌다고 하자.
$I \subset R$를 아이디얼이라 하자.
그러면 모든 $n \geq c$에 대해
$$
f^{-1}(I^nN) = K + I^{n-c}f^{-1}(I^cN)
\quad\text{and}\quad
f(M) \cap I^nN \subset f(I^{n - c}M)
$$
이 되는 $c$가 존재한다.
\end{lemma}

\begin{proof}
$\Im(f) \subset N$에 보조정리 \ref{algebra-lemma-Artin-Rees}를 적용하고,
$f : I^{n-c}M \to I^{n-c}f(M)$이 전사임에 유의하라.
\end{proof}

\begin{lemma}[크룰의 교집합 정리]
\label{algebra-lemma-intersect-powers-ideal-module-zero}
$R$을 뇌터 국소환이라 하자. $I \subset R$를 진아이디얼이라 하자.
$M$을 유한 $R$-모듈이라 하자. 그러면
$\bigcap_{n \geq 0} I^nM = 0$이다.
\end{lemma}

\begin{proof}
$N = \bigcap_{n \geq 0} I^nM$으로 두자.
그러면 모든 $n \geq 0$에 대해 $N = I^nM \cap N$이다.
아르틴--리스 보조정리 \ref{algebra-lemma-Artin-Rees}에 의해,
충분히 큰 어떤 $n$에 대해 $N = I^nM \cap N \subset IN$이다.
나카야마의 보조정리 \ref{algebra-lemma-NAK}에 의해 $N = 0$이다.
\end{proof}

\begin{lemma}
\label{algebra-lemma-intersection-powers-ideal-module}
$R$을 뇌터 환이라 하자. $I \subset R$를 아이디얼이라 하자.
$M$을 유한 $R$-모듈이라 하고 $N = \bigcap_n I^n M$으로 두자.
\begin{enumerate}
\item 모든 소 아이디얼 $\mathfrak p$, $I \subset \mathfrak p$에 대해
% P01-E0363
$N_f = 0$이 되는 $f \in R$, $f \not \in \mathfrak p$가 존재한다.
\item $I$가 $R$의 야코브슨 근기에 포함되면 $N = 0$이다.
\end{enumerate}
\end{lemma}

\begin{proof}
(1)의 증명. 모듈 $N$의 생성원들을 $x_1, \ldots, x_n$이라 하자.
보조정리 \ref{algebra-lemma-Noetherian-basic}를 참조하라.
모든 소 아이디얼 $\mathfrak p$, $I \subset \mathfrak p$에 대해,
보조정리 \ref{algebra-lemma-intersect-powers-ideal-module-zero}에 의해
$N$의 국소화 $M_{\mathfrak p}$ 안에서의 상은 영이다.
따라서 각 $x_i$가 $N_{g_i}$에서 영으로 사상되도록 하는
$g_i \in R$, $g_i \not \in \mathfrak p$를 찾을 수 있다.
그러므로 $N_{g_1g_2\ldots g_n} = 0$이다.

\medskip\noindent
(2)는 (1)과 보조정리 \ref{algebra-lemma-characterize-zero-local}에서 따른다.
\end{proof}

\begin{remark}
\label{algebra-remark-intersection-powers-ideal}
특히 보조정리 \ref{algebra-lemma-intersect-powers-ideal-module-zero}는
$R$이 뇌터 국소환이고 $I \subset R$가 비가역원 아이디얼이면
$\bigcap_n I^n = (0)$임을 함의한다.
더 일반적으로 $R$을 뇌터 환, $I \subset R$를 아이디얼이라 하고,
$f \in \bigcap_{n \in \mathbf{N}} I^n$이라고 하자.
그러면 보조정리 \ref{algebra-lemma-intersection-powers-ideal-module}에 의해
모든 소 아이디얼 $I \subset \mathfrak p$에 대해
$f$가 $R_g$에서 영으로 사상되도록 하는
$g \in R$, $g \not \in \mathfrak p$가 존재한다.
대수기하에서는 이를 ``$\Spec(R)$의 닫힌집합 $V(I)$의
어떤 열린 근방에서 $f$가 영이다''라고 표현한다.
\end{remark}

\begin{lemma}[아르틴--테이트]
\label{algebra-lemma-Artin-Tate}
$R$을 뇌터 환이라 하자. $S$를 유한 생성 $R$-대수라 하자.
$T \subset S$가 $R$-부분대수이고 $S$가 유한 생성 $T$-모듈이면,
$T$는 $R$ 위에서 유한형이다.
\end{lemma}

\begin{proof}
$R$-대수로서 $S$를 생성하는 원소들을
$x_1, \ldots, x_n \in S$라 하자.
$T$-모듈로서 $S$를 생성하는 원소들을
$S$ 안에서 $y_1, \ldots, y_m$이라 하자.
그러면 $x_i = \sum a_{ij} y_j$인 $a_{ij} \in T$가 존재한다.
또한 $y_i y_j = \sum b_{ijk} y_k$인 $b_{ijk} \in T$가 존재한다.
$T' \subset T$를 $a_{ij}$와 $b_{ijk}$로 생성되는
$R$-부분대수라 하자. 이는 유한 생성 $R$-대수이므로 뇌터 환이다.
대수
$$
S' = T'[Y_1, \ldots, Y_m]/(Y_i Y_j - \sum b_{ijk} Y_k).
$$
을 생각하자. $S'$는 $T'$ 위에서 유한이다.
실제로 $T'$-모듈로서 $1, Y_1, \ldots, Y_m$의 류들로 생성된다.
$Y_i$를 $y_i$로 보내는 $T'$-대수 준동형 $S' \to S$를 생각하자.
% P01-E0364
$a_{ij} \in T'$이므로 $x_j$는 이 사상의 상에 속한다.
따라서 $S' \to S$는 전사이다.
그러므로 $S$도 $T'$ 위에서 유한이다.
$T'$가 뇌터 환이므로 $T \subset S$가 $T'$ 위에서 유한임을 얻고,
증명이 끝난다.
\end{proof}




\section{길이}
\label{algebra-section-length}

\begin{definition}
\label{algebra-definition-length}
$R$을 환이라 하자. 임의의 $R$-모듈 $M$에 대해
$R$ 위에서 $M$의 {\it 길이}를 다음 공식으로 정의한다.
$$
\text{length}_R(M)
=
\sup
\{
n
\mid
\exists\ 0 = M_0 \subset M_1 \subset \ldots \subset M_n = M,
\text{ }M_i \not = M_{i + 1}
\}.
$$
\end{definition}

\noindent
다르게 말하면 부분모듈 사슬들의 길이의 상한이다.
한 부분모듈 사슬이 다른 사슬의 세분이라는 개념은 자명하다.
이에 따라 모든 부분모듈 사슬의 모음에는 부분순서가 주어진다.
$M$이 영이 아니면 가장 작은 사슬은
$0 = M_0 \subset M_1 = M$ 꼴이다.
$M$의 길이가 유한하면 모든 사슬을 극대 사슬로 세분할 수 있다는
명백한 사실에 유의하자. 그러나 모든 극대 사슬의 길이가 같다는
사실은 그만큼 명백하지 않다(뒤에서 보게 될 것이다).

\begin{lemma}
\label{algebra-lemma-finite-length-finite}
\begin{slogan}
유한 길이 모듈은 유한 모듈이다.
\end{slogan}
$R$을 환이라 하자.
$M$을 $R$-모듈이라 하자.
$\text{length}_R(M) < \infty$이면 $M$은 유한 $R$-모듈이다.
\end{lemma}

\begin{proof}
생략한다.
\end{proof}

\begin{lemma}
\label{algebra-lemma-length-additive}
\begin{slogan}
길이는 짧은 완전열에서 가법적이다.
\end{slogan}
$0 \to M' \to M \to M'' \to 0$이 $R$ 위의 모듈들의
짧은 완전열이면, $M$의 길이는 $M'$과 $M''$의 길이의 합이다.
\end{lemma}

\begin{proof}
$M'$과 $M''$의 길이가 각각 $n', n''$인 여과들이 주어지면,
이에 대응하는 길이 $n' + n''$의 $M$의 여과를 쉽게 만들 수 있다.
따라서
$\text{length}_R M \geq \text{length}_R M' + \text{length}_R M''$이다.
역으로 $M$의 여과
$M_0 \subset M_1 \subset \ldots \subset M_n$이 주어졌다고 하자.
유도된 여과 $M_i' = M_i \cap M'$와
$M_i'' = \Im(M_i \to M'')$을 생각하자.
$\{M'_i\}$의 단계 수를 $n'$, $\{M''_i\}$의 단계 수를 $n''$이라 하자.
$M_i' = M_{i + 1}'$이고 $M_i'' = M_{i + 1}''$이면
$M_i = M_{i + 1}$이다. 따라서 $n' + n'' \geq n$이다.
앞의 결과와 합치면 증명이 끝난다.
\end{proof}

\begin{lemma}
\label{algebra-lemma-length-infinite}
$R$을 극대 아이디얼 $\mathfrak m$을 갖는 국소환이라 하자.
$M$이 $R$-모듈이고 모든 $n \geq 0$에 대해
$\mathfrak m^n M \not = 0$이면 $\text{length}_R(M) = \infty$이다.
다르게 말하면 $M$의 길이가 유한하면 어떤 $n$에 대해
$\mathfrak m^nM = 0$이다.
\end{lemma}

\begin{proof}
모든 $n\geq 0$에 대해 $\mathfrak m^n M \not = 0$이라고 가정하자.
$f_1f_2 \ldots f_n x \not = 0$이 되도록
$x \in M$과 $f_1, \ldots, f_n \in \mathfrak m$을 택하자.
여과
$$
0 \subset R f_1 \ldots f_n x \subset R f_1 \ldots f_{n - 1} x
\subset \ldots \subset R x \subset M
$$
의 처음 $n$개 단계는 서로 다르다. 예를 들어
$R f_1 x = R f_1 f_2 x$이면 어떤 $g$에 대해
$f_1 x = g f_1 f_2 x$이고, 따라서 $(1 - gf_2) f_1 x = 0$이다.
$1 - gf_2$가 가역원이므로 $f_1 x = 0$인데, 이는
$x$와 $f_1, \ldots, f_n$의 선택에 모순이다.
따라서 길이는 무한이다.
\end{proof}

\begin{lemma}
\label{algebra-lemma-length-independent}
$R \to S$를 환 준동형이라 하자. $M$을 $S$-모듈이라 하자.
항상 $\text{length}_R(M) \geq \text{length}_S(M)$이다.
$R \to S$가 전사이면 등호가 성립한다.
\end{lemma}

\begin{proof}
% P01-E0366
$S$-부분모듈들에 의한 $M$의 여과는
$R$-부분모듈들에 의한 $M$의 여과를 준다. 이로써 부등식이 증명된다.
$R \to S$가 전사이면 $M$의 모든 $R$-부분모듈은 자동으로
$S$-부분모듈이다. 따라서 이 경우에는 등호가 성립한다.
\end{proof}

\begin{lemma}
\label{algebra-lemma-dimension-is-length}
$R$을 극대 아이디얼 $\mathfrak m$을 갖는 환이라 하자.
$M$이 $\mathfrak m M  =  0$을 만족하는 $R$-모듈이라고 하자.
그러면 $R$-모듈로서 $M$의 길이는
% P01-E1371
$R/\mathfrak m$-벡터공간으로서 $M$의 차원과 같다.
길이가 유한일 필요충분조건은 $M$이 유한 $R$-모듈인 것이다.
\end{lemma}

\begin{proof}
첫째 부분은 보조정리 \ref{algebra-lemma-length-independent}의 특수한 경우이다.
따라서 길이가 유한일 필요충분조건은 $M$이
$R/\mathfrak m$-벡터공간으로서 유한 기저를 갖는 것이고,
이는 다시 $M$이 $R$-모듈로서 유한 생성원 집합을 갖는 것과 동치이다.
\end{proof}

\begin{lemma}
\label{algebra-lemma-length-localize}
$R$을 환이라 하자. $M$을 $R$-모듈이라 하자.
$S \subset R$를 곱셈적 부분집합이라 하자. 그러면
$\text{length}_R(M) \geq \text{length}_{S^{-1}R}(S^{-1}M)$이다.
\end{lemma}

\begin{proof}
보조정리 \ref{algebra-lemma-submodule-localization}에 의해
모든 부분모듈 $N' \subset S^{-1}M$은 어떤 $R$-부분모듈
$N \subset M$에 대해 $S^{-1}N$ 꼴이다. 보조정리가 따른다.
\end{proof}

\begin{lemma}
\label{algebra-lemma-length-finite}
$R$을 유한 생성 극대 아이디얼 $\mathfrak m$을 갖는 환이라 하자.
% P01-E1372
(예를 들어 $R$이 뇌터 환인 경우이다.)
$M$이 유한 $R$-모듈이고 어떤 $n$에 대해
$\mathfrak m^n M  =  0$이라고 하자.
그러면 $\text{length}_R(M) < \infty$이다.
\end{lemma}

\begin{proof}
여과
$0 = \mathfrak m^n M \subset
\mathfrak m^{n-1} M \subset
\ldots \subset \mathfrak m M \subset M$을 생각하자.
모든 부분몫은 보조정리 \ref{algebra-lemma-dimension-is-length}가 적용되는
유한 생성 $R$-모듈이다. 가법성, 즉 보조정리
\ref{algebra-lemma-length-additive}에 의해 결론을 얻는다.
\end{proof}

\begin{definition}
\label{algebra-definition-simple-module}
$R$을 환이라 하자. $M$을 $R$-모듈이라 하자.
$M \not = 0$이고 $M$의 모든 부분모듈이 $M$ 또는 $0$이면
$M$을 {\it 단순}이라고 한다.
\end{definition}

\begin{lemma}
\label{algebra-lemma-characterize-length-1}
$R$을 환이라 하자. $M$을 $R$-모듈이라 하자.
다음은 서로 동치이다.
\begin{enumerate}
\item $M$은 단순이다.
\item $\text{length}_R(M) = 1$이다.
\item 어떤 극대 아이디얼 $\mathfrak m \subset R$에 대해
$M \cong R/\mathfrak m$이다.
\end{enumerate}
\end{lemma}

\begin{proof}
$\mathfrak m$을 $R$의 극대 아이디얼이라 하자.
보조정리 \ref{algebra-lemma-dimension-is-length}에 의해 모듈
$R/\mathfrak m$의 길이는 $1$이다.
첫째 주장과 둘째 주장의 동치는 정의에서 바로 따른다.
$M$이 단순이라고 하자. $x \in M$, $x \not = 0$을 택하자.
$M$이 단순이므로 $M = R \cdot x$이다.
$I \subset R$를 $x$의 소멸자, 즉
$I = \{f \in R \mid fx = 0\}$라 하자.
사상 $R/I \to M$, $f \bmod I \mapsto fx$는 동형사상이므로
$R/I$는 단순 $R$-모듈이다. $R/I \not = 0$이므로 $I \not = R$이다.
% P01-E0367
$I \subset \mathfrak m$이고 $I$를 포함하는 극대 아이디얼을 택하자.
$I \not = \mathfrak m$이면 $\mathfrak m /I \subset R/I$는
$R/I$의 단순성에 모순되는 자명하지 않은 부분모듈이다.
따라서 원하는 대로 $I = \mathfrak m$이다.
\end{proof}

\begin{lemma}
\label{algebra-lemma-simple-pieces}
$R$을 환이라 하자. $M$을 유한 길이 $R$-모듈이라 하자.
부분모듈들의 임의의 극대 사슬
$$
0 = M_0 \subset M_1 \subset M_2 \subset \ldots \subset M_n = M
$$
을 택하되 $i = 1, \ldots, n$에 대해
$M_i \not = M_{i-1}$이라고 하자. 그러면 다음이 성립한다.
\begin{enumerate}
\item $n = \text{length}_R(M)$이다.
\item 각 $M_i/M_{i-1}$는 단순이다.
\item 각 $M_i/M_{i-1}$는 어떤 극대 아이디얼 $\mathfrak m_i$에 대해
$R/\mathfrak m_i$ 꼴이다.
\item 극대 아이디얼 $\mathfrak m \subset R$가 주어지면
$$
\# \{i \mid \mathfrak m_i = \mathfrak m\}
=
\text{length}_{R_{\mathfrak m}} (M_{\mathfrak m}).
$$
\end{enumerate}
\end{lemma}

\begin{proof}
$M_i/M_{i-1}$가 단순이 아니면 여과를 더 세분할 수 있으므로
이 여과는 극대가 아니다. 따라서 $M_i/M_{i-1}$는 단순이다.
보조정리 \ref{algebra-lemma-characterize-length-1}에 의해 모듈
$M_i/M_{i-1}$의 길이는 $1$이고, 어떤 극대 아이디얼
$\mathfrak m_i$에 대해 $R/\mathfrak m_i$ 꼴이다.
길이의 가법성, 즉 보조정리 \ref{algebra-lemma-length-additive}에 의해
$n = \text{length}_R(M)$이다.
국소화는 완전하므로
$$
0 = (M_0)_{\mathfrak m}
\subset (M_1)_{\mathfrak m}
\subset (M_2)_{\mathfrak m}
\subset \ldots
\subset (M_n)_{\mathfrak m} = M_{\mathfrak m}
$$
은 연속된 몫들이 $(M_i/M_{i-1})_{\mathfrak m}$인
$M_{\mathfrak m}$의 여과이다.
따라서 마지막 주장은 $R$의 극대 아이디얼
$\mathfrak m$, $\mathfrak m'$에 대해
$$
(R/\mathfrak m')_{\mathfrak m}
\cong
\left\{
\begin{matrix}
0 &
\text{if } \mathfrak m \not = \mathfrak m', \\
R_{\mathfrak m}/\mathfrak m R_{\mathfrak m} &
\text{if } \mathfrak m  = \mathfrak m'
\end{matrix}
\right.
$$
임에서 바로 따른다. 이는 독자에게 맡긴다.
\end{proof}

\begin{lemma}
\label{algebra-lemma-pushdown-module}
$A$를 극대 아이디얼 $\mathfrak m$을 갖는 국소환이라 하자.
$B$를 극대 아이디얼들이 $\mathfrak m_i$, $i = 1, \ldots, n$인
반국소환이라 하자.
$A \to B$가 각 $\mathfrak m_i$가 $\mathfrak m$ 위에 놓이고
$$
[\kappa(\mathfrak m_i) : \kappa(\mathfrak m)] < \infty.
$$
를 만족하는 준동형이라고 하자.
$M$을 유한 길이 $B$-모듈이라 하자. 그러면
$$
\text{length}_A(M) = \sum\nolimits_{i = 1, \ldots, n}
[\kappa(\mathfrak m_i) : \kappa(\mathfrak m)]
\text{length}_{B_{\mathfrak m_i}}(M_{\mathfrak m_i}),
$$
이고, 특히 $\text{length}_A(M) < \infty$이다.
\end{lemma}

\begin{proof}
% P01-E1373
보조정리 \ref{algebra-lemma-simple-pieces}에서와 같이 $B$-부분모듈들의 극대 사슬
$$
0 = M_0
\subset M_1
\subset M_2
\subset \ldots
\subset M_m = M
$$
을 택하자. 그러면 각 몫 $M_j/M_{j - 1}$는 어떤
$i(j) \in \{1, \ldots, n\}$에 대해
$\kappa(\mathfrak m_{i(j)})$와 동형이다.
더욱이 보조정리 \ref{algebra-lemma-dimension-is-length}에 의해
$\text{length}_A(\kappa(\mathfrak m_i)) =
[\kappa(\mathfrak m_i) : \kappa(\mathfrak m)]$이다.
길이의 가법성(보조정리 \ref{algebra-lemma-length-additive})에서 결과가 따른다.
\end{proof}

\begin{lemma}
\label{algebra-lemma-pullback-module}
$A \to B$를 국소환들의 평탄 국소 준동형이라 하자.
그러면 임의의 $A$-모듈 $M$에 대해
$$
\text{length}_A(M) \text{length}_B(B/\mathfrak m_AB)
=
\text{length}_B(M \otimes_A B).
$$
이다. 특히 $\text{length}_B(B/\mathfrak m_AB) < \infty$이면,
$M$의 길이가 유한일 필요충분조건은
$M \otimes_A B$의 길이가 유한인 것이다.
\end{lemma}

\begin{proof}
보조정리 \ref{algebra-lemma-local-flat-ff}에 의해 환 준동형
$A \to B$는 충실히 평탄하다.
따라서 $0 = M_0 \subset M_1 \subset \ldots \subset M_n = M$이
$M$ 안의 길이 $n$인 사슬이면, 이에 대응하는 사슬
$0 = M_0 \otimes_A B \subset M_1 \otimes_A B \subset
\ldots \subset M_n \otimes_A B = M \otimes_A B$도 길이가 $n$이다.
이는
$\text{length}_A(M) = \infty \Rightarrow
\text{length}_B(M \otimes_A B) = \infty$를 증명한다.
이제 $\text{length}_A(M) < \infty$라고 가정하자.
이 경우 $M$은 길이가 $\ell = \text{length}_A(M)$이고
몫들이 $A/\mathfrak m_A$인 여과를 갖는다.
위와 같이 논증하면 $M \otimes_A B$는 길이가 $\ell$이고
몫들이 $B \otimes_A A/\mathfrak m_A =  B/\mathfrak m_AB$와
동형인 여과를 갖는다. 따라서 보조정리가 따른다.
\end{proof}

\begin{lemma}
\label{algebra-lemma-pullback-transitive}
$A \to B \to C$를 국소환들의 평탄 국소 준동형들이라 하자. 그러면
$$
\text{length}_B(B/\mathfrak m_A B)
\text{length}_C(C/\mathfrak m_B C)
=
\text{length}_C(C/\mathfrak m_A C)
$$
이다.
\end{lemma}

\begin{proof}
환 준동형 $B \to C$와 $B$-모듈
$M = B/\mathfrak m_A B$에 보조정리
\ref{algebra-lemma-pullback-module}를 적용하면 따른다.
\end{proof}




\section{아르틴 환}
\label{algebra-section-artinian}

\noindent
아르틴 환, 특히 아르틴 국소환은 대수기하에서 중요한 역할을 한다.
예를 들면 변형 이론에서 그러하다.

\begin{definition}
\label{algebra-definition-artinian}
환 $R$이 아이디얼들에 대한 하강 사슬 조건을 만족하면
이를 {\it 아르틴 환}이라고 한다.
\end{definition}

\begin{lemma}
\label{algebra-lemma-finite-dimensional-algebra}
$R$이 체 위의 유한 차원 대수라고 하자.
그러면 $R$은 아르틴 환이다.
\end{lemma}

\begin{proof}
아이디얼들에 대한 하강 사슬 조건은 명백히 성립한다.
\end{proof}

\begin{lemma}
\label{algebra-lemma-artinian-finite-nr-max}
$R$이 아르틴 환이면 $R$은 극대 아이디얼을 유한 개만 갖는다.
\end{lemma}

\begin{proof}
$\mathfrak m_i$, $i = 1, 2, 3, \ldots$가 서로 다른 극대 아이디얼들이라고 하자.
그러면
$\mathfrak m_1 \supset \mathfrak m_1\cap \mathfrak m_2
\supset \mathfrak m_1 \cap \mathfrak m_2 \cap \mathfrak m_3 \supset \ldots$
는 무한 하강 사슬이다
(중국인의 나머지 정리에 의해 모든 사상
$R \to \oplus_{i = 1}^n R/\mathfrak m_i$가 전사이기 때문이다).
\end{proof}

\begin{lemma}
\label{algebra-lemma-artinian-radical-nilpotent}
$R$을 아르틴 환이라 하자. $R$의 야코브슨 근기는 멱영 아이디얼이다.
\end{lemma}

\begin{proof}
$I \subset R$를 야코브슨 근기라 하자.
$I \supset I^2 \supset I^3 \supset \ldots$는 하강 사슬임에 유의하라.
따라서 어떤 $n$에 대해 $I^n = I^{n + 1}$이다.
$J = \{ x\in R \mid xI^n = 0\}$로 두자. $J = R$임을 보여야 한다.
그렇지 않다면 $J \subset J'$이고 $J' \not = J$인 아이디얼 가운데
극소인 것을 택하자(아르틴성에 의해 가능하다).
그러면 $J'/J$는 단순 $R$-모듈이므로 어떤 극대 아이디얼
$\mathfrak m$에 대해 $R/\mathfrak m$과 동형이다.
보조정리 \ref{algebra-lemma-characterize-length-1}을 참조하라.
그러면 $\mathfrak m I^n$은 $J'$를 영으로 만든다.
$I \subset \mathfrak m$이므로 $I^{n + 1} = I^n$도 $J'$를 영으로 만든다.
따라서 $J' = J$이고, 이는 모순이다.
\end{proof}

\begin{lemma}
\label{algebra-lemma-product-local}
극대 아이디얼이 유한 개이고 야코브슨 근기가 국소 멱영인 모든 환은
그 극대 아이디얼들에서의 국소화들의 곱이다.

또한 모든 소 아이디얼은 극대 아이디얼이다.
\end{lemma}

\begin{proof}
$R$을 극대 아이디얼들이
$\mathfrak m_1, \ldots, \mathfrak m_n$인 환이라 하자.
$I = \bigcap_{i = 1}^n \mathfrak m_i$를 $R$의 야코브슨 근기라 하자.
$I$가 국소 멱영이라고 가정하자.
$\mathfrak p$를 $R$의 소 아이디얼이라 하자.
모든 소 아이디얼은 $R$의 모든 멱영원을 포함하므로
$ \mathfrak p \supset \mathfrak m_1 \cap \ldots \cap \mathfrak m_n$이다.
또한
$\mathfrak m_1 \cap \ldots \cap \mathfrak m_n \supset
\mathfrak m_1 \ldots \mathfrak m_n$이므로
$\mathfrak p \supset \mathfrak m_1 \ldots \mathfrak m_n$이다.
따라서 어떤 $i$에 대해 $\mathfrak p \supset \mathfrak m_i$이고,
그러므로 $\mathfrak p = \mathfrak m_i$이다.
이에 따라 $R$의 스펙트럼은 이산 위상공간
$\{\mathfrak m_1, \ldots, \mathfrak m_n\}$이다.
보조정리 \ref{algebra-lemma-disjoint-implies-product}를 $n - 1$번 적용하면,
각 $i$에 대해 $R_i$의 스펙트럼이 한원소집합인
$R = R_1 \times \ldots \times R_n$을 얻는다.
따라서 $R_i$는 국소환이고, (그 보조정리에 의해) $R$의 국소화이므로
원하는 대로 $R$의 국소환들 가운데 하나이다.
\end{proof}

\begin{lemma}
\label{algebra-lemma-artinian-finite-length}
환 $R$이 아르틴 환일 필요충분조건은 자기 자신 위의 모듈로서
유한 길이를 갖는 것이다. 그러한 모든 환 $R$은 아르틴 환이자
뇌터 환이고, $R$의 모든 소 아이디얼은 극대 아이디얼이며,
$R$은 그 극대 아이디얼들에서의 국소화들의 (유한) 곱과 같다.
\end{lemma}

\begin{proof}
$R$이 자기 자신 위에서 유한 길이를 가지면 아이디얼들에 대한
상승 사슬 조건과 하강 사슬 조건을 모두 만족한다.
따라서 뇌터 환이자 아르틴 환이다.
% P01-E0368
모든 아르틴 환은 보조정리 \ref{algebra-lemma-artinian-finite-nr-max},
\ref{algebra-lemma-artinian-radical-nilpotent}, 그리고
\ref{algebra-lemma-product-local}에 의해 극대 아이디얼들에서의
국소화들의 곱과 같다.

\medskip\noindent
$R$이 아르틴 환이라고 하자. $R$이 자기 자신 위에서
유한 길이를 가짐을 보이겠다.
연속된 몫들이 유한 길이를 갖는 부분모듈 사슬 하나를 제시하면 충분하다.
앞에서 말한 바에 의해 $R$이 극대 아이디얼 $\mathfrak m$을 갖는
국소환이라고 가정해도 된다.
보조정리 \ref{algebra-lemma-artinian-radical-nilpotent}에 의해
어떤 $n$에 대해 $\mathfrak m^n =0$이다.
사슬
$0 = \mathfrak m^n \subset \mathfrak m^{n-1} \subset
\ldots \subset \mathfrak m \subset R$을 생각하자.
보조정리 \ref{algebra-lemma-dimension-is-length}에 의해 각 부분몫
$\mathfrak m^j/\mathfrak m^{j + 1}$의 길이는
$\kappa(\mathfrak m)$ 위의 벡터공간으로서 그 차원이다.
% P01-E0369
이 차원은 유한이어야 한다. 그렇지 않으면 벡터 부분공간들의
무한 하강 사슬을 얻게 되고, 이는 $R$의 아이디얼들의
무한 하강 사슬에 대응하기 때문이다.
\end{proof}




\section{본질적으로 유한형인 준동형}
\label{algebra-section-essentially-of-finite-type}

\noindent
유한형 환 준동형의 국소화에 관한 몇 가지 간단한 설명을 제시한다.

\begin{definition}
\label{algebra-definition-essentially-finite-p-t}
$R \to S$를 환 준동형이라 하자.
\begin{enumerate}
\item $S$가 유한형 $R$-대수의 국소화이면
$R \to S$가 {\it 본질적으로 유한형}이라고 한다.
\item $S$가 유한 표시 $R$-대수의 국소화이면
$R \to S$가 {\it 본질적으로 유한 표시}라고 한다.
\end{enumerate}
\end{definition}

\begin{lemma}
\label{algebra-lemma-composition-essentially-of-finite-type}
본질적으로 유한형인 환 준동형들의 클래스는 합성 아래에서 보존된다.
본질적으로 유한 표시인 경우도 마찬가지이다.
\end{lemma}

\begin{proof}
생략한다.
\end{proof}

\begin{lemma}
\label{algebra-lemma-base-change-essentially-of-finite-type}
본질적으로 유한형인 환 준동형들의 클래스는 밑변환 아래에서 보존된다.
본질적으로 유한 표시인 경우도 마찬가지이다.
\end{lemma}

\begin{proof}
생략한다.
\end{proof}

\begin{lemma}
\label{algebra-lemma-essentially-of-finite-type-into-artinian-local}
$R \to S$를 환 준동형이라 하자.
$S$가 극대 아이디얼 $\mathfrak m$을 갖는 아르틴 국소환이라고 가정하자.
그러면 다음이 성립한다.
\begin{enumerate}
\item $R \to S$가 유한일 필요충분조건은
$R \to S/\mathfrak m$이 유한인 것이다.
\item $R \to S$가 유한형일 필요충분조건은
$R \to S/\mathfrak m$이 유한형인 것이다.
\item $R \to S$가 본질적으로 유한형일 필요충분조건은
합성 $R \to S/\mathfrak m$이 본질적으로 유한형인 것이다.
\end{enumerate}
\end{lemma}

\begin{proof}
$R \to S$가 유한이면 보조정리 \ref{algebra-lemma-finite-transitive}에 의해
$R \to S/\mathfrak m$도 유한이다.
역으로 $R \to S/\mathfrak m$이 유한이라고 가정하자.
$S$는 자기 자신 위에서 유한 길이를 가지므로
(보조정리 \ref{algebra-lemma-artinian-finite-length}),
아이디얼들에 의한 여과
$$
0 \subset I_1 \subset \ldots \subset I_n = S
$$
를 택하여 $S$-모듈로서
$I_i/I_{i - 1} \cong S/\mathfrak m$이 되게 할 수 있다.
따라서 $S$는 각 연속된 몫이 유한 $R$-모듈인
$R$-부분모듈 $I_i$들에 의한 여과를 갖는다.
보조정리 \ref{algebra-lemma-extension}에 의해 $S$는 유한 $R$-모듈이다.

\medskip\noindent
$R \to S$가 유한형이면 보조정리
\ref{algebra-lemma-compose-finite-type}에 의해
$R \to S/\mathfrak m$도 유한형이다.
역으로 $R \to S/\mathfrak m$이 유한형이라고 가정하자.
$S/\mathfrak m$의 생성원들로 사상되는
$f_1, \ldots, f_n \in S$를 택하자.
그러면 $A = R[x_1, \ldots, x_n] \to S$, $x_i \mapsto f_i$는
$A \to S/\mathfrak m$이 전사(특히 유한)가 되는 환 준동형이다.
따라서 (1)에 의해 $A \to S$는 유한이고,
보조정리 \ref{algebra-lemma-compose-finite-type}에 의해
$R \to S$가 유한형임을 알 수 있다.

\medskip\noindent
$R \to S$가 본질적으로 유한형이면 보조정리
\ref{algebra-lemma-composition-essentially-of-finite-type}에 의해
$R \to S/\mathfrak m$도 본질적으로 유한형이다.
역으로 $R \to S/\mathfrak m$이 본질적으로 유한형이라고 가정하자.
$S/\mathfrak m$이 $R[x_1, \ldots, x_n]/I$의 국소화라고 하자.
$\mathfrak m$을 법으로 한 합동류가 $I$를 법으로 한
$x_1, \ldots, x_n$의 합동류에 대응하는
$f_1, \ldots, f_n \in S$를 택하자.
핵이 $J$인 사상
$R[x_1, \ldots, x_n] \to S$, $x_i \mapsto f_i$를 생각하자.
$A = R[x_1, \ldots, x_n]/J \subset S$와
$\mathfrak p = A \cap \mathfrak m$으로 두자.
$A/\mathfrak p \subset S/\mathfrak m$은
$S/\mathfrak m$ 안에서 $R[x_1, \ldots, x_n]/I$의 상과 같음에 유의하라.
따라서 $\kappa(\mathfrak p) = S/\mathfrak m$이다.
그러므로 (1)에 의해 $A_\mathfrak p \to S$는 유한이다.
보조정리 \ref{algebra-lemma-composition-essentially-of-finite-type}에 의해
$S$가 본질적으로 유한형이라는 결론을 얻는다.
\end{proof}

\noindent
다음 보조정리는 여기서 제시하는 대수적 논증 대신
사영공간의 고유성을 사용하여 증명할 수도 있다.

\begin{lemma}
\label{algebra-lemma-localization-at-closed-point-special-fibre}
$\varphi : R \to S$를 본질적으로 유한형인 준동형이라 하고,
$R$과 $S$가 국소환이라고 하자
(단, $\varphi$가 국소 준동형일 필요는 없다).
그러면 어떤 $n$과 $\mathfrak m_R$ 위에 놓이는 극대 아이디얼
$\mathfrak m \subset R[x_1, \ldots, x_n]$이 존재하여,
$S$는 $R[x_1, \ldots, x_n]_\mathfrak m$의 한 몫의 국소화가 된다.
\end{lemma}

\begin{proof}
$S$를 $R[x_1, \ldots, x_n]$의 한 몫의 국소화로 쓸 수 있다.
따라서 어떤 소 아이디얼
$\mathfrak q \subset R[x_1, \ldots, x_n]$에 대해
$S = R[x_1, \ldots, x_n]_\mathfrak q$인 경우를 증명하면 충분하다.
$\mathfrak q + \mathfrak m_R R[x_1, \ldots, x_n] \not = R[x_1, \ldots, x_n]$
이면, 보조정리의 명제에 나오는 것처럼
$\mathfrak q \subset \mathfrak m$인 극대 아이디얼
$\mathfrak m$을 찾을 수 있고 결과는 명백하다.

\medskip\noindent
$R \to \kappa(\mathfrak q)$의 상을 지배하는 값매김환
$A \subset \kappa(\mathfrak q)$를 택하자
(보조정리 \ref{algebra-lemma-dominate}).
% P01-E0370
$x_i$의 상 $\lambda_i \in \kappa(\mathfrak q)$가 모두
$A$에 포함되면, $\mathfrak q$는
$R[x_1, \ldots, x_n] \to A$ 아래에서 $\mathfrak m_A$의
역상에 포함되므로 앞의 경우로 돌아간다.
따라서 어떤 $i$가 존재하여 $\lambda_i^{-1} \in A$이고,
모든 $j = 1, \ldots, n$에 대해 $\lambda_j/\lambda_i \in A$이다
($A$의 값군이 전순서이기 때문이다. 보조정리
\ref{algebra-lemma-valuation-group}를 참조하라).
이제 사상
$$
R[y_0, y_1, \ldots, \hat{y_i}, \ldots, y_n]
\to R[x_1, \ldots, x_n]_\mathfrak q,\quad
y_0 \mapsto 1/x_i,\quad y_j \mapsto x_j/x_i
$$
을 생각하자. $\mathfrak q' \subset R[y_0, \ldots, \hat{y_i}, \ldots, y_n]$
를 $\mathfrak q$의 역상이라 하자.
$y_0 \not \in \mathfrak q'$이므로 위 사상이 국소화들 사이의
동형사상을 정의함을 쉽게 알 수 있다.
한편 $y_j$가 $A$의 원소로 사상되므로
$R[y_0, \ldots, \hat{y_i}, \ldots, y_n]$에는
첫째 문단의 결과를 적용할 수 있다.
이로써 증명이 끝난다.
\end{proof}
 
\section{K-군}
\label{algebra-section-K-groups}

\noindent
$R$을 환이라 하자. $R$에 연관된 두 아벨군을 도입하겠다.
첫째 군은 $K'_0(R)$로 나타내며 다음 성질을 갖는다%
\footnote{이 정의는 임의의 환에 대해 의미가 있지만,
$R$이 뇌터 환이 아니면 거의 사용되지 않는다.}.
\begin{enumerate}
\item 모든 유한 $R$-모듈 $M$마다 $K'_0(R)$의 원소 $[M]$이 주어진다.
\item 유한 $R$-모듈들의 모든 짧은 완전열
$0 \to M' \to M \to M'' \to 0$에 대해 관계
$[M] = [M'] + [M'']$이 성립한다.
\item 군 $K'_0(R)$은 원소들 $[M]$로 생성된다.
\item 생성원들 $[M]$ 사이에서 $K'_0(R)$에 성립하는 모든 관계는
위와 같은 완전열에서 오는 관계들의 $\mathbf{Z}$-선형결합이다.
\end{enumerate}
실제 구성은 유한 생성 $R$-모듈 전체의 모음이 진클래스라는 점을
처리해야 하므로 조금 더 번거롭다.
그러나 군 $K'_0(R)$의 생성원 집합으로, $n$이 모든 정수를 달리고
% P01-E0371
$K$가 모든 부분모듈 $K \subset R^n$을 달릴 때의 원소
$[R^n/K]$를 취하면 이 문제를 극복할 수 있다.
이 원소들에 부과되는 관계 부분군의 생성원은 항들이
$R^n/K$ 꼴인 짧은 완전열들에서 오는 관계들이다.
원소 $[M]$은 $M \cong R^n/K$가 되도록 $n$과 $K$를 택하고
$[M] = [R^n/K]$로 두어 정의한다. 세부사항은 독자에게 맡긴다.

\begin{lemma}
\label{algebra-lemma-length-K0}
$R$이 아르틴 국소환이면 길이 함수는 자연스러운 아벨군 준동형
$\text{length}_R : K'_0(R) \to \mathbf{Z}$를 정의한다.
\end{lemma}

\begin{proof}
모든 유한 $R$-모듈의 길이는 유한하다.
실제로 이는 보조정리 \ref{algebra-lemma-artinian-finite-length}에 의해
유한 길이를 갖는 $R^n$의 몫이다.
또한 길이 함수는 가법적이다. 보조정리
\ref{algebra-lemma-length-additive}를 참조하라.
\end{proof}

\noindent
둘째 군은 $K_0(R)$로 나타내며 다음 성질을 갖는다.
\begin{enumerate}
\item 모든 유한 사영 $R$-모듈 $M$마다
$K_0(R)$의 원소 $[M]$이 주어진다.
\item 유한 사영 $R$-모듈들의 모든 짧은 완전열
$0 \to M' \to M \to M'' \to 0$에 대해
관계 $[M] = [M'] + [M'']$이 성립한다.
\item 군 $K_0(R)$은 원소들 $[M]$로 생성된다.
\item $K_0(R)$의 모든 관계는 위와 같은 완전열에서 오는 관계들의
$\mathbf{Z}$-선형결합이다.
\end{enumerate}
이 군도 위와 같이 구성한다.

\medskip\noindent
일반적으로 동형사상은 아닌 명백한 사상
$K_0(R) \to K'_0(R)$가 있음에 유의하라.

\begin{example}
\label{algebra-example-K0-field}
$R = k$가 체이면 차원 함수(동시에 길이 함수이기도 하다)가 주는
동형사상 아래에서
$K_0(k) = K'_0(k) \cong \mathbf{Z}$임이 명백하다.
\end{example}

\begin{example}
\label{algebra-example-K0-PID}
$R$을 주아이디얼 정역이라 하자.
$K_0(R) = K'_0(R) = \mathbf{Z}$라고 주장한다.
실제로 모든 유한 사영 $R$-모듈은 유한 자유이다.
보조정리 \ref{algebra-lemma-rank}에 의해 유한 자유 모듈의 랭크는 잘 정의된다.
유한 자유 모듈들의 짧은 완전열
$$
0 \to M' \to M \to M'' \to 0
$$
이 주어지면
$\text{rank}(M) = \text{rank}(M') + \text{rank}(M'')$이다.
% P01-E0372
이 경우 $M \cong M' \oplus M'$이기 때문이다
(예를 들어 보조정리 \ref{algebra-lemma-lift-map}에 의해 분할이 존재한다).
따라서 $K_0(R) = \mathbf{Z}$이다.

\medskip\noindent
주아이디얼 정역 위의 모듈에 대한 구조정리에 따르면,
모든 유한 생성 $R$-모듈은
$M = R^{\oplus r} \oplus R/(d_1) \oplus \ldots \oplus R/(d_k)$ 꼴이다.
짧은 완전열
$$
0 \to (d_i) \to R \to R/(d_i) \to 0
$$
을 생각하자. 아이디얼 $(d_i)$는 모듈로서 $R$과 동형이므로
(생성원 $d_i$를 갖는 자유 모듈이다),
$K'_0(R)$에서 $[(d_i)] = [R]$이다.
% P01-E1374
그러면 $[R/(d_i)] = [(d_i)]-[R] = 0$이다.
따라서 꼬임 모듈의 $K'_0(R)$에서의 류는 영이다.
자유 부분의 랭크를 사용하면 식별 $K'_0(R) = \mathbf{Z}$를 얻고,
표준 준동형 $K_0(R) \to K'_0(R)$는 동형사상이다.
\end{example}

\begin{example}
\label{algebra-example-K0-polynomial-ring}
$k$를 체라 하자. 그러면
$K_0(k[x]) = K'_0(k[x]) = \mathbf{Z}$이다.
$R = k[x]$가 주아이디얼 정역이므로 이는 예
\ref{algebra-example-K0-PID}에서 따른다.
\end{example}

\begin{example}
\label{algebra-example-K0-node}
$k$를 체라 하자.
$R = \{f \in k[x] \mid f(0) = f(1)\}$로 두자.
예 \ref{algebra-example-affine-open-not-standard}와 비교하라.
이 경우 $K_0(R) \cong k^* \oplus \mathbf{Z}$이지만
$K'_0(R) = \mathbf{Z}$이다.
\end{example}

\begin{lemma}
\label{algebra-lemma-K0-product}
$R = R_1 \times R_2$라 하자. 그러면
$K_0(R) = K_0(R_1) \times K_0(R_2)$이고
$K'_0(R) = K'_0(R_1) \times K'_0(R_2)$이다.
\end{lemma}

\begin{proof}
생략한다.
\end{proof}

\begin{lemma}
\label{algebra-lemma-K0prime-Artinian}
$R$을 아르틴 국소환이라 하자.
보조정리 \ref{algebra-lemma-length-K0}의 사상
$\text{length}_R : K'_0(R) \to \mathbf{Z}$는 동형사상이다.
\end{lemma}

\begin{proof}
생략한다.
\end{proof}

\begin{lemma}
\label{algebra-lemma-K0-local}
$(R, \mathfrak m)$을 국소환이라 하자.
모든 유한 사영 $R$-모듈은 유한 자유이다.
$[M] \to \text{rank}_R(M)$으로 정의되는 사상
$\text{rank}_R : K_0(R) \to \mathbf{Z}$는 잘 정의되며 동형사상이다.
\end{lemma}

\begin{proof}
$P$를 유한 사영 $R$-모듈이라 하자.
$P/\mathfrak m P$의 기저로 사상되는 원소들
$x_1, \ldots, x_n \in P$를 택하자.
나카야마의 보조정리 \ref{algebra-lemma-NAK}에 의해 이 원소들은 $P$를 생성한다.
이에 대응하는 전사 $u : R^{\oplus n} \to P$는
$P$가 사영이므로 분할을 갖는다.
따라서 $Q = \Ker(u)$에 대해 $R^{\oplus n} = P \oplus Q$이다.
그러므로 $Q/\mathfrak m Q = 0$이고, 나카야마의 보조정리에 의해
$Q$는 영이다. 이로써 모든 유한 사영 $R$-모듈이 유한 자유임을 알 수 있다.
보조정리 \ref{algebra-lemma-rank}에 의해 유한 자유 모듈의 랭크는 잘 정의된다.
유한 자유 $R$-모듈들의 짧은 완전열
$$
0 \to M' \to M \to M'' \to 0
$$
이 주어지면
$\text{rank}(M) = \text{rank}(M') + \text{rank}(M'')$이다.
% P01-E0373
이 경우 $M \cong M' \oplus M'$이기 때문이다
(예를 들어 보조정리 \ref{algebra-lemma-lift-map}에 의해 분할이 존재한다).
따라서 $K_0(R) = \mathbf{Z}$이다.
\end{proof}

\begin{lemma}
\label{algebra-lemma-K0-and-K0prime-Artinian-local}
$R$을 아르틴 국소환이라 하자. 다음 가환 그림이 존재한다.
$$
\xymatrix{
K_0(R) \ar[rr] \ar[d]_{\text{rank}_R} & &
K'_0(R) \ar[d]^{\text{length}_R} \\
\mathbf{Z} \ar[rr]^{\text{length}_R(R)} & &
\mathbf{Z}
}
$$
여기서 수직 사상들은 보조정리
\ref{algebra-lemma-K0prime-Artinian}과 \ref{algebra-lemma-K0-local}에 의해 동형사상이다.
\end{lemma}

\begin{proof}
$P$를 유한 사영 $R$-모듈이라 하자.
$\text{length}_R(P) = \text{rank}_R(P) \text{length}_R(R)$임을 보여야 한다.
보조정리 \ref{algebra-lemma-K0-local}에 의해 모듈 $P$는 유한 자유이다.
따라서 어떤 $n \geq 0$에 대해 $P \cong R^{\oplus n}$이다.
그러면 $\text{rank}_R(P) = n$이고, 길이의 가법성
(보조정리 \ref{algebra-lemma-length-additive})에 의해
$\text{length}_R(R^{\oplus n}) = n \text{length}_R(R)$이다.
따라서 결과가 성립한다.
\end{proof}




\section{등급환}
\label{algebra-section-graded}

\noindent
여기서 {\it 등급환}이란 기저 아벨군의 직합 분해
$S = \bigoplus_{d \geq 0} S_d$가 주어진 환 $S$로서
$S_d \cdot S_e \subset S_{d + e}$를 만족하는 것을 말한다.
음의 차수의 영이 아닌 원소는 허용하지 않는다는 점에 유의하라.
{\it 무관 아이디얼}은 아이디얼
$S_{+} = \bigoplus_{d > 0} S_d$이다.
{\it 등급 모듈}이란 기저 아벨군의 직합 분해
$M = \bigoplus_{n\in \mathbf{Z}} M_n$가 주어진 $S$-모듈 $M$로서
$S_d \cdot M_e \subset M_{d + e}$를 만족하는 것을 말한다.
모듈에서는 음의 차수의 영이 아닌 원소를 허용한다는 점에 유의하라.
$k > 0$에 대해 $S_{-k} = (0)$으로 두어 $S$를 등급 $S$-모듈로 생각한다.
$M$의 원소 $x$ (각각 $S$의 원소 $f$)가 어떤 $d$에 대해
$x \in M_d$ (각각 $f \in S_d$)이면 이를 {\it 동차}라고 한다.
{\it 등급 $S$-모듈의 사상}이란
$\varphi(M_d) \subset M'_d$를 만족하는 $S$-모듈 사상
$\varphi : M \to M'$을 말한다.
차수를 이동시키는 사상은 허용하지 않는다.
$M$에서 $N$으로 가는 등급 모듈 준동형들의 $S_0$-모듈을
$\text{GrHom}_0(M, N)$로 나타내자.

\medskip\noindent
이제 등급 아이디얼, 등급 몫환, 등급 부분모듈, 등급 몫모듈,
등급 텐서곱 등의 개념이 있다.
% P01-E1375
관련 정의와 보조정리를 찾는 일은 독자에게 맡긴다.
예를 들어 등급 모듈들의 짧은 완전열은 각 차수에서 짧은 완전열이다.

\medskip\noindent
등급환 $S$, 등급 $S$-모듈 $M$, 그리고
$n \in \mathbf{Z}$가 주어지면,
$M(n)_d = M_{n + d}$인 등급 $S$-모듈을 $M(n)$으로 나타낸다.
이를 {\it $n$에 의한 $M$의 트위스트}라고 한다.
특히 사영 스킴의 연구에서 중요한 역할을 할 모듈들
$S(n)$, $n \in \mathbf{Z}$를 얻는다.
다음과 같은 자명한 함자적 동형사상들이 있다.
$(M \oplus N)(n) = M(n) \oplus N(n)$,
$(M \otimes_S N)(n) = M \otimes_S N(n) = M(n) \otimes_S N$.
또한 $S_0$-모듈
$$
\text{GrHom}(M, N) =
\bigoplus\nolimits_{n \in \mathbf{Z}} \text{GrHom}_n(M, N),
\quad
\text{GrHom}_n(M, N) = \text{GrHom}_0(M, N(n)).
$$
에 등급 $S$-모듈 구조를 정의할 수 있다.
% P01-E0374
곱셈의 정의는 생략한다.

\begin{lemma}
\label{algebra-lemma-graded-NAK}
$S$를 등급환이라 하자. $M$을 등급 $S$-모듈이라 하자.
\begin{enumerate}
\item $S_+M = M$이고 $M$이 유한이면 $M = 0$이다.
\item $N, N' \subset M$이 등급 부분모듈들이고,
$M = N + S_+N'$이며 $N'$이 유한이면 $M = N$이다.
\item $N \to M$이 등급 모듈의 사상이고,
$N/S_+N \to M/S_+M$이 전사이며 $M$이 유한이면
$N \to M$은 전사이다.
\item $x_1, \ldots, x_n \in M$이 동차이고 $M/S_+M$을 생성하며
$M$이 유한이면 $x_1, \ldots, x_n$은 $M$을 생성한다.
\end{enumerate}
\end{lemma}

\begin{proof}
(1)의 증명. $S$ 위에서 $M$의 생성원들을
$y_1, \ldots, y_r$라 하자.
$y_i$가 차수 $d_i$의 동차 원소라고 가정해도 된다.
번호를 다시 매겨 $d_r = \min(d_i)$라고 가정해도 된다.
그러면 조건 $S_+M = M$은 $y_r = 0$을 함의한다.
$r$에 관한 귀납법에 의해 $M = 0$이다.
(2)는 $M/N$에 (1)을 적용하면 따른다.
(3)은 부분모듈들 $\Im(N \to M)$과 $M$에 (2)를 적용하면 따른다.
(4)는 모듈 사상
$\bigoplus S(-d_i) \to M$, $(s_1, \ldots, s_n) \mapsto \sum s_i x_i$
에 (3)을 적용하면 따른다.
\end{proof}

\noindent
$S$를 등급환이라 하자. $d \geq 1$을 정수라 하자.
$S^{(d)} = \bigoplus_{n \geq 0} S_{nd}$로 두자.
$S^{(d)}$를 차수 $n$ 성분이
$(S^{(d)})_n = S_{nd}$인 등급환으로 생각한다.
등급 $S$-모듈 $M$이 주어지면 마찬가지로
$M^{(d)} = \bigoplus_{n \in \mathbf{Z}} M_{nd}$를 생각할 수 있고,
이는 등급 $S^{(d)}$-모듈이다.

\begin{lemma}
\label{algebra-lemma-uple-generated-degree-1}
$S$를 $S_0$ 위에서 유한 생성인 등급환이라 하자.
그러면 충분히 가분적인 모든 $d$에 대해
대수 $S^{(d)}$는 $S_0$ 위에서 차수 $1$의 원소들로 생성된다.
\end{lemma}

\begin{proof}
$S$가 $S_0$ 위에서 $f_1, \ldots, f_r \in S$로 생성된다고 하자.
$f_i$를 그 동차 성분들로 대체하여, $f_i$가 차수
$d_i > 0$의 동차 원소라고 가정해도 된다.
그러면 $S_n$의 모든 원소는
$\sum e_i d_i = n$을 만족하는 단항식
$f_1^{e_1} \ldots f_r^{e_r}$들의 $S_0$ 계수 선형결합이다.
$m$을 $\text{lcm}(d_i)$의 배수라 하자.
임의의 $N \geq r$에 대해
$$
\sum e_i d_i = N m
$$
이면 초등적인 논증에 의해 어떤 $i$에 대해 $e_i \geq m/d_i$이다.
따라서 차수 $N m$의 모든 단항식은 차수 $m$의 단항식
$f_i^{m/d_i}$와 차수 $(N - 1)m$의 단항식의 곱이다.
그러므로 $n \geq 2$일 때 차수 $nrm$의 모든 단항식은
차수 $rm$인 단항식들의 곱이다.
따라서 $S^{(rm)}$은 $S_0$ 위에서 차수 $1$의 원소들로 생성된다.
\end{proof}

\begin{lemma}
\label{algebra-lemma-integral-closure-graded}
\begin{reference}
\cite[Theorem 2.3.2]{Huneke-Swanson}
\end{reference}
$R \to S$를 등급환들의 준동형이라 하자.
$S' \subset S$를 $S$ 안에서 $R$의 정수적 폐포라 하자.
그러면
$$
S' = \bigoplus\nolimits_{d \geq 0} S' \cap S_d,
$$
이다. 즉 $S'$는 $S$의 등급 $R$-부분대수이다.
\end{lemma}

\begin{proof}
다음을 보여야 한다.
$s = s_n + s_{n + 1} + \ldots + s_m \in S'$이면
각 동차 성분 $s_j \in S'$이다.
등급환들의 모든 준동형 $R \to S$에 대해
$m - n$에 관한 귀납법으로 이를 증명하겠다.
먼저 $R \to R_0$와 호환되는 환 준동형 $S \to S_0$가 있으므로,
$s_0$이 $R_0$ 위에서 정수적이고 따라서 $R$ 위에서도
정수적인 것은 바로 알 수 있다.
그러므로 $s$를 $s - s_0$으로 대체하여 $n > 0$이라고 가정해도 된다.
$t$의 차수가 $0$인 등급환의 확대
$R[t, t^{-1}] \to S[t, t^{-1}]$을 생각하자.
다음 가환 그림이 있다.
$$
\xymatrix{
S[t, t^{-1}] \ar[rr]_{s \mapsto t^{\deg(s)}s} & & S[t, t^{-1}] \\
R[t, t^{-1}] \ar[u] \ar[rr]^{r \mapsto t^{\deg(r)}r} & &  R[t, t^{-1}] \ar[u]
}
$$
여기서 수평 사상들은 환의 자기동형사상이다.
% P01-E1376
따라서 $R[t, t^{-1}]$ 위에서 $S[t, t^{-1}]$의 정수적 폐포 $C$는
자기 자신 안으로 사상된다. 그러므로
$$
t^m(s_n + s_{n + 1} + \ldots + s_m) -
(t^ns_n + t^{n + 1}s_{n + 1} + \ldots + t^ms_m) \in C
$$
이고, 귀납가정에 의해 $i = n, \ldots, m - 1$에 대해
각 $(t^m - t^i)s_i \in C$이다.
임의의 환 $A$와 $m > i \geq n > 0$에 대해
$A[t, t^{-1}]/(t^m - t^i - 1) \cong A[t]/(t^m - t^i - 1) \supset A$임에
유의하라. 실제로 $A[t]/(t^m - t^i - 1)$에서
$t(t^{m - 1} - t^{i - 1}) = 1$이다.
$t^m - t^i$가 $1$로 사상되므로,
$i = n, \ldots, m - 1$에 대해 환
$S[t]/(t^m - t^i - 1)$에서의 $s_i$의 상은
$R[t]/(t^m - t^i - 1)$ 위에서 정수적이다.
$R \to R[t]/(t^m - t^i - 1)$이 유한이므로
추이성에 의해 $s_i$는 $R$ 위에서 정수적이다.
보조정리 \ref{algebra-lemma-integral-transitive}를 참조하라.
마지막으로
$s_m = s - \sum_{i = n, \ldots, m - 1} s_i$도
$R$ 위에서 정수적임을 얻는다.
\end{proof}




\section{등급환의 Proj}
\label{algebra-section-proj}

\noindent
$S$를 등급환이라 하자.
{\it 동차 아이디얼}이란 $S$의 등급 부분모듈이기도 한
아이디얼 $I \subset S$를 말한다.
동치로, 동차 원소들로 생성되는 아이디얼이다.
또 동치로, $f \in I$이고
$$
f = f_0 + f_1 + \ldots + f_n
$$
가 $S$에서 $f$를 동차 성분들로 분해한 것이면
각 $i$에 대해 $f_i \in I$이다.
동차 아이디얼 $\mathfrak p$가 소 아이디얼인지 확인하려면,
$a, b$가 동차이고 $ab \in \mathfrak p$일 때
$a \in \mathfrak p$이거나 $b \in \mathfrak p$임을 확인하면 충분하다.

\begin{definition}
\label{algebra-definition-proj}
$S$를 등급환이라 하자.
$S_{+} \not \subset \mathfrak p$를 만족하는 $S$의 동차 소 아이디얼
$\mathfrak p$들의 집합을 $\text{Proj}(S)$로 정의한다.
집합 $\text{Proj}(S)$은 $\Spec(S)$의 부분집합이며,
이에 유도 위상을 준다.
위상공간 $\text{Proj}(S)$을 등급환 $S$의
{\it 동차 스펙트럼}이라고 한다.
\end{definition}

\noindent
구성에 의해 연속 사상
$$
\text{Proj}(S) \longrightarrow \Spec(S_0).
$$
이 있음에 유의하라.

\medskip\noindent
$S = \oplus_{d \geq 0} S_d$를 등급환이라 하자.
$f\in S_d$이고 $d \geq 1$이라고 가정하자.
$S_{(f)}$를 $S_f$의 부분환으로 정의하되,
그 원소는 $r$이 동차이고 $\deg(r) = nd$인 $r/f^n$ 꼴의 원소로 한다.
$M$이 등급 $S$-모듈이면, $S_{(f)}$-모듈 $M_{(f)}$를
% P01-E1377
$M_f$의 부분모듈로 정의하되 그 원소는 $x$가 차수 $nd$인
동차 원소인 $x/f^n$ 꼴의 원소로 한다.

\begin{lemma}
\label{algebra-lemma-Z-graded}
$S$를 양의 차수인 동차 가역원을 포함하는
$\mathbf{Z}$-등급환이라 하자.
$S$의 $\mathbf{Z}$-등급 소 아이디얼들의 집합
$G \subset \Spec(S)$는 (유도 위상 아래에서)
$\Spec(S_0)$와 위상동형으로 사상된다.
\end{lemma}

\begin{proof}
먼저 역을 구성하여 이 사상이 전단사임을 보이자.
$f \in S_d$, $d > 0$을 $S$의 가역원이라 하자.
$\mathfrak p_0$가 $S_0$의 소 아이디얼이면,
$\mathfrak p_0S$는 $S$의 $\mathbf{Z}$-등급 아이디얼이고
$\mathfrak p_0S \cap S_0 = \mathfrak p_0$이다.
또한 $a$, $b$가 동차이고 $ab \in \mathfrak p_0S$이면
$a^db^d/f^{\deg(a) + \deg(b)} \in \mathfrak p_0$이다.
따라서 $a^d/f^{\deg(a)} \in \mathfrak p_0$이거나
$b^d/f^{\deg(b)} \in \mathfrak p_0$이다.
다르게 말하면 $a^d \in \mathfrak p_0S$이거나
$b^d \in \mathfrak p_0S$이다.
그러므로 $\sqrt{\mathfrak p_0S}$는 $S$의 $\mathbf{Z}$-등급
소 아이디얼이고 $S_0$와의 교집합은 $\mathfrak p_0$이다.

\medskip\noindent
사상이 위상동형임을 보이기 위해 $G \cap D(g)$의 상이 열림을 보이자.
$g = \sum g_i$이고 $g_i \in S_i$이면, 위의 결과에 의해
$G \cap D(g)$는 열린집합 $\bigcup D(g_i^d/f^i)$ 위로 사상된다.
\end{proof}

\noindent
차수가 $> 0$인 동차 원소 $f \in S$에 대해
$$
D_{+}(f) = \{ \mathfrak p \in \text{Proj}(S) \mid f \not\in \mathfrak p \}.
$$
로 정의한다. 마지막으로 동차 아이디얼 $I \subset S$에 대해
$$
V_{+}(I) = \{ \mathfrak p \in \text{Proj}(S) \mid I \subset \mathfrak p \}.
$$
로 정의한다. 더 일반적으로 동차 원소들의 임의의 집합 $E$에 대해,
$E \subset S$이면 표기 $V_{+}(E)$를 사용하겠다.

\begin{lemma}[Proj의 위상]
\label{algebra-lemma-topology-proj}
$S = \oplus_{d \geq 0} S_d$를 등급환이라 하자.
\begin{enumerate}
\item 집합들 $D_{+}(f)$는 $\text{Proj}(S)$에서 열린집합이다.
\item $D_{+}(ff') = D_{+}(f) \cap D_{+}(f')$이다.
\item $g = g_0 + \ldots + g_m$을 $g_i \in S_i$인 $S$의 원소라 하자.
그러면
$$
D(g) \cap \text{Proj}(S) =
(D(g_0) \cap \text{Proj}(S))
\cup
\bigcup\nolimits_{i \geq 1} D_{+}(g_i).
$$
\item 차수 $0$의 동차 원소 $g_0\in S_0$에 대해
$$
D(g_0) \cap \text{Proj}(S)
=
\bigcup\nolimits_{f \in S_d, \ d\geq 1} D_{+}(g_0 f).
$$
\item 열린집합들 $D_{+}(f)$는 $\text{Proj}(S)$의 위상의 기저를 이룬다.
\item $f \in S$를 양의 차수인 동차 원소라 하자.
환 $S_f$는 자연스러운 $\mathbf{Z}$-등급을 갖는다.
환 준동형들 $S \to S_f \leftarrow S_{(f)}$는 위상동형사상들
$$
D_{+}(f)
\leftarrow
\{\mathbf{Z}\text{-graded primes of }S_f\}
\to
\Spec(S_{(f)}).
$$
을 유도한다.
\item $\text{Proj}(S)$가 준콤팩트가 아닌 $S$가 존재한다.
\item 집합들 $V_{+}(I)$는 닫힌집합이다.
\item 임의의 닫힌 부분집합 $T \subset \text{Proj}(S)$은
어떤 동차 아이디얼 $I \subset S$에 대해 $V_{+}(I)$ 꼴이다.
\item 임의의 등급 아이디얼 $I \subset S$에 대해
$V_{+}(I) = \emptyset$일 필요충분조건은
$S_{+} \subset \sqrt{I}$인 것이다.
\end{enumerate}
\end{lemma}

\begin{proof}
$D_{+}(f) = \text{Proj}(S) \cap D(f)$이므로 이 집합들은 열린집합이다.
이로써 (1)이 증명된다.
또한 $D(ff') = D(f) \cap D(f')$이므로 (2)가 따른다.
마찬가지로 집합들
$V_{+}(I) = \text{Proj}(S) \cap V(I)$는 닫힌집합이다.
이로써 (8)이 증명된다.

\medskip\noindent
$T \subset \text{Proj}(S)$가 닫혔다고 하자.
그러면 어떤 아이디얼 $J \subset S$에 대해
$T = \text{Proj}(S) \cap V(J)$로 쓸 수 있다.
동차 아이디얼의 정의에 의해, $g \in J$이고
$g = g_0 + \ldots + g_m$이며 $g_d \in S_d$이면
모든 $\mathfrak p \in T$에 대해 $g_d \in \mathfrak p$이다.
$I \subset S$를 $J$의 원소들의 동차 성분들로 생성되는
아이디얼이라 두면 $T = V_{+}(I)$이다.
이로써 (9)가 증명된다.

\medskip\noindent
$g \in S$일 때 $\text{Proj}(S) \cap D(g)$에 대한 공식은
정의에서 바로 따른다. 이로써 (3)이 증명된다.
$\text{Proj}(S) \cap D(g_0)$에 대한 공식을 생각하자.
우변이 좌변에 포함되는 것은 명백하다.
반대 포함을 위해 $\mathfrak p \in \text{Proj}(S)$이고
$g_0 \not \in \mathfrak p$라고 하자.
양의 차수인 모든 동차 $f$에 대해
$g_0f \in \mathfrak p$이면 $S_{+} \subset \mathfrak p$가 되어 모순이다.
이로써 반대 포함을 얻고 (4)가 증명된다.

\medskip\noindent
표준 열린집합들 $D(g) \subset \Spec(S)$이 $\Spec(S)$의
위상의 기저를 이루므로, 열린집합들
$D(g) \cap \text{Proj}(S)$의 모음은 위상의 기저를 이룬다.
위 공식들에 의해 $D(g) \cap \text{Proj}(S)$을
열린집합들 $D_{+}(f)$의 합집합으로 나타낼 수 있다.
따라서 열린집합들 $D_{+}(f)$의 모음도 위상의 기저를 이룬다.
이로써 (5)가 증명된다.

\medskip\noindent
(6)의 증명. 먼저 보조정리 \ref{algebra-lemma-standard-open}에 의해
$D_{+}(f)$를 $D(f) = \Spec(S_f)$의 부분집합(유도 위상 포함)으로
식별할 수 있음에 유의하자.
환 $S_f$는 $\mathbf{Z}$-등급을 갖는다.
동차 원소들은 동차 $r \in S$에 대해 $r/f^n$ 꼴이고,
차수는 $\deg(r/f^n) = \deg(r) - n\deg(f)$이다.
부분집합 $D_{+}(f)$는 $\mathbf{Z}$-등급 아이디얼인
(즉 동차 원소들로 생성되는) 소 아이디얼
$\mathfrak p \subset S_f$들에 정확히 대응한다.
따라서 $S_f$의 $\mathbf{Z}$-등급 소 아이디얼들의 집합이
$\Spec(S_{(f)})$와 위상동형으로 사상됨을 보이면 된다.
이는 보조정리 \ref{algebra-lemma-Z-graded}에서 따른다.

\medskip\noindent
$S = \mathbf{Z}[X_1, X_2, X_3, \ldots]$라 하고,
각 $X_i$의 차수가 $1$인 등급을 주자. 그러면
$$
\text{Proj}(S) = \bigcup\nolimits_{i = 1}^\infty D_{+}(X_i)
$$
가 유한 세분을 갖지 않음을 쉽게 알 수 있다.
이로써 (7)이 증명된다.

\medskip\noindent
$I \subset S$를 등급 아이디얼이라 하자.
$\sqrt{I} \supset S_{+}$이면, 정의에 의해 모든 소 아이디얼
$\mathfrak p \in \text{Proj}(S)$은 $S_{+}$를 포함하지 않으므로
$V_{+}(I) = \emptyset$이다.
역으로 $S_{+} \not \subset \sqrt{I}$라고 하자.
그러면 원소 $f \in S_{+}$를 찾을 수 있으며, 이때 $f$는
$I$을 법으로 멱영이 아니다.
이는 분명 $f$의 동차 성분 가운데 하나가 $I$을 법으로
멱영이 아니라는 뜻이므로, $f$가 동차라고 가정해도 되고 그렇게 한다.
이는 $I S_f \not = S_f$, 즉 $(S/I)_f$가 영이 아님을 함의한다.
따라서 $(S/I)_{(f)} \not = 0$이다.
이는 $(S/I)_f$ 안으로 사상되는 환이기 때문이다.
소 아이디얼 $\mathfrak q \subset (S/I)_{(f)}$를 택하자.
이는 무관 아이디얼 $(S/I)_{+}$를 포함하지 않는
$S/I$의 등급 소 아이디얼에 대응한다.
다시 이는 $I$를 포함하지만 $S_{+}$를 포함하지 않는
$S$의 등급 소 아이디얼 $\mathfrak p$에 대응한다.
이로써 (10)이 증명되고 전체 증명이 끝난다.
\end{proof}

\begin{example}
\label{algebra-example-proj-polynomial-ring-1-variable}
$R$을 환이라 하자. $\deg(X) = 1$인 $S = R[X]$이면
자연스러운 사상 $\text{Proj}(S) \to \Spec(R)$은
전단사이고, 사실 위상동형사상이다.
실제로 $\mathfrak p \in \text{Proj}(S)$라고 하자.
$S_{+} \not \subset \mathfrak p$이므로 $X \not \in \mathfrak p$이다.
따라서 $a \in R$, $n > 0$이고 $aX^n \in \mathfrak p$이면
$a \in \mathfrak p$이다.
그러므로 $\mathfrak p_0 = \mathfrak p \cap R$에 대해
$\mathfrak p = \mathfrak p_0S$이다.
\end{example}

\noindent
$\mathfrak p \in \text{Proj}(S)$이면 $S_{(\mathfrak p)}$를
다음 환으로 정의한다.
그 원소는 같은 차수의 동차 원소 $r, f \in S$에 대해
$f \not\in \mathfrak p$인 분수 $r/f$이다.
통상과 같이, 어떤 동차 원소 $f'' \in S$가 존재하여
$f'' \not \in \mathfrak p$이고 $f''(rf' - r'f) = 0$일
필요충분조건이 $r/f = r'/f'$인 것으로 정한다.
등급 $S$-모듈 $M$이 주어지면 $M_{(\mathfrak p)}$를
다음 $S_{(\mathfrak p)}$-모듈로 둔다.
그 원소는 같은 차수의 동차 원소 $x \in M$, $f \in S$에 대해
$f \not \in \mathfrak p$인 분수 $x/f$이다.
어떤 동차 원소 $f'' \in S$가 존재하여
$f'' \not \in \mathfrak p$이고 $f''(xf' - x'f) = 0$일
필요충분조건이 $x/f = x'/f'$인 것으로 정한다.

\begin{lemma}
\label{algebra-lemma-proj-prime}
$S$를 등급환이라 하자. $M$을 등급 $S$-모듈이라 하자.
$\mathfrak p$를 $\text{Proj}(S)$의 원소라 하자.
$f \in S$를 $f \not \in \mathfrak p$, 즉
$\mathfrak p \in D_{+}(f)$를 만족하는 양의 차수의 동차 원소라 하자.
$\mathfrak p' \subset S_{(f)}$를 보조정리
\ref{algebra-lemma-topology-proj}에서처럼 $\mathfrak p$에 대응하는
$\Spec(S_{(f)})$의 원소라 하자. 그러면
$S_{(\mathfrak p)} = (S_{(f)})_{\mathfrak p'}$이고, 이와 호환되게
$M_{(\mathfrak p)} = (M_{(f)})_{\mathfrak p'}$이다.
\end{lemma}

\begin{proof}
사상 $\psi : M_{(\mathfrak p)} \to (M_{(f)})_{\mathfrak p'}$를 정의하자.
$x/g \in M_{(\mathfrak p)}$라 하고
$$
\psi(x/g) = (x g^{\deg(f) - 1}/f^{\deg(x)})/(g^{\deg(f)}/f^{\deg(g)}).
$$
로 둔다. $\deg(x) = \deg(g)$이고
$g^{\deg(f)}/f^{\deg(g)} \not \in \mathfrak p'$이므로 이는 의미가 있다.
$\psi$가 잘 정의되고 모듈 사상이며 동형사상이라는 확인은 생략한다.
힌트: 역은 $(x/f^n)/(g/f^m)$을 $(xf^m)/(g f^n)$으로 보낸다.
\end{proof}

\noindent
다음은 보조정리 \ref{algebra-lemma-silly}의 등급판이다.

\begin{lemma}
\label{algebra-lemma-graded-silly}
$S$가 등급환이고, $\mathfrak p_i$, $i = 1, \ldots, r$가
동차 소 아이디얼들이며, $I \subset S_{+}$가 등급 아이디얼이라고 하자.
모든 $i$에 대해 $I \not\subset \mathfrak p_i$라고 가정하자.
그러면 모든 $i$에 대해 $x\not\in \mathfrak p_i$인
양의 차수의 동차 원소 $x\in I$가 존재한다.
\end{lemma}

\begin{proof}
$\mathfrak p_i$들 사이에 포함관계가 없다고 가정해도 된다.
$r = 1$인 경우 결과는 참이다.
$r - 1$에 대해 결과가 성립한다고 하자.
모든 $i = 1, \ldots, r - 1$에 대해
$x \not \in \mathfrak p_i$인 양의 차수의 동차 원소 $x \in I$를 택하자.
$x \not\in \mathfrak p_r$이면 끝났다.
그러므로 $x \in \mathfrak p_r$이라고 가정하자.
$I \mathfrak p_1 \ldots \mathfrak p_{r-1} \subset \mathfrak p_r$이면
$I \subset \mathfrak p_r$이 되어 모순이다.
동차 원소 $y \in I\mathfrak p_1 \ldots \mathfrak p_{r-1}$를 택하되,
$y \not \in \mathfrak p_r$라 하자.
그러면 $x^{\deg(y)} + y^{\deg(x)}$가 원하는 원소이다.
\end{proof}

\begin{lemma}
\label{algebra-lemma-smear-out}
$S$를 등급환이라 하자.
$\mathfrak p \subset S$를 소 아이디얼이라 하자.
$\mathfrak q$를 $\mathfrak p$의 동차 원소들로 생성되는
$S$의 동차 아이디얼이라 하자.
그러면 $\mathfrak q$는 $S$의 소 아이디얼이다.
\end{lemma}

\begin{proof}
$\mathfrak q$가 소 아이디얼임을 증명하려면,
$f, g \in S$가 동차이고 $fg \in \mathfrak q$이면
% P01-E1378
$f$나 $g$ 가운데 하나가 $\mathfrak q$에 속함을 확인하면 충분하다.
% P01-E0375
$\mathfrak p \subset \mathfrak q$이므로 $fg \in \mathfrak p$이다.
$\mathfrak p$가 소 아이디얼이므로
$f \in \mathfrak p$이거나 $g \in \mathfrak p$이다.
$f$와 $g$가 동차이므로
$f \in \mathfrak q$이거나 $g \in \mathfrak q$임이 명백하다.
\end{proof}

\begin{lemma}
\label{algebra-lemma-graded-ring-minimal-prime}
$S$를 등급환이라 하자.
\begin{enumerate}
\item $S$의 모든 극소 소 아이디얼은 $S$의 동차 아이디얼이다.
\item 동차 아이디얼 $I \subset S$가 주어지면,
$I$ 위의 모든 극소 소 아이디얼은 동차이다.
\end{enumerate}
\end{lemma}

\begin{proof}
첫째 주장은 보조정리 \ref{algebra-lemma-smear-out}에서 구성한 소 아이디얼
$\mathfrak q$가 $\mathfrak q \subset \mathfrak p$를 만족하므로 성립한다.
% P01-E1379
둘째 주장은 $S/I$를 생각하고 첫째 부분을 적용하면 따른다.
\end{proof}

\begin{lemma}
\label{algebra-lemma-dehomogenize-finite-type}
$R$을 환이라 하자. $S$를 등급 $R$-대수라 하자.
$f \in S_{+}$를 동차 원소라 하자.
$S$가 $R$ 위에서 유한형이라고 가정하자. 그러면
\begin{enumerate}
\item 환 $S_{(f)}$는 $R$ 위에서 유한형이다.
\item 임의의 유한 등급 $S$-모듈 $M$에 대해
모듈 $M_{(f)}$는 유한 $S_{(f)}$-모듈이다.
\end{enumerate}
\end{lemma}

\begin{proof}
$R$-대수로서 $S$를 생성하는
$f_1, \ldots, f_n \in S$를 택하자.
각 $f_i$를 그 동차 성분들로 분해하여,
각 $f_i$가 동차라고 가정해도 된다.
$S_{(f)}$의 원소는
$$
\sum\nolimits_{e\deg(f) =
\sum e_i\deg(f_i)} \lambda_{e_1 \ldots e_n} f_1^{e_1} \ldots f_n^{e_n}/f^e
$$
꼴의 합이며 $\lambda_{e_1 \ldots e_n} \in R$이다.
따라서 $S_{(f)}$는
$e\deg(f) = \sum e_i\deg(f_i)$를 만족하는 원소들
$f_1^{e_1} \ldots f_n^{e_n} /f^e$로 $R$-대수로서 생성된다.
$e_i \geq \deg(f)$이면 이를
$$
f_1^{e_1} \ldots f_n^{e_n}/f^e =
f_i^{\deg(f)}/f^{\deg(f_i)} \cdot
f_1^{e_1} \ldots f_i^{e_i - \deg(f)} \ldots f_n^{e_n}/f^{e - \deg(f_i)}
$$
로 쓸 수 있다.
따라서 원소들 $f_i^{\deg(f)}/f^{\deg(f_i)}$와,
$e \deg(f) = \sum e_i \deg(f_i)$ 및 $e_i < \deg(f)$를 만족하는
원소들 $f_1^{e_1} \ldots f_n^{e_n} /f^e$만 있으면 된다.
이는 유한 목록이므로 (1)이 참이다.

\medskip\noindent
(2)를 보이기 위해 $M$이 동차 원소들
$x_1, \ldots, x_m$로 생성된다고 하자.
그러면 위와 같이 논증하여 $M_{(f)}$가
$e \deg(f) = \sum e_i \deg(f_i) + \deg(x_j)$와
$e_i < \deg(f)$를 만족하는 유한 목록의 원소들
$f_1^{e_1} \ldots f_n^{e_n} x_j /f^e$로
$S_{(f)}$-모듈로서 생성됨을 알 수 있다.
\end{proof}

\begin{lemma}
\label{algebra-lemma-homogenize}
$R$을 환이라 하자.
$R'$을 유한형 $R$-대수라 하고, $M$을 유한 $R'$-모듈이라 하자.
차수 $1$의 동차 원소 $f \in S$를 갖는 등급 $R$-대수 $S$와
등급 $S$-모듈 $N$이 존재하여 다음을 만족한다.
\begin{enumerate}
\item $R' \cong S_{(f)}$이고 $M \cong N_{(f)}$이다(모듈로서).
\item $S_0 = R$이고 $S$는 $R$ 위에서 차수 $1$인
유한 개의 원소로 생성된다.
\item $N$은 유한 $S$-모듈이다.
\end{enumerate}
\end{lemma}

\begin{proof}
어떤 아이디얼 $I$에 대해
$R' = R[x_1, \ldots, x_n]/I$로 쓸 수 있다.
% P01-E1310
$g \in R[x_1, \ldots, x_n]$에 대해
$g = \tilde g(1, x_1, \ldots, x_n)$를 만족하는
최소 차수의 동차 원소
$\tilde g \in R[X_0, \ldots, X_n]$를 나타내자.
% P01-E1380
$\tilde I \subset R[X_0, \ldots, X_n]$를
모든 원소 $\tilde g$, $g \in I$로 생성되는 아이디얼이라 하자.
$S = R[X_0, \ldots, X_n]/\tilde I$로 두고,
% P01-E1381
$f$를 $S$에서 $X_0$의 상이라 하자.
구성에 의해 동형사상
$$
S_{(f)} \longrightarrow R', \quad
X_i/X_0 \longmapsto x_i.
$$
을 얻는다.
모듈 $M$에 대해 같은 일을 하기 위해 표시
$$
M = (R')^{\oplus r}/\sum\nolimits_{j \in J} R'k_j
$$
를 택하자. 여기서 $k_j = (k_{1j}, \ldots, k_{rj})$이다.
% P01-E0376
$d_{ij} = \deg(\tilde k_{ij})$로 두고
$d_j = \max\{d_{ij}\}$로 두자.
$K_{ij} = X_0^{d_j - d_{ij}}\tilde k_{ij}$로 두면
이는 차수 $d_j$의 동차 원소이다. 이 표기를 사용하여
$$
N = \Coker\Big(
\bigoplus\nolimits_{j \in J} S(-d_j) \xrightarrow{(K_{ij})} S^{\oplus r}
\Big)
$$
로 두면 된다. 몇 가지 세부사항은 생략한다.
\end{proof}





\section{뇌터 등급환}
\label{algebra-section-noetherian-graded}

\noindent
힐베르트 다항식에 관한 내용을 포함하여
뇌터 등급환의 이론을 조금 다룬다.

\begin{lemma}
\label{algebra-lemma-S-plus-generated}
$S$를 등급환이라 하자. 동차 원소들의 집합
$f_i \in S_{+}$가 $S$를 $S_0$ 위의 대수로서 생성할 필요충분조건은
그 집합이 $S_{+}$를 $S$의 아이디얼로서 생성하는 것이다.
\end{lemma}

\begin{proof}
$f_i$들이 $S$를 $S_0$ 위의 대수로서 생성하면,
$S_{+}$의 모든 원소는 $f_i$들에 관한 상수항이 없는 다항식이다.
따라서 $S_{+}$는 아이디얼로서 $f_i$들에 의해 생성된다.
역으로 $S_{+} = \sum Sf_i$라고 하자.
$S$의 임의의 원소 $f$를 $S_0$에 계수를 둔
$f_i$들의 다항식으로 쓸 수 있음을 보이겠다.
동차 원소에 대해서만 보이면 충분하다. $f$의 차수를 $d$라 하자.
$d$에 관한 귀납법을 쓸 수 있다. $d = 0$인 경우는 자명하다.
$d > 0$이면 $f \in S_{+}$이므로 어떤 $g_i \in S$에 대해
$f = \sum g_i f_i$로 쓸 수 있다.
$S$가 등급환이므로 $g_i$를 차수 $d - \deg(f_i)$인
동차 성분으로 대체할 수 있다.
귀납가정에 의해 각 $g_i$는 $f_i$들의 다항식이므로 결론을 얻는다.
\end{proof}

\begin{lemma}
\label{algebra-lemma-graded-Noetherian}
등급환 $S$가 뇌터일 필요충분조건은
$S_0$가 뇌터이고 $S_{+}$가 $S$의 아이디얼로서 유한 생성되는 것이다.
\end{lemma}

\begin{proof}
$S$가 뇌터이면 $S_0 = S/S_{+}$가 뇌터이고
$S_{+}$가 유한 생성된다는 것은 명백하다.
역으로 $S_0$가 뇌터이고 $S_{+}$가 $S$의 아이디얼로서
유한 생성된다고 하자. 생성원
$S_{+} = (f_1, \ldots, f_n)$을 택하자.
$f_i$들을 동차 성분들로 분해하여 각 $f_i$가 동차라고 가정해도 된다.
보조정리 \ref{algebra-lemma-S-plus-generated}에 의해,
% P01-E0377
$X_i$를 $f_i$로 보내는 $S_0[X_1, \ldots X_n] \to S$는 전사이다.
따라서 보조정리 \ref{algebra-lemma-Noetherian-permanence}에 의해 $S$는 뇌터이다.
\end{proof}

\begin{definition}
\label{algebra-definition-numerical-polynomial}
$A$를 아벨 군이라 하자.
충분히 큰 모든 정수 $n$에 대해 정의된 함수
$f : n \mapsto f(n) \in A$가 {\it 수치 다항식}이라는 것은,
$r \geq 0$과 원소 $a_0, \ldots, a_r\in A$가 존재하여
모든 $n \gg 0$에 대해
$$
f(n) = \sum\nolimits_{i = 0}^r \binom{n}{i} a_i
$$
가 성립한다는 뜻이다.
\end{definition}

\noindent
이항계수를 쓰는 이유는 정수점에서의 값이 모두 정수인 임의의 다항식
$P \in \mathbf{Q}[T]$가 $a_i \in \mathbf{Z}$에 대해
$P(T) = \sum a_i \binom{T}{i}$ 꼴의 합과 같다는
초등적 사실 때문이다. 특히 식 $\binom{T + 1}{i + 1}$도
이 꼴임에 유의하자.

\begin{lemma}
\label{algebra-lemma-numerical-polynomial-functorial}
$A \to A'$가 아벨 군의 준동형이고
$f : n \mapsto f(n) \in A$가 수치 다항식이면,
그 합성도 수치 다항식이다.
\end{lemma}

\begin{proof}
이는 정의에서 바로 따른다.
\end{proof}

\begin{lemma}
\label{algebra-lemma-numerical-polynomial}
$f: n \mapsto f(n) \in A$가 충분히 큰 모든 $n$에 대해 정의되고,
$n \mapsto f(n) - f(n-1)$이 수치 다항식이라고 하자.
그러면 $f$는 수치 다항식이다.
\end{lemma}

\begin{proof}
모든 $n \gg 0$에 대해
$f(n) - f(n-1) = \sum\nolimits_{i = 0}^r \binom{n}{i} a_i$라 하자.
$g(n) = f(n) - \sum\nolimits_{i = 0}^r \binom{n + 1}{i + 1} a_i$로 두자.
그러면 모든 $n \gg 0$에 대해 $g(n) - g(n-1) = 0$이다.
따라서 $g$는 결국 상수이며, 그 값을 $a_{-1}$이라 하자.
$a_{-1} + \sum\nolimits_{i = 0}^r \binom{n + 1}{i + 1} a_i$가
요구되는 꼴임을 보이는 것은 독자에게 맡긴다
(보조정리 앞의 설명을 보라).
\end{proof}

\begin{lemma}
\label{algebra-lemma-graded-module-fg}
$M$이 유한 생성 등급 $S$-모듈이고
$S$가 $S_0$ 위에서 유한 생성되면,
각 $M_n$은 유한 $S_0$-모듈이다.
\end{lemma}

\begin{proof}
$M$의 생성원들을 $m_i$, $S$의 생성원들을 $f_i$라 하자.
동차 성분들을 취하여 $m_i$와 $f_i$가 동차라고 가정할 수 있고,
$f_i \in S_{+}$라고 가정할 수도 있다.
이 경우 각 $M_n$은 차수가 $n$인 “단항식”
$\prod f_i^{e_i} m_j$들에 의해 $S_0$ 위에서 생성됨이 명백하다.
\end{proof}

\begin{proposition}
\label{algebra-proposition-graded-hilbert-polynomial}
$S$를 뇌터 등급환, $M$을 유한 등급 $S$-모듈이라 하자.
함수
$$
\mathbf{Z} \longrightarrow K'_0(S_0), \quad
n \longmapsto [M_n]
$$
를 생각하자. 보조정리 \ref{algebra-lemma-graded-module-fg}를 보라.
$S_{+}$가 차수 $1$인 원소들로 생성되면 이 함수는 수치 다항식이다.
\end{proposition}

\begin{proof}
$S_1$의 생성원들의 최소 개수에 관한 귀납법으로 증명한다.
이 개수가 $0$이면 모든 $n \gg 0$에 대해 $M_n = 0$이므로
결론이 성립한다.
귀납 단계를 증명하기 위해 최소 생성집합의 한 원소
$x\in S_1$을 택하여 등급환 $S/(x)$에 귀납가정을 적용할 수 있게 하자.

\medskip\noindent
먼저 $x$가 $M$ 위에서 멱영이면 결론이 성립함을 보인다.
이는 $x^r M  = 0$이 되는 최소 정수 $r$에 관한 귀납법으로 증명한다.
$r = 1$이면 $M$은 $S/xS$ 위의 모듈이므로,
다른 귀납가정에 의해 결론이 성립한다.
% P01-E0378
$r > 1$이면 정수 $r', r''$가 $r$보다 모두 작도록 하는 짧은 완전열
$0 \to M' \to M \to M'' \to 0$을 찾을 수 있다.
따라서 $M''$와 $M'$에 대해서는 결론을 알고 있다.
그러므로 $K'_0(S_0)$에서의 관계
$
[M_d]  = [M'_d] + [M''_d]
$
에 의해 $M$에 대한 결론을 얻는다.

\medskip\noindent
$x$가 $M$ 위에서 멱영이 아니면,
$x$가 멱영으로 작용하는 가장 큰 부분모듈을 $M' \subset M$이라 하자.
완전열 $0 \to M' \to M \to M/M' \to 0$을 생각하면,
다시 $M/M'$에 대해 결론을 증명하면 충분함을 알 수 있다.
즉 $x$에 의한 곱셈이 단사라고 가정해도 된다.

\medskip\noindent
$\overline{M} = M/xM$으로 두자.
사상 $x : M \to M$은 $M_d$를 $M_d$로 보내지 않으므로
등급 $S$-모듈의 사상이 {\it 아니다}.
구체적으로 각 $d$에 대해 다음 짧은 완전열이 있다.
$$
0 \to M_d \xrightarrow{x} M_{d + 1} \to \overline{M}_{d + 1} \to 0
$$
따라서 $[M_{d + 1}] - [M_d] = [\overline{M}_{d + 1}]$이다.
보조정리 \ref{algebra-lemma-numerical-polynomial}에 의해 결론을 얻는다.
\end{proof}

\begin{remark}
\label{algebra-remark-period-polynomial}
$S$가 여전히 뇌터이지만 $S$가 차수 $1$에서 생성되지 않으면,
등급 $S$-모듈에 딸린 함수는 주기 다항식이다
(즉 어떤 $n$에 대해 정수들의 법 $n$ 합동류 각각에서 수치 다항식이다).
\end{remark}

\begin{example}
\label{algebra-example-hilbert-function}
$S = k[X_1, \ldots, X_d]$라 하자.
예제 \ref{algebra-example-K0-field}에 의해
$K_0(k) = K'_0(k) = \mathbf{Z}$로 식별할 수 있다.
따라서 임의의 유한 생성 등급 $k[X_1, \ldots, X_d]$-모듈은
수치 다항식 $n \mapsto \dim_k(M_n)$을 준다.
\end{example}

\begin{lemma}
\label{algebra-lemma-quotient-smaller-d}
$k$를 체라 하자. $I \subset k[X_1, \ldots, X_d]$를
0이 아닌 등급 아이디얼이라 하자.
$M = k[X_1, \ldots, X_d]/I$로 두자.
그러면 수치 다항식 $n \mapsto \dim_k(M_n)$
(예제 \ref{algebra-example-hilbert-function}를 보라)의 차수는
$ < d - 1$이다($d = 1$이면 영이다).
\end{lemma}

\begin{proof}
등급 모듈 $k[X_1, \ldots, X_d]$에 딸린 수치 다항식은
$n \mapsto \binom{n - 1 + d}{d - 1}$이다.
% P01-E0379
차수 $e$인 임의의 0이 아닌 동차 원소 $f \in I$와
임의의 차수 $n >> e$에 대해
$I_n \supset f \cdot k[X_1, \ldots, X_d]_{n-e}$이므로
$\dim_k(I_n) \geq \binom{n - e - 1 + d}{d - 1}$이다.
따라서
$\dim_k(M_n) \leq \binom{n - 1 + d}{d - 1} - \binom{n - e - 1 + d}{d - 1}$이다.
마지막 식의 차수는 $ < d - 1$이므로 결론을 얻는다
($d = 1$이면 영이다).
\end{proof}






\section{뇌터 국소환}
\label{algebra-section-Noetherian-local}

\noindent
이 절 전체에서 $(R, \mathfrak m, \kappa)$는 뇌터 국소환이다.
이 절에서는 모듈의 힐베르트 함수에 관한 이론을 전개한다.
$M$을 유한 $R$-모듈이라 하자. $M$의 {\it 힐베르트 함수}를
모든 정수 $n \geq 0$에 대해 정의된 함수
$$
\varphi_M : n
\longmapsto
\text{length}_R(\mathfrak m^nM/{\mathfrak m}^{n + 1}M)
$$
로 정의한다. 또 하나의 중요한 불변량은 모든 정수 $n \geq 0$에 대해
정의된 함수
$$
\chi_M : n
\longmapsto
\text{length}_R(M/{\mathfrak m}^{n + 1}M)
$$
이다. 보조정리 \ref{algebra-lemma-length-additive}에 의해
$$
\chi_M(n) = \sum\nolimits_{i = 0}^n \varphi_M(i).
$$
가 성립함에 유의하자.
이 구성에는 정의 아이디얼을 사용하는 변형이 있다.

\begin{definition}
\label{algebra-definition-ideal-definition}
$(R, \mathfrak m)$을 뇌터 국소환이라 하자.
$\sqrt{I} = \mathfrak m$을 만족하는 아이디얼 $I \subset R$을
{\it $R$의 정의 아이디얼}이라 한다.
\end{definition}

\noindent
$I \subset R$을 정의 아이디얼이라 하자.
$R$이 뇌터이므로 이는 어떤 $r$에 대해
$\mathfrak m^r \subset I$라는 뜻이다.
보조정리 \ref{algebra-lemma-Noetherian-power}를 보라.
% P01-E0380
따라서 $I$의 한 거듭제곱에 의해 소멸되는 모든 유한 $R$-모듈은
유한 길이를 갖는다. 보조정리 \ref{algebra-lemma-length-finite}를 보라.
그러므로 모든 $n \geq 0$에 대해
$$
\varphi_{I, M}(n) = \text{length}_R(I^nM/I^{n + 1}M)
\quad\text{and}\quad
\chi_{I, M}(n) = \text{length}_R(M/I^{n + 1}M)
$$
로 정의하는 것이 타당하다. 다시
$$
\chi_{I, M}(n) = \sum\nolimits_{i = 0}^n \varphi_{I, M}(i).
$$
가 성립한다.

\begin{lemma}
\label{algebra-lemma-differ-finite}
$M' \subset M$을 몫이 유한 길이를 갖는 유한 $R$-모듈들이라 하자.
% P01-E0381
그러면 상수 $c_1, c_2$가 존재하여 모든 $n \geq c_2$에 대해
$$
c_1 + \chi_{I, M'}(n - c_2) \leq \chi_{I, M}(n) \leq
c_1 + \chi_{I, M'}(n)
$$
가 성립한다.
\end{lemma}

\begin{proof}
$M/M'$가 유한 길이를 가지므로
$I^{c_2}M \subset M'$가 되는 $c_2 \geq 0$이 존재한다.
$c_1 = \text{length}_R(M/M')$로 두자.
$n \geq c_2$에 대해
\begin{eqnarray*}
\chi_{I, M}(n)
& = &
\text{length}_R(M/I^{n + 1}M) \\
& = &
c_1 + \text{length}_R(M'/I^{n + 1}M) \\
& \leq &
c_1 + \text{length}_R(M'/I^{n + 1}M') \\
& = &
c_1 + \chi_{I, M'}(n)
\end{eqnarray*}
이다. 한편 $I^{c_2}M \subset M'$이므로,
$n \geq c_2$에 대해 $I^nM \subset I^{n - c_2}M'$이다.
따라서 $n \geq c_2$에 대해
\begin{eqnarray*}
\chi_{I, M}(n)
& = &
\text{length}_R(M/I^{n + 1}M) \\
& = &
c_1 + \text{length}_R(M'/I^{n + 1}M) \\
& \geq &
c_1 + \text{length}_R(M'/I^{n + 1 - c_2}M') \\
& = &
c_1 + \chi_{I, M'}(n - c_2)
\end{eqnarray*}
를 얻으므로 증명이 끝난다.
\end{proof}

\begin{lemma}
\label{algebra-lemma-hilbert-ses}
$0 \to M' \to M \to M'' \to 0$을 유한 $R$-모듈들의 짧은 완전열이라 하자.
그러면 유한 여길이 $l$을 갖는 부분모듈 $N \subset M'$와
$c \geq 0$이 존재하여 모든 $n \geq c$에 대해
$$
\chi_{I, M}(n) = \chi_{I, M''}(n) + \chi_{I, N}(n - c) + l
$$
및
$$
\varphi_{I, M}(n) = \varphi_{I, M''}(n) + \varphi_{I, N}(n - c)
$$
가 성립한다.
\end{lemma}

\begin{proof}
$M/I^nM \to M''/I^nM''$은 전사이고 그 핵은
$M' / M' \cap I^nM$임에 유의하자.
아르틴--리스 보조정리 \ref{algebra-lemma-Artin-Rees}에 의해
$M' \cap I^nM = I^{n - c}(M' \cap I^cM)$이 되는 상수 $c$가 존재한다.
$N = M' \cap I^cM$이라 쓰자.
$I^c M' \subset N \subset M'$임에 유의하자.
따라서 $n \geq c$에 대해
$\text{length}_R(M' / M' \cap I^nM)
= \text{length}_R(M'/N) + \text{length}_R(N/I^{n - c}N)$이다.
짧은 완전열
$$
0 \to M' / M' \cap I^nM \to M/I^nM \to M''/I^nM'' \to 0
$$
과 길이의 가법성(보조정리 \ref{algebra-lemma-length-additive})에서
$n \geq c$에 대해 등식
$$
\chi_{I, M}(n - 1)
=
\chi_{I, M''}(n - 1)
+
\chi_{I, N}(n - c - 1)
+
\text{length}_R(M'/N)
$$
을 얻는다.
$\varphi_{I, M}(n) = \chi_{I, M}(n) - \chi_{I, M}(n - 1)$이고
모듈 $M''$와 $N$에 대해서도 마찬가지이다.
따라서 모든 $n \geq c$에 대해
$\varphi_{I, M}(n) = \varphi_{I, M''}(n) + \varphi_{I, N}(n-c)$을 얻는다.
\end{proof}

\begin{lemma}
\label{algebra-lemma-hilbert-change-I}
$I$, $I'$을 뇌터 국소환 $R$의 두 정의 아이디얼이라 하자.
$M$을 유한 $R$-모듈이라 하자. 그러면 상수 $a$가 존재하여
$n \geq 1$에 대해 $\chi_{I, M}(n) \leq \chi_{I', M}(an)$이 성립한다.
\end{lemma}

\begin{proof}
$(I')^c \subset I$인 정수 $c \geq 1$이 존재한다.
따라서 전사 $M/(I')^{c(n + 1)}M \to M/I^{n + 1}M$을 얻는다.
그러므로 $a = 2c - 1$로 두면 결과가 따른다.
\end{proof}

\begin{proposition}
\label{algebra-proposition-hilbert-function-polynomial}
$R$을 뇌터 국소환이라 하자. $M$을 유한 $R$-모듈이라 하자.
$I \subset R$을 정의 아이디얼이라 하자.
힐베르트 함수 $\varphi_{I, M}$와 함수
$\chi_{I, M}$은 수치 다항식이다.
\end{proposition}

\begin{proof}
등급환 $S = R/I \oplus I/I^2 \oplus I^2/I^3 \oplus
\ldots = \bigoplus_{d \geq 0} I^d/I^{d + 1}$과 등급
$S$-모듈 $N = M/IM \oplus IM/I^2M \oplus \ldots =
\bigoplus_{d \geq 0} I^dM/I^{d + 1}M$을 생각하자.
이 쌍 $(S, N)$은 명제 \ref{algebra-proposition-graded-hilbert-polynomial}의
가정들을 만족한다.
따라서 $\varphi_{I, M}$에 관한 결과는 그 명제와
보조정리 \ref{algebra-lemma-length-K0}에서 따른다.
$\chi_{I, M}$에 관한 결과는 이 사실과
보조정리 \ref{algebra-lemma-numerical-polynomial}에서 따른다.
\end{proof}

\begin{definition}
\label{algebra-definition-hilbert-polynomial}
$R$을 뇌터 국소환이라 하자. $M$을 유한 $R$-모듈이라 하자.
$R$ 위에서 $M$의 {\it 힐베르트 다항식}은
모든 $n \gg 0$에 대해 $P(n) = \varphi_M(n)$을 만족하는 원소
$P(t) \in \mathbf{Q}[t]$이다.
\end{definition}

\noindent
명제 \ref{algebra-proposition-hilbert-function-polynomial}에 의해
힐베르트 다항식이 존재함을 알 수 있다.

\begin{lemma}
\label{algebra-lemma-d-independent}
$R$을 뇌터 국소환이라 하자. $M$을 유한 $R$-모듈이라 하자.
\begin{enumerate}
\item 수치 다항식 $\varphi_{I, M}$의 차수는
정의 아이디얼 $I$에 무관하다.
\item 수치 다항식 $\chi_{I, M}$의 차수는
정의 아이디얼 $I$에 무관하다.
\end{enumerate}
\end{lemma}

\begin{proof}
(2)는 보조정리 \ref{algebra-lemma-hilbert-change-I}에서 바로 따른다.
$n \geq 1$에 대해
$\varphi_{I, M}(n) = \chi_{I, M}(n) - \chi_{I, M}(n - 1)$이므로
(1)은 (2)에서 따른다.
\end{proof}

\begin{definition}
\label{algebra-definition-d}
$R$을 뇌터 국소환, $M$을 유한 $R$-모듈이라 하자.
다음과 같이 정의되는 $\{-\infty, 0, 1, 2, \ldots \}$의 원소를
{\it $d(M)$}으로 나타낸다.
\begin{enumerate}
\item $M = 0$이면 $d(M) = -\infty$로 둔다.
\item $M \not = 0$이면 $d(M)$은 수치 다항식 $\chi_M$의 차수이다.
\end{enumerate}
\end{definition}

\noindent
모든 $n$에 대해 $\mathfrak m^nM \not = 0$이면
$d(M)$은 $M$의 힐베르트 다항식의 차수 $+1$이다.

\begin{lemma}
\label{algebra-lemma-differ-finite-chi}
$R$을 뇌터 국소환이라 하자. $I \subset R$을 정의 아이디얼이라 하자.
$M$을 유한 길이를 갖지 않는 유한 $R$-모듈이라 하자.
% P01-E0382
$M' \subset M$이 유한 여길이를 갖는 부분모듈이면
$\chi_{I, M} - \chi_{I, M'}$은 그 차수가 두 다항식 각각의 차수보다
$<$인 다항식이다.
\end{lemma}

\begin{proof}
보조정리 \ref{algebra-lemma-differ-finite}에서 초등 계산으로 따른다.
\end{proof}

\begin{lemma}
\label{algebra-lemma-hilbert-ses-chi}
$R$을 뇌터 국소환이라 하자. $I \subset R$을 정의 아이디얼이라 하자.
$0 \to M' \to M \to M'' \to 0$을 유한 $R$-모듈들의 짧은 완전열이라 하자.
그러면
\begin{enumerate}
\item $M'$이 유한 길이를 갖지 않으면
$\chi_{I, M} - \chi_{I, M''} - \chi_{I, M'}$은
차수가 $\chi_{I, M'}$의 차수보다 $<$인 수치 다항식이다.
\item $\max\{ \deg(\chi_{I, M'}), \deg(\chi_{I, M''}) \} = \deg(\chi_{I, M})$이고,
\item $\max\{d(M'), d(M'')\} = d(M)$이다.
\end{enumerate}
\end{lemma}

\begin{proof}
먼저 (1)을 증명한다.
$N \subset M'$을 보조정리 \ref{algebra-lemma-hilbert-ses}에서처럼 택하자.
보조정리 \ref{algebra-lemma-differ-finite-chi}에 의해 수치 다항식
$\chi_{I, M'} - \chi_{I, N}$의 차수는
$\chi_{I, M'}$와 $\chi_{I, N}$의 공통 차수보다 $<$이다.
보조정리 \ref{algebra-lemma-hilbert-ses}에 의해 차
$$
\chi_{I, M}(n) - \chi_{I, M''}(n) - \chi_{I, N}(n - c)
$$
는 $n \gg 0$에 대해 상수이다. 초등 계산에 의해 차
$\chi_{I, N}(n) - \chi_{I, N}(n - c)$의 차수는
$\chi_{I, N}$의 차수보다 $<$이고, 후자는 0보다 크다(위를 보라).
이들을 모두 합치면 (1)을 얻는다.

\medskip\noindent
$\chi_{I, M'}$와 $\chi_{I, M''}$의 최고차항 계수가 음이 아님에 유의하자.
따라서 $\chi_{I, M'} + \chi_{I, M''}$의 차수는
두 차수의 최댓값과 같다.
$M'$이 유한 길이를 갖지 않으면 (2)는 (1)에서 따른다.
$M'$이 유한 길이를 가지면 아르틴--리스
(보조정리 \ref{algebra-lemma-Artin-Rees})에 의해 모든 $n \gg 0$에 대해
$I^nM \to I^nM''$은 동형사상이다.
따라서 $n \gg 0$에 대해 $M/I^nM \to M''/I^nM''$은
핵이 $M'$인 전사이고, 모든 $n \gg 0$에 대해
$\chi_{I, M}(n) - \chi_{I, M''}(n) = \text{length}(M')$임을 알 수 있다.
그러므로 이 경우에도 (2)가 성립한다.

\medskip\noindent
(3)의 증명. $M$, $M'$, $M''$ 가운데 하나가 영인 경우를 제외하면
이는 (2)에서 따른다. 이 특별한 경우들의 증명은 생략한다.
\end{proof}

\section{차원}
\label{algebra-section-dimension}

\noindent
위상수학, 절 \ref{topology-section-krull-dimension}과 비교하라.

\begin{definition}
\label{algebra-definition-chain-primes}
$R$을 환이라 하자. {\it 소 아이디얼 사슬}은 $R$의 소 아이디얼들의 열
$\mathfrak p_0 \subset \mathfrak p_1 \subset \ldots \subset \mathfrak p_n$으로서,
$i = 0, \ldots, n - 1$에 대해
$\mathfrak p_i \not = \mathfrak p_{i + 1}$인 것이다.
이 소 아이디얼 사슬의 {\it 길이}는 $n$이다.
\end{definition}

\noindent
환 $R$의 소 아이디얼들과 $\Spec(R)$의 기약 닫힌 부분집합들 사이에는
포함관계를 뒤집는 전단사가 있음을 상기하자.
보조정리 \ref{algebra-lemma-irreducible}을 보라.

\begin{definition}
\label{algebra-definition-Krull}
환 $R$의 {\it 크룰 차원}은 위상공간 $\Spec(R)$의 크룰 차원이다.
위상수학, 정의 \ref{topology-definition-Krull}을 보라.
즉 이는 $R$이
$$
\mathfrak p_0
\subset
\mathfrak p_1
\subset
\ldots
\subset
\mathfrak p_n, \quad
\mathfrak p_i \not = \mathfrak p_{i + 1}.
$$
와 같은 길이 $n$의 소 아이디얼 사슬을 갖게 하는
정수 $n\geq 0$들의 상한이다.
\end{definition}

\begin{definition}
\label{algebra-definition-height}
환 $R$의 소 아이디얼 $\mathfrak p$의 {\it 높이}는
국소환 $R_{\mathfrak p}$의 차원이다.
\end{definition}

\begin{lemma}
\label{algebra-lemma-dimension-height}
$R$의 크룰 차원은 그 (극대) 소 아이디얼들의 높이의 상한이다.
\end{lemma}

\begin{proof}
소 아이디얼 사슬의 끝에는 언제나 극대 아이디얼을 하나 덧붙일 수 있기 때문이다.
\end{proof}

\begin{lemma}
\label{algebra-lemma-Noetherian-dimension-0}
차원 $0$인 뇌터 환은 아르틴 환이다.
역으로 모든 아르틴 환은 차원 0인 뇌터 환이다.
\end{lemma}

\begin{proof}
$R$이 차원 $0$인 뇌터 환이라고 하자.
보조정리 \ref{algebra-lemma-Noetherian-topology}에 의해 공간 $\Spec(R)$은 뇌터이다.
위상수학, 보조정리 \ref{topology-lemma-Noetherian}에 의해
$\Spec(R)$은 유한 개의 기약 성분을 가지며,
$\Spec(R) = Z_1 \cup \ldots \cup Z_r$라 쓰자.
보조정리 \ref{algebra-lemma-irreducible}에 따르면,
% P01-E0383
각 $Z_i = V(\mathfrak p_i)$이고 $\mathfrak p_i$는 극소 아이디얼이다.
차원이 $0$이므로 이 $\mathfrak p_i$들은 극대이기도 하다.
따라서 $\Spec(R)$은 원소들이 $\mathfrak p_i$인 이산 위상공간이다.
야코브슨 근기 $\bigcap \mathfrak p_i$의 모든 원소 $f$는 멱영이다.
그렇지 않으면 $R_f$가 영환이 아니어서 다른 소 아이디얼이 존재하기 때문이다.
보조정리 \ref{algebra-lemma-product-local}에 의해
$R$은 $\prod R_{\mathfrak p_i}$와 같다.
$R_{\mathfrak p_i}$도 차원 $0$인 뇌터 환이므로,
앞의 논증은 그 근기 $\mathfrak p_iR_{\mathfrak p_i}$가
국소적으로 멱영임을 보여 준다.
보조정리 \ref{algebra-lemma-Noetherian-power}에 의해
어떤 $n \geq 1$에 대해
$\mathfrak p_i^nR_{\mathfrak p_i} = 0$이다.
% P01-E0384
보조정리 \ref{algebra-lemma-length-finite}에 의해
$R_{\mathfrak p_i}$는 $R$ 위에서 유한 길이를 갖는다.
따라서 보조정리 \ref{algebra-lemma-artinian-finite-length}에 의해
$R$은 아르틴 환이다.

\medskip\noindent
$R$이 아르틴 환이면 보조정리 \ref{algebra-lemma-artinian-finite-length}에 의해
뇌터이다. 보조정리 \ref{algebra-lemma-artinian-finite-nr-max},
\ref{algebra-lemma-artinian-radical-nilpotent}, \ref{algebra-lemma-product-local}을
결합하면 그 모든 소 아이디얼은 극대이다.
\end{proof}

\noindent
이하에서는 정의 \ref{algebra-definition-d}에서 정의한 불변량 $d(-)$를 사용한다.
먼저 준비 보조정리를 하나 제시한다.

\begin{lemma}
\label{algebra-lemma-dimension-0-d-0}
$R$을 뇌터 국소환이라 하자.
그러면 $\dim(R) = 0 \Leftrightarrow d(R) = 0$이다.
\end{lemma}

\begin{proof}
$d(R) = 0$일 필요충분조건은 $R$이 $R$-모듈로서
유한 길이를 갖는 것이기 때문이다.
보조정리 \ref{algebra-lemma-artinian-finite-length}을 보라.
\end{proof}

\begin{proposition}
\label{algebra-proposition-dimension-zero-ring}
$R$을 환이라 하자. 다음 조건들은 서로 동치이다.
\begin{enumerate}
\item $R$은 아르틴 환이다.
\item $R$은 뇌터이고 $\dim(R) = 0$이다.
\item $R$은 자기 자신 위의 모듈로서 유한 길이를 갖는다.
\item $R$은 아르틴 국소환들의 유한 곱이다.
\item $R$은 뇌터이고 $\Spec(R)$은 유한 이산 위상공간이다.
\item $R$은 차원 $0$인 뇌터 국소환들의 유한 곱이다.
\item $R$은 $d(R_i) = 0$인 뇌터 국소환
$R_i$들의 유한 곱이다.
\item $R$은 극대 아이디얼들이 멱영인 뇌터 국소환
$R_i$들의 유한 곱이다.
\item $R$은 뇌터이고 극대 아이디얼이 유한 개이며
그 야코브슨 근기 아이디얼은 멱영이다.
\item $R$은 뇌터이고 그 소 아이디얼들 사이에는
진정한 포함관계가 없다.
\end{enumerate}
\end{proposition}

\begin{proof}
이는 보조정리
\ref{algebra-lemma-product-local},
\ref{algebra-lemma-artinian-finite-length},
\ref{algebra-lemma-Noetherian-dimension-0}, \ref{algebra-lemma-dimension-0-d-0}을
결합하면 얻는다.
\end{proof}

\begin{lemma}
\label{algebra-lemma-height-1}
$R$을 뇌터 국소환이라 하자.
다음 조건들은 서로 동치이다.
\begin{enumerate}
\item
\label{algebra-item-dim-1}
$\dim(R) = 1$,
\item
\label{algebra-item-d-1}
$d(R) = 1$,
\item
\label{algebra-item-Vx}
어떤 $x \in \mathfrak m$가 존재하고, $x$는 멱영이 아니며
$V(x) = \{\mathfrak m\}$이고,
\item
\label{algebra-item-x}
어떤 $x \in \mathfrak m$가 존재하고, $x$는 멱영이 아니며
$\mathfrak m = \sqrt{(x)}$이며,
\item
\label{algebra-item-ideal-1}
$1$개 원소로 생성되는 정의 아이디얼이 존재하고,
$0$개 원소로 생성되는 정의 아이디얼은 없다.
\end{enumerate}
\end{lemma}

\begin{proof}
먼저 $\dim(R) = 1$이라고 하자.
$\mathfrak p_i$들을 $R$의 극소 소 아이디얼들이라 하자.
차원이 $1$이므로 $R$의 다른 소 아이디얼은 $\mathfrak m$뿐이다.
보조정리 \ref{algebra-lemma-Noetherian-irreducible-components}에 따르면
그 수는 유한하다. 따라서 보조정리 \ref{algebra-lemma-silly}에 의해
$x \in \mathfrak m$, $x \not \in \mathfrak p_i$인 원소를 찾을 수 있다.
그러면 $x$를 포함하는 유일한 소 아이디얼은 $\mathfrak m$이므로
(\ref{algebra-item-Vx})를 얻는다.

\medskip\noindent
(\ref{algebra-item-Vx})이면 보조정리 \ref{algebra-lemma-Zariski-topology}에 의해
$\mathfrak m = \sqrt{(x)}$이므로 (\ref{algebra-item-x})가 성립한다.
그 역도 명백하다.
% P01-E0385
(\ref{algebra-item-x})와 (\ref{algebra-item-ideal-1})의 동치는
정의에서 바로 따른다.

\medskip\noindent
(\ref{algebra-item-ideal-1})을 가정하자.
$I = (x)$를 정의 아이디얼이라 하자.
$I^n/I^{n + 1}$은 $x^n$에 의한 곱셈을 통해 $R/I$의 몫이므로,
$\text{length}_R(I^n/I^{n + 1})$은 유계이다.
따라서 $d(R) = 0$ 또는 $d(R) = 1$인데,
$0$이 정의 아이디얼이 아니라는 가정에 의해 $d(R) = 0$은 배제된다.

\medskip\noindent
(\ref{algebra-item-d-1})을 가정하자. 모순을 얻기 위해
두 포함이 모두 진정한 소 아이디얼 사슬
$\mathfrak p \subset \mathfrak q \subset \mathfrak m$이
존재한다고 하자. 어떤 정의 아이디얼 $I \subset R$을 택하자.
보조정리 \ref{algebra-lemma-hilbert-ses-chi}를 거듭 사용하겠다.
우선 이 보조정리를 완전열
$0 \to \mathfrak p \to R \to R/\mathfrak p \to 0$에 적용하면
$d(R/\mathfrak p) \leq 1$을 얻는다. 그러나 이는 분명 0일 수 없다.
$x\in \mathfrak q$, $x\not \in \mathfrak p$인 원소를 택하자.
짧은 완전열
% P01-E0386
$$
0 \to R/\mathfrak p \to R/\mathfrak p \to R/(xR + \mathfrak p) \to 0.
$$
을 생각하자. 이는
$\chi_{I, R/\mathfrak p} - \chi_{I, R/\mathfrak p}
- \chi_{I, R/(xR + \mathfrak p)} = - \chi_{I, R/(xR + \mathfrak p)}$의
차수가 $ < 1$임을 함의한다.
즉 $d(R/(xR + \mathfrak p)) = 0$이고,
보조정리 \ref{algebra-lemma-dimension-0-d-0}에 의해
$\dim(R/(xR + \mathfrak p)) = 0$이다.
그러나 $R/(xR + \mathfrak p)$은 서로 다른 소 아이디얼
$\mathfrak q/(xR + \mathfrak p)$와
$\mathfrak m/(xR + \mathfrak p)$을 가지므로 원하는 모순을 얻는다.
\end{proof}

\begin{proposition}
\label{algebra-proposition-dimension}
$R$을 뇌터 국소환이라 하자. $d \geq 0$을 정수라 하자.
다음 조건들은 서로 동치이다.
\begin{enumerate}
\item
\label{algebra-item-dim-d}
$\dim(R) = d$,
\item
\label{algebra-item-d-d}
$d(R) = d$,
\item
\label{algebra-item-ideal-d}
$d$개 원소로 생성되는 정의 아이디얼이 존재하고,
$d$개보다 적은 원소로 생성되는 정의 아이디얼은 없다.
\end{enumerate}
\end{proposition}

\begin{proof}
이 증명은 사실 보조정리 \ref{algebra-lemma-height-1}의 증명과 같다.
$d$에 관한 귀납법으로 명제를 증명하겠다.
보조정리 \ref{algebra-lemma-dimension-0-d-0}과 \ref{algebra-lemma-height-1}에 의해
$d > 1$이라고 가정해도 된다.
$R$의 정의 아이디얼의 생성원 최소 개수를 $d'(R)$로 나타내자.
부등식
$\dim(R) \geq d'(R) \geq d(R) \geq \dim(R)$
을 증명하여 이들이 모두 같음을 보이겠다.

\medskip\noindent
먼저 $\dim(R) = d$라고 하자.
$\mathfrak p_i$들을 $R$의 극소 소 아이디얼들이라 하자.
보조정리 \ref{algebra-lemma-Noetherian-irreducible-components}에 따르면
그 수는 유한하다. 따라서 보조정리 \ref{algebra-lemma-silly}에 의해
$x \in \mathfrak m$, $x \not \in \mathfrak p_i$인 원소를 찾을 수 있다.
소 아이디얼의 모든 극대 사슬은 어떤 $\mathfrak p_i$에서 시작하므로
$R/xR$의 차원은 기껏해야 $d-1$이다.
귀납가정에 의해 $R/xR$에서 정의 아이디얼을 생성하는
$x_2, \ldots, x_d$가 존재한다.
따라서 $R$에는 (기껏해야) $d$개 원소로 생성되는 정의 아이디얼이 있다.

\medskip\noindent
$d'(R) = d$라고 하자.
$I = (x_1, \ldots, x_d)$를 정의 아이디얼이라 하자.
$I^n/I^{n + 1}$은 차수 $n$인 $x_1, \ldots, x_d$의 모든 단항식에 의한
곱셈을 통해 $\binom{d + n - 1}{d - 1}$개의 $R/I$의 직합의 몫이다.
따라서 $\text{length}_R(I^n/I^{n + 1})$은
차수 $d-1$인 다항식에 의해 유계이다.
그러므로 $d(R) \leq d$이다.

\medskip\noindent
$d(R) = d$라고 하자.
모든 포함이 진정하고 $e \geq 2$인 소 아이디얼 사슬
$\mathfrak p \subset \mathfrak q \subset
\mathfrak q_2 \subset \ldots \subset \mathfrak q_e = \mathfrak m$
을 생각하자.
어떤 정의 아이디얼 $I \subset R$을 택하자.
보조정리 \ref{algebra-lemma-hilbert-ses-chi}를 거듭 사용하겠다.
우선 이를 완전열
$0 \to \mathfrak p \to R \to R/\mathfrak p \to 0$에 적용하면
$d(R/\mathfrak p) \leq d$를 얻는다. 그러나 이는 분명 0일 수 없다.
$x\in \mathfrak q$, $x\not \in \mathfrak p$인 원소를 택하자.
짧은 완전열
% P01-E0386
$$
0 \to R/\mathfrak p \to R/\mathfrak p \to R/(xR + \mathfrak p) \to 0.
$$
을 생각하자. 이는
$\chi_{I, R/\mathfrak p} - \chi_{I, R/\mathfrak p}
- \chi_{I, R/(xR + \mathfrak p)} = - \chi_{I, R/(xR + \mathfrak p)}$의
차수가 $ < d$임을 함의한다.
즉 $d(R/(xR + \mathfrak p)) \leq d - 1$이고,
귀납가정에 의해
$\dim(R/(xR + \mathfrak p)) \leq d - 1$이다.
이제 $R/(xR + \mathfrak p)$에는 소 아이디얼 사슬
$\mathfrak q/(xR + \mathfrak p) \subset \mathfrak q_2/(xR + \mathfrak p)
\subset \ldots \subset \mathfrak q_e/(xR + \mathfrak p)$이 있으므로
$e - 1 \leq d - 1$을 얻는다.
임의의 소 아이디얼 사슬에서 시작했으므로
$\dim(R) \leq d(R)$임이 증명된다.

\medskip\noindent
앞의 논증을 되짚어 보면 원하는 순환 부등식들을
증명했음을 알 수 있다.
\end{proof}

\noindent
$(R, \mathfrak m)$을 뇌터 국소환이라 하자.
위의 결과에서 $\mathfrak m$은 $\dim(R)$개보다 적은 변수로
생성될 수 없음이 명백하다.
% P01-E0387
나카야마 보조정리 \ref{algebra-lemma-NAK}에 의해
$\mathfrak m$의 생성원 최소 개수는
$\dim_{\kappa(\mathfrak m)} \mathfrak m/\mathfrak m^2$와 같다.
따라서 다음 기본 부등식을 얻는다.
$$
\dim(R) \leq \dim_{\kappa(\mathfrak m)} \mathfrak m/\mathfrak m^2.
$$
등호가 성립하는 환은 좋은 성질을 많이 갖는다는 것이 밝혀진다.
그러한 환을 정칙 국소환이라 한다.

\begin{definition}
\label{algebra-definition-regular-local}
$(R, \mathfrak m)$을 차원 $d$인 뇌터 국소환이라 하자.
\begin{enumerate}
\item {\it $R$의 매개변수계}는
$R$의 정의 아이디얼을 생성하는 원소들의 열
$x_1, \ldots, x_d \in \mathfrak m$이다.
\item $\mathfrak m = (x_1, \ldots, x_d)$가 되게 하는
$x_1, \ldots, x_d \in \mathfrak m$가 존재하면
$R$을 {\it 정칙 국소환}이라 하고,
$x_1, \ldots, x_d$를 {\it 정칙 매개변수계}라 한다.
\end{enumerate}
\end{definition}

\noindent
다음 보조정리들은 위의 보조정리와 명제의 증명에서 명백하지만,
편리하게 참조할 수 있도록 명시한다.

\begin{lemma}
\label{algebra-lemma-minimal-over-1}
$R$을 뇌터 환이라 하자. $x \in R$이라 하자.
\begin{enumerate}
\item $\mathfrak p$가 $(x)$ 위의 극소 소 아이디얼이면
$\mathfrak p$의 높이는 $0$ 또는 $1$이다.
\item $\mathfrak p, \mathfrak q \in \Spec(R)$이고
$\mathfrak q$가 $(\mathfrak p, x)$ 위의 극소 소 아이디얼이면,
$\mathfrak p$와 $\mathfrak q$ 사이에는 소 아이디얼이 엄밀히 존재하지 않는다.
\end{enumerate}
\end{lemma}

\begin{proof}
(1)의 증명. $\mathfrak p$가 $x$ 위에서 극소이면
$x$를 포함하는 $R_\mathfrak p$의 유일한 소 아이디얼은
극대 아이디얼 $\mathfrak p R_\mathfrak p$이다.
이는 $R_\mathfrak p$의 소 아이디얼들이
$\mathfrak p$에 포함되는 $R$의 소 아이디얼들과
$1$-대-$1$로 대응하기 때문이다.
보조정리 \ref{algebra-lemma-spec-localization}을 보라.
따라서 $x$가 $R_\mathfrak p$에서 멱영이 아니면
보조정리 \ref{algebra-lemma-height-1}에 의해
$\dim(R_\mathfrak p) = 1$이다.
물론 $x$가 $R_\mathfrak p$에서 멱영이면,
이 논증은 $\mathfrak pR_\mathfrak p$가 유일한 소 아이디얼임을 보여
높이가 $0$임을 알 수 있다.

\medskip\noindent
(2)의 증명. (1)에 의해 $\mathfrak q/\mathfrak p$는
$R/\mathfrak p$에서 높이가 $1$ 또는 $0$인 소 아이디얼이다.
이는 $\mathfrak p$와 $\mathfrak q$ 사이에
소 아이디얼이 엄밀히 존재할 수 없음을 즉시 함의한다.
\end{proof}

\begin{lemma}
\label{algebra-lemma-minimal-over-r}
$R$을 뇌터 환이라 하자. $f_1, \ldots, f_r \in R$이라 하자.
\begin{enumerate}
\item $\mathfrak p$가 $(f_1, \ldots, f_r)$ 위의 극소 소 아이디얼이면
$\mathfrak p$의 높이는 $\leq r$이다.
\item $\mathfrak p, \mathfrak q \in \Spec(R)$이고
$\mathfrak q$가 $(\mathfrak p, f_1, \ldots, f_r)$ 위의
극소 소 아이디얼이면, $\mathfrak p$와 $\mathfrak q$ 사이의
모든 소 아이디얼 사슬의 길이는 기껏해야 $r$이다.
\end{enumerate}
\end{lemma}

\begin{proof}
(1)의 증명. $\mathfrak p$가 $f_1, \ldots, f_r$ 위에서 극소이면,
$f_1, \ldots, f_r$을 포함하는 $R_\mathfrak p$의 유일한 소 아이디얼은
극대 아이디얼 $\mathfrak p R_\mathfrak p$이다.
이는 $R_\mathfrak p$의 소 아이디얼들이
$\mathfrak p$에 포함되는 $R$의 소 아이디얼들과
$1$-대-$1$로 대응하기 때문이다.
보조정리 \ref{algebra-lemma-spec-localization}을 보라.
따라서 명제 \ref{algebra-proposition-dimension}에 의해
$\dim(R_\mathfrak p) \leq r$이다.

\medskip\noindent
(2)의 증명. (1)에 의해 $\mathfrak q/\mathfrak p$는
높이가 $\leq r$인 소 아이디얼이다.
이는 $\mathfrak p$와 $\mathfrak q$ 사이의 소 아이디얼 사슬에 관한
주장을 즉시 함의한다.
\end{proof}

\begin{lemma}
\label{algebra-lemma-one-equation}
$R$이 뇌터 국소환이고 $x\in \mathfrak m$가
그 극대 아이디얼의 원소라고 하자.
그러면 $\dim R \leq \dim R/xR + 1$이다.
$x$가 $R$의 어느 극소 소 아이디얼에도 포함되지 않으면 등호가 성립한다.
(예를 들어 $x$가 영인자가 아니면 그렇다.)
\end{lemma}

\begin{proof}
$$
x_1, \ldots, x_{\dim R/xR} \in R
$$
의 $R/xR$에서의 상들이
$R/xR$의 정의 아이디얼을 생성하면,
$x, x_1, \ldots, x_{\dim R/xR}$은 $R$의 정의 아이디얼을 생성한다.
따라서 명제 \ref{algebra-proposition-dimension}에 의해 부등식이 성립한다.
한편 $x$가 $R$의 어느 극소 소 아이디얼에도 포함되지 않으면,
$R/xR$의 모든 소 아이디얼 사슬은 $R$에서
극대가 되기까지 적어도 한 단계가 남은 사슬을 준다.
\end{proof}

\begin{lemma}
\label{algebra-lemma-elements-generate-ideal-definition}
$(R, \mathfrak m)$을 뇌터 국소환이라 하자.
$x_1, \ldots, x_d \in \mathfrak m$가 정의 아이디얼을 생성하고
$d = \dim(R)$이라고 하자.
그러면 모든 $i = 1, \ldots, d$에 대해
$\dim(R/(x_1, \ldots, x_i)) = d - i$이다.
\end{lemma}

\begin{proof}
명제 \ref{algebra-proposition-dimension}의 증명에서 바로 따르거나,
$d$에 관한 귀납법과 보조정리 \ref{algebra-lemma-one-equation}을 사용하면 된다.
\end{proof}

\section{차원론의 응용}
\label{algebra-section-applications-dimension-theory}

\noindent
차원에 관한 결과들을 사용하여 어떤 환들의 스펙트럼이 무한함을 증명하고
더 많은 야코브슨 환을 만들어 낼 수 있다.

\begin{lemma}
\label{algebra-lemma-Noetherian-local-domain-dim-2-infinite-opens}
$R$을 차원이 $\geq 2$인 뇌터 국소 정역이라 하자.
공집합이 아닌 열린 부분집합 $U \subset \Spec(R)$은 무한하다.
\end{lemma}

\begin{proof}
모순을 얻기 위해 $U \subset \Spec(R)$이 유한하다고 하자.
이 경우 $(0) \in U$이고 $\{(0)\}$은 $U$의 열린 부분집합이다
($\{(0)\}$의 여집합은 다른 점들의 폐포들의 합집합이기 때문이다).
따라서 $U = \{(0)\}$이라고 가정해도 된다.
$\mathfrak m \subset R$을 극대 아이디얼이라 하자.
$V(x) \cup U = \Spec(R)$가 되게 하는
$x \in \mathfrak m$, $x \not = 0$인 원소를 찾을 수 있다.
즉 $D(x) = \{(0)\}$이다. 특히 보조정리 \ref{algebra-lemma-one-equation}에 의해
$\dim(R/xR) = \dim(R) - 1 \geq 1$임을 알 수 있다.
명제 \ref{algebra-proposition-dimension}에 따라
$\overline{y}_2, \ldots, \overline{y}_{\dim(R)} \in R/xR$가
$R/xR$의 정의 아이디얼을 생성하게 하자.
그 올림 $y_2, \ldots, y_{\dim(R)} \in R$을 택하여
$x, y_2, \ldots, y_{\dim(R)}$가 $R$에서 정의 아이디얼을
생성하게 하자.
보조정리 \ref{algebra-lemma-elements-generate-ideal-definition}에 의해
$\dim(R/(y_2)) = \dim(R) - 1$ 및
$\dim(R/(y_2, x)) = \dim(R) - 2$이다.
따라서 $y_2$를 포함하지만 $x$를 포함하지 않는
소 아이디얼 $\mathfrak p$가 존재한다.
이는 $D(x) = \{(0)\}$라는 사실과 모순이다.
\end{proof}

\noindent
% P01-E0388
$k$가 체일 때의 환 $k[[t]]$ 또는 $p$-진수의 환은
소 아이디얼을 유한 개만 갖는 차원 $1$의 뇌터 환이다.
이 현상이 일어날 수 있는 차원의 최댓값은 이것이다.

\begin{lemma}
\label{algebra-lemma-Noetherian-finite-nr-primes}
소 아이디얼이 유한 개인 뇌터 환의 차원은 $\leq 1$이다.
\end{lemma}

\begin{proof}
$R$을 소 아이디얼이 유한 개인 뇌터 환이라 하자.
$R$이 국소 정역이면 보조정리
\ref{algebra-lemma-Noetherian-local-domain-dim-2-infinite-opens}에서 결론이 따른다.
$R$이 정역이면 국소 경우에 의해 모든 극대 아이디얼
$\mathfrak m$에 대해 $R_\mathfrak m$의 차원이 $\leq 1$이다.
따라서 보조정리 \ref{algebra-lemma-dimension-height}에 의해
$\dim(R) \leq 1$이다.
$R$이 일반적인 경우, $R$의 모든 극소 소 아이디얼
$\mathfrak q$에 대해 $\dim(R/\mathfrak q) \leq 1$이다.
모든 소 아이디얼은 극소 소 아이디얼을 포함하므로
(보조정리 \ref{algebra-lemma-Zariski-topology}),
이는 $\dim(R) \leq 1$을 함의한다.
\end{proof}

\begin{lemma}
\label{algebra-lemma-finite-type-algebra-finite-nr-primes}
$S$를 체 $k$ 위의 0이 아닌 유한형 대수라 하자.
다음 조건들은 서로 동치이다.
\begin{enumerate}
\item $\dim(S) = 0$,
\item $S$는 소 아이디얼을 유한 개 갖는다.
\item $S$는 극대 아이디얼을 유한 개 갖는다.
\item $\Spec(S)$은 보조정리
\ref{algebra-lemma-ring-with-only-minimal-primes}의 동치 조건 중 하나를 만족한다.
\item $\dim_k(S) < \infty$,
\item $S$는 아르틴 환이다.
\item $\Spec(S)$은 이산 위상공간이다.
% P01-E0389
\item 여기에 더 추가한다.
\end{enumerate}
\end{lemma}

\begin{proof}
보조정리 \ref{algebra-lemma-ring-with-only-minimal-primes}의 (5)를 보면
(1)과 (4)가 동치임은 정의에서 바로 따른다.
$\Spec(S)$은 소버이고, 뇌터이며, 야코브슨임을 상기하자.
보조정리 \ref{algebra-lemma-spec-spectral}, \ref{algebra-lemma-Noetherian-topology},
\ref{algebra-lemma-finite-type-field-Jacobson}, \ref{algebra-lemma-jacobson}을 보라.
$S$의 차원이 $0$이면 각 점은 기약 성분 하나를 정하며,
기약 성분은 유한 개뿐이다
(위상수학, 보조정리 \ref{topology-lemma-Noetherian}).
따라서 (1)은 (2)를 함의한다. (2)가 (3)을 함의하는 것은 자명하다.
(3)이 성립하면 위상수학, 보조정리
\ref{topology-lemma-finite-jacobson}에 의해 $\Spec(S)$은 이산이고,
따라서 $S$의 차원은 $0$이다.

\medskip\noindent
이제 (1) -- (4)가 서로 동치임을 안다.
함의 (5) $\Rightarrow$ (6)은 보조정리
\ref{algebra-lemma-finite-dimensional-algebra}이다.
함의 (6) $\Rightarrow$ (7)은 명제
\ref{algebra-proposition-dimension-zero-ring}에서 따르고,
함의 (7) $\Rightarrow$ (4)는 자명하다.
역으로 $S$가 (1) -- (4)를 만족하면
$S$는 모두 극대인 유한 개의 소 아이디얼
$\mathfrak m_1, \ldots, \mathfrak m_r$을 갖는다.
힐베르트 영점정리(정리 \ref{algebra-theorem-nullstellensatz})에 의해
$\kappa(\mathfrak m_i)$는 $k$의 유한 확대이다.
명제 \ref{algebra-proposition-dimension-zero-ring}에 의해
$S$가 아르틴 환임도 알 수 있다.
다음으로 보조정리 \ref{algebra-lemma-artinian-finite-length}는
$\text{length}_S(S) < \infty$임을 알려 준다.
따라서 보조정리 \ref{algebra-lemma-pushdown-module}에 의해
$\dim_k(S) < \infty$이다.
결론적으로 (1) -- (7)은 서로 동치이다.
(주의: $\dim(S) = 0 \Rightarrow \dim_k(S) < \infty$를
증명하는 또 하나의 더 표준적인 방법은 뇌터 정규화를 사용하는 것이지만,
이 단계에서는 아직 쓸 수 없다.)
\end{proof}

\begin{lemma}
\label{algebra-lemma-noetherian-dim-1-Jacobson}
뇌터 야코브슨 환에 관하여 다음이 성립한다.
\begin{enumerate}
\item 차원이 $1$이고 소 아이디얼이 무한 개인 임의의 뇌터 정역
$R$은 야코브슨 환이다.
\item 모든 소 아이디얼 $\mathfrak p$가 극대이거나
무한히 많은 소 아이디얼에 포함되는 임의의 뇌터 환은 야코브슨 환이다.
\end{enumerate}
\end{lemma}

\begin{proof}
(1)은 보조정리 \ref{algebra-lemma-pid-jacobson}을 달리 표현한 것이다.

\medskip\noindent
$R$을 모든 비극대 소 아이디얼 $\mathfrak p$가
무한히 많은 소 아이디얼에 포함되는 뇌터 환이라 하자.
모순을 얻기 위해 $\Spec(R)$이 야코브슨이 아니라고 하자.
보조정리 \ref{algebra-lemma-irreducible}과
\ref{algebra-lemma-Noetherian-topology}에 의해
$\Spec(R)$은 소버 뇌터 위상공간이다.
위상수학, 보조정리
\ref{topology-lemma-non-jacobson-Noetherian-characterize}에 의해
$\{\mathfrak p\}$가 $\Spec(R)$의 국소 닫힌 부분집합이 되는
% P01-E0390
비극대 아이디얼 $\mathfrak p \subset R$이 존재한다.
즉 $\mathfrak p$는 극대가 아니고
$\{\mathfrak p\}$는 $V(\mathfrak p)$의 열린 부분집합이다.
% P01-E0391
$\mathfrak p \subset \mathfrak q$인 소 아이디얼
$\mathfrak q \subset R$을 생각하자.
$(R/\mathfrak p)_{\mathfrak q} =
R_{\mathfrak q}/\mathfrak pR_{\mathfrak q}$의 스펙트럼의 위상은
$\Spec(R)$의 위상에서 유도됨을 상기하자.
보조정리 \ref{algebra-lemma-spec-localization}과
\ref{algebra-lemma-spec-closed}를 보라.
따라서 $\{(0)\}$은
$\Spec((R/\mathfrak p)_{\mathfrak q})$의 국소 닫힌 부분집합이다.
보조정리 \ref{algebra-lemma-Noetherian-local-domain-dim-2-infinite-opens}에 의해
$\dim((R/\mathfrak p)_{\mathfrak q}) = 1$이다.
이것이 모든 $\mathfrak q \supset \mathfrak p$에 대해 성립하므로
$\dim(R/\mathfrak p) = 1$이다.
이제 $\mathfrak p$가 무한히 많은 소 아이디얼에 포함된다는
가정에 의해 $\Spec(R/\mathfrak p)$이 무한임을 알 수 있다.
따라서 보조정리의 (1)에 의해
$V(\mathfrak p) \cong \Spec(R/\mathfrak p)$은
그 닫힌 점들의 폐포이다.
이는 $\{\mathfrak p\} \subset V(\mathfrak p)$가
열릴 수 없다는 뜻이므로 원하는 모순이다.
\end{proof}

\section{가군의 지지집합과 차원}
\label{algebra-section-support}

\noindent
가군의 지지집합과 차원에 관한 몇 가지 기본 결과를 제시한다.

\begin{lemma}
\label{algebra-lemma-filter-Noetherian-module}
$R$을 뇌터 환이라 하고 $M$을 유한 $R$-가군이라 하자.
다음과 같은 $R$-부분가군에 의한 여과가 존재한다.
$$
0 = M_0 \subset M_1 \subset \ldots \subset M_n = M
$$
각 몫 $M_i/M_{i-1}$은 $R$의 어떤 소 아이디얼 $\mathfrak p_i$에 대한
$R/\mathfrak p_i$와 동형이다.
\end{lemma}

\begin{proof}[첫째 증명]
보조정리 \ref{algebra-lemma-trivial-filter-finite-module}에 의해
어떤 아이디얼 $I$에 대해 $M = R/I$인 경우만 다루면 충분하다.
가군 $R/J$에 대해 이 보조정리가 성립하지 않는 아이디얼 $J$들의
집합을 $S$라 하고 포함관계로 순서를 주자.
모순을 얻기 위해 $S$가 공집합이 아니라고 가정하자.
$R$이 뇌터 환이므로 $S$에는 극대원 $J$가 있다.
$S$의 정의에 의해 아이디얼 $J$는 소 아이디얼일 수 없다.
$ab \in J$이지만 $a \in J$도 아니고 $b\in J$도 아닌
$a, b\in R$을 택하자. 여과
$0 \subset aR/(J \cap aR) \subset R/J$를 생각하자.
부분가군 $aR/(J \cap aR)$과 몫가군
$(R/J)/(aR/(J \cap aR))$은 모두 순환가군이다.
이를 각각 $R/J'$와 $R/J''$로 쓰면 짧은 완전열
$0 \to R/J' \to R/J \to R/J'' \to 0$을 얻는다.
$b \in J'$이므로 포함 $J \subset J'$은 진정하고,
$a \in J''$이므로 포함 $J \subset J''$도 진정하다.
따라서 $J$의 극대성에 의해 $R/J'$과 $R/J''$은 모두
위와 같은 여과를 가지며, 그러므로 $R/J$도 그러하다. 모순이다.
\end{proof}

\begin{proof}[둘째 증명]
$R$-가군 $M$에 대해 보조정리의 명제와 같은 여과가 존재할 때
$P(M)$이 성립한다고 하자. $P$는 확대에 대해 안정적이고
$0$에 대해 성립한다. 보조정리
\ref{algebra-lemma-trivial-filter-finite-module}에 의해 모든 아이디얼 $I$에 대해
$P(R/I)$가 성립함을 보이면 충분하다.
그렇지 않다면 $R$이 뇌터 환이므로 극대
% P01-E0392
반례 $J$가 존재한다. 예 \ref{algebra-example-oka-family-property-modules}와
명제 \ref{algebra-proposition-oka}에 의해 아이디얼 $J$는 소 아이디얼이고,
이는 모순이다.
\end{proof}

\begin{lemma}
\label{algebra-lemma-filter-primes-in-support}
% P01-E1382
$R$, $M$, $M_i$, $\mathfrak p_i$를 보조정리
\ref{algebra-lemma-filter-Noetherian-module}에서와 같이 두자.
그러면 $\text{Supp}(M) = \bigcup V(\mathfrak p_i)$이고,
특히 $\mathfrak p_i \in \text{Supp}(M)$이다.
\end{lemma}

\begin{proof}
이는 보조정리 \ref{algebra-lemma-support-closed}와
\ref{algebra-lemma-support-quotient}에서 따른다.
\end{proof}

\begin{lemma}
\label{algebra-lemma-support-point}
$R$을 극대 아이디얼 $\mathfrak m$을 갖는 뇌터 국소환이라 하자.
$M$을 영이 아닌 유한 $R$-가군이라 하자.
그러면 $\text{Supp}(M) = \{ \mathfrak m\}$인 것과
$M$이 $R$ 위에서 유한 길이를 갖는 것은 동치이다.
\end{lemma}

\begin{proof}
$\text{Supp}(M) = \{ \mathfrak m\}$이라고 하자.
보조정리 \ref{algebra-lemma-filter-Noetherian-module}의 여과에 나타나는
모든 소 아이디얼 $\mathfrak p_i$가 극대 아이디얼임을 보이면 충분하다.
이는 보조정리 \ref{algebra-lemma-filter-primes-in-support}에 의해 자명하다.

\medskip\noindent
$M$이 $R$ 위에서 유한 길이를 갖는다고 하자.
보조정리 \ref{algebra-lemma-length-infinite}에 의해 $\mathfrak m^n M = 0$이다.
$R$의 임의의 소 아이디얼 $\mathfrak p \not = \mathfrak m$에 대해
$\mathfrak m$의 어떤 원소가 $R_{\mathfrak p}$에서 가역원으로
보내지므로 $M_{\mathfrak p} = 0$이다.
\end{proof}

\begin{lemma}
\label{algebra-lemma-Noetherian-power-ideal-kills-module}
$R$을 뇌터 환이라 하자.
$I \subset R$을 아이디얼이라 하자.
$M$을 유한 $R$-가군이라 하자.
그러면 어떤 $n \geq 0$에 대해 $I^nM = 0$인 것과
$\text{Supp}(M) \subset V(I)$인 것은 동치이다.
\end{lemma}

\begin{proof}
실제로 $I^nM = 0$인 것은 $I^n \subset \text{Ann}(M)$인 것과 동치이다.
$R$이 뇌터 환이므로 이는 보조정리
\ref{algebra-lemma-Noetherian-power}에 의해
$I \subset \sqrt{\text{Ann}(M)}$인 것과 동치이다.
다시 보조정리 \ref{algebra-lemma-Zariski-topology}에 의해 이는
$V(I) \supset V(\text{Ann}(M))$인 것과 동치이다.
보조정리 \ref{algebra-lemma-support-closed}에 의해 후자는
$V(I) \supset \text{Supp}(M)$인 것과 동치이다.
\end{proof}

\begin{lemma}
\label{algebra-lemma-filter-minimal-primes-in-support}
% P01-E1382
$R$, $M$, $M_i$, $\mathfrak p_i$를 보조정리
\ref{algebra-lemma-filter-Noetherian-module}에서와 같이 두자.
집합 $\{\mathfrak p_i\}$의 극소원들은
$\text{Supp}(M)$의 극소원들이다.
극소 소 아이디얼 $\mathfrak p$가 나타나는 횟수는
$$
\#\{i \mid \mathfrak p_i = \mathfrak p\}
=
\text{length}_{R_\mathfrak p} M_{\mathfrak p}.
$$
이다.
\end{lemma}

\begin{proof}
첫째 명제는
$\text{Supp}(M) = \bigcup V(\mathfrak p_i)$에서 따른다.
보조정리 \ref{algebra-lemma-filter-primes-in-support}를 보라.
$\mathfrak p \in \text{Supp}(M)$이 극소라고 하자.
$M_{\mathfrak p}$의 지지집합은 극대 아이디얼
$\mathfrak p R_{\mathfrak p}$ 하나로 이루어진 집합이다.
따라서 보조정리 \ref{algebra-lemma-support-point}에 의해
$M_{\mathfrak p}$의 길이는 유한하고 $> 0$이다.
다음으로 $M_{\mathfrak p}$에는 다음 부분몫들을 갖는 여과가 있다.
$
(R/\mathfrak p_i)_{\mathfrak p}
=
R_{\mathfrak p}/{\mathfrak p_i}R_{\mathfrak p}
$.
$\mathfrak p_i \not \subset \mathfrak p$이면 이들은 영이고,
$\mathfrak p_i \subset \mathfrak p$이면
$\mathfrak p$의 극소성에 의해 이 경우
$\mathfrak p_i = \mathfrak p$이므로 $\kappa(\mathfrak p)$와 같다.
$\kappa(\mathfrak p)$의 길이가 $1$이므로 결론이 따른다.
\end{proof}

\begin{lemma}
\label{algebra-lemma-support-dimension-d}
$R$을 뇌터 국소환이라 하자.
$M$을 유한 $R$-가군이라 하자.
그러면 $d(M) = \dim(\text{Supp}(M))$이다.
여기서 $d(M)$은 정의 \ref{algebra-definition-d}에서와 같다.
\end{lemma}

\begin{proof}
$M_i, \mathfrak p_i$를 보조정리
\ref{algebra-lemma-filter-Noetherian-module}에서와 같이 두자.
보조정리 \ref{algebra-lemma-hilbert-ses-chi}에 의해 등식
$d(M) = \max \{ d(R/\mathfrak p_i) \}$을 얻는다.
명제 \ref{algebra-proposition-dimension}에 의해
$d(R/\mathfrak p_i) = \dim(R/\mathfrak p_i)$이다.
$\dim(R/\mathfrak p_i) = \dim V(\mathfrak p_i)$임은 자명하다.
$\text{Supp}(M)$의 모든 극소 소 아이디얼이
$\mathfrak p_i$들 가운데 나타나므로
(보조정리 \ref{algebra-lemma-filter-minimal-primes-in-support}) 결론이 따른다.
\end{proof}

\begin{lemma}
\label{algebra-lemma-ses-dimension}
$R$을 뇌터 환이라 하자. 유한 $R$-가군들의 짧은 완전열
$0 \to M' \to M \to M'' \to 0$을 두자. 그러면
$\max\{\dim(\text{Supp}(M')), \dim(\text{Supp}(M''))\} =
\dim(\text{Supp}(M))$이다.
\end{lemma}

\begin{proof}
$R$이 국소환이면 이는 보조정리
\ref{algebra-lemma-support-dimension-d}와 \ref{algebra-lemma-hilbert-ses-chi}에서
곧바로 따른다.
$R$이 국소환이 아닐 때에도 통하는 더 초등적인 논증은
$\text{Supp}(M')$, $\text{Supp}(M'')$, $\text{Supp}(M)$가
닫혀 있다는 사실(보조정리 \ref{algebra-lemma-support-closed})과
$\text{Supp}(M) = \text{Supp}(M') \cup \text{Supp}(M'')$라는
사실(보조정리 \ref{algebra-lemma-support-quotient})을 사용하는 것이다.
\end{proof}

\section{연관 소 아이디얼}
\label{algebra-section-ass}

\noindent
다음은 표준 정의이다. 뇌터가 아닌 환이나 유한이 아닌 가군에는
절 \ref{algebra-section-weakly-ass}의 정의가 더 적절할 수 있다.

\begin{definition}
\label{algebra-definition-associated}
$R$을 환이라 하고 $M$을 $R$-가군이라 하자.
$R$의 소 아이디얼 $\mathfrak p$가 $M$에 {\it 연관된다}는 것은
소멸자가 $\mathfrak p$인 원소 $m \in M$이 존재한다는 뜻이다.
그러한 모든 소 아이디얼의 집합을 $\text{Ass}_R(M)$ 또는
$\text{Ass}(M)$로 나타낸다.
\end{definition}

\begin{lemma}
\label{algebra-lemma-ass-support}
$R$을 환이라 하고 $M$을 $R$-가군이라 하자.
그러면 $\text{Ass}(M) \subset \text{Supp}(M)$이다.
\end{lemma}

\begin{proof}
$m \in M$의 소멸자가 $\mathfrak p$이면, 특히
$R \setminus \mathfrak p$의 어떤 원소도 $m$을 소멸시키지 않는다.
따라서 $m$은 $M_{\mathfrak p}$의 영이 아닌 원소이다.
즉 $\mathfrak p \in \text{Supp}(M)$이다.
\end{proof}

\begin{lemma}
\label{algebra-lemma-ass}
$R$을 환이라 하자. $R$-가군들의 짧은 완전열
$0 \to M' \to M \to M'' \to 0$을 두자. 그러면
$\text{Ass}(M') \subset \text{Ass}(M)$이고
$\text{Ass}(M) \subset \text{Ass}(M') \cup \text{Ass}(M'')$이다.
또한 $\text{Ass}(M' \oplus M'') = \text{Ass}(M') \cup \text{Ass}(M'')$이다.
\end{lemma}

\begin{proof}
$m' \in M'$이면 $m'$을 $M'$의 원소로 볼 때의 소멸자는
$m'$을 $M$의 원소로 볼 때의 소멸자와 같다.
따라서 포함 $\text{Ass}(M') \subset \text{Ass}(M)$을 얻는다.
소멸자가 소 아이디얼 $\mathfrak p$인 원소 $m \in M$을 택하자.
$m' = gm \in M'$가 되는 $g \in R$, $g \not \in \mathfrak p$가
존재하면 $m'$의 소멸자는 $\mathfrak p$이다.
그러한 $g \in R$, $g \not \in \mathfrak p$ 가운데
$gm \in M'$인 것이 존재하지 않으면,
% P01-E0393
$m$의 상 $m'' \in M''$의 소멸자는 $\mathfrak p$이다.
이는 포함
$\text{Ass}(M) \subset \text{Ass}(M') \cup \text{Ass}(M'')$을 증명한다.
마지막 명제의 증명은 생략한다.
\end{proof}

\begin{lemma}
\label{algebra-lemma-ass-filter}
$R$을 환, $M$을 $R$-가군이라 하자.
다음과 같은 $R$-부분가군에 의한 여과가 존재한다고 하자.
$$
0 = M_0 \subset M_1 \subset \ldots \subset M_n = M
$$
각 몫 $M_i/M_{i-1}$은 $R$의 어떤 소 아이디얼 $\mathfrak p_i$에 대해
$R/\mathfrak p_i$와 동형이라고 하자.
그러면 $\text{Ass}(M) \subset \{\mathfrak p_1, \ldots, \mathfrak p_n\}$이다.
\end{lemma}

\begin{proof}
여과 $\{ M_i \}$의 길이 $n$에 관해 귀납한다.
소멸자가 소 아이디얼 $\mathfrak p$인 $m \in M$을 택하자.
$m \in M_{n-1}$이면 귀납가정으로 끝난다.
그렇지 않으면 $m$은
$M/M_{n-1} \cong R/\mathfrak p_n$의 영이 아닌 원소로 보내진다.
따라서 $\mathfrak p \subset \mathfrak p_n$이다.
등호가 성립하지 않으면
$f \in \mathfrak p_n$, $f \not\in \mathfrak p$인 원소를 택할 수 있다.
이 경우 $fm$의 소멸자는 여전히 $\mathfrak p$이고
$fm \in M_{n-1}$이다. 따라서 귀납가정으로 결론을 얻는다.
\end{proof}

\begin{lemma}
\label{algebra-lemma-finite-ass}
$R$을 뇌터 환이라 하자.
$M$을 유한 $R$-가군이라 하자.
그러면 $\text{Ass}(M)$은 유한집합이다.
\end{lemma}

\begin{proof}
보조정리 \ref{algebra-lemma-ass-filter}와
보조정리 \ref{algebra-lemma-filter-Noetherian-module}에서 곧바로 따른다.
\end{proof}

\begin{proposition}
\label{algebra-proposition-minimal-primes-associated-primes}
$R$을 뇌터 환이라 하자.
$M$을 유한 $R$-가군이라 하자.
다음 소 아이디얼들의 집합은 모두 같다.
\begin{enumerate}
\item $M$의 지지집합에 있는 극소 소 아이디얼들.
\item $\text{Ass}(M)$의 극소 소 아이디얼들.
\item $M_i/M_{i-1} \cong R/\mathfrak p_i$인 임의의 여과
$0 = M_0 \subset M_1 \subset \ldots
\subset M_{n-1} \subset M_n = M$에 대해,
집합 $\{\mathfrak p_i\}$의 극소 소 아이디얼들.
\end{enumerate}
\end{proposition}

\begin{proof}
(3)과 같은 여과를 택하자.
보조정리 \ref{algebra-lemma-filter-minimal-primes-in-support}에서
(1)과 (3)의 집합이 같음을 보였다.

\medskip\noindent
$\mathfrak p$를 집합 $\{\mathfrak p_i\}$의 극소원이라 하자.
$\mathfrak p = \mathfrak p_i$가 되는 최소의 $i$를 택하자.
$m \in M_i$, $m \not \in M_{i-1}$을 택하자.
$m$의 소멸자는 $\mathfrak p_i = \mathfrak p$에 포함되고
$\mathfrak p_1 \mathfrak p_2 \ldots \mathfrak p_i$를 포함한다.
$i$와 $\mathfrak p$의 선택에 의해 $j < i$이면
$\mathfrak p_j \not \subset \mathfrak p$이고, 따라서
$\mathfrak p_1 \mathfrak p_2 \ldots \mathfrak p_{i - 1}
\not \subset \mathfrak p_i$이다.
$f \in \mathfrak p_1 \mathfrak p_2 \ldots \mathfrak p_{i - 1}$,
$f \not \in \mathfrak p$를 택하자.
그러면 $fm$의 소멸자는 $\mathfrak p$이다.
이로부터 $\mathfrak p$가 $M$의 연관 소 아이디얼임을 알 수 있다.
보조정리 \ref{algebra-lemma-ass-support}에 의해
$\text{Ass}(M) \subset \text{Supp}(M)$이고, 따라서
$\mathfrak p$는 $\text{Ass}(M)$에서 극소이다.
그러므로 (1)의 소 아이디얼 집합은 (2)의 소 아이디얼 집합에 포함된다.

\medskip\noindent
$\mathfrak p$를 $\text{Ass}(M)$의 극소원이라 하자.
$\text{Ass}(M) \subset \text{Supp}(M)$이므로
$\mathfrak q \subset \mathfrak p$인
$\text{Supp}(M)$의 극소원 $\mathfrak q$가 존재한다.
방금 $\mathfrak q \in \text{Ass}(M)$임을 보였다.
$\mathfrak p$의 극소성에 의해 $\mathfrak q = \mathfrak p$이다.
따라서 (2)의 소 아이디얼 집합은 (1)의 소 아이디얼 집합에 포함된다.
\end{proof}

\begin{lemma}
\label{algebra-lemma-ass-zero}
\begin{slogan}
뇌터 환 위의 영이 아닌 각 가군은 연관 소 아이디얼을 갖는다.
\end{slogan}
$R$을 뇌터 환이라 하고 $M$을 $R$-가군이라 하자.
그러면
$$
M = (0) \Leftrightarrow \text{Ass}(M) = \emptyset.
$$
\end{lemma}

\begin{proof}
$M = (0)$이면 정의에 의해 $\text{Ass}(M) = \emptyset$이다.
$M \not = 0$이면 영이 아닌 임의의 유한 생성 부분가군
$M' \subset M$을 택하자. 예를 들어 영이 아닌 원소 하나로 생성되는
부분가군을 택하면 된다.
보조정리 \ref{algebra-lemma-support-zero}에 의해
$\text{Supp}(M')$은 공집합이 아니다.
명제 \ref{algebra-proposition-minimal-primes-associated-primes}에 의해
$\text{Ass}(M')$도 공집합이 아니다.
보조정리 \ref{algebra-lemma-ass}에 의해
$\text{Ass}(M) \not = \emptyset$이다.
\end{proof}

\begin{lemma}
\label{algebra-lemma-ass-minimal-prime-support}
$R$을 뇌터 환이라 하자.
$M$을 $R$-가군이라 하자.
$\text{Supp}(M)$의 원소들 가운데 극소인
임의의 $\mathfrak p \in \text{Supp}(M)$은
$\text{Ass}(M)$의 원소이다.
\end{lemma}

\begin{proof}
$M$이 유한 $R$-가군이면 이는 명제
\ref{algebra-proposition-minimal-primes-associated-primes}의 귀결이다.
일반적으로 이를 그 유한 부분가군들의 합집합
$M = \bigcup M_\lambda$로 쓰고,
$\text{Supp}(M) = \bigcup \text{Supp}(M_\lambda)$ 및
$\text{Ass}(M) = \bigcup \text{Ass}(M_\lambda)$를 사용하면 된다.
\end{proof}

\begin{lemma}
\label{algebra-lemma-ass-zero-divisors}
$R$을 뇌터 환이라 하자.
$M$을 $R$-가군이라 하자.
합집합 $\bigcup_{\mathfrak q \in \text{Ass}(M)} \mathfrak q$는
$M$ 위에서 영인자인 $R$의 원소들의 집합이다.
\end{lemma}

\begin{proof}
어떤 연관 소 아이디얼에 속한 원소든 자명하게 $M$ 위에서 영인자이다.
역으로 $x \in R$이 $M$ 위에서 영인자라고 하자.
부분가군 $N = \{m \in M \mid xm = 0\}$을 생각하자.
$N$은 영이 아니므로 보조정리 \ref{algebra-lemma-ass-zero}에 의해
연관 소 아이디얼 $\mathfrak q$를 갖는다.
그러면 $x \in \mathfrak q$이고, 보조정리 \ref{algebra-lemma-ass}에 의해
$\mathfrak q$는 $M$의 연관 소 아이디얼이다.
\end{proof}

\begin{lemma}
\label{algebra-lemma-one-equation-module}
% P01-E0394
$R$을 뇌터 국소환, $M$을 유한 $R$-가군,
$f \in \mathfrak m$을 $R$의 극대 아이디얼의 원소라 하자. 그러면
$$
\dim(\text{Supp}(M/fM)) \leq
\dim(\text{Supp}(M)) \leq
\dim(\text{Supp}(M/fM)) + 1
$$
$f$가 $M$의 지지집합의 어느 극소 소 아이디얼에도 속하지 않으면
(예를 들어 $f$가 $M$ 위에서 영인자가 아니면)
오른쪽 부등식에서 등호가 성립한다.
\end{lemma}

\begin{proof}
(괄호 안의 명제는 보조정리 \ref{algebra-lemma-ass-zero-divisors}에서 따른다.)
첫째 부등식은
$\text{Supp}(M/fM) \subset \text{Supp}(M)$에서 따른다.
보조정리 \ref{algebra-lemma-support-quotient}를 보라.
둘째 부등식을 위해
$\text{Supp}(M/fM) = \text{Supp}(M) \cap V(f)$임을 주목하자.
보조정리 \ref{algebra-lemma-support-quotient}를 보라.
예를 들어 보조정리 \ref{algebra-lemma-filter-primes-in-support}와
차원의 초등적 성질들에 의해, $R$의 소 아이디얼 $\mathfrak p$에 대해
$\dim V(\mathfrak p) \leq \dim (V(\mathfrak p) \cap V(f)) + 1$임을
보이면 충분하다. 이는 보조정리 \ref{algebra-lemma-one-equation}의 귀결이다.
끝으로 $f$가 $M$의 지지집합의 어느 극소 소 아이디얼에도
포함되지 않으면, $\text{Supp}(M/fM)$ 안의 모든 소 아이디얼 사슬은
$\text{Supp}(M)$에서 극대 사슬이 되기까지 적어도 한 단계가 남은
사슬을 낳는다.
\end{proof}

\begin{lemma}
\label{algebra-lemma-ass-functorial}
$\varphi : R \to S$를 환 준동형사상이라 하자.
$M$을 $S$-가군이라 하자.
그러면 $\Spec(\varphi)(\text{Ass}_S(M)) \subset \text{Ass}_R(M)$이다.
\end{lemma}

\begin{proof}
$\mathfrak q \in \text{Ass}_S(M)$이면, $S$에서 $m$의 소멸자가
$\mathfrak q$인 $M$의 원소 $m$이 존재한다.
그러면 $R$에서 $m$의 소멸자는 $\mathfrak q \cap R$이다.
\end{proof}

\begin{remark}
\label{algebra-remark-ass-reverse-functorial}
$\varphi : R \to S$를 환 준동형사상이라 하자.
$M$을 $S$-가군이라 하자.
이때
$\Spec(\varphi)(\text{Ass}_S(M)) \supset \text{Ass}_R(M)$가
항상 성립하는 것은 아니다.
예를 들어 환 준동형사상
$R = k \to S = k[x_1, x_2, x_3, \ldots]/(x_i^2)$와 $M = S$를 생각하자.
그러면 $\text{Ass}_R(M)$은 공집합이 아니지만
$\text{Ass}_S(S)$는 공집합이다.
\end{remark}

\begin{lemma}
\label{algebra-lemma-ass-functorial-Noetherian}
$\varphi : R \to S$를 환 준동형사상이라 하자.
$M$을 $S$-가군이라 하자. $S$가 뇌터 환이면
$\Spec(\varphi)(\text{Ass}_S(M)) = \text{Ass}_R(M)$이다.
\end{lemma}

\begin{proof}
보조정리 \ref{algebra-lemma-ass-functorial}에서 이미
$\Spec(\varphi)(\text{Ass}_S(M)) \subset \text{Ass}_R(M)$임을 보였다.
역포함을 위해 소 아이디얼
$\mathfrak p \in \text{Ass}_R(M)$을 택하자.
$R$에서 $m$의 소멸자가 $\mathfrak p$인 원소 $m \in M$을 택하자.
$I = \{g \in S \mid gm = 0\}$를 $S$에서 $m$의 소멸자라 하자.
그러면 $R/\mathfrak p \subset S/I$는 단사이다.
보조정리 \ref{algebra-lemma-injective-minimal-primes-in-image}와
\ref{algebra-lemma-minimal-prime-image-minimal-prime}를 결합하면
$I$ 위에서 극소이고 $\mathfrak p$로 가는 소 아이디얼
$\mathfrak q \subset S$가 존재함을 알 수 있다.
명제 \ref{algebra-proposition-minimal-primes-associated-primes}에 의해
$\mathfrak q$는 $S/I$의 연관 소 아이디얼이고,
따라서 보조정리 \ref{algebra-lemma-ass}에 의해
$\mathfrak q$는 $M$의 연관 소 아이디얼이다. 이로써 결론을 얻는다.
\end{proof}

\begin{lemma}
\label{algebra-lemma-ass-quotient-ring}
$R$을 환이라 하자.
$I$를 아이디얼이라 하자.
$M$을 $R/I$-가군이라 하자.
표준 단사
$\Spec(R/I) \to \Spec(R)$를 통해
$\text{Ass}_{R/I}(M) = \text{Ass}_R(M)$이다.
\end{lemma}

\begin{proof}
생략한다.
\end{proof}

\begin{lemma}
\label{algebra-lemma-associated-primes-localize}
$R$을 환이라 하자.
$M$을 $R$-가군이라 하자.
$\mathfrak p \subset R$을 소 아이디얼이라 하자.
\begin{enumerate}
\item $\mathfrak p \in \text{Ass}(M)$이면
$\mathfrak pR_{\mathfrak p} \in \text{Ass}(M_{\mathfrak p})$이다.
\item $\mathfrak p$가 유한 생성이면 그 역도 성립한다.
\end{enumerate}
\end{lemma}

\begin{proof}
$\mathfrak p \in \text{Ass}(M)$이면 소멸자가 $\mathfrak p$인
원소 $m \in M$이 존재한다.
국소화가 완전하므로(명제 \ref{algebra-proposition-localization-exact}),
$M_{\mathfrak p}$에서 $m/1$의 소멸자는
$\mathfrak pR_{\mathfrak p}$이다. 따라서 (1)이 성립한다.
$\mathfrak pR_{\mathfrak p} \in \text{Ass}(M_{\mathfrak p})$이고
$\mathfrak p = (f_1, \ldots, f_n)$이라고 하자.
$M_{\mathfrak p}$에서 소멸자가
$\mathfrak pR_{\mathfrak p}$인 원소 $m/g$를 택하자.
이는 $m$의 소멸자가 $\mathfrak p$에 포함됨을 뜻한다.
$M_{\mathfrak p}$에서 $f_i m/g = 0$이므로
$g_i f_i m = 0$가 $M$에서 성립하게 하는
$g_i \in R$, $g_i \not \in \mathfrak p$가 존재한다.
이를 종합하면 $g_1\ldots g_nm$의 소멸자는 $\mathfrak p$이다.
따라서 $\mathfrak p \in \text{Ass}(M)$이다.
\end{proof}

\begin{lemma}
\label{algebra-lemma-localize-ass}
$R$을 환이라 하고 $M$을 $R$-가군이라 하자.
$S \subset R$을 곱셈적 부분집합이라 하자.
표준 단사 $\Spec(S^{-1}R) \to \Spec(R)$를 통해 다음이 성립한다.
\begin{enumerate}
\item $\text{Ass}_R(S^{-1}M) = \text{Ass}_{S^{-1}R}(S^{-1}M)$,
\item
$\text{Ass}_R(M) \cap \Spec(S^{-1}R) \subset \text{Ass}_R(S^{-1}M)$,
\item $R$이 뇌터 환이면 이 포함은 등호이다.
\end{enumerate}
\end{lemma}

\begin{proof}
$m \in S^{-1}M$에 대해 $I \subset R$과 $J \subset S^{-1}R$를
$m$의 소멸자들이라 하자.
그러면 $I$는 준동형사상 $R \to S^{-1}R$에 의한 $J$의 역상이고
$J = S^{-1}I$이다.
(1)의 등식은 보조정리 \ref{algebra-lemma-spec-localization}에 있는
$\Spec(S^{-1}R) \to \Spec(R)$의 기술에서 따른다.
$m \in M$에 대해 $I \subset R$을 $R$에서 $m$의 소멸자라 하고,
$J \subset S^{-1}R$을 $m/1 \in S^{-1}M$의 소멸자라 하자.
$J = S^{-1}I$이므로 (2)가 따른다.
뇌터인 경우의 등식은 보조정리
\ref{algebra-lemma-associated-primes-localize}에서 따른다. 실제로
% P01-E0395
$\mathfrak p \in R$, $S \cap \mathfrak p = \emptyset$이면
$M_{\mathfrak p} = (S^{-1}M)_{S^{-1}\mathfrak p}$이다.
\end{proof}

\begin{lemma}
\label{algebra-lemma-localize-ass-nonzero-divisors}
$R$을 환이라 하고 $M$을 $R$-가군이라 하자.
$S \subset R$을 곱셈적 부분집합이라 하자.
모든 $s \in S$가 $M$ 위에서 영인자가 아니라고 가정하자.
그러면
$$
\text{Ass}_R(M) = \text{Ass}_R(S^{-1}M).
$$
\end{lemma}

\begin{proof}
가정에 의해 $M \subset S^{-1}M$이므로 보조정리
\ref{algebra-lemma-ass}에서 포함
$\text{Ass}(M) \subset \text{Ass}(S^{-1}M)$을 얻는다.
역으로 소멸자가 소 아이디얼 $\mathfrak p$인
원소 $n/s \in S^{-1}M$이 있다고 하자.
그러면 $n \in M$의 소멸자도 $\mathfrak p$이다.
\end{proof}

\begin{lemma}
\label{algebra-lemma-ideal-nonzerodivisor}
$R$을 극대 아이디얼 $\mathfrak m$을 갖는 뇌터 국소환이라 하자.
$I \subset \mathfrak m$을 아이디얼이라 하고
$M$을 유한 $R$-가군이라 하자.
다음은 동치이다.
\begin{enumerate}
\item $M$ 위에서 영인자가 아닌 $x \in I$가 존재한다.
\item 모든 $\mathfrak q \in \text{Ass}(M)$에 대해
$I \not \subset \mathfrak q$이다.
\end{enumerate}
\end{lemma}

\begin{proof}
$I$ 안에 영인자가 아닌 $x$가 존재하면,
$x$는 자명하게 $M$의 어느 연관 소 아이디얼에도 속할 수 없다.
역으로 모든 $\mathfrak q \in \text{Ass}(M)$에 대해
$I \not \subset \mathfrak q$라고 하자.
이 경우 보조정리 \ref{algebra-lemma-finite-ass}와 \ref{algebra-lemma-silly}에 의해
모든 $\mathfrak q \in \text{Ass}(M)$에 대해
$x \in I$, $x \not \in \mathfrak q$인 원소를 택할 수 있다.
보조정리 \ref{algebra-lemma-ass-zero-divisors}에 의해
원소 $x$는 $M$ 위에서 영인자가 아니다.
\end{proof}

\begin{lemma}
\label{algebra-lemma-zero-at-ass-zero}
$R$을 환이라 하고 $M$을 $R$-가군이라 하자.
$R$이 뇌터 환이면 준동형사상
$$
M
\longrightarrow
\prod\nolimits_{\mathfrak p \in \text{Ass}(M)} M_{\mathfrak p}
$$
은 단사이다.
\end{lemma}

\begin{proof}
$x \in M$을 이 준동형사상의 핵에 속하는 원소라 하자.
$Rx \subset M$의 연관 소 아이디얼이 $\mathfrak p$이면,
한편으로 보조정리 \ref{algebra-lemma-ass}에 의해
$\mathfrak p \in \text{Ass}(M)$이고,
다른 한편으로 $(Rx)_{\mathfrak p} \subset M_{\mathfrak p}$는
영이 아니다. 이 모순은 $\text{Ass}(Rx) = \emptyset$임을 보인다.
따라서 보조정리 \ref{algebra-lemma-ass-zero}에 의해 $Rx = 0$이다.
\end{proof}

\noindent
% P01-E0396
이 보조정리는 아마 다른 곳에 놓아야 할 것이다.

\begin{lemma}
\label{algebra-lemma-dim-not-zero-exists-nonzerodivisor-nonunit}
$k$를 체라 하고 $S$를 유한형 $k$-대수라 하자.
$\dim(S) > 0$이면 영인자가 아니고 가역원도 아닌
원소 $f \in S$가 존재한다.
\end{lemma}

\begin{proof}
보조정리 \ref{algebra-lemma-finite-ass}에 의해 환 $S$에는
연관 소 아이디얼이 유한 개 있다.
보조정리 \ref{algebra-lemma-finite-type-algebra-finite-nr-primes}에 의해
환 $S$에는 극대 아이디얼이 무한히 많다.
따라서 $S$의 연관 소 아이디얼이 아닌 극대 아이디얼
$\mathfrak m \subset S$를 택할 수 있다.
소 아이디얼 회피(보조정리 \ref{algebra-lemma-silly})에 의해
$S$의 어느 연관 소 아이디얼에도 속하지 않는
영이 아닌 $f \in \mathfrak m$을 택할 수 있다.
보조정리 \ref{algebra-lemma-ass-zero-divisors}에 의해 원소 $f$는
영인자가 아니고, $f \in \mathfrak m$이므로 $f$는 가역원이 아니다.
\end{proof}

\section{상징적 거듭제곱}
\label{algebra-section-symbolic-power}

\noindent
다음과 같이 정의한다.

\begin{definition}
\label{algebra-definition-symbolic-power}
$R$을 환이라 하고 $\mathfrak p$를 소 아이디얼이라 하자.
$n \geq 0$에 대해 $\mathfrak p$의 $n$번째
{\it 상징적 거듭제곱}은 아이디얼
$\mathfrak p^{(n)} = \Ker(R \to R_\mathfrak p/\mathfrak p^nR_\mathfrak p)$이다.
\end{definition}

\noindent
$\mathfrak p^n \subset \mathfrak p^{(n)}$이지만
등호가 항상 성립하는 것은 아님을 주목하자.

\begin{lemma}
\label{algebra-lemma-symbolic-power-associated}
$R$을 뇌터 환이라 하자.
$\mathfrak p$를 소 아이디얼이라 하자.
$n > 0$이라 하자. 그러면
$\text{Ass}(R/\mathfrak p^{(n)}) = \{\mathfrak p\}$이다.
\end{lemma}

\begin{proof}
$\mathfrak q$가 $R/\mathfrak p^{(n)}$의 연관 소 아이디얼이면
자명하게 $\mathfrak p \subset \mathfrak q$이다.
한편 임의의 $x \in R$, $x \not \in \mathfrak p$는
$R/\mathfrak p^{(n)}$ 위에서 영인자가 아니다.
실제로 $y \in R$이고
$xy \in \mathfrak p^{(n)} = R \cap \mathfrak p^nR_{\mathfrak p}$이면
$y \in \mathfrak p^nR_{\mathfrak p}$이고, 따라서
$y \in \mathfrak p^{(n)}$이다. 그러므로 결론이 따른다.
\end{proof}

\begin{lemma}
\label{algebra-lemma-symbolic-power-flat-extension}
% P01-E0397
$R \to S$를 평탄한 환 준동형사상이라 하자.
$\mathfrak p \subset R$을 $\mathfrak q = \mathfrak p S$가
$S$의 소 아이디얼이 되는 소 아이디얼이라 하자.
그러면 $\mathfrak p^{(n)} S = \mathfrak q^{(n)}$이다.
\end{lemma}

\begin{proof}
$\mathfrak p^{(n)} = \Ker(R \to R_\mathfrak p/\mathfrak p^nR_\mathfrak p)$이므로,
평탄성을 사용하면 $\mathfrak p^{(n)} S$가 준동형사상
$S \to S_\mathfrak p/\mathfrak p^nS_\mathfrak p$의 핵임을 알 수 있다.
한편 $\mathfrak q^{(n)}$은 준동형사상
$S \to S_\mathfrak q/\mathfrak q^nS_\mathfrak q =
S_\mathfrak q/\mathfrak p^nS_\mathfrak q$의 핵이다.
따라서
$$
S_\mathfrak p/\mathfrak p^nS_\mathfrak p
\longrightarrow
S_\mathfrak q/\mathfrak p^nS_\mathfrak q
$$
가 단사임을 보이면 충분하다.
오른쪽 가군은 왼쪽 가군을
$f \in S$, $f \not \in \mathfrak q$인 원소들로 국소화한 것이다.
따라서 이 원소들이
$S_\mathfrak p/\mathfrak p^nS_\mathfrak p$ 위에서
영인자가 아님을 보이면 충분하다.
평탄성에 의해 가군
$S_\mathfrak p/\mathfrak p^nS_\mathfrak p$에는 부분몫들이
$$
\mathfrak p^iS_\mathfrak p/\mathfrak p^{i + 1}S_\mathfrak p
\cong \mathfrak p^iR_\mathfrak p/\mathfrak p^{i + 1}R_\mathfrak p
\otimes_{R_\mathfrak p} S_\mathfrak p \cong
V \otimes_{\kappa(\mathfrak p)} (S/\mathfrak q)_\mathfrak p
$$
인 유한 여과가 있다.
여기서 $V$는 $\kappa(\mathfrak p)$-벡터공간이다.
% P01-E0398
따라서 원하는 대로 $f$는 가역적으로 작용한다.
\end{proof}

\section{상대 연관 소 아이디얼 집합}
\label{algebra-section-relative-assassin}

\noindent
상대 연관 소 아이디얼 집합을 논의하자. $R \to S$를 환 준동형사상이라 하고,
$N$을 $S$-가군이라 하자. 이 상황에서 다음과 같은 $S$의 소 아이디얼
$\mathfrak q$들의 집합을 도입할 수 있다.
\begin{enumerate}
\item $A$: $\mathfrak p = R \cap \mathfrak q$라 두면
$\mathfrak q \in \text{Ass}_S(N \otimes_R \kappa(\mathfrak p))$인 것,
\item $A'$: $\mathfrak p = R \cap \mathfrak q$라 두면
% P01-E0400
$\mathfrak q$가 표준 사상
$\Spec(S \otimes_R \kappa(\mathfrak p)) \to \Spec(S)$에 의한
$\text{Ass}_{S \otimes \kappa(\mathfrak p)}(N \otimes_R \kappa(\mathfrak p))$의
상에 속하는 것,
\item $A_{fin}$: $\mathfrak p = R \cap \mathfrak q$라 두면
$\mathfrak q \in \text{Ass}_S(N/\mathfrak pN)$인 것,
\item $A'_{fin}$: 어떤 소 아이디얼 $\mathfrak p' \subset R$에 대하여
$\mathfrak q \in \text{Ass}_S(N/\mathfrak p'N)$인 것,
\item $B$: 어떤 $R$-가군 $M$에 대하여
$\mathfrak q \in \text{Ass}_S(N \otimes_R M)$인 것, 그리고
\item $B_{fin}$: 어떤 유한 $R$-가군 $M$에 대하여
$\mathfrak q \in \text{Ass}_S(N \otimes_R M)$인 것.
\end{enumerate}
이 집합들 사이의 관계 몇 가지를 알아보자.

\begin{lemma}
\label{algebra-lemma-compare-relative-assassins}
$R \to S$를 환 준동형사상이라 하고 $N$을 $S$-가군이라 하자.
% P01-E0399
$A$, $A'$, $A_{fin}$, $B$, $B_{fin}$을 위에서 도입한
$\Spec(S)$의 부분집합들이라 하자.
\begin{enumerate}
\item 항상 $A = A'$이다.
\item 항상 $A_{fin} \subset A$,
$B_{fin} \subset B$, $A_{fin} \subset A'_{fin} \subset B_{fin}$
및 $A \subset B$이다.
\item $S$가 뇌터 환이면 $A = A_{fin}$이고 $B = B_{fin}$이다.
\item $N$이 $R$ 위에서 평탄이면
$A = A_{fin} = A'_{fin}$이고 $B = B_{fin}$이다.
\item $R$이 뇌터 환이고 $N$이 $R$ 위에서 평탄이면 모든 집합이 같으며,
즉 $A = A' = A_{fin} = A'_{fin} = B = B_{fin}$이다.
\end{enumerate}
\end{lemma}

\begin{proof}
이 증명에 나오는 몇몇 논증은 뒤에서 더 정밀한 보조정리들의 증명에
되풀이된다. 그 보조정리들은 모든 가군 또는 모든 소 아이디얼의 모임이 아니라
주어진 가군 $M$이나 주어진 소 아이디얼 $\mathfrak p$를 다루기 때문이다.

\medskip\noindent
(1)의 증명. $\mathfrak p$를 $R$의 소 아이디얼이라 하자. 그러면
$$
\text{Ass}_S(N \otimes_R \kappa(\mathfrak p)) =
\text{Ass}_{S/\mathfrak pS}(N \otimes_R \kappa(\mathfrak p)) =
\text{Ass}_{S \otimes_R \kappa(\mathfrak p)}(N \otimes_R \kappa(\mathfrak p))
$$
이다. 첫째 등식은 보조정리 \ref{algebra-lemma-ass-quotient-ring}에서,
둘째 등식은 보조정리 \ref{algebra-lemma-localize-ass}의 (1)에서 따른다.
% P01-E0401
이는 $A = A'$임을 증명한다. 포함관계
$A_{fin} \subset A'_{fin}$은 명백하다.

\medskip\noindent
(2)의 증명. $A_{fin} \subset A$일 가능성을 제외하면 각 포함관계는
정의에서 바로 따른다. 남은 포함관계는 보조정리
\ref{algebra-lemma-localize-ass}와 $A_{fin}$의 서술에서
$\mathfrak p = R \cap \mathfrak q$를 요구한다는 사실에서 따른다.

\medskip\noindent
(3)의 증명. $S$가 뇌터 환이면 등식 $A = A_{fin}$은
보조정리 \ref{algebra-lemma-localize-ass}의 (3)에서 따른다.
$\mathfrak q = (g_1, \ldots, g_m)$을 $S$의 유한 생성 소 아이디얼이라 하자.
$z \in N \otimes_R M$을 소멸자가 $\mathfrak q$인 원소라 하자.
$z$가 어떤 $z' \in N \otimes_R M'$의 상이 되도록 유한 부분가군
$M' \subset M$을 택할 수 있다. 그러면
$\text{Ann}_S(z') \subset \mathfrak q = \text{Ann}_S(z)$이다.
$N \otimes_R -$는 여극한과 가환하고 $M$은 유한 $R$-가군들의
유향 여극한이므로, $M' \subset M'' \subset M$이고 상
$z'' \in N \otimes_R M''$가 $g_1, \ldots, g_m$에 의해 소멸되는
경우를 찾을 수 있다. 따라서 $\text{Ann}_S(z'') = \mathfrak q$이다.
이는 $S$가 뇌터 환이면 $B = B_{fin}$임을 증명한다.

\medskip\noindent
(4)의 증명. $N$이 평탄이면 함자 $N \otimes_R -$는 완전하다.
특히 $M' \subset M$이면 $N \otimes_R M' \subset N \otimes_R M$이다.
따라서 $z \in N \otimes_R M$의 소멸자
$\mathfrak q = \text{Ann}_S(z)$가 소 아이디얼이면,
$z \in N \otimes_R M'$이 되는 임의의 유한 $R$-부분가군
$M' \subset M$을 택할 수 있고, $N \otimes_R M'$의 원소로 본
$z$의 소멸자도 $\mathfrak q$와 같다. 따라서 $B = B_{fin}$이다.
$\mathfrak p'$을 $R$의 소 아이디얼이라 하고, $\mathfrak q$를
$N/\mathfrak p'N$의 연관 소 아이디얼인 $S$의 소 아이디얼이라 하자.
이는 $\mathfrak p'S \subset \mathfrak q$를 함의한다. $N$이 $R$ 위에서
평탄이므로 $N/\mathfrak p'N$은 정역 $R/\mathfrak p'$ 위에서 평탄하다.
따라서 $R/\mathfrak p'$의 모든 영이 아닌 원소는
% P01-E0402
$N/\mathfrak p'$ 위의 영인자가 아니다. 그러므로 이 원소들 가운데 어느 것도
$\mathfrak q$의 원소로 갈 수 없고, $\mathfrak p' = R \cap \mathfrak q$이다.
따라서 $A_{fin} = A'_{fin}$이다. 마지막으로 보조정리
\ref{algebra-lemma-localize-ass-nonzero-divisors}에 의해
$\text{Ass}_S(N/\mathfrak p'N) =
\text{Ass}_S(N \otimes_R \kappa(\mathfrak p'))$, 즉 $A'_{fin} = A$이다.

\medskip\noindent
(5)의 증명. 다른 등식들은 (4)에서 증명했으므로
$A'_{fin} = B_{fin}$만 보이면 된다. 이를 위해 $M$을 유한
$R$-가군이라 하자. 보조정리 \ref{algebra-lemma-filter-Noetherian-module}에 의해
$R$-부분가군들로 이루어진 여과
$$
0 = M_0 \subset M_1 \subset \ldots \subset M_n = M
$$
가 존재하며, 각 몫 $M_i/M_{i-1}$은 $R$의 어떤 소 아이디얼
$\mathfrak p_i$에 대한 $R/\mathfrak p_i$와 동형이다.
$N$이 평탄이므로 $S$-부분가군들로 이루어진 여과
$$
0 = N \otimes_R M_0 \subset N \otimes_R M_1 \subset \ldots \subset
N \otimes_R M_n = N \otimes_R M
$$
를 얻으며, 각 부분몫은 $N/\mathfrak p_iN$과 동형이다.
보조정리 \ref{algebra-lemma-ass}에 의해
$\text{Ass}_S(N \otimes_R M) \subset \bigcup \text{Ass}_S(N/\mathfrak p_iN)$이다.
그러므로 $B_{fin} \subset A'_{fin}$이다. 반대쪽 포함관계는 (2)의 일부이므로
결론을 얻는다.
\end{proof}

\noindent
$S/R$ 위의 $N$의 상대 연관 소 아이디얼 집합을 위의 집합
$A = A'$으로 정의한다. 이 집합은 섬유환 위의 섬유가군
$N \otimes_R \kappa(\mathfrak p)$에만 의존한다는 것이 그 동기이다.
가군의 연관 소 아이디얼 집합과 마찬가지로, 이 개념은 섬유환
$S \otimes_R \kappa(\mathfrak p)$이 뇌터 환일 때, 예를 들어
$R \to S$가 유한형일 때 가장 의미가 있음을 주의한다.

\begin{definition}
\label{algebra-definition-relative-assassin}
$R \to S$를 환 준동형사상이라 하고 $N$을 $S$-가군이라 하자.
{\it $S/R$ 위의 $N$의 상대 연관 소 아이디얼 집합}은 집합
$$
\text{Ass}_{S/R}(N)
=
\{ \mathfrak q \subset S \mid
\mathfrak q \in \text{Ass}_S(N \otimes_R \kappa(\mathfrak p))
\text{ with }\mathfrak p = R \cap \mathfrak q\}.
$$
이다. 이는 보조정리 \ref{algebra-lemma-compare-relative-assassins}에서
$A$라고 이름 붙인 집합이다.
\end{definition}

\noindent
다음 몇 결과의 취지는 겉보기에는 분명하지 않더라도
상대 연관 소 아이디얼 집합에 관한 것이라는 데 있다.

\begin{lemma}
\label{algebra-lemma-bourbaki}
$R \to S$를 환 준동형사상이라 하자.
$M$을 $R$-가군이라 하고 $N$을 $S$-가군이라 하자.
% P01-E1383
$N$이 $R$-가군으로서 평탄이면
$$
\text{Ass}_S(M \otimes_R N)
\supset
\bigcup\nolimits_{\mathfrak p \in \text{Ass}_R(M)} \text{Ass}_S(N/\mathfrak pN)
$$
이고, $R$이 뇌터 환이면 등호가 성립한다.
\end{lemma}

\begin{proof}
$\mathfrak p \in \text{Ass}_R(M)$이면 단사사상
$R/\mathfrak p \to M$이 존재한다. $N$이 $R$ 위에서 평탄이므로 단사사상
$R/\mathfrak p \otimes_R N \to M \otimes_R N$을 얻는다.
$R/\mathfrak p \otimes_R N = N/\mathfrak pN$이므로 보조정리
\ref{algebra-lemma-ass}에 의해
$\text{Ass}_S(N/\mathfrak pN) \subset \text{Ass}_S(M \otimes_R N)$이다.
그러므로 우변은 좌변에 포함된다.

\medskip\noindent
$M = \bigcup M_\lambda$를 유한 생성 $R$-부분가군들의 합집합으로 쓰자.
그러면 $N$이 $R$-평탄이므로
$N \otimes_R M = \bigcup N \otimes_R M_\lambda$이기도 하다.
연관 소 아이디얼의 정의에 의해
$\text{Ass}_S(N \otimes_R M) = \bigcup \text{Ass}_S(N \otimes_R M_\lambda)$이고
$\text{Ass}_R(M) = \bigcup \text{Ass}(M_\lambda)$이다.
따라서 $M$이 유한 생성이라고 가정해도 된다.

\medskip\noindent
$\mathfrak q \in \text{Ass}_S(M \otimes_R N)$이라 하고,
$R$이 뇌터 환이며 $M$이 유한 $R$-가군이라고 가정하자.
증명을 끝내려면 $\mathfrak q$가 우변의 원소임을 보여야 한다.
먼저 보조정리 \ref{algebra-lemma-associated-primes-localize}에 의해
$\mathfrak qS_{\mathfrak q} \in
\text{Ass}_{S_{\mathfrak q}}((M \otimes_R N)_{\mathfrak q})$임을 관찰한다.
$\mathfrak p$를 $R$의 대응하는 소 아이디얼이라 하자. 다음에 주목하자.
$$
(M \otimes_R N)_{\mathfrak q} = M \otimes_R N_{\mathfrak q}
= M_{\mathfrak p} \otimes_{R_{\mathfrak p}} N_{\mathfrak q}
$$
만일
$\mathfrak pR_{\mathfrak p} \not \in
\text{Ass}_{R_{\mathfrak p}}(M_{\mathfrak p})$이면,
$M_{\mathfrak p}$에서 영인자가 아닌 원소
$x \in \mathfrak pR_{\mathfrak p}$가 존재한다(보조정리
\ref{algebra-lemma-ideal-nonzerodivisor} 참조). $N_{\mathfrak q}$가
$R_{\mathfrak p}$ 위에서 평탄이므로, $\mathfrak qS_{\mathfrak q}$에서
$x$의 상은 $(M \otimes_R N)_{\mathfrak q}$ 위에서 영인자가 아니다.
이는 다음 가정과 모순이다.
% P01-E0403
$\mathfrak qS_{\mathfrak q} \in \text{Ass}_S((M \otimes_R N)_{\mathfrak q})$.
따라서 $\mathfrak p$는 $M$의 연관 소 아이디얼 가운데 하나이다.

\medskip\noindent
논증을 계속하여 여과
$$
0 = M_0 \subset M_1 \subset \ldots \subset M_n = M
$$
를 택하자. 각 몫 $M_i/M_{i-1}$은 $R$의 어떤 소 아이디얼
$\mathfrak p_i$에 대한 $R/\mathfrak p_i$와 동형이다. 보조정리
\ref{algebra-lemma-filter-Noetherian-module}를 참조하라.
(보조정리 \ref{algebra-lemma-ass-filter}에 의해 적어도 하나의 $i$에 대하여
$\mathfrak p_i = \mathfrak p$이다.) 이로부터 여과
$$
0 = M_0 \otimes_R N \subset M_1 \otimes_R N \subset \ldots
\subset M_n \otimes_R N = M \otimes_R N
$$
를 얻고, 그 부분몫들은 $N/\mathfrak p_iN$과 동형이다.
$\mathfrak p_i \not = \mathfrak p$이면 앞 단락의 결과에 의해
$\mathfrak q$는 가군 $N/\mathfrak p_iN$의 연관 소 아이디얼일 수 없다
($\text{Ass}_R(R/\mathfrak p_i) = \{\mathfrak p_i\}$이기 때문이다).
그러므로 원하는 대로 $\mathfrak q$는 $N/\mathfrak pN$의
연관 소 아이디얼이다.
\end{proof}

\begin{lemma}
\label{algebra-lemma-post-bourbaki}
$R \to S$를 환 준동형사상이라 하고 $N$을 $S$-가군이라 하자.
$N$이 $R$-가군으로서 평탄이고 $R$이 분수체 $K$를 갖는 정역이라고 가정하자.
그러면
$$
\text{Ass}_S(N) =
\text{Ass}_S(N \otimes_R K) =
\text{Ass}_{S \otimes_R K}(N \otimes_R K)
$$
이다. 여기서는 표준 포함
$\Spec(S \otimes_R K) \subset \Spec(S)$를 사용한다.
\end{lemma}

\begin{proof}
$S \otimes_R K = (R \setminus \{0\})^{-1}S$이고
$N \otimes_R K = (R \setminus \{0\})^{-1}N$임에 주목하자.
모든 영이 아닌 $x \in R$에 대하여 $N$이 $R$ 위에서 평탄이므로
$N$ 위에서 $x$에 의한 곱셈은 단사이다. 따라서 보조정리
\ref{algebra-lemma-localize-ass-nonzero-divisors}와 보조정리
\ref{algebra-lemma-localize-ass}의 (1)을 결합하면 결론이 따른다.
\end{proof}

\begin{lemma}
\label{algebra-lemma-bourbaki-fibres}
$R \to S$를 환 준동형사상이라 하자.
$M$을 $R$-가군이라 하고 $N$을 $S$-가군이라 하자.
% P01-E1383
$N$이 $R$-가군으로서 평탄이라고 가정하자. 그러면
$$
\text{Ass}_S(M \otimes_R N)
\supset
\bigcup\nolimits_{\mathfrak p \in \text{Ass}_R(M)}
\text{Ass}_{S \otimes_R \kappa(\mathfrak p)}(N \otimes_R \kappa(\mathfrak p))
$$
이다. 여기서는 주석 \ref{algebra-remark-fundamental-diagram}을 사용하여
섬유환의 스펙트럼을 $\Spec(S)$의 부분집합으로 생각한다.
$R$이 뇌터 환이면 이 포함관계는 등식이다.
\end{lemma}

\begin{proof}
이는 보조정리 \ref{algebra-lemma-bourbaki}와 동치이다. 그 동치는 보조정리
\ref{algebra-lemma-ass-quotient-ring}, \ref{algebra-lemma-flat-base-change},
\ref{algebra-lemma-post-bourbaki}에서 따른다.
\end{proof}

\begin{remark}
\label{algebra-remark-bourbaki}
$R \to S$를 환 준동형사상이라 하고 $N$을 $S$-가군이라 하자.
$\mathfrak p$를 $R$의 소 아이디얼이라 하자. 그러면
$$
\text{Ass}_S(N \otimes_R \kappa(\mathfrak p)) =
\text{Ass}_{S/\mathfrak pS}(N \otimes_R \kappa(\mathfrak p)) =
\text{Ass}_{S \otimes_R \kappa(\mathfrak p)}(N \otimes_R \kappa(\mathfrak p)).
$$
% P01-E0404
첫째 등식은 보조정리 \ref{algebra-lemma-ass-quotient-ring}에서,
둘째 등식은 보조정리 \ref{algebra-lemma-localize-ass}의 (1)에서 따른다.
\end{remark}

\section{약하게 연관된 소 아이디얼}
\label{algebra-section-weakly-ass}

\noindent
% P01-E1384
이는 연관 소 아이디얼 개념의 한 변형으로, 뇌터가 아닌 환과
유한이 아닌 가군에 유용하다.

\begin{definition}
\label{algebra-definition-weakly-associated}
$R$을 환이라 하고 $M$을 $R$-가군이라 하자.
$R$의 소 아이디얼 $\mathfrak p$가 $M$에 {\it 약하게 연관된다}고 함은
어떤 원소 $m \in M$가 존재하여 $\mathfrak p$가 소멸자
$\text{Ann}(m) = \{f \in R \mid fm = 0\}$를 포함하는 소 아이디얼들 가운데
극소인 것을 말한다. 이러한 모든 소 아이디얼의 집합을
$\text{WeakAss}_R(M)$ 또는 $\text{WeakAss}(M)$으로 나타낸다.
\end{definition}

\noindent
따라서 연관 소 아이디얼은 약하게 연관된 소 아이디얼이다.
다음은 그 소 아이디얼에서의 국소화로 나타낸 특성화이다.

\begin{lemma}
\label{algebra-lemma-weakly-ass-local}
$R$을 환이라 하고 $M$을 $R$-가군이라 하자.
$\mathfrak p$를 $R$의 소 아이디얼이라 하자.
다음 조건들은 동치이다.
\begin{enumerate}
\item $\mathfrak p$는 $M$에 약하게 연관된다,
\item $\mathfrak pR_{\mathfrak p}$는 $M_{\mathfrak p}$에 약하게 연관된다,
그리고
\item $M_{\mathfrak p}$에는 소멸자의 근기가
$\mathfrak pR_{\mathfrak p}$와 같은 원소가 들어 있다.
\end{enumerate}
\end{lemma}

\begin{proof}
(1)을 가정하자. 그러면 어떤 원소 $m \in M$가 존재하여
$\mathfrak p$는 $m$의 소멸자
$I = \{x \in R \mid xm = 0\}$를 포함하는 소 아이디얼들 가운데 극소이다.
국소화는 완전하므로 $M_{\mathfrak p}$에서 $m$의 소멸자는
$I_{\mathfrak p}$이다. 따라서 $\mathfrak pR_{\mathfrak p}$는
$R_{\mathfrak p}$에서 $M_{\mathfrak p}$의 원소 $m$의 소멸자
$I_{\mathfrak p}$를 포함하는 극소 소 아이디얼이다. 이는 (2)를 함의하며,
$\sqrt{I_{\mathfrak p}} = \mathfrak pR_{\mathfrak p}$를 함의하므로
(3)도 함의한다.

\medskip\noindent
$R_{\mathfrak p}$ 위의 $M_{\mathfrak p}$에 함의
(1) $\Rightarrow$ (3)을 적용하면 (2) $\Rightarrow$ (3)을 얻는다.

\medskip\noindent
마지막으로 (3)을 가정하자. 이는 소멸자의 근기가
$\mathfrak pR_{\mathfrak p}$와 같은 원소
$m/f \in M_{\mathfrak p}$가 존재한다는 뜻이다. $M$에서 $m$의 소멸자를
$I = \{x \in R \mid xm = 0\}$라 하면
$\sqrt{I_{\mathfrak p}} = \mathfrak pR_{\mathfrak p}$이다.
이는 분명히 $\mathfrak p$가 $I$를 포함하고 $I$를 포함하는
소 아이디얼들 가운데 극소라는 뜻이다. 즉 (1)이 성립한다.
\end{proof}

\begin{lemma}
\label{algebra-lemma-reduced-weakly-ass-minimal}
환원환에서 그 환에 약하게 연관된 소 아이디얼들은
극소 소 아이디얼들이다.
\end{lemma}

\begin{proof}
$(R, \mathfrak m)$을 환원 국소환이라 하자.
$x \in R$을 소멸자의 근기가 $\mathfrak m$인 원소라고 하자.
$\mathfrak m \not = 0$이면 $x$는 가역원일 수 없으므로
$x \in \mathfrak m$이다. 특히 어떤 $n \geq 0$에 대하여
$x^{1 + n} = 0$이다. 따라서 $x = 0$이다.
% P01-E0405
이는 $x$의 소멸자가 $\mathfrak m$에 포함된다는 가정과 모순이다.
그러므로 $\mathfrak m = 0$, 즉 $R$은 체이다.
보조정리 \ref{algebra-lemma-weakly-ass-local}에 의해 보조정리의 명제가 따른다.
\end{proof}

\begin{lemma}
\label{algebra-lemma-weakly-ass}
$R$을 환이라 하자.
$0 \to M' \to M \to M'' \to 0$을 $R$-가군들의 짧은 완전열이라 하자.
그러면 $\text{WeakAss}(M') \subset \text{WeakAss}(M)$이고
$\text{WeakAss}(M) \subset \text{WeakAss}(M') \cup \text{WeakAss}(M'')$이다.
\end{lemma}

\begin{proof}
보조정리 \ref{algebra-lemma-weakly-ass-local}의 약하게 연관된 소 아이디얼에 대한
특성화를 사용하겠다. $\mathfrak p$를 $R$의 소 아이디얼이라 하자.
국소화는 완전하므로 짧은 완전열
$0 \to M'_{\mathfrak p} \to M_{\mathfrak p} \to M''_{\mathfrak p} \to 0$
을 얻는다. $m \in M_{\mathfrak p}$를 소멸자의 근기가
$\mathfrak pR_{\mathfrak p}$인 원소라 하자. 그러면
$M''_{\mathfrak p}$에서 $m$의 상 $\overline{m}$이 영이고
$m \in M'_{\mathfrak p}$이거나, $\overline{m}$의 소멸자의 근기가
$\mathfrak pR_{\mathfrak p}$이다. 이는
$\text{WeakAss}(M) \subset \text{WeakAss}(M') \cup \text{WeakAss}(M'')$
임을 증명한다. 포함관계
$\text{WeakAss}(M') \subset \text{WeakAss}(M)$은 정의에서 바로 따른다.
\end{proof}

\begin{lemma}
\label{algebra-lemma-weakly-ass-zero}
\begin{slogan}
영이 아닌 모든 가군에는 약하게 연관된 소 아이디얼이 존재한다.
\end{slogan}
$R$을 환이라 하고 $M$을 $R$-가군이라 하자. 그러면
$$
M = (0) \Leftrightarrow \text{WeakAss}(M) = \emptyset
$$
\end{lemma}

\begin{proof}
$M = (0)$이면 정의에 의해 $\text{WeakAss}(M) = \emptyset$이다.
역으로 $M \not = 0$이라 하자. 영이 아닌 원소 $m \in M$를 택하자.
% P01-E1385
$m$의 소멸자를 $I = \{x \in R \mid xm = 0\}$라 쓰자.
그러면 $R/I \subset M$이다. 따라서 보조정리
\ref{algebra-lemma-weakly-ass}에 의해
$\text{WeakAss}(R/I) \subset \text{WeakAss}(M)$이다.
% P01-E1386
$I \not = R$이므로 $V(I) = \Spec(R/I)$에는 극소 소 아이디얼이 들어 있다.
보조정리 \ref{algebra-lemma-Zariski-topology}와
\ref{algebra-lemma-spec-closed}를 참조하라. 이로써 결론을 얻는다.
\end{proof}

\begin{lemma}
\label{algebra-lemma-weakly-ass-support}
$R$을 환이라 하고 $M$을 $R$-가군이라 하자. 그러면
$$
\text{Ass}(M) \subset \text{WeakAss}(M) \subset \text{Supp}(M).
$$
\end{lemma}

\begin{proof}
첫째 포함관계는 정의에서 바로 따른다.
$\mathfrak p \in \text{WeakAss}(M)$이면 보조정리
\ref{algebra-lemma-weakly-ass-local}에 의해 $M_{\mathfrak p} \not = 0$이고,
따라서 $\mathfrak p \in \text{Supp}(M)$이다.
\end{proof}

\begin{lemma}
\label{algebra-lemma-weakly-ass-zero-divisors}
$R$을 환이라 하고 $M$을 $R$-가군이라 하자.
합집합 $\bigcup_{\mathfrak q \in \text{WeakAss}(M)} \mathfrak q$는
$M$ 위에서 영인자인 $R$의 원소들의 집합이다.
\end{lemma}

\begin{proof}
$f \in \mathfrak q \in \text{WeakAss}(M)$이라 하자.
그러면 어떤 원소 $m \in M$가 존재하여 $\mathfrak q$는
$I = \{x \in R \mid xm = 0\}$ 위에서 극소이다.
따라서 $g \in R$, $g \not \in \mathfrak q$ 및 $n > 0$이 존재하여
$f^ngm = 0$이다. $g \not \in I$이므로 $gm \not = 0$임에 주목하자.
위 조건을 만족하는 $n$을 극소로 택하면
$f (f^{n - 1}gm) = 0$이고 $f^{n - 1}gm \not = 0$이므로
$f$는 $M$ 위에서 영인자이다.
역으로 $f \in R$이 $M$ 위에서 영인자라고 하자.
부분가군 $N = \{m \in M \mid fm = 0\}$을 생각하자.
$N$이 영이 아니므로 보조정리 \ref{algebra-lemma-weakly-ass-zero}에 의해
약하게 연관된 소 아이디얼 $\mathfrak q$를 갖는다.
분명히 $f \in \mathfrak q$이고, 보조정리
\ref{algebra-lemma-weakly-ass}에 의해 $\mathfrak q$는
$M$에 약하게 연관된 소 아이디얼이다.
\end{proof}

\begin{lemma}
\label{algebra-lemma-weakly-ass-minimal-prime-support}
$R$을 환이라 하고 $M$을 $R$-가군이라 하자.
$\text{Supp}(M)$의 원소들 가운데 극소인 임의의
$\mathfrak p \in \text{Supp}(M)$는 $\text{WeakAss}(M)$의 원소이다.
\end{lemma}

\begin{proof}
$\Spec(R_{\mathfrak p})$ 안에서
$\text{Supp}(M_{\mathfrak p}) = \{\mathfrak pR_{\mathfrak p}\}$임에
주목하자. 특히 $M_{\mathfrak p}$는 영이 아니므로 보조정리
\ref{algebra-lemma-weakly-ass-zero}에 의해
$\text{WeakAss}(M_{\mathfrak p}) \not = \emptyset$이다.
보조정리 \ref{algebra-lemma-weakly-ass-support}에 의해
$\text{WeakAss}(M_{\mathfrak p}) \subset \text{Supp}(M_{\mathfrak p})$이므로
$\text{WeakAss}(M_{\mathfrak p}) = \{\mathfrak pR_{\mathfrak p}\}$이고,
따라서 보조정리 \ref{algebra-lemma-weakly-ass-local}에 의해
$\mathfrak p \in \text{WeakAss}(M)$이다.
\end{proof}

\begin{lemma}
\label{algebra-lemma-ass-weakly-ass}
$R$을 환이라 하고 $M$을 $R$-가군이라 하자.
$\mathfrak p$를 유한 생성인 $R$의 소 아이디얼이라 하자. 그러면
$$
\mathfrak p \in \text{Ass}(M) \Leftrightarrow
\mathfrak p \in \text{WeakAss}(M).
$$
특히 $R$이 뇌터 환이면 $\text{Ass}(M) = \text{WeakAss}(M)$이다.
\end{lemma}

\begin{proof}
어떤 $g_i \in R$들에 대하여 $\mathfrak p = (g_1, \ldots, g_n)$이라 쓰자.
% P01-E0406
반대쪽 함의는 일반적으로 성립하므로 함의 ``$\Leftarrow$''만 증명하면
충분하다. 보조정리 \ref{algebra-lemma-weakly-ass-support}를 참조하라.
$\mathfrak p \in \text{WeakAss}(M)$이라 가정하자.
보조정리 \ref{algebra-lemma-weakly-ass-local}에 의해 어떤 원소
$m \in M_{\mathfrak p}$가 존재하여
$I = \{x \in R_{\mathfrak p} \mid xm = 0\}$의 근기가
$\mathfrak pR_{\mathfrak p}$이다. 따라서 각 $i$에 대해
$M_{\mathfrak p}$에서 $g_i^{e_i}m = 0$이 되는 가장 작은
$e_i > 0$이 존재한다. 어떤 $i$에 대해 $e_i > 1$이면 $m$을
$g_i^{e_i - 1} m \not = 0$으로 바꾸어 $\sum e_i$를 줄일 수 있다.
따라서 $m \in M_{\mathfrak p}$의 소멸자가
$(g_1, \ldots, g_n)R_{\mathfrak p} = \mathfrak p R_{\mathfrak p}$
이라고 가정해도 된다. 보조정리
\ref{algebra-lemma-associated-primes-localize}에 의해
$\mathfrak p \in \text{Ass}(M)$임을 알 수 있다.
\end{proof}

\begin{remark}
\label{algebra-remark-weakly-ass-not-functorial}
$\varphi : R \to S$를 환 준동형사상이라 하고 $M$을 $S$-가군이라 하자.
연관 소 아이디얼의 경우와 달리(보조정리
\ref{algebra-lemma-ass-functorial} 참조), 항상
$\Spec(\varphi)(\text{WeakAss}_S(M)) \subset \text{WeakAss}_R(M)$인 것은
아니다. 한 예로 환 준동형사상
$$
R = k[x_1, x_2, x_3, \ldots] \to
S = k[x_1, x_2, x_3, \ldots, y_1, y_2, y_3, \ldots]/
(x_1y_1, x_2y_2, x_3y_3, \ldots)
$$
과 $M = S$를 생각하자. 이 경우
$\mathfrak q = \sum x_iS$는 $S$의 극소 소 아이디얼이고,
따라서 $M = S$에 약하게 연관된 소 아이디얼이다
(보조정리 \ref{algebra-lemma-weakly-ass-minimal-prime-support} 참조).
그러나 한편 $S$의 영이 아닌 임의의 원소의 $R$에서의 소멸자는
유한 생성이므로, 그 근기는
$R \cap \mathfrak q = (x_1, x_2, x_3, \ldots)$와 같을 수 없다
(세부사항은 생략한다).
\end{remark}

\begin{lemma}
\label{algebra-lemma-weakly-ass-reverse-functorial}
$\varphi : R \to S$를 환 준동형사상이라 하고 $M$을 $S$-가군이라 하자.
그러면
$$
\Spec(\varphi)(\text{WeakAss}_S(M)) \supset \text{WeakAss}_R(M)
$$
이다.
\end{lemma}

\begin{proof}
$\mathfrak p$를 $\text{WeakAss}_R(M)$의 원소라 하자.
그러면 소멸자
$I = \{x \in R_{\mathfrak p} \mid xm = 0\}$의 근기가
$\mathfrak pR_{\mathfrak p}$인 어떤 $m \in M_{\mathfrak p}$가 존재한다.
$S_{\mathfrak p}$에서 $m$의 소멸자
$J = \{x \in S_{\mathfrak p} \mid xm = 0 \}$를 생각하자.
$IS_{\mathfrak p} \subset J$이므로 $J$ 위의 임의의 극소 소 아이디얼
$\mathfrak q \subset S_{\mathfrak p}$는 $\mathfrak p$ 위에 놓인다.
더욱이 예를 들어 보조정리 \ref{algebra-lemma-weakly-ass-local}에 의해
그러한 $\mathfrak q$는 $M$에 약하게 연관된 소 아이디얼에 대응한다.
\end{proof}

\begin{remark}
\label{algebra-remark-ass-functorial}
$\varphi : R \to S$를 환 준동형사상이라 하고 $M$을 $S$-가군이라 하자.
% P01-E0407
스펙트럼 위의 대응하는 사상을
$f : \Spec(S) \to \Spec(R)$로 나타내자. 그러면
$$
f(\text{Ass}_S(M)) \subset
\text{Ass}_R(M) \subset
\text{WeakAss}_R(M) \subset
f(\text{WeakAss}_S(M))
$$
이다. 보조정리 \ref{algebra-lemma-ass-functorial},
\ref{algebra-lemma-weakly-ass-reverse-functorial},
\ref{algebra-lemma-weakly-ass-support}를 참조하라.
일반적으로 모든 포함관계가 진포함일 수 있다. 주석
\ref{algebra-remark-ass-reverse-functorial}과
\ref{algebra-remark-weakly-ass-not-functorial}을 참조하라.
$S$가 뇌터 환이면 바깥쪽 두 포함관계가 보조정리
\ref{algebra-lemma-ass-weakly-ass}에 의해 등식이므로 모든 포함관계가 등식이다.
\end{remark}

\begin{lemma}
\label{algebra-lemma-weakly-ass-finite-ring-map}
$\varphi : R \to S$를 환 준동형사상이라 하고 $M$을 $S$-가군이라 하자.
% P01-E0407
스펙트럼 위의 대응하는 사상을
$f : \Spec(S) \to \Spec(R)$로 나타내자.
$\varphi$가 유한 환 준동형사상이면
$$
\text{WeakAss}_R(M) = f(\text{WeakAss}_S(M)).
$$
\end{lemma}

\begin{proof}
포함관계 가운데 하나는 이미 증명하였다. 주석
\ref{algebra-remark-ass-functorial}을 참조하라. 다른 포함관계를 증명하기 위해
$\mathfrak q \in \text{WeakAss}_S(M)$이라 하고 $\mathfrak p$를
이에 대응하는 $R$의 소 아이디얼이라 하자.
$\mathfrak q$가
$J = \{g \in S \mid gm = 0\}$ 위의 극소 소 아이디얼이 되는
원소 $m \in M$를 택하자. 따라서
$JS_{\mathfrak q}$의 근기는 $\mathfrak qS_{\mathfrak q}$이다.
$R \to S$가 유한이므로 $\mathfrak p$ 위에는 유한 개의 소 아이디얼
$\mathfrak q = \mathfrak q_1, \mathfrak q_2, \ldots, \mathfrak q_l$이
있다. 보조정리 \ref{algebra-lemma-finite-finite-fibres}를 참조하라.
$i > 1$에 대해 $x \not \in \mathfrak q_i$인
$x \in \mathfrak q$를 택하자. 보조정리 \ref{algebra-lemma-silly}를 참조하라.
위에서 본 바에 따라 어떤 원소 $y \in S$, $y \not \in \mathfrak q$와
어떤 정수 $t > 0$이 존재하여 $y x^t m = 0$이다.
따라서 원소 $ym \in M$은 $x^t$에 의해 소멸되므로
$ym$은 $M_{\mathfrak q_i}$, $i = 2, \ldots, l$에서 영으로 간다.
% P01-E0408
분명히 $ym$은 $S_{\mathfrak q}$에서는 영으로 가지 않는다.

\medskip\noindent
유한 환 준동형사상에 대한 상향 정리에 의해 환 $S_{\mathfrak p}$는
극대 아이디얼들이 $\mathfrak q_i S_{\mathfrak p}$인 반국소환이다.
보조정리 \ref{algebra-lemma-integral-going-up}을 참조하라.
$f \in \mathfrak pR_{\mathfrak p}$이면 $f$의 어떤 거듭제곱이
$JS_{\mathfrak q}$에 들어가므로, 어떤 $t > 0$에 대해
$f^t ym$은 $M_{\mathfrak q}$에서 영으로 간다.
$ym$은 $S_{\mathfrak p}$의 다른 극대 아이디얼들에서 영이 되므로
$f^t ym$은 $M_{\mathfrak p}$에서 영이다. 보조정리
\ref{algebra-lemma-characterize-zero-local}을 참조하라.
이로써 $\mathfrak p$가 $R$에서 $ym$의 소멸자 위의 극소 소 아이디얼임을
알 수 있고, 결론을 얻는다.
\end{proof}

\begin{lemma}
\label{algebra-lemma-weakly-ass-quotient-ring}
$R$을 환이라 하고 $I$를 아이디얼이라 하자.
$M$을 $R/I$-가군이라 하자. 표준 단사사상
$\Spec(R/I) \to \Spec(R)$을 통해
$\text{WeakAss}_{R/I}(M) = \text{WeakAss}_R(M)$이다.
\end{lemma}

\begin{proof}
보조정리 \ref{algebra-lemma-weakly-ass-finite-ring-map}의 특수한 경우이다.
\end{proof}

\begin{lemma}
\label{algebra-lemma-localize-weakly-ass}
$R$을 환이라 하고 $M$을 $R$-가군이라 하자.
$S \subset R$을 곱셈적 부분집합이라 하자.
표준 단사사상 $\Spec(S^{-1}R) \to \Spec(R)$을 통해
$\text{WeakAss}_R(S^{-1}M) = \text{WeakAss}_{S^{-1}R}(S^{-1}M)$이고
$$
\text{WeakAss}(M) \cap \Spec(S^{-1}R) = \text{WeakAss}(S^{-1}M).
$$
\end{lemma}

\begin{proof}
$m \in S^{-1}M$이라 하자.
$I = \{x \in R \mid xm = 0\}$ 및
$I' = \{x' \in S^{-1}R \mid x'm = 0\}$라 하자.
그러면 $I' = S^{-1}I$이고, $I = R$이 아닌 한
$I \cap S = \emptyset$이다(확인은 생략한다).
따라서 $I'$ 위의 $S^{-1}R$의 극소 소 아이디얼들은
$I$ 위에서 극소이고 $S$를 피하는 $R$의 소 아이디얼들과 일대일로 대응한다.
이는 등식
$\text{WeakAss}_R(S^{-1}M) = \text{WeakAss}_{S^{-1}R}(S^{-1}M)$을 증명한다.
둘째 등식은 보조정리 \ref{algebra-lemma-weakly-ass-local}과
% P01-E0409
$\mathfrak p \in R$, $S \cap \mathfrak p = \emptyset$일 때의 등식
$M_{\mathfrak p} = (S^{-1}M)_{S^{-1}\mathfrak p}$에서 따른다.
\end{proof}

\begin{lemma}
\label{algebra-lemma-localize-weakly-ass-nonzero-divisors}
$R$을 환이라 하고 $M$을 $R$-가군이라 하자.
$S \subset R$을 곱셈적 부분집합이라 하자.
모든 $s \in S$가 $M$ 위에서 영인자가 아니라고 가정하자. 그러면
$$
\text{WeakAss}(M) = \text{WeakAss}(S^{-1}M).
$$
\end{lemma}

\begin{proof}
가정에 의해 $M \subset S^{-1}M$이므로 보조정리
\ref{algebra-lemma-weakly-ass}에서
$\text{WeakAss}(M) \subset \text{WeakAss}(S^{-1}M)$을 얻는다.
역으로 $n/s \in S^{-1}M$을 소멸자가 $I$인 원소라 하고,
$\mathfrak p$를 $I$ 위의 극소 소 아이디얼이라 하자.
그러면 원소 $n \in M$의 소멸자는 $I$이고,
$\mathfrak p$는 $I$ 위의 극소 소 아이디얼이다.
\end{proof}

\begin{lemma}
\label{algebra-lemma-zero-at-weakly-ass-zero}
$R$을 환이라 하고 $M$을 $R$-가군이라 하자. 사상
$$
M
\longrightarrow
\prod\nolimits_{\mathfrak p \in \text{WeakAss}(M)} M_{\mathfrak p}
$$
은 단사이다.
\end{lemma}

\begin{proof}
$x \in M$을 이 사상의 핵에 속하는 원소라 하자.
$N = Rx \subset M$이라 두자. $\mathfrak p$가 $N$에 약하게 연관된
소 아이디얼이면 한편으로
$\mathfrak p \in \text{WeakAss}(M)$이고
(보조정리 \ref{algebra-lemma-weakly-ass}), 다른 한편으로
$N_{\mathfrak p} \subset M_{\mathfrak p}$는 영이 아니다.
이 모순은 $\text{WeakAss}(N) = \emptyset$임을 보인다.
따라서 보조정리 \ref{algebra-lemma-weakly-ass-zero}에 의해
$N = 0$, 즉 $x = 0$이다.
\end{proof}

\begin{lemma}
\label{algebra-lemma-weak-post-bourbaki}
$R \to S$를 환 준동형사상이라 하고 $N$을 $S$-가군이라 하자.
$N$이 $R$-가군으로서 평탄이고 $R$이 분수체 $K$를 갖는
정역이라고 가정하자. 그러면
$$
\text{WeakAss}_S(N) = \text{WeakAss}_{S \otimes_R K}(N \otimes_R K)
$$
이다. 여기서는 표준 포함
$\Spec(S \otimes_R K) \subset \Spec(S)$를 사용한다.
\end{lemma}

\begin{proof}
$S \otimes_R K = (R \setminus \{0\})^{-1}S$이고
$N \otimes_R K = (R \setminus \{0\})^{-1}N$임에 주목하자.
영이 아닌 임의의 $x \in R$에 대해 $N$이 $R$ 위에서 평탄이므로
$N$ 위에서 $x$에 의한 곱셈은 단사이다. 따라서 보조정리
\ref{algebra-lemma-localize-weakly-ass-nonzero-divisors}에서 결론이 따른다.
\end{proof}

\begin{lemma}
\label{algebra-lemma-weakly-ass-change-fields}
$K/k$를 체확대라 하고 $R$을 $k$-대수라 하자.
$M$을 $R$-가군이라 하자.
$\mathfrak q \subset R \otimes_k K$를
$\mathfrak p \subset R$ 위에 놓이는 소 아이디얼이라 하자.
$\mathfrak q$가 $M \otimes_k K$에 약하게 연관되면
$\mathfrak p$는 $M$에 약하게 연관된다.
\end{lemma}

\begin{proof}
$z \in M \otimes_k K$를 $\mathfrak q$가 $z$의 소멸자
$J \subset R \otimes_k K$ 위에서 극소가 되게 하는 원소라 하자.
$z \in M \otimes_k L$이 되는 유한 생성 부분확대 $K/L/k$를 택하자.
$R \otimes_k L \to R \otimes_k K$가 평탄이므로
$J = I(R \otimes_k K)$이다. 여기서 $I \subset R \otimes_k L$은
더 작은 환에서의 $z$의 소멸자이다
(보조정리 \ref{algebra-lemma-annihilator-flat-base-change}).
따라서 하향 정리에 의해
$\mathfrak q \cap (R \otimes_k L)$은 $I$ 위에서 극소이다
(보조정리 \ref{algebra-lemma-flat-going-down}).
이렇게 하여 다음 단락에서 서술하는 경우로 환원한다.

\medskip\noindent
$K/k$가 유한 생성 체확대라고 가정하자.
$x_1, \ldots, x_r \in K$를 $k$ 위의 $K$의 초월기저라 하자.
Fields, Section \ref{fields-section-transcendence}를 참조하라.
$L = k(x_1, \ldots, x_r)$이라 두자.
$[K : L] = n$이라 하자. 그러면
$R \otimes_k L \to R \otimes_k K$는 유한 환 준동형사상이다.
따라서 보조정리 \ref{algebra-lemma-weakly-ass-finite-ring-map}에 의해
$\mathfrak q \cap (R \otimes_k L)$은 $M \otimes_k K$를
$R \otimes_k L$-가군으로 보았을 때 약하게 연관된 소 아이디얼이다.
$R \otimes_k L$-가군으로서
$M \otimes_k K \cong (M \otimes_k L)^{\oplus n}$이므로,
$\mathfrak q \cap (R \otimes_k L)$은 $M \otimes_k L$에 약하게 연관된
소 아이디얼이다(예를 들어 보조정리
\ref{algebra-lemma-weakly-ass}와 귀납법을 사용하라).
이렇게 하여 다음 단락에서 논의하는 경우로 환원한다.

\medskip\noindent
$K = k(x_1, \ldots, x_r)$가 순수 초월 체확대라고 가정하자.
$R$을 $R_\mathfrak p$로, $M$을 $M_\mathfrak p$로,
$\mathfrak q$를 $\mathfrak q(R_\mathfrak p \otimes_k K)$로 바꾸어도 된다.
보조정리 \ref{algebra-lemma-localize-weakly-ass}를 참조하라.
이렇게 하여 다음 단락에서 논의하는 경우로 환원한다.

\medskip\noindent
$K = k(x_1, \ldots, x_r)$가 순수 초월 체확대이고
$R$이 극대 아이디얼 $\mathfrak p$를 갖는 국소환이라고 가정하자.
임의의 $f \in R \otimes_k K$,
$f \not \in \mathfrak p(R \otimes_k K)$가 $M \otimes_k K$ 위에서
영인자가 아니라고 주장한다. 실제로 $z \in M \otimes_k K$라 하자.
$z \in M' \otimes_k K$인 유한 $R$-부분가군
$M' \subset M$이 존재하며, $M'$은 이 성질에 관하여 극소이다.
$k$-벡터공간으로서 $K$의 기저 $\{t_\alpha\}$를 택하여
$z = \sum m_\alpha \otimes t_\alpha$로 쓰고,
$M'$을 $m_\alpha$들이 생성하는 $R$-부분가군으로 택하면 된다.
$z \in \mathfrak p(M' \otimes_k K) = \mathfrak p M' \otimes_k K$이면
$\mathfrak pM' = M'$이고 보조정리 \ref{algebra-lemma-NAK}에 의해
% P01-E1387
$M' = 0$이므로 모순이다. 따라서 $z$는
$M'/\mathfrak p M' \otimes_k K$에서 영이 아닌 상
$\overline{z}$를 갖는다.
% P01-E1388
% P01-E0410
그러나 $R/\mathfrak p \otimes_k K$는
$\kappa(\mathfrak p)[x_1, \ldots, x_n]$의 국소화로서 정역이고,
$M'/\mathfrak p M' \otimes_k K$는 자유가군이므로
$f\overline{z} \not = 0$이다. 이는 주장을 증명한다.

\medskip\noindent
마지막으로 $z \in M \otimes_k K$를 $\mathfrak q$가 $z$의 소멸자
$J \subset R \otimes_k K$ 위에서 극소가 되게 하는 원소라 하자.
$f \in \mathfrak p$에 대하여 어떤 $n \geq 1$과
$g \in R \otimes_k K$, $g \not \in \mathfrak q$가 존재하여
% P01-E0411
$g f^n z \in J$, 즉 $g f^n z = 0$이다.
(이는 $\mathfrak q$가 $\mathfrak p$ 위에 놓이고
$\mathfrak q$가 $J$ 위에서 극소이기 때문이다.)
위에서 $g$가 영인자가 아님을 보였으므로 $f^n z = 0$이다.
이는 $M \otimes_k K$를 $R$-가군으로 보았을 때
$\mathfrak p$가 약하게 연관된 소 아이디얼임을 뜻한다.
$M \otimes_k K$는 $M$의 복사본들의 직합이므로 앞에서와 같이
$\mathfrak p$가 $M$에 약하게 연관된 소 아이디얼이라고 결론짓는다.
\end{proof}

\section{매립 소 아이디얼}
\label{algebra-section-embedded-primes}

\noindent
정의는 다음과 같다.

\begin{definition}
\label{algebra-definition-embedded-primes}
$R$을 환이라 하고 $M$을 $R$-가군이라 하자.
\begin{enumerate}
\item $M$의 연관 소 아이디얼 가운데 $M$의 연관 소 아이디얼들 사이에서
극소가 아닌 것을 $M$의 {\it 매립 연관 소 아이디얼}이라 한다.
\item {\it $R$의 매립 소 아이디얼}은 $R$을 $R$-가군으로 보았을 때의
매립 연관 소 아이디얼이다.
\end{enumerate}
\end{definition}

\noindent
다음은 이들을 제거하는 한 방법이다.

\begin{lemma}
\label{algebra-lemma-remove-embedded-primes}
$R$을 뇌터 환이라 하고 $M$을 유한 $R$-가군이라 하자.
$R$-부분가군들의 집합
$$
\{
K \subset M
\mid
\text{Supp}(K)
\text{ nowhere dense in }
\text{Supp}(M)
\}.
$$
을 생각하자. 이 집합에는 극대원 $K$가 있고, 몫
$M' = M/K$는 다음 성질들을 갖는다.
\begin{enumerate}
\item $\text{Supp}(M) = \text{Supp}(M')$,
\item $M'$에는 매립 연관 소 아이디얼이 없다,
\item $M$의 모든 매립 연관 소 아이디얼에 속하는 임의의
$f \in R$에 대하여 $M_f \cong M'_f$이다.
\end{enumerate}
\end{lemma}

\begin{proof}
이후 별도의 언급 없이 보조정리 \ref{algebra-lemma-finite-ass}와
명제 \ref{algebra-proposition-minimal-primes-associated-primes}를 사용하겠다.
$\mathfrak q_1, \ldots, \mathfrak q_t$를 $M$의 지지집합에 속하는
극소 소 아이디얼들이라 하자.
$\mathfrak p_1, \ldots, \mathfrak p_s$를 $M$의 매립 연관
소 아이디얼들이라 하자. 그러면
$\text{Ass}(M) = \{\mathfrak q_j, \mathfrak p_i\}$이다. 다음과 같이 두자.
$$
K = \{m \in M \mid \text{Supp}(Rm) \subset \bigcup V(\mathfrak p_i)\}
$$
이것이 부분가군임은 바로 알 수 있다. $M$은 뇌터 환 위의
유한 가군이므로 $K$도 유한이다. 따라서
$\text{Supp}(K)$는 $\text{Supp}(M)$에서 어디에서도 조밀하지 않다.
지지집합이 $\text{Supp}(M)$에서 어디에서도 조밀하지 않은 다른 부분가군
$K' \subset M$을 택하자. 이는
% P01-E0412
$K_{\mathfrak q_j} = 0$이라는 뜻이다. 그러므로 $m \in K'$이면
$m$은 $M_{\mathfrak q_j}$에서 영으로 가고, 다시
$(Rm)_{\mathfrak q_j} = 0$을 함의한다.
한편 $\text{Ass}(Rm) \subset \text{Ass}(M)$이다.
따라서 $Rm$의 지지집합은 $\bigcup V(\mathfrak p_i)$에 포함된다.
그러므로 $m \in K$이고, $m$은 $K'$의 임의의 원소였으므로
$K' \subset K$이다.

\medskip\noindent
$M' = M/K$라 하자. $K_{\mathfrak q_j}=0$이므로 모든 $j$에 대해
$M'_{\mathfrak q_j} = M_{\mathfrak q_j}$이다.
따라서 $M$과 $M'$은 같은 지지집합을 갖는다.

\medskip\noindent
$\overline{m} \in M'$을 $m \in M$의 상이라 하고
$\mathfrak q = \text{Ann}(\overline{m}) \in \text{Ass}(M')$라 하자.
그러면 $m \not \in K$이므로 $Rm$의 지지집합은
$\mathfrak q_j$ 가운데 하나를 포함해야 한다.
$M_{\mathfrak q_j} = M'_{\mathfrak q_j}$이므로
$\overline{m}$은 $M'_{\mathfrak q_j}$에서 영으로 가지 않는다.
따라서 $\mathfrak q \subset \mathfrak q_j$이다
(실제로는 등식이다). 이는 $M'$의 모든 연관 소 아이디얼이
매립된 것이 아님을 뜻한다.

\medskip\noindent
$f$를 모든 $\mathfrak p_i$에 속하는 원소라 하자. 그러면
% P01-E0413
$D(f) \cap \text{supp}(K) = 0$이다.
$K_f = 0$이므로 $M_f = M'_f$이다.
\end{proof}

\begin{lemma}
\label{algebra-lemma-remove-embedded-primes-localize}
$R$을 뇌터 환이라 하고 $M$을 유한 $R$-가군이라 하자.
임의의 $f \in R$에 대하여 $(M')_f = (M_f)'$이다.
여기서 $M \to M'$과 $M_f \to (M_f)'$은 보조정리
\ref{algebra-lemma-remove-embedded-primes}에서 구성한 몫들이다.
\end{lemma}

\begin{proof}
생략한다.
\end{proof}

\begin{lemma}
\label{algebra-lemma-no-embedded-primes-endos}
$R$을 뇌터 환이라 하자.
$M$을 매립 연관 소 아이디얼이 없는 유한 $R$-가군이라 하자.
$I = \{x \in R \mid xM = 0\}$라 하자. 그러면 환 $R/I$에는
매립 소 아이디얼이 없다.
\end{lemma}

\begin{proof}
$R$을 $R/I$로 바꾸어도 된다. 따라서 $R$의 영이 아닌 모든 원소가
$M$ 위에서 영이 아니게 작용한다고 가정해도 된다.
보조정리 \ref{algebra-lemma-support-closed}에 의해 이는
$\Spec(R)$이 $M$의 지지집합과 같음을 함의한다.
$\mathfrak p$가 $R$의 매립 소 아이디얼이라고 가정하자.
$x \in R$을 소멸자가 $\mathfrak p$인 원소라 하자.
영이 아닌 가군 $N = xM \subset M$을 생각하자.
이는 $\mathfrak p$에 의해 소멸된다. 따라서 $N$의 임의의 연관
소 아이디얼 $\mathfrak q$는 $\mathfrak p$를 포함하고
$M$의 연관 소 아이디얼이기도 하다. 그러면 $\mathfrak q$는
$M$의 매립 연관 소 아이디얼이 되어 보조정리의 가정과 모순된다.
\end{proof}

\section{정칙열}
\label{algebra-section-regular-sequences}

\noindent
이 절에서는 정칙열의 몇 가지 기본 성질을 전개한다.

\begin{definition}
\label{algebra-definition-regular-sequence}
$R$을 환이라 하고 $M$을 $R$-가군이라 하자. $R$의 원소들의 열
$f_1, \ldots, f_r$가 다음 조건들을 만족하면 이를
{\it $M$-정칙열}이라 한다.
\begin{enumerate}
\item 각 $i = 1, \ldots, r$에 대하여 $f_i$는
$M/(f_1, \ldots, f_{i - 1})M$ 위의 비영인자이고,
\item 가군 $M/(f_1, \ldots, f_r)M$은 영이 아니다.
\end{enumerate}
$I$가 $R$의 아이디얼이고 $f_1, \ldots, f_r \in I$이면
$f_1, \ldots, f_r$를 {\it $I$ 안의 $M$-정칙열}이라 한다.
$M = R$이면 $f_1, \ldots, f_r$를 간단히
({\it $I$ 안의}) {\it 정칙열}이라 한다.
\end{definition}

\noindent
이 정의는 원소들 $f_1, \ldots, f_r$의 순서에 의존한다는 점에
주의해야 한다(아래의 예들을 보라). 어떤 논문이나 책에서는 가군
$M/(f_1, \ldots, f_r)M$이 영이 아니어야 한다는 조건을 생략한다.
그러면 정칙열이라는 성질이 국소화 아래에서 보존된다는 장점이 있다.
그러나 우리는 주로 $R$이 국소환일 때 가군의 깊이를 정의하기 위해
이 정의를 사용할 것이다. 이 경우 $f_i$들은 극대 아이디얼에 속해야 하는데,
이 조건은 $R$에서 국소화 $R_\mathfrak p$로 옮길 때 보존되지 않는다.

\begin{example}
\label{algebra-example-global-regular}
$k$를 체라 하자. 환 $k[x, y, z]$에서 열
$x, y(1-x), z(1-x)$는 정칙이지만 열
$y(1-x), z(1-x), x$는 정칙이 아니다.
\end{example}

\begin{example}
\label{algebra-example-local-regular}
$k$를 체라 하자. 환
$k[x, y, w_0, w_1, w_2, \ldots]/I$를 생각하자.
여기서 $I$는 $yw_i$, $i = 0, 1, 2, \ldots$와
$w_i - xw_{i + 1}$, $i = 0, 1, 2, \ldots$로 생성된다.
열 $x, y$는 정칙이지만 $y$는 영인자이다.
더 나아가 극대 아이디얼 $(x, y, w_i)$에서 국소화해도
여전히 예를 얻는다.
\end{example}

\begin{lemma}
\label{algebra-lemma-permute-xi}
$R$을 뇌터 국소환이라 하자.
$M$을 유한 $R$-가군이라 하자.
$x_1, \ldots, x_c$를 $M$-정칙열이라 하자. 그러면
% P01-E0414
$x_i$들의 임의의 순열도 정칙열이다.
\end{lemma}

\begin{proof}
먼저 $c = 2$인 경우를 다룬다.
$K \subset M$을 $x_2 : M \to M$의 핵이라 하자.
임의의 $z \in K$에 대하여 $x_2$가 $M/x_1M$ 위의 비영인자이므로
어떤 $z' \in M$에 대해 $z = x_1 z'$임을 안다.
$x_1$이 $M$ 위의 비영인자이므로 $x_2 z' = 0$이기도 하다.
따라서 $x_1 : K \to K$는 전사이다.
그러므로 나카야마 보조정리 \ref{algebra-lemma-NAK}에 의해 $K = 0$이다.
다음으로 $M/x_2M$ 위의 $x_1$에 의한 곱셈을 생각하자.
$z \in M$이 이 사상의 핵에 속하는 원소
$\overline{z} \in M/x_2M$로 가면, 어떤 $y \in M$에 대해
$x_1 z = x_2 y$이다. 그런데 $x_1, x_2$가 정칙열이므로
어떤 $y' \in M$에 대해 $y = x_1 y'$이다. 따라서
$x_1 ( z - x_2 y' ) =0$이고, 이에 따라 $z = x_2 y'$이며,
원하는 대로 $\overline{z} = 0$이다.

\medskip\noindent
일반적인 경우에는 임의의 순열이 인접한 첨자들의 전치들의 합성임을
관찰하자. 그러므로
$$
x_1, \ldots, x_{i-2}, x_i, x_{i-1}, x_{i + 1}, \ldots, x_c
$$
가 $M$-정칙열임을 보이면 충분하다. 이는 방금 다룬 경우를 가군
% P01-E0415
$M/(x_1, \ldots, x_{i-2})$와 길이 $2$의 정칙열
$x_{i-1}, x_i$에 적용하면 따른다.
\end{proof}

\begin{lemma}
\label{algebra-lemma-flat-increases-depth}
\begin{slogan}
평탄 국소환 준동형은 정칙열을 보존하고 반영한다.
\end{slogan}
$R, S$를 국소환이라 하자. $R \to S$를 평탄 국소환 준동형이라 하자.
$x_1, \ldots, x_r$를 $R$ 안의 열이라 하고 $M$을 $R$-가군이라 하자.
다음은 서로 동치이다.
\begin{enumerate}
\item $x_1, \ldots, x_r$는 $R$ 안의 $M$-정칙열이고,
\item $S$에서 $x_1, \ldots, x_r$의 상들은
$M \otimes_R S$-정칙열을 이룬다.
\end{enumerate}
\end{lemma}

\begin{proof}
보조정리 \ref{algebra-lemma-local-flat-ff}에 의해 $R \to S$가 충실히 평탄하기
때문이다.
\end{proof}

\begin{lemma}
\label{algebra-lemma-regular-sequence-in-neighbourhood}
$R$을 뇌터 환이라 하고 $M$을 유한 $R$-가군이라 하자.
$\mathfrak p$를 소 아이디얼이라 하자. $x_1, \ldots, x_r$를
$R$ 안의 열이라 하고, $R_{\mathfrak p}$에서의 그 상들이
$M_{\mathfrak p}$-정칙열을 이룬다고 하자. 그러면
$g \in R$, $g \not \in \mathfrak p$가 존재하여 $R_g$에서
$x_1, \ldots, x_r$의 상들이 $M_g$-정칙열을 이룬다.
\end{lemma}

\begin{proof}
다음과 같이 두자.
$$
K_i = \Ker\left(x_i : M/(x_1, \ldots, x_{i - 1})M \to
M/(x_1, \ldots, x_{i - 1})M\right).
$$
이는 가정에 의해 $\mathfrak p$에서의 국소화가 영인 유한
$R$-가군이다. 따라서 모든 $i = 1, \ldots, r$에 대해
$(K_i)_g = 0$이 되게 하는 $g \in R$, $g \not \in \mathfrak p$가
존재한다. 이 $g$가 원하는 원소이다.
\end{proof}

\begin{lemma}
\label{algebra-lemma-join-regular-sequences}
$A$를 환이라 하자. $I$를 $A$ 안의 정칙열
$f_1, \ldots, f_n$으로 생성되는 아이디얼이라 하자.
$g_1, \ldots, g_m \in A$를 그 상들
$\overline{g}_1, \ldots, \overline{g}_m$이 $A/I$에서 정칙열을 이루는
원소들이라 하자. 그러면 $f_1, \ldots, f_n, g_1, \ldots, g_m$은
$A$ 안의 정칙열이다.
\end{lemma}

\begin{proof}
이는 정의에서 바로 따른다.
\end{proof}

\begin{lemma}
\label{algebra-lemma-regular-sequence-short-exact-sequence}
$R$을 환이라 하자. $0 \to M_1 \to M_2 \to M_3 \to 0$을
$R$-가군들의 짧은 완전열이라 하자. $f_1, \ldots, f_r \in R$라 하자.
$f_1, \ldots, f_r$가 $M_1$-정칙이고 $M_3$-정칙이면
$f_1, \ldots, f_r$는 $M_2$-정칙이다.
\end{lemma}

\begin{proof}
보조정리 \ref{algebra-lemma-snake}에 의해 $f_1 : M_1 \to M_1$과
$f_1 : M_3 \to M_3$가 단사이면 $f_1 : M_2 \to M_2$도 단사이고,
다음 짧은 완전열을 얻는다.
$$
0 \to M_1/f_1M_1 \to M_2/f_1M_2 \to M_3/f_1M_3 \to 0
$$
보조정리는 이것과 $r$에 대한 귀납법에서 따른다. 몇몇 세부사항은
생략한다.
\end{proof}

\begin{lemma}
\label{algebra-lemma-regular-sequence-powers}
$R$을 환이라 하고 $M$을 $R$-가군이라 하자.
$f_1, \ldots, f_r \in R$라 하고 $e_1, \ldots, e_r > 0$을 정수들이라 하자.
그러면 $f_1, \ldots, f_r$가 $M$-정칙열인 것과
$f_1^{e_1}, \ldots, f_r^{e_r}$가 $M$-정칙열인 것은 서로 동치이다.
\end{lemma}

\begin{proof}
$r$에 대한 귀납법으로 이를 증명하겠다. $r = 1$이면 다음 두 가지
쉬운 사실에서 따른다. (a) $M$ 위의 비영인자의 거듭제곱은 $M$ 위의
비영인자이고, (b) $M$ 위의 비영인자의 약수는 $M$ 위의 비영인자이다.
$r > 1$이면 $M/f_1M$에 귀납가정을 적용하여
$f_1, f_2, \ldots, f_r$가 $M$-정칙열인 것과
$f_1, f_2^{e_2}, \ldots, f_r^{e_r}$가 $M$-정칙열인 것이 동치임을 얻는다.
따라서 주어진 $e > 0$에 대해 $f_1^e, f_2, \ldots, f_r$가
$M$-정칙열인 것과 $f_1, \ldots, f_r$가 $M$-정칙열인 것이 동치임을
보이면 충분하다. 이를 $e$에 대한 귀납법으로 증명하겠다.
$e = 1$인 경우는 자명하다. 두 가정 아래에서 모두 $f_1$은
비영인자이므로($r = 1$인 경우에 의해) 다음 짧은 완전열이 있다.
$$
0 \to M/f_1M \xrightarrow{f_1^{e - 1}} M/f_1^eM \to M/f_1^{e - 1}M \to 0
$$
$f_1, f_2, \ldots, f_r$가 $M$-정칙열이라고 하자.
그러면 귀납법에 의해 원소들 $f_2, \ldots, f_r$는
$M/f_1M$-정칙열이면서 $M/f_1^{e - 1}M$-정칙열이다.
보조정리 \ref{algebra-lemma-regular-sequence-short-exact-sequence}에 의해
$f_2, \ldots, f_r$는 $M/f_1^eM$-정칙이다. 따라서
$f_1^e, f_2, \ldots, f_r$는 $M$-정칙이다. 역으로
$f_1^e, f_2, \ldots, f_r$가 $M$-정칙열이라고 하자. 그러면
$f_2 : M/f_1^eM \to M/f_1^eM$가 단사이고, 따라서
$f_2 : M/f_1M \to M/f_1M$가 단사이며, 따라서 귀납법(!)에 의해
$f_2 : M/f_1^{e - 1}M \to M/f_1^{e - 1}M$가 단사이다. 따라서
$$
0 \to
M/(f_1, f_2)M \xrightarrow{f_1^{e - 1}}
M/(f_1^e, f_2)M \to
M/(f_1^{e - 1}, f_2)M \to 0
$$
은 보조정리 \ref{algebra-lemma-snake}에 의해 짧은 완전열이다. 이는
$r = 2$일 때 역을 증명한다. $r > 2$이면
$f_3 : M/(f_1^e, f_2)M \to M/(f_1^e, f_2)M$가 단사이고, 따라서
$f_3 : M/(f_1, f_2)M \to M/(f_1, f_2)M$가 단사이며, 이런 식으로
계속한다. 몇몇 세부사항은 생략한다.
\end{proof}

\begin{lemma}
\label{algebra-lemma-regular-sequence-in-polynomial-ring}
$R$을 환이라 하자. $f_1, \ldots, f_r \in R$가 단위 아이디얼을
생성하지 않는다고 하자. 다음은 서로 동치이다.
\begin{enumerate}
\item $f_1, \ldots, f_r$의 임의의 순열은 정칙열이다.
\item $f_1, \ldots, f_r$의 임의의 부분열(주어진 순서로)은
정칙열이다.
\item $f_1x_1, \ldots, f_rx_r$는 다항식환
$R[x_1, \ldots, x_r]$ 안의 정칙열이다.
\end{enumerate}
\end{lemma}

\begin{proof}
(1)이 (2)를 함의함은 명백하다. (2)가 (1)을 함의함을 $r$에 대한
귀납법으로 증명하겠다. $r = 1$인 경우는 자명하다.
$r = 2$인 경우는 $a, b \in R$가 정칙열이고 $b$가 비영인자이면
$b, a$가 정칙열이라는 말이다. $a$와 $b$가 모두 비영인자이면
$a : R/(b) \to R/(b)$의 핵은 $b : R/(a) \to R/(a)$의 핵과
동형이므로 이는 명백하다. 이제 $r > 2$인 경우를 생각하자.
(2)가 성립한다고 하고, 어떤 순열 $\sigma$에 대해
$f_{\sigma(1)}, \ldots, f_{\sigma(r)}$가 정칙열임을 증명하려 한다고 하자.
귀납법에 의해 $f_{\sigma(1)}, \ldots, f_{\sigma(r - 1)}$가
정칙열임을 이미 안다. 그러므로 $s = \sigma(r)$라 할 때 $f_s$가
$f_1, \ldots, \hat f_s, \ldots, f_r$로 나눈 몫에서 비영인자임을
보이면 충분하다. $s = r$이면 끝났다. $s < r$이면 $f_s$와 $f_r$가
환 $R/(f_1, \ldots, \hat f_s, \ldots, f_{r - 1})$에서 모두
비영인자임에 주의하자(다시 귀납가정에 의해).
이 환에서 $f_s, f_r$가 정칙열임을 아니, 길이 $2$인 열의 경우에 의해
$f_r, f_s$도 그러하다고 결론짓는다.

\medskip\noindent
$R[x_1, \ldots, x_r]/(f_1x_1, \ldots, f_ix_i)$는 $R$-가군으로서
다음 가군들의 직합임에 주의하자.
$$
R/I_E \cdot x_1^{e_1} \ldots x_r^{e_r}
$$
이들은 다중첨자 $E = (e_1, \ldots, e_r)$로 첨자화되며, 여기서
$I_E$는 $1 \leq j \leq i$이고 $e_j > 0$인 첨자에 대한 $f_j$들로
생성되는 아이디얼이다. 따라서
% P01-E0416
$f_{i + 1}x_i$가 이것 위의 비영인자인 것과 모든 $E$에 대해
$f_{i + 1}$가 $R/I_E$ 위의 비영인자인 것은 서로 동치이다.
모든 성분이 양수인 $E$를 택하면 $f_{i + 1}$가
$R/(f_1, \ldots, f_i)$ 위의 비영인자임을 알 수 있다. 따라서
(3)은 (2)를 함의한다. 역으로 (2)가 성립하면
$f_1, \ldots, f_i, f_{i + 1}$의 임의의 부분열이 정칙열이고,
특히 $f_{i + 1}$는 모든 $R/I_E$ 위의 비영인자이다.
이와 같이 (2)가 (3)을 함의함을 알 수 있다.
\end{proof}

\section{준정칙열}
\label{algebra-section-quasi-regular}

\noindent
정칙열보다 약간 약하고 사용하기 더 쉬운 준정칙열의 개념을 도입한다.
$R$을 환이라 하고 $f_1, \ldots, f_c \in R$라 하자.
$J = (f_1, \ldots, f_c)$라 두자. $M$을 $R$-가군이라 하자.
그러면 다음과 같은 표준 사상이 있다.
\begin{equation}
\label{algebra-equation-quasi-regular}
M/JM \otimes_{R/J} R/J[X_1, \ldots, X_c]
\longrightarrow
\bigoplus\nolimits_{n \geq 0} J^nM/J^{n + 1}M
\end{equation}
이는 등급 $R/J[X_1, \ldots, X_c]$-가군들의 사상이며 다음 규칙으로
정의된다.
$$
\overline{m} \otimes X_1^{e_1} \ldots X_c^{e_c} \longmapsto
f_1^{e_1} \ldots f_c^{e_c} m \bmod J^{e_1 + \ldots + e_c + 1}M.
$$
(\ref{algebra-equation-quasi-regular})은 항상 전사임에 주의하자.

\begin{definition}
\label{algebra-definition-quasi-regular-sequence}
$R$을 환이라 하자.
$M$을 $R$-가군이라 하자.
$R$의 원소들의 열 $f_1, \ldots, f_c$에 대하여
(\ref{algebra-equation-quasi-regular})이 동형이면 이를
{\it $M$-준정칙}이라 한다. $M = R$이면
$f_1, \ldots, f_c$를 간단히 {\it 준정칙열}이라 한다.
\end{definition}

\noindent
따라서 $f_1, \ldots, f_c$가 준정칙열이면
$$
R/J[X_1, \ldots, X_c] = \bigoplus\nolimits_{n \geq 0} J^n/J^{n + 1}
$$
이다. 여기서 $J = (f_1, \ldots, f_c)$이다. 준정칙열이라는 성질이
$f_1, \ldots, f_c$의 순서와 무관함은 명백하다.

\begin{lemma}
\label{algebra-lemma-regular-quasi-regular}
$R$을 환이라 하자.
\begin{enumerate}
\item $R$의 정칙열 $f_1, \ldots, f_c$는 준정칙열이다.
\item $M$이 $R$-가군이고 $f_1, \ldots, f_c$가
$M$-정칙열이라고 하자. 그러면 $f_1, \ldots, f_c$는
$M$-준정칙열이다.
\end{enumerate}
\end{lemma}

\begin{proof}
$J = (f_1, \ldots, f_c)$라 두자.
$c$에 대한 귀납법으로 첫 번째 주장을 증명한다.
임의의 관계식 $\sum_{|I| = n} a_I f^I \in J^{n + 1}$가
$a_I \in R$와 함께 주어졌을 때 실제로 모든 다중첨자 $I$에 대해
$a_I \in J$임을 보여야 한다. $J^{n + 1}$의 임의의 원소는
$b_I \in J$인 $\sum_{|I| = n} b_I f^I$ 꼴이므로, $a_I$를
$a_I - b_I$로 바꾸어 관계식이 $\sum_{|I| = n} a_I f^I = 0$이라고
가정해도 된다. 이를 다음과 같이 다시 쓸 수 있다.
$$
\sum\nolimits_{e = 0}^n
\left(
\sum\nolimits_{|I'| = n - e}
a_{I', e} f^{I'}
\right)
f_c^e
=
0
$$
여기와 아래에서 “프라임이 붙은” 다중첨자 $I'$는
$I' = (i_1, \ldots, i_{c - 1}, 0)$ 꼴이어야 한다.
% P01-E0418
$l \in \{0, \ldots, n\}$에 대한 귀납법으로, 관계식
$$
\sum\nolimits_{e = 0}^l
\left(
\sum\nolimits_{|I'| = n - e}
a_{I', e} f^{I'}
\right)
f_c^e
=
0
$$
이 주어지면 모든 $I', e$에 대해 $a_{I', e} \in J$임을 보이겠다.
구체적으로 $J' = (f_1, \ldots, f_{c-1})$라 두자.
$\sum\nolimits_{|I'| = n - l} a_{I', l} f^{I'}$는
$f_c^{l}$에 의해 $(J')^{n - l + 1}$ 안으로 보내짐을 관찰하자.
귀납가정($c$에 대한 귀납가정)에 의해
$f_c^l a_{I', l} \in J'$임을 알 수 있다.
$f_1, \ldots, f_c$가 정칙열이므로 $f_c$는 $R/J'$ 위의 영인자가
아니다. 따라서 $a_{I', l} \in J'$라고 결론짓는다.
이에 따라 항
$(\sum\nolimits_{|I'| = n - l} a_{I', l} f^{I'})f_c^l$을
$(\sum\nolimits_{|I'| = n - l + 1} f_c b_{I', l - 1}
f^{I'})f_c^{l-1}$ 꼴로 다시 쓸 수 있다. 그 결과 다음 꼴의 새
관계식을 얻는다.
$$
\left(\sum\nolimits_{|I'| = n - l + 1}
(a_{I', l-1} + f_c b_{I', l - 1}) f^{I'}\right)f_c^{l-1}
+
\sum\nolimits_{e = 0}^{l - 2}
\left(
\sum\nolimits_{|I'| = n - e}
a_{I', e} f^{I'}
\right)
f_c^e
=
0
$$
이제 귀납가정(이번에는 $l$에 대한 귀납가정)에 의해 모든
$a_{I', l-1} + f_c b_{I', l - 1} \in J$이고,
$e \leq l - 2$인 모든 $a_{I', e} \in J$이다.
이를 위에서 본 $a_{I', l} \in J' \subset J$와 합치면
귀납 단계의 증명이 끝난다.

\medskip\noindent
두 번째 주장은 임의의 형식적 표현
$F = \sum_{|I| = n} m_I X^I$, $m_I \in M$가 주어지고
$\sum m_I f^I \in J^{n + 1}M$이면 모든 계수 $m_I$가
% P01-E0417
$J$에 속한다는 뜻이다. 이는 위의 첫 번째 주장에 해당하는 결과를
증명한 것과 정확히 같은 방법으로 증명된다.
\end{proof}

\begin{lemma}
\label{algebra-lemma-flat-base-change-quasi-regular}
$R \to R'$을 평탄 환 사상이라 하자. $M$을 $R$-가군이라 하자.
$f_1, \ldots, f_r \in R$가 $M$-준정칙열을 이룬다고 하자.
그러면 $R'$에서 $f_1, \ldots, f_r$의 상들은
$M \otimes_R R'$-준정칙열을 이룬다.
\end{lemma}

\begin{proof}
$J = (f_1, \ldots, f_r)$, $J' = JR'$ 및 $M' = M \otimes_R R'$라 두자.
표준 사상
$\mu : R'/J'[X_1, \ldots X_r] \otimes_{R'/J'} M'/J'M' \to
\bigoplus (J')^nM'/(J')^{n + 1}M'$가 동형임을 보여야 한다.
$R \to R'$이 평탄하므로 완전열
$0 \to J^nM \to M$과
$0 \to J^{n + 1}M \to J^nM \to J^nM/J^{n + 1}M \to 0$은
$R'$과 텐서곱한 뒤에도 완전하다. 첫 번째 사실로부터
$J^nM \otimes_R R' = (J')^nM'$이고, 이어서
$(J')^nM'/(J')^{n + 1}M' = J^nM/J^{n + 1}M \otimes_R R'$임을 얻는다.
따라서 $\mu$는 가정에 의해 동형인
(\ref{algebra-equation-quasi-regular})과 $\text{id}_{R'}$의 텐서곱이며,
결론이 따른다.
\end{proof}

\begin{lemma}
\label{algebra-lemma-quasi-regular-sequence-in-neighbourhood}
$R$을 뇌터 환이라 하고 $M$을 유한 $R$-가군이라 하자.
$\mathfrak p$를 소 아이디얼이라 하자. $x_1, \ldots, x_c$를
$R$ 안의 열이라 하고, $R_{\mathfrak p}$에서의 그 상들이
$M_{\mathfrak p}$-준정칙열을 이룬다고 하자. 그러면
$g \in R$, $g \not \in \mathfrak p$가 존재하여 $R_g$에서
$x_1, \ldots, x_c$의 상들이 $M_g$-준정칙열을 이룬다.
\end{lemma}

\begin{proof}
(\ref{algebra-equation-quasi-regular})의 핵 $K$를 생각하자.
$M/JM \otimes_{R/J} R/J[X_1, \ldots, X_c]$는 유한
$R/J[X_1, \ldots, X_c]$-가군이고 $R/J[X_1, \ldots, X_c]$는
뇌터 환이므로, $K$도 유한 $R/J[X_1, \ldots, X_c]$-가군이다.
동차 생성원들 $k_1, \ldots, k_t \in K$를 택하자.
가정에 의해 각 $i = 1, \ldots, t$에 대해
$g_i \in R$, $g_i \not \in \mathfrak p$이고 $g_i k_i = 0$인
원소가 존재한다. 따라서 $g = g_1 \ldots g_t$가 원하는 원소이다.
\end{proof}

\begin{lemma}
\label{algebra-lemma-truncate-quasi-regular}
$R$을 환이라 하고 $M$을 $R$-가군이라 하자.
$f_1, \ldots, f_c \in R$를 $M$-준정칙열이라 하자.
임의의 $i$에 대하여
$\overline{R} = R/(f_1, \ldots, f_i)$의 원소들의 열
$\overline{f}_{i + 1}, \ldots, \overline{f}_c$는
$\overline{M} = M/(f_1, \ldots, f_i)M$-준정칙열이다.
\end{lemma}

\begin{proof}
$i = 1$인 경우를 증명하면 충분하다.
$\overline{J} = (\overline{f}_2, \ldots, \overline{f}_c) \subset \overline{R}$라 두자.
그러면
\begin{align*}
\overline{J}^n\overline{M}/\overline{J}^{n + 1}\overline{M}
& =
(J^nM + f_1M)/(J^{n + 1}M + f_1M) \\
& = J^nM / (J^{n + 1}M + J^nM \cap f_1M).
\end{align*}
따라서 보조정리를 증명하려면
% P01-E0419
$J^{n + 1}M + J^nM \cap f_1M = J^{n + 1}M + f_1J^{n - 1}M$임을
보이면 충분하다. 이 등식으로부터
$\bigoplus_{n \geq 0}
\overline{J}^n\overline{M}/\overline{J}^{n + 1}\overline{M}$이
$\bigoplus_{n \geq 0} J^nM/J^{n + 1}M \cong M/JM[X_1, \ldots, X_c]$을
% P01-E0420
$X_1$로 나눈 몫임을 알 수 있기 때문이다. 실제로
$J^nM \cap f_1M = f_1J^{n - 1}M$이다.
실제로 $m \not \in J^{n - 1}M$이면 $f_1m \not \in J^nM$이다.
이는 가정에 의해 $\bigoplus J^nM/J^{n + 1}M$이 다항식 대수
$M/J[X_1, \ldots, X_c]$이기 때문이다.
\end{proof}

\begin{lemma}
\label{algebra-lemma-quasi-regular-regular}
$(R, \mathfrak m)$을 뇌터 국소환이라 하자.
$M$을 영이 아닌 유한 $R$-가군이라 하자.
$f_1, \ldots, f_c \in \mathfrak m$을 $M$-준정칙열이라 하자.
그러면 $f_1, \ldots, f_c$는 $M$-정칙열이다.
\end{lemma}

\begin{proof}
$J = (f_1, \ldots, f_c)$라 두자.
$f_1$이 $M$ 위의 비영인자임을 보이자.
$x \in M$이 영이 아니라고 하자.
크룰 교집합 정리에 의해 $x \in J^rM$이지만
$x \not \in J^{r + 1}M$이 되게 하는 정수 $r$가 존재한다.
보조정리 \ref{algebra-lemma-intersect-powers-ideal-module-zero}를 보라.
그러면 $f_1 x \in J^{r + 1}M$이고, 가정한
$\bigoplus J^nM/J^{n + 1}M$의 구조에 의해 그 유는
$J^{r + 1}M/J^{r + 2}M$에서 영이 아니다.
따라서 $f_1x \not = 0$이다.

\medskip\noindent
이제 보조정리 \ref{algebra-lemma-truncate-quasi-regular}를 사용하여
$c$에 대한 귀납법으로 증명을 끝낼 수 있다.
\end{proof}

\begin{remark}[다른 종류의 정칙열]
\label{algebra-remark-koszul-regular}
논문 \cite{Kabele}에서 저자는 환 $R$의 원소들의 열
$x_1, \ldots, x_r$에 대한 두 가지 정칙성 조건을 더 논한다.
구체적으로 $i \geq 1$일 때
$H_i(K_{\bullet}(R, x_{\bullet})) = 0$이면 그 열을
{\it 코쥘-정칙}이라 한다. 여기서
$K_{\bullet}(R, x_{\bullet})$은 코쥘 복합체이다.
$H_1(K_{\bullet}(R, x_{\bullet})) = 0$이면 그 열을
{\it $H_1$-정칙}이라 한다.
다음 함의들이 성립한다. 정칙 $\Rightarrow$ 코쥘-정칙 $\Rightarrow$
$H_1$-정칙 $\Rightarrow$ 준정칙. 저자는 예들을 통해 $R$이
(비뇌터) 국소환이고 그 열이 $R$의 극대 아이디얼을 생성하는 경우에도
일반적으로 이 함의들을 역으로 돌릴 수 없음을 보인다.
이 개념들은 More on Algebra의 절
\ref{more-algebra-section-koszul-regular}에서 더 자세히 도입한다.
\end{remark}

\begin{remark}
\label{algebra-remark-join-quasi-regular-sequences}
$k$를 체라 하자. 환
$$
A = k[x, y, w, z_0, z_1, z_2, \ldots]/
(y^2z_0 - wx, z_0 - yz_1, z_1 - yz_2, \ldots)
$$
을 생각하자. 이 환에서 $x$는 비영인자이고 $A/xA$에서 $y$의 상은
준정칙열을 이룬다. 그러나 $x, y$가 $A$ 안의 준정칙열인 것은 아니다.
왜냐하면 $(x, y)/(x, y)^2$에서 $wx = 0$이지만
$A/(x, y)$에서 $w$는 영이 아니므로, $(x, y)/(x, y)^2$는
$A/(x, y)$ 위의 계수 2인 자유가군이 아니기 때문이다. 따라서
보조정리 \ref{algebra-lemma-join-regular-sequences}와 유사한 명제는
준정칙열에 대해서는 성립하지 않는다.
\end{remark}

\begin{lemma}
\label{algebra-lemma-quasi-regular-on-quotient}
$R$을 환이라 하자. $J = (f_1, \ldots, f_r)$를 $R$의 아이디얼이라 하자.
$M$을 $R$-가군이라 하자.
$\overline{R} = R/\bigcap_{n \geq 0} J^n$,
$\overline{M} = M/\bigcap_{n \geq 0} J^nM$라 두고,
$\overline{f}_i$로 $\overline{R}$에서 $f_i$의 상을 나타내자.
그러면 $f_1, \ldots, f_r$가 $M$-준정칙인 것과
$\overline{f}_1, \ldots, \overline{f}_r$가
$\overline{M}$-준정칙인 것은 서로 동치이다.
\end{lemma}

\begin{proof}
이는
% P01-E0421
$J^nM/J^{n + 1}M \cong
\overline{J}^n\overline{M}/\overline{J}^{n + 1}\overline{M}$이기
때문이다.
\end{proof}

\section{블로업 대수}
\label{algebra-section-blow-up}

\noindent
이 절에서는 블로업에 관한 몇 가지 초등적인 사실을 살펴본다.

\begin{definition}
\label{algebra-definition-blow-up}
$R$을 환이라 하자.
$I \subset R$를 아이디얼이라 하자.
\begin{enumerate}
\item 쌍 $(R, I)$에 딸린 {\it 블로업 대수}, 또는 {\it 리스 대수}는
다음 등급 $R$-대수이다.
$$
\text{Bl}_I(R) =
\bigoplus\nolimits_{n \geq 0} I^n =
R \oplus I \oplus I^2 \oplus \ldots
$$
여기서 직합항 $I^n$은 차수 $n$에 놓인다.
\item $a \in I$를 원소라 하자. 리스 대수에서 $a$를 차수 $1$의
원소로 본 것을 $a^{(1)}$로 나타내자. 그러면
{\it 아핀 블로업 대수} $R[\frac{I}{a}]$는
절 \ref{algebra-section-proj}에서 구성한 대수
$(\text{Bl}_I(R))_{(a^{(1)})}$이다.
\end{enumerate}
\end{definition}

\noindent
달리 말하면 $R[\frac{I}{a}]$의 원소는 $x \in I^n$인
$x/a^n$ 꼴의 식으로 표현된다. 두 표현 $x/a^n$과 $y/a^m$이
같은 원소를 정의할 필요충분조건은 어떤 $k \geq 0$에 대해
$a^k(a^mx - a^ny) = 0$인 것이다.

\begin{lemma}
\label{algebra-lemma-affine-blowup}
$R$을 환, $I \subset R$를 아이디얼, $a \in I$를 원소라 하자.
$R' = R[\frac{I}{a}]$를 아핀 블로업 대수라 하자. 그러면
\begin{enumerate}
\item $R'$에서 $a$의 상은 비영인자이고,
\item $IR' = aR'$이며,
\item $(R')_a = R_a$이다.
\end{enumerate}
\end{lemma}

\begin{proof}
위의 $R[\frac{I}{a}]$에 대한 설명에서 바로 따른다.
\end{proof}

\begin{lemma}
\label{algebra-lemma-blowup-base-change}
$R \to S$를 환 사상이라 하자. $I \subset R$를 아이디얼,
$a \in I$를 원소라 하자. $J = IS$라 두고 $b \in J$를 $a$의
상이라 하자. 그러면 $S[\frac{J}{b}]$는
$S \otimes_R R[\frac{I}{a}]$를 $b$의 어떤 거듭제곱에 의해
소멸되는 원소들의 아이디얼로 나눈 몫이다.
\end{lemma}

\begin{proof}
$S'$을 $S \otimes_R R[\frac{I}{a}]$를 그 $b$-거듭제곱 꼬임
원소들로 나눈 몫이라 하자. 환 사상
$$
S \otimes_R R[\textstyle{\frac{I}{a}}]
\longrightarrow
S[\textstyle{\frac{J}{b}}]
$$
은 전사이고, $S[\frac{J}{b}]$에서 $b$가 비영인자이므로
% P01-E0424
$a$-거듭제곱 꼬임을 소멸시킨다. 따라서 전사 사상
$S' \to S[\frac{J}{b}]$를 얻는다. 핵이 자명함을 보이기 위해
역사상을 구성한다. 구체적으로 $z = y/b^n$을
$S[\frac{J}{b}]$의 원소라 하자. 즉, $y \in J^n$이다.
$x_i \in I^n$이고 $s_i \in S$이도록 $y = \sum x_is_i$라 쓰자.
$z$를 $S'$에서 $\sum s_i \otimes x_i/a^n$의 유로 보낸다.
이 사상은 잘 정의된다. 실제로 사상
$S \otimes_R I^n \to J^n$의 핵의 원소는
% P01-E0424
$a^n$에 의해 소멸되므로 $S'$에서 영으로 가기 때문이다.
\end{proof}

\begin{example}
\label{algebra-example-rees-algebra-polynomial}
$R$을 환이라 하자. $P = R[t_1, \ldots, t_n]$을 다항식 대수라 하자.
$I = (t_1, \ldots, t_n) \subset P$라 하자. 정의
\ref{algebra-definition-blow-up}의 표기 아래에서 다음 동형이 있다.
$$
P[T_1, \ldots, T_n]/(t_iT_j - t_jT_i) \longrightarrow \text{Bl}_I(P)
$$
이는 $T_i$를 $t_i^{(1)}$로 보낸다. 이 사상이 잘 정의됨을 보이는 일은
독자에게 맡긴다. $I$는 $t_1, \ldots, t_n$으로 생성되므로
이 사상은 전사이다. 이 사상이 단사임을 보이려면 다음을 보여야 한다.
각 $e \geq 1$에 대해 $P$-가군 $I^e$는 다중첨자
$E = (e_1, \ldots, e_n)$가 차수 $|E| = e$를 달릴 때의 단항식
% P01-E0422
$t^E = t_1^{e_1} \ldots x_n^{e_n}$들로 생성되고, 이들 사이의 관계는
$|E| = |E'| = e$이고
$e_a + \delta_{a i} = e'_a + \delta_{a j},\ a = 1, \ldots, n$일 때의
$t_i t^E = t_j t^{E'}$뿐이다(크로네커 델타).
세부사항은 생략한다.
\end{example}

\begin{example}
\label{algebra-example-affine-blowup-algebra-polynomial}
$R$을 환이라 하자. $P = R[t_1, \ldots, t_n]$을 다항식 대수라 하자.
$I = (t_1, \ldots, t_n) \subset P$라 하자. $a = t_1$이라 하자.
정의 \ref{algebra-definition-blow-up}의 표기 아래에서 다음 동형이 있다.
$$
P[x_2, \ldots, x_n]/(t_1x_2 - t_2, \ldots, t_1x_n - t_n)
\longrightarrow
\textstyle{P[\frac{I}{a}] = P[\frac{I}{t_1}]}
$$
이는 $x_i$를 $t_i/t_1$로 보낸다. 이 사상이 잘 정의됨을 보이는 일은
독자에게 맡긴다. $I$는 $t_1, \ldots, t_n$으로 생성되므로
이 사상은 전사이다. 이 사상이 단사임을 보이기 위해 독자는 사상의
정의역과 공역이 $t_1$-꼬임이 없고 $t_1$을 가역화한 뒤 사상이
동형임을 논증할 수 있다. 보조정리 \ref{algebra-lemma-affine-blowup}을 보라.
또는 예 \ref{algebra-example-rees-algebra-polynomial}의 리스 대수에 대한
설명을 사용할 수 있다. 세부사항은 생략한다.
\end{example}

\begin{lemma}
\label{algebra-lemma-affine-blowup-quotient-description}
$R$을 환이라 하자. $I = (a_1, \ldots, a_n)$을 $R$의 아이디얼이라 하자.
$a = a_1$이라 하자. 그러면 다음 전사 사상이 있다.
$$
R[x_2, \ldots, x_n]/(a x_2 - a_2, \ldots, a x_n - a_n)
\longrightarrow
\textstyle{R[\frac{I}{a}]}
$$
그 핵은 정의역의 $a$-거듭제곱 꼬임이다.
\end{lemma}

\begin{proof}
$t_i$를 $a_i$로 보내는 환 사상
$P = \mathbf{Z}[t_1, \ldots, t_n] \to R$를 생각하자.
$J = (t_1, \ldots, t_n)$라 두자. 예
\ref{algebra-example-affine-blowup-algebra-polynomial}에 의해
$P[\frac{J}{t_1}] =
P[x_2, \ldots, x_n]/(t_1 x_2 - t_2, \ldots, t_1 x_n - t_n)$이다.
% P01-E0423
사상 $P \to A$에 보조정리 \ref{algebra-lemma-blowup-base-change}를 적용하면
결론이 따른다.
\end{proof}

\begin{lemma}
\label{algebra-lemma-blowup-in-principal}
$R$을 환, $I \subset R$를 아이디얼, $a \in I$를 원소라 하자.
$R' = R[\frac{I}{a}]$라 두자. $f \in R$가 $V(f) = V(I)$를 만족하면,
$f$는 $R'$에서 비영인자로 가고 $R'_f = R'_a = R_a$이다.
\end{lemma}

\begin{proof}
이후 별도의 언급 없이 보조정리 \ref{algebra-lemma-affine-blowup}의 결과들을
사용하겠다. 가정 $V(f) = V(I)$는
$V(fR') = V(IR') = V(aR')$를 함의한다. 따라서 어떤
$b, c \in R'$에 대해 $a^n = fb$이고 $f^m = ac$이다.
보조정리가 따른다.
\end{proof}

\begin{lemma}
\label{algebra-lemma-blowup-add-principal}
$R$을 환, $I \subset R$를 아이디얼, $a \in I$ 및 $f \in R$를
원소들이라 하자. $R' = R[\frac{I}{a}]$ 및
$R'' = R[\frac{fI}{fa}]$라 두자. 그러면 전사 $R$-대수 사상
$R' \to R''$가 있으며, 그 핵은 $R'$의 $f$-거듭제곱 꼬임
원소들의 집합이다.
\end{lemma}

\begin{proof}
사상은 $x \in I^n$인 $x/a^n$을 $f^nx/(fa)^n$으로 보내어 주어진다.
이 사상이 잘 정의되고 전사임은 바로 확인할 수 있다.
$R''$에서 $af$가 비영인자이므로
(보조정리 \ref{algebra-lemma-affine-blowup}), $f$-거듭제곱 꼬임 원소들의
집합은 영으로 보내진다. 역으로 $x \in R'$이고 모든 $n > 0$에 대해
$f^n x \not = 0$이라 하자. $R'$에서 $a$가 비영인자이므로 모든
$n$에 대해 $(af)^n x \not = 0$이다. 정의
\ref{algebra-definition-blow-up} 뒤의 $R''$에 대한 설명에 의해 $R''$에서
$x$의 상은 영이 아니다.
\end{proof}

\begin{lemma}
\label{algebra-lemma-blowup-reduced}
\begin{slogan}
기약성은 블로업 아래에서 불변이다
\end{slogan}
$R$이 기약환이면 $R$의 모든 (아핀) 블로업 대수는 기약환이다.
\end{lemma}

\begin{proof}
$I \subset R$를 아이디얼, $a \in I$를 원소라 하자.
$x \in I^n$인 $x/a^n$이 $R[\frac{I}{a}]$의 멱영원이라고 하자.
그러면 $(x/a^n)^m = 0$이다. 따라서 어떤 $N \geq 0$에 대해
$R$에서 $a^N x^m = 0$이다. 필요하다면 $N$을 키워서 어떤
$e \geq 0$에 대해 $N = me$라고 가정해도 된다.
그러면 $(a^e x)^m = 0$이고 $R$이 기약환이므로 $a^e x = 0$이다.
이는 $R[\frac{I}{a}]$에서 $x/a^n = 0$이라는 뜻이다.
\end{proof}

\begin{lemma}
\label{algebra-lemma-blowup-domain}
$R$을 정역, $I \subset R$를 아이디얼, $a \in I$를 영이 아닌
원소라 하자. 그러면 아핀 블로업 대수 $R[\frac{I}{a}]$는 정역이다.
\end{lemma}

\begin{proof}
$x \in I^n$, $y \in I^m$인 $x/a^n$, $y/a^m$을 곱이 영인
$R[\frac{I}{a}]$의 원소들이라 하자. 그러면 $R$에서
$a^N x y = 0$이다. $R$이 정역이므로 $x = 0$이거나 $y = 0$이라고
결론짓는다.
\end{proof}

\begin{lemma}
\label{algebra-lemma-blowup-dominant}
$R$을 환이라 하자. $I \subset R$를 아이디얼이라 하자.
$a \in I$라 하자. $a$가 $R$의 어느 극소 소 아이디얼에도
속하지 않으면 $\Spec(R[\frac{I}{a}]) \to \Spec(R)$의 상은
조밀하다.
\end{lemma}

\begin{proof}
$x \in R$에 대하여 $a^k x = 0$이면 $x$는 $R$의 모든 극소
소 아이디얼에 속하며, 따라서 멱영이다. 보조정리
\ref{algebra-lemma-Zariski-topology}를 보라. 그러므로
$R \to R[\frac{I}{a}]$의 핵은 멱영원들로 이루어진다.
이제 보조정리 \ref{algebra-lemma-image-dense-generic-points}로부터
결론이 따른다.
\end{proof}

\begin{lemma}
\label{algebra-lemma-valuation-ring-colimit-affine-blowups}
% P01-E0425
$(R, \mathfrak m)$을 분수체가 $K$인 국소 정역이라 하자.
$R$을 지배하는 부치환환 $R \subset A \subset K$를 택하자. 그러면
$$
A = \colim R[\textstyle{\frac{I}{a}}]
$$
는 다음 성질들을 만족하는 아핀 블로업
$R \to R[\frac{I}{a}]$들의 유향 여극한이다.
\begin{enumerate}
\item $a \in I \subset \mathfrak m$,
\item $I$는 유한 생성이고,
\item $\mathfrak m$에서 $R \to R[\frac{I}{a}]$의 섬유환은
영이 아니다.
\end{enumerate}
\end{lemma}

\begin{proof}
임의의 블로업 대수 $R[\frac{I}{a}]$는 $K$에 포함되는 정역이다.
보조정리 \ref{algebra-lemma-blowup-domain}을 보라. 보조정리의 내용은
$A$가 $a \in I$가 성질 (1), (2), (3)을 만족하는 것들의 유향
합집합이라는 것이다. $R[\frac{I}{a}] \subset A$이고
$R[\frac{J}{b}] \subset A$이면
$$
R[\textstyle{\frac{I}{a}}] \cup R[\textstyle{\frac{J}{b}}]  \subset
R[\textstyle{\frac{IJ}{ab}}] \subset A
$$
이다. 첫 번째 포함은 $x/a^n = b^nx/(ab)^n$이기 때문이고,
두 번째 포함은 $z \in (IJ)^n$이면 $x_i \in I^n$ 및
$y_i \in J^n$인 $z = \sum x_iy_i$로 쓸 수 있으며, 따라서
$z/(ab)^n = \sum (x_i/a^n)(y_i/b^n)$가 $A$에 속하기 때문이다.

\medskip\noindent
% P01-E0426
유한 부분집합 $E \subset A$를 생각하자.
$E = \{e_1, \ldots, e_n\}$이라 하자. 모든
$i = 1, \ldots, n$에 대해 $e_i = f_i/a$로 쓸 수 있게 하는
영이 아닌 $a \in R$를 택하자. $I = (f_1, \ldots, f_n, a)$라 두자.
$R[\frac{I}{a}] \subset A$라고 주장한다. $R[\frac{I}{a}]$의
원소는 원소들 $e_i$의 다항식으로 표현될 수 있으므로 이는 명백하다.
이 관찰로부터 보조정리가 바로 따른다.
\end{proof}

\section{Ext 군}
\label{algebra-section-ext}

\noindent
이 절에서는 뇌터 국소환 위의 깊이에 관한 몇 가지 기본 성질을
확립하기 위해 호몰로지 대수학을 조금 다룬다.

\begin{lemma}
\label{algebra-lemma-resolution-by-finite-free}
$R$을 환이라 하고 $M$을 $R$-가군이라 하자.
\begin{enumerate}
\item 다음과 같은 완전 복합체가 존재한다.
$$
\ldots \to F_2 \to F_1 \to F_0 \to M \to 0.
$$
여기서 $F_i$들은 자유 $R$-가군이다.
\item $R$이 뇌터 환이고 $M$이 유한 $R$-가군이면,
$F_i$가 유한 자유가군이 되도록 이 복합체를 택할 수 있다.
달리 말하면 다음과 같은 완전 복합체를 찾을 수 있다.
$$
\ldots \to R^{\oplus n_2} \to R^{\oplus n_1} \to R^{\oplus n_0} \to M \to 0.
$$
\end{enumerate}
\end{lemma}

\begin{proof}
뇌터인 경우만 설명하겠다.
첫 단계로 전사 $R^{n_0} \to M$을 택한다.
그다음 길이 $e$인 완전 복합체를 구성했다면
% P01-E0430
단지 전사 $R^{n_{e + 1}} \to
\Ker(R^{n_e} \to R^{n_{e-1}})$를 택하면 되며,
$R$이 뇌터 환이므로 이것이 가능하다.
\end{proof}

\begin{definition}
\label{algebra-definition-finite-free-resolution}
$R$을 환이라 하고 $M$을 $R$-가군이라 하자.
\begin{enumerate}
\item $M$의 (왼쪽) {\it 분해} $F_\bullet \to M$은 다음과 같은
$R$-가군들의 완전 복합체이다.
$$
\ldots \to F_2 \to F_1 \to F_0 \to M \to 0
$$
\item {\it 자유 $R$-가군들에 의한 $M$의 분해}는 각 $F_i$가
자유 $R$-가군인 분해 $F_\bullet \to M$이다.
\item {\it 유한 자유 $R$-가군들에 의한 $M$의 분해}는 각 $F_i$가
유한 자유 $R$-가군인 분해 $F_\bullet \to M$이다.
\end{enumerate}
\end{definition}

\noindent
우리는 흔히 $F_{\bullet}$라는 표기로 다음과 같은 $R$-가군들의
복합체를 나타낸다.
$$
\ldots \to F_i \to F_{i-1} \to \ldots
$$
이 경우 사상 $F_i \to F_{i-1}$을 흔히 $d_i$ 또는 $d_{F, i}$로
나타낸다. 이 절에서는 항상 $F_0$가 복합체의 마지막 영이 아닌
항이라고 가정한다. {\it 복합체 $F_{\bullet}$의 $i$번째 호몰로지군}은
군 $H_i = \Ker(d_{F, i})/\Im(d_{F, i + 1})$이다.
{\it 복합체의 사상 $\alpha : F_{\bullet} \to G_{\bullet}$}은
다음을 만족하는 사상들 $\alpha_i : F_i \to G_i$로 주어진다.
% P01-E0427
$\alpha_{i-1} \circ d_{F, i} = d_{G, i-1} \circ \alpha_i$.
그러한 사상은 호몰로지 위의 사상 $H_i(\alpha) :
H_i(F_{\bullet}) \to H_i(G_{\bullet})$을 유도한다.
$\alpha, \beta :  F_{\bullet} \to G_{\bullet}$가 복합체의
사상들이면, $\alpha$와 $\beta$ 사이의 {\it 호모토피}는
다음을 만족하는 사상들의 모음 $h_i : F_i \to G_{i + 1}$로 주어진다.
$\alpha_i - \beta_i = d_{G, i + 1} \circ h_i +
h_{i-1} \circ d_{F, i}$.
$\alpha$와 $\beta$ 사이에 호모토피가 존재하면 두 사상
$\alpha, \beta : F_{\bullet} \to G_{\bullet}$를
{\it 호모토픽}이라고 한다.

\medskip\noindent
다음과 같은 모양의 복합체 $F^{\bullet}$에 대해서도 매우 비슷한
표기를 사용할 것이다.
$$
\ldots \to F^i \xrightarrow{d^i} F^{i + 1} \to \ldots
$$
복합체의 사상, 호모토피 등도 있다. 이 경우
$H^i(F^{\bullet}) =
\Ker(d^i)/\Im(d^{i - 1})$로 두고 이를
{\it $i$번째 코호몰로지군}이라 한다.

\begin{lemma}
\label{algebra-lemma-homotopic-equal-homology}
호모토픽인 임의의 두 복합체 사상은 (코)호몰로지군 위에 같은 사상을
유도한다.
\end{lemma}

\begin{proof}
생략한다.
\end{proof}

\begin{lemma}
\label{algebra-lemma-compare-resolutions}
$R$을 환이라 하자. $M \to N$을 $R$-가군의 사상이라 하자.
$N_\bullet \to N$을 임의의 분해라 하자. 다음이 각 $F_i$가
자유 $R$-가군인 $R$-가군들의 복합체라고 하자.
$$
\ldots \to F_2 \to F_1 \to F_0 \to M
$$
그러면 다음이 성립한다.
\begin{enumerate}
\item 다음 그림을 가환으로 만드는 복합체의 사상
$F_\bullet \to N_\bullet$이 존재하고,
$$
\xymatrix{
F_0 \ar[r] \ar[d] & M \ar[d] \\
N_0 \ar[r] & N
}
$$
\item (1)과 같은 임의의 두 사상
$\alpha, \beta : F_\bullet \to N_\bullet$는 호모토픽이다.
\end{enumerate}
\end{lemma}

\begin{proof}
(1)의 증명. $F_0$가 자유가군이므로 사상 $F_0 \to M \to N$을
올리는 사상 $F_0 \to N_0$를 찾을 수 있다. 이로부터
$F_1 \to F_0 \to N_0$인 사상이 유도되며, 그 상은
$N_1 \to N_0$의 상 안에 들어간다. $F_1$이 자유가군이므로 이를
사상 $F_1 \to N_1$로 올릴 수 있다. 다시 이로부터
$F_2 \to F_1 \to N_1$인 사상이 유도되며, 이는 $N_0$에서 영으로
간다. $N_\bullet$가 완전하므로 이 사상의 상은
$N_2 \to N_1$의 상에 포함된다. 따라서 올려서 사상
$F_2 \to N_2$를 얻을 수 있다. 이 과정을 반복한다.

\medskip\noindent
(2)의 증명. $\alpha, \beta$가 호모토픽임을 보이려면 그 차
$\gamma = \alpha - \beta$가 영과 호모토픽임을 보이면 충분하다.
$\gamma_0 : F_0 \to N_0$의 상은 $N_1 \to N_0$의 상에 포함된다.
따라서 $\gamma_0$를 사상 $h_0 : F_0 \to N_1$로 올릴 수 있다.
사상 $\gamma_1' = \gamma_1 - h_0 \circ d_{F, 1}$을 생각하자.
$h_0$의 선택에 의해 $\gamma_1'$의 상은
$N_1 \to N_0$의 핵에 포함된다. $N_\bullet$가 완전하므로
$\gamma_1'$를 사상 $h_1 : F_1 \to N_2$로 올릴 수 있다.
이 단계에서 $\gamma_1 = h_0 \circ d_{F, 1}
+ d_{N, 2} \circ h_1$이다. 이 과정을 반복한다.
\end{proof}

\noindent
이제 군 $\Ext^i_R(M, N)$을 정의할 준비가 되었다.
즉, 보조정리 \ref{algebra-lemma-resolution-by-finite-free}를 사용하여
$M$의 자유 $R$-가군들에 의한 분해 $F_{\bullet}$을 택한다.
다음 (코호몰로지) 복합체를 생각하자.
$$
\Hom_R(F_\bullet, N) :
\Hom_R(F_0, N) \to
\Hom_R(F_1, N) \to
\Hom_R(F_2, N) \to \ldots
$$
$i \geq 0$에 대해 $\Ext^i_R(M, N)$을 이 복합체의 $i$번째
% P01-E1389
코호몰로지군으로 정의한다\footnote{이 단계에서는 아마 “그” Ext 군보다
“하나의” Ext 군이라고 하는 편이 더 적절할 것이다.}.
$i < 0$에 대해서는 $\Ext^i_R(M, N) = 0$으로 둔다.
계속하기 전에 다음을 지적하자.
$$
\Ext^0_R(M, N) = \Ker(\Hom_R(F_0, N) \to \Hom_R(F_1, N)) =
\Hom_R(M, N)
$$
이는 완전열 $F_1 \to F_0 \to M \to 0$에 보조정리
\ref{algebra-lemma-hom-exact}의 (1)을 적용할 수 있기 때문이다.
다음 보조정리는 이 정의가 어떤 의미에서 잘 정의되는지 설명한다.

\begin{lemma}
\label{algebra-lemma-ext-welldefined}
$R$을 환이라 하자. $M_1, M_2, N$을 $R$-가군이라 하자.
$F_{\bullet}$이 가군 $M_1$의 자유 분해이고,
$G_{\bullet}$이 가군 $M_2$의 자유 분해라고 하자.
$\varphi : M_1 \to M_2$를 가군 사상이라 하자.
$\alpha : F_{\bullet} \to G_{\bullet}$를
$M_1 = \Coker(d_{F, 1}) \to M_2 = \Coker(d_{G, 1})$ 위에서
$\varphi$를 유도하는 복합체 사상이라 하자. 보조정리
\ref{algebra-lemma-compare-resolutions}를 보라. 그러면 유도된 사상들
% P01-E0428
$$
H^i(\alpha) :
H^i(\Hom_R(F_{\bullet}, N))
\longrightarrow
H^i(\Hom_R(G_{\bullet}, N))
$$
은 $\alpha$의 선택과 무관하다. $\varphi$가 동형이면 모든 사상
$H^i(\alpha)$도 동형이다. $M_1 = M_2$, $F_\bullet = G_\bullet$이고
$\varphi$가 항등사상이면 모든 사상
% P01-E0429
$H_i(\alpha)$도 항등사상이다.
\end{lemma}

\begin{proof}
$\varphi$를 유도하는 다른 사상 $\beta : F_{\bullet} \to G_{\bullet}$는
보조정리 \ref{algebra-lemma-compare-resolutions}에 의해 $\alpha$와
호모토픽이다. 따라서 사상들
% P01-E0428
$\Hom_R(F_\bullet, N) \to
\Hom_R(G_\bullet, N)$은 호모토픽이다.
따라서 독립성에 관한 결과는 보조정리
\ref{algebra-lemma-homotopic-equal-homology}에서 따른다.

\medskip\noindent
$\varphi$가 동형이라고 하자.
$\psi : M_2 \to M_1$을 그 역사상이라 하자.
$\psi :
M_2 = \Coker(d_{G, 1}) \to M_1 = \Coker(d_{F, 1})$을 유도하는
% P01-E1390
사상 $\beta : G_{\bullet} \to F_{\bullet}$을 택하자.
보조정리 \ref{algebra-lemma-compare-resolutions}를 보라.
좋다. 이제 사상
$H^i(\alpha) \circ H^i(\beta) =
H^i(\alpha \circ \beta)$를 생각하자. 위에서 보인 바에 의해
사상 $H^i(\alpha \circ \beta)$는 사상
$H^i(\text{id}_{G_{\bullet}}) = \text{id}$와 {\it 같다}.
합성 $H^i(\beta) \circ H^i(\alpha)$에 대해서도 마찬가지이다.
따라서 $H^i(\alpha)$와 $H^i(\beta)$는 서로의 역사상이다.
\end{proof}

\begin{lemma}
\label{algebra-lemma-long-exact-seq-ext}
$R$을 환이라 하자. $M$을 $R$-가군이라 하자.
$0 \to N' \to N \to N'' \to 0$을 짧은 완전열이라 하자.
그러면 다음 긴 완전열을 얻는다.
$$
\begin{matrix}
0
\to \Hom_R(M, N')
\to \Hom_R(M, N)
\to \Hom_R(M, N'')
\\
\phantom{0\ }
\to \Ext^1_R(M, N')
\to \Ext^1_R(M, N)
\to \Ext^1_R(M, N'')
\to \ldots
\end{matrix}
$$
\end{lemma}

\begin{proof}
자유 분해 $F_{\bullet} \to M$을 택하자.
% P01-E1392
각 $F_i$가 자유가군이므로 다음 복합체들의 짧은 완전열을 얻는다.
$$
0 \to
\Hom_R(F_{\bullet}, N') \to
\Hom_R(F_{\bullet}, N) \to
\Hom_R(F_{\bullet}, N'') \to
0
$$
따라서 여기에 뱀 보조정리를 적용하여 긴 완전열을 얻는다.
\end{proof}

\begin{lemma}
\label{algebra-lemma-reverse-long-exact-seq-ext}
$R$을 환이라 하자. $N$을 $R$-가군이라 하자.
$0 \to M' \to M \to M'' \to 0$을 짧은 완전열이라 하자.
그러면 다음 긴 완전열을 얻는다.
$$
\begin{matrix}
0
\to \Hom_R(M'', N)
\to \Hom_R(M, N)
\to \Hom_R(M', N)
\\
\phantom{0\ }
\to \Ext^1_R(M'', N)
\to \Ext^1_R(M, N)
\to \Ext^1_R(M', N)
\to \ldots
\end{matrix}
$$
\end{lemma}

\begin{proof}
$M'$와 $M''$의 생성원 집합
$\{m'_{i'}\}_{i' \in I'}$와
$\{m''_{i''}\}_{i'' \in I''}$을 택하자.
각 $i'' \in I''$에 대해 원소 $m''_{i''} \in M''$의 올림
$\tilde m''_{i''} \in M$을 택하자.
$F' = \bigoplus_{i' \in I'} R$,
$F'' = \bigoplus_{i'' \in I''} R$ 및 $F = F' \oplus F''$로 두자.
이 자유가군들의 생성원들을 대응하는 선택된 생성원들로 보내면
전사 $R$-가군 사상 $F' \to M'$, $F'' \to M''$ 및
$F \to M$을 얻는다. 다음 짧은 완전열들의 사상을 얻는다.
$$
\begin{matrix}
0 & \to & M' & \to & M & \to & M'' & \to & 0 \\
& & \uparrow & & \uparrow & & \uparrow \\
0 & \to & F' & \to & F & \to & F'' & \to & 0 \\
\end{matrix}
$$
뱀 보조정리에 의해 핵들의 열
% P01-E1391
$0 \to K' \to K \to K'' \to 0$은 $R$-가군들의 짧은 완전열이다.
따라서 이 과정을 무한히 계속할 수 있다. 달리 말하면 다음 그림에
들어맞는 분해들의 짧은 완전열을 얻는다.
$$
\begin{matrix}
0 & \to & M' & \to & M & \to & M'' & \to & 0 \\
& & \uparrow & & \uparrow & & \uparrow \\
0 & \to & F_\bullet' & \to & F_\bullet & \to & F_\bullet'' & \to & 0 \\
\end{matrix}
$$
각 열 $0 \to F'_n \to F_n \to F''_n \to 0$은 구성에 의해
분할 완전하므로, 함자 $\Hom_R(-, N)$을 적용하여 다음 복합체들의
짧은 완전열을 얻는다.
$$
0 \to
\Hom_R(F''_{\bullet}, N) \to
\Hom_R(F_{\bullet}, N) \to
\Hom_R(F'_{\bullet}, N) \to
0
$$
따라서 여기에 뱀 보조정리를 적용하여 긴 완전열을 얻는다.
\end{proof}

\begin{lemma}
\label{algebra-lemma-annihilate-ext}
$R$을 환이라 하자. $M$, $N$을 $R$-가군이라 하자.
$xN = 0$ 또는 $xM = 0$을 만족하는 임의의 $x\in R$는
각 가군 $\Ext^i_R(M, N)$을 소멸시킨다.
\end{lemma}

\begin{proof}
$M$의 자유 분해 $F_{\bullet}$을 택하자.
$\Ext^i_R(M, N)$은 복합체 $\Hom_R(F_{\bullet}, N)$의
코호몰로지로 정의되므로 $xN = 0$인 경우에는 보조정리가 명백하다.
$xM = 0$이면 $F_{\bullet}$ 위에서 $x$에 의한 곱셈은
$M$ 위의 영사상을 올린다. 따라서 보조정리
\ref{algebra-lemma-ext-welldefined}에 의해 이것이 Ext 군 위에 유도하는
사상은 영사상이 유도하는 사상과 같다.
\end{proof}

\begin{lemma}
\label{algebra-lemma-ext-noetherian}
$R$을 뇌터 환이라 하자. $M$, $N$을 유한 $R$-가군이라 하자.
그러면 모든 $i$에 대해 $\Ext^i_R(M, N)$은 유한 $R$-가군이다.
\end{lemma}

\begin{proof}
보조정리 \ref{algebra-lemma-resolution-by-finite-free}를 보라.
$\Ext^i_R(M, N)$은 각 $F_n$이 유한 자유 $R$-가군인 복합체
$\Hom_R(F_\bullet, N)$의 코호몰로지군으로 계산되므로 성립한다.
\end{proof}

\section{깊이}
\label{algebra-section-depth}

\noindent
다음은 우리가 사용할 정의이다.

\begin{definition}
\label{algebra-definition-depth}
$R$을 환, $I \subset R$를 아이디얼이라 하자.
$M$을 유한 $R$-가군이라 하자. $M$의 {\it $I$-깊이}를
$\text{depth}_I(M)$로 나타내고 다음과 같이 정의한다.
\begin{enumerate}
\item $IM \not = M$이면 $\text{depth}_I(M)$은
$I$ 안의 $M$-정칙열들의 길이들의 집합
$\{0, 1, 2, \ldots, \infty\}$에서의 상한이고,
\item $IM = M$이면 $\text{depth}_I(M) = \infty$로 둔다.
\end{enumerate}
$(R, \mathfrak m)$이 국소환이면
$\text{depth}_{\mathfrak m}(M)$을 간단히 $M$의 {\it 깊이}라 한다.
\end{definition}

\noindent
설명. 정의 \ref{algebra-definition-regular-sequence}에 따르면 빈 열은 영가군
위의 정칙열이 아니지만, 실제 목적을 위해서는 $0$ 가군의 깊이를
$+\infty$로 두는 것이 편리하다는 사실을 알게 될 것이다.
$I = R$이면 모든 유한 $R$-가군 $M$에 대해
$\text{depth}_I(M) = \infty$임을 주목하자.
$I$가 $R$의 야코브슨 근기에 포함되면
(예를 들어 $R$이 국소환이고 $I \subset \mathfrak m_R$이면),
나카야마 보조정리에 의해 $M \not = 0 \Rightarrow IM \not = M$이다.
가군 $M$의 $I$-깊이가 $0$일 필요충분조건은 $M$이 영이 아니고
$I$가 $M$ 위의 비영인자를 포함하지 않는 것이다.

\medskip\noindent
예 \ref{algebra-example-global-regular}는 환이 뇌터 환이어도 깊이가 잘
거동하지 않음을 보이고, 예 \ref{algebra-example-local-regular}는 환이
국소환이지만 비뇌터 환이면 깊이가 잘 거동하지 않음을 보인다.
환이 뇌터 국소환이면 깊이가 잘 거동함을 보게 될 것이다.

\begin{lemma}
\label{algebra-lemma-depth-weak-sequence}
$R$을 환, $I \subset R$를 아이디얼, $M$을 유한 $R$-가군이라 하자.
그러면 $\text{depth}_I(M)$은 $f_i$가
$M/(f_1, \ldots, f_{i - 1})M$ 위의 비영인자인 열
$f_1, \ldots, f_r \in I$들의 길이들의 상한과 같다.
\end{lemma}

\begin{proof}
$IM = M$이라고 하자. 그러면 보조정리 \ref{algebra-lemma-NAK}에 의해
$f : M \to M$이 $\text{id}_M$이 되게 하는 $f \in I$가 존재한다.
따라서 $f, 0, 0, 0, \ldots$는 연속적인 비영인자들의 무한열이고,
이 경우 두 값이 일치한다. $IM \not =  M$이면 보조정리에서 말한
열이 $M$-정칙열이므로 역시 두 값이 일치한다.
\end{proof}

\begin{lemma}
\label{algebra-lemma-bound-depth}
$(R, \mathfrak m)$을 뇌터 국소환이라 하자.
$M$을 영이 아닌 유한 $R$-가군이라 하자. 그러면
$\dim(\text{Supp}(M)) \geq \text{depth}(M)$이다.
\end{lemma}

\begin{proof}
$\dim(\text{Supp}(M))$에 대해 귀납한다.
$\dim(\text{Supp}(M)) = 0$이면
$\text{Supp}(M) = \{\mathfrak m\}$이고, 따라서
$\text{Ass}(M) = \{\mathfrak m\}$이다
(보조정리 \ref{algebra-lemma-ass-support}와 \ref{algebra-lemma-ass-zero}에 의해).
그러므로 예를 들어 보조정리 \ref{algebra-lemma-ideal-nonzerodivisor}에 의해
$M$의 깊이는 영이다. 귀납 단계에서는
$\dim(\text{Supp}(M)) > 0$이라고 가정한다.
$f_i$가 $M/(f_1, \ldots, f_{i - 1})M$ 위의 비영인자가 되도록
$\mathfrak m$의 원소들로 이루어진 열 $f_1, \ldots, f_d$를 택하자.
보조정리 \ref{algebra-lemma-depth-weak-sequence}에 의해
$\dim(\text{Supp}(M)) \geq d$임을 보이면 충분하다.
그렇지 않으면 보조정리가 성립하므로 $d > 0$이라고 가정해도 된다.
보조정리 \ref{algebra-lemma-one-equation-module}에 의해
$\dim(\text{Supp}(M/f_1M)) = \dim(\text{Supp}(M)) - 1$이다.
귀납법에 의해 원하는 대로
$\dim(\text{Supp}(M/f_1M)) \geq d - 1$이다.
\end{proof}

\begin{lemma}
\label{algebra-lemma-depth-finite-noetherian}
$R$을 뇌터 환, $I \subset R$를 아이디얼이라 하고,
$M$을 $IM \not = M$을 만족하는 유한하고 영이 아닌 $R$-가군이라 하자.
그러면 $\text{depth}_I(M) < \infty$이다.
\end{lemma}

\begin{proof}
$M/IM$이 영이 아니므로 보조정리 \ref{algebra-lemma-support-zero}에 의해
$\mathfrak p \in \text{Supp}(M/IM)$을 택할 수 있다.
그러면 $(M/IM)_\mathfrak p \not = 0$이고, 이는
$I \subset \mathfrak p$를 함의하며 국소화가 완전하므로
$M_\mathfrak p \not = IM_\mathfrak p$도 함의한다.
$f_1, \ldots, f_r \in I$를 $M$-정칙열이라 하자.
$(f_1, \ldots, f_r) \subset I$이므로
$M_\mathfrak p/(f_1, \ldots, f_r)M_\mathfrak p$는 영이 아니다.
국소화는 평탄하므로 $f_1, \ldots, f_r$의 상들은
$I_\mathfrak p$ 안의 $M_\mathfrak p$-정칙열을 이룬다.
이것이 $I$ 안의 모든 $M$-정칙열에 대해 성립하므로
$\text{depth}_I(M) \leq \text{depth}_{I_\mathfrak p}(M_\mathfrak p)$이다.
후자는 $\leq \text{depth}(M_\mathfrak p)$이고,
보조정리 \ref{algebra-lemma-bound-depth}에 의해 $< \infty$이다.
\end{proof}

\begin{lemma}
\label{algebra-lemma-depth-ext}
$R$을 극대 아이디얼이 $\mathfrak m$인 뇌터 국소환이라 하자.
$M$을 영이 아닌 유한 $R$-가군이라 하자. 그러면
$\text{depth}(M)$은 $\Ext^i_R(R/\mathfrak m, M)$이 영이 아닌
가장 작은 정수 $i$와 같다.
\end{lemma}

\begin{proof}
$M$의 깊이를 $\delta(M)$로 나타내고,
$\Ext^i_R(R/\mathfrak m, M)$이 영이 아닌 가장 작은 정수 $i$를
$i(M)$으로 나타내자. 잠시 뒤에 $i(M) < \infty$임을 보게 될 것이다.
보조정리 \ref{algebra-lemma-ideal-nonzerodivisor}에 의해
$\delta(M) = 0$일 필요충분조건은 $i(M) = 0$이다. 실제로
$\mathfrak m \in \text{Ass}(M)$이라는 것은 정확히
$i(M) = 0$이라는 뜻이다. 따라서 $\delta(M)$ 또는 $i(M)$이
$> 0$이면, (a) $x$가 $M$ 위의 비영인자이고
(b) $\text{depth}(M/xM) = \delta(M) - 1$이 되게 하는
$x \in \mathfrak m$을 택할 수 있다. 보조정리
\ref{algebra-lemma-long-exact-seq-ext}에 의해 짧은 완전열
$0 \to M \to M \to M/xM \to 0$에 딸린 Ext 군들의 긴 완전열을
생각하자.
$$
\begin{matrix}
0
\to \Hom_R(\kappa, M)
\to \Hom_R(\kappa, M)
\to \Hom_R(\kappa, M/xM)
\\
\phantom{0\ }
\to \Ext^1_R(\kappa, M)
\to \Ext^1_R(\kappa, M)
\to \Ext^1_R(\kappa, M/xM)
\to \ldots
\end{matrix}
$$
$x \in \mathfrak m$이므로 모든 사상
$\Ext^i_R(\kappa, M) \to \Ext^i_R(\kappa, M)$은 영사상이다.
보조정리 \ref{algebra-lemma-annihilate-ext}를 보라. 따라서
$i(M/xM) = i(M) - 1$임은 명백하다. $\delta(M)$에 대한 귀납법으로
증명이 끝난다.
\end{proof}

\begin{lemma}
\label{algebra-lemma-depth-in-ses}
$R$을 뇌터 국소환이라 하자. $0 \to N' \to N \to N'' \to 0$을
영이 아닌 유한 $R$-가군들의 짧은 완전열이라 하자.
\begin{enumerate}
\item
$\text{depth}(N) \geq \min\{\text{depth}(N'), \text{depth}(N'')\}$
\item
$\text{depth}(N'') \geq \min\{\text{depth}(N), \text{depth}(N') - 1\}$
\item
$\text{depth}(N') \geq \min\{\text{depth}(N), \text{depth}(N'') + 1\}$
\end{enumerate}
\end{lemma}

\begin{proof}
보조정리 \ref{algebra-lemma-depth-ext}를 보라.
Ext 군 $\Ext^i(\kappa, N)$을 사용한 깊이의 특징지음을 적용하고,
보조정리 \ref{algebra-lemma-long-exact-seq-ext}에서 오는 다음 긴
코호몰로지 완전열을 사용한다.
$$
\begin{matrix}
0
\to \Hom_R(\kappa, N')
\to \Hom_R(\kappa, N)
\to \Hom_R(\kappa, N'')
\\
\phantom{0\ }
\to \Ext^1_R(\kappa, N')
\to \Ext^1_R(\kappa, N)
\to \Ext^1_R(\kappa, N'')
\to \ldots
\end{matrix}
$$
\end{proof}

\begin{lemma}
\label{algebra-lemma-depth-drops-by-one}
$R$을 뇌터 국소환이라 하고 $M$을 영이 아닌 유한 $R$-가군이라 하자.
\begin{enumerate}
\item $x \in \mathfrak m$이 $M$ 위의 비영인자이면
$\text{depth}(M/xM) = \text{depth}(M) - 1$이다.
\item 임의의 $M$-정칙열 $x_1, \ldots, x_r$는 길이가
$\text{depth}(M)$인 $M$-정칙열로 연장될 수 있다.
\end{enumerate}
\end{lemma}

\begin{proof}
(2)는 (1)의 형식적 귀결이다. $x \in R$를 (1)과 같다고 하자.
짧은 완전열 $0 \to M \to M \to M/xM \to 0$과 보조정리
\ref{algebra-lemma-depth-in-ses}에 의해 깊이는 많아야 1만큼 감소한다.
다른 한편, $x_1, \ldots, x_r \in \mathfrak m$이
$M/xM$의 정칙열이면 $x, x_1, \ldots, x_r$는 $M$의 정칙열이다.
따라서 깊이는 적어도 1만큼 감소한다.
\end{proof}

\begin{lemma}
\label{algebra-lemma-inherit-minimal-primes}
$(R, \mathfrak m)$을 뇌터 국소환, $M$을 유한 $R$-가군이라 하자.
$x \in \mathfrak m$, $\mathfrak p \in \text{Ass}(M)$라 하고,
$\mathfrak q$가 $\mathfrak p + (x)$ 위의 극소 소 아이디얼이라고 하자.
그러면 어떤 $n \geq 1$에 대해
$\mathfrak q \in \text{Ass}(M/x^nM)$이다.
\end{lemma}

\begin{proof}
$N \cong R/\mathfrak p$인 부분가군 $N \subset M$을 택하자.
아르틴--리스 보조정리(보조정리 \ref{algebra-lemma-Artin-Rees})에 의해
$N \cap x^nM \subset xN$이 되게 하는 $n > 0$을 택할 수 있다.
$N \to M \to M/x^nM$의 상을
$\overline{N} \subset M/x^nM$으로 나타내자.
보조정리 \ref{algebra-lemma-ass}에 의해
$\mathfrak q \in \text{Ass}(\overline{N})$임을 보이면 충분하다.
$n$의 선택에 의해 전사
% P01-E0431
$\overline{N} \to N/xN = R/\mathfrak p + (x)$가 있고,
따라서 $\mathfrak q$는 $\overline{N}$의 지지집합에 속한다.
$\overline{N}$은 $x^n$과 $\mathfrak p$에 의해 소멸되므로
$\mathfrak q$는 $\overline{N}$의 지지집합에 속하는 소 아이디얼들
가운데 극소이다. 따라서 보조정리
\ref{algebra-lemma-ass-minimal-prime-support}에 의해 $\mathfrak q$는
$\overline{N}$의 연관 소 아이디얼이다.
\end{proof}

\begin{lemma}
\label{algebra-lemma-depth-dim-associated-primes}
$(R, \mathfrak m)$을 뇌터 국소환, $M$을 유한 $R$-가군이라 하자.
$\mathfrak p \in \text{Ass}(M)$에 대해
$\dim(R/\mathfrak p) \geq \text{depth}(M)$이다.
\end{lemma}

\begin{proof}
$\mathfrak m \in \text{Ass}(M)$이면 $\mathfrak m$의 모든 원소에 의해
소멸되는 영이 아닌 원소 $x \in M$이 존재한다.
따라서 $\text{depth}(M) = 0$이다. 특히 이 경우 보조정리가 성립한다.

\medskip\noindent
$\text{depth}(M) = 1$이면 첫 문단에 의해
$\mathfrak m \not \in \text{Ass}(M)$이다. 따라서 모든
$\mathfrak p \in \text{Ass}(M)$에 대해
$\dim(R/\mathfrak p) \geq 1$이고, 이 경우에도 보조정리가 성립한다.

\medskip\noindent
일반적으로 $\text{depth}(M)$에 대해 귀납하여 보조정리를 증명하겠다.
이를 $> 1$이라고 가정해도 되고 그렇게 하자.
$M$ 위의 비영인자 $x \in \mathfrak m$을 택하자.
$x \not \in \mathfrak p$임을 주목하자
(보조정리 \ref{algebra-lemma-ass-zero-divisors}).
보조정리 \ref{algebra-lemma-one-equation}에 의해
% P01-E0432
$\dim(R/\mathfrak p + (x)) = \dim(R/\mathfrak p) - 1$이다.
따라서 $\mathfrak p + (x)$ 위의 극소 소 아이디얼
$\mathfrak q$ 가운데
$\dim(R/\mathfrak q) = \dim(R/\mathfrak p) - 1$을 만족하는 것이
존재한다(짧은 논증은 생략한다. 힌트: 뇌터 국소환 $A$의 차원은
$A$의 극소 소 아이디얼 $\mathfrak r$들을 달리는
$A/\mathfrak r$의 차원들의 최댓값이다).
보조정리 \ref{algebra-lemma-inherit-minimal-primes}에서와 같은 $n$을 택하여
$\mathfrak q$가 $M/x^nM$의 연관 소 아이디얼이 되게 하자.
보조정리 \ref{algebra-lemma-depth-drops-by-one}에 의해
$\text{depth}(M/x^nM) = \text{depth}(M) - 1$이므로
$M/x^nM$과 $\mathfrak q$에 귀납 가정을 적용할 수 있다.
그 결과 $\dim(R/\mathfrak q) \geq \text{depth}(M/x^nM)$을 얻으며,
이것으로 증명이 끝난다.
\end{proof}

\begin{lemma}
\label{algebra-lemma-depth-localization}
$R$을 뇌터 국소환, $M$을 유한 $R$-가군이라 하자.
소 아이디얼 $\mathfrak p \subset R$에 대해
$\text{depth}(M_\mathfrak p) + \dim(R/\mathfrak p) \geq \text{depth}(M)$이다.
\end{lemma}

\begin{proof}
$M_\mathfrak p = 0$이면 $\text{depth}(M_\mathfrak p) = \infty$이므로
보조정리가 성립한다.
$\text{depth}(M) \leq \dim(R/\mathfrak p)$이면 보조정리가 성립한다.
$\text{depth}(M) > \dim(R/\mathfrak p)$이면 보조정리
\ref{algebra-lemma-depth-dim-associated-primes}에 의해 $\mathfrak p$는
$M$의 어느 연관 소 아이디얼 $\mathfrak q$에도 포함되지 않는다.
따라서 보조정리 \ref{algebra-lemma-silly}와 \ref{algebra-lemma-finite-ass}에 의해
$M$의 어느 연관 소 아이디얼에도 속하지 않는
$x \in \mathfrak p$를 찾을 수 있다. 그러면 보조정리
\ref{algebra-lemma-ass-zero-divisors}에 의해 $x$는 $M$ 위의 비영인자이다.
따라서 $\text{depth}(M/xM) = \text{depth}(M) - 1$이고,
$M_\mathfrak p$가 영이 아니라면
$\text{depth}(M_\mathfrak p / x M_\mathfrak p) =
\text{depth}(M_\mathfrak p) - 1$이다.
보조정리 \ref{algebra-lemma-depth-drops-by-one}를 보라.
그러므로 $\text{depth}(M)$에 대한 귀납법으로 결론을 얻는다.
\end{proof}

\begin{lemma}
\label{algebra-lemma-depth-goes-down-finite}
$(R, \mathfrak m)$을 뇌터 국소환이라 하자.
$R \to S$를 유한 환 사상이라 하자.
$\mathfrak m_1, \ldots, \mathfrak m_n$을 $S$의 극대 아이디얼들이라
하자. $N$을 유한 $S$-가군이라 하자. 그러면
$$
\min\nolimits_{i = 1, \ldots, n} \text{depth}(N_{\mathfrak m_i}) =
\text{depth}_\mathfrak m(N)
$$
\end{lemma}

\begin{proof}
보조정리 \ref{algebra-lemma-integral-no-inclusion},
\ref{algebra-lemma-integral-going-up} 및
\ref{algebra-lemma-finite-finite-fibres}에 의해 $S$의 극대 아이디얼들은
정확히 $\mathfrak m$ 위에 놓이는 $S$의 소 아이디얼들이며 그 수는
유한하다. 따라서 보조정리의 명제는 의미가 있다.
% P01-E0433
$k = \min\nolimits_{i = 1, \ldots, n} \text{depth}(N_{\mathfrak m_i})$에
대해 귀납하여 보조정리를 증명하겠다. $k = 0$이면 어떤 $i$에 대해
$\text{depth}(N_{\mathfrak m_i}) = 0$이다. 보조정리
\ref{algebra-lemma-depth-ext}에 의해 이는
$\mathfrak m_i S_{\mathfrak m_i}$가 $N_{\mathfrak m_i}$의 연관
소 아이디얼이라는 뜻이고, 따라서 $\mathfrak m_i$는 $N$의 연관
소 아이디얼이다(보조정리 \ref{algebra-lemma-localize-ass}).
보조정리 \ref{algebra-lemma-ass-functorial-Noetherian}에 의해
$\mathfrak m$은 $R$-가군으로 본 $N$의 연관 소 아이디얼이다.
그러므로 $\text{depth}_\mathfrak m(N) = 0$이다.
이로써 기저 단계가 증명되었다.
$k > 0$이면 $\mathfrak m_i \not \in \text{Ass}_S(N)$이다.
따라서 다시 보조정리 \ref{algebra-lemma-ass-functorial-Noetherian}에 의해
$\mathfrak m \not \in \text{Ass}_R(N)$이다.
그러므로 보조정리 \ref{algebra-lemma-ideal-nonzerodivisor}에 의해
$N$ 위의 영인자가 아닌 $f \in \mathfrak m$을 찾을 수 있다.
보조정리 \ref{algebra-lemma-depth-drops-by-one}에 의해 $N$에서 $N/fN$으로
넘어갈 때 모든 깊이가 정확히 $1$만큼 감소하며,
나머지는 귀납 가정에서 따른다.
\end{proof}

\section{Ext의 함자성}
\label{algebra-section-functoriality-ext}

\noindent
이 절에서는 환의 변경 등에 대한 $\Ext$의 함자성을 간단히 논한다.
확인해야 할 항목들의 목록은 다음과 같다.
\begin{enumerate}
\item $R \to R'$, $R$-가군 $M$ 및 $R'$-가군 $N'$가 주어지면,
$R$-가군 $\Ext^i_R(M, N')$에는 자연스러운 $R'$-가군 구조가 있다.
더욱이 표준적인 $R'$-선형 사상
$\Ext^i_{R'}(M \otimes_R R', N') \to
\Ext^i_R(M, N')$가 있다.
\item $R \to R'$와 $R$-가군들 $M$, $N$이 주어지면 자연스러운
$R$-가군 사상
% P01-E0434
$\Ext^i_R(M, N) \to \text{Ext}^i_R(M, N \otimes_R R')$가 있다.
\end{enumerate}

\begin{lemma}
\label{algebra-lemma-flat-base-change-ext}
평탄 환 사상 $R \to R'$, $R$-가군 $M$ 및 $R'$-가군 $N'$가
주어지면 자연스러운 사상
% P01-E0434
$$
\Ext^i_{R'}(M \otimes_R R', N') \to \text{Ext}^i_R(M, N')
$$
은 $i \geq 0$에 대해 동형이다.
\end{lemma}

\begin{proof}
$M$의 자유 분해 $F_\bullet$을 택하자.
$R \to R'$가 평탄하므로 $F_\bullet \otimes_R R'$는
$R'$ 위에서 $M \otimes_R R'$의 자유 분해이다.
명제는 사상
$$
\Hom_{R'}(F_\bullet \otimes_R R', N') \to
\Hom_R(F_\bullet, N')
$$
이 호몰로지군 위에 동형을 유도한다는 것이다.
% P01-E0435
이는 보조정리 \ref{algebra-lemma-adjoint-tensor-restrict}에 의해 이 사상이
복합체들의 동형이므로 참이다.
\end{proof}

\section{Ext 군의 한 응용}
\label{algebra-section-ext-application}

\noindent
그 응용은 다음과 같다.

\begin{lemma}
\label{algebra-lemma-split-injection-after-completion}
$R$을 뇌터 환이라 하자. $I \subset R$를 $R$의 야코브슨 근기에
포함되는 아이디얼이라 하자. $N \to M$을 유한 $R$-가군들의
준동형이라 하자. $N/I^nN \to M/I^nM$이 분할 단사인 임의로 큰
$n$들이 존재한다고 하자. 그러면 $N \to M$은 분할 단사이다.
\end{lemma}

\begin{proof}
$\varphi : N \to M$이 보조정리의 가정들을 만족한다고 하자.
이는 임의로 큰 $n$에 대해 $\Ker(\varphi) \subset I^nN$임을
함의한다. 따라서 보조정리
\ref{algebra-lemma-intersection-powers-ideal-module}에 의해 $\varphi$는
% P01-E1393
단사이다. $Q = M/N$라 두면 다음 짧은 완전열을 얻는다.
$$
0 \to N \to M \to Q \to 0.
$$
다음이 $Q$의 유한 자유 분해라고 하자.
$$
F_2 \xrightarrow{d_2} F_1 \xrightarrow{d_1} F_0 \to Q \to 0
$$
사상 $F_0 \to Q$를 올리는 사상 $\alpha : F_0 \to M$을 택할 수 있다.
이로부터 $\beta \circ d_2 = 0$인 사상
$\beta : F_1 \to N$이 유도된다. 위의 확장이 분할될 필요충분조건은
$\beta = \gamma \circ d_1$인 사상 $\gamma : F_0 \to N$이
존재하는 것이다. 달리 말하면 $\Ext^1_R(Q, N)$에서 $\beta$의 유는
위의 짧은 완전열이 분할되는 데 대한 장애물이다.

\medskip\noindent
$N/I^nN \to M/I^nM$이 분할 단사인 큰 정수 $n$을 택하자.
이는 다음 열이 여전히 짧은 완전열임을 함의한다.
$$
0 \to N/I^nN \to M/I^nM \to Q/I^nQ \to 0.
$$
또한 다음 열도 여전히 완전하다.
$$
F_1/I^nF_1 \xrightarrow{d_1} F_0/I^nF_0 \to Q/I^nQ \to 0
$$
위와 같이 논증하면 $\beta$가 유도하는 사상
$\overline{\beta} : F_1/I^nF_1 \to N/I^nN$은 어떤 사상
$\overline{\gamma_n} : F_0/I^nF_0 \to N/I^nN$에 대해
$\overline{\gamma_n} \circ d_1$과 같다.
$F_0$가 자유가군이므로 $\overline{\gamma_n}$을 사상
$\gamma_n : F_0 \to N$으로 올릴 수 있고, 그러면
$\beta - \gamma_n \circ d_1$은 $F_1$에서 $I^nN$으로 가는 사상이다.
달리 말하면 이 $n$에 대해 다음을 얻는다.
$$
\beta \in
\Im\Big(\Hom_R(F_0, N) \to \Hom_R(F_1, N)\Big) + I^n\Hom_R(F_1, N).
$$

\medskip\noindent
가정에 의해 이 성질이 임의로 큰 $n$에 대해 성립하므로,
보조정리 \ref{algebra-lemma-intersection-powers-ideal-module}에 의해
$\Hom_R(F_0, N) \to \Hom_R(F_1, N)$의 여핵에서 $\beta$의 상은
영이다. 따라서 원하는 대로 $\beta$는 사상
$\Hom_R(F_0, N) \to \Hom_R(F_1, N)$의 상에 속한다.
\end{proof}

\section{Tor 군과 평탄성}
\label{algebra-section-tor}

\noindent
이 절에서는 Tor 군이 무엇인지 설명하기 위해 앞 절에서 전개한
호몰로지 대수학의 일부를 사용한다. 즉, $R$이 환이고 $M$, $N$이
두 $R$-가군이라고 하자. $M$의 자유 $R$-가군들에 의한 분해
$F_\bullet$을 택한다. 보조정리
\ref{algebra-lemma-resolution-by-finite-free}를 보라.
호몰로지 복합체
$$
F_\bullet \otimes_R N
:
\ldots
\to F_2 \otimes_R N
\to F_1 \otimes_R N
\to F_0 \otimes_R N
$$
를 생각하자. $\text{Tor}^R_i(M, N)$을 이 복합체의 $i$번째
호몰로지군으로 정의한다. 다음 보조정리는 이것이 어떤 의미에서
잘 정의되는지 설명한다.

\begin{lemma}
\label{algebra-lemma-tor-welldefined}
$R$을 환이라 하자. $M_1, M_2, N$을 $R$-가군이라 하자.
$F_\bullet$이 가군 $M_1$의 자유 분해이고 $G_\bullet$이 가군
$M_2$의 자유 분해라고 하자. $\varphi : M_1 \to M_2$를
가군 사상이라 하자. $\alpha : F_\bullet \to G_\bullet$를
$M_1 = \Coker(d_{F, 1}) \to M_2 = \Coker(d_{G, 1})$ 위에서
$\varphi$를 유도하는 복합체 사상이라 하자.
보조정리 \ref{algebra-lemma-compare-resolutions}를 보라.
그러면 유도된 사상들
$$
H_i(\alpha) :
H_i(F_\bullet \otimes_R N)
\longrightarrow
H_i(G_\bullet \otimes_R N)
$$
은 $\alpha$의 선택과 무관하다. $\varphi$가 동형이면 모든 사상
$H_i(\alpha)$도 동형이다. $M_1 = M_2$, $F_\bullet = G_\bullet$이고
$\varphi$가 항등사상이면 모든 사상 $H_i(\alpha)$도 항등사상이다.
\end{lemma}

\begin{proof}
이 보조정리의 증명은 보조정리 \ref{algebra-lemma-ext-welldefined}의
증명과 동일하다.
\end{proof}

\noindent
이 보조정리는 Tor 가군들이 잘 정의된다는 것뿐 아니라 첫 번째
변수에 대한 구성 $(M, N) \mapsto \text{Tor}_i^R(M, N)$의 함자성도
제공한다. 물론 두 번째 변수에 대한 함자성은 명백하다.
각 $\text{Tor}_i^R$가 실제로 다음 함자임을 확인하는 일은 독자에게
맡긴다.
$$
\text{Mod}_R \times \text{Mod}_R \to \text{Mod}_R.
$$
여기서 $\text{Mod}_R$는 $R$-가군들의 범주를 나타내며,
곱범주의 정의는 Categories, 정의
\ref{categories-definition-product-category}를 보라.
즉, $R$-가군들의 사상 $M_1 \to M_2$와 $N_1 \to N_2$가 주어지면
다음 가환 그림을 얻는다.
$$
\xymatrix{
\text{Tor}_i^R(M_1, N_1) \ar[r] \ar[d] &
\text{Tor}_i^R(M_1, N_2) \ar[d] \\
\text{Tor}_i^R(M_2, N_1) \ar[r] &
\text{Tor}_i^R(M_2, N_2) \\
}
$$

\begin{lemma}
\label{algebra-lemma-long-exact-sequence-tor}
$R$을 환, $M$을 $R$-가군이라 하자.
$0 \to N' \to N \to N'' \to 0$이 $R$-가군들의 짧은 완전열이라고
하자. 다음 긴 완전열이 존재한다.
$$
\text{Tor}_1^R(M, N')
\to \text{Tor}_1^R(M, N)
\to \text{Tor}_1^R(M, N'')
\to
M \otimes_R N'
\to M \otimes_R N
\to M \otimes_R N''
\to 0
$$
\end{lemma}

\begin{proof}
이 증명은 보조정리 \ref{algebra-lemma-long-exact-seq-ext}의 증명과 같다.
\end{proof}

\noindent
다음 $R$-가군들의 호몰로지 이중복합체를 생각하자.
$$
\xymatrix{
\ldots \ar[r]^d &
A_{2, 0} \ar[r]^d &
A_{1, 0} \ar[r]^d &
A_{0, 0} \\
\ldots \ar[r]^d &
A_{2, 1} \ar[r]^d \ar[u]^\delta &
A_{1, 1} \ar[r]^d \ar[u]^\delta &
A_{0, 1} \ar[u]^\delta \\
\ldots \ar[r]^d &
A_{2, 2} \ar[r]^d \ar[u]^\delta &
A_{1, 2} \ar[r]^d \ar[u]^\delta &
A_{0, 2} \ar[u]^\delta \\
&
\ldots \ar[u]^\delta &
\ldots \ar[u]^\delta &
\ldots \ar[u]^\delta \\
}
$$
이는 $d_{i, j} : A_{i, j} \to A_{i-1, j}$와
$\delta_{i, j} : A_{i, j} \to A_{i, j-1}$가 다음 성질들을
만족한다는 뜻이다.
\begin{enumerate}
\item 임의의 두 $d_{i, j}$의 합성은 영이다.
달리 말하면 이중복합체의 행들은 복합체이다.
\item 임의의 두 $\delta_{i, j}$의 합성은 영이다.
달리 말하면 이중복합체의 열들은 복합체이다.
\item 임의의 쌍 $(i, j)$에 대해
$\delta_{i-1, j} \circ d_{i, j}
= d_{i, j-1} \circ \delta_{i, j}$이다.
달리 말하면 모든 사각형이 가환이다.
\end{enumerate}
올바른 방법은 그러한 모든 이중복합체에 스펙트럼 열을 대응시키는
것이다. 그러나 지금은 약간 더 쉬운 일을 해도 충분하다.

\medskip\noindent
즉, 이 절에서만 이중복합체
$(A_{\bullet, \bullet}, d, \delta)$가 주어졌을 때
$R(A)_j = \Coker(A_{1, j} \to A_{0, j})$ 및
$U(A)_i = \Coker(A_{i, 1} \to A_{i, 0})$로 둔다.
(문자 $R$과 $U$는 각각 Right와 Up을 암시한다.)
사상들 $\delta$를 사용하여 $R(A)_\bullet$에 복합체 구조를 준다.
마찬가지로 사상들 $d$를 사용하여 $U(A)_\bullet$에 복합체 구조를
준다. 달리 말하면 다음의 거대한 가환 그림을 얻는다.
$$
\xymatrix{
\ldots \ar[r]^d &
U(A)_2 \ar[r]^d &
U(A)_1 \ar[r]^d &
U(A)_0 &
\\
\ldots \ar[r]^d &
A_{2, 0} \ar[r]^d \ar[u] &
A_{1, 0} \ar[r]^d \ar[u] &
A_{0, 0} \ar[r] \ar[u] &
R(A)_0 \\
\ldots \ar[r]^d &
A_{2, 1} \ar[r]^d \ar[u]^\delta &
A_{1, 1} \ar[r]^d \ar[u]^\delta &
A_{0, 1} \ar[r] \ar[u]^\delta &
R(A)_1 \ar[u]^\delta \\
\ldots \ar[r]^d &
A_{2, 2} \ar[r]^d \ar[u]^\delta &
A_{1, 2} \ar[r]^d \ar[u]^\delta &
A_{0, 2} \ar[r] \ar[u]^\delta &
R(A)_2 \ar[u]^\delta \\
&
\ldots \ar[u]^\delta &
\ldots \ar[u]^\delta &
\ldots \ar[u]^\delta &
\ldots \ar[u]^\delta \\
}
$$
(물론 이것은 더 이상 이중복합체가 아니다.)
이중복합체의 사상
$\Phi : (A_{\bullet, \bullet}, d, \delta)
\to (B_{\bullet, \bullet}, d, \delta)$가 무엇인지는 명백하며,
이것이 복합체의 사상
$R(\Phi) : R(A)_\bullet \to R(B)_\bullet$ 및
$U(\Phi) : U(A)_\bullet \to U(B)_\bullet$을 유도함도 명백하다.

\begin{lemma}
\label{algebra-lemma-no-spectral-sequence}
$(A_{\bullet, \bullet}, d, \delta)$가 다음을 만족하는
이중복합체라고 하자.
\begin{enumerate}
\item 각 행 $A_{\bullet, j}$는 $R(A)_j$의 분해이다.
\item 각 열 $A_{i, \bullet}$은 $U(A)_i$의 분해이다.
\end{enumerate}
그러면 표준적인 동형
$$
H_i(R(A)_\bullet)
\cong
H_i(U(A)_\bullet).
$$
들이 있다. 이 동형들은 위의 성질들을 만족하는 이중복합체의
사상에 대해 함자적이다.
\end{lemma}

\begin{proof}
% P01-E0436
$H_i(R(A)_\bullet))$와 $H_i(U(A)_\bullet)$이 제3의 군과 표준적으로
동형임을 보이겠다. 즉,
$$
\mathbf{H}_i(A) :=
\frac{
\{
(a_{i, 0}, a_{i-1, 1}, \ldots, a_{0, i})
\mid
d(a_{i, 0}) = \delta(a_{i-1, 1}), \ldots,
d(a_{1, i-1}) = \delta(a_{0, i})
\}}
{
\{
d(a_{i + 1, 0}) + \delta(a_{i, 1}),
d(a_{i, 1}) + \delta(a_{i-1, 2}),
\ldots,
d(a_{1, i}) + \delta(a_{0, i + 1})
\}
}
$$
이다. 여기서는 $a_{i, j}$가 $A_{i, j}$의 원소를 나타낸다는
표기 관례를 사용한다. 달리 말하면 $\mathbf{H}_i$의 원소는
지그재그로 표현된다. $i = 2$인 경우 이는 다음과 같다.
$$
\xymatrix{
a_{2, 0} \ar@{|->}[r] & d(a_{2, 0}) = \delta(a_{1, 1}) & \\
& a_{1, 1} \ar@{|->}[u] \ar@{|->}[r] &
d(a_{1, 1}) = \delta(a_{0, 2}) \\
& & a_{0, 2} \ar@{|->}[u] \\
}
$$
물론 ‘자명한’ 지그재그들, 즉 다음 꼴의 원소들로 생성되는
부분가군으로 나눈다.
% P01-E0437
$$
(0, \ldots, 0, -\delta(a_{t + 1, t-i}),
d(a_{t + 1, t-i}), 0, \ldots, 0)
$$
이들이다.
다음 표준적인 준동형들이 있음을 주목하자.
$$
\mathbf{H}_i(A) \to H_i(R(A)_\bullet), \quad
(a_{i, 0}, a_{i-1, 1}, \ldots, a_{0, i}) \mapsto
\text{class of image of }a_{0, i}
$$
및
$$
\mathbf{H}_i(A) \to H_i(U(A)_\bullet), \quad
(a_{i, 0}, a_{i-1, 1}, \ldots, a_{0, i}) \mapsto
\text{class of image of }a_{i, 0}
$$

\medskip\noindent
먼저 이 사상들이 전사임을 보이자.
$\overline{r} \in H_i(R(A)_\bullet)$이라고 하자.
% P01-E0438
$r \in R(A)_i$를 $\overline{r}$의 유를 나타내는 코사이클이라 하자.
$r$로 가는 원소 $a_{0, i} \in A_{0, i}$를 택하자.
$\delta(r) = 0$이므로 $\delta(a_{0, i})$는 $d$의 상에 속한다.
따라서 $d(a_{1, i-1}) = \delta(a_{0, i})$인 원소
$a_{1, i-1} \in A_{1, i-1}$가 존재한다. 이는 다시
$\delta(a_{1, i-1})$가 $d$의 핵에 속함을 함의한다
% P01-E1394
($d(\delta(a_{1, i-1})) = \delta(d(a_{1, i-1}))
= \delta(\delta(a_{0, i})) = 0$이기 때문이다.
행들의 완전성에 의해
$d(a_{2, i-2}) = \delta(a_{1, i-1})$인 원소
$a_{2, i-2}$를 찾을 수 있다. 완전한 지그재그를 찾을 때까지
이 과정을 계속한다. 물론 $\mathbf{H}_i \to H_i(U(A))$의 전사성도
같은 방법으로 보인다.

\medskip\noindent
단사성을 증명할 때도 정확히 같은 방식으로 논증한다.
$H_i(R(A)_\bullet)$에서 영으로 가는 지그재그
$(a_{i, 0}, a_{i-1, 1}, \ldots, a_{0, i})$가 주어졌다고 하자.
이는 $a_{0, i}$가
% P01-E0439
$\Coker(A_{i, 1} \to A_{i, 0})$에서 다음 사상의 상에 속하는
원소로 간다는 뜻이다.
$\delta : \Coker(A_{i + 1, 1} \to A_{i + 1, 0}) \to
\Coker(A_{i, 1} \to A_{i, 0})$.
달리 말하면 $a_{0, i}$는 사상
$\delta \oplus d : A_{0, i + 1} \oplus A_{1, i} \to A_{0, i}$의
상에 속한다. 자명한 지그재그의 정의에 의해 우리 지그재그를
자명한 것으로 수정하여 $a_{0, i} = 0$이라고 가정할 수 있다.
이는 곧바로 $d(a_{1, i-1}) = 0$을 함의한다. 행들이 완전하므로
$a_{1, i-1}$은 $d : A_{2, i-1} \to A_{1, i-1}$의 상에 속한다.
따라서 우리 지그재그를 다시 자명한 지그재그로 수정하여
$(a_{i, 0}, a_{i-1, 1}, \ldots, a_{2, i-2}, 0, 0)$ 꼴이라고
가정할 수 있다. 이와 같이 계속하면 원하는 단사성을 얻는다.

\medskip\noindent
$\Phi : (A_{\bullet, \bullet}, d, \delta)
\to (B_{\bullet, \bullet}, d, \delta)$가 보조정리의 조건들을
모두 만족하는 이중복합체들의 사상이면, 명백히 다음 가환 그림을
얻는다.
$$
\xymatrix{
H_i(U(A)_\bullet) \ar[d] &
\mathbf{H}_i(A) \ar[r] \ar[l] \ar[d] &
H_i(R(A)_\bullet) \ar[d] \\
H_i(U(B)_\bullet) &
\mathbf{H}_i(B) \ar[r] \ar[l] &
H_i(R(B)_\bullet) \\
}
$$
이로써 함자성이 증명된다.
\end{proof}

\begin{remark}
\label{algebra-remark-signs-double-complex}
위에서 구성한 동형은 부호를 무시하면 ‘올바른’ 동형이다.
호몰로지 대수학의 상당 부분은 여러 사상의 부호를 선택하고 적절한
부호를 개입시켜 그림의 가환성을 보이는 일과 관련된다. 지금은 단지
위 증명에서 주어진 동형을 사용하고 부호는 나중에 다루겠다.
\end{remark}

\begin{lemma}
\label{algebra-lemma-tor-left-right}
$R$을 환이라 하자. 임의의 $i \geq 0$에 대해 함자들
$\text{Mod}_R \times \text{Mod}_R \to \text{Mod}_R$,
$(M, N) \mapsto \text{Tor}_i^R(M, N)$ 및
$(M, N) \mapsto \text{Tor}_i^R(N, M)$은 표준적으로 동형이다.
\end{lemma}

\begin{proof}
$F_\bullet$을 가군 $M$의 자유 분해라 하고
$G_\bullet$을 가군 $N$의 자유 분해라 하자.
이중복합체 $(A_{i, j}, d, \delta)$를 다음과 같이 정의한다.
\begin{enumerate}
\item $A_{i, j} = F_i \otimes_R G_j$로 두고,
% P01-E0440
\item $d_{i, j} : F_i \otimes_R G_j \to F_{i-1} \otimes G_j$를
$d_{F, i} \otimes \text{id}$와 같게 두며,
\item $\delta_{i, j} : F_i \otimes_R G_j \to F_i \otimes G_{j-1}$를
$\text{id} \otimes d_{G, j}$와 같게 둔다.
\end{enumerate}
이 이중복합체를 보통 간단히 $F_\bullet \otimes_R G_\bullet$로
나타낸다.

\medskip\noindent
각 $G_j$가 자유가군이고 따라서 평탄하므로, 이중복합체의 각 행은
호몰로지 차수 $0$을 제외하고 완전하다. 각 $F_i$가 자유가군이고
따라서 평탄하므로, 이중복합체의 각 열은 호몰로지 차수 $0$을
제외하고 완전하다. 따라서 이중복합체는 보조정리
\ref{algebra-lemma-no-spectral-sequence}의 조건들을 만족한다.

\medskip\noindent
보조정리가 무엇을 말하는지 알아보기 위해
$R(A)_\bullet$과 $U(A)_\bullet$을 계산하자. 즉,
\begin{eqnarray*}
R(A)_i & = & \Coker(A_{1, i} \to A_{0, i}) \\
& = & \Coker(F_1 \otimes_R G_i \to F_0 \otimes_R G_i) \\
& = & \Coker(F_1 \to F_0) \otimes_R G_i \\
& = & M \otimes_R G_i
\end{eqnarray*}
이다. 실제로 이 동형들은 미분 $\delta$와 양립하며, 호몰로지
복합체로서 $R(A)_\bullet = M \otimes_R G_\bullet$임을 알 수 있다.
정확히 같은 방법으로 $U(A)_\bullet = F_\bullet \otimes_R N$임을
알 수 있다. 따라서
\begin{eqnarray*}
\text{Tor}_i^R(M, N)
& = & H_i(F_\bullet \otimes_R N) \\
& = & H_i(U(A)_\bullet) \\
& = & H_i(R(A)_\bullet) \\
& = & H_i(M \otimes_R G_\bullet) \\
& = & H_i(G_\bullet \otimes_R M) \\
& = & \text{Tor}_i^R(N, M)
\end{eqnarray*}
을 얻는다. 여기서 세 번째 등식은 보조정리
\ref{algebra-lemma-no-spectral-sequence}이고, 다섯 번째 등식은
텐서곱의 동형 $V \otimes W = W \otimes V$를 사용한다.

\medskip\noindent
함자성. $R$-가군들 $M_\nu$, $N_\nu$, $\nu = 1, 2$가 있다고 하자.
$\varphi : M_1 \to M_2$와 $\psi : N_1 \to N_2$를
$R$-가군들의 사상이라 하자. $M_\nu$의 자유 분해
$F_{\nu, \bullet}$과 $N_\nu$의 자유 분해
$G_{\nu, \bullet}$이 있다고 하자.
보조정리 \ref{algebra-lemma-compare-resolutions}에 의해 $\varphi$ 및
$\psi$와 양립하는 복합체의 사상
$\alpha : F_{1, \bullet} \to F_{2, \bullet}$ 및
$\beta : G_{1, \bullet} \to G_{2, \bullet}$를 택할 수 있다.
쌍 $(\alpha, \beta)$가 이중복합체의 사상
$$
\alpha \otimes \beta :
F_{1, \bullet} \otimes_R G_{1, \bullet}
\longrightarrow
F_{2, \bullet} \otimes_R G_{2, \bullet}
$$
을 유도한다고 주장한다. 이는
$F_{1, i} \otimes_R G_{1, j} \to F_{2, i} \otimes_R G_{2, j}$가
$\alpha_i \otimes \beta_j$로 주어진다는 규칙을 사용하여 바로
확인된다. 여기서 $\alpha_i$와 $\beta_j$는 각각 $\alpha$의
차수 $i$ 성분과 $\beta$의 차수 $j$ 성분이다.
유도된 사상
$R(F_{1, \bullet} \otimes_R G_{1, \bullet})_\bullet
\to R(F_{2, \bullet} \otimes_R G_{2, \bullet})_\bullet$이
$\varphi \otimes \beta$에 의해 유도되는 사상
$M_1 \otimes_R G_{1, \bullet}
\to M_2 \otimes_R G_{2, \bullet}$과 일치함도 독자는 바로 확인할
수 있다. $U(-)_\bullet$ 복합체들 위에 유도되는 사상에 대해서도
마찬가지이다. 따라서 함자성에 관한 명제는 보조정리
\ref{algebra-lemma-no-spectral-sequence}의 함자성 명제에서 따른다.
\end{proof}

\begin{remark}
\label{algebra-remark-curiosity-signs-swap}
위에서 $M = N$인 경우에는 흥미로운 일이 생긴다.
이 경우 표준적인 사상
$\text{Tor}_i^R(M, M) \to \text{Tor}_i^R(M, M)$을 얻는다.
이 사상은 항등사상이 아님을 주목하자. $i = 0$일 때조차
항등사상이 아니다! 예를 들어 $V$가 어떤 체 위의 차원 $n$인
벡터공간이면 교환 사상 $V \otimes_k V \to V \otimes_k V$는
고윳값 $+1$을 $(n^2 + n)/2$개, 고윳값 $-1$을
% P01-E1395
$(n^2-n)/2$개 가진다. 표수가 $2$이면 대각화 가능하지조차 않다.
사상의 부호를 바꾸어도 이 현상은 없어지지 않는다.
\end{remark}

\begin{lemma}
\label{algebra-lemma-tor-noetherian}
$R$을 뇌터 환이라 하자. $M$, $N$을 유한 $R$-가군이라 하자.
그러면 모든 $p$에 대해 $\text{Tor}_p^R(M, N)$은 유한
$R$-가군이다.
\end{lemma}

\begin{proof}
% P01-E0441
$\text{Tor}_p^R(M, N)$은 각 $F_n$이 유한 자유 $R$-가군인
복합체 $F_\bullet \otimes_R N$의 코호몰로지군으로 계산되므로
성립한다. 보조정리 \ref{algebra-lemma-resolution-by-finite-free}를 보라.
\end{proof}

\begin{lemma}
\label{algebra-lemma-characterize-flat}
$R$을 환이라 하고 $M$을 $R$-가군이라 하자.
다음은 서로 동치이다.
\begin{enumerate}
\item 가군 $M$은 $R$ 위에서 평탄하다.
\item 모든 $i > 0$에 대해 함자 $\text{Tor}_i^R(M, -)$는 영이다.
\item 함자 $\text{Tor}_1^R(M, -)$는 영이다.
\item 모든 아이디얼 $I \subset R$에 대해
$\text{Tor}_1^R(M, R/I) = 0$이다.
\item 모든 유한 생성 아이디얼 $I \subset R$에 대해
$\text{Tor}_1^R(M, R/I) = 0$이다.
\end{enumerate}
\end{lemma}

\begin{proof}
$M$이 평탄하다고 하자. $N$을 $R$-가군이라 하자.
$F_\bullet$을 $N$의 자유 분해라 하자.
$M$의 평탄성에 의해 $F_\bullet \otimes_R M$은
$N \otimes_R M$의 분해이다. 따라서 모든 고차 Tor 군이 사라진다.

\medskip\noindent
이제 마지막 조건이 $M$의 평탄성을 함의함을 보이면 충분하다.
$I \subset R$를 아이디얼이라 하자. 짧은 완전열
$0 \to I \to R \to R/I \to 0$을 생각하자.
보조정리 \ref{algebra-lemma-long-exact-sequence-tor}를 적용하면 다음
완전열을 얻는다.
$$
\text{Tor}_1^R(M, R/I) \to
M \otimes_R I \to
M \otimes_R R \to
M \otimes_R R/I \to
0
$$
명백히 $M \otimes_R R = M$이므로 마지막 가정은 모든 유한 생성
아이디얼 $I$에 대해 $M \otimes_R I \to M$이 단사임을 함의한다.
따라서 보조정리 \ref{algebra-lemma-flat}에 의해 $M$은 평탄하다.
\end{proof}

\begin{remark}
\label{algebra-remark-Tor-ring-mod-ideal}
보조정리 \ref{algebra-lemma-characterize-flat}의 증명은 실제로
$$
\text{Tor}_1^R(M, R/I)
=
\Ker(I \otimes_R M \to M).
$$
임을 보인다.
\end{remark}

\section{Tor의 함자성}
\label{algebra-section-functoriality-tor}

\noindent
이 절에서는 환의 변경 등에 관한 $\text{Tor}$의 함자성을 간단히
논한다. 다음은 구체적으로 확인해야 할 사항들의 목록이다.
\begin{enumerate}
\item 환 사상 $R \to R'$, $R$-가군 $M$, 그리고 $R'$-가군 $N'$이
주어지면 $R$-가군 $\text{Tor}_i^R(M, N')$은 자연스러운
$R'$-가군 구조를 갖는다.
\item 환 사상 $R \to R'$과 $R$-가군 $M$, $N$이 주어지면 자연스러운
$R$-가군 사상
$\text{Tor}_i^R(M, N) \to \text{Tor}_i^{R'}(M \otimes_R R', N \otimes_R R')$이
있다.
\item 환 사상 $R \to R'$, $R$-가군 $M$, 그리고 $R'$-가군 $N'$이
주어지면 자연스러운 $R'$-가군 사상
$\text{Tor}_i^R(M, N') \to \text{Tor}_i^{R'}(M \otimes_R R', N')$이
존재한다.
\end{enumerate}

\begin{lemma}
\label{algebra-lemma-flat-base-change-tor}
평탄 환 사상 $R \to R'$과 $R$-가군 $M$, $N$이 주어졌다고 하자.
그러면 자연스러운 $R$-가군 사상
$\text{Tor}_i^R(M, N)\otimes_R R'
\to \text{Tor}_i^{R'}(M \otimes_R R', N \otimes_R R')$은
모든 $i$에 대해 동형사상이다.
\end{lemma}

\begin{proof}
생략한다. 실제로 $R$ 위에서 $M$의 자유 분해 $F_\bullet$은
$R$ 위에서 $R'$과 텐서곱을 취해도 완전한 채로 남으므로,
$(F_\bullet \otimes_R N)\otimes_R R'$은 $R'$ 위의 Tor 군을 계산한다.
\end{proof}

\noindent
다음 보조정리는 달리 들어맞는 곳이 없는 듯하다.

\begin{lemma}
\label{algebra-lemma-tor-commutes-filtered-colimits}
$R$을 환이라 하자. $M = \colim M_i$를 $R$-가군들의 여과 여극한이라
하고, $N$을 $R$-가군이라 하자. 그러면 모든 $n$에 대해
$\text{Tor}_n^R(M, N) = \colim \text{Tor}_n^R(M_i, N)$이다.
\end{lemma}

\begin{proof}
$N$의 자유 분해 $F_\bullet$을 택하자. 보조정리
\ref{algebra-lemma-tensor-products-commute-with-limits}에 의해 복합체로서
$F_\bullet \otimes_R M = \colim F_\bullet \otimes_R M_i$이다.
% P01-E1396
따라서 보조정리 \ref{algebra-lemma-directed-colimit-exact}에 의해 결론이 따른다.
\end{proof}








\section{사영 가군}
\label{algebra-section-projective}

\noindent
사영 가군에 관한 몇 가지 보조정리를 제시한다.

\begin{definition}
\label{algebra-definition-projective}
$R$을 환이라 하자. $R$-가군 $P$가 {\it 사영적}이라는 것은 함자
$\Hom_R(P, -) : \text{Mod}_R \to \text{Mod}_R$가 완전 함자라는 것과
필요충분이다.
\end{definition}

\noindent
임의의 $R$-가군 $M$에 대해 함자 $\Hom_R(M, - )$는 왼쪽 완전이다.
보조정리 \ref{algebra-lemma-hom-exact}를 보라. 따라서 $P$가 사영적이라는
조건이 실제로 뜻하는 바는, $R$-가군들의 전사 $N \to N'$이 주어질 때
사상 $\Hom_R(P, N) \to \Hom_R(P, N')$이 전사라는 것이다.

\begin{lemma}
\label{algebra-lemma-characterize-projective}
$R$을 환이라 하고 $P$를 $R$-가군이라 하자. 다음은 서로 동치이다.
\begin{enumerate}
\item $P$는 사영적이다.
\item $P$는 자유 $R$-가군의 직합인자이다.
\item 모든 $R$-가군 $M$에 대해 $\Ext^1_R(P, M) = 0$이다.
\end{enumerate}
\end{lemma}

\begin{proof}
$P$가 사영적이라고 하자. $F$가 자유 $R$-가군인 전사
$\pi : F \to P$를 택하자. $P$가 사영적이므로
$\pi \circ i = \text{id}_P$를 만족하는 $i \in \Hom_R(P, F)$가
존재한다. 다시 말해 $F \cong \Ker(\pi) \oplus i(P)$이므로,
$P$가 $F$의 직합인자임을 알 수 있다.

\medskip\noindent
역으로 $P \oplus Q = F$가 자유 $R$-가군이라고 하자. 자유 가군
$F = \bigoplus_{i \in I} R$는 사영적임을 주목하자. 실제로
$\Hom_R(F, M) = \prod_{i \in I} M$이고 함자
$M \mapsto \prod_{i \in I} M$은 완전하다. 함자로서
$\Hom_R(F, -) = \Hom_R(P, -) \times \Hom_R(Q, -)$이므로
$P$와 $Q$는 모두 사영적이다.

\medskip\noindent
$P \oplus Q = F$가 자유 $R$-가군이라고 하자. 그러면 다음 꼴의
자유 분해 $F_\bullet$이 있다.
% P01-E0442
$$
\ldots F \xrightarrow{a} F \xrightarrow{b} F \to P \to 0
$$
% P01-E1397
여기서 사상 $a, b$는 번갈아 나타나며 $P$와 $Q$로의 사영자이다.
따라서 복합체 $\Hom_R(F_\bullet, M)$은 차수 $\geq 1$에서
분할 완전이고, 이로부터 (3)의 소멸을 얻는다.

\medskip\noindent
모든 $R$-가군 $M$에 대해 $\Ext^1_R(P, M) = 0$이라고 하자.
자유 분해 $F_\bullet \to P$를 택하고
$M = \Im(F_1 \to F_0) = \Ker(F_0 \to P)$라 두자.
몫사상 $\pi : F_1 \to M$의 류가 주는 원소
$\xi \in \Ext^1_R(P, M)$를 생각하자. $\xi$가 영이므로
$\pi = s \circ (F_1 \to F_0)$를 만족하는 사상 $s : F_0 \to M$이
존재한다. 이는 명백히
$$
F_0 = \Ker(s) \oplus \Ker(F_0 \to P) =
P \oplus \Ker(F_0 \to P)
$$
를 뜻하며, 결론을 얻는다.
\end{proof}

\begin{lemma}
\label{algebra-lemma-characterize-finite-projective-noetherian}
$R$을 뇌터 환이라 하고 $P$를 유한 $R$-가군이라 하자.
모든 유한 $R$-가군 $M$에 대해 $\Ext^1_R(P, M) = 0$이면
$P$는 사영적이다.
\end{lemma}

\noindent
이 보조정리는 더 강화할 수 있다. $R$이 뇌터라는 가정이 없을 때에는
유한 표시 $R$-가군에 대한 판본이 있고, 뇌터인 경우에는 $M$이 모든
유한 길이 가군을 달리도록 하는 판본이 있다.

\begin{proof}
핵이 $M$인 전사 $R^{\oplus n} \to P$를 택하자.
$\Ext^1_R(P, M) = 0$이므로 이 전사는 분할되며, 보조정리
\ref{algebra-lemma-characterize-projective}에 의해 결론이 따른다.
\end{proof}

\begin{lemma}
\label{algebra-lemma-direct-sum-projective}
사영 가군들의 직합은 사영적이다.
\end{lemma}

\begin{proof}
보조정리 \ref{algebra-lemma-characterize-projective}에서 사영 가군을 자유
가군의 직합인자로 특징지은 것으로부터 따른다.
\end{proof}

\begin{lemma}
\label{algebra-lemma-lift-projective-module}
$R$을 환이라 하고 $I \subset R$를 멱영 아이디얼이라 하자.
$\overline{P}$를 사영 $R/I$-가군이라 하자. 그러면
$P/IP \cong \overline{P}$를 만족하는 사영 $R$-가군 $P$가 존재한다.
\end{lemma}

\begin{proof}
보조정리 \ref{algebra-lemma-characterize-projective}에 의해 어떤 $R/I$-가군
$\overline{K}$, 어떤 집합 $A$에 대하여 직합분해
$\bigoplus_{\alpha \in A} R/I = \overline{P} \oplus \overline{K}$를
택할 수 있다. $A$ 위의 자유 $R$-가군을
$F = \bigoplus_{\alpha \in A} R$라 쓰자.
$\bigoplus_{\alpha \in A} R/I$의 직합인자 $\overline{P}$에 대응하는
사영자 $\overline{p}$의 올림 $p : F \to F$를 택하자.
$p^2 - p \in \text{End}_R(F)$는 $F$의 멱영 자기준동형임을 주목하자.
실제로 $I$가 멱영이고 $p^2 - p$의 행렬 성분들이 $I$에 들어간다.
더 정확히 말해 $I^n = 0$이면 $(p^2 - p)^n = 0$이다.
따라서 보조정리 \ref{algebra-lemma-lift-idempotents-noncommutative}에 의해
$p$의 선택을 수정하여 $p$가 사영자라고 가정할 수 있다.
$P = \Im(p)$라 두자.
\end{proof}

\begin{lemma}
\label{algebra-lemma-lift-finite-projective-module}
$R$을 환이라 하고 $I \subset R$를 국소 멱영 아이디얼이라 하자.
$\overline{P}$를 유한 사영 $R/I$-가군이라 하자. 그러면
$P/IP \cong \overline{P}$를 만족하는 유한 사영 $R$-가군 $P$가
존재한다.
\end{lemma}

\begin{proof}
보조정리 \ref{algebra-lemma-characterize-projective}에 의해 $\overline{P}$는
자유 $R/I$-가군 $\bigoplus_{\alpha \in A} R/I$의 직합인자임을
상기하자. $\overline{P}$가 유한이므로, 어떤 유한 부분집합
$A' \subset A$에 대해 $\overline{P}$는
$\bigoplus_{\alpha \in A'} R/I$에 포함된다. 따라서 어떤 $n$과
어떤 $R/I$-가군 $\overline{K}$에 대해 직합분해
$(R/I)^{\oplus n} = \overline{P} \oplus \overline{K}$가 있다고
가정해도 된다. $(R/I)^{\oplus n}$의 직합인자 $\overline{P}$에
대응하는 사영자 $\overline{p}$의 올림
$p \in \text{Mat}(n \times n, R)$를 택하자.
$p^2 - p \in \text{Mat}(n \times n, R)$가 멱영임을 주목하자.
$I$가 국소 멱영이고 $p^2 - p$의 행렬 성분 $c_{ij}$가 $I$에 속하므로
어떤 $t > 0$에 대해 $c_{ij}^t = 0$이고, 따라서 행렬 성분을 살펴보면
$(p^2 - p)^{tn^2} = 0$이다. 그러므로 보조정리
\ref{algebra-lemma-lift-idempotents-noncommutative}에 의해 $p$의 선택을
수정하여 $p$가 사영자라고 가정할 수 있다. $P = \Im(p)$라 두자.
\end{proof}

\begin{lemma}
\label{algebra-lemma-lift-projective}
$R$을 환이라 하자. $I \subset R$를 아이디얼이라 하고
$M$을 $R$-가군이라 하자. 다음을 가정하자.
\begin{enumerate}
\item $I$는 멱영이다.
\item $M/IM$은 사영 $R/I$-가군이다.
\item $M$은 평탄 $R$-가군이다.
\end{enumerate}
그러면 $M$은 사영 $R$-가군이다.
\end{lemma}

\begin{proof}
보조정리 \ref{algebra-lemma-lift-projective-module}에 의해 사영 $R$-가군
$P$와 동형사상 $P/IP \to M/IM$을 찾을 수 있다. 증명을 끝내기 위해
$M$이 $P$와 동형임을 보이겠다. $P$가 사영적이므로 사상
$P \to P/IP \to M/IM$을 $R$-가군 사상 $P \to M$으로 올릴 수 있고,
이 사상은 $I$로 나눈 뒤 동형사상이다. 어떤 $n$에 대해 $I^n = 0$이므로
여과
\begin{align*}
0 = I^nM \subset I^{n - 1}M \subset \ldots \subset IM \subset M \\
0 = I^nP \subset I^{n - 1}P \subset \ldots \subset IP \subset P
\end{align*}
를 사용하면, 유도된 사상
$I^aP/I^{a + 1}P \to I^aM/I^{a + 1}M$이 전단사임을 보이는 것으로
충분함을 알 수 있다. $P$와 $M$은 모두 평탄 $R$-가군이므로 이 사상을
다음 사상과 동일시할 수 있다.
$$
I^a/I^{a + 1} \otimes_{R/I} P/IP
\longrightarrow
I^a/I^{a + 1} \otimes_{R/I} M/IM
$$
이는 $P \to M$에 의해 유도된다. 유도된 사상 $P/IP \to M/IM$이
동형사상이 되도록 $P \to M$을 택했으므로 결론을 얻는다.
\end{proof}

\begin{lemma}
\label{algebra-lemma-projective-modulo-two-ideals}
$R$을 환이라 하자. $I \cap J = 0$을 만족하는 아이디얼
$I, J \subset R$를 택하자. $P/IP$가 사영 $R/I$-가군이고
$P/JP$가 사영 $R/J$-가군인 $R$-가군 $P$를 택하자.
그러면 $P$는 사영 $R$-가군이다.
\end{lemma}

\begin{proof}
$F$가 자유 $R$-가군인 전사 $p : F \to P$를 택하자.
$P/IP$가 사영 $R/I$-가군이므로 $p \bmod I$의 오른쪽 역인 사상
$f : P/IP \to F/IF$를 택할 수 있다. 몫사상
$F/JF \to F/(I + J)F$와 $p$에 의해 유도되는 사상
$q : F/JF \to F/(I + J)F \times_{P/(I + J)P} P/JP$를 생각하자.
$q$는 $R/J$-가군의 전사 사상임을 주목하자. 작은 세부사항은 생략한다.
$f$와 항등사상에 의해 유도되는 사상
$f' : P/JP \to F/(I + J)F \times_{P/(I + J)P} P/JP$를 생각하자.
% P01-E1398
$P/JP$가 사영 $R/J$-가군이고 $q$가 전사이므로,
$q \circ g = f'$를 만족하는 사상 $g : P/JP \to F/JF$를 택할 수 있다.
그러면 $f \bmod I + J = g \bmod I + J$이다.
$I \cap J = 0$에 의해
$F = F/JF \times_{F/(I + J)F} F/IF$이므로, 잘 정의된 $R$-가군
준동형 $h = (f, g) : P \to F$를 얻는다. 사상
$a = p \circ h : P \to P$는 $I$로 나눈 뒤에도 $J$로 나눈 뒤에도
항등사상으로 내려간다. 그러면 $a$는 부분가군 $IP$ 위에서도
항등사상을 유도한다. 실제로 $x = \sum i_\alpha x_\alpha$이고
$i_\alpha \in I$, $x_\alpha \in P$이면
$a(x) = \sum i_\alpha a(x_\alpha) = \sum i_\alpha x_\alpha = x$이다.
이는 $a(x_\alpha) = x_\alpha + j_\alpha$이고
% P01-E1399
$j_\alpha \in J$이며 $i_\alpha j_\alpha = 0$이기 때문이다.
따라서 $a : P \to P$는 $IP$ 위에서도 $P/IP$ 위에서도 항등으로
작용하므로, $a$가 $P$의 자기동형사상임을 얻는다. 작은 세부사항은
생략한다. 따라서 $P$는 $F$의 직합인자이고, 그러므로 사영적이다.
\end{proof}















\section{유한 사영 가군}
\label{algebra-section-finite-projective-modules}

\begin{definition}
\label{algebra-definition-locally-free}
$R$을 환이라 하고 $M$을 $R$-가군이라 하자.
\begin{enumerate}
\item $\Spec(R)$을 표준 열린집합 $D(f_i)$, $i \in I$로 덮을 수 있고
모든 $i \in I$에 대해 $M_{f_i}$가 자유 $R_{f_i}$-가군이면,
$M$을 {\it 국소 자유}라고 한다.
\item 각 $M_{f_i}$가 유한 자유가 되도록 덮개를 택할 수 있으면
$M$을 {\it 유한 국소 자유}라고 한다.
\item 각 $M_{f_i}$가 $R_{f_i}^{\oplus r}$와 동형이 되도록 덮개를
택할 수 있으면 $M$을 {\it 랭크 $r$인 유한 국소 자유}라고 한다.
\end{enumerate}
\end{definition}

\noindent
유한 국소 자유 $R$-가군은 보조정리 \ref{algebra-lemma-cover}에 의해 자동으로
유한 표시임을 주목하자. 또한 $M$이 환 $R$ 위의 랭크 $r$인 유한
국소 자유 가군이고 $R$이 영환이 아니면, 보조정리
\ref{algebra-lemma-rank}에 의해 $r$은 유일하게 정해진다. 실제로 국소화들 중
% P01-E1311
적어도 하나인 $R_{f_i}$는 영환이 아니다.

\begin{lemma}
\label{algebra-lemma-finite-projective}
$R$을 환이라 하고 $M$을 $R$-가군이라 하자. 다음은 서로 동치이다.
\begin{enumerate}
\item $M$은 유한 표시이고 $R$-평탄하다.
\item $M$은 유한 사영 가군이다.
\item $M$은 유한 자유 $R$-가군의 직합인자이다.
\item $M$은 유한 표시이고, 모든 $\mathfrak p \in \Spec(R)$에 대해
국소화 $M_{\mathfrak p}$가 자유이다.
\item $M$은 유한 표시이고, 모든 극대 아이디얼 $\mathfrak m \subset R$에
대해 국소화 $M_{\mathfrak m}$가 자유이다.
\item $M$은 유한이고 국소 자유이다.
\item $M$은 유한 국소 자유이다.
\item $M$은 유한이고, 모든 소 아이디얼 $\mathfrak p$에 대해 가군
$M_{\mathfrak p}$가 자유이며, 함수
$$
\rho_M : \Spec(R) \to \mathbf{Z}, \quad
\mathfrak p
\longmapsto
\dim_{\kappa(\mathfrak p)} M \otimes_R \kappa(\mathfrak p)
$$
는 자리스키 위상에서 국소 상수이다.
\end{enumerate}
\end{lemma}

\begin{proof}
먼저 $M$이 유한 사영 가군, 즉 (2)가 성립한다고 하자.
전사 $R^n \to M$을 택하고 그 핵을 $K$라 하자.
$M$이 사영적이므로 $0 \to K \to R^n \to M \to 0$은 분할된다.
따라서 (2) $\Rightarrow$ (3)이다.
(3) $\Rightarrow$ (2)는 사영 가군의 직합인자가 사영적이라는
사실에서 따른다. 보조정리 \ref{algebra-lemma-characterize-projective}를 보라.

\medskip\noindent
(3)을 가정하자. 그러면 $K \oplus M \cong R^{\oplus n}$으로 쓸 수 있다.
따라서 $K$는 $R^n$의 직합인자이고, 그러므로 유한 생성이다.
이는 $M = R^{\oplus n}/K$가 유한 표시임을 보인다.
다시 말해 (3) $\Rightarrow$ (1)이다.

\medskip\noindent
$M$이 유한 표시이고 평탄하다고, 즉 (1)이 성립한다고 하자.
(7)이 성립함을 증명하겠다. 임의의 소 아이디얼 $\mathfrak p$와
$M \otimes_R \kappa(\mathfrak p)$의 기저로 사상되는
$x_1, \ldots, x_r \in M$을 택하자.
나카야마 보조정리(보조정리 \ref{algebra-lemma-NAK-localization}의 형태)에
의해, 어떤 $g \in R$, $g \not \in \mathfrak p$에 대해 이 원소들은
$M_g$를 생성한다. 대응하는 전사
$\varphi : R_g^{\oplus r} \to M_g$는 다음 두 성질을 갖는다.
(a) $\Ker(\varphi)$는 유한 $R_g$-가군이다(보조정리
\ref{algebra-lemma-extension} 참조).
(b) $M_g$가 $R_g$ 위에서 평탄하므로
$\Ker(\varphi) \otimes \kappa(\mathfrak p) = 0$이다
(보조정리 \ref{algebra-lemma-flat-tor-zero} 참조).
따라서 나카야마 보조정리를 다시 적용하면
$\Ker(\varphi)_{g'} = 0$인 $g' \in R_g$가 존재한다.
% P01-E0443
다시 말해 $M_{gg'}$는 자유이다.

\medskip\noindent
유한 국소 자유 가군은 유한 가군이다. 보조정리 \ref{algebra-lemma-cover}를
보라. 따라서 (7) $\Rightarrow$ (6)이다.
(6) $\Rightarrow$ (7)과 (7) $\Rightarrow$ (8)은 명백하다.

\medskip\noindent
유한 국소 자유 가군은 유한 표시 가군이다. 보조정리
\ref{algebra-lemma-cover}를 보라. 따라서 (7) $\Rightarrow$ (4)이다.
물론 (4)는 (5)를 함의한다.
평탄성은 국소적으로 확인할 수 있으므로(보조정리
\ref{algebra-lemma-flat-localization} 참조), (5)는 (1)을 함의한다.
이 시점에서 다음을 얻었다.
$$
\xymatrix{
(2) \ar@{<=>}[r] & (3) \ar@{=>}[r] & (1) \ar@{=>}[r] &
(7)  \ar@{<=>}[r] \ar@{=>}[rd] \ar@{=>}[d] & (6) \\
& & (5) \ar@{=>}[u] & (4) \ar@{=>}[l] & (8)
}
$$

\medskip\noindent
$M$이 (1), (4), (5), (6), (7)을 만족한다고 하자. (3)이 성립함을
증명하겠다. $M$이 사영적임을 보이면 충분하다. 즉
$\Hom_R(M, -)$가 완전함을 보여야 한다.
% P01-E1400
$R$-가군의 짧은 완전열 $0 \to N'' \to N \to N'\to 0$이
주어졌다고 하자. 다음 열이 완전함을 보여야 한다.
$0 \to \Hom_R(M, N'') \to \Hom_R(M, N) \to
\Hom_R(M, N') \to 0$.
$M$이 유한 국소 자유이므로 $M_{f_i}$가 유한 자유가 되도록 하는
덮개 $\Spec(R) = \bigcup D(f_i)$가 존재한다. 보조정리
\ref{algebra-lemma-hom-from-finitely-presented}에 의해
$$
0 \to \Hom_R(M, N'')_{f_i} \to \Hom_R(M, N)_{f_i} \to
\Hom_R(M, N')_{f_i} \to 0
$$
은
$0 \to \Hom_{R_{f_i}}(M_{f_i}, N''_{f_i}) \to
\Hom_{R_{f_i}}(M_{f_i}, N_{f_i}) \to
\Hom_{R_{f_i}}(M_{f_i}, N'_{f_i}) \to 0$과 같다.
$M_{f_i}$가 자유이고 국소화
$0 \to N''_{f_i} \to N_{f_i} \to N'_{f_i} \to 0$가
완전하므로(국소화가 완전하기 때문이다), 이 열은 완전하다.
그러므로 보조정리 \ref{algebra-lemma-cover}에 의해
$0 \to \Hom_R(M, N'') \to \Hom_R(M, N) \to
\Hom_R(M, N') \to 0$은 완전하다.

\medskip\noindent
마지막으로 (8)이 성립한다고 하자. 극대 아이디얼
$\mathfrak m \subset R$를 택하자.
$M \otimes_R \kappa(\mathfrak m) = M/\mathfrak mM$의
$\kappa(\mathfrak m)$-기저로 사상되는
$x_1, \ldots, x_r \in M$을 택하자. 특히
$\rho_M(\mathfrak m) = r$이다. 나카야마 보조정리
\ref{algebra-lemma-NAK}에 의해, 어떤 $f \in R$, $f \not \in \mathfrak m$에
대해 $x_1, \ldots, x_r$은 $R_f$ 위에서 $M_f$를 생성한다.
$\rho_M$이 국소 상수라는 가정에 의해, 어떤
$g \in R$, $g \not \in \mathfrak m$가 존재하여 $\rho_M$은
$D(g)$ 위에서 상수 $r$이다. 다음 사상이 동형사상이라고 주장한다.
$$
\Psi : R_{fg}^{\oplus r} \longrightarrow M_{fg}, \quad
(a_1, \ldots, a_r) \longmapsto \sum a_i x_i
$$
이 주장은 $M$이 유한 국소 자유, 즉 (7)이 성립함을 보인다.
주장을 보이기 위해서는 모든 $\mathfrak p \in D(fg)$에 대해 국소화에
유도된 사상
$\Psi_{\mathfrak p} : R_{\mathfrak p}^{\oplus r} \to M_{\mathfrak p}$이
동형사상임을 보이면 충분하다. 보조정리
\ref{algebra-lemma-characterize-zero-local}을 보라.
$f$의 선택에 의해 사상 $\Psi_{\mathfrak p}$은 전사이다.
가정 (8)에 의해
$M_{\mathfrak p} \cong R_{\mathfrak p}^{\oplus \rho_M(\mathfrak p)}$이고,
$g$의 선택에 의해 $\rho_M(\mathfrak p) = r$이다.
따라서 $\Psi_{\mathfrak p}$은 전사
$R_{\mathfrak p}^{\oplus r} \to
M_{\mathfrak p} \cong R_{\mathfrak p}^{\oplus r}$를 정한다.
% P01-E1401
그러므로 보조정리 \ref{algebra-lemma-fun}에 의해 동형사상이다.
(물론 이 마지막 사실은 간단한 행렬 논증으로도 따른다.)
\end{proof}

\begin{lemma}
\label{algebra-lemma-finite-projective-reduced}
$R$을 축소환이라 하고 $M$을 $R$-가군이라 하자. 그러면 보조정리
\ref{algebra-lemma-finite-projective}의 동치조건들은 다음 조건과도 동치이다.
\begin{enumerate}
\item[(9)] $M$은 유한이고, 함수
$\rho_M : \Spec(R) \to \mathbf{Z}$,
$\mathfrak p \mapsto
\dim_{\kappa(\mathfrak p)} M \otimes_R \kappa(\mathfrak p)$는
자리스키 위상에서 국소 상수이다.
\end{enumerate}
\end{lemma}

\begin{proof}
극대 아이디얼 $\mathfrak m \subset R$를 택하자.
$M \otimes_R \kappa(\mathfrak m) = M/\mathfrak mM$의
$\kappa(\mathfrak m)$-기저로 사상되는
$x_1, \ldots, x_r \in M$을 택하자. 특히
$\rho_M(\mathfrak m) = r$이다. 나카야마 보조정리
\ref{algebra-lemma-NAK}에 의해, 어떤 $f \in R$, $f \not \in \mathfrak m$에
대해 $x_1, \ldots, x_r$은 $R_f$ 위에서 $M_f$를 생성한다.
$\rho_M$이 국소 상수라는 가정에 의해, 어떤
$g \in R$, $g \not \in \mathfrak m$가 존재하여 $\rho_M$은
$D(g)$ 위에서 상수 $r$이다. 다음 사상이 동형사상이라고 주장한다.
$$
\Psi : R_{fg}^{\oplus r} \longrightarrow M_{fg}, \quad
(a_1, \ldots, a_r) \longmapsto \sum a_i x_i
$$
이 주장은 $M$이 유한 국소 자유, 즉 (7)이 성립함을 보인다.
$\Psi$가 전사이므로 $\Psi$가 단사임을 보이면 충분하다.
$R_{fg}$가 축소환이므로, $R_{fg}$의 모든 극소 소 아이디얼
$\mathfrak p$에서 국소화한 뒤 $\Psi$가 단사임을 보이면 충분하다.
보조정리 \ref{algebra-lemma-reduced-ring-sub-product-fields}를 보라.
그런데 보조정리 \ref{algebra-lemma-minimal-prime-reduced-ring}에 의해
$R_\mathfrak p = \kappa(\mathfrak p)$이고,
$\rho_M(\mathfrak p) = r$이다. 따라서
$\Psi_\mathfrak p :
R_\mathfrak p^{\oplus r} \to M \otimes_R \kappa(\mathfrak p)$은
차원이 같은 유한 차원 벡터공간 사이의 전사이므로 동형사상이다.
\end{proof}

\begin{remark}
\label{algebra-remark-warning}
유한 $R$-가군이 $R$-평탄이면 자동으로 사영적이라는 주장은 거짓이다.
% P01-E1402
반례는 다음과 같다. $R = \mathcal{C}^\infty(\mathbf{R})$를
$\mathbf{R}$ 위의 무한 번 미분 가능 함수들의 환이라 하고,
$M = R_{\mathfrak m} = R/I$라 하자. 여기서
$\mathfrak m = \{f \in R \mid f(0) = 0\}$이고
$I = \{f \in R \mid \exists \epsilon, \epsilon > 0 :
f(x) = 0\ \forall x, |x| < \epsilon\}$이다.
\end{remark}

\begin{lemma}
\label{algebra-lemma-finite-flat-local}
(주의: 주석 \ref{algebra-remark-warning}를 보라.)
$R$이 국소환이고 $M$이 유한 평탄 $R$-가군이라고 하자.
그러면 $M$은 유한 자유이다.
\end{lemma}

\begin{proof}
평탄성의 방정식 판정법에서 따른다. 보조정리
\ref{algebra-lemma-flat-eq}를 보라. 구체적으로
$x_1, \ldots, x_r \in M$이 $M/\mathfrak mM$의 기저로 사상된다고 하자.
나카야마 보조정리 \ref{algebra-lemma-NAK}에 의해 이 원소들은 $M$을 생성한다.
$x_i$들 사이에 관계식이 없음을 보이고자 한다. 대신 $n$에 대한
귀납법으로 다음을
보이겠다. $x_1, \ldots, x_n \in M$이 벡터공간
$M/\mathfrak mM$에서 선형독립이면 이들은 $R$ 위에서도 독립이다.

\medskip\noindent
귀납법의 첫 경우에는 $x \in M$, $x \not\in \mathfrak mM$이고
관계식 $fx = 0$이 주어져 있다. 방정식 판정법에 의해 모든 $j$에 대해
$x = \sum a_j y_j$와 $fa_j = 0$을 만족하는
$y_j \in M$ 및 $a_j \in R$가 존재한다.
$x \not\in \mathfrak mM$이므로 적어도 하나의 $a_j$는 단원이고,
따라서 $f = 0 $이다.

\medskip\noindent
$\sum f_i x_i$가 $x_1, \ldots, x_n$ 사이의 관계식이라고 하자.
$x_i$의 선택에 의해 $f_i \in \mathfrak m$이다.
평탄성의 방정식 판정법에 따라
$x_i = \sum a_{ij} y_j$와 $\sum f_i a_{ij} = 0$을 만족하는
$a_{ij} \in R$ 및 $y_j \in M$이 존재한다.
$x_n \not \in \mathfrak mM$이므로 적어도 하나의 $j$에 대해
$a_{nj}\not\in \mathfrak m$이다.
$\sum f_i a_{ij} = 0$이므로
$f_n = \sum_{i = 1}^{n-1} (-a_{ij}/a_{nj}) f_i$를 얻는다.
관계식 $\sum f_i x_i = 0$은 이제
$\sum_{i = 1}^{n-1} f_i( x_i + (-a_{ij}/a_{nj}) x_n) = 0$으로
다시 쓸 수 있다. 원소
$x_i + (-a_{ij}/a_{nj}) x_n$은 $M/\mathfrak mM$에서
$n-1$개의 선형독립 원소로 사상됨을 주목하자.
귀납가정에 의해 모든 $f_i$, $i \leq n-1$는 영이어야 하고,
또한 $f_n = \sum_{i = 1}^{n-1} (-a_{ij}/a_{nj}) f_i$이다.
이로써 귀납 단계를 증명했다.
\end{proof}

\begin{lemma}
\label{algebra-lemma-finite-projective-descends}
$R \to S$를 국소환들의 평탄 국소 준동형이라 하자.
$M$을 유한 $R$-가군이라 하자. 그러면 $M$이 $R$ 위의 유한 사영
가군인 것과 $M \otimes_R S$가 $S$ 위의 유한 사영 가군인 것은
필요충분이다.
\end{lemma}

\begin{proof}
보조정리 \ref{algebra-lemma-finite-projective}에 의해 국소환 위에서 유한
사영이라는 것은 유한 자유라는 것과 같다.
$M \otimes_R S$가 유한 자유 $S$-가군이라고 하자.
$M/\mathfrak m_RM$에서의 상들이
% P01-E0444
$\kappa(\mathfrak m)$ 위의 기저를 이루는
$x_1, \ldots, x_r \in M$을 택하자. 그러면
$x_1 \otimes 1, \ldots, x_r \otimes 1$은 $M \otimes_R S$의
기저임을 알 수 있다. 이는 사상
$R^{\oplus r} \to M, (a_i) \mapsto \sum a_i x_i$가
$S$와 텐서곱을 취한 뒤 동형사상이 됨을 함의한다.
$R \to S$의 충실 평탄성에 의해 이 사상은 동형사상이다.
보조정리 \ref{algebra-lemma-local-flat-ff}를 보라.
\end{proof}

\begin{lemma}
\label{algebra-lemma-locally-free-semi-local-free}
$R$을 반국소환이라 하고 $M$을 유한 국소 자유 가군이라 하자.
$M$의 랭크가 상수이면 $M$은 자유이다. 특히 $R$의 스펙트럼이
연결이면 $M$은 자유이다.
\end{lemma}

\begin{proof}
생략한다. 힌트는 다음과 같다. 먼저 랭크가 상수임을 이용하여
$R$의 모든 극대 아이디얼 $\mathfrak m_1, \ldots, \mathfrak m_n$에 대해
% P01-E1403
$M/\mathfrak m_iM$의 차원이 같은 $d$임을 보인다.
다음으로 중국인의 나머지 정리를 사용하여, 각
$M/\mathfrak m_iM$에서 기저를 이루는 원소
$x_1, \ldots, x_d \in M$이 존재함을 보인다.
마지막으로 $x_1, \ldots, x_d$가 $M$의 기저임을 보인다.
\end{proof}

\noindent
다음은 군대상에 관한 장에서 사용하는 기술적 보조정리이다.

\begin{lemma}
\label{algebra-lemma-semi-local-module-basis-in-submodule}
$R$을 극대 아이디얼 $\mathfrak m$과 무한 잉여체를 갖는 국소환이라
하자. $R \to S$를 환 사상이라 하자. $M$을 $S$-가군이라 하고
$N \subset M$을 $R$-부분가군이라 하자. 다음을 가정하자.
\begin{enumerate}
\item $S$는 반국소이고 $\mathfrak mS$는 $S$의 야코브슨 근기에
포함된다.
\item $M$은 유한 자유 $S$-가군이다.
\item $N$은 $S$-가군으로서 $M$을 생성한다.
\end{enumerate}
그러면 $N$은 $M$의 $S$-기저를 포함한다.
\end{lemma}

\begin{proof}
$M$이 랭크 $n$인 자유가군이라고 하자. $I \subset S$를 야코브슨
근기라 하자. 나카야마 보조정리 \ref{algebra-lemma-NAK}에 의해 원소들의 열
$m_1, \ldots, m_n$이 $M$의 기저일 필요충분조건은
$\overline{m}_i \in M/IM$이 $M/IM$을 생성하는 것이다.
따라서 $M$을 $M/IM$으로, $N$을 $N/(N \cap IM)$으로,
$R$을 $R/\mathfrak m$으로, $S$를 $S/IS$로 바꾸어도 된다.
이 경우 $S$는 체들의 유한곱
$S = k_1 \times \ldots \times k_r$이고
$M = k_1^{\oplus n} \times \ldots \times k_r^{\oplus n}$임을 알 수 있다.
$N \subset M$이 $S$-가군으로서 $M$을 생성한다는 것은
$a_j \in S$를 계수로 하는 선형결합 $\sum a_j x_j$가 각 인자
$k_i^{\oplus n}$에서 영이 아닌 성분을 갖도록 하는
$x_j \in N$이 존재한다는 뜻이다.
$R = k$가 무한체이므로, 어떤 $c_j \in k$에 대해 선형결합
$y = \sum c_j x_j$ 역시 각 인자에서 영이 아닌 성분을 갖는다.
따라서 $y \in N$은 자유 직합인자 $Sy \subset M$을 생성한다.
$M/Sy$와 부분가군 $\overline{N} = N/(N \cap Sy)$에 대해
$n$에 관한 귀납법을 적용하면 결론이 따른다.
다시 말해 $M/Sy$를 자유롭게 생성하는
$\overline{y}_2, \ldots, \overline{y}_n$이 $\overline{N}$ 안에 존재한다.
그러면 $y, y_2, \ldots, y_n$은 $M$을 자유롭게 생성하며, 결론을 얻는다.
\end{proof}

\begin{lemma}
\label{algebra-lemma-evaluation-map-iso-finite-projective}
% P01-E1404
$R$을 환이라 하고 $L$, $M$, $N$을 $R$-가군이라 하자.
표준 사상
$$
\Hom_R(M, N) \otimes_R L \to \Hom_R(M, N \otimes_R L)
$$
은 $M$이 유한 사영 가군이면 동형사상이다.
\end{lemma}

\begin{proof}
보조정리 \ref{algebra-lemma-finite-projective}에 의해 $M$은 유한 표시이면서
유한 국소 자유이다. 보조정리
\ref{algebra-lemma-hom-from-finitely-presented}와
\ref{algebra-lemma-tensor-product-localization}에 의해 화살표의 왼쪽과
오른쪽을 만드는 것은 국소화와 가환한다. 이 사상이 동형인지 여부는
국소화한 뒤 확인할 수 있다. 보조정리 \ref{algebra-lemma-cover}를 보라.
따라서 $M$이 유한 자유라고 가정해도 된다. 이 경우 보조정리는
즉시 성립한다.
\end{proof}


 
\section{가군 사상이 정의하는 열린 자취}
\label{algebra-section-loci-maps}

\noindent
주어진 가군 사상이 전사이거나 동형사상이 되는 소 아이디얼들의 집합은
때로 열린집합이다. 유한 사영 가군의 경우에는 사상의 랭크를 살펴볼
수 있다.

\begin{lemma}
\label{algebra-lemma-map-between-finite}
$R$을 환이라 하자. $\varphi : M \to N$을 $R$-가군들의 사상이라 하고
$N$을 유한 $R$-가군이라 하자. 그러면 다음 등식이 성립한다.
\begin{align*}
U & = \{\mathfrak p \subset R \mid
\varphi_{\mathfrak p} : M_{\mathfrak p} \to N_{\mathfrak p}
\text{ is surjective}\} \\
& = \{\mathfrak p \subset R \mid
\varphi \otimes \kappa(\mathfrak p) :
M \otimes \kappa(\mathfrak p) \to N \otimes \kappa(\mathfrak p)
\text{ is surjective}\}
\end{align*}
또한 $U$는 $\Spec(R)$의 열린 부분집합이다. 더욱이
$D(f) \subset U$를 만족하는 임의의 $f \in R$에 대해 사상
$M_f \to N_f$는 전사이다.
\end{lemma}

\begin{proof}
표시된 공식의 등식은 나카야마 보조정리에서 따른다.
나카야마 보조정리는 $U$가 열려 있음도 함의한다.
보조정리 \ref{algebra-lemma-NAK}, 특히 (3)을 보라.
$D(f) \subset U$이면 $M_f \to N_f$는 $R_f$의 모든 소
아이디얼에서 국소화한 뒤 전사이므로, 보조정리
\ref{algebra-lemma-characterize-zero-local}에 의해 전사이다.
\end{proof}

\begin{lemma}
\label{algebra-lemma-map-between-finitely-presented}
$R$을 환이라 하자. $\varphi : M \to N$을 $R$-가군들의 사상이라
하고, $M$은 유한이며 $N$은 유한 표시라고 하자. 그러면
$$
U = \{\mathfrak p \subset R \mid
\varphi_{\mathfrak p} : M_{\mathfrak p} \to N_{\mathfrak p}
\text{ is an isomorphism}\}
$$
는 $\Spec(R)$의 열린 부분집합이다.
\end{lemma}

\begin{proof}
$\mathfrak p \in U$라 하자. 표시
$N = R^{\oplus n}/\sum_{j = 1, \ldots, m} R k_j$를 택하자.
$R^{\oplus n}$의 $i$번째 기저벡터가 $N$에서 갖는 상을 $e_i$라
쓰자. 각 $i \in \{1, \ldots, n\}$에 대해, 어떤
$f_i \in R$, $f_i \not \in \mathfrak p$에 대하여
$\varphi(m_i) = f_i e_i$를 만족하는 원소
$m_i \in M_{\mathfrak p}$를 택하자.
$\varphi_{\mathfrak p}$가 동형사상이므로 이는 가능하다.
% P01-E0445
$f = f_1 \ldots f_n$이라 두고, $i$번째 기저벡터를 $m_i/f_i$로
보내는 사상을 $\psi : R_f^{\oplus n} \to M_f$라 하자.
$\varphi_f \circ \psi$는 주어진 사상
$R^{\oplus n} \to N$을 $f$에서 국소화한 것임을 주목하자.
$\varphi_{\mathfrak p}$가 동형사상이므로
% P01-E0446
$\psi(k_j)$는 $M_{\mathfrak p}$에서 영으로 사상되는 $M$의 원소임을
알 수 있다. 따라서 $g_j \psi(k_j) = 0$을 만족하는
$g_j \in R$, $g_j \not \in \mathfrak p$가 존재한다.
$g = g_1 \ldots g_m$이라 두면, $\psi_g$는 $N_{fg}$를 통해
인수분해되어 사상 $\chi : N_{fg} \to M_{fg}$를 준다.
구성에 의해 $\chi$는 $\varphi_{fg}$의 오른쪽 역이다.
따라서 $\chi_\mathfrak p$는 동형사상이다. 보조정리
\ref{algebra-lemma-map-between-finite}에 의해, 어떤
$h \in R$, $h \not \in \mathfrak p$에 대해
$\chi_h : N_{fgh} \to M_{fgh}$는 전사이다.
그러므로 $\varphi_{fgh}$와 $\chi_h$는 서로 역인 사상이고,
원하는 대로 $D(fgh) \subset U$이다.
\end{proof}

\begin{lemma}
\label{algebra-lemma-finitely-presented-localization-free}
$R$을 환이라 하고 $\mathfrak p \subset R$를 소 아이디얼이라 하자.
$M$을 유한 표시 $R$-가군이라 하자. $M_\mathfrak p$가 자유이면,
어떤 $f \in R$, $f \not \in \mathfrak p$가 존재하여
$M_f$는 자유 $R_f$-가군이다.
\end{lemma}

\begin{proof}
기저 $x_1, \ldots, x_n \in M_\mathfrak p$를 택하자.
각 $x_i$가 어떤 $y_i \in M_f$의 상이 되도록 하는
$f \in R$, $f \not \in \mathfrak p$를 택할 수 있다.
$m \gg 0$에 대해 $y_i$를 $f^m y_i$로 바꾸면
$y_i \in M$이라고 가정할 수 있다. 실제로 이는
$x_1, \ldots, x_n$을 $f^mx_1, \ldots, f^mx_n$으로 바꾸는데,
$f$는 $R_\mathfrak p$에서 단원으로 사상되므로 이들도 여전히
기저이다. 따라서 $\mathfrak p$에서의 국소화가 동형사상인
$R$-가군 준동형
$\varphi = (y_1, \ldots, y_n) : R^{\oplus n} \to M$을 얻는다.
보조정리 \ref{algebra-lemma-map-between-finitely-presented}에 의해,
$f \not \in \mathfrak q$인 모든 소 아이디얼
$\mathfrak q \subset R$에 대해 $\varphi_\mathfrak q$가
동형사상이 되도록 하는 어떤 $f \in R$, $f \not \in \mathfrak p$를
찾을 수 있다. 그러면 보조정리 \ref{algebra-lemma-characterize-zero-local}에
의해 $\varphi_f$는 동형사상이고, 증명이 끝난다.
\end{proof}

\begin{lemma}
\label{algebra-lemma-cokernel-flat}
$R$을 환이라 하자. $\varphi : P_1 \to P_2$를 유한 사영 가군들
사이의 사상이라 하자. 그러면 다음이 성립한다.
\begin{enumerate}
\item $\varphi \otimes \kappa(\mathfrak p)$가 단사인 소 아이디얼
$\mathfrak p \in \Spec(R)$들의 집합 $U$는 열려 있고,
$D(f) \subset U$를 만족하는 임의의 $f\in R$에 대해
\begin{enumerate}
\item $P_{1, f} \to P_{2, f}$는 단사이고,
\item 가군 $\Coker(\varphi)_f$는 $R_f$ 위에서 유한 사영이다.
\end{enumerate}
\item $\varphi \otimes \kappa(\mathfrak p)$가 전사인 소 아이디얼
$\mathfrak p \in \Spec(R)$들의 집합 $W$는 열려 있고,
$D(f) \subset W$를 만족하는 임의의 $f\in R$에 대해
\begin{enumerate}
\item $P_{1, f} \to P_{2, f}$는 전사이고,
\item 가군 $\Ker(\varphi)_f$는 $R_f$ 위에서 유한 사영이다.
\end{enumerate}
\item $\varphi \otimes \kappa(\mathfrak p)$가 동형사상인 소 아이디얼
$\mathfrak p \in \Spec(R)$들의 집합 $V$는 열려 있고,
$D(f) \subset V$를 만족하는 임의의 $f\in R$에 대해 사상
$\varphi : P_{1, f} \to P_{2, f}$는 $R_f$ 위 가군들의 동형사상이다.
\end{enumerate}
\end{lemma}

\begin{proof}
$U$가 열린집합임을 증명하기 위해 $\Spec(R)$ 위에서 국소적으로
논해도 된다. 따라서 $R$을 적당한 국소화로 바꾸고
$P_1 = R^{n_1}$ 및 $P_2 = R^{n_2}$라고 가정할 수 있다.
보조정리 \ref{algebra-lemma-finite-projective}를 보라.
이 경우 $\varphi \otimes \kappa(\mathfrak p)$가 단사일
필요충분조건은 $n_1 \leq n_2$이고 $\varphi$의 행렬의 어떤
$n_1 \times n_1$ 소행렬식 $f$가 $\kappa(\mathfrak p)$에서
가역인 것이다. 따라서 $D(f) \subset U$이다.
이 논증은 $\mathfrak p \in U$에 대해
$P_{1, \mathfrak p} \to P_{2, \mathfrak p}$가 단사임도 보인다.

\medskip\noindent
이제 $D(f) \subset U$라고 하자. 앞 문단의 주의와 보조정리
\ref{algebra-lemma-characterize-zero-local}에 의해
$P_{1, f} \to P_{2, f}$는 단사, 즉 (1)(a)가 성립한다.
보조정리 \ref{algebra-lemma-finite-projective}에 의해 (1)(b)를
증명하려면 $\Coker(\varphi)$가 $D(f)$ 위에서 국소적으로 유한
사영임을 보이면 충분하다. 따라서 위에서 본 것처럼
$P_1 = R^{n_1}$ 및 $P_2 = R^{n_2}$이고 $\varphi$의 행렬의 어떤
소행렬식이 $R$에서 가역이라고 가정할 수 있다.
그 소행렬식이 $R^{n_2}$의 처음 $n_1$개 기저벡터에 대응하면,
마지막 $n_2 - n_1$개 기저벡터를 사용하여 사상
$R^{n_2 - n_1} \to R^{n_2} \to \Coker(\varphi)$을 얻는데,
이는 쉽게 동형사상임을 알 수 있다.

\medskip\noindent
$W$의 열린성과 $D(f) \subset W$일 때의 (2)(a)는 보조정리
\ref{algebra-lemma-map-between-finite}에서 따른다. $P_{2, f}$는 $R_f$ 위에서
사영적이므로 $\varphi_f : P_{1, f} \to P_{2, f}$는 단면을 가지며,
% P01-E0447
따라서 $\Ker(\varphi)_f$는 $P_{2, f}$의 직합인자이다.
그러므로 $\Ker(\varphi)_f$는 유한 사영이다. 따라서 (2)(b)도
성립한다.

\medskip\noindent
$V = U \cap W$가 열려 있음은 명백하고, (3)의 나머지 명제는
(1)(a)와 (2)(a)에서 따른다.
\end{proof}





\section{가군의 사영성에 대한 충실 평탄 하강}
\label{algebra-section-ffdescent-projectivity-introduction}

\medskip\noindent
다음 몇 절에서는 Raynaud와 Gruson \cite{GruRay}을 따라,
충실 평탄 환 사상을 따라 가군의 사영성이 하강함을 증명한다.
증명의 발상은 Kaplansky식 d\'evissage \cite{Kaplansky}를 사용하여
가산 생성 가군의 경우로 환원하는 것이다. 가군 $M$에 잘 거동하는
여과가 주어지면, d\'evissage를 통해 $M$을 여과 부분가군들의 연속
몫들의 직합으로 나타낼 수 있다(절 \ref{algebra-section-transfinite-devissage}를
보라). 이 기법을 사용하여 사영 가군은 가산 생성 가군들의 직합임을
증명한다(정리 \ref{algebra-theorem-projective-direct-sum}). 가산 생성
가군에 대한 사영성의 하강을 증명하기 위해 가군에 “Mittag-Leffler”
조건을 도입하고, 가산 생성 가군이 사영일 필요충분조건은 평탄하고
Mittag-Leffler인 것임을 증명한 뒤(정리
\ref{algebra-theorem-projectivity-characterization}), Mittag-Leffler
가군이라는 성질이 하강함을 보인다(보조정리
\ref{algebra-lemma-ffdescent-ML}). 마지막으로, 충실 평탄 환 사상에 의한
기저변환이 사영인 임의의 가군 $M$이 주어지면, 연속 몫들이 가산
생성 사영 가군인 부분가군들로 $M$을 여과한다. 그러고 나서
d\'evissage에 의해 $M$이 사영 가군들의 직합이고, 따라서 그 자체로
사영임을 결론짓는다(정리 \ref{algebra-theorem-ffdescent-projectivity}).

\medskip\noindent
\cite{GruRay}에 제시된 사영성의 충실 평탄 하강 증명에는 오류가
있음을 지적한다. 그곳에서는 더 일반적인 종류의 환 사상을 따른
사영성의 하강으로부터 충실 평탄 환 사상을 따른 사영성의 하강을
도출한다(\cite[Example 3.1.4(1) of Part II]{GruRay}).
그러나 이 더 일반적인 종류의 사상을 따른 하강의 증명은 옳지 않다.
\cite{G}에서 Gruson은 무엇이 잘못되었는지 설명하지만, 여기서
관심 있는 경우에 대한 수정은 제시하지 않는다. 사영성의 충실 평탄
하강 증명에서 이 빈틈을 메우는 일은 Mittag-Leffler 가군이라는
성질이 충실 평탄 환 사상을 따라 하강함을 증명하는 것으로 귀착된다.
이를 보조정리 \ref{algebra-lemma-ffdescent-ML}에서 수행한다.





\section{평탄성의 특징짓기}
\label{algebra-section-characterize-flatness}

\noindent
이 절에서는 평탄성의 판정법을 논한다. 이 절의 주된 결과는
Lazard 정리(아래 정리 \ref{algebra-theorem-lazard})로, 평탄 가군은
유한 자유 가군들의 유향계의 여극한임을 말한다.
평탄성의 “등식 판정법”은 보조정리 \ref{algebra-lemma-flat-eq}를 보라.
이를 가공하면 겉보기에는 훨씬 강한 다음 성질을 얻을 수 있다.

\begin{lemma}
\label{algebra-lemma-flat-factors-free}
$M$을 $R$-가군이라 하자. 다음 조건들은 동치이다.
\begin{enumerate}
\item $M$은 평탄하다.
\item $f: R^n \to M$이 가군 사상이고 $x \in \Ker(f)$이면,
$f = g \circ h$ 및 $x \in \Ker(h)$를 만족하는 가군 사상
$h: R^n \to R^m$과 $g: R^m \to M$이 존재한다.
\item $f: R^n \to M$이 가군 사상이고,
$N \subset \Ker(f)$가 임의의 부분가군이며,
$h: R^n \to R^{m}$가 $N \subset \Ker(h)$를 만족하고
$f$가 $h$를 통해 인수분해되는 사상이라고 하자. 그러면 임의의
$x \in \Ker(f)$에 대해 $N + Rx \subset \Ker(h')$이고
$f$가 $h'$를 통해 인수분해되도록 하는 사상
$h': R^n \to R^{m'}$을 찾을 수 있다.
\item $f: R^n \to M$이 가군 사상이고
$N \subset \Ker(f)$가 유한 생성 부분가군이면,
$f = g \circ h$ 및 $N \subset \Ker(h)$를 만족하는 가군 사상
$h: R^n \to R^m$과 $g: R^m \to M$이 존재한다.
\end{enumerate}
\end{lemma}

\begin{proof}
(1)과 (2)의 동치는 평탄성의 등식 판정법을 다시 표현한 것에
불과하다\footnote{실제로 가군 사상 $f : R^n \to M$은 $M$의
원소들 $x_1, x_2, \ldots, x_n$의 선택에 대응한다(즉 표준
기저 원소들 $e_1, e_2, \ldots, e_n$의 상들이다). 또한
원소 $x \in \Ker(f)$는 이 $x_1, x_2, \ldots, x_n$ 사이의
관계에 대응한다(즉 $x$의 좌표를 $f_i$라 할 때의 관계
$\sum_i f_i x_i = 0$이다). 가군 사상 $h$는
$m \times n$ 행렬로 나타내며 보조정리 \ref{algebra-lemma-flat-eq}의
행렬 $(a_{ij})$에 대응하고, 보조정리 \ref{algebra-lemma-flat-eq}의
$y_j$들은 $g$에 의한 $R^m$의 표준 기저벡터들의 상이다.}.
(2)가 (3)을 함의함을 보이기 위해, $f$가
$f = g \circ h$로 인수분해되게 하는 사상 $g: R^m \to M$을
택하자. (2)에 의해 사상 $h'': R^m \to R^{m'}$을 찾되,
$h''$가 $h(x)$를 영으로 보내고 $g: R^m \to M$이
$h''$를 통해 인수분해되도록 한다.
그러면 $h' = h'' \circ h$로 두면 된다. (3)이 (4)를 함의함은
(4)의 $N \subset \Ker(f)$의 생성원 개수에 관한 귀납법으로
따른다. (4)가 (2)를 함의함은 명백하다.
\end{proof}

\begin{lemma}
\label{algebra-lemma-flat-factors-fp}
$M$을 $R$-가군이라 하자. 그러면 $M$이 평탄일 필요충분조건은
다음 조건이 성립하는 것이다. $P$가 유한 표시 $R$-가군이고
$f: P \to M$이 가군 사상이면, 유한 자유 $R$-가군 $F$와
$f = g \circ h$를 만족하는 가군 사상
$h: P \to F$ 및 $g: F \to M$이 존재한다.
\end{lemma}

\begin{proof}
이는 보조정리 \ref{algebra-lemma-flat-factors-free}의 조건 (4)를 다시
표현한 것에 불과하다.
\end{proof}

\begin{lemma}
\label{algebra-lemma-flat-surjective-hom}
$M$을 $R$-가군이라 하자. 그러면 $M$이 평탄일 필요충분조건은
다음 조건이 성립하는 것이다. 모든 유한 표시 $R$-가군 $P$에 대해,
$N \to M$이 전사 $R$-가군 사상이면 유도된 사상
$\Hom_R(P, N) \to \Hom_R(P, M)$은 전사이다.
\end{lemma}

\begin{proof}
먼저 $M$이 평탄하다고 하자. $P$가 유한 표시일 때, 주어진 사상
$f: P \to M$이 사상 $N \to M$을 통해 인수분해됨을 보여야 한다.
보조정리 \ref{algebra-lemma-flat-factors-fp}에 의해 사상 $f$는
$F$가 유한 자유인 어떤 사상 $F \to M$을 통해 인수분해된다.
$F$가 자유이므로 이 사상은 $N \to M$을 통해 인수분해된다.
따라서 $f$도 $N \to M$을 통해 인수분해된다.

\medskip\noindent
역으로 보조정리의 조건이 성립한다고 하자. 유한 표시 가군
$P$에서 오는 사상 $f: P \to M$을 택하자. $M$ 위로의 전사
$N \to M$을 갖는 자유가군 $N$을 택한다. 그러면 $f$는
$N \to M$을 통해 인수분해되고, $P$가 유한 생성이므로 $f$는
$N$의 유한 자유 부분가군을 통해 인수분해된다. 따라서 $M$은
보조정리 \ref{algebra-lemma-flat-factors-fp}의 조건을 만족하며, 평탄하다.
\end{proof}

\begin{theorem}[Lazard 정리]
\label{algebra-theorem-lazard}
$M$을 $R$-가군이라 하자. 그러면 $M$이 평탄일 필요충분조건은
유한 자유 $R$-가군들의 유향계의 여극한인 것이다.
\end{theorem}

\begin{proof}
평탄 가군들의 유향계의 여극한은 평탄하다. 유향 여극한을 취하는
것은 완전하고 텐서곱과 가환하기 때문이다. 따라서 $M$이 유한 자유
가군들의 유향계의 여극한이면 $M$은 평탄하다.

\medskip\noindent
역을 위해 먼저, 다음과 같은 방식으로 임의의 가군 $M$을 유한 표시
가군들의 유향계의 여극한으로 쓸 수 있음을 상기하자. 어떤 집합
$I$에 대해 전사 $f: R^I \to M$을 택하고 그 핵을 $K$라 하자.
$E$를 $J$가 $I$의 유한 부분집합이고 $N$이
$R^J \cap K$의 유한 생성 부분가군인 순서쌍 $(J, N)$들의 집합이라
하자. $(J, N) \leq (J', N')$를 $J \subset J'$이고
$N \subset N'$인 것과 동치로 정의하면 $E$는 유향 부분순서집합이
된다. $e = (J, N)$에 대해 $M_e = R^J/N$이라 정의하고,
$e \leq e'$일 때 $f_{ee'}: M_e \to M_{e'}$를 자연스러운 사상으로
정의하자. 그러면 $(M_e, f_{ee'})$는 유향계이고, 자연스러운 사상들
$f_e: M_e \to M$은 동형사상
$\colim_{e \in E} M_e \xrightarrow{\cong} M$을 유도한다.

\medskip\noindent
이제 $M$이 평탄하다고 하자. $I = M \times \mathbf{Z}$라 두고
$R^{I}$의 표준 기저를 $(x_i)$라 쓰며, 위 논의의
$f: R^I \to M$을 $x_i$를 $i$의 $M$으로의 사영에 보내는
사상으로 택하자. 정리를 증명하려면 $M_e$가 자유인 $e \in E$들이
$E$의 공종 부분집합을 이룸을 보이는 것으로 충분하다.
따라서 임의의 $e = (J, N) \in E$를 택하자.
보조정리 \ref{algebra-lemma-flat-factors-fp}에 의해 유한 자유 가군 $F$와
사상 $h: R^J/N \to F$ 및 $g: F \to M$이 존재하여 자연스러운
사상 $f_e: R^J/N \to M$은
$R^J/N \xrightarrow{h} F \xrightarrow{g} M$으로 인수분해된다.
어떤 $e' \geq e$에 대해 $F$를 $M_{e'}$로 실현하고자 한다.

\medskip\noindent
$F$의 유한 기저 $\{ b_1, \ldots, b_n \}$을 택하자.
모든 $\ell$에 대해 $i_{\ell} \notin J$이고,
$f: R^I \to M$에 의한 $x_{i_{\ell}}$의 상이
$g: F \to M$에 의한 $b_{\ell}$의 상과 같도록 하는 서로 다른
$n$개의 원소 $i_1, \ldots, i_n \in I$를 택한다.
$I$의 선택에 의해 $M$의 모든 원소는 서로 다른 무한히 많은
$i \in I$에 대해 $f(x_i)$로 쓸 수 있으므로 이는 가능하다.
이제 $J' = J \cup \{i_1, \ldots , i_n \}$라 두고,
$i \in J$이면 $x_i \mapsto h(x_i)$로,
$\ell = 1, \ldots, n$이면 $x_{i_{\ell}} \mapsto b_{\ell}$로
보내어 $R^{J'} \to F$를 정의하자.
$N' = \Ker(R^{J'} \to F)$라 두자. 다음을 관찰하자.
\begin{enumerate}
\item 정사각형
$$
\xymatrix{
R^{J'} \ar[r] \ar@{^{(}->}[d] & F \ar[d]^{g} \\
R^{I} \ar[r]_{f} & M
}
$$
은 가환하고,
%$R^{J'} \to F$ factors $f: R^I \to M$,
따라서 $N' \subset K = \Ker(f)$이다.
\item $R^{J'} \to F$는 유한 자유 가군 위로의 전사이므로 분할된다.
따라서 $N'$은 유한 생성이다.
\item $J \subset J'$이고 $N \subset N'$이다.
\end{enumerate}
(1)과 (2)에 의해 $e' = (J', N')$은 $E$에 속하고,
(3)에 의해 $e' \geq e$이며, 구성에 의해
$M_{e'} = R^{J'}/N' \cong F$는 자유이다.
\end{proof}

\section{보편 단사 가군 사상}
\label{algebra-section-universally-injective}

\noindent
이제 평탄 가군과 어떤 의미에서 상보적인 보편 단사 가군 사상을
논한다(보조정리 \ref{algebra-lemma-flat-universally-injective}를 보라).
Lazard의 학위논문 \cite{Autour}을 따르며, \cite{Lam}도 보라.

\begin{definition}
\label{algebra-definition-universally-injective}
$f: M \to N$을 $R$-가군들의 사상이라 하자. 모든 $R$-가군 $Q$에
대해 사상
$f \otimes_R \text{id}_Q: M \otimes_R Q \to N \otimes_R Q$가
단사이면 $f$를 {\it 보편 단사}라고 한다. $R$-가군들의 열
$0 \to M_1 \to M_2 \to M_3 \to 0$이 완전하고
$M_1 \to M_2$가 보편 단사이면 이를 {\it 보편 완전}이라고 한다.
\end{definition}

\begin{example}
\label{algebra-example-universally-exact}
보편 완전열의 예는 다음과 같다.
\begin{enumerate}
\item 분할 짧은 완전열은 보편 완전이다. 텐서곱은 직합을 취하는
것과 가환하기 때문이다.
\item 보편 완전열들의 유향계의 여극한은 보편 완전이다. 이는 유향
여극한을 취하는 것이 완전하고 텐서곱이 여극한을 취하는 것과
가환한다는 사실에서 따른다. 특히 분할 완전열들의 유향계의 여극한은
보편 완전이다. 역으로 모든 보편 완전열이 이 방식으로 생긴다는 것을
아래에서 보게 된다.
\end{enumerate}
\end{example}

\noindent
이제 짧은 완전열이 보편 완전일 판정법들을 나열한다. 이들은 앞에서
준 평탄성 판정법들의 유사물이다. 아래 (3)--(6)은 각각 보조정리
\ref{algebra-lemma-flat-eq}, \ref{algebra-lemma-flat-factors-free},
\ref{algebra-lemma-flat-surjective-hom} 및 정리 \ref{algebra-theorem-lazard}의
평탄성 판정법에 대응한다.

\begin{theorem}
\label{algebra-theorem-universally-exact-criteria}
$R$-가군들의 완전열
$$
0 \to M_1 \xrightarrow{f_1} M_2 \xrightarrow{f_2} M_3 \to 0
$$
이 주어졌다고 하자. 다음 조건들은 동치이다.
\begin{enumerate}
\item 열 $0 \to M_1 \to M_2 \to M_3 \to 0$은 보편 완전이다.
\item 모든 유한 표시 $R$-가군 $Q$에 대해 열
$$
0 \to M_1 \otimes_R Q \to M_2 \otimes_R Q \to
M_3 \otimes_R Q \to 0
$$
은 완전하다.
\item 모든 $i$에 대해
$$
f_1(x_i) = \sum\nolimits_j a_{ij} y_j,
$$
를 만족하는 원소 $x_i \in M_1$ $(i = 1, \ldots, n)$,
$y_j \in M_2$ $(j = 1, \ldots, m)$ 및
$a_{ij} \in R$ $(i = 1, \ldots, n, j = 1, \ldots, m)$가
주어지면, 모든 $i$에 대해
$$
x_i = \sum\nolimits_j a_{ij} z_j .
$$
를 만족하는 $z_j \in M_1$ $(j =1, \ldots, m)$가 존재한다.
\item $R$-가군 사상들의 가환 그림
$$
\xymatrix{
R^n \ar[r] \ar[d] &  R^m \ar[d] \\
M_1 \ar[r]^{f_1}        &  M_2
}
$$
이 주어졌다고 하자. 여기서 $m$과 $n$은 정수이다. 그러면 위쪽
삼각형을 가환하게 하는 사상 $R^m \to M_1$이 존재한다.
\item 모든 유한 표시 $R$-가군 $P$에 대해 $R$-가군 사상
$\Hom_R(P, M_2) \to \Hom_R(P, M_3)$은 전사이다.
\item 열 $0 \to M_1 \to M_2 \to M_3 \to 0$은 다음 꼴의 분할
완전열들의 유향계의 여극한이다.
$$
0 \to M_{1} \to M_{2, i} \to M_{3, i} \to 0
$$
여기서 $M_{3, i}$들은 유한 표시이다.
\end{enumerate}
\end{theorem}

\begin{proof}
(1)이 (2)를 함의함은 명백하다.

\medskip\noindent
이제 (2)가 (3)을 함의함을 보이자. (3)에서와 같은 관계
$f_1(x_i) = \sum_j a_{ij} y_j$가 주어졌다고 하자.
$(d_j)$를 $R^m$의 기저, $(e_i)$를 $R^n$의 기저라 하고,
$R^m \to R^n$을 $d_j \mapsto \sum_i a_{ij} e_i$로 주어지는
사상이라 하자. $Q$를 $R^m \to R^n$의 여핵이라 하자.
그러면 $R^m \to R^n \to Q \to 0$에 사상
$f_1: M_1 \to M_2$를 텐서하여 다음 가환 그림을 얻는다.
$$
\xymatrix{
M_1^{\oplus m} \ar[r] \ar[d] & M_1^{\oplus n} \ar[r] \ar[d] & M_1 \otimes_R Q
\ar[r] \ar[d] & 0 \\
M_2^{\oplus m} \ar[r] & M_2^{\oplus n} \ar[r] & M_2 \otimes_R Q \ar[r] & 0
}
$$
여기서 $M_1^{\oplus m} \to M_1^{\oplus n}$은
$$
(z_1, \ldots, z_m) \mapsto
(\sum\nolimits_j a_{1j} z_j, \ldots, \sum\nolimits_j a_{nj} z_j),
$$
로 주어지고, $M_2^{\oplus m} \to M_2^{\oplus n}$도 마찬가지로
주어진다. $x = (x_1, \ldots, x_n) \in M_1^{\oplus n}$이
$M_1^{\oplus m} \to M_1^{\oplus n}$의 상에 속함을 보이고자 한다.
(2)에 의해 사상 $M_1 \otimes Q \to M_2 \otimes Q$는 단사이므로,
윗줄의 완전성에 의해 $x$가 $M_2 \otimes Q$에서 $0$으로
사상됨을 보이는 것으로 충분하다. 아랫줄의 완전성에 의해 이는
$M_2^{\oplus n}$에서의 $x$의 상이
$M_2^{\oplus m} \to M_2^{\oplus n}$의 상에 속함을 보이는 것으로
충분하다. 이는 가정에 의해 참이다.

\medskip\noindent
조건 (4)는 (3)을 그림의 형태로 옮겨 쓴 것에 불과하다.

\medskip\noindent
이제 (4)가 (5)를 함의함을 보이자.
$\varphi : P \to M_3$를 유한 표시 $R$-가군 $P$에서 오는
사상이라 하자. $\varphi$가 사상 $P \to M_2$로 올려짐을
보여야 한다. $P$의 표시
$$
R^n \xrightarrow{g_1} R^m \xrightarrow{g_2} P \to 0.
$$
를 택하자. $R^n$과 $R^m$의 자유성을 이용하여 다음 그림이
가환하도록 먼저 $h_2: R^m \to M_2$를, 이어서
$h_1: R^n \to M_1$을 구성할 수 있다.
$$
\xymatrix{
         & R^n \ar[r]^{g_1} \ar[d]^{h_1} & R^m \ar[r]^{g_2} \ar[d]^{h_2} & P
\ar[r] \ar[d]^{\varphi} & 0 \\
0 \ar[r] & M_1 \ar[r]^{f_1} & M_2 \ar[r]^{f_2} & M_3 \ar[r] & 0 .
}
$$
(4)에 의해 $k_1 \circ g_1 = h_1$을 만족하는 사상
$k_1: R^m \to M_1$이 존재한다. 이제 $h_2' = h_2 - f_1 \circ k_1$로
$h'_2: R^m \to M_2$를 정의하자. 그러면
$$
h'_2 \circ g_1 = h_2 \circ g_1 - f_1 \circ k_1 \circ g_1 =
h_2 \circ g_1 - f_1 \circ h_1 = 0 .
$$
따라서 몫으로 내려가면 $h'_2$는
$\varphi' \circ g_2 = h_2'$를 만족하는 사상
$\varphi': P \to M_2$를 정의한다. 그림으로 쓰면 다음과 같다.
$$
\xymatrix{
R^m \ar[r]^{g_2} \ar[d]_{h'_2} & P \ar[d]^{\varphi} \ar[dl]_{\varphi'} \\
M_2 \ar[r]^{f_2} & M_3.
}
$$
여기서 위쪽 삼각형은 가환한다. $\varphi'$가 원하는 올림, 즉
$f_2 \circ \varphi' = \varphi$를 만족한다고 주장한다.
정의에서 다음을 얻는다.
$$
f_2 \circ \varphi' \circ g_2 = f_2 \circ h'_2 =
f_2 \circ h_2 - f_2 \circ f_1 \circ k_1 = f_2 \circ h_2 =
\varphi \circ g_2.
$$
$g_2$가 전사이므로 증명이 끝난다.

\medskip\noindent
이제 (5)가 (6)을 함의함을 보이자. $M_{3}$을 유한 표시 가군들
$M_{3, i}$의 유향계의 여극한으로 쓰자. 보조정리
\ref{algebra-lemma-module-colimit-fp}를 보라. $M_{2, i}$를 $M_{3}$ 위에서
$M_{3, i}$와 $M_{2}$의 섬유곱이라 하자. 정의에 의해 이는 두
$M_3$으로의 사영이 같은 원소들로 이루어진 $M_2 \times M_{3, i}$의
부분가군이다. $M_{1, i}$를 사영 $M_{2, i} \to M_{3, i}$의 핵이라
하자. 그러면 완전열들의 유향계
$$
0 \to M_{1, i} \to M_{2, i} \to M_{3, i} \to 0,
$$
와 각 $i$에 대해 다음 완전열들의 사상을 얻는다.
$$
\xymatrix{
0  \ar[r] & M_{1, i} \ar[d] \ar[r]  &  M_{2, i}  \ar[r] \ar[d] & M_{3, i} \ar[d]
\ar[r] & 0 \\
0  \ar[r] & M_{1} \ar[r]  &  M_{2}  \ar[r] & M_{3} \ar[r] & 0
}
$$
이들은 유향계와 양립한다. 섬유곱 $M_{2, i}$의 정의로부터 사상
$M_{1, i} \to M_1$은 동형사상이다. (5)에 의해
$M_{3, i} \to M_3$의 올림인 사상 $M_{3, i} \to M_{2}$가 존재하고,
섬유곱의 보편 성질에 의해 이는 $M_{2, i} \to M_{3, i}$의 단면을
준다. 따라서 열
$$
0 \to M_{1, i} \to M_{2, i} \to M_{3, i} \to 0
$$
은 분할된다. 여극한을 취하면 다음 가환 그림을 얻는다.
$$
\xymatrix{
0  \ar[r] & \colim M_{1, i} \ar[d]^{\cong} \ar[r]  &  \colim M_{2, i}
 \ar[r]
\ar[d] & \colim M_{3, i} \ar[d]^{\cong} \ar[r] & 0 \\
0  \ar[r] & M_{1} \ar[r]  &  M_{2}  \ar[r] & M_{3} \ar[r] & 0
}
$$
두 행은 완전하고 양 끝의 수직 사상은 동형사상이다. 따라서
$\colim M_{2, i} \to M_2$ 역시 동형사상이고 (6)이 성립한다.

\medskip\noindent
조건 (6)이 (1)을 함의함은 예 \ref{algebra-example-universally-exact} (2)에서
따른다.
\end{proof}

\noindent
앞 정리는 보편 완전열이 언제나 분할 짧은 완전열들의 여극한임을
보인다. 보편 단사 사상의 여핵이 유한 표시이면 실제로 사상 자체가
분할된다.

\begin{lemma}
\label{algebra-lemma-universally-exact-split}
$R$-가군들의 완전열
$$
0 \to M_1 \to M_2 \to M_3 \to 0
$$
이 주어졌다고 하자. $M_3$가 유한 표시라고 가정하자. 그러면
$$
0 \to M_1 \to M_2 \to M_3 \to 0
$$
이 보편 완전일 필요충분조건은 분할되는 것이다.
\end{lemma}

\begin{proof}
분할 짧은 완전열은 언제나 보편 완전이다. 예
\ref{algebra-example-universally-exact}를 보라. 역으로 이 열이 보편
완전이면, 정리 \ref{algebra-theorem-universally-exact-criteria} (5)를
$P = M_3$에 적용하여 사상 $M_2 \to M_3$가 단면을 가짐을 얻는다.
\end{proof}

\noindent
다음 보조정리는 보편 단사 사상과 평탄 가군이 어떻게 상보적인지
보여 준다.

\begin{lemma}
\label{algebra-lemma-flat-universally-injective}
$M$을 $R$-가군이라 하자. 그러면 $M$이 평탄일 필요충분조건은
$R$-가군들의 임의의 완전열
$$
0 \to M_1 \to M_2 \to M \to 0
$$
이 보편 완전인 것이다.
\end{lemma}

\begin{proof}
이는 보조정리 \ref{algebra-lemma-flat-surjective-hom}와 정리
\ref{algebra-theorem-universally-exact-criteria} (5)에서 따른다.
\end{proof}

\begin{example}
\label{algebra-example-universally-exact-non-split-non-flat}
분할되지 않거나 평탄하지 않은 보편 완전열의 예는 다음과 같다.
\begin{enumerate}
\item 보조정리 \ref{algebra-lemma-universally-exact-split}에도 불구하고,
보편 완전이지만 분할되지 않는 $R$-가군들의 짧은 완전열
$$
0 \to M_1 \to M_2 \to M_3 \to 0
$$
이 존재할 수 있다. 예를 들어 $R = \mathbf{Z}$로 두고,
$M_1 = \bigoplus_{n=1}^{\infty} \mathbf{Z}$,
$M_{2} = \prod_{n = 1}^{\infty} \mathbf{Z}$라 하며,
$M_{3}$를 포함 사상 $M_1 \to M_2$의 여핵이라 하자.
$M_1, M_2, M_3$는 모두 꼬임이 없으므로 평탄하다
(More on Algebra, 보조정리
\ref{more-algebra-lemma-dedekind-torsion-free-flat}). 따라서 보조정리
\ref{algebra-lemma-flat-universally-injective}에 의해
$$
0 \to M_1 \to M_2 \to M_3 \to 0
$$
은 보편 완전이다. 그러나 단면 $s: M_3 \to M_2$는 존재할 수 없다.
실제로 $x$를 원소 $(2, 2^2, 2^3, \ldots) \in M_2$가
$M_3$에서 갖는 상이라 하면, 임의의 가군 사상 $s: M_3 \to M_2$는
$x$를 영으로 보내야 한다. 이는 모든 $n \geq 1$에 대해
$x \in 2^n M_3$이고, 따라서 모든 $n \geq 1$에 대해 $s(x)$가
$2^n$으로 나누어지므로 $0$이어야 하기 때문이다.
\item 보조정리 \ref{algebra-lemma-flat-universally-injective}에도 불구하고,
$M_1, M_2, M_3$가 모두 평탄하지 않으면서 보편 완전인
$R$-가군들의 짧은 완전열
$$
0 \to M_1 \to M_2 \to M_3 \to 0
$$
이 존재할 수 있다. 실제로 $M$이 임의의 비평탄 가군이면 분할 완전열
$$
0 \to M \to M \oplus M \to M \to 0.
$$
을 택하면 된다. 예를 들어 $R = \mathbf{Z}$ 위에서는 $M$을 임의의
꼬임 가군으로 택한다.
\item (1)에서와 같은 완전열과 (2)에서와 같은 완전열의 직합을
취하면, 보편 완전이고 분할되지 않으며
$M_1, M_2, M_3$가 모두 평탄하지 않은 $R$-가군들의 짧은 완전열
$$
0 \to M_1 \to M_2 \to M_3 \to 0
$$
을 얻는다.
\end{enumerate}
\end{example}

\begin{lemma}
\label{algebra-lemma-ui-flat-domain}
$0 \to M_1 \to M_2 \to M_3 \to 0$을 $R$-가군들의 보편 완전열이라
하고, $M_2$가 평탄하다고 가정하자. 그러면 $M_1$과 $M_3$는
평탄하다.
\end{lemma}

\begin{proof}
$0 \to N \to N' \to N'' \to 0$을 $R$-가군들의 짧은 완전열이라
하자. 다음 가환 그림을 생각하자.
$$
\xymatrix{
M_1 \otimes_R N \ar[r] \ar[d] &
M_2 \otimes_R N \ar[r] \ar[d] &
M_3 \otimes_R N \ar[d] \\
M_1 \otimes_R N' \ar[r] \ar[d] &
M_2 \otimes_R N' \ar[r] \ar[d] &
M_3 \otimes_R N' \ar[d] \\
M_1 \otimes_R N'' \ar[r] &
M_2 \otimes_R N'' \ar[r] &
M_3 \otimes_R N''
}
$$
(경계의 $0$들은 생략하였다.) 가정에 의해 각 행은 짧은 완전열이고
화살표 $M_2 \otimes N \to M_2 \otimes N'$은 단사이다.
이는 명백히 $M_1 \otimes N \to M_1 \otimes N'$이 단사임을
함의하므로 $M_1$이 평탄임을 알 수 있다. 특히 왼쪽과 가운데 열에서
짧은 완전열을 얻는다. 그림 추적에 의해 화살표
$M_3 \otimes N \to M_3 \otimes N'$이 단사임이 따른다.
따라서 $M_3$는 평탄하다.
\end{proof}

\begin{lemma}
\label{algebra-lemma-universally-injective-tensor}
$R$을 환이라 하자. $M \to M'$을 보편 단사 $R$-가군 사상이라
하자. 그러면 임의의 $R$-가군 $N$에 대해 사상
$M \otimes_R N \to M' \otimes_R N$은 보편 단사이다.
\end{lemma}

\begin{proof}
생략한다.
\end{proof}

\begin{lemma}
\label{algebra-lemma-composition-universally-injective}
$R$을 환이라 하자. 보편 단사 $R$-가군 사상들의 합성은 보편
단사이다.
\end{lemma}

\begin{proof}
생략한다.
\end{proof}

\begin{lemma}
\label{algebra-lemma-universally-injective-permanence}
$R$을 환이라 하자. $M \to M'$과 $M' \to M''$을 $R$-가군 사상이라
하자. 그 합성 $M \to M''$이 보편 단사이면 $M \to M'$도 보편
단사이다.
\end{lemma}

\begin{proof}
생략한다.
\end{proof}

\begin{lemma}
\label{algebra-lemma-universally-injective-check-stalks}
$R \to S$를 환 사상이라 하고 $M \to M'$을 $S$-가군들의 사상이라
하자. 다음 조건들은 동치이다.
\begin{enumerate}
\item $M \to M'$은 $R$-가군들의 사상으로서 보편 단사이다.
\item $S$의 각 소 아이디얼 $\mathfrak q$에 대해 사상
$M_{\mathfrak q} \to M'_{\mathfrak q}$은 $R$-가군들의 사상으로서
보편 단사이다.
\item $S$의 각 극대 아이디얼 $\mathfrak m$에 대해 사상
$M_{\mathfrak m} \to M'_{\mathfrak m}$은 $R$-가군들의 사상으로서
보편 단사이다.
\item $S$의 각 소 아이디얼 $\mathfrak q$에 대해 사상
$M_{\mathfrak q} \to M'_{\mathfrak q}$은
$R_{\mathfrak p}$-가군들의 사상으로서 보편 단사이다. 여기서
$\mathfrak p$는 $R$에서 $\mathfrak q$의 역상이다.
\item $S$의 각 극대 아이디얼 $\mathfrak m$에 대해 사상
$M_{\mathfrak m} \to M'_{\mathfrak m}$은
$R_{\mathfrak p}$-가군들의 사상으로서 보편 단사이다. 여기서
$\mathfrak p$는 $R$에서 $\mathfrak m$의 역상이다.
\end{enumerate}
\end{lemma}

\begin{proof}
$N$을 $R$-가군이라 하자. $\mathfrak q$를 $R$의 소 아이디얼
$\mathfrak p$ 위에 놓인 $S$의 소 아이디얼이라 하자. 그러면
$$
(M \otimes_R N)_{\mathfrak q} =
M_{\mathfrak q} \otimes_R N =
M_{\mathfrak q} \otimes_{R_{\mathfrak p}} N_{\mathfrak p}.
$$
또한 $M'$에 대해서도 같은 것이 성립하고, 국소화는 완전하다.
그리고 $N$이 $R_{\mathfrak p}$-가군이면
$N_{\mathfrak p} = N$이다. 이를 사용하면 동치들을 곧바로
증명할 수 있다.

\medskip\noindent
예를 들어 (5)가 성립한다고 하자.
$K = \Ker(M \otimes_R N \to M' \otimes_R N)$라 두자. 위의
관찰에 의해 $S$의 각 극대 아이디얼 $\mathfrak m$에 대해
$K_{\mathfrak m} = 0$이다. 따라서 보조정리
\ref{algebra-lemma-characterize-zero-local}에 의해 $K = 0$이다.
그러므로 (1)이 성립한다. 역으로 (1)이 성립한다고 하자.
$\mathfrak p \subset R$ 위에 놓인 임의의
$\mathfrak q \subset S$를 택하고, $R_{\mathfrak p}$ 위의 임의의
가군 $N$을 택하자. 가정에 의해
$\Ker(M \otimes_R N \to M' \otimes_R N) = 0$이다.
따라서 위 공식들과 $N = N_{\mathfrak p}$라는 사실에 의해
$\Ker(M_{\mathfrak q} \otimes_{R_{\mathfrak p}} N \to
M'_{\mathfrak q} \otimes_{R_{\mathfrak p}} N) = 0$임을 알 수 있다.
다시 말해 (4)가 성립한다. 물론
(4) $\Rightarrow$ (5)는 즉시 성립한다. 따라서 (1), (4), (5)는
모두 동치이다. 다른 동치들의 증명은 생략한다.
\end{proof}

\begin{lemma}
\label{algebra-lemma-universally-injective-localize}
$\varphi : A \to B$를 환 사상이라 하자. $S \subset A$와
$S' \subset B$를 $\varphi(S) \subset S'$를 만족하는 곱셈적
부분집합이라 하자. $M \to M'$을 $B$-가군들의 사상이라 하자.
\begin{enumerate}
\item $M \to M'$이 $A$-가군들의 사상으로서 보편 단사이면,
$(S')^{-1}M \to (S')^{-1}M'$은 $A$-가군들의 사상으로서도,
$S^{-1}A$-가군들의 사상으로서도 보편 단사이다.
\item $M$과 $M'$이 $(S')^{-1}B$-가군이면, $M \to M'$이
$A$-가군들의 사상으로서 보편 단사일 필요충분조건은
$S^{-1}A$-가군들의 사상으로서 보편 단사인 것이다.
\end{enumerate}
\end{lemma}

\begin{proof}
보조정리 \ref{algebra-lemma-universally-injective-check-stalks}를 사용하여
증명할 수 있지만, 다음과 같이 직접 증명할 수도 있다.
$M \to M'$이 $A$-보편 단사라고 가정하자. $Q$를 $A$-가군이라
하자. 그러면 $Q \otimes_A M \to Q \otimes_A M'$은 단사이다.
국소화는 완전하므로
$(S')^{-1}(Q \otimes_A M) \to (S')^{-1}(Q \otimes_A M')$은
단사이다. $(S')^{-1}(Q \otimes_A M) = Q \otimes_A (S')^{-1}M$이고
$M'$에 대해서도 마찬가지이므로,
$Q \otimes_A (S')^{-1}M \to Q \otimes_A (S')^{-1}M'$이 단사임을
알 수 있다. 따라서 $(S')^{-1}M \to (S')^{-1}M'$은
$A$-가군들의 사상으로서 보편 단사이다.
% P01-E0448
이는 (1)의 첫 부분을 증명한다. (2)를 보기 위해 다음 두 사실을
사용할 수 있다. (a) $Q$가 $S^{-1}A$-가군이면
$Q \otimes_A S^{-1}A = Q$이다. 즉 $A$ 위에서 $Q$와 텐서하는 것은
$S^{-1}A$ 위에서 $Q$와 텐서하는 것과 같다. (b) $M$이 $S$의
원소들이 가역적으로 작용하는 임의의 $A$-가군이면
$M \otimes_A Q = M \otimes_{S^{-1}A} S^{-1}Q$이다.
(2)는 이로부터 즉시 따른다.
\end{proof}

\begin{lemma}
\label{algebra-lemma-check-universally-injective-into-flat}
$R$을 환이라 하고 $M \to M'$을 $R$-가군들의 사상이라 하자.
$M'$이 평탄이면, $M \to M'$이 보편 단사일 필요충분조건은
$R$의 모든 유한 생성 아이디얼 $I$에 대해
$M/IM \to M'/IM'$이 단사인 것이다.
\end{lemma}

\begin{proof}
모든 유한 $R$-가군 $Q$에 대해
$M \otimes_R Q \to M' \otimes_R Q$가 단사임을 보이는 것으로
충분하다. 정리 \ref{algebra-theorem-universally-exact-criteria}를 보라.
그러면 $Q$는 연속 몫들이 순환 가군 $R/I_i$와 동형인 부분가군들에
의한 유한 여과
$0 = Q_0 \subset Q_1 \subset \ldots \subset Q_n = Q$를 갖는다.
보조정리 \ref{algebra-lemma-trivial-filter-finite-module}를 보라.
$M'$이 평탄하므로, 연속 몫들이
$M' \otimes_R R/I_i = M'/I_iM'$인 부분가군
$M' \otimes_R Q_i$에 의한 $M' \otimes_R Q$의 여과
$$
\xymatrix{
M \otimes Q_1 \ar[r] \ar[d] &
M \otimes Q_2 \ar[r] \ar[d] &
\ldots \ar[r] &
M \otimes Q \ar[d] \\
M' \otimes Q_1 \ar@{^{(}->}[r] &
M' \otimes Q_2 \ar@{^{(}->}[r] &
\ldots \ar@{^{(}->}[r] &
M' \otimes Q
}
$$
를 얻는다. 간단한 귀납 논증에 의해
$M/I_i M \to M'/I_i M'$이 단사인지 확인하는 것으로 충분하다.
$I'_i \subset I_i$인 유한 생성 아이디얼들의 모임은 유향집합임을
주목하자. 따라서 $M/I_iM = \colim M/I'_iM$은 여과 여극한이고,
$M'$에 대해서도 마찬가지이다. 가정에 의해 사상
$M/I'_iM \to M'/I'_i M'$은 단사이며, 여과 여극한이 완전하므로
(보조정리 \ref{algebra-lemma-directed-colimit-exact}) 결론을 얻는다.
\end{proof}

\begin{lemma}
\label{algebra-lemma-base-change-along-universally-injective-ring-map}
$R \to S$를 $R$-가군들의 사상으로서 보편 단사인 환 사상이라 하자.
그러면 $R$-가군 위의 함자 $M \mapsto M \otimes_R S$는 단사,
전사 및 동형사상을 반영한다.
\end{lemma}

\begin{proof}
$M \to N$을 핵 $K$와 여핵 $Q$를 갖는 $R$-가군들의 사상이라 하자.
$M \otimes_R S \to N \otimes_R S$가 단사이면
$K \otimes_R S \to M \otimes_R S$의 상은 영이다.
$K \subset K \otimes_R S$이고 $M \subset M \otimes_R S$이므로
$K = 0$임을 얻는다. 한편
$M \otimes_R S \to N \otimes_R S$가 전사이면,
텐서곱의 오른쪽 완전성에 의해 $Q \otimes_R S$는 영이다.
따라서 $Q \subset Q \otimes_R S$도 영이다.
\end{proof}

\begin{lemma}
\label{algebra-lemma-faithfully-flat-universally-injective}
$R \to S$를 충실 평탄 환 사상이라 하자. 그러면 $R \to S$는
$R$-가군들의 사상으로서 보편 단사이다. 특히 임의의 아이디얼
$I \subset R$에 대해 $R \cap IS = I$이다.
\end{lemma}

\begin{proof}
$N$을 $R$-가군이라 하자. $N \to N \otimes_R S$가 단사임을
보여야 한다. $S$가 $R$-가군으로서 충실 평탄이므로, $S$와 텐서한
뒤 이를 증명하면 충분하다. 따라서
$N \otimes_R S \to N \otimes_R S \otimes_R S$,
$n \otimes s \mapsto n \otimes 1 \otimes s$가 단사임을 보이면
충분하다. 실제로 다음 퇴축이 존재하므로 이는 참이다.
$n \otimes s \otimes s' \mapsto n \otimes ss'$.
\end{proof}




\section{유한 사영 가군에 대한 하강}
\label{algebra-section-finite-projective}

\noindent
이 절에서는 {\it 유한} 사영 가군이라는 성질이 충실 평탄 환 사상을
따라 하강한다는 사실의 초등적 증명을 제시한다. 유한성 조건을
없애면 이 증명은 적용되지 않는다. 그러나 이 방법은 {\it 가산 생성}
사영 가군이라는 성질의 하강을 증명할 때 사용할 방법을 시사한다.
이 절 끝의 논평을 보라.

\begin{lemma}
\label{algebra-lemma-finite-projective-again}
$M$을 $R$-가군이라 하자. 그러면 $M$이 유한 사영일 필요충분조건은
$M$이 유한 표시이고 평탄인 것이다.
\end{lemma}

\begin{proof}
이는 보조정리 \ref{algebra-lemma-finite-projective}의 일부이다.
그러나 이 시점에서는 그 보조정리의 함의 (1) $\Rightarrow$ (2)를
다음과 같이 더 우아하게 증명할 수 있다. $M$이 유한 표시이고
평탄이면 전사 $R^n \to M$을 택하자. 보조정리
\ref{algebra-lemma-flat-surjective-hom}을 $P = M$에 적용하면 사상
$R^n \to M$은 단면을 갖는다. 따라서 $M$은 자유가군의
직합인자이고, 그러므로 사영이다.
\end{proof}

\noindent
다음은 하강하는 가군의 성질들이다.

\begin{lemma}
\label{algebra-lemma-descend-properties-modules}
$R \to S$를 충실 평탄 환 사상이라 하고 $M$을 $R$-가군이라 하자.
그러면 다음이 성립한다.
\begin{enumerate}
\item $S$-가군 $M \otimes_R S$가 유한형이면 $M$은 유한형이다.
\item $S$-가군 $M \otimes_R S$가 유한 표시이면 $M$은 유한
표시이다.
\item $S$-가군 $M \otimes_R S$가 평탄이면 $M$은 평탄이다.
\item 필요에 따라 여기에 더 추가한다.
% P01-E0449
\end{enumerate}
\end{lemma}

\begin{proof}
$M \otimes_R S$가 유한형이라고 가정하자.
$y_1, \ldots, y_m$을 $S$ 위에서 $M \otimes_R S$의 생성원들이라
하자. 어떤 $x_1, \ldots, x_n \in M$에 대해
% P01-E0450
$y_j = \sum x_i \otimes f_i$로 쓰자. 그러면 사상
$\varphi : R^{\oplus n} \to M$은
$\varphi \otimes \text{id}_S : S^{\oplus n} \to M \otimes_R S$가
전사라는 성질을 가짐을 알 수 있다. $R \to S$가 충실 평탄이므로
$\varphi$는 전사이고 $M$은 유한 생성이다.

\medskip\noindent
$M \otimes_R S$가 유한 표시라고 가정하자. (1)에 의해 $M$은
유한형이다. 전사 $R^{\oplus n} \to M$을 택하고 그 핵을 $K$라
쓰자. $R \to S$가 평탄이므로 $K \otimes_R S$는 기저변환
$S^{\oplus n} \to M \otimes_R S$의 핵이다.
$M \otimes_R S$가 유한 표시이므로 $K \otimes_R S$가 유한형임을
얻는다. 따라서 (1)에 의해 $K$는 유한형이고, 그러므로 $M$은 유한
표시이다.

\medskip\noindent
(3)은 보조정리 \ref{algebra-lemma-flatness-descends}이다.
\end{proof}

\begin{proposition}
\label{algebra-proposition-ffdescent-finite-projectivity}
$R \to S$를 충실 평탄 환 사상이라 하고 $M$을 $R$-가군이라 하자.
$S$-가군 $M \otimes_R S$가 유한 사영이면 $M$은 유한 사영이다.
\end{proposition}

\begin{proof}
보조정리 \ref{algebra-lemma-finite-projective-again}과
\ref{algebra-lemma-descend-properties-modules}에서 따른다.
\end{proof}

\noindent
다음 몇 절에서는 d\'evissage를 사용하여 가산 생성인 경우로
환원함으로써 유한성 가정을 없앤다. 가산 생성인 경우의 전략은
보조정리 \ref{algebra-lemma-finite-projective-again}과 유사한 가산 생성
사영 가군의 특징짓기를 찾고, 이 특징짓기가 하강함을 직접 증명하는
것이다. 이를 위해 Mittag-Leffler 가군의 개념을 도입하고, 가군
$M$이 가산 생성이면 사영일 필요충분조건은 평탄하고
Mittag-Leffler인 것임을 증명한다(정리
\ref{algebra-theorem-projectivity-characterization}). $M$이 유한
생성이면 이 명제는 보조정리 \ref{algebra-lemma-finite-projective-again}으로
환원된다. 예 \ref{algebra-example-ML} (1)에 따르면 유한 생성 가군이
Mittag-Leffler일 필요충분조건은 유한 표시이기 때문이다.

\section{가군의 초한 d\'evissage}
\label{algebra-section-transfinite-devissage}

\noindent
이 절에서는 가군을 직합으로 분해하기 위한 d\'evissage 기법을
도입한다. 주된 결과는 사영 가군이 가산 생성 가군들의 직합이라는
것이다(아래 정리 \ref{algebra-theorem-projective-direct-sum}).
\cite{Kaplansky}를 따른다.


\begin{definition}
\label{algebra-definition-devissage}
$M$을 $R$-가군이라 하자. $M$의 {\it 직합 d\'evissage}란,
순서수 $S$로 첨자화되고 포함관계에 관해 증가하며 다음을 만족하는
부분가군들의 족 $(M_{\alpha})_{\alpha \in S}$이다.
\begin{enumerate}
\item[(0)] $M_0 = 0$;
\item[(1)] $M = \bigcup_{\alpha} M_{\alpha}$;
\item[(2)] $\alpha \in S$가 극한 순서수이면
$M_{\alpha} = \bigcup_{\beta < \alpha} M_{\beta}$;
\item[(3)] $\alpha + 1 \in S$이면 $M_{\alpha}$는
$M_{\alpha + 1}$의 직합인자이다.
\end{enumerate}
더욱이
\begin{enumerate}
\item[(4)] $\alpha + 1 \in S$에 대해
$M_{\alpha + 1}/M_{\alpha}$가 가산 생성이다
\end{enumerate}
까지 성립하면 $(M_{\alpha})_{\alpha \in S}$를 $M$의
{\it Kaplansky d\'evissage}라고 한다.
\end{definition}

\noindent
이 용어는 다음 보조정리에 의해 정당화된다.

\begin{lemma}
\label{algebra-lemma-direct-sum-devissage}
$M$을 $R$-가군이라 하자.
$(M_{\alpha})_{\alpha \in S}$가 $M$의 직합 d\'evissage이면
$M \cong \bigoplus_{\alpha + 1 \in S} M_{\alpha + 1}/M_{\alpha}$이다.
\end{lemma}

\begin{proof}
직합 d\'evissage의 성질 (3)에 의해 각 $\alpha \in S$에 대해
포함 사상 $M_{\alpha + 1}/M_{\alpha} \to M$이 존재한다.
이 포함 사상들의 합으로 주어지는 사상
$$
f : \bigoplus\nolimits_{\alpha + 1\in S} M_{\alpha + 1}/M_{\alpha} \to M
$$
을 생각하자. 또한 각 $\beta\in S$에 대해 제한 사상
$$
f_{\beta} :
\bigoplus\nolimits_{\alpha + 1 \leq \beta} M_{\alpha + 1}/M_{\alpha}
\longrightarrow
M
$$
을 생각하자. $S$에 관한 초한 귀납법으로 $f_{\beta}$의 상이
$M_{\beta}$임을 알 수 있다. $\beta=0$이면 이는 $(0)$에 의해
참이다. $\beta+1$이 후계 순서수이고 $\beta$에 대해 참이면
(3)에 의해 $\beta + 1$에 대해서도 참이다. $\beta$가 극한
순서수이고 모든 $\alpha < \beta$에 대해 참이면 (2)에 의해
$\beta$에 대해서도 참이다. 따라서 (1)에 의해 $f$는 전사이다.

\medskip\noindent
$S$에 관한 초한 귀납법은 제한 사상 $f_{\beta}$들이 단사임도
보인다. $\beta = 0$이면 참이다. $\beta+1$이 후계 순서수이고
$f_{\beta}$가 단사라고 하자. 핵의 원소 $x$를 성분
$x_{\alpha + 1} \in M_{\alpha + 1}/M_{\alpha}$들을 사용하여
$x = (x_{\alpha + 1})_{\alpha + 1 \leq \beta + 1}$로 쓰자.
성질 (3)과 $f_{\beta}$의 상이 $M_{\beta}$라는 사실에 의해
$(x_{\alpha + 1})_{\alpha + 1 \leq \beta}$와
$x_{\beta + 1}$은 모두 $0$으로 사상된다. 따라서
$x_{\beta+1} = 0$이고, 제한 사상 $f_{\beta}$가 단사라는
가정에 의해 모든 $\alpha + 1 \leq \beta$에 대해
$x_{\alpha + 1} = 0$이다. 그러므로 $x = 0$이고
$f_{\beta+1}$은 단사이다. $\beta$가 극한 순서수이면 핵의 원소
$x$를 생각하자. 그러면 $x$는 이미 어떤 $\alpha < \beta$에 대한
$f_{\alpha}$의 정의역에 포함된다. 따라서 $x = 0$이고 귀납법이
끝난다. 각 $\beta \in S$에 대해 $f_{\beta}$가 단사이므로
$f$가 단사임을 결론짓는다.
\end{proof}

\begin{lemma}
\label{algebra-lemma-Kaplansky-devissage}
$M$을 $R$-가군이라 하자. 그러면 $M$이 가산 생성 $R$-가군들의
직합일 필요충분조건은 Kaplansky d\'evissage를 갖는 것이다.
\end{lemma}

\begin{proof}
앞 보조정리는 “충분” 방향을 처리한다. 역으로
$M = \bigoplus_{i \in I} N_i$이고 각 $N_i$가 가산 생성
$R$-가군이라고 하자. $I$를 정렬하여 순서수로 생각할 수 있게 하자.
% P01-E0451
그러면 $M_i = \bigoplus_{j < i} N_j$로 두어 $M$의 Kaplansky
d\'evissage $(M_i)_{i \in I}$를 얻는다.
\end{proof}

\begin{theorem}
\label{algebra-theorem-kaplansky-direct-sum}
$M$이 가산 생성 $R$-가군들의 직합이라고 하자. $P$가 $M$의
직합인자이면 $P$ 역시 가산 생성 $R$-가군들의 직합이다.
\end{theorem}

\begin{proof}
$M = P \oplus Q$로 쓰자. 정의의 성질 (0)--(4)에 더하여 다음을
만족하는 $M$의 Kaplansky d\'evissage
$(M_{\alpha})_{\alpha \in S}$를 구성하고자 한다.
\begin{enumerate}
\item[(5)] 각 $M_{\alpha}$는 $M$의 직합인자이다.
\item[(6)] $M_{\alpha} = P_{\alpha} \oplus Q_{\alpha}$이다.
여기서 $P_{\alpha} =P \cap M_{\alpha}$이고
$Q_\alpha = Q \cap M_{\alpha}$이다.
\end{enumerate}
(주의: 성질 (0)--(2)가 성립하면, 실제로 성질 (3)과 성질 (5)는
동치이다.)

\medskip\noindent
이것이 정리를 함의함을 보기 위해,
$(P_{\alpha})_{\alpha \in S}$가 $P$의 Kaplansky d\'evissage를
이룸을 보이면 충분하다. 성질 (0), (1), (2)는 명백하다.
$(M_{\alpha})$에 대한 (5)와 (6)에 의해 각 $P_{\alpha}$는
$M$의 직합인자이다. $P_{\alpha} \subset P_{\alpha + 1}$이므로
이는 $P_{\alpha}$가 $P_{\alpha + 1}$의 직합인자임을 함의한다.
따라서 $(P_{\alpha})$에 대해 (3)이 성립한다. (4)를 위해
다음을 주목하자.
$$
M_{\alpha + 1}/M_{\alpha} \cong P_{\alpha + 1}/P_{\alpha} \oplus
Q_{\alpha + 1}/Q_{\alpha},
$$
$M_{\alpha + 1}/M_{\alpha}$가 가산 생성이므로
$P_{\alpha + 1}/P_{\alpha}$도 가산 생성이다.

\medskip\noindent
$M_{\alpha}$들을 구성하는 일만 남았다.
$M = \bigoplus_{i \in I} N_i$이고 각 $N_i$가 가산 생성
$R$-가군이라고 쓰자. $I$의 정렬을 하나 택하자. 초한 재귀법으로
각 순서수 $\alpha$에 대해 $M$의 증가하는 부분가군
$M_{\alpha}$를 정의하되, $M_{\alpha}$가 $N_i$들 가운데 일부의
직합이 되도록 할 것이다.

\medskip\noindent
$\alpha = 0$이면 $M_{0} = 0$으로 둔다. $\alpha$가 극한 순서수이고
모든 $\beta < \alpha$에 대해 $M_{\beta}$가 정의되어 있으면
$M_{\alpha} = \bigcup_{\beta < \alpha} M_{\beta}$로 정의한다.
각 $\beta < \alpha$에 대한 $M_{\beta}$가 $N_i$들 가운데 일부의
직합이므로 $M_{\alpha}$도 그러하다. $\alpha + 1$이 후계
순서수이고 $M_{\alpha}$가 정의되어 있으면 다음과 같이
$M_{\alpha + 1}$을 정의한다. $M_{\alpha} = M$이면
$M_{\alpha + 1} = M$으로 둔다. 그렇지 않으면 $N_j$가
$M_{\alpha}$에 포함되지 않는 가장 작은 $j \in I$를 택한다.
다음을 만족하는 무한 행렬 $(x_{mn}), m, n = 1, 2, 3, \ldots$을
구성할 것이다.
\begin{enumerate}
\item $N_j$는 성분 $x_{mn}$들이 생성하는 $M$의 부분가군에
포함된다.
\item 임의의 성분 $x_{k\ell}$을 그 $P$-성분과 $Q$-성분으로
$x_{k\ell} = y_{k\ell} + z_{k\ell}$이라 쓸 때,
행렬 $(x_{mn})$은 $y_{k\ell}$ 또는 $z_{k\ell}$이 영이 아닌
성분을 갖는 각 $N_i$의 생성원 집합을 포함한다.
\end{enumerate}
그러면 $M_{\alpha + 1}$을 $M_{\alpha}$와 모든 $x_{mn}$이 생성하는
$M$의 부분가군으로 정의한다. 행렬 $(x_{mn})$의 성질 (2)에 의해
$M_{\alpha + 1}$은 $N_i$들 가운데 일부의 직합이다.
행렬 $(x_{mn})$을 구성하기 위해,
$x_{11}, x_{12}, x_{13}, \ldots$를 $N_j$의 가산 생성원 집합이라
하자. $x_{11} = y_{11} + z_{11}$이 $P$-성분과 $Q$-성분으로의
분해이면, $x_{21}, x_{22}, x_{23}, \ldots$를 $y_{11}$ 또는
$z_{11}$이 영이 아닌 성분을 갖는 $N_i$들의 합의 가산 생성원
집합이라 하자. 이 과정을 $x_{12}$에 반복하여 행렬의 셋째 행인
$x_{31}, x_{32}, \ldots$를 얻는다. $x_{21}$에 반복하여 넷째
행을, $x_{13}$에 반복하여 다섯째 행을 얻는 식으로 진행하며,
아래에 나타낸 것처럼 연속하는 역대각선을 따라 내려간다.
$$
\left(
\vcenter{
\xymatrix@R=2mm@C=2mm{
x_{11} & x_{12} \ar[dl] & x_{13} \ar[dl] & x_{14} \ar[dl] & \ldots  \\
x_{21} & x_{22} \ar[dl] & x_{23} \ar[dl] & \ldots  \\
x_{31} & x_{32} \ar[dl] & \ldots \\
x_{41} & \ldots \\
\ldots
}
}
\right).
$$

\medskip\noindent
$I$에 관한 초한 귀납법으로 각 $i \in I$에 대해 $N_i$가 어떤
$M_{\alpha}$에 포함됨을 알 수 있다. 여기서는 $N_j$가
$M_{\alpha}$에 포함되지 않는 가장 작은 $j$에 대해
$M_{\alpha + 1}$이 $N_j$를 포함하도록 구성했다는 사실을 사용한다.
따라서 각 $i \in I$에 대해 $N_i$가 $M_{\alpha}$에 포함되는
$\alpha \in S$가 존재하도록 하는 충분히 큰 순서수 $S$가 존재한다.
이는 $(M_{\alpha})_{\alpha \in S}$가 $M$의 Kaplansky
d\'evissage의 성질 (1)을 만족한다는 뜻이다. 또한 족
$(M_{\alpha})_{\alpha \in S}$는 다른 정의 성질들과 위의 (5), (6)도
만족한다. 성질 (0), (2), (4), (6)은 구성으로부터 명백하다.
각 $M_{\alpha}$는 구성상 $N_i$들 가운데 일부의 직합이므로 성질
(5)가 성립한다. 그리고 (5)와
$M_{\alpha} \subset M_{\alpha + 1}$이라는 사실은 (3)을 함의한다.
\end{proof}

\noindent
따름정리로 이 절 처음에 말한 사영 가군에 대한 결과를 얻는다.

\begin{theorem}
\label{algebra-theorem-projective-direct-sum}
\begin{slogan}
모든 사영 가군은 가산 생성 사영 가군들의 직합이다.
\end{slogan}
$P$가 사영 $R$-가군이면 $P$는 가산 생성 사영 $R$-가군들의
직합이다.
\end{theorem}

\begin{proof}
가군이 사영일 필요충분조건은 자유가군의 직합인자인 것이다.
따라서 이는 정리 \ref{algebra-theorem-kaplansky-direct-sum}에서 따른다.
\end{proof}
\section{국소환 위의 사영 가군}
\label{algebra-section-projective-local-ring}

\noindent
이 절에서는 아주 멋진 결과를 증명한다. 즉, 국소환 위의 사영 가군
$M$은 자유이다(아래의 정리
\ref{algebra-theorem-projective-free-over-local-ring}). $M$이 유한이라는
가정을 더하면 이 결과는 보조정리 \ref{algebra-lemma-finite-flat-local}이다.
일반적으로 다음이 성립한다.

\begin{lemma}
\label{algebra-lemma-projective-free}
$R$을 환이라 하자. 모든 사영 $R$-가군이 자유일 필요충분조건은
모든 가산 생성 사영 $R$-가군이 자유인 것이다.
\end{lemma}

\begin{proof}
정리 \ref{algebra-theorem-projective-direct-sum}에서 곧바로 따른다.
\end{proof}

\noindent
다음은 가산 생성 가군이 자유이기 위한 판정법이다.

\begin{lemma}
\label{algebra-lemma-freeness-criteria}
$M$을 다음 성질을 갖는 가산 생성 $R$-가군이라 하자. $N'$이 유한
자유 $R$-가군이고 $M = N \oplus N'$이면, $N$의 임의의 원소는
$N$의 어떤 자유 직합인자에 포함된다. 그러면 $M$은 자유이다.
\end{lemma}

\begin{proof}
$x_1, x_2, \ldots$를 $M$의 가산 생성원 집합이라 하자. 모든 $n$에
대해 $F_1 \oplus \ldots \oplus F_n$이 $x_1, \ldots, x_n$을 포함하는
$M$의 직합인자가 되도록 $M$의 유한 자유 직합인자
$F_1, F_2, \ldots$를 귀납적으로 구성한다. 실제로 원하는 성질을
갖는 $F_1, \ldots, F_n$이 주어졌다고 하고
$$
M = F_1 \oplus \ldots \oplus F_n \oplus N
$$
이라 쓰며, $x \in N$을 $x_{n + 1}$의 상이라 하자. 보조정리의
가정에 의해 $x$를 포함하는 자유 직합인자 $F_{n + 1} \subset N$을
찾을 수 있다. 물론 $F_{n + 1}$을 $F_{n + 1}$의 유한 자유
직합인자로 바꿀 수 있으므로 귀납 단계가 끝난다. 이제
$M = \bigoplus_{i = 1}^{\infty} F_i$는 자유이다.
\end{proof}

\begin{lemma}
\label{algebra-lemma-projective-freeness-criteria}
$P$를 국소환 $R$ 위의 사영 가군이라 하자. 그러면 $P$의 임의의
원소는 $P$의 어떤 자유 직합인자에 포함된다.
\end{lemma}

\begin{proof}
$P$는 사영이므로 어떤 자유 $R$-가군 $F$의 직합인자이다. 이를
$F = P \oplus Q$라 쓰자. $x \in P$를 $P$의 자유 직합인자에
포함됨을 보이고자 하는 원소라 하자. $x$를 나타내는 데 필요한
기저 원소의 수가 최소가 되도록 $F$의 기저 $B$를 택하자. 즉, 어떤
$e_i \in B$와 $a_i \in R$에 대해
$x = \sum_{i=1}^n a_i e_i$라 하자. 그러면 어느 $a_j$도 다른
$a_i$들의 선형결합으로 나타낼 수 없다. 실제로 어떤 $b_i \in R$에
대해 $a_j = \sum_{i \neq  j} a_i b_i$라면, $i \neq j$일 때 $e_i$를
$e_i + b_ie_j$로 바꾸고 $B$의 다른 원소는 그대로 두어 $F$의 새
기저를 얻으며, 이 기저로는 $x$를 더 짧게 나타낼 수 있기 때문이다.

\medskip\noindent
$e_i$의 $P$-성분과 $Q$-성분으로의 분해를
$e_i = y_i + z_i, y_i \in P, z_i \in Q$라 하자. $t_i$가
$e_1, \ldots, e_n$ 이외의 $B$의 원소들의 선형결합이 되도록
$y_i = \sum_{j=1}^{n} b_{ij} e_j + t_i$라 쓰자. 증명을 끝내려면
행렬 $(b_{ij})$가 가역임을 보이면 충분하다. 실제로 그러면
$i=1, \ldots, n$에 대해 $e_i \mapsto y_i$로 보내고
$B \setminus \{e_1, \ldots, e_n\}$을 고정하는 사상 $F \to F$는
동형이다. 따라서 $y_1, \ldots, y_n$과
$B \setminus \{e_1, \ldots, e_n\}$은 함께 $F$의 기저를 이룬다.
이때 $y_1, \ldots, y_n$으로 생성되는 부분가군 $N$은 $P$의 자유
부분가군이다. $N \subset P$이고 $N$과 $P$가 모두 $F$의
직합인자이므로 $N$은 $P$의 직합인자이다. 또한 $x \in P$이므로
$x = \sum_{i=1}^n a_i e_i = \sum_{i=1}^n a_i y_i$이고, 따라서
$x \in N$이다.

\medskip\noindent
이제 $(b_{ij})$가 가역임을 증명하자.
$y_i = \sum_{j=1}^{n} b_{ij} e_j + t_i$를
$\sum_{i=1}^n a_i e_i = \sum_{i=1}^n a_i y_i$에 대입하고 $e_j$의
계수를 비교하면 $a_j = \sum_{i=1}^n a_i b_{ij}$를 얻는다. 위에서
보았듯이 $B$의 선택에 의해 어느 $a_j$도 다른 $a_i$들의 선형결합으로
쓸 수 없다. 따라서 $i \neq j$이면 $b_{ij}$는 비가역원이고,
모든 $i$에 대해 $1-b_{ii}$는 비가역원이다. 특히 $b_{ii}$는
가역원이다. 그런데 국소환 위에서 대각 원소가 가역원이고 나머지
원소가 비가역원인 행렬은 그 행렬식이 가역원이므로 가역이다.
\end{proof}

\begin{theorem}
\label{algebra-theorem-projective-free-over-local-ring}
\begin{slogan}
국소환 위의 사영 가군은 자유이다.
\end{slogan}
$P$가 국소환 $R$ 위의 사영 가군이면 $P$는 자유이다.
\end{theorem}

\begin{proof}
보조정리 \ref{algebra-lemma-projective-free}, \ref{algebra-lemma-freeness-criteria},
\ref{algebra-lemma-projective-freeness-criteria}에서 따른다.
\end{proof}



\section{Mittag-Leffler 계}
\label{algebra-section-mittag-leffler}

\noindent
이 절의 목적은 Mittag-Leffler 계를 정의하고 이 개념이 왜 유용한지
설명하는 것이다.

\medskip\noindent
이하에서 $I$는 유향집합이다. 범주론, 정의
\ref{categories-definition-directed-set}를 보라.
$(A_i, \varphi_{ji}: A_j \to A_i)$를 $I$로 지표화된 집합 또는 가군의
역계라 하자. 범주론, 정의
\ref{categories-definition-directed-system}를 보라. $I$가 유향이라고
가정했으므로 이는 유향 역계이다(범주론, 정의
\ref{categories-definition-directed-system}). 각 $i \in I$에 대해
$j \geq i$일 때의 상들 $\varphi_{ji}(A_j) \subset A_i$는 $A_i$의
부분집합(또는 부분가군)으로 이루어진 감소 유향족을 이룬다.
$A'_i = \bigcap_{j \geq i} \varphi_{ji}(A_j)$라 하자. 그러면
$j \geq i$일 때 $\varphi_{ji}(A'_j) \subset A'_i$이므로 제한하여
유향 역계 $(A'_i, \varphi_{ji}|_{A'_j})$를 얻는다. 집합 또는 가군의
범주에서 역계의 극한을 구성하는 방식으로부터
$\lim A_i = \lim A'_i$이다. $(A_i, \varphi_{ji})$에 대한
Mittag-Leffler 조건은 어떤 $j \geq i$에 대해 $A'_i$가
$\varphi_{ji}(A_j)$와 같은 것이다(따라서 모든 $k \geq j$에 대해
$\varphi_{ki}(A_k)$와도 같다).

\begin{definition}
\label{algebra-definition-ML-system}
$(A_i, \varphi_{ji})$를 $I$ 위의 집합들의 유향 역계라 하자. 각
$i \in I$에 대해 $j \geq i$일 때의 족
$\varphi_{ji}(A_j) \subset A_i$가 안정화되면
$(A_i, \varphi_{ji})$를 {\it Mittag-Leffler}라고 한다. 명시적으로,
이는 각 $i \in I$에 대해 어떤 $j \geq i$가 존재하여 $k \geq j$이면
$\varphi_{ki}(A_k) = \varphi_{ji}( A_j)$가 된다는 뜻이다.
$(A_i, \varphi_{ji})$가 환 $R$ 위의 가군들의 유향 역계일 때, 그
밑바탕 집합들의 역계가 Mittag-Leffler이면 이 가군의 역계도
Mittag-Leffler라고 한다.
\end{definition}

\begin{example}
\label{algebra-example-ML-surjective-maps}
$(A_i, \varphi_{ji})$가 집합 또는 가군의 유향 역계이고 사상
$\varphi_{ji}$가 전사이면, 이 계가 Mittag-Leffler임은 명백하다.
역으로 $(A_i, \varphi_{ji})$가 Mittag-Leffler라고 하자.
$j \geq i$일 때 $\varphi_{ji}(A_j)$의 안정된 상을
$A'_i \subset A_i$라 하자. 그러면 $j \geq i$일 때
$\varphi_{ji}|_{A'_j}: A'_j \to A'_i$는 전사이고
$\lim A_i = \lim A'_i$이다. 따라서 Mittag-Leffler 계
$(A_i, \varphi_{ji})$의 극한은 전사 사상들로 이루어진 $I$ 위의 유향
역계의 극한으로도 쓸 수 있다.
\end{example}

\begin{lemma}
\label{algebra-lemma-ML-limit-nonempty}
$(A_i, \varphi_{ji})$를 $I$ 위의 유향 역계라 하자. $I$가 가산이라고
가정하자. $(A_i, \varphi_{ji})$가 Mittag-Leffler이고 $A_i$들이
공집합이 아니면, $\lim A_i$도 공집합이 아니다.
\end{lemma}

\begin{proof}
$i_1, i_2, i_3, \ldots$를 $I$의 원소들의 열거라 하자. $n = 1, 2,
3, \ldots$에 대해 $j_n \in I$인 원소의 열을 다음 조건으로
귀납적으로 정의한다. $j_1 = i_1$이고, $j_n \geq i_n$이며 $m < n$이면
$j_n \geq j_m$이다. 그러면 열 $j_n$은 증가하며 $I$의 공종 부분집합을
이룬다. 따라서 $I =\{1, 2, 3, \ldots \}$라고 가정해도 된다. 이제
예 \ref{algebra-example-ML-surjective-maps}에 의해 양의 정수로 지표화되고
전사 사상들로 이루어진 공집합이 아닌 집합들의 역계의 극한이
공집합이 아님을 보이는 문제로 환원된다. 이는 선택공리에서 따른다.
\end{proof}

\noindent
Mittag-Leffler 조건은 다음 완전성 성질 때문에 우리에게 중요하다.

\begin{lemma}
\label{algebra-lemma-ML-exact-sequence}
$$
0 \to A_i \xrightarrow{f_i} B_i \xrightarrow{g_i} C_i \to 0
$$
을 $I$ 위의 아벨군 유향 역계들의 완전열이라 하자. $I$가 가산이라고
가정하자. $(A_i)$가 Mittag-Leffler이면
$$
0 \to \lim A_i \to \lim B_i \to \lim C_i\to 0
$$
은 완전이다.
\end{lemma}

\begin{proof}
유향 역계의 극한을 취하는 것은 좌완전이므로
$\lim B_i \to \lim C_i$의 전사성만 증명하면 된다.
$(c_i) \in \lim C_i$라 하자. 각 $i \in I$에 대해
$E_i = g_i^{-1}(c_i)$라 하자. $g_i: B_i \to C_i$가 전사이므로 이는
공집합이 아니다. $(B_i)$의 사상계 $\varphi_{ji}: B_j \to B_i$는
$E_j \to E_i$인 사상들로 제한되며, 이로써 $(E_i)$는 공집합이 아닌
집합들의 역계가 된다. $(E_i)$가 Mittag-Leffler임을 보이면 충분하다.
실제로 그러면 보조정리 \ref{algebra-lemma-ML-limit-nonempty}에 의해
$\lim E_i$가 공집합이 아니고, $\lim E_i$의 임의의 원소는
$(c_i)$로 가는 $\lim B_i$의 원소를 준다.

\medskip\noindent
단사 $f_i: A_i \to B_i$를 통해 $A_i$를 $B_i$의 부분집합으로
간주하자. $(A_i)$가 Mittag-Leffler이므로 $i \in I$이면 어떤
$j \geq i$가 존재하여 $k \geq j$에 대해
$\varphi_{ki}(A_k) = \varphi_{ji}(A_j)$이다. 이때 또한
$k \geq j$에 대해 $\varphi_{ki}(E_k) = \varphi_{ji}(E_j)$라고
주장한다. $k \geq j$일 때
$\varphi_{ki}(E_k) \subset \varphi_{ji}(E_j)$는 항상 성립한다.
역포함을 위해 $e_j \in E_j$라 하고,
$\varphi_{ki}(x_k) = \varphi_{ji}(e_j)$를 만족하는
$x_k \in E_k$를 찾아야 한다. $e'_k \in E_k$를 임의로 택하고
$e'_j = \varphi_{kj}(e'_k)$라 하자. 그러면
$g_j(e_j - e'_j) = c_j - c_j = 0$이므로
$e_j - e'_j = a_j \in A_j$이다.
$\varphi_{ki}(A_k) = \varphi_{ji}(A_j)$이므로
$\varphi_{ki}(a_k) = \varphi_{ji}(a_j)$인 $a_k \in A_k$가 존재한다.
따라서
$$
\varphi_{ki}(e'_k + a_k) = \varphi_{ji}(e'_j) + \varphi_{ji}(a_j) =
\varphi_{ji}(e_j),
$$
이므로 $x_k = e'_k + a_k$로 택할 수 있다.
\end{proof}




\section{역계}
\label{algebra-section-inverse-systems}

\noindent
많은 논문에서(그리고 이 절에서) {\it 역계}라는 말은 부분순서집합
$(\mathbf{N}, \geq)$ 위의 역계를 가리키는 데 쓰인다. 이 절에서는
그러한 계를 간단히 논한다. 이 내용은 호몰로지, 절
\ref{homology-section-inverse-systems}에서 더 폭넓게 다룰 것이다.
환 $R$과 표시된 사상들을 갖는 $R$-가군의 열
$$
M_1 \xleftarrow{\varphi_2} M_2 \xleftarrow{\varphi_3} M_3 \leftarrow \ldots
$$
이 주어졌다고 하자. 연속한 사상들을 합성하면 $i \geq i'$일 때마다
사상 $\varphi_{ii'} : M_i \to M_{i'}$를 얻고, 이들은
$i \geq i' \geq i''$일 때
$\varphi_{ii''} = \varphi_{i'i''} \circ \varphi_{i i'}$를 만족한다.
역으로 사상계 $\varphi_{ii'}$가 주어지면
$\varphi_i = \varphi_{i(i-1)}$로 놓아 위에 표시한 사상들을 복원할
수 있다. 이 경우
$$
\lim M_i
=
\{(x_i) \in \prod M_i \mid \varphi_i(x_i) = x_{i - 1}, \ i = 2, 3, \ldots\}
$$
이다. 범주론, 절 \ref{categories-section-limit-sets}와 비교하라.
호몰로지, 절 \ref{homology-section-inverse-systems}에서 설명하듯이,
이는 실제로 범주론, 절 \ref{categories-section-limits}에서 정의한
$R$-가군의 범주에서의 극한이다.

\begin{lemma}
\label{algebra-lemma-Mittag-Leffler}
$R$을 환이라 하자. $i \geq 1$에 대해 $R$-가군들의 짧은 완전열
$0 \to K_i \to L_i \to M_i \to 0$들이 주어지고, 이들이 다음 짧은
완전열들의 사상들에 들어맞는다고 하자.
$$
\xymatrix{
0 \ar[r] &
K_i \ar[r] &
L_i \ar[r] &
M_i \ar[r] &
0 \\
0 \ar[r] &
K_{i + 1} \ar[r] \ar[u] &
L_{i + 1} \ar[r] \ar[u] &
M_{i + 1} \ar[r] \ar[u] &
0}
$$
모든 $i$에 대해 어떤 $c = c(i) \geq i$가 존재하여 모든
$j \geq c$에 대해
$\Im(K_c \to K_i) = \Im(K_j \to K_i)$이면, 열
$$
0 \to \lim K_i \to \lim L_i \to \lim M_i \to 0
$$
은 완전이다.
\end{lemma}

\begin{proof}
이는 더 일반적인 보조정리 \ref{algebra-lemma-ML-exact-sequence}의 특수한
경우이다.
\end{proof}







\section{Mittag-Leffler 가군}
\label{algebra-section-mittag-leffler-modules}

\noindent
Mittag-Leffler 가군은 (매우 거칠게 말해) 그 쌍대가 Mittag-Leffler
계인 유향극한으로 쓸 수 있는 가군이다. 정확한 정의를 내리려면
약간의 준비가 필요하다.

\begin{definition}
\label{algebra-definition-ML-inductive-system}
$(M_i, f_{ij})$를 $R$-가군들의 유향계라 하자. 각 $M_i$가 유한 표시
$R$-가군이고 모든 $R$-가군 $N$에 대해 역계
$$
(\Hom_R(M_i, N), \Hom_R(f_{ij}, N))
$$
가 Mittag-Leffler이면 $(M_i, f_{ij})$를 가군들의
{\it Mittag-Leffler 유향계}라고 한다.
\end{definition}

\noindent
가군들의 Mittag-Leffler 유향계의 여극한이 되는 $R$-가군들을
특징짓고자 한다.

\begin{definition}
\label{algebra-definition-domination}
$f: M \to N$과 $g: M \to M'$을 $R$-가군들의 사상이라 하자. 임의의
$R$-가군 $Q$에 대해
$\Ker(f \otimes_R \text{id}_Q) \subset \Ker(g \otimes_R \text{id}_Q)$
이면 $g$가 $f$를 {\it 지배한다}고 한다.
\end{definition}

\noindent
이 조건은 유한 표시 가군들에 대해서만 확인해도 충분하다.

\begin{lemma}
\label{algebra-lemma-domination-fp}
$f: M \to N$과 $g: M \to M'$을 $R$-가군들의 사상이라 하자. 임의의
유한 표시 $R$-가군 $Q$에 대해
$\Ker(f \otimes_R \text{id}_Q) \subset \Ker(g \otimes_R
\text{id}_Q)$일 필요충분조건은 $g$가 $f$를 지배하는 것이다.
\end{lemma}

\begin{proof}
모든 유한 표시 가군 $Q$에 대해
$\Ker(f \otimes_R \text{id}_Q) \subset \Ker(g \otimes_R
\text{id}_Q)$라고 가정하자. $Q$가 임의의 가군이면 이를 유한 표시
가군 $Q_i$들의 유향계의 여극한
$Q = \colim_{i \in I} Q_i$로 쓰자. 그러면 모든 $i$에 대해
$\Ker(f \otimes_R \text{id}_{Q_i}) \subset \Ker(g \otimes_R
\text{id}_{Q_i})$이다. 유향 여극한을 취하는 것은 완전하고 텐서곱과
가환하므로
$\Ker(f \otimes_R \text{id}_Q) \subset \Ker(g
\otimes_R \text{id}_Q)$가 따른다.
\end{proof}

\begin{lemma}
\label{algebra-lemma-domination-universally-injective}
$f : M \to N$과 $g : M \to M'$을 $R$-가군들의 사상이라 하자.
$f$와 $g$의 밀어내기
$$
\xymatrix{
M  \ar[r]_f \ar[d]_g & N \ar[d]^{g'} \\
M' \ar[r]^{f'} & N'
}
$$
를 생각하자. $g$가 $f$를 지배할 필요충분조건은 $f'$이 보편
단사인 것이다.
\end{lemma}

\begin{proof}
$N'$은 $M' \oplus N$을 $x \in M$에 대한 원소
$(g(x), -f(x))$들로 이루어진 부분가군으로 나눈 몫임을 상기하라.
$N'$의 구성으로부터 짧은 완전열
$$
0 \to \Ker(f) \cap \Ker(g) \to \Ker(f) \to \Ker(f')
\to 0.
$$
을 얻는다. 텐서곱을 취하는 것은 밀어내기를 취하는 것과 가환하므로,
모든 $R$-가군 $Q$에 대해 다음과 같은 짧은 완전열이 있다.
$$
0 \to \Ker(f \otimes \text{id}_Q ) \cap \Ker(g \otimes
\text{id}_Q) \to \Ker(f \otimes \text{id}_Q)
\to \Ker(f' \otimes \text{id}_Q) \to 0
$$
따라서 $f'$이 보편 단사일 필요충분조건은 모든 $Q$에 대해
$\Ker(f \otimes \text{id}_Q ) \subset \Ker(g \otimes
\text{id}_Q)$인 것이고, 이는 다시 $g$가 $f$를 지배하는 것과
동치이다.
\end{proof}

\noindent
다음 보조정리가 보여 주듯이 위의 지배 정의는 때때로 사상의
통상적인 지배 개념과 관련된다.

\begin{lemma}
\label{algebra-lemma-domination}
$f: M \to N$과 $g: M \to M'$을 $R$-가군들의 사상이라 하자.
$\Coker(f)$가 유한 표시라고 가정하자. $g$가 $f$를 지배할
필요충분조건은 $g$가 $f$를 통하여 인수분해되는 것이다. 즉,
$g = h \circ f$인 가군 사상 $h: N \to M'$이 존재하는 것이다.
\end{lemma}

\begin{proof}
보조정리 \ref{algebra-lemma-domination-universally-injective}의 명제에서와 같이
$f$와 $g$의 밀어내기를 생각하자. 밀어내기의 구성으로부터
$\Coker(f') = \Coker(f)$이고, 따라서 $\Coker(f')$는 유한 표시이다.
그러면 보조정리 \ref{algebra-lemma-universally-exact-split}에 의해 $f'$이
보편 단사일 필요충분조건은
$$
0 \to M' \xrightarrow{f'} N' \to \Coker(f') \to 0
$$
이 분할되는 것이다. 이는 $h' \circ f' = \text{id}_{M'}$인 사상
$h' : N' \to M'$이 존재하는 것과 동치이다. 밀어내기의 보편성에
의해 그러한 $h'$의 존재는 $g$가 $f$를 통하여 인수분해되는 것과
동치이다.
\end{proof}


\begin{proposition}
\label{algebra-proposition-ML-characterization}
$M$을 $R$-가군이라 하자. $(M_i, f_{ij})$를 $I$로 지표화된 유한 표시
$R$-가군들의 유향계라 하고 $M = \colim M_i$라고 하자.
$f_i: M_i \to M$을 표준 사상이라 하자. 다음은 서로 동치이다.
\begin{enumerate}
\item 모든 유한 표시 $R$-가군 $P$와 가군 사상 $f: P \to M$에 대해,
어떤 유한 표시 $R$-가군 $Q$와 가군 사상 $g: P \to Q$가 존재하여
$g$와 $f$가 서로를 지배한다. 즉, 모든 $R$-가군 $N$에 대해
$\Ker(f \otimes_R \text{id}_N) = \Ker(g \otimes_R \text{id}_N)$이다.
\item 각 $i \in I$에 대해 어떤 $j \geq i$가 존재하여
$f_{ij}: M_i \to M_j$가 $f_i: M_i \to M$을 지배한다.
\item 각 $i \in I$에 대해 어떤 $j \geq i$가 존재하여 모든
$k \geq i$에 대해 $f_{ij}: M_i \to M_j$가
$f_{ik}: M_i \to M_k$를 통하여 인수분해된다.
\item 모든 $R$-가군 $N$에 대해 역계
$(\Hom_R(M_i, N), \Hom_R(f_{ij}, N))$가 Mittag-Leffler이다.
\item $N = \prod_{s \in I} M_s$에 대해 역계
$(\Hom_R(M_i, N), \Hom_R(f_{ij}, N))$가 Mittag-Leffler이다.
\end{enumerate}
\end{proposition}

\begin{proof}
먼저 (1)과 (2)의 동치를 증명하자. (1)이 성립한다고 하고
$i \in I$라 하자. 사상 $f_i: M_i \to M$에 대응하여 (1)에서와 같은
$g: M_i \to Q$를 택할 수 있다. $M_i$와 $Q$가 유한 표시이므로
$\Coker(g)$도 유한 표시이다. 따라서 보조정리
\ref{algebra-lemma-domination}에 의해 $f_i : M_i \to M$은
$g: M_i \to Q$를 통하여 인수분해된다. 즉, 어떤 $h: Q \to M$에
대해 $f_i = h \circ g$이다. 이제 $Q$가 유한 표시이므로 어떤
$j \geq i$에 대해 $h$는 $M_j \to M$을 통하여 인수분해된다. 즉,
어떤 $h': Q \to M_j$에 대해 $h = f_j \circ h'$이다. 모두 합치면
다음 가환 그림을 얻는다.
$$
\xymatrix{
  &  M  & \\
M_i \ar[dr]_g \ar[ur]^{f_i} \ar[rr]^{f_{ij}} &
& M_j \ar[ul]_{f_j} \\
  & Q  \ar[ur]_{h'} &
}
$$
따라서 $f_{ij}$는 $g$를 지배한다. 그런데 $g$는 $f_i$를 지배하므로
$f_{ij}$는 $f_i$를 지배한다.

\medskip\noindent
역으로 (2)가 성립한다고 하자. $P$를 유한 표시 가군이라 하고
$f: P \to M$을 가군 사상이라 하자. 그러면 어떤 $i \in I$에 대해
$f$는 $f_i: M_i \to M$을 통하여 인수분해된다. 즉, 어떤
$g': P \to M_i$에 대해 $f = f_i \circ g'$이다. (2)에 의해
$f_{ij}$가 $f_i$를 지배하도록 하는 $j \geq i$를 택하자. 다음
가환 그림이 있다.
$$
\xymatrix{
P \ar[d]_{g'} \ar[r]^{f}           & M  \\
M_i \ar[ur]^{f_i} \ar[r]_{f_{ij}} & M_j \ar[u]_{f_j}
}
$$
그림과 $f_{ij}$가 $f_i$를 지배한다는 사실로부터 $f$와
$f_{ij} \circ g'$가 서로를 지배함을 안다. 따라서
$g = f_{ij} \circ g' : P \to M_j$로 택하면 된다.

\medskip\noindent
다음으로 (2)와 (3)이 동치임을 증명하자. $i \in I$라 하자.
$f_i$는 $f_{ik}$를 통하여 인수분해되므로, $k \geq i$이면 $f_i$가
$f_{ik}$를 지배한다. (2)가 성립하면 $f_{ij}$가 $f_i$를 지배하도록
하는 $j \geq i$를 택하자. 지배 관계는 추이적이므로, 모든
$k \geq i$에 대해 $f_{ij}$는 $f_{ik}$를 지배한다. 모든 $M_i$가
유한 표시이므로 $k \geq i$일 때 $\Coker(f_{ik})$는 유한 표시이다.
보조정리 \ref{algebra-lemma-domination}에 의해 모든 $k \geq i$에 대해
$f_{ij}$는 $f_{ik}$를 통하여 인수분해된다. 따라서 (2)는 (3)을
함의한다. 한편 (3)이 성립하면 임의의 $R$-가군 $N$에 대해
$k \geq i$일 때 $f_{ij} \otimes_R \text{id}_N$은
$f_{ik} \otimes_R \text{id}_N$을 통하여 인수분해된다. 따라서
$k \geq i$일 때
$\Ker(f_{ik} \otimes_R \text{id}_N) \subset \Ker(f_{ij} \otimes_R
\text{id}_N)$이다. 그런데
$\Ker(f_i \otimes_R \text{id}_N: M_i \otimes_R N
\to M \otimes_R N)$은 $k \geq i$일 때의
$\Ker(f_{ik} \otimes_R \text{id}_N)$들의 합집합이다. 그러므로 임의의
$R$-가군 $N$에 대해
$\Ker(f_i \otimes_R \text{id}_N) \subset \Ker(f_{ij} \otimes_R
\text{id}_N)$이고, 정의에 의해 이는 $f_{ij}$가 $f_i$를 지배한다는
뜻이다.

\medskip\noindent
(3)이 (4)를, (4)가 (5)를 함의함은 자명하다. (5)가 (3)을
함의함을 보이자. $N = \prod_{s \in I} M_s$라 하자. (5)가 성립하면
주어진 $i \in I$에 대해 어떤 $j \geq i$를 택하여
$$
\Im( \Hom(M_j, N) \to  \Hom(M_i, N)) =
\Im( \Hom(M_k, N) \to  \Hom(M_i, N))
$$
이 모든 $k \geq j$에 대해 성립하도록 할 수 있다. $s \in I$에 대한
곱을 $\Hom$ 밖으로 옮기고 곱의 각 성분에서 사상을 살펴보면, 이는
$$
\Im( \Hom(M_j, M_s) \to  \Hom(M_i, M_s)) =
\Im( \Hom(M_k, M_s) \to   \Hom(M_i, M_s))
$$
이 모든 $k \geq j$와 $s \in I$에 대해 성립한다는 뜻이다.
$s = j$로 두면
$$
\Im( \Hom(M_j, M_j) \to  \Hom(M_i, M_j)) =
\Im( \Hom(M_k, M_j) \to   \Hom(M_i, M_j))
$$
이 모든 $k \geq j$에 대해 성립한다. $f_{ij}$는
$\Hom(M_j, M_j) \to  \Hom(M_i, M_j)$ 아래에서
$\text{id} \in  \Hom(M_j, M_j)$의 상이므로, 이는 임의의
$k \geq j$에 대해 $f_{ij} = h \circ f_{ik}$인
$h \in \Hom(M_k, M_j)$가 존재함을 보인다. $j \geq k$이면
$h = f_{kj}$로 택할 수 있다. 따라서 (3)이 성립한다.
\end{proof}

\begin{definition}
\label{algebra-definition-mittag-leffler-module}
$M$을 $R$-가군이라 하자. 명제
\ref{algebra-proposition-ML-characterization}의 동치 조건들이 성립하면
$M$을 {\it Mittag-Leffler}라고 한다.
\end{definition}

\noindent
특히 유한 표시 가군은 Mittag-Leffler이다.

\begin{remark}
\label{algebra-remark-flat-ML}
$M$을 평탄 $R$-가군이라 하자. Lazard의 정리(정리
\ref{algebra-theorem-lazard})에 의해 유한 자유 $R$-가군 $M_i$들로
이루어진 유향계 $(M_i, f_{ij})$의 여극한으로
$M = \colim M_i$를 쓸 수 있다. $M$이 Mittag-Leffler임을 보이려면
쌍대의 역계 $(\Hom_R(M_i, R), \Hom_R(f_{ij}, R))$가
Mittag-Leffler임을 보이면 충분하다. 이는 명제
\ref{algebra-proposition-ML-characterization}의 판정법 (4)와, 유한 자유
$R$-가군 $F$ 및 임의의 $R$-가군 $N$에 대해 함자적 동형
$\Hom_R(F, R) \otimes_R N \cong \Hom_R(F, N)$이 존재한다는 사실에서
따른다.
\end{remark}

\begin{lemma}
\label{algebra-lemma-tensor-ML-modules}
$R$이 환이고 $M$, $N$이 $R$ 위의 Mittag-Leffler 가군이면,
$M \otimes_R N$도 Mittag-Leffler 가군이다.
\end{lemma}

\begin{proof}
$M = \colim_{i \in I} M_i$와 $N = \colim_{j \in J} N_j$를 유한 표시
$R$-가군들의 유향 여극한으로 쓰자. 전이 사상을
$f_{ii'} : M_i \to M_{i'}$와 $g_{jj'} : N_j \to N_{j'}$로 표기하자.
그러면 $M_i \otimes_R N_j$는 유한 표시 $R$-가군이고(보조정리
\ref{algebra-lemma-tensor-finiteness}를 보라),
% P01-E0452
$M \otimes_R N = \colim_{(i, j) \in I \times J} M_i \otimes_R M_j$이다.
$(i, j) \in I \times J$를 택하자. Mittag-Leffler 가군의 정의에 의해
두 계 모두에 명제 \ref{algebra-proposition-ML-characterization}의 (3)을
적용할 수 있다. 다시 말해, 어떤 $i' \geq i$와 $j' \geq j$가
존재하여 $i'' \geq i$와 $j'' \geq j$를 어떻게 택하든 사상
% P01-E0453
$a : M_{i''} \to M_{i'}$와 $b : M_{j''} \to M_{j'}$가 존재하고
$f_{ii'} = a \circ f_{ii''}$와 $g_{jj'} = b \circ g_{jj''}$를 만족한다.
그러면
$a \otimes b : M_{i''} \otimes_R N_{j''} \to M_{i'} \otimes_R N_{j'}$가
계 $(M_i \otimes_R N_j, f_{ii'} \otimes g_{jj'})$에 대해 같은 역할을
함은 명백하다. 따라서 명제
\ref{algebra-proposition-ML-characterization}의 특징짓기 (3)에 의해
$M \otimes_R N$이 Mittag-Leffler임을 결론내린다.
\end{proof}

\begin{lemma}
\label{algebra-lemma-ML-also}
$R$을 환이라 하고 $M$을 $R$-가군이라 하자. $M$이
Mittag-Leffler일 필요충분조건은 모든 유한 자유 $R$-가군 $F$와
가군 사상 $f: F \to M$에 대해 어떤 유한 표시 $R$-가군 $Q$와 가군
사상 $g : F \to Q$가 존재하여 $g$와 $f$가 서로를 지배하는 것이다.
즉, 모든 $R$-가군 $N$에 대해
$\Ker(f \otimes_R \text{id}_N) = \Ker(g \otimes_R \text{id}_N)$이다.
\end{lemma}

\begin{proof}
% P01-E0454
이 조건은 명제 \ref{algebra-proposition-ML-characterization}의 조건 (1)보다
명백히 약하므로 Mittag-Leffler 가군은 이 조건을 만족한다.
역으로 $M$이 이 조건을 만족하고, $f : P \to M$이 유한 표시
$R$-가군 $P$에서 $M$으로 가는 $R$-가군 사상이라고 하자. $F$가
유한 자유 $R$-가군인 전사 $F \to P$를 택하자. 가정에 의해 유한 표시
$R$-가군 $Q$로 가는 사상 $F \to Q$를 찾아 $F \to Q$와
$F \to M$이 서로를 지배하도록 할 수 있다. 특히 $F \to Q$의 핵은
$F \to P$의 핵을 포함하므로, $F \to Q$가 합성
$F \to P \to Q$와 같도록 하는 $R$-가군 사상 $g : P \to Q$를 얻는다.
$N$을 임의의 $R$-가군이라 하고 다음 가환 그림을 생각하자.
$$
\xymatrix{
F \otimes_R N \ar[d] \ar[r] & Q \otimes_R N \\
P \otimes_R N \ar[ru] \ar[r] & M \otimes_R N
}
$$
가정에 의해 $F \otimes_R N \to Q \otimes_R N$과
$F \otimes_R N \to M \otimes_R N$의 핵은 같다. 따라서
$F \otimes_R N \to P \otimes_R N$이 전사이므로
$P \otimes_R N \to Q \otimes_R N$과
$P \otimes_R N \to M \otimes_R N$의 핵도 같다.
\end{proof}

\begin{lemma}
\label{algebra-lemma-restrict-ML-modules}
$R \to S$를 유한이며 유한 표시인 환 사상이라 하자. $M$을
$S$-가군이라 하자. $M$이 $S$ 위의 Mittag-Leffler 가군이면 $M$은
$R$ 위의 Mittag-Leffler 가군이다.
\end{lemma}

\begin{proof}
$M$이 $S$ 위의 Mittag-Leffler 가군이라고 가정하자. 이를 유한 표시
$S$-가군 $M_i$들의 유향 여극한 $M = \colim M_i$로 쓰자. $M$이
$S$ 위에서 Mittag-Leffler이므로 각 $i$에 대해 어떤 지표
$j \geq i$가 존재하여 모든 $k \geq j$에 대해
$f_{ij} = h \circ f_{ik}$인 인수분해가 있다(여기서 $h$는 $i$와
$j$의 선택 및 $k$에 의존한다). 보조정리
\ref{algebra-lemma-finite-finitely-presented-extension}에 의해 $M_i$들은
$R$-가군으로서도 유한 표시임에 유의하라. 또한 사상
$f_{ij}, f_{ik}, h$는 모두 $R$-가군의 사상이다. 따라서
$(M_i, f_{ij})$를 $R$-가군의 계로 보아도 같은 조건을 만족함을 안다.
그러므로 $M$은 $R$-가군으로서 Mittag-Leffler이다.
\end{proof}

\begin{lemma}
\label{algebra-lemma-mod-ideal-ML-modules}
$R$을 환이라 하자. 어떤 유한 생성 아이디얼 $I$에 대해
$S = R/I$라 하자. $M$을 $S$-가군이라 하자. $M$이 $R$ 위의
Mittag-Leffler 가군일 필요충분조건은 $M$이 $S$ 위의
Mittag-Leffler 가군인 것이다.
\end{lemma}

\begin{proof}
한쪽 함의는 보조정리 \ref{algebra-lemma-restrict-ML-modules}에서 따른다.
다른 쪽을 증명하기 위해 $M$이 $R$-가군으로서 Mittag-Leffler라고
가정하자. 이를 유한 표시 $S$-가군들의 유향 여극한
$M = \colim M_i$로 쓰자. $I$가 유한 생성이므로 환 $S$는
$R$-대수로서 유한이며 유한 표시이고, 따라서 $M_i$들은
$R$-가군으로서 유한 표시이다. 보조정리
\ref{algebra-lemma-finite-finitely-presented-extension}를 보라. 다음으로
$N$을 임의의 $S$-가군이라 하자. $R \to S$가 전사이므로 각 $i$에
대해 $\Hom_R(M_i, N) = \Hom_S(M_i, N)$임에 유의하라. 따라서 역계
$(\Hom_R(M_i, N))_i$가 Mittag-Leffler 조건을 만족한다는 사실은
$(\Hom_S(M_i, N))_i$도 Mittag-Leffler 조건을 만족함을 함의한다.
그러므로 정의에 의해 $M$은 $S$ 위에서 Mittag-Leffler이다.
\end{proof}

\begin{remark}
\label{algebra-remark-go-up-ML-modules}
$R \to S$를 유한이며 유한 표시인 환 사상이라 하자. $M$을
$R$-가군으로서 Mittag-Leffler인 $S$-가군이라 하자. 일반적으로
$M$이 $S$-가군으로서 Mittag-Leffler인 것은 아니다. 예를 들어
$S$가 $R$ 위의 쌍대수 환, 즉 $\epsilon^2 = 0$인
$S = R \oplus R\epsilon$이라고 하자. 그러면 $S$-가군은 제곱영
$R$-선형 자기준동형 $\epsilon : M \to M$이 주어진 $R$-가군 $M$으로
이루어진다. $M_0$가 Mittag-Leffler가 아닌 $R$-가군이라고 하자.
$F_1$과 $F_0$가 자유 $R$-가군인 표시
$F_1 \xrightarrow{u} F_0 \to M_0 \to 0$을 택하자.
$M = F_1 \oplus F_0$로 놓고
$$
\epsilon =
\left(
\begin{matrix}
0 & 0 \\
u & 0
\end{matrix}
\right) : M \longrightarrow M.
$$
로 정의하자. 그러면 $M/\epsilon M \cong F_1 \oplus M_0$는
$R = S/\epsilon S$ 위에서 Mittag-Leffler가 아니고, 따라서 $S$
위에서도 Mittag-Leffler가 아니다(보조정리
\ref{algebra-lemma-mod-ideal-ML-modules}를 보라). 한편
$M/\epsilon M = M \otimes_S S/\epsilon S$이고, $M$이
Mittag-Leffler라면 보조정리 \ref{algebra-lemma-tensor-ML-modules}에 의해
이는 $S$ 위에서 Mittag-Leffler일 것이다.
\end{remark}


\section{직접곱과 텐서곱의 교환}
\label{algebra-section-products-tensor}

\noindent
$M$을 $R$-가군이라 하고 $(Q_{\alpha})_{\alpha \in A}$를
$R$-가군들의 족이라 하자. 그러면 순수 텐서에서
$x \otimes (q_{\alpha}) \mapsto (x \otimes q_{\alpha})$로 주어지는
표준 사상
$M \otimes_R \left( \prod_{\alpha \in A} Q_{\alpha} \right) \to
\prod_{\alpha \in A} ( M \otimes_R Q_{\alpha})$가 있다. 다음 예가
보여 주듯이 이 사상은 반드시 단사이거나 전사인 것은 아니다.

\begin{example}
\label{algebra-example-Q-not-ML}
$R = \mathbf{Z}$, $M = \mathbf{Q}$로 두고 $n \geq 1$에 대해 족
$Q_n = \mathbf{Z}/n$을 생각하자. 그러면
$\prod_n (M \otimes Q_n) = 0$이다. 그러나 단사
$\mathbf{Z} \to \prod_n Q_n$에 $M$을 텐서하여 얻는 단사
$\mathbf{Q} \to M \otimes (\prod_n Q_n)$가 있으므로
$M \otimes (\prod_n Q_n)$은 영이 아니다. 따라서
$M \otimes (\prod_n Q_n) \to \prod_n (M \otimes Q_n)$은 단사가 아니다.

\medskip\noindent
한편 다시 $R = \mathbf{Z}$, $M = \mathbf{Q}$로 두고 $n \geq 1$에 대해
$Q_n = \mathbf{Z}$라 하자. 사상
$M \otimes (\prod_n Q_n) \to \prod_n (M \otimes Q_n) = \prod_n M$의 상은
$a_n \in \mathbf{Z}$이고 $m$이 어떤 영이 아닌 정수일 때의 꼴
$(a_n/m)_{n \geq 1}$인 열들로 정확히 이루어진다. 따라서 이 사상은
전사가 아니다.
\end{example}

\noindent
아래에서 $(Q_{\alpha})_{\alpha \in A}$의 모든 선택에 대해 사상
$M \otimes_R \left( \prod_{\alpha} Q_{\alpha} \right) \to
\prod_{\alpha} (M \otimes_R Q_{\alpha})$가 각각 전사, 전단사 또는
단사가 되기 위해 $M$에 필요한 정확한 조건을 결정한다. 단사인
경우의 가군들이 정확히 Mittag-Leffler 가군으로 드러나므로 이는
중요하다(명제 \ref{algebra-proposition-ML-tensor}). 이하에서 $M$이
$R$-가군이고 $A$가 집합이면 $M^A$는 곱
$\prod_{\alpha \in A} M$을 뜻한다.

\begin{proposition}
\label{algebra-proposition-fg-tensor}
$M$을 $R$-가군이라 하자. 다음은 서로 동치이다.
\begin{enumerate}
\item $M$은 유한 생성이다.
\item 모든 $R$-가군의 족 $(Q_{\alpha})_{\alpha \in A}$에 대해 표준
사상 $M \otimes_R \left( \prod_{\alpha} Q_{\alpha} \right)
\to \prod_{\alpha} (M \otimes_R Q_{\alpha})$는 전사이다.
\item 모든 $R$-가군 $Q$와 모든 집합 $A$에 대해 표준 사상
$M \otimes_R Q^{A} \to (M \otimes_R Q)^{A}$는 전사이다.
\item 모든 집합 $A$에 대해 표준 사상
$M \otimes_R R^{A} \to M^{A}$는 전사이다.
\end{enumerate}
\end{proposition}

\begin{proof}
먼저 (1)이 (2)를 함의함을 증명하자. 전사 $R^n \to M$을 택하고
가환 그림
$$
\xymatrix{
R^n \otimes_R (\prod_{\alpha} Q_{\alpha})  \ar[r]^{\cong} \ar[d] &
\prod_{\alpha} (R^n \otimes_R Q_{\alpha}) \ar[d] \\
M \otimes_R (\prod_{\alpha} Q_{\alpha})  \ar[r] & \prod_{\alpha} ( M
\otimes_R Q_{\alpha}).
}
$$
을 생각하자. 위쪽 화살표는 동형이고 수직 화살표들은 전사이다.
따라서 아래쪽 화살표도 전사이다.

\medskip\noindent
(2)가 (3)을, (3)이 (4)를 함의함은 명백하므로 (4)가 (1)을
함의함을 증명하면 된다. 사실 (1)이 성립하려면 $M^M$의 원소
$d = (x)_{x \in M}$이 사상 $f: M \otimes_R R^{M} \to M^M$의 상에
속하는 것으로 충분하다. 이 경우 어떤 $x_i \in M$과
$a_i \in R^M$에 대해
$d = \sum_{i = 1}^{n} f(x_i \otimes a_i)$이다. $x \in M$에 대해
% P01-E0455
$p_x: M^M \to M$을 $x$번째 인자로의 사영이라 쓰면
$$
x = p_x(d) = \sum\nolimits_{i = 1}^{n} p_x(f(x_i \otimes a_i)) =
\sum\nolimits_{i=1}^{n} p_x(a_i) x_i.
$$
따라서 $x_1, \ldots, x_n$은 $M$을 생성한다.
\end{proof}

\begin{proposition}
\label{algebra-proposition-fp-tensor}
$M$을 $R$-가군이라 하자. 다음은 서로 동치이다.
\begin{enumerate}
\item $M$은 유한 표시이다.
\item 모든 $R$-가군의 족 $(Q_{\alpha})_{\alpha \in A}$에 대해 표준
사상 $M \otimes_R \left( \prod_{\alpha} Q_{\alpha} \right)
\to \prod_{\alpha} (M \otimes_R Q_{\alpha})$는 전단사이다.
\item 모든 $R$-가군 $Q$와 모든 집합 $A$에 대해 표준 사상
$M \otimes_R Q^{A} \to (M \otimes_R Q)^{A}$는 전단사이다.
\item 모든 집합 $A$에 대해 표준 사상
$M \otimes_R R^{A} \to M^{A}$는 전단사이다.
\end{enumerate}
\end{proposition}

\begin{proof}
먼저 (1)이 (2)를 함의함을 증명하자. 표시
$R^m \to R^n \to M$을 택하고 가환 그림
$$
\xymatrix{
R^m \otimes_R (\prod_{\alpha} Q_{\alpha}) \ar[r] \ar[d]^{\cong} & R^n
\otimes_R (\prod_{\alpha} Q_{\alpha}) \ar[r] \ar[d]^{\cong} & M \otimes_R
(\prod_{\alpha} Q_{\alpha}) \ar[r] \ar[d] & 0 \\
\prod_{\alpha} (R^m \otimes_R Q_{\alpha}) \ar[r] & \prod_{\alpha} (R^n
\otimes_R Q_{\alpha}) \ar[r] & \prod_{\alpha} (M \otimes_R Q_{\alpha})
\ar[r] & 0.
}
$$
을 생각하자. 처음 두 수직 화살표는 동형이고 두 행은 완전이다.
따라서 사상
$M \otimes_R (\prod_{\alpha} Q_{\alpha}) \to
\prod_{\alpha} ( M \otimes_R Q_{\alpha})$는 전사이고, 그림 추적에
의해 단사이기도 하다. 그러므로 (2)가 성립한다.

\medskip\noindent
(2)가 (3)을, (3)이 (4)를 함의함은 명백하므로 (4)가 (1)을
함의함을 증명하면 된다. 명제 \ref{algebra-proposition-fg-tensor}에 의해
(4)가 성립하면 $M$이 유한 생성임을 이미 안다. 따라서 $F$가 유한
자유인 전사 $F \to M$을 택할 수 있다. 그 핵을 $K$라 하자. $K$가
유한 생성임을 보여야 한다. 임의의 집합 $A$에 대해 가환 그림
$$
\xymatrix{
& K \otimes_R R^A \ar[r] \ar[d]_{f_3} & F \otimes_R R^A \ar[r]
\ar[d]_{f_2}^{\cong} & M \otimes_R R^A \ar[r] \ar[d]_{f_1}^{\cong} & 0 \\
0 \ar[r] & K^A \ar[r] & F^A \ar[r] & M^A \ar[r] & 0 .
}
$$
이 있다. 가정에 의해 $f_1$은 동형이고, $F$가 유한 자유이므로
$f_2$도 동형이며, 두 행은 완전이다. 그림 추적으로 $f_3$가 전사임을
알 수 있으므로 명제 \ref{algebra-proposition-fg-tensor}에 의해 $K$는 유한
생성이다.
\end{proof}

\noindent
다음 명제에 아래 보조정리가 필요하다.

\begin{lemma}
\label{algebra-lemma-kernel-tensored-fp}
$M$을 $R$-가군, $P$를 유한 표시 $R$-가군, $f: P \to M$을 사상이라
하자. $Q$를 $R$-가군이라 하고
$x \in \Ker(P \otimes Q \to M \otimes Q)$라고 하자. 그러면 유한
표시 $R$-가군 $P'$과 사상 $f': P \to P'$가 존재하여 $f$가 $f'$을
통하여 인수분해되고
$x \in \Ker(P \otimes Q \to P' \otimes Q)$이다.
\end{lemma}

\begin{proof}
$M$을 유한 표시 가군 $M_i$들의 유향계의 여극한
$M = \colim_{i \in I} M_i$로 쓰자. $P$가 유한 표시이므로 사상
$f: P \to M$은 어떤 $j \in I$에 대해 $M_j \to M$을 통하여
인수분해된다. $Q$를 텐서하면 가환 그림
$$
\xymatrix{
& M_j \otimes Q \ar[dr] & \\
P \otimes Q \ar[ur] \ar[rr] & & M \otimes Q .
}
$$
을 얻는다. $M_j \otimes Q$에서 $x$의 상 $y$는
$M_j \otimes Q \to M \otimes Q$의 핵에 속한다.
$M \otimes Q = \colim_{i \in I} (M_i \otimes Q)$이므로, 이는 어떤
$j' \geq j$에 대해 $y$가 $M_{j'} \otimes Q$에서 $0$으로 간다는
뜻이다. 따라서 $P' = M_{j'}$로, $f'$은 합성
$P \to M_j \to M_{j'}$로 택하면 된다.
\end{proof}

\begin{proposition}
\label{algebra-proposition-ML-tensor}
$M$을 $R$-가군이라 하자. 다음은 서로 동치이다.
\begin{enumerate}
\item $M$은 Mittag-Leffler이다.
\item 모든 $R$-가군의 족 $(Q_{\alpha})_{\alpha \in A}$에 대해 표준
사상 $M \otimes_R \left( \prod_{\alpha} Q_{\alpha} \right)
\to \prod_{\alpha} (M \otimes_R Q_{\alpha})$는 단사이다.
\end{enumerate}
\end{proposition}

\begin{proof}
먼저 (1)이 (2)를 함의함을 증명하자. $M$이 Mittag-Leffler이고
$x$가 $M \otimes_R (\prod_{\alpha} Q_{\alpha}) \to
\prod_{\alpha} (M \otimes_R Q_{\alpha})$의 핵에 속한다고 하자.
$M$을 유한 표시 가군 $M_i$들의 유향계의 여극한
$M = \colim_{i \in I} M_i$로 쓰자. 그러면
$M \otimes_R (\prod_{\alpha} Q_{\alpha})$는
$M_i \otimes_R (\prod_{\alpha} Q_{\alpha})$들의 여극한이다. 따라서
$x$는 어떤 원소
$x_i \in M_i \otimes_R (\prod_{\alpha} Q_{\alpha})$의 상이다. 어떤
$j \geq i$에 대해 $x_i$가
$M_j \otimes_R (\prod_{\alpha} Q_{\alpha})$에서 $0$으로 감을 보여야
한다. $M$이 Mittag-Leffler이므로 $M_i \to M_j$와 $M_i \to M$이
서로를 지배하도록 하는 $j \geq i$를 택할 수 있다. 이제 가환 그림
$$
\xymatrix{
M \otimes_R (\prod_{\alpha} Q_{\alpha}) \ar[r] & \prod_{\alpha} (M
\otimes_R Q_{\alpha}) \\
M_i \otimes_R (\prod_{\alpha} Q_{\alpha}) \ar[r]^{\cong} \ar[d] \ar[u] &
\prod_{\alpha} (M_i \otimes_R Q_{\alpha}) \ar[d] \ar[u] \\
M_j \otimes_R (\prod_{\alpha} Q_{\alpha}) \ar[r]^{\cong} & \prod_{\alpha}
(M_j \otimes_R Q_{\alpha})
}
$$
을 생각하자. 명제 \ref{algebra-proposition-fp-tensor}에 의해 아래쪽 두 수평
사상은 동형이다. $x_i$는
$\prod_{\alpha} (M \otimes_R Q_{\alpha})$에서 $0$으로 가므로,
$\prod_{\alpha} (M_i \otimes_R Q_{\alpha})$에서 그 상은 사상
$\prod_{\alpha} (M_i \otimes_R Q_{\alpha}) \to
\prod_{\alpha} (M \otimes_R Q_{\alpha})$의 핵에 속한다. 그런데 $j$의
선택에 의해 이 핵은
$\prod_{\alpha} (M_i \otimes_R Q_{\alpha}) \to
\prod_{\alpha} (M_j \otimes_R Q_{\alpha})$의 핵과 같다. 따라서
$x_i$는 $\prod_{\alpha} (M_j \otimes_R Q_{\alpha})$에서 $0$으로 가고,
그러므로 $M_j \otimes_R (\prod_{\alpha} Q_{\alpha})$에서도 $0$으로
간다.

\medskip\noindent
이제 (2)가 성립한다고 하자. 명제
\ref{algebra-proposition-ML-characterization}에서 Mittag-Leffler 조건의
정식화 (1)을 $M$이 만족함을 증명하자. $f: P \to M$을 유한 표시
가군 $P$에서 $M$으로 가는 사상이라 하자. 유한 표시 $R$-가군들의
동형류를 대표하는 집합 $B$를 택하자. $Q \in B$이고
$x \in \Ker(P \otimes Q \to M \otimes Q)$인 쌍 $(Q, x)$들의 집합을
$A$라 하자. $\alpha = (Q, x) \in A$에 대해 $Q_{\alpha}$는 $Q$를,
$x_{\alpha}$는 $x$를 뜻한다고 하자. 가환 그림
$$
\xymatrix{
M \otimes_R (\prod_{\alpha} Q_{\alpha}) \ar[r] &
\prod_{\alpha} (M \otimes_R Q_{\alpha}) \\
P \otimes_R (\prod_{\alpha} Q_{\alpha}) \ar[r]^{\cong} \ar[u] &
\prod_{\alpha} (P \otimes_R Q_{\alpha}) \ar[u] .
}
$$
을 생각하자. 가정에 의해 위쪽 화살표는 단사이고, 명제
\ref{algebra-proposition-fp-tensor}에 의해 아래쪽 화살표는 동형이다. 이
동형 아래에서 $(x_{\alpha}) \in
\prod_{\alpha} (P \otimes_R Q_{\alpha})$에 대응하는 원소를
$x \in P \otimes_R (\prod_{\alpha} Q_{\alpha})$라 하자. 그림의 위쪽
화살표가 단사이므로
$x \in \Ker( P \otimes_R (\prod_{\alpha} Q_{\alpha})
\to M \otimes_R (\prod_{\alpha} Q_{\alpha}))$이다. 보조정리
\ref{algebra-lemma-kernel-tensored-fp}에 의해 유한 표시 가군 $P'$과 사상
$f': P \to P'$가 존재하여 $f: P \to M$이 $f'$을 통하여
인수분해되고
$x \in \Ker(P \otimes_R (\prod_{\alpha} Q_{\alpha})
\to P' \otimes_R (\prod_{\alpha} Q_{\alpha}))$이다. 가환 그림
$$
\xymatrix{
P' \otimes_R (\prod_{\alpha} Q_{\alpha}) \ar[r]^{\cong} &
\prod_{\alpha} (P' \otimes_R Q_{\alpha}) \\
P \otimes_R (\prod_{\alpha} Q_{\alpha}) \ar[r]^{\cong} \ar[u] &
\prod_{\alpha} (P \otimes_R Q_{\alpha}) \ar[u] .
}
$$
이 있다. 여기서 위쪽과 아래쪽 화살표는 모두 명제
\ref{algebra-proposition-fp-tensor}에 의해 동형이다. 따라서 $x$가 왼쪽 수직
사상의 핵에 속하므로 $(x_{\alpha})$는 오른쪽 수직 사상의 핵에
속한다. 이는 모든 $\alpha \in A$에 대해
$x_{\alpha} \in \Ker(P \otimes_R Q_{\alpha} \to P' \otimes_R
Q_{\alpha})$라는 뜻이다. $A$의 정의에 의해 모든 유한 표시 $Q$에
대해
$\Ker(P \otimes_R Q \to P' \otimes_R Q) \supset
\Ker(P \otimes_R Q \to M \otimes_R Q)$이고, $f: P \to M$이
$f': P \to P'$을 통하여 인수분해되므로 실제로 등호가 성립한다.
보조정리 \ref{algebra-lemma-domination-fp}에 의해 $f$와 $f'$은 서로를
지배한다.
\end{proof}

\begin{lemma}
\label{algebra-lemma-minimal-contains}
$M$을 $R$ 위의 평탄 Mittag-Leffler 가군이라 하자. $F$를
$R$-가군이라 하고 $x \in F \otimes_R M$이라 하자. 그러면
$x \in F' \otimes_R M$이 되게 하는 가장 작은 부분가군
$F' \subset F$가 존재한다. 또한 $F'$은 유한 $R$-가군이다.
\end{lemma}

\begin{proof}
$M$이 평탄이므로 $F' \subset F$가 부분가군이면
$F' \otimes_R M \subset F \otimes_R M$이고, 따라서 명제의 뜻이
명확하다. $I = \{F' \subset F \mid x \in F' \otimes_R M\}$라 하고,
$i \in I$에 해당하는 부분가군을 $F_i \subset F$로 표기하자.
그러면 $x$는 사상
$$
F \otimes_R M \longrightarrow \prod (F/F_i \otimes_R M)
$$
아래에서 영으로 가므로, 명제 \ref{algebra-proposition-ML-tensor}에 의해
사상
$$
F \otimes_R M \longrightarrow \left(\prod F/F_i\right) \otimes_R M
$$
아래에서도 $x$는 영으로 간다. $M$이 평탄이므로 이 화살표의 핵은
$(\bigcap F_i) \otimes_R M$이고, 이로써 $F' = \bigcap F_i$임이
증명된다. $F'$이 유한 가군임을 보이기 위해 어떤 $f_j \in F'$와
$m_j \in M$에 대해
$x = \sum_{j = 1, \ldots, m} f_j \otimes m_j$라고 하자. 이때
$F'' \subset F'$이 $f_1, \ldots, f_m$으로 생성되는 부분가군이면
$x \in F'' \otimes_R M$이다. 그러면 물론 $F'' = F'$이고, 마지막
명제가 따른다.
\end{proof}

\begin{lemma}
\label{algebra-lemma-pure-submodule-ML}
$0 \to M_1 \to M_2 \to M_3 \to 0$을 $R$-가군들의 보편 완전열이라
하자. 그러면 다음이 성립한다.
\begin{enumerate}
\item $M_2$가 Mittag-Leffler이면 $M_1$도 Mittag-Leffler이다.
\item $M_1$과 $M_3$가 Mittag-Leffler이면 $M_2$도 Mittag-Leffler이다.
\end{enumerate}
\end{lemma}

\begin{proof}
임의의 $R$-가군의 족 $(Q_{\alpha})_{\alpha \in A}$에 대해 완전한
행들을 갖는 가환 그림
$$
\xymatrix{
0 \ar[r] & M_1 \otimes_R (\prod_{\alpha} Q_{\alpha}) \ar[r] \ar[d] & M_2
\otimes_R (\prod_{\alpha} Q_{\alpha}) \ar[r] \ar[d] & M_3 \otimes_R
(\prod_{\alpha} Q_{\alpha}) \ar[r] \ar[d] & 0 \\
0 \ar[r] & \prod_{\alpha}(M_1 \otimes Q_{\alpha}) \ar[r] & \prod_{\alpha}(M_2
\otimes Q_{\alpha}) \ar[r] & \prod_{\alpha}(M_3 \otimes Q_{\alpha})\ar[r] & 0
}
$$
이 있다. 따라서 (1)과 (2)는 명제 \ref{algebra-proposition-ML-tensor}에서
따른다.
\end{proof}

\begin{lemma}
\label{algebra-lemma-quotient-module-ML}
$M_1 \to M_2 \to M_3 \to 0$을 $R$-가군들의 완전열이라 하자.
$M_1$이 유한 생성이고 $M_2$가 Mittag-Leffler이면 $M_3$도
Mittag-Leffler이다.
\end{lemma}

\begin{proof}
임의의 $R$-가군의 족 $(Q_{\alpha})_{\alpha \in A}$에 대해, 텐서곱은
우완전이므로 완전한 행들을 갖는 가환 그림
$$
\xymatrix{
M_1 \otimes_R (\prod_{\alpha} Q_{\alpha}) \ar[r] \ar[d] & M_2
\otimes_R (\prod_{\alpha} Q_{\alpha}) \ar[r] \ar[d] & M_3 \otimes_R
(\prod_{\alpha} Q_{\alpha}) \ar[r] \ar[d] & 0 \\
\prod_{\alpha}(M_1 \otimes Q_{\alpha}) \ar[r] & \prod_{\alpha}(M_2
\otimes Q_{\alpha}) \ar[r] & \prod_{\alpha}(M_3 \otimes Q_{\alpha})\ar[r] & 0
}
$$
이 있다. 명제 \ref{algebra-proposition-fg-tensor}에 의해 왼쪽 수직 화살표는
전사이다. 명제 \ref{algebra-proposition-ML-tensor}에 의해 가운데 수직
화살표는 단사이다. 그림 추적으로 오른쪽 수직 화살표가 단사임을
알 수 있다. 그러므로 명제 \ref{algebra-proposition-ML-tensor}에 의해 $M_3$는
Mittag-Leffler이다.
\end{proof}


\begin{lemma}
\label{algebra-lemma-colimit-universally-injective-ML}
$M = \colim M_i$가 보편 단사 전이 사상을 갖는 Mittag-Leffler
$R$-가군 $M_i$들의 유향계의 여극한이면 $M$은 Mittag-Leffler이다.
\end{lemma}

\begin{proof}
$(Q_{\alpha})_{\alpha \in A}$를 $R$-가군들의 족이라 하자. 사상
$M \otimes_R (\prod Q_\alpha) \to \prod M \otimes_R Q_\alpha$가
단사임을 보여야 하고, 각 $i$에 대해 사상
$M_i \otimes_R (\prod Q_\alpha) \to \prod M_i \otimes_R Q_\alpha$가
단사임은 명제 \ref{algebra-proposition-ML-tensor}에서 안다. $\otimes$는 여과
여극한과 가환하므로 사상
$\prod M_i \otimes_R Q_\alpha \to \prod M \otimes_R Q_\alpha$가
단사임을 보이면 충분하다. 전이 사상이 보편 단사라는 가정에 의해
각 사상 $M_i \otimes_R Q_\alpha \to M \otimes_R Q_\alpha$가
단사이므로 이는 명백하다.
\end{proof}

\begin{lemma}
\label{algebra-lemma-direct-sum-ML}
$M = \bigoplus_{i \in I} M_i$가 $R$-가군들의 직합이면, $M$이
Mittag-Leffler일 필요충분조건은 각 $M_i$가 Mittag-Leffler인 것이다.
\end{lemma}

\begin{proof}
“필요” 방향은 보조정리 \ref{algebra-lemma-pure-submodule-ML} (1)과 분할되는
짧은 완전열이 보편 완전이라는 사실에서 따른다. 역은 보조정리
\ref{algebra-lemma-colimit-universally-injective-ML}에서 따르지만, 다음과 같이
직접 논증할 수도 있다. 먼저 $I$가 유한이면 보조정리
\ref{algebra-lemma-pure-submodule-ML} (2)에서 따른다. 일반적인 $I$에 대해
모든 $M_i$가 Mittag-Leffler이면, 명제
\ref{algebra-proposition-ML-characterization}의 조건 (1)을 확인하여 $M$도
그러함을 증명한다. $f: P \to M$을 유한 표시 가군 $P$에서 오는
사상이라 하자. 그러면 $I$의 어떤 유한 부분집합 $I'$에 대해 $f$는
$P \xrightarrow{f'} \bigoplus_{i' \in I'} M_{i'} \hookrightarrow
\bigoplus_{i \in I} M_i$로 인수분해된다. 유한인 경우에 의해
$\bigoplus_{i' \in I'} M_{i'}$는 Mittag-Leffler이므로, 어떤 유한 표시
가군 $Q$와 사상 $g: P \to Q$가 존재하여 $g$와 $f'$이 서로를
지배한다. 그러면 $g$와 $f$도 서로를 지배한다.
\end{proof}

\begin{lemma}
\label{algebra-lemma-flat-ML-over-ML-ring}
$R \to S$를 환 사상이라 하고 $M$을 $S$-가군이라 하자. $S$가
$R$-가군으로서 Mittag-Leffler이고 $M$이 $S$-가군으로서 평탄하고
Mittag-Leffler이면, $M$은 $R$-가군으로서 Mittag-Leffler이다.
\end{lemma}

\begin{proof}
명제 \ref{algebra-proposition-ML-tensor}의 특징짓기에서 이를 이끌어 낸다.
실제로 $Q_\alpha$가 $R$-가군들의 족이라고 하자. 합성
$$
\xymatrix{
M \otimes_R \prod_\alpha Q_\alpha =
M \otimes_S S \otimes_R \prod_\alpha Q_\alpha \ar[d] \\
M \otimes_S \prod_\alpha (S \otimes_R Q_\alpha) \ar[d] \\
\prod_\alpha (M \otimes_S S \otimes_R Q_\alpha) =
\prod_\alpha (M \otimes_R Q_\alpha)
}
$$
을 생각하자. $M$이 $S$ 위에서 평탄이고 $S$가 $R$ 위에서
Mittag-Leffler이므로 첫 번째 화살표는 단사이며, $M$이 $S$ 위에서
Mittag-Leffler이므로 두 번째 화살표도 단사이다. 따라서 $M$은
$R$ 위에서 Mittag-Leffler이다.
\end{proof}


\section{코히런트 환}
\label{algebra-section-coherent}

\noindent
$\prod$와 $\otimes$의 교환에 관한 논의를 사용하여 평탄 가군들의
곱이 평탄해지는 환들을 결정한다. 이들은 이른바 코히런트 환임이
드러난다. 환 달린 공간 위의 코히런트 $\mathcal{O}_X$-가군이라는
개념이 더 익숙할 수도 있다. 가군론, 절
\ref{modules-section-coherent}를 보라.

\begin{definition}
\label{algebra-definition-coherent}
$R$을 환, $M$을 $R$-가군이라 하자.
\begin{enumerate}
\item $M$이 유한 생성이고 $M$의 모든 유한 생성 부분가군이 $R$
위에서 유한 표시이면 이를 {\it 코히런트 가군}이라고 한다.
\item $R$이 자기 자신 위의 가군으로서 코히런트이면 이를
{\it 코히런트 환}이라고 한다.
\end{enumerate}
\end{definition}

\noindent
따라서 환이 코히런트일 필요충분조건은 모든 유한 생성 아이디얼이
가군으로서 유한 표시인 것이다.

\begin{example}
\label{algebra-example-valuation-ring-coherent}
부치환은 코히런트 환이다. 실제로 영이 아닌 모든 유한 생성
아이디얼은 주 아이디얼이고(보조정리
\ref{algebra-lemma-characterize-valuation-ring}), 부치환은 정역이므로
자유이며, 따라서 유한 표시이다.
\end{example}

\noindent
코히런트 가군들의 범주는 아벨 범주이다.

\begin{lemma}
\label{algebra-lemma-coherent}
$R$을 환이라 하자.
\begin{enumerate}
\item 코히런트 가군의 유한 부분가군은 코히런트이다.
\item $\varphi : N \to M$을 유한 가군에서 코히런트 가군으로 가는
준동형이라 하자. 그러면 $\Ker(\varphi)$는 유한이고,
$\Im(\varphi)$와 $\Coker(\varphi)$는 코히런트이다.
\item $\varphi : N \to M$을 코히런트 가군들의 준동형이라 하자.
그러면 $\Ker(\varphi)$와 $\Coker(\varphi)$는 코히런트 가군이다.
\item $R$-가군들의 짧은 완전열
$0 \to M_1 \to M_2 \to M_3 \to 0$에서 셋 가운데 둘이 코히런트이면
나머지 하나도 코히런트이다.
\end{enumerate}
\end{lemma}

\begin{proof}
첫 번째 명제는 정의에서 곧바로 따른다.

\medskip\noindent
$\varphi : N \to M$이 (2)의 가정을 만족한다고 하자. 먼저
$\Im(\varphi)$는 유한이므로 (1)에 의해 코히런트이다. 특히
$\Im(\varphi)$는 유한 표시이다. 따라서 완전열
$0 \to \Ker(\varphi) \to N \to \Im(\varphi) \to 0$에 보조정리
\ref{algebra-lemma-extension}을 적용하면 $\Ker(\varphi)$가 유한임을 안다.
$\Coker(\varphi)$가 코히런트임을 보이기 위해
$E \subset \Coker(\varphi)$를 유한 부분가군이라 하고, $E'$을
$M$ 안에서 그 역상이라 하자. 완전열
% P01-E0456
$0 \to \Im(\varphi) \to E' \to E \to 0$과 $\Ker(\varphi)$가
유한이라는 사실에 의해 보조정리 \ref{algebra-lemma-extension}으로부터
$E' \subset M$이 유한임을 결론내린다. $M$이 코히런트이므로 이는
유한 표시이다. 이제 같은 완전열로부터 $E$가 유한 표시임을 얻고,
주장이 증명된다.

\medskip\noindent
(3)은 (1)과 (2)에서 곧바로 따른다.

\medskip\noindent
$0 \to M_1 \xrightarrow{i} M_2 \xrightarrow{p} M_3 \to 0$을 (4)에서와
같은 $R$-가군들의 짧은 완전열이라 하자. $M_1$과 $M_3$가
코히런트이면 $M_2$도 그러함을 증명하는 일만 남는다. 보조정리
\ref{algebra-lemma-extension}에 의해 $M_2$가 유한임을 안다.
$N_2 \subset M_2$를 유한 부분가군이라 하자.
$N_3 = p(N_2) \subset M_3$와 $N_1 = i^{-1}(N_2) \subset M_1$로 놓자.
완전열 $0 \to N_1 \to N_2 \to N_3 \to 0$이 있다. $N_3$는
$N_2$의 몫이므로 유한이고, 따라서 $M_3$의 유한 부분가군이므로
유한 표시이다. 보조정리 \ref{algebra-lemma-extension} (5)에 의해 $N_1$이
유한이고, 따라서 $M_1$의 유한 부분가군으로서 유한 표시이다.
보조정리 \ref{algebra-lemma-extension} (2)에 의해 $N_2$가 유한 표시임을
결론내린다.
\end{proof}

\begin{lemma}
\label{algebra-lemma-coherent-ring}
$R$을 환이라 하자. $R$이 코히런트이면, 가군이 코히런트일
필요충분조건은 유한 표시인 것이다.
\end{lemma}

\begin{proof}
코히런트 가군이 유한 표시임은 (어느 환 위에서도) 명백하다. 역으로
$R$이 코히런트이면 $R^{\oplus n}$도 코히런트이고, 임의의 사상
$R^{\oplus m} \to R^{\oplus n}$의 여핵도 코히런트이다. 보조정리
\ref{algebra-lemma-coherent}를 보라.
\end{proof}

\begin{lemma}
\label{algebra-lemma-Noetherian-coherent}
뇌터 환은 코히런트 환이다.
\end{lemma}

\begin{proof}
보조정리 \ref{algebra-lemma-Noetherian-finite-type-is-finite-presentation}에 의해
모든 유한 $R$-가군은 유한 표시이다. 특히 $R$의 모든 아이디얼은
유한 표시이다.
\end{proof}

\begin{proposition}
\label{algebra-proposition-characterize-coherent}
\begin{reference}
이는 \cite[Theorem 2.1]{Chase}이다.
\end{reference}
$R$을 환이라 하자. 다음은 서로 동치이다.
\begin{enumerate}
\item $R$은 코히런트이다.
\item 평탄 $R$-가군들의 임의의 곱은 평탄이다.
\item 모든 집합 $A$에 대해 가군 $R^A$는 평탄이다.
\end{enumerate}
\end{proposition}

\begin{proof}
$R$이 코히런트라고 가정하고 $Q_\alpha$, $\alpha \in A$를 평탄
$R$-가군들의 집합이라 하자. $R$의 모든 유한 생성 아이디얼 $I$에
대해 $I \otimes_R \prod_\alpha Q_\alpha \to \prod Q_\alpha$가
단사임을 보여야 한다. 보조정리 \ref{algebra-lemma-flat}을 보라. $R$이
코히런트이므로 $I$는 유한 표시 $R$-가군이다. 따라서 명제
\ref{algebra-proposition-fp-tensor}에 의해
$I \otimes_R \prod_\alpha Q_\alpha = \prod I \otimes_R Q_\alpha$이다.
$Q_\alpha$의 평탄성에 의해 $I \otimes_R Q_\alpha \to Q_\alpha$가
단사이므로 원하는 단사성이 따른다.

\medskip\noindent
(2) $\Rightarrow$ (3)은 자명하다.

\medskip\noindent
모든 집합 $A$에 대해 $R$-가군 $R^A$가 평탄이라고 가정하자.
$I$를 $R$의 유한 생성 아이디얼이라 하자. 가정에 의해
$I \otimes_R R^A \to R^A$는 단사이다. 명제
\ref{algebra-proposition-fg-tensor}와 $I$의 유한성에 의해 그 상은 $I^A$와
같다. 따라서 모든 집합 $A$에 대해 $I \otimes_R R^A = I^A$이고,
명제 \ref{algebra-proposition-fp-tensor}에 의해 $I$가 유한 표시임을
결론내린다.
\end{proof}







\section{Mittag-Leffler 가군의 예와 반례}
\label{algebra-section-examples}

\noindent
Mittag-Leffler 가군의 몇 가지 예와 반례로 이 논의를 마친다.

\begin{example}
\label{algebra-example-ML}
Mittag-Leffler 가군들.
\begin{enumerate}
\item 모든 유한 표시 가군은 Mittag-Leffler이다. 예를 들어 이는
명제 \ref{algebra-proposition-ML-characterization} (1)에서 따른다. 일반적으로
% P01-E0457
유한 생성 가군이 Mittag-Leffler일 필요충분조건은 유한 표시인
것이다. 이는 명제 \ref{algebra-proposition-fg-tensor},
\ref{algebra-proposition-fp-tensor}, \ref{algebra-proposition-ML-tensor}에서 따른다.
\item 자유가군은 명제 \ref{algebra-proposition-ML-characterization}의 조건
(1)을 만족하므로 Mittag-Leffler이다.
\item 앞의 예와 보조정리 \ref{algebra-lemma-direct-sum-ML}에 의해 사영
가군은 Mittag-Leffler이다.
\end{enumerate}
\end{example}

\noindent
예의 목록에 뇌터 환 $R$ 위의 멱급수환도 추가하고자 한다.
% P01-E0458
이는 다음 보조정리의 따름결과가 될 것이다.

\begin{lemma}
\label{algebra-lemma-flat-ML-criterion}
$M$을 평탄 $R$-가군이라 하자. 다음은 서로 동치이다.
\begin{enumerate}
\item $M$은 Mittag-Leffler이다.
\item $F$가 유한 자유 $R$-가군이고 $x \in F \otimes_R M$이면,
$x \in F' \otimes_R M$이 되게 하는 $F$의 가장 작은 부분가군
$F'$이 존재한다.
\end{enumerate}
\end{lemma}

\begin{proof}
(1) $\Rightarrow$ (2)는 보조정리 \ref{algebra-lemma-minimal-contains}의 특수한
경우이다. (2)를 가정하자. 정리 \ref{algebra-theorem-lazard}에 의해 $M$을
유한 자유 $R$-가군들의 유향계 $(M_i, f_{ij})$의 여극한
$M = \colim_{i \in I} M_i$로 쓸 수 있다. 주의
\ref{algebra-remark-flat-ML}에 의해 역계
$(\Hom_R(M_i, R), \Hom_R(f_{ij}, R))$가 Mittag-Leffler임을 보이면
충분하다. 다시 말해 $i \in I$를 고정하고, $j \geq i$에 대해
$Q_j$를 $\Hom_R(M_j, R) \to \Hom_R(M_i, R)$의 상이라 할 때,
$Q_j$들이 안정화됨을 보여야 한다.

\medskip\noindent
$M_i$가 유한 자유이므로 모든 $j$에 대해
$\Hom_R(M_i, M_j) = \Hom_R(M_i, R) \otimes_R  M_j$로 식별할 수 있다.
$M_j$들이 자유라는 사실을 사용하면, $j \geq i$에 대해 $Q_j$는
$f_{ij} \in Q_j \otimes_R M_j$가 되게 하는
$\Hom_R(M_i, R)$의 가장 작은 부분가군이다. 식별
$\Hom_R(M_i, M) = \Hom_R(M_i, R) \otimes_R M$ 아래에서 표준 사상
$f_i: M_i \to M$은 $\Hom_R(M_i, R) \otimes_R M$에 속한다. $M$에
대한 가정에 의해 $f_i \in Q \otimes_R M$이 되게 하는
$\Hom_R(M_i, R)$의 가장 작은 부분가군 $Q$가 존재한다. $Q_j$들이
$Q$로 안정화됨을 보이겠다.

\medskip\noindent
$j \geq i$에 대해 가환 그림
$$
\xymatrix{
Q_j \otimes_R M_j \ar[r] \ar[d] & \Hom_R(M_i, R) \otimes_R M_j
\ar[d] \\
Q_j \otimes_R M \ar[r] & \Hom_R(M_i, R) \otimes_R M.
}
$$
이 있다. $f_{ij} \in Q_j \otimes_R M_j$는
$f_i \in \Hom_R(M_i, R) \otimes_R M$으로 가므로
$f_i \in Q_j \otimes_R M$이다. 따라서 $Q$의 선택에 의해 모든
$j \geq i$에 대해 $Q \subset Q_j$이다.

\medskip\noindent
$Q_j$들이 감소하고 모든 $j \geq i$에 대해 $Q \subset Q_j$이므로,
$Q_j$들이 $Q$로 안정화됨을 보이려면 $Q_j \subset Q$인 어떤
$j \geq i$를 찾으면 충분하다. 다음 가군의 원소로서
% P01-E0459
$$
\Hom_R(M_i, R) \otimes_R M = \colim_{j \in J}
(\Hom_R(M_i, R) \otimes_R M_j),
$$
$f_i$는 $j \geq i$일 때의 $f_{ij}$들의 여극한이고, 또한 $f_i$는
부분가군
$$
\colim_{j \in J} (Q \otimes_R M_j) \subset \colim_{j \in J}
(\Hom_R(M_i, R) \otimes_R M_j).
$$
에 속한다. 따라서 어떤 $j \geq i$에 대해 $f_{ij}$는
$Q \otimes_R M_j$에 속한다. $Q_j$는
$f_{ij} \in Q_j \otimes_R M_j$가 되게 하는 $\Hom_R(M_i, R)$의 가장
작은 부분가군이므로 $Q_j\subset Q$임을 결론내린다.
\end{proof}

\begin{lemma}
\label{algebra-lemma-product-over-Noetherian-ring}
$R$을 뇌터 환, $A$를 집합이라 하자. 그러면 $M = R^A$는 평탄
Mittag-Leffler $R$-가군이다.
\end{lemma}

\begin{proof}
보조정리 \ref{algebra-lemma-Noetherian-coherent}와 명제
\ref{algebra-proposition-characterize-coherent}를 결합하면 $M$이 $R$ 위에서
평탄임을 안다. $M$이 보조정리 \ref{algebra-lemma-flat-ML-criterion}의 조건을
만족함을 보이자. $F$를 유한 자유 $R$-가군이라 하자. $F'$이 $F$의
임의의 부분가군이면 $R$이 뇌터이므로 그것은 유한 표시이다. 따라서
명제 \ref{algebra-proposition-fp-tensor}에 의해 가환 그림
$$
\xymatrix{
F' \otimes_R M \ar[r] \ar[d]^{\cong} & F \otimes_R M \ar[d]^{\cong} \\
(F')^A \ar[r] & F^A
}
$$
이 있다. 이를 통해 사상 $F' \otimes_R M \to F \otimes_R M$을
$(F')^A \to F^A$와 식별할 수 있다. 따라서
$x \in F \otimes_R M$이 $(x_\alpha) \in F^A$에 대응하면,
% P01-E0460
$x_\alpha$들로 생성되는 $F$의 부분가군 $F'$은
$x \in F' \otimes_R M$이 되게 하는 $F$의 가장 작은 부분가군이다.
\end{proof}

\begin{lemma}
\label{algebra-lemma-power-series-ML}
$R$을 뇌터 환, $n$을 양의 정수라 하자. 그러면 $R$-가군
$M = R[[t_1, \ldots, t_n]]$은 평탄이고 Mittag-Leffler이다.
\end{lemma}

\begin{proof}
$R$-가군으로서 어떤 (가산) 집합 $A$에 대해 $M = R^A$이다. 따라서
이 보조정리는 보조정리 \ref{algebra-lemma-product-over-Noetherian-ring}의
특수한 경우이다.
\end{proof}

\begin{example}
\label{algebra-example-not-ML}
Mittag-Leffler가 아닌 가군들.
\begin{enumerate}
\item 예 \ref{algebra-example-Q-not-ML}과 명제 \ref{algebra-proposition-ML-tensor}에
의해 $\mathbf{Q}$는 Mittag-Leffler $\mathbf{Z}$-가군이 아니다.
\item 아래에서(정리 \ref{algebra-theorem-projectivity-characterization}) 평탄
가산 생성 가군이 사영일 필요충분조건은 Mittag-Leffler인 것임을
증명한다. 따라서 사영이 아닌 임의의 평탄 가산 생성 가군 $M$은
Mittag-Leffler가 아닌 가군의 예이다. 그러한 예는 주의
\ref{algebra-remark-warning}를 보라.
\item $k$를 체라 하자. $R = k[[x]]$라 하자. $R$-가군
$M = \prod_{n \in \mathbf{N}} R/(x^n)$은 Mittag-Leffler가 아니다.
실제로 $\xi_{2^m} = x^{2^{m - 1}}$이고 그 밖에는 $\xi_n = 0$으로
정의되는 원소 $\xi = (\xi_1, \xi_2, \xi_3, \ldots)$를 생각하자. 즉,
$$
\xi = (0, x, 0, x^2, 0, 0, 0, x^4, 0, 0, 0, 0, 0, 0, 0, x^8, \ldots)
$$
이다. 그러면 $m \gg 0$에 대해 $M/x^{2^m}M$에서 $\xi$의 소멸자는
% P01-E0461
$x^{2^{m - 1}}$로 생성된다. 그러나 $M$이 Mittag-Leffler라면,
유한 $R$-가군 $Q$와 원소 $\xi' \in Q$가 존재하여 모든
$l \geq 1$에 대해 $Q/x^l Q$에서 $\xi'$의 소멸자가 $M/x^l M$에서
$\xi$의 소멸자와 같을 것이다. 명제
\ref{algebra-proposition-ML-characterization} (1)을 보라. 이제 어떤 정수
$a \geq 0$가 존재하여 모든 $l \gg 0$에 대해 $Q/x^l Q$에서
$\xi'$의 소멸자는 $x^a$ 또는 $x^{l - a}$로 생성됨을 증명할 수 있다
($\xi' \in Q$가 꼬임인지 아닌지에 따라 달라진다). 위의 사실들을
결합하면 모든
% P01-E0462
$l = 2^m >> 0$에 대해 $a = l/2$ 또는 $l - a = l/2$라는 등식을
얻게 되는데, 이는 불합리하다.
\item 같은 논증은 $\bigoplus_{n \in \mathbf{N}} R/(x^n)$의
$(x)$-진 완비화가 $R = k[[x]]$ 위에서 Mittag-Leffler가 아님을
보인다(힌트: $\xi$는 실제로 이 완비화의 원소이다).
\item $R = k[a, b]/(a^2, ab, b^2)$라 하자. $S$를 표시
$S = R[t]/(at - b)$를 갖는 유한 표시 $R$-대수라 하자. 그러면
$R$-가군으로서 $S$는 가산 생성이고 분해 불가능이다(자세한 내용은
생략한다). 한편 $R$은 아르틴 국소환이므로 완비 국소환이고, 따라서
헨젤 국소환이다. 보조정리 \ref{algebra-lemma-complete-henselian}을 보라.
$S$가 $R$-가군으로서 Mittag-Leffler라면 보조정리
\ref{algebra-lemma-split-ML-henselian}에 의해 유한 $R$-가군들의 직합일
것이다. 따라서 $S$는 $R$-가군으로서 Mittag-Leffler가 아니다.
\end{enumerate}
\end{example}


\section{가산 생성 Mittag-Leffler 가군}
\label{algebra-section-ML-countable}

\noindent
가산 생성 Mittag-Leffler 가군은 특히 단순한 구조를 갖는 것으로
드러난다.

\begin{lemma}
\label{algebra-lemma-ML-countable-colimit}
$M$을 $R$-가군이라 하자. $(M_i, f_{ij})$가 유한 표시 $R$-가군들의
유향계일 때 $M = \colim_{i \in I} M_i$라 쓰자. $M$이
Mittag-Leffler이고 가산 생성이면, 어떤 유향 가산 부분집합
$I' \subset I$가 존재하여 $M \cong \colim_{i \in I'} M_i$이다.
\end{lemma}

\begin{proof}
$x_1, x_2, \ldots$를 $M$의 가산 생성원 집합이라 하자. 각 $x_n$에
대해 $x_n$이 표준 사상 $f_i: M_i \to M$의 상에 속하도록 하는
$i \in I$를 택하고, 이 모든 $i$의 집합을 $I'_{0} \subset I$라 하자.
이제 $M$이 Mittag-Leffler이므로 각 $i \in I'_{0}$에 대해
$j \geq i$이고 모든 $k \geq i$에 대해
$f_{ij}: M_i \to M_j$가 $f_{ik}: M_i \to M_k$를 통하여
인수분해되도록 하는 $j \in I$를 택할 수 있다(명제
\ref{algebra-proposition-ML-characterization}의 조건 (3)). 이 모든 $j$와
$I'_0$의 합집합을 $I'_1$이라 하자.
% P01-E0463
$I'_1$이 가산이므로 이를 가산 유향집합 $I'_{2} \subset I$로
확대할 수 있다. 이제 $I'_{0}$에 적용한 것과 같은 절차를
$I'_{2}$에 적용하여 새 가산 집합 $I'_{3} \subset I$를 얻는다.
그런 다음 $I'_{3}$을 가산 유향집합 $I'_{4}$로 확대한다. 이와 같이
계속한다. 즉,
% P01-E0464
$\ell$이 홀수이면 각 $ i \in I'_{\ell}$에 대해 명제
\ref{algebra-proposition-ML-characterization} (3)에서와 같은 $j$를 추가하고,
$\ell$이 짝수이면 $I'_{\ell}$을 유향집합으로 확대하여,
$\ell \geq 0$에 대한 부분집합의 열 $I'_{\ell} \subset I$를 얻는다.
합집합 $I' = \bigcup I'_{\ell}$은 다음을 만족한다.
\begin{enumerate}
\item $I'$은 가산이고 유향이다.
\item 각 $x_n$은 어떤 $i \in I'$에 대해 $f_i: M_i \to M$의 상에
속한다.
\item $i \in I'$이면 어떤 $j \in I'$가 존재하여 $j \geq i$이고,
모든 $k \in I$ 중 $k \geq i$인 것에 대해 $f_{ij}: M_i \to M_j$가
$f_{ik}: M_i \to M_k$를 통하여 인수분해된다. 특히 $k \geq i$이면
$\Ker(f_{ik}) \subset \Ker(f_{ij})$이다.
\end{enumerate}
표준 사상
$\colim_{i \in I'} M_i \to \colim_{i \in I} M_i = M$이 동형이라고
주장한다. (2)에 의해 이는 전사이다. 단사성을 위해
$x \in \colim_{i \in I'} M_i$가 $\colim_{i \in I} M_i$에서 $0$으로
간다고 하자. $x$를 어떤 $i \in I'$에 대한 원소
$\tilde{x} \in M_i$로 나타내면, 이는 어떤 $k \in I, k \geq i$에
대해 $f_{ik}(\tilde{x}) = 0$이라는 뜻이다. 그런데 (3)에 의해 어떤
$j \in I', j \geq i,$가 존재하여 $f_{ij}(\tilde{x}) = 0$이다.
따라서 $\colim_{i \in I'} M_i$에서 $x = 0$이다.
\end{proof}

\noindent
보조정리 \ref{algebra-lemma-ML-countable-colimit}에 의해 $R$ 위의 가산 생성
Mittag-Leffler 가군 $M$은 각 $M_n$이 유한 표시 $R$-가군인 계
$$
M_1 \to M_2 \to M_3 \to M_4 \to \ldots
$$
의 여극한이다. 이를 보려면 보조정리
\ref{algebra-lemma-ML-limit-nonempty}의 증명에서와 같이 논증하여 가산
유향집합이
% P01-E0465
$(\mathbf{N}, \geq)$와 동형인 공종 부분집합을 가짐을 보라.
$R = k[x_1, x_2, x_3, \ldots]$이고 $M = R/(x_i)$라고 하자. 그러면
$M$은 유한 생성이지만 유한 표시가 아니므로 Mittag-Leffler가 아니다
(예 \ref{algebra-example-ML}의 (1)을 보라). 그러나
$M_n = R/(x_1, \ldots, x_n)$으로 놓으면 물론
$M = \colim_n M_n$으로 쓸 수 있다. 따라서 그러한 극한으로 쓸 수
$M$을 그러한 극한으로 쓸 수 있다는 조건은 $M$이 Mittag-Leffler임을
함의하지 않는다.

\begin{lemma}
\label{algebra-lemma-ML-countable}
$R$을 환이라 하자. $M$을 $R$-가군이라 하자. $M$이
Mittag-Leffler이고 가산 생성이라고 가정하자. $P$가 유한 생성일 때
임의의 $R$-가군 사상 $f : P \to M$에 대해 다음을 만족하는
자기준동형 $\alpha : M \to M$이 존재한다.
\begin{enumerate}
\item $\alpha : M \to M$은 유한 표시 $R$-가군을 통하여
인수분해된다.
\item $\alpha \circ f = f$이다.
\end{enumerate}
\end{lemma}

\begin{proof}
$I$가 가산인 유한 표시 $R$-가군들의 유향 여극한으로
$M = \colim_{i \in I} M_i$라 쓰자. 보조정리
\ref{algebra-lemma-ML-countable-colimit}를 보라. 전이 사상을 $f_{ij}$로
표기하고 $M$으로 가는 표준 사상을 $f_i : M_i \to M$으로 표기한다.
$N = \prod_{s \in I} M_s$로 놓자. 다음과 같이 표기하자.
$$
M_i^* = \Hom_R(M_i, N) = \prod\nolimits_{s \in I} \Hom_R(M_i, M_s)
$$
그러면 $(M_i^*)$는 $I$ 위의 $R$-가군의 역계이다. 또한
$\Hom_R(M, N) = \lim M_i^*$임에 유의하라. $M$이 Mittag-Leffler이므로
모든 $i \in I$에 대해 어떤 지표 $k(i) \geq i$를 찾아
$$
E_i := \bigcap\nolimits_{i' \geq i} \Im(M_{i'}^* \to M_i^*)
=
\Im(M_{k(i)}^* \to M_i^*)
$$
가 되게 할 수 있다. $\Im(P \to M) \subset \Im(M_j \to M)$이 되도록
$j \in I$를 하나 택하여 고정하자. $P$가 유한 생성이므로 이는
가능하다. $k = k(j)$로 놓자. 다음 원소를 생각하자.
$x = (0, \ldots, 0, \text{id}_{M_k}, 0, \ldots, 0) \in M_k^*$
이 원소는
$y = (0, \ldots, 0, f_{jk}, 0, \ldots, 0) \in M_j^*$로 간다.
$k$의 선택에 의해 $y \in E_j$이다. 예
\ref{algebra-example-ML-surjective-maps}에 의해 각 $i \geq j$에 대해 전이
사상 $E_i \to E_j$는 전사이고
$\lim E_i = \lim M_i^* = \Hom_R(M, N)$이다. 따라서 보조정리
\ref{algebra-lemma-ML-limit-nonempty}에 의해 $E_j \subset M_j^*$에서 $y$로
가는 원소 $z \in \Hom_R(M, N)$가 존재한다. $z_k$를 $z$의 $k$번째
성분이라 하자. 그러면 $z_k : M \to M_k$는 그림
$$
\xymatrix{
M \ar[r]_{z_k} & M_k \\
M_j \ar[ru]_{f_{jk}} \ar[u]^{f_j}
}
$$
을 가환하게 하는 준동형이다. $\alpha : M \to M$를 합성
$f_k \circ z_k : M \to M_k \to M$이라 하자. 구성에 의해 $\alpha$는
유한 표시 가군을 통하여 인수분해되고
$\alpha \circ f_j = f_j$이다. $f$의 상이 $f_j$의 상에 포함되므로
이로부터 $\alpha \circ f = f$도 따른다.
\end{proof}

\noindent
나중에(보조정리 \ref{algebra-lemma-split-ML-henselian}를 보라) 보조정리
\ref{algebra-lemma-ML-countable}가 다음을 뜻함을 볼 것이다. 헨젤 국소환
위의 가산 생성 Mittag-Leffler 가군은 유한 표시 가군들의 직합이다.


\section{사영 가군의 특징짓기}
\label{algebra-section-characterize-projective}

\noindent
이 절의 목표는 가군이 사영 가군일 필요충분조건이 평탄하고
Mittag-Leffler이며 가산 생성 가군들의 직합인 것임을 증명하는
것이다(아래의 정리 \ref{algebra-theorem-projectivity-characterization}).

\begin{lemma}
\label{algebra-lemma-countgen-projective}
$M$을 $R$-가군이라 하자. $M$이 평탄하고 Mittag-Leffler이며
가산 생성이면 $M$은 사영 가군이다.
\end{lemma}

\begin{proof}
Lazard 정리(정리 \ref{algebra-theorem-lazard})에 의해 $M =
\colim_{i
\in I} M_i$를 집합 $I$로 지표화된 유한 자유 $R$-가군의 유향계
$(M_i, f_{ij})$로 쓸 수 있다. 보조정리
\ref{algebra-lemma-ML-countable-colimit}에 의해 $I$가 가산이라고 가정해도
된다. 이제
$$
0 \to N_1 \to N_2 \to N_3 \to 0
$$
을 $R$-가군의 완전열이라 하자. $\Hom_R(M, -)$을 적용해도
완전성이 보존됨을 보여야 한다. $M_i$가 유한 자유 가군이므로 각
$i$에 대해
$$
0 \to \Hom_R(M_i, N_1) \to \Hom_R(M_i, N_2) \to
\Hom_R(M_i, N_3) \to 0
$$
은 완전열이다. $M$이 Mittag-Leffler이므로
$(\Hom_R(M_i, N_{1}))$은 Mittag-Leffler 역계이다. 따라서 보조정리
\ref{algebra-lemma-ML-exact-sequence}에 의해
$$
0 \to \lim_{i \in I} \Hom_R(M_i, N_1) \to
\lim_{i \in I} \Hom_R(M_i, N_2) \to
\lim_{i \in I} \Hom_R(M_i, N_3) \to 0
$$
은 완전열이다. 한편 임의의 $R$-가군 $N$에 대해 함자적 동형사상
$\Hom_R(M, N) \cong \lim_{i \in I} \Hom_R(M_i, N)$이 있으므로
$$
0 \to \Hom_R(M, N_1) \to \Hom_R(M, N_2) \to
\Hom_R(M, N_3) \to 0
$$
도 완전열이다.
\end{proof}

\begin{remark}
\label{algebra-remark-characterize-projective}
보조정리 \ref{algebra-lemma-countgen-projective}는 가산 생성이라는 가정이
없으면 성립하지 않는다. 예를 들어 $\mathbf Z$-가군 $M =
\mathbf{Z}[[x]]$는 평탄하고 Mittag-Leffler이지만 사영 가군이
아니다. 보조정리 \ref{algebra-lemma-power-series-ML}에 의해 이 가군은
Mittag-Leffler이다. 자유 아벨군의 부분군은 자유이므로 사영
$\mathbf Z$-가군은 사실 자유 가군이고, 그 부분가군도 자유
가군이다. 따라서 $M$이 사영 가군이 아님을 보이려면 자유가 아닌
부분가군 하나를 구성하면 충분하다. 소수 $p$를 고정하고, 모든 정수
$m \geq 1$에 대해 유한 개를 제외한 모든 $i$에서 $p^m$이 $a_i$를
나누는 멱급수 $f(x) = \sum a_i x^i$들로 이루어진 부분가군 $N$을
생각하자. 그러면 모든 $a_i \in \mathbf{Z}$에 대해
$\sum a_i p^i x^i$가 $N$에 속하므로 $N$은 비가산이다. 따라서
$N$이 자유 가군이라면 그 랭크는 비가산이고, $\mathbf{Z}/p$ 위의
$N/pN$의 차원도 비가산이어야 한다. 그러나 $i \geq 0$인 원소
$x^i \in N/pN$들이 $N/pN$을 생성하므로 이는 사실이 아니다.
\end{remark}

\begin{theorem}
\label{algebra-theorem-projectivity-characterization}
$M$을 $R$-가군이라 하자. $M$이 사영 가군일 필요충분조건은 다음과
같다.
\begin{enumerate}
\item $M$은 평탄하다.
\item $M$은 Mittag-Leffler이다.
\item $M$은 가산 생성 $R$-가군들의 직합이다.
\end{enumerate}
\end{theorem}

\begin{proof}
먼저 $M$이 사영 가군이라고 하자. 그러면 $M$은 자유 가군의
직합인자이다. 평탄성과 Mittag-Leffler 성질은 직합인자에
내려가므로 $M$은 평탄하고 Mittag-Leffler이다. Kaplansky
정리(정리 \ref{algebra-theorem-projective-direct-sum})에 의해 $M$은
(3)을 만족한다.

\medskip\noindent
역으로 $M$이 (1)--(3)을 만족한다고 하자. 평탄성과 Mittag-Leffler
성질은 직합인자에 내려가므로 $M$은 평탄하고 Mittag-Leffler인
가산 생성 $R$-가군들의 직합이다. 보조정리
\ref{algebra-lemma-countgen-projective}에 의해 $M$은 사영 가군들의
직합이다. 따라서 $M$은 사영 가군이다.
\end{proof}

\begin{lemma}
\label{algebra-lemma-ML-ui-descent}
% P01-E0466: frozen English omits an article before “universally injective map”.
$f: M \to N$을 $R$-가군의 보편 단사 사상이라 하자. $M$이 가산
생성 $R$-가군들의 직합이고 $N$이 평탄하고 Mittag-Leffler라고
가정하자. 그러면 $M$은 사영 가군이다.
\end{lemma}

\begin{proof}
보조정리 \ref{algebra-lemma-ui-flat-domain}과
\ref{algebra-lemma-pure-submodule-ML}에 의해 $M$은 평탄하고
Mittag-Leffler이므로, 정리
\ref{algebra-theorem-projectivity-characterization}에서 결론이 따른다.
\end{proof}

\begin{lemma}
\label{algebra-lemma-universally-injective-submodule-powerseries}
% P01-E0467: frozen English uses “a R-module” instead of “an R-module”.
$R$을 뇌터 환이라 하고 $M$을 $R$-가군이라 하자. $M$이 가산 생성
$R$-가군들의 직합이고, 어떤 $n$에 대해 보편 단사 사상
$M \to R[[t_1, \ldots, t_n]]$이 존재한다고 가정하자. 그러면 $M$은
사영 가군이다.
\end{lemma}

\begin{proof}
보조정리 \ref{algebra-lemma-ML-ui-descent}와
\ref{algebra-lemma-power-series-ML}에서 따른다.
\end{proof}



\section{가군 성질의 상승}
\label{algebra-section-ascent}

\noindent
정리 \ref{algebra-theorem-projectivity-characterization}에 나오는 가군의
모든 성질은 임의의 환 사상을 따라 상승한다.

\begin{lemma}
\label{algebra-lemma-ascend-properties-modules}
$R \to S$를 환 사상이라 하고 $M$을 $R$-가군이라 하자. 그러면
다음이 성립한다.
\begin{enumerate}
\item $M$이 평탄하면 $S$-가군 $M \otimes_R S$도 평탄하다.
\item $M$이 Mittag-Leffler이면 $S$-가군 $M \otimes_R S$도
Mittag-Leffler이다.
\item $M$이 가산 생성 $R$-가군들의 직합이면 $S$-가군
$M \otimes_R S$도 가산 생성 $S$-가군들의 직합이다.
\item $M$이 사영 가군이면 $S$-가군 $M \otimes_R S$도 사영
가군이다.
\end{enumerate}
\end{lemma}

\begin{proof}
(2)를 제외하면 모두 자명하다. (2)에 대해서는 Mittag-Leffler
성질에 관한 명제 \ref{algebra-proposition-ML-characterization}의 조건
(3)과 텐서곱이 여극한을 취하는 것과 교환한다는 사실을 사용한다.
또는 명제 \ref{algebra-proposition-ML-tensor}의 특징짓기를 사용할 수도
있다.
\end{proof}



\section{가군 성질의 하강}
\label{algebra-section-descent}

\noindent
사영성을 특징짓는 정리
\ref{algebra-theorem-projectivity-characterization}의 성질들이 충실 평탄
사상에 대해 하강하는지를 다룬다. 평탄성이 있을 때
Mittag-Leffler 가군이라는 성질은 하강한다.

\begin{lemma}
\label{algebra-lemma-ffdescent-ML}
\begin{reference}
2014년 12월 20일 Juan Pablo Acosta Lopez가 보낸 전자우편.
\end{reference}
$R \to S$를 충실 평탄 환 사상이라 하고 $M$을 $R$-가군이라 하자.
$S$-가군 $M \otimes_R S$가 Mittag-Leffler이면 $M$도
Mittag-Leffler이다.
\end{lemma}

\begin{proof}
$M = \colim_{i\in I} M_i$를 유한 표시 $R$-가군 $M_i$들의 유향
여극한으로 쓰자. 명제 \ref{algebra-proposition-ML-characterization}를
사용하면, 각 $i \in I$에 대해 $M_i\rightarrow M_j$가
$M_i\rightarrow M$을 지배하도록 하는 $i \leq j$, $j\in I$가
존재함을 증명하면 됨을 알 수 있다.

\medskip\noindent
% P01-E0468: frozen English omits “to be” in the pushout construction.
$N$을 다음 밀어내기로 잡자.
$$
\xymatrix{
M_i  \ar[r] \ar[d] & M_j \ar[d] \\
M \ar[r] & N
}
$$
그러면 보조정리
\ref{algebra-lemma-domination-universally-injective}에 의해 이 보조정리는
$M_j\rightarrow N$이 보편 단사가 되는 $j$가 존재한다는 것과
동치이다.
% P01-E0469: frozen English resumes this sentence with capitalized “Is”.
$S$로 텐서화한 그림
$$
\xymatrix{
M_i\otimes_R S  \ar[r] \ar[d] & M_j\otimes_R S \ar[d] \\
M\otimes_R S \ar[r] & N\otimes_R S
}
$$
도 밀어내기 그림이다. 이제
$M \otimes_R S = \colim_{i\in I} M_i \otimes_R S$는
$M\otimes_R S$를 유한 표시 $S$-가군들의 여극한으로 나타내고,
$M\otimes_R S$는 Mittag-Leffler이므로
$M_j\otimes_R S\rightarrow N\otimes_R S$가 보편 단사가 되는
$j \geq i$가 존재한다. $R\rightarrow S$가 충실 평탄이므로
$M_j\rightarrow N$도 보편 단사라고 결론 내린다.
\end{proof}

\begin{lemma}
\label{algebra-lemma-ffdescent-countable}
$R \to S$를 충실 평탄 환 사상이라 하고 $M$을 $R$-가군이라 하자.
$S$-가군 $M \otimes_R S$가 가산 생성이면 $M$도 가산 생성이다.
\end{lemma}

\begin{proof}
$M \otimes_R S$가 원소 $y_i$, $i = 1, 2, 3, \ldots$에 의해
생성된다고 하자. 어떤 $n_i \geq 0$, $x_{ij} \in M$,
$s_{ij} \in S$에 대해
$y_i = \sum_{j = 1, \ldots, n_i} x_{ij} \otimes s_{ij}$로 쓰자.
원소 $x_{ij}$들의 가산 모음이 생성하는 부분가군을
$M' \subset M$이라 하자. $M' \otimes_R S \to M \otimes_R S$의
상은 생성원 $y_i$들을 포함하므로 이 사상은 전사이다. $S$가
$R$ 위에서 충실 평탄이므로 원하는 대로 $M' = M$이라고 결론
내린다.
\end{proof}

\noindent
이제 가산 생성 사영 가군에 대한 충실 평탄 하강은 쉽게 따라온다.

\begin{lemma}
\label{algebra-lemma-ffdescent-countable-projectivity}
$R \to S$를 충실 평탄 환 사상이라 하고 $M$을 $R$-가군이라 하자.
$S$-가군 $M \otimes_R S$가 가산 생성 사영 가군이면 $M$도 가산
생성 사영 가군이다.
\end{lemma}

\begin{proof}
보조정리 \ref{algebra-lemma-descend-properties-modules},
\ref{algebra-lemma-ffdescent-ML}, \ref{algebra-lemma-ffdescent-countable}과 정리
\ref{algebra-theorem-projectivity-characterization}에서 따른다.
\end{proof}

\noindent
이제 남은 일은 d\'evissage를 사용하여 일반적인 사영성의 하강을
가산 생성인 경우로 환원하는 것이다. 먼저 간단한 두 보조정리를
제시한다.

\begin{lemma}
\label{algebra-lemma-lift-countably-generated-submodule}
$R \to S$를 환 사상, $M$을 $R$-가군, $Q$를
$M \otimes_R S$의 가산 생성 $S$-부분가군이라 하자. 그러면
$\Im(P \otimes_R S \to M \otimes_R S)$가 $Q$를 포함하도록 하는
$M$의 가산 생성 $R$-부분가군 $P$가 존재한다.
\end{lemma}

\begin{proof}
$y_1, y_2, \ldots$를 $Q$의 생성원들이라 하고, 어떤
$x_{jk} \in M$과 $s_{jk} \in S$에 대해
$y_j = \sum_k x_{jk} \otimes s_{jk}$로 쓰자.
% P01-E0470: frozen English omits “to” in “take P be”.
그러면 $x_{jk}$들이 생성하는 $M$의 부분가군을 $P$로 잡으면 된다.
\end{proof}

\begin{lemma}
\label{algebra-lemma-adapted-submodule}
$R \to S$를 환 사상이라 하고 $M$을 $R$-가군이라 하자.
$M \otimes_R S = \bigoplus_{i \in I} Q_i$가 가산 생성
$S$-가군 $Q_i$들의 직합이라고 가정하자. $N$이 $M$의 가산 생성
부분가군이면, $N' \supset N$이고 어떤 부분집합 $I' \subset I$에
대해 $\Im(N' \otimes_R S \to M \otimes_R S) =
\bigoplus_{i \in I'} Q_i$가 되는 $M$의 가산 생성 부분가군 $N'$가
존재한다.
\end{lemma}

\begin{proof}
$N'_0 = N$이라 하자. $\ell = 0, 1, 2, \ldots$에 대해 가산 생성
부분가군 $N'_{\ell} \subset M$들의 증가열을 귀납적으로 구성한다.
$I'_{\ell}$을 $\Im(N'_{\ell} \otimes_R S \to M \otimes_R S)$의
$Q_i$ 위로의 사영이 영이 아닌 $i \in I$들의 집합이라 하면,
$\Im(N'_{\ell + 1} \otimes_R S \to M \otimes_R
S)$가 모든 $i \in I'_{\ell}$에 대해 $Q_i$를 포함하도록 구성한다.
$N'_\ell$로부터 $N'_{\ell + 1}$을 구성하려면
$i \in I'_{\ell}$인 (가산 개의) $Q_i$들의 합을 $Q$라 하고,
보조정리 \ref{algebra-lemma-lift-countably-generated-submodule}의 $P$를
택한 다음 $N'_{\ell + 1} = N'_{\ell} + P$로 놓는다.
$N'_{\ell}$들을 구성한 뒤에는
$N' = \bigcup_{\ell} N'_{\ell}$ 및
$I' = \bigcup_{\ell} I'_{\ell}$로 잡으면 된다.
\end{proof}

\begin{theorem}
\label{algebra-theorem-ffdescent-projectivity}
$R \to S$를 충실 평탄 환 사상이라 하고 $M$을 $R$-가군이라 하자.
$S$-가군 $M \otimes_R S$가 사영 가군이면 $M$도 사영 가군이다.
\end{theorem}

\begin{proof}
사영 가군들의 직합이고 따라서 사영 가군임을 보이기 위해
$M$의 Kaplansky d\'evissage를 구성한다. 정리
\ref{algebra-theorem-projective-direct-sum}에 의해
$M \otimes_R S = \bigoplus_{i \in I} Q_i$를 가산 생성
$S$-가군 $Q_i$들의 직합으로 쓸 수 있다. $M$ 위의 정렬순서를 하나
택한다. 초한 재귀를 사용하여 각 순서수 $\alpha$에 대해 $M$의
부분가군 $M_{\alpha}$들의 증가족을 정의하되,
$M_{\alpha} \otimes_R S$가 $Q_i$들 중 일부의 직합이 되게 한다.

\medskip\noindent
$\alpha = 0$이면 $M_0 = 0$으로 놓는다. $\alpha$가 극한
순서수이고 모든 $\beta < \alpha$에 대해 $M_{\beta}$가
정의되었다면 $M_\alpha = \bigcup_{\beta < \alpha} M_{\beta}$로
정의한다. 각 $\beta < \alpha$에 대해
$M_{\beta} \otimes_R S$가 $Q_i$들 중 일부의 직합이므로
$M_{\alpha} \otimes_R S$도 그러하다. $\alpha + 1$이 후계
순서수이고 $M_{\alpha}$가 정의되었다면 다음과 같이
$M_{\alpha + 1}$을 정의한다. $M_{\alpha} = M$이면
$M_{\alpha +1} = M$으로 놓는다. 그렇지 않으면 고정한
정렬순서에 관해 $x \notin M_{\alpha}$인 가장 작은 $x \in M$을
택한다. $S$가 $R$ 위에서 평탄하므로
$(M/M_{\alpha}) \otimes_R S = M
\otimes_R S/M_{\alpha} \otimes_R S$이고,
$M_{\alpha} \otimes_R S$가 일부 $Q_i$들의 직합이므로
$(M/M_{\alpha}) \otimes_R S$도 그러하다. 보조정리
\ref{algebra-lemma-adapted-submodule}에 의해, $M/M_{\alpha}$에서 $x$의
상을 포함하고 $P \otimes_R S$가 일부 $Q_i$들의 직합이 되는
$M/M_{\alpha}$의 가산 생성 $R$-부분가군 $P$를 찾을 수 있다.
% P01-E0471: frozen source sends this image to M tensor S rather than
% (M/M_alpha) tensor S; the frozen formula is preserved below.
(여기서 이 가군은 $S$가 $R$ 위에서 평탄하므로
$\Im(P \otimes_R S \to M \otimes_R S)$와 같다.)
$M \otimes_R S = \bigoplus_{i \in I} Q_i$가 사영 가군이고 사영성은
직합인자에 내려가므로 $P \otimes_R S$도 사영 가군이다. 따라서
보조정리 \ref{algebra-lemma-ffdescent-countable-projectivity}에 의해
$P$도 사영 가군이다. 마지막으로 $M_{\alpha + 1}$을 $M$ 안에서
$P$의 역상으로 정의한다. 그러면
$M_{\alpha + 1}/M_{\alpha} = P$는 가산 생성 사영 가군이다. 특히
$M_{\alpha + 1}/M_{\alpha}$의 사영성 때문에 열
$0 \to M_{\alpha} \to M_{\alpha + 1} \to
M_{\alpha + 1}/M_{\alpha} \to 0$이 분할되므로
$M_{\alpha}$는 $M_{\alpha + 1}$의 직합인자이다.

\medskip\noindent
% P01-E0472: frozen English says “transfinite induction on M” although
% the induction runs over the ordinal stages.
$M$에 대한 초한 귀납법(구성에서 $M_{\alpha}$에 들어 있지 않은
가장 작은 $x \in M$을 $M_{\alpha + 1}$에 넣었다는 사실을
사용한다)에 의해 각 $x \in M$은 어떤 $M_{\alpha}$에 들어간다.
% P01-E0473: frozen source reuses S, already the target ring, as an ordinal.
따라서 충분히 큰 어떤 순서수 $S$가 존재하여 모든 $x \in M$에
대해 $x \in M_{\alpha}$가 되는 $\alpha \in S$가 존재한다.
이는 $(M_{\alpha})_{\alpha \in S}$가 $M$의 Kaplansky
d\'evissage의 성질 (1)을 만족한다는 뜻이다. 다른 성질들은 구성상
명백하다. 따라서
$M = \bigoplus_{\alpha + 1 \in S} M_{\alpha + 1}/M_{\alpha}$이다.
각 $M_{\alpha + 1}/M_{\alpha}$는 구성상 사영 가군이므로
$M$도 사영 가군이다.
\end{proof}


\section{완비화}
\label{algebra-section-completion}

\noindent
$R$을 환, $I$를 아이디얼이라 하자. {\it $I$에 관한 $R$의
완비화}를 다음 극한으로 정의한다.
$$
R^\wedge = \lim_n R/I^n.
$$
$R^\wedge$의 원소는 모든 $n$에 대해
$f_n \equiv f_{n + 1} \bmod I^n$을 만족하는 원소
$f_n \in R/I^n$들의 열로 주어진다. $R^\wedge$를 $R$-대수로
간주한다. 마찬가지로 $M$이 $R$-가군이면 {\it $I$에 관한
$M$의 완비화}를 다음 극한으로 정의한다.
$$
M^\wedge = \lim_n M/I^nM.
$$
$M^\wedge$의 원소는 모든 $n$에 대해
$m_n \equiv m_{n + 1} \bmod I^nM$을 만족하는 원소
$m_n \in M/I^nM$들의 열로 주어진다. $M^\wedge$를
$R^\wedge$-가군으로 간주한다. 이 설명에서 항상 표준 사상
$$
M \longrightarrow M^\wedge
\quad\text{and}\quad
M \otimes_R R^\wedge \longrightarrow M^\wedge.
$$
이 있음을 알 수 있다. 또한 가군 사상 $\varphi : M \to N$이
주어지면 완비화 사이의 유도 사상
$\varphi^\wedge : M^\wedge \to N^\wedge$을 얻고, 이 사상은 그림
$$
\xymatrix{
M \ar[r] \ar[d] & N \ar[d] \\
M^\wedge \ar[r] & N^\wedge
}
$$
을 가환하게 한다. 일반적으로 완비화는 완전 함자가 아니다.
Examples, Section \ref{examples-section-completion-not-exact}를 보라.
먼저 몇 가지 긍정적인 결과를 제시한다.

\begin{lemma}
\label{algebra-lemma-completion-generalities}
$R$을 환, $I \subset R$을 아이디얼이라 하자.
$\varphi : M \to N$을 $R$-가군의 사상이라 하자.
\begin{enumerate}
\item $M/IM \to N/IN$이 전사이면 $M^\wedge \to N^\wedge$도
전사이다.
\item $M \to N$이 전사이면 $M^\wedge \to N^\wedge$도 전사이다.
\item $0 \to K \to M \to N \to 0$이 $R$-가군의 짧은 완전열이고
$N$이 평탄하면
$0 \to K^\wedge \to M^\wedge \to N^\wedge \to 0$도 짧은
완전열이다.
\item 임의의 유한 $R$-가군 $M$에 대해 사상
$M \otimes_R R^\wedge \to M^\wedge$은 전사이다.
\end{enumerate}
\end{lemma}

\begin{proof}
$M/IM \to N/IN$이 전사라고 가정하자. Nakayama 보조정리에 의해
각 $n \geq 1$에 대해 사상 $M/I^nM \to N/I^nN$은 전사이다.
좀 더 정확히 말하면, 환 $R/I^n$과 멱영 아이디얼 $I/I^n$ 위의
사상 $M/I^nM \to N/I^nN$에 보조정리 \ref{algebra-lemma-NAK}의 (11)을
적용한다. $K_n = \{x \in M \mid \varphi(x) \in I^nN\}$으로
놓자. 그러면 짧은 완전열
$$
0 \to K_n/I^nM \to M/I^nM \to N/I^nN \to 0
$$
을 얻는다. 표준 사상
$K_{n + 1}/I^{n + 1}M \to K_n/I^nM$이 전사임을 주장한다.
실제로 $x \in K_n$이면 $z_j \in I^n$, $n_j \in N$인 어떤
원소들에 대해 $\varphi(x) = \sum z_j n_j$로 쓴다. 가정에 의해
$m_j \in M$, $z_{jk} \in I$, $n_{jk} \in N$인 어떤 원소들에 대해
$n_j = \varphi(m_j) + \sum z_{jk}n_{jk}$로 쓸 수 있다. 따라서
$$
\varphi(x - \sum z_j m_j) = \sum z_jz_{jk} n_{jk}.
$$
이는 $x' = x - \sum z_j m_j \in K_{n + 1}$이
$x \bmod I^nM$으로 사상된다는 뜻이고, 주장을 증명한다. 이제
위의 짧은 완전열들의 역계에 보조정리
\ref{algebra-lemma-Mittag-Leffler}를 적용하여 (1)을 얻는다.
(2)는 (1)의 특수한 경우이다.
(3)의 가정들이 성립하면 각 $n$에 대해
$$
0 \to K/I^nK \to M/I^nM \to N/I^nN \to 0
$$
은 보조정리 \ref{algebra-lemma-flat-tor-zero}에 의해 짧은 완전열이다.
따라서 보조정리 \ref{algebra-lemma-Mittag-Leffler}를 직접 적용하여
(3)을 얻는다.
(4)를 보이기 위해 생성원 $x_i \in M$, $i = 1, \ldots, n$을
택하자. 사상
$R^{\oplus n} \to M$, $(a_1, \ldots, a_n) \mapsto \sum a_ix_i$는
전사이다. 따라서 (2)에 의해
$(R^\wedge)^{\oplus n} \to M^\wedge$,
$(a_1, \ldots, a_n) \mapsto \sum a_ix_i$는 전사이다.
이로부터 (4)가 따른다.
\end{proof}

\begin{definition}
\label{algebra-definition-complete}
$R$을 환, $I \subset R$을 아이디얼, $M$을 $R$-가군이라 하자.
사상
$$
M \longrightarrow M^\wedge = \lim_n M/I^nM
$$
이 동형사상이면 $M$이 {\it $I$-진 완비}라고 한다\footnote{여기에는
$\bigcap I^nM = 0$이라는 조건이 포함된다.}. $R$이
$R$-가군으로서 $I$-진 완비이면 $R$이 {\it $I$-진 완비}라고
한다.
\end{definition}

\noindent
$R$-가군 $M$의 $I$에 관한 완비화가 $I$-진 완비라는 명제는
일반적으로 참이 아니다. 예는 Examples, Section
\ref{examples-section-noncomplete-completion}을 보라.
아이디얼이 유한 생성이면 완비화는 완비이다.

\begin{lemma}
\label{algebra-lemma-hathat-finitely-generated}
\begin{reference}
\cite[Theorem 15]{Matlis}. 여기에 제시한 간결한 증명은 2016년
11월 5일 Bjorn Poonen이 보낸 전자우편에서 가져왔다.
\end{reference}
$R$을 환, $I$를 $R$의 유한 생성 아이디얼, $M$을 $R$-가군이라
하자. 그러면
\begin{enumerate}
\item 완비화 $M^\wedge$은 $I$-진 완비이고,
\item 모든 $n \geq 1$에 대해
$I^nM^\wedge = \Ker(M^\wedge \to M/I^nM) = (I^nM)^\wedge$이다.
\end{enumerate}
특히 $R^\wedge$은 $I$-진 완비이고,
$I^nR^\wedge = (I^n)^\wedge$이며
$R^\wedge/I^nR^\wedge = R/I^n$이다.
\end{lemma}

\begin{proof}
$I$가 유한 생성이므로 $I^n$도 유한 생성이다. 그 생성원을
$f_1, \ldots, f_r$이라 하자. 전사
$(f_1, \ldots, f_r) : M^{\oplus r} \to I^n M$에 보조정리
\ref{algebra-lemma-completion-generalities}의 (2)를 적용하면 전사
$$
(M^\wedge)^{\oplus r} \xrightarrow{(f_1, \ldots, f_r)} (I^n M)^\wedge =
\lim_{m \geq n} I^n M/I^m M = \Ker(M^\wedge \to M/I^n M).
$$
을 얻는다. 한편 사상
$(f_1, \ldots, f_r) : (M^\wedge)^{\oplus r} \to M^\wedge$의
상은 $I^n M^\wedge$이다. 따라서
$M^\wedge / I^n M^\wedge \simeq M/I^n M$이다. 역극한을 취하면
$(M^\wedge)^\wedge \simeq M^\wedge$을 얻는다. 즉,
$M^\wedge$은 $I$-진 완비이다.
\end{proof}

\begin{lemma}
\label{algebra-lemma-completion-differ-by-torsion}
$R$을 환, $I \subset R$을 아이디얼이라 하자.
$0 \to M \to N \to Q \to 0$을 $R$-가군의 완전열이라 하고,
$Q$가 $I$의 어떤 거듭제곱에 의해 소멸한다고 하자.
그러면 완비화는 완전열
$0 \to M^\wedge \to N^\wedge \to Q \to 0$을 만든다.
\end{lemma}

\begin{proof}
$I^c Q = 0$이라 하자. 그러면 $n \geq c$에 대해
$Q/I^nQ = Q$이다. 한편 $n \geq c$에 대해
$I^nM \subset M \cap I^nN \subset I^{n - c}M$임은 명백하다.
따라서 $M^\wedge = \lim M/(M \cap I^n N)$이다.
$n \geq c$인 완전열의 계
$$
0 \to M/(M \cap I^n N) \to N/I^n N \to Q \to 0
$$
에 보조정리 \ref{algebra-lemma-Mittag-Leffler}를 적용하면 결론을 얻는다.
\end{proof}

\begin{lemma}
\label{algebra-lemma-hathat}
\begin{reference}
Lenstra와 de Smit의 미출판 노트에서 가져왔다.
\end{reference}
$R$을 환, $I \subset R$을 아이디얼, $M$을 $R$-가군이라 하자.
$K_n = \Ker(M^\wedge \to M/I^nM)$으로 표기하자.
$M^\wedge$이 $I$-진 완비일 필요충분조건은 모든 $n \geq 1$에
대해 $K_n$이 $I^nM^\wedge$과 같은 것이다.
\end{lemma}

\begin{proof}
가군 $I^n M^\wedge$은 $K_n$에 포함된다. 따라서 각 $n \geq 1$에
대해 표준 완전열
$$
0 \to K_n/I^nM^\wedge \to M^\wedge/I^nM^\wedge \to M/I^nM \to 0.
$$
이 있다. $I^nM^\wedge$이 $I^nM/I^{n + 1}M$ 위로 전사되므로
$K_{n + 1} + I^n M^\wedge = K_n$임을 알 수 있다. 따라서 역계
$\{K_n/I^n M^\wedge\}_{n \geq 1}$의 전이 사상들은 전사이다.
보조정리 \ref{algebra-lemma-Mittag-Leffler}에 의해 짧은 완전열
$$
0 \to
\lim_n K_n/I^n M^\wedge \to
(M^\wedge)^\wedge \to
M^\wedge \to 0
$$
이 있음을 알 수 있다. 따라서 $M^\wedge$이 완비일 필요충분조건은
모든 $n \geq 1$에 대해 $K_n/I^n M^\wedge = 0$인 것이다.
\end{proof}

\begin{lemma}
\label{algebra-lemma-radical-completion}
$R$을 환, $I \subset R$을 아이디얼이라 하고
$R^\wedge = \lim R/I^n$으로 놓자.
\begin{enumerate}
\item $R/I$의 단원으로 사상되는 $R^\wedge$의 모든 원소는
단원이다.
\item $1 + I$의 모든 원소는 $R^\wedge$의 가역원으로 사상된다.
\item $1 + IR^\wedge$의 모든 원소는 $R^\wedge$에서 가역이다.
\item 아이디얼 $IR^\wedge$과 $\Ker(R^\wedge \to R/I)$는
$R^\wedge$의 Jacobson 근기에 포함된다.
\end{enumerate}
\end{lemma}

\begin{proof}
$x \in R^\wedge$이 $R/I$의 단원 $x_1$으로 사상된다고 하자.
보조정리 \ref{algebra-lemma-locally-nilpotent-unit}에 의해 모든 $n$에
대해 $x$는 $R/I^n$의 단원 $x_n$으로 사상된다. 따라서
$y = (x_n^{-1}) \in \lim R/I^n = R^\wedge$은 $x$의 역원이다.
(2)와 (3)은 (1)에서 즉시 따른다. (4)는 (1)과 보조정리
\ref{algebra-lemma-contained-in-radical}에서 따른다.
\end{proof}

\begin{lemma}
\label{algebra-lemma-when-surjective-to-completion}
% P01-E0474: frozen source never introduces M as an A-module.
$A$를 환이라 하자. $I = (f_1, \ldots, f_r)$를 유한 생성
아이디얼이라 하자. 각 $i$에 대해
$M \to \lim M/f_i^nM$이 전사이면 $M \to \lim M/I^nM$도
전사이다.
\end{lemma}

\begin{proof}
$$
I^n \supset (f_1^n, \ldots, f_r^n) \supset I^{rn}
$$
이므로
$$
\lim M/I^nM = \lim M/(f_1^n, \ldots, f_r^n)M
$$
임에 유의하라.
$\lim M/(f_1^n, \ldots, f_r^n)M$의 원소 $\xi$는 기호적으로
$$
\xi = \sum\nolimits_{n \geq 0} \sum\nolimits_i f_i^n x_{n, i}
$$
로 쓸 수 있으며, 여기서 $x_{n, i} \in M$이다.
$M \to \lim M/f_i^nM$이 전사이면
$\lim M/f_i^nM$에서 $\sum x_{n, i} f_i^n$으로 사상되는
$x_i \in M$이 존재한다. 그러면 $x = \sum x_i$는
$\lim M/I^nM$에서 $\xi$로 사상된다.
\end{proof}

\begin{lemma}
\label{algebra-lemma-complete-by-sub}
% P01-E0475: frozen source never introduces M as an A-module.
$A$를 환이라 하고 $I \subset J \subset A$를 아이디얼이라 하자.
$M$이 $J$-진 완비이고 $I$가 유한 생성이면 $M$은 $I$-진
완비이다.
\end{lemma}

\begin{proof}
$M$이 $J$-진 완비이고 $I$가 유한 생성이라고 가정하자.
$\bigcap J^nM = 0$이므로 $\bigcap I^nM = 0$이다.
보조정리 \ref{algebra-lemma-when-surjective-to-completion}에 의해
$I$가 한 원소로 생성되는 경우에
$M \to \lim M/I^nM$의 전사성만 증명하면 충분하다.
$I = (f)$라 하자. $x_{n + 1} - x_n \in f^nM$인 $x_n \in M$이
주어졌다고 하자. 모든 $n$에 대해 $x_n - x \in f^nM$이 되는
$x \in M$이 존재함을 보여야 한다.
$x_{n + 1} - x_n \in J^nM$이고 $M$이 $J$-진 완비이므로
$x_n - x \in J^nM$이 되는 원소 $x \in M$이 존재한다.
$x_n$을 $x_n - x$로 바꾸어 $x_n \in J^nM$이라고 가정해도 된다.
증명을 끝내기 위해 이 조건에서 $x_n \in I^nM$이 따름을 보인다.
실제로 $x_n - x_{n + 1} = f^nz_n$으로 쓰면
$$
x_n = f^n(z_n + fz_{n + 1} + f^2z_{n + 2} + \ldots)
$$
이다. $f^c \in J^c$이므로 합
$z_n + fz_{n + 1} + f^2z_{n + 2} + \ldots$는 $M$에서
수렴한다. 그 부분합은 $x_n - x_{n + c}$이고
$x_{n + c} \in J^{n + c}M$이므로 합
$f^n(z_n + fz_{n + 1} + f^2z_{n + 2} + \ldots)$는 $M$에서
$x_n$으로 수렴한다.
\end{proof}

\begin{lemma}
\label{algebra-lemma-change-ideal-completion}
$R$을 환, $I$, $J$를 $R$의 아이디얼들이라 하자.
$I^c \subset J$ 및 $J^d \subset I$를 만족하는 정수
$c, d > 0$이 존재한다고 가정하자. 그러면 임의의 $R$-가군에 대해
$I$에 관한 완비화와 $J$에 관한 완비화가 일치한다.
특히 $R$-가군 $M$이 $I$-진 완비일 필요충분조건은
$J$-진 완비인 것이다.
\end{lemma}

\begin{proof}
사상들의 계
$M/I^nM \to M/J^{\lfloor n/c \rfloor}M$과 사상들의 계
$M/J^mM \to M/I^{\lfloor m/d \rfloor}M$을 생각하여 두 완비화
사이의 서로 역인 사상들을 얻는다.
\end{proof}

\begin{lemma}
\label{algebra-lemma-quotient-complete}
$R$을 환, $I$를 $R$의 아이디얼이라 하자. $M$을 $I$-진 완비
$R$-가군이라 하고 $K \subset M$을 $R$-부분가군이라 하자.
다음은 서로 동치이다.
\begin{enumerate}
\item $K = \bigcap (K + I^nM)$이다.
\item $M/K$는 $I$-진 완비이다.
\end{enumerate}
\end{lemma}

\begin{proof}
$N = M/K$로 놓자. 보조정리
\ref{algebra-lemma-completion-generalities}에 의해 사상
$M = M^\wedge \to N^\wedge$은 전사이다. 따라서
$N \to N^\wedge$은 전사이다. $N \to N^\wedge$의 핵이
가군 $\bigcap (K + I^nM) / K$임은 쉽게 알 수 있다.
\end{proof}

\begin{lemma}
\label{algebra-lemma-when-finite-module-complete-over-complete-ring}
$R$을 환, $I$를 $R$의 아이디얼, $M$을 $R$-가군이라 하자.
(a) $R$이 $I$-진 완비이고, (b) $M$이 유한 $R$-가군이며,
(c) $\bigcap I^nM = (0)$이면 $M$은 $I$-진 완비이다.
\end{lemma}

\begin{proof}
보조정리 \ref{algebra-lemma-completion-generalities}에 의해 사상
$M = M \otimes_R R = M \otimes_R R^\wedge \to M^\wedge$은
전사이다. 이 사상의 핵은 $\bigcap I^nM$이고 가정에 의해
영이다. 따라서 $M \cong M^\wedge$이고 $M$은 완비이다.
\end{proof}

\begin{lemma}
\label{algebra-lemma-finite-over-complete-ring}
$R$을 환, $I \subset R$을 아이디얼, $M$을 $R$-가군이라 하자.
다음을 가정하자.
\begin{enumerate}
\item $R$은 $I$-진 완비이다.
\item $\bigcap_{n \geq 1} I^nM = (0)$이다.
\item $M/IM$은 유한 $R/I$-가군이다.
\end{enumerate}
그러면 $M$은 유한 $R$-가군이다.
\end{lemma}

\begin{proof}
% P01-E0476/P01-E0477: frozen English has the wrong article before R/I-module
% and omits “by” in the naming construction.
$x_1, \ldots, x_n \in M$을 그 상들이 $M/IM$에서 $R/I$-가군으로서
$M/IM$을 생성하는 원소들이라 하자. $x_1, \ldots, x_n$이
생성하는 $R$-부분가군을 $M' \subset M$이라 하자. 보조정리
\ref{algebra-lemma-completion-generalities}에 의해 사상
$(M')^\wedge \to M^\wedge$은 전사이다.
$\bigcap I^nM = 0$이므로 특히 $\bigcap I^nM' = (0)$이다.
따라서 보조정리
\ref{algebra-lemma-when-finite-module-complete-over-complete-ring}에 의해
$M'$은 완비이고, $M' \to M^\wedge$이 전사라고 결론 내린다.
마지막으로 $M \to M^\wedge$의 핵은
$\bigcap I^nM = (0)$과 같으므로 영이다.
따라서 $M \cong M' \cong M^\wedge$은 유한 생성이다.
\end{proof}


\section{뇌터 환에서의 완비화}
\label{algebra-section-completion-noetherian}

\noindent
이 절에서는 뇌터 환의 아이디얼에 관한 완비화를 다룬다.

\begin{lemma}
\label{algebra-lemma-completion-tensor}
$I$를 뇌터 환 $R$의 아이디얼이라 하자.
% P01-E0478: frozen English omits “by” in the notation declaration.
${}^\wedge$으로 $I$에 관한 완비화를 표기하자.
\begin{enumerate}
\item $K \to N$이 유한 $R$-가군들의 단사 사상이면 완비화 사이의
사상 $K^\wedge \to N^\wedge$도 단사이다.
\item $0 \to K \to N \to M \to 0$이 유한 $R$-가군들의 짧은
완전열이면 $0 \to K^\wedge \to N^\wedge \to M^\wedge \to 0$도
짧은 완전열이다.
\item $M$이 유한 $R$-가군이면
$M^\wedge = M \otimes_R R^\wedge$이다.
\end{enumerate}
\end{lemma}

\begin{proof}
$M = N/K$로 놓으면 (1)이 (2)에서 따름을 알 수 있다.
$0 \to K \to N \to M \to 0$을 (2)에서와 같이 잡자. 각 $n$에
대해 짧은 완전열
$$
0 \to K/(I^nN \cap K) \to N/I^nN \to M/I^nM \to 0.
$$
을 얻는다. 보조정리 \ref{algebra-lemma-Mittag-Leffler}에 의해 완전열
$$
0 \to \lim K/(I^nN \cap K) \to N^\wedge \to M^\wedge \to 0.
$$
을 얻는다. Artin--Rees 보조정리 \ref{algebra-lemma-Artin-Rees}에 의해
$n \geq c$일 때
$I^nK \subset I^n N \cap K \subset I^{n-c} K$가 되도록 $c$를
택할 수 있다. 따라서
$K^\wedge = \lim K/I^nK = \lim K/(I^nN \cap K)$이고 (2)가
성립한다고 결론 내린다.

\medskip\noindent
$M$을 (3)에서와 같이 잡고
$0 \to K \to R^{\oplus t} \to M \to 0$을 $M$의 표시라 하자.
가환 그림
$$
\xymatrix{
&
K \otimes_R R^\wedge \ar[r] \ar[d] &
R^{\oplus t} \otimes_R R^\wedge \ar[r] \ar[d] &
M \otimes_R R^\wedge \ar[r] \ar[d] &
0 \\
0 \ar[r] &
K^\wedge \ar[r] &
(R^{\oplus t})^\wedge \ar[r] &
M^\wedge \ar[r] & 0
}
$$
을 얻는다. 윗행은 Section \ref{algebra-section-flat}에 의해 완전열이다.
아랫행은 (2)에 의해 완전열이다. 보조정리
\ref{algebra-lemma-completion-generalities}에 의해 수직 화살표들은
전사이다. 가운데 수직 화살표는 동형사상이다. 뱀 보조정리
\ref{algebra-lemma-snake}에 의해 (3)이 성립한다고 결론 내린다.
\end{proof}

\begin{lemma}
\label{algebra-lemma-completion-flat}
$I$를 뇌터 환 $R$의 아이디얼이라 하자.
% P01-E0479: frozen English repeats the malformed notation declaration.
${}^\wedge$으로 $I$에 관한 완비화를 표기하자.
\begin{enumerate}
\item 환 사상 $R \to R^\wedge$은 평탄하다.
\item 함자 $M \mapsto M^\wedge$은 유한 생성 $R$-가군들의
범주에서 완전하다.
\end{enumerate}
\end{lemma}

\begin{proof}
$J$를 $R$의 임의의 아이디얼이라 하고
$J \otimes_R R^\wedge \to R \otimes_R R^\wedge = R^\wedge$을
생각하자. 보조정리 \ref{algebra-lemma-completion-tensor}에 따르면 이
사상은 $J^\wedge \to R^\wedge$과 동일시되고,
$J^\wedge \to R^\wedge$은 단사이다. (1)은 보조정리
\ref{algebra-lemma-flat}에서 따른다. (2)는 보조정리
\ref{algebra-lemma-completion-tensor}의 (2)를 다시 표현한 것이다.
\end{proof}

\begin{lemma}
\label{algebra-lemma-completion-faithfully-flat}
$I$를 뇌터 환 $R$의 아이디얼이라 하자. $R^\wedge$으로
$I$에 관한 $R$의 완비화를 표기하자. $I$가 $R$의 Jacobson 근기에
포함되면 환 사상 $R \to R^\wedge$은 충실 평탄이다. 특히
$(R, \mathfrak m)$이 뇌터 국소환이면 완비화
$\lim_n R/\mathfrak m^n$은 충실 평탄이다.
\end{lemma}

\begin{proof}
보조정리 \ref{algebra-lemma-completion-flat}에 의해 이 사상은 평탄하다.
마지막 사상이 사영 사상 $R^\wedge \to R/I$인 합성
$R \to R^\wedge \to R/I$을 보면, $R$의 모든 극대 아이디얼이
$\Spec(R^\wedge) \to \Spec(R)$의 상에 들어감을 알 수 있다.
따라서 보조정리 \ref{algebra-lemma-ff}에 의해 이 사상은 충실 평탄이다.
\end{proof}

\begin{lemma}
\label{algebra-lemma-completion-complete}
$R$을 뇌터 환, $I$를 $R$의 아이디얼, $M$을 $R$-가군이라 하자.
$I$에 관한 $M$의 완비화 $M^\wedge$은 $I$-진 완비이고,
$I^n M^\wedge = (I^nM)^\wedge$이며
$M^\wedge/I^nM^\wedge = M/I^nM$이다.
\end{lemma}

\begin{proof}
$I$가 유한 생성 아이디얼이므로 이는 보조정리
\ref{algebra-lemma-hathat-finitely-generated}의 특수한 경우이다.
\end{proof}

\begin{lemma}
\label{algebra-lemma-completion-Noetherian}
$I$를 환 $R$의 아이디얼이라 하자. 다음을 가정하자.
\begin{enumerate}
\item $R/I$는 뇌터 환이다.
\item $I$는 유한 생성이다.
\end{enumerate}
그러면 $I$에 관한 $R$의 완비화 $R^\wedge$은
$IR^\wedge$에 관해 완비인 뇌터 환이다.
\end{lemma}

\begin{proof}
보조정리 \ref{algebra-lemma-hathat-finitely-generated}에 의해
$R^\wedge$은 $I$-진 완비이다. 따라서 $IR^\wedge$-진
완비이기도 하다. $R^\wedge/IR^\wedge = R/I$가 뇌터 환이므로,
$R$을 $R^\wedge$으로 바꾼 뒤에는 (1), (2)뿐 아니라 $R$도
$I$-진 완비라고 가정할 수 있다.

\medskip\noindent
$f_1, \ldots, f_t$를 $I$의 생성원들이라 하자. 그러면 $T_i$를
원소 $\overline{f}_i \in I/I^2$으로 보내는 환의 전사
$R/I[T_1, \ldots, T_t] \to \bigoplus I^n/I^{n + 1}$이 있다.
따라서 $\bigoplus I^n/I^{n + 1}$은 뇌터 환이다.
$J \subset R$을 아이디얼이라 하자. 다음 아이디얼을 생각하자.
$$
\bigoplus J \cap I^n/J \cap I^{n + 1} \subset \bigoplus I^n/I^{n + 1}.
$$
$\overline{g}_1, \ldots, \overline{g}_m$을 이 아이디얼의
생성원들이라 하자. $\overline{g}_j$를 차수 $d_j$인 동차원소로
택할 수 있고, $\overline{g}_j \in J \cap I^{d_j}/J \cap
I^{d_j + 1}$으로 사상되는
$g_j \in J \cap I^{d_j}$를 택할 수 있다.
$g_1, \ldots, g_m$이 $J$를 생성한다고 주장한다.

\medskip\noindent
$x \in J \cap I^n$이라 하자.
$x - \sum a_j g_j \in J \cap I^{n + 1}$이 되도록 하는
$a_j \in I^{\max(0, n - d_j)}$가 존재한다. 그 이유는
% P01-E0480: frozen source uses unrestricted negative ideal exponents;
% preserve the formula and record the source defect only in the sidecar.
$g_1, \ldots, g_m$의 선택에 의해
$J \cap I^n/J \cap I^{n + 1}$이
$\sum \overline{g}_j I^{n - d_j}/I^{n - d_j + 1}$와 같기 때문이다.
따라서 $x \in J$에서 시작하여
$a_{j, n} \in I^{\max(0, n - d_j)}$이고
$$
x =
\sum\nolimits_{n = 0, \ldots, N}
\sum\nolimits_{j = 1, \ldots, m} a_{j, n} g_j \bmod I^{N + 1}
$$
을 만족하는 벡터열
$(a_{1, n}, \ldots, a_{m, n})_{n \geq 0}$을 찾을 수 있다.
$A_j = \sum_{n \geq 0} a_{j, n}$으로 놓으면 $R$이 완비이므로
$x = \sum A_j g_j$임을 알 수 있다. 따라서 $J$는 유한 생성이고
결론이 따른다.
\end{proof}

\begin{lemma}
\label{algebra-lemma-completion-Noetherian-Noetherian}
$R$을 뇌터 환, $I$를 $R$의 아이디얼이라 하자.
$I$에 관한 $R$의 완비화 $R^\wedge$은 뇌터 환이다.
\end{lemma}

\begin{proof}
이는 보조정리 \ref{algebra-lemma-completion-Noetherian}의 결과이다.
다음과 같이 직접 볼 수도 있다. $f_1, \ldots, f_n$을 $I$의
생성원들로 택하자. 사상
$$
R[[x_1, \ldots, x_n]] \longrightarrow R^\wedge,
\quad
x_i \longmapsto f_i.
$$
을 생각하자. 이는 잘 정의된 전사 환 사상이다(세부사항은
생략한다). $R[[x_1, \ldots, x_n]]$이 뇌터 환이므로(보조정리
\ref{algebra-lemma-Noetherian-power-series}를 보라) 결론이 따른다.
\end{proof}

\noindent
$R \to S$가 국소환 $(R, \mathfrak m)$과 $(S, \mathfrak n)$ 사이의
국소 준동형이라고 하자. $S^\wedge$을 $\mathfrak n$에 관한 $S$의
완비화라 하자. 일반적으로 $S^\wedge$은 $S$의
$\mathfrak m$-진 완비화가 아니다. 어떤 $t \geq 1$에 대해
$\mathfrak n^t \subset \mathfrak mS$이면 보조정리
\ref{algebra-lemma-change-ideal-completion}에 의해
$S^\wedge = \lim S/\mathfrak m^nS$가 성립한다. 어떤 경우에는
이로부터 더 나아가 $S^\wedge$이 $R^\wedge$ 위에서 유한임이
따른다.

\begin{lemma}
\label{algebra-lemma-finite-after-completion}
$R \to S$가 국소환 $(R, \mathfrak m)$과 $(S, \mathfrak n)$ 사이의
국소 준동형이라 하자. $R^\wedge$, resp.\ $S^\wedge$을
$\mathfrak m$, resp.\ $\mathfrak n$에 관한 $R$, resp.\ $S$의
완비화라 하자. $\mathfrak m$과 $\mathfrak n$이 유한 생성이고
$\dim_{\kappa(\mathfrak m)} S/\mathfrak mS < \infty$이면
\begin{enumerate}
\item $S^\wedge$은 $S$의 $\mathfrak m$-진 완비화와 같고,
\item $S^\wedge$은 유한 $R^\wedge$-가군이다.
\end{enumerate}
\end{lemma}

\begin{proof}
$R \to S$가 국소환 사상이므로
$\mathfrak mS \subset \mathfrak n$이다.
가정 $\dim_{\kappa(\mathfrak m)} S/\mathfrak mS < \infty$에서
$S/\mathfrak mS$가 아르틴 환임이 따른다. 보조정리
\ref{algebra-lemma-finite-dimensional-algebra}를 보라.
% P01-E0481: frozen English omits the subject of “Hence has dimension 0”.
따라서 보조정리 \ref{algebra-lemma-Noetherian-dimension-0}에 의해 이 환의
차원은 $0$이고, 그러므로
$\mathfrak n = \sqrt{\mathfrak mS}$이다.
이 사실과 $\mathfrak n$이 유한 생성이라는 사실에 의해 어떤
$t \geq 1$에 대해 $\mathfrak n^t \subset \mathfrak mS$이다.
보조정리 \ref{algebra-lemma-change-ideal-completion}에 의해
$S^\wedge$을 $S$의 $\mathfrak m$-진 완비화와 동일시할 수 있다.
$\mathfrak m$이 유한 생성이므로 보조정리
\ref{algebra-lemma-hathat-finitely-generated}에 의해 $S^\wedge$과
$R^\wedge$은 $\mathfrak m$-진 완비이다. 이제
$S^\wedge$을 $R^\wedge$-가군으로 보고 보조정리
\ref{algebra-lemma-finite-over-complete-ring}를 적용하면 결론을 얻는다.
\end{proof}

\begin{lemma}
\label{algebra-lemma-completion-finite-extension}
$R$을 뇌터 환, $R \to S$를 유한 환 사상이라 하자.
$\mathfrak p \subset R$을 소아이디얼이라 하고,
$\mathfrak q_1, \ldots, \mathfrak q_m$을 $\mathfrak p$ 위에 놓인
$S$의 소아이디얼들이라 하자(보조정리
\ref{algebra-lemma-finite-finite-fibres}). 그러면
$$
R_\mathfrak p^\wedge \otimes_R S =
(S_\mathfrak p)^\wedge =
S_{\mathfrak q_1}^\wedge \times \ldots \times S_{\mathfrak q_m}^\wedge
$$
이다. 여기서 $(S_\mathfrak p)^\wedge$은 $\mathfrak p$에 관한
완비화이고, 국소환 $R_\mathfrak p$와
$S_{\mathfrak q_i}$는 각각의 극대 아이디얼에 관해 완비화한다.
\end{lemma}

\begin{proof}
$R$을 국소화 $R_\mathfrak p$로 바꾸고, $S$를
$S_\mathfrak p = S \otimes_R R_\mathfrak p$로 바꾸어도 된다.
따라서 $R$이 국소 뇌터 환이고 $\mathfrak p = \mathfrak m$이 그
극대 아이디얼이라고 가정할 수 있다.
보조정리 \ref{algebra-lemma-finite-after-completion}에 의해
$\mathfrak q_iS_{\mathfrak q_i}$-진 완비화
$S_{\mathfrak q_i}^\wedge$은 $\mathfrak m$-진 완비화와 같다.
모든 $n \geq 1$에 대해 $S/\mathfrak m^nS$의 소아이디얼들은
$S$의 극대 아이디얼 $\mathfrak q_1, \ldots, \mathfrak q_m$과
일대일로 대응한다($R$ 위의 $S$에 대한 상승 정리, 보조정리
\ref{algebra-lemma-integral-going-up}를 보라). 따라서 보조정리
\ref{algebra-lemma-artinian-finite-length}에 의해
$S/\mathfrak m^nS = \prod S_{\mathfrak q_i}/\mathfrak m^nS_{\mathfrak q_i}$
이다(예를 들어 $S/\mathfrak m^nS$가 아르틴 환임을 보기 위해
명제 \ref{algebra-proposition-dimension-zero-ring}을 사용한다).
그러므로 $S$의 $\mathfrak m$-진 완비화 $S^\wedge$은
$\prod S_{\mathfrak q_i}^\wedge$과 같다. 마지막으로 보조정리
\ref{algebra-lemma-completion-tensor}에 의해
$R^\wedge \otimes_R S = S^\wedge$이다.
\end{proof}

\begin{lemma}
\label{algebra-lemma-split-completed-sequence}
$R$을 환, $I \subset R$을 아이디얼이라 하자.
$0 \to K \to P \to M \to 0$을 $R$-가군의 짧은 완전열이라 하자.
$M$이 $R$ 위에서 평탄하고 $M/IM$이 사영 $R/I$-가군이면
$I$-진 완비화들의 열
$$
0 \to K^\wedge \to P^\wedge \to M^\wedge \to 0
$$
은 분할 완전열이다.
\end{lemma}

\begin{proof}
$M$이 평탄하므로 각 열
$$
0 \to K/I^nK \to P/I^nP \to M/I^nM \to 0
$$
은 짧은 완전열이다. 보조정리
\ref{algebra-lemma-flat-tor-zero}를 보라. 또한 보조정리
\ref{algebra-lemma-completion-generalities}에 의해
$0 \to K^\wedge \to P^\wedge \to M^\wedge \to 0$은 짧은
완전열이다. $s_{n + 1} \bmod I^n = s_n$을 만족하는 분할 사상
$s_n : M/I^nM \to P/I^nP$들을 찾을 수 있음을 보이면 충분하다.
$n$에 대한 귀납법으로 이 $s_n$들을 구성한다.
$M/IM$이 사영 $R/I$-가군이므로 존재하는 임의의 분할 사상
$s_1$을 택한다. 어떤 $n > 0$에 대해 $s_n$이 주어졌다고 하자.
$P_{n + 1} = \{x \in P \mid x \bmod I^nP \in \Im(s_n)\}$로
놓자. 사상
$\pi : P_{n + 1}/I^{n + 1}P_{n + 1} \to M/I^{n + 1}M$은
전사이다(세부사항은 생략한다). 보조정리
\ref{algebra-lemma-lift-projective}에 의해
$M/I^{n + 1}M$은 사영 $R/I^{n + 1}$-가군이므로 $\pi$의 절단
$t : M/I^{n + 1}M \to P_{n + 1}/I^{n + 1}P_{n + 1}$을
택할 수 있다. $s_{n + 1}$을 $t$와 표준 사상
$P_{n + 1}/I^{n + 1}P_{n + 1} \to P/I^{n + 1}P$의 합성으로
놓으면 된다.
\end{proof}

\begin{lemma}
\label{algebra-lemma-complete-modulo-nilpotent}
$A$를 뇌터 환, $I, J \subset A$를 아이디얼들이라 하자.
$A$가 $I$-진 완비이고 $A/I$가 $J$-진 완비이면 $A$는
$J$-진 완비이다.
\end{lemma}

\begin{proof}
$B$를 $A$의 $(I + J)$-진 완비화라 하자. 보조정리
\ref{algebra-lemma-completion-flat}에 의해 $B/IB$는 $A/I$의 $J$-진
완비화이고, 따라서 가정에 의해 $A/I$와 동형이다. 또한 보조정리
\ref{algebra-lemma-complete-by-sub}에 의해 $B$는 $I$-진 완비이다.
그러므로 보조정리 \ref{algebra-lemma-finite-over-complete-ring}에 의해
$B$는 유한 $A$-가군이다. Nakayama 보조정리(보조정리
\ref{algebra-lemma-NAK})를 적용하되, $I$가 $A$의 Jacobson 근기에
속한다는 사실(보조정리 \ref{algebra-lemma-radical-completion})을
사용하면 $A \to B$가 전사임을 알 수 있다.
% P01-E0482: frozen English omits “that” in the Nakayama parenthesis.
보조정리 \ref{algebra-lemma-completion-flat}에 의해 $A \to B$는
평탄하다. $\Spec(B) \to \Spec(A)$의 상은 $V(I)$를 포함하고,
$I$가 $A$의 Jacobson 근기에 포함되므로 $A \to B$는 충실
평탄이다(보조정리 \ref{algebra-lemma-ff-rings}). 따라서 $A \to B$는
단사이다. 그러므로 $A$는 $I + J$에 관해 완비이고, 따라서 더욱
강하게 $J$에 관해 완비이다.
\end{proof}


\section{가군의 극한 취하기}
\label{algebra-section-limits}

\noindent
이 절에서는 가군들의 극한을 취할 때 어떤 일이 일어나는지 다룬다.

\begin{lemma}
\label{algebra-lemma-limit-complete-pre}
$I \subset A$를 어떤 환의 유한 생성 아이디얼이라 하자.
$(M_n)$을 $I^n M_n = 0$인 $A$-가군의 역계라 하자.
그러면 $M = \lim M_n$은 $I$-진 완비이다.
\end{lemma}

\begin{proof}
$M \to M/I^nM \to M_n$이 있다. 극한을 취하면
$M \to M^\wedge \to M$을 얻는다. 따라서 $M$은 $M^\wedge$의
직합인자이다. 보조정리
\ref{algebra-lemma-hathat-finitely-generated}에 의해 $M^\wedge$이
$I$-진 완비이므로 $M$도 그러하다.
\end{proof}

\begin{lemma}
\label{algebra-lemma-limit-complete}
$I \subset A$를 어떤 환의 유한 생성 아이디얼이라 하자.
$(M_n)$을 $M_n = M_{n + 1}/I^nM_{n + 1}$인 $A$-가군의
역계라 하자. $M = \lim M_n$으로 놓자. 그러면
$M/I^nM = M_n$이고 $M$은 $I$-진 완비이다.
\end{lemma}

\begin{proof}
보조정리 \ref{algebra-lemma-limit-complete-pre}에 의해 $M$은 $I$-진
완비이다. 전이 사상들이 전사이므로 사상 $M \to M_n$들도
전사이다. $N_n$을 정의하는 짧은 완전열들의 역계
$$
0 \to N_n \to M \to M_n \to 0
$$
을 생각하자. $M_n = M_{n + 1}/I^nM_{n + 1}$이므로 사상
$N_{n + 1} + I^nM \to N_n$은 전사이다. 따라서
$N_{n + 1}/(N_{n + 1} \cap I^{n + 1}M) \to
N_n/(N_n \cap I^nM)$은 전사이다. 짧은 완전열
$$
0 \to N_n/(N_n \cap I^nM) \to M/I^nM \to M_n \to 0
$$
의 역극한을 취하면 완전열
$$
0 \to \lim N_n/(N_n \cap I^nM) \to M^\wedge \to M
$$
을 얻는다. $M$은 $I$-진 완비이므로
$\lim N_n/(N_n \cap I^nM) = 0$이라고 결론 내린다. 전이
사상들의 전사성에 의해 모든 $n$에 대해
$N_n/(N_n \cap I^nM) = 0$이다. 따라서 원하는 대로
$M_n = M/I^nM$이다.
\end{proof}

\begin{lemma}
\label{algebra-lemma-finiteness-graded}
$A$를 뇌터 등급 환, $I \subset A_+$를 동차 아이디얼이라 하자.
$(N_n)$을 $N_n = N_{n + 1}/I^n N_{n + 1}$인 유한 등급
$A$-가군의 역계라 하자. 그러면 모든 $n$에 대해 등급 가군으로
$N_n = N/I^nN$이 되는 유한 등급 $A$-가군 $N$이 존재한다.
\end{lemma}

\begin{proof}
$r$과 $N_1$을 생성하는 차수 $d_1, \ldots, d_r$의 동차원소
$x_{1, 1}, \ldots, x_{1, r} \in N_1$을 택하자. 전이 사상들이
전사이므로 $x_{1, i}$를 올리는 동차원소들의 양립하는 계
$x_{n, i} \in N_n$을 택할 수 있다. 등급 Nakayama
보조정리(보조정리 \ref{algebra-lemma-graded-NAK})에 의해 $N_n$은
차수 $d_1, \ldots, d_r$에 놓인 원소
$x_{n, 1}, \ldots, x_{n, r}$에 의해 생성된다.
따라서 $m \leq n$이면 사상 $N_n \to N_n/I^m N_n$은
차수 $< \min(d_i) + m$에서 동형사상이다(그 차수들에서는
$I^mN_n$이 영이기 때문이다). 따라서 차수 $d$ 부분들의 역계
% P01-E0483: frozen stabilization includes an unjustified equality and
% can invoke N_0; preserve its indices and record the defect in the sidecar.
$$
\ldots
= N_{2 + d - \min(d_i), d}
= N_{1 + d - \min(d_i), d}
= N_{d - \min(d_i), d} \to N_{-1 + d - \min(d_i), d} \to \ldots
$$
은 표시된 것처럼 안정화한다. $N$을 그 $d$차 등급 부분이 이
안정화된 값인 등급 $A$-가군이라 하자. 특히 $N$에는 원소
$x_i = \lim x_{n, i}$가 있다. $x_i$들이 $N$을 생성한다고
주장한다. 임의의 $x \in N_d$는 $x_1, \ldots, x_r$의
선형결합이다. 실제로 $x_{d - \min(d_i), i}$들이
$N_{d - \min(d_i)}$을 생성하므로
$N_{d - \min(d_i), d}$에서 이를 확인할 수 있다.
마지막으로 독자는 전사
$N/I^nN \to N_n$이 동형사상임을 앞에서와 같이 각 차수에서
무슨 일이 일어나는지 확인하여 알 수 있다. 세부사항은 생략한다.
\end{proof}

\begin{lemma}
\label{algebra-lemma-daniel-litt}
$A$를 등급 환, $I \subset A_+$를 동차 아이디얼이라 하자.
$A' = \lim A/I^n$으로 표기하자. $(G_n)$을 $I^n$에 의해 소멸하는
등급 $A$-가군 $G_n$들의 역계라 하자. $M$을 등급 $A$-가군이라
하고 $\varphi_n : M \to G_n$을 양립하는 등급 $A$-가군 사상들의
계라 하자. 유도 사상
$$
\varphi : M \otimes_A A' \longrightarrow \lim G_n
$$
이 동형사상이면 모든 $d \in \mathbf{Z}$에 대해
$M_d \to \lim G_{n, d}$는 동형사상이다.
\end{lemma}

\begin{proof}
관례에 따라 등급 환은 차수 $\geq 0$에 놓이고, 등급 가군은
임의의 차수에서 영이 아닌 부분을 가질 수 있다. Section
\ref{algebra-section-graded}를 보라. $G_n$은 $I^n$에 의해 소멸하므로
$\lim G_n$이 $A'$ 위의 가군이고, 따라서 사상 $\varphi$가
존재한다. 다음 포함도 기억해 두면 유용하다.
$$
\bigoplus\nolimits_{d \in \mathbf{Z}} \lim G_{n, d} \subset
\lim G_n \subset
\prod\nolimits_{d \in \mathbf{Z}} \lim G_{n, d}
$$
여기서 아래첨자 ${\ }_d$는 $d$차 등급 부분을 뜻한다.

\medskip\noindent
단사성. $x \in M_d$라 하자. $x \mapsto 0$이 $\lim G_{n, d}$에서
성립하면 $M \otimes_A A'$에서 $x \otimes 1 = 0$이다.
그러면 $x \in M'$이고 $M' \otimes_A A'$에서
$x \otimes 1$이 영이 되는 유한 생성 부분가군
$M' \subset M$을 찾을 수 있다. $M'$이 차수
$d_1, \ldots, d_r$에 놓인 동차원소들에 의해 생성된다고 하자.
$n = d - \min(d_i) + 1$로 놓자. $A'$에서 $A/I^n$으로 가는
사상이 있고 $A \to A/I^n$은 차수 $\leq n - 1$에서
동형사상이므로,
% P01-E0484: frozen source transfers this coefficient-degree bound to
% M' without the generator shifts; preserve the stated bound.
$M' \to M' \otimes_A A'$은 차수 $\leq n - 1$에서 단사이다.
따라서 원하는 대로 $x = 0$이다.

\medskip\noindent
전사성. $y \in \lim G_{n, d}$라 하자.
$y$로 사상되는 $M \otimes_A A'$의 유한합
$\sum x_i \otimes f'_i$를 택하자. $x_i$가 동차이고 그 차수가
$d_i$라고 가정해도 된다. $A'$은 등급 환이 아니지만 등급 환
$A/I^nA$들의 극한이고, 더욱이 주어진 각 차수에서 전이
사상들은 결국 동형사상이 됨에 유의하라(위를 보라). 이로부터
$$
A = \bigoplus\nolimits_{d \geq 0} A_d \subset A' \subset
\prod\nolimits_{d \geq 0} A_d
$$
을 얻는다. 따라서
$$
f'_i = \sum\nolimits_{j = 0, \ldots, d - d_i - 1} f_{i, j} + f_i + g'_i
$$
로 쓸 수 있다. 여기서 $f_{i, j} \in A_j$,
$f_i \in A_{d - d_i}$이고, $g'_i \in A'$은
$\prod_{j \leq d - d_i} A_j$에서 영으로 사상된다.
이제 $\varphi_n(\sum_{i, j} f_{i, j}x_i) \in G_n$을 계산하면
차수가 $< d$인 동차원소들의 합을 얻는다. 따라서
$\varphi(\sum x_i \otimes f_{i, j})$는
$\lim G_{n, d}$에서 영으로 사상된다. 마찬가지로 계산하면
원소 $\varphi(\sum x_i \otimes g'_i)$가
$\prod_{d' \leq d} \lim G_{n, d'}$에서 영으로 사상됨을 알 수
있다. $\varphi(\sum x_i \otimes f'_i)$가 $y$임을 알고 있으므로
원하는 대로 $\sum f_ix_i \in M_d$가 $y$로 사상된다고 결론
내린다.
\end{proof}

\section{평탄성 판정법}
\label{algebra-section-criteria-flatness}

\noindent
이 절에서는 뇌터인 경우에 필요한 몇 가지 중요한 기술적 보조정리를
증명한다. 이 결과들은 절 \ref{algebra-section-more-flatness-criteria}에서
뇌터가 아닌 경우로 (부분적으로) 일반화한다.

\begin{lemma}
\label{algebra-lemma-mod-injective}
$R \to S$가 국소환들의 국소 준동형이고 $S$가 뇌터 환이라고 하자.
% P01-E0485: frozen source omits “by” in three notation declarations.
$\mathfrak m$을 $R$의 극대 아이디얼이라 하자. $M$을 평탄
$R$-가군, $N$을 유한 $S$-가군이라 하자. $u : N \to M$을
$R$-가군의 사상이라 하자.
$\overline{u} : N/\mathfrak m N \to M/\mathfrak m M$이
단사이면 $u$도 단사이다. 이 경우 $M/u(N)$은 $R$ 위에서 평탄하다.
\end{lemma}

\begin{proof}
먼저 모든 $n \geq 1$에 대해
$u_n : N/{\mathfrak m}^nN \to M/{\mathfrak m}^nM$이
단사임을 주장한다. 귀납법을 사용한다. 시작 단계에서는
$\overline{u} = u_1$이 단사이다. $M$이 $R$ 위에서 평탄하다는
가정에 의해 짧은 완전열
$0 \to M \otimes_R {\mathfrak m}^n/{\mathfrak m}^{n + 1}
\to M/{\mathfrak m}^{n + 1}M \to M/{\mathfrak m}^n M \to 0$을
얻는다. 또한
$M \otimes_R {\mathfrak m}^n/{\mathfrak m}^{n + 1}
= M/{\mathfrak m}M \otimes_{R/{\mathfrak m}}
{\mathfrak m}^n/{\mathfrak m}^{n + 1}$이다. $N$에 대해서도
왼쪽의 영을 제외하면 이와 비슷한 완전열
$N \otimes_R {\mathfrak m}^n/{\mathfrak m}^{n + 1}
\to N/{\mathfrak m}^{n + 1}N \to N/{\mathfrak m}^n N \to 0$이
있다. 또한
$N \otimes_R {\mathfrak m}^n/{\mathfrak m}^{n + 1}
= N/{\mathfrak m}N \otimes_{R/{\mathfrak m}}
{\mathfrak m}^n/{\mathfrak m}^{n + 1}$이다. 따라서 $u_n$과
$\overline{u} \otimes \text{id}_{{\mathfrak m}^n/{\mathfrak m}^{n + 1}}$가
모두 단사이므로 $u_{n + 1}$도 단사이다.

\medskip\noindent
$S$ 위의 $N$과 아이디얼 $\mathfrak mS$에 Krull의 교차 정리
(보조정리 \ref{algebra-lemma-intersect-powers-ideal-module-zero})를
적용하면 $\bigcap \mathfrak m^nN = 0$을 얻는다. 따라서 모든
$n$에 대한 $u_n$의 단사성은 $u$의 단사성을 함의한다.

\medskip\noindent
$M/u(N)$이 $R$ 위에서 평탄함을 보이려면 모든 아이디얼
$I \subset R$에 대해
$\text{Tor}_1^R(M/u(N), R/I) = 0$임을 보이면 충분하다.
보조정리 \ref{algebra-lemma-characterize-flat}를 보라. 짧은 완전열
$$
0 \to N \xrightarrow{u} M \to M/u(N) \to 0
$$
과 $M$의 평탄성으로부터 Tor의 완전열
$$
0 \to \text{Tor}_1^R(M/u(N), R/I) \to N/IN \to M/IM
$$
을 얻는다. 보조정리 \ref{algebra-lemma-long-exact-sequence-tor}를 보라.
따라서 $N/IN$이 $M/IM$으로 단사임을 보이면 충분하다.
$R/I \to S/IS$는 국소환들의 국소 준동형이고 $S/IS$는 뇌터
환이다. $N/IN \to M/IM$은 $R/I$-가군의 사상이고, $N/IN$은
$S/IS$ 위에서 유한이며, $M/IM$은 $R/I$ 위에서 평탄하고,
$u \bmod I : N/IN \to M/IM$은 $\mathfrak m$을 법으로 단사이다.
그러므로 증명의 첫 부분을 $u \bmod I$에 적용하면 결론을 얻는다.
\end{proof}

\begin{lemma}
\label{algebra-lemma-grothendieck}
$R \to S$가 뇌터 국소환들의 평탄 국소 준동형이라고 하자.
$\mathfrak m$을 $R$의 극대 아이디얼이라 하자. $f \in S$가
$S/{\mathfrak m}S$에서 비영인자라고 하자. 그러면 $S/fS$는
$R$ 위에서 평탄하고, $f$는 $S$에서 비영인자이다.
\end{lemma}

\begin{proof}
보조정리 \ref{algebra-lemma-mod-injective}에서 곧바로 따른다.
\end{proof}

\begin{lemma}
\label{algebra-lemma-grothendieck-regular-sequence}
$R \to S$가 뇌터 국소환들의 평탄 국소 준동형이라고 하자.
$\mathfrak m$을 $R$의 극대 아이디얼이라 하자.
$f_1, \ldots, f_c$를 그 상
$\overline{f}_1, \ldots, \overline{f}_c$가
$S/{\mathfrak m}S$에서 정칙열을 이루는 $S$의 원소들의 열이라 하자.
그러면 $f_1, \ldots, f_c$는 $S$에서 정칙열이고, 각 몫
$S/(f_1, \ldots, f_i)$는 $R$ 위에서 평탄하다.
\end{lemma}

\begin{proof}
귀납법과 보조정리 \ref{algebra-lemma-grothendieck}을 적용한다.
\end{proof}

\begin{lemma}
\label{algebra-lemma-free-fibre-flat-free}
$R \to S$를 뇌터 국소환들의 국소 준동형이라 하자.
$\mathfrak m$을 $R$의 극대 아이디얼이라 하자.
$M$을 영이 아닌 유한 $S$-가군이라 하자.
(a) $M/\mathfrak mM$이 자유 $S/\mathfrak mS$-가군이고
(b) $M$이 $R$ 위에서 평탄하다고 하자.
그러면 $M$은 자유가군이고 $S$는 $R$ 위에서 평탄하다.
\end{lemma}

\begin{proof}
$\overline{x}_1, \ldots, \overline{x}_n$을 자유가군
$M/\mathfrak mM$의 기저라 하자. $x_i$가 $\overline{x}_i$로
사상되도록 $x_1, \ldots, x_n \in M$을 택하자. $i$번째 표준
기저벡터를 $x_i$로 보내는 사상을
$u : S^{\oplus n} \to M$이라 하자. 보조정리
\ref{algebra-lemma-mod-injective}에 의해 $u$는 단사이다. 한편 Nakayama의
보조정리 \ref{algebra-lemma-NAK}에 의해 이 사상은 전사이다. 이제 결론이
따른다.
\end{proof}

\begin{lemma}
\label{algebra-lemma-complex-exact-mod}
$R \to S$를 국소 뇌터 환들의 국소 준동형이라 하자.
$\mathfrak m$을 $R$의 극대 아이디얼이라 하자.
$0 \to F_e \to F_{e-1} \to \ldots \to F_0$를 유한
$S$-가군들의 유한 복합체라 하자. 각 $F_i$가 $R$-평탄이고 복합체
$0 \to F_e/\mathfrak m F_e \to F_{e-1}/\mathfrak m F_{e-1}
\to \ldots \to F_0 / \mathfrak m F_0$가 완전하다고 가정하자.
그러면 $0 \to F_e \to F_{e-1} \to \ldots \to F_0$는
완전하고, 또한 가군 $\Coker(F_1 \to F_0)$는 $R$-평탄이다.
\end{lemma}

\begin{proof}
$e$에 대한 귀납법을 사용한다. $e = 1$이면 이는 정확히 보조정리
\ref{algebra-lemma-mod-injective}이다. $e > 1$이면 보조정리
\ref{algebra-lemma-mod-injective}에 의해 $F_e \to F_{e-1}$은 단사이고
$C = \Coker(F_e \to F_{e-1})$는 $R$ 위에서 평탄한 유한
$S$-가군이다. 따라서 복합체
$0 \to C \to F_{e-2} \to \ldots \to F_0$에 귀납가설을
적용할 수 있다. 그러면 $C \to F_{e-2}$가 단사이고 복합체가
완전임을 얻으며, $F_1 \to F_0$의 여핵이 평탄함도 얻는다.
\end{proof}

\noindent
이 절의 나머지 부분에서는 이른바 {\it 평탄성의 국소 판정법}의
두 가지 형태를 증명한다. 아래의 흥미로운 보조정리
\ref{algebra-lemma-CM-over-regular-flat}도 주목하라.

\begin{lemma}
\label{algebra-lemma-prepare-local-criterion-flatness}
$R$을 극대 아이디얼이 $\mathfrak m$이고 잉여체가
$\kappa = R/\mathfrak m$인 국소환이라 하자.
$M$을 $R$-가군이라 하자.
$\text{Tor}_1^R(\kappa, M) = 0$이면 모든 유한 길이
$R$-가군 $N$에 대해 $\text{Tor}_1^R(N, M) = 0$이다.
\end{lemma}

\begin{proof}
% P01-E0486: frozen source misdescribes this as descending induction.
$N$의 길이에 대한 귀납법을 사용한다. $N$의 길이가 $1$이면
$N \cong \kappa$이므로 끝난다. $N$의 길이가 $1$보다 크면
$N'$과 $N''$이 더 작은 유한 길이 $R$-가군인 짧은 완전열
$0 \to N' \to N \to N'' \to 0$에 $N$을 넣을 수 있다.
$\text{Tor}_1^R(N', M)$과 $\text{Tor}_1^R(N'', M)$의 소멸
(귀납가설), 그리고 Tor 군들의 긴 완전열로부터
$\text{Tor}_1^R(N, M)$의 소멸이 따른다. 보조정리
\ref{algebra-lemma-long-exact-sequence-tor}를 보라.
\end{proof}

\begin{lemma}[평탄성의 국소 판정법]
\label{algebra-lemma-local-criterion-flatness}
$R \to S$를 국소 뇌터 환들의 국소 준동형이라 하자.
$\mathfrak m$을 $R$의 극대 아이디얼이라 하고
$\kappa = R/\mathfrak m$이라 하자. $M$을 유한 $S$-가군이라 하자.
$\text{Tor}_1^R(\kappa, M) = 0$이면 $M$은 $R$ 위에서 평탄하다.
\end{lemma}

\begin{proof}
$I \subset R$을 아이디얼이라 하자. 보조정리 \ref{algebra-lemma-flat}에
의해 $I \otimes_R M \to M$이 단사임을 보이면 충분하다.
주의 \ref{algebra-remark-Tor-ring-mod-ideal}에 의해 그 핵은
$\text{Tor}_1^R(M, R/I)$와 같다. 보조정리
\ref{algebra-lemma-prepare-local-criterion-flatness}에 의해 유한 여길이를
갖는 모든 아이디얼에 대해 $J \otimes_R M \to M$이 단사이다.
% P01-E0487: frozen source leaves J unquantified; Korean supplies the evident range.

\medskip\noindent
$n >> 0$을 택하고 다음 짧은 완전열을 생각하자.
$$
0
\to I \cap \mathfrak m^n
\to I \oplus \mathfrak m^n
\to I + \mathfrak m^n
\to 0
$$
이는 짧은 완전열 $0 \to R \to R^{\oplus 2} \to R \to 0$의
부분열이다. 따라서 다음 도표를 얻는다.
$$
\xymatrix{
(I\cap \mathfrak m^n) \otimes_R M \ar[r] \ar[d] &
I \otimes_R M \oplus \mathfrak m^n \otimes_R M \ar[r] \ar[d] &
(I + \mathfrak m^n) \otimes_R M \ar[d] \\
M \ar[r] &
M \oplus M \ar[r] &
M
}
$$
$I + \mathfrak m^n$과 $\mathfrak m^n$은 유한 여길이를 갖는
아이디얼들이다. 따라서 도표 추적을 하면
$\Ker((I \cap \mathfrak m^n)\otimes_R M \to M)
\to \Ker(I \otimes_R M \to M)$이 전사임을 알 수 있다.
특히 $K = \Ker(I \otimes_R M \to M)$은
$I \otimes_R M$ 안에서 $(I \cap \mathfrak m^n) \otimes_R M$의
상에 포함된다. 아르틴--리스 보조정리, 즉 보조정리
\ref{algebra-lemma-Artin-Rees}에 의해 어떤 $c > 0$과 모든 $n >> 0$에
대해 $K$는 $\mathfrak m^{n-c}(I \otimes_R M)$에 포함된다.
$I \otimes_R M$은 유한 $S$-가군이고(!) $S$는 뇌터 환이므로,
이는 $K = 0$을 함의한다. 실제로 위의 사실은 $K$가
$I \otimes_R M$의 $\mathfrak mS$-진 완비화에서 영으로
사상됨을 뜻한다. 그런데 보조정리
\ref{algebra-lemma-completion-faithfully-flat}에 의해 $S$에서 그
$\mathfrak mS$-진 완비화로 가는 사상은 충실 평탄하다.
따라서 원하는 대로 $K = 0$이다.
\end{proof}

\noindent
다음과 같은 조건을 자주 만나게 된다.
“$M/IM$이 $R/I$ 위에서 평탄하고
$\text{Tor}_1^R(R/I, M) = 0$이다.”
다음 보조정리는 이 조건들의 몇 가지 귀결을 준다
(보조정리 \ref{algebra-lemma-prepare-local-criterion-flatness}의
일반화이다).

\begin{lemma}
\label{algebra-lemma-what-does-it-mean}
$R$을 환, $I \subset R$을 아이디얼, $M$을 $R$-가군이라 하자.
$M/IM$이 $R/I$ 위에서 평탄하고
$\text{Tor}_1^R(R/I, M) = 0$이면
\begin{enumerate}
\item 모든 $n \geq 1$에 대해 $M/I^nM$은 $R/I^n$ 위에서
평탄하고,
\item 어떤 $m \geq 0$에 대해 $I^m$에 의해 소멸하는 임의의
가군 $N$에 대해 $\text{Tor}_1^R(N, M) = 0$이다.
\end{enumerate}
특히 $I$가 멱영이면 $M$은 $R$ 위에서 평탄하다.
\end{lemma}

\begin{proof}
$M/IM$이 $R/I$ 위에서 평탄하고
$\text{Tor}_1^R(R/I, M) = 0$이라고 가정하자.
$N$을 $R/I$-가군이라 하자. 짧은 완전열
% P01-E0488: frozen source reuses the ideal I as the free-module index set.
$$
0 \to K \to \bigoplus\nolimits_{i \in I} R/I \to N \to 0
$$
을 택하자. $\text{Tor}$의 긴 완전열과
$\text{Tor}_1^R(R/I, M)$의 소멸로부터
$$
0 \to \text{Tor}_1^R(N, M) \to K \otimes_R M \to
(\bigoplus\nolimits_{i \in I} R/I) \otimes_R M \to N \otimes_R M \to 0
$$
을 얻는다. 그런데 $K$, $\bigoplus_{i \in I} R/I$, $N$은 모두
$I$에 의해 소멸하므로
\begin{align*}
K \otimes_R M & = K \otimes_{R/I} M/IM, \\
(\bigoplus\nolimits_{i \in I} R/I) \otimes_R M & =
(\bigoplus\nolimits_{i \in I} R/I) \otimes_{R/I} M/IM, \\
N \otimes_R M & = N \otimes_{R/I} M/IM.
\end{align*}
$M/IM$이 $R/I$ 위에서 평탄하므로
% P01-E0489: frozen source has malformed tensor base R/; retained exactly.
$$
0 \to K \otimes_{R/I} M/IM \to
(\bigoplus\nolimits_{i \in I} R/I) \otimes_{R/I} M/IM \to
N \otimes_{R/} M/IM \to 0
$$
은 완전하다. 이를 위의 열과 결합하면 $I$에 의해 소멸하는 임의의
$R$-가군 $N$에 대해 $\text{Tor}_1^R(N, M) = 0$임을 얻는다.

\medskip\noindent
$m$에 대한 귀납법으로 (2)를 증명하자. $m = 1$인 경우는 앞
문단에서 증명했다. $m > 1$이고 $I^m$에 의해 소멸하는 $N$에
대해서는 $N'$과 $N''$이 $I^{m - 1}$에 의해 소멸하는 완전열
$0 \to N' \to N \to N'' \to 0$을 택할 수 있다. 예를 들어
$N' = IN$과 $N'' = N/IN$을 택할 수 있다. 그러면 완전열
$$
\text{Tor}_1^R(N', M) \to
\text{Tor}_1^R(N, M) \to
\text{Tor}_1^R(N'', M)
$$
과 귀납가설이 원하는 소멸을 증명한다.

\medskip\noindent
마지막으로 (1)을 증명하자. $n \geq 1$이 주어졌을 때
$M/I^nM$이 $R/I^n$ 위에서 평탄함을 보여야 한다. 다시 말해
$I^n$에 의해 소멸하는 $R$-가군 $N$들의 범주에서 함자
$N \mapsto N \otimes_{R/I^n} M/I^nM$이 완전함을 보여야 한다.
그러한 $N$에 대해서는
$N \otimes_{R/I^n} M/I^nM = N \otimes_R M$이다. (2)의
$\text{Tor}_1$ 소멸에 의해 어떤 $I$의 거듭제곱에 의해 소멸하는
$N$들의 범주에서 함자 $N \mapsto N \otimes_R M$은 완전하다.
따라서 결론을 얻는다.
\end{proof}

\begin{lemma}
\label{algebra-lemma-what-does-it-mean-again}
$R$을 환, $I \subset R$을 아이디얼, $M$을 $R$-가군이라 하자.
\begin{enumerate}
\item $M/IM$이 $R/I$ 위에서 평탄하고
$M \otimes_R I/I^2 \to IM/I^2M$이 단사이면,
$M/I^2M$은 $R/I^2$ 위에서 평탄하다.
\item $M/IM$이 $R/I$ 위에서 평탄하고
$n = 1, \ldots, k$에 대해
$M \otimes_R I^n/I^{n + 1} \to I^nM/I^{n + 1}M$이 단사이면,
$M/I^{k + 1}M$은 $R/I^{k + 1}$ 위에서 평탄하다.
\end{enumerate}
\end{lemma}

\begin{proof}
첫 번째 명제는 $R$을 $R/I^2$로, $M$을 $M/I^2M$으로 바꾸어
보조정리 \ref{algebra-lemma-what-does-it-mean}을 적용하고
$$
\text{Tor}_1^{R/I^2}(M/I^2M, R/I) =
\Ker(M \otimes_R I/I^2 \to IM/I^2M),
$$
을 사용하면 따른다. 주의 \ref{algebra-remark-Tor-ring-mod-ideal}를 보라.
두 번째 명제도 같은 방식으로, $n = 1, \ldots, k$에 대해
$M/I^{n + 1}M$이 $R/I^{n + 1}$ 위에서 평탄함을 보이는
$n$에 대한 귀납법을 사용하여 따른다. 여기서는 모든 $n$에 대해
$$
\text{Tor}_1^{R/I^{n + 1}}(M/I^{n + 1}M, R/I^n) =
\Ker(M \otimes_R I^n/I^{n + 1} \to I^nM/I^{n + 1}M)
$$
임을 사용한다.
\end{proof}

\begin{lemma}[국소 판정법의 변형]
\label{algebra-lemma-variant-local-criterion-flatness}
$R \to S$를 뇌터 국소환들의 국소 준동형이라 하자.
$I \not = R$을 $R$의 아이디얼이라 하자. $M$을 유한
$S$-가군이라 하자. $\text{Tor}_1^R(M, R/I) = 0$이고
$M/IM$이 $R/I$ 위에서 평탄하면 $M$은 $R$ 위에서 평탄하다.
\end{lemma}

\begin{proof}
첫 번째 증명: 보조정리 \ref{algebra-lemma-what-does-it-mean}에 의해,
$\kappa$가 $R$의 잉여체일 때
$\text{Tor}_1^R(\kappa, M)$은 영이다. 따라서 보조정리
\ref{algebra-lemma-local-criterion-flatness}에 의해 $M$은 $R$ 위에서
평탄하다.

\medskip\noindent
두 번째 증명: $\mathfrak m$을 $R$의 극대 아이디얼이라 하자.
$\mathfrak m \otimes_R M \to M$이 단사임을 보인 뒤 보조정리
\ref{algebra-lemma-local-criterion-flatness}을 적용하겠다.
$\sum f_i \otimes x_i \in \mathfrak m \otimes_R M$이고
$M$에서 $\sum f_i x_i = 0$이라고 하자. $R/I$ 위의 $M/IM$에
평탄성의 방정식 판정법, 보조정리 \ref{algebra-lemma-flat-eq}를 적용하면
$\overline{a}_{ij} \in R/I$와 $\overline{y}_j \in M/IM$가
존재하여
$x_i \bmod IM = \sum_j \overline{a}_{ij} \overline{y}_j $이고
$0 = \sum_i (f_i \bmod I) \overline{a}_{ij}$이다.
$a_{ij} \in R$을 $\overline{a}_{ij}$의 올림으로 택하고, 마찬가지로
$y_j \in M$을 $\overline{y}_j$의 올림으로 택하자. 그러면
\begin{eqnarray*}
\sum f_i \otimes x_i
& = &
\sum f_i \otimes x_i +
\sum f_ia_{ij} \otimes y_j -
\sum f_i \otimes a_{ij} y_j
\\
& = &
\sum f_i \otimes (x_i - \sum a_{ij} y_j) +
\sum (\sum f_i a_{ij}) \otimes y_j
\end{eqnarray*}
임을 알 수 있다. $x_i - \sum a_{ij} y_j \in IM$이고
$\sum f_i a_{ij} \in I$이므로, $I \otimes_R M$의 어떤 원소가
$\mathfrak m \otimes_R M$에서 주어진 원소
$\sum f_i \otimes x_i$로 사상된다. 그런데 가정에 의해
$I \otimes_R M \to M$은 단사이다(주의
\ref{algebra-remark-Tor-ring-mod-ideal}를 보라). 이로써 증명이 끝난다.
\end{proof}

\noindent
특히 보조정리 \ref{algebra-lemma-variant-local-criterion-flatness}의
상황에서 $I = (x)$가 $R$의 비영인자 $x$ 하나에 의해
생성되었다고 하자. 그러면
$\text{Tor}_1^R(M, R/(x)) = (0)$일 필요충분조건은 $x$가
$M$ 위의 비영인자인 것이다.

\begin{lemma}
\label{algebra-lemma-flat-module-powers}
$R \to S$를 환 사상, $I \subset R$을 아이디얼, $M$을
$S$-가군이라 하자. 다음을 가정하자.
\begin{enumerate}
\item $R$은 뇌터 환이다.
\item $S$는 뇌터 환이다.
\item $M$은 유한 $S$-가군이다.
\item 각 $n \geq 1$에 대해 가군 $M/I^n M$은 $R/I^n$ 위에서
평탄하다.
\end{enumerate}
그러면 모든 $\mathfrak q \in V(IS)$에 대해 국소화
$M_{\mathfrak q}$는 $R$ 위에서 평탄하다.
특히 $S$가 국소환이고 $IS$가 그 극대 아이디얼에 포함되면
$M$은 $R$ 위에서 평탄하다.
\end{lemma}

\begin{proof}
보조정리 \ref{algebra-lemma-variant-local-criterion-flatness}을 사용하겠다.
가정에 의해 $M/IM$은 $R/I$ 위에서 평탄하다. 따라서
$\text{Tor}_1^R(M, R/I)$이 $\mathfrak q$에서 국소화한 뒤
영임을 확인하면 충분하다. 주의
\ref{algebra-remark-Tor-ring-mod-ideal}에 의해 이 Tor 군은
$K = \Ker(I \otimes_R M \to M)$과 같다. 모든 $n \geq 1$에 대해
$I/I^n \otimes_{R/I^n} M/I^nM \to M/I^nM$의 핵이 영임을 안다.
따라서 $K$의 원소는
$I/I^n \otimes_{R/I^n} M/I^nM$에서 영으로 사상된다.
다음 등식
$$
I/I^n \otimes_{R/I^n} M/I^nM = I/I^n \otimes_R M =
(I \otimes_R M)/I^{n - 1}(I \otimes_R M)
$$
에 의해 모든 $n \geq 1$에 대해
$K \subset I^{n - 1}(I \otimes_R M)$이다.
아르틴--리스 보조정리, 더 정확히 보조정리
\ref{algebra-lemma-intersection-powers-ideal-module}에 의해 원하는 대로
$K_{\mathfrak q} = 0$을 얻는다.
\end{proof}

\begin{lemma}
\label{algebra-lemma-surjective-on-tor-one}
$R \to R' \to R''$을 환 사상들이라 하자. $M$을 $R$-가군이라 하자.
$M \otimes_R R'$이 $R'$ 위에서 평탄하다고 하자. 그러면 자연스러운
사상
$\text{Tor}_1^R(M, R') \otimes_{R'} R'' \to
\text{Tor}_1^R(M, R'')$은 전사이다.
\end{lemma}

\begin{proof}
$F_\bullet$을 $R$ 위의 $M$의 자유 분해라 하자. 복합체
$F_2 \otimes_R R' \to F_1\otimes_R R' \to F_0 \otimes_R R'$는
$\text{Tor}_1^R(M, R')$을 계산한다. 복합체
$F_2 \otimes_R R'' \to F_1\otimes_R R'' \to F_0 \otimes_R R''$는
$\text{Tor}_1^R(M, R'')$을 계산한다.
$F_i \otimes_R R' \otimes_{R'} R'' = F_i \otimes_R R''$임을
주목하라.
$K' = \Ker(F_1\otimes_R R' \to F_0 \otimes_R R')$로 놓고,
마찬가지로
$K'' = \Ker(F_1\otimes_R R'' \to F_0 \otimes_R R'')$로 놓자.
그러면 완전열
$$
0 \to K' \to F_1\otimes_R R' \to F_0 \otimes_R R' \to M \otimes_R R' \to 0.
$$
을 얻는다. $M \otimes_R R'$이 $R'$ 위에서 평탄하다는 가정에
의해 다음 열도 여전히 완전하다.
$$
K' \otimes_{R'} R'' \to
F_1 \otimes_R R'' \to
F_0 \otimes_R R'' \to
M \otimes_R R'' \to 0
$$
이는 $K' \otimes_{R'} R'' \to K''$이 전사임을 뜻한다.
$\text{Tor}_1^R(M, R')$은 $K'$의 몫이고
$\text{Tor}_1^R(M, R'')$은 $K''$의 몫이므로 결론을 얻는다.
\end{proof}

\begin{lemma}
\label{algebra-lemma-surjective-on-tor-one-trivial}
$R \to R'$을 환 사상이라 하자. $I \subset R$을 아이디얼이라 하고
$I' = IR'$이라 하자. $M$을 $R$-가군이라 하고
$M' = M \otimes_R R'$으로 놓자. 자연스러운 사상
$\text{Tor}_1^R(R'/I', M) \to \text{Tor}_1^{R'}(R'/I', M')$은
전사이다.
\end{lemma}

\begin{proof}
$F_2 \to F_1 \to F_0 \to M \to 0$을 $R$ 위의 $M$의 자유 분해라
하자. $F_i' = F_i \otimes_R R'$으로 놓자. 열
$F_2' \to F_1' \to F_0' \to M' \to 0$은 이제 $F_1'$에서
완전하지 않을 수 있다. 따라서 $R'$ 위의 $M'$의 자유 분해는
적당한 자유 $R'$-가군 $F_2''$에 대해 다음과 같은 꼴이다.
$$
F_2' \oplus F_2'' \to F_1' \to F_0' \to M' \to 0
$$
다음으로
$F_i \otimes_R R'/I' = F_i' / IF_i' = F_i'/I'F_i'$임을
주목하라. 그러므로 복합체
$F_2'/I'F_2' \to F_1'/I'F_1' \to F_0'/I'F_0'$는
$\text{Tor}_1^R(M, R'/I')$을 계산한다. 한편
$F_i' \otimes_{R'} R'/I' = F_i'/I'F_i'$이고 $F_2''$에 대해서도
마찬가지이다. 따라서 복합체
$F_2'/I'F_2' \oplus F_2''/I'F_2'' \to F_1'/I'F_1' \to F_0'/I'F_0'$는
$\text{Tor}_1^{R'}(M', R'/I')$을 계산한다. 복합체들의 수직 사상
$$
\xymatrix{
F_2'/I'F_2' \ar[r] \ar[d] &
F_1'/I'F_1' \ar[r] \ar[d] &
F_0'/I'F_0' \ar[d] \\
F_2'/I'F_2' \oplus F_2''/I'F_2'' \ar[r] &
F_1'/I'F_1' \ar[r] &
F_0'/I'F_0'
}
$$
% P01-E0490: frozen source says cohomology for these homological Tor complexes.
이 호몰로지 위에 전사를 유도함은 명백하므로 결론을 얻는다.
\end{proof}

\begin{lemma}
\label{algebra-lemma-another-variant-local-criterion-flatness}
국소 뇌터 환들의 국소 준동형들로 이루어진 교환 도표
$$
\xymatrix{
S \ar[r] & S' \\
R \ar[r] \ar[u] & R' \ar[u]
}
$$
가 주어졌다고 하자. $I \subset R$을 진 아이디얼이라 하자.
$M$을 유한 $S$-가군이라 하자.
% P01-E0491: frozen source omits “by” in the notation declaration.
$I' = IR'$ 및 $M' = M \otimes_S S'$으로 표기하자.
다음을 가정하자.
\begin{enumerate}
\item $S'$은 텐서곱 $S \otimes_R R'$의 한 국소화이다.
\item $M/IM$은 $R/I$ 위에서 평탄하다.
\item $\text{Tor}_1^R(M, R/I) \to \text{Tor}_1^{R'}(M', R'/I')$은
영사상이다.
\end{enumerate}
그러면 $M'$은 $R'$ 위에서 평탄하다.
\end{lemma}

\begin{proof}
$S'$은 $S \otimes_R R'$의 국소화이므로 $M'$은
$M \otimes_R R'$의 국소화이다. 보조정리
\ref{algebra-lemma-flat-base-change}에 의해 가군
$M/IM \otimes_{R/I} R'/I'
= M \otimes_R R' /I'(M \otimes_R R')$은 $R'/I'$ 위에서
평탄하다. 따라서 $M'/I'M'$은 평탄 가군의 국소화이므로
$R'/I'$ 위에서 평탄하다. 보조정리
\ref{algebra-lemma-variant-local-criterion-flatness}에 의해
$\text{Tor}_1^{R'}(M', R'/I')$이 영임을 보이면 충분하다.
$M'$은 $M \otimes_R R'$의 국소화이므로 마지막 가정에 의해
다음 사상이 전사임을 보이면 충분하다.
$\text{Tor}_1^R(M, R/I) \otimes_R R'
\to
\text{Tor}_1^{R'}(M \otimes_R R', R'/I')$

\medskip\noindent
보조정리 \ref{algebra-lemma-surjective-on-tor-one-trivial}에 의해
$\text{Tor}_1^R(M, R'/I') \to \text{Tor}_1^{R'}(M \otimes_R R', R'/I')$은
전사이다. 따라서 이제 다음 사상이 전사임을 보이면 충분하다.
$\text{Tor}_1^R(M, R/I) \otimes_R R'
\to
\text{Tor}_1^R(M, R'/I')$
이는 환 사상 $R \to R/I \to R'/I'$과 가군 $M$을 사용하여
보조정리 \ref{algebra-lemma-surjective-on-tor-one}을 적용하면 따른다.
\end{proof}

\noindent
다음 보조정리를 보조정리
\ref{algebra-lemma-criterion-flatness-fibre-nilpotent}
(멱영 아이디얼의 경우) 및 보조정리
\ref{algebra-lemma-criterion-flatness-fibre}
(유한 표시 대수의 경우)와 비교하라.

\begin{lemma}[섬유별 평탄성 판정법; 뇌터인 경우]
\label{algebra-lemma-criterion-flatness-fibre-Noetherian}
$R$, $S$, $S'$을 뇌터 국소환들이라 하고
$R \to S \to S'$을 국소환 준동형들이라 하자.
$\mathfrak m \subset R$을 극대 아이디얼이라 하자.
$M$을 $S'$-가군이라 하자. 다음을 가정한다.
\begin{enumerate}
\item 가군 $M$은 $S'$ 위에서 유한이다.
\item 가군 $M$은 영이 아니다.
\item 가군 $M/\mathfrak m M$은 평탄 $S/\mathfrak m S$-가군이다.
\item 가군 $M$은 평탄 $R$-가군이다.
\end{enumerate}
그러면 $S$는 $R$ 위에서 평탄하고 $M$은 평탄 $S$-가군이다.
\end{lemma}

\begin{proof}
$I = \mathfrak mS \subset S$로 놓자. (3)에 의해 $M/IM$은 평탄
$S/I$-가군이다.
$\mathfrak m \otimes_R S' \to I \otimes_S S'$이 전사이므로
$\mathfrak m \otimes_R M \to I \otimes_S M$도 전사이다.
다음을 생각하자.
$$
\mathfrak m \otimes_R M \to I \otimes_S M \to M.
$$
$M$은 $R$ 위에서 평탄하므로 합성은 단사이고, 따라서 두 화살은
모두 단사이다. 특히
$\text{Tor}_1^S(S/I, M) = 0$이다. 주의
\ref{algebra-remark-Tor-ring-mod-ideal}를 보라. 보조정리
\ref{algebra-lemma-variant-local-criterion-flatness}에 의해 $M$은
$S$ 위에서 평탄하다. Nakayama의 보조정리 \ref{algebra-lemma-NAK}에
의해 $M/\mathfrak m_{S'}M$이 영이 아니므로, 실제로 보조정리
\ref{algebra-lemma-ff}에 의해 $M$은 $S$ 위에서 충실 평탄하다
(이는 $M/\mathfrak m_SM \not = 0$을 강제하기 때문이다).

\medskip\noindent
완전열 $0 \to \mathfrak m \to R \to \kappa \to 0$을 생각하자.
이 열은 완전열
$0 \to \text{Tor}_1^R(\kappa, S) \to \mathfrak m \otimes_R S \to I \to 0$을
준다. $M$은 $S$ 위에서 평탄하므로 이는 완전열
$0 \to \text{Tor}_1^R(\kappa, S)\otimes_S M \to
\mathfrak m \otimes_R M \to I \otimes_S M \to 0$을 준다.
위의 결과에 의해
$\text{Tor}_1^R(\kappa, S)\otimes_S M = 0$이다.
$M$은 $S$ 위에서 충실 평탄하므로
$\text{Tor}_1^R(\kappa, S) = 0$이다. 따라서 보조정리
\ref{algebra-lemma-local-criterion-flatness}에 의해 $S$는 $R$ 위에서
평탄하다.
\end{proof}

\begin{lemma}
\label{algebra-lemma-flatness-scallop}
$A$를 환, $M$을 $A$-가군, $f \in A$를 원소라 하자.
다음을 가정하자.
\begin{enumerate}
\item $f$는 $A$와 $M$ 위의 비영인자이다.
\item $M_f$는 평탄 $A_f$-가군이다.
\item $M/fM$은 평탄 $A/fA$-가군이다.
\end{enumerate}
그러면 $M$은 평탄 $A$-가군이다. “평탄”을 “충실 평탄”으로
바꾸어도 마찬가지이다.
\end{lemma}

\begin{proof}
$0 \to A \to A \to A/fA \to 0$과
$0 \to M \to M \to M/fM \to 0$이 완전하므로
$i = 1$ 및 $i = 2$에 대해
$\text{Tor}_i^A(M, A/fA) = 0$이다.
보조정리 \ref{algebra-lemma-what-does-it-mean}에 의해 $f$에 의해
소멸하는 모든 $A$-가군 $N$에 대해
$\text{Tor}_1^A(M, N) = 0$이다(여기서는 $M/fM$이
$A/fA$ 위에서 평탄함을 사용한다). $f$에 의해 소멸하는
$A$-가군 $N$이 주어지면 $F$가 $A/fA$의 복사본들의 직합인
$A$-가군들의 짧은 완전열
$0 \to N' \to F \to N \to 0$을 택할 수 있다. 완전열
$$
\text{Tor}_2^A(M, F) \to
\text{Tor}_2^A(M, N) \to
\text{Tor}_1^A(M, N')
$$
로부터 $f$에 의해 소멸하는 모든 $A$-가군 $N$에 대해
$\text{Tor}_2^A(M, N) = 0$을 얻는다.
이제 $K$를 임의의 $A$-가군이라 하자. 사상 $f : K \to K$를
두 짧은 완전열
$$
0 \to K[f] \to K \to K' \to 0
\quad\text{and}\quad
0 \to K' \to K \to K/fK \to 0
$$
로 나눌 수 있다. 여기서 $K' = K/K[f] \cong fK$이다. Tor의
완전열을 적용하면 완전열
$$
\text{Tor}_1^A(K[f], M) \to
\text{Tor}_1^A(K, M) \to
\text{Tor}_1^A(K', M)
$$
과
$$
\text{Tor}_2^A(K/fK, M) \to
\text{Tor}_1^A(K', M) \to
\text{Tor}_1^A(K, M)
$$
을 얻는다. $\text{Tor}_1^A(K[f], M)$과
$\text{Tor}_2^A(K/fK, M)$의 소멸을 사용하면
$f : K \to K$가 $\text{Tor}^A_1(M, K)$ 위에 단사 사상을
유도함을 얻는다. 다시 말해 $\text{Tor}^A_1(M, K)$가 영일
필요충분조건은
$\text{Tor}^A_1(M, K) \otimes_A A_f$가 영인 것이다.
보조정리 \ref{algebra-lemma-flat-base-change-tor}에 의해
$$
\text{Tor}^A_1(M, K) \otimes_A A_f = \text{Tor}^{A_f}_1(M_f, K_f)  = 0
$$
% P01-E0492: frozen source has this sentence fragment twice.
마지막 소멸은 $M_f$가 $A_f$ 위에서 평탄함에서 따른다
(보조정리 \ref{algebra-lemma-characterize-flat}).
따라서 모든 $A$-가군 $K$에 대해 $\text{Tor}_1^A(M, K) = 0$이다.
그러므로 보조정리 \ref{algebra-lemma-characterize-flat}에 의해
$M$은 $A$ 위에서 평탄하다.

\medskip\noindent
충실 평탄 가군의 경우에 대한 논증은 생략한다.
\end{proof}

\begin{lemma}
\label{algebra-lemma-flatness-scallop-pre}
$A$를 환, $M$을 $A$-가군이라 하자.
$I = (f_1, \ldots, f_r)$를 $r \geq 1$개의 원소로 생성되는
$A$의 아이디얼이라 하자. 다음을 가정하자.
\begin{enumerate}
\item $i = 1, \ldots, r$에 대해 $M_{f_i}$는 평탄
$A_{f_i}$-가군이다.
\item $M/IM$은 평탄 $A/I$-가군이다.
\item $i = 1, \ldots, r + 1$에 대해
$\text{Tor}_i^A(M, A/I) = 0$이다.
\end{enumerate}
그러면 $M$은 평탄 $A$-가군이다. “평탄”을 “충실 평탄”으로
바꾸어도 마찬가지이다.
\end{lemma}

\begin{proof}
보조정리 \ref{algebra-lemma-what-does-it-mean}에 의해 $I$에 의해
소멸하는 모든 $A$-가군 $K$에 대해
$\text{Tor}_1^A(M, K) = 0$이다.
$I$에 의해 소멸하는 $A$-가군 $K$가 주어지면, $F$가
$A/I$의 복사본들의 직합인 $A$-가군들의 짧은 완전열
$0 \to N \to F \to K \to 0$을 택할 수 있다. 완전열
$$
\text{Tor}_i^A(M, F) \to
\text{Tor}_i^A(M, K) \to
\text{Tor}_{i - 1}^A(M, N)
$$
을 얻는다. 따라서 가정 (3)과 $i$에 대한 귀납법을 사용하면
$I$에 의해 소멸하는 모든 $A$-가군 $K$와
$i = 1, \ldots, r + 1$에 대해
$\text{Tor}_i^A(M, K) = 0$임을 얻는다.

\medskip\noindent
어떤 $1 \leq j \leq r$에 대해 $f_1, \ldots, f_j$에 의해
소멸하는 모든 $A$-가군 $K$와 $i = 1, \ldots, j + 1$에 대해
$\text{Tor}_i^A(M, K) = 0$임을 보였다고 하자.
$K$를 $f_1, \ldots, f_{j - 1}$에 의해 소멸하는
$A$-가군이라 하자. 사상 $f_j : K \to K$를 두 짧은 완전열
$$
0 \to K[f_j] \to K \to K' \to 0
\quad\text{and}\quad
0 \to K' \to K \to K/f_jK \to 0
$$
로 나눌 수 있다. 여기서 $K' = K/K[f_j] \cong f_jK$이다.
$1 \leq i \leq j$라 하자. Tor의 완전열을 적용하면
$$
\text{Tor}_i^A(K[f_j], M) \to
\text{Tor}_i^A(K, M) \to
\text{Tor}_i^A(K', M)
$$
과
$$
\text{Tor}_{i + 1}^A(K/f_jK, M) \to
\text{Tor}_i^A(K', M) \to
\text{Tor}_i^A(K, M)
$$
을 얻는다. $\text{Tor}_i^A(K[f_j], M)$과
$\text{Tor}_{i + 1}^A(K/f_jK, M)$의 소멸을 사용하면
$f_j : K \to K$가 $\text{Tor}^A_i(M, K)$ 위에 단사 사상을
유도함을 얻는다. 다시 말해 $\text{Tor}^A_i(M, K)$가 영일
필요충분조건은
$\text{Tor}^A_i(M, K) \otimes_A A_{f_j}$가 영인 것이다.
보조정리 \ref{algebra-lemma-flat-base-change-tor}에 의해
$$
\text{Tor}^A_i(M, K) \otimes_A A_{f_j} =
\text{Tor}^{A_{f_j}}_i(M_{f_j}, K_{f_j})  = 0
$$
마지막 소멸은 $M_{f_j}$가 $A_{f_j}$ 위에서 평탄함에서 따른다
(보조정리 \ref{algebra-lemma-characterize-flat}).
따라서 $f_1, \ldots, f_{j - 1}$에 의해 소멸하는 모든
$A$-가군 $K$와 $i = 1, \ldots, j$에 대해
$\text{Tor}_i^A(M, K) = 0$이다.
$j$에 대한 내림 귀납법을 사용하면 이는 $j = 0$에도 성립한다.
즉 모든 $A$-가군 $K$에 대해
$\text{Tor}_1^A(M, K) = 0$이다.
따라서 보조정리 \ref{algebra-lemma-characterize-flat}에 의해
$M$은 $A$ 위에서 평탄하다.

\medskip\noindent
충실 평탄 가군의 경우에 대한 논증은 생략한다.
\end{proof}

\section{밑바꿈과 평탄성}
\label{algebra-section-base-change-flat}

\noindent
밑바꿈을 할 때 평탄성에 어떤 일이 일어나는지를 다루는
보조정리들을 모아 둔다.

\begin{lemma}
\label{algebra-lemma-base-change-flat-up-down}
국소환들의 국소 준동형들로 이루어진 교환 도표
$$
\xymatrix{
S \ar[r] & S' \\
R \ar[r] \ar[u] & R' \ar[u]
}
$$
가 주어졌다고 하자. $S'$이 텐서곱 $S \otimes_R R'$의
한 국소화라고 가정하자. $M$을 $S$-가군이라 하고
$M' = S' \otimes_S M$으로 놓자.
\begin{enumerate}
\item $M$이 $R$ 위에서 평탄하면 $M'$은 $R'$ 위에서 평탄하다.
\item $M'$이 $R'$ 위에서 평탄하고 $R \to R'$이 평탄하면
$M$은 $R$ 위에서 평탄하다.
\end{enumerate}
특히 다음이 성립한다.
\begin{enumerate}
\item[(3)] $S$가 $R$ 위에서 평탄하면 $S'$은 $R'$ 위에서 평탄하다.
\item[(4)] $R' \to S'$과 $R \to R'$이 평탄하면 $S$는
$R$ 위에서 평탄하다.
\end{enumerate}
\end{lemma}

\begin{proof}
(1)의 증명. $M$이 $R$ 위에서 평탄하면 보조정리
\ref{algebra-lemma-flat-base-change}에 의해 $M \otimes_R R'$은
$R'$ 위에서 평탄하다. $W \subset S \otimes_R R'$가
$W^{-1}(S \otimes_R R') = S'$을 만족하는 곱셈적 부분집합이면
$M' = W^{-1}(M \otimes_R R')$이다. 따라서 $M'$은 평탄 가군의
국소화이므로 $R'$ 위에서 평탄하다. 보조정리
\ref{algebra-lemma-flat-localization}의 (5)를 보라. 이로써 (1)을
증명했으며, 특히 (3)도 성립한다.

\medskip\noindent
(2)의 증명. $M'$이 $R'$ 위에서 평탄하고 $R \to R'$이
평탄하다고 하자. 북서 대각선에 관해 반사한 도표에 (3)을
적용하면 $S \to S'$이 평탄함을 알 수 있다. 따라서 보조정리
\ref{algebra-lemma-local-flat-ff}에 의해 $S \to S'$은 충실 평탄하다.
$M$이 평탄함을 보이기 위해 보조정리 \ref{algebra-lemma-flat}의
판정법 (\ref{algebra-item-f-ideal})을 사용하겠다.
$I \subset R$을 아이디얼이라 하자.
% P01-E0493: frozen source says merely “has a kernel”; descent requires a nonzero kernel.
$I \otimes_R M \to M$의 핵이 영이 아니면
$(I \otimes_R M) \otimes_S S' \to M \otimes_S S' = M'$의
핵도 영이 아니다. $R \to R'$이 평탄하므로
$I \otimes_R R' = IR'$이고,
$$
(I \otimes_R M) \otimes_S S' =
(I \otimes_R R') \otimes_{R'} (M \otimes_S S') =
IR' \otimes_{R'} M'.
$$
임을 주목하라. $M'$이 $R'$ 위에서 평탄하므로 이 가군은
$M'$으로 단사적으로 사상된다. 이로써 (2)의 증명이 끝나고,
(4)도 성립한다.
\end{proof}

\noindent
평탄성의 국소 판정법을 응용한 또 하나의 결과는 다음과 같다.

\begin{lemma}
\label{algebra-lemma-yet-another-variant-local-criterion-flatness}
국소환과 국소 준동형들의 교환 도표
$$
\xymatrix{
S \ar[r] & S' \\
R \ar[r] \ar[u] & R' \ar[u]
}
$$
를 생각하자. $M$을 유한 $S$-가군이라 하자. 다음을 가정하자.
\begin{enumerate}
\item 수평 화살들은 평탄 환 사상이다.
\item $M$은 $R$ 위에서 평탄하다.
\item $\mathfrak m_R R' = \mathfrak m_{R'}$이다.
\item $R'$과 $S'$은 뇌터 환이다.
\end{enumerate}
그러면 $M' = M \otimes_S S'$은 $R'$ 위에서 평탄하다.
\end{lemma}

\begin{proof}
$\mathfrak m_R \subset R$이고 $R \to R'$이 평탄하므로 가정 (3)에
의해
$\mathfrak m_R \otimes_R R' = \mathfrak m_R R' = \mathfrak m_{R'}$이다.
$M'$은 유한 $S'$-가군이고 보조정리
\ref{algebra-lemma-flatness-descends-more-general}에 의해 $R$ 위에서
평탄하다. 따라서 $\mathfrak m_R \otimes_R M' \to M'$은
단사이다. 이제
$$
\mathfrak m_R \otimes_R M' =
\mathfrak m_R \otimes_R R' \otimes_{R'} M' =
\mathfrak m_{R'} \otimes_{R'} M'
$$
을 얻는다. 그러므로
$\mathfrak m_{R'} \otimes_{R'} M' \to M'$은 단사이다.
이는 $\text{Tor}_1^{R'}(\kappa_{R'}, M') = 0$임을 보인다
(주의 \ref{algebra-remark-Tor-ring-mod-ideal}).
따라서 보조정리 \ref{algebra-lemma-local-criterion-flatness}에 의해
$M'$은 $R'$ 위에서 평탄하다.
\end{proof}

\section{아르틴 환 위의 평탄성 판정법}
\label{algebra-section-flatness-artinian}

\noindent
아르틴 환 위의 가군에 대한 몇 가지 평탄성 판정법을 논한다.
아르틴 국소환의 극대 아이디얼은 멱영이므로 다음 두 보조정리를
아르틴 국소환에 적용할 수 있음에 주목하라.

\begin{lemma}
\label{algebra-lemma-local-artinian-basis-when-flat}
$(R, \mathfrak m)$을 극대 아이디얼 $\mathfrak m$이 멱영인
국소환이라 하자. $M$을 평탄 $R$-가군이라 하자.
$A$가 집합이고 $x_\alpha \in M$, $\alpha \in A$가 $M$의
원소들의 모임이면 다음은 동치이다.
\begin{enumerate}
\item $\{\overline{x}_\alpha\}_{\alpha \in A}$는
$R/\mathfrak m$ 위의 벡터공간 $M/\mathfrak mM$의 기저를 이루고,
\item $\{x_\alpha\}_{\alpha \in A}$는 $R$ 위의 $M$의 기저를 이룬다.
\end{enumerate}
\end{lemma}

\begin{proof}
(2) $\Rightarrow$ (1)은 자명하다. (1)을 가정하자. Nakayama의
보조정리 \ref{algebra-lemma-NAK}에 의해 원소들 $x_\alpha$는 $M$을
생성한다. 따라서 짧은 완전열
$$
0 \to K \to \bigoplus\nolimits_{\alpha \in A} R \to M \to 0
$$
을 얻는다. $R/\mathfrak m$과 텐서하고 보조정리
\ref{algebra-lemma-flat-tor-zero}를 사용하면
$K/\mathfrak mK = 0$을 얻는다. Nakayama의 보조정리
\ref{algebra-lemma-NAK}에 의해 $K = 0$이다.
\end{proof}

\begin{lemma}
\label{algebra-lemma-local-artinian-characterize-flat}
$R$을 극대 아이디얼이 멱영인 국소환이라 하자.
$M$을 $R$-가군이라 하자. 다음은 동치이다.
\begin{enumerate}
\item $M$은 $R$ 위에서 평탄하다.
\item $M$은 자유 $R$-가군이다.
\item $M$은 사영 $R$-가군이다.
\end{enumerate}
\end{lemma}

\begin{proof}
모든 사영 가군은 자유가군의 직합인자로서 평탄하고, 모든
자유가군은 사영이므로 평탄 가군이 자유임을 보이면 충분하다.
$M$을 평탄 가군이라 하자. $A$를 집합이라 하고,
$\overline{x_\alpha} \in M/\mathfrak m M$가 $R$의 잉여체 위에서
기저를 이루도록 원소 $x_\alpha \in M$, $\alpha \in A$를 택하자.
보조정리 \ref{algebra-lemma-local-artinian-basis-when-flat}에 의해
$x_\alpha$들은 $R$ 위의 $M$의 기저이다. 이로써 끝난다.
\end{proof}

\begin{lemma}
\label{algebra-lemma-lift-basis}
$R$을 환, $I \subset R$을 아이디얼, $M$을 $R$-가군이라 하자.
$A$를 집합이라 하고 $x_\alpha \in M$, $\alpha \in A$를
$M$의 원소들의 모임이라 하자. 다음을 가정하자.
\begin{enumerate}
\item $I$는 멱영이다.
\item $\{\overline{x}_\alpha\}_{\alpha \in A}$는 $R/I$ 위의
$M/IM$의 기저를 이룬다.
\item $\text{Tor}_1^R(R/I, M) = 0$이다.
\end{enumerate}
그러면 $M$은 $R$ 위에서 $\{x_\alpha\}_{\alpha \in A}$를
기저로 하는 자유가군이다.
\end{lemma}

\begin{proof}
$R$, $I$, $M$, $\{x_\alpha\}_{\alpha \in A}$가 보조정리에서와
같고 가정 (1), (2), (3)을 만족한다고 하자. Nakayama의 보조정리
\ref{algebra-lemma-NAK}에 의해 원소들 $x_\alpha$는 $R$ 위에서 $M$을
생성한다. 가정 $\text{Tor}_1^R(R/I, M) = 0$은 짧은 완전열
$$
0 \to I \otimes_R M \to M \to M/IM \to 0.
$$
이 있음을 함의한다. $\sum f_\alpha x_\alpha = 0$을 $M$에서의
관계식이라 하자. $x_\alpha$들의 선택에 의해
$f_\alpha \in I$이다. 따라서 $I \otimes_R M$에서
$\sum f_\alpha \otimes x_\alpha = 0$이다.
사상 $I \otimes_R M \to I/I^2 \otimes_{R/I} M/IM$과
$\{x_\alpha\}_{\alpha \in A}$가 $M/IM$의 기저를 이룬다는
사실로부터 $f_\alpha \in I^2$임을 얻는다! 따라서 $M/I^2M$에서
$x_\alpha$들의 상들 사이에는 관계식이 없다. 다시 말해
$M/I^2M$은 $x_\alpha$들의 상들을 기저로 하는 자유가군이다.
이제 사상 $I \otimes_R M \to I/I^3 \otimes_{R/I^2} M/I^2M$을
사용하면 $f_\alpha \in I^3$임을 얻는다! 이 과정을 계속한다.
가정 (1)에 의해 어떤 $n$에 대해 $I^n = 0$이므로 끝난다.
\end{proof}

\begin{lemma}
\label{algebra-lemma-prepare-lift-flatness}
$\varphi : R \to R'$을 환 사상, $I \subset R$을 아이디얼,
$M$을 $R$-가군이라 하자. 다음을 가정하자.
\begin{enumerate}
\item $M/IM$은 $R/I$ 위에서 평탄하다.
\item $R' \otimes_R M$은 $R'$ 위에서 평탄하다.
\end{enumerate}
$I_2 = \varphi^{-1}(\varphi(I^2)R')$로 놓자.
그러면 $M/I_2M$은 $R/I_2$ 위에서 평탄하다.
\end{lemma}

\begin{proof}
$R$, $M$, $R'$을 각각 $R/I_2$, $M/I_2M$,
$R'/\varphi(I)^2R'$으로 바꾸어도 된다. 그러면 $I^2 = 0$이고
$\varphi$는 단사이다. 보조정리 \ref{algebra-lemma-what-does-it-mean}과
$I^2 = 0$이라는 사실에 의해
$\text{Tor}^R_1(R/I, M) = K = \Ker(I \otimes_R M \to M)$가
영임을 보이면 충분하다.
$M' = M \otimes_R R'$ 및 $I' = IR'$으로 놓자.
가정에 의해 사상 $I' \otimes_{R'} M' \to M'$은 단사이다.
따라서 $K$는
$$
I' \otimes_{R'} M' = I' \otimes_R M = I' \otimes_{R/I} M/IM.
$$
에서 영으로 사상된다. 이제 $I \to I'$은 $R/I$-가군의 단사
사상이다. $M/IM$은 $R/I$ 위에서 평탄하므로 사상
$$
I \otimes_{R/I} M/IM \longrightarrow I' \otimes_{R/I} M/IM
$$
은 단사이다. 따라서 원하는 대로
$I \otimes_R M = I \otimes_{R/I} M/IM$에서 $K$는 영이다.
\end{proof}

\begin{lemma}
\label{algebra-lemma-lift-flatness}
$\varphi : R \to R'$을 환 사상, $I \subset R$을 아이디얼,
$M$을 $R$-가군이라 하자. 다음을 가정하자.
\begin{enumerate}
\item $I$는 멱영이다.
\item $R \to R'$은 단사이다.
\item $M/IM$은 $R/I$ 위에서 평탄하다.
\item $R' \otimes_R M$은 $R'$ 위에서 평탄하다.
\end{enumerate}
그러면 $M$은 $R$ 위에서 평탄하다.
\end{lemma}

\begin{proof}
귀납적으로 $I_1 = I$로 놓고, $n \geq 1$에 대해
$I_{n + 1} = \varphi^{-1}(\varphi(I_n)^2R')$로 정의하자.
보조정리 \ref{algebra-lemma-prepare-lift-flatness}에 의해 각 $n \geq 1$에
대해 $M/I_nM$은 $R/I_n$ 위에서 평탄하다.
$\varphi(I_{n + 1}) \subset \varphi(I)^{2^n}R'$임은 명백하다.
$I$가 멱영이므로 어떤 $n$에 대해 $\varphi(I_n) = 0$이다.
$\varphi$가 단사이므로 어떤 $n$에 대해 $I_n = 0$이다.
이로써 끝난다.
\end{proof}

\noindent
다음은 평탄성의 국소 판정법에 대한 아르틴 국소환 버전이다.

\begin{lemma}
\label{algebra-lemma-artinian-variant-local-criterion-flatness}
$R$을 아르틴 국소환, $M$을 $R$-가군,
$I \subset R$을 진 아이디얼이라 하자. 다음은 동치이다.
\begin{enumerate}
\item $M$은 $R$ 위에서 평탄하다.
\item $M/IM$은 $R/I$ 위에서 평탄하고
$\text{Tor}_1^R(R/I, M) = 0$이다.
\end{enumerate}
\end{lemma}

\begin{proof}
(1) $\Rightarrow$ (2)는 정의에서 곧바로 따른다.
$M/IM$이 $R/I$ 위에서 평탄하고
$\text{Tor}_1^R(R/I, M) = 0$이라고 하자.
보조정리 \ref{algebra-lemma-local-artinian-characterize-flat}에 의해
$M/IM$은 $R/I$ 위에서 자유이다. $M/IM$에서의 상들이 기저를
이루도록 집합 $A$와 원소들 $x_\alpha \in M$을 택하자.
보조정리 \ref{algebra-lemma-lift-basis}에 의해 $M$은 자유이고,
특히 평탄하다.
\end{proof}

\noindent
시작환이 아르틴 환인 단사 준동형을 따라 평탄성이 하강한다.

\begin{lemma}
\label{algebra-lemma-descent-flatness-injective-map-artinian-rings}
$R \to S$를 환 사상, $M$을 $R$-가군이라 하자. 다음을 가정하자.
\begin{enumerate}
\item $R$은 아르틴 환이다.
\item $R \to S$는 단사이다.
\item $M \otimes_R S$는 평탄 $S$-가군이다.
\end{enumerate}
그러면 $M$은 평탄 $R$-가군이다.
\end{lemma}

\begin{proof}
첫 번째 증명: $I \subset R$을 $R$의 Jacobson 근기라 하자.
그러면 $I$는 멱영이고, 절 \ref{algebra-section-artinian}을 보면
$R/I$는 체들의 곱이므로 $M/IM$은 $R/I$ 위에서 평탄하다.
따라서 보조정리 \ref{algebra-lemma-lift-flatness}을 적용하면
$M$이 평탄함을 얻는다.

\medskip\noindent
두 번째 증명: 보조정리 \ref{algebra-lemma-artinian-finite-length}에 의해
$R = \prod R_i$를 국소 아르틴 환들의 유한 곱으로 쓸 수 있다.
% P01-E0496: frozen source redundantly says both R and S, omitting M.
이는 $R$과 $S$ 양쪽에 비슷한 곱 분해를 유도한다. 따라서
$R$이 국소 아르틴 환인 경우로 환원된다(세부사항은 생략한다).

\medskip\noindent
$R \to S$와 $M$이 보조정리의 (1), (2), (3)을 만족하고,
추가로 $R$이 극대 아이디얼 $\mathfrak m$을 갖는 국소환이라고
하자. $A$를 집합이라 하고
% P01-E0494: frozen source types x_alpha as an element of the index set A.
$\overline{x}_\alpha$가 $R/\mathfrak m$ 위의
$M/\mathfrak mM$의 기저를 이루도록 원소들
$x_\alpha \in A$를 택하자. Nakayama의 보조정리
\ref{algebra-lemma-NAK}에 의해 원소들 $x_\alpha$는 $R$-가군으로서
$M$을 생성한다.
$N = S \otimes_R M$ 및 $I = \mathfrak mS$로 놓자.
그러면 $\{1 \otimes x_\alpha\}_{\alpha \in A}$는 $N$의
원소들의 모임이고 $N/IN$의 기저를 이룬다. 또한 $N$은
$S$ 위에서 평탄하므로
$\text{Tor}_1^S(S/I, N) = 0$이다. 따라서 보조정리
\ref{algebra-lemma-lift-basis}에 의해 $N$은
$\{1 \otimes x_\alpha\}_{\alpha \in A}$를 기저로 하는
자유가군이다. 이제 $R \to S$의 단사성은 $x_\alpha$들 사이에
$R$의 계수를 갖는 자명하지 않은 관계식이 있을 수 없음을
보장한다.
\end{proof}

\noindent
다음 보조정리를 보조정리
\ref{algebra-lemma-criterion-flatness-fibre-Noetherian}
(뇌터 국소환의 경우), 보조정리
\ref{algebra-lemma-criterion-flatness-fibre}
(유한 표시 대수의 경우), 그리고 보조정리
\ref{algebra-lemma-criterion-flatness-fibre-locally-nilpotent}
(국소 멱영 아이디얼의 경우)와 비교하라.

\begin{lemma}[섬유별 평탄성 판정법: 멱영인 경우]
\label{algebra-lemma-criterion-flatness-fibre-nilpotent}
환들의 범주에서 교환 도표
$$
\xymatrix{
S \ar[rr] & & S' \\
& R \ar[lu] \ar[ru]
}
$$
가 주어졌다고 하자. $I \subset R$을 멱영 아이디얼,
$M$을 $S'$-가군이라 하자. 다음을 가정하자.
\begin{enumerate}
\item 가군 $M/IM$은 평탄 $S/IS$-가군이다.
\item 가군 $M$은 평탄 $R$-가군이다.
\end{enumerate}
그러면 $M$은 평탄 $S$-가군이고,
$M \otimes_S \kappa(\mathfrak q)$가 영이 아닌 모든
$\mathfrak q \subset S$에 대해 $S_{\mathfrak q}$는
$R$ 위에서 평탄하다.
\end{lemma}

\begin{proof}
$M$이 $R$ 위에서 평탄하므로 짧은 완전열
$0 \to I \to R \to R/I \to 0$과 텐서하면 짧은 완전열
$$
0 \to I \otimes_R M \to M \to M/IM \to 0.
$$
을 얻는다. $I \otimes_R M \to IS \otimes_S M$은 전사임을
주목하라. 위의 사실과 결합하면
$$
I \otimes_R M \to IS \otimes_S M \to M
$$
의 두 사상이 모두 단사임을 얻는다. 따라서
% P01-E0495: frozen source has Tor_1^S(IS,M), not the quotient required by the criterion.
$\text{Tor}_1^S(IS, M) = 0$이다(주의
\ref{algebra-remark-Tor-ring-mod-ideal}를 보라). 보조정리
\ref{algebra-lemma-what-does-it-mean}에 의해 $M$은 평탄 $S$-가군이다.
마지막으로 $M \otimes_S \kappa(\mathfrak q)$가 영이 아닌 임의의
소 아이디얼 $\mathfrak q \subset S$에 대해 $S_{\mathfrak q}$가
$R$ 위에서 평탄함을 보여야 한다. 이는 보조정리
\ref{algebra-lemma-ff}와 \ref{algebra-lemma-flat-permanence}에서 따른다.
\end{proof}

\section{복합체를 완전하게 만드는 것은 무엇인가?}
\label{algebra-section-complex-exact}

\noindent
이 자료의 일부는 Buchsbaum과 Eisenbud의 논문 \cite{WhatExact}에서
찾을 수 있다.

\begin{situation}
\label{algebra-situation-complex}
여기서 $R$은 환이고, 다음 복합체가 주어져 있다.
$$
0
\to
R^{n_e}
\xrightarrow{\varphi_e}
R^{n_{e-1}}
\xrightarrow{\varphi_{e-1}}
\ldots
\xrightarrow{\varphi_{i + 1}}
R^{n_i}
\xrightarrow{\varphi_i}
R^{n_{i-1}}
\xrightarrow{\varphi_{i-1}}
\ldots
\xrightarrow{\varphi_1}
R^{n_0}
$$
다시 말해 $i = 1, \ldots, e - 1$에 대해
$\varphi_i \circ \varphi_{i + 1} = 0$을 요구한다.
\end{situation}

\begin{lemma}
\label{algebra-lemma-add-trivial-complex}
$R$이 환이라고 하자. 유한 자유 $R$-가군들의 복합체
$$
\ldots
\xrightarrow{\varphi_{i + 1}}
R^{n_i}
\xrightarrow{\varphi_i}
R^{n_{i-1}}
\xrightarrow{\varphi_{i-1}}
\ldots
$$
가 주어졌다고 하자. 어떤 $i$에 대해 사상 $\varphi_i$의 행렬 계수
가운데 하나가 가역이라고 하자. 그러면 위 복합체는 복합체
$$
\ldots \to
R^{n_{i + 2}} \xrightarrow{\varphi_{i + 2}}
R^{n_{i + 1}} \to
R^{n_i - 1} \to
R^{n_{i - 1} - 1} \to
R^{n_{i - 2}} \xrightarrow{\varphi_{i - 2}}
R^{n_{i - 3}} \to
\ldots
$$
와 복합체 $\ldots \to 0 \to R \to R \to 0 \to \ldots$의 직합과
동형이다. 여기서 사상 $R \to R$은 항등사상이다.
\end{lemma}

\begin{proof}
가정은 $R^{n_i}$와 $R^{n_{i-1}}$의 기저를 바꾼 뒤에는
$R^{n_i}$의 첫째 기저벡터가 $\varphi_i$에 의해
$R^{n_{i-1}}$의 첫째 기저벡터로 사상된다는 뜻이다.
$e_j$를 $R^{n_i}$의 $j$번째 기저벡터, $f_k$를
$R^{n_{i-1}}$의 $k$번째 기저벡터라 하자.
$\varphi_i(e_j) = \sum a_{jk} f_k$라 쓰자. 그러면 $k = 1$이 아니면
$a_{1k} = 0$이고 $a_{11} = 1$이다. $j > 1$에 대해
$e'_j = e_j - a_{j1} e_1$로 놓아 $R^{n_i}$의 기저를 다시 바꾸자.
이 좌표변환 뒤에는 $j > 1$에 대해 $a_{j1} = 0$이다.
$R^{n_{i + 1}} \to R^{n_i}$의 상은 $e_j$, $j > 1$이 생성하는
% P01-E0497: frozen source says “subspace” although R is an arbitrary ring.
부분가군에 포함됨을 주목하라. 또한 그 첫째 기저벡터가 상에 속하므로
$R^{n_{i-1}} \to R^{n_{i-2}}$는 $f_1$을 영으로 보내야 한다.
이 조건들과 $\varphi_i$의 행렬 $(a_{jk})$의 모양으로부터 결론이
따른다.
\end{proof}

\noindent
상황 \ref{algebra-situation-complex}에서 다음 꼴의 복합체
$$
0 \to \ldots \to 0 \to R \xrightarrow{1} R \to 0 \to \ldots \to 0
$$
또는 다음 꼴의 복합체
$$
0 \to \ldots \to 0 \to R
$$
를 {\it 자명하다}고 한다. 더 정확히 말해,
$0 \to R^{n_e} \to R^{n_{e-1}} \to \ldots \to R^{n_0}$가
자명하다는 것은, $e \geq i \geq 1$인 어떤 첨자에 대해
$n_i = n_{i - 1} = 1$, $\varphi_i = \text{id}_R$이고
$j \not \in \{i, i - 1\}$인 모든 첨자에 대해 $n_j = 0$이거나,
$n_0 = 1$이고 $i > 0$인 모든 첨자에 대해 $n_i = 0$이라는 뜻이다.
위 보조정리는 국소환 위의 유한 자유가군들의 모든 유한 복합체가
자명 복합체들과의 직합을 무시하면, 모든 사상의 모든 행렬 계수가
극대 아이디얼에 속하는 복합체와 같다는 것을 분명히 말해 준다.

\begin{lemma}
\label{algebra-lemma-exact-depth-zero-local}
상황 \ref{algebra-situation-complex}에서 $R$이 극대 아이디얼 $\mathfrak m$을
갖는 뇌터 국소환이라고 하자. $\mathfrak m \in \text{Ass}(R)$, 즉
$R$의 깊이가 $0$이라고 가정하자.
$0 \to R^{n_e} \to R^{n_{e-1}} \to \ldots \to R^{n_0}$가
$R^{n_e}, \ldots, R^{n_1}$에서 완전하다고 하자. 그러면 이 복합체는
자명 복합체들의 직합과 동형이다.
\end{lemma}

\begin{proof}
$x \in R$, $x \not = 0$, $\mathfrak m x = 0$인 원소를 택하자.
$n_i > 0$인 가장 큰 첨자를 $i$라 하자. $i = 0$이면 주장은 참이다.
$i > 0$이면 $R^{n_i}$의 첫째 기저벡터를 $f_1$이라 하자.
복합체의 완전성에 의해 $xf_1$은 영으로 사상되지 않으므로,
사상 $R^{n_i} \to R^{n_{i - 1}}$의 어떤 행렬 계수는
$\mathfrak m$에 속하지 않는다. 이제 보조정리
\ref{algebra-lemma-add-trivial-complex}를 사용하면
$n_e + \ldots + n_1$을 줄일 수 있다. 귀납법으로 증명이 끝난다.
\end{proof}

\begin{lemma}
\label{algebra-lemma-exact-artinian-local}
상황 \ref{algebra-situation-complex}에서 $R$이 아르틴 국소환이라고 하자.
% P01-E0498: frozen source has “a Artinian”.
$0 \to R^{n_e} \to R^{n_{e-1}} \to \ldots \to R^{n_0}$가
$R^{n_e}, \ldots, R^{n_1}$에서 완전하다고 하자. 그러면 이 복합체는
자명 복합체들의 직합과 동형이다.
\end{lemma}

\begin{proof}
아르틴 국소환의 깊이는 $0$이므로 이는 보조정리
\ref{algebra-lemma-exact-depth-zero-local}의 특수한 경우이다.
\end{proof}

\noindent
이제 유한 자유가군 사이 사상의 계수를 정의한다. 이는 계수의 가능한
정의 가운데 하나일 뿐이다. 이 절에 알맞은 정의일 뿐이며, 다른
상황에서는 다른 정의가 더 편리할 수 있다.

\begin{definition}
\label{algebra-definition-rank}
$R$이 환이고 $\varphi : R^m \to R^n$이 유한 자유가군들의 사상이라고
하자.
\begin{enumerate}
\item $\varphi$의 {\it 계수}는
$\wedge^r \varphi : \wedge^r R^m \to \wedge^r R^n$이 영이 아닌
가장 큰 $r$이다.
\item 위에서 정의한 계수가 $r$일 때, $\varphi$의 행렬의
$r \times r$ 소행렬식들이 생성하는 아이디얼을
$I(\varphi) \subset R$이라 한다.
\end{enumerate}
\end{definition}

\noindent
$\varphi : R^m \to R^n$의 계수가 $0$인 것과 $\varphi = 0$인 것은
동치이며, 이 경우 $I(\varphi) = R$이다.

\begin{lemma}
\label{algebra-lemma-trivial-case-exact}
상황 \ref{algebra-situation-complex}에서 복합체가 자명 복합체들의 직합과
동형이라고 하자. 그러면 다음이 성립한다.
\begin{enumerate}
\item 사상 $\varphi_i$의 계수는
$r_i = n_i - n_{i + 1} + \ldots + (-1)^{e-i-1} n_{e-1} + (-1)^{e-i} n_e$이다.
\item 모든 $i$, $1 \leq i \leq e - 1$에 대해
$\text{rank}(\varphi_{i + 1}) + \text{rank}(\varphi_i) = n_i$이다.
\item 각 $I(\varphi_i) = R$이다.
\end{enumerate}
\end{lemma}

\begin{proof}
복합체 자체가 자명 복합체들의 직합이라고 가정해도 된다. 각 $i$에
대해 $R^{n_i}$의 표준 기저벡터들을 $R^{n_{i-1}}$의 기저벡터로
사상되는 것들과 영으로 사상되는 것들로 나눌 수 있다(그리고 $i > 0$이면
후자는 $R^{n_{i + 1}}$의 기저벡터들의 상이다). $i = e$에서 시작하는
하강 귀납법을 사용하면, $R^{n_i}$의 기저벡터 가운데 영으로 사상되는
것이 $r_{i + 1}$개이고 $R^{n_{i-1}}$의 기저벡터로 사상되는 것이
% P01-E0500: frozen source has the malformed compound “r_{i+1}-basis elements”.
$r_i$개임을 쉽게 보일 수 있다. 이로부터 결과가 따른다.
\end{proof}

\begin{lemma}
\label{algebra-lemma-div-x-exact-one-less}
상황 \ref{algebra-situation-complex}에서 $R$이 극대 아이디얼 $\mathfrak m$을
갖는 국소환이라고 하자.
$0 \to R^{n_e} \to R^{n_{e-1}} \to \ldots \to R^{n_0}$가
$R^{n_e}, \ldots, R^{n_1}$에서 완전하다고 하자.
$x \in \mathfrak m$이 비영인자라고 하자. 그러면 복합체
$0 \to (R/xR)^{n_e} \to \ldots \to (R/xR)^{n_1}$은
$(R/xR)^{n_e}, \ldots, (R/xR)^{n_2}$에서 완전하다.
\end{lemma}

\begin{proof}
항이 $F_i = R^{n_i}$이고 미분이 $\varphi_i$로 주어지는 복합체를
% P01-E0499: frozen source omits “by” in notation declarations.
$F_\bullet$이라 하자. 그러면 복합체들의 짧은 완전열
$$
0 \to F_\bullet \xrightarrow{x} F_\bullet \to F_\bullet/xF_\bullet \to 0
$$
이 있다. 뱀 보조정리를 적용하면 긴 완전열
$$
H_i(F_\bullet) \xrightarrow{x} H_i(F_\bullet) \to
H_i(F_\bullet/xF_\bullet) \to H_{i - 1}(F_\bullet)
\xrightarrow{x} H_{i - 1}(F_\bullet)
$$
을 얻는다. 보조정리는 여기서 따른다.
\end{proof}

\begin{lemma}[비순환성 보조정리]
\label{algebra-lemma-acyclic}
\begin{reference}
\cite[Lemma 1.8]{Peskine-Szpiro}
\end{reference}
$R$이 뇌터 국소환이라고 하자.
$0 \to M_e \to M_{e-1} \to \ldots \to M_0$를 유한 $R$-가군들의
복합체라 하자. $\text{depth}(M_i) \geq i$라고 가정하자.
복합체가 $M_i$에서 완전하지 않은 가장 큰 첨자를 $i$라 하자.
$i > 0$이면
$\Ker(M_i \to M_{i-1})/\Im(M_{i + 1} \to M_i)$의 깊이는
$\geq 1$이다.
\end{lemma}

\begin{proof}
$H = \Ker(M_i \to M_{i-1})/\Im(M_{i + 1} \to M_i)$를 문제의
% P01-E0501: frozen source calls H a cohomology group although this is a homological complex.
호몰로지군이라 하자. 복합체를 다음 완전열들로 나눌 수 있다.
% P01-E0502: frozen source calls every listed sequence short exact, but one ends at M_{i-1}.
$0 \to M_e \to M_{e-1} \to K_{e-2} \to 0$,
$0 \to K_j \to M_j \to K_{j-1} \to 0$, $i + 2 \leq j \leq e-2 $,
$0 \to K_{i + 1} \to M_{i + 1} \to B_i \to 0$,
$0 \to K_i \to M_i \to M_{i-1}$, 그리고
$0 \to B_i \to K_i \to H \to 0$.
보조정리 \ref{algebra-lemma-depth-in-ses}를 반복해서 사용하며 이 열들을
위쪽으로 따라가 깊이에 관한 주장들을 증명한다. 먼저
$\text{depth}(M_e) \geq e$이고 $\text{depth}(M_{e-1}) \geq e-1$이므로
$\text{depth}(K_{e-2}) \geq e - 1$을 얻는다. 이제
$i + 2 \leq j \leq e-2 $에 대한 열
$0 \to K_j \to M_j \to K_{j-1} \to 0$로부터 마찬가지로
$i + 2 \leq j \leq e-2$에 대해
$\text{depth}(K_{j-1}) \geq j$임을 얻는다. 이어서 열
$0 \to K_{i + 1} \to M_{i + 1} \to B_i \to 0$로부터
$\text{depth}(B_i) \geq i + 1$을 얻는다. 가정에 의해 $M_i$의
깊이는 $\geq i \geq 1$이므로, 열
$0 \to K_i \to M_i \to M_{i-1}$은
$\text{depth}(K_i) \geq 1$임을 보여 준다. 마지막으로 열
$0 \to B_i \to K_i \to H \to 0$로부터 결과가 따른다.
\end{proof}

\begin{proposition}
\label{algebra-proposition-what-exact}
\begin{reference}
\cite[Corollary 1]{WhatExact}
\end{reference}
상황 \ref{algebra-situation-complex}에서 $R$이 뇌터 국소환이라고 하자.
다음은 서로 동치이다.
\begin{enumerate}
\item $0 \to R^{n_e} \to R^{n_{e-1}} \to \ldots \to R^{n_0}$가
$R^{n_e}, \ldots, R^{n_1}$에서 완전하다.
\item 모든 $i$, $1 \leq i \leq e$에 대해 다음 두 조건이 성립한다.
\begin{enumerate}
\item $r_i = n_i - n_{i + 1} + \ldots + (-1)^{e-i-1} n_{e-1} + (-1)^{e-i} n_e$일 때
$\text{rank}(\varphi_i) = r_i$이다.
\item $I(\varphi_i) = R$이거나 $I(\varphi_i)$가 길이 $i$인
정칙열을 포함한다.
\end{enumerate}
\end{enumerate}
\end{proposition}

\begin{proof}
어떤 $i$에 대해 $\varphi_i$의 어떤 행렬 계수가 $\mathfrak m$에
속하지 않으면 보조정리 \ref{algebra-lemma-add-trivial-complex}를 적용한다.
주어진 복합체에 대한 명제와 거기에 자명 복합체 하나를 더한 복합체에
대한 명제가 동치임은 쉽게 알 수 있다. 따라서 각 $\varphi_i$의 모든
행렬 성분이 극대 아이디얼의 원소라고 가정해도 된다. 또한
$e \geq 1$이라고 가정해도 된다.

\medskip\noindent
복합체가 $R^{n_e}, \ldots, R^{n_1}$에서 완전하다고 가정하자.
$\mathfrak q \in \text{Ass}(R)$을 택하자. 환 $R_{\mathfrak q}$의
깊이는 $0$이고 $\mathfrak q$에서 국소화한 뒤에도 복합체는
완전하다. 보조정리 \ref{algebra-lemma-exact-depth-zero-local}와
\ref{algebra-lemma-trivial-case-exact}를 $R_{\mathfrak q}$ 위의 국소화된
복합체에 적용한다. 그러면 모든 $i$에 대해
$\varphi_{i, \mathfrak q}$의 계수가 $r_i$임을 얻는다.
$R \to \bigoplus_{\mathfrak q \in \text{Ass}(R)} R_\mathfrak q$가
단사이므로(보조정리 \ref{algebra-lemma-zero-at-ass-zero}), 정의
\ref{algebra-definition-rank}의 계수 정의에 의해 $\varphi_i$의 $R$ 위의
계수도 $r_i$이다. 따라서 계수가 변하지 않으므로
$I(\varphi_i)_\mathfrak q = I(\varphi_{i, \mathfrak q})$이다.
위에서 인용한 보조정리들에 의해 모든 아이디얼
$I(\varphi_i)_{\mathfrak q}$, $e \geq i \geq 1$은
$R_{\mathfrak q}$와 같으므로, 아이디얼 $I(\varphi_i)$ 가운데 어느
것도 $\mathfrak q$에 포함되지 않는다. 따라서
$I(\varphi_e)I(\varphi_{e-1})\ldots I(\varphi_1)$은 $R$의 어떤
연관 소 아이디얼에도 포함되지 않는다. 보조정리 \ref{algebra-lemma-silly}에
의해 모든 $\mathfrak q \in \text{Ass}(R)$에 대해
$x \not \in \mathfrak q$인
$x \in I(\varphi_e)I(\varphi_{e - 1})\ldots I(\varphi_1)$을
택할 수 있다. $x$는 비영인자임을 주목하라(보조정리
\ref{algebra-lemma-ass-zero-divisors}). 보조정리
\ref{algebra-lemma-div-x-exact-one-less}에 따르면 복합체
$0 \to (R/xR)^{n_e} \to \ldots \to (R/xR)^{n_1}$은
$(R/xR)^{n_e}, \ldots, (R/xR)^{n_2}$에서 완전하다. $e$에 대한
귀납법으로 모든 아이디얼 $I(\varphi_i)/xR$은 길이 $i - 1$인
정칙열을 갖는다. 이는 $I(\varphi_i)$가 길이 $i$인 정칙열을
포함함을 증명한다.

\medskip\noindent
(2)(a)와 (2)(b)가 성립한다고 가정하자. $\dim(R)$에 대한 귀납법으로
(1)을 증명한다. $\dim(R) = 0$이면 (2)(b)에 의해
$1 \leq i \leq e$에 대해 $I(\varphi_i) = R$이어야 한다.
$\varphi_i$의 계수들이 극대 아이디얼에 포함되어 있으므로, 이는 모든
$i$에 대해 $r_i = 0$일 때만 가능하다. (2)(a)에 의해 $e = 0$이고
(1)이 성립한다. 이제 $\dim(R) > 0$이라고 하자. 임의의 소 아이디얼
$\mathfrak p \subset R$에 대해, $R_\mathfrak p$ 위에서 사상이
$\varphi_{i, \mathfrak p}$인 복합체
$0 \to R_\mathfrak p^{n_e} \to R_\mathfrak p^{n_{e - 1}} \to \ldots \to
R_\mathfrak p^{n_0}$도 조건 (2)(a)와 (2)(b)를 만족한다고 주장한다.
실제로 $I(\varphi_i)$가 비영인자를 포함하므로 $R_\mathfrak p$에서
$I(\varphi_i)$의 상은 영이 아니다. 따라서
$\varphi_{i, \mathfrak p}$의 계수는 $\varphi_i$의 계수와 같다.
환 $R$ 위의 행렬 $\varphi$에 대해 위와 같이 정의한 계수는
$R$-대수 $R'$로 옮길 때에만 감소할 수 있고, 이는
$I(\varphi)$가 $R'$에서 영으로 사상되는 것과 동치이다. 그러므로
(2)(a)가 성립한다. 이제
$I(\varphi_{i, \mathfrak p}) = I(\varphi_i)_\mathfrak p$이고,
(2)(b)도 국소화 아래 보존됨을 알 수 있다. $R$의 차원에 대한
귀납법으로, $R$의 극대가 아닌 임의의 소 아이디얼 $\mathfrak p$에서
국소화한 복합체가 완전하다고 가정할 수 있다. 따라서
$\Ker(\varphi_i)/\Im(\varphi_{i + 1})$의 지지집합은
$\{\mathfrak m\}$에 포함되고, 영이 아니면 깊이가 $0$이다. 증명의
첫째 문단에서 말한 이유로 모든 $i$에 대해
% P01-E0503: frozen source overlooks rank-zero maps, for which I(varphi_i)=R.
$I(\varphi_i) \subset \mathfrak m$이므로 (2)(b)는
$\text{depth}(R) \geq e$를 함의한다. 보조정리
\ref{algebra-lemma-acyclic}에 의해 복합체는
$R^{n_e}, \ldots, R^{n_1}$에서 완전하며, 증명이 끝난다.
\end{proof}

\begin{remark}
\label{algebra-remark-what-exact}
명제 \ref{algebra-proposition-what-exact}에서 동치 조건 (1)과 (2)가 성립하면,
$I(\varphi_i) = R$인 것과 $i \geq j$인 것이 동치가 되는 $j$가
존재한다. 명제의 증명에서와 마찬가지로, 모든 행렬의 계수가 $R$의
극대 아이디얼 $\mathfrak m$에 속하는 경우만 확인하면 충분하다.
이 경우 $I(\varphi_j) = R$인 것과 $\varphi_j = 0$인 것은 동치이다.
그러나 $\varphi_j = 0$이면 임의로 긴 완전 복합체
$0 \to R^{n_e} \to R^{n_{e-1}} \to \ldots
\to R^{n_j} \to 0 \to 0 \to \ldots \to 0$을 얻는다. 따라서 명제에
의해 $i > j$에 대해 $I(\varphi_i)$는 $R$이어야 한다(그렇지 않으면
뇌터 국소환의 진 아이디얼이 임의로 긴 정칙열을 포함하게 되는데,
이는 불가능하다).
\end{remark}

\section{Cohen--Macaulay 가군}
\label{algebra-section-CM}

\noindent
여기서는 Cohen--Macaulay 가군들이 좋은 성질을 가짐을 보인다.
쌍대성 등과의 관계를 확립하기 위한 Ext 군의 사용은 뒤로 미룬다.

\begin{definition}
\label{algebra-definition-CM}
$R$을 뇌터 국소환, $M$을 유한 $R$-가군이라 하자.
$\dim(\text{Supp}(M)) = \text{depth}(M)$이면 $M$을
{\it Cohen--Macaulay}라고 한다.
\end{definition}

\noindent
첫째 목표는 명제 \ref{algebra-proposition-CM-module}을 확립하는 것이다.
이를 위해 깊이에 관한 앞선 결과를 거의 사용하지 않는, (어쩌면
표준적이지 않은) 초등적 보조정리들의 열을 이용한다. 몇 가지
표기를 도입하자.

\medskip\noindent
$R$을 뇌터 국소환이라 하자. $M$을 Cohen--Macaulay 가군이라 하고,
$f_1, \ldots, f_d$를 $d = \dim(\text{Supp}(M))$인 $M$-정칙열이라 하자.
모든 $i = 0, 1, \ldots, d-1$에 대해
$\dim (\text{Supp}(M) \cap V(g, f_1, \ldots, f_i))
= d - i - 1$이면 $g \in \mathfrak m$을
{\it $(M, f_1, \ldots, f_d)$에 관해 좋은 원소}라고 한다.
이는 $i = 0, 1, \ldots, d - 1$에 대해
$\dim(\text{Supp}(M/(f_1, \ldots, f_i)M) \cap V(g)) =
d - i - 1$이라는 조건과 동치이다.

\begin{lemma}
\label{algebra-lemma-good-element}
표기와 가정은 위와 같다. $g$가 $(M, f_1, \ldots, f_d)$에 관해
좋은 원소이면, (a) $g$는 $M$ 위의 비영인자이고, (b) $M/gM$은
$f_1, \ldots, f_{d - 1}$을 극대 정칙열로 갖는 Cohen--Macaulay
가군이다.
\end{lemma}

\begin{proof}
$d$에 대한 귀납법으로 증명한다.
% P01-E0504: the frozen d=0 claim conflicts with vacuous goodness and can make the conclusion false.
$d = 0$이면 $M$은 유한이고 보조정리가 적용되는 경우가 없다.
$d = 1$이면 $g : M \to M$이 단사임을 보여야 한다. 가정에 의해
$\dim \text{Supp}(M) \cap V(g) = 0$이므로 핵 $K$의 지지집합은
$\{\mathfrak m\}$이다. 따라서 $K$의 길이는 유한하다.
$f_1 : K \to K$가 단사이면 그 상의 길이는 $K$의 길이와 같고,
따라서 $f_1 K = K$이다. Nakayama의 보조정리 \ref{algebra-lemma-NAK}에 의해
$K = 0$이다. 또한 $\dim \text{Supp}(M/gM) = 0$이므로 $M/gM$은
깊이 $0$인 Cohen--Macaulay 가군이다.

\medskip\noindent
$d > 1$이라고 하자. 정의에서 쉽게 알 수 있듯이 $g$는
$(M/f_1M, f_2, \ldots, f_d)$에 관해 좋은 원소이다. 귀납가정에 의해
(a) $g$는 $M/f_1M$ 위의 비영인자이고,
(b) $M/(g, f_1)M$은 $f_2, \ldots, f_{d - 1}$을 극대 정칙열로 갖는
Cohen--Macaulay 가군이다. 보조정리 \ref{algebra-lemma-permute-xi}에 의해
$g, f_1$은 $M$-정칙열이다. 따라서 $g$는 $M$ 위의 비영인자이고
$f_1, \ldots, f_{d - 1}$은 $M/gM$-정칙열이다.
\end{proof}

\begin{lemma}
\label{algebra-lemma-CM-one-g}
$R$을 뇌터 국소환이라 하고, $M$을 $R$ 위의 Cohen--Macaulay
가군이라 하자. $g \in \mathfrak m$이
$\dim(\text{Supp}(M) \cap V(g)) = \dim(\text{Supp}(M)) - 1$을
만족한다고 하자. 그러면 (a) $g$는 $M$ 위의 비영인자이고,
(b) $M/gM$은 깊이가 하나 작은 Cohen--Macaulay 가군이다.
\end{lemma}

\begin{proof}
$d = \dim(\text{Supp}(M))$인 $M$-정칙열
% P01-E0506: frozen source has the wrong English article before M-regular.
$f_1, \ldots, f_d$를 택하자. $g$가
$(M, f_1, \ldots, f_d)$에 관해 좋은 원소이면 보조정리
\ref{algebra-lemma-good-element}에 의해 끝난다. 특히 $d = 1$이면 보조정리가
성립한다. ($d = 0$인 경우는 일어나지 않는다.) $d > 1$이라고 하자.
(\romannumeral1) $h$가 $(M, f_1, \ldots, f_d)$에 관해 좋고,
(\romannumeral2) $\dim(\text{Supp}(M) \cap V(h, g)) = d - 2$인
원소 $h \in R$을 택하자. 이러한 $h$가 존재함을 보이기 위해,
닫힌집합 $\text{Supp}(M)$,
$\text{Supp}(M)\cap V(f_1, \ldots, f_i)$, $i = 1, \ldots, d - 1$,
그리고 $\text{Supp}(M) \cap V(g)$의 극소 소 아이디얼들로 이루어진
(유한) 집합을 $\{\mathfrak q_j\}$라 하자. 이 $\mathfrak q_j$들은
어느 것도 $\mathfrak m$과 같지 않으므로, 보조정리
\ref{algebra-lemma-silly}에 의해 $h \in \mathfrak m$,
$h \not \in \mathfrak q_j$인 원소를 찾을 수 있다. 이 $h$가
(\romannumeral1)과 (\romannumeral2)를 만족함은 분명하다.
보조정리 \ref{algebra-lemma-good-element}에 의해 $M/hM$은
Cohen--Macaulay 가군이다. (\romannumeral2)에 의해 쌍
$(M/hM, g)$는 귀납가정을 만족한다. 따라서 $M/(h, g)M$은
Cohen--Macaulay 가군이고 $g : M/hM \to M/hM$은 단사이다.
보조정리 \ref{algebra-lemma-permute-xi}에 의해 $g : M \to M$과
$h : M/gM \to M/gM$은 단사이다. $M/(g, h)M$이 Cohen--Macaulay
가군이라는 사실과 결합하면 증명이 끝난다.
\end{proof}

\begin{proposition}
\label{algebra-proposition-CM-module}
$R$을 극대 아이디얼 $\mathfrak m$을 갖는 뇌터 국소환이라 하자.
$M$을 지지집합의 차원이 $d$인 $R$ 위의 Cohen--Macaulay 가군이라
하자. $g_1, \ldots, g_c$가 $\mathfrak m$의 원소들이고
$\dim(\text{Supp}(M/(g_1, \ldots, g_c)M)) = d - c$라고 하자.
그러면 $g_1, \ldots, g_c$는 $M$-정칙열이며, 극대 $M$-정칙열로
연장할 수 있다.
\end{proposition}

\begin{proof}
$Z = \text{Supp}(M) \subset \Spec(R)$라 하자. 보조정리
\ref{algebra-lemma-one-equation}에 의해 사슬
$Z \supset Z \cap V(g_1) \supset \ldots \supset Z \cap V(g_1, \ldots, g_c)$의
각 단계에서 차원은 기껏해야 $1$만큼 감소한다. 따라서 가정에 의해
각 단계에서 정확히 $1$만큼 감소한다. 그러므로 가군
% P01-E0505: frozen source omits the terminal module M in this quotient.
$M/(g_1, \ldots, g_i)$와 원소 $g_{i + 1}$에 보조정리
\ref{algebra-lemma-CM-one-g}를 차례로 적용할 수 있다.

\medskip\noindent
$c < d$일 때 $g_1, \ldots, g_c$에 원소 하나를 더 붙여 연장하려면,
보조정리 \ref{algebra-lemma-silly}를 사용하여
$Z \cap V(g_1, \ldots, g_c)$의 유한 개의 극소 소 아이디얼 가운데
어느 것에도 속하지 않는 원소 $g_{c + 1} \in \mathfrak m$을 택하면
된다.
\end{proof}

\noindent
명제 \ref{algebra-proposition-CM-module}을 증명했으므로 이제 표준 이론을
계속 전개한다.

\begin{lemma}
\label{algebra-lemma-nonzerodivisor-on-CM}
$R$을 극대 아이디얼 $\mathfrak m$을 갖는 뇌터 국소환이라 하자.
$M$을 유한 $R$-가군, $x \in \mathfrak m$을 $M$ 위의 비영인자라
하자. 그러면 $M$이 Cohen--Macaulay인 것과 $M/xM$이
Cohen--Macaulay인 것은 동치이다.
\end{lemma}

\begin{proof}
보조정리 \ref{algebra-lemma-depth-drops-by-one}에 의해
$$
\text{depth}(M/xM) = \text{depth}(M)-1
$$
이다. 보조정리
\ref{algebra-lemma-one-equation-module}에 의해
$$
\dim(\text{Supp}(M/xM)) = \dim(\text{Supp}(M)) - 1
$$
이다.
\end{proof}

\begin{lemma}
\label{algebra-lemma-CM-over-quotient}
$R \to S$를 뇌터 국소환들의 전사 준동형이라 하자. $N$을 유한
$S$-가군이라 하자. 그러면 $N$이 $S$-가군으로서 Cohen--Macaulay인
것과 $N$이 $R$-가군으로서 Cohen--Macaulay인 것은 동치이다.
\end{lemma}

\begin{proof}
생략한다.
\end{proof}

\begin{lemma}
\label{algebra-lemma-CM-ass-minimal-support}
\begin{reference}
\cite[Chapter 0, Proposition 16.5.4]{EGA}
\end{reference}
$R$을 뇌터 국소환, $M$을 유한 Cohen--Macaulay $R$-가군이라 하자.
$\mathfrak p \in \text{Ass}(M)$이면
$\dim(R/\mathfrak p) = \dim(\text{Supp}(M))$이고 $\mathfrak p$는
$M$의 지지집합 안의 극소 소 아이디얼이다. 특히 $M$은 매장된 연관
소 아이디얼을 갖지 않는다.
\end{lemma}

\begin{proof}
보조정리 \ref{algebra-lemma-depth-dim-associated-primes}에 의해
$\text{depth}(M) \leq \dim(R/\mathfrak p)$이다. 물론
$\mathfrak p \in \text{Supp}(M)$이므로(보조정리
\ref{algebra-lemma-ass-support})
$\dim(R/\mathfrak p) \leq \dim(\text{Supp}(M))$이다. $M$이
Cohen--Macaulay이므로 두 부등식 모두 등식이다. 따라서
$\mathfrak p$는 $\text{Supp}(M)$에서 극소여야 한다. 그렇지 않으면
$\dim(R/\mathfrak p) < \dim(\text{Supp}(M))$이기 때문이다.
마지막으로, $M$의 지지집합의 극소 소 아이디얼들은
$\text{Ass}(M)$의 극소 원소들과 같다(명제
\ref{algebra-proposition-minimal-primes-associated-primes}). 따라서 $M$은
매장된 연관 소 아이디얼을 갖지 않는다(정의
\ref{algebra-definition-embedded-primes}).
\end{proof}

\begin{definition}
\label{algebra-definition-maximal-CM}
$R$을 뇌터 국소환이라 하자. $R$ 위의 유한 가군 $M$이
$\text{depth}(M) = \dim(R)$을 만족하면 이를 {\it 극대
Cohen--Macaulay} 가군이라고 한다.
\end{definition}

\noindent
다시 말해, 뇌터 국소환 위의 극대 Cohen--Macaulay 가군은 그 환
위에서 가능한 가장 큰 깊이를 갖는 유한 가군이다. 동치로, 뇌터
국소환 $R$ 위의 극대 Cohen--Macaulay 가군은 차원이 그 환의 차원과
같은 Cohen--Macaulay 가군이다. 특히 $M$이
$\Spec(R) = \text{Supp}(M)$인 Cohen--Macaulay $R$-가군이면 $M$은
극대 Cohen--Macaulay이다. 따라서 다음 두 보조정리는 극대
Cohen--Macaulay 가군에 관한 것이다.

\begin{lemma}
\label{algebra-lemma-maximal-chain-maximal-CM}
\begin{slogan}
국소 Cohen--Macaulay 환에서 소 아이디얼의 모든 극대 사슬의 길이는
그 차원과 같다.
\end{slogan}
$R$을 뇌터 국소환이라 하자. $\Spec(R) = \text{Supp}(M)$인
Cohen--Macaulay 가군 $M$이 존재한다고 가정하자. 그러면 소 아이디얼의
모든 극대 사슬 $\mathfrak p_0 \subset
\mathfrak p_1 \subset \ldots \subset \mathfrak p_n$의 길이는
$n = \dim(R)$이다.
\end{lemma}

\begin{proof}
$\dim(R)$에 대한 귀납법으로 증명한다. $\dim(R) = 0$이면 주장은
분명하다. $\dim(R) > 0$이라고 하자. 그러면 $n > 0$이다.
원소 $x \in \mathfrak p_1$을 택하되, $x$는 $R$의 어떤 극소
소 아이디얼에도 속하지 않고 $x \not \in \mathfrak p_0$이게 하자. (보조정리
\ref{algebra-lemma-silly}를 보라.) 보조정리 \ref{algebra-lemma-one-equation}에 의해
$\dim(R/xR) = \dim(R) - 1$이다. 명제
\ref{algebra-proposition-CM-module}과 보조정리
\ref{algebra-lemma-CM-over-quotient}에 의해 $M/xM$은 $R/xR$ 위의
Cohen--Macaulay 가군이다. 보조정리 \ref{algebra-lemma-support-quotient}에
의해 $M/xM$의 지지집합은 $\Spec(R/xR)$이다. 어떤 $n$에 대해
$x$를 $x^n$으로 바꾼 뒤에는 $\mathfrak p_1$이 $M/xM$의 연관
소 아이디얼이라고 가정할 수 있다. 보조정리
\ref{algebra-lemma-inherit-minimal-primes}를 보라. 보조정리
\ref{algebra-lemma-CM-ass-minimal-support}에 의해
$\mathfrak p_1/(x)$은 $R/xR$의 극소 소 아이디얼이다. 따라서 사슬
$\mathfrak p_1/(x) \subset \ldots \subset \mathfrak p_n/(x)$은
$R/xR$에서 소 아이디얼들의 극대 사슬이다. 귀납법에 의해 이 사슬의
길이는 원하는 대로 $\dim(R/xR) = \dim(R) - 1$이다.
\end{proof}

\begin{lemma}
\label{algebra-lemma-dim-formula-maximal-CM}
$R$이 뇌터 국소환이라고 하자. $\Spec(R) = \text{Supp}(M)$인
Cohen--Macaulay 가군 $M$이 존재한다고 가정하자. 그러면
소 아이디얼 $\mathfrak p \subset R$에 대해
$$
\dim(R) = \dim(R_{\mathfrak p}) + \dim(R/\mathfrak p).
$$
\end{lemma}

\begin{proof}
보조정리 \ref{algebra-lemma-maximal-chain-maximal-CM}에서 바로 따른다.
\end{proof}

\begin{lemma}
\label{algebra-lemma-localize-CM-module}
$R$이 뇌터 국소환이라고 하고, $M$을 $R$ 위의 Cohen--Macaulay
가군이라 하자. 임의의 소 아이디얼 $\mathfrak p \subset R$에 대해
$M_{\mathfrak p}$는 $R_\mathfrak p$ 위의 Cohen--Macaulay 가군이다.
\end{lemma}

\begin{proof}
% P01-E0507: frozen source omits the verb in “M not zero”.
$\mathfrak p \not = \mathfrak m$이고 $M$이 영이 아니라고 가정해도
된다. 소 아이디얼의 극대 사슬
$\mathfrak p = \mathfrak p_c \subset
\mathfrak p_{c - 1} \subset \ldots \subset \mathfrak p_1 \subset \mathfrak m$을
택하자. $R_{\mathfrak p_1}$ 위의 $M_{\mathfrak p_1}$에 대해 결과를
증명하면 $c$에 대한 귀납법으로 보조정리가 따른다. 따라서
$\mathfrak p$와 $\mathfrak m$ 사이에 소 아이디얼이 없다고 가정해도
된다. $M_\mathfrak p$의 지지집합 안의 모든 소 아이디얼 사슬에는
$M$의 지지집합 안에서 소 아이디얼 하나(곧 $\mathfrak m$)를 더
붙일 수 있으므로
$\dim(\text{Supp}(M_\mathfrak p)) \leq \dim(\text{Supp}(M)) - 1$이다.
한편 보조정리 \ref{algebra-lemma-depth-localization}와 $\mathfrak p$의 선택에
의해
$\text{depth}(M_\mathfrak p) \geq \text{depth}(M) - \dim(R/\mathfrak p) =
\text{depth}(M) - 1$이다. 따라서 원하는 대로
$\text{depth}(M_\mathfrak p) \geq \dim(\text{Supp}(M_\mathfrak p))$이다
(반대 부등식은 보조정리 \ref{algebra-lemma-bound-depth}이다).
\end{proof}

\begin{definition}
\label{algebra-definition-module-CM}
$R$을 뇌터 환, $M$을 유한 $R$-가군이라 하자. $R$의 모든 소 아이디얼
$\mathfrak p$에 대해 $M_\mathfrak p$가 $R_\mathfrak p$ 위의
Cohen--Macaulay 가군이면 $M$을 {\it Cohen--Macaulay}라고 한다.
\end{definition}

\noindent
보조정리 \ref{algebra-lemma-localize-CM-module}에 의해 $R$의 극대
아이디얼들에서만 이를 확인하면 충분하다.

\begin{lemma}
\label{algebra-lemma-maximal-CM-polynomial-algebra}
$R$을 뇌터 환이라 하고, $M$을 $R$ 위의 Cohen--Macaulay 가군이라
하자. 그러면 $M \otimes_R R[x_1, \ldots, x_n]$은
$R[x_1, \ldots, x_n]$ 위의 Cohen--Macaulay 가군이다.
\end{lemma}

\begin{proof}
변수의 개수에 대한 귀납법으로 $R[x]$ 위의
$M[x] = M \otimes_R R[x]$인 경우만 증명하면 충분하다.
$\mathfrak m \subset R[x]$을 극대 아이디얼이라 하고
$\mathfrak p = R \cap \mathfrak m$이라 하자. $R_{\mathfrak p}$의
극대 아이디얼에서 길이가
$d = \dim(\text{Supp}(M_{\mathfrak p}))$인 $M_\mathfrak p$-정칙열
% P01-E0506: frozen source has the wrong English article before M_p-regular.
$f_1, \ldots, f_d$를 택하자. $R[x]$가 $R$ 위에서 평탄하므로
$R[x]_{\mathfrak m}$은 $R_{\mathfrak p}$ 위에서 평탄하다.
따라서 보조정리 \ref{algebra-lemma-flat-increases-depth}에 의해
$f_1, \ldots, f_d$는 $R[x]_{\mathfrak m}$ 안에서 길이 $d$인
$M[x]_{\mathfrak m}$-정칙열이다. 몫
$$
Q = M[x]_{\mathfrak m}/(f_1, \ldots, f_d)M[x]_{\mathfrak m} =
M_{\mathfrak p}/(f_1, \ldots, f_d)M_{\mathfrak p}
\otimes_{R_\mathfrak p} R[x]_{\mathfrak m}
$$
의 지지집합은 $\mathfrak p$ 위에 놓이는 소 아이디얼들과 같다.
실제로 $R_\mathfrak p \to R[x]_\mathfrak m$은 평탄하고
$M_{\mathfrak p}/(f_1, \ldots, f_d)M_{\mathfrak p}$의 지지집합은
$\{\mathfrak p\}$와 같다(자세한 내용은 생략한다. 도움말: 보조정리
\ref{algebra-lemma-annihilator-flat-base-change}와
\ref{algebra-lemma-support-closed}에서 따른다). 따라서 그 차원은 $1$이다.
증명을 끝내려면 $Q$ 위의 비영인자인 $f \in \mathfrak m$을 찾으면
충분하다. $\mathfrak m$은 극대 아이디얼이므로 체 확장
$\kappa(\mathfrak m)/\kappa(\mathfrak p)$는 유한이다(정리
\ref{algebra-theorem-nullstellensatz}). 따라서 $x$에 관한 다항식으로 보았을
때 최고차항의 계수가 $\mathfrak p$에 속하지 않는
$f \in \mathfrak m$을 찾을 수 있다. 이러한 $f$는
$$
M_{\mathfrak p}/(f_1, \ldots, f_d)M_{\mathfrak p} \otimes_R R[x] =
\bigoplus\nolimits_{n \geq 0}
M_{\mathfrak p}/(f_1, \ldots, f_d)M_{\mathfrak p} \cdot x^n
$$
위의 비영인자로 작용하고, 따라서 $Q$ 위의 비영인자로 작용한다.
\end{proof}

\section{Cohen--Macaulay 환}
\label{algebra-section-CM-ring}

\noindent
이 절의 결과 대부분은 절 \ref{algebra-section-CM}의 결과의 특수한 경우이다.

\begin{definition}
\label{algebra-definition-local-ring-CM}
뇌터 국소환 $R$이 자기 자신 위의 가군으로서 Cohen--Macaulay이면
이를 {\it Cohen--Macaulay}라고 한다.
\end{definition}

\noindent
이는 극대 아이디얼 안에 $R$-정칙열 $x_1, \ldots, x_d$가 존재하여
% P01-E0508: frozen source has the wrong English article before R-regular.
$R/(x_1, \ldots, x_d)$의 차원이 $0$일 것을 요구하는 것과 동치임을
주목하라. 보통은 ``$R$-정칙열''이라 하지 않고 단지 ``정칙열''이라고
말한다.

\begin{lemma}
\label{algebra-lemma-reformulate-CM}
\begin{slogan}
Cohen--Macaulay 국소환의 정칙열은 올바른 차원의 대상을 잘라 낸다는
성질로 특징지어진다.
\end{slogan}
$R$을 극대 아이디얼 $\mathfrak m $을 갖는 뇌터 Cohen--Macaulay
국소환이라 하자. $x_1, \ldots, x_c \in \mathfrak m$을 원소들이라
하자. 그러면
$$
x_1, \ldots, x_c
\text{ is a regular sequence }
\Leftrightarrow
\dim(R/(x_1, \ldots, x_c)) = \dim(R) - c
$$
이다. 이 조건들이 성립하면 $x_1, \ldots, x_c$는 길이
$\dim(R)$인 정칙열로 연장될 수 있고, 각 몫
$R/(x_1, \ldots, x_i)$는 차원이 $\dim(R) - i$인
Cohen--Macaulay 환이다.
\end{lemma}

\begin{proof}
명제 \ref{algebra-proposition-CM-module}의 특수한 경우이다.
\end{proof}

\begin{lemma}
\label{algebra-lemma-maximal-chain-CM}
% P01-E0509: frozen source omits both article and head noun after “Noetherian local”.
$R$을 뇌터 국소환이라 하자. $R$이 차원 $d$인 Cohen--Macaulay라고
하자.
% P01-E0510: frozen theorem says arbitrary ideals although dimension concerns prime ideals.
아이디얼의 모든 극대 사슬 $\mathfrak p_0 \subset
\mathfrak p_1 \subset \ldots \subset \mathfrak p_n$의 길이는
$n = d$이다.
\end{lemma}

\begin{proof}
보조정리 \ref{algebra-lemma-maximal-chain-maximal-CM}의 특수한 경우이다.
\end{proof}

\begin{lemma}
\label{algebra-lemma-CM-dim-formula}
$R$이 차원 $d$인 뇌터 Cohen--Macaulay 국소환이라고 하자. 임의의
소 아이디얼 $\mathfrak p \subset R$에 대해
$$
\dim(R) = \dim(R_{\mathfrak p}) + \dim(R/\mathfrak p).
$$
\end{lemma}

\begin{proof}
보조정리 \ref{algebra-lemma-maximal-chain-CM}에서 바로 따른다. (또한 이는
보조정리 \ref{algebra-lemma-dim-formula-maximal-CM}의 특수한 경우이다.)
\end{proof}

\begin{lemma}
\label{algebra-lemma-localize-CM}
$R$이 Cohen--Macaulay 국소환이라고 하자. 임의의 소 아이디얼
$\mathfrak p \subset R$에 대해 환 $R_{\mathfrak p}$도
Cohen--Macaulay이다.
\end{lemma}

\begin{proof}
보조정리 \ref{algebra-lemma-localize-CM-module}의 특수한 경우이다.
\end{proof}

\begin{definition}
\label{algebra-definition-ring-CM}
뇌터 환 $R$의 모든 국소환이 Cohen--Macaulay이면 이를
{\it Cohen--Macaulay}라고 한다.
\end{definition}

\begin{lemma}
\label{algebra-lemma-CM-polynomial-algebra}
$R$이 뇌터 Cohen--Macaulay 환이라고 하자. $R$ 위의 모든
다항식 대수는 Cohen--Macaulay이다.
\end{lemma}

\begin{proof}
보조정리 \ref{algebra-lemma-maximal-CM-polynomial-algebra}의 특수한 경우이다.
\end{proof}

\begin{lemma}
\label{algebra-lemma-dimension-shift}
$R$을 차원 $d$인 뇌터 Cohen--Macaulay 국소환이라 하자.
$0 \to K \to R^{\oplus n} \to M \to 0$를 $R$-가군들의 완전열이라
하자. 그러면 $M = 0$이거나 $\text{depth}(K) > \text{depth}(M)$이거나
$\text{depth}(K) = \text{depth}(M) = d$이다.
\end{lemma}

\begin{proof}
이는 보조정리 \ref{algebra-lemma-depth-in-ses}의 특수한 경우이다.
\end{proof}

\begin{lemma}
\label{algebra-lemma-mcm-resolution}
$R$을 차원 $d$인 뇌터 Cohen--Macaulay 국소환이라 하자. $M$을
깊이가 $e$인 유한 $R$-가군이라 하자. 각 $F_i$가 유한 자유이고
$K$가 극대 Cohen--Macaulay인 완전 복합체
% P01-E0511: frozen displayed index d-e-1 is undefined in the case e=d.
$$
0 \to K \to F_{d-e-1} \to \ldots \to F_0 \to M \to 0
$$
가 존재한다.
\end{lemma}

\begin{proof}
정의와 보조정리 \ref{algebra-lemma-dimension-shift}에서 바로 따른다.
\end{proof}

\begin{lemma}
\label{algebra-lemma-find-sequence-image-regular}
$\varphi : A \to B$를 국소환들의 사상이라 하자. $B$가 뇌터이자
Cohen--Macaulay이고
$\mathfrak m_B = \sqrt{\varphi(\mathfrak m_A) B}$라고 가정하자.
그러면 $\varphi(f_1), \ldots, \varphi(f_{\dim(B)})$가 $B$에서
정칙열이 되는 원소들의 열 $f_1, \ldots, f_{\dim(B)}$가 $A$에
존재한다.
\end{lemma}

\begin{proof}
$\dim(B)$에 대한 귀납법으로 다음을 증명하면 충분하다.
$\dim(B) \geq 1$이면 $B$에서 비영인자로 사상되는 원소 $f$를
$A$에서 찾을 수 있다. 보조정리 \ref{algebra-lemma-reformulate-CM}에 의해,
$B$의 유한 개의 극소 소 아이디얼
$\mathfrak q_1, \ldots, \mathfrak q_r$ 가운데 어느 것에도 그 상이
속하지 않는 $f \in A$를 찾으면 충분하다. 여기서 그 상은 $B$에서
취한다. 가정
$\mathfrak m_B = \sqrt{\varphi(\mathfrak m_A) B}$에 의해
$\mathfrak m_A \not \subset \varphi^{-1}(\mathfrak q_i)$이다.
그러므로 보조정리 \ref{algebra-lemma-silly}에 의해 $f$를 찾을 수 있다.
\end{proof}

\section{카테너리 환}
\label{algebra-section-catenary}

\noindent
Topology, 절 \ref{topology-section-catenary-spaces}와 비교하라.

\begin{definition}
\label{algebra-definition-catenary}
환 $R$이 다음 조건을 만족하면 {\it 카테너리}라고 한다. 소 아이디얼의
임의의 쌍 $\mathfrak p \subset \mathfrak q$에 대해 소 아이디얼들의
모든 유한 사슬
$\mathfrak p = \mathfrak p_0 \subset \mathfrak p_1 \subset \ldots \subset
\mathfrak p_e = \mathfrak q$의 길이를 제한하는 정수가 존재하고,
그러한 모든 극대 사슬의 길이가 같다.
\end{definition}

\begin{lemma}
\label{algebra-lemma-catenary}
환 $R$이 카테너리인 것과 위상공간 $\Spec(R)$이 카테너리인 것은
동치이다(Topology, 정의 \ref{topology-definition-catenary}를 보라).
\end{lemma}

\begin{proof}
정의와 보조정리 \ref{algebra-lemma-irreducible}의 기약 닫힌 부분집합의
특징짓기에서 바로 따른다.
\end{proof}

\noindent
일반적으로 $R$이 카테너리라고 해서 유한 생성 $R$-대수가
카테너리인 것은 아니다. 따라서 다음 정의를 둔다.

\begin{definition}
\label{algebra-definition-universally-catenary}
뇌터 환 $R$ 위의 모든 유한형 $R$-대수가 카테너리이면 그 환을
{\it 보편적으로 카테너리}라고 한다.
\end{definition}

\noindent
뇌터가 아닌 환에 대해 이 정의가 올바른 것인지는 분명하지 않으므로
뇌터 환으로 한정한다. 보조정리 \ref{algebra-lemma-quotient-catenary}에 의해
뇌터 환 $R$이 보편적으로 카테너리임을 확인하려면 각 다항식 대수
$R[x_1, \ldots, x_n]$이 카테너리임을 확인하면 충분하다.

\begin{lemma}
\label{algebra-lemma-localization-catenary}
카테너리 환의 모든 국소화는 카테너리이다. 뇌터 보편적 카테너리
환의 모든 국소화는 보편적으로 카테너리이다.
\end{lemma}

\begin{proof}
$A$를 환, $S \subset A$를 곱셈적 부분집합이라 하자. 보조정리
\ref{algebra-lemma-spec-localization}의 $\Spec(S^{-1}A)$에 대한 서술은
$A$가 카테너리이면 $S^{-1}A$도 카테너리임을 보여 준다.
$S^{-1}A \to C$가 유한형이면, 어떤 유한형 환 사상 $A \to B$에
대해 $C = S^{-1}B$이다. 따라서 $A$가 뇌터이자 보편적으로
카테너리이면 $B$는 카테너리이고, $C$도 카테너리이다. 이로써
보조정리가 증명된다.
\end{proof}

\begin{lemma}
\label{algebra-lemma-universally-catenary}
$A$를 뇌터 보편적 카테너리 환이라 하자. $A$ 위에서 본질적으로
유한형인 모든 $A$-대수는 보편적으로 카테너리이다.
\end{lemma}

\begin{proof}
$B$가 유한형 $A$-대수이면 보조정리
\ref{algebra-lemma-Noetherian-permanence}에 의해 $B$는 뇌터이다. 모든 유한형
$B$-대수는 유한형 $A$-대수이고, $A$가 보편적으로 카테너리라는
가정에 의해 카테너리이다. 따라서 $B$는 보편적으로 카테너리이다.
보조정리 \ref{algebra-lemma-localization-catenary}에 의해 $B$의 모든
국소화는 보편적으로 카테너리이며, 이로써 증명이 끝난다.
\end{proof}

\begin{lemma}
\label{algebra-lemma-catenary-check-local}
$R$을 환이라 하자. 다음은 서로 동치이다.
\begin{enumerate}
\item $R$은 카테너리이다.
\item 모든 소 아이디얼 $\mathfrak p$에 대해 $R_\mathfrak p$는
카테너리이다.
\item 모든 극대 아이디얼 $\mathfrak m$에 대해 $R_\mathfrak m$은
카테너리이다.
\end{enumerate}
$R$이 뇌터라고 가정하자. 다음은 서로 동치이다.
\begin{enumerate}
\item $R$은 보편적으로 카테너리이다.
\item 모든 소 아이디얼 $\mathfrak p$에 대해 $R_\mathfrak p$는
보편적으로 카테너리이다.
\item 모든 극대 아이디얼 $\mathfrak m$에 대해 $R_\mathfrak m$은
보편적으로 카테너리이다.
\end{enumerate}
\end{lemma}

\begin{proof}
두 경우 모두 (1) $\Rightarrow$ (2)는 보조정리
\ref{algebra-lemma-localization-catenary}에서 따르고, (2) $\Rightarrow$ (3)은
바로 따른다. $R$의 모든 극대 아이디얼 $\mathfrak m$에 대해
$R_\mathfrak m$이 카테너리라고 하자. $R$ 안의 소 아이디얼
$\mathfrak p \subset \mathfrak q$가 주어지면 극대 아이디얼
$\mathfrak q \subset \mathfrak m$을 택한다.
% P01-E0512: frozen source says “chains of primes ideals”.
$\mathfrak p$와 $\mathfrak q$ 사이의 소 아이디얼 사슬들은
$\mathfrak pR_\mathfrak m$과 $\mathfrak qR_\mathfrak m$ 사이의
소 아이디얼 사슬들과 일대일 대응하므로 $R$은 카테너리이다.
이제 $R$이 뇌터이고 $R$의 모든 극대 아이디얼 $\mathfrak m$에 대해
$R_\mathfrak m$이 보편적으로 카테너리라고 하자. $R \to S$를
유한형 환 사상이라 하고, $S$의 소 아이디얼 $\mathfrak q$가
소 아이디얼 $\mathfrak p \subset R$ 위에 놓인다고 하자. $R$에서
극대 아이디얼 $\mathfrak p \subset \mathfrak m$을 택하자.
$R_\mathfrak p$는 $R_\mathfrak m$의 국소화이므로 보조정리
\ref{algebra-lemma-localization-catenary}에 의해 보편적으로 카테너리이다.
그러면 $S_\mathfrak p$는 $R_\mathfrak p$ 위의 유한형 환으로서
카테너리이다. 따라서 $S_\mathfrak q$는 국소화로서 카테너리이다.
위에서 다룬 첫째 경우에 의해 $S$는 카테너리이다.
\end{proof}

\begin{lemma}
\label{algebra-lemma-quotient-catenary}
카테너리 환의 모든 몫은 카테너리이다. 뇌터 보편적 카테너리 환의
모든 몫은 보편적으로 카테너리이다.
\end{lemma}

\begin{proof}
$A$를 환, $I \subset A$를 아이디얼이라 하자. 보조정리
\ref{algebra-lemma-spec-closed}의 $\Spec(A/I)$에 대한 서술은 $A$가
카테너리이면 $A/I$도 카테너리임을 보여 준다. 둘째 주장은
보조정리 \ref{algebra-lemma-universally-catenary}의 특수한 경우이다.
\end{proof}

\begin{lemma}
\label{algebra-lemma-catenary-check-irreducible}
$R$을 뇌터 환이라 하자.
\begin{enumerate}
\item $R$이 카테너리인 것과 모든 극소 소 아이디얼 $\mathfrak p$에
대해 $R/\mathfrak p$가 카테너리인 것은 동치이다.
\item $R$이 보편적으로 카테너리인 것과 모든 극소 소 아이디얼
$\mathfrak p$에 대해 $R/\mathfrak p$가 보편적으로 카테너리인 것은
동치이다.
\end{enumerate}
\end{lemma}

\begin{proof}
$\mathfrak a \subset \mathfrak b$가 $R$의 소 아이디얼들의 포함이면
극소 소 아이디얼 $\mathfrak p \subset \mathfrak a$를 찾을 수 있으므로
첫째 주장은 분명하다. 둘째 주장의 증명은 생략한다.
\end{proof}

\begin{lemma}
\label{algebra-lemma-CM-ring-catenary}
뇌터 Cohen--Macaulay 환은 보편적으로 카테너리이다. 더 일반적으로,
$R$이 뇌터 환이고 $M$이
$\text{Supp}(M) = \Spec(R)$인 Cohen--Macaulay $R$-가군이면
$R$은 보편적으로 카테너리이다.
\end{lemma}

\begin{proof}
보조정리 \ref{algebra-lemma-CM-polynomial-algebra}에 의해 $R$ 위의 다항식
대수는 Cohen--Macaulay이므로, Cohen--Macaulay 환이 카테너리임을
보이면 충분하다. $R$을 Cohen--Macaulay 환이라 하고
$\mathfrak p \subset \mathfrak q$를 $R$의 소 아이디얼들이라 하자.
정의에 의해 $R_{\mathfrak q}$와 $R_{\mathfrak p}$는
Cohen--Macaulay이다. 소 아이디얼들의 극대 사슬
$\mathfrak p = \mathfrak p_0 \subset \mathfrak p_1 \subset
\ldots \subset \mathfrak p_n = \mathfrak q$를 택하자. 이어서
소 아이디얼들의 극대 사슬
$\mathfrak q_0 \subset \mathfrak q_1 \subset \ldots \subset
\mathfrak q_m = \mathfrak p$를 택하자. 보조정리
\ref{algebra-lemma-maximal-chain-CM}에 의해
$n + m = \dim(R_{\mathfrak q})$이고, 같은 보조정리에 의해
$m = \dim(R_{\mathfrak p})$이다. 따라서
$n = \dim(R_{\mathfrak q}) - \dim(R_{\mathfrak p})$은 선택과
무관하다.

\medskip\noindent
더 일반적인 주장을 증명하려면 위와 똑같이 논하되, 보조정리
\ref{algebra-lemma-maximal-CM-polynomial-algebra}와
\ref{algebra-lemma-maximal-chain-maximal-CM}을 사용한다.
\end{proof}

\begin{lemma}
\label{algebra-lemma-catenary-Noetherian-local}
$(A, \mathfrak m)$을 뇌터 국소환이라 하자. 다음은 서로 동치이다.
\begin{enumerate}
\item $A$는 카테너리이다.
\item $\mathfrak p \mapsto \dim(A/\mathfrak p)$는 $\Spec(A)$ 위의
차원 함수이다.
\end{enumerate}
\end{lemma}

\begin{proof}
$A$가 카테너리이면 Topology, 보조정리
\ref{topology-lemma-locally-dimension-function}와 보조정리
\ref{algebra-lemma-catenary}에 의해 $\Spec(A)$는 차원 함수 $\delta$를 갖는다.
$\delta(\mathfrak m) = 0$이라고 가정해도 된다. 그러면 Topology,
보조정리 \ref{topology-lemma-dimension-function-catenary}에 의해
$$
\delta(\mathfrak p) = \text{codim}(V(\mathfrak m), V(\mathfrak p)) =
\dim(A/\mathfrak p)
$$
이다. 이렇게 하여 (1)이 (2)를 함의함을 알 수 있다. 역함의도
Topology, 보조정리
\ref{topology-lemma-dimension-function-catenary}에서 따른다.
\end{proof}

\section{정칙 국소환}
\label{algebra-section-regular}

\noindent
정칙 국소환은 정의 \ref{algebra-definition-regular-local}에서 정의했다.
정칙 국소환의 모든 소 아이디얼 국소화가 정칙임을 보이는 일은 그리
쉽지 않다. 사실 지금까지 전개한 자료의 상당 부분은 이 사실의 증명을
위한 것이다. 명제 \ref{algebra-proposition-finite-gl-dim-regular}을 보고,
그곳의 인용들을 거슬러 올라가라.

\begin{lemma}
\label{algebra-lemma-regular-graded}
$(R, \mathfrak m, \kappa)$를 차원 $d$인 정칙 국소환이라 하자.
등급환 $\bigoplus \mathfrak m^n / \mathfrak m^{n + 1}$은 등급
다항식 대수 $\kappa[X_1, \ldots, X_d]$와 동형이다.
\end{lemma}

\begin{proof}
$x_1, \ldots, x_d$를 극대 아이디얼 $\mathfrak m$의 극소
생성집합이라 하자. 정의 \ref{algebra-definition-regular-local}을 보라.
$X_i$를 $\mathfrak m/\mathfrak m^2$에서 $x_i$의 동치류로 보내는
전사
$\kappa[X_1, \ldots, X_d] \to \bigoplus \mathfrak m^n/\mathfrak m^{n + 1}$가
있다. 명제 \ref{algebra-proposition-dimension}에 의해 $d(R) = d$이므로
수치다항식 $n \mapsto \dim_\kappa \mathfrak m^n/\mathfrak m^{n + 1}$의
차수는 $d - 1$이다. 보조정리 \ref{algebra-lemma-quotient-smaller-d}에 의해
전사 $\kappa[X_1, \ldots, X_d]
\to \bigoplus \mathfrak m^n/\mathfrak m^{n + 1}$는 동형이다.
\end{proof}

\begin{lemma}
\label{algebra-lemma-regular-domain}
모든 정칙 국소환은 정역이다.
\end{lemma}

\begin{proof}
보조정리 \ref{algebra-lemma-intersect-powers-ideal-module-zero}에 의해
$\bigcap \mathfrak m^n = 0$임을 사용한다. $fg = 0$인
$f, g \in R$을 택하자.
% P01-E0513: frozen proof omits the necessary reduction to nonzero f and g.
둘 중 하나가 영이면 자명하므로 둘 다 영이 아니라고 가정한다.
$a, b$가 극대가 되도록 $f \in \mathfrak m^a$이고
$g \in \mathfrak m^b$라고 하자. $fg = 0 \in \mathfrak m^{a + b + 1}$이므로
보조정리 \ref{algebra-lemma-regular-graded}의 결과에 의해
$f \in \mathfrak m^{a + 1}$이거나 $g \in \mathfrak m^{b + 1}$이다.
이는 모순이다.
\end{proof}

\begin{lemma}
\label{algebra-lemma-regular-ring-CM}
$R$을 정칙 국소환이라 하고, $x_1, \ldots, x_d$를 극대 아이디얼
$\mathfrak m$의 극소 생성집합이라 하자. 그러면
$x_1, \ldots, x_d$는 정칙열이고, 각 $R/(x_1, \ldots, x_c)$는
차원 $d - c$인 정칙 국소환이다. 특히 $R$은 Cohen--Macaulay이다.
\end{lemma}

\begin{proof}
보조정리 \ref{algebra-lemma-one-equation}에 의해 $R/x_1R$은 차원이
$\geq d - 1$인 뇌터 국소환이고, $x_2, \ldots, x_d$가 그 극대
아이디얼을 생성한다. 따라서 정의에 의해 정칙 국소환이다.
보조정리 \ref{algebra-lemma-regular-domain}에 의해 $R$은 정역이므로
$x_1$은 비영인자이다.
\end{proof}

\begin{lemma}
\label{algebra-lemma-regular-quotient-regular}
$R$을 정칙 국소환이라 하자. $I \subset R$을 $R/I$도 정칙
국소환이 되게 하는 아이디얼이라 하자. 그러면 어떤
$0 \leq c \leq d$에 대해 $I = (x_1, \ldots, x_c)$이고
$x_1, \ldots, x_d$가 $R$의 극대 아이디얼 $\mathfrak m$의 극소
생성집합이 되도록 택할 수 있다.
\end{lemma}

\begin{proof}
$\dim(R) = d$이고 $\dim(R/I) = d - c$라고 하자.
$R/I$의 극대 아이디얼을
$\overline{\mathfrak m} = \mathfrak m/I$로 나타내고,
$\kappa = R/\mathfrak m$이라 하자. 정칙 국소환의 정의에 의해
$$
\dim_\kappa((I + \mathfrak m^2)/\mathfrak m^2) =
\dim_\kappa(\mathfrak m/\mathfrak m^2)
- \dim(\overline{\mathfrak m}/\overline{\mathfrak m}^2) = d - (d - c) = c
$$
이다. 따라서 $\mathfrak m/\mathfrak m^2$에서 상들이 선형독립인
$x_1, \ldots, x_c \in I$를 택하고, $x_{c + 1}, \ldots, x_d$를
더하여 $\mathfrak m$의 극소 생성집합을 얻을 수 있다. 유도된 사상
$R/(x_1, \ldots, x_c) \to R/I$는 같은 차원의 정칙 국소환들 사이의
전사이다(보조정리 \ref{algebra-lemma-regular-ring-CM}). 따라서 그 핵은
영이다. 즉 $I = (x_1, \ldots, x_c)$이다. 그렇지 않으면 보조정리
\ref{algebra-lemma-regular-domain}과 \ref{algebra-lemma-one-equation}에 의해
$\dim(R/I) < \dim(R/(x_1, \ldots, x_c))$이기 때문이다.
\end{proof}

\begin{lemma}
\label{algebra-lemma-free-mod-x}
$R$을 뇌터 국소환, $x \in \mathfrak m$을 원소라 하자. $M$을 유한
$R$-가군이라 하되, $x$가 $M$ 위의 비영인자이고 $M/xM$이 $R/xR$
위에서 자유라고 하자. 그러면 $M$은 $R$ 위에서 자유이다.
\end{lemma}

\begin{proof}
$M/xM$의 $R/xR$-기저로 사상되는 $M$의 원소들을
$m_1, \ldots, m_r$라 하자. Nakayama의 보조정리 \ref{algebra-lemma-NAK}에
의해 $m_1, \ldots, m_r$는 $M$을 생성한다. 관계
$\sum a_i m_i = 0$이 있으면 모든 $i$에 대해 $a_i \in xR$이다.
따라서 어떤 $b_i \in R$에 대해 $a_i = b_i x$이다. 그러므로
$R^r \to M$의 핵 $K$는 $xK = K$를 만족하고, Nakayama의 보조정리에
의해 영이다.
\end{proof}

\begin{lemma}
\label{algebra-lemma-regular-mcm-free}
$R$을 정칙 국소환이라 하자. $R$ 위의 모든 극대 Cohen--Macaulay
가군은 자유이다.
\end{lemma}

\begin{proof}
$M$을 $R$ 위의 극대 Cohen--Macaulay 가군이라 하자.
$x \in \mathfrak m$을 $\mathfrak m$을 생성하는 정칙열의 일부라
하자. 명제 \ref{algebra-proposition-CM-module}에 의해 $x$는 $M$ 위의
비영인자이고, $M/xM$은 $R/xR$ 위의 극대 Cohen--Macaulay 가군이다.
$\dim(R)$에 대한 귀납법으로 $M/xM$이 자유임을 알 수 있다.
보조정리 \ref{algebra-lemma-free-mod-x}에 의해 끝난다.
\end{proof}

\begin{lemma}
\label{algebra-lemma-regular-mod-x}
$R$이 뇌터 국소환이라고 하자. $x \in \mathfrak m$이 비영인자이고
$R/xR$이 정칙 국소환이라고 하자. 그러면 $R$은 정칙 국소환이다.
더 일반적으로 $x_1, \ldots, x_r$가 $R$의 정칙열이고
$R/(x_1, \ldots, x_r)$가 정칙 국소환이면 $R$은 정칙 국소환이다.
\end{lemma}

\begin{proof}
$x$와 $R/xR$의 극대 아이디얼의 극소 생성계의 원상들을 합치면
$\mathfrak m$의 생성원 $\dim(R)$개를 얻기 때문이다. 보조정리
\ref{algebra-lemma-one-equation}을 사용하라. 마지막 주장은 첫째 주장과
귀납법에서 따른다.
\end{proof}

\begin{lemma}
\label{algebra-lemma-colimit-regular}
$(R_i, \varphi_{ii'})$를 전이사상이 국소환 사상인 국소환들의
유향계라 하자. 각 $R_i$가 정칙 국소환이고
$R = \colim R_i$가 뇌터이면 $R$은 정칙 국소환이다.
\end{lemma}

\begin{proof}
$\mathfrak m \subset R$을 극대 아이디얼이라 하자. 이는 극대
아이디얼 $\mathfrak m_i \subset R_i$들의 여극한이다.
$d = \dim \mathfrak m/\mathfrak m^2$에 대한 귀납법으로 증명한다.
$d = 0$이면 $R = R/\mathfrak m$은 체이고 $R$은 정칙 국소환이다.
$d > 0$이면 $x \in \mathfrak m$, $x \not \in \mathfrak m^2$인
원소를 택하자. 어떤 $i$에 대해 $x$로 사상되는
$x_i \in \mathfrak m_i$를 찾을 수 있다. 환
$R/xR = \colim_{i' \geq i} R_{i'}/x_iR_{i'}$은 뇌터 국소환이다.
보조정리 \ref{algebra-lemma-regular-ring-CM}에 의해
$R_{i'}/x_iR_{i'}$는 정칙 국소환이다. 따라서 귀납법에 의해
$R/xR$은 정칙 국소환이다. 각 $R_i$가 정역이므로(보조정리
\ref{algebra-lemma-regular-graded}) $R$은 정역이다. 따라서 $x$는
비영인자이고, 보조정리 \ref{algebra-lemma-regular-mod-x}에 의해 $R$은
정칙 국소환이다.
\end{proof}





\section{환의 epimorphism}
\label{algebra-section-epimorphism}

\noindent
모든 범주에는 {\it epimorphism}이라는 개념이 있다.
이 절의 일부 내용은 \cite{Autour}와 \cite{Mazet}에서 가져왔다.

\begin{lemma}
\label{algebra-lemma-epimorphism}
$R \to S$를 환 사상이라 하자. 다음 조건들은 동치이다.
\begin{enumerate}
\item $R \to S$는 epimorphism이다.
\item 두 환 사상 $S \to S \otimes_R S$는 같다.
\item 두 환 사상 $S \to S \otimes_R S$ 중 어느 하나는 동형사상이다.
\item 환 사상 $S \otimes_R S \to S$는 동형사상이다.
\end{enumerate}
\end{lemma}

\begin{proof}
생략한다.
\end{proof}

\begin{lemma}
\label{algebra-lemma-composition-epimorphism}
환의 두 epimorphism의 합성은 epimorphism이다.
\end{lemma}

\begin{proof}
생략한다. 힌트: 이는 모든 범주에서 참이다.
\end{proof}

\begin{lemma}
\label{algebra-lemma-base-change-epimorphism}
$R \to S$가 환의 epimorphism이고 $R \to R'$이 임의의 환 사상이면,
$R' \to R' \otimes_R S$는 epimorphism이다.
\end{lemma}

\begin{proof}
생략한다. 힌트: 밂이 존재하는 모든 범주에서 참이다.
\end{proof}

\begin{lemma}
\label{algebra-lemma-permanence-epimorphism}
$A \to B \to C$가 환 사상들이고 $A \to C$가 epimorphism이면
$B \to C$도 epimorphism이다.
\end{lemma}

\begin{proof}
생략한다. 힌트: 이는 모든 범주에서 참이다.
\end{proof}

\noindent
특히 이는 $R \to S$가 epimorphism이고 그 상이
$\overline{R} \subset S$이면 $\overline{R} \to S$가 epimorphism임을 뜻한다.
따라서 epimorphism에 관한 결과를 증명할 때에는 흔히 그 사상이 단사라고
가정해도 된다. 다음 보조정리는 특히 모든 국소화가 epimorphism임을 뜻한다.

\begin{lemma}
\label{algebra-lemma-epimorphism-local}
$R \to S$를 환 사상이라 하자. 다음 조건들은 동치이다.
\begin{enumerate}
\item $R \to S$는 epimorphism이다.
\item $R$의 각 소 아이디얼 $\mathfrak p$에 대해
$R_{\mathfrak p} \to S_{\mathfrak p}$는 epimorphism이다.
\end{enumerate}
\end{lemma}

\begin{proof}
$S_{\mathfrak p} = R_{\mathfrak p} \otimes_R S$이므로(보조정리
\ref{algebra-lemma-tensor-localization} 참조), 보조정리
\ref{algebra-lemma-base-change-epimorphism}에 의해 (1)은 (2)를 함의한다.
역으로 (2)가 성립한다고 하자. $a, b : S \to A$를 환 $S$에서 환 $A$로
가는 두 환 사상으로서 $R \to S$를 등화한다고 하자. 가정에 의해
$R$의 모든 소 아이디얼 $\mathfrak p$에 대하여 유도된 사상
$a_{\mathfrak p}, b_{\mathfrak p} : S_{\mathfrak p} \to A_{\mathfrak p}$는
같다. 보조정리 \ref{algebra-lemma-characterize-zero-local}에서 보듯이
$A \subset \prod_{\mathfrak p} A_{\mathfrak p}$이므로 $a = b$이다.
\end{proof}

\begin{lemma}
\label{algebra-lemma-finite-epimorphism-surjective}
\begin{slogan}
환 사상이 전사일 필요충분조건은 그것이 유한 epimorphism인 것이다.
\end{slogan}
$R \to S$를 환 사상이라 하자. 다음 조건들은 동치이다.
\begin{enumerate}
\item $R \to S$는 epimorphism이고 유한이다.
\item $R \to S$는 전사이다.
\end{enumerate}
\end{lemma}

\begin{proof}
(이 보조정리는 문헌에서 여러 번 재증명된 듯하며 서로 다른 증명도 많다.)
전사인 환 사상이 epimorphism이라는 것은 명백하다.
$R \to S$가 유한 환 사상이고 $S \otimes_R S \to S$가 동형사상이라고
하자. 목표는 $R \to S$가 전사임을 보이는 것이다. $S/R$이 영이 아니라고
가정하자. 완전열 $R \to S \to S/R \to 0$에서 완전열
$$
R \otimes_R S \to S \otimes_R S \to S/R \otimes_R S \to 0.
$$
을 얻는다. 가정에 의해 첫째 화살표는 동형사상이므로
$S/R \otimes_R S = 0$이다. 따라서 $S/R \otimes_R S/R = 0$이기도 하다.
보조정리 \ref{algebra-lemma-trivial-filter-finite-module}에 의해 어떤 진 아이디얼
$I \subset R$에 대해 $R$-가군의 전사 $S/R \to R/I$가 존재한다.
따라서 전사
$S/R \otimes_R S/R \to R/I \otimes_R R/I = R/I \not = 0$이 존재하여
모순이다.
\end{proof}

\begin{lemma}
\label{algebra-lemma-faithfully-flat-epimorphism}
충실 평탄 epimorphism은 동형사상이다.
\end{lemma}

\begin{proof}
보조정리 \ref{algebra-lemma-epimorphism}의 (3)에서 명백하다. 실제로 사상
$S \to S \otimes_R S$는 $R \to S$에 $S$를 텐서한 사상이다.
\end{proof}

\begin{lemma}
\label{algebra-lemma-epimorphism-over-field}
$k \to S$가 epimorphism이고 $k$가 체이면 $S = k$ 또는 $S = 0$이다.
\end{lemma}

\begin{proof}
이는 보조정리 \ref{algebra-lemma-faithfully-flat-epimorphism}의 결과에서
명백하다($k$ 위의 영이 아닌 모든 대수는 충실 평탄이다). 또는
$\dim_k(R) \leq 1$이 아니면 $R \to R \otimes_k R$이 epimorphism일 수 없음을
직접 논증해도 된다.
\end{proof}

\begin{lemma}
\label{algebra-lemma-epimorphism-injective-spec}
$R \to S$를 환의 epimorphism이라 하자. 그러면
\begin{enumerate}
\item $\Spec(S) \to \Spec(R)$는 단사이고,
\item $\mathfrak p \subset R$ 위에 놓이는 $\mathfrak q \subset S$에 대해
$\kappa(\mathfrak p) = \kappa(\mathfrak q)$이다.
\end{enumerate}
\end{lemma}

\begin{proof}
$\mathfrak p$를 $R$의 소 아이디얼이라 하자. 사상의 섬유는 섬유환
$S \otimes_R \kappa(\mathfrak p)$의 스펙트럼이다. 보조정리
\ref{algebra-lemma-base-change-epimorphism}에 의해 사상
$\kappa(\mathfrak p) \to S \otimes_R \kappa(\mathfrak p)$는
epimorphism이고, 따라서 보조정리 \ref{algebra-lemma-epimorphism-over-field}에 의해
$S \otimes_R \kappa(\mathfrak p) = 0$ 또는
$S \otimes_R \kappa(\mathfrak p) = \kappa(\mathfrak p)$이다. 이는 (1)과
(2)를 증명한다.
\end{proof}

\begin{lemma}
\label{algebra-lemma-relations}
$R$을 환이라 하자. $M$, $N$을 $R$-가군이라 하자.
$\{x_i\}_{i \in I}$를 $M$의 생성원 집합이라 하고,
$\{y_j\}_{j \in J}$를 $N$의 생성원 집합이라 하자.
$\{m_j\}_{j \in J}$를 유한 개의 $j$를 제외하고 모두 $m_j = 0$인
$M$의 원소족이라 하자. 그러면
$$
\sum\nolimits_{j \in J} m_j \otimes y_j = 0 \text{ in } M \otimes_R N
$$
인 것은 다음과 동치이다. 유한 개의 쌍 $(i, j)$를 제외하고 모두
$a_{i, j} = 0$인 $a_{i, j} \in R$가 존재하여
\begin{align*}
m_j & = \sum\nolimits_{i \in I} a_{i, j} x_i \quad\text{for all } j \in J, \\
0 & = \sum\nolimits_{j \in J} a_{i, j} y_j \quad\text{for all } i \in I.
\end{align*}
를 만족한다.
\end{lemma}

\begin{proof}
충분성은 즉시 알 수 있다.
$\sum_{j \in J} m_j \otimes y_j = 0$이라고 하자. 짧은 완전열
$$
0 \to K \to \bigoplus\nolimits_{j \in J} R \to N \to 0
$$
을 생각하자. 여기서 $\bigoplus\nolimits_{j \in J} R$의 $j$번째 기저
벡터는 $y_j$로 사상된다. 여기에 $M$을 텐서하여 완전열
$$
K \otimes_R M \to \bigoplus\nolimits_{j \in J} M \to N \otimes_R M \to 0.
$$
을 얻는다. 가정에 의해 $\sum k_i \otimes x_i$가 가운데 항의 원소
$(m_j)_{j \in J}$로 사상되게 하는 원소 $k_i \in K$들이 존재한다.
% P01-E1405
$k_i = (a_{i, j})_{j \in J}$로 쓰면 원하는 것을 얻는다.
\end{proof}

\begin{lemma}
\label{algebra-lemma-kernel-difference-projections}
$\varphi : R \to S$를 환 사상이라 하고 $g \in S$라 하자.
다음 조건들은 동치이다.
\begin{enumerate}
\item $S \otimes_R S$에서 $g \otimes 1 = 1 \otimes g$이다.
\item 어떤 $n \geq 0$과 $1 \leq i, j \leq n$에 대해 원소
$y_i, z_j \in S$ 및 $x_{i, j} \in R$가 존재하여 다음을 만족한다.
\begin{enumerate}
\item $g = \sum_{i, j \leq n} x_{i, j} y_i z_j$,
\item 각 $j$에 대해 $\sum x_{i, j}y_i \in \varphi(R)$이고,
\item 각 $i$에 대해 $\sum x_{i, j}z_j \in \varphi(R)$이다.
\end{enumerate}
\end{enumerate}
\end{lemma}

\begin{proof}
(2)가 (1)을 함의함은 명백하다. 역으로
$g \otimes 1 = 1 \otimes g$라고 하자. $0, 1 \in I$이고
$s_0 = 1$, $s_1 = g$인 $R$-가군 $S$의 생성원
$\{s_i\}_{i \in I}$를 택하자. 관계
$g \otimes s_0 + (-1) \otimes s_1 = 0$에 보조정리
\ref{algebra-lemma-relations}를 적용한다. 그러면 $a_{i, j} \in R$들이 존재하여
$g = \sum_i a_{i, 0} s_i$, $-1 = \sum_i a_{i, 1} s_i$이고,
$j \not = 0, 1$이면 $0 = \sum_i a_{i, j} s_i$이며, 또한 모든 $i$에 대해
$\sum_j a_{i, j}s_j = 0$이다. 따라서
$$
\sum\nolimits_{i, j \not = 0} a_{i, j} s_i s_j = -g + a_{0, 0}
$$
이고 각 $j \not = 0$에 대해
$\sum_{i \not = 0} a_{i, j}s_i \in R$이다. 이는
$-g + a_{0, 0}$을 (2)와 같이 쓸 수 있음을 증명한다. 따라서 $g$도
(2)와 같이 쓸 수 있다. 세부 사항은 생략한다.
힌트: (2)와 같은 표현을 갖는 $S$의 원소들의 집합이 $S$의
% P01-E1406
$R$-부분대수를 이룸을 보여라.
\end{proof}

\begin{remark}
\label{algebra-remark-matrices-associated-to-elements-epicenter}
$R \to S$를 환 사상이라 하자. $g \otimes 1 = 1 \otimes g$를 만족하는
원소 $g \in S$들의 집합을 때때로 $S$의 {\it 에피중심}이라고 한다.
이는 $R$-대수이다. 보조정리
\ref{algebra-lemma-kernel-difference-projections}의 구성에 의해 에피중심의 각
$g$에 대해 행렬 분해
$$
(g) = Y X Z
$$
를 얻는다. 여기서 $X \in \text{Mat}(n \times n, R)$,
$Y \in \text{Mat}(1 \times n, S)$이고
$Z \in \text{Mat}(n \times 1, S)$이다. 실제로 보조정리의 (2)에서와 같은
$x_{i, j}, y_i, z_j$를 택하자. $X = (x_{i, j})$로 두고, $y_i$들을
성분으로 하는 행벡터를 $y$, $z_j$들을 성분으로 하는 열벡터를 $z$라
하자. 이 표기에서 보조정리
\ref{algebra-lemma-kernel-difference-projections}의 조건 (b)와 (c)는 정확히
$Y X \in \text{Mat}(1 \times n, R)$,
$X Z \in \text{Mat}(n \times 1, R)$임을 뜻한다.
행렬의 삼중항 $(X, YX, XZ)$을 고려하는 것이 매우 편리하다.
$n \in \mathbf{N}$과 삼중항 $(P, U, V)$가 주어졌을 때, 위와 같은 행렬
분해가 존재하여 $P = X$, $U = YX$, $V = XZ$이면
$(P, U, V)$를 {\it $g$에 결부된 $n$-삼중항}이라고 한다.
\end{remark}

\begin{lemma}
\label{algebra-lemma-epimorphism-cardinality}
$R \to S$를 환의 epimorphism이라 하자. 그러면 $S$의 기수는 $R$의 기수
이하이다. 식으로 쓰면 $|S| \leq |R|$이다.
\end{lemma}

\begin{proof}
$R \to S$가 epimorphism이라는 조건은 각 $g \in S$가
$g \otimes 1 = 1 \otimes g$를 만족한다는 뜻이다. 보조정리
\ref{algebra-lemma-epimorphism}를 보라. 비고
\ref{algebra-remark-matrices-associated-to-elements-epicenter}에서 도입한 표기를
사용하자. $g, g' \in S$라 하고 $(P, U, V)$가 $g$와 $g'$ 모두에 결부된
$n$-삼중항이라고 하자. 이때 $g = g'$임을 주장한다. 실제로 $g$의 행렬
분해 $(g) = YXZ$에 대해 $(P, U, V) = (X, YX, XZ)$라 쓰고, $g'$의 행렬
분해 $(g') = Y'X'Z'$에 대해
$(P, U, V) = (X', Y'X', X'Z')$라 쓰자. 그러면
$$
(g) = YXZ = UZ = Y'X'Z = Y'PZ = Y'XZ = Y'V = Y'X'Z' = (g')
$$
이므로 $g = g'$이다. 따라서 $S$의 기수는 가능한 삼중항의 수로
유계이고, 이는 기껏해야 $\sup_{n \in \mathbf{N}} |R|^n$이다.
$R$이 무한이면 이는 기껏해야 $|R|$이다. \cite[Ch. I, 10.13]{Kunen}을
보라.

\medskip\noindent
$R$이 유한환이면 위의 논증은 $S$가 기껏해야 가산임만을 증명한다.
사실 이 경우 $R$은 아르틴환이고 사상 $R \to S$는 전사이다. 이 경우의
증명은 생략한다.
\end{proof}

\begin{lemma}
\label{algebra-lemma-epimorphism-modules}
환 준동형 $R \to S$에 대하여 다음 조건들은 동치이다.
\begin{enumerate}
\item $R \to S$는 환의 epimorphism이다.
\item 임의의 $S$-가군 $N_1, N_2$에 대해
$\Hom_S(N_1, N_2) = \Hom_R(N_1, N_2)$이다.
\item 제한 함자 $\text{Mod}_S \to \text{Mod}_R$는 충만충실하다.
\end{enumerate}
\end{lemma}

\begin{proof}
(2)와 (3)은 정의에 의해 동치이다.

\medskip\noindent
(1)을 가정하고, $N_1, N_2$를 $S$-가군이라 하며
$\varphi : N_1 \to N_2$를 $R$-선형 사상이라 하자. 각 $x \in N_1$에
대하여 규칙 $g \otimes g' \mapsto g\varphi(g'x)$로 정의된 사상
$S \otimes_R S \to N_2$를 생각하자. 두 사상 $S \to S \otimes_R S$가
모두 동형사상이므로(보조정리 \ref{algebra-lemma-epimorphism}),
$g \varphi(g'x) = gg'\varphi(x) = \varphi(gg' x)$이다. 따라서 $\varphi$는
$S$-선형이다.

\medskip\noindent
(2)를 가정하자. $N_1 = S \otimes_R S$를 왼쪽 $S$-가군 작용에 의한
$S$-가군으로 보고, $N_2 = S \otimes_R S$에는 오른쪽 $S$-가군 작용을
주자. $N_1 = N_2$가 $R$-가군으로서 성립하므로,
(2)에 의해 $S \otimes_R S$ 위의 두 $S$-가군 구조는 일치한다. 이는
보조정리 \ref{algebra-lemma-epimorphism}에 의해 (1)을 함의한다.
\end{proof}



\section{순수 아이디얼}
\label{algebra-section-pure-ideals}

\noindent
이 절의 내용은 많은 논문에서 논의된다. 예를 들어 \cite{Lazard},
\cite{Bkouche}, \cite{DeMarco}를 보라.

\begin{definition}
\label{algebra-definition-pure-ideal}
$R$을 환이라 하자. 몫환 $R/I$가 $R$ 위에서 평탄이면
$I \subset R$를 {\it 순수}하다고 한다.
\end{definition}

\begin{lemma}
\label{algebra-lemma-pure}
$R$을 환, $I \subset R$를 아이디얼이라 하자. 다음 조건들은 동치이다.
\begin{enumerate}
\item $I$는 순수하다.
\item 모든 아이디얼 $J \subset R$에 대해 $J \cap I = IJ$이다.
\item 모든 유한 생성 아이디얼 $J \subset R$에 대해
$J \cap I = JI$이다.
\item 모든 $x \in R$에 대해 $(x) \cap I = xI$이다.
\item 모든 $x \in I$에 대해 어떤 $y \in I$가 존재하여 $x = yx$이다.
\item 모든 $x_1, \ldots, x_n \in I$에 대해 어떤 $y \in I$가 존재하여
모든 $i = 1, \ldots, n$에 대해 $x_i = yx_i$이다.
\item $R$의 모든 소 아이디얼 $\mathfrak p$에 대해
$IR_{\mathfrak p} = 0$ 또는 $IR_{\mathfrak p} = R_{\mathfrak p}$이다.
\item $\text{Supp}(I) = \Spec(R) \setminus V(I)$이다.
\item $I$는 사상 $R \to (1 + I)^{-1}R$의 핵이다.
\item $R$의 어떤 곱셈 부분집합 $S$에 대해 $R$-대수로서
$R/I \cong S^{-1}R$이다.
\item $R$-대수로서 $R/I \cong (1 + I)^{-1}R$이다.
\end{enumerate}
\end{lemma}

\begin{proof}
$R$의 임의의 아이디얼 $J$에 대해 짧은 완전열
$0 \to J \to R \to R/J \to 0$이 있다. 여기에 $R/I$를 텐서하면
% P01-E0523
완전열 $J \otimes_R R/I \to R/I \to R/I + J \to 0$을 얻고,
$J \otimes_R R/I = J/JI$이다. 따라서 보조정리 \ref{algebra-lemma-flat}에서
(1), (2), (3)의 동치가 따른다. 또한 이 조건들은 (4)를 함의한다.

\medskip\noindent
(4) $\Rightarrow$ (5)는 자명하다. (5)를 가정하고
$x_1, \ldots, x_n \in I$라 하자. $x_i = y_ix_i$를 만족하는
$y_i \in I$를 택하자. 등식
$1 - y = \prod_{i = 1, \ldots, n} (1 - y_i)$로 정해지는 원소
$y \in I$를 택하자. 그러면 모든 $i = 1, \ldots, n$에 대해
$x_i = yx_i$이다. 따라서 (6)이 성립하고, (5) $\Leftrightarrow$ (6)이다.

\medskip\noindent
(5)를 가정하자. $x \in I$라 하자. 어떤 $y \in I$에 대해
$x = yx$이다. 따라서 $x(1 - y) = 0$이고, 이는 $x$가
$(1 + I)^{-1}R$에서 영으로 사상됨을 보인다. 물론 사상
$R \to (1 + I)^{-1}R$의 핵은 항상 $I$에 포함된다. 따라서 (5)는
(9)를 함의한다. (9)를 가정하자. 임의의 $x \in I$에 대해 어떤
$y \in I$가 존재하여 $x(1 - y) = 0$이다. 다시 말해 $x = yx$이다.
따라서 (5)와 (9)는 동치이다.

\medskip\noindent
(5)를 가정하고 $\mathfrak p$를 $R$의 소 아이디얼이라 하자.
$\mathfrak p \not \in V(I)$이면 $IR_{\mathfrak p} = R_{\mathfrak p}$이다.
$\mathfrak p \in V(I)$, 다시 말해 $I \subset \mathfrak p$이면,
% P01-E0524
$x \in I$에 대해 어떤 $y \in I$가 존재하여 $x(1 - y) = 0$이고,
따라서 $x$는 $R_{\mathfrak p}$에서 영으로 사상된다. 즉
$IR_{\mathfrak p} = 0$이다. 그러므로 (7)이 성립한다.

\medskip\noindent
(7)을 가정하자. 그러면 $R$의 임의의 소 아이디얼 $\mathfrak p$에 대해
$(R/I)_{\mathfrak p}$는 $0$ 또는 $R_{\mathfrak p}$이다. 따라서 보조정리
\ref{algebra-lemma-flat-localization}에 의해 (1)이 성립한다. 이 시점에서
(1) -- (7)과 (9)가 모두 동치임을 알 수 있다.

\medskip\noindent
$IR_{\mathfrak p} = I_{\mathfrak p}$이므로 (7)은 (8)을 함의한다.
마지막으로 (8)이 성립하면, 이는 정확히 $I_{\mathfrak p}$가 영 가군일
필요충분조건이 $\mathfrak p \in V(I)$라는 뜻이며, 이는 명백히 (7)을
뜻한다. 이제 (1) -- (9)는 동치이다.

\medskip\noindent
(1) -- (9)가 성립한다고 하자. (9)에 의해
$R/I \subset (1 + I)^{-1}R$이고, (7)이 주는 소 아이디얼에서의 국소화
기술에 의해 사상 $R/I \to (1 + I)^{-1}R$도 전사이다. 따라서 (11)이
성립한다.

\medskip\noindent
(11) $\Rightarrow$ (10)은 자명하다. 또한 $R$의 국소화는 $R$ 위에서
평탄하므로 (10)은 (1)을 함의한다. 보조정리
\ref{algebra-lemma-flat-localization}을 보라.
\end{proof}

\begin{lemma}
\label{algebra-lemma-pure-ideal-determined-by-zero-set}
\begin{slogan}
순수 아이디얼은 그 영점 자취로 결정된다.
\end{slogan}
$R$을 환이라 하자. $I, J \subset R$가 순수 아이디얼이면
$V(I) = V(J)$는 $I = J$를 함의한다.
\end{lemma}

\begin{proof}
예를 들어 보조정리 \ref{algebra-lemma-pure}의 성질 (7)에 의해
$I = \Ker(R \to \prod_{\mathfrak p \in V(I)} R_{\mathfrak p})$이므로,
결부된 닫힌 부분집합으로부터 이를 복원할 수 있다.
\end{proof}

\begin{lemma}
\label{algebra-lemma-pure-open-closed-specializations}
$R$을 환이라 하자. 규칙 $I \mapsto V(I)$는 전단사
$$
\{I \subset R \text{ pure}\}
\leftrightarrow
\{Z \subset \Spec(R)\text{ closed and closed under generalizations}\}
$$
를 정한다.
\end{lemma}

\begin{proof}
$I$를 순수 아이디얼이라 하자. $R \to R/I$가 평탄하므로, 하강 성질에
의해 사상 $\Spec(R/I) \to \Spec(R)$을 따라 일반화를 올릴 수 있다.
따라서 $V(I)$는 일반화에 대해 닫혀 있다. 이는 사상이 잘 정의됨을
보인다. 보조정리 \ref{algebra-lemma-pure-ideal-determined-by-zero-set}에 의해
사상은 단사이다. $Z \subset \Spec(R)$가 닫혀 있고 일반화에 대해서도
닫혀 있다고 하자. $Z = V(J)$를 만족하는 근기 아이디얼 $J \subset R$를
택하자. $I = \{x \in R : x \in xJ\}$로 두자. $I$는 아이디얼이다.
실제로 $x, y \in I$이면 $x = xf$와 $y = yg$를 만족하는
$f, g \in J$가 존재하고,
$$
x + y = (x + y)(f + g - fg)
$$
이다. 확인은 독자에게 맡긴다. $I$가 순수이고 $V(I) = V(J)$라고
주장한다. 이 주장이 참이면 보조정리의 사상은 전사이고 보조정리가
성립한다.

\medskip\noindent
$I \subset J$이므로 $V(J) \subset V(I)$이다. $I \subset \mathfrak p$인
소 아이디얼을 택하자. 곱셈 부분집합
$S = (R \setminus \mathfrak p)(1 + J)$를 생각하자. $I$의 정의와
$I \subset \mathfrak p$에 의해 $0 \not \in S$이다. 따라서 $S$와
서로소인 $R$의 소 아이디얼 $\mathfrak q$가 존재한다. 보조정리
\ref{algebra-lemma-localization-zero}와 \ref{algebra-lemma-spec-localization}을 보라.
따라서 $\mathfrak q \subset \mathfrak p$이고
$\mathfrak q \cap (1 + J) = \emptyset$이다. 이는
$\mathfrak q + J$가 $R$의 진 아이디얼임을 함의한다.
$\mathfrak q + J$를 포함하는 극대 아이디얼 $\mathfrak m$을 택하자.
그러면 $\mathfrak m \in V(J)$이고, $Z$가 일반화에 대해 닫혀 있다고
가정했으므로 $\mathfrak q \in V(J) = Z$이다. 다시
$\mathfrak q \subset \mathfrak p$이므로 $\mathfrak p \in V(J)$이다.
따라서 $V(I) = V(J)$이다.

\medskip\noindent
마지막으로 $V(I) = V(J)$이고 $J$가 근기적이므로 $J = \sqrt{I}$이다.
$x \in I$를 택하자. 정의에 의해 어떤 $y \in J$에 대해 $x = xy$이다.
그러면 $x = xy = xy^2 = \ldots = xy^n$이다. 어떤 $n > 0$에 대해
$y^n \in I$이므로 보조정리 \ref{algebra-lemma-pure}의 성질 (5)가 성립하고,
$I$가 실제로 순수임을 알 수 있다.
\end{proof}

\begin{lemma}
\label{algebra-lemma-finitely-generated-pure-ideal}
$R$을 환, $I \subset R$를 아이디얼이라 하자. 다음 조건들은 동치이다.
\begin{enumerate}
\item $I$는 순수하고 유한 생성이다.
\item $I$는 멱등원 하나로 생성된다.
\item $I$는 순수하고 $V(I)$는 열린집합이다.
\item $R/I$는 사영 $R$-가군이다.
\end{enumerate}
\end{lemma}

\begin{proof}
(1)이 성립하면 보조정리 \ref{algebra-lemma-pure}에 의해
$I = I \cap I = I^2$이다. 따라서 보조정리
\ref{algebra-lemma-ideal-is-squared-union-connected}에 의해 $I$는 멱등원으로
생성된다. 그러므로 (1) $\Rightarrow$ (2)이다. (2)가 성립하면
$I = (e)$이고 $R$-가군으로서 $R = (1 - e) \oplus (e)$이다. 따라서
$R/I$는 평탄하고 $I$는 순수하며 $V(I) = D(1 - e)$는 열린집합이다.
그러므로 (2) $\Rightarrow$ (1) $+$ (3)이다. 마지막으로 (3)을 가정하자.
$V(I)$는 열린집합이자 닫힌집합이므로, $R$의 어떤 멱등원 $e$에 대해
$V(I) = D(1 - e)$이다. 보조정리 \ref{algebra-lemma-disjoint-decomposition}을
보라. 아이디얼 $J = (e)$는 $V(J) = V(I)$인 순수 아이디얼이므로,
보조정리 \ref{algebra-lemma-pure-ideal-determined-by-zero-set}에 의해 $I = J$이다.
따라서 (3) $\Rightarrow$ (2)임을 알 수 있다. 보조정리
\ref{algebra-lemma-finite-projective}에 의해 (4)는 $I$가 순수하고 $R/I$가 유한
표현이라는 명제와 동치이다. 또한 보조정리 \ref{algebra-lemma-extension}에서
보듯이 $R/I$가 유한 표현일 필요충분조건은 $I$가 유한 생성인 것이다.
따라서 (4)는 (1)과 동치이다.
\end{proof}

\noindent
위의 결과를 사용하여 모든 유한 평탄 가군이 유한 표현인 환들을
특징지을 수 있다.

\begin{lemma}
\label{algebra-lemma-finite-flat-module-finitely-presented}
$R$을 환이라 하자. 다음 조건들은 동치이다.
\begin{enumerate}
\item 닫혀 있고 일반화에 대해서도 닫혀 있는 모든
$Z \subset \Spec(R)$는 열린집합이기도 하다.
\item 모든 유한 평탄 $R$-가군은 유한 국소 자유이다.
\end{enumerate}
\end{lemma}

\begin{proof}
모든 유한 평탄 $R$-가군이 유한 국소 자유라면, $I$가 순수 아이디얼일
때 $R/I$의 지지집합은 열린집합이다. 따라서
% P01-E0525
보조정리 \ref{algebra-lemma-pure-ideal-determined-by-zero-set}에서
(2) $\Rightarrow$ (1)이 따른다.

\medskip\noindent
역으로 $R$이 (1)을 만족한다고 하자. $M$을 유한 평탄 $R$-가군이라
하자. $M$의 지지집합 $Z = \text{Supp}(M)$는 닫혀 있다. 보조정리
\ref{algebra-lemma-support-closed}를 보라. 한편 $\mathfrak p \subset \mathfrak p'$이면
보조정리 \ref{algebra-lemma-finite-flat-local}에 의해 가군
$M_{\mathfrak p'}$는 자유이고,
% P01-E0526
$M_{\mathfrak p} = M_{\mathfrak p'} \otimes_{R_{\mathfrak p'}} R_{\mathfrak p}$이다.
따라서
$\mathfrak p' \in \text{Supp}(M) \Rightarrow \mathfrak p \in \text{Supp}(M)$,
즉 지지집합은 일반화에 대해 닫혀 있다. $R$이 (1)을 만족하므로 $M$의
지지집합은 열린집합이자 닫힌집합이다. $M$이 $r$개의 원소
$m_1, \ldots, m_r$로 생성된다고 하자. 가군
$\wedge^i(M)$, $i = 1, \ldots, r$도 유한 평탄 $R$-가군이다. 실제로
$\wedge^i(M)_{\mathfrak p} = \wedge^i(M_{\mathfrak p})$는
$R_{\mathfrak p}$ 위에서 자유이다. 또한
$\text{Supp}(\wedge^{i + 1}(M)) \subset \text{Supp}(\wedge^i(M))$이다.
따라서 열린닫힌 부분집합들에 의한 분해
$$
\Spec(R) = U_0 \amalg U_1 \amalg \ldots \amalg U_r
$$
가 존재하여, 모든 $i = 0, \ldots, r$에 대해 $\wedge^i(M)$의 지지집합은
$U_r \cup \ldots \cup U_i$이다. $\mathfrak p$를 $R$의 소 아이디얼이라
하고 $\mathfrak p \in U_i$라 하자. 등식
$\wedge^i(M) \otimes_R \kappa(\mathfrak p) =
\wedge^i(M \otimes_R \kappa(\mathfrak p))$에 주목하자. 필요하면
$m_1, \ldots, m_r$의 번호를 다시 매겨 $m_1, \ldots, m_i$가
$M \otimes_R \kappa(\mathfrak p)$를 생성한다고 해도 된다. Nakayama의
보조정리 \ref{algebra-lemma-NAK}에 의해 어떤 $f \in R$,
$f \not \in \mathfrak p$에 대해 전사
$$
R_f^{\oplus i} \longrightarrow M_f, \quad
(a_1, \ldots, a_i) \longmapsto \sum a_im_i
$$
를 얻는다. 또한 $D(f) \subset U_i$라고 가정해도 된다. 이는
$\wedge^i(M_f) = \wedge^i(M)_f$가 그 지지집합이
$\Spec(R_f)$ 전체인 평탄 $R_f$ 가군임을 뜻한다. 위에서 보았듯이 이는
단일 원소 $m_1 \wedge \ldots \wedge m_i$로 생성된다. 따라서 어떤 순수
아이디얼 $J \subset R_f$에 대해 $\wedge^i(M)_f \cong R_f/J$이고
$V(J) = \Spec(R_f)$이다. 이는 명백히 $J = (0)$을 뜻한다. 보조정리
\ref{algebra-lemma-pure-ideal-determined-by-zero-set}을 보라. 그러므로
$m_1 \wedge \ldots \wedge m_i$는 $\wedge^i(M_f)$의 기저이고, 따라서
위에 표시한 사상은 전사일 뿐 아니라 단사이기도 하다. 이는 원하는 대로
$M$이 유한 국소 자유임을 증명한다.
\end{proof}







\section{유한 전역 차원인 환}
\label{algebra-section-ring-finite-gl-dim}

\noindent
다음 보조정리는 주어진 가군의 서로 다른 사영 분해를 비교할 때 자주
사용된다.

\begin{lemma}[Schanuel의 보조정리]
\label{algebra-lemma-Schanuel}
$R$을 환, $M$을 $R$-가군이라 하자.
$$
0 \to K \xrightarrow{c_1} P_1 \xrightarrow{p_1} M \to 0
\quad\text{and}\quad
0 \to L \xrightarrow{c_2} P_2 \xrightarrow{p_2} M \to 0
$$
이 $P_i$가 사영인 두 짧은 완전열이라고 하자. 그러면
$K \oplus P_2 \cong L \oplus P_1$이다. 더 정확히 말하면, 수직 화살표가
% P01-E1407
동형사상인 가환 그림
$$
\xymatrix{
0 \ar[r] &
K \oplus P_2 \ar[r]_{(c_1, \text{id})} \ar[d] &
P_1 \oplus P_2 \ar[r]_{(p_1, 0)} \ar[d] &
M \ar[r] \ar@{=}[d] &
0 \\
0 \ar[r] &
P_1 \oplus L \ar[r]^{(\text{id}, c_2)} &
P_1 \oplus P_2 \ar[r]^{(0, p_2)} &
M \ar[r] &
0
}
$$
이 존재한다.
\end{lemma}

\begin{proof}
짧은 완전열 $0 \to N \to P_1 \oplus P_2 \to M \to 0$로 정의되는
가군 $N$을 생각하자. 여기서 마지막 사상은 두 사상 $P_i \to M$의
합이다. 사영 $N \to P_1$은 핵이 $L$인 전사이고, $N \to P_2$는 핵이
$K$인 전사임을 쉽게 알 수 있다. $P_i$들이 사영이므로
$N \cong K \oplus P_2 \cong L \oplus P_1$이다. 이는 첫째 명제를
증명한다.

\medskip\noindent
둘째 명제를 증명하고 첫째 명제도 다시 증명하기 위해,
$p_1 = p_2 \circ a$와 $p_2 = p_1 \circ b$를 만족하도록
$a : P_1 \to P_2$와 $b : P_2 \to P_1$을 택하자. $P_1$과 $P_2$가
사영이므로 이는 가능하다. 그러면 가환 그림
$$
\xymatrix{
0 \ar[r] &
K \oplus P_2 \ar[r]_{(c_1, \text{id})} &
P_1 \oplus P_2 \ar[r]_{(p_1, 0)} &
M \ar[r] &
0 \\
0 \ar[r] &
N \ar[r] \ar[d] \ar[u] &
P_1 \oplus P_2 \ar[r]_{(p_1, p_2)}
\ar[d]_S \ar[u]^T &
M \ar[r] \ar@{=}[d] \ar@{=}[u] &
0 \\
0 \ar[r] &
P_1 \oplus L \ar[r]^{(\text{id}, c_2)} &
P_1 \oplus P_2 \ar[r]^{(0, p_2)} &
M \ar[r] &
0
}
$$
을 얻는다. 여기서 $T$와 $S$는 행렬
$$
S = \left(
\begin{matrix}
\text{id} & 0 \\
a & \text{id}
\end{matrix}
\right)
\quad\text{and}\quad
T = \left(
\begin{matrix}
\text{id} & b \\
0 & \text{id}
\end{matrix}
\right)
$$
로 주어진다. 그러면 원하는 대로 $S$, $T$ 및 사상
$N \to P_1 \oplus L$과 $N \to K \oplus P_2$는 동형사상이다.
\end{proof}

\begin{definition}
\label{algebra-definition-finite-proj-dim}
$R$을 환, $M$을 $R$-가군이라 하자. 사영 $R$-가군들에 의한 유한 길이
분해가 존재하면 $M$이 {\it 유한 사영 차원}을 갖는다고 한다. 그러한
분해의 최소 길이를 $M$의 {\it 사영 차원}이라고 한다.
\end{definition}

\noindent
$M$의 사영 차원이 $0$일 필요충분조건은 $M$이 사영 가군인 것임이
명백하다. 다음 보조정리는 사영 차원이 사영 분해의 선택에 어느 정도
독립적인지를 설명한다.

\begin{lemma}
\label{algebra-lemma-independent-resolution}
$R$을 환이라 하자. $M$이 사영 차원 $d$인 $R$-가군이라고 하자.
$F_i$들이 사영이고 $e \geq d - 1$인 완전열
$F_e \to F_{e-1} \to \ldots \to F_0 \to M \to 0$이 있다고 하자.
그러면 $F_e \to F_{e-1}$의 핵은 사영이다($e = 0$인 경우에는
$F_0 \to M$의 핵이 사영이다).
\end{lemma}

\begin{proof}
$d$에 대한 귀납법으로 증명한다. $d = 0$이면 $M$은 사영이다. 이 경우
분해 $F_0 = \Ker(F_0 \to M) \oplus M$가 있으므로
$\Ker(F_0 \to M)$은 사영이다. 이는 $e = 0$일 때 증명을 끝낸다.
$e > 0$이면 $M$을 $\Ker(F_0 \to M)$으로 바꾸어 $e$를 줄인다.

\medskip\noindent
이제 $d > 0$이라고 하자. $P_i$들이 사영인 최소 길이의 유한 분해
$0 \to P_d \to P_{d-1} \to \ldots \to P_0 \to M \to 0$을 택하자.
Schanuel의 보조정리 \ref{algebra-lemma-Schanuel}에 의해
$P_0 \oplus \Ker(F_0 \to M) \cong F_0 \oplus \Ker(P_0 \to M)$이다.
$d = 1$, $e = 0$이면 우변은 사영인 $F_0 \oplus P_1$이므로 이는 그
경우를 증명한다. 따라서 이제 $e > 0$이라고 해도 된다. 가군
$F_0 \oplus \Ker(P_0 \to M)$은 길이 $d - 1$의 유한 사영 분해
$$
0 \to P_d \to P_{d-1} \to \ldots \to
P_2 \to P_1 \oplus F_0 \to \Ker(P_0 \to M) \oplus F_0 \to 0
$$
를 갖는다. 길이 $e - 1$인 완전열
$$
F_e \to F_{e-1} \to \ldots \to F_2 \to P_0 \oplus F_1 \to
P_0 \oplus \Ker(F_0 \to M) \to 0
$$
에 귀납 가정을 적용하면, $e \geq 2$일 때
$\Ker(F_e \to F_{e - 1})$이 사영이거나
$\Ker(F_1 \oplus P_0 \to F_0 \oplus P_0)$이 사영이다. 이는 보조정리를
함의한다.
\end{proof}

\begin{lemma}
\label{algebra-lemma-what-kind-of-resolutions}
$R$을 환, $M$을 $R$-가군이라 하고 $d \geq 0$이라 하자.
다음 조건들은 동치이다.
\begin{enumerate}
\item $M$의 사영 차원은 $\leq d$이다.
\item $P_i$들이 사영인 분해
$0 \to P_d \to P_{d - 1} \to \ldots \to P_0 \to M \to 0$이 존재한다.
\item $P_i$들이 사영인 어떤 분해
$\ldots \to P_2 \to P_1 \to P_0 \to M \to 0$에 대하여,
$d \geq 2$이면 $\Ker(P_{d - 1} \to P_{d - 2})$가 사영이고,
$d = 1$이면 $\Ker(P_0 \to M)$이 사영이며, $d = 0$이면 $M$이 사영이다.
\item $P_i$들이 사영인 임의의 분해
$\ldots \to P_2 \to P_1 \to P_0 \to M \to 0$에 대하여,
$d \geq 2$이면 $\Ker(P_{d - 1} \to P_{d - 2})$가 사영이고,
$d = 1$이면 $\Ker(P_0 \to M)$이 사영이며, $d = 0$이면 $M$이 사영이다.
\end{enumerate}
\end{lemma}

\begin{proof}
(1)과 (2)의 동치는 사영 차원의 정의이다. 정의
\ref{algebra-definition-finite-proj-dim}을 보라. 보조정리
\ref{algebra-lemma-independent-resolution}에 의해 (2) $\Rightarrow$ (4)이다.
(4) $\Rightarrow$ (3)과 (3) $\Rightarrow$ (2)는 즉시 알 수 있다.
\end{proof}

\begin{lemma}
\label{algebra-lemma-what-kind-of-resolutions-local}
$R$을 국소환, $M$을 $R$-가군이라 하고 $d \geq 0$이라 하자.
보조정리 \ref{algebra-lemma-what-kind-of-resolutions}의 동치인 조건 (1) -- (4)는
다음 조건과도 동치이다.
\begin{enumerate}
\item[(5)] $P_i$들이 자유인 분해
$0 \to P_d \to P_{d - 1} \to \ldots \to P_0 \to M \to 0$이 존재한다.
\end{enumerate}
\end{lemma}

\begin{proof}
보조정리 \ref{algebra-lemma-what-kind-of-resolutions}와 정리
\ref{algebra-theorem-projective-free-over-local-ring}에서 따른다.
\end{proof}

\begin{lemma}
\label{algebra-lemma-what-kind-of-resolutions-Noetherian}
$R$을 뇌터환, $M$을 유한 $R$-가군이라 하고 $d \geq 0$이라 하자.
보조정리 \ref{algebra-lemma-what-kind-of-resolutions}의 동치인 조건 (1) -- (4)는
다음 조건과도 동치이다.
\begin{enumerate}
\item[(6)] $P_i$들이 유한 사영인 분해
$0 \to P_d \to P_{d - 1} \to \ldots \to P_0 \to M \to 0$이 존재한다.
\end{enumerate}
\end{lemma}

\begin{proof}
$F_i$들이 유한 자유인 분해
$\ldots \to F_2 \to F_1 \to F_0 \to M \to 0$을 택하자(보조정리
\ref{algebra-lemma-resolution-by-finite-free}). 보조정리
\ref{algebra-lemma-what-kind-of-resolutions}에 의해 적어도 $d \geq 2$이면
% P01-E0527
$P_d = \Ker(F_{d - 1} \to F_{d - 2})$는 사영이다. $R$이 뇌터이고
$P_d \subset F_{d - 1}$이며 후자가 유한 자유이므로 $P_d$는 유한
$R$-가군이다. 그러므로
$0 \to P_d \to F_{d - 1} \to \ldots \to F_1 \to F_0 \to M \to 0$은
원하는 분해이다.
\end{proof}

\begin{lemma}
\label{algebra-lemma-what-kind-of-resolutions-Noetherian-local}
$R$을 뇌터 국소환, $M$을 유한 $R$-가군이라 하고 $d \geq 0$이라 하자.
보조정리 \ref{algebra-lemma-what-kind-of-resolutions}의 동치인 조건 (1) -- (4),
보조정리 \ref{algebra-lemma-what-kind-of-resolutions-local}의 조건 (5), 그리고
보조정리 \ref{algebra-lemma-what-kind-of-resolutions-Noetherian}의 조건 (6)은
다음 조건과도 동치이다.
\begin{enumerate}
\item[(7)] $F_i$들이 유한 자유인 분해
$0 \to F_d \to F_{d - 1} \to \ldots \to F_0 \to M \to 0$이 존재한다.
\end{enumerate}
\end{lemma}

\begin{proof}
이는 보조정리 \ref{algebra-lemma-what-kind-of-resolutions},
\ref{algebra-lemma-what-kind-of-resolutions-local},
\ref{algebra-lemma-what-kind-of-resolutions-Noetherian} 및 국소환 위의 유한 사영
가군이 유한 자유라는 사실에서 따른다. 보조정리
\ref{algebra-lemma-finite-projective}를 보라.
\end{proof}

\begin{lemma}
\label{algebra-lemma-projective-dimension-ext}
$R$을 환, $M$을 $R$-가군이라 하고 $n \geq 0$이라 하자.
다음 조건들은 동치이다.
\begin{enumerate}
\item $M$의 사영 차원은 $\leq n$이다.
\item 모든 $R$-가군 $N$과 모든 $i \geq n + 1$에 대해
$\Ext^i_R(M, N) = 0$이다.
\item 모든 $R$-가군 $N$에 대해 $\Ext^{n + 1}_R(M, N) = 0$이다.
\end{enumerate}
\end{lemma}

\begin{proof}
(1)을 가정하자. $M$의 자유 분해 $F_\bullet \to M$을 택하고
$d_e : F_e \to F_{e - 1}$로 쓰자. 보조정리
\ref{algebra-lemma-independent-resolution}에 의해 $e \geq n - 1$이면
$P_e = \Ker(d_e)$가 사영임을 알 수 있다. 따라서 $e \geq n$이면
% P01-E0528
$F_e \cong P_e \oplus P_{e - 1}$이고, 여기서 $d_e$는 직합 인자
$P_{e - 1}$을 $F_{e - 1}$ 안의 $P_{e - 1}$에 동형으로 사상한다.
그러므로 임의의 $R$-가군 $N$에 대해 복합체
$\Hom_R(F_\bullet, N)$은 차수 $\geq n + 1$에서 분할 완전이다.
따라서 (2)가 성립한다. (2) $\Rightarrow$ (3)은 자명하다.

\medskip\noindent
(3)이 성립한다고 가정하자. $n = 0$이면 보조정리
\ref{algebra-lemma-characterize-projective}에 의해 $M$은 사영이므로 (1)이 성립한다.
$n > 0$이면 자유 $R$-가군 $F$와 핵이 $K$인 전사
$F \to M$을 택하자. 보조정리
\ref{algebra-lemma-reverse-long-exact-seq-ext}와 (1)에 의해 모든 $i > 0$에 대해
$\Ext_R^i(F, N)$이 소멸한다는 사실에서, 모든 $R$-가군 $N$에 대해
$\Ext_R^n(K, N) = 0$임을 알 수 있다. 따라서 귀납법에 의해 $K$의 사영
차원은 $\leq n - 1$이다. $K$의 임의의 유한 사영 분해 앞에 $F$를
붙이면 길이가 하나 더 긴 $M$의 사영 분해를 얻으므로, $M$의 사영
차원은 $\leq n$이다.
\end{proof}

\begin{lemma}
\label{algebra-lemma-exact-sequence-projective-dimension}
$R$을 환이라 하고
$0 \to M' \to M \to M'' \to 0$을 $R$-가군들의 짧은 완전열이라 하자.
\begin{enumerate}
\item $M$의 사영 차원이 $\leq n$이고 $M''$의 사영 차원이
$\leq n + 1$이면, $M'$의 사영 차원은 $\leq n$이다.
\item $M'$과 $M''$의 사영 차원이 모두 $\leq n$이면, $M$의 사영
차원은 $\leq n$이다.
\item $M'$의 사영 차원이 $\leq n$이고 $M$의 사영 차원이
$\leq n + 1$이면, $M''$의 사영 차원은 $\leq n + 1$이다.
\end{enumerate}
\end{lemma}

\begin{proof}
보조정리 \ref{algebra-lemma-projective-dimension-ext}의 사영 차원에 대한 특성화와
보조정리 \ref{algebra-lemma-reverse-long-exact-seq-ext}의 Ext 군의 긴 완전열을
결합하면 된다.
\end{proof}

\begin{definition}
\label{algebra-definition-finite-gl-dim}
$R$을 환이라 하자. 모든 $R$-가군이 길이 $n$ 이하인 사영 $R$-가군들에
의한 분해를 갖도록 하는 정수 $n$이 존재하면, 환 $R$이
{\it 유한 전역 차원}을 갖는다고 한다. 그러한 $n$ 중 최소인 것을
$R$의 {\it 전역 차원}이라고 한다.
\end{definition}

\noindent
다음 보조정리의 증명에 나오는 논증은 Auslander의 논문
\cite{Auslander}에서 찾을 수 있다.

\begin{lemma}
\label{algebra-lemma-colimit-projective-dimension}
$R$을 환이라 하자. $M_e$들이 포함관계에 의해 잘 정렬된 부분가군들이고
$M = \bigcup_{e \in E} M_e$라고 가정하자. 몫가군
$M_e/\bigcup\nolimits_{e' < e} M_{e'}$들의 사영 차원이 $\leq n$이라고
가정하자. 그러면 $M$의 사영 차원은 $\leq n$이다.
\end{lemma}

\begin{proof}
$n$에 대한 귀납법으로 증명한다.

\medskip\noindent
기본 경우는 $n = 0$이다. 이때
$P_e = M_e/\bigcup_{e' < e} M_{e'}$는 사영이다. 따라서 사영
$M_e \to P_e$의 절단 $P_e \to M_e$를 택할 수 있다. 유도된 사상
$\psi : \bigoplus_{e \in E} P_e \to M$이 동형이라고 주장한다.
실제로 $x = \sum x_e \in \bigoplus P_e$가 영이 아니면,
$x_{e_{max}}$가 영이 아닌 가장 큰 $e_{max}$를 택한다. 그러면
$y = \psi(x) = \psi(\sum x_e)$는 영이 아니다. 왜냐하면
$y \in M_{e_{max}}$의 $P_{e_{max}}$에서의 상 $x_{e_{max}}$가 영이
아니기 때문이다. 한편 $y \in M$이라 하자. 그러면 어떤 $e$에 대해
$y \in M_e$이다. $e$에 대한 초한 귀납법으로 $y \in \Im(\psi)$임을
보이자. $x_e \in P_e$를 $y$의 상이라 하자. 그러면
$y - \psi(x_e) \in \bigcup_{e' < e} M_{e'}$이다. 귀납 가정에 의해
$y - \psi(x_e) \in \Im(\psi)$이므로 $y \in \Im(\psi)$이다.
따라서 주장은 참이고 $\psi$는 동형이다. 보조정리
\ref{algebra-lemma-direct-sum-projective}에 의해 $M$은 사영가군들의 직합으로서
사영이다.

\medskip\noindent
$n > 0$이면, 각 $e \in E$에 대해 $F_e$를 $M_e$의 원소들의 집합을
기저로 하는 자유 $R$-가군이라 하자. 그러면 잘 정렬된 집합 $E$ 위의
짧은 완전열들의 계
$$
0 \to K_e \to F_e \to M_e \to 0
$$
를 얻는다. 전이사상 $F_{e'} \to F_e$와 $K_{e'} \to K_e$도
단사임에 주의하자. $F = \bigcup F_e$와 $K = \bigcup K_e$로 두자.
그러면
$$
0 \to
K_e/\bigcup\nolimits_{e' < e} K_{e'} \to
F_e/\bigcup\nolimits_{e' < e} F_{e'} \to
M_e/\bigcup\nolimits_{e' < e} M_{e'} \to 0
$$
도 $R$-가군들의 짧은 완전열이고,
$F_e/\bigcup_{e' < e} F_{e'}$는 $M_e$에서
$\bigcup_{e' < e} M_{e'}$에 포함되지 않는 원소들의 집합을 기저로 하는
자유 $R$-가군이다. 따라서 보조정리
\ref{algebra-lemma-exact-sequence-projective-dimension}에 의해
$K_e/\bigcup_{e' < e} K_{e'}$의 사영 차원은 $n - 1$ 이하이다.
귀납법에 의해 $K$의 사영 차원은 $n - 1$ 이하이다. 따라서 $M$의
사영 차원은 $n$ 이하이고, 증명이 끝난다.
\end{proof}

\begin{lemma}
\label{algebra-lemma-finite-gl-dim}
$R$을 환이라 하자. 다음 조건들은 동치이다.
\begin{enumerate}
\item $R$의 유한 전역 차원은 $\leq n$이다.
\item 모든 유한 $R$-가군의 사영 차원은 $\leq n$이다.
\item 모든 순환 $R$-가군 $R/I$의 사영 차원은 $\leq n$이다.
\end{enumerate}
\end{lemma}

\begin{proof}
(1) $\Rightarrow$ (2)와 (2) $\Rightarrow$ (3)은 명백하다. (3)을
가정하자. $M$의 생성원 집합 $E \subset M$을 택하고 $E$ 위에 하나의
정렬순서를 택하자. $e \in E$에 대해 $M_e$를 $e' \leq e$를 만족하는
원소 $e' \in E$들이 생성하는 $M$의 부분가군이라 하자. 그러면
$M = \bigcup_{e \in E} M_e$이다. 각 $e \in E$에 대해 몫가군
$$
M_e/\bigcup\nolimits_{e' < e} M_{e'}
$$
은 영이거나 하나의 원소로 생성되므로, (3)에 의해 사영 차원이
$\leq n$이다. 보조정리 \ref{algebra-lemma-colimit-projective-dimension}에 의해
이는 $M$의 사영 차원이 $\leq n$임을 뜻한다.
\end{proof}

\begin{lemma}
\label{algebra-lemma-localize-finite-gl-dim}
$R$을 환, $M$을 $R$-가군이라 하고 $S \subset R$을 곱셈적 부분집합이라
하자.
\begin{enumerate}
\item $M$의 사영 차원이 $\leq n$이면 $S^{-1}M$의 $S^{-1}R$ 위의
사영 차원은 $\leq n$이다.
\item $R$의 유한 전역 차원이 $\leq n$이면 $S^{-1}R$의 유한 전역
차원은 $\leq n$이다.
\end{enumerate}
\end{lemma}

\begin{proof}
$0 \to P_n \to P_{n - 1} \to \ldots \to P_0 \to M \to 0$을 사영
분해라 하자. 명제 \ref{algebra-proposition-localization-exact}에 의해 국소화는
완전이고, 보조정리 \ref{algebra-lemma-ascend-properties-modules}에 의해 각
$S^{-1}P_i$는 사영 $S^{-1}R$-가군이므로,
$0 \to S^{-1}P_n \to \ldots \to S^{-1}P_0 \to S^{-1}M \to 0$은
$S^{-1}M$의 사영 분해이다. 이는 (1)을 증명한다. $M'$을
$S^{-1}R$-가군이라 하자. $M' = S^{-1}M'$임에 주의하면 (2)는
(1)에서 따른다.
\end{proof}










\section{정칙환과 전역 차원}
\label{algebra-section-regular-finite-gl-dim}

\noindent
유한 전역 차원인 환에 관한 결과를 이용하여 정칙 국소환의 또 다른
특성화를 줄 수 있다.

\begin{proposition}
\label{algebra-proposition-regular-finite-gl-dim}
$R$을 차원 $d$인 정칙 국소환이라 하자. 깊이가 $e$인 모든 유한
$R$-가군 $M$은 유한 자유 분해
$$
0 \to F_{d-e} \to \ldots \to F_0 \to M \to 0.
$$
를 갖는다. 특히 정칙 국소환의 전역 차원은 $\leq d$이다.
\end{proposition}

\begin{proof}
첫째 명제는 보조정리 \ref{algebra-lemma-regular-mcm-free}와 보조정리
\ref{algebra-lemma-mcm-resolution}에서 따른다. 마지막 명제는 이 결과와 보조정리
\ref{algebra-lemma-finite-gl-dim}에서 따른다.
\end{proof}

\begin{lemma}
\label{algebra-lemma-finite-gl-dim-primes}
$R$을 뇌터환이라 하고 $n \geq 0$을 정수라 하자. $R$의 유한 전역 차원이
$\leq n$일 필요충분조건은 $R$의 모든 극대 아이디얼 $\mathfrak m$에 대해
환 $R_{\mathfrak m}$의 전역 차원이 $\leq n$인 것이다.
\end{lemma}

\begin{proof}
% P01-E0529
보조정리 \ref{algebra-lemma-localize-finite-gl-dim}에서 $R$의 유한 전역 차원이
$n$이면 모든 국소화 $R_{\mathfrak m}$의 유한 전역 차원이 $n$ 이하임을
보았다. 역으로 모든 $R_{\mathfrak m}$의 전역 차원이 $\leq n$이라고
가정하자. $M$을 유한 $R$-가군이라 하고
% P01-E0530
$0 \to K_n \to F_{n-1} \to \ldots \to F_0 \to M \to 0$을
$F_i$들이 유한 자유인 분해라 하자. 그러면 $K_n$은 유한 $R$-가군이다.
보조정리 \ref{algebra-lemma-independent-resolution}와 가정에 의해 모든 가군
$K_n \otimes_R R_{\mathfrak m}$은 사영이다. 따라서 보조정리
\ref{algebra-lemma-finite-projective}에 의해 가군 $K_n$은 유한 사영이다.
\end{proof}

\begin{lemma}
\label{algebra-lemma-length-resolution-residue-field}
$R$이 극대 아이디얼이 $\mathfrak m$이고 잉여체가 $\kappa$인 뇌터
국소환이라고 하자. 이 경우 $\kappa$의 사영 차원은
$\geq \dim_\kappa \mathfrak m / \mathfrak m^2$이다.
\end{lemma}

\begin{proof}
$x_1 , \ldots, x_n$을 $\mathfrak m$의 원소들로서 그 상들이
$\mathfrak m / \mathfrak m^2$의 기저를 이룬다고 하자.
$x_1, \ldots, x_n$에 대한 {\it Koszul 복합체}를 생각하자. 이는 사상들이
$$
0 \to \wedge^n R^n \to \wedge^{n-1} R^n \to \wedge^{n-2} R^n \to
\ldots \to \wedge^i R^n \to \ldots \to R^n \to R
$$
이고
$$
e_{j_1} \wedge \ldots \wedge e_{j_i}
\longmapsto
\sum_{a = 1}^i (-1)^{a + 1} x_{j_a} e_{j_1} \wedge \ldots
\wedge \hat e_{j_a} \wedge \ldots \wedge e_{j_i}
$$
로 주어진 복합체이다. 이것이 복합체 $K_{\bullet}(R, x_{\bullet})$임은
쉽게 알 수 있다. $K_{\bullet}(R, x_{\bullet})$의 마지막 사상의 여핵은
보조정리 \ref{algebra-lemma-NAK}의 (8)에 의해 $\kappa$임에 주의하자.

\medskip\noindent
$\kappa$의 유한 사영 차원이 $d$이면 유한 자유 $R$-가군들에 의한 길이
$d$의 분해 $F_{\bullet} \to \kappa$를 찾을 수 있다(보조정리
\ref{algebra-lemma-what-kind-of-resolutions-Noetherian-local}). 자명한 직합 인자를
분해에서 제거하면 기껏해야 분해가 짧아질 뿐이므로, 보조정리
\ref{algebra-lemma-add-trivial-complex}에 의해 복합체 $F_{\bullet}$의 모든 사상이
$\Im(F_i \to F_{i-1}) \subset \mathfrak m F_{i-1}$을 만족한다고 가정해도
된다. 보조정리 \ref{algebra-lemma-compare-resolutions}에 의해 $\kappa$ 위에서
항등사상을 유도하는 복합체의 사상
$\alpha : K_{\bullet}(R, x_{\bullet}) \to F_{\bullet}$을 찾을 수 있다.
사상 $\alpha_i : \wedge^i R^n = K_i(R, x_{\bullet}) \to F_i$가
$\alpha_i \otimes \kappa : \wedge^i \kappa^n \to F_i \otimes \kappa$가
단사라는 성질을 가짐을 귀납법으로 증명하겠다. 이로부터
$F_n \not = 0$이고 따라서 원하는 대로 $d \geq n$임을 알 수 있다.

\medskip\noindent
$R \xrightarrow{\alpha_0} F_0 \to \kappa$의 합성이 영이 아니므로
$i = 0$인 경우는 명백하다. $F_0$의 계수는 $1$이어야 함에 주의하자.
그렇지 않으면 여핵이 $\kappa$ 한 개뿐인 사상 $F_1 \to F_0$의 상이
$\mathfrak m F_0$에 포함될 수 없기 때문이다.

\medskip\noindent
다음으로 $i = 1$인 경우를 확인하자. 이 경우가 유익하다고 생각하기
때문이다. 귀납 단계가 $i = 0$인 경우로부터 $i = 1$인 경우를 이끌어
낼 것이므로 독자는 이 부분을 건너뛰어도 된다. 위에서
$F_0 = R$이고 $F_1 \to F_0 = R$의 상이 $\mathfrak m$임을 보았다.
가환 그림
$$
\begin{matrix}
R^n & = & K_1(R, x_{\bullet}) & \to & K_0(R, x_{\bullet}) & = & R \\
& & \downarrow & & \downarrow & & \downarrow \\
& & F_1 & \to & F_0 & = & R
\end{matrix}
$$
이 있고, 여기서 가장 오른쪽 수직 화살표는 단원에 의한 곱셈으로
주어진다. 따라서 합성 $R^n \to F_1 \to F_0 = R$의 상도
$\mathfrak m$과 같다. 그러므로
$\dim_\kappa (\mathfrak m / \mathfrak m^2) = n$이므로 사상
$R^n \otimes \kappa \to F_1 \otimes \kappa$는 단사여야 한다.

\medskip\noindent
$i \geq 1$이라 하고 모든 $j \leq i - 1$에 대해
$\alpha_j \otimes \kappa$의 단사성이 증명되었다고 가정하자. 가환 그림
$$
\begin{matrix}
\wedge^i R^n & = & K_i(R, x_{\bullet}) & \to & K_{i-1}(R, x_{\bullet})
& = & \wedge^{i-1} R^n \\
& & \downarrow & & \downarrow & & \\
& & F_i & \to & F_{i-1} & &
\end{matrix}
$$
을 생각하자. $\wedge^{i-1} \kappa^n \to F_{i-1} \otimes \kappa$가
단사임을 알고 있다. 따라서
$\wedge^{i-1} \kappa^n \otimes_{\kappa} \mathfrak m/\mathfrak m^2
\to F_{i-1} \otimes \mathfrak m/\mathfrak m^2$도 단사이다. 또한 복합체를
택한 방식에 의해 $F_i$는 $\mathfrak mF_{i-1}$로 사상되며 Koszul
복합체도 마찬가지이다. 따라서 가환 그림
$$
\begin{matrix}
\wedge^i \kappa^n & \to &
\wedge^{i-1} \kappa^n \otimes \mathfrak m/\mathfrak m^2 \\
\downarrow & & \downarrow \\
F_i \otimes \kappa & \to & F_{i-1} \otimes \mathfrak m/\mathfrak m^2
\end{matrix}
$$
을 얻는다. 이제 사상
$\wedge^i \kappa^n \to
\wedge^{i-1} \kappa^n \otimes \mathfrak m/\mathfrak m^2$가 단사임을
확인하는 것으로 충분하며, 이는 직접 계산할 수 있다.
\end{proof}

\begin{lemma}
\label{algebra-lemma-dim-gl-dim}
$R$을 뇌터 국소환이라 하자. 잉여체 $\kappa$의 $R$ 위의 유한 사영
차원이 $n$이라고 가정하자. 이 경우 $\dim(R) \geq n$이다.
\end{lemma}

\begin{proof}
$F_{\bullet}$을 유한 자유 $R$-가군들에 의한 $\kappa$의 유한 분해라 하자
(보조정리 \ref{algebra-lemma-what-kind-of-resolutions-Noetherian-local}). 자명한
직합 인자를 분해에서 제거하면 기껏해야 분해가 짧아질 뿐이므로,
보조정리 \ref{algebra-lemma-add-trivial-complex}에 의해 복합체 $F_{\bullet}$의
% P01-E0531
모든 사상이 $\Im(F_i \to F_{i-1}) \subset \mathfrak m F_{i-1}$을
만족한다고 가정해도 된다. $F_n \not = 0$이고 $i > n$에 대해
$F_i = 0$이라고 하자. 그러면 $\kappa$의 사영 차원은 $n$이다.
명제 \ref{algebra-proposition-what-exact}에 의해 $I(\varphi_n)$이 복합체를 택한
방식상 $R$과 같을 수 없으므로
$\text{depth}_{I(\varphi_n)}(R) \geq n$임을 알 수 있다. 따라서 보조정리
\ref{algebra-lemma-bound-depth}에 의해 $\dim(R) \geq n$이기도 하다.
\end{proof}

\begin{proposition}
\label{algebra-proposition-finite-gl-dim-regular}
$(R, \mathfrak m, \kappa)$를 뇌터 국소환이라 하자. 다음 조건들은
동치이다.
\begin{enumerate}
\item $\kappa$는 $R$-가군으로서 유한 사영 차원을 갖는다.
\item $R$은 유한 전역 차원을 갖는다.
\item $R$은 정칙 국소환이다.
\end{enumerate}
더욱이 이 경우 $R$의 전역 차원은
$\dim(R) = \dim_\kappa(\mathfrak m/\mathfrak m^2)$와 같다.
\end{proposition}

\begin{proof}
명제 \ref{algebra-proposition-regular-finite-gl-dim}에 의해
(3) $\Rightarrow$ (2)이고, (2) $\Rightarrow$ (1)은 자명하다. (1)을
가정하자. 보조정리 \ref{algebra-lemma-length-resolution-residue-field}와
\ref{algebra-lemma-dim-gl-dim}에 의해
$\dim(R) \geq \dim_\kappa(\mathfrak m /\mathfrak m^2)$임을 알 수 있다.
따라서 정의 \ref{algebra-definition-regular-local}과 그 앞의 논의에 의해 $R$은
정칙이다. 동치인 조건 (1) -- (3)이 성립한다고 가정하자. 명제
\ref{algebra-proposition-regular-finite-gl-dim}에 의해 $R$의 전역 차원은
$\dim(R)$ 이하이고, 보조정리
\ref{algebra-lemma-length-resolution-residue-field}에 의해
$\dim_\kappa(\mathfrak m/\mathfrak m^2)$ 이상이다. 따라서 명시한 등식이
성립한다.
\end{proof}

\begin{lemma}
\label{algebra-lemma-localization-of-regular-local-is-regular}
뇌터 국소환 $R$이 정칙 국소환일 필요충분조건은 유한 전역 차원을 갖는
것이다. 이 경우 모든 소 아이디얼 $\mathfrak p$에 대해
$R_{\mathfrak p}$는 정칙 국소환이다.
\end{lemma}

\begin{proof}
명제 \ref{algebra-proposition-finite-gl-dim-regular}와
\ref{algebra-proposition-regular-finite-gl-dim}에 의해 뇌터 국소환이 정칙
국소환일 필요충분조건은 유한 전역 차원을 갖는 것임을 알 수 있다.
또한 모든 국소화 $R_{\mathfrak p}$는 유한 전역 차원을 가지므로(보조정리
\ref{algebra-lemma-localize-finite-gl-dim}) 정칙 국소환이다.
\end{proof}

\noindent
보조정리 \ref{algebra-lemma-localization-of-regular-local-is-regular}에 의해 다음
정의는 앞의 정칙 국소환의 정의와 충돌하지 않으므로 의미가 있다.

\begin{definition}
\label{algebra-definition-regular}
뇌터환 $R$의 소 아이디얼에서의 모든 국소화 $R_{\mathfrak p}$가 정칙
국소환이면 이를 {\it 정칙}이라고 한다.
\end{definition}

\noindent
극대 아이디얼에서의 국소환들이 정칙일 것만 요구해도 충분하다. 이는
$R$이 유한 전역 차원을 갖는다고 요구하는 것과 같지 않음에 주의하자.
$R$이 뇌터라고 가정해도 마찬가지이다. 실제로 유한 전역 차원을 갖지
않는 정칙 뇌터환의 예가 있는데, 그 환은 유한 차원을 갖지 않기 때문이다.

\begin{lemma}
\label{algebra-lemma-finite-gl-dim-finite-dim-regular}
$R$을 뇌터환이라 하자. 다음 조건들은 동치이다.
\begin{enumerate}
\item $R$의 유한 전역 차원은 $n$이다.
\item $R$은 차원이 $n$인 정칙환이다.
\item 모든 극대 아이디얼에서의 국소화 $R_{\mathfrak m}$이 차원
$\leq n$인 정칙환이고 적어도 하나의 $\mathfrak m$에서 등호가 성립하도록
하는 정수 $n$이 존재한다.
\item 모든 소 아이디얼에서의 국소화 $R_{\mathfrak p}$가 차원
$\leq n$인 정칙환이고 적어도 하나의 $\mathfrak p$에서 등호가 성립하도록
하는 정수 $n$이 존재한다.
\end{enumerate}
\end{lemma}

\begin{proof}
이는 위의 논의에서 따른다. 더 정확히 말하면 정의
\ref{algebra-definition-regular}, 보조정리 \ref{algebra-lemma-finite-gl-dim-primes}, 명제
\ref{algebra-proposition-finite-gl-dim-regular}를 결합하면 된다.
\end{proof}

\begin{lemma}
\label{algebra-lemma-flat-under-regular}
$R \to S$를 뇌터 국소환들의 국소 준동형사상이라 하자. $R \to S$가
평탄이고 $S$가 정칙이라고 가정하자. 그러면 $R$은 정칙이다.
\end{lemma}

\begin{proof}
$\mathfrak m \subset R$을 극대 아이디얼이라 하고
$\kappa = R/\mathfrak m$을 잉여체라 하자. $d = \dim S$로 두자.
각 $F_i$가 유한 자유 $R$-가군인 임의의 분해
$F_\bullet \to \kappa$를 택하고
% P01-E0532
$K_d = \Ker(F_{d - 1} \to F_{d - 2})$로 두자. $R \to S$의 평탄성에
의해 복합체
$0 \to K_d \otimes_R S \to F_{d - 1} \otimes_R S \to \ldots
\to F_0 \otimes_R S \to \kappa \otimes_R S \to 0$은 여전히 완전이다.
% P01-E0533
$S$의 전역 차원은 $d$이므로(명제
\ref{algebra-proposition-regular-finite-gl-dim}),
$K_d \otimes_R S$는 유한 자유 $S$-가군이다(보조정리
\ref{algebra-lemma-independent-resolution}도 보라). 보조정리
\ref{algebra-lemma-finite-projective-descends}에 의해 $K_d$는 유한 자유
$R$-가군이다. 따라서 $\kappa$는 유한 사영 차원을 가지며, 명제
\ref{algebra-proposition-finite-gl-dim-regular}에 의해 $R$은 정칙이다.
\end{proof}








\section{Auslander--Buchsbaum}
\label{algebra-section-Auslander-Buchsbaum}

\noindent
다음 결과는 \cite{Auslander-Buchsbaum}에서 찾을 수 있다.

\begin{proposition}
\label{algebra-proposition-Auslander-Buchsbaum}
$R$을 뇌터 국소환이라 하고 $M$을 유한 사영 차원
$\text{pd}_R(M)$을 갖는 영이 아닌 유한 $R$-가군이라 하자. 그러면
$$
\text{depth}(R) = \text{pd}_R(M) + \text{depth}(M)
$$
이다.
\end{proposition}

\begin{proof}
$\text{depth}(M)$에 대한 귀납법으로 증명한다. 가장 흥미로운 경우는
$\text{depth}(M) = 0$인 경우이다. 이 경우
$$
0 \to R^{n_e} \to R^{n_{e-1}} \to \ldots \to R^{n_0} \to M \to 0
$$
을 최소 유한 자유 분해라 하자. 그러면 $e = \text{pd}_R(M)$이다.
보조정리 \ref{algebra-lemma-add-trivial-complex}에 의해 복합체의 사상들을 나타내는
모든 행렬 계수가 $R$의 극대 아이디얼에 포함된다고 가정해도 된다.
한편 명제 \ref{algebra-proposition-what-exact}에 의해
$\text{depth}(R) \geq e$임을 알 수 있다. 다른 한편 긴 완전열을 짧은
완전열들
% P01-E0534
\begin{align*}
0 \to R^{n_e} \to R^{n_{e - 1}} \to K_{e - 2} \to 0,\\
0 \to K_{e - 2} \to R^{n_{e - 2}} \to K_{e - 3} \to 0,\\
\ldots,\\
0 \to K_0 \to R^{n_0} \to M \to 0
\end{align*}
로 나누고 보조정리 \ref{algebra-lemma-depth-in-ses}를 사용하면
\begin{align*}
\text{depth}(K_{e - 2}) \geq \text{depth}(R) - 1,\\
\text{depth}(K_{e - 3}) \geq \text{depth}(R) - 2,\\
\ldots,\\
\text{depth}(K_0) \geq \text{depth}(R) - (e - 1),\\
\text{depth}(M) \geq \text{depth}(R) - e
\end{align*}
임을 알 수 있다. $\text{depth}(M) = 0$이므로
$\text{depth}(R) \leq e$임을 얻는다. 이로써
$\text{depth}(M) = 0$인 경우의 증명이 끝난다.

\medskip\noindent
귀납 단계. $\text{depth}(M) > 0$이면 $M$과 $R$ 모두에서 영인자가 아닌
$x \in \mathfrak m$을 택한다. 실제로 $\text{pd}_R(M) > 0$이면 앞서
언급한 명제 \ref{algebra-proposition-what-exact}에 의해
$\text{depth}(R) > 0$이고, $\text{pd}_R(M) = 0$이면 $M$은 유한 자유이므로
$\text{depth}(R) = \text{depth}(M) > 0$이어서 이는 가능하다. 따라서
보조정리 \ref{algebra-lemma-depth-in-ses} 등에 의해
$\text{depth}(R \oplus M) > 0$이고, $R$과 $M$ 모두에서 영인자가 아닌
% P01-E0535
$x \in \mathfrak m$을 찾을 수 있다. 위와 같이
$$
0 \to R^{n_e} \to R^{n_{e-1}} \to \ldots \to R^{n_0} \to M \to 0
$$
을 최소 분해라 하자. 뱀 보조정리를 적용하면
$$
0 \to (R/xR)^{n_e} \to (R/xR)^{n_{e-1}} \to \ldots \to (R/xR)^{n_0} \to
M/xM \to 0
$$
도 최소 분해임을 알 수 있다. 따라서
$\text{pd}_R(M) = \text{pd}_{R/xR}(M/xM)$이다. 보조정리
\ref{algebra-lemma-depth-drops-by-one}에 의해
$\text{depth}(R/xR) = \text{depth}(R) - 1$이고
$\text{depth}(M/xM) = \text{depth}(M) - 1$이다. 지금까지 모든 깊이는
$R$-가군으로서의 깊이였지만,
$\text{depth}_R(M/xM) = \text{depth}_{R/xR}(M/xM)$이고 $R/xR$에 대해서도
마찬가지임을 관찰하자. 귀납 가정에 의해 $R/xR$ 위의 $M/xM$에 대해
Auslander--Buchsbaum 공식이 성립한다. $R/xR$과 $M/xM$의 깊이가 모두
하나씩 줄고 사영 차원은 변하지 않았으므로 결론이 따른다.
\end{proof}









\section{준동형사상과 차원}
\label{algebra-section-homomorphism-dimension}

\noindent
이 절에서는 환들 사이에 사상이 있을 때 그 차원들을 연관시키는 쉬운
결과들을 모은다.

\begin{lemma}
\label{algebra-lemma-dimension-going-up}
$R \to S$가 정의 \ref{algebra-definition-going-up-down}의 상승 성질 또는 정의
\ref{algebra-definition-going-up-down}의 하강 성질을 만족하는 환 준동형사상이라고
하자. 또한
$\Spec(S) \to \Spec(R)$가 전사라고 가정하자. 그러면
$\dim(R) \leq \dim(S)$이다.
\end{lemma}

\begin{proof}
상승 성질을 가정하자. $R$의 소 아이디얼들의 임의의 사슬
$\mathfrak p_0 \subset \mathfrak p_1 \subset \ldots
\subset \mathfrak p_e$를 택하자. 전사성에 의해 $\mathfrak p_0$로
사상되는 소 아이디얼 $\mathfrak q_0$을 택할 수 있다. 상승 성질에 의해
이를 $\mathfrak p_i$ 위에 놓이는 소 아이디얼 $\mathfrak q_i$들의 길이
$e$인 사슬로 확장할 수 있다. 따라서 $\dim(S) \geq \dim(R)$이다.
하강 성질의 경우도 완전히 같다. 순수 위상적 판본은 위상수학의 보조정리
\ref{topology-lemma-dimension-specializations-lift}도 보라.
\end{proof}

\begin{lemma}
\label{algebra-lemma-going-up-maximal-on-top}
$R \to S$가 정의 \ref{algebra-definition-going-up-down}의 상승 성질을 갖는 환
준동형사상이라고 하자. 만일
% P01-E0536
$\mathfrak q \subset S$가 극대 아이디얼이면, $R$에서 $\mathfrak q$의
역상도 극대 아이디얼이다.
\end{lemma}

\begin{proof}
자명하다.
\end{proof}

\begin{lemma}
\label{algebra-lemma-integral-dim-up}
$R \to S$가 환 준동형사상이고 $S$가 $R$ 위에서 정수적이라고 하자.
그러면 $\dim (R) \geq \dim(S)$이고, $\Spec(S)$의 모든 닫힌점은
$\Spec(R)$의 닫힌점으로 사상된다.
\end{lemma}

\begin{proof}
보조정리 \ref{algebra-lemma-integral-no-inclusion},
\ref{algebra-lemma-going-up-maximal-on-top} 및 정의들에서 즉시 따른다.
\end{proof}

\begin{lemma}
\label{algebra-lemma-integral-sub-dim-equal}
$R \subset S$이고 $S$가 $R$ 위에서 정수적이라고 하자. 그러면
$\dim(R) = \dim(S)$이다.
\end{lemma}

\begin{proof}
보조정리 \ref{algebra-lemma-integral-going-up},
\ref{algebra-lemma-integral-overring-surjective},
\ref{algebra-lemma-dimension-going-up}, \ref{algebra-lemma-integral-dim-up}을 결합하면 된다.
\end{proof}

\begin{definition}
\label{algebra-definition-fibre}
$R \to S$가 환 준동형사상이라고 하자. $\mathfrak q \subset S$를
$R$의 소 아이디얼 $\mathfrak p$ 위에 놓이는 소 아이디얼이라 하자.
{\it $\mathfrak q$에서 올의 국소환}은 국소환
$$
S_{\mathfrak q}/\mathfrak pS_{\mathfrak q}
=
(S/\mathfrak pS)_{\mathfrak q}
=
(S \otimes_R \kappa(\mathfrak p))_{\mathfrak q}
$$
이다.
\end{definition}

\begin{lemma}
\label{algebra-lemma-dimension-base-fibre-total}
$R \to S$를 뇌터환들의 준동형사상이라 하자. $\mathfrak q \subset S$를
소 아이디얼 $\mathfrak p$ 위에 놓이는 소 아이디얼이라 하자. 그러면
$$
\dim(S_{\mathfrak q})
\leq
\dim(R_{\mathfrak p})
+
\dim(S_{\mathfrak q}/\mathfrak pS_{\mathfrak q}).
$$
\end{lemma}

\begin{proof}
명제 \ref{algebra-proposition-dimension}의 차원 특성화를 사용한다.
$x_1, \ldots, x_d$를 $\mathfrak p$의 원소들로서
$d = \dim(R_{\mathfrak p})$이고 $R_{\mathfrak p}$의 정의 아이디얼을
생성하도록 택하자. $y_1, \ldots, y_e$를 $\mathfrak q$의 원소들로서
$e = \dim(S_{\mathfrak q}/\mathfrak pS_{\mathfrak q})$이고
$S_{\mathfrak q}/\mathfrak pS_{\mathfrak q}$의 정의 아이디얼을 생성하도록
택하자. $S_{\mathfrak q}/(x_1, \ldots, x_d, y_1, \ldots, y_e)$의 극대
아이디얼이 멱영임은 명백하다. 따라서
$x_1, \ldots, x_d, y_1, \ldots, y_e$는 $S_{\mathfrak q}$의 정의
아이디얼을 생성한다.
\end{proof}

\begin{lemma}
\label{algebra-lemma-dimension-base-fibre-equals-total}
$R \to S$를 뇌터환들의 준동형사상이라 하자. $\mathfrak q \subset S$를
소 아이디얼 $\mathfrak p$ 위에 놓이는 소 아이디얼이라 하자.
$R \to S$가 하강 성질을 만족한다고 가정하자(예를 들어
$R \to S$가 평탄이면 보조정리 \ref{algebra-lemma-flat-going-down}에 의해 그렇다).
그러면
$$
\dim(S_{\mathfrak q})
=
\dim(R_{\mathfrak p})
+
\dim(S_{\mathfrak q}/\mathfrak pS_{\mathfrak q}).
$$
\end{lemma}

\begin{proof}
보조정리 \ref{algebra-lemma-dimension-base-fibre-total}에 의해 부등식
$\dim(S_{\mathfrak q}) \leq
\dim(R_{\mathfrak p}) + \dim(S_{\mathfrak q}/\mathfrak pS_{\mathfrak q})$를
얻는다. 등식을 얻기 위해
$\mathfrak pS \subset \mathfrak q_0 \subset \mathfrak q_1 \subset \ldots
\subset \mathfrak q_d = \mathfrak q$를
$d = \dim(S_{\mathfrak q}/\mathfrak pS_{\mathfrak q})$인 소 아이디얼들의
사슬로 택하자. 한편 소 아이디얼들의 사슬
$\mathfrak p_0 \subset \mathfrak p_1 \subset \ldots \subset \mathfrak p_e
= \mathfrak p$를 $e = \dim(R_{\mathfrak p})$가 되도록 택하자. 하강 정리에
의해 $\mathfrak p_{e-1}$ 위에 놓이는
$\mathfrak q_{-1} \subset \mathfrak q_0$을 택할 수 있다. 이어서
$\mathfrak p_{e-2}$ 위에 놓이는
$\mathfrak q_{-2} \subset \mathfrak q_{-1}$을 택할 수 있다. 귀납적으로
계속하면 길이 $e + d$인 사슬
$\mathfrak q_{-e} \subset \ldots \subset \mathfrak q_d$를 얻는다.
\end{proof}

\begin{lemma}
\label{algebra-lemma-flat-over-regular-with-regular-fibre}
$R \to S$를 뇌터 국소환들의 국소 준동형사상이라 하자. 다음을 가정하자.
\begin{enumerate}
\item $R$은 정칙이다.
\item $S/\mathfrak m_RS$는 정칙이다.
\item $R \to S$는 평탄이다.
\end{enumerate}
그러면 $S$는 정칙이다.
\end{lemma}

\begin{proof}
보조정리 \ref{algebra-lemma-dimension-base-fibre-equals-total}에 의해
$\dim(S) = \dim(R) + \dim(S/\mathfrak m_RS)$이다.
$d = \dim(R)$이고 그 극대 아이디얼을 생성하도록
$x_1, \ldots, x_d \in \mathfrak m_R$을 택하자. 또한
$e = \dim(S/\mathfrak m_RS)$이고 $S/\mathfrak m_RS$의 극대 아이디얼을
생성하도록 $y_1, \ldots, y_e \in \mathfrak m_S$를 택하자. 그러면
$x_1, \ldots, x_d, y_1, \ldots, y_e$는 $S$의 극대 아이디얼을 생성하고
$e + d = \dim(S)$임을 알 수 있다.
\end{proof}

\noindent
아래 보조정리는 나중에 체 위 유한형 환이 Cohen--Macaulay일 필요충분조건이
다항식환 위에서 준유한 평탄인 것임을 보이는 데 쓰인다. 이는 보조정리
\ref{algebra-lemma-CM-over-regular-flat}의 부분적 역이다.

\begin{lemma}
\label{algebra-lemma-finite-flat-over-regular-CM}
$R \to S$를 뇌터 국소환들의 국소 준동형사상이라 하고 $R$이
Cohen--Macaulay라고 가정하자. $S$가 $R$ 위 유한 평탄이거나,
$S$가 $R$ 위 평탄이고 $\dim(S) \leq \dim(R)$이면, $S$는
Cohen--Macaulay이고 $\dim(R) = \dim(S)$이다.
\end{lemma}

\begin{proof}
$x_1, \ldots, x_d \in \mathfrak m_R$을 길이가 $d = \dim(R)$인 정칙열로
택하자. 보조정리 \ref{algebra-lemma-flat-increases-depth}에 의해 이는 $S$에서
정칙열로 사상된다. 따라서 $\dim(S) \leq d$이면 $S$는
Cohen--Macaulay이다. $S$가 $R$ 위 유한 평탄인 경우 이는 보조정리
\ref{algebra-lemma-integral-sub-dim-equal}에 의해 참이고, 둘째 경우에는 가정했다.
\end{proof}









\section{차원 공식}
\label{algebra-section-dimension-formula}

\noindent
카테너리(정의 \ref{algebra-definition-catenary})와 보편적으로 카테너리
(정의 \ref{algebra-definition-universally-catenary})의 정의를 상기하자.

\begin{lemma}
\label{algebra-lemma-dimension-formula}
$R \to S$를 환 준동형사상이라 하자. $\mathfrak q$를 $R$의 소 아이디얼
$\mathfrak p$ 위에 놓이는 $S$의 소 아이디얼이라 하자. 다음을 가정하자.
\begin{enumerate}
\item $R$은 뇌터이다.
\item $R \to S$는 유한형이다.
\item $R$, $S$는 정역이다.
\item $R \subset S$이다.
\end{enumerate}
그러면
$$
\text{height}(\mathfrak q)
\leq
\text{height}(\mathfrak p) + \text{trdeg}_R(S)
- \text{trdeg}_{\kappa(\mathfrak p)} \kappa(\mathfrak q)
$$
이고, $R$이 보편적으로 카테너리이면 등호가 성립한다.
\end{lemma}

\begin{proof}
$R \subset S' \subset S$이고 $S'$이 $S$의 유한 생성 $R$-부분대수라고
하자. 이 경우 $\mathfrak q' = S' \cap \mathfrak q$로 두자. 체의 탑에서
초월차수의 가법성(체, 보조정리
\ref{fields-lemma-transcendence-degree-tower})에 의해 환 준동형사상
$R \to S'$과 $S' \to S$에 대한 보조정리가 $R \to S$에 대한 보조정리를
함의한다. 따라서 $R$ 위에서 $S$의 생성원 수에 대해 귀납하여 $S$가
$R$ 위에서 하나의 원소로 생성되는 경우로 환원할 수 있다.

\medskip\noindent
경우 I: $S = R[x]$가 $R$ 위 다항식대수라고 하자. 이 경우
$\text{trdeg}_R(S) = 1$이다. 또한 $R \to S$는 평탄이므로
$$
\dim(S_{\mathfrak q}) =
\dim(R_{\mathfrak p}) + \dim(S_{\mathfrak q}/\mathfrak pS_{\mathfrak q})
$$
이다. 보조정리 \ref{algebra-lemma-dimension-base-fibre-equals-total}을 보라.
$\mathfrak r = \mathfrak pS$로 두자. 그러면
$\text{trdeg}_{\kappa(\mathfrak p)} \kappa(\mathfrak q) = 1$일
필요충분조건은 $\mathfrak q = \mathfrak r$인 것이고, 이는
$\dim(S_{\mathfrak q}/\mathfrak pS_{\mathfrak q}) = 0$을 함의한다.
마찬가지로
$\text{trdeg}_{\kappa(\mathfrak p)} \kappa(\mathfrak q) = 0$일
필요충분조건은 $\mathfrak r \subset \mathfrak q$가 진포함인 것이고,
이는 $\dim(S_{\mathfrak q}/\mathfrak pS_{\mathfrak q}) = 1$을 함의한다.
따라서 경우 I에서는 모든 경우에 등호가 성립하여 증명이 끝난다.

\medskip\noindent
경우 II: $S = R[x]/\mathfrak n$이고 $\mathfrak n \not = 0$이라고 하자.
이 경우 $\text{trdeg}_R(S) = 0$이다.
% P01-E0537
$\mathfrak q$에 대응하는 소 아이디얼을 $\mathfrak q' \subset R[x]$라
하자. 그러면
$$
S_{\mathfrak q} = (R[x])_{\mathfrak q'}/\mathfrak n(R[x])_{\mathfrak q'}
$$
이다. 앞의 경우에 의해
$\dim((R[x])_{\mathfrak q'}) =
\dim(R_{\mathfrak p}) + 1
- \text{trdeg}_{\kappa(\mathfrak p)} \kappa(\mathfrak q)$이다.
$\mathfrak n \not = 0$이므로 $S_{\mathfrak q}$의 차원은 적어도 하나
줄어든다. 보조정리 \ref{algebra-lemma-one-equation}을 보라. 이는 보조정리의
부등식을 증명한다.
% P01-E0538
$R$이 보편적으로 카테너리인 경우 등호를 보이기 위해,
$\mathfrak n \subset R[x]$가 높이 하나인 소 아이디얼임에 주의하자.
실제로 이는 $R$의 분수체 $F$에 대해 $F[x]$의 영이 아닌 소 아이디얼에
대응한다. 따라서
$S_\mathfrak q = R[x]_{\mathfrak q'}/\mathfrak nR[x]_{\mathfrak q'}$의
임의의 극대 소 아이디얼 사슬은 $R[x]$ 안에서
% P01-E0539
$\mathfrak q'$와 $(0)$ 사이의 길이가 하나 더 큰 극대 소 아이디얼 사슬에
대응한다. $R$이 보편적으로 카테너리이면 이 사슬들은 모두 같은 길이를
가지며 그 길이는 $\mathfrak q'$의 높이와 같다. 이로부터
$\dim(S_\mathfrak q) = \dim(R[x]_{\mathfrak q'}) - 1$이고 원하는 대로
등호가 성립함을 알 수 있다.
\end{proof}

\noindent
다음 보조정리는 일반적으로 유한인 사상이 여차원 $1$에서 준유한인
경향이 있음을 말한다.

\begin{lemma}
\label{algebra-lemma-finite-in-codim-1}
$A \to B$를 환 준동형사상이라 하자. 다음을 가정하자.
\begin{enumerate}
\item $A \subset B$는 정역들의 확대이다.
\item 유도된 분수체 확대는 유한이다.
\item $A$는 뇌터이다.
\item $A \to B$는 유한형이다.
\end{enumerate}
$\mathfrak p \subset A$를 높이 $1$인 소 아이디얼이라 하자. 그러면
$\mathfrak p$ 위에 놓이는 $B$의 소 아이디얼은 유한 개 이하이고 모두
높이 $1$이다.
\end{lemma}

\begin{proof}
차원 공식(보조정리 \ref{algebra-lemma-dimension-formula})에 의해
$\mathfrak p$ 위에 놓이는 임의의 소 아이디얼 $\mathfrak q$에 대해
$$
\dim(B_{\mathfrak q}) \leq
\dim(A_{\mathfrak p}) - \text{trdeg}_{\kappa(\mathfrak p)} \kappa(\mathfrak q).
$$
이다. 정역 $B_\mathfrak q$는 적어도 $2$개의 소 아이디얼을 가지므로
$\dim(B_{\mathfrak q}) \geq 1$이다. 따라서
$\dim(B_{\mathfrak q}) = 1$이고 확대
$\kappa(\mathfrak p) \subset \kappa(\mathfrak q)$는 대수적이다.
그러므로 $\mathfrak q$는 올
$\Spec(B \otimes_A \kappa(\mathfrak p))$의 닫힌점을 정의한다. 보조정리
\ref{algebra-lemma-finite-residue-extension-closed}을 보라.
$B \otimes_A \kappa(\mathfrak p)$가 뇌터환이므로 올
$\Spec(B \otimes_A \kappa(\mathfrak p))$는 뇌터 위상공간이다. 보조정리
\ref{algebra-lemma-Noetherian-topology}을 보라. 닫힌점들로 이루어진 뇌터
위상공간은 유한하다. 예를 들어 위상수학, 보조정리
\ref{topology-lemma-Noetherian}을 보라.
\end{proof}









\section{체 위 유한형 대수의 차원}
\label{algebra-section-dimension-finite-type-algebras}

\noindent
이 절에서는 체 위 다항식환의 차원을 계산한다. 또한 체 위 유한형 정역의
차원이 극대 아이디얼에서의 국소환들의 차원임을 증명한다. 기초체 위의
초월차수와의 관계는 절
\ref{algebra-section-dimension-finite-type-algebras-reprise}에서 확립한다.

\begin{lemma}
\label{algebra-lemma-dim-affine-space}
$\mathfrak m$을 $k[x_1, \ldots, x_n]$의 극대 아이디얼이라 하자.
아이디얼 $\mathfrak m$은 $n$개의 원소로 생성된다.
$k[x_1, \ldots, x_n]_{\mathfrak m}$의 차원은 $n$이다. 따라서
$k[x_1, \ldots, x_n]_{\mathfrak m}$은 차원 $n$인 정칙 국소환이다.
\end{lemma}

\begin{proof}
Hilbert 영점정리(정리 \ref{algebra-theorem-nullstellensatz})에 의해 잉여체
$\kappa = \kappa(\mathfrak m)$가 $k$의 유한 확대임을 안다.
% P01-E0540
$\alpha_i \in \kappa$를 $x_i$의 상이라 하자.
% P01-E0541
$i = 1, \ldots, n$에 대해
$\kappa_i = k(\alpha_1, \ldots, \alpha_i) \subset \kappa$로 두고
$\kappa_0 = k$로 두자. 체론에 의해
$\kappa_i = k[\alpha_1, \ldots, \alpha_i]$임에 주의하자. 다음과 같이
$f_i \in \mathfrak m \cap k[x_1, \ldots, x_i]$를 귀납적으로 정의한다.
$P_i(T) \in \kappa_{i-1}[T]$를 $\kappa_{i-1}$ 위에서 $\alpha_i$의
일계수 최소다항식이라 하자. $Q_i(T) \in k[x_1, \ldots, x_{i-1}][T]$를
$P_i(T)$의 같은 차수를 갖는 일계수 끌어올림이라 하고
$f_i = Q_i(x_i)$로 두자.
$d_i = \deg_T(P_i) = \deg_T(Q_i) = \deg_{x_i}(f_i)$이면 체, 보조정리
\ref{fields-lemma-multiplicativity-degrees}와
\ref{fields-lemma-degree-minimal-polynomial}에 의해
$d_1d_2\ldots d_i = [\kappa_i : k]$임에 주의하자.

\medskip\noindent
모든 $i = 0, 1, \ldots, n$에 대해 동형사상
$$
\psi_i : k[x_1, \ldots, x_i] /(f_1, \ldots, f_i) \cong \kappa_i.
$$
가 있다고 주장한다. 구성에 의해 합성
$k[x_1, \ldots, x_i] \to k[x_1, \ldots, x_n] \to \kappa$의 상은
$\kappa_i$이고 $f_1, \ldots, f_i$는 핵에 속한다. 따라서 전사
준동형사상을 얻는다. $\psi_i$가 단사임을 귀납법으로 증명하자.
$i = 0$인 경우는 명백하다. $i$에 대한 명제가 성립한다고 하고
$i + 1$에 대해 증명하자. 환 확대
$k[x_1, \ldots, x_i]/(f_1, \ldots, f_i) \to
k[x_1, \ldots, x_{i + 1}]/(f_1, \ldots, f_{i + 1})$는 체 위에서 하나의
$1$개 원소와 하나의 기약 방정식으로 생성된다. 초등 체론에 의해
$k[x_1, \ldots, x_{i + 1}]/(f_1, \ldots, f_{i + 1})$는 체이고, 따라서
% P01-E0542
$\psi_i$는 단사이다.

\medskip\noindent
이로부터 $\mathfrak m = (f_1, \ldots, f_n)$임을 알 수 있다. 또한
$$
k[x_1, \ldots, x_n]/(f_1, \ldots, f_i)
\cong
\kappa_i[x_{i + 1}, \ldots, x_n].
$$
임도 알 수 있다. 따라서 $(f_1, \ldots, f_i)$는 소 아이디얼이다.
그러므로
$$
(0) \subset (f_1) \subset (f_1, f_2) \subset \ldots \subset
(f_1, \ldots, f_n) = \mathfrak m
$$
은 길이 $n$인 소 아이디얼들의 사슬이다. 보조정리가 따른다.
\end{proof}

\begin{proposition}
\label{algebra-proposition-finite-gl-dim-polynomial-ring}
체 위 $n$변수 다항식대수는 정칙환이고 전역 차원은 $n$이다.
극대 아이디얼에서의 모든 국소화는 차원 $n$인 정칙 국소환이다.
\end{proposition}

\begin{proof}
보조정리 \ref{algebra-lemma-dim-affine-space}에 의해 극대 아이디얼에서의 모든
국소화 $k[x_1, \ldots, x_n]_{\mathfrak m}$은 차원 $n$인 정칙
국소환이다. 따라서 보조정리
\ref{algebra-lemma-finite-gl-dim-finite-dim-regular}에 의해 결론이 따른다.
\end{proof}

\begin{lemma}
\label{algebra-lemma-dimension-height-polynomial-ring}
$k$를 체라 하고
$\mathfrak p \subset \mathfrak q \subset k[x_1, \ldots, x_n]$를 소
아이디얼들의 쌍이라 하자. $\mathfrak p$와 $\mathfrak q$ 사이의 임의의
극대 소 아이디얼 사슬의 길이는
$\text{height}(\mathfrak q) - \text{height}(\mathfrak p)$이다.
\end{lemma}

\begin{proof}
명제 \ref{algebra-proposition-finite-gl-dim-polynomial-ring}에 의해
$k[x_1, \ldots, x_n]$의 모든 국소환은 정칙이다. 따라서 보조정리
\ref{algebra-lemma-regular-ring-CM}에 의해 모든 국소환은 Cohen--Macaulay이다.
극대 아이디얼에서의 국소환들은 차원 $n$이므로
$k[x_1, \ldots, x_n]$의 모든 극대 소 아이디얼 사슬은 길이 $n$이다.
보조정리 \ref{algebra-lemma-maximal-chain-CM}을 보라. 그러므로 예를 들어 보조정리
\ref{algebra-lemma-CM-dim-formula}에 의해 $(0)$과 $\mathfrak p$ 사이의 모든
극대 소 아이디얼 사슬은 길이 $\text{height}(\mathfrak p)$이다.
이들을 결합하면 보조정리의 주장이 따른다.
\end{proof}

\begin{lemma}
\label{algebra-lemma-dimension-spell-it-out}
$k$를 체라 하고 $S$를 정역인 유한형 $k$-대수라 하자. 그러면 $S$의
임의의 극대 아이디얼 $\mathfrak m$에 대해
$\dim(S) = \dim(S_{\mathfrak m})$이다. 다시 말해 모든 극대 소 아이디얼
사슬의 길이는 $S$의 차원과 같다.
\end{lemma}

\begin{proof}
$S = k[x_1, \ldots, x_n]/\mathfrak p$로 쓰자. 명제
\ref{algebra-proposition-finite-gl-dim-polynomial-ring}와 보조정리
\ref{algebra-lemma-dimension-height-polynomial-ring}에 의해 $S$의 모든 극대 소
아이디얼 사슬은(이들은 반드시 극대 아이디얼에서 끝난다) 길이
$n - \text{height}(\mathfrak p)$를 갖는다. 따라서 이 수는 $S$의
차원이며 $S$의 임의의 극대 아이디얼 $\mathfrak m$에 대한
$S_{\mathfrak m}$의 차원이기도 하다.
\end{proof}

\noindent
위상공간 $X$의 점 $x$에서 차원 $\dim_x(X)$을 위상수학, 정의
\ref{topology-definition-Krull}에서 정의했음을 상기하자.

\begin{lemma}
\label{algebra-lemma-dimension-at-a-point-finite-type-over-field}
$k$를 체, $S$를 유한형 $k$-대수, $X = \Spec(S)$라 하자.
$\mathfrak p \subset S$를 소 아이디얼이라 하고 $x \in X$를 대응하는
점이라 하자. 다음 수들은 서로 같다.
\begin{enumerate}
\item $\dim_x(X)$.
\item $x$를 지나는 $X$의 기약성분 $Z$들에 대한 최댓값
$\max \dim(Z)$.
\item $\mathfrak p \subset \mathfrak m$인 극대 아이디얼
$\mathfrak m$들에 대한 최솟값 $\min \dim(S_{\mathfrak m})$.
\end{enumerate}
\end{lemma}

\begin{proof}
$X = \bigcup_{i \in I} Z_i$를 $X$의 기약성분들로의 분해라 하자.
그 수는 유한하다(보조정리 \ref{algebra-lemma-obvious-Noetherian}과
\ref{algebra-lemma-Noetherian-topology}을 보라). $I' = \{i \mid x \in Z_i\}$로
두고 $T = \bigcup_{i \not \in I'} Z_i$로 두자. 그러면
$U = X \setminus T$는 점 $x$를 포함하는 $X$의 열린 부분집합이다.
수 (2)는 $\max_{i \in I'} \dim(Z_i)$이다. $x \in W$인 임의의 열린집합
$W \subset U$에 대해 $W$의 기약성분은 $i \in I'$에 대한 기약집합
$W_i = Z_i \cap W$이고, $x$는 이들 각각에 포함된다. $X$가 Jacobson이므로
각 $W_i$, $i \in I'$는 닫힌점을 포함한다. 절
\ref{algebra-section-ring-jacobson}을 보라. $W_i \subset Z_i$이므로
$\dim(W_i) \leq \dim(Z_i)$이다. 닫힌점의 존재와 보조정리
\ref{algebra-lemma-dimension-spell-it-out}에 의해 열린집합 $W_i$ 안에는 길이가
$\dim(Z_i)$인 기약 닫힌 부분집합들의 사슬이 있다. 따라서 모든
$i \in I'$에 대해 $\dim(W_i) = \dim(Z_i)$이다. 그러므로 $\dim(W)$는
수 (2)와 같다. 이는 (1) $ = $ (2)를 증명한다.

\medskip\noindent
$\mathfrak p$를 포함하는 임의의 극대 아이디얼
$\mathfrak m \supset \mathfrak p$를 택하고, $x_0 \in X$를 대응하는
점이라 하자. 우선 $x_0$은 모든 기약성분 $Z_i$, $i \in I'$에 포함된다.
$\mathfrak q_i$를 기약성분 $Z_i$에 대응하는 $S$의 극소 소 아이디얼이라
하자. $x_0 \in Z_i$인 각 $i$에 대해(이는
$\mathfrak m \supset \mathfrak q_i$와 동치이다) 전사
$$
S_{\mathfrak m} \longrightarrow
S_\mathfrak m/\mathfrak q_i S_\mathfrak m =(S/\mathfrak q_i)_{\mathfrak m}
$$
가 있다. 또한 이렇게 얻은 소 아이디얼
$\mathfrak q_i S_\mathfrak m$들은 뇌터 국소환 $S_{\mathfrak m}$의 극소
소 아이디얼을 모두 준다. 보조정리
\ref{algebra-lemma-irreducible-components-containing-x}를 보라. 보조정리
\ref{algebra-lemma-dimension-spell-it-out}을 사용하면 $S_{\mathfrak m}$의 차원은
$x_0$을 지나는 $Z_i$들의 차원 중 최댓값임을 알 수 있다. 보조정리의
증명을 끝내려면 $x_0 \in Z_i \Rightarrow i \in I'$가 되도록 $x_0$을
택할 수 있음을 보이는 것으로 충분하다. 위에서 보았듯이 $S$는
Jacobson이므로 위의 $T$에 대해 $V(\mathfrak p) \setminus T$가 공집합이
아님을 보이는 것으로 충분하다. 이는 점 $x$(즉 $\mathfrak p$)를
포함하므로 명백하다.
\end{proof}

\begin{lemma}
\label{algebra-lemma-dimension-closed-point-finite-type-field}
$k$를 체, $S$를 유한형 $k$-대수, $X = \Spec(S)$라 하자.
$\mathfrak m \subset S$를 극대 아이디얼이라 하고 $x \in X$를 대응하는
닫힌점이라 하자. 그러면 $\dim_x(X) = \dim(S_{\mathfrak m})$이다.
\end{lemma}

\begin{proof}
이는 보조정리 \ref{algebra-lemma-dimension-at-a-point-finite-type-over-field}의
특수한 경우이다.
\end{proof}

\begin{lemma}
\label{algebra-lemma-disjoint-decomposition-CM-algebra}
$k$를 체라 하고 $S$를 유한형 $k$ 대수라 하자.
% P01-E0543
$S$가 Cohen--Macaulay라고 가정하자. 그러면
$\Spec(S) = \coprod T_d$는 열린닫힌 부분집합 $T_d$들의 유한 서로소
합이고, 각 $T_d$는 차원이 $d$인 등차원이다(위상수학, 정의
\ref{topology-definition-equidimensional}). 동치로, $S$는 환들
$S_d$, $d = 0, \ldots, \dim(S)$의 곱이고 $S_d$의 모든 극대 아이디얼
$\mathfrak m$의 높이는 $d$이다.
\end{lemma}

\begin{proof}
두 명제의 동치는 보조정리 \ref{algebra-lemma-disjoint-implies-product}에서 따른다.
$\mathfrak m \subset S$를 극대 아이디얼이라 하자. 보조정리
\ref{algebra-lemma-maximal-chain-CM}에 의해 $S_{\mathfrak m}$의 모든 극대 소
아이디얼 사슬은 같은 길이 $\dim(S_{\mathfrak m})$를 갖는다. 따라서
보조정리 \ref{algebra-lemma-dimension-spell-it-out}에 의해 $\mathfrak m$에
대응하는 점을 지나는 기약성분들의 차원은 모두
% P01-E0544
$\dim(S_{\mathfrak m})$와 같다. $\Spec(S)$가 Jacobson 위상공간이므로
그 두 기약성분의 교집합이 공집합이 아니면 닫힌점을 포함한다. 보조정리
\ref{algebra-lemma-finite-type-field-Jacobson}과 \ref{algebra-lemma-jacobson}을 보라.
따라서 서로 만나는 두 기약성분의 차원이 같음을 보였다. 이 사실과
$\Spec(S)$가 유한 개의 기약성분을 갖는다는 사실(보조정리
\ref{algebra-lemma-obvious-Noetherian}과 \ref{algebra-lemma-Noetherian-topology}을 보라)에서
보조정리가 쉽게 따른다.
\end{proof}
















\section{Noether 정규화}
\label{algebra-section-Noether-normalization}

\noindent
이 절에서는 Noether 정규화 보조정리의 여러 판본을 증명한다. 사용할
핵심 성분은 다음 두 보조정리에 담겨 있다.

\begin{lemma}
\label{algebra-lemma-helper}
$n \in \mathbf{N}$이라 하자. $N$을 다중지표
$\nu = (\nu_1, \ldots, \nu_n)$들의 유한하고 공집합이 아닌 집합이라
하자. $e = (e_1, \ldots, e_n)$이 주어지면
$e \cdot \nu = \sum e_i\nu_i$로 두자. 그러면
$e_1 \gg e_2 \gg \ldots \gg e_{n-1} \gg e_n$일 때 다음이 성립한다.
$\nu, \nu' \in N$이면
$$
(e \cdot \nu = e \cdot \nu')
\Leftrightarrow
(\nu = \nu')
$$
이다.
\end{lemma}

\begin{proof}
$N = \{\nu_j\}$이고 $\nu_j = (\nu_{j1}, \ldots, \nu_{jn})$이라고 하자.
$A_i = \max_j \nu_{ji} - \min_j \nu_{ji}$로 두자.
% P01-E0545
각 $i$에 대해
$e_{i - 1} > A_ie_i + A_{i + 1}e_{i + 1} + \ldots + A_ne_n$이면
보조정리가 성립한다. 실제로 $e \cdot (\nu - \nu') = 0$이라고 하자.
$n \ge 2$이면
$$
e_1(\nu_1 - \nu'_1) = \sum\nolimits_{i = 2}^n e_i(\nu'_i - \nu_i).
$$
이다. $(\nu_1 - \nu'_1) \ge 0$이라고 가정해도 된다. 만일
$(\nu_1 - \nu'_1) > 0$이면
$$
e_1(\nu_1 - \nu'_1) \ge e_1 >
A_2e_2 + \ldots + A_ne_n \ge
\sum\nolimits_{i = 2}^n e_i|\nu'_i - \nu_i| \ge
\sum\nolimits_{i = 2}^n e_i(\nu'_i - \nu_i).
$$
이다. 이 모순으로부터 $\nu'_1 = \nu_1$을 얻는다. 귀납법에 의해
$2 \le i \le n$인 모든 지표에 대해 $\nu'_i = \nu_i$이다.
\end{proof}

\begin{lemma}
\label{algebra-lemma-helper-polynomial}
$R$을 환이라 하고 $g \in R[x_1, \ldots, x_n]$을 상수가 아닌 원소,
즉 $g \not \in R$인 원소라 하자.
$e_1 \gg e_2 \gg \ldots \gg e_{n-1} \gg e_n = 1$이면 다항식
$$
g(x_1 + x_n^{e_1}, x_2 + x_n^{e_2}, \ldots, x_{n - 1} + x_n^{e_{n - 1}}, x_n)
=
ax_n^d + \text{lower order terms in }x_n
$$
이고, 여기서 $d > 0$이며 $a \in R$은 $g$의 영이 아닌 계수 중 하나이다.
\end{lemma}

\begin{proof}
$g = \sum_{\nu \in N} a_\nu x^\nu$로 쓰자. 여기서
% P01-E0546
$a_\nu \in R$은 영이 아니며, $N$은 보조정리 \ref{algebra-lemma-helper}에서와
같은 다중지표들의 유한 집합이고
$x^\nu = x_1^{\nu_1} \ldots x_n^{\nu_n}$이다. 다음 식의 최고차항은
$$
(x_1 + x_n^{e_1})^{\nu_1} \ldots (x_{n-1} + x_n^{e_{n-1}})^{\nu_{n-1}}
x_n^{\nu_n}
\quad\text{is}\quad
x_n^{e_1\nu_1 + \ldots + e_{n-1}\nu_{n-1} + \nu_n}.
$$
이다. 따라서 보조정리 \ref{algebra-lemma-helper}에서 결론이 따른다. 그 보조정리는
$g$의 영이 아닌 항 $a_\nu x^\nu$ 중 정확히 하나만이
$g(x_1 + x_n^{e_1}, x_2 + x_n^{e_2}, \ldots, x_{n - 1} + x_n^{e_{n - 1}},
x_n)$의 최고차항을 준다는 것을 보장한다. 즉 $e \cdot \nu$가 최대가 되는
유일한 $\nu \in N$에 대해 $a = a_\nu$이다.
\end{proof}

\begin{lemma}
\label{algebra-lemma-one-relation}
$k$를 체라 하고 어떤 진아이디얼 $I$에 대해
$S = k[x_1, \ldots, x_n]/I$라고 하자. $I \not = 0$이면
$S$가 $k[y_1, \ldots, y_{n-1}]$ 위에서 유한이 되도록 하는
$y_1, \ldots, y_{n-1} \in k[x_1, \ldots, x_n]$이 존재한다. 더욱이
$y_i$들을 $x_1, \ldots, x_n$이 생성하는
$k[x_1, \ldots, x_n]$의 $\mathbf{Z}$-부분대수 안에서 택할 수 있다.
\end{lemma}

\begin{proof}
$f \in I$, $f\not = 0$을 택하자. $S$가
$k[x_1, \ldots, x_n]/(f)$의 몫환이므로 이 환에 대해 보조정리를 보이는
것으로 충분하다. 적당한 정수 $e_i$에 대해
$y_i = x_i - x_n^{e_i}$, $i = 1, \ldots, n-1$로 택할 것이다. 언제 이것이
성립하는가? $\overline{x_n} \in k[x_1, \ldots, x_n]/(f)$가 환
$k[y_1, \ldots, y_{n-1}]$ 위에서 정수적임을 보이는 것으로 충분하다.
이 환 위에서 $\overline{x_n}$이 만족하는 방정식은
$$
f(y_1 + x_n^{e_1}, \ldots, y_{n-1} + x_n^{e_{n-1}}, x_n) = 0.
$$
이다. 따라서 위 방정식을 $x_n$에 관한 방정식으로 보았을 때 최고차항의
계수가 $k$의 영이 아닌 원소가 되게 하는 정수 $e_i$들이 존재함을 보이면
증명이 끝난다. 이는 예를 들어
% P01-E0547
$e_1 \gg e_2 \gg \ldots \gg e_{n-1}$로 택하여 얻을 수 있다. 보조정리
\ref{algebra-lemma-helper-polynomial}을 보라.
\end{proof}

\begin{lemma}
\label{algebra-lemma-Noether-normalization}
\begin{slogan}
Noether 정규화
\end{slogan}
$k$를 체라 하고 어떤 아이디얼 $I$에 대해
$S = k[x_1, \ldots, x_n]/I$라고 하자. $I \neq (1)$이면 다음 조건을
만족하는 $r\geq 0$과
$y_1, \ldots, y_r \in k[x_1, \ldots, x_n]$이 존재한다.
(a) $k[y_1, \ldots, y_r] \to S$는 단사이며, 여기서 출발환은
$y_1, \ldots, y_r$에 대한 다항식환이다. (b) 사상
$k[y_1, \ldots, y_r] \to S$는 유한이다. 이 경우 정수 $r$은 $S$의
차원이다. 더욱이 $y_i$들을 $x_1, \ldots, x_n$이 생성하는
$k[x_1, \ldots, x_n]$의 $\mathbf{Z}$-부분대수 안에서 택할 수 있다.
\end{lemma}

\begin{proof}
$n$에 대한 귀납법으로 증명한다. $n = 0$인 경우는 자명하다.
$I = 0$이면 $r = n$과 $y_i = x_i$로 택한다. $I \not = 0$이면 보조정리
\ref{algebra-lemma-one-relation}에서와 같이 $y_1, \ldots, y_{n-1}$을 택한다.
$S' \subset S$를 $y_i$들의 상이 생성하는 부분환이라 하자. 귀납법에 의해
$k[z_1, \ldots, z_r] \to S'$에 대해 (a), (b)가 성립하도록 하는
$r$과 $z_1, \ldots, z_r \in k[y_1, \ldots, y_{n-1}]$을 택할 수 있다.
$S' \to S$가 단사이고 유한이므로 $k[z_1, \ldots, z_r] \to S$에
대해서도 (a), (b)가 성립함을 알 수 있다.
% P01-E0548
마지막 주장은 보조정리 \ref{algebra-lemma-integral-sub-dim-equal}에서 따른다.
\end{proof}

\begin{lemma}
\label{algebra-lemma-Noether-normalization-at-point}
$k$를 체, $S$를 유한형 $k$ 대수라 하고 $X = \Spec(S)$로 쓰자.
$\mathfrak q$를 $S$의 소 아이디얼, $x \in X$를 대응하는 점이라 하자.
$\dim(S_g) = \dim_x(X) =: d$이고 유한 단사사상
$k[y_1, \ldots, y_d] \to S_g$가 존재하도록 하는
$g \in S$, $g \not \in \mathfrak q$가 존재한다.
\end{lemma}

\begin{proof}
정의에 의해 $\dim_x(X)$는 $g \in S$, $g \not \in \mathfrak q$에 대한
$S_g$들의 차원 중 최솟값, 즉 실제로 달성되는 최솟값임에 주의하자.
따라서 보조정리 \ref{algebra-lemma-Noether-normalization}에서 결론이 따른다.
\end{proof}

\begin{lemma}
\label{algebra-lemma-refined-Noether-normalization}
$k$를 체라 하고
$\mathfrak q \subset k[x_1, \ldots, x_n]$를 소 아이디얼이라 하자.
$r = \text{trdeg}_k\ \kappa(\mathfrak q)$로 두자. 그러면
$\varphi^{-1}(\mathfrak q) = (y_{r + 1}, \ldots, y_n)$을 만족하는 유한
환 준동형사상
$\varphi : k[y_1, \ldots, y_n] \to k[x_1, \ldots, x_n]$이 존재한다.
\end{lemma}

\begin{proof}
$n$에 대한 귀납법으로 증명한다. $n = 0$인 경우는 명백하다.
$n > 0$이라고 가정하자. $r = n$이면 $\mathfrak q = (0)$이고 결과는
명백하다. 영이 아닌 $f \in \mathfrak q$를 택하자. 물론 $f$는 상수가
아니다. 다음 꼴의 자기동형사상
$$
k[x_1, \ldots, x_n] \longrightarrow k[x_1, \ldots, x_n],
\quad
x_n \mapsto x_n,
\quad
x_i \mapsto x_i + x_n^{e_i}\ (i < n)
$$
을 적용한 뒤 $f$가 $k[x_1, \ldots, x_n]$ 위에서 $x_n$에 관한
일계수 다항식이라고 가정해도 된다. 보조정리
\ref{algebra-lemma-helper-polynomial}을 보라. 따라서 환 준동형사상
$$
k[y_1, \ldots, y_n] \longrightarrow k[x_1, \ldots, x_n],
\quad
y_n \mapsto f,
\quad
y_i \mapsto x_i\ (i < n)
$$
은 유한이다. 또한 구성에 의해
$y_n \in \mathfrak q \cap k[y_1, \ldots, y_n]$이다. 따라서
$\mathfrak q \cap k[y_1, \ldots, y_n] = \mathfrak pk[y_1, \ldots, y_n] + (y_n)$
이고, 여기서 $\mathfrak p \subset k[y_1, \ldots, y_{n - 1}]$는 소
아이디얼이다. $\kappa(\mathfrak p) \subset \kappa(\mathfrak q)$가
유한이므로 $r = \text{trdeg}_k\ \kappa(\mathfrak p)$임에 주의하자.
쌍 $(k[y_1, \ldots, y_{n - 1}], \mathfrak p)$에 귀납 가정을 적용하면
$\mathfrak p \cap k[z_1, \ldots, z_{n - 1}] =
(z_{r + 1}, \ldots, z_{n - 1})$을 만족하는 유한 환 준동형사상
$k[z_1, \ldots, z_{n - 1}] \to k[y_1, \ldots, y_{n - 1}]$을 얻는다.
환 준동형사상
$k[z_1, \ldots, z_{n - 1}] \to k[y_1, \ldots, y_{n - 1}]$을
$z_n$을 $y_n$으로 사상시켜 환 준동형사상
$k[z_1, \ldots, z_n] \to k[y_1, \ldots, y_n]$으로 확장한다. 환
준동형사상들의 합성
$$
k[z_1, \ldots, z_n] \to k[y_1, \ldots, y_n] \to k[x_1, \ldots, x_n]
$$
이 문제를 푼다.
\end{proof}

\begin{lemma}
\label{algebra-lemma-Noether-normalization-over-a-domain}
$R \to S$를 단사인 유한형 환 준동형사상이라 하고 $R$이 정역이라고
가정하자. 그러면 단사사상들에 의한 분해
$$
R \to R[y_1, \ldots, y_d] \to S' \to S
$$
가 존재하도록 하는 정수 $d$가 존재한다. 여기서 $S'$은
$R[y_1, \ldots, y_d]$ 위에서 유한이고, 어떤 영이 아닌 $f \in R$에
대해 $S'_f \cong S_f$이다.
\end{lemma}

\begin{proof}
$x_1, \ldots, x_n \in S$를 $R$ 위에서 $S$를 생성하도록 택하자.
$K$를 $R$의 분수체라 하고 $S_K = S \otimes_R K$로 두자. 보조정리
\ref{algebra-lemma-Noether-normalization}에 의해
$K[y_1, \ldots, y_d] \to S_K$가 유한 단사사상이 되도록 하는
$y_1, \ldots, y_d \in S$를 찾을 수 있다. $y_j$들을
$x_1, \ldots, x_n$이 생성하는 $\mathbf{Z}$-대수 안에서 택할 수 있으므로
$y_i \in S$임에 주의하자. 유한 환 준동형사상은 정수적이므로(보조정리
\ref{algebra-lemma-finite-is-integral}을 보라), $S_K$에서 $P_i(x_i) = 0$을
만족하는 일계수 다항식
$P_i \in K[y_1, \ldots, y_d][T]$를 찾을 수 있다. 모든 $i$에 대해
$fP_i \in R[y_1, \ldots, y_d][T]$가 되도록 하는 영이 아닌
$f \in R$을 택하자. 그러면 $fP_i(x_i)$는 $S_K$에서 영으로 사상된다.
따라서 $f$를 $R$의 다른 영이 아닌 원소로 바꾸어 $fP_i(x_i)$가
$S$에서도 영이라고 가정해도 된다. $x_i' = fx_i$로 두고
$S' \subset S$를 $y_1, \ldots, y_d$와 $x'_1, \ldots, x'_n$이 생성하는
$R$-부분대수라 하자. $Q_i(x_i') = 0$이고
$Q_i = f^{\deg_T(P_i)}P_i(T/f)$는 $f$를 택한 방식에 의해
$R[y_1, \ldots, y_d]$에 계수를 갖는 $T$의 일계수 다항식이므로
$x'_i$는 $R[y_1, \ldots, y_d]$ 위에서 정수적임에 주의하자.
따라서 보조정리 \ref{algebra-lemma-characterize-finite-in-terms-of-integral}에 의해
$R[y_1, \ldots, y_d] \subset S'$은 유한이다. $S' \subset S$이므로
$S'_f \subset S_f$이다(국소화는 완전이다). 다른 한편 $S'_f$의 원소
$x_i = x'_i/f$들은 $R_f$ 위에서 $S_f$를 생성하므로
$S'_f \to S_f$는 전사이다. 따라서 $S'_f \cong S_f$이고 증명이 끝난다.
\end{proof}







\section{체 위 유한형 대수의 차원, 다시 보기}
\label{algebra-section-dimension-finite-type-algebras-reprise}

\noindent
이 절은 절 \ref{algebra-section-dimension-finite-type-algebras}의 연속이다.
이 절에서는 체 위 유한형 정역에 대해 차원과 기초체 위의 초월차수 사이의
관계를 확립한다.

\begin{lemma}
\label{algebra-lemma-dimension-prime-polynomial-ring}
$k$를 체라 하자.
$S$를 정역인 유한형 $k$-대수라 하자.
$K$를 $S$의 분수체라 하자.
$r = \text{trdeg}(K/k)$를 $k$ 위에서 $K$의 초월차수라 하자.
그러면 $\dim(S) = r$이다. 더욱이 $S$의 각 극대 아이디얼에서의
국소환의 차원은 $r$이다.
\end{lemma}

\begin{proof}
$S = k[x_1, \ldots, x_n]/\mathfrak p$로 쓸 수 있다.
보조정리 \ref{algebra-lemma-dimension-height-polynomial-ring}에 의해
$S$의 극대 아이디얼에서의 모든 국소환은 같은 차원을 갖는다.
보조정리 \ref{algebra-lemma-Noether-normalization}을 적용하자. 그러면
$$
k[y_1, \ldots, y_d] \to S
$$
가 유한 단사 환 준동형사상이 되도록 할 수 있으며, 여기서
$d = \dim(S)$이다. 명백히 $k(y_1, \ldots, y_d) \subset K$는
유한 확대이므로 결론이 따른다.
\end{proof}

\begin{lemma}
\label{algebra-lemma-tr-deg-specialization}
$k$를 체라 하자. $S$를 유한형 $k$-대수라 하자.
$\mathfrak q \subset \mathfrak q' \subset S$를 서로 다른 소 아이디얼이라
하자. 그러면
$\text{trdeg}_k\ \kappa(\mathfrak q') < \text{trdeg}_k\ \kappa(\mathfrak q)$
이다.
\end{lemma}

\begin{proof}
보조정리 \ref{algebra-lemma-dimension-prime-polynomial-ring}에 의해
$\dim V(\mathfrak q) = \text{trdeg}_k\ \kappa(\mathfrak q)$이고
$\mathfrak q'$에 대해서도 마찬가지이다. 따라서 진포함
$V(\mathfrak q') \subset V(\mathfrak q)$이 차원의 엄격한 부등식을
함의하므로 결과가 따른다.
\end{proof}

\noindent
다음 보조정리는 보조정리
\ref{algebra-lemma-dimension-closed-point-finite-type-field}를 일반화한다.

\begin{lemma}
\label{algebra-lemma-dimension-at-a-point-finite-type-field}
$k$를 체라 하자.
$S$를 유한형 $k$-대수라 하자.
$X = \Spec(S)$로 두자.
$\mathfrak p \subset S$를 소 아이디얼이라 하고,
$x \in X$를 이에 대응하는 점이라 하자.
그러면
$$
\dim_x(X) = \dim(S_{\mathfrak p}) + \text{trdeg}_k\ \kappa(\mathfrak p).
$$
이다.
\end{lemma}

\begin{proof}
보조정리 \ref{algebra-lemma-dimension-prime-polynomial-ring}에 의해
$r = \text{trdeg}_k\ \kappa(\mathfrak p)$가
$V(\mathfrak p)$의 차원과 같음을 안다.
$S$에서 $\mathfrak p$로 시작하는 소 아이디얼들의 임의의 극대 사슬
$\mathfrak p \subset \mathfrak p_1 \subset \ldots \subset \mathfrak p_r$를
택하자. 보조정리 \ref{algebra-lemma-dimension-spell-it-out}에 의해 이 사슬의
길이는 $r$이다. $\mathfrak q_j$, $j \in J$를 $\mathfrak p$에 포함되는
$S$의 극소 소 아이디얼들이라 하자. 이들은 대응
$\mathfrak q_j \mapsto \mathfrak q_jS_{\mathfrak p}$를 통해
$S_{\mathfrak p}$의 극소 소 아이디얼들과 $1-1$로 대응한다.
보조정리 \ref{algebra-lemma-dimension-at-a-point-finite-type-over-field}에 의해
$\dim_x(X)$는 환 $S/\mathfrak q_j$들의 차원의 최댓값과 같다.
각 $j$에 대해 소 아이디얼들의 극대 사슬
$\mathfrak q_j \subset \mathfrak p'_1 \subset \ldots \subset
\mathfrak p'_{s(j)} = \mathfrak p$를 택하자. 그러면
$\dim(S_{\mathfrak p}) = \max_{j \in J} s(j)$이다.
이제 각 사슬
$$
\mathfrak q_j \subset \mathfrak p'_1 \subset \ldots \subset
\mathfrak p'_{s(j)} = \mathfrak p \subset
\mathfrak p_1 \subset \ldots \subset \mathfrak p_r
$$
은 $S/\mathfrak q_j$에서 극대 사슬이고, 앞에서 말한 바에 의해
$\dim_x(X) = \max_{j \in J} r + s(j)$이다.
따라서 보조정리가 따른다.
\end{proof}

\noindent
다음 보조정리는 한 유한형 스펙트럼의 다른 유한형 스펙트럼 안에서의
여차원이 높이들의 차임을 말한다.

\begin{lemma}
\label{algebra-lemma-codimension}
$k$를 체라 하자.
$S' \to S$를 유한형 $k$-대수들의 전사사상이라 하자.
$\mathfrak p \subset S$를 소 아이디얼이라 하고,
$\mathfrak p'$를 이에 대응하는 $S'$의 소 아이디얼이라 하자.
$X = \Spec(S)$ 및 $X' = \Spec(S')$로 두고,
$x \in X$ 및 $x'\in X'$를 각각 $\mathfrak p$ 및 $\mathfrak p'$에
대응하는 점이라 하자. 그러면
$$
\dim_{x'} X' - \dim_x X =
\text{height}(\mathfrak p') - \text{height}(\mathfrak p).
$$
이다.
\end{lemma}

\begin{proof}
보조정리 \ref{algebra-lemma-dimension-at-a-point-finite-type-field}에서 즉시 따른다.
\end{proof}

\begin{lemma}
\label{algebra-lemma-dimension-preserved-field-extension}
$k$를 체라 하자.
$S$를 유한형 $k$-대수라 하자.
$K/k$를 체 확대라 하자.
그러면 $\dim(S) = \dim(K \otimes_k S)$이다.
\end{lemma}

\begin{proof}
보조정리 \ref{algebra-lemma-Noether-normalization}에 의해
$d = \dim(S)$이고 유한 단사사상
$k[y_1, \ldots, y_d] \to S$가 존재한다.
$K$는 $k$ 위에서 평탄이므로 유한 단사사상
$K[y_1, \ldots, y_d] \to K \otimes_k S$도 얻는다.
결과는 보조정리 \ref{algebra-lemma-integral-sub-dim-equal}에서 따른다.
\end{proof}

\begin{lemma}
\label{algebra-lemma-dimension-at-a-point-preserved-field-extension}
$k$를 체라 하자.
$S$를 유한형 $k$-대수라 하자.
$X = \Spec(S)$로 두자.
$K/k$를 체 확대라 하자.
$S_K = K \otimes_k S$ 및 $X_K = \Spec(S_K)$로 두자.
$\mathfrak q \subset S$를 $x \in X$에 대응하는 소 아이디얼이라 하고,
$\mathfrak q_K \subset S_K$를 $\mathfrak q$ 위에 놓이며
$x_K \in X_K$에 대응하는 소 아이디얼이라 하자.
그러면 $\dim_x X = \dim_{x_K} X_K$이다.
\end{lemma}

\begin{proof}
$S = k[x_1, \ldots, x_n]/I$인 표시를 택하자.
이는 표시
$K \otimes_k S = K[x_1, \ldots, x_n]/(K \otimes_k I)$를 준다.
$\mathfrak q_K' \subset K[x_1, \ldots, x_n]$ 및
$\mathfrak q' \subset k[x_1, \ldots, x_n]$를 각각 대응하는
소 아이디얼이라 하자. 다음의 뇌터 국소환들의 가환 그림을 생각하자.
$$
\xymatrix{
K[x_1, \ldots, x_n]_{\mathfrak q_K'} \ar[r] &
(K \otimes_k S)_{\mathfrak q_K} \\
k[x_1, \ldots, x_n]_{\mathfrak q'} \ar[r] \ar[u] &
S_{\mathfrak q} \ar[u]
}
$$
두 수직 화살표는 평탄 환 준동형사상
$S \to S_K$ 및
$k[x_1, \ldots, x_n] \to K[x_1, \ldots, x_n]$의 국소화이므로
평탄이다. 더욱이 두 수직 화살표의 올환은 서로 같다. 따라서 보조정리
\ref{algebra-lemma-dimension-base-fibre-equals-total}에 의해
$\text{height}(\mathfrak q') - \text{height}(\mathfrak q)
= \text{height}(\mathfrak q_K') - \text{height}(\mathfrak q_K)$임을 알 수
있다.
% P01-E1408
$x' \in X' = \Spec(k[x_1, \ldots, x_n])$ 및
$x'_K \in X'_K = \Spec(K[x_1, \ldots, x_n])$를 각각
$\mathfrak q'$ 및 $\mathfrak q_K'$에 대응하는 점이라 하자.
보조정리 \ref{algebra-lemma-codimension}와 위에서 보인 바에 의해
\begin{eqnarray*}
n - \dim_x X & = & \dim_{x'} X' - \dim_x X \\
& = & \text{height}(\mathfrak q') - \text{height}(\mathfrak q) \\
& = & \text{height}(\mathfrak q_K') - \text{height}(\mathfrak q_K) \\
& = & \dim_{x'_K} X'_K - \dim_{x_K} X_K \\
& = & n - \dim_{x_K} X_K
\end{eqnarray*}
이고 보조정리가 따른다.
\end{proof}

\begin{lemma}
\label{algebra-lemma-inequalities-under-field-extension}
$k$를 체라 하자. $S$를 유한형 $k$-대수라 하자.
$K/k$를 체 확대라 하자. $S_K = K \otimes_k S$로 두자.
$\mathfrak q \subset S$를 소 아이디얼이라 하고 $\mathfrak q_K \subset S_K$를
$\mathfrak q$ 위에 놓이는 소 아이디얼이라 하자. 그러면
$$
\dim (S_K \otimes_S \kappa(\mathfrak q))_{\mathfrak q_K} =
\dim (S_K)_{\mathfrak q_K} - \dim S_\mathfrak q =
\text{trdeg}_k \kappa(\mathfrak q) - \text{trdeg}_K \kappa(\mathfrak q_K)
$$
이다. 더욱이 주어진 $\mathfrak q$에 대해 위 수가 영이 되도록
$\mathfrak q_K$를 항상 택할 수 있다.
\end{lemma}

\begin{proof}
$S_\mathfrak q \to (S_K)_{\mathfrak q_K}$는 특수올이
$(S_K \otimes_S \kappa(\mathfrak q))_{\mathfrak q_K}$인
뇌터 국소환들의 평탄 국소 준동형사상이다. 따라서 첫째 등식은 보조정리
\ref{algebra-lemma-dimension-base-fibre-equals-total}에서 따른다.
둘째 등식은 보조정리
\ref{algebra-lemma-dimension-at-a-point-preserved-field-extension}의 표기를 사용하면
$\dim_x X = \dim_{x_K} X_K$라는 사실과, 보조정리
\ref{algebra-lemma-dimension-at-a-point-finite-type-field}에 의해
$\dim_x X = \dim S_\mathfrak q + \text{trdeg}_k \kappa(\mathfrak q)$이고
$\dim_{x_K} X_K$에 대해서도 마찬가지라는 사실에서 따른다.
$\mathfrak q_K$를 $\mathfrak q S_K$ 위의 극소 소 아이디얼로 택하면
올환의 차원은 영이다.
\end{proof}







\section{체 위 등급 대수의 차원}
\label{algebra-section-dimension-graded}

\noindent
다음은 기본적인 결과이다.

\begin{lemma}
\label{algebra-lemma-dimension-graded}
$k$를 체라 하자. $S$를 차수 $1$인 유한 개의 원소로 $k$ 위에서
생성되는 등급 $k$-대수라 하자.
$S_0 = k$라고 가정하자. 모든 $d \gg 0$에 대해
$\dim(S_d) = P(d)$가 되도록 하는 다항식
$P(T) \in \mathbf{Q}[T]$를 택하자. 명제
\ref{algebra-proposition-graded-hilbert-polynomial}을 보라. 그러면
\begin{enumerate}
\item 무관 아이디얼 $S_{+}$는 극대 아이디얼 $\mathfrak m$이다.
\item $S$의 임의의 극소 소 아이디얼은 동차 아이디얼이고
$S_{+} = \mathfrak m$에 포함된다.
\item $\dim(S) = \deg(P) + 1 = \dim_x\Spec(S)$이다
($\deg(0) = -1$이라는 관례를 쓴다). 여기서 $x$는 극대 아이디얼
$S_{+} = \mathfrak m$에 대응하는 점이다.
\item 국소환 $R = S_{\mathfrak m}$의 힐베르트 함수는 $S$의
힐베르트 함수와 같다.
\end{enumerate}
\end{lemma}

\begin{proof}
첫째 명제는 명백하다.
둘째 명제는 보조정리 \ref{algebra-lemma-graded-ring-minimal-prime}에서 따른다.
(2)에 의해 모든 기약 성분은 $x$를 지난다.
따라서 보조정리
\ref{algebra-lemma-dimension-at-a-point-finite-type-over-field}에 의해
$\dim(S) = \dim_x\Spec(S) = \dim(S_\mathfrak m)$이다.
$\mathfrak m^d/\mathfrak m^{d + 1} \cong
\mathfrak m^dS_\mathfrak m/\mathfrak m^{d + 1}S_\mathfrak m$이므로
국소환 $S_\mathfrak m$의 힐베르트 함수가 $S$의 힐베르트 함수와
같음을 알 수 있는데, 이것이 (4)이다. 명제 \ref{algebra-proposition-dimension}에
의해 (3)의 마지막 등식을 얻는다.
\end{proof}











\section{일반적 평탄성}
\label{algebra-section-generic-flatness}

\noindent
요지는 정역 위의 유한형 대수가 그 정역의 원소 하나를 가역화하면
평탄해진다는 것이다. 이 결과에는 강도가 점점 높아지는 여러 판본이 있다.

\begin{lemma}
\label{algebra-lemma-generic-flatness-Noetherian}
$R \to S$를 환 사상이라 하자.
$M$을 $S$-가군이라 하자.
다음을 가정하자.
\begin{enumerate}
\item $R$은 뇌터이다.
\item $R$은 정역이다.
\item $R \to S$는 유한형이다.
\item $M$은 유한형 $S$-가군이다.
\end{enumerate}
그러면 $M_f$가 자유 $R_f$-가군이 되도록 하는 영이 아닌
$f \in R$이 존재한다.
\end{lemma}

\begin{proof}
$K$를 $R$의 분수체라 하고 $S_K = K \otimes_R S$로 두자. 이는
$K$ 위의 유한형 대수이다. $d = \dim(S_K)$에 대한 귀납법을 쓴다
(이 수가 유한인 것은 예를 들어 뇌터 정규화에서 따르며, 절
\ref{algebra-section-Noether-normalization}을 보라).
$d \geq 0$을 고정하자.
$\dim(S_K) < d$인 모든 경우에 보조정리가 성립한다고 가정하자.

\medskip\noindent
보조정리의 조건을 만족하고 $\dim(S_K) = d$인 $R \to S$와 $M$이
주어졌다고 하자. 보조정리 \ref{algebra-lemma-filter-Noetherian-module}에 의해
$M_i/M_{i - 1}$가 $S$의 어떤 소 아이디얼 $\mathfrak q$에 대해
$S/\mathfrak q$와 동형이 되도록 하는 여과
$0 \subset M_1 \subset M_2 \subset \ldots \subset M_n = M$이 존재한다.
$\dim((S/\mathfrak q)_K) \leq \dim(S_K)$임에 주의하자. 또한 자유
가군들의 확대는 자유이다(기본 개념 \ref{algebra-item-extension-free}을 보라).
따라서 $M = S$이고 $S$가 $R$ 위의 유한형 정역이라고 가정해도 된다.

\medskip\noindent
$R \to S$의 핵이 영이 아니면, 그 핵에서 영이 아닌 $f \in R$을
택하자. 이 경우 $S_f = 0$이므로 보조정리가 성립한다. (이는 실제로
$d = -1$인 경우이며 귀납법의 시작이다.) 따라서 $R \to S$가
뇌터 정역들의 유한형 확대라고 가정해도 된다.

\medskip\noindent
보조정리 \ref{algebra-lemma-Noether-normalization-over-a-domain}을 적용하고
$R$을 $R_f$로 바꾸면(여기서 $f$는 그 보조정리의 원소이다) 둘째 확대가
유한인 분해
$$
R \subset R[y_1, \ldots, y_d] \subset S
$$
를 얻는다. $S$의 분수체가 $R[y_1, \ldots, y_d]$의 분수체 위에서 갖는
기저를 이루는 $z_1, \ldots, z_r \in S$를 택하자. 이는 짧은 완전열
$$
0 \to
R[y_1, \ldots, y_d]^{\oplus r} \xrightarrow{(z_1, \ldots, z_r)}
S \to N \to 0
$$
을 준다. 구성에 의해 $N$은 유한 $R[y_1, \ldots, y_d]$-가군이고 그
지지집합은 $\Spec(R[y_1, \ldots, y_d])$의 일반점 $(0)$을 포함하지 않는다.
보조정리 \ref{algebra-lemma-support-closed}에 의해 영이 아닌
$g \in R[y_1, \ldots, y_d]$가 존재하고, 이 $g$가 $N$을 영멸한다.
따라서 $N$을
$S' = R[y_1, \ldots, y_d]/(g)$ 위의 유한 가군으로 볼 수 있다.
$\dim(S'_K) < d$이므로 귀납 가정에 의해 $N_f$가 자유
$R_f$-가군이 되도록 하는 영이 아닌 $f \in R$이 존재한다.
또한 $(R[y_1, \ldots, y_d])_f \cong R_f[y_1, \ldots, y_d]$가 자유이므로,
앞에서 말한 자유 가군들의 확대가 자유라는 사실로부터 결론이 따른다.
\end{proof}

\begin{lemma}
\label{algebra-lemma-generic-flatness-finitely-presented}
\begin{slogan}
일반적 자유성.
\end{slogan}
$R \to S$를 환 사상이라 하자.
$M$을 $S$-가군이라 하자.
다음을 가정하자.
\begin{enumerate}
\item $R$은 정역이다.
\item $R \to S$는 유한 표시이다.
\item $M$은 유한 표시 $S$-가군이다.
\end{enumerate}
그러면 $M_f$가 자유 $R_f$-가군이 되도록 하는 영이 아닌
$f \in R$이 존재한다.
\end{lemma}

\begin{proof}
$S = R[x_1, \ldots, x_n]/(g_1, \ldots, g_m)$로 쓰자.
$g \in R[x_1, \ldots, x_n]$에 대해 $\overline{g}$로 그 $S$에서의 상을
나타내자. 어떤 $n_i \in S^{\oplus t}$들에 대해
$M = S^{\oplus t}/\sum Sn_i$로 쓸 수 있다. 어떤
$g_{ij} \in R[x_1, \ldots, x_n]$들에 대해
$n_i = (\overline{g}_{i1}, \ldots, \overline{g}_{it})$로 쓰자.
$R_0 \subset R$를 모든 원소
$g_i, g_{ij} \in R[x_1, \ldots, x_n]$의 모든 계수로 생성되는
부분환이라 하자.
$S_0 = R_0[x_1, \ldots, x_n]/(g_1, \ldots, g_m)$으로 정의하고
$M_0 = S_0^{\oplus t}/\sum S_0n_i$로 정의하자.
그러면 $R_0$은 $\mathbf{Z}$ 위의 유한형 정역이므로 뇌터이다
(보조정리 \ref{algebra-lemma-Noetherian-permanence}을 보라).
더욱이 단사사상 $R_0 \to R$을 통해
$S \cong R \otimes_{R_0} S_0$이고
$M \cong R \otimes_{R_0} M_0$이다. 보조정리
\ref{algebra-lemma-generic-flatness-Noetherian}을 적용하면 $(M_0)_f$가 자유
$(R_0)_f$-가군이 되도록 하는 영이 아닌 $f \in R_0$을 얻는다.
따라서 $M_f = R_f \otimes_{(R_0)_f} (M_0)_f$는 자유 $R_f$-가군이다.
\end{proof}

\begin{lemma}
\label{algebra-lemma-generic-flatness}
$R \to S$를 환 사상이라 하자.
$M$을 $S$-가군이라 하자.
다음을 가정하자.
\begin{enumerate}
\item $R$은 정역이다.
\item $R \to S$는 유한형이다.
\item $M$은 유한형 $S$-가군이다.
\end{enumerate}
그러면 다음이 성립하도록 하는 영이 아닌 $f \in R$이 존재한다.
\begin{enumerate}
\item[(a)] $M_f$와 $S_f$는 자유 $R_f$-가군이다.
\item[(b)] $S_f$는 유한 표시 $R_f$-대수이고 $M_f$는 유한 표시
$S_f$-가군이다.
\end{enumerate}
\end{lemma}

\begin{proof}
먼저 $S = R[x_1, \ldots, x_n]$인 경우에 보조정리를 증명한 다음
일반적인 결과를 도출한다.

\medskip\noindent
$S = R[x_1, \ldots, x_n]$이라고 가정하자.
$M$을 생성하는 원소 $m_1, \ldots, m_t$를 택하자. 이는 짧은 완전열
$$
0 \to N \to S^{\oplus t} \xrightarrow{(m_1, \ldots, m_t)} M \to 0.
$$
을 준다. $K$로 $R$의 분수체를 나타내자.
$S_K = K \otimes_R S = K[x_1, \ldots, x_n]$로 나타내고, 마찬가지로
$N_K = K \otimes_R N$, $M_K = K \otimes_R M$으로 나타내자.
$R \to K$가 평탄이므로 $K$와 텐서한 뒤에도 열은 완전하다.
$S_K = K[x_1, \ldots, x_n]$가 뇌터환이므로(보조정리
\ref{algebra-lemma-Noetherian-permanence}을 보라) 이를 생성하는 유한 개의 원소
$n'_1, \ldots, n'_s \in N_K$를 찾을 수 있다. 어떤 $a_{ij} \in K$에
대해 $n'_i = \sum a_{ij}n_j$가 되도록 하는
$n_1, \ldots, n_r \in N$을 택하자. 다음과 같이 두자.
$$
M' = S^{\oplus t}/\sum\nolimits_{i = 1, \ldots, r} Sn_i
$$
구성에 의해 $M'$은 유한 표시 $S$-가군이고, 전사사상
$M' \to M$은 동형사상 $M'_K \cong M_K$를 유도한다.
$R \to S$와 $M'$에 보조정리
\ref{algebra-lemma-generic-flatness-finitely-presented}를 적용하면 $M'_f$가 자유
$R_f$-가군이 되도록 하는 $f \in R$을 얻는다. 따라서
$M'_f \to M_f$는 정역 $R_f$ 위 가군들의 전사사상이며 그 원천은 자유
가군이고, $K$와 텐서하면 동형사상이 된다. $M'_f$가 자유이고 $R_f$가
분수체 $K$를 갖는 정역이므로 $M'_f \subset M'_K$이다. 따라서 이 사상은
단사이다. 그러므로 $M'_f \to M_f$는 동형사상이고 결과가 증명되었다.

\medskip\noindent
일반적인 경우에는 전사사상 $R[x_1, \ldots, x_n] \to S$를 택하자.
$S$와 $M$을 모두 $R[x_1, \ldots, x_n]$ 위의 유한 가군으로 보자.
앞에서 증명한 특수한 경우에 의해 $S_f$와 $M_f$가 모두 자유
$R_f$-가군이고 유한 표시 $R_f[x_1, \ldots, x_n]$-가군이 되도록 하는
영이 아닌 $f \in R$이 존재한다. 이는 명백히 $S_f$가 유한 표시
$R_f$-대수이고 $M_f$가 유한 표시 $S_f$-가군임을 함의한다.
\end{proof}

\noindent
$R \to S$를 환 사상이라 하고 $M$을 $S$-가군이라 하자.
원소 $f \in R$에 대한 다음 조건을 생각하자.
\begin{equation}
\label{algebra-equation-flat-and-finitely-presented}
\left\{
\begin{matrix}
S_f & \text{is of finite presentation over }R_f\\
M_f & \text{is of finite presentation as }S_f\text{-module}\\
S_f, M_f & \text{are free as }R_f\text{-modules}
\end{matrix}
\right.
\end{equation}
다음과 같이 정의하자.
\begin{equation}
\label{algebra-equation-good-locus}
U(R \to S, M)
=
\bigcup\nolimits_{f \in
R\text{ with }(\ref{algebra-equation-flat-and-finitely-presented})}
D(f)
\end{equation}
이는 $\Spec(R)$의 열린 부분집합이다.

\begin{lemma}
\label{algebra-lemma-generic-flatness-locus-extension}
$R \to S$를 환 사상이라 하자.
$0 \to M_1 \to M_2 \to M_3 \to 0$을 $S$-가군들의 짧은 완전열이라 하자.
그러면
$$
U(R \to S, M_1) \cap U(R \to S, M_3) \subset U(R \to S, M_2).
$$
이다.
\end{lemma}

\begin{proof}
$u \in U(R \to S, M_1) \cap U(R \to S, M_3)$라 하자.
$u \in D(f_1)$ 및 $u \in D(f_3)$이고
(\ref{algebra-equation-flat-and-finitely-presented})가 각각 $f_1$과 $M_1$ 및
$f_3$과 $M_3$에 대해 성립하도록 하는 $f_1, f_3 \in R$을 택하자.
$f = f_1f_3$으로 두자. 그러면 $u \in D(f)$이고
(\ref{algebra-equation-flat-and-finitely-presented})는 $f$와 $M_1$, $M_3$에 대해
모두 성립한다. 자유 가군들의 확대는 자유이고,
유한 표시 가군들의 확대는 유한 표시이다(보조정리
\ref{algebra-lemma-extension}). 따라서
(\ref{algebra-equation-flat-and-finitely-presented})가 $f$와 $M_2$에 대해
성립함을 알 수 있다. 그러므로 $u \in U(R \to S, M_2)$이고 결론이 따른다.
\end{proof}

\begin{lemma}
\label{algebra-lemma-generic-flatness-locus-localize}
$R \to S$를 환 사상이라 하자.
$M$을 $S$-가군이라 하자.
$f \in R$이라 하자.
동일시 $\Spec(R_f) = D(f)$를 사용하면
$U(R_f \to S_f, M_f) = D(f) \cap U(R \to S, M)$이다.
\end{lemma}

\begin{proof}
$u \in U(R_f \to S_f, M_f)$라고 하자. 그러면 $u \in D(g)$이고
(\ref{algebra-equation-flat-and-finitely-presented})가 순서쌍
$((R_f)_g \to (S_f)_g, (M_f)_g)$에 대해 성립하도록 하는
$g \in R_f$가 존재한다. 어떤 $a \in R$에 대해 $g = a/f^n$으로 쓰고
$h = af$로 두자. 그러면 $R_h = (R_f)_g$, $S_h = (S_f)_g$ 및
$M_h = (M_f)_g$이다. 더욱이 $u \in D(h)$이다. 따라서
$u \in U(R \to S, M)$이다. 역으로
$u \in D(f) \cap U(R \to S, M)$라고 하자. 그러면 $u \in D(g)$이고
(\ref{algebra-equation-flat-and-finitely-presented})가 순서쌍
$(R_g \to S_g, M_g)$에 대해 성립하도록 하는 $g \in R$이 존재한다.
그러면 (\ref{algebra-equation-flat-and-finitely-presented})가 순서쌍
$(R_{fg} \to S_{fg}, M_{fg}) = ((R_f)_g \to (S_f)_g, (M_f)_g)$에
대해서도 성립함이 명백하다. 따라서 $u \in U(R_f \to S_f, M_f)$이고
결론이 따른다.
\end{proof}

\begin{lemma}
\label{algebra-lemma-generic-flatness-locus-reduce}
$R \to S$를 환 사상이라 하자.
$M$을 $S$-가군이라 하자.
$U \subset \Spec(R)$를 조밀한 열린집합이라 하자.
각 $i \in I$에 대해 $U(R_{f_i} \to S_{f_i}, M_{f_i})$가
$D(f_i)$에서 조밀하도록 하는 열린집합들의 덮개
$U = \bigcup_{i \in I} D(f_i)$가 존재한다고 가정하자.
그러면 $U(R \to S, M)$은 $\Spec(R)$에서 조밀하다.
\end{lemma}

\begin{proof}
보조정리 \ref{algebra-lemma-generic-flatness-locus-localize}에 비추어 보면
이는 순전히 위상적인 명제이다. 실제로 그 보조정리에 의해 각
$i \in I$에 대해 $U(R \to S, M) \cap D(f_i)$가 $D(f_i)$에서 조밀함을
알 수 있다. 위상수학, 보조정리
\ref{topology-lemma-nowhere-dense-local}에 의해
$U(R \to S, M) \cap U$는 $U$에서 조밀하다.
$U$가 $\Spec(R)$에서 조밀하므로 $U(R \to S, M)$도
$\Spec(R)$에서 조밀하다.
\end{proof}

\begin{lemma}
\label{algebra-lemma-generic-flatness-reduced}
$R \to S$를 환 사상이라 하고 $M$을 $S$-가군이라 하자.
다음을 가정하자.
\begin{enumerate}
\item $R \to S$는 유한형이다.
\item $M$은 유한 $S$-가군이다.
\item $R$은 기약환이다.
\end{enumerate}
그러면 다음을 만족하는 부분집합 $U \subset \Spec(R)$가 존재한다.
\begin{enumerate}
\item $U$는 $\Spec(R)$에서 열리고 조밀하다.
\item 모든 $u \in U$에 대해 $u \in D(f) \subset U$이고 다음을 만족하는
$f \in R$이 존재한다.
\begin{enumerate}
\item $M_f$와 $S_f$는 $R_f$ 위에서 자유이다.
\item $S_f$는 유한 표시 $R_f$-대수이다.
\item $M_f$는 유한 표시 $S_f$-가군이다.
\end{enumerate}
\end{enumerate}
\end{lemma}

\begin{proof}
이 보조정리는 식 (\ref{algebra-equation-good-locus})의 열린집합
$U(R \to S, M)$이 $\Spec(R)$에서 조밀하다는 명제와 동치임에 주의하자.
먼저 $S = R[x_1, \ldots, x_n]$인 경우에 보조정리를 증명한 다음
일반적인 결과를 도출한다.

\medskip\noindent
$S = R[x_1, \ldots, x_n]$이고 $M$이 $S$ 위의 임의의 유한 가군인 경우를
증명하자. 이 경우 $S_f = R_f[x_1, \ldots, x_n]$는 $R_f$ 위에서 자유이고
유한 표시이므로, $S$에 관한 조건은 고려할 필요가 없고 $M$에 관한
조건만 고려하면 된다. $n$에 대한 귀납법을 쓴다.

\medskip\noindent
어떤 아이디얼 $J_i \subset S$들에 대해
$M_i/M_{i - 1} \cong S/J_i$가 되도록 하는 유한 여과
$$
0 \subset M_1 \subset M_2 \subset \ldots \subset M_t = M
$$
이 존재한다. 보조정리 \ref{algebra-lemma-trivial-filter-finite-module}을 보라.
조밀한 열린집합들의 유한 교집합은 조밀한 열린집합이므로, 보조정리
\ref{algebra-lemma-generic-flatness-locus-extension}에 의해 각 가군
% P01-E1409
$R/J_i$에 대해 보조정리를 증명하면 충분함을 알 수 있다. 따라서
$S = R[x_1, \ldots, x_n]$의 어떤 아이디얼 $J$에 대해
$M = S/J$라고 가정해도 된다.

\medskip\noindent
$I \subset R$을 $J$의 원소들의 계수로 생성되는 아이디얼이라 하자.
$U_1 = \Spec(R) \setminus V(I)$로 두고
$$
U_2 = \Spec(R) \setminus \overline{U_1}.
$$
로 두자. 그러면 $U = U_1 \cup U_2$가 $\Spec(R)$에서 조밀함은 명백하다.
(a) $D(f) \subset U_1$이거나 (b) $D(f) \subset U_2$인 원소
$f \in R$을 택하자. 이러한 모든 $f$에 대해 순서쌍
$(R_f \to R_f[x_1, \ldots, x_n], M_f)$에 보조정리가 성립한다면,
보조정리 \ref{algebra-lemma-generic-flatness-locus-reduce}에 의해
$U(R \to S, M)$이 $\Spec(R)$에서 조밀함을 알 수 있다.
따라서 (a) $I = R$이거나 (b) $V(I) = \Spec(R)$이라고 가정해도 된다.

\medskip\noindent
(b)의 경우 $R$이 기약환이므로 실제로 $I = 0$이다! 따라서
$J = 0$이고 $M = S$이므로 이 경우 보조정리가 성립한다.

\medskip\noindent
(a)의 경우에는 조금 더 작업해야 한다. $J$의 원소들의 계수들의 집합이
아이디얼을 이루므로(세부 사항은 생략한다), $I$의 모든 원소는 실제로
$J$의 어떤 원소의 어떤 단항식의 계수임에 주의하자. 따라서
$$
g = \sum\nolimits_{K \in E} a_K x^K \in J
$$
이고 적어도 하나의 계수 $a_{K_0}$가 $R$의 단원인 원소를 찾을 수 있다.
여기서 $E$는 다중 첨자 $K = (k_1, \ldots, k_n)$들의 유한 집합이다.
실제로 (a)의 경우 $1 \in I$이므로 계수 하나가 $1$인 원소를 찾을 수 있다.
$m = \#\{K \in E \mid a_K \text{ is not a unit}\}$로 두자.
$0 \leq m \leq \# E - 1$임에 주의하자.
$m$에 대한 귀납법을 쓴다.

\medskip\noindent
$m = 0$인 경우를 보자. 이 경우 $g$의 모든 계수
$a_K$, $K \in E$는 단원이고 $E \not = \emptyset$이다.
$E = \{K_0\}$가 한원소 집합이고 $K_0 = (0, \ldots, 0)$이면 $g$는
단원이고 $J = S$이므로 결과는 확실히 성립한다. (이는 특히 $n = 0$일
때 일어나며 $n$에 대한 귀납법의 기초 단계를 제공한다.)
$E = \{(0, \ldots, 0)\}$이 아니면 적어도 하나의 $K$는
$(0, \ldots, 0)$과 다르다. 즉 $g \not \in R$이다. 이제 통상적인
뇌터 정규화 기법을 사용한다. 즉 다음을 생각하자.
$$
G(y_1, \ldots, y_n) =
g(y_1 + y_n^{e_1}, y_2 + y_n^{e_2}, \ldots, y_{n - 1} + y_n^{e_{n - 1}}, y_n)
$$
여기서 $0 \ll e_{n -1} \ll e_{n - 2} \ll \ldots \ll e_1$이다.
보조정리 \ref{algebra-lemma-helper-polynomial}에 의해
$G(y_1, \ldots, y_n)$은 $y_n$의 다항식으로서 다음과 같은 꼴이다.
$$
a_K y_n^{k_n + \sum_{i = 1, \ldots, n - 1} e_i k_i} +
\text{lower order terms in }y_n
$$
$a_K$가 단원이므로 $M = R[x_1, \ldots, x_n]/J$는
$R[y_1, \ldots, y_{n - 1}]$ 위에서 유한하다. 따라서
$$
U(R \to R[x_1, \ldots, x_n], M) =
U(R \to R[y_1, \ldots, y_{n - 1}], M)
$$
이고 $n$에 대한 귀납 가정으로부터 결론이 따른다.

\medskip\noindent
$m > 0$인 경우를 보자. $a_K$가 단원이 아닌 다중 첨자
$K \in E$를 택하자. 앞에서와 같이
$U_1 = \Spec(R_{a_K}) = \Spec(R) \setminus V(a_K)$로 두고
$$
U_2 = \Spec(R) \setminus \overline{U_1}.
$$
로 두자. 그러면 $U = U_1 \cup U_2$가 $\Spec(R)$에서 조밀함은 명백하다.
(a) $D(f) \subset U_1$이거나 (b) $D(f) \subset U_2$인 원소
$f \in R$을 택하자. 이러한 모든 $f$에 대해 순서쌍
$(R_f \to R_f[x_1, \ldots, x_n], M_f)$에 보조정리가 성립한다면,
보조정리 \ref{algebra-lemma-generic-flatness-locus-reduce}에 의해
$U(R \to S, M)$이 $\Spec(R)$에서 조밀함을 알 수 있다.
따라서 (a) $a_KR = R$이거나 (b) $V(a_K) = \Spec(R)$이라고
가정해도 된다. (a)의 경우 $a_K$가 단원이 되었으므로 수 $m$이 감소한다.
(b)의 경우 $R$이 기약환이므로 $a_K = 0$이라고 결론 내릴 수 있다.
따라서 집합 $E$가 작아지고 이 경우에도 수 $m$이 감소한다.
두 경우 모두 $m$에 대한 귀납 가정으로부터 결론이 따른다.

\medskip\noindent
이제 $S = R[x_1, \ldots, x_n]$인 경우의 보조정리를 증명했다.
$(R \to S, M)$이 보조정리의 조건을 만족하는 임의의 순서쌍이라고 하자.
전사사상 $R[x_1, \ldots, x_n] \to S$를 택하자. 식
(\ref{algebra-equation-good-locus})에서 도입한 표기를 사용하면
$$
U(R \to S, M) =
U(R \to R[x_1, \ldots, x_n], S)
\cap
U(R \to R[x_1, \ldots, x_n], M)
$$
이다. 방금 우변의 두 열린집합이 조밀함을 증명했으므로 좌변의
열린집합도 조밀하다.
\end{proof}






\section{Krull--Akizuki 정리 주변}
\label{algebra-section-krull-akizuki}

\noindent
Krull--Akizuki 정리의 한 응용은 이산 값매김환이 풍부하게 존재함을
보이는 것이다. 더 일반적으로 이 절에서는 뇌터 국소환을 지배하는
이산 값매김환을 구성하는 방법을 보인다.

\medskip\noindent
먼저 극대 아이디얼을 블로업하여 뇌터 국소 정역을 $1$차원 뇌터 국소
정역으로 지배하는 방법을 보인다.

\begin{lemma}
\label{algebra-lemma-dominate-by-dimension-1}
$R$을 분수체 $K$를 갖는 뇌터 국소 정역이라 하자.
$R$이 체가 아니라고 가정하자.
그러면 다음을 만족하는 $R \subset R' \subset K$가 존재한다.
\begin{enumerate}
\item $R'$은 차원 $1$인 뇌터 국소환이다.
\item $R \to R'$은 국소환 사상이다. 즉 $R'$은 $R$을 지배한다.
\item $R \to R'$은 본질적으로 유한형이다.
\end{enumerate}
\end{lemma}

\begin{proof}
$R$을 지배하는 임의의 값매김환 $A \subset K$를 택하자(보조정리
% P01-E1410
\ref{algebra-lemma-dominate}에 의해 그러한 환이 존재한다).
대응하는 값매김을 $v$로 나타내자.
$x_1, \ldots, x_r$를 $R$의 극대 아이디얼 $\mathfrak m$의 극소 생성계라
하자. $v(x_r) = \min\{v(x_1), \ldots, v(x_r)\}$이라고 가정할 수 있고
그렇게 하자. 다음 환을 생각하자.
$$
S = R[x_1/x_r, x_2/x_r, \ldots, x_{r - 1}/x_r] \subset K.
$$
$\mathfrak mS = x_rS$가 주 아이디얼임에 주의하자.
$S \subset A$이고 $v(x_r) > 0$이므로 $x_rS \not = S$임을 알 수 있다.
$x_rS$ 위의 극소 소 아이디얼 $\mathfrak q$를 택하자. 그러면 보조정리
\ref{algebra-lemma-minimal-over-1}에 의해 $\text{height}(\mathfrak q) = 1$이고
$\mathfrak q$는 $\mathfrak m$ 위에 놓인다. 따라서
$R' = S_{\mathfrak q}$가 답임을 알 수 있다.
\end{proof}

\begin{lemma}[Koll\'ar]
\label{algebra-lemma-hart-serre-loc-thm}
\begin{reference}
이는 J\'anos Koll\'ar의 출간 예정 논문
``Variants of normality for Noetherian schemes''에서 가져온 것이다.
\end{reference}
$(R, \mathfrak m)$을 뇌터 국소환이라 하자.
그러면 다음 가운데 정확히 하나가 성립한다.
\begin{enumerate}
\item $(R, \mathfrak m)$은 아르틴이다.
\item $(R, \mathfrak m)$은 차원 $1$인 정칙환이다.
\item $\text{depth}(R) \geq 2$이다.
\item 동형사상이 아닌 유한 환 사상 $R \to R'$이 존재하여, 그 핵과
여핵은 $\mathfrak m$의 어떤 거듭제곱에 의해 영멸되고,
$\mathfrak m$은 $R'$의 연관 소 아이디얼이 아니며 $R' \not = 0$이다.
\end{enumerate}
\end{lemma}

\begin{proof}
$(R, \mathfrak m)$이 아르틴환이 아닐 필요충분조건은
$V(\mathfrak m) \subset \Spec(R)$가 조밀한 곳이 전혀 없는 것이다.
명제 \ref{algebra-proposition-dimension-zero-ring}을 보라. 이제부터 이를
가정한다.

\medskip\noindent
$J \subset R$을 $\mathfrak m$의 어떤 거듭제곱에 의해 영멸되는 가장 큰
아이디얼이라 하자. $J \not = 0$이면 $R \to R/J$에 의해
$(R, \mathfrak m)$은 (4)의 경우이다.

\medskip\noindent
그렇지 않으면 $J = 0$이다. 특히 $\mathfrak m$은 $R$의 연관 소
아이디얼이 아니며, 보조정리 \ref{algebra-lemma-ideal-nonzerodivisor}에 의해
비영인자 $x \in \mathfrak m$이 존재함을 알 수 있다.
$\mathfrak m$이 $R/xR$의 연관 소 아이디얼이 아니면 같은 보조정리에 의해
$\text{depth}(R) \geq 2$이다. 따라서 이제 어떤
$y \in R$, $y \not \in xR$가 존재하여 $y \mathfrak m \subset xR$인 경우만
남는다.

\medskip\noindent
$y \mathfrak m \subset x \mathfrak m$이면 사상
$\varphi : \mathfrak m \to \mathfrak m$, $f \mapsto yf/x$를 생각할 수 있다
($x$가 비영인자이므로 잘 정의된다). 보조정리
\ref{algebra-lemma-charpoly-module}의 행렬식 기법에 의해 계수가 $R$에 속하고
$P(\varphi) = 0$인 일계수 다항식 $P$가 존재한다. 따라서
$R_x$에서 $P(y/x) = 0$이다.
$R' \subset R_x$를 $R$과 $y/x$로 생성되는 환이라 하자. 그러면
$R \subset R'$이고 $R'/R$은 $\mathfrak m$의 어떤 거듭제곱에 의해
영멸되는 유한 $R$-가군이다. 따라서 $R$은 (4)의 경우이다.

\medskip\noindent
그렇지 않으면 $R$의 어떤 단원 $u$와 어떤 $t \in \mathfrak m$이 존재하여
$y t = u x$이다. $t$를 $u^{-1}t$로 바꾸면 $yt = x$를 얻는다.
특히 $y$는 비영인자이다. 임의의 $t' \in \mathfrak m$에 대해 어떤
$s \in R$가 존재하여 $y t' = x s$이다. 따라서
$y (t' - s t ) = x s - x s = 0$이다. $y$가 비영인자이므로 이는
$t' = ts$를 함의하고, 따라서 $\mathfrak m = (t)$이다.
그러므로 $(R, \mathfrak m)$은 차원 1인 정칙환이다.

\medskip\noindent
지금까지의 논증은 모든 $R$이 네 경우 중 하나에 속함을 보인다.
증명을 끝내려면 성질들의 여섯 가지 조합 각각에 대해 두 성질이 동시에
성립할 수 없음을 보여야 한다. 여기서는 (2)와 (3)이 각각 (4)를
배제함만 보이고, 나머지 경우들은 독자에게 맡긴다.

\medskip\noindent
$R$이 차원 $1$인 정칙환이고 $R \to R'$이 유한 환 사상이며, 그 핵과
여핵은 $\mathfrak m$의 어떤 거듭제곱에 의해 영멸되고,
$\mathfrak m$은 $R'$의 연관 소 아이디얼이 아니라고 가정하자.
그러면 $R'$은 깊이가 적어도 $1$인 유한 $R$-가군이므로 보조정리
\ref{algebra-lemma-regular-mcm-free}에 의해 자유이다. $R \to R'$이 일반점에서
동형사상이므로 $R'$은 $R$-가군으로서 계수 $1$의 자유 가군이어야 한다.
그러면 $R' \to \text{End}_R(R') \cong R$은 사상 $R \to R'$의
역사상이다. 따라서 (4)는 성립할 수 없다.

\medskip\noindent
$R$의 깊이가 $\geq 2$이고 $R \to R'$이 유한 환 사상이며, 그 핵과
여핵은 $\mathfrak m$의 어떤 거듭제곱에 의해 영멸되고,
$\mathfrak m$은 $R'$의 연관 소 아이디얼이 아니라고 가정하자.
그러면 $R \to R'$은 반드시 단사이다(그렇지 않으면 핵의 깊이가
$0$인 동시에 $\geq 1$이어야 한다). 또한 여핵의 깊이는 적어도
$1$이다(보조정리 \ref{algebra-lemma-depth-in-ses}와 가정에 의해 $R'$의 깊이가
$\geq 1$이라는 사실을 사용한다). 따라서 여핵이
$\{\mathfrak m\}$ 위에 지지되려면 여핵 자체도 $0$이어야 한다.
그러므로 (4)는 성립할 수 없다.
\end{proof}

\begin{lemma}
\label{algebra-lemma-nonregular-dimension-one}
$R$을 극대 아이디얼 $\mathfrak m$을 갖는 국소환이라 하자.
$R$이 뇌터이고 차원이 $1$이며
$\dim(\mathfrak m/\mathfrak m^2) > 1$이라고 가정하자. 그러면 다음을
만족하는 환 사상 $R \to R'$이 존재한다.
\begin{enumerate}
\item $R \to R'$은 유한이다.
\item $R \to R'$은 동형사상이 아니다.
\item $R \to R'$의 핵과 여핵은 $\mathfrak m$의 어떤 거듭제곱에 의해
영멸된다.
\item $\mathfrak m$은 $R'$의 연관 소 아이디얼이 아니다.
\end{enumerate}
\end{lemma}

\begin{proof}
이는 보조정리 \ref{algebra-lemma-hart-serre-loc-thm}과 $R$이 아르틴환도,
정칙환도 아니고 깊이가 $\geq 2$이지도 않다는 사실에서 따른다
(마지막 사실은 보조정리 \ref{algebra-lemma-bound-depth}에 의해 깊이가 차원을
넘지 않는다는 데서 따른다).
\end{proof}

\begin{example}
\label{algebra-example-nonreduced}
다음 뇌터 국소환을 생각하자.
$$
R = k[[x, y]]/(y^2)
$$
이는 차원이 1이고 Cohen--Macaulay이다. 보조정리
\ref{algebra-lemma-nonregular-dimension-one}에서와 같은 확대의 한 예는
$$
k[[x, y]]/(y^2) \subset k[[x, z]]/(z^2), \ \ y \mapsto xz
$$
이다. 다시 말해 $R$에 $y/x$를 첨가하여 얻는다. 이 구성을
$n > 1$번 반복하면 원소 $y/x^n$을 첨가하게 된다.
\end{example}

\begin{example}
\label{algebra-example-bad-dvr-char-p}
$k$를 특성 $p > 0$인 체라 하되, 그 $p$제곱들의 부분체 $k^p$ 위에서
$k$의 차수가 무한이라고 하자. 예를 들어
$k = \mathbf{F}_p(t_1, t_2, t_3, \ldots)$로 둘 수 있다.
다음 환을 생각하자.
$$
A =
\left\{
\sum a_i x^i \in k[[x]] \text{ such that }
[k^p(a_0, a_1, a_2, \ldots) : k^p] < \infty
\right\}
$$
그러면 $A$는 이산 값매김환이고 그 완비화는
$A^\wedge = k[[x]]$이다. $A \subset k[[x]]$에서 유도되는 분수체들의
확대는 무한 순수 비분리 확대임에 주의하자. 임의의
$f \in k[[x]]$, $f \not \in A$를 택하고
$R = A[f] \subset k[[x]]$로 두자. 그러면 $R$은 차원이 $1$인 뇌터 국소
정역이고 그 완비화 $R^\wedge$은 기약환이 아니다(생각해 보라!).
\end{example}

\begin{remark}
\label{algebra-remark-resolution-dim-1}
$R$이 $1$차원 반국소 뇌터 정역이라고 하자. 어떤 극대 아이디얼
$\mathfrak m \subset R$에 대해 $R_{\mathfrak m}$이 정칙환이 아니면,
$(R, \mathfrak m)$에 보조정리 \ref{algebra-lemma-nonregular-dimension-one}을
적용하여 유한 환 확대 $R \subset R_1$을 얻을 수 있다.
(예를 들어 $\Spec(R_1) \to \Spec(R)$이 아이디얼 $\mathfrak m$에서의
$\Spec(R)$의 블로업이 되도록 할 수 있다.) 물론 $R_1$은 $R$과 같은
분수체를 갖는 $1$차원 반국소 뇌터 정역이다. $R_1$이 정칙 반국소환이
아니면 구성을 반복하여 $R_1 \subset R_2$를 얻을 수 있다. 따라서 유한
환 확대들의 열
$$
R \subset R_1 \subset R_2 \subset R_3 \subset \ldots
$$
을 얻으며, 어떤 $n$에 대해 $R_n$이 정칙이면 이 열은 멈출 수 있다.
특이점 해소는 결국 $R_n$이 실제로 정칙이 된다는 주장일 것이다.
실제로는 그렇지 않다. 즉 완비화가 기약환이 아닌, 특성 $0$이고 차원
$1$인 뇌터 국소 정역 $A$가 존재한다. \cite[Proposition 3.1]{Ferrand-Raynaud}
또는 우리의 예, 절 \ref{examples-section-local-completion-nonreduced}을
보라. 특성 $p > 0$인 예는 예
\ref{algebra-example-bad-dvr-char-p}를 보라. 블로업 구성은 완비화와 가환하므로
이 열이 결코 안정화되지 않음을 쉽게 알 수 있다. 논의는
\cite{Bennett}를 보라(주로 양의 특성의 경우를 다룬다). 다른 한편
$R$의 모든 극대 아이디얼에서의 완비화가 기약환이면 이 절차는 멈춘다
(여기에 나중 참고문헌 삽입).
\end{remark}

\begin{lemma}
\label{algebra-lemma-characterize-dvr}
$A$를 환이라 하자. 다음은 서로 동치이다.
\begin{enumerate}
\item 환 $A$는 이산 값매김환이다.
\item 환 $A$는 값매김환이고 뇌터이지만 체는 아니다.
\item 환 $A$는 차원 $1$인 정칙 국소환이다.
\item 환 $A$는 하나의 영이 아닌 원소로 생성되는 극대 아이디얼
$\mathfrak m$을 갖는 뇌터 국소 정역이다.
\item 환 $A$는 차원 $1$인 뇌터 국소 정규 정역이다.
\end{enumerate}
이 경우 $\pi$가 $A$의 극대 아이디얼의 생성원이면 $A$의 모든 영이 아닌
원소는 $u \in A$가 단원인 $u\pi^n$의 꼴로 유일하게 쓸 수 있다.
\end{lemma}

\begin{proof}
(1)과 (2)의 동치는 보조정리
\ref{algebra-lemma-valuation-ring-Noetherian-discrete}이다. 더욱이 보조정리
\ref{algebra-lemma-valuation-ring-Noetherian-discrete}의 증명에서 $A$가 이산
값매김환이면 $A$가 주 아이디얼 정역임을 보았으므로 (3)이 성립한다.
정칙 국소환은 정역임에 주의하자(보조정리
\ref{algebra-lemma-regular-domain}을 보라). 이를 사용하면 (3)과 (4)의 동치는
차원론에서 따른다. 절 \ref{algebra-section-dimension}을 보라.

\medskip\noindent
(3)을 가정하고 $\pi$를 극대 아이디얼 $\mathfrak m$의 생성원이라 하자.
모든 $n \geq 0$에 대해
$\dim_{A/\mathfrak m} \mathfrak m^n/\mathfrak m^{n + 1} = 1$이다.
실제로 이는 $\pi^n$으로 생성되며 영일 수 없다. 특히
$\mathfrak m^n = (\pi^n)$이고 등급 환
$\bigoplus \mathfrak m^n/\mathfrak m^{n + 1}$은 다항식환
$A/\mathfrak m[T]$와 동형이다.
$x \in A \setminus \{0\}$에 대해
$v(x) = \max\{n \mid x \in \mathfrak m^n\}$로 정의하자.
다시 말해 $u \in A^*$인 $x = u \pi^{v(x)}$로 쓸 수 있다.
위의 설명에 의해 모든 $x, y \in A \setminus \{0\}$에 대해
$v(xy) = v(x) + v(y)$이다. 이를 $v(a/b) = v(a) - v(b)$로 두어
$A$의 분수체 $K$로 확장한다(위에서 보인 곱셈성에 의해 잘 정의된다).
그러면 $A$는 값매김이 $\geq 0$인 $K$의 원소들의 집합임이 명백하다.
따라서 보조정리 \ref{algebra-lemma-valuation-valuation-ring}에 의해 $A$는
값매김환이다.

\medskip\noindent
보조정리 \ref{algebra-lemma-valuation-ring-normal}에 의해 값매김환은 정규
정역이다. 따라서 동치인 조건 (1) -- (3)이 (5)를 함의한다.
(5)를 가정하자. $\mathfrak m$이 $1$개의 원소로 생성될 수 없다고
가정하여 모순을 얻자. 그러면 보조정리
\ref{algebra-lemma-nonregular-dimension-one}에 의해, $\mathfrak m$의 임의의
영이 아닌 원소를 가역화하면 동형사상이 되지만 그 자체는 동형사상이
아닌 유한 환 사상 $A \to A'$이 존재한다. 특히 $A'$을 $A$의 분수체의
부분집합과 동일시할 수 있다. $A \to A'$은 유한이므로 정수적이다
(보조정리 \ref{algebra-lemma-finite-is-integral}을 보라).
$A$가 정규이므로 $A = A'$을 얻어 모순이다.
\end{proof}

\begin{definition}
\label{algebra-definition-uniformizer}
$A$를 이산 값매김환이라 하자. $A$의 극대 아이디얼을 생성하는 원소
$\pi \in A$를 {\it 균일화원}이라고 한다.
\end{definition}

\noindent
보조정리 \ref{algebra-lemma-characterize-dvr}에 의해 이산 값매김환의 임의의
두 균일화원은 서로 동반원이다.

\begin{lemma}
\label{algebra-lemma-finite-length}
$R$을 분수체 $K$를 갖는 정역이라 하자.
$M$을 $K^{\oplus r}$의 $R$-부분가군이라 하자.
$R$이 차원 $1$인 뇌터 국소환이라고 가정하자.
임의의 영이 아닌 $x \in R$에 대해
$\text{length}_R(R/xR) < \infty$이고
$$
\text{length}_R(M/xM) \leq r \cdot \text{length}_R(R/xR).
$$
이다.
\end{lemma}

\begin{proof}
$x$가 단원이면 결과가 성립한다. 따라서 $R$의 극대 아이디얼에 대해
$x \in \mathfrak m$이라고 가정해도 된다.
$x$는 영이 아니고 $R$은 정역이므로 $\dim(R/xR) = 0$이며, 따라서
$R/xR$은 유한 길이를 갖는다. 보조정리에서처럼
$M \subset K^{\oplus r}$를 생각하자. 필요하면 $K^{\oplus r}$를 더 작은
부분공간으로 바꾸어, $M$의 원소들이 $K$-벡터공간으로서
$K^{\oplus r}$를 생성한다고 가정해도 된다.

\medskip\noindent
먼저 $M$이 유한 $R$-가군이라고 하자. 이 경우 분모들을 제거하여
$M \subset R^{\oplus r}$이라고 가정할 수 있다. $M$이 벡터공간으로서
% P01-E1411
$K^{\oplus r}$를 생성하므로 $R^{\oplus r}/M$은 유한 길이를 갖는다.
특히 $x^cR^{\oplus r} \subset M$이 되도록 하는 정수 $c \geq 0$이
존재한다. $M \supset xM \supset x^2M \supset \ldots$은 연속한 몫들이
각각 $M/xM$과 동형인 가군들의 열임에 주의하자. 따라서
$$
n \text{length}_R(M/xM) = \text{length}_R(M/x^nM).
$$
임을 알 수 있다. $M = R^{\oplus r}$에 대한 같은 논증은
$$
n \text{length}_R(R^{\oplus r}/xR^{\oplus r}) =
\text{length}_R(R^{\oplus r}/x^nR^{\oplus r}).
$$
임을 보인다. 위에서 택한 $c$에 의해 $x^nM$은
$x^n R^{\oplus r}$와 $x^{n + c}R^{\oplus r}$ 사이에 놓인다.
이로부터 쉽게
$$
r(n + c) \text{length}_R(R/xR)
\geq
n \text{length}_R(M/xM)
\geq
r (n - c) \text{length}_R(R/xR)
$$
을 얻는다. 따라서 유한 가군인 경우에는 실제로 보조정리의 결과가
등호와 함께 성립한다.

\medskip\noindent
이제 $M$이 유한 가군이 아니라고 하자. 어떤 자연수 $k$에 대해
$M/xM$의 길이가 $\geq k$라고 하자. 그러면
% P01-E1412
$$
0 \subset N_0 \subset N_1 \subset N_2 \subset \ldots N_k \subset M/xM
$$
이고 $i = 0, \ldots k - 1$에 대해 $N_i \not = N_{i + 1}$인 사슬을
찾을 수 있다. 각 $i = 1, \ldots, k$에 대해 $xM$을 법으로 한 합동류가
$N_i$에는 속하지만 $N_{i - 1}$에는 속하지 않는 원소 $m_i \in M$을
택하자. 유한 $R$-가군
$M' = Rm_1 + \ldots + Rm_k \subset M$을 생각하자.
$N'_i \subset M'/xM'$를 $N_i$의 역상이라 하자. $m_i$를 택한 방식에 의해
$N'_i \not =N'_{i + 1}$임은 명백하다. 따라서
$\text{length}_R(M'/xM') \geq k$임을 알 수 있다. 유한 가군인 경우에
의해 원하는 대로 $k \leq r\text{length}_R(R/xR)$라고 결론 내린다.
\end{proof}

\noindent
다음은 첫 번째 응용이다.

\begin{lemma}
\label{algebra-lemma-finite-extension-residue-fields-dimension-1}
$R \to S$를 분수체들의 단사사상 $K \subset L$을 유도하는 정역들의
준동형사상이라 하자. $R$이 차원 $1$인 뇌터 국소환이고
$[L : K] < \infty$이면 다음이 성립한다.
\begin{enumerate}
\item $R$의 극대 아이디얼 $\mathfrak m$ 위에 놓이는 $S$의 각 소
아이디얼 $\mathfrak n_i$는 극대이다.
\item 이러한 소 아이디얼은 유한 개이다.
\item 각 $i$에 대해
$[\kappa(\mathfrak n_i) : \kappa(\mathfrak m)] < \infty$이다.
\end{enumerate}
\end{lemma}

\begin{proof}
영이 아닌 $x \in \mathfrak m$을 택하자. 여기서 $n = [L : K]$로 두고
부분가군 $S \subset L \cong K^{\oplus n}$에 보조정리
\ref{algebra-lemma-finite-length}를 적용하자. 그러면 환 $S/xS$는 $R$ 위에서
유한 길이를 갖는다. 따라서 $S/\mathfrak m S$는
$\kappa(\mathfrak m)$ 위에서 유한 길이를 갖는다. 다시 말해
$\dim_{\kappa(\mathfrak m)} S/\mathfrak m S$는 유한하다
(보조정리 \ref{algebra-lemma-dimension-is-length}). 따라서
$S/\mathfrak mS$는 아르틴환이다(보조정리
\ref{algebra-lemma-finite-dimensional-algebra}). 아르틴환에 관한 구조 정리들은
% P01-E1413
(1)과 (2)를 함의한다. 예를 들어 보조정리
\ref{algebra-lemma-artinian-finite-length}를 보라. (3)은 위에서 확립한
유한성에서 따른다.
\end{proof}

\begin{lemma}
\label{algebra-lemma-finite-length-global}
$R$을 분수체 $K$를 갖는 정역이라 하자.
$M$을 $K^{\oplus r}$의 $R$-부분가군이라 하자.
$R$이 차원 $1$인 뇌터환이라고 가정하자.
임의의 영이 아닌 $x \in R$에 대해
$\text{length}_R(M/xM) < \infty$이다.
\end{lemma}

\begin{proof}
$R$의 차원이 $1$이므로 $x$는 유한 개의 소 아이디얼
$\mathfrak m_i$, $i = 1, \ldots, n$에 포함되며, 이들은 모두 극대이다.
$R$이 뇌터이므로 $R/xR$은 아르틴환이고, 명제
\ref{algebra-proposition-dimension-zero-ring}과 보조정리
\ref{algebra-lemma-artinian-finite-length}에 의해
$R/xR = \prod_{i = 1, \ldots, n} (R/xR)_{\mathfrak m_i}$이다.
따라서 $M/xM$도 마찬가지로 $\mathfrak m_i$에서의 국소화들의 곱
$M/xM = \prod (M/xM)_{\mathfrak m_i}$로 분해된다.
$R_{\mathfrak m_i}$ 위의 $M_{\mathfrak m_i}$에 보조정리
\ref{algebra-lemma-finite-length}를 적용하면 각각의
$M_{\mathfrak m_i}/xM_{\mathfrak m_i} = (M/xM)_{\mathfrak m_i}$가
$R_{\mathfrak m_i}$ 위에서 유한 길이를 가짐을 알 수 있다.
따라서 $M/xM$은 $R$ 위에서 유한 길이를 갖는다. 실제로 위의 사실에
의해 $M/xM$은 연속한 몫들이 잉여체 $\kappa(\mathfrak m_i)$와 동형인
$R$-부분가군들의 유한 여과를 갖는다.
\end{proof}

\begin{lemma}[Krull-Akizuki]
\label{algebra-lemma-krull-akizuki}
$R$을 분수체 $K$를 갖는 정역이라 하자.
$L/K$를 체들의 유한 확대라 하자.
$R$이 뇌터이고 $\dim(R) = 1$이라고 가정하자.
이 경우 $R \subset A \subset L$을 만족하는 임의의 환 $A$는 뇌터이다.
\end{lemma}

\begin{proof}
$I \subset A$를 영이 아닌 아이디얼이라 하자. 보조정리
\ref{algebra-lemma-intersection-not-zero}에 의해 영이 아닌 원소
$x \in I \cap R$을 찾을 수 있다. 그러면
$I/xA \subset A/xA$를 얻는다. 보조정리
\ref{algebra-lemma-finite-length-global}에 의해 $R$-가군 $A/xA$는
$R$-가군으로서 유한 길이를 갖는다. 따라서 $I/xA$도 $R$-가군으로서
유한 길이를 갖는다. 그러므로 $I$는 $A$의 아이디얼로서 유한 생성이다.
\end{proof}

\begin{lemma}
\label{algebra-lemma-exists-dvr}
$R$을 분수체 $K$를 갖는 뇌터 국소 정역이라 하자.
$R$이 체가 아니라고 가정하자.
$L/K$를 유한 생성 체 확대라 하자.
그러면 분수체가 $L$이고 $R$을 지배하는 이산 값매김환 $A$가 존재한다.
\end{lemma}

\begin{proof}
$L$이 $K$ 위에서 유한하지 않으면 $K$ 위에서 $L$의 초월기저
$x_1, \ldots, x_r$를 택하고, $R$을 $\mathfrak m_R$과
$x_1, \ldots, x_r$로 생성되는 극대 아이디얼에서 국소화한
$R[x_1, \ldots, x_r]$로 바꾸자. 따라서 $K \subset L$이 유한이라고
가정해도 된다.

\medskip\noindent
보조정리 \ref{algebra-lemma-dominate-by-dimension-1}에 의해
$\dim(R) = 1$이라고 가정해도 된다.

\medskip\noindent
$A \subset L$을 $L$ 안에서 $R$의 정수적 폐포라 하자.
보조정리 \ref{algebra-lemma-krull-akizuki}에 의해 이는 뇌터이다.
보조정리 \ref{algebra-lemma-integral-overring-surjective}에 의해 $R$의 극대
아이디얼 위에 놓이는 소 아이디얼 $\mathfrak q \subset A$가 존재한다.
보조정리 \ref{algebra-lemma-characterize-dvr}에 의해 환 $A_{\mathfrak q}$는
원하는 대로 $R$을 지배하는 이산 값매김환이다.
\end{proof}

\section{인수분해}
\label{algebra-section-factoring}

\noindent
다음은 보통 학부 1학년 대수학 과정에서 배우는 몇 가지 개념과
그 사이의 관계이다.

\begin{definition}
\label{algebra-definition-irreducible-prime-element}
$R$을 정역이라 하자.
\begin{enumerate}
\item 단원 $u \in R^*$가 존재하여 $x = uy$이면 원소
$x, y \in R$를 {\it 동반원}이라고 한다.
\item 원소 $x \in R$가 영이 아니고 단원이 아니며, $x = yz$이고
$y, z \in R$일 때마다 $y$가 단원이거나 $x$의 동반원이면
{\it 기약원}이라고 한다.
\item $x$가 생성하는 아이디얼이 소 아이디얼이면 원소
$x \in R$를 {\it 소원}이라고 한다.
\end{enumerate}
\end{definition}

\begin{lemma}
\label{algebra-lemma-easy-divisibility}
$R$을 정역이라 하자. $x, y \in R$라 하자.
$x$, $y$가 동반원일 필요충분조건은 $(x) = (y)$인 것이다.
\end{lemma}

\begin{proof}
어떤 단원 $u \in R$에 대하여 $x = uy$이면 $(x) \subset (y)$이고,
$y = u^{-1}x$이므로 $(y) \subset (x)$이기도 하다. 역으로
$(x) = (y)$라고 하자.
% P01-E1414
그러면 어떤 $f, g \in A$에 대하여 $x = fy$이고 $y = gx$이다.
따라서 $x = fg x$이고 $R$이 정역이므로 $fg = 1$이다. 그러므로
$x$와 $y$는 동반원이다.
\end{proof}

\begin{lemma}
\label{algebra-lemma-factorization-exists}
$R$을 정역이라 하자. 다음 조건들을 생각하자.
\begin{enumerate}
\item 환 $R$은 주 아이디얼에 대한 오름 사슬 조건을 만족한다.
\item 영이 아니고 단원이 아닌 모든 원소 $a \in R$에는 각 $b_i$가
$R$의 기약원인 인수분해 $a = b_1 \ldots b_k$가 존재한다.
\end{enumerate}
그러면 (1)은 (2)를 함의한다.
\end{lemma}

\begin{proof}
영이 아니고 단원이 아니며 기약원들로 인수분해되지 않는 원소를
$x$라 하자. $x_1 = x$로 두자. $y$와 $z$가 모두 기약원도 단원도
아닌 방식으로 $x = yz$라고 쓸 수 있다.
$y$가 기약원들로 인수분해되지 않으면 $x_2 = y$로 두고,
$z$가 기약원들로 인수분해되지 않으면 $x_2 = z$로 둔다.
이 과정을 계속하면 $R$의 원소들로 이루어진 수열
$$
\ldots | x_3 | x_2 | x_1
$$
을 얻으며, 여기서 $x_n/x_{n + 1}$은 단원이 아니다.
이는 주 아이디얼들의 엄밀한 증가 수열
$(x_1) \subset (x_2) \subset (x_3) \subset \ldots$을 주므로 증명이
끝난다.
\end{proof}

\begin{definition}
\label{algebra-definition-UFD}
정역 $R$이 다음 성질을 가지면 {\it 유일 인수분해 정역}, 줄여서
{\it UFD}라고 한다. $x \in R$가 영이 아니고 단원이 아니면
$x$는 기약원들로 인수분해되고, 기약원들로의 인수분해
$$
x = a_1 \ldots a_m = b_1 \ldots b_n
$$
이 주어지면 $n = m$이며, 각 $a_i$와 $b_{\sigma(i)}$가 동반원이 되게
하는 순열 $\sigma : \{1, \ldots, n\} \to \{1, \ldots, n\}$이 존재한다.
\end{definition}

\begin{lemma}
\label{algebra-lemma-characterize-UFD}
$R$을 정역이라 하자. 영이 아니고 단원이 아닌 모든 원소가
기약원들로 인수분해된다고 가정하자. 그러면 $R$이 UFD일
필요충분조건은 모든 기약원이 소원인 것이다.
\end{lemma}

\begin{proof}
$R$이 UFD라고 가정하고 $x \in R$를 기약원이라 하자.
$ab \in (x)$, 즉 $ab = cx$라고 하자. 인수분해
$a = a_1 \ldots a_n$, $b = b_1 \ldots b_m$, 그리고
$c = c_1 \ldots c_r$를 택하자. 인수분해
$$
a_1 \ldots a_n b_1 \ldots b_m = c_1 \ldots c_r x
$$
의 유일성에 의해 $x$는 원소들 $a_1, \ldots, b_m$ 중 하나의
동반원이다. 다시 말해 $a \in (x)$ 또는 $b \in (x)$이고, 따라서
$x$는 소원이다.

\medskip\noindent
모든 기약원이 소원이라고 가정하자. 기약원들로의 인수분해가
순열과 동반원으로 바꾸는 것까지 유일함을 증명해야 한다.
$a_i$와 $b_j$가 기약원이고
$a_1 \ldots a_m = b_1 \ldots b_n$이라고 하자.
$a_1$이 소원이므로 어떤 $j$에 대하여 $b_j \in (a_1)$이다.
번호를 다시 매겨 $b_1 \in (a_1)$이라고 가정해도 된다.
그러면 $b_1 = a_1 u$이고, $b_1$이 기약원이므로 $u$는 단원이다.
따라서 $a_1$과 $b_1$은 동반원이며,
% P01-E1415
$a_2 \ldots a_n = ub_2\ldots b_m$이다. $n + m$에 관한 귀납법으로
$n = m$이고, 원하는 대로 $i = 2, \ldots, n$에 대하여
$a_i$는 $b_{\sigma(i)}$의 동반원임을 알 수 있다.
\end{proof}

\begin{lemma}
\label{algebra-lemma-characterize-UFD-height-1}
$R$을 뇌터 정역이라 하자. 그러면 $R$이 UFD일 필요충분조건은
높이 $1$인 모든 소 아이디얼이 주 아이디얼인 것이다.
\end{lemma}

\begin{proof}
$R$이 UFD라고 가정하고 $\mathfrak p$를 높이 1인 소 아이디얼이라
하자. 영이 아닌 $x \in \mathfrak p$를 택하고
$x = a_1 \ldots a_n$을 기약원들로의 인수분해라 하자.
$\mathfrak p$가 소 아이디얼이므로 어떤 $i$에 대하여
$a_i \in \mathfrak p$이다. 보조정리 \ref{algebra-lemma-characterize-UFD}에
의해 아이디얼 $(a_i)$는 소 아이디얼이다. $\mathfrak p$의 높이가
$1$이므로 $(a_i) = \mathfrak p$라고 결론 내린다.

\medskip\noindent
높이 $1$인 모든 소 아이디얼이 주 아이디얼이라고 가정하자.
$R$이 뇌터이므로 영이 아니고 단원이 아닌 모든 원소 $x$는
기약원들로 인수분해된다. 보조정리
\ref{algebra-lemma-factorization-exists}를 보라. 보조정리
\ref{algebra-lemma-characterize-UFD}에 의해 기약원 $x$가 소원임을 보이면
충분하다. $(x)$ 위에 극소인 소 아이디얼
$(x) \subset \mathfrak p$를 택하자. 보조정리
\ref{algebra-lemma-minimal-over-1}에 의해 $\mathfrak p$의 높이는 $1$이다.
가정에 의해 $\mathfrak p = (y)$이다. 따라서 $x = yz$이고,
$x$가 기약원이므로 $z$는 단원이다. 그러므로 $(x) = (y)$이고
$x$가 소원임을 알 수 있다.
\end{proof}

\begin{lemma}[Nagata의 팩토리얼성 판정법]
\label{algebra-lemma-invert-prime-elements}
\begin{reference}
\cite[Lemma 2]{Nagata-UFD}
\end{reference}
$A$를 정역이라 하자. $S \subset A$를 소원들로 생성되는 곱셈적
부분집합이라 하자. $x \in A$를 기약원이라 하자. 그러면
\begin{enumerate}
\item $S^{-1}A$에서 $x$의 상은 기약원이거나 단원이고,
\item $x$가 소원일 필요충분조건은 $S^{-1}A$에서 $x$의 상이
$S^{-1}A$의 소원이거나 단원인 것이다.
\end{enumerate}
더 나아가 이 경우 $A$가 UFD일 필요충분조건은 $A$의 영이 아니고
단원이 아닌 모든 원소가 기약원들로 인수분해되고 $S^{-1}A$가
UFD인 것이다.
\end{lemma}

\begin{proof}
$\alpha, \beta \in S^{-1}A$에 대하여 $x = \alpha \beta$라고 하자.
그러면 어떤 $a, b \in A$, $s, s' \in S$에 대하여
$\alpha = a/s$이고 $\beta = b/s'$이다. 따라서 $ss'x = ab$를 얻는다.
가정에 의해 어떤 소원들 $p_i$에 대하여
$ss' = p_1 \ldots p_r$라고 쓸 수 있다. 각 $i$에 대하여
$p_i$는 $a$ 또는 $b$를 나눈다. 나누어 나가면 어떤
$s'', s''' \in S$에 대하여 $x = a' b'$이고 $a = s'' a'$,
$b = s''' b'$인 인수분해를 얻는다. $x$가 기약원이므로
$a'$ 또는 $b'$은 단원이다. 이를 거슬러 올라가면
$\alpha$ 또는 $\beta$가 단원임을 알 수 있다. 이로써 (1)을
증명했다.

\medskip\noindent
$x$가 소원이라고 하자. 그러면 $A/(x)$는 정역이다. 따라서
$S^{-1}A/xS^{-1}A = S^{-1}(A/(x))$는 정역이거나 영환이다.
그러므로 $x$는 소원 또는 단원으로 사상된다.

\medskip\noindent
$S^{-1}A$에서 $x$의 상이 단원이라고 하자. 그러면 어떤
$s \in S$와 $y \in A$에 대하여 $y x = s$이다. 가정에 의해
$p_i$들이 소원인 $s = p_1 \ldots p_r$를 얻는다. 각 $i$에 대하여
$p_i$는 $y$를 나누거나 $p_i$는 $x$를 나눈다. 두 번째 경우에는
$p_i$와 $x$가 동반원이므로($x$가 기약원이기 때문이다) 증명이
끝난다. 그러나 모든 $i = 1, \ldots, r$에 대하여 첫 번째 경우가
일어나면 $x$가 단원이라는 모순을 얻는다.

\medskip\noindent
$S^{-1}A$에서 $x$의 상이 소원이라고 하자. $a, b \in A$이고
$ab \in (x)$라고 가정하자. 그러면 어떤 $s \in S$와 $y \in A$에
대하여 $sa = xy$ 또는 $sb = xy$이다. 첫 번째 경우라고 하자.
가정에 의해 $p_i$들이 소원인 $s = p_1 \ldots p_r$를 얻는다.
각 $i$에 대하여 $p_i$는 $y$를 나누거나 $p_i$는 $x$를 나눈다.
두 번째 경우에는 $p_i$와 $x$가 동반원이므로($x$가 기약원이기
때문이다) 증명이 끝난다. 모든 $i = 1, \ldots, r$에 대하여
첫 번째 경우가 일어나면 원하는 대로 $a \in (x)$이다.
이로써 (2)의 증명이 끝난다.

\medskip\noindent
보조정리의 마지막 명제는 (1), (2), 그리고 보조정리
\ref{algebra-lemma-characterize-UFD}에서 따른다.
\end{proof}

\begin{lemma}
\label{algebra-lemma-UFD-ascending-chain-condition-principal-ideals}
UFD는 주 아이디얼에 대한 오름 사슬 조건을 만족한다.
\end{lemma}

\begin{proof}
$R$의 주 아이디얼들의 오름 사슬
$(a_1) \subset (a_2) \subset (a_3) \subset \ldots$을 생각하자.
$p_i$들이 소원이도록 $a_1 = p_1^{e_1} \ldots p_r^{e_r}$라고 쓰자.
그러면 어떤 $0 \leq c_i \leq e_i$에 대하여 $a_n$이
$p_1^{c_1} \ldots p_r^{c_r}$의 동반원임을 알 수 있다. 가능한
경우가 유한 개뿐이므로 결론을 얻는다.
\end{proof}

\begin{lemma}
\label{algebra-lemma-factoring-in-polynomial}
$R$을 정역이라 하자. $R$이 주 아이디얼에 대한 오름 사슬 조건을
만족한다고 가정하자. 그러면 $R$ 위의 다항식환도 같은 성질을
만족한다.
\end{lemma}

\begin{proof}
$R[x]$의 주 아이디얼들의 오름 사슬
$(f_1) \subset (f_2) \subset (f_3) \subset \ldots$을 생각하자.
$f_{n + 1}$이 $f_n$을 나누므로 이 수열에서 차수들은 감소한다.
따라서 모든 $n \gg 0$에 대하여 $f_n$의 차수는 어떤 고정된
$d \geq 0$이다. $a_n$을 $f_n$의 최고차항 계수라 하자.
조건 $f_n \in (f_{n + 1})$은 모든 $n$에 대하여 $a_{n + 1}$이
$a_n$을 나눔을 함의한다. $R$에 관한 가정에 의해 충분히 큰 모든
$n$에 대하여 $a_{n + 1}$과 $a_n$은 동반원이다
(보조정리 \ref{algebra-lemma-easy-divisibility}). 따라서 큰 $n$에 대하여
$f_n = u f_{n + 1}$임을 알 수 있는데, 여기서 $u \in R$는
(차수 때문에) 환의 원소이며 ($a_n$과 $a_{n + 1}$이 동반원이므로)
단원이다.
\end{proof}

\begin{lemma}
\label{algebra-lemma-polynomial-ring-UFD}
UFD 위의 다항식환은 UFD이다. 특히 $k$가 체이면
$k[x_1, \ldots, x_n]$은 UFD이다.
\end{lemma}

\begin{proof}
$R$을 UFD라 하자. 그러면 $R$은 주 아이디얼에 대한 오름 사슬
조건을 만족하고
(보조정리 \ref{algebra-lemma-UFD-ascending-chain-condition-principal-ideals}),
따라서 $R[x]$도 주 아이디얼에 대한 오름 사슬 조건을 만족하며
(보조정리 \ref{algebra-lemma-factoring-in-polynomial}), 따라서
$R[x]$의 모든 원소는 기약원들로 인수분해된다
(보조정리 \ref{algebra-lemma-factorization-exists}).
$S \subset R$를 소원들로 생성되는 곱셈적 부분집합이라 하자.
$R$의 모든 비단원은 소원들의 곱이므로
$K = S^{-1}R$은 $R$의 분수체이다. $R$의 모든 소원은 $R[x]$의
소원으로 사상되고, $S^{-1}(R[x]) = S^{-1}R[x] = K[x]$는 UFD
(더 나아가 PID)이기도 함을 관찰하자. 따라서 보조정리
\ref{algebra-lemma-invert-prime-elements}를 적용하여 결론을 얻을 수 있다.
\end{proof}

\begin{lemma}
\label{algebra-lemma-UFD-normal}
유일 인수분해 정역은 정규이다.
\end{lemma}

\begin{proof}
$R$을 UFD라 하자. $x$를 $R$ 위에서 정수적인 $R$의 분수체의
원소라 하자. $a_i \in R$인
$x^d - a_1 x^{d - 1} - \ldots - a_d = 0$이 성립한다고 하자.
$u$가 단원이고 $e_i \in \mathbf{Z}$이며
$p_1, \ldots, p_r$가 서로 동반원이 아닌 기약원들이도록
$x = u p_1^{e_1} \ldots p_r^{e_r}$라고 쓸 수 있다.
보조정리를 증명하려면 $e_i \geq 0$임을 보여야 한다. 그렇지 않고
$e_1 < 0$이라고 하자. 그러면 $N \gg 0$에 대하여
$$
u^d p_2^{de_2 + N} \ldots p_r^{de_r + N} =
p_1^{-de_1}p_2^N \ldots p_r^N(
\sum\nolimits_{i = 1, \ldots, d} a_i x^{d - i} ) \in (p_1)
$$
을 얻는데, 이는 $R$에서의 인수분해의 유일성과 모순이다.
\end{proof}

\begin{definition}
\label{algebra-definition-PID}
모든 아이디얼이 주 아이디얼인 정역 $R$을
{\it 주 아이디얼 정역}, 줄여서 {\it PID}라고 한다.
\end{definition}

\begin{lemma}
\label{algebra-lemma-PID-UFD}
주 아이디얼 정역은 유일 인수분해 정역이다.
\end{lemma}

\begin{proof}
PID는 뇌터이므로 보조정리
\ref{algebra-lemma-characterize-UFD-height-1}에서 따른다.
\end{proof}

\begin{definition}
\label{algebra-definition-dedekind-domain}
정역 $R$의 영이 아닌 모든 아이디얼 $I \subset R$가
영이 아닌 소 아이디얼들의 곱
$$
I = \mathfrak p_1 \ldots \mathfrak p_r
$$
으로 $\mathfrak p_i$들의 순열까지 유일하게 표현되면 이를
{\it 데데킨트 정역}이라고 한다.
\end{definition}

\begin{lemma}
\label{algebra-lemma-PID-dedekind}
PID는 데데킨트 정역이다.
\end{lemma}

\begin{proof}
$R$을 PID라 하자. $R$의 영이 아닌 모든 아이디얼은 주 아이디얼이고
$R$은 UFD이므로(보조정리 \ref{algebra-lemma-PID-UFD}), $R$의 모든 기약원이
소원이라는 사실(보조정리 \ref{algebra-lemma-characterize-UFD})에서 결론이
따른다. 실제로 원소들의 인수분해가 소 아이디얼들로의 인수분해를
준다.
\end{proof}

\begin{lemma}
\label{algebra-lemma-product-ideals-principal}
\begin{slogan}
아이디얼들의 곱이 가역 가군일 필요충분조건은 두 인수가 모두
가역 가군인 것이다.
\end{slogan}
$A$를 환이라 하자. $I$와 $J$를 $A$의 영이 아닌 아이디얼들이라
하고, 어떤 영인자가 아닌 원소 $f \in A$에 대하여 $IJ = (f)$라고
하자. 그러면 $I$와 $J$는 유한 생성 아이디얼들이며, $A$-가군으로서
계수 $1$의 유한 국소 자유 가군이다.
\end{lemma}

\begin{proof}
보조정리 \ref{algebra-lemma-finite-projective}에 의해 $I$와 $J$가
계수 $1$의 유한 국소 자유 $A$-가군임을 보이면 충분하다.
이를 위해 $x_i \in I$이고 $y_i \in J$인
$f = \sum_{i = 1, \ldots, n} x_i y_i$라고 쓰자. 또한 어떤
$a_i \in A$에 대하여 $x_i y_i = a_i f$라고 쓸 수 있다.
$f$가 영인자가 아니므로 $\sum a_i = 1$이다. 따라서 각
$I_{a_i}$와 $J_{a_i}$가 $A_{a_i}$ 위의 계수 $1$ 자유 가군임을
보이면 충분하다. $A$를 $A_{a_i}$로 바꾸고 나면 어떤 $x \in I$와
$y \in J$에 대하여 $f = xy$라고 결론 내린다. $x$와 $y$는 모두
영인자가 아님을 주목하자. $I = (x)$이고 $J = (y)$라고 주장한다.
그러면 증명이 끝난다. 실제로 $x' \in I$이면 어떤 $a \in A$에
대하여 $x'y = af = axy$이다. 따라서 $x' = ax$이고 결론을 얻는다.
\end{proof}

\begin{lemma}
\label{algebra-lemma-characterize-Dedekind}
$R$을 환이라 하자. 다음 조건들은 서로 동치이다.
\begin{enumerate}
\item $R$은 데데킨트 정역이다.
\item $R$은 뇌터 정역이고, 영이 아닌 모든 극대 아이디얼
$\mathfrak m$에 대하여 국소환 $R_{\mathfrak m}$은 이산 값매김환이다.
\item $R$은 뇌터 정규 정역이고 $\dim(R) \leq 1$이다.
\end{enumerate}
\end{lemma}

\begin{proof}
(1)을 가정하자. 데데킨트 정역의 정의에서 $R$이 뇌터라고 가정하지
않았으므로 논증은 자명하지 않다. $\mathfrak p \subset R$를
소 아이디얼이라 하자. 정의에 나오는 인수분해의 유일성에 의해
$\mathfrak p \not = \mathfrak p^2$임을 관찰하자.
$x \not \in \mathfrak p^2$인 $x \in \mathfrak p$를 택하자.
$y \in \mathfrak p$를 두 번째 원소로 택하자(예를 들어 $y = 0$).
$(x, y) = \mathfrak p_1 \ldots \mathfrak p_r$라고 쓰자.
$(x, y) \subset \mathfrak p$이므로 소 아이디얼들
$\mathfrak p_i$ 중 적어도 하나는 $\mathfrak p$에 포함된다.
그러나 $x \not \in \mathfrak p^2$이므로 많아야 하나만 포함된다.
따라서 $\mathfrak p_1, \ldots, \mathfrak p_r$ 중 정확히 하나가
$\mathfrak p$에 포함되며, 이를 $\mathfrak p_1 \subset \mathfrak p$라
하자. 모든 $y$의 선택에 대하여
$(x, y)R_\mathfrak p = \mathfrak p_1R_\mathfrak p$가 소 아이디얼이라고
결론 내린다. $(x)R_\mathfrak p = \mathfrak pR_\mathfrak p$라고
주장한다. 실제로 $y \in \mathfrak p$를 택하자. 위의 논증을
$y^2$에 적용하면 $(x, y^2)R_\mathfrak p$가 소 아이디얼이다.
따라서 $y \in (x, y^2)R_\mathfrak p$, 즉 $R_\mathfrak p$에서
$y = ax + by^2$이다. 그러므로
$(1 - by)y = ax \in (x)R_\mathfrak p$, 즉 원하는 대로
$y \in (x)R_\mathfrak p$이다.

\medskip\noindent
$(x) = \mathfrak p_1 \ldots \mathfrak p_r$이고
$\mathfrak p_1 \subset \mathfrak p$이도록 새로 쓰면,
$\mathfrak p_1 R_\mathfrak p = \mathfrak p R_\mathfrak p$, 즉
$\mathfrak p_1 = \mathfrak p$라고 결론 내린다. 더 나아가
보조정리 \ref{algebra-lemma-product-ideals-principal}에 의해
$\mathfrak p_1 = \mathfrak p$는 $R$의 유한 생성 아이디얼이다.
보조정리 \ref{algebra-lemma-cohen}에 의해 $R$이 뇌터라고 결론 내린다.
% P01-E1416
더 나아가 모든 소 아이디얼 $\mathfrak p$에 대하여
$R_\mathfrak m$이 이산 값매김환임이 따른다. 보조정리
\ref{algebra-lemma-characterize-dvr}를 보라.

\medskip\noindent
(2)와 (3)의 동치는 보조정리 \ref{algebra-lemma-normality-is-local}과
\ref{algebra-lemma-characterize-dvr}에서 따른다. (2)와 (3)이 성립한다고
가정하자. $I \subset R$를 아이디얼이라 하자. $I$의 인수분해를
구성하겠다. $I$가 소 아이디얼이면 증명할 것이 없다. 그렇지 않으면
$\mathfrak p \subset R$가 극대이도록 $I \subset \mathfrak p$를
택하자. $J = \{x \in R \mid x \mathfrak p \subset I\}$라 하자.
$J \mathfrak p = I$라고 주장한다. $R$의 극대 아이디얼
$\mathfrak m$들에서 국소화한 뒤 이를 확인하면 충분하다
($J$의 형성은 국소화와 가환하며 보조정리
\ref{algebra-lemma-characterize-zero-local}을 사용한다).
그러면 $\mathfrak p R_\mathfrak m = R_\mathfrak m$이고 결론이
분명하거나, $\mathfrak p R_\mathfrak m = \mathfrak m R_\mathfrak m$이다.
뒤의 경우 $\mathfrak p R_\mathfrak m = (\pi)$이고,
$\mathfrak p$가 주 아이디얼인 경우의 결론은 즉시 따른다.
뇌터 귀납법에 의해 아이디얼 $J$에는 인수분해가 존재하며,
따라서 $I$의 원하는 인수분해를 얻는다. 인수분해의 유일성 증명은
생략한다.
\end{proof}

\noindent
다음은 Krull--Akizuki 보조정리의 한 변형이다.

\begin{lemma}
\label{algebra-lemma-integral-closure-Dedekind}
$A$를 분수체 $K$를 갖는 차원 $1$의 뇌터 정역이라 하자.
$L/K$를 유한 확대라 하자. $B$를 $L$ 안에서 $A$의 정수적 폐포라
하자. 그러면 $B$는 데데킨트 정역이고,
$\Spec(B) \to \Spec(A)$는 전사이며 유한한 섬유들을 갖고
유한 잉여체 확대들을 유도한다.
\end{lemma}

\begin{proof}
Krull--Akizuki(보조정리 \ref{algebra-lemma-krull-akizuki})에 의해 환
$B$는 뇌터이다. 보조정리 \ref{algebra-lemma-integral-sub-dim-equal}에 의해
$\dim(B) = 1$이다. 따라서 보조정리
\ref{algebra-lemma-characterize-Dedekind}에 의해 $B$는 데데킨트 정역이다.
스펙트럼 사이 사상의 전사성은 보조정리
\ref{algebra-lemma-integral-overring-surjective}에서 따른다.
마지막 두 명제는 보조정리
\ref{algebra-lemma-finite-extension-residue-fields-dimension-1}에서 따른다.
\end{proof}

\section{소멸 차수}
\label{algebra-section-orders-of-vanishing}

\begin{lemma}
\label{algebra-lemma-ord-additive}
$R$을 차원 $1$인 반국소 뇌터환이라 하자.
$a, b \in R$가 영인자가 아니면
$$
\text{length}_R(R/(ab)) =
\text{length}_R(R/(a)) +
\text{length}_R(R/(b))
$$
이고 이 길이들은 유한하다.
\end{lemma}

\begin{proof}
유한성은 보조정리 \ref{algebra-lemma-finite-length-global}에서 보았다.
첫 번째 사상이 $b$를 곱하는 사상으로 주어지는 짧은 완전열
$0 \to R/(a) \to R/(ab) \to R/(b) \to 0$이 있으므로 가법성이
성립한다. (길이의 가법성을 사용하라. 보조정리
\ref{algebra-lemma-length-additive}를 보라.)
\end{proof}

\begin{definition}
\label{algebra-definition-ord}
$K$를 체라 하고, $R \subset K$를 분수체가 $K$인 차원 $1$의
국소\footnote{$R$이 반국소이기만 한 경우에도 이를 정의할 수 있지만,
아마 실제로 원하는 개념은 결코 아닐 것이다!} 뇌터 부분환이라 하자.
이 경우 {\it $R$을 따른 소멸 차수}
$$
\text{ord}_R : K^* \longrightarrow \mathbf{Z}
$$
를 다음 규칙으로 정의한다.
$$
\text{ord}_R(x) = \text{length}_R(R/(x))
$$
단, $x \in R$인 경우이고, 영이 아닌 $x, y \in R$에 대해서는
$\text{ord}_R(x/y) = \text{ord}_R(x) - \text{ord}_R(y)$로 둔다.
\end{definition}

\noindent
소멸 차수를 사용하여 벡터공간 안의 격자들을 비교할 수 있다.
다음이 그 정의이다.

\begin{definition}
\label{algebra-definition-lattice}
$R$을 분수체 $K$를 갖는 차원 $1$의 뇌터 국소 정역이라 하자.
$V$를 유한 차원 $K$-벡터공간이라 하자. $V = K \otimes_R M$을
만족하는 유한 $R$-부분가군 $M \subset V$를 {\it $V$의 격자}라고
한다.
\end{definition}

\noindent
조건 $V = K \otimes_R M$은 $M$이 벡터공간 $V$의 한 기저를
포함한다는 뜻이다. 문헌의 여러 곳에서는 환 $R$이 이산 값매김환인
경우에만 격자라는 개념을 정의하기도 한다. $R$이 이산 값매김환이면
모든 격자는 자유 $R$-가군이지만, 일반적으로는 그렇지 않을 수 있다.

\begin{lemma}
\label{algebra-lemma-compare-lattices}
$R$을 분수체 $K$를 갖는 차원 $1$의 뇌터 국소 정역이라 하자.
$V$를 유한 차원 $K$-벡터공간이라 하자.
\begin{enumerate}
\item $M$이 $V$의 격자이고 $M \subset M' \subset V$가 $M$을
포함하는 $V$의 $R$-부분가군이면 다음 조건들은 서로 동치이다.
\begin{enumerate}
\item $M'$은 격자이다.
\item $\text{length}_R(M'/M)$은 유한하다.
\item $M'$은 유한 생성이다.
\end{enumerate}
\item $M$이 $V$의 격자이고 $M' \subset M$이 $M$의
$R$-부분가군이면, $M'$이 격자일 필요충분조건은
$\text{length}_R(M/M')$가 유한한 것이다.
\item $M$, $M'$이 $V$의 격자이면 $M \cap M'$과 $M + M'$도
그렇다.
\item $M \subset M' \subset M'' \subset V$가 $V$의 격자들이면
$$
\text{length}_R(M''/M) =
\text{length}_R(M'/M) +
\text{length}_R(M''/M').
$$
\item $M$, $M'$, $N$, $N'$이 $V$의 격자들이고
$N \subset M \cap M'$, $M + M' \subset N'$이면
\begin{eqnarray*}
& & \text{length}_R(M/M \cap M') - \text{length}_R(M'/M \cap M')\\
& = &
\text{length}_R(M/N) - \text{length}_R(M'/N) \\
& = &
\text{length}_R(M + M' / M') - \text{length}_R(M + M'/M) \\
& = &
\text{length}_R(N' / M') - \text{length}_R(N'/M)
\end{eqnarray*}
\end{enumerate}
\end{lemma}

\begin{proof}
(1)을 증명하자. (1)(a)를 가정하자.
$y_1, \ldots, y_m$이 $M'$을 생성한다고 하자. 그러면 각
$y_i = x_i/f_i$이며, 여기서 $x_i \in M$이고 $f_i \in R$는 영이
아니다. 따라서 $f_1 \ldots f_m M' \subset M$임을 알 수 있다.
$R$은 차원 $1$의 뇌터 국소환이므로 어떤 $n$에 대하여
$\mathfrak m^n \subset (f_1 \ldots f_m)$이다(예를 들어 보조정리
\ref{algebra-lemma-one-equation}과 명제
\ref{algebra-proposition-dimension-zero-ring}을 함께 사용하거나,
보조정리 \ref{algebra-lemma-finite-length}와
\ref{algebra-lemma-length-infinite}를 함께 사용하라).
다시 말해 어떤 $n$에 대하여 $\mathfrak m^nM' \subset M$이다.
% P01-E1417
따라서 보조정리 \ref{algebra-lemma-length-finite}에 의해
$\text{length}(M'/M) < \infty$이고, 다시 말해 (1)(b)가 성립한다.
(1)(b)를 가정하자. 그러면 $M'/M$은 유한 $R$-가군이다
(보조정리 \ref{algebra-lemma-finite-length-finite}를 보라).
따라서 $M'$은 유한 $R$-가군들의 확대이므로 유한 $R$-가군이다.
그러므로 (1)(c)가 성립한다. 함의 (1)(c) $\Rightarrow$ (1)(a)는
정의 \ref{algebra-definition-lattice} 뒤의 설명에서 따른다.

\medskip\noindent
(2)를 증명하자. $M$이 $V$의 격자이고 $M' \subset M$이
$R$-부분가군이라고 하자. (1)에서 $M'$이 격자이면
$\text{length}_R(M/M') < \infty$임을 보았다. 역으로
$\text{length}_R(M/M') < \infty$라고 가정하자. $R$이 뇌터이므로
$M'$은 유한 생성이고, 어떤 $n$에 대하여
$\mathfrak m^n M \subset M'$이다
(보조정리 \ref{algebra-lemma-length-infinite}). 따라서 $M'$은 $V$의
한 기저를 포함하며 $M'$은 격자이다.

\medskip\noindent
(3)을 증명하자. $M$, $M'$이 $V$의 격자라고 하자. $R$이
뇌터이므로 $M$의 부분가군 $M \cap M'$은 유한하다.
$M$이 격자이므로 $V$의 $K$-기저를 이루는
$x_1, \ldots, x_n \in M$을 찾을 수 있다. $M'$이 격자이므로
$y_i \in M'$이고 $f_i \in R$인 $x_i = y_i/f_i$라고 쓸 수 있다.
따라서 $f_ix_i \in M \cap M'$이다. 그러므로 $M \cap M'$도
격자이다. $M + M'$이 격자라는 사실은 (1)에서 따른다.

\medskip\noindent
(4)는 길이의 가법성(보조정리 \ref{algebra-lemma-length-additive})과
완전열
$$
0 \to M'/M \to M''/M \to M''/M' \to 0
$$
에서 따른다. (5)는 (4)를 반복하여 적용하면 따른다.
\end{proof}

\begin{definition}
\label{algebra-definition-distance}
$R$을 분수체 $K$를 갖는 차원 $1$의 뇌터 국소 정역이라 하자.
$V$를 유한 차원 $K$-벡터공간이라 하자. $M$, $M'$을 $V$의 두
격자라 하자. {\it $M$과 $M'$ 사이의 거리}는 보조정리
\ref{algebra-lemma-compare-lattices}의 (5)에 나오는 정수
$$
d(M, M') = \text{length}_R(M/M \cap M') - \text{length}_R(M'/M \cap M')
$$
이다.
\end{definition}

\noindent
특히 $M' \subset M$이면
$d(M, M') = \text{length}_R(M/M')$이다.

\begin{lemma}
\label{algebra-lemma-properties-distance-function}
$R$을 분수체 $K$를 갖는 차원 $1$의 뇌터 국소 정역이라 하자.
$V$를 유한 차원 $K$-벡터공간이라 하자. 이 거리 함수는
$V$의 세 격자 $M$, $M'$, $M''$이 주어질 때마다
$$
d(M, M'') = d(M, M') + d(M', M'')
$$
을 만족한다. 특히 $d(M, M') = - d(M', M)$이다.
\end{lemma}

\begin{proof}
생략한다.
\end{proof}

\begin{lemma}
\label{algebra-lemma-order-vanishing-determinant}
$R$을 분수체 $K$를 갖는 차원 $1$의 뇌터 국소 정역이라 하자.
$V$를 유한 차원 $K$-벡터공간이라 하자.
$\varphi : V \to V$를 $K$-선형 동형사상이라 하자.
임의의 격자 $M \subset V$에 대하여
$$
d(M, \varphi(M)) = \text{ord}_R(\det(\varphi))
$$
이다.
\end{lemma}

\begin{proof}
정수 $d(M, \varphi(M))$이 격자 $M$에 의존하지 않음을 다음과 같이
알 수 있다. $M'$을 또 하나의 그러한 격자라 하자. 그러면
\begin{eqnarray*}
d(M, \varphi(M)) & = & d(M, M') + d(M', \varphi(M)) \\
& = & d(M, M') + d(\varphi(M'), \varphi(M)) + d(M', \varphi(M'))
\end{eqnarray*}
이다. $\varphi$가 동형사상이므로
$d(\varphi(M'), \varphi(M)) = d(M', M) = -d(M, M')$이고, 따라서
$d(M, \varphi(M)) = d(M', \varphi(M'))$이다. 또한
(보조정리의) 등식 양변은 $\varphi$에 대해 가법적이다. 즉,
$$
\text{ord}_R(\det(\varphi \circ \psi))
=
\text{ord}_R(\det(\varphi))
+
\text{ord}_R(\det(\psi))
$$
이고, 위에서 보인 독립성에 의해
% P01-E1418
\begin{eqnarray*}
d(M, \varphi(\psi((M))) & = &
d(M, \psi(M)) + d(\psi(M), \varphi(\psi(M))) \\
& = & d(M, \psi(M)) + d(M, \varphi(M))
\end{eqnarray*}
이기도 하다. 따라서 $\text{GL}(V)$의 생성원들에 대하여 보조정리를
증명하면 충분하다. 동형사상 $K^{\oplus n} \cong V$를 택하자.
그러면 $\text{GL}(V) = \text{GL}_n(K)$는 기본행렬 $E$들로
생성된다. $E$가 단위행렬인 경우 결론은 분명하다.
$i \not = j$, $\lambda \in K$, $\lambda \not = 0$인
$E = E_{ij}(\lambda)$인 경우, 예를 들어
$$
E_{12}(\lambda) =
\left(
\begin{matrix}
1 & \lambda & \ldots \\
0 & 1 & \ldots \\
\ldots & \ldots & \ldots
\end{matrix}
\right)
$$
인 경우에는 다른 기저에 관하여 $E_{12}(1)$을 얻는다.
$R^{\oplus n} \subset K^{\oplus n}$을 격자로 택하면
$E = E_{12}(1)$인 경우의 결론은 분명하다. 마지막으로
$a \in K^*$이고 $E = E_i(a)$인 경우, 예를 들어
$$
E_1(a) =
\left(
\begin{matrix}
a & 0 & \ldots \\
0 & 1 & \ldots \\
\ldots & \ldots & \ldots
\end{matrix}
\right)
$$
인 경우에는
% P01-E1419
$E_1(a)(R^{\oplus b}) = aR \oplus R^{\oplus n - 1}$이고,
원하는 대로
$d(R^{\oplus n}, aR \oplus R^{\oplus n - 1})
= \text{ord}_R(a)$임이 분명하다.
\end{proof}

\begin{lemma}
\label{algebra-lemma-finite-extension-dim-1}
$A \to B$를 환 준동형사상이라 하자. 다음을 가정하자.
\begin{enumerate}
\item $A$는 차원 $1$의 뇌터 국소 정역이다.
\item $A \subset B$는 정역들의 유한 확대이다.
\end{enumerate}
$L/K$를 이에 대응하는 분수체들의 유한 확대라 하자.
$y \in L^*$이고 $x = \text{Nm}_{L/K}(y)$라 하자.
이 상황에서 $B$는 반국소이다. $\mathfrak m_i$,
$i = 1, \ldots, n$을 $B$의 극대 아이디얼들이라 하자. 그러면
$$
\text{ord}_A(x) =
\sum\nolimits_i
[\kappa(\mathfrak m_i) : \kappa(\mathfrak m_A)]
\text{ord}_{B_{\mathfrak m_i}}(y)
$$
이다. 여기서 $\text{ord}$는 정의 \ref{algebra-definition-ord}에서처럼
정의한다.
\end{lemma}

\begin{proof}
보조정리 \ref{algebra-lemma-finite-in-codim-1}에 의해 환 $B$는 반국소이다.
어떤 $b, b' \in B$에 대하여 $y = b/b'$라고 쓰자.
$\text{ord}$의 가법성과 $\text{Nm}$의 곱셈성에 의해
$y = b$ 또는 $y = b'$인 경우에 보조정리를 증명하면 충분하다.
다시 말해 $y \in B$라고 가정해도 된다. 이 경우 공식의 우변은
$$
\sum [\kappa(\mathfrak m_i) : \kappa(\mathfrak m_A)]
\text{length}_{B_{\mathfrak m_i}}((B/yB)_{\mathfrak m_i})
$$
이다. 보조정리 \ref{algebra-lemma-pushdown-module}에 의해 이는
$\text{length}_A(B/yB)$와 같다. 보조정리
\ref{algebra-lemma-order-vanishing-determinant}에 의해
$$
\text{length}_A(B/yB) = d(B, yB) =
\text{ord}_A(\det\nolimits_K(L \xrightarrow{y} L)).
$$
정의에 의해
$x = \text{Nm}_{L/K}(y) = \det\nolimits_K(L \xrightarrow{y} L)$이므로
보조정리가 증명되었다.
\end{proof}

\section{준유한 사상}
\label{algebra-section-quasi-finite}

\noindent
유한형 환 준동형사상 $R \to S$를 생각하자.
사상 $\Spec(S) \to \Spec(R)$는 한 점이 그 섬유에서 고립점일 때
그 점에서 준유한이라고 한다. 이는 그 섬유가 그 점에서
영차원이라는 뜻이다. 이 절에서는 중요하지만 기술적인 이 개념의
기본 성질을 연구한다. 더 고급 내용은 다음 절에서 찾을 수 있다.

\begin{lemma}
\label{algebra-lemma-isolated-point}
$k$를 체라 하자. $S$를 유한형 $k$-대수라 하자.
$\mathfrak q$를 $S$의 소 아이디얼이라 하자. 다음 조건들은
서로 동치이다.
\begin{enumerate}
\item $\mathfrak q$는 $\Spec(S)$의 고립점이다.
\item $S_{\mathfrak q}$는 $k$ 위에서 유한하다.
\item $D(g) = \{ \mathfrak q \}$가 되는
$g \in S$, $g \not\in \mathfrak q$가 존재한다.
\item $\dim_{\mathfrak q} \Spec(S) = 0$이다.
\item $\mathfrak q$는 $\Spec(S)$의 닫힌점이고
$\dim(S_{\mathfrak q}) = 0$이다.
\item 체확대 $\kappa(\mathfrak q)/k$가 유한하고
$\dim(S_{\mathfrak q}) = 0$이다.
\end{enumerate}
이 경우 어떤 유한형 $k$-대수 $S'$에 대하여
$S = S_{\mathfrak q} \times S'$이다. 또한 (3)의 원소 $g$는
$S_{\mathfrak q} = S_g$라는 성질을 갖는다.
\end{lemma}

\begin{proof}
$\mathfrak q$가 $\Spec(S)$의 고립점, 즉
$\{\mathfrak q\}$가 $\Spec(S)$에서 열려 있다고 하자.
$\Spec(S)$는 야콥슨 공간이므로(보조정리
\ref{algebra-lemma-finite-type-field-Jacobson}과
\ref{algebra-lemma-jacobson}을 보라) $\mathfrak q$는 닫힌점이다.
따라서 $\{\mathfrak q\}$는 $\Spec(S)$에서 열린닫힌집합이다.
보조정리 \ref{algebra-lemma-disjoint-decomposition}과
\ref{algebra-lemma-disjoint-implies-product}에 의해, $\mathfrak q$가
$\Spec(S_1)$의 유일한 점에 대응하도록
$S = S_1 \times S_2$라고 쓸 수 있다. 따라서
$S_1 = S_{\mathfrak q}$는 $k$ 위의 유한형 영차원 환이다.
그러므로 예를 들어 보조정리 \ref{algebra-lemma-Noether-normalization}에
의해 $k$ 위에서 유한하다. (1)이 (2)를 함의함을 증명했다.

\medskip\noindent
$S_{\mathfrak q}$가 $k$ 위에서 유한하다고 하자. 그러면 보조정리
\ref{algebra-lemma-finite-dimensional-algebra}에 의해
$S_{\mathfrak q}$는 아르틴 국소환이다. 따라서 보조정리
\ref{algebra-lemma-artinian-finite-length}에 의해
$\Spec(S_{\mathfrak q}) = \{\mathfrak qS_{\mathfrak q}\}$이다.
완전열 $0 \to K \to S \to S_{\mathfrak q}
\to Q \to 0$을 생각하자. $K_{\mathfrak q} = Q_{\mathfrak q} = 0$임은
분명하다. 또한 $S$가 뇌터이므로 $K$는 유한 $S$-가군이고,
$S_{\mathfrak q}$가 $k$ 위에서 유한하므로 $Q$도 유한
$S$-가군이다. 따라서 $K_g = Q_g = 0$이 되게 하는
$g \in S$, $g \not \in \mathfrak q$가 존재한다. 그러므로
$S_{\mathfrak q} = S_g$이고 $D(g) = \{ \mathfrak q \}$이다.
(2)가 (3)을 함의함을 증명했다.

\medskip\noindent
$D(g) =  \{ \mathfrak q \}$라고 하자. $\Spec(S)$의 위상을 구성한
방식에 의해 $D(g)$는 열려 있으므로 $\mathfrak q$는
$\Spec(S)$의 고립점이다. (3)이 (1)을 함의함을 증명했다.
다시 말해 (1), (2), (3)은 서로 동치이다.

\medskip\noindent
$\dim_{\mathfrak q} \Spec(S) = 0$이라고 가정하자. 이는
$\Spec(S)$에서 $\mathfrak q$의 어떤 열린 근방이 영차원이라는
뜻이다. 그러면 영차원인 $D(g)$ 꼴의 열린 근방이 존재한다.
$S_g$가 뇌터이므로 $S_g$는 아르틴이고
$D(g) = \Spec(S_g)$는 유한 이산집합이다. 명제
\ref{algebra-proposition-dimension-zero-ring}을 보라. 따라서
$\mathfrak q$는 $D(g)$의 고립점이고, 위의 (1)과 (2)의 동치를
$\mathfrak qS_g \subset S_g$에 적용하면
$S_{\mathfrak q} = (S_g)_{\mathfrak qS_g}$가 $k$ 위에서 유한함을
알 수 있다. 그러므로 (4)는 (2)를 함의한다. (1)이 (4)를 함의함은
분명하다. 이로써 (1)--(4)는 모두 서로 동치이다.

\medskip\noindent
보조정리 \ref{algebra-lemma-dimension-closed-point-finite-type-field}은
함의 (5) $\Rightarrow$ (4)를 준다.
함의 (4) $\Rightarrow$ (6)은 보조정리
\ref{algebra-lemma-dimension-at-a-point-finite-type-field}에서 따른다.
함의 (6) $\Rightarrow$ (5)는 보조정리
\ref{algebra-lemma-finite-residue-extension-closed}에서 따른다.
이제 (1)--(6)이 서로 동치임을 안다.

\medskip\noindent
보조정리 끝의 두 명제는 위에서 (1), (2), (3)의 동치를 증명하는
과정에서 보았다.
\end{proof}

\begin{lemma}
\label{algebra-lemma-isolated-point-fibre}
\begin{slogan}
섬유의 고립점에 대한 동치 조건
\end{slogan}
$R \to S$를 유한형 환 준동형사상이라 하자.
$\mathfrak q \subset S$를 $\mathfrak p \subset R$ 위에 놓이는
소 아이디얼이라 하자. 비고 \ref{algebra-remark-fundamental-diagram}에서처럼
$F = \Spec(S \otimes_R \kappa(\mathfrak p))$를
$\Spec(S) \to \Spec(R)$의 섬유라 하자. $\mathfrak q$에 대응하는
점을 $\overline{\mathfrak q} \in F$로 나타내자. 다음 조건들은
서로 동치이다.
\begin{enumerate}
\item $\overline{\mathfrak q}$는 $F$의 고립점이다.
\item $S_{\mathfrak q}/\mathfrak pS_{\mathfrak q}$는
$\kappa(\mathfrak p)$ 위에서 유한하다.
\item $\mathfrak p$로 사상되는 $D(g)$의 유일한 소 아이디얼이
$\mathfrak q$가 되게 하는
$g \in S$, $g \not \in \mathfrak q$가 존재한다.
\item $\dim_{\overline{\mathfrak q}}(F) = 0$이다.
\item $\overline{\mathfrak q}$는 $F$의 닫힌점이고
$\dim(S_{\mathfrak q}/\mathfrak pS_{\mathfrak q}) = 0$이다.
\item 체확대 $\kappa(\mathfrak q)/\kappa(\mathfrak p)$가 유한하고
$\dim(S_{\mathfrak q}/\mathfrak pS_{\mathfrak q}) = 0$이다.
\end{enumerate}
\end{lemma}

\begin{proof}
$S_{\mathfrak q}/\mathfrak pS_{\mathfrak q} =
(S \otimes_R \kappa(\mathfrak p))_{\overline{\mathfrak q}}$임을
주목하자. 또한 $S \otimes_R \kappa(\mathfrak p)$는
$\kappa(\mathfrak p)$ 위의 유한형 대수이다.
각 조건은 $\kappa(\mathfrak p)$-대수
$S \otimes_R \kappa(\mathfrak p)$와 소 아이디얼
$\overline{\mathfrak q}$에 관한 보조정리
\ref{algebra-lemma-isolated-point}의 조건들에 정확히 대응하므로 서로
동치이다.
\end{proof}

\begin{definition}
\label{algebra-definition-quasi-finite}
$R \to S$를 유한형 환 준동형사상이라 하자.
$\mathfrak q \subset S$를 소 아이디얼이라 하자.
\begin{enumerate}
\item 보조정리 \ref{algebra-lemma-isolated-point-fibre}의 동치 조건들이
성립하면 $R \to S$가 {\it $\mathfrak q$에서 준유한}이라고 한다.
\item 환 준동형사상 $A \to B$가 유한형이고 $B$의 모든 소
아이디얼에서 준유한이면 이를 {\it 준유한}이라고 한다.
\end{enumerate}
\end{definition}

\begin{lemma}
\label{algebra-lemma-quasi-finite-above-p}
$R \to S$를 유한형 환 준동형사상이라 하고
$\mathfrak p$를 $R$의 소 아이디얼이라 하자. 다음 조건들은
서로 동치이다.
\begin{enumerate}
\item $R \to S$는 $\mathfrak p$ 위에 놓이는 $S$의 모든 소
아이디얼에서 준유한이다.
\item $S \otimes_R \kappa(\mathfrak p)$는 유한
$\kappa(\mathfrak p)$-대수이다.
\item $\Spec(S \otimes_R \kappa(\mathfrak p))$는 유한집합이다.
\end{enumerate}
\end{lemma}

\begin{proof}
조건 (1)은 $\Spec(S \otimes_R \kappa(\mathfrak p))$의 모든 점이
열려 있다는 것, 즉 그 위상이 이산이라는 것을 말한다. 따라서
(1), (2), (3)의 동치는 보조정리
\ref{algebra-lemma-finite-type-algebra-finite-nr-primes}에서 따른다.
\end{proof}

\begin{lemma}
\label{algebra-lemma-quasi-finite}
$R \to S$를 유한형 환 준동형사상이라 하자.
$R \to S$가 준유한일 필요충분조건은 모든 소 아이디얼
$\mathfrak p \subset R$에 대하여 환
$S \otimes_R \kappa(\mathfrak p)$가 $\kappa(\mathfrak p)$ 위에서
유한한 것이다.
\end{lemma}

\begin{proof}
더 일반적인 보조정리 \ref{algebra-lemma-quasi-finite-above-p}에서 즉시
따른다.
\end{proof}

\begin{lemma}
\label{algebra-lemma-quasi-finite-local}
$R \to S$를 유한형 환 준동형사상이라 하자.
$\mathfrak q \subset S$를 $\mathfrak p \subset R$ 위에 놓이는
소 아이디얼이라 하자.
$f \in R$, $f \not \in \mathfrak p$와
$g \in S$, $g \not \in \mathfrak q$라 하자.
$R \to S$가 $\mathfrak q$에서 준유한일 필요충분조건은
$R_f \to S_{fg}$가 $\mathfrak qS_{fg}$에서 준유한인 것이다.
\end{lemma}

\begin{proof}
$\Spec(S_{fg}) \to \Spec(R_f)$의 섬유는
$\Spec(S) \to \Spec(R)$의 섬유의 한 열린 부분집합과 위상동형이다.
따라서 보조정리 \ref{algebra-lemma-isolated-point-fibre}의 동치 조건 중
(1)에서 결론이 따른다.
\end{proof}

\begin{lemma}
\label{algebra-lemma-four-rings}
소 아이디얼들이 표시된 환들의 가환 그림
$$
\xymatrix{
S \ar[r] & S' & &
\mathfrak q \ar@{-}[r] & \mathfrak q' \\
R \ar[u] \ar[r] &  R' \ar[u] & &
\mathfrak p \ar@{-}[r] \ar@{-}[u] & \mathfrak p' \ar@{-}[u]
}
$$
이 주어졌다고 하자. $R \to S$가 유한형이고
$S \otimes_R R' \to S'$이 전사라고 가정하자.
$R \to S$가 $\mathfrak q$에서 준유한이면
$R' \to S'$은 $\mathfrak q'$에서 준유한이다.
\end{lemma}

\begin{proof}
보조정리 \ref{algebra-lemma-isolated-point}에서처럼
$S_1$은 $\kappa(\mathfrak p)$ 위에서 유한하고 $\mathfrak q$는
$S_1$의 한 점에 대응하도록
$S \otimes_R \kappa(\mathfrak p) = S_1 \times S_2$라고 쓰자.
이 곱분해는 임의의 $S \otimes_R \kappa(\mathfrak p)$-대수에
대응하는 곱분해를 유도한다. 특히
$S' \otimes_{R'} \kappa(\mathfrak p') = S'_1 \times S'_2$를 얻는다.
$S \otimes_R R' \to S'$이 전사이므로 표준 사상
$(S \otimes_R \kappa(\mathfrak p)) \otimes_{\kappa(\mathfrak p)}
\kappa(\mathfrak p') \to S' \otimes_{R'} \kappa(\mathfrak p')$은
전사이고, 따라서
$S_i \otimes_{\kappa(\mathfrak p)} \kappa(\mathfrak p') \to S'_i$도
전사이다. 그러므로 $S'_1$은 $\kappa(\mathfrak p')$ 위에서 유한하다.
사상 $S' \otimes_{R'} \kappa(\mathfrak p') \to
\kappa(\mathfrak q')$은 $S_1'$을 통하여 인수분해된다
(즉 인수 $S_2'$을 영으로 보낸다). 이는 사상
$S \otimes_R \kappa(\mathfrak p) \to
\kappa(\mathfrak q)$이 $S_1$을 통하여 인수분해되기 때문이다
(즉 인수 $S_2$를 영으로 보낸다). 따라서 보조정리
\ref{algebra-lemma-spec-product}에서처럼 $\mathfrak q'$는 섬유의 분리합집합
분해
$\Spec(S' \otimes_{R'} \kappa(\mathfrak p'))
= \Spec(S_1') \amalg \Spec(S_2')$에서 $\Spec(S_1')$의 한 점에
대응한다. $S_1'$은 체 위에서 유한하므로
% P01-E1420
아르틴환이고, 따라서 $\Spec(S_1')$은 유한 이산집합이다.
(명제 \ref{algebra-proposition-dimension-zero-ring}을 보라.)
그러므로 원하는 대로 $\mathfrak q'$는 그 섬유에서 고립되어 있다.
\end{proof}

\begin{lemma}
\label{algebra-lemma-quasi-finite-composition}
준유한 환 준동형사상들의 합성은 준유한이다.
\end{lemma}

\begin{proof}
$A \to B$와 $B \to C$가 준유한 환 준동형사상이라고 하자.
보조정리 \ref{algebra-lemma-compose-finite-type}에 의해
$A \to C$가 유한형임을 알 수 있다. $\mathfrak r \subset C$를
$\mathfrak q \subset B$와 $\mathfrak p \subset A$ 위에 놓이는
$C$의 소 아이디얼이라 하자.
% P01-E1421
$A \to B$와 $B \to C$가 각각 $\mathfrak q$와 $\mathfrak r$에서
준유한이므로, $\mathfrak p$로 사상되는 $D(b)$의 유일한 소
아이디얼이 $\mathfrak q$이고 마찬가지로 $\mathfrak q$로 사상되는
$D(c)$의 유일한 소 아이디얼이 $\mathfrak r$이 되게 하는
$b \in B$와 $c \in C$가 존재한다. $c' \in C$가 $b \in B$의
상이면 $\mathfrak r$은 $\mathfrak p$로 사상되는 $D(cc')$의
유일한 소 아이디얼이다. 그러므로 $A \to C$는
$\mathfrak r$에서 준유한이다.
\end{proof}

\begin{lemma}
\label{algebra-lemma-quasi-finite-base-change}
$R \to S$를 유한형 환 준동형사상이라 하자.
$R \to R'$을 임의의 환 준동형사상이라 하자.
$S' = R' \otimes_R S$로 두자.
\begin{enumerate}
\item 집합
$\{\mathfrak q' \mid R' \to S' \text{ quasi-finite at }\mathfrak q'\}$
은 표준 사상 $\Spec(S') \to \Spec(S)$ 아래에서
$\Spec(S)$의 대응하는 집합의 역상이다.
\item $\Spec(R') \to \Spec(R)$가 전사이면
$R \to S$가 준유한일 필요충분조건은 $R' \to S'$이 준유한인 것이다.
\item 준유한 환 준동형사상의 모든 밑변환은 준유한이다.
\end{enumerate}
\end{lemma}

\begin{proof}
$\mathfrak p' \subset R'$을 $\mathfrak p \subset R$ 위에 놓이는
소 아이디얼이라 하자. 그러면 섬유환
$S' \otimes_{R'} \kappa(\mathfrak p')$은 체확대
$\kappa(\mathfrak p) \to \kappa(\mathfrak p')$에 의한 섬유환
$S \otimes_R \kappa(\mathfrak p)$의 밑변환이다. 따라서 첫 번째
명제는 체확대 아래에서 차원이 불변이라는 사실
(보조정리 \ref{algebra-lemma-dimension-at-a-point-preserved-field-extension})과
보조정리 \ref{algebra-lemma-isolated-point}에서 따른다.
밑변환 아래에서 준유한 사상이 안정적이라는 사실은 이 결과와
유한형 성질이 밑변환 아래에서 안정적이라는 사실에서 따른다.
두 번째 명제는 가정에 의해 임의의 소 아이디얼
$\mathfrak q \subset S$ 위에 놓이는 소 아이디얼
$\mathfrak q' \subset S'$을 찾을 수 있기 때문에 따른다.
\end{proof}

\begin{lemma}
\label{algebra-lemma-quasi-finite-permanence}
$A \to B$와 $B \to C$를 $A \to C$가 유한형이 되게 하는
환 준동형사상들이라 하자. $\mathfrak r$을
$\mathfrak q \subset B$와 $\mathfrak p \subset A$ 위에 놓이는
$C$의 소 아이디얼이라 하자. $A \to C$가 $\mathfrak r$에서
준유한이면 $B \to C$도 $\mathfrak r$에서 준유한이다.
\end{lemma}

\begin{proof}
보조정리 \ref{algebra-lemma-compose-finite-type}에 의해 $B \to C$가
유한형이므로 이 명제는 의미가 있다. 보조정리
\ref{algebra-lemma-isolated-point-fibre}의 조건 (3)을 사용하자.
$A \to C$가 $\mathfrak r$에서 준유한이면 어떤 $c \in C$가
존재하여
$$
\{\mathfrak r' \subset C \text{ lying over }\mathfrak p\} \cap D(c) =
\{\mathfrak{r}\}.
$$
$\mathfrak q$ 위에 놓이는 소 아이디얼들 $\mathfrak r' \subset C$은
$\mathfrak p$ 위에 놓이는 소 아이디얼들
$\mathfrak r' \subset C$의 부분집합이므로, $B \to C$가
$\mathfrak r$에서 준유한이라고 결론 내린다.
\end{proof}

\begin{lemma}
\label{algebra-lemma-generically-finite}
$R \to S$를 유한형 환 준동형사상이라 하자.
$\mathfrak p \subset R$를 극소 소 아이디얼이라 하자.
$\mathfrak p$ 위에 놓이는 $S$의 소 아이디얼이 유한 개 이하라고
가정하자. 그러면 $R_g \to S_g$가 유한이 되게 하는
$g \in R$, $g \not \in \mathfrak p$가 존재한다.
\end{lemma}

\begin{proof}
$x_1, \ldots, x_n$을 $R$ 위에서 $S$의 생성원들이라 하자.
보조정리 \ref{algebra-lemma-quasi-finite-above-p}에 의해 가정은
$S \otimes_R \kappa(\mathfrak p) = S_{\mathfrak p}/\mathfrak pS_{\mathfrak p}$가
유한 $\kappa(\mathfrak p)$-대수라는 뜻이다. 따라서
$S_{\mathfrak p}/\mathfrak pS_{\mathfrak p}$에서 $P_i(x_i)$가
영으로 사상되게 하는 모닉 다항식 $P_i \in R_{\mathfrak p}[X]$를
찾을 수 있다. $\mathfrak p$가 극소 소 아이디얼이므로
$\mathfrak pR_{\mathfrak p}$는 국소 멱영 아이디얼이다.
보조정리 \ref{algebra-lemma-minimal-prime-reduced-ring}을 보라. 따라서
$\mathfrak pS_{\mathfrak p}$도 국소 멱영 아이디얼이다.
보조정리 \ref{algebra-lemma-locally-nilpotent}를 보라. 그러므로
% P01-E1422
$S_{\mathfrak p}$에서 $P(x_i)^{e_i} = 0$이 되게 하는
$e_i \geq 1$들이 존재한다. 모든 $i$에 대하여 $P_i$의 계수들이
$R[1/g_1]$에 속하게 하는
$g_1 \in R$, $g_1 \not \in \mathfrak p$를 택하자. 다음으로
$S_{g_1g_2}$에서 $P(x_i)^{e_i} = 0$이 되게 하는
$g_2 \in R$, $g_2 \not \in \mathfrak p$를 택하자.
$g = g_1g_2$로 두면 결론을 얻는다.
\end{proof}

\section{Zariski의 주정리}
\label{algebra-section-Zariski}

\noindent
이 절의 목표는 Zariski의 주정리의 대수적 판본을 증명하는 것이다.
이 정리는 이 Stacks Project의 뒤쪽에서 전개할 스킴 및 스킴의
사상 이론에 관한 많은 추가 논의의 토대가 된다.

\medskip\noindent
$R \to S$를 유한형 환 준동형사상이라 하자. 이 절에서는 사상이
준유한인 $\Spec(S)$의 점들의 집합이 {\it 열린집합}임을 보이려 한다
(정리 \ref{algebra-theorem-main-theorem}). 실제로 어떤 의미에서
$S/R$의 준유한 자취가 $\Spec(S')$에서 열리도록 하는 유한 환
준동형사상 $R \to S'$이 존재한다는 것이 드러날 것이다
(그러나 여기서는 스킴의 언어를 전개하지 않으므로 대수학 장에서
이를 증명하지 않는다. $R \to S$가 준유한인 경우에는 보조정리
\ref{algebra-lemma-quasi-finite-open-integral-closure}를 보라).
이 명제들을 증명하는 일은 다소 까다롭고, 정수적 환확대와 유한
환확대에 관한 긴 보조정리 목록을 통하여 증명한다. 이 내용은
\cite{Henselian}과 \cite{Peskine}에서 찾을 수 있다. Thierry
Coquand의 노트도 도움이 되었다.

\begin{lemma}
\label{algebra-lemma-make-integral-trivial}
$\varphi : R \to S$를 환 준동형사상이라 하자.
$t \in S$가 관계
$\varphi(a_0) + \varphi(a_1)t + \ldots + \varphi(a_n) t^n = 0$을
만족한다고 하자. 그러면 $\varphi(a_n)t$는 $R$ 위에서 정수적이다.
\end{lemma}

\begin{proof}
실제로 등식
$$
\varphi(a_0) + \varphi(a_1)t + \ldots + \varphi(a_n) t^n = 0
$$
에
$\varphi(a_n)^{n-1}$을 곱하고
$$
\varphi(a_0 a_n^{n-1}) +
\varphi(a_1 a_n^{n-2}) (\varphi(a_n)t) +
\ldots +
(\varphi(a_n) t)^n = 0
$$
으로 쓰면 된다.
\end{proof}

\noindent
다음 보조정리는 어떤 의미에서 이 절의 핵심 보조정리이다.

\begin{lemma}
\label{algebra-lemma-make-integral-trick}
$R$을 환이라 하자. $\varphi : R[x] \to S$를 환 준동형사상이라
하자. $t \in S$라 하자. (a) $t$가 $R[x]$ 위에서 정수적이고,
(b) $t \varphi(p) \in \Im(\varphi)$가 되는 모닉 다항식
$p \in R[x]$가 존재한다고 가정하자. 그러면
$t - \varphi(q)$가 $R$ 위에서 정수적이 되게 하는
$q \in R[x]$가 존재한다.
\end{lemma}

\begin{proof}
어떤 $r \in R[x]$에 대하여 $t \varphi(p) = \varphi(r)$라고 쓰자.
유클리드 나눗셈을 사용하여 $q, r' \in R[x]$이고
$\deg(r') < \deg(p)$이도록 $r = qp + r'$라고 쓰자.
$t$를 여전히 $R[x]$ 위에서 정수적인 $t - \varphi(q)$로
바꾸어도 되며, 그러면 $t \varphi(p) = \varphi(r')$을 얻는다.
환 $S_t$에서는 이를
$\varphi(p) - (1/t) \varphi(r') = 0$으로 쓸 수 있다.
이는 $\varphi(x)$가 국소화 $S_t$의 원소이며 부분환
$\varphi(R)[1/t] \subset S_t$ 위에서 정수적임을 함의한다.
한편 $t$는 부분환 $\varphi(R)[\varphi(x)] \subset S$ 위에서
정수적이다. 두 사실을 합치고 보조정리
\ref{algebra-lemma-integral-transitive}를 사용하면 $t$가 부분환
$\varphi(R)[1/t] \subset S_t$ 위에서 정수적이라고 결론 내린다.
다시 말해 $r_{i, j} \in R$인 다음 꼴의 등식이 $S_t$에서
성립한다.
$$
t^d + \sum\nolimits_{i < d}
\left(\sum\nolimits_{j = 0, \ldots, n_i} \varphi(r_{i, j})/t^j\right) t^i = 0
$$
이는 충분히 큰 어떤 $N$에 대하여 $S$에서
$t^{d + N} +
\sum_{i < d} \sum_{j = 0, \ldots, n_i} \varphi(r_{i, j}) t^{i + N - j} = 0$임을
뜻한다. 다시 말해 $t$는 $R$ 위에서 정수적이다.
\end{proof}

\begin{lemma}
\label{algebra-lemma-combine-lemmas}
$R$을 환이라 하자. $\varphi : R[x] \to S$를 환 준동형사상이라
하자. $t \in S$라 하자. $t$가 $R[x]$ 위에서 정수적이라고 가정하자.
$p \in R[x]$, $p = a_0 + a_1x + \ldots +
a_k x^k$이고 $t \varphi(p) \in \Im(\varphi)$라고 하자.
그러면 $\varphi(a_k)^n t - \varphi(q)$가 $R$ 위에서 정수적이 되게
하는 $q \in R[x]$와 $n \geq 0$이 존재한다.
\end{lemma}

\begin{proof}
$R'$과 $S'$을 각각 $a_k$에서 $R$과 $S$를 국소화한 환이라 하자.
$\varphi' : R'[x] \to S'$을 $\varphi$의 국소화라 하자.
$t' \in S'$을 $t$의 상이라 하자. $p' = p/a_k \in R'[x]$로 두자.
$t \varphi(p) \in \Im(\varphi)$이므로
$t' \varphi'(p') \in \Im(\varphi')$이다. $p'$은 모닉이므로
보조정리 \ref{algebra-lemma-make-integral-trick}에 의해
$t' - \varphi'(q')$가 $R'$ 위에서 정수적이 되게 하는
$q' \in R'[x]$가 존재한다. $a_k^n q'$이 어떤
$q \in R[x]$의 상이 되게 하는 $n \geq 0$과 $q$를 택할 수 있다.
그러면 $\varphi(a_k)^n t - \varphi(q)$는 $S$의 원소이고,
$S'$에서의 그 상은 $R'$ 위에서 정수적이다. 보조정리
\ref{algebra-lemma-integral-closure-localize}에 의해
$\varphi(a_k)^m(\varphi(a_k)^n t - \varphi(q))$가 $R$ 위에서
정수적이 되게 하는 $m \geq 0$이 존재한다. 따라서 원하는 대로
$\varphi(a_k)^{m + n}t - \varphi(a_k^m q)$는 $R$ 위에서
정수적이다.
\end{proof}

\begin{situation}
\label{algebra-situation-one-transcendental-element}
$R$을 환이라 하자. $\varphi : R[x] \to S$를 유한 환
준동형사상이라 하자.
$$
J = \{ g \in S \mid gS \subset \Im(\varphi)\}
$$
를 $\varphi$의 ‘도체 아이디얼’이라 하자.
$\varphi(R) \subset S$가 $S$에서 정수적으로 닫혀 있다고 가정하자.
\end{situation}

\begin{lemma}
\label{algebra-lemma-leading-coefficient-in-J}
상황 \ref{algebra-situation-one-transcendental-element}에 있다고 하자.
$u \in S$, $a_0, \ldots, a_k \in R$이고
$u \varphi(a_0 + a_1x + \ldots + a_k x^k) \in J$라고 하자.
그러면 $u \varphi(a_k)^m \in J$가 되게 하는 $m \geq 0$이 존재한다.
\end{lemma}

\begin{proof}
$S$가 $R[x]$-가군으로서 $t_1, \ldots, t_n$에 의해 생성된다고
가정하자. 이 경우
$J = \{ g \in S \mid gt_i \in \Im(\varphi)\text{ for all }i\}$이다.
각 원소 $u t_i$가 $R[x]$ 위에서 정수적임을 주목하자.
보조정리 \ref{algebra-lemma-finite-is-integral}을 보라.
또한
$\varphi(a_0 + a_1x + \ldots + a_k x^k) u t_i \in
\Im(\varphi)$이다. 보조정리 \ref{algebra-lemma-combine-lemmas}에 의해
각 $i$에 대하여
$\varphi(a_k^{n_i}) u t_i - \varphi(q_i)$가 $R$ 위에서
정수적이 되게 하는 정수 $n_i$와 원소 $q_i \in R[x]$가 존재한다.
가정에 의해 이 원소는 $\varphi(R)$에 속하므로
$\varphi(a_k^{n_i}) u t_i \in \Im(\varphi)$이다.
따라서 $m = \max\{n_1, \ldots, n_n\}$을 택하면 된다.
\end{proof}

\begin{lemma}
\label{algebra-lemma-all-coefficients-in-J}
상황 \ref{algebra-situation-one-transcendental-element}에 있다고 하자.
$u \in S$, $a_0, \ldots, a_k \in R$이고
$u \varphi(a_0 + a_1x + \ldots + a_k x^k) \in \sqrt{J}$라고 하자.
그러면 모든 $i$에 대하여 $u \varphi(a_i) \in \sqrt{J}$이다.
\end{lemma}

\begin{proof}
보조정리의 가정 아래에서 어떤 $n \geq 1$에 대하여
$u^n \varphi(a_0 + a_1x + \ldots + a_k x^k)^n \in J$이다.
보조정리 \ref{algebra-lemma-leading-coefficient-in-J}에 의해 어떤
$m \geq 1$에 대하여 $u^n \varphi(a_k^{nm}) \in J$임을 얻는다.
따라서 $u \varphi(a_k) \in \sqrt{J}$이고, 그러므로
$u \varphi(a_0 + a_1x + \ldots + a_k x^k) - u \varphi(a_k x^k) =
u \varphi(a_0 + a_1x + \ldots + a_{k-1} x^{k-1}) \in \sqrt{J}$이다.
$k$에 관한 귀납법으로 결론을 얻는다.
\end{proof}

\noindent
이 보조정리는 다음 정의를 시사한다.

\begin{definition}
\label{algebra-definition-strongly-transcendental}
환들의 포함 $R \subset S$와 원소 $x \in S$가 주어졌다고 하자.
$u \in S$와 $a_i \in R$에 대하여
$u(a_0 + a_1 x + \ldots + a_k x^k) = 0$일 때마다 모든 $i$에 대하여
$ua_i = 0$이면 $x$가 {\it $R$ 위에서 강하게 초월적}이라고 한다.
\end{definition}

\noindent
$S$가 정역이면 이는 $S$의 분수체의 원소 $x$가 $R$의 분수체
위에서 초월적이라고 말하는 것과 같다.

\begin{lemma}
\label{algebra-lemma-reduced-strongly-transcendental-minimal-prime}
$R \subset S$가 축소환들의 포함이고 $x \in S$가 $R$ 위에서
강하게 초월적이라고 하자. $\mathfrak q \subset S$를 극소 소
아이디얼이라 하고 $\mathfrak p = R \cap \mathfrak q$라 하자.
그러면 $S/\mathfrak q$에서 $x$의 상은 부분환
$R/\mathfrak p$ 위에서 강하게 초월적이다.
\end{lemma}

\begin{proof}
$u(a_0 + a_1x + \ldots + a_k x^k) \in \mathfrak q$라고 하자.
보조정리 \ref{algebra-lemma-minimal-prime-reduced-ring}에 의해
국소환 $S_{\mathfrak q}$는 체이고, 따라서
$u(a_0 + a_1x + \ldots + a_k x^k) $는 $S_{\mathfrak q}$에서
영이다. 그러므로 어떤 $u' \in S$, $u' \not\in \mathfrak q$에 대하여
$uu'(a_0 + a_1x + \ldots + a_k x^k) = 0$이다.
$x$가 $R$ 위에서 강하게 초월적이므로 모든 $i$에 대하여
$uu'a_i = 0$을 얻는다. 이는 다시 $ua_i \in \mathfrak q$를
함의한다.
\end{proof}

\begin{lemma}
\label{algebra-lemma-domains-transcendental-not-quasi-finite}
$R\subset S$가 정역들의 포함이고 $x \in S$라 하자.
$x$가 $R$ 위에서 (강하게) 초월적이고 $S$가 $R[x]$ 위에서
유한하다고 가정하자. 그러면 $R \to S$는 $S$의 어느 소
아이디얼에서도 준유한이 아니다.
\end{lemma}

\begin{proof}
먼저 $R$이 정규라고 가정하자. 정의
\ref{algebra-definition-ring-normal}을 보라. 보조정리
\ref{algebra-lemma-polynomial-ring-normal}에 의해 $R[x]$는 정규이다.
소 아이디얼 $\mathfrak q \subset S$를 택하고
$\mathfrak p = R \cap \mathfrak q$로 두자. 확대
$\kappa(\mathfrak p) \subset \kappa(\mathfrak q)$가 유한하다고
가정하자. $R \to S$가 $\mathfrak q$에서 준유한이면 이 조건이
성립할 것이다. $\mathfrak r = R[x] \cap \mathfrak q$라 하자.
$\kappa(\mathfrak p)
\subset \kappa(\mathfrak r) \subset \kappa(\mathfrak q)$이므로
확대 $\kappa(\mathfrak p)
\subset \kappa(\mathfrak r)$도 유한함을 알 수 있다.
따라서 포함 $\mathfrak r \supset \mathfrak p R[x]$은 엄밀하다.
$R[x] \subset S$에 대한 하강 성질, 즉 명제
\ref{algebra-proposition-going-down-normal-integral}에 의해
$\mathfrak pR[x]$ 위에 놓이는 소 아이디얼
$\mathfrak q' \subset \mathfrak q$를 찾는다. 따라서 섬유
$\Spec(S \otimes_R \kappa(\mathfrak p))$에는 $\mathfrak q$와
다른 점, 즉 $\mathfrak q'$가 들어 있고, 그 폐포는
$\mathfrak q$를 포함한다. 그러므로 $\mathfrak q$는 그 섬유의
고립점이 아니다.

\medskip\noindent
$R$이 정규가 아니면, $R \subset R' \subset K$를 $R$의 분수체
$K$ 안에서의 정수적 폐포 $R'$이라 하자. $S \subset S' \subset L$을
$S$의 분수체 $L$에서 $R'$과 $S$로 생성되는 부분환 $S'$이라 하자.
구성에 의해
사상 $S \otimes_R R' \to S'$이 전사임을 주목하자.
이는 $R'[x] \subset S'$이 유한임을 함의한다. 또한
$S \subset S'$은 $\Spec$ 위의 전사를 유도한다. 보조정리
\ref{algebra-lemma-integral-overring-surjective}를 보라.
보조정리 \ref{algebra-lemma-four-rings}와 방금 논의한 정규인 경우에 의해
결론을 얻는다.
\end{proof}

\begin{lemma}
\label{algebra-lemma-reduced-strongly-transcendental-not-quasi-finite}
$R \subset S$가 축소환들의 포함이라고 하자.
% P01-E1423
$x \in S$가 $R$ 위에서 강하게 초월적이고 $S$가 $R[x]$ 위에서
유한하다고 가정하자. 그러면 $R \to S$는 $S$의 어느 소
아이디얼에서도 준유한이 아니다.
\end{lemma}

\begin{proof}
$\mathfrak q \subset S$를 임의의 소 아이디얼이라 하자.
극소 소 아이디얼 $\mathfrak q' \subset \mathfrak q$를 택하자.
보조정리
\ref{algebra-lemma-reduced-strongly-transcendental-minimal-prime}과
\ref{algebra-lemma-domains-transcendental-not-quasi-finite}에 의해
확대 $R/(R \cap \mathfrak q') \subset
S/\mathfrak q'$은 $\mathfrak q$에 대응하는 소 아이디얼에서
준유한이 아니다. 보조정리 \ref{algebra-lemma-four-rings}에 의해
확대 $R \to S$는 $\mathfrak q$에서 준유한이 아니다.
\end{proof}

\begin{lemma}
\label{algebra-lemma-quasi-finite-monogenic}
$R$을 환이라 하자. $S = R[x]/I$라 하자.
$\mathfrak q \subset S$를 소 아이디얼이라 하자.
$R \to S$가 $\mathfrak q$에서 준유한이라고 가정하자.
$S' \subset S$을 $S$에서 $R$의 정수적 폐포라 하자.
그러면 $S'_g \cong S_g$가 되게 하는
$g \in S'$, $g \not\in \mathfrak q$가 존재한다.
\end{lemma}

\begin{proof}
$\mathfrak p$를 $\Spec(R)$에서 $\mathfrak q$의 상이라 하자.
어떤 $i$에 대하여 $a_i \not \in \mathfrak p$인
$f \in I$, $f = a_nx^n + \ldots + a_0$가 존재한다.
실제로 그렇지 않으면 섬유환
$S \otimes_R \kappa(\mathfrak p)$가 $\kappa(\mathfrak p)[x]$이고
사상은 $\mathfrak p$ 위에 놓이는 어느 소 아이디얼에서도
준유한이지 않을 것이다. 따라서 $b_j \in S'$,
$j = 0, \ldots, m$이고 어떤 $j$에 대하여
$b_j \not \in \mathfrak q \cap S'$인 관계
$b_m x^m + \ldots + b_0 = 0$이 존재한다고 결론 내린다.
$m$에 관한 귀납법으로 보조정리를 증명한다.
% P01-E1424
기초 단계 $m = 0$은 공허하다(명제들 $b_0 = 0$과
$b_0 \not \in \mathfrak q$가 서로 모순이기 때문이다).

\medskip\noindent
$b_m \not \in \mathfrak q$인 경우. 이 경우 $x$는
$S'_{b_m}$ 위에서 정수적이고, 실제로 보조정리
\ref{algebra-lemma-make-integral-trivial}에 의해 $b_mx \in S'$이다.
따라서 단사사상 $S'_{b_m} \to S_{b_m}$은 전사이기도 하며,
즉 원하는 동형사상이다.

\medskip\noindent
$b_m \in \mathfrak q$인 경우. 이 경우에도 보조정리
\ref{algebra-lemma-make-integral-trivial}에 의해 $b_mx \in S'$이다.
$b'_{m - 1} = b_mx + b_{m - 1}$로 두자. 그러면
$$
b'_{m - 1}x^{m - 1} + b_{m - 2}x^{m - 2} + \ldots + b_0 = 0
$$
이다. $b'_{m - 1}$은 $S' \cap \mathfrak q$를 법으로
$b_{m - 1}$과 합동이므로
$b'_{m - 1}, b_{m - 2}, \ldots, b_0$ 중 하나는 여전히
$S' \cap \mathfrak q$에 속하지 않는다. 따라서 $m$에 관한
귀납법으로 결론을 얻는다.
\end{proof}

\begin{theorem}[Zariski의 주정리]
\label{algebra-theorem-main-theorem}
$R$을 환이라 하자.
% P01-E1425
$R \to S$를 유한형 $R$-대수라 하자.
$S' \subset S$을 $S$에서 $R$의 정수적 폐포라 하자.
$\mathfrak q \subset S$를 $S$의 소 아이디얼이라 하자.
$R \to S$가 $\mathfrak q$에서 준유한이면
$S'_g \cong S_g$가 되게 하는
$g \in S'$, $g \not \in \mathfrak q$가 존재한다.
\end{theorem}

\begin{proof}
$S$가 $x_1, \ldots, x_n$으로 생성되는 $R$-부분대수 위에서
유한이 되게 하는 유한 개의 원소
$x_1, \ldots, x_n \in S$가 존재한다. (예를 들어 $R$ 위에서
$S$의 생성원들을 택하라.)
% P01-E1426
그러한 최소 수 $n$에 관한 귀납법으로 명제를 증명한다.

\medskip\noindent
$n = 0$인 경우에는 $S' = S$이므로 자명하다.
보조정리 \ref{algebra-lemma-finite-is-integral}을 보라.

\medskip\noindent
$n = 1$인 경우. $R$을 $S$에서 그 정수적 폐포로 바꾸어도 된다
(보조정리 \ref{algebra-lemma-quasi-finite-permanence}는 $R \to S$가
여전히 $\mathfrak q$에서 준유한임을 보장한다).
따라서 $R \subset S$가 $S$에서 정수적으로 닫혀 있다고,
다시 말해 $R = S'$이라고 가정해도 된다.
사상 $\varphi : R[x] \to S$, $x \mapsto x_1$을 생각하자.
(아래에서 $\varphi$가 단사가 아님을 볼 것이다.)
가정에 의해 $\varphi$는 유한이다. 따라서 상황
\ref{algebra-situation-one-transcendental-element}에 있다.
$J \subset S$를 상황 \ref{algebra-situation-one-transcendental-element}에서
정의한 ‘도체 아이디얼’이라 하자. 그림
$$
\xymatrix{
R[x] \ar[r] & S \ar[r] & S/\sqrt{J} & R/(R \cap \sqrt{J})[x] \ar[l]
\\
& R \ar[lu] \ar[r] \ar[u] & R/(R \cap \sqrt{J}) \ar[u] \ar[ru] &
}
$$
을 생각하자. 보조정리 \ref{algebra-lemma-all-coefficients-in-J}에 의해
몫 $S/\sqrt{J}$에서 $x$의 상은
$R/ (R \cap \sqrt{J})$ 위에서 강하게 초월적이다. 따라서 보조정리
\ref{algebra-lemma-reduced-strongly-transcendental-not-quasi-finite}에 의해
환 준동형사상 $R/ (R \cap \sqrt{J}) \to S/\sqrt{J}$은
$S/\sqrt{J}$의 어느 소 아이디얼에서도 준유한이 아니다.
보조정리 \ref{algebra-lemma-four-rings}에 의해
$\mathfrak q$가 $V(J) \subset \Spec(S)$에 속하지 않는다고
결론 내린다. 따라서 $s \in J$,
$s \not\in \mathfrak q$인 원소가 존재한다. $J$의 정의에 의해
어떤 다항식 $f \in R[x]$에 대하여 $s = \varphi(f)$라고 쓸 수 있다.
$I = \Ker(\varphi : R[x] \to S)$라 하자. $\varphi(f) \in J$이므로
$(R[x]/I)_f \cong S_{\varphi(f)}$를 얻는다. 또한
$s \not \in \mathfrak q$라는 것은
$f \not \in \varphi^{-1}(\mathfrak q)$라는 뜻이다. 따라서
$\varphi^{-1}(\mathfrak q)/I$는 $R \to R[x]/I$가 준유한인
$R[x]/I$의 소 아이디얼이다. 보조정리
\ref{algebra-lemma-quasi-finite-local}을 보라.
$R$이 $S$에서 정수적으로 닫혀 있으므로 $R$은 $R[x]/I$에서도
정수적으로 닫혀 있음을 주목하자. 보조정리
\ref{algebra-lemma-quasi-finite-monogenic}에 의해
$R_h \cong (R[x]/I)_h$가 되게 하는
$h \in R$, $h \not \in R \cap \mathfrak q$가 존재한다.
따라서 $(R[x]/I)_{fh} = S_{\varphi(fh)}$는 어떤
$h' \in R$, $h' \not \in \mathfrak q$에서 $R$을 주국소화한
$R_{h'}$과 동형이다.

\medskip\noindent
$n > 1$인 경우. $R' \subset S$을 $S$에서
$R[x_1, \ldots, x_{n-1}]$의 정수적 폐포라 하자.
보조정리 \ref{algebra-lemma-quasi-finite-permanence}에 의해 확대
$S/R'$은 $\mathfrak q$에서 준유한이다. 또한 $S$가
$R'[x_n]$ 위에서 유한함을 주목하자. 위의 $n = 1$인 경우에 의해
$(R')_{g'} \cong S_{g'}$가 되게 하는
$g' \in R'$, $g' \not \in \mathfrak q$가 존재한다.
이 단계에서는 $R'$이 $R$ 위에서 유한형일 필요가 없으므로
$R \to R'$에 귀납가정을 적용할 수 없다. $S$가 $R$ 위에서
유한 생성이므로 특히 $(R')_{g'}$가 $R$ 위에서 유한 생성임을
얻는다. $y_i \in R'$이고, 원소
$g'$, 그리고 $y_1/(g')^{n_1}, \ldots, y_N/(g')^{n_N}$이
$R$ 위에서 $(R')_{g'}$을 생성한다고 하자.
$R''$을 $x_1, \ldots, x_{n-1}, y_1, \ldots, y_N, g'$으로
생성되는 $R'$의 $R$-부분대수라 하자. 이는
$(R'')_{g'} \cong S_{g'}$라는 성질을 갖는다.
전사성은 $y_i$들을 택한 방식에서, 단사성은
$R'' \subset R'$이라는 사실과 국소화의 완전성에서 따른다.
$R'$의 선택과 보조정리 \ref{algebra-lemma-characterize-integral}에 의해
$R''$은 $R[x_1, \ldots, x_{n-1}]$ 위에서 유한임을 주목하자.
$\mathfrak q'' = R'' \cap \mathfrak q$라 하자.
$(R'')_{\mathfrak q''} = S_{\mathfrak q}$이므로 보조정리
\ref{algebra-lemma-isolated-point-fibre}에 의해 $R \to R''$은
$\mathfrak q''$에서 준유한이다. $R \to R''$,
$\mathfrak q''$, 그리고 $x_1, \ldots, x_{n-1} \in R''$에
귀납가정을 적용한다. 그러면 $R$ 위에서 정수적인 부분환
$R''' \subset R''$과, $(R''')_{g''} \cong (R'')_{g''}$이 되게 하는
$g'' \in R'''$, $g'' \not \in \mathfrak q''$을 얻는다.
$(R'')_{g''}$에서 $g'$의 상을 어떤 $g''' \in  R'''$에 대하여
$g'''/(g'')^n$으로 쓰자.
$g = g''g''' \in R'''$로 두자. 그러면
$g \not\in \mathfrak q$이고 $(R''')_g \cong S_g$임이 분명하다.
구성에 의해 $R''' \subset S'$이므로 원하는 대로
$S'_g \cong S_g$이기도 하다.
\end{proof}

\begin{lemma}
\label{algebra-lemma-quasi-finite-open}
$R \to S$를 유한형 환 준동형사상이라 하자.
$S/R$이 준유한인 $\Spec(S)$의 점들 $\mathfrak q$의 집합은
$\Spec(S)$에서 열려 있다.
\end{lemma}

\begin{proof}
$\mathfrak q \subset S$를 환 준동형사상이 준유한인 점이라 하자.
정리 \ref{algebra-theorem-main-theorem}에 의해 정수적 환 준동형사상
$R \to S'$, $S' \subset S$과
$S'_g \cong S_g$가 되게 하는 $g \in S'$,
$g\not \in \mathfrak q$가 존재한다. $S$, 따라서 $S_g$가
$R$ 위에서 유한형이므로, $S'$의 유한 개의 원소
$y_1, \ldots, y_N$을 찾아 $g, y_1, \ldots, y_N$으로 생성되는
$R$-부분대수 $S'' \subset S'$에 대해 $S''_g \cong S_g$가 되게
할 수 있다. $S''$은 $R$ 위에서
유한하므로(보조정리 \ref{algebra-lemma-characterize-integral}을 보라)
$S''$은 $R$ 위에서 준유한이다(보조정리
\ref{algebra-lemma-quasi-finite}를 보라). 이는 $S''_g$가 $R$ 위에서
준유한임을 함의함을 쉽게 알 수 있다. 예를 들어 한 소
아이디얼에서 준유한이라는 성질은 그 소 아이디얼에서의 국소환에만
의존하기 때문이다. 따라서 $S_g$가 $R$ 위에서 준유한임을 안다.
마찬가지로 이는 $R \to S$가 $D(g)$에 속하는 $S$의 모든 소
아이디얼에서 준유한임을 함의한다.
\end{proof}

\begin{lemma}
\label{algebra-lemma-quasi-finite-open-integral-closure}
$R \to S$를 유한형 환 준동형사상이라 하자.
$S$가 $R$ 위에서 준유한이라고 가정하자.
$S' \subset S$을 $S$에서 $R$의 정수적 폐포라 하자. 그러면
\begin{enumerate}
\item $\Spec(S) \to \Spec(S')$은 한 열린 부분집합 위로의
위상동형사상이다.
\item $g \in S'$이고 $D(g)$가 이 사상의 상에 포함되면
$S'_g \cong S_g$이다.
\item 환 준동형사상 $S'' \to S$에 대하여 (1)과 (2)가 성립하도록
하는 유한 $R$-대수 $S'' \subset S'$이 존재한다.
\end{enumerate}
\end{lemma}

\begin{proof}
$S/R$이 준유한이므로 $\Spec(S)$의 각 점
$\mathfrak q$에 정리 \ref{algebra-theorem-main-theorem}을 적용할 수 있다.
보조정리 \ref{algebra-lemma-quasi-compact}에 의해 $\Spec(S)$는
준콤팩트이므로, $S'_{g_i} = S_{g_i}$이고
$g_1, \ldots, g_n$이 $S$의 단위 아이디얼을 생성하도록 하는
유한 개의 $g_i \in S'$, $i = 1, \ldots, n$을 택할 수 있다
(다시 말해 $g_1, \ldots, g_n$에 대응하는 $\Spec(S)$의 표준
열린집합들이 $\Spec(S)$ 전체를 덮는다).

\medskip\noindent
$D(g) \subset \Spec(S')$이 이 사상의 상에 포함된다고 하자.
그러면 $D(g) \subset \bigcup D(g_i)$이다. 다시 말해
$g_1, \ldots, g_n$은 $S'_g$의 단위 아이디얼을 생성한다.
$g_i$를 택한 방식에 의해 $S'_{gg_i} \cong S_{gg_i}$임을
주목하자. 따라서 보조정리 \ref{algebra-lemma-cover}에 의해
$S'_g \cong S_g$이다.

\medskip\noindent
(3)의 유한 대수 $S'' \subset S'$을 구성하자.
이를 위해 각 $S'_{g_i} \cong S_{g_i}$가 유한형 $R$-대수임을
주목하자. 각 $i$에 대하여 $S'_{g_i}$가 $1/g_i$와 원소들
$y_{ij}$로 $R$-대수로서 생성되게 하는 원소
$y_{ij} \in S'$을 택하자. 이제 모든 $g_i$와 모든 $y_{ij}$로
생성되는 $S'$의 $R$-부분대수를 $S''$으로 둔다. 세부사항은
생략한다.
\end{proof}

\section{Zariski의 주정리의 응용}
\label{algebra-section-apply-main-theorem}

\noindent
다음은 $1$차원 반국소환들 사이의 준유한 사상 중 유한인 것들을
보조정리 \ref{algebra-lemma-finite-extension-dim-1}의 공식에서 언제나
등호가 성립하는 것으로 특징짓는 즉각적인 응용이다.

\begin{lemma}
\label{algebra-lemma-quasi-finite-extension-dim-1}
$A \subset B$를 정역들의 확대라 하자. 다음을 가정하자.
\begin{enumerate}
\item $A$는 차원 $1$의 뇌터 국소환이다.
\item $A \to B$는 유한형이다.
\item 유도된 분수체 확대 $L/K$는 유한하다.
\end{enumerate}
그러면 $B$는 반국소이다. $x \in \mathfrak m_A$,
$x \not = 0$이라 하자. $\mathfrak m_i$,
$i = 1, \ldots, n$을 $B$의 극대 아이디얼들이라 하자. 그러면
$$
[L : K]\text{ord}_A(x)
\geq
\sum\nolimits_i
[\kappa(\mathfrak m_i) : \kappa(\mathfrak m_A)]
\text{ord}_{B_{\mathfrak m_i}}(x)
$$
이다. 여기서 $\text{ord}$는 정의 \ref{algebra-definition-ord}에서처럼
정의한다. 등호가 성립할 필요충분조건은 $A \to B$가 유한인 것이다.
\end{lemma}

\begin{proof}
보조정리 \ref{algebra-lemma-finite-in-codim-1}에 의해 환 $B$는 반국소이다.
$B'$을 $B$에서 $A$의 정수적 폐포라 하자. 보조정리
\ref{algebra-lemma-quasi-finite-open-integral-closure}에 의해
$\mathfrak n_i = C \cap \mathfrak m_i$로 두면
$C_{\mathfrak n_i} \cong B_{\mathfrak m_i}$이고 소 아이디얼들
$\mathfrak n_1, \ldots, \mathfrak n_n$이 쌍마다 서로 다르도록 하는
유한 $A$-부분대수 $C \subset B'$을 찾을 수 있다.
보조정리 \ref{algebra-lemma-finite-in-codim-1}에 의해 환 $C$는 반국소이다.
$\mathfrak p_j$, $j = 1, \ldots, m$을 $C$의 나머지 극대
아이디얼들(‘빠진 점들’)이라 하자. 보조정리
\ref{algebra-lemma-finite-extension-dim-1}에 의해
$$
\text{ord}_A(x^{[L : K]}) =
\sum\nolimits_i
[\kappa(\mathfrak n_i) : \kappa(\mathfrak m_A)]
\text{ord}_{C_{\mathfrak n_i}}(x)
+
\sum\nolimits_j
[\kappa(\mathfrak p_j) : \kappa(\mathfrak m_A)]
\text{ord}_{C_{\mathfrak p_j}}(x)
$$
이고, 따라서 부등식이 따른다. 등호가 성립하면
$m = 0$이라고(‘빠진 점’이 없다고) 결론 내린다. 그러므로
$C \subset B$는 극대 아이디얼들 위의 전단사와 모든 극대
아이디얼에서의 국소화들 사이의 동형사상을 유도하는 반국소환들의
포함이다. 따라서 $b \in B$이면
$I = \{x \in C \mid xb \in C\}$는 $C$의 어느 극대 아이디얼에도
포함되지 않는 $C$의 아이디얼이다. 그러므로 $I = C$이고
$b \in C$이다. 따라서 $B = C$이고 $B$는 $A$ 위에서 유한하다.
\end{proof}

\noindent
다음은 국소환들의 국소 준동형사상 구조에 대한 Zariski의 주정리의
더 표준적인 응용이다.

\begin{lemma}
\label{algebra-lemma-essentially-finite-type-fibre-dim-zero}
$(R, \mathfrak m_R) \to (S, \mathfrak m_S)$를 국소환들의 국소
준동형사상이라 하자. 다음을 가정하자.
\begin{enumerate}
\item $R \to S$는 본질적으로 유한형이다.
\item $\kappa(\mathfrak m_R) \subset \kappa(\mathfrak m_S)$는
유한하다.
\item $\dim(S/\mathfrak m_RS) = 0$이다.
\end{enumerate}
그러면 $S$는 유한 $R$-대수의 국소화이다.
\end{lemma}

\begin{proof}
$S = S'_{\mathfrak q'}$가 되는 유한형 $R$-대수 $S'$과
$S'$의 소 아이디얼 $\mathfrak q'$을 택하자. 정의
\ref{algebra-definition-quasi-finite}에 의해 $R \to S'$은
$\mathfrak q'$에서 준유한이다. 어떤
$g' \in S'$, $g' \not \in \mathfrak q'$에 대하여 $S'$을
$S'_{g'}$로 바꾼 뒤, 보조정리 \ref{algebra-lemma-quasi-finite-open}에 의해
$R \to S'$이 준유한이라고 가정해도 된다. 그러면 보조정리
\ref{algebra-lemma-quasi-finite-open-integral-closure}에 의해
$R$-대수로서 $S'_{g'} \cong S''_{g''}$이 되게 하는 유한
$R$-대수 $S''$과 원소들
$g' \in S'$, $g' \not \in \mathfrak q'$ 및 $g'' \in S''$이
존재한다. 이로써 보조정리가 증명된다.
\end{proof}

\begin{lemma}
\label{algebra-lemma-completion-at-quasi-finite-prime}
$R \to S$를 환 준동형사상이라 하고, $\mathfrak q$를
$R$의 $\mathfrak p$ 위에 놓이는 $S$의 소 아이디얼이라 하자.
다음을 가정하자.
\begin{enumerate}
\item $R$은 뇌터이다.
\item $R \to S$는 유한형이다.
\item $R \to S$는 $\mathfrak q$에서 준유한이다.
\end{enumerate}
그러면 어떤 $R_\mathfrak p^\wedge$-대수 $B$에 대하여
$R_\mathfrak p^\wedge \otimes_R S = S_\mathfrak q^\wedge \times B$이다.
\end{lemma}

\begin{proof}
$S'_g = S_g$이고 특히
$S'_{\mathfrak q'} = S_\mathfrak q$가 되게 하는 유한
$R$-대수 $S' \subset S$과
$g \in S'$, $g \not \in \mathfrak q' = S' \cap \mathfrak q$가
존재한다. 보조정리
\ref{algebra-lemma-quasi-finite-open-integral-closure}를 보라.
보조정리 \ref{algebra-lemma-completion-finite-extension}에 의해
$$
R_\mathfrak p^\wedge \otimes_R S' = (S'_{\mathfrak q'})^\wedge \times B'
$$
이다. 이 곱분해 아래에서 $g$는 쌍 $(u, b')$으로 사상되는데,
$g \not \in \mathfrak q'$이므로
$u \in (S'_{\mathfrak q'})^\wedge$는 단원임을 관찰하자.
$R_\mathfrak p^\wedge \otimes_R S'$의 곱분해는 곱분해
$$
R_\mathfrak p^\wedge \otimes_R S = A \times B
$$
를 유도한다. $S'_g = S_g$이므로
$(R_\mathfrak p^\wedge \otimes_R S')_g = (R_\mathfrak p^\wedge \otimes_R S)_g$
이기도 하다. $g \mapsto (u, b')$이고 $u$가 단원이므로
$(S'_{\mathfrak q'})^\wedge = A$임을 알 수 있다.
동형사상 $S'_{\mathfrak q'} = S_\mathfrak q$는 완비화들 사이의
동형사상을 결정하므로 $A = S_\mathfrak q^\wedge$이기도 하다.
이로써 증명이 끝나지만, 유도된 사상
$\phi : R_\mathfrak p^\wedge \otimes_R S \to A = S_\mathfrak q^\wedge$이
자연스러운 사상인지 건전성 검사를 해야 한다.
$x \in R_\mathfrak p^\wedge$인 $x \otimes 1$ 꼴의 원소에 대해서는
자연스러운 사상 $R_\mathfrak p^\wedge \to S_\mathfrak q^\wedge$이
$(S'_{\mathfrak q'})^\wedge$을 통하여 인수분해되므로 분명하다.
$y \in S$인 $1 \otimes y$ 꼴의 원소에 대해서는, 어떤
$n \geq 1$에 대하여 원소 $g^ny$가 어떤 $y' \in S'$의 상이라고
논증할 수 있다. 따라서 $\phi(1 \otimes g^ny)$는 합성
$S' \to (S'_{\mathfrak q'})^\wedge \to S_\mathfrak q^\wedge$에
의한 $y'$의 상이고, 이는 사상 $S \to S_\mathfrak q^\wedge$에
의한 $g^ny$의 상과 같다. $g$가 단원으로 사상되므로
$\phi(1 \otimes y)$도 올바른 값, 즉
$S \to S_\mathfrak q^\wedge$에 의한 $y$의 상을 갖는다.
\end{proof}

\section{섬유의 차원}
\label{algebra-section-dimension-fibres}

\noindent
Zariski의 주정리를 이용하여 섬유들의 차원이 어떻게 변하는지
연구한다. 위상공간 $X$의 점 $x$에서의 차원 $\dim_x(X)$을
Topology, 정의 \ref{topology-definition-Krull}에서 정의했음을
상기하자.

\begin{definition}
\label{algebra-definition-relative-dimension}
$R \to S$가 유한형이고, $\mathfrak q \subset S$가 $R$의 소
아이디얼 $\mathfrak p$ 위에 놓이는 소 아이디얼이라고 하자.
$\mathfrak q$에서의 {\it $S/R$의 상대 차원}을
$\dim_{\mathfrak q}(S/R)$로 나타내며, 이를 $\mathfrak q$에
대응하는 점에서의
$\Spec(S \otimes_R \kappa(\mathfrak p))$의 차원으로 정의한다.
모든 $\mathfrak q$에 대한 $\dim_{\mathfrak q}(S/R)$의 상한을
$\dim(S/R)$로 둔다. 이를 {\it $S/R$의 상대 차원}이라고 한다.
\end{definition}

\noindent
특히 $R \to S$가 $\mathfrak q$에서 준유한일 필요충분조건은
$\dim_{\mathfrak q}(S/R) = 0$인 것이다. 다음 보조정리는
Zariski의 주정리를 다소 다르게 서술한 것이다.

\begin{lemma}
\label{algebra-lemma-quasi-finite-over-polynomial-algebra}
$R \to S$를 유한형 환 준동형사상이라 하고,
$\mathfrak q \subset S$를 소 아이디얼이라 하자.
$\dim_{\mathfrak q}(S/R) = n$이라고 가정하자.
$S_g$가 다항식 대수 $R[t_1, \ldots, t_n]$ 위에서 준유한이
되게 하는 $g \in S$, $g \not\in \mathfrak q$가 존재한다.
\end{lemma}

\begin{proof}
환 $\overline{S} = S \otimes_R \kappa(\mathfrak p)$는
$\kappa(\mathfrak p)$ 위에서 유한형이다.
$\overline{\mathfrak q}$를 $\mathfrak q$에 대응하는
$\overline{S}$의 소 아이디얼이라 하자. 한 점에서의 위상공간의
차원의 정의에 의해 $\dim(U) = n$인 열린집합
$U \subset \Spec(\overline{S})$가 존재한다.
% P01-E1427
여기서 $\overline{q} \in U$이다.
$\Spec(\overline{S})$의 위상은 $\Spec(S)$의 위상에서
유도되므로(주의 \ref{algebra-remark-fundamental-diagram}을 보라), 상
$\overline{g} \in \overline{S}$가
$D(\overline{g}) \subset U$를 만족하게 하는
$g \in S$, $g \not \in \mathfrak q$를 찾을 수 있다.
따라서 $S$를 $S_g$로 바꾸고 나면
$\dim(\overline{S}) = n$이라고 해도 된다.

\medskip\noindent
다음으로 $R$-대수 $S$의 생성원 $x_1, \ldots, x_N$을 택하자.
보조정리 \ref{algebra-lemma-Noether-normalization}에 의해
$x_1, \ldots, x_N$으로 생성되는 $S$의 $\mathbf{Z}$-부분대수
안에 원소 $y_1, \ldots, y_n$이 존재하여, 사상
$R[t_1, \ldots, t_n] \to S$, $t_i \mapsto y_i$는
% P01-E1428
$\kappa(\mathfrak p)[t_1\ldots, t_n] \to \overline{S}$가
유한이라는 성질을 갖는다. 특히 $S$는 $\mathfrak q$에서
$R[t_1, \ldots, t_n]$ 위에 준유한이다. 따라서 보조정리
\ref{algebra-lemma-quasi-finite-open}에 의해 $S$를 어떤
$g\in S$, $g \not \in \mathfrak q$에 대한 $S_g$로 바꾸어
$R[t_1, \ldots, t_n] \to S$가 준유한이 되게 할 수 있다.
\end{proof}

\begin{lemma}
\label{algebra-lemma-refined-quasi-finite-over-polynomial-algebra}
$R \to S$를 환 준동형사상이라 하자. $\mathfrak q \subset S$가
$R$의 소 아이디얼 $\mathfrak p$ 위에 놓이는 소 아이디얼이라
하자. 다음을 가정하자.
\begin{enumerate}
\item $R \to S$는 유한형이고,
\item $\dim_{\mathfrak q}(S/R) = n$이며,
\item $\text{trdeg}_{\kappa(\mathfrak p)}\kappa(\mathfrak q) = r$이다.
\end{enumerate}
그러면 $f \in R$, $f \not \in \mathfrak p$,
$g \in S$, $g \not\in \mathfrak q$와 준유한 환 준동형사상
$$
\varphi : R_f[x_1, \ldots, x_n] \longrightarrow S_g
$$
가 존재하여
$\varphi^{-1}(\mathfrak qS_g) =
(\mathfrak p, x_{r + 1}, \ldots, x_n)R_f[x_1, \ldots, x_n]$이다.
% P01-E1429
\end{lemma}

\begin{proof}
$S$를 주 국소화로 바꾸어, 준유한 환 준동형사상
$\varphi : R[t_1, \ldots, t_n] \to S$가 존재한다고 해도 된다.
보조정리 \ref{algebra-lemma-quasi-finite-over-polynomial-algebra}를 보라.
$\mathfrak q' = \varphi^{-1}(\mathfrak q)$로 두자.
$\overline{\mathfrak q}' \subset \kappa(\mathfrak p)[t_1, \ldots, t_n]$을
$\mathfrak q'$에 대응하는 소 아이디얼이라 하자. 보조정리
\ref{algebra-lemma-refined-Noether-normalization}에 의해
$\overline{\mathfrak q}'$의 역상이
$(x_{r + 1}, \ldots, x_n)$인 유한 환 준동형사상
$\kappa(\mathfrak p)[x_1, \ldots, x_n] \to
\kappa(\mathfrak p)[t_1, \ldots, t_n]$이 존재한다.
$\overline{h}_i \in \kappa(\mathfrak p)[t_1, \ldots, t_n]$를
$x_i$의 상이라 하자. $\kappa(\mathfrak p)[t_1, \ldots, t_n]$에서
$\overline{h}_i$로 사상되는
$h_i \in R_f[t_1, \ldots, t_n]$와
$f \in R$, $f \not \in \mathfrak p$를 찾을 수 있다.
그러면 환 준동형사상
$$
R_f[x_1, \ldots, x_n] \longrightarrow R_f[t_1, \ldots, t_n]
$$
은 $\kappa(\mathfrak p)$와 텐서곱한 뒤 유한이 된다. 특히
$R_f[t_1, \ldots, t_n]$은 소 아이디얼
$\mathfrak q'R_f[t_1, \ldots, t_n]$에서
$R_f[x_1, \ldots, x_n]$ 위에 준유한이다. 따라서 보조정리
\ref{algebra-lemma-quasi-finite-open}에 의해
$g \in R_f[t_1, \ldots, t_n]$,
$g \not \in \mathfrak q'R_f[t_1, \ldots, t_n]$가 존재하여
$R_f[x_1, \ldots, x_n] \to R_f[t_1, \ldots, t_n, 1/g]$는
준유한이다. 따라서 합성
$$
R_f[x_1, \ldots, x_n] \longrightarrow
R_f[t_1, \ldots, t_n, 1/g] \longrightarrow S_{\varphi(g)}
$$
은 준유한이고 결론을 얻는다.
\end{proof}

\begin{lemma}
\label{algebra-lemma-dimension-inequality-quasi-finite}
$R \to S$를 유한형 환 준동형사상이라 하자.
$\mathfrak q \subset S$가 $\mathfrak p \subset R$ 위에 놓이는
소 아이디얼이라 하자. $R \to S$가 $\mathfrak q$에서 준유한이면
$\dim(S_{\mathfrak q}) \leq \dim(R_{\mathfrak p})$이다.
\end{lemma}

\begin{proof}
$R_{\mathfrak p}$가 뇌터이면
($S_{\mathfrak q}$는 $R_{\mathfrak p}$ 위에서 본질적 유한형이므로
역시 뇌터이다), 이는 보조정리
\ref{algebra-lemma-dimension-base-fibre-total}와 정의들에서 즉시 따른다.
일반적인 경우 $S'$을 $S_\mathfrak p$ 안에서
$R_\mathfrak p$의 정수적 폐포라 하자. Zariski의 주정리
\ref{algebra-theorem-main-theorem}에 의해
% P01-E1430
$\mathfrak q$ 위에 놓이는 어떤 $\mathfrak q' \subset S'$에 대하여
$S_{\mathfrak q} = S'_{\mathfrak q'}$이다.
보조정리 \ref{algebra-lemma-integral-dim-up}에 의해
$\dim(S') \leq \dim(R_\mathfrak p)$이고, 따라서 더구나
$\dim(S_\mathfrak q) = \dim(S'_{\mathfrak q'}) \leq \dim(R_\mathfrak p)$이다.
\end{proof}

\begin{lemma}
\label{algebra-lemma-dimension-quasi-finite-over-polynomial-algebra}
\begin{slogan}
아핀 n-공간의 준유한 덮개의 차원은 n 이하이다.
\end{slogan}
$k$를 체라 하자. $S$를 유한형 $k$-대수라 하자.
준유한 $k$-대수 준동형사상
$k[t_1, \ldots, t_n] \subset S$가 존재한다고 하자.
그러면 $\dim(S) \leq n$이다.
\end{lemma}

\begin{proof}
보조정리 \ref{algebra-lemma-dim-affine-space}에 의해
$k[t_1, \ldots, t_n]$의 모든 국소환의 차원은 $n$ 이하이다.
따라서 보조정리 \ref{algebra-lemma-dimension-inequality-quasi-finite}에서
결론이 따른다.
\end{proof}

\begin{lemma}
\label{algebra-lemma-dimension-fibres-bounded-open-upstairs}
$R \to S$를 유한형 환 준동형사상이라 하자.
$\mathfrak q \subset S$를 소 아이디얼이라 하자.
$\dim_{\mathfrak q}(S/R) = n$이라고 가정하자.
모든 $\mathfrak q' \in V$에 대하여
$\dim_{\mathfrak q'}(S/R) \leq n$인
$\Spec(S)$ 안의 $\mathfrak q$의 열린 근방 $V$가 존재한다.
\end{lemma}

\begin{proof}
보조정리 \ref{algebra-lemma-quasi-finite-over-polynomial-algebra}에 의해
$S$가 다항식 대수 $R[t_1, \ldots, t_n]$ 위에서 준유한이라고
가정해도 된다. 섬유들을 생각하면 보조정리
\ref{algebra-lemma-dimension-quasi-finite-over-polynomial-algebra}로 환원된다.
\end{proof}

\noindent
달리 말하면, 이 보조정리는 섬유의 차원이 $\leq n$인 점들의
집합이 $\Spec(S)$에서 열린다는 뜻이다. 다음 보조정리는 이
열린집합의 형성이 밑변환과 가환함을 말한다. 환 준동형사상이
유한 표시이면 이 집합은 준콤팩트 열린집합이다(아래를 보라).

\begin{lemma}
\label{algebra-lemma-dimension-fibres-bounded-open-upstairs-base-change}
$R \to S$를 유한형 환 준동형사상이라 하자.
$R \to R'$을 임의의 환 준동형사상이라 하자.
$S' = R' \otimes_R S$로 두고, 스펙트럼 사이의 대응하는 사상을
% P01-E1431
$f : \Spec(S') \to \Spec(S)$로 나타내자.
$n \geq 0$이라 하자. 역상
$f^{-1}(\{\mathfrak q \in \Spec(S) \mid
\dim_{\mathfrak q}(S/R) \leq n\})$
은
$\{\mathfrak q' \in \Spec(S') \mid
\dim_{\mathfrak q'}(S'/R') \leq n\}$
과 같다.
\end{lemma}

\begin{proof}
이 조건은 체 위에서 유한형인 섬유환들의 차원으로 표현된다.
이를 보조정리
\ref{algebra-lemma-dimension-at-a-point-preserved-field-extension}과
결합하면 결론을 얻는다.
\end{proof}

\begin{lemma}
\label{algebra-lemma-dimension-fibres-bounded-quasi-compact-open-upstairs}
$R \to S$를 유한 표시 환 준동형사상이라 하자.
$n \geq 0$이라 하자. 집합
$$
V_n = \{\mathfrak q \in \Spec(S) \mid \dim_{\mathfrak q}(S/R) \leq n\}
$$
은 $\Spec(S)$의 준콤팩트 열린부분집합이다.
\end{lemma}

\begin{proof}
보조정리 \ref{algebra-lemma-dimension-fibres-bounded-open-upstairs}에 의해
열려 있다. $S$의 표시
$S = R[x_1, \ldots, x_n]/(f_1, \ldots, f_m)$을 택하자.
$R_0$를 다항식들 $f_i$의 계수들로 생성되는 $R$의
$\mathbf{Z}$-부분대수라 하자.
$S_0 = R_0[x_1, \ldots, x_n]/(f_1, \ldots, f_m)$로 두자.
그러면 $S = R \otimes_{R_0} S_0$이다. 보조정리
\ref{algebra-lemma-dimension-fibres-bounded-open-upstairs-base-change}에
의해 $V_n$은 준콤팩트 연속사상
$\Spec(S) \to \Spec(S_0)$ 아래에서 어떤 열린집합 $V_{0, n}$의
역상이다. $S_0$가 뇌터이므로 $V_{0, n}$은 준콤팩트이고,
결론을 얻는다.
\end{proof}

\begin{lemma}
\label{algebra-lemma-finite-type-domain-over-valuation-ring-dim-fibres}
$R$을 잉여체가 $k$이고 분수체가 $K$인 값매김환이라 하자.
$S$를 $R$을 포함하는 정역이라 하자. $S$는 $R$ 위에서
유한형이라고 가정하자.
$S \otimes_R k$가 영환이 아니면
$$
\dim(S \otimes_R k) = \dim(S \otimes_R K)
$$
이다. 실제로 $\Spec(S \otimes_R k)$는 등차원이다.
% P01-E1432
\end{lemma}

\begin{proof}
% P01-E1433
$\mathfrak m_RS$를 포함하는 $S$의 모든 소 아이디얼
$\mathfrak q$에 대하여 $\dim_{\mathfrak q}(S/k)$가
$\dim(S \otimes_R K)$와 같음을 보이면 충분하다. 그러한 소
아이디얼을 택하자. 보조정리
\ref{algebra-lemma-dimension-fibres-bounded-open-upstairs}에 의해 부등식
$\dim_{\mathfrak q}(S/k) \geq \dim(S \otimes_R K)$가 성립한다.
$n = \dim_{\mathfrak q}(S/k)$로 두자. 보조정리
\ref{algebra-lemma-quasi-finite-over-polynomial-algebra}에 의해
$S$를 어떤 $g \in S$, $g \not \in \mathfrak q$에 대한
$S_g$로 바꾸고 나면 준유한 환 준동형사상
$R[t_1, \ldots, t_n] \to S$가 존재한다.
$\dim(S \otimes_R K) < n$이면
$K[t_1, \ldots, t_n] \to S \otimes_R K$는 영이 아닌 핵을 갖는다.
$f = \sum a_I t_1^{i_1}\ldots t_n^{i_n}$라고 하자.
$f$를 값매김이 최소인 $f$의 영이 아닌 계수로 나누어
$f\in R[t_1, \ldots, t_n]$이고 어떤 $a_I$가 $k$에서 영으로
사상되지 않는다고 해도 된다. 따라서 환 준동형사상
$k[t_1, \ldots, t_n] \to S \otimes_R k$는 영이 아닌 핵을 가지며,
이는 $\dim(S \otimes_R k) < n$을 함의한다. 모순이다.
\end{proof}

\section{유한 표시 대수와 가군}
\label{algebra-section-finite-presentation}

\noindent
이 절에서는 가정에 유한형 또는 유한 표시 조건이 들어가는 것이
핵심인 몇 가지 표준 결과를 논의한다.

\begin{lemma}
\label{algebra-lemma-finite-type-descends}
$R \to S$를 환 준동형사상이라 하자.
$R \to R'$을 충실 평탄 환 준동형사상이라 하자.
$S' = R'\otimes_R S$로 두자.
$R \to S$가 유한형일 필요충분조건은 $R' \to S'$이
유한형인 것이다.
\end{lemma}

\begin{proof}
$R \to S$가 유한형이면 $R' \to S'$이 유한형임은 분명하다.
$R' \to S'$이 유한형이라고 가정하자.
$y_1, \ldots, y_m$이 $R'$ 위에서 $S'$을 생성한다고 하자.
어떤 $a_{ij} \in R'$와 $x_{ji} \in S$에 대하여
$y_j = \sum_i a_{ij} \otimes x_{ji}$라고 쓰자.
% P01-E1434
$A \subset S$를 $x_{ij}$들로 생성되는 $R$-부분대수라 하자.
평탄성에 의해 $A' := R' \otimes_R A \subset S'$이고,
구성에 의해 $y_j \in A'$이다. 따라서 $A' = S'$이다.
충실 평탄성에 의해 $A = S$이다.
\end{proof}

\begin{lemma}
\label{algebra-lemma-finite-presentation-descends}
$R \to S$를 환 준동형사상이라 하자.
$R \to R'$을 충실 평탄 환 준동형사상이라 하자.
$S' = R'\otimes_R S$로 두자.
$R \to S$가 유한 표시일 필요충분조건은 $R' \to S'$이
유한 표시인 것이다.
\end{lemma}

\begin{proof}
$R \to S$가 유한 표시이면 $R' \to S'$이 유한 표시임은 분명하다.
$R' \to S'$이 유한 표시라고 가정하자. 보조정리
\ref{algebra-lemma-finite-type-descends}에 의해 $R \to S$는 유한형이다.
$S = R[x_1, \ldots, x_n]/I$라고 쓰자. 평탄성에 의해
$S' = R'[x_1, \ldots, x_n]/R'\otimes I$이다.
$g_1, \ldots, g_m$이 $R'[x_1, \ldots, x_n]$ 위에서
$R'\otimes I$를 생성한다고 하자. 어떤 $a_{ij} \in R'$와
$f_{ji} \in I$에 대하여
$g_j = \sum_i a_{ij} \otimes f_{ji}$라고 쓰자.
% P01-E1435
$J \subset I$를 $f_{ij}$들로 생성되는 아이디얼이라 하자.
평탄성에 의해 $R' \otimes_R J \subset R'\otimes_R I$이고,
둘 다 $R'[x_1, \ldots, x_n]$ 위의 아이디얼이다.
구성에 의해 $g_j \in R' \otimes_R J$이다. 따라서
$R' \otimes_R J = R'\otimes_R I$이다.
충실 평탄성에 의해 $J = I$이다.
\end{proof}

\begin{lemma}
\label{algebra-lemma-construct-fp-module}
$R$을 환이라 하자. $I \subset R$을 아이디얼이라 하고,
$S \subset R$을 곱셈적 부분집합이라 하자.
$R' = S^{-1}(R/I) = S^{-1}R/S^{-1}I$로 두자.
\begin{enumerate}
\item 임의의 유한 $R'$-가군 $M'$에 대하여
$S^{-1}(M/IM) \cong M'$인 유한 $R$-가군 $M$이 존재한다.
\item 임의의 유한 표시 $R'$-가군 $M'$에 대하여
$S^{-1}(M/IM) \cong M'$인 유한 표시 $R$-가군 $M$이 존재한다.
\end{enumerate}
\end{lemma}

\begin{proof}
(1)의 증명. 짧은 완전열
$0 \to K' \to (R')^{\oplus n} \to M' \to 0$을 택하자.
$K \subset R^{\oplus n}$를 사상
$R^{\oplus n} \to (R')^{\oplus n}$ 아래에서 $K'$의 역상이라 하자.
그러면 $M = R^{\oplus n}/K$가 원하는 가군이다.

\medskip\noindent
(2)의 증명. 표시
$(R')^{\oplus m} \to (R')^{\oplus n} \to M' \to 0$을 택하자.
첫 번째 사상이 행렬 $A' = (a'_{ij})$로 주어지고 두 번째 사상이
생성원 $x'_i \in M'$, $i = 1, \ldots, n$으로 결정된다고 하자.
$R' = S^{-1}(R/I)$이므로 $s \in S$와 $R$에 계수를 갖는 행렬
$A = (a_{ij})$를 택하여
$a'_{ij} = a_{ij} / s \bmod S^{-1}I$이 되게 할 수 있다.
첫 번째 사상이 행렬 $A$로 주어지고 두 번째 사상이 생성원
$x_i \in M$, $i = 1, \ldots, n$으로 결정되는 표시
$R^{\oplus m} \to R^{\oplus n} \to M \to 0$을 갖는
유한 표시 $R$-가군을 $M$이라 하자.
그러면 사상 $M \to M'$, $x_i \mapsto x'_i$는 동형사상
$S^{-1}(M/IM) \cong M'$을 유도한다.
\end{proof}

\begin{lemma}
\label{algebra-lemma-construct-fp-module-from-localization}
$R$을 환이라 하자. $S \subset R$을 곱셈적 부분집합이라 하고,
$M$을 $R$-가군이라 하자.
\begin{enumerate}
\item $S^{-1}M$이 유한 $S^{-1}R$-가군이면, 유한 $R$-가군 $M'$과
동형사상 $S^{-1}M' \to S^{-1}M$을 유도하는 사상
$M' \to M$이 존재한다.
\item $S^{-1}M$이 유한 표시 $S^{-1}R$-가군이면, 유한 표시
$R$-가군 $M'$과 동형사상 $S^{-1}M' \to S^{-1}M$을 유도하는
사상 $M' \to M$이 존재한다.
\end{enumerate}
\end{lemma}

\begin{proof}
(1)의 증명. $S^{-1}R$-가군으로서 $S^{-1}M$을 생성하는 원소
$x_1, \ldots, x_n \in M$을 택하자. $M'$을
$x_1, \ldots, x_n$으로 생성되는 $M$의 $R$-부분가군이라 하자.

\medskip\noindent
(2)의 증명. $S^{-1}R$-가군으로서 $S^{-1}M$을 생성하는 원소
$x_1, \ldots, x_n \in M$을 택하자. 규칙
$(a_1, \ldots, a_n) \mapsto \sum a_i x_i$로 주어지는 사상에 대하여
$K = \Ker(R^{\oplus n} \to M)$로 두자. 보조정리
\ref{algebra-lemma-extension}에 의해 $S^{-1}K$는 유한
$S^{-1}R$-가군이다. (1)에 의해 $S^{-1}K' = S^{-1}K$인
유한 부분가군 $K' \subset K$를 찾을 수 있다.
$M' = \Coker(K' \to R^{\oplus n})$로 두자.
\end{proof}

\begin{lemma}
\label{algebra-lemma-construct-fp-module-from-stalk}
$R$을 환이라 하자. $\mathfrak p \subset R$을 소 아이디얼이라 하고,
$M$을 $R$-가군이라 하자.
\begin{enumerate}
\item $M_{\mathfrak p}$가 유한 $R_{\mathfrak p}$-가군이면,
유한 $R$-가군 $M'$과 동형사상
$M'_{\mathfrak p} \to M_{\mathfrak p}$을 유도하는 사상
$M' \to M$이 존재한다.
\item $M_{\mathfrak p}$가 유한 표시 $R_{\mathfrak p}$-가군이면,
유한 표시 $R$-가군 $M'$과 동형사상
$M'_{\mathfrak p} \to M_{\mathfrak p}$을 유도하는 사상
$M' \to M$이 존재한다.
\end{enumerate}
\end{lemma}

\begin{proof}
이는 보조정리
\ref{algebra-lemma-construct-fp-module-from-localization}의 특수한 경우이다.
% P01-E1436
\end{proof}

\begin{lemma}
\label{algebra-lemma-local-isomorphism}
$\varphi : R \to S$를 환 준동형사상이라 하자.
$\mathfrak q \subset S$가 $\mathfrak p \subset R$ 위에 놓이는
소 아이디얼이라 하자. 다음을 가정하자.
\begin{enumerate}
\item $S$는 $R$ 위에서 유한 표시이고,
\item $\varphi$는 동형사상 $R_\mathfrak p \cong S_\mathfrak q$를
유도한다.
\end{enumerate}
그러면 $f \in R$, $f \not \in \mathfrak p$와
$R_f$-대수 $C$가 존재하여 $R_f$-대수로서
$S_f \cong R_f \times C$이다.
\end{lemma}

\begin{proof}
$S = R[x_1, \ldots, x_n]/(g_1, \ldots, g_m)$이라고 쓰자.
$a_i \in R_\mathfrak p$를 $S_\mathfrak q$에서 $x_i$의 상으로
사상되는 원소라 하자. 어떤 $f \in R$, $f \not \in \mathfrak p$에
대하여 $a_i = b_i/f$라고 쓰자. $R$을 $R_f$로, $x_i$를
$x_i - a_i$로 바꾸어, $x_i$가 $S_\mathfrak q$에서 영으로 사상되는
$S = R[x_1, \ldots, x_n]/(g_1, \ldots, g_m)$이라고 해도 된다.
그러면 $c_j$가 $g_j$의 상수항일 때 $c_j$가
$R_\mathfrak p$에서 영으로 사상된다고 결론 내린다.
$R$을 한 번 더 바꾸어 $g_j$들의 상수 계수 $c_j$들이 영이라고
해도 된다. 따라서 핵이 아이디얼 $(x_1, \ldots, x_n)$인
$R$-대수 준동형사상 $S \to R$, $x_i \mapsto 0$을 얻는다.

\medskip\noindent
동형사상들
$$
R_\mathfrak p \to S_\mathfrak q \to R_\mathfrak p
$$
이 있고 $S \to R$은 $x_i$를 영으로 보낸다. 따라서
$$
S_\mathfrak q = R_\mathfrak p[x_1, \ldots, x_n]/(x_1, \ldots, x_n)
$$
이고,
더구나
$$
S_\mathfrak q = S_\mathfrak p/(x_1, \ldots, x_n)S_\mathfrak p
$$
이다.
이는 유한 생성 아이디얼
$(x_1, \ldots, x_n)S_\mathfrak p$가 $S_\mathfrak p$에서
순수하다는 뜻이다. 정의 \ref{algebra-definition-pure-ideal}을 보라.
따라서 보조정리 \ref{algebra-lemma-finitely-generated-pure-ideal}에 의해
$$
(x_1, \ldots, x_n)S_\mathfrak p
$$
는 $S_\mathfrak p$의 어떤
멱등원 $e$로 생성된다. 어떤 $f \in R$, $f \not \in \mathfrak p$에
대하여 $R \to S$를 $R_f \to S_f$로 바꾸어
$e$로 사상되는 멱등원 $e' \in S$를 찾을 수 있다.
그러면 $e'S$와 $(x_1, \ldots, x_n)S$는 $S_\mathfrak p$에서
같아지는 유한 생성 아이디얼들이다. 따라서 어떤
$f \in R$, $f \not \in \mathfrak p$에 대하여 다시
$R \to S$를 $R_f \to S_f$로 바꾸어
$e'S = (x_1, \ldots, x_n)S$라고 해도 된다.
$C = e'S$로 두면 증명이 끝난다.
\end{proof}

\begin{lemma}
\label{algebra-lemma-isomorphic-local-rings}
$R$을 환이라 하자. $S$, $S'$은 $R$ 위에서 유한 표시라고 하자.
$\mathfrak q \subset S$와 $\mathfrak q' \subset S'$을
소 아이디얼들이라 하자. $R$-대수로서
$S_{\mathfrak q} \cong S'_{\mathfrak q'}$이면,
$R$-대수로서 $S_g \cong S'_{g'}$이 되게 하는
$g \in S$, $g \not \in \mathfrak q$와
$g' \in S'$, $g' \not \in \mathfrak q'$이 존재한다.
\end{lemma}

\begin{proof}
$\psi : S_{\mathfrak q} \to S'_{\mathfrak q'}$를 보조정리의
가정에 나오는 동형사상이라 하자.
$S = R[x_1, \ldots, x_n]/(f_1, \ldots, f_r)$와
$S' = R[y_1, \ldots, y_m]/J$라고 쓰자.
각 $i = 1, \ldots, n$에 대하여, $\psi$에 의한 $x_i$의 상을
나타내는 분수 $h_i/g_i$를 택하자. 여기서
$h_i, g_i \in R[y_1, \ldots, y_m]$이고
$g_i \bmod J$는 $\mathfrak q'$에 속하지 않는다.
$S'$을 $S'_{g_1 \ldots g_n}$로 바꾸고,
$R[y_1, \ldots, y_m, y_{m + 1}]$에서 $y_{m + 1}$을
$1/(g_1\ldots g_n)$으로 사상시켜, $\psi(x_i)$가 어떤
$h_i \in R[y_1, \ldots, y_m]$의 상이라고 해도 된다.
원소 $f_j(h_1, \ldots, h_n) \in R[y_1, \ldots, y_m]$들을
생각하자. $\psi$가 각 $f_j$를 죽이므로,
각 $j = 1, \ldots, r$에 대하여
$g f_j(h_1, \ldots, h_n) \in J$가 되게 하는
$g \in R[y_1, \ldots, y_m]$, $g \bmod J \not \in \mathfrak q'$이
존재한다. 앞에서처럼 $S'$을 $S'_g$로,
$R[y_1, \ldots, y_m, y_{m + 1}]$로 바꾸어
$f_j(h_1, \ldots, h_n) \in J$라고 해도 된다.
따라서 국소환들 위에서 $\psi$를 유도하는 환 준동형사상
$S \to S'$, $x_i \mapsto h_i$를 얻는다.
보조정리 \ref{algebra-lemma-compose-finite-type}에 의해 $S \to S'$은
유한 표시이다. 보조정리 \ref{algebra-lemma-local-isomorphism}에 의해
$S' = S \times C$라고 해도 된다.
% P01-E1437
따라서 인자 $C$에 대응하는 멱등원에서 $S'$을 국소화하면
결론을 얻는다.
\end{proof}

\begin{lemma}
\label{algebra-lemma-finite-type-mod-nilpotent}
$R$을 환이라 하고 $I \subset R$을 멱영 아이디얼이라 하자.
$R/I \to S/IS$가 유한형인 $R$-대수를 $S$라 하자.
그러면 $R \to S$는 유한형이다.
\end{lemma}

\begin{proof}
$S/IS$에서의 상들이 $R/I$ 위의 대수로서 $S/IS$를 생성하는
$s_1, \ldots, s_n \in S$를 택하자. 보조정리
\ref{algebra-lemma-NAK}의 (11)에 의해 $R$-대수 준동형사상
$R[x_1, \ldots, x_n] \to S$, $x_i \mapsto s_i$는 전사이고
결론을 얻는다.
\end{proof}

\begin{lemma}
\label{algebra-lemma-surjective-mod-locally-nilpotent}
$R$을 환이라 하고 $I \subset R$을 국소 멱영 아이디얼이라 하자.
$S \to S'/IS'$이 전사이고 $S'$이 $R$ 위에서 유한형인
$R$-대수 준동형사상 $S \to S'$을 생각하자.
그러면 $S \to S'$은 전사이다.
\end{lemma}

\begin{proof}
어떤 아이디얼 $K$에 대하여
$S' = R[x_1, \ldots, x_m]/K$라고 쓰자. 가정에 의해
% P01-E1438
$g_j = x_j + \sum \delta_{j, J} x^J \in R[x_1, \ldots, x_n]$,
$\delta_{j, J} \in I$이고
$g_j \bmod K \in \Im(S \to S')$인 원소들이 존재한다.
따라서 $g_1, \ldots, g_m$이 $R[x_1, \ldots, x_n]$을 생성함을
보이면 충분하다. 적어도 $\delta_{j, J}$들을 포함하는
$R$의 유한 생성 $\mathbf{Z}$-부분대수를 $R_0 \subset R$라 하자.
그러면 보조정리 \ref{algebra-lemma-Noetherian-power}에 의해
$R_0 \cap I$는 멱영 아이디얼이다. 따라서
$R_0[x_1, \ldots, x_n]$은 $g_1, \ldots, g_m$으로 생성된다.
(실제로 $x_j \mapsto g_j$는
$R_0[x_1, \ldots, x_m]$의 자기동형사상을 정의한다. 자세한 내용은
생략한다.) $R$은 부분환들 $R_0$의 합집합이므로 결론을 얻는다.
\end{proof}

\begin{lemma}
\label{algebra-lemma-isomorphism-modulo-ideal}
$R$을 환이라 하고 $I \subset R$을 아이디얼이라 하자.
$S \to S'$을 $R$-대수 준동형사상이라 하자.
$IS \subset \mathfrak q \subset S$를 소 아이디얼이라 하자.
다음을 가정하자.
\begin{enumerate}
\item $S \to S'$은 전사이고,
\item $S_\mathfrak q/IS_\mathfrak q \to S'_\mathfrak q/IS'_\mathfrak q$는
동형사상이며,
\item $S$는 $R$ 위에서 유한형이고,
% P01-E1439
\item $S'$은 $R$ 위에서 유한 표시이며,
\item $S'_\mathfrak q$는 $R$ 위에서 평탄하다.
\end{enumerate}
그러면 어떤 $g \in S$, $g \not \in \mathfrak q$에 대하여
$S_g \to S'_g$는 동형사상이다.
\end{lemma}

\begin{proof}
$J = \Ker(S \to S')$로 두자. 보조정리
\ref{algebra-lemma-compose-finite-type}에 의해 $J$는 유한 생성
아이디얼이다. $S'_\mathfrak q$가 $R$ 위에서 평탄하므로
$J_\mathfrak q/IJ_\mathfrak q \subset S_\mathfrak q/IS_{\mathfrak q}$이다.
($0 \to J \to S \to S' \to 0$에 보조정리
\ref{algebra-lemma-flat-tor-zero}를 적용하라.) 가정 (2)에 의해
$J_\mathfrak q/IJ_\mathfrak q$는 영이다. Nakayama의 보조정리
(보조정리 \ref{algebra-lemma-NAK})에 의해 $J_g = 0$이 되게 하는
$g \in S$, $g \not \in \mathfrak q$가 존재한다. 따라서 원하는 대로
$S_g \cong S'_g$이다.
\end{proof}

\begin{lemma}
\label{algebra-lemma-isomorphism-modulo-locally-nilpotent}
$R$을 환이라 하고 $I \subset R$을 아이디얼이라 하자.
$S \to S'$을 $R$-대수 준동형사상이라 하자. 다음을 가정하자.
\begin{enumerate}
\item $I$는 국소 멱영이고,
\item $S/IS \to S'/IS'$은 동형사상이며,
\item $S$는 $R$ 위에서 유한형이고,
% P01-E1440
\item $S'$은 $R$ 위에서 유한 표시이며,
\item $S'$은 $R$ 위에서 평탄하다.
\end{enumerate}
그러면 $S \to S'$은 동형사상이다.
\end{lemma}

\begin{proof}
보조정리 \ref{algebra-lemma-surjective-mod-locally-nilpotent}에 의해 사상
$S \to S'$은 전사이다. $I$가 국소 멱영이므로 아이디얼들
$IS$와 $IS'$도 국소 멱영이다(보조정리
\ref{algebra-lemma-locally-nilpotent}). 따라서 $S$의 모든 소 아이디얼
$\mathfrak q$는 $IS$를 포함하고, (자명하게)
$S_\mathfrak q/IS_\mathfrak q \cong S'_\mathfrak q/IS'_\mathfrak q$이다.
그러므로 보조정리 \ref{algebra-lemma-isomorphism-modulo-ideal}을 적용하면
모든 소 아이디얼 $\mathfrak q \subset S$에 대하여
$S_\mathfrak q \to S'_\mathfrak q$가 동형사상임을 알 수 있다.
예를 들어 보조정리 \ref{algebra-lemma-characterize-zero-local}에 의해
$S \to S'$은 단사이고, 결론을 얻는다.
\end{proof}

\section{여극한과 유한 표시 사상}
\label{algebra-section-colimits-flat}

\noindent
% P01-E1441
이 절에서는 나중에 절대 Noetherian 축소를 이용하여 결과를 증명하는 데
도움이 될 몇 가지 예비 보조정리를 증명한다.
Categories, Section \ref{categories-section-directed-colimits}에서
일반적인 여과 여극한을 논의한다.
다음은 이 매우 일반적인 개념의 한 예이다.

\begin{lemma}
\label{algebra-lemma-ring-colimit-fp-category}
$R \to A$를 환 준동형사상이라 하자. $R$ 위에서 유한 표시인
$A'$을 갖는 모든 $R$-대수 준동형사상의 도식 $A' \to A$들로 이루어진
범주 $\mathcal{I}$를 생각하자. 그러면 $\mathcal{I}$는 여과 범주이고,
$\mathcal{I}$ 위에서 $A'$들의 여극한은 $A$와 동형이다.
\end{lemma}

\begin{proof}
범주\footnote{집합론적 어려움을 피하기 위하여, $A'$이
$R[x_1, x_2, x_3, \ldots]$의 몫인 $A' \to A$만을 생각한다.}
$\mathcal{I}$는 $R \to R$이 그 대상이므로 공집합이 아니다.
$\mathcal{I}$의 두 대상 $A' \to A$, $A'' \to A$를 생각하자.
그러면 $A' \otimes_R A'' \to A$는 $\mathcal{I}$에 속한다
(보조정리 \ref{algebra-lemma-compose-finite-type}과
\ref{algebra-lemma-base-change-finiteness}를 사용하라). 환 준동형사상
$A' \to A' \otimes_R A''$와 $A'' \to A' \otimes_R A''$는
$\mathcal{I}$의 사상을 정의하므로, Categories, Definition
\ref{categories-definition-directed}을 보아 여과 범주의 두 번째
정의 성질이 성립한다. 마지막으로 $\mathcal{I}$에서 두 사상
$\sigma, \tau : A' \to A''$가 주어졌다고 하자.
$x_1, \ldots, x_r \in A'$이 $R$-대수로서 $A'$의 생성원이라면
$A''' = A''/(\sigma(x_i) - \tau(x_i))$를 생각할 수 있다.
이는 유한 표시 $R$-대수이고, 주어진 $R$-대수 준동형사상
$A'' \to A$는 전사 $\nu : A'' \to A'''$를 거쳐 분해된다.
따라서 $\nu$는 원하는 대로 $\sigma$와 $\tau$를 같게 만드는
$\mathcal{I}$의 사상이다.

\medskip\noindent
% P01-E1442
우리 지표 범주가 공여과 범주라는 사실로부터 집합의 범주에서
$B = \colim_{A' \to A} A'$의 값을 계산할 수 있다
(일부 세부 사항은 생략한다. Categories, Section
\ref{categories-section-directed-colimits}의 논의와 비교하라).
$B \to A$가 전사임을 보이기 위해, 모든 $a \in A$에 대하여
$R[x] \to A$, $x \mapsto a$를 사용하면 $a$가 $B \to A$의
상에 속함을 알 수 있다. 반대로 $b \in B$가 $A$의 영으로
사상된다고 하자. 그러면 $\mathcal{I}$의 어떤 $A' \to A$와
$b$로 사상되는 $a' \in A'$을 찾을 수 있다. 이때
$A'/(a') \to A$도 $\mathcal{I}$에 속하고, 사상 $A' \to B$는
$A' \to A'/(a') \to B$로 분해되므로 원하는 대로 $b = 0$이다.
\end{proof}

\noindent
여극한을 준순서집합 위에서 생각하는 편이 더 쉬울 때가 많다.
% P01-E1443
$(\Lambda, \geq)$를 준순서집합이라 하자.
$\Lambda$ 위의 환의 계란, 각 $\lambda \in \Lambda$에 대한 환
$R_\lambda$와 $\lambda \leq \mu$일 때마다 주어진 준동형사상
$R_\lambda \to R_\mu$로 이루어진다. 이 준동형사상들은 모든
$\lambda \leq \mu \leq \nu$에 대하여
$R_\lambda \to R_\mu \to R_\nu$가 사상
$R_\lambda \to R_\nu$와 같다는 규칙을 만족해야 한다.
Categories, Section \ref{categories-section-posets-limits}을 보라.
% P01-E1444
우리는 흔히 $(I, \leq)$가 {\it 유향}이라고 가정할 것이다.
이는 $\Lambda$가 공집합이 아니고, $\lambda, \mu \in \Lambda$가
주어지면 $\lambda \leq \nu$와 $\mu \leq \nu$를 만족하는
$\nu \in \Lambda$가 존재한다는 뜻이다. 이 경우 여극한
$\colim_\lambda R_\lambda$를 때때로 ``직접 극한''이라고 부른다는
것을 상기하자(그러나 우리는 이 용어를 사용하지 않을 것이다).

\medskip\noindent
Categories, Lemma \ref{categories-lemma-directed-category-system}에 의하면
여과 지표 범주 위의 여극한은 유향 집합 위의 여극한과 같은 것임에
유의하라.

\begin{lemma}
\label{algebra-lemma-ring-colimit-fp}
$R \to A$를 환 준동형사상이라 하자. 유한 표시 $R$-대수들의 유향계
$A_\lambda$로서 $A = \colim_\lambda A_\lambda$인 것이 존재한다.
$A$가 $R$ 위에서 유한형이면, $A_\lambda$의 계에 있는 모든 전이
사상이 전사가 되도록 정할 수 있다.
\end{lemma}

\begin{proof}
첫 번째 증명은 이것이 보조정리 \ref{algebra-lemma-ring-colimit-fp-category}와
Categories, Lemma \ref{categories-lemma-directed-category-system}에서
따른다는 것이다.

\medskip\noindent
두 번째 증명.
보조정리 \ref{algebra-lemma-module-colimit-fp}의 증명과 비교하라.
임의의 유한 부분집합 $S \subset A$와 $S$의 원소들 사이의 다항식
관계들로 이루어진 임의의 유한 모음 $E$를 생각하자. 따라서 각
$s \in S$는 $x_s \in A$에 대응하고, 각 $e \in E$는
$f_e(x_s) = 0$을 만족하는 다항식
$f_e \in R[X_s; s\in S]$로 이루어진다.
$A_{S, E} = R[X_s; s\in S]/(f_e; e\in E)$라 두면 이는 유한 표시
$R$-대수이다. 표준 사상 $A_{S, E} \to A$가 있다.
$S \subset S'$이고, $E$의 원소들이 사상
$R[X_s; s \in S] \to R[X_s; s\in S']$을 통해 $E'$의 부분집합에
대응한다면, $A$로 가는 사상들과 가환하는 자명한 사상
$A_{S, E} \to A_{S', E'}$이 있다. 따라서
% P01-E1445
$\Lambda$를 위와 같은 포함관계로 순서가 주어진 쌍 $(S, E)$들의
집합과 같게 두면, 유향 부분순서집합을 얻는다.
이 유향계의 여극한이 $A$임은 명백하다.

\medskip\noindent
마지막 명제를 위하여 $A = R[x_1, \ldots, x_n]/I$라고 하자.
이 경우 위의 계들 $(S, E)$ 가운데 $S = \{x_1, \ldots, x_n\}$인
것들로 이루어진 부분집합 $\Lambda' \subset \Lambda$를 생각하자.
여전히 $A = \colim_{\lambda' \in \Lambda'} A_{\lambda'}$임을
쉽게 알 수 있다. 더욱이 전이 사상들은 명백히 전사이다.
\end{proof}

\noindent
유한 표시 환 준동형사상은 다음과 같이 특징지을 수 있다.
이는 어떤 의미에서 유한 표시 대수들이 $R$-대수의 범주에서
``콤팩트'' 대상이라는 것을 말해 준다.

\begin{lemma}
\label{algebra-lemma-characterize-finite-presentation}
환 준동형사상 $\varphi : R \to S$에 대하여 다음은 동치이다.
\begin{enumerate}
\item $\varphi$는 유한 표시이고,
\item 모든 $R$-대수의 유향계 $A_\lambda$에 대하여 사상
$$
\colim_\lambda \Hom_R(S, A_\lambda) \longrightarrow
\Hom_R(S, \colim_\lambda A_\lambda)
$$
은 전단사이며,
\item 모든 $R$-대수의 유향계 $A_\lambda$에 대하여 사상
$$
\colim_\lambda \Hom_R(S, A_\lambda) \longrightarrow
\Hom_R(S, \colim_\lambda A_\lambda)
$$
은 전사이다.
\end{enumerate}
\end{lemma}

\begin{proof}
(1)을 가정하고 $S = R[x_1, \ldots, x_n] / (f_1, \ldots, f_m)$이라
쓰자. $A = \colim A_\lambda$라 두자. $R$-대수 준동형사상
$S \to A$ 또는 $S \to A_\lambda$는 $x_1, \ldots, x_n$의 상으로
결정됨에 유의하라. 따라서
$\colim_\lambda \Hom_R(S, A_\lambda) \to \Hom_R(S, A)$가
단사임은 명백하다. 전사임을 보이기 위하여 $R$-대수 준동형사상
$\chi : S \to A$를 택하자. 그러면 각 $x_i$는 어떤
$A_{\lambda_i}$의 상에 있는 원소로 사상된다.
$\mu \geq \lambda_i$, $i = 1, \ldots, n$을 택하여
$i = 1, \ldots, n$에 대하여 $\chi(x_i)$가
$y_i \in A_\mu$의 상이라고 가정할 수 있다.
$z_j = f_j(y_1, \ldots, y_n) \in A_\mu$를 생각하자.
$\chi$가 준동형사상이므로 $z_j$의 $A = \colim_\lambda A_\lambda$에서의
상은 영이다. 따라서 $z_j$가 $A_{\mu_j}$에서 영으로 사상되게 하는
$\mu_j \geq \mu$가 존재한다. $\nu \geq \mu_j$,
$j = 1, \ldots, m$을 택하자. 그러면 $z_1, \ldots, z_m$의 상은
$A_\nu$에서 영이다. 이는 정확히 $y_i$가 관계
$f_j(y'_1, \ldots, y'_n) = 0$을 만족하는 원소
$y'_i \in A_\nu$로 사상된다는 뜻이다. 따라서 환 준동형사상
$S \to A_\nu$를 얻는다. 이로써 (1)이 (2)를 함의함을 보였다.

\medskip\noindent
(2)가 (3)을 함의함은 명백하다. (3)을 가정하자.
보조정리 \ref{algebra-lemma-ring-colimit-fp}에 의해 $R$ 위에서 유한 표시인
$S_\lambda$들을 써서 $S = \colim_\lambda S_\lambda$라 쓸 수 있다.
그러면 어떤 $\lambda$에 대하여 항등사상은
$$
S \to S_\lambda \to S
$$
로 분해된다. $S \to S_\lambda \to S$에 보조정리
\ref{algebra-lemma-compose-finite-type}의 (4)를 적용하면 $S$가
$S_\lambda$ 위에서 유한 표시임을 알 수 있다. 같은 보조정리의
(2)를 $R \to S_\lambda \to S$에 적용하면 $S$가 $R$ 위에서
유한 표시임을 얻는다.
\end{proof}

\noindent
위의 기본 내용을 이용하면, 대수 $A$가 주어진 유형의 대수들의
여과 여극한이 되는 기준을 다음과 같이 줄 수 있다.
% P01-E1446

\begin{lemma}
\label{algebra-lemma-when-colimit}
$R \to \Lambda$를 환 준동형사상이라 하자. $\mathcal{E}$를
$R$-대수들의 집합으로서 각 $A \in \mathcal{E}$가 $R$ 위에서
유한 표시인 것이라 하자. 그러면 다음 두 명제는 동치이다.
\begin{enumerate}
\item $\Lambda$는 $\mathcal{E}$의 원소들의 여과 여극한이고,
% P01-E1447
\item $R$-대수 준동형사상 $A \to \Lambda$로서 $A$가 $R$ 위에서
유한 표시인 모든 것에 대하여, $B \in \mathcal{E}$인 분해
$A \to B \to \Lambda$를 찾을 수 있다.
\end{enumerate}
\end{lemma}

\begin{proof}
$\mathcal{I} \to \mathcal{E}$, $i \mapsto A_i$가
$\Lambda = \colim_i A_i$를 만족하는 여과 도식이라고 하자.
$A \to \Lambda$를 $R$ 위에서 유한 표시인 $A$를 갖는
$R$-대수 준동형사상이라 하자. 보조정리
\ref{algebra-lemma-characterize-finite-presentation}을 적용하면 분해
$A \to A_i \to \Lambda$를 얻는다. 따라서 (1)은 (2)를 함의한다.

\medskip\noindent
보조정리 \ref{algebra-lemma-ring-colimit-fp-category}의 범주
$\mathcal{I}$를 생각하자. Categories, Lemma
\ref{categories-lemma-cofinal-in-filtered}에 의해
$A \in \mathcal{E}$인 $A \to \Lambda$들로 이루어진 충만 부분범주
$\mathcal{J}$는 $\mathcal{I}$에서 공종이고 여과 범주이다.
그러면 Categories, Lemma \ref{categories-lemma-cofinal}에 의해
$\Lambda$는 $\mathcal{J}$ 위의 여극한이기도 하다.
\end{proof}

\noindent
그러나 더 강한 사실이 성립한다. 즉 $R = \colim_\lambda R_\lambda$가
주어지면 유한 표시 $R$-가군들의 범주는 유한 표시
$R_\lambda$-가군들의 범주들의 극한과 동치임을 알 수 있다.
% P01-E1448
유한 표시 $R$-대수들의 범주에 대해서도 마찬가지이다.

\begin{lemma}
\label{algebra-lemma-module-map-property-in-colimit}
$A$를 환이라 하고 $M, N$을 $A$-가군이라 하자.
$R = \colim_{i \in I} R_i$가 $A$-대수들의 유향 여극한이라고 하자.
\begin{enumerate}
\item $M$이 유한 $A$-가군이고 $u, u' : M \to N$이
$A$-가군 준동형사상으로서
$u \otimes 1 = u' \otimes 1 : M \otimes_A R \to N \otimes_A R$
이면, 어떤 $i$에 대하여
$u \otimes 1 = u' \otimes 1 : M \otimes_A R_i \to N \otimes_A R_i$이다.
\item $N$이 유한 $A$-가군이고 $u : M \to N$이 $A$-가군
준동형사상으로서
$u \otimes 1 : M \otimes_A R \to N \otimes_A R$이 전사이면,
어떤 $i$에 대하여 사상
$u \otimes 1 : M \otimes_A R_i \to N \otimes_A R_i$가 전사이다.
\item $N$이 유한 표시 $A$-가군이고
$v : N \otimes_A R \to M \otimes_A R$이 $R$-가군 준동형사상이면,
$v = v_i \otimes 1$을 만족하는 어떤 $i$와 $R_i$-가군 준동형사상
$v_i : N \otimes_A R_i \to M \otimes_A R_i$가 존재한다.
\item $M$이 유한 $A$-가군이고 $N$이 유한 표시 $A$-가군이며
$u : M \to N$이 $A$-가군 준동형사상으로서
$u \otimes 1 : M \otimes_A R \to N \otimes_A R$이 동형사상이면,
어떤 $i$에 대하여 사상
$u \otimes 1 : M \otimes_A R_i \to N \otimes_A R_i$가 동형사상이다.
\end{enumerate}
\end{lemma}

\begin{proof}
(1)을 증명하기 위하여 $u$가 (1)과 같다고 하고
$x_1, \ldots, x_m \in M$을 생성원이라 하자.
$N \otimes_A R = \colim_i N \otimes_A R_i$이므로,
% P01-E1449
$M \otimes_A R_i$에서 $u(x_j) \otimes 1 = u'(x_j) \otimes 1$,
$j = 1, \ldots, m$이 되게 하는 $i \in I$를 택할 수 있다.
그러한 $i$에 대하여
$u \otimes 1 = u' \otimes 1 : M \otimes_A R_i \to N \otimes_A R_i$이다.

\medskip\noindent
(2)를 증명하기 위하여 $u \otimes 1$이 전사라고 하고
$y_1, \ldots, y_m \in N$을 생성원이라 하자.
$N \otimes_A R = \colim_i N \otimes_A R_i$이므로,
$N \otimes_A R$에서의 상이 $y_j \otimes 1$인
$i \in I$와 $z_j \in M \otimes_A R_i$, $j = 1, \ldots, m$을
택할 수 있다. 그러한 $i$에 대하여 사상
$u \otimes 1 : M \otimes_A R_i \to N \otimes_A R_i$는 전사이다.

\medskip\noindent
(3)을 증명하기 위하여 $y_1, \ldots, y_m \in N$을 생성원이라 하자.
% P01-E1450
$K = \Ker(A^{\oplus m} \to N)$라 하되, 여기서 사상은 규칙
$(a_1, \ldots, a_m) \mapsto \sum a_j x_j$로 주어진다.
$k_1, \ldots, k_t$를 $K$의 생성원이라 하자.
$k_s = (k_{s1}, \ldots, k_{sm})$라 하자.
$M \otimes_A R = \colim_i M \otimes_A R_i$이므로,
$M \otimes_A R$에서의 상이 $v(y_j \otimes 1)$인
$i \in I$와 $z_j \in M \otimes_A R_i$, $j = 1, \ldots, m$을
택할 수 있다. 이 $z_j$들을 사용하여 사상
$v_i : N \otimes_A R_i \to M \otimes_A R_i$를 정의하고자 한다.
$K \otimes_A R_i \to R_i^{\oplus m} \to N \otimes_A R_i \to 0$은
표시이므로, 각 $s = 1, \ldots, t$에 대하여
$\xi_s = \sum_j k_{sj}z_j$가 $M \otimes_A R_i$에서 영임을 확인하면
충분하다. 그렇지 않을 수도 있지만, $\xi_s$의 $M \otimes_A R$에서의
상이 영이므로 $i$를 조금 증가시키면 그렇게 됨을 알 수 있다.

\medskip\noindent
(4)를 증명하기 위하여 $u \otimes 1$이 동형사상이고,
$M$이 유한이며, $N$이 유한 표시라고 하자.
$v : N \otimes_A R \to M \otimes_A R$을 $u \otimes 1$의
역사상이라 하자. (3)을 적용하여 어떤 $i$에 대한 사상
$v_i : N \otimes_A R_i \to M \otimes_A R_i$를 얻는다.
(1)을 적용하면 $i$를 증가시킨 뒤
% P01-E1451
$v_i \circ (u \otimes 1) = \text{id}_{M \otimes_R R_i}$ 및
$(u \otimes 1) \circ v_i = \text{id}_{N \otimes_R R_i}$이다.
\end{proof}

\begin{lemma}
\label{algebra-lemma-colimit-category-fp-modules}
$R = \colim_{\lambda \in \Lambda} R_\lambda$가 환들의 유향 여극한이라고
하자. 그러면 유한 표시 $R$-가군들의 범주는 유한 표시
$R_\lambda$-가군들의 범주들의 여극한이다. 더 정확히 말하면 다음과 같다.
\begin{enumerate}
\item 유한 표시 $R$-가군 $M$이 주어지면,
$M \cong M_\lambda \otimes_{R_\lambda} R$을 만족하는
$\lambda \in \Lambda$와 유한 표시 $R_\lambda$-가군 $M_\lambda$가
존재한다.
\item $\lambda \in \Lambda$, 유한 표시 $R_\lambda$-가군
$M_\lambda, N_\lambda$, 그리고 $R$-가군 준동형사상
$\varphi : M_\lambda \otimes_{R_\lambda} R \to N_\lambda \otimes_{R_\lambda} R$
이 주어지면, $\varphi = \varphi_\mu \otimes 1_R$을 만족하는
$\mu \geq \lambda$와 $R_\mu$-가군 준동형사상
$\varphi_\mu : M_\lambda \otimes_{R_\lambda} R_\mu \to
N_\lambda \otimes_{R_\lambda} R_\mu$가 존재한다.
\item $\lambda \in \Lambda$, 유한 표시 $R_\lambda$-가군
$M_\lambda, N_\lambda$, 그리고 $R_\lambda$-가군 준동형사상
$\varphi, \psi : M_\lambda \to N_\lambda$가 주어지고
$\varphi \otimes 1_R = \psi \otimes 1_R$이면, 어떤
$\mu \geq \lambda$에 대하여
$\varphi \otimes 1_{R_\mu} = \psi \otimes 1_{R_\mu}$이다.
\end{enumerate}
\end{lemma}

\begin{proof}
(1)을 증명하기 위하여 표시
$R^{\oplus m} \to R^{\oplus n} \to M \to 0$을 택하자.
첫 번째 사상이 행렬 $A = (a_{ij})$로 주어진다고 하자.
$R$에서 $A$로 사상되는, $R_\lambda$에 계수를 갖는 행렬
$A_\lambda = (a_{\lambda, ij})$와 $\lambda \in \Lambda$를
택할 수 있다. 이제 첫 번째 화살표가 $A_\lambda$로 주어지는 표시
$R_\lambda^{\oplus m} \to R_\lambda^{\oplus n} \to M_\lambda \to 0$
를 갖는 $R_\lambda$-가군을 $M_\lambda$로 두기만 하면 된다.

\medskip\noindent
(2)와 (3)은 보조정리 \ref{algebra-lemma-module-map-property-in-colimit}에서
따른다.
\end{proof}

\begin{lemma}
\label{algebra-lemma-algebra-map-property-in-colimit}
$A$를 환이라 하고 $B, C$를 $A$-대수라 하자.
$R = \colim_{i \in I} R_i$가 $A$-대수들의 유향 여극한이라고 하자.
\begin{enumerate}
\item $B$가 유한형 $A$-대수이고 $u, u' : B \to C$가
$A$-대수 준동형사상으로서
$u \otimes 1 = u' \otimes 1 : B \otimes_A R \to C \otimes_A R$
이면, 어떤 $i$에 대하여
$u \otimes 1 = u' \otimes 1 : B \otimes_A R_i \to C \otimes_A R_i$이다.
\item $C$가 유한형 $A$-대수이고 $u : B \to C$가 $A$-대수
준동형사상으로서
$u \otimes 1 : B \otimes_A R \to C \otimes_A R$이 전사이면,
어떤 $i$에 대하여 사상
$u \otimes 1 : B \otimes_A R_i \to C \otimes_A R_i$는 전사이다.
\item $C$가 $A$ 위에서 유한 표시이고
$v : C \otimes_A R \to B \otimes_A R$이 $R$-대수 준동형사상이면,
$v = v_i \otimes 1$을 만족하는 어떤 $i$와 $R_i$-대수 준동형사상
$v_i : C \otimes_A R_i \to B \otimes_A R_i$가 존재한다.
\item $B$가 유한형 $A$-대수이고 $C$가 유한 표시 $A$-대수이며
$u \otimes 1 : B \otimes_A R \to C \otimes_A R$이 동형사상이면,
어떤 $i$에 대하여 사상
$u \otimes 1 : B \otimes_A R_i \to C \otimes_A R_i$는 동형사상이다.
\end{enumerate}
\end{lemma}

\begin{proof}
(1)을 증명하기 위하여 $u$가 (1)과 같다고 하고
$x_1, \ldots, x_m \in B$를 생성원이라 하자.
$B \otimes_A R = \colim_i B \otimes_A R_i$이므로,
$B \otimes_A R_i$에서 $u(x_j) \otimes 1 = u'(x_j) \otimes 1$,
$j = 1, \ldots, m$이 되게 하는 $i \in I$를 택할 수 있다.
그러한 $i$에 대하여
$u \otimes 1 = u' \otimes 1 : B \otimes_A R_i \to C \otimes_A R_i$이다.

\medskip\noindent
(2)를 증명하기 위하여 $u \otimes 1$이 전사라고 하고
$y_1, \ldots, y_m \in C$를 생성원이라 하자.
% P01-E1452
$B \otimes_A R = \colim_i B \otimes_A R_i$이므로,
$C \otimes_A R$에서의 상이 $y_j \otimes 1$인
$i \in I$와 $z_j \in B \otimes_A R_i$, $j = 1, \ldots, m$을
택할 수 있다. 그러한 $i$에 대하여 사상
$u \otimes 1 : B \otimes_A R_i \to C \otimes_A R_i$는 전사이다.

\medskip\noindent
(3)을 증명하기 위하여 $c_1, \ldots, c_m \in C$를 생성원이라 하자.
% P01-E1453
$K = \Ker(A[x_1, \ldots, x_m] \to N)$라 하되, 여기서 사상은
규칙 $x_j \mapsto \sum c_j$로 주어진다. $f_1, \ldots, f_t$를
$A[x_1, \ldots, x_m]$의 아이디얼로서 $K$의 생성원이라 하자.
$f_j = f_j(x_1, \ldots, x_m)$를 다항식으로 생각한다.
$B \otimes_A R = \colim_i B \otimes_A R_i$이므로,
$B \otimes_A R$에서의 상이 $v(c_j \otimes 1)$인
$i \in I$와 $z_j \in B \otimes_A R_i$, $j = 1, \ldots, m$을
택할 수 있다. 이 $z_j$들을 사용하여 사상
$v_i : C \otimes_A R_i \to B \otimes_A R_i$를 정의하고자 한다.
$K \otimes_A R_i \to R_i[x_1, \ldots, x_m] \to C \otimes_A R_i \to 0$
은 표시이므로,
% P01-E1454
각 $s = 1, \ldots, t$에 대하여 $\xi_s = f_j(z_1, \ldots, z_m)$가
$B \otimes_A R_i$에서 영임을 확인하면 충분하다. 그렇지 않을 수도
있지만, $\xi_s$의 $B \otimes_A R$에서의 상이 영이므로 $i$를 조금
증가시키면 그렇게 됨을 알 수 있다.

\medskip\noindent
(4)를 증명하기 위하여 $u \otimes 1$이 동형사상이고,
$B$가 유한형 $A$-대수이며, $C$가 유한 표시 $A$-대수라고 하자.
% P01-E1455
$v : B \otimes_A R \to C \otimes_A R$을 $u \otimes 1$의
역사상이라 하자. (3)과 같은
$v_i : C \otimes_A R_i \to B \otimes_A R_i$를 택하자.
(1)을 적용하면 $i$를 증가시킨 뒤
% P01-E1456
$v_i \circ (u \otimes 1) = \text{id}_{B \otimes_R R_i}$ 및
$(u \otimes 1) \circ v_i = \text{id}_{C \otimes_R R_i}$이다.
\end{proof}

\begin{lemma}
\label{algebra-lemma-colimit-category-fp-algebras}
$R = \colim_{\lambda \in \Lambda} R_\lambda$가 환들의 유향 여극한이라고
하자. 그러면 유한 표시 $R$-대수들의 범주는 유한 표시
$R_\lambda$-대수들의 범주들의 여극한이다. 더 정확히 말하면 다음과 같다.
\begin{enumerate}
\item 유한 표시 $R$-대수 $A$가 주어지면,
$A \cong A_\lambda \otimes_{R_\lambda} R$을 만족하는
$\lambda \in \Lambda$와 유한 표시 $R_\lambda$-대수 $A_\lambda$가
존재한다.
\item $\lambda \in \Lambda$, 유한 표시 $R_\lambda$-대수
$A_\lambda, B_\lambda$, 그리고 $R$-대수 준동형사상
$\varphi : A_\lambda \otimes_{R_\lambda} R \to B_\lambda \otimes_{R_\lambda} R$
이 주어지면, $\varphi = \varphi_\mu \otimes 1_R$을 만족하는
$\mu \geq \lambda$와 $R_\mu$-대수 준동형사상
$\varphi_\mu : A_\lambda \otimes_{R_\lambda} R_\mu \to
B_\lambda \otimes_{R_\lambda} R_\mu$가 존재한다.
\item $\lambda \in \Lambda$, 유한 표시 $R_\lambda$-대수
$A_\lambda, B_\lambda$, 그리고 $R_\lambda$-대수 준동형사상
$\varphi_\lambda, \psi_\lambda : A_\lambda \to B_\lambda$가 주어지고,
% P01-E1457
$\varphi \otimes 1_R = \psi \otimes 1_R$이면 어떤
$\mu \geq \lambda$에 대하여
$\varphi \otimes 1_{R_\mu} = \psi \otimes 1_{R_\mu}$이다.
\end{enumerate}
\end{lemma}

\begin{proof}
(1)을 증명하기 위하여 표시
$A = R[x_1, \ldots, x_n]/(f_1, \ldots, f_m)$을 택하자.
$f_j \in R[x_1, \ldots, x_n]$로 사상되는 원소
$f_{\lambda, j} \in R_\lambda[x_1, \ldots, x_n]$와
$\lambda \in \Lambda$를 택할 수 있다. 이제 단순히
$A_\lambda =
R_\lambda[x_1, \ldots, x_n]/(f_{\lambda, 1}, \ldots, f_{\lambda, m})$
라 두면 된다.

\medskip\noindent
(2)와 (3)은 보조정리 \ref{algebra-lemma-algebra-map-property-in-colimit}에서
따른다.
\end{proof}

\begin{lemma}
\label{algebra-lemma-limit-no-condition-local}
$R \to S$가 국소환들의 국소 준동형사상이라고 하자.
그러면 유향 집합 $(\Lambda, \leq)$와 국소환들의 국소 준동형사상
$R_\lambda \to S_\lambda$들의 계로서 다음을 만족하는 것이 존재한다.
\begin{enumerate}
\item 계 $R_\lambda \to S_\lambda$의 여극한은 $R \to S$와 같다.
\item 각 $R_\lambda$는 $\mathbf{Z}$ 위에서 본질적으로 유한형이다.
\item 각 $S_\lambda$는 $R_\lambda$ 위에서 본질적으로 유한형이다.
\end{enumerate}
\end{lemma}

\begin{proof}
% P01-E1458
환 준동형사상을 $\varphi : R  \to S$로 나타내자.
$\mathfrak m \subset R$을 $R$의 극대 아이디얼이라 하고
$\mathfrak n \subset S$를 $S$의 극대 아이디얼이라 하자. 다음과 같이 두자.
$$
\Lambda = \{
(A, B)
\mid
A \subset R, B \subset S, \# A < \infty, \# B < \infty, \varphi(A) \subset B
\}.
$$
부분순서로 포함관계를 취한다. 각
$\lambda = (A, B) \in \Lambda$에 대하여 $R'_\lambda$를
$a \in A$가 생성하는 부분 $\mathbf{Z}$-대수라 하고,
$S'_\lambda$를 $b$, $b \in B$가 생성하는 부분
$\mathbf{Z}$-대수라 하자. $R_\lambda$를 소 아이디얼
$R'_\lambda \cap \mathfrak m$에서 $R'_\lambda$를 국소화한 환이라 하고,
$S_\lambda$를 소 아이디얼 $S'_\lambda \cap \mathfrak n$에서
$S'_\lambda$를 국소화한 환이라 하자. 그림으로 쓰면
$$
\xymatrix{
B \ar[r] &
S'_\lambda \ar[r] &
S_\lambda \ar[r] &
S \\
A \ar[r] \ar[u] &
R'_\lambda \ar[r] \ar[u] &
R_\lambda \ar[r] \ar[u] &
R \ar[u]
}.
$$
전이 사상들은 명백하다. 나머지 명제들의 증명은 독자에게 맡긴다.
\end{proof}

\begin{lemma}
\label{algebra-lemma-limit-essentially-finite-type}
$R \to S$가 국소환들의 국소 준동형사상이라고 하자.
$S$가 $R$ 위에서 본질적으로 유한형이라고 가정하자.
그러면 유향 집합 $(\Lambda, \leq)$와 국소환들의 국소 준동형사상
$R_\lambda \to S_\lambda$들의 계로서 다음을 만족하는 것이 존재한다.
\begin{enumerate}
\item 계 $R_\lambda \to S_\lambda$의 여극한은 $R \to S$와 같다.
\item 각 $R_\lambda$는 $\mathbf{Z}$ 위에서 본질적으로 유한형이다.
\item 각 $S_\lambda$는 $R_\lambda$ 위에서 본질적으로 유한형이다.
\item 각 $\lambda \leq \mu$에 대하여 사상
$S_\lambda \otimes_{R_\lambda} R_\mu \to S_\mu$는 $S_\mu$를
$S_\lambda \otimes_{R_\lambda} R_\mu$의 한 몫의 국소화로 표시한다.
\end{enumerate}
\end{lemma}

\begin{proof}
% P01-E1458
환 준동형사상을 $\varphi : R  \to S$로 나타내자.
$\mathfrak m \subset R$을 $R$의 극대 아이디얼이라 하고
$\mathfrak n \subset S$를 $S$의 극대 아이디얼이라 하자.
$S$가 $x_1, \ldots, x_n$이 생성하는 $S$의 부분 $R$-대수를
국소화한 환이 되게 하는 원소 $x_1, \ldots, x_n \in S$를 택하자.
달리 말해 $S$는 다항식환 $R[x_1, \ldots, x_n]$의 한 국소화의
몫이다.

\medskip\noindent
$\Lambda = \{ A \subset R \mid \# A < \infty\}$를 $R$의 유한
부분집합들로 이루어진 집합이라 하자. 부분순서로 포함관계를 취한다.
각 $\lambda = A \in \Lambda$에 대하여 $R'_\lambda$를
$a \in A$가 생성하는 부분 $\mathbf{Z}$-대수라 하고,
$S'_\lambda$를 $\varphi(a)$, $a \in A$와 원소
$x_1, \ldots, x_n$이 생성하는 부분 $\mathbf{Z}$-대수라 하자.
$R_\lambda$를 소 아이디얼 $R'_\lambda \cap \mathfrak m$에서
$R'_\lambda$를 국소화한 환이라 하고, $S_\lambda$를 소 아이디얼
$S'_\lambda \cap \mathfrak n$에서 $S'_\lambda$를 국소화한 환이라 하자.
그림으로 쓰면
$$
\xymatrix{
\varphi(A) \amalg \{x_i\} \ar[r] &
S'_\lambda \ar[r] &
S_\lambda \ar[r] &
S \\
A \ar[r] \ar[u] &
R'_\lambda \ar[r] \ar[u] &
R_\lambda \ar[r] \ar[u] &
R \ar[u]
}
$$
$A \subset B$가 $\Lambda$에서 $\lambda \leq \mu$에 대응하면
표준 사상 $R_\lambda \to R_\mu$와 $S_\lambda \to S_\mu$가
존재함은 명백하고, 따라서 유향 집합 $\Lambda$ 위의 계를 얻는다.

\medskip\noindent
모든 사상 $R_\lambda \to R$이 단사이고 $R$의 모든 원소는 결국
상에 들어가므로 $R = \colim R_\lambda$라는 명제는 명백하다.
$S = \colim S_\lambda$에도 같은 논증이 적용된다.
(2), (3)은 구성에 의해 참이다. 마지막 명제는 사상
$S'_\lambda \otimes_{R'_\lambda} R'_\mu
\to S'_\mu$가 명백히 전사이므로 성립한다.
\end{proof}

\begin{lemma}
\label{algebra-lemma-limit-essentially-finite-presentation}
$R \to S$가 국소환들의 국소 준동형사상이라고 하자.
$S$가 $R$ 위에서 본질적으로 유한 표시라고 가정하자.
그러면 유향 집합 $(\Lambda, \leq)$와
% P01-E1459
국소환들의 국소 준동형사상 $R_\lambda \to S_\lambda$들의 계로서
다음을 만족하는 것이 존재한다.
\begin{enumerate}
\item 계 $R_\lambda \to S_\lambda$의 여극한은 $R \to S$와 같다.
\item 각 $R_\lambda$는 $\mathbf{Z}$ 위에서 본질적으로 유한형이다.
\item 각 $S_\lambda$는 $R_\lambda$ 위에서 본질적으로 유한형이다.
\item 각 $\lambda \leq \mu$에 대하여 사상
$S_\lambda \otimes_{R_\lambda} R_\mu \to S_\mu$는 $S_\mu$를
$S_\lambda \otimes_{R_\lambda} R_\mu$의 한 소 아이디얼에서의
국소화로 표시한다.
\end{enumerate}
\end{lemma}

\begin{proof}
가정에 의해 동형사상
$\Phi : (R[x_1, \ldots, x_n]/I)_{\mathfrak q} \to S$를 택할 수 있다.
여기서 $I \subset R[x_1, \ldots, x_n]$은 유한 생성 아이디얼이고
$\mathfrak q \subset R[x_1, \ldots, x_n]/I$는 소 아이디얼이다.
($R \cap \mathfrak q$는 $R$의 극대 아이디얼 $\mathfrak m$과 같음에
유의하라.) 또한 $f_1, \ldots, f_m \in I$를 아이디얼 $I$의
생성원으로 택한다. 각 $R_\lambda$가 국소이고 $\mathbf{Z}$ 위에서
본질적으로 유한형인 유향 집합 $(\Lambda, \leq )$ 위의 여극한
$R = \colim R_\lambda$로 $R$을 임의의 방식으로 쓰자.
어떤 $\lambda_0 \in \Lambda$가 존재하여 $f_j$는 어떤
$f_{j, \lambda_0} \in R_{\lambda_0}[x_1, \ldots, x_n]$의 상이다.
모든 $\lambda \geq \lambda_0$에 대하여
% P01-E1460
$f_{j, \lambda} \in R_{\lambda}[x_1, \ldots, x_n]$를
$f_{j, \lambda_0}$의 상으로 나타내자. 따라서 환 준동형사상들의 계
$$
R_\lambda[x_1, \ldots, x_n]/(f_{1, \lambda}, \ldots, f_{m, \lambda})
\to
R[x_1, \ldots, x_n]/(f_1, \ldots, f_m) \to S
$$
를 얻는다.
% P01-E1461
$\mathfrak q_\lambda$를 $\mathfrak q$의 역상으로 두자.
$S_\lambda = (R_\lambda[x_1, \ldots, x_n]/
(f_{1, \lambda}, \ldots, f_{m, \lambda}))_{\mathfrak q_\lambda}$라 두자.
이것이 성립함을 확인하는 일은 독자에게 맡긴다.
\end{proof}

\begin{remark}
\label{algebra-remark-suitable-systems-limits}
$R \to S$가 본질적으로 유한 표시인 국소환들의 국소 준동형사상이라고
하자. 보조정리 \ref{algebra-lemma-limit-essentially-finite-type}에 열거된
성질들을 갖는 임의의 계 $(\Lambda, \leq)$,
$R_\lambda \to S_\lambda$를 택하자. 이 계가 ``잘못된'' 계일 수도
있다. 즉 보조정리 \ref{algebra-lemma-limit-essentially-finite-presentation}의
성질 (4)를 만족하지 않을 수도 있다. 다음은 한 예이다.
$k = \mathbf{F}_2$라 하자. 환
$$
R = \text{localization of }
k[z, y_1, y_2, \ldots]/(y_i^2 - zy_{i + 1})
\text{ at }(z, y_1, y_2, \ldots)
$$
을 생각하자. $S = R/zR$라 두자. 계로서 $\Lambda = \mathbf{N}$과
$$
R_n = \text{localization of }
k[z, y_1, \ldots, y_n]/(\{y_i^2 - zy_{i + 1}\}_{i \leq n-1})
\text{ at }(z, y_1, \ldots, y_n)
$$
을 취하고 $S_n = R_n/(z, y_n^2)$라 하자.
% P01-E1462
모든 사상 $S_n \otimes_{R_n} R_{n + 1} \to S_{n + 1}$은
$1 \otimes y_{n + 1}^2$이 영으로 사상되므로 국소화가 아니다
(즉, 이 경우에는 동형사상이 아니다).
그 대신 $S_n' = R_n/zR_n$을 취하면 사상
$S'_n \otimes_{R_n} R_{n + 1} \to S'_{n + 1}$은 동형사상이다.
이 주석의 요지는 계를 고를 때 실제로 조금 주의해야 한다는 것이다.
\end{remark}

\begin{lemma}
\label{algebra-lemma-limit-module-essentially-finite-presentation}
$R \to S$가 국소환들의 국소 준동형사상이라고 하자.
$S$가 $R$ 위에서 본질적으로 유한 표시라고 가정하자.
$M$을 유한 표시 $S$-가군이라 하자. 그러면 유향 집합
$(\Lambda, \leq)$와 국소환들의 국소 준동형사상
$R_\lambda \to S_\lambda$들의 계 및 $S_\lambda$-가군
$M_\lambda$들로서 다음을 만족하는 것이 존재한다.
\begin{enumerate}
\item 계 $R_\lambda \to S_\lambda$의 여극한은 $R \to S$와 같다.
계 $M_\lambda$의 여극한은 $M$이다.
\item 각 $R_\lambda$는 $\mathbf{Z}$ 위에서 본질적으로 유한형이다.
\item 각 $S_\lambda$는 $R_\lambda$ 위에서 본질적으로 유한형이다.
\item 각 $M_\lambda$는 $S_\lambda$ 위에서 유한이다.
\item 각 $\lambda \leq \mu$에 대하여 사상
$S_\lambda \otimes_{R_\lambda} R_\mu \to S_\mu$는 $S_\mu$를
$S_\lambda \otimes_{R_\lambda} R_\mu$의 한 소 아이디얼에서의
국소화로 표시한다.
\item 각 $\lambda \leq \mu$에 대하여 사상
$M_\lambda \otimes_{S_\lambda} S_\mu \to M_\mu$는 동형사상이다.
\end{enumerate}
\end{lemma}

\begin{proof}
보조정리 \ref{algebra-lemma-limit-essentially-finite-presentation}의 증명에서처럼,
먼저 $R = \colim R_\lambda$를 본질적으로 유한형인 국소
$\mathbf{Z}$-대수들의 유향 여극한으로 쓸 수 있다.
다음으로 어떤 $\lambda_1 \in \Lambda$에 대하여
$f_{j, \lambda_1} \in R_{\lambda_1}[x_1, \ldots, x_n]$들이 존재하여
$$
S =
\colim_{\lambda \geq \lambda_1} S_\lambda, \text{ with }
S_\lambda =
(R_\lambda[x_1, \ldots, x_n]/
(f_{1, \lambda}, \ldots, f_{m, \lambda}))_{\mathfrak q_\lambda}
$$
라고 가정할 수 있다. $S$ 위의 $M$의 표시
$$
S^{\oplus s} \to S^{\oplus t} \to M \to 0
$$
을 택하자. $A \in \text{Mat}(t \times s, S)$를 이 표시의 행렬이라
하자. 어떤 $\lambda_2 \in \Lambda$, $\lambda_2 \geq \lambda_1$에
대하여 $A$로 사상되는 행렬
$A_{\lambda_2} \in \text{Mat}(t \times s, S_{\lambda_2})$를 찾을 수
있다. 모든 $\lambda \geq \lambda_2$에 대하여
$M_\lambda = \Coker(S_\lambda^{\oplus s} \xrightarrow{A_\lambda}
S_\lambda^{\oplus t})$라 두자. 이것이 성립함을 확인하는 일은
독자에게 맡긴다.
\end{proof}

\begin{lemma}
\label{algebra-lemma-limit-no-condition}
$R \to S$를 환 준동형사상이라 하자. 그러면 유향 집합
$(\Lambda, \leq)$와 환 준동형사상 $R_\lambda \to S_\lambda$들의
계로서 다음을 만족하는 것이 존재한다.
\begin{enumerate}
\item 계 $R_\lambda \to S_\lambda$의 여극한은 $R \to S$와 같다.
\item 각 $R_\lambda$는 $\mathbf{Z}$ 위에서 유한형이다.
\item 각 $S_\lambda$는 $R_\lambda$ 위에서 유한형이다.
\end{enumerate}
\end{lemma}

\begin{proof}
이는 보조정리 \ref{algebra-lemma-limit-no-condition-local}의 비국소 버전이다.
증명은 비슷하므로 독자에게 맡긴다.
\end{proof}

\begin{lemma}
\label{algebra-lemma-limit-integral}
$R \to S$를 환 준동형사상이라 하자.
$S$가 $R$ 위에서 정수적이라고 가정하자. 그러면 유향 집합
$(\Lambda, \leq)$와 환 준동형사상 $R_\lambda \to S_\lambda$들의
계로서 다음을 만족하는 것이 존재한다.
\begin{enumerate}
\item 계 $R_\lambda \to S_\lambda$의 여극한은 $R \to S$와 같다.
\item 각 $R_\lambda$는 $\mathbf{Z}$ 위에서 유한형이다.
\item 각 $S_\lambda$는 $R_\lambda$ 위에서 유한이다.
\end{enumerate}
\end{lemma}

\begin{proof}
$E \subset R$이 유한 부분집합이고 $F \subset S$가 유한 부분집합이며,
모든 원소 $f \in F$가 계수들이 $E$에 속하는 어떤 일계수 다항식
$P(X) \in R[X]$의 근인 쌍 $(E, F)$들의 집합 $\Lambda$를 생각하자.
$E \subset E'$ 및 $F \subset F'$일 때
$(E, F) \leq (E', F')$라 하자.
$\lambda = (E, F) \in \Lambda$가 주어지면 $R_\lambda \subset R$을
$E$가 생성하는 $R$의 $\mathbf{Z}$-부분대수와 같게 두고,
$S_\lambda \subset S$를 $F$ 및 $S$ 안에서의 $E$의 상이 생성하는
$\mathbf{Z}$-부분대수와 같게 두자. $R = \colim R_\lambda$임은
명백하다.
% P01-E1463
$S$의 모든 원소가 $S$ 위에서 정수적이므로
$S = \colim S_\lambda$이다. 환 준동형사상
$R_\lambda \to S_\lambda$가 유한인 것은 보조정리
\ref{algebra-lemma-characterize-finite-in-terms-of-integral}와, 쌍 $(E, F)$에
대한 조건에 의해 $R_\lambda$ 위에서 정수적인 $F$의 원소들이
$S_\lambda$를 $R_\lambda$ 위에서 생성한다는 사실에서 따른다.
따라서 보조정리가 성립한다.
\end{proof}

\begin{lemma}
\label{algebra-lemma-limit-finite-type}
$R \to S$를 환 준동형사상이라 하자.
$S$가 $R$ 위에서 유한형이라고 가정하자. 그러면 유향 집합
$(\Lambda, \leq)$와 환 준동형사상 $R_\lambda \to S_\lambda$들의
계로서 다음을 만족하는 것이 존재한다.
\begin{enumerate}
\item 계 $R_\lambda \to S_\lambda$의 여극한은 $R \to S$와 같다.
\item 각 $R_\lambda$는 $\mathbf{Z}$ 위에서 유한형이다.
\item 각 $S_\lambda$는 $R_\lambda$ 위에서 유한형이다.
\item 각 $\lambda \leq \mu$에 대하여 사상
$S_\lambda \otimes_{R_\lambda} R_\mu \to S_\mu$는 $S_\mu$를
$S_\lambda \otimes_{R_\lambda} R_\mu$의 몫으로 표시한다.
\end{enumerate}
\end{lemma}

\begin{proof}
이는 보조정리 \ref{algebra-lemma-limit-essentially-finite-type}의 비국소
버전이다. 증명은 비슷하므로 독자에게 맡긴다.
\end{proof}

\begin{lemma}
\label{algebra-lemma-limit-finite-presentation}
$R \to S$를 환 준동형사상이라 하자.
$S$가 $R$ 위에서 유한 표시라고 가정하자. 그러면 유향 집합
$(\Lambda, \leq)$와 환 준동형사상 $R_\lambda \to S_\lambda$들의
계로서 다음을 만족하는 것이 존재한다.
\begin{enumerate}
\item 계 $R_\lambda \to S_\lambda$의 여극한은 $R \to S$와 같다.
\item 각 $R_\lambda$는 $\mathbf{Z}$ 위에서 유한형이다.
\item 각 $S_\lambda$는 $R_\lambda$ 위에서 유한형이다.
\item 각 $\lambda \leq \mu$에 대하여 사상
$S_\lambda \otimes_{R_\lambda} R_\mu \to S_\mu$는 동형사상이다.
\end{enumerate}
\end{lemma}

\begin{proof}
이는 보조정리 \ref{algebra-lemma-limit-essentially-finite-presentation}의
비국소 버전이다. 증명은 비슷하므로 독자에게 맡긴다.
\end{proof}

\begin{lemma}
\label{algebra-lemma-limit-module-finite-presentation}
$R \to S$를 환 준동형사상이라 하자.
$S$가 $R$ 위에서 유한 표시라고 가정하자.
$M$을 유한 표시 $S$-가군이라 하자. 그러면 유향 집합
$(\Lambda, \leq)$와 환 준동형사상 $R_\lambda \to S_\lambda$들의
계 및 $S_\lambda$-가군 $M_\lambda$들로서 다음을 만족하는 것이
존재한다.
\begin{enumerate}
\item 계 $R_\lambda \to S_\lambda$의 여극한은 $R \to S$와 같다.
계 $M_\lambda$의 여극한은 $M$이다.
\item 각 $R_\lambda$는 $\mathbf{Z}$ 위에서 유한형이다.
\item 각 $S_\lambda$는 $R_\lambda$ 위에서 유한형이다.
\item 각 $M_\lambda$는 $S_\lambda$ 위에서 유한이다.
\item 각 $\lambda \leq \mu$에 대하여 사상
$S_\lambda \otimes_{R_\lambda} R_\mu \to S_\mu$는 동형사상이다.
\item 각 $\lambda \leq \mu$에 대하여 사상
$M_\lambda \otimes_{S_\lambda} S_\mu \to M_\mu$는 동형사상이다.
\end{enumerate}
특히 모든 $\lambda \in \Lambda$에 대하여
$$
M = M_\lambda \otimes_{S_\lambda} S
= M_\lambda \otimes_{R_\lambda} R.
$$
\end{lemma}

\begin{proof}
이는 보조정리 \ref{algebra-lemma-limit-module-essentially-finite-presentation}의
비국소 버전이다. 증명은 비슷하므로 독자에게 맡긴다.
\end{proof}

\section{평탄성의 추가 기준}
\label{algebra-section-more-flatness-criteria}

\noindent
다음 보조정리는 대수기하학에서 정규 곡면으로부터 매끄러운 곡면으로
가는 유한 사상이 평탄함을 보일 때 흔히 사용된다. 이는 보조정리
\ref{algebra-lemma-finite-flat-over-regular-CM}의 부분적인 역이다. 실제로 단사인
유한 국소환 준동형사상은 분명히 조건 (3)을 만족한다.

\begin{lemma}
\label{algebra-lemma-CM-over-regular-flat}
\begin{slogan}
기적의 평탄성
\end{slogan}
$R \to S$를 뇌터 국소환들의 국소 준동형사상이라 하자. 다음을 가정하자.
\begin{enumerate}
\item $R$은 정칙이고,
% P01-E0638: frozen assumption omits the finite verb "is".
\item $S$ Cohen--Macaulay이며,
\item $\dim(S) = \dim(R) + \dim(S/\mathfrak m_R S)$이다.
\end{enumerate}
그러면 $R \to S$는 평탄하다.
\end{lemma}

\begin{proof}
$\dim(R)$에 관하여 귀납한다. $\dim(R) = 0$인 경우에는 $R$이 체이므로
자명하다. $\dim(R) > 0$이라고 하자. (3)에 의해
$\dim(S) > 0$이다. $\mathfrak q_1, \ldots, \mathfrak q_r$를 $S$의
극소 소 아이디얼들이라 하자. 다음이 성립하므로
$$
\dim(S/\mathfrak q_i) = \dim(S) > \dim(S/\mathfrak m_R S)
$$
% P01-E0639: frozen prose lacks punctuation at this display boundary.
$\mathfrak q_i \not \supset \mathfrak m_R S$임에 유의하라. 첫 번째
등식은 보조정리 \ref{algebra-lemma-maximal-chain-CM}에 의하고, 부등식은 (3)에
의한다. 따라서 $\mathfrak p_i = R \cap \mathfrak q_i$는
$\mathfrak m_R$와 같지 않다. 보조정리 \ref{algebra-lemma-silly}를 사용하여
$x \in \mathfrak m_R$, $x \not \in \mathfrak m_R^2$, 그리고
$x \not \in \mathfrak p_i$인 원소를 택하자. 그러면 $x$는 $S$의
어떤 극소 소 아이디얼에도 들어 있지 않다. 따라서 (2)와 보조정리
\ref{algebra-lemma-reformulate-CM}에 의해 $x$는 $S$ 위의 비영인자이고,
$S/xS$는 $\dim(S/xS) = \dim(S) - 1$인 Cohen--Macaulay 환이다.
(1)과 보조정리 \ref{algebra-lemma-regular-ring-CM}에 의해 환 $R/xR$은
$\dim(R/xR) = \dim(R) - 1$인 정칙환이다. 귀납에 의해
$R/xR \to S/xS$는 평탄하다. 그러므로 보조정리
\ref{algebra-lemma-variant-local-criterion-flatness}와 그 뒤의 주석을 사용하여
결론을 얻는다.
\end{proof}

\begin{lemma}
\label{algebra-lemma-flat-over-regular}
$R \to S$를 뇌터 국소환들의 준동형사상이라 하자. $R$이 정칙
국소환이고, 한 정칙 매개변수계가 $S$의 정칙열로 사상된다고 하자.
그러면 $R \to S$는 평탄하다.
\end{lemma}

\begin{proof}
$x_1, \ldots, x_d$가 $R$의 매개변수계이고 $S$의 정칙열로 사상된다고
하자. $R/(x_1, \ldots, x_d)$가 체이므로
$S/(x_1, \ldots, x_d)S$는 그 위에서 평탄함에 유의하라.
이제 $x_d$는 $S/(x_1, \ldots, x_{d - 1})S$의 비영인자이므로,
평탄성의 국소 기준에 의해 $S/(x_1, \ldots, x_{d - 1})S$는
$R/(x_1, \ldots, x_{d - 1})$ 위에서 평탄하다(보조정리
\ref{algebra-lemma-variant-local-criterion-flatness}와 그 뒤의 주석을 보라).
그다음 $x_{d - 1}$은 $S/(x_1, \ldots, x_{d - 2})S$의 비영인자이므로,
평탄성의 국소 기준에 의해 $S/(x_1, \ldots, x_{d - 2})S$는
$R/(x_1, \ldots, x_{d - 2})$ 위에서 평탄하다(보조정리
\ref{algebra-lemma-variant-local-criterion-flatness}와 그 뒤의 주석을 보라).
이 과정을 계속하면 $S$가 $R$ 위에서 평탄하다는 결론에 이른다.
\end{proof}

\noindent
다음 보조정리는, 어떤 기초환 위의 유한 표시 환들 위의 유한 표시
가군에 관한 결과가 뇌터인 경우의 유한 가군에 관한 대응하는 결과에서
따름을 증명하는 핵심이다.

\begin{lemma}
\label{algebra-lemma-colimit-eventually-flat}
$R \to S$, $M$, $\Lambda$, $R_\lambda \to S_\lambda$, $M_\lambda$가
보조정리 \ref{algebra-lemma-limit-module-essentially-finite-presentation}과
같다고 하자. $M$이 $R$ 위에서 평탄하다고 가정하자.
그러면 어떤 $\lambda \in \Lambda$에 대하여 가군 $M_\lambda$는
$R_\lambda$ 위에서 평탄하다.
\end{lemma}

\begin{proof}
어떤 $\lambda \in \Lambda$를 택하고 다음을 생각하자.
$$
\text{Tor}_1^{R_\lambda}(M_\lambda, R_\lambda/\mathfrak m_\lambda)
=
\Ker(\mathfrak m_\lambda \otimes_{R_\lambda} M_\lambda
\to M_\lambda).
$$
Remark \ref{algebra-remark-Tor-ring-mod-ideal}을 보라. 우변을 보면 이것이 유한
생성 $S_\lambda$-가군임을 알 수 있다($S_\lambda$가 뇌터이고 관련된
가군들이 유한이기 때문이다). $\xi_1, \ldots, \xi_n$을 생성원이라
하자. $M$이 $R$ 위에서 평탄하므로
$0 = \Ker(\mathfrak m_\lambda R \otimes_R M \to M)$이다.
$\otimes$는 여극한과 가환하므로 각 $\xi_i$가
$\mathfrak m_{\lambda}R_{\lambda'} \otimes_{R_{\lambda'}} M_{\lambda'}$
의 영으로 사상되게 하는 $\lambda' \geq \lambda$가 존재한다.
따라서
$$
\text{Tor}_1^{R_\lambda}(M_\lambda, R_\lambda/\mathfrak m_\lambda)
\longrightarrow
\text{Tor}_1^{R_{\lambda'}}(M_{\lambda'},
R_{\lambda'}/\mathfrak m_{\lambda}R_{\lambda'})
$$
는 영 사상이다. 마지막 환이 체이므로
$M_\lambda \otimes_{R_\lambda} R_\lambda/\mathfrak m_\lambda$는
$R_\lambda/\mathfrak m_\lambda$ 위에서 평탄함에 유의하라. 따라서
보조정리 \ref{algebra-lemma-another-variant-local-criterion-flatness}를 적용하면
$M_{\lambda'}$가 $R_{\lambda'}$ 위에서 평탄함을 얻는다.
\end{proof}

\noindent
위 보조정리를 이용하면 뇌터가 아닐 때에도 Section
\ref{algebra-section-criteria-flatness}의 결과들을 다시 증명하기 시작할 수 있다.

\begin{lemma}
\label{algebra-lemma-mod-injective-general}
$R \to S$가 국소환들의 국소 준동형사상이라고 하자.
% P01-E0640: three frozen declarations omit "by/as".
$\mathfrak m$을 $R$의 극대 아이디얼로 나타내자.
$u : M \to N$을 $S$-가군 준동형사상이라 하자. 다음을 가정하자.
\begin{enumerate}
\item $S$는 $R$ 위에서 본질적으로 유한 표시이고,
\item $M$, $N$은 $S$ 위에서 유한 표시이며,
\item $N$은 $R$ 위에서 평탄하고,
\item $\overline{u} : M/\mathfrak mM \to N/\mathfrak mN$은 단사이다.
\end{enumerate}
그러면 $u$는 단사이고 $N/u(M)$은 $R$ 위에서 평탄하다.
\end{lemma}

\begin{proof}
보조정리 \ref{algebra-lemma-limit-module-essentially-finite-presentation}과 그 증명에
의해, 국소환 준동형사상들의 계 $R_\lambda \to S_\lambda$ 및
$S_\lambda$-가군 준동형사상
$u_\lambda : M_\lambda \to N_\lambda$를 찾을 수 있다. 이들은
% P01-E0641: frozen final modifier is attached to the wrong noun.
그 보조정리의 결론 (1)--(6)을 $N$과 $M$ 각각에 대하여 만족하고,
준동형사상 $u_\lambda$들의 여극한은 $u$이다. 보조정리
\ref{algebra-lemma-colimit-eventually-flat}에 의해 충분히 큰 모든
$\lambda$에 대하여 $N_\lambda$가 $R_\lambda$ 위에서 평탄하다고
가정할 수 있다. $\mathfrak m_\lambda \subset R_\lambda$를 극대
아이디얼이라 하고 $\kappa_\lambda = R_\lambda / \mathfrak m_\lambda$,
resp.\ $\kappa = R/\mathfrak m$를 잉여체라 하자.

\medskip\noindent
사상
$$
\Psi_\lambda :
M_\lambda/\mathfrak m_\lambda M_\lambda \otimes_{\kappa_\lambda} \kappa
\longrightarrow
M/\mathfrak m M.
$$
을 생각하자. $S_\lambda/\mathfrak m_\lambda S_\lambda$는 체
$\kappa_\lambda$ 위에서 본질적으로 유한형이므로 텐서곱
$S_\lambda/\mathfrak m_\lambda S_\lambda \otimes_{\kappa_\lambda} \kappa$
는 $\kappa$ 위에서 본질적으로 유한형이다. 따라서 이는 뇌터 환이고,
$\Psi_\lambda$의 핵은 유한 생성임을 얻는다.
$M/\mathfrak m M$은 계 $M_\lambda/\mathfrak m_\lambda M_\lambda$의
여극한이고 $\kappa$는 체 $\kappa_\lambda$들의 여극한이므로,
% P01-E0642: a directed preorder guarantees only a later index >= lambda.
$\Psi_\lambda$의 핵이 다음 사상의 핵에 의해 생성되게 하는
$\lambda' > \lambda$가 존재한다.
$$
\Psi_{\lambda, \lambda'} :
M_\lambda/\mathfrak m_\lambda M_\lambda
\otimes_{\kappa_\lambda}
\kappa_{\lambda'}
\longrightarrow
M_{\lambda'}/\mathfrak m_{\lambda'} M_{\lambda'}.
$$
구성에 의해 곱셈적 부분집합
$W \subset S_\lambda \otimes_{R_\lambda} R_{\lambda'}$가 존재하여
$S_{\lambda'} = W^{-1}(S_\lambda \otimes_{R_\lambda} R_{\lambda'})$이고
$$
W^{-1}(M_\lambda/\mathfrak m_\lambda M_\lambda
\otimes_{\kappa_\lambda}
\kappa_{\lambda'})
=
M_{\lambda'}/\mathfrak m_{\lambda'} M_{\lambda'}.
$$
이제 $x$가 다음 사상의 핵의 원소라고 하자.
$$
\Psi_{\lambda'} :
M_{\lambda'}/\mathfrak m_{\lambda'} M_{\lambda'}
\otimes_{\kappa_{\lambda'}} \kappa
\longrightarrow
M/\mathfrak m M.
$$
그러면 어떤 $w \in W$에 대하여
% P01-E0643: frozen membership omits its tensor base and representative map.
$wx \in M_\lambda/\mathfrak m_\lambda M_\lambda \otimes \kappa$이다.
따라서 $wx \in \Ker(\Psi_\lambda)$이다. 그러므로 $wx$는
$\Psi_{\lambda, \lambda'}$의 핵에 속하는 원소들의 선형결합이다.
따라서 $M_{\lambda'}/\mathfrak m_{\lambda'} M_{\lambda'}
\otimes_{\kappa_{\lambda'}} \kappa$에서 $wx = 0$이고,
$w$는 $S_{\lambda'}$에서 가역이므로 $x = 0$이다.
충분히 큰 모든 $\lambda'$에 대하여 $\Psi_{\lambda'}$의 핵이 영임을
얻었다!

\medskip\noindent
앞 문단의 결과에 의해 충분히 큰 모든 $\lambda$에 대하여
$\Psi_\lambda$의 핵이 영이라고 가정할 수 있다. 이는 사상
$M_\lambda/\mathfrak m_\lambda M_\lambda \to M/\mathfrak m M$이
단사임을 함의한다. 여기에 $\overline{u}$가 단사라는 사실을 결합하면
$\overline{u_\lambda} : M_\lambda/\mathfrak m_\lambda M_\lambda
\to N_\lambda/\mathfrak m_\lambda N_\lambda$도 단사임을 형식적으로
얻는다. 보조정리 \ref{algebra-lemma-mod-injective}에 의해 (충분히 큰 모든
$\lambda$에 대하여) 사상 $u_\lambda$는 단사이고
$N_\lambda/u_\lambda(M_\lambda)$는 $R_\lambda$ 위에서 평탄하다.
따라서 보조정리가 성립한다.
\end{proof}

\begin{lemma}
\label{algebra-lemma-grothendieck-general}
$R \to S$가 국소환들의 국소환 준동형사상이라고 하자.
% P01-E0640
$\mathfrak m$을 $R$의 극대 아이디얼로 나타내자. 다음을 가정하자.
\begin{enumerate}
\item $S$는 $R$ 위에서 본질적으로 유한 표시이고,
\item $S$는 $R$ 위에서 평탄하며,
\item $f \in S$는 $S/{\mathfrak m}S$의 비영인자이다.
\end{enumerate}
그러면 $S/fS$는 $R$ 위에서 평탄하고 $f$는 $S$의 비영인자이다.
\end{lemma}

\begin{proof}
% P01-E0645: frozen proof is a sentence fragment without a subject.
이는 보조정리 \ref{algebra-lemma-mod-injective-general}에서 바로 따른다.
\end{proof}

\begin{lemma}
\label{algebra-lemma-grothendieck-regular-sequence-general}
$R \to S$가 국소환들의 국소환 준동형사상이라고 하자.
% P01-E0640
$\mathfrak m$을 $R$의 극대 아이디얼로 나타내자. 다음을 가정하자.
\begin{enumerate}
\item $R \to S$는 본질적으로 유한 표시이고,
\item $R \to S$는 평탄하며,
\item $f_1, \ldots, f_c$는 $S$의 원소들의 열로서, 그 상
$\overline{f}_1, \ldots, \overline{f}_c$가
$S/{\mathfrak m}S$에서 정칙열을 이룬다.
\end{enumerate}
그러면 $f_1, \ldots, f_c$는 $S$의 정칙열이고, 각 몫
$S/(f_1, \ldots, f_i)$는 $R$ 위에서 평탄하다.
\end{lemma}

\begin{proof}
% P01-E0646: frozen proof is a noun phrase rather than a sentence.
귀납과 보조정리 \ref{algebra-lemma-grothendieck-general}을 사용한다.
\end{proof}

\noindent
다음은 국소적으로 유한 표시인 국소환 준동형사상의 경우에 대한
평탄성의 국소 기준이다.

\begin{lemma}
\label{algebra-lemma-variant-local-criterion-flatness-general}
$R \to S$를 국소환들의 국소 준동형사상이라 하자.
$I \not = R$을 $R$의 아이디얼이라 하고 $M$을 $S$-가군이라 하자.
다음을 가정하자.
\begin{enumerate}
\item $S$는 $R$ 위에서 본질적으로 유한 표시이고,
\item $M$은 $S$ 위에서 유한 표시이며,
\item $\text{Tor}_1^R(M, R/I) = 0$이고,
\item $M/IM$은 $R/I$ 위에서 평탄하다.
\end{enumerate}
그러면 $M$은 $R$ 위에서 평탄하다.
\end{lemma}

\begin{proof}
$\Lambda$, $R_\lambda \to S_\lambda$, $M_\lambda$가 보조정리
\ref{algebra-lemma-limit-module-essentially-finite-presentation}과 같다고 하자.
% P01-E0647: frozen designation omits "by/as".
$I_\lambda \subset R_\lambda$를 $I$의 역상으로 나타내자.
% P01-E0648: frozen quotient-system map has the unquotiented source.
이 경우 계 $R/I \to S/IS$, $M/IM$,
$R_\lambda \to S_\lambda/I_\lambda S_\lambda$, 그리고
$M_\lambda/I_\lambda M_\lambda$도 보조정리
\ref{algebra-lemma-limit-module-essentially-finite-presentation}의 결론을 만족한다.
따라서 보조정리 \ref{algebra-lemma-colimit-eventually-flat}에 의해 (지표집합
$\Lambda$를 축소한 뒤) 모든 $\lambda$에 대하여
% P01-E0649: frozen sentence does not name the flatness base.
$M_\lambda/I_\lambda M_\lambda$가 평탄하다고 가정할 수 있다.
어떤 $\lambda$를 택하고 다음을 생각하자.
$$
\text{Tor}_1^{R_\lambda}(M_\lambda, R_\lambda/I_\lambda)
=
\Ker(I_\lambda \otimes_{R_\lambda} M_\lambda
\to M_\lambda).
$$
Remark \ref{algebra-remark-Tor-ring-mod-ideal}을 보라. 우변을 보면 이것이 유한
생성 $S_\lambda$-가군임을 알 수 있다($S_\lambda$가 뇌터이고 관련된
가군들이 유한이기 때문이다). $\xi_1, \ldots, \xi_n$을 생성원이라
하자. $\text{Tor}_1^R(M, R/I) = 0$이고 $\otimes$가 여극한과
가환하므로, 각 $\xi_i$가
$\text{Tor}_1^{R_{\lambda'}}(M_{\lambda'}, R_{\lambda'}/I_{\lambda'})$
의 영으로 사상되게 하는 $\lambda' \geq \lambda$가 존재한다.
다음 사상들의 합성
$$
\xymatrix{
R_{\lambda'} \otimes_{R_\lambda}
\text{Tor}_1^{R_\lambda}(M_\lambda, R_\lambda/I_\lambda)
\ar[d]^{\text{surjective by Lemma \ref{algebra-lemma-surjective-on-tor-one}}} \\
\text{Tor}_1^{R_\lambda}(M_\lambda, R_{\lambda'}/I_\lambda R_{\lambda'})
\ar[d]^{\text{surjective up to localization by
Lemma \ref{algebra-lemma-surjective-on-tor-one-trivial}}} \\
\text{Tor}_1^{R_{\lambda'}}(M_{\lambda'}, R_{\lambda'}/I_\lambda R_{\lambda'})
\ar[d]^{\text{surjective by Lemma \ref{algebra-lemma-surjective-on-tor-one}}} \\
\text{Tor}_1^{R_{\lambda'}}(M_{\lambda'}, R_{\lambda'}/I_{\lambda'}).
}
$$
은 표시된 이유들에 의해 국소화까지 전사이다. 국소화가 필요한 까닭은
$M_{\lambda'}$가 $M_\lambda \otimes_{R_\lambda} R_{\lambda'}$와 같지
않기 때문이다. 즉 이는 $M_\lambda \otimes_{S_\lambda} S_{\lambda'}$와
같고, $S_{\lambda'}$는
$S_{\lambda} \otimes_{R_\lambda} R_{\lambda'}$의 국소화이므로
% P01-E0650: frozen consequence clause is grammatically malformed.
국소화까지의 명제(또는 $S_{\lambda'}$와 텐서한 뒤의 명제)를 얻는다.
보조정리 \ref{algebra-lemma-surjective-on-tor-one}은 첫째와 셋째 화살표에
적용된다. 실제로 $M_\lambda/I_\lambda M_\lambda$는
$R_\lambda/I_\lambda$ 위에서 평탄하고,
$M_{\lambda'}/I_\lambda M_{\lambda'}$는 평탄 가군
$M_\lambda/I_\lambda M_\lambda$의 기저변환이므로
$R_{\lambda'}/I_\lambda R_{\lambda'}$ 위에서 평탄하다.
위에서 설명했듯이 합성은 생성원 $\xi_i$들을 영으로 보낸다. 마침내
$\text{Tor}_1^{R_{\lambda'}}(M_{\lambda'}, R_{\lambda'}/I_{\lambda'})$가
영임을 얻는다. 보조정리 \ref{algebra-lemma-variant-local-criterion-flatness}에
의해 이는 $M_{\lambda'}$가 $R_{\lambda'}$ 위에서 평탄함을 함의한다.
\end{proof}

\noindent
아래 보조정리를 보조정리 \ref{algebra-lemma-criterion-flatness-fibre-Noetherian}
(뇌터 국소환인 경우) 및 보조정리
\ref{algebra-lemma-criterion-flatness-fibre-nilpotent}(기초의 아이디얼이 멱영인
경우)과 비교하라.

\begin{lemma}[Crit\`ere de platitude par fibres]
\label{algebra-lemma-criterion-flatness-fibre}
$R$, $S$, $S'$을 국소환이라 하고
$R \to S \to S'$을 국소환 준동형사상이라 하자.
$M$을 $S'$-가군이라 하자. $\mathfrak m \subset R$을 극대 아이디얼이라
하자. 다음을 가정하자.
\begin{enumerate}
\item 환 준동형사상 $R \to S$와 $R \to S'$은 본질적으로 유한
표시이다.
\item 가군 $M$은 $S'$ 위에서 유한 표시이다.
\item 가군 $M$은 영이 아니다.
\item 가군 $M/\mathfrak mM$은 평탄 $S/\mathfrak mS$-가군이다.
\item 가군 $M$은 평탄 $R$-가군이다.
\end{enumerate}
그러면 $S$는 $R$ 위에서 평탄하고 $M$은 평탄 $S$-가군이다.
\end{lemma}

\begin{proof}
보조정리 \ref{algebra-lemma-limit-essentially-finite-presentation}의 증명에서처럼,
먼저 $R = \colim R_\lambda$를 본질적으로 유한형인 국소
$\mathbf{Z}$-대수들의 유향 여극한으로 쓸 수 있다.
% P01-E0651: frozen designation omits "by/as".
$\mathfrak p_\lambda$를 $R_\lambda$의 극대 아이디얼로 나타내자.
다음으로 어떤 $\lambda_1 \in \Lambda$에 대하여
$f_{j, \lambda_1} \in R_{\lambda_1}[x_1, \ldots, x_n]$들이 존재하여
$$
S =
\colim_{\lambda \geq \lambda_1} S_\lambda, \text{ with }
S_\lambda =
(R_\lambda[x_1, \ldots, x_n]/
(f_{1, \lambda}, \ldots, f_{u, \lambda}))_{\mathfrak q_\lambda}
$$
라고 가정할 수 있다. 어떤 $\lambda_2 \in \Lambda$,
$\lambda_2 \geq \lambda_1$에 대하여
$g_{j, \lambda_2} \in R_{\lambda_2}[x_1, \ldots, x_n, y_1, \ldots, y_m]$
들이 존재하고, 그 상
$\overline{g}_{j, \lambda_2} \in S_{\lambda_2}[y_1, \ldots, y_m]$에
대하여
$$
S' =
\colim_{\lambda \geq \lambda_2} S'_\lambda, \text{ with }
S'_\lambda =
(S_\lambda[y_1, \ldots, y_m]/
(\overline{g}_{1, \lambda}, \ldots,
\overline{g}_{v, \lambda}))_{\overline{\mathfrak q}'_\lambda}
$$
라고 할 수 있다. 이는 또한 다음을 함의함에 유의하라.
% P01-E0652: frozen presentation omits the defining f-relations.
$$
S'_\lambda =
(R_\lambda[x_1, \ldots, x_n, y_1, \ldots, y_m]/
(g_{1, \lambda}, \ldots, g_{v, \lambda}))_{\mathfrak q'_\lambda}
$$
$S'$ 위의 $M$의 표시
$$
(S')^{\oplus s} \to (S')^{\oplus t} \to M \to 0
$$
를 택하자. $A \in \text{Mat}(t \times s, S')$를 이 표시의 행렬이라
하자. 어떤 $\lambda_3 \in \Lambda$, $\lambda_3 \geq \lambda_2$에
대하여 $A$로 사상되는 행렬
% P01-E0653: frozen matrix ring uses S_lambda instead of S'_lambda.
$A_{\lambda_3} \in \text{Mat}(t \times s, S_{\lambda_3})$를 찾을 수
있다. 모든 $\lambda \geq \lambda_3$에 대하여
$M_\lambda = \Coker((S'_\lambda)^{\oplus s} \xrightarrow{A_\lambda}
(S'_\lambda)^{\oplus t})$라 두자.

\medskip\noindent
이 선택들에 대하여 각 $\lambda_3 \leq \lambda \leq \mu$마다
$S_\lambda \otimes_{R_{\lambda}} R_\mu \to S_\mu$는 국소화이고,
$S'_\lambda \otimes_{S_{\lambda}} S_\mu \to S'_\mu$는 국소화이며,
사상 $M_\lambda \otimes_{S'_\lambda} S'_\mu \to M_\mu$는
동형사상이다. 이는 또한
$S'_\lambda \otimes_{R_{\lambda}} R_\mu \to S'_\mu$가 국소화임을
함의한다. 따라서 $M$이 $R$ 위에서 평탄하므로 보조정리
\ref{algebra-lemma-colimit-eventually-flat}에 의해 충분히 큰 모든
$\lambda$에 대하여 가군 $M_\lambda$는 $R_\lambda$ 위에서
평탄하다. 더욱이 다음에 유의하라.
$
\mathfrak m = \colim \mathfrak p_\lambda
$,
$
S/\mathfrak mS = \colim S_\lambda/\mathfrak p_\lambda S_\lambda
$,
$
S'/\mathfrak mS' = \colim S'_\lambda/\mathfrak p_\lambda S'_\lambda
$,
그리고
$
M/\mathfrak mM = \colim M_\lambda/\mathfrak p_\lambda M_\lambda
$이다. 또한 각 $\lambda_3 \leq \lambda \leq \mu$에 대하여 위에
열거된 성질들로부터
$$
S'_\lambda/\mathfrak p_\lambda S'_\lambda
\otimes_{S_{\lambda}/\mathfrak p_\lambda S_\lambda}
S_\mu/\mathfrak p_\mu S_\mu
\longrightarrow
S'_\mu/\mathfrak p_\mu S'_\mu
$$
가 국소화이고, 사상
$$
M_\lambda / \mathfrak p_\lambda M_\lambda
\otimes_{S'_\lambda/\mathfrak p_\lambda S'_\lambda}
S'_\mu /\mathfrak p_\mu S'_\mu
\longrightarrow
M_\mu/\mathfrak p_\mu M_\mu
$$
가 동형사상임을 알 수 있다. 따라서 계
$(S_\lambda/\mathfrak p_\lambda S_\lambda \to
S'_\lambda/\mathfrak p_\lambda S'_\lambda,
M_\lambda/\mathfrak p_\lambda M_\lambda)$도 보조정리
\ref{algebra-lemma-limit-module-essentially-finite-presentation}과 같은 계이다.
$M/\mathfrak m M$이 $S/\mathfrak mS$ 위에서 평탄하다고 가정했으므로
보조정리 \ref{algebra-lemma-colimit-eventually-flat}을 다시 적용할 수 있고,
충분히 큰 모든 $\lambda$에 대하여
$M_\lambda/\mathfrak p_\lambda M_\lambda$가
$S_\lambda/\mathfrak p_\lambda S_\lambda$ 위에서 평탄함을 얻는다.
따라서 충분히 큰 $\lambda$에 대하여 자료
$R_\lambda \to S_\lambda \to S'_\lambda, M_\lambda$는 보조정리
\ref{algebra-lemma-criterion-flatness-fibre-Noetherian}의 가정을 만족한다.
그러한 $\lambda$를 택하자. 그러면
% P01-E0654: frozen proof asserts equality with an unlocalized base change.
$S = S_\lambda \otimes_{R_\lambda} R$은 $R$ 위에서 평탄하고,
$M = M_\lambda \otimes_{S_\lambda} S$는 $S$ 위에서 평탄하다
(평탄 가군의 기저변환은 평탄하기 때문이다).
\end{proof}

\noindent
다음은 ``crit\`ere de platitude par fibres'' 보조정리
\ref{algebra-lemma-criterion-flatness-fibre}의 쉬운 따름정리이다. 이런 종류의
더 많은 결과는 More on Flatness, Section
\ref{flat-section-introduction}을 보라.

\begin{lemma}
\label{algebra-lemma-criterion-flatness-fibre-fp-over-ft}
$R$, $S$, $S'$을 국소환이라 하고
$R \to S \to S'$을 국소환 준동형사상이라 하자.
$M$을 $S'$-가군이라 하자. $\mathfrak m \subset R$을 극대 아이디얼이라
하자. 다음을 가정하자.
\begin{enumerate}
\item $R \to S'$은 본질적으로 유한 표시이고,
\item $R \to S$는 본질적으로 유한형이며,
\item $M$은 $S'$ 위에서 유한 표시이고,
\item $M$은 영이 아니며,
\item $M/\mathfrak mM$은 평탄 $S/\mathfrak mS$-가군이고,
\item $M$은 평탄 $R$-가군이다.
\end{enumerate}
그러면 $S$는 $R$ 위에서 본질적으로 유한 표시이고 평탄하며,
$M$은 평탄 $S$-가군이다.
\end{lemma}

\begin{proof}
% P01-E0655: frozen proof assumes finite presentation instead of finite type.
$S$가 $R$ 위에서 본질적으로 유한 표시이므로,
$S = C_{\overline{\mathfrak q}}$가 되게 하는 유한형 $R$-대수
$C$가 존재한다.
$C = R[x_1, \ldots, x_n]/I$라 쓰자.
% P01-E0656: frozen declaration mixes "denote" with "be".
$\mathfrak q \subset R[x_1, \ldots, x_n]$를
$\overline{\mathfrak q}$에 대응하는 소 아이디얼이라 하자.
그러면 $B = R[x_1, \ldots, x_n]_{\mathfrak q}$는 $R$ 위에서
본질적으로 유한 표시이고 $J = IB$일 때 $S = B/J$임을 알 수 있다.
$B/\mathfrak mB$ 안에서의 $J$의 상 $\overline{J}$를 생성하는 상
$\overline{f}_i \in B/\mathfrak mB$을 갖는
$f_1, \ldots, f_k \in J$를 찾을 수 있다. 따라서
% P01-E0657: frozen plural existential is paired with one ideal symbol.
$B/J' \to B/J$가 동형사상
$$
(B/J') \otimes_R R/\mathfrak m \longrightarrow
B/J \otimes_R R/\mathfrak m = S/\mathfrak mS.
$$
을 유도하게 하는 유한 생성 아이디얼 $J' \subset J$가 존재한다.
위와 같은 임의의 $J'$에 대하여 보조정리
\ref{algebra-lemma-criterion-flatness-fibre}를 환 준동형사상
$$
R \longrightarrow B/J' \longrightarrow S'
$$
과 가군 $M$에 적용할 수 있다. 따라서 위와 같은 어떤 선택
$J'$에 대해서도 $B/J'$은 $R$ 위에서 평탄하다. 이제
% P01-E0658: frozen inclusion chain repeats J' instead of introducing J''.
$J' \subset J' \subset J$가 위와 같은 두 유한 생성 아이디얼이면,
$B/J' \to B/J''$은 평탄하고 본질적으로 유한 표시인
$R$-대수들 사이의 전사로서 $\mathfrak m$을 법으로 동형사상이다.
% P01-E0659: frozen consecutive relative clauses have conflicting antecedents.
따라서 보조정리 \ref{algebra-lemma-mod-injective-general}에 의해
$B/J' = B/J''$, 즉 $J' = J''$이다. 이는 명백히 $J$가 유한
생성임을 뜻한다. 즉 $S$는 $R$ 위에서 본질적으로 유한 표시이다.
따라서 보조정리 \ref{algebra-lemma-criterion-flatness-fibre}를
$R \to S \to S'$에 적용할 수 있고, 결론을 얻는다.
\end{proof}

\begin{lemma}[Crit\`ere de platitude par fibres: locally nilpotent case]
\label{algebra-lemma-criterion-flatness-fibre-locally-nilpotent}
환의 범주에서 다음 가환 도식이 주어졌다고 하자.
$$
\xymatrix{
S \ar[rr] & & S' \\
& R \ar[lu] \ar[ru]
}
$$
$I \subset R$을 국소 멱영 아이디얼이라 하고 $M$을 $S'$-가군이라
하자. 다음을 가정하자.
\begin{enumerate}
\item $R \to S$는 유한형이고,
\item $R \to S'$은 유한 표시이며,
\item $M$은 유한 표시 $S'$-가군이고,
\item $M/IM$은 $S/IS$-가군으로서 평탄하며,
\item $M$은 $R$-가군으로서 평탄하다.
\end{enumerate}
그러면 $M$은 평탄 $S$-가군이고,
$M \otimes_S \kappa(\mathfrak q)$가 영이 아닌 모든
$\mathfrak q \subset S$에 대하여 $S_\mathfrak q$는 $R$ 위에서
평탄하고 본질적으로 유한 표시이다.
\end{lemma}

\begin{proof}
$M \otimes_S \kappa(\mathfrak q)$가 영이 아니면
$S' \otimes_S \kappa(\mathfrak q)$도 영이 아니므로,
$\mathfrak q$ 위에 놓이는 소 아이디얼 $\mathfrak q' \subset S'$이
존재한다(보조정리 \ref{algebra-lemma-in-image}).
$\mathfrak p \subset R$을 $\Spec(R)$에서 $\mathfrak q$의 상이라 하자.
$I$가 국소 멱영이므로 $I \subset \mathfrak p$이고, 따라서
$M/\mathfrak p M$은 $S/\mathfrak pS$ 위에서 평탄하다.
그러므로 보조정리 \ref{algebra-lemma-criterion-flatness-fibre-fp-over-ft}를
$R_\mathfrak p \to S_\mathfrak q \to S'_{\mathfrak q'}$ 및
$M_{\mathfrak q'}$에 적용할 수 있다. $M_{\mathfrak q'}$가 $S$ 위에서
평탄하고 $S_\mathfrak q$가 $R$ 위에서 평탄하고 본질적으로 유한
표시임을 얻는다. $\mathfrak q'$은 $S'$의 임의의 소 아이디얼이었으므로
$M$도 $S$ 위에서 평탄함을 알 수 있다(보조정리
\ref{algebra-lemma-flat-localization}).
\end{proof}

\section{평탄 궤적의 열린성}
\label{algebra-section-open-flat}

\noindent
보조정리 \ref{algebra-lemma-colimit-eventually-flat}을 이용하여 뇌터인 경우로
축소한다. 뇌터인 경우는 Section \ref{algebra-section-complex-exact}에서 주어진
완전 복합체의 특징짓기를 사용하여 다룬다.

\begin{lemma}
\label{algebra-lemma-CM-dim-finite-type}
$k$를 체라 하자. $S$를 유한형 $k$-대수라 하자.
$f_1, \ldots, f_i$를 $S$의 원소라 하자. $S$가 차원 $d$인
Cohen--Macaulay 등차원 환이고
$\dim V(f_1, \ldots, f_i) \leq d - i$라고 가정하자. 그러면 등식이
성립하고, 모든 $V(f_1, \ldots, f_i)$의 소 아이디얼
$\mathfrak q$에 대하여
% P01-E0660: frozen plural subject takes the singular verb "forms".
% P01-E0661: frozen point/set preposition is malformed.
$f_1, \ldots, f_i$는 $S_{\mathfrak q}$의 정칙열을 이룬다.
\end{lemma}

\begin{proof}
$S$가 차원 $d$인 Cohen--Macaulay 등차원 환이면, 보조정리
\ref{algebra-lemma-disjoint-decomposition-CM-algebra}에 의해 $S$의 모든 극대
아이디얼 $\mathfrak m$에 대하여 $\dim(S_{\mathfrak m}) = d$이다.
Proposition \ref{algebra-proposition-CM-module}에 의해 모든 극대 아이디얼
$\mathfrak m \in V(f_1, \ldots, f_i)$에 대하여 이 열은
$S_{\mathfrak m}$의 정칙열이고, 국소환
$S_{\mathfrak m}/(f_1, \ldots, f_i)$는 차원 $d - i$인
Cohen--Macaulay 환이다. 이는 실제로 $S/(f_1, \ldots, f_i)$가
차원 $d - i$인 Cohen--Macaulay 등차원 환이라는 뜻이다.
\end{proof}

\begin{lemma}
\label{algebra-lemma-open-regular-sequence}
$R \to S$를 유한형 환 준동형사상이라 하자. 모든 섬유
$S \otimes_R \kappa(\mathfrak p)$가 차원 $d$인 Cohen--Macaulay
등차원 환이 되게 하는 정수 $d$가 있다고 하자.
$f_1, \ldots, f_i$를 $S$의 원소라 하자. 집합
$$
\{ \mathfrak q \in V(f_1, \ldots, f_i)
\mid f_1, \ldots, f_i
\text{ are a regular sequence in }
S_{\mathfrak q}/\mathfrak p S_{\mathfrak q}
\text{ where }\mathfrak p = R \cap \mathfrak q
\}
$$
은 $V(f_1, \ldots, f_i)$에서 열린집합이다.
\end{lemma}

\begin{proof}
$\overline{S} = S/(f_1, \ldots, f_i)$라 쓰자.
$\mathfrak q$가 보조정리에서 정의된 집합의 원소이고
$\mathfrak p$가 이에 대응하는 $R$의 소 아이디얼이라고 하자.
Definition \ref{algebra-definition-relative-dimension}에서 정의된 상대 차원을
사용할 것이다. 먼저 보조정리
\ref{algebra-lemma-dimension-at-a-point-finite-type-field}에 의해
$d = \dim_{\mathfrak q}(S/R) =
\dim(S_{\mathfrak q}/\mathfrak pS_{\mathfrak q}) +
\text{trdeg}_{\kappa(\mathfrak p)}\ \kappa(\mathfrak q)$임에 유의하라.
$f_1, \ldots, f_i$가 뇌터 국소환
$S_{\mathfrak q}/\mathfrak pS_{\mathfrak q}$에서 정칙열을 이루므로
보조정리 \ref{algebra-lemma-one-equation}에 의해
$\dim(\overline{S}_{\mathfrak q}/\mathfrak p\overline{S}_{\mathfrak q})
= \dim(S_{\mathfrak q}/\mathfrak pS_{\mathfrak q}) - i$이다.
보조정리 \ref{algebra-lemma-dimension-at-a-point-finite-type-field}에 의해
$\dim_{\mathfrak q}(\overline{S}/R)
= \dim(\overline{S}_{\mathfrak q}/\mathfrak p\overline{S}_{\mathfrak q}) +
\text{trdeg}_{\kappa(\mathfrak p)}\ \kappa(\mathfrak q)
= d - i$라는 결론을 얻는다. 보조정리
\ref{algebra-lemma-dimension-fibres-bounded-open-upstairs}에 의해
$\mathfrak q$의 한 근방에 속하는 모든
$\mathfrak q' \in V(f_1, \ldots, f_i) = \Spec(\overline{S})$에 대하여
$\dim_{\mathfrak q'}(\overline{S}/R) \leq d - i$이다.
따라서 어떤 $g \in S$, $g \not \in \mathfrak q$에 대하여 $S$를
$S_g$로 바꾸고 나면 이 부등식이 모든 $\mathfrak q'$에 대하여
성립한다고 가정할 수 있다. 결과는 보조정리
\ref{algebra-lemma-CM-dim-finite-type}에서 따른다.
\end{proof}

\begin{lemma}
\label{algebra-lemma-exact-on-fibres-open}
$R \to S$를 환 준동형사상이라 하자. 유한 자유 $S$-가군들의 유한
호몰로지 복합체를 생각하자.
$$
F_{\bullet} :
0
\to
S^{n_e}
\xrightarrow{\varphi_e}
S^{n_{e-1}}
\xrightarrow{\varphi_{e-1}}
\ldots
\xrightarrow{\varphi_{i + 1}}
S^{n_i}
\xrightarrow{\varphi_i}
S^{n_{i-1}}
\xrightarrow{\varphi_{i-1}}
\ldots
\xrightarrow{\varphi_1}
S^{n_0}
$$
$S$의 모든 소 아이디얼 $\mathfrak q$에 대하여 복합체
$\overline{F}_{\bullet, \mathfrak q} =
F_{\bullet, \mathfrak q} \otimes_R \kappa(\mathfrak p)$를 생각하자.
여기서 $\mathfrak p$는 $R$에서 $\mathfrak q$의 역상이다.
$R$이 뇌터이고 어떤 정수 $d$가 존재하여 $R \to S$가 유한형이고
평탄하며, 섬유 $S \otimes_R \kappa(\mathfrak p)$가 차원 $d$인
Cohen--Macaulay 환이라고 가정하자.
% P01-E0662: frozen coordinated hypothesis omits finite verbs.
그러면 집합
$$
\{\mathfrak q \in \Spec(S) \mid
\overline{F}_{\bullet, \mathfrak q}\text{ is exact}\}
$$
은 $\Spec(S)$에서 열린집합이다.
\end{lemma}

\begin{proof}
$\mathfrak q$를 보조정리에서 정의된 집합의 원소라 하자.
Proposition \ref{algebra-proposition-what-exact}을 이용하여
$D(g)$가 보조정리에서 정의된 집합에 포함되게 하는
$g \in S$, $g \not \in \mathfrak q$가 존재함을 보일 것이다.
달리 말해 $S$를 $S_g$로 바꾸고 나면 보조정리의 집합이
$\Spec(S)$ 전체가 됨을 보일 것이다. 따라서 증명 도중 유한 번
$S$를 이러한 국소화로 바꿀 것이다. 환이 국소환일 때 Proposition
\ref{algebra-proposition-what-exact}은 복합체의 완전성을 사상 $\varphi_i$의
계수와 아이디얼 $I(\varphi_i)$를 이용하여 특징짓는다는 것을
상기하자. 먼저 계수 조건을 다룬다.
$r_i = n_i - n_{i + 1} + \ldots + (-1)^{e - i} n_e$라 두자.
$r_i + r_{i + 1} = n_i$이고, (완전한 경우) $r_i$는
$\varphi_i$의 기대 계수임에 유의하라.

\medskip\noindent
보조정리 \ref{algebra-lemma-complex-exact-mod}에 의해
$\overline{F}_{\bullet, \mathfrak q}$가 완전하면 국소화
$F_{\bullet, \mathfrak q}$도 완전하다. 특히 어떤
$g \in S$, $g \not \in \mathfrak q$로 국소화하면 복합체
$F_\bullet$이 완전해진다. 이 경우 $S$를 모든 소 아이디얼에서
국소화한 환들에 Proposition \ref{algebra-proposition-what-exact}을 적용하면
$\varphi_i$의 모든 $(r_i + 1) \times (r_i + 1)$-소행렬식이 영임을
알 수 있다. 따라서 $\varphi_i$의 계수는 $r_i$ 이하이다.

\medskip\noindent
$I_i \subset S$를 $\varphi_i$의 행렬의
$r_i \times r_i$-소행렬식들이 생성하는 아이디얼이라 하자.
Proposition \ref{algebra-proposition-what-exact}에 의해 복합체
$\overline{F}_{\bullet, \mathfrak q}$가 완전일 필요충분조건은 모든
$1 \leq i \leq e$에 대하여 $(I_i)_{\mathfrak q} = S_{\mathfrak q}$이거나
$(I_i)_{\mathfrak q}$가 길이 $i$인
$S_{\mathfrak q}/\mathfrak p S_{\mathfrak q}$-정칙열을 포함하는 것이다.
실제로 위의 $r_i$의 선택과 $\varphi_i$의 계수에 대한 경계에 의해
Proposition \ref{algebra-proposition-what-exact}의 조건이 성립하는 방법은
이것뿐이다.

\medskip\noindent
$(I_i)_{\mathfrak q} = S_{\mathfrak q}$이면
$g \not\in \mathfrak q$인 어떤 원소에서 $S$를 국소화한 뒤
$I_i = S$라고 가정할 수 있다. 이는 명백히 열린 조건이다.

\medskip\noindent
$(I_i)_{\mathfrak q} \not = S_{\mathfrak q}$이면
$S_{\mathfrak q}/\mathfrak pS_{\mathfrak q}$에서 정칙열을 이루는 열
% P01-E0663: frozen proof does not explicitly clear localized denominators.
$f_1, \ldots, f_i \in (I_i)_{\mathfrak q}$가 있다.
$(f_1, \ldots, f_i) \not \subset \mathfrak q'$인 모든 소 아이디얼
$\mathfrak q' \subset S$에 대하여
$(I_i)_{\mathfrak q'} = S_{\mathfrak q'}$임에 유의하라.
따라서 결과는 보조정리 \ref{algebra-lemma-open-regular-sequence}에서 따른다.
\end{proof}

\begin{theorem}
\label{algebra-theorem-openness-flatness}
$R$을 환이라 하자. $R \to S$를 유한 표시 환 준동형사상이라 하자.
$M$을 유한 표시 $S$-가군이라 하자. 집합
$$
\{ \mathfrak q \in \Spec(S) \mid
M_{\mathfrak q}\text{ is flat over }R\}
$$
은 $\Spec(S)$에서 열린집합이다.
\end{theorem}

\begin{proof}
$\mathfrak q \in \Spec(S)$를 소 아이디얼이라 하자.
$\mathfrak p \subset R$을 $R$에서 $\mathfrak q$의 역상이라 하자.
$M_{\mathfrak q}$가 $R$ 위에서 평탄일 필요충분조건은
$R_{\mathfrak p}$ 위에서 평탄인 것임에 유의하라.
$M_{\mathfrak q}$가 $R$ 위에서 평탄하다고 가정하자.
$M_g$가 $R$ 위에서 평탄하게 하는
$g \in S$, $g \not \in \mathfrak q$가 존재한다고 주장한다.

\medskip\noindent
먼저 $R$과 $S$가 $\mathbf{Z}$ 위에서 유한형인 경우로 축소한다.
보조정리 \ref{algebra-lemma-limit-module-finite-presentation}과 같은 유향 집합
$\Lambda$와 계 $(R_\lambda \to S_\lambda, M_\lambda)$를 택하자.
$\mathfrak p_\lambda$를 $R_\lambda$에서 $\mathfrak p$의 역상과 같게
두고, $\mathfrak q_\lambda$를 $S_\lambda$에서 $\mathfrak q$의
역상과 같게 두자. 그러면 계
$$
((R_\lambda)_{\mathfrak p_\lambda},
(S_\lambda)_{\mathfrak q_\lambda},
(M_\lambda)_{\mathfrak q_{\lambda}})
$$
는 보조정리 \ref{algebra-lemma-limit-module-essentially-finite-presentation}과
같은 계이다. 따라서 보조정리 \ref{algebra-lemma-colimit-eventually-flat}에 의해
어떤 $\lambda$에 대하여 가군 $M_\lambda$는 소 아이디얼
$\mathfrak q_{\lambda}$에서 $R_\lambda$ 위에 평탄하다.
계 $(R_\lambda \to S_\lambda, M_\lambda, \mathfrak q_{\lambda})$에
대하여 주장을 증명할 수 있다고 하자. 달리 말해
$(M_\lambda)_g$가 $R_\lambda$ 위에서 평탄하게 하는
$g \in S_\lambda$, $g \not\in \mathfrak q_\lambda$를 찾을 수 있다고
하자. 보조정리 \ref{algebra-lemma-limit-module-finite-presentation}에 의해
$M = M_\lambda \otimes_{R_\lambda} R$이고, 따라서
$M_g = (M_\lambda)_g \otimes_{R_\lambda} R$이기도 하다. 그러므로
보조정리 \ref{algebra-lemma-flat-base-change}에 의해 계
$(R \to S, M, \mathfrak q)$에 대한 주장을 얻는다.

\medskip\noindent
이제 $R$과 $S$가 $\mathbf{Z}$ 위에서 유한형이라고 가정할 수 있다.
$S$를 다항식환 $R[x_1, \ldots, x_n]$의 몫으로 쓸 수 있다.
물론 $S$를 $R[x_1, \ldots, x_n]$으로 바꾸어 $S$가 $R$ 위의
다항식환이라고 가정할 수 있다. 특히 $R \to S$는 평탄하고,
% P01-E0664: frozen phrase uses a plural noun attributively.
모든 섬유 환 $S \otimes_R \kappa(\mathfrak p)$의 전역 차원은
$n$임을 알 수 있다.

\medskip\noindent
보조정리 \ref{algebra-lemma-resolution-by-finite-free}를 사용하여, 각 $F_i$가
유한 자유인 $S$ 위의 $M$의 분해 $F_\bullet$을 택하자.
$K_n = \Ker(F_{n-1} \to F_{n-2})$라 하자. 각 $F_i$가 $R$ 위에서
평탄하고 $M$에 대한 가정이 있으므로, 보조정리
\ref{algebra-lemma-flat-ses}에 의해 $(K_n)_{\mathfrak q}$는 $R$ 위에서
평탄함에 유의하라. 또한 열
$$
0 \to
K_n/\mathfrak p K_n \to
F_{n-1}/ \mathfrak p F_{n-1} \to
\ldots \to
F_0 / \mathfrak p F_0 \to
M/\mathfrak p M \to
0
$$
은 $\mathfrak q$에서 국소화하면 완전하다. 이는
$\text{Tor}_i^{R_\mathfrak p}(\kappa(\mathfrak p), M_{\mathfrak q})$가
소멸하기 때문이다.
% P01-E0665: frozen proof claims exact global dimension n rather than a bound.
$S_\mathfrak q/\mathfrak p S_{\mathfrak q}$의 전역 차원이 $n$이므로,
$\mathfrak q$에서 국소화한 $K_n / \mathfrak p K_n$은
$S_\mathfrak q/\mathfrak p S_{\mathfrak q}$ 위의 유한 자유 가군임을
얻는다. 보조정리 \ref{algebra-lemma-free-fibre-flat-free}에 의해
$(K_n)_{\mathfrak q}$는 $S_{\mathfrak q}$ 위에서 자유이다.
특히 $(K_n)_g$가 $S_g$ 위에서 유한 자유가 되게 하는
$g \in S$, $g \not \in \mathfrak q$가 존재한다.

\medskip\noindent
보조정리 \ref{algebra-lemma-exact-on-fibres-open}에 의해 더 국소화한
$S_g$가 존재하여 복합체
$$
0 \to K_n \to F_{n-1} \to \ldots \to F_0
$$
는 $R \to S$의 {\it 모든 섬유에서} 완전하다. 보조정리
\ref{algebra-lemma-complex-exact-mod}에 의해 이는 $F_1 \to F_0$의 여핵이
평탄함을 함의한다. 이로써 뇌터인 경우의 정리를 증명하였다.
\end{proof}

\section{Cohen--Macaulay 궤적의 열린성}
\label{algebra-section-CM-open}

\noindent
이 절에서는 유한형 대수의 Cohen--Macaulay 성질을 평탄성으로 특징짓는다.
% P01-E0666: the frozen introduction omits “that” before its object clause.
그런 다음 이를 사용하여 이러한 대수가 Cohen--Macaulay가 되는 점들의
집합이 열려 있음을 증명한다.

\begin{lemma}
\label{algebra-lemma-where-CM}
$S$를 체 $k$ 위의 유한형 대수라 하자.
$\varphi : k[y_1, \ldots, y_d] \to S$를 준유한 환 준동형사상이라 하자.
$\Spec(S)$의 부분집합으로서 다음 등식이 성립한다.
% P01-E0667: the frozen left predicate omits “is” before “flat”.
% P01-E0668: the frozen right predicate omits “is” before “CM”.
$$
\{ \mathfrak q \mid
S_{\mathfrak q} \text{ flat over }k[y_1, \ldots, y_d]\}
=
\{ \mathfrak q \mid
S_{\mathfrak q} \text{ CM and }\dim_{\mathfrak q}(S/k) = d\}
$$
기호는 정의 \ref{algebra-definition-relative-dimension}을 보라.
\end{lemma}

\begin{proof}
$\mathfrak q \subset S$를 소 아이디얼이라 하자.
% P01-E0669: the frozen designation construction lacks a grammatical complement.
$\mathfrak p = k[y_1, \ldots, y_d] \cap \mathfrak q$라 쓰자.
예를 들어 보조정리 \ref{algebra-lemma-dimension-inequality-quasi-finite}에 의해 항상
$\dim(S_{\mathfrak q}) \leq \dim(k[y_1, \ldots, y_d]_{\mathfrak p})$
임에 유의하라. 또한 체확대
$\kappa(\mathfrak q)/\kappa(\mathfrak p)$는 유한이므로
$\text{trdeg}_k(\kappa(\mathfrak p)) = \text{trdeg}_k(\kappa(\mathfrak q))$
이다.

\medskip\noindent
$\mathfrak q$가 왼쪽 집합의 원소라고 하자. 그러면 보조정리
\ref{algebra-lemma-finite-flat-over-regular-CM}를 적용할 수 있으므로,
$S_{\mathfrak q}$는 Cohen--Macaulay이고
$\dim(S_{\mathfrak q}) = \dim(k[y_1, \ldots, y_d]_{\mathfrak p})$임을
얻는다. 이를 위의 초월차수 등식 및 보조정리
\ref{algebra-lemma-dimension-at-a-point-finite-type-field}과 결합하면
$\dim_{\mathfrak q}(S/k) = d$가 따른다. 따라서 $\mathfrak q$는
오른쪽 집합의 원소이다.

\medskip\noindent
$\mathfrak q$가 오른쪽 집합의 원소라고 하자. 위의 초월차수 등식,
가정 $\dim_{\mathfrak q}(S/k) = d$, 그리고 보조정리
\ref{algebra-lemma-dimension-at-a-point-finite-type-field}에 의해
$\dim(S_{\mathfrak q}) = \dim(k[y_1, \ldots, y_d]_{\mathfrak p})$를 얻는다.
따라서 보조정리 \ref{algebra-lemma-CM-over-regular-flat}를 적용할 수 있고,
$\mathfrak q$가 왼쪽 집합의 원소임을 알 수 있다.
\end{proof}

\begin{lemma}
\label{algebra-lemma-finite-type-over-field-CM-open}
$S$를 체 $k$ 위의 유한형 대수라 하자. $S_{\mathfrak q}$가
Cohen--Macaulay인 소 아이디얼 $\mathfrak q$들의 집합은 $S$에서 열린다.
% P01-E0670: the frozen statement names the ring S rather than Spec(S).
\end{lemma}

\noindent
이 보조정리는 아래의 보조정리
\ref{algebra-lemma-finite-presentation-flat-CM-locus-open}의 특수한 경우이다.
원한다면 곧바로 그 보조정리의 증명으로 건너뛰어도 된다.

\begin{proof}
$\mathfrak q \subset S$를 $S_{\mathfrak q}$가 Cohen--Macaulay인
소 아이디얼이라 하자. 환 $S_g$가 Cohen--Macaulay가 되게 하는
$g \in S$, $g \not \in \mathfrak q$가 존재함을 보여야 한다.
임의의 $g \in S$, $g \not \in \mathfrak q$에 대하여
$S$를 $S_g$로, $\mathfrak q$를 $\mathfrak qS_g$로 바꾸어도 된다.
이를 보조정리 \ref{algebra-lemma-Noether-normalization-at-point} 및
\ref{algebra-lemma-dimension-at-a-point-finite-type-field}과 결합하면
$d = \dim(S_{\mathfrak q}) + \text{trdeg}_k(\kappa(\mathfrak q))$이고
유한 단사인 환 준동형사상 $k[y_1, \ldots, y_d] \to S$가
존재한다고 가정할 수 있다.
$\mathfrak p = k[y_1, \ldots, y_d] \cap \mathfrak q$라 두자.
구성에 의해 $\mathfrak q$는 보조정리 \ref{algebra-lemma-where-CM}의
표시된 등식 오른쪽 집합의 원소이고, 따라서 왼쪽 집합의 원소이기도 하다.

\medskip\noindent
정리 \ref{algebra-theorem-openness-flatness}에 의해 어떤
$g \in S$, $g \not \in \mathfrak q$에 대하여 환 $S_g$는
$k[y_1, \ldots, y_d]$ 위에서 평탄하다. 따라서 보조정리
\ref{algebra-lemma-where-CM}의 등식을 다시 적용하면, 원하는 대로
$S_g$의 모든 국소환이 Cohen--Macaulay임을 얻는다.
\end{proof}

\begin{lemma}
\label{algebra-lemma-generic-CM}
$k$를 체라 하고, $S$를 유한형 $k$-대수라 하자.
% P01-E0671: the frozen finite-type compound lacks standard hyphenation.
Cohen--Macaulay인 소 아이디얼들의 집합은
조밀 열린집합 $U \subset \Spec(S)$를 이룬다.
\end{lemma}

\begin{proof}
보조정리 \ref{algebra-lemma-finite-type-over-field-CM-open}에 의해 이 집합은 열린다.
이는 모든 극소 소 아이디얼 $\mathfrak q \subset S$를 포함한다.
극소 소 아이디얼에서의 국소환 $S_{\mathfrak q}$는 차원이 영이고,
따라서 Cohen--Macaulay이기 때문이다.
\end{proof}

\begin{lemma}
\label{algebra-lemma-finite-presentation-flat-CM-locus-open}
$R$을 환이라 하자. $R \to S$가 유한 표시이고 평탄하다고 하자.
임의의 $d \geq 0$에 대하여 집합
$$
\left\{
\begin{matrix}
\mathfrak q \in \Spec(S)
\text{ such that setting }\mathfrak p = R \cap \mathfrak q
\text{ the fibre ring}\\
S_{\mathfrak q}/\mathfrak pS_{\mathfrak q}
\text{ is Cohen-Macaulay}
\text{ and } \dim_{\mathfrak q}(S/R) = d
\end{matrix}
\right\}
$$
은 $\Spec(S)$에서 열린다.
\end{lemma}

\begin{proof}
$\mathfrak q$를 표시된 집합의 원소라 하고, $\mathfrak p$를 이에 대응하는
$R$의 소 아이디얼이라 하자.
% P01-E0672: the frozen local goal omits the fixed relative-dimension condition.
$R \to S_g$의 모든 섬유환이 Cohen--Macaulay가 되게 하는
$g \in S$, $g \not \in \mathfrak q$를 찾아야 한다.
증명하는 동안 유한 번에 한하여, 어떤
$g \in S$, $g \not \in \mathfrak q$에 대해 $S$를 $S_g$로
바꾸어도 된다. 따라서 보조정리
\ref{algebra-lemma-quasi-finite-over-polynomial-algebra}에 의해
$d = \dim_{\mathfrak q}(S/R)$인 준유한 환 준동형사상
$R[t_1, \ldots, t_d] \to S$가 존재한다고 가정할 수 있다.
$\mathfrak q' = R[t_1, \ldots, t_d] \cap \mathfrak q$라 두자.
보조정리 \ref{algebra-lemma-where-CM}에 의해 환 준동형사상
$$
R[t_1, \ldots, t_d]_{\mathfrak q'} /
\mathfrak p R[t_1, \ldots, t_d]_{\mathfrak q'}
\longrightarrow
S_{\mathfrak q}/\mathfrak p S_{\mathfrak q}
$$
은 평탄하다. 따라서 섬유에 의한 평탄성 판정법인 보조정리
\ref{algebra-lemma-criterion-flatness-fibre}에 의해
$R[t_1, \ldots, t_d]_{\mathfrak q'} \to S_{\mathfrak q}$는 평탄하다.
그러므로 정리 \ref{algebra-theorem-openness-flatness}에 의해 어떤
% P01-E0673: the frozen sentence lacks punctuation after the point condition.
$g \in S$, $g \not \in \mathfrak q$에 대하여 환 준동형사상
$R[t_1, \ldots, t_d] \to S_g$는 평탄하다. $S$를 $S_g$로 바꾸면,
모든 소 아이디얼 $\mathfrak r \subset S$에 대하여
$\mathfrak r' = R[t_1, \ldots, t_d] \cap \mathfrak r$ 및
$\mathfrak p' = R \cap \mathfrak r$라 두었을 때 국소환 준동형사상
$R[t_1, \ldots, t_d]_{\mathfrak r'} \to S_{\mathfrak r}$는 평탄하다.
따라서 그 기저변환
$$
R[t_1, \ldots, t_d]_{\mathfrak r'} /
\mathfrak p' R[t_1, \ldots, t_d]_{\mathfrak r'}
\longrightarrow
S_{\mathfrak r}/\mathfrak p' S_{\mathfrak r}
$$
도 평탄하다. 그러므로 $k = \kappa(\mathfrak p')$로 놓고 보조정리
\ref{algebra-lemma-where-CM}를 적용하면 원하는 대로
$\mathfrak r$가 이 보조정리의 집합에 속함을 알 수 있다.
\end{proof}

\begin{lemma}
\label{algebra-lemma-generic-CM-flat-finite-presentation}
$R$을 환이라 하자. $R \to S$가 평탄하고 유한 표시라고 하자.
$\mathfrak p = R \cap \mathfrak q$일 때 섬유환
$S_{\mathfrak q} \otimes_R \kappa(\mathfrak p)$가 Cohen--Macaulay인
소 아이디얼 $\mathfrak q$들의 집합은
$\Spec(S) \to \Spec(R)$의 모든 섬유에서 열리고 조밀하다.
\end{lemma}

\begin{proof}
이 집합을 $W$라 하자. 보조정리
\ref{algebra-lemma-finite-presentation-flat-CM-locus-open}에 의해 이는 열린다.
$W$와 한 섬유의 교집합은 보조정리 \ref{algebra-lemma-generic-CM}를
적용할 수 있는 그 섬유의 해당 집합이므로, 이는 각 섬유에서 조밀하다.
\end{proof}

\begin{lemma}
\label{algebra-lemma-extend-field-CM-locus}
$k$를 체라 하자. $S$를 유한형 $k$-대수라 하자.
$K/k$를 체확대라 하고 $S_K = K \otimes_k S$라 두자.
$\mathfrak q \subset S$를 $S$의 소 아이디얼이라 하자.
$\mathfrak q_K \subset S_K$를 $\mathfrak q$ 위에 놓인
$S_K$의 소 아이디얼이라 하자. 그러면 $S_{\mathfrak q}$가
Cohen--Macaulay일 필요충분조건은
$(S_K)_{\mathfrak q_K}$가 Cohen--Macaulay인 것이다.
\end{lemma}

\begin{proof}
증명하는 동안 유한 번에 한하여, 임의의
$g \in S$, $g \not \in \mathfrak q$에 대해 $S$를 $S_g$로
바꾸어도 된다. 따라서 보조정리
\ref{algebra-lemma-Noether-normalization-at-point}를 사용하여
$\dim(S) = \dim_{\mathfrak q}(S/k) =: d$라고 가정하고
유한 단사 사상 $k[x_1, \ldots, x_d] \to S$를 찾을 수 있다.
기저변환에 의해 이것은 유한 단사 사상
$K[x_1, \ldots, x_d] \to S_K$도 유도함에 유의하라.
보조정리 \ref{algebra-lemma-dimension-at-a-point-preserved-field-extension}에 의해
$\dim_{\mathfrak q_K}(S_K/K) = d$이다.
$\mathfrak p = k[x_1, \ldots, x_d] \cap \mathfrak q$ 및
$\mathfrak p_K = K[x_1, \ldots, x_d] \cap \mathfrak q_K$라 두자.
다음의 뇌터 국소환들의 가환 그림을 생각하자.
$$
\xymatrix{
S_{\mathfrak q} \ar[r] &
(S_K)_{\mathfrak q_K} \\
k[x_1, \ldots, x_d]_{\mathfrak p} \ar[r] \ar[u] &
K[x_1, \ldots, x_d]_{\mathfrak p_K} \ar[u]
}
$$
보조정리 \ref{algebra-lemma-where-CM}에 의해 왼쪽 수직 화살표가 평탄일
필요충분조건이 오른쪽 수직 화살표가 평탄인 것임을 보여야 한다.
아래쪽 화살표가 평탄하므로 이 동치는 보조정리
\ref{algebra-lemma-base-change-flat-up-down}에 의해 성립한다.
\end{proof}

\begin{lemma}
\label{algebra-lemma-CM-locus-commutes-base-change}
$R$을 환이라 하자. $R \to S$를 유한형이라 하자.
$R \to R'$를 임의의 환 준동형사상이라 하자.
$S' = R' \otimes_R S$라 두자.
% P01-E0674: the frozen “Denote f ... the map” construction omits “by/as”.
환 준동형사상 $S \to S'$에 연관된 사상을
$f : \Spec(S') \to \Spec(S)$라 쓰자.
$\mathfrak p = R \cap \mathfrak q$일 때 섬유환
$S_{\mathfrak q} \otimes_R \kappa(\mathfrak p)$가 Cohen--Macaulay인
소 아이디얼 $\mathfrak q$들의 집합을 $W$라 하자.
$S'/R'$에 대한 유사한 집합을 $W'$라 하자. 그러면
$W' = f^{-1}(W)$이다.
\end{lemma}

\begin{proof}
% P01-E0675: the frozen one-sentence proof omits its subject.
이는 보조정리 \ref{algebra-lemma-extend-field-CM-locus}와 정의에서 자명하다.
\end{proof}

\begin{lemma}
\label{algebra-lemma-relative-dimension-CM}
$R$을 환이라 하자. $R \to S$가 (a) 평탄하고,
(b) 유한 표시이며, (c) Cohen--Macaulay 섬유를 갖는 환 준동형사상이라 하자.
% P01-E0676: the frozen predicate list is grammatically nonparallel.
그러면 $S = S_0 \times \ldots \times S_n$을 $R$-대수 $S_d$들의
곱으로 쓸 수 있고, 각 $S_d$는 (a), (b), (c)를 만족하며
모든 섬유가 차원 $d$로 등차원이다.
\end{lemma}

\begin{proof}
% P01-E0677: the frozen designation construction omits “by/as”.
각 정수 $d$에 대하여, 보조정리
\ref{algebra-lemma-finite-presentation-flat-CM-locus-open}에서 정의한 집합을
$W_d \subset \Spec(S)$라 쓰자. 명백히
$\Spec(S) = \coprod W_d$이고, 방금 인용한 보조정리에 의해
각 $W_d$는 열린다. 따라서 결론은 보조정리
\ref{algebra-lemma-disjoint-implies-product}에서 따른다.
\end{proof}

\section{미분}
\label{algebra-section-differentials}

\noindent
이 절에서는 환 준동형사상의 미분 가군을 정의한다.

\begin{definition}
\label{algebra-definition-derivation}
$\varphi : R \to S$를 환 준동형사상이라 하고, $M$을 $S$-가군이라 하자.
$M$으로 가는 {\it 미분}, 더 정확히 말해 {\it $R$-미분}이란
가법이고 $\varphi(R)$의 원소들을 소멸시키며
{\it 라이프니츠 법칙} $D(ab) = aD(b) + bD(a)$를 만족하는
사상 $D : S \to M$을 말한다.
\end{definition}

\noindent
$r \in R$이고 $a \in S$이면 $D(ra) = rD(a)$임에 유의하라.
동치인 정의로, $R$-미분은 라이프니츠 법칙을 만족하는
$R$-선형 사상 $D : S \to M$이다.
모든 $R$-미분의 집합은 $S$-가군을 이룬다.
두 $R$-미분 $D, D'$가 주어지면 그 합
$D + D' : S \to M$, $a \mapsto D(a)+D'(a)$는 $R$-미분이고,
$R$-미분 $D$와 원소 $c\in S$가 주어지면 스칼라배
$cD : S \to M$, $a \mapsto cD(a)$도 $R$-미분이다.
이 $S$-가군을
$$
\text{Der}_R(S, M).
$$
로 나타낸다. 또한 $\alpha : M \to N$이 $S$-가군 준동형사상이면
합성 $\alpha \circ D$는 $N$으로 가는 $R$-미분이다.
따라서 대응 $M \mapsto \text{Der}_R(S, M)$은 공변 함자이다.

\medskip\noindent
다음의 자유 $S$-가군 준동형사상을 생각하자.
$$
\bigoplus\nolimits_{(a, b)\in S^2} S[(a, b)] \oplus
\bigoplus\nolimits_{(f, g)\in S^2} S[(f, g)] \oplus
\bigoplus\nolimits_{r\in R} S[r]
\longrightarrow
\bigoplus\nolimits_{a\in S} S[a]
$$
이는 다음 규칙들로 정의된다.
$$
[(a, b)] \longmapsto [a + b] - [a] - [b],\quad
[(f, g)] \longmapsto [fg] -f[g] - g[f],\quad
[r]      \longmapsto [\varphi(r)]
$$
여기서 기호의 뜻은 명백하다. $\Omega_{S/R}$를 이 사상의 여핵이라 하자.
$a$를 여핵 안의 $[a]$의 동치류 $\text{d}a$로 보내는
사상 $\text{d} : S \to \Omega_{S/R}$가 있다.
$\Omega_{S/R}$에 부과된 관계에 의해 이는 $R$-미분이다.
달리 말해
$$
\text{d}(a + b) = \text{d}a + \text{d}b, \quad
\text{d}(fg) = f\text{d}g + g\text{d}f, \quad
\text{d}\varphi(r) = 0
$$
% P01-E0678: the frozen display ends with a period before the following where-clause.
이며, 여기서 $a,b,f,g \in S$이고 $r \in R$이다.

\begin{definition}
\label{algebra-definition-differentials}
쌍 $(\Omega_{S/R}, \text{d})$를 $R$ 위의 $S$의
{\it K\"ahler 미분 가군} 또는 {\it 미분 가군}이라 한다.
\end{definition}

\begin{lemma}
\label{algebra-lemma-universal-omega}
\begin{slogan}
미분 가군에서 나가는 사상은 미분과 같다.
\end{slogan}
$R$ 위의 $S$의 미분 가군은 다음 보편 성질을 갖는다. 사상
$$
\Hom_S(\Omega_{S/R}, M)
\longrightarrow
\text{Der}_R(S, M), \quad
\alpha
\longmapsto
\alpha \circ \text{d}
$$
은 함자들의 동형사상이다.
\end{lemma}

\begin{proof}
정의상 $R$-미분은 각 $a \in S$에 원소 $D(a) \in M$을
대응시키는 규칙이다. 따라서 $D$는 사상
$[D] : \bigoplus S[a] \to M$을 유도한다.
한편 $R$-미분이라는 조건들은 정확히 $[D]$가 위에서 주어진
$\Omega_{S/R}$의 표시 사상의 상을 소멸시킨다는 뜻이다.
\end{proof}

\begin{lemma}
\label{algebra-lemma-trivial-differential-surjective}
$R \to S$가 전사라고 하자. 그러면 $\Omega_{S/R} = 0$이다.
\end{lemma}

\begin{proof}
모든 $R$-미분이 명백히 영이어야 한다는 것으로도 알 수 있고,
$\Omega_{S/R}$의 표시에서 나오는 사상이 전사라는 것으로도 알 수 있다.
\end{proof}

\noindent
환들의 가환 그림
\begin{equation}
\label{algebra-equation-functorial-omega}
\vcenter{
\xymatrix{
S \ar[r]_\varphi
&
S'
\\
R \ar[r]^\psi \ar[u]^\alpha
&
R' \ar[u]_\beta
}
}
\end{equation}
이 주어졌다고 하자. 이 경우 미분 가군들의 자연스러운 사상이 존재하며,
이는 가환 그림
$$
\xymatrix{
\Omega_{S/R} \ar[r] &
\Omega_{S'/R'}
\\
S \ar[u]^{\text{d}} \ar[r]^{\varphi}
&
S' \ar[u]_{\text{d}}
}
$$
에 들어맞는다. 이 사상을 구성하려면
$\Omega_{S/R}$와 $\Omega_{S'/R'}$의 표시 사이의 명백한 사상을 사용하면 된다.
구체적으로,
\begin{equation}
\label{algebra-equation-map-presentations}
\vcenter{
\xymatrix{
\bigoplus S'[(a', b')]
\oplus
\bigoplus S'[(f', g')]
\oplus
\bigoplus S'[r'] \ar[r]
&
\bigoplus S' [a'] \\
 \\
\bigoplus S[(a, b)]
\oplus
\bigoplus S[(f, g)]
\oplus
\bigoplus S[r] \ar[r]
\ar[uu]^{
\begin{matrix}
[(a, b)] \mapsto [(\varphi(a), \varphi(b))] \\
[(f, g)] \mapsto [(\varphi(f), \varphi(g))] \\
[r]\mapsto [\psi(r)]
\end{matrix}
} &
\bigoplus S[a] \ar[uu]_{[a] \mapsto [\varphi(a)]}
}
}
\end{equation}
이다. 결과적으로 $f\text{d}g \in \Omega_{S/R}$는
$\varphi(f)\text{d}\varphi(g)$로 보내진다.


\begin{lemma}
\label{algebra-lemma-colimit-differentials}
$I$를 유향 집합이라 하자. 범주론의 절
\ref{categories-section-posets-limits}을 참고하여,
$(R_i \to S_i, \varphi_{ii'})$를 $I$ 위의 환 준동형사상들의 계라 하자.
그러면
$$
\Omega_{S/R} =
\colim_i \Omega_{S_i/R_i}.
$$
이고, 여기서 $R \to S = \colim (R_i \to S_i)$이다.
\end{lemma}

\begin{proof}
이는 $\Omega_{S/R}$의 정의 표시와 위에서 설명한 이 구성의 함자성으로부터
명백하다.
% P01-E0679: the frozen phrase “the functoriality of this described above” lacks a noun.
\end{proof}

\begin{lemma}
\label{algebra-lemma-differential-surjective}
그림 (\ref{algebra-equation-functorial-omega})에서
$S \to S'$가 전사이고 그 핵이 $I \subset S$라고 하자.
그러면 $\Omega_{S/R} \to \Omega_{S'/R'}$는 전사이고,
그 핵은 $\varphi(a) \in \beta(R')$인 $a \in S$에 대한 원소
$\text{d}a$들에 의해 $S$-가군으로 생성된다.
(특히 원소 $\text{d}(i)$, $i \in I$도 여기에 포함된다.)
\end{lemma}

\begin{proof}[첫째 증명]
표시들의 사상 (\ref{algebra-equation-map-presentations})을 생각하자.
자유 가군들의 오른쪽 수직 사상은 명백히 전사이다.
따라서 이 사상은 전사이다.
$\Omega_{S/R}$의 어떤 원소 $\eta$가 $\Omega_{S'/R'}$에서
영으로 보내진다고 하자. 어떤 $s_i, a_i \in S$에 대하여
$\eta$를 $\sum s_i[a_i]$의 상으로 쓰자.
그러면 $\sum \varphi(s_i)[\varphi(a_i)]$는 그림의 왼쪽 위 꼭짓점에 있는 원소
$$
\theta = \sum s_j'[a_j', b_j'] + \sum s_k'[f_k', g_k'] + \sum s_l'[r_l']
$$
의 상임을 알 수 있다. $\varphi$가 전사이므로 항
$s_j'[a_j', b_j']$와 $s_k'[f_k', g_k']$는 오른쪽 아래 꼭짓점에
있는 원소들의 상이다.
% P01-E0680: the frozen proof names the lower right corner, but these generators lift from the lower left.
따라서 이 원소들의 상을 빼서 $\eta$와 $\theta$를 바꾸면
$\theta = \sum s_l'[r_l']$이라고 가정할 수 있다.
달리 말해 $\sum \varphi(s_i)[\varphi(a_i)]$는
$\sum s'_l [\beta(r'_l)]$ 꼴이다.
다음으로 모든 $i$에 대하여 $a' = \varphi(a_i)$이고
모든 $l$에 대하여 $a' = \beta(r_l')$인 어떤 $a' \in S'$가
있다고 가정할 수 있다.
% P01-E0681: the frozen proof suppresses the componentwise reduction behind this common a'.
이는 그림의 오른쪽 위 꼭짓점의 직합 분해로부터 명백하다.
$\varphi(a) = a'$인 $a \in S$를 택하자.
그러면 어떤 $x_i \in I$에 대하여 $a_i = a + x_i$로 쓸 수 있다.
따라서 허용된 관계를 사용하면 모든 $a_i$가 $a$와 같다고
가정할 수 있다. 그러면 우리 원소가 $s[a]$ 꼴이라고 가정할 수 있다.
여전히
$\varphi(s)[a'] = \sum \varphi(s_l')[\beta(r_l')]$
임을 알고 있다.
% P01-E0682: the frozen formula applies phi to coefficients already lying in S'.
따라서 $\varphi(s) = 0$이면 결론을 얻고, 그렇지 않으면
$a' = \varphi(a)$가 $\beta$의 상에 속하므로 역시 결론을 얻는다.
\end{proof}

\begin{proof}[둘째 증명]
더 언급하지 않고 보조정리 \ref{algebra-lemma-universal-omega}에 주어진
미분 가군의 보편 성질을 사용하겠다.

\medskip\noindent
(\ref{algebra-equation-functorial-omega})에서
$R'' = S \times_{S'} R'$라 두자. 그러면 다음 그림을 얻는다.
% P01-E0683: the frozen phrase “we have following diagram” omits the article.
$$
\xymatrix{
S \ar[r] & S \ar[r] & S' \\
R \ar[r] \ar[u] & R'' \ar[r] \ar[u] & R' \ar[u]
}
$$
$M$을 $S$-가군이라 하자. 정의로부터 곧바로,
$R$-미분 $D : S \to M$이 $R''$-미분일 필요충분조건은
$R'' \to S$의 상에 속하는 원소들을 소멸시키는 것임을 알 수 있다.
보편 성질에 의해 이는 자연스러운 사상
$\Omega_{S/R} \to \Omega_{S/R''}$가 전사이고 그 핵이
$R''$의 상에 의해 $S$-가군으로 생성된다는 명제로 옮겨진다.
% P01-E0684: the frozen kernel description names elements of R'' rather than their differentials.

\medskip\noindent
앞 문단에 의해, $S \to S'$가 전사이고
$R = S \times_{S'} R'$일 때
$\Omega_{S/R} \to \Omega_{S'/R'}$가 동형사상임을 보이면 충분하다.
$M'$를 $S'$-가군이라 하자. 임의의 $R'$-미분
$D' : S' \to M'$에 $S \to S'$를 앞에 합성하면
$R$-미분을 얻음에 유의하라. 역으로, $M$을 $S$-가군이라 하고
$D : S \to M$을 $R$-미분이라 하자. $i \in I$이면
($R = S \times_{S'} R'$이므로) $\alpha(a) = i$인
어떤 $a \in R$이 존재한다. 따라서 $D(i) = 0$이고,
모든 $s \in S$에 대하여 $0 = D(is) = iD(s)$이다.
그러므로 $D$의 상은 $I$에 의해 소멸되는 원소들로 이루어진
부분가군 $M' \subset M$에 포함되고, 더욱이 유도된 사상
$S \to M'$는 $R'$-미분 $S' \to M'$를 거쳐 인수분해된다.
보편 성질을 사용하여 이것이
$\Omega_{S/R} \to \Omega_{S'/R'}$가 동형사상이라는 뜻임을
확인하는 일은 연습문제로 남긴다. 세부사항은 생략한다.
\end{proof}

\begin{lemma}
\label{algebra-lemma-exact-sequence-differentials}
$A \to B \to C$를 환 준동형사상들이라 하자.
그러면 $C$-가군들의 표준 완전열
$$
C \otimes_B \Omega_{B/A} \to
\Omega_{C/A} \to
\Omega_{C/B} \to 0
$$
이 존재한다.
\end{lemma}

\begin{proof}
$R = A$, $S = C$, $R' = B$, $S' = C$로 두면
그림 (\ref{algebra-equation-functorial-omega})을 얻는다.
보조정리 \ref{algebra-lemma-differential-surjective}에 의해 사상
$\Omega_{C/A} \to \Omega_{C/B}$는 전사이고, 그 핵은
$B \to C$의 상에 속하는 $c \in C$에 대한 원소
$\text{d}(c)$들로 생성된다. 따라서 보조정리가 따른다.
\end{proof}

\begin{lemma}
\label{algebra-lemma-differentials-localize}
$\varphi : A \to B$를 환 준동형사상이라 하자.
\begin{enumerate}
\item $S \subset A$가 $B$의 가역원들로 보내지는 곱셈적 부분집합이면
$\Omega_{B/A} = \Omega_{B/S^{-1}A}$이다.
\item $S \subset B$가 곱셈적 부분집합이면
$S^{-1}\Omega_{B/A} = \Omega_{S^{-1}B/A}$이다.
\end{enumerate}
\end{lemma}

\begin{proof}
(1)의 등식을 보이려면 임의의 $A$-미분
$D : B \to M$이 원소 $\varphi(s)^{-1}$들을 소멸시킴을 보이면 충분하다.
이는 라이프니츠 법칙을
$1 = \varphi(s) \varphi(s)^{-1}$에 적용하면 명백하다.
(2)를 보이기 위해 명백한 사상
$S^{-1}\Omega_{B/A} \to \Omega_{S^{-1}B/A}$가 있음에 유의하라.
% P01-E0697: the frozen introductory phrase lacks a comma after “(2)”.
이것이 동형사상임을 보이려면 $S^{-1}B$에서
$S^{-1}\Omega_{B/A}$로 가는 $A$-미분 $\text{d}'$가
존재함을 보이면 충분하다.
% P01-E0685: the frozen source says “a A-derivation”.
이를 정의하기 위해 단순히
$\text{d}'(b/s) = (1/s)\text{d}b - (1/s^2)b\text{d}s$로 둔다.
세부사항은 생략한다.
\end{proof}

\begin{lemma}
\label{algebra-lemma-differential-seq}
그림 (\ref{algebra-equation-functorial-omega})에서
$S \to S'$가 전사이고 그 핵이 $I \subset S$라고 하며,
$R' = R$이라고 가정하자.
그러면 $S'$-가군들의 표준 완전열
$$
I/I^2
\longrightarrow
\Omega_{S/R} \otimes_S S'
\longrightarrow
\Omega_{S'/R}
\longrightarrow
0
$$
% P01-E0692: the frozen source omits terminal punctuation after this display.
이 존재한다. 가장 왼쪽 사상은 $f \in I$가
$\text{d}f \otimes 1$로 보내진다는 규칙으로 특징지어진다.
\end{lemma}

\begin{proof}
가운데 항은 $\Omega_{S/R} \otimes_S S/I$이다.
$f \in I$에 대하여 $I/I^2$에서의 $f$의 상을
$\overline{f}$라 쓰자.
% P01-E0686: the frozen designation construction omits “by/as”.
사상 $\overline{f} \mapsto \text{d}f \otimes 1$이 잘 정의됨을
보이려면, $f_1, f_2 \in I$일 때
$\text{d} f_1f_2 \otimes 1 = 0$임을 확인하면 된다.
이는 라이프니츠 법칙
$\text{d} f_1f_2 \otimes 1 = (f_1 \text{d}f_2 + f_2 \text{d} f_1 )\otimes 1 =
\text{d}f_2 \otimes f_1 + \text{d}f_1 \otimes f_2 = 0$
으로부터 명백하다. 비슷한 계산으로 이 사상이
$S' = S/I$-선형임을 알 수 있다.
% P01-E0687: the frozen source has subject-verb disagreement in “A similar computation show”.

\medskip\noindent
사상 $\Omega_{S/R} \otimes_S S' \to \Omega_{S'/R}$는
$S$-선형 사상 $\Omega_{S/R} \to \Omega_{S'/R}$에 연관된
표준 $S'$-선형 사상이다.
보조정리 \ref{algebra-lemma-differential-surjective}에 의해
$\Omega_{S/R} \to \Omega_{S'/R}$가 전사이므로 이 사상도 전사이다.

\medskip\noindent
$f \in I$에 대하여 $\text{d}f$가 $\Omega_{S'/R}$에서 영으로
보내지므로 두 사상의 합성은 영이다.
완전성은 $\Omega_{S/R} \to \Omega_{S'/R}$의 핵이
부분가군 $I\Omega_{S/R}$와 원소 $\text{d}f$, $f \in I$에 의해
$S$-부분가군으로 생성된다는 뜻일 뿐임에 유의하라.
보조정리 \ref{algebra-lemma-differential-surjective}에 의해 이 핵은
어떤 $r \in R$에 대하여 $\varphi(a) = \beta(r)$인 원소
$\text{d}(a)$들로 생성됨을 알고 있다.
그런데 $a = \alpha(r) + a - \alpha(r)$이므로
$\text{d}(a) = \text{d}(a - \alpha(r))$이다. 또한
$\varphi(a - \alpha(r)) =
\varphi(a) - \varphi(\alpha(r)) = \beta(r) - \beta(r) = 0$이므로
$a - \alpha(r) \in I$이다.
따라서 $f \in I$인 원소 $\text{d}f$들이 이미
원하는 대로 그 핵을 $S$-가군으로 생성한다는 결론을 얻는다.
% P01-E0696: the frozen sentence omits “that” after “We conclude”.
\end{proof}

\begin{lemma}
\label{algebra-lemma-differential-seq-split}
그림 (\ref{algebra-equation-functorial-omega})에서
$S \to S'$가 전사이고 그 핵이 $I \subset S$라고 하며
$R' = R$이라고 가정하자. 또한 $S \to S'$의 오른쪽 역인
$R$-대수 준동형사상 $S' \to S$가 존재한다고 가정하자.
그러면 보조정리 \ref{algebra-lemma-differential-seq}의 $S'$-가군 완전열은
짧은 완전열
$$
0 \longrightarrow
I/I^2
\longrightarrow
\Omega_{S/R} \otimes_S S'
\longrightarrow
\Omega_{S'/R}
\longrightarrow
0
$$
이 되며, 나아가 분할 짧은 완전열이다.
\end{lemma}

\begin{proof}
$\beta : S' \to S$를 전사 $\alpha : S \to S'$의 오른쪽 역이라 하자.
사상
$$
D : S \longrightarrow I/I^2,
\quad
x \longmapsto x - \beta(\alpha(x))
$$
을 생각하자. $D$가 $R$-미분임은 쉽게 보인다(생략한다).
또한 $x \in I, s \in S$이면 $x D(s) = 0$이다.
따라서 보편 성질에 의해 $D$는 사상
$\tau : \Omega_{S/R} \otimes_S S' \to I/I^2$를 유도한다.
이것이 $\text{d} : I/I^2  \to \Omega_{S/R} \otimes_S S'$의
왼쪽 역임을 확인하는 일은 생략한다. 따라서 결론을 얻는다.
\end{proof}

\begin{lemma}
\label{algebra-lemma-differential-mod-power-ideal}
$R \to S$를 환 준동형사상이라 하자. $I \subset S$를 아이디얼이라 하자.
$n \geq 1$을 정수라 하자. $S' = S/I^{n + 1}$이라 두자.
사상 $\Omega_{S/R} \to \Omega_{S'/R}$는 동형사상
$$
\Omega_{S/R} \otimes_S S/I^n
\longrightarrow
\Omega_{S'/R} \otimes_{S'} S/I^n.
$$
% P01-E0693: the frozen source omits terminal punctuation after this displayed map.
을 유도한다.
\end{lemma}

\begin{proof}
이는 보조정리 \ref{algebra-lemma-differential-seq}와
$\text{d}$에 대한 라이프니츠 법칙으로부터 얻는 사실
$\text{d}(I^{n + 1}) \subset I^n\Omega_{S/R}$에서 나온다.
\end{proof}

\begin{lemma}
\label{algebra-lemma-differentials-base-change}
환 준동형사상 $R \to R'$ 및 $R \to S$가 있다고 하자.
$S' = S \otimes_R R'$라 두면 그림
(\ref{algebra-equation-functorial-omega})을 얻는다.
그러면 위에서 정의한 표준 사상은 동형사상
$\Omega_{S/R} \otimes_R R' = \Omega_{S'/R'}$을 유도한다.
\end{lemma}

\begin{proof}
$\text{d}' : S' = S \otimes_R R' \to \Omega_{S/R} \otimes_R R'$를
사상
$\text{d}'( \sum a_i \otimes x_i ) = \sum \text{d}(a_i) \otimes x_i$로
정의하자. 이는 사상
$S \times R' \to \Omega_{S/R} \otimes_R R'$,
$(a, x)\mapsto \text{d}a \otimes_R x$가 $R$-쌍선형이므로 존재한다.
간단한 계산으로 이것이 $R'$-미분임을 확인할 수 있다.
쌍 $(\Omega_{S/R} \otimes_R R', \text{d}')$가 보편 성질을
만족함을 보이겠다.
$D : S' \to M'$를 $S'$-가군으로 가는 $R'$-미분이라 하자.
합성 $S \to S' \to M'$은 $R$-미분이므로
$S$-선형 사상 $\varphi_D : \Omega_{S/R} \to M'$를 얻는다.
이를 $R'$와 텐서하여 사상
$\varphi'_D : \Omega_{S/R} \otimes_R R' \to M'$,
$\varphi'_D(\eta \otimes x) = x\varphi_D(\eta)$를 얻는다.
$D = \varphi'_D \circ \text{d}'$임은 명백하다.
\end{proof}

\noindent
곱셈 사상 $S \otimes_R S \to S$는 $a \otimes b$를 $S$의 $ab$로
보내는 $R$-대수 준동형사상이다.
$S \otimes_R S$를 두 사상 $S \to S \otimes_R S$ 가운데 어느 하나를 통해
$S$-대수로 생각하면, 이는 또한 $S$-대수 준동형사상이다.

\begin{lemma}
\label{algebra-lemma-differentials-diagonal}
$R \to S$를 환 준동형사상이라 하자.
$J = \Ker(S \otimes_R S \to S)$를 곱셈 사상의 핵이라 하자.
$S$-가군들의 표준 동형사상 $\Omega_{S/R} \to J/J^2$가 존재하며,
이는 $a \text{d} b \mapsto a \otimes b - ab \otimes 1$로 주어진다.
\end{lemma}

\begin{proof}[첫째 증명]
보조정리 \ref{algebra-lemma-differential-seq-split}을 가환 그림
$$
\xymatrix{
S \otimes_R S \ar[r] & S \\
S \ar[r] \ar[u] & S \ar[u]
}
$$
에 적용하자. 여기서 왼쪽 수직 화살표는 $a \mapsto a \otimes 1$이다.
완전열
$0 \to J/J^2 \to
\Omega_{S \otimes_R S/S} \otimes_{S \otimes_R S} S \to \Omega_{S/S} \to 0$
을 얻는다. 보조정리 \ref{algebra-lemma-trivial-differential-surjective}에 의해
항 $\Omega_{S/S}$는 $0$이고, 따라서 나머지 두 항 사이의
동형사상을 얻는다. $S \to S \otimes_R S$는
$R \to S$의 기저변환이므로, 보조정리
\ref{algebra-lemma-differentials-base-change}에 의해
$\Omega_{S \otimes_R S/S} = \Omega_{S/R} \otimes_S (S \otimes_R S)$이다.
따라서
$$
\Omega_{S \otimes_R S/S} \otimes_{S \otimes_R S} S =
\Omega_{S/R} \otimes_S (S \otimes_R S) \otimes_{S \otimes_R S} S =
\Omega_{S/R}
$$
% P01-E0694: the frozen source omits terminal punctuation after this displayed chain.
이다. 사상이 보조정리의 규칙으로 주어짐을 확인하는 일은 생략한다.
\end{proof}

\begin{proof}[둘째 증명]
먼저 규칙 $a \text{d} b \mapsto a \otimes b - ab \otimes 1$이
잘 정의됨을 보이자. 이를 위해 $\text{d}r$과
$a\text{d}b + b \text{d}a - d(ab)$가 영으로 보내짐을 보여야 한다.
% P01-E0688: the frozen relation uses d(ab) rather than the established \text{d}(ab).
첫째 것은 텐서곱의 정의에 의해
$r \otimes 1 - 1 \otimes r = 0$이기 때문이다.
둘째 것은
% P01-E0689: the frozen “The first because...” clause is a sentence fragment.
$$
(a \otimes b - ab \otimes 1) +
(b \otimes a - ba \otimes 1) -
(1 \otimes ab - ab \otimes 1)
=
(a \otimes 1 - 1\otimes a)(1\otimes b - b \otimes 1)
$$
이 $J^2$에 속하기 때문이다.

\medskip\noindent
반대 방향의 사상을 구성하자.
$S \to S \otimes_R S$, $a \mapsto a \otimes 1$을
$R \to S$의 기저변환으로 생각할 수 있다.
따라서 보조정리 \ref{algebra-lemma-differentials-base-change}에 의해
$\Omega_{S \otimes_R S/S} = \Omega_{S/R} \otimes_S (S \otimes_R S)$이다.
이제 보조정리 \ref{algebra-lemma-differential-seq}의 열은 사상
$$
J/J^2  \to \Omega_{S \otimes_R S/ S} \otimes_{S \otimes_R S} S
= (\Omega_{S/R} \otimes_S (S \otimes_R S))\otimes_{S \otimes_R S} S
= \Omega_{S/R}.
$$
을 준다. 이것이 위 사상의 역임을 확인하는 일은 독자에게 맡긴다.
% P01-E0690: the frozen phrase “see it is the inverse” omits “that”.
\end{proof}

\begin{lemma}
\label{algebra-lemma-differentials-polynomial-ring}
$S = R[x_1, \ldots, x_n]$이면 $\Omega_{S/R}$는
$\text{d}x_1, \ldots, \text{d}x_n$을 기저로 하는
유한 자유 $S$-가군이다.
\end{lemma}

\begin{proof}
먼저 $\text{d}x_1, \ldots, \text{d}x_n$이
$\Omega_{S/R}$를 $S$-가군으로 생성함을 보이자.
이를 위해 임의의 $g \in R[x_1, \ldots, x_n]$에 대하여
$\text{d}g$가 합 $\sum g_i \text{d}x_i$로 표현됨을 보인다.
$g$의 (전체) 차수에 대한 귀납법을 사용하자.
$g$의 차수가 $0$이면 $\text{d}g = 0$이므로 명백하다.
$g$의 차수가 $> 0$이면, $c\in R$이고
$\deg(g_i) < \deg(g)$인 $c + \sum g_i x_i$로
$g$를 쓸 수 있다. 라이프니츠 법칙에 의해
$\text{d}g = \sum g_i \text{d} x_i + \sum x_i \text{d}g_i$이므로
귀납법으로 결론을 얻는다.

\medskip\noindent
$R$-미분 $\partial / \partial x_i :
R[x_1, \ldots, x_n] \to R[x_1, \ldots, x_n]$을 생각하자.
(그 정의는 독자에게 맡긴다. 정의 성질은
$\partial / \partial x_i (x_j) = \delta_{ij}$이다.)
보편 성질에 의해 이는 $\text{d}x_i$를 $1$로,
$j \not = i$인 경우 $\text{d}x_j$를 $0$으로 보내는
$S$-가군 준동형사상
$l_i : \Omega_{S/R} \to R[x_1, \ldots, x_n]$에 대응한다.
따라서 원소 $\text{d}x_1, \ldots, \text{d}x_n$ 사이에는
$S$-선형 관계가 없다는 것이 명백하다.
\end{proof}

\begin{lemma}
\label{algebra-lemma-differentials-finitely-presented}
$R \to S$가 유한 표시라고 하자.
그러면 $\Omega_{S/R}$는 유한 표시 $S$-가군이다.
\end{lemma}

\begin{proof}
$S = R[x_1, \ldots, x_n]/(f_1, \ldots, f_m)$로 쓰자.
$I = (f_1, \ldots, f_m)$라 쓰자.
보조정리 \ref{algebra-lemma-differential-seq}에 의해
$S$-가군들의 완전열
$$
I/I^2
\to
\Omega_{R[x_1, \ldots, x_n]/R} \otimes_{R[x_1, \ldots, x_n]} S
\to
\Omega_{S/R}
\to
0
$$
% P01-E0695: the frozen source omits terminal punctuation after this displayed sequence.
이 존재한다. $I/I^2$가 $f_i$들의 상으로 생성되는 유한
$S$-가군이고, 가운데 항이 보조정리
\ref{algebra-lemma-differentials-polynomial-ring}에 의해 유한 자유이므로
결론이 따른다.
\end{proof}

\begin{lemma}
\label{algebra-lemma-differentials-finitely-generated}
$R \to S$가 유한형이라고 하자. 그러면
$\Omega_{S/R}$는 유한 생성 $S$-가군이다.
% P01-E0691: the frozen source omits the article in “is finitely generated S-module”.
\end{lemma}

\begin{proof}
이는 보조정리 \ref{algebra-lemma-differentials-finitely-presented}의 증명과
매우 비슷하지만 더 쉽다.
\end{proof}

\section{드람 복합체}
\label{algebra-section-de-rham-complex}

\noindent
$A \to B$를 환 준동형사상이라 하자.
% P01-E0698: the frozen designation grammatically identifies d as the module rather than its universal derivation.
절 \ref{algebra-section-differentials}에서 구성한 보편 $A$-미분을
미분 가군과 함께 $\text{d} : B \to \Omega_{B/A}$로 쓰자.
$i \geq 0$에 대하여
$\Omega_{B/A}^i = \wedge^i_B(\Omega_{B/A})$를
절 \ref{algebra-section-tensor-algebra}에서와 같은 $i$번째 외승이라 하자.
{\it $A$ 위의 $B$의 드람 복합체}란 아래에서 구성하고 설명할
$A$-선형 미분을 갖춘 복합체
$$
\Omega_{B/A}^0 \to \Omega_{B/A}^1 \to \Omega_{B/A}^2 \to \ldots
$$
를 말한다.

\medskip\noindent
사상 $\text{d} : \Omega^0_{B/A} \to \Omega^1_{B/A}$는
보편 미분 $\text{d} : B \to \Omega_{B/A}$이다.
이는 실제로 $A$-선형임에 유의하라.

\medskip\noindent
$p \geq 1$에 대하여 유일한 $A$-선형 사상
$\text{d} : \Omega_{B/A}^p \to \Omega_{B/A}^{p + 1}$이 존재하여
\begin{equation}
\label{algebra-equation-rule}
\text{d}\left(b_0\text{d}b_1 \wedge \ldots \wedge \text{d}b_p\right) =
\text{d}b_0 \wedge \text{d}b_1 \wedge \ldots \wedge \text{d}b_p
\end{equation}
% P01-E0699: the frozen numbered equation lacks terminal punctuation.
을 만족한다고 주장한다.
$\Omega_{B/A}$는 원소 $\text{d}b$들에 의해 $B$-가군으로
생성됨을 상기하라. 따라서 $\Omega^p_{B/A}$는 원소
$b_0\text{d}b_1 \wedge \ldots \wedge \text{d}b_p$들에 의해
$A$-가군으로 생성된다. 그러므로 사상
$\text{d} : \Omega^p_{B/A} \to \Omega^{p + 1}_{B/A}$가
존재한다면 유일하다.

\medskip\noindent
$\text{d} : \Omega_{B/A}^1 \to \Omega_{B/A}^2$의 구성.
정의 \ref{algebra-definition-differentials}에 의해 원소 $\text{d}b$들은
$a \in A$에 대한 관계 $\text{d}a = 0$ 및
$\text{d}(b' + b'') = \text{d}b' + \text{d}b''$와
$b', b'' \in B$에 대한 관계
$\text{d}(b'b'') = b'\text{d}b'' + b''\text{d}b'$를 제외하면
$\Omega_{B/A}$를 $B$-가군으로 자유롭게 생성한다.
따라서 규칙
$$
\sum b'_i \text{d}b_i \longmapsto \sum \text{d}b'_i \wedge \text{d}b_i
$$
이 잘 정의됨을 보이려면, $a \in A$ 및 $b, b', b'' \in B$에 대한 원소
$$
b\text{d}a,
\quad\text{and}\quad
b\text{d}(b' + b'') - b\text{d}b' - b\text{d}b''
\quad\text{and}\quad
b\text{d}(b'b'') - bb'\text{d}b'' - bb''\text{d}b'
$$
가 영으로 보내짐을 보이면 된다.
이는 $\text{d}$에 대한 라이프니츠 법칙을 사용해 직접 계산하면 명백하다.

\medskip\noindent
합성 $\Omega^0_{B/A} \to \Omega^1_{B/A} \to \Omega^2_{B/A}$는 영이다.
실제로
$\text{d}(\text{d}(b)) = \text{d}(1 \text{d}b) =
\text{d}(1) \wedge \text{d}(b) = 0 \wedge \text{d}b = 0$이다.
$1 \in B$는 $A \to B$의 상에 속하므로 여기서
$\text{d}(1) = 0$이다. 이를 아래에서 사용하겠다.

\medskip\noindent
$p \geq 2$일 때
$\text{d} : \Omega_{B/A}^p \to \Omega_{B/A}^{p + 1}$의 구성.
$A$-선형 사상
$$
\gamma : \Omega^1_{B/A} \otimes_A \ldots \otimes_A \Omega^1_{B/A}
\longrightarrow
\Omega_{B/A}^{p + 1}
$$
을 공식
$$
\omega_1 \otimes \ldots \otimes \omega_p
\longmapsto
\sum (-1)^{i + 1}
\omega_1 \wedge \ldots \wedge \text{d}(\omega_i) \wedge \ldots \wedge \omega_p
$$
로 정의하자.
% P01-E0700: the frozen alternating sum omits the bounds i = 1,...,p.
이 사상이 자연스러운 전사
$\Omega^1_{B/A} \otimes_A \ldots \otimes_A \Omega^1_{B/A} \to \Omega^p_{B/A}$
를 거쳐 인수분해되어 원하는 사상
$\text{d} : \Omega^p_{B/A} \to \Omega^{p + 1}_{B/A}$를 준다는 것을 보이자.
보조정리 \ref{algebra-lemma-present-wedge}에 따르면
$\Omega^1_{B/A} \otimes_A \ldots \otimes_A \Omega^1_{B/A} \to \Omega^p_{B/A}$
의 핵은 어떤 $i \not = j$에 대해 $\omega_i = \omega_j$인 원소
$\omega_1 \otimes \ldots \otimes \omega_p$들과, 어떤 $f \in B$에 대한 원소
$\omega_1 \otimes \ldots \otimes f\omega_i \otimes \ldots \otimes \omega_p -
\omega_1 \otimes \ldots \otimes f\omega_j \otimes \ldots \otimes \omega_p$
들에 의해 $A$-가군으로 생성된다.
직접 계산하면 첫째 종류의 원소는 $\gamma$에 의해 $0$으로 보내진다.
달리 말해 $\gamma$는 교대적이다. 마무리하려면
$$
\gamma(
\omega_1 \otimes \ldots \otimes f\omega_i \otimes \ldots \otimes \omega_p)
=
\gamma(
\omega_1 \otimes \ldots \otimes f\omega_j \otimes \ldots \otimes \omega_p)
$$
가 $f \in B$에 대하여 성립함을 보여야 한다.
$A$-선형성과 교대성에 의해
$p = 2$, $i = 1$, $j = 2$, $\omega_1 = b \text{d}b'$,
$\omega_2 = c \text{d} c'$이고 $b, b', c, c' \in B$인 경우만
보이면 충분하다. 따라서 다음을 보여야 한다.
\begin{align*}
& \text{d}(fb) \wedge \text{d}b' \wedge c \text{d}c'
- fb \text{d}b' \wedge \text{d}c \wedge \text{d}c' \\
& =
\text{d}b \wedge \text{d}b' \wedge fc\text{d}c'
- b \text{d}b' \wedge \text{d}(fc) \wedge \text{d}c'
\end{align*}
% P01-E0701: the frozen align display lacks punctuation before “in other words”.
달리 말해 다음을 보여야 한다.
$$
(c \text{d}(fb) + fb \text{d}c - fc \text{d}b - b \text{d}(fc))
\wedge \text{d}b' \wedge \text{d}c' = 0.
$$
이는 라이프니츠 법칙으로부터 따른다. 원소
$b_0\text{d}b_1 \otimes \text{d}b_2 \otimes \ldots \otimes \text{d}b_p$에서
$\gamma$의 값은
$\text{d}b_0 \wedge \text{d}b_1 \wedge \ldots \wedge \text{d}b_p$임에
유의하라. 따라서 이렇게 얻은 사상
$\text{d} : \Omega^p_{B/A} \to \Omega^{p + 1}_{B/A}$는
(\ref{algebra-equation-rule})을 만족한다.

\medskip\noindent
끝으로 $\Omega^p_{B/A}$는 원소
$b_0\text{d}b_1 \wedge \ldots \wedge \text{d}b_p$들에 의해
가법적으로 생성되고,
$\text{d}(b_0\text{d}b_1 \wedge \ldots \wedge \text{d}b_p) =
\text{d}b_0 \wedge \ldots \wedge \text{d}b_p$이므로,
정확히 같은 방식으로 합성
$
\Omega^p_{B/A} \to
\Omega^{p + 1}_{B/A} \to
\Omega^{p + 2}_{B/A}
$
이 $p \geq 1$에 대하여 영임을 알 수 있다.
따라서 드람 복합체는 실제로 복합체이다.

\medskip\noindent
환 $R$만 주어졌을 때에는 $\Omega_R = \Omega_{R/\mathbf{Z}}$라 둔다.
이를 때때로 $R$의 절대 미분 가군이라고 한다. 이 명칭은 타당하다.
$\Omega_R$가 라이프니츠 법칙만 가정하고 $\text{d}1$의 소멸은
가정하지 않은 미분 가군이라면, 라이프니츠 법칙으로부터
$\text{d}1 = \text{d}(1 \cdot 1) =
1 \text{d}1 + 1 \text{d}1 = 2 \text{d}1$을 얻고, 따라서
$\Omega_R$에서 $\text{d}1 = 0$이다.
이 경우 {\it $R$의 절대 드람 복합체}는 이에 대응하는 복합체
$$
\Omega_R^0 \to \Omega_R^1 \to \Omega_R^2 \to \ldots
$$
이며, 여기서 $\Omega^i_R = \Omega^i_{R/\mathbf{Z}}$ 등으로 둔다.

\medskip\noindent
환들의 가환 그림
$$
\xymatrix{
B \ar[r] & B' \\
A \ar[r] \ar[u] & A' \ar[u]
}
$$
% P01-E0703: the frozen diagram lacks terminal punctuation before the next sentence.
이 주어졌다고 하자. 드람 복합체들의 자연스러운 사상
$$
\Omega^\bullet_{B/A} \longrightarrow \Omega^\bullet_{B'/A'}
$$
% P01-E0704: the frozen displayed map lacks terminal punctuation before “Namely”.
이 존재한다. 구체적으로 차수 $0$에서는 사상 $B \to B'$이고,
차수 $1$에서는 절 \ref{algebra-section-differentials}에서 구성한 사상
$\Omega_{B/A} \to \Omega_{B'/A'}$이며, $p \geq 2$에서는 유도된 사상
$\Omega^p_{B/A} = \wedge^p_B(\Omega_{B/A}) \to
\wedge^p_{B'}(\Omega_{B'/A'}) = \Omega^p_{B'/A'}$이다.
미분들과의 호환성은 공식 (\ref{algebra-equation-rule})에 의한 미분의
특징지음으로부터 따른다.

\begin{lemma}
\label{algebra-lemma-de-rham-complex-base-change}
환 준동형사상 $A \to A'$ 및 $A \to B$가 있다고 하자.
$B' = B \otimes_A A'$라 두면 위와 같은 그림을 얻는다.
그러면 위에서 정의한 표준 사상은 복합체들의 동형사상
$\Omega^\bullet_{B/A} \otimes_A A' = \Omega^\bullet_{B'/A'}$을 유도한다.
\end{lemma}

\begin{proof}
이는 보조정리 \ref{algebra-lemma-differentials-base-change}와
외승을 취하는 것이 기저변환과 가환한다는 사실에서 따른다.
\end{proof}

\begin{lemma}
\label{algebra-lemma-de-rham-complex}
$A \to B$를 환 준동형사상이라 하자.
$\pi : \Omega_{B/A} \to \Omega$를 전사 $B$-가군 준동형사상이라 하자.
% P01-E0705: the frozen designation omits “by/as”.
$\pi$와 보편 미분 $\text{d}_{B/A} : B \to \Omega_{B/A}$의 합성을
$\text{d} : B \to \Omega$라 쓰자. $\Omega^i = \wedge_B^i(\Omega)$라 두자.
$\pi$의 핵이 다음 조건을 만족하는 원소
$\omega \in \Omega_{B/A}$들에 의해 $B$-가군으로 생성된다고 가정하자.
즉 $\text{d}_{B/A}(\omega) \in \Omega_{B/A}^2$가
$\Omega^2$에서 영으로 보내진다고 하자.
그러면 드람 복합체
$$
\Omega^0 \to \Omega^1 \to \Omega^2 \to \ldots
$$
가 존재하고, 그 미분은 규칙
$$
\text{d} : \Omega^p \to \Omega^{p + 1},\quad
\text{d}\left(f_0\text{d}f_1 \wedge \ldots \wedge \text{d}f_p\right) =
\text{d}f_0 \wedge \text{d}f_1 \wedge \ldots \wedge \text{d}f_p
$$
% P01-E0706: the frozen defining display lacks terminal punctuation at the end of the lemma.
으로 정의된다.
\end{lemma}

\begin{proof}
가환 그림
$$
\xymatrix{
\Omega_{B/A}^0 \ar[d] \ar[r]_{\text{d}_{B/A}} &
\Omega_{B/A}^1 \ar[d]_\pi \ar[r]_{\text{d}_{B/A}} &
\Omega_{B/A}^2 \ar[d]_{\wedge^2\pi} \ar[r]_{\text{d}_{B/A}} &
\ldots \\
\Omega^0 \ar[r]^{\text{d}} &
\Omega^1 \ar[r]^{\text{d}} &
\Omega^2 \ar[r]^{\text{d}} &
\ldots
}
$$
% P01-E0707: the frozen source runs the diagram directly into a lowercase independent clause.
이 존재함을 보이겠다. 위에서 주어진 $A$ 위의 $B$의 드람 복합체의 미분
$\text{d}_{B/A} : \Omega^p_{B/A} \to \Omega^{p + 1}_{B/A}$의
구성으로부터 사상 $\text{d}$의 설명이 따른다.
가장 왼쪽 수직 화살표는 동형사상이므로 첫째 정사각형을 얻는다.
% P01-E0708: the frozen source writes the closed compound “left most”.
$\pi$가 전사이므로 둘째 정사각형을 얻기 위해서는
$\text{d}_{B/A}$가 $\pi$의 핵을 $\wedge^2\pi$의 핵으로
보냄을 보이면 충분하다. $\pi$의 핵의 임의의 원소는
$\pi(\omega_i) = 0$이고
$\wedge^2\pi(\text{d}_{B/A}(\omega_i)) = 0$인
$\sum b_i\omega_i$ 꼴이라는 것이 주어져 있다.
$\text{d}_{B/A}$에 대한 라이프니츠 법칙으로부터
$\text{d}_{B/A}(\sum b_i\omega_i) = \sum b_i \text{d}_{B/A}(\omega_i) +
\sum \text{d}_{B/A}(b_i) \wedge \omega_i$를 얻는다.
따라서 이는 $\wedge^2\pi$에 의해 영으로 보내진다.

\medskip\noindent
$i > 1$에 대하여 $\wedge^i \pi$는 전사이고, 그 핵은 사상
$\Ker(\pi) \wedge \Omega^{i - 1}_{B/A}
\to \Omega_{B/A}^i$의 상이다. $\omega_1 \in \Ker(\pi)$ 및
$\omega_2 \in \Omega^{i - 1}_{B/A}$에 대하여
$$
\text{d}_{B/A}(\omega_1 \wedge \omega_2) =
\text{d}_{B/A}(\omega_1) \wedge \omega_2 -
\omega_1 \wedge \text{d}_{B/A}(\omega_2)
$$
% P01-E0709: the frozen display lacks a comma before the following which-clause.
이며, 위에서 방금 증명한 것에 의해 이는 $\wedge^{i + 1}\pi$의
핵에 속한다. 따라서 위 그림의 $(i + 1)$번째 정사각형을 얻는다.
이로써 증명이 끝난다.
\end{proof}

\section{유한 차수 미분 연산자}
\label{algebra-section-differential-operators}

\noindent
이 절에서는 유한 차수의 미분 연산자를 도입한다.

\begin{definition}
\label{algebra-definition-differential-operators}
$R \to S$를 환 사상이라 하자. $M$, $N$을 $S$-가군이라 하자.
$k \geq 0$을 정수라 하자. {\it 차수가 $k$인 미분 연산자
$D : M \to N$}를 다음과 같이 귀납적으로 정의한다. 이는 모든
$g \in S$에 대하여 사상 $m \mapsto D(gm) - gD(m)$이 차수
$k - 1$의 미분 연산자인 $R$-선형 사상이다. 기초 경우 $k = 0$에서는
차수 $0$의 미분 연산자를 $S$-선형 사상으로 정의한다.
\end{definition}

\noindent
$D : M \to N$이 차수 $k$의 미분 연산자이면, 모든 $g \in S$에
대하여 사상 $gD$도 차수 $k$의 미분 연산자이다. 차수 $k$인 두
미분 연산자의 합도 역시 그러하다. 따라서 이들 전체의 집합
$$
\text{Diff}^k(M, N) = \text{Diff}^k_{S/R}(M, N)
$$
은 $S$-가군이다. 또한
$$
\text{Diff}^0(M, N) \subset
\text{Diff}^1(M, N) \subset
\text{Diff}^2(M, N) \subset \ldots
$$
% P01-E0710: the frozen inclusion display lacks terminal punctuation at the end of the paragraph.

\begin{lemma}
\label{algebra-lemma-composition-differential-operators}
$R \to S$를 환 사상이라 하자. $L, M, N$을 $S$-가군이라 하자.
$D : L \to M$과 $D' : M \to N$이 각각 차수 $k$와 $k'$의
미분 연산자이면, $D' \circ D$는 차수 $k + k'$의 미분 연산자이다.
\end{lemma}

\begin{proof}
$g \in S$라 하자. 그러면 $x \in L$을
$$
D'(D(gx)) - gD'(D(x)) = D'(D(gx)) - D'(gD(x)) + D'(gD(x)) - gD'(D(x))
$$
로 보내는 사상은 차수가 더 낮은 미분 연산자들의 두 합성의 합이다.
따라서 $k + k'$에 대한 귀납법으로 보조정리가 따른다.
\end{proof}

\begin{lemma}
\label{algebra-lemma-module-principal-parts}
$R \to S$를 환 사상이라 하자. $M$을 $S$-가군이라 하자.
$k \geq 0$이라 하자. $S$-가군 $P^k_{S/R}(M)$과 표준적 동형
$$
\text{Diff}^k_{S/R}(M, N) = \Hom_S(P^k_{S/R}(M), N)
$$
이 존재하며, 이는 $S$-가군 $N$에 대해 함자적이다.
\end{lemma}

\noindent
요네다 보조정리에 의해 이는 차수 $k$의 ``보편'' 미분 연산자
$D_{univ} : M \to P^k_{S/R}(M)$이 존재하여, 차수 $k$인 모든
미분 연산자 $D : M \to N$에 대하여
$D = D_{univ} \circ \alpha$를 만족하는 유일한 $S$-선형 사상
$\alpha : P^k_{S/R}(M) \to N$이 존재함을 뜻한다.
% P01-E0711: the frozen universal factorization reverses the composition order; the typed composite is alpha after D_univ.
쌍 $(P^k_{S/R}(M), D_{univ})$은 유일한 동형을 제외하고 유일하다.

\begin{proof}
$P^k_{S/R}(M)$의 존재는 일반적인 범주론적 논증
(추후 이곳에 참고문헌을 삽입)으로부터 따르지만, 우리도 하나의
구성을 제시한다.
% P01-E0712: the frozen source contains the editorial placeholder “insert future reference here”.
$F = \bigoplus_{m \in M} S[m]$으로 두자. 여기서 $[m]$은 $m$에
대응하는 직합 성분의 기저원을 나타내는 기호이다. 임의의 미분 연산자
$D : M \to N$이 주어지면 $[m]$을 $D(m)$으로 보내는 $S$-선형 사상
$L_D : F \to N$을 얻는다. $D$의 차수가 $0$이면 $L_D$는 다음 원소들을
영으로 보낸다.
$$
[m + m'] - [m] - [m'],\quad
g_0[m] - [g_0m]
$$
여기서 $g_0 \in S$이고 $m, m' \in M$이다.
$D$의 차수가 $1$이면 $L_D$는 다음 원소들을 영으로 보낸다.
$$
[m + m'] - [m] - [m'],\quad
f[m] - [fm], \quad
g_0g_1[m] - g_0[g_1m] - g_1[g_0m] + [g_1g_0m]
$$
여기서 $f \in R$, $g_0, g_1 \in S$, 그리고 $m \in M$이다.
$D$의 차수가 $k$이면 $L_D$는 원소
$[m + m'] - [m] - [m']$, $f[m] - [fm]$ 및 다음 원소들을 영으로 보낸다.
$$
g_0g_1\ldots g_k[m] - \sum g_0 \ldots \hat g_i \ldots g_k[g_im] + \ldots
+(-1)^{k + 1}[g_0\ldots g_km]
$$
% P01-E0713: the frozen general relation display lacks terminal punctuation before “Conversely”.
반대로, 앞 문장에 열거된 모든 원소를 영으로 보내는 $S$-선형 사상
$L : F \to N$이 주어지면 $m \mapsto L([m])$은 차수 $k$의 미분
연산자이다. 따라서 $P^k_{S/R}(M)$은 이 원소들이 생성하는 부분가군으로
$F$를 나눈 몫임을 알 수 있다.
\end{proof}

\begin{definition}
\label{algebra-definition-module-principal-parts}
$R \to S$를 환 사상이라 하자. $M$을 $S$-가군이라 하자. 보조정리
\ref{algebra-lemma-module-principal-parts}에서 구성한 가군 $P^k_{S/R}(M)$을
$M$의 {\it 차수 $k$의 주부 가군}이라 한다.
\end{definition}

\noindent
포함관계
$$
\text{Diff}^0(M, N) \subset
\text{Diff}^1(M, N) \subset
\text{Diff}^2(M, N) \subset \ldots
$$
는 요네다 보조정리(Categories, Lemma \ref{categories-lemma-yoneda})에 의해
전사
$$
\ldots \to P^2_{S/R}(M) \to P^1_{S/R}(M) \to P^0_{S/R}(M) = M
$$
에 대응한다.
% P01-E0714: the frozen surjection display lacks terminal punctuation at the end of the paragraph.

\begin{example}
\label{algebra-example-derivations-and-differential-operators}
$R \to S$를 환 사상이라 하고 $N$을 $S$-가군이라 하자.
$\text{Diff}^1(S, N) = \text{Der}_R(S, N) \oplus N$임을 보자.
실제로 $D : S \to N$이 차수 $1$의 미분 연산자이면,
$\sigma_D(g) := D(g) - gD(1)$로 정의되는 $\sigma_D : S \to N$은
$R$-미분이고, $D = \sigma_D + \lambda_{D(1)}$이다. 여기서
$\lambda_x : S \to N$은 $g$를 $gx$로 보내는 선형 사상이다.
따라서 $\Omega_{S/R}$의 보편 성질에 의해
$P^1_{S/R} = \Omega_{S/R} \oplus S$이다.
\end{example}

\begin{lemma}
\label{algebra-lemma-sequence-of-principal-parts}
$R \to S$를 환 사상이라 하자. $M$을 $S$-가군이라 하자. $M$에 대해
함자적인 표준적 짧은 완전열
$$
0 \to \Omega_{S/R} \otimes_S M \to P^1_{S/R}(M) \to M \to 0
$$
이 존재하며, 이를 {\it 주부열}이라 한다.
\end{lemma}

\begin{proof}
사상 $P^1_{S/R}(M) \to M$은 위에서 주어졌다.
$N$을 $S$-가군이라 하고 $D : M \to N$을 차수 $1$의 미분 연산자라 하자.
$m \in M$에 대하여 사상
$$
g \longmapsto D(gm) - gD(m)
$$
은 차수 $1$ 미분 연산자의 공리에 의해 $R$-미분 $S \to N$이다.
따라서 이는 규칙 $a\text{d}b \mapsto aD(bm) - abD(m)$으로 결정되는
선형 사상 $D_m : \Omega_{S/R} \to N$에 대응한다
(보조정리 \ref{algebra-lemma-universal-omega} 참조). 사상
$$
\Omega_{S/R} \times M \longrightarrow N,\quad
(\eta, m) \longmapsto D_m(\eta)
$$
은 $S$-쌍선형이고(세부사항은 생략한다), 따라서 $S$-선형 사상
$$
\sigma_D : \Omega_{S/R} \otimes_S M \to N
$$
을 정한다.
% P01-E0715: the frozen sigma_D display lacks terminal punctuation before the next sentence.
이와 같이 $N$에 대해 함자적인 사상
$\text{Diff}^1(M, N) \to \Hom_S(\Omega_{S/R} \otimes_S M, N)$,
$D \mapsto \sigma_D$를 얻는다. 요네다 보조정리에 의해 이는 사상
$\Omega_{S/R} \otimes_S M \to P^1_{S/R}(M)$에 대응한다.
% P01-E0716: the frozen phrase “this corresponds a map” omits “to”.
구성으로부터 이 사상이 $M$에 대해 함자적임은 즉시 알 수 있다. 열
$$
\Omega_{S/R} \otimes_S M \to P^1_{S/R}(M) \to M \to 0
$$
은 모든 가군 $N$에 대하여 열
$$
0 \to \Hom_S(M, N) \to
\text{Diff}^1(M, N) \to
\Hom_S(\Omega_{S/R} \otimes_S M, N)
$$
이 직접 확인으로 완전하므로 완전하다.

\medskip\noindent
$\Omega_{S/R} \otimes_S M \to P^1_{S/R}(M)$이 단사임을 보이기 위해
다음과 같이 논증한다. $F$가 자유 $S$-가군인 완전열
$$
0 \to M' \to F \to M \to 0
$$
을 택하자. 이는 모든 $N$에 대하여 완전열
$$
0 \to \text{Diff}^1(M, N) \to \text{Diff}^1(F, N) \to \text{Diff}^1(M', N)
$$
을 유도한다. 따라서 가환 그림
$$
\xymatrix{
0 \ar[r] &
\Omega_{S/R} \otimes_S M' \ar[r] \ar[d] &
P^1_{S/R}(M') \ar[r] \ar[d] &
M' \ar[r] \ar[d] &
0 \\
0 \ar[r] &
\Omega_{S/R} \otimes_S F \ar[r] \ar[d] &
P^1_{S/R}(F) \ar[r] \ar[d] &
F \ar[r] \ar[d] &
0 \\
0 \ar[r] &
\Omega_{S/R} \otimes_S M \ar[r] \ar[d] &
P^1_{S/R}(M) \ar[r] \ar[d] &
M \ar[r] \ar[d] &
0 \\
& 0 & 0 & 0
}
$$
에서 가운데 열은 완전하다. 왼쪽 열은
$\Omega_{S/R} \otimes_S -$의 오른쪽 완전성에 의해 완전하다.
뱀 보조정리(절 \ref{algebra-section-snake} 참조)에 의해 자유 가군 $F$에 대해
왼쪽에서 완전함을 보이면 충분하다. $P^1_{S/R}(-)$가 직합과 가환한다는
것을 이용하면 $M = S$인 경우로 환원된다. 이 경우는 예
\ref{algebra-example-derivations-and-differential-operators}의 논의에서 따른다.
\end{proof}

\begin{remark}
\label{algebra-remark-functoriality-principal-parts}
환들의 가환 그림
$$
\xymatrix{
B \ar[r] & B' \\
A \ar[u] \ar[r] & A' \ar[u]
}
$$
과 $B$-가군 $M$, $B'$-가군 $M'$, 그리고 $B$-선형 사상 $M \to M'$이
주어졌다고 하자.
% P01-E0717: the frozen multi-object “Suppose given ...” sentence lacks a comma after the displayed diagram.
그러면 가군 사상들의 양립하는 계
$$
\xymatrix{
\ldots \ar[r] &
P^2_{B'/A'}(M') \ar[r] &
P^1_{B'/A'}(M') \ar[r] &
P^0_{B'/A'}(M') \\
\ldots \ar[r] &
P^2_{B/A}(M) \ar[r] \ar[u] &
P^1_{B/A}(M) \ar[r] \ar[u] &
P^0_{B/A}(M) \ar[u]
}
$$
을 얻는다.
% P01-E0718: the frozen system diagram lacks terminal punctuation before the next sentence.
이 사상들은 이 유형의 사상들을 더 합성하는 것과 양립한다. 이를 보는 가장
쉬운 방법은 보조정리 \ref{algebra-lemma-module-principal-parts}의 증명에서
$P^k_{B/A}(M)$을 생성원과 관계식으로 나타낸 기술을 이용하는 것이지만,
이 가군들의 보편 성질로부터 직접 볼 수도 있다.
% P01-E0719: the frozen compound sentence lacks a comma before “but”.
더욱이 이 사상들은 보조정리 \ref{algebra-lemma-sequence-of-principal-parts}의
짧은 완전열들과 양립한다.
\end{remark}

\begin{lemma}
\label{algebra-lemma-principal-parts-base-change}
환 사상 $A \to A'$와 $A \to B$ 및 $B$-가군 $M$이 주어졌다고 하자.
$B' = B \otimes_A A'$으로 두고 $M' = M \otimes_A A'$을 $B'$-가군으로
보자. 주석 \ref{algebra-remark-functoriality-principal-parts}의 사상은 동형
$P^k_{B/A}(M) \otimes_A A' = P^k_{B'/A'}(M')$을 유도한다.
\end{lemma}

\begin{proof}
$D_{univ} : M \to P^k_{B/A}(M)$을 보편 미분 연산자라 하자. 유도된
$A'$-선형 사상
$D_{univ} \otimes 1 : M' = M \otimes_A A' \to P^k_{B/A}(M) \otimes_A A'$은
$M'/B'/A'$에 대한 차수 $k$의 미분 연산자임을 쉽게 확인할 수 있다.
$D_{univ} \otimes 1$이 보편 성질을 만족함을 보이겠다. $N'$을 $B'$-가군이라
하고 $D' : M' \to N'$을 차수 $k$의 미분 연산자라 하자. 그러면 합성
$D : M \to M' \to N'$은 $N'$을 $B$-가군으로 보았을 때 그 안으로 가는
미분 연산자이다. 따라서 $D = \gamma \circ D_{univ}$을 만족하는
$B$-선형 사상 $\gamma : P^k_{B/A}(M) \to N'$이 존재한다. 유도된
$B'$-선형 사상 $\gamma' : P^k_{B/A}(M) \otimes_A A' \to N'$이
$\gamma' \circ (D_{univ} \otimes 1) = D'$을 만족함은 독자가 쉽게 확인할
수 있다. $\gamma'$의 유일성에 대한 증명은 생략한다.
\end{proof}

\begin{lemma}
\label{algebra-lemma-principal-parts-diagonal}
$R \to S$를 환 사상이라 하자. $M$을 $S$-가군이라 하자.
$J = \Ker(S \otimes_R S \to S)$를 곱셈 사상의 핵이라 하자.
$S$-가군의 표준적 동형
$$
P^k_{S/R}(M) \longrightarrow (S \otimes_R M)/J^{k + 1}(S \otimes_R M)
$$
이 존재한다. 여기서 $s \in S$는 $s \otimes 1$을 곱하는 것으로 공역에
작용하며, 차수 $k$의 보편 미분 연산자는
$m \mapsto \text{class of }1 \otimes m$으로 주어지는 사상에 대응한다.
% P01-E0720: the frozen statement omits the verb “corresponds” and mismatches singular/plural in its universal-operator clause.
\end{lemma}

\begin{proof}
사상 $T : M \to S \otimes M$, $m \mapsto 1 \otimes m$을 생각하자.
% P01-E0721: the frozen map writes S tensor M without the base-ring subscript R used by every surrounding tensor product.
% P01-E0722: the frozen sentence also omits punctuation after the displayed map before “Since”.
$T$는 $R$-선형이고
\begin{align*}
g_0g_1\ldots g_k T(m)
- \sum g_0 \ldots \hat g_i \ldots g_k T(g_im) + \ldots
+(-1)^{k + 1}T(g_0\ldots g_km) \\
=
\prod_{i = 0,\ldots,k} (g_i \otimes 1 - 1 \otimes g_i) 1 \otimes m
\end{align*}
이 $m \in M$ 및 $g_0, \ldots, g_k \in S$에 대하여 성립하므로, 보조정리의
서술에 있는 규칙 $m \mapsto \text{class of }1 \otimes m$은 차수 $k$의
미분 연산자라고 결론내린다(보조정리
\ref{algebra-lemma-module-principal-parts}의 증명 참조). $P^k_{S/R}(M)$의 보편
성질에 의해 보조정리의 서술에 있는 화살표를 얻는다. 한편
$D : M \to N$이 차수 $k$의 미분 연산자이면 $S$-선형 사상
$L : S \otimes_R M \to N$, $g \otimes m \mapsto gD(m)$을 생각할 수 있다.
독자는 위와 비슷한 계산과 $J$가 $g \otimes 1 - 1 \otimes g$ 꼴의 원소들로
생성된다는 사실을 이용하여 $L$이 $J^{k + 1}(S \otimes_R M)$을 영으로
보냄을 확인할 수 있다. 특히 이를 보편 미분 연산자
$D_{univ} : M \to P^k_{S/R}(M)$에 적용하면 $S$-선형 사상
$(S \otimes_R M)/J^{k + 1}(S \otimes_R M) \to P^k_{S/R}(M)$을 얻는다.
이 사상이 보조정리의 서술에 있는 사상의 역임은 독자에게 맡긴다.
\end{proof}

\begin{lemma}
\label{algebra-lemma-differentials-de-rham-complex-order-1}
$A \to B$를 환 사상이라 하자. 미분
$\text{d} : \Omega^i_{B/A}  \to \Omega^{i + 1}_{B/A}$은 차수 $1$의
미분 연산자이다.
\end{lemma}

\begin{proof}
$b \in B$가 주어졌을 때 $\text{d} \circ b - b \circ \text{d}$가 선형
연산자임을 보여야 한다. 따라서
$$
\text{d} \circ b \circ b' - b \circ \text{d} \circ b' -
b' \circ \text{d} \circ b + b' \circ b \circ \text{d} = 0
$$
임을 보여야 한다.
% P01-E0723: the frozen commutator display lacks terminal punctuation before the next sentence.
이를 위해서는 $\Omega^i_{B/A}$의 가법적 생성원들에 대해 확인하면 충분하다.
따라서
$$
\text{d}(bb'b_0\text{d}b_1 \wedge \ldots \wedge \text{d}b_i) -
b\text{d}(b'b_0\text{d}b_1 \wedge \ldots \wedge \text{d}b_i) -
b'\text{d}(bb_0\text{d}b_1 \wedge \ldots \wedge \text{d}b_i) +
bb'\text{d}(b_0\text{d}b_1 \wedge \ldots \wedge \text{d}b_i)
$$
가 영임을 보이면 충분하다. 이는 라이프니츠 법칙을 이용하는 유쾌한 계산이며,
독자에게 맡긴다.
\end{proof}

\begin{lemma}
\label{algebra-lemma-check-differential-operators}
$A \to B$를 환 사상이라 하자. $g_i \in B$, $i \in I$를 $A$-대수로서
$B$의 생성원들의 집합이라 하자. $M, N$을 $B$-가군이라 하자.
$D : M \to N$을 $A$-선형 사상이라 하자. $D$가 차수 $k$의 미분
연산자임을 보이기 위해서는 $i \in I$에 대하여
$D \circ g_i - g_i \circ D$가 차수 $k - 1$의 미분 연산자임을 보이면 충분하다.
% P01-E0724: the frozen criterion omits k >= 1, although order k - 1 is undefined by the preceding definition when k = 0.
\end{lemma}

\begin{proof}
실제로 $D \circ g - g \circ D$가 차수 $k - 1$의 미분 연산자가 되게 하는
원소 $g \in B$들의 집합은 $B$의 $A$-부분대수라고 주장한다. 이는 관계식
$$
D \circ (g + g') - (g + g') \circ D =
(D \circ g - g \circ D) + (D \circ g' - g' \circ D)
$$
및
$$
D \circ gg' - gg' \circ D =
(D \circ g - g \circ D) \circ g' + g \circ (D \circ g' - g' \circ D)
$$
에서 따른다.
% P01-E0725: the frozen product display lacks terminal punctuation before “Strictly speaking”.
엄밀히 말하면 곱에 대해 결론을 내릴 때 보조정리
\ref{algebra-lemma-composition-differential-operators}도 사용한다.
\end{proof}

\begin{lemma}
\label{algebra-lemma-invert-system-differential-operators}
$A \to B$를 환 사상이라 하자. $M, N$을 $B$-가군이라 하자.
$S \subset B$를 곱셈적 부분집합이라 하자. 차수 $k$의 임의의 미분 연산자
$D : M \to N$은 차수 $k$의 미분 연산자
$E : S^{-1}M \to S^{-1}N$으로 유일하게 확장된다.
\end{lemma}

\begin{proof}
$k$에 대한 귀납법을 쓴다. $k = 0$이면 $D$는 $B$-선형이므로 국소화의
함자성에 의해 확장을 얻는다. $b \in B$가 주어졌을 때 연산자
$L_b : m \mapsto D(bm) - bD(m)$은 차수 $k - 1$이다. 따라서 귀납법에
의해 이는 차수 $k - 1$의 미분 연산자
$E_b : S^{-1}M \to S^{-1}N$으로 유일하게 확장된다. 또한 계산하면
$L_{b'b} = L_{b'} \circ b + b' \circ L_b$이고, 따라서 유일성에 의해
$E_{b'b} = E_{b'} \circ b + b' \circ E_b$를 얻는다.
% P01-E0728: the frozen sentence lacks punctuation before “hence”.
마찬가지로
$E_{b'} \circ b - b \circ E_{b'} =
E_b \circ b' - b' \circ E_b$를 얻는다.
이제 $m \in M$ 및 $g \in S$에 대하여
$$
E(m/g) = (1/g)(D(m) - E_g(m/g))
$$
로 둔다.
% P01-E0729: the frozen defining display lacks terminal punctuation before the next sentence.
이것이 잘 정의됨을 보이기 위해서는 $g' \in S$에 대하여 대표원
$g'm/g'g$를 사용해도 같은 결과를 얻음을 보이면 충분하다. 계산하면
\begin{align*}
(1/g'g)(D(g'm) - E_{g'g}(g'm/gg'))
& =
(1/gg')(g'D(m) + E_{g'}(m) - E_{g'g}(g'm/gg')) \\
& =
(1/g'g)(g'D(m) - g' E_g(m/g))
\end{align*}
% P01-E0730: the frozen align display lacks a comma before the following “which” clause.
이며, 이는 앞의 결과와 같다. $D$와 $E_g$가 $R$-선형이므로 $E$도
$R$-선형임은 분명하다.
% P01-E0726: the frozen proof uses an undefined ring R; the lemma's base ring is A.
$g = 1$로 놓고 $E_1 = 0$을 사용하면 $E$가 $D$를 확장함을 알 수 있다.
보조정리 \ref{algebra-lemma-check-differential-operators}에 의해 $E$가 차수 $k$의
미분 연산자임을 보이려면, $b \in B$에 대한
$E \circ b - b \circ E$와 $g' \in S$에 대한
$E \circ 1/g' - 1/g' \circ E$가 차수 $k - 1$의 미분 연산자임을 보이면
충분하다. 첫 번째 것에 대해서는 $S^{-1}M$의 원소 $m/g$를 택하고
다음을 관찰한다.
\begin{align*}
E(b m/g) - bE(m/g)
& =
(1/g)(D(bm) - bD(m) - E_g(bm/g) + bE_g(m/g)) \\
& =
(1/g)(L_b(m) - E_b(m) + gE_b(m/g)) \\
& =
E_b(m/g)
\end{align*}
% P01-E0731: the frozen align display lacks a comma before the following “which” clause.
이며, 이는 차수 $k - 1$의 미분 연산자이다. 마지막으로
\begin{align*}
E(m/g'g) - (1/g')E(m/g)
& =
(1/g'g)(D(m) - E_{g'g}(m/g'g)) -
(1/g'g)(D(m) - E_g(m/g)) \\
& =
-(1/g')E_{g'}(m/g'g)
\end{align*}
% P01-E0732: the frozen align display lacks a comma before the following “which” clause.
이며, 이 역시 선형 사상들($1/g'$를 곱하는 사상과 부호)과 $E_{g'}$의
합성이므로 차수 $k - 1$의 미분 연산자이다. 유일성의 증명은 생략한다.
\end{proof}

\begin{lemma}
\label{algebra-lemma-base-change-differential-operators}
$R \to A$와 $R \to B$를 환 사상이라 하자. $M$과 $M'$을 $A$-가군이라
하자. $D : M \to M'$을 $R \to A$에 관한 차수 $k$의 미분 연산자라 하자.
$N$을 임의의 $B$-가군이라 하자. 그러면 사상
$$
D \otimes \text{id}_N : M \otimes_R N \to M' \otimes_R N
$$
은 $B \to A \otimes_R B$에 관한 차수 $k$의 미분 연산자이다.
\end{lemma}

\begin{proof}
$D' = D \otimes \text{id}_N$이 $B$-선형임은 분명하다. 보조정리
\ref{algebra-lemma-check-differential-operators}에 의해
$$
D' \circ a \otimes 1 -  a \otimes 1 \circ D' =
(D \circ a - a \circ D) \otimes \text{id}_N
$$
이 차수 $k - 1$의 미분 연산자임을 보이면 충분하며, 이는 $k$에 대한
귀납법으로부터 따른다.
% P01-E0727: the frozen proof invokes the k - 1 criterion without treating k = 0, where negative order was never defined.
\end{proof}
\section{소박한 여접 복합체}
\label{algebra-section-netherlander}

\noindent
$R \to S$를 환 사상이라 하자. 변수들이 원소 $s \in S$인 다항식환을
$R[S]$로 나타내자.
% P01-E0733: the frozen designation “Denote R[S] the polynomial ring” omits “by/as”.
$s \in S$에 대응하는 변수를 $[s] \in R[S]$로 나타내자. 따라서 $R[S]$는
$s_1, \ldots, s_n$이 $S$의 원소들로 이루어진 모든 순서 없는 열을
달릴 때 기저원 $[s_1] \ldots [s_n]$ 위의 자유 $R$-가군이다.
% P01-E0734: the frozen source has subject-verb disagreement in “s_1,...,s_n ranges”.
표준적 전사
\begin{equation}
\label{algebra-equation-canonical-presentation}
R[S] \longrightarrow S,\quad [s] \longmapsto s
\end{equation}
가 존재하며, 그 핵을 $I \subset R[S]$로 나타낸다. $I$가 원소
$[s + s'] - [s] - [s']$, $[s][s'] - [ss']$ 및 $[r] - r$로
생성된다는 것은 간단한 관찰이다. 보조정리
\ref{algebra-lemma-differential-seq}에 따르면 표준적 사상
\begin{equation}
\label{algebra-equation-naive-cotangent-complex}
I/I^2 \longrightarrow \Omega_{R[S]/R} \otimes_{R[S]} S
\end{equation}
이 존재하며, 그 여핵은 표준적으로 $\Omega_{S/R}$와 동형이다.
$S$-가군 $\Omega_{R[S]/R} \otimes_{R[S]} S$는 생성원
$\text{d}[s]$ 위에서 자유임에 유의하라.

\begin{definition}
\label{algebra-definition-naive-cotangent-complex}
$R \to S$를 환 사상이라 하자. {\it 소박한 여접 복합체}
$\NL_{S/R}$는 사슬 복합체 (\ref{algebra-equation-naive-cotangent-complex})
$$
\NL_{S/R} = \left(I/I^2 \longrightarrow \Omega_{R[S]/R} \otimes_{R[S]} S\right)
$$
이다. 여기서 $I/I^2$는 (호몰로지) 차수 $1$에 놓이고,
$\Omega_{R[S]/R} \otimes_{R[S]} S$는 차수 $0$에 놓인다. 차수 $1$의
호몰로지를
$H_1(L_{S/R}) = H_1(\NL_{S/R})$\footnote{문헌에서는 이 가군을
때때로 $\Gamma_{S/R}$로 나타낸다.}로 나타낸다.
% P01-E0735: the frozen notation sentence “denote ... the homology” omits “by/as”.
\end{definition}

\noindent
계속하기 전에 실제 여접 복합체에 관해 몇 마디 하자
(Cotangent, Section \ref{cotangent-section-cotangent-ring-map}).
환 사상 $R \to S$가 주어지면, 그 항들이 다항식 대수인 표준적 단체
$R$-대수 $P_\bullet$이 존재하며, 이는 표준적 호모토피 동치
$$
P_\bullet \longrightarrow S
$$
를 갖춘다.
% P01-E0736: the frozen display lacks terminal punctuation before the next sentence.
$R$ 위의 $S$의 여접 복합체 $L_{S/R}$는 단체 가군
$$
\Omega_{P_\bullet/R} \otimes_{P_\bullet} S
$$
에 연관된 사슬 복합체로 정의된다.
% P01-E0737: the frozen display lacks terminal punctuation before the next sentence.
위에서 정의한 소박한 여접 복합체는 절단
$\tau_{\leq 1}L_{S/R}$와 표준적으로 동형이다
(Homology, Section \ref{homology-section-truncations} 및
Cotangent, Section \ref{cotangent-section-surjections} 참조).
따라서 실제로 $H_1(\NL_{S/R}) = H_1(L_{S/R})$이므로 우리의 정의는
여접 복합체를 사용하는 정의와 양립한다. 더욱이 위에서 본 바와 같이
$H_0(L_{S/R}) = H_0(\NL_{S/R}) = \Omega_{S/R}$이다.

\medskip\noindent
$R \to S$를 환 사상이라 하자. {\it $R$ 위의 $S$의 표시}란 $P$가
(어떤 변수 집합 위의) 다항식 대수인 $R$-대수의 전사
$\alpha : P \to S$를 말한다. $S$가 $R$ 위에서 유한형일 때에는 흔히
이를 “$R[x_1, \ldots, x_n] \to S$를 $S/R$의 표시라 하자”라고 쓰거나,
표시의 핵이 $I$임을 나타내고 싶으면
“$0 \to I \to R[x_1, \ldots, x_n] \to S \to 0$을 $S/R$의
표시라 하자”라고 쓴다. 소박한 여접 복합체를 정의할 때 사용한
사상 $R[S] \to S$도 표시의 한 예임에 유의하라.

\medskip\noindent
모든 표시 $\alpha$에 대하여 $S$-가군의 두 항 사슬 복합체
% P01-E0738: the frozen compound modifier “two term” should be hyphenated.
$$
\NL(\alpha) : I/I^2 \longrightarrow \Omega_{P/R} \otimes_P S.
$$
를 얻음에 유의하라. 여기서 항 $I/I^2$는 차수 $1$에 놓이고 항
$\Omega_{P/R} \otimes S$는 차수 $0$에 놓인다. $I/I^2$에서
$f \in I$의 동치류는 $\Omega_{P/R} \otimes S$의
$\text{d}f \otimes 1$로 보내진다.
% P01-E0739: the frozen two subsequent tensor products omit the base P shown in the defining display.
이 복합체의 여핵은 표준적으로 $\Omega_{S/R}$이다. 보조정리
\ref{algebra-lemma-differential-seq}를 보라. 복합체 $\NL(\alpha)$를
{\it $S/R$의 표시 $\alpha : P \to S$에 연관된 소박한 여접 복합체}라
한다. $P = R[S]$이고 $S$로 가는 표준적 전사가 주어지면
$\NL_{S/R}$를 다시 얻는다. $P = R[x_1, \ldots, x_n]$이면 때때로
이 복합체를
$I/I^2 \to \bigoplus_{i = 1, \ldots, n} S\text{d}x_i$로 나타낸다.
% P01-E0740: the frozen phrase “then will sometimes use” omits the subject “we”.

\medskip\noindent
환들의 가환 그림
\begin{equation}
\label{algebra-equation-functoriality-NL}
\vcenter{
\xymatrix{
S \ar[r]_{\phi} & S' \\
R \ar[r] \ar[u] & R' \ar[u]
}
}
\end{equation}
이 주어졌다고 하자. $\alpha : P \to S$를 $R$ 위의 $S$의 표시라 하고,
$\alpha' : P' \to S'$을 $R'$ 위의 $S'$의 표시라 하자.
{\it 표시 $\alpha : P \to S$에서 표시 $\alpha' : P' \to S'$으로 가는
표시의 사상}은 $\phi \circ \alpha = \alpha' \circ \varphi$를 만족하는
$R$-대수 사상
$$
\varphi : P \to P'
$$
로 정의된다. 이 경우 $\varphi(I) \subset I'$임에 유의하라. 여기서
$I = \Ker(\alpha)$이고 $I' = \Ker(\alpha')$이다. 따라서 $\varphi$는
$S$-가군의 사상 $I/I^2 \to I'/(I')^2$을 유도하고, 미분의 함자성에
의해 $S$-가군의 사상
$\Omega_{P/R} \otimes S \to \Omega_{P'/R'} \otimes S'$도 유도한다.
% P01-E0741: the frozen functoriality formula omits the tensor bases P and P'.
이 사상들은 $\NL(\alpha)$와 $\NL(\alpha')$의 미분과 양립하며,
소박한 여접 복합체의 사상
$$
\NL(\alpha) \longrightarrow \NL(\alpha').
$$
을 얻는다. 유도된 사상
$\NL(\alpha) \otimes_S S' \to \NL(\alpha')$을 생각하는 것이 흔히 편리하다.

\medskip\noindent
특수한 경우 $P = R[S]$이고 $P' = R'[S']$이면, 사상
$\phi : S \to S'$는 규칙 $[s] \mapsto [\phi(s)]$에 의해 표준적 환 사상
$\varphi : P \to P'$을 유도한다. 따라서 위 구성은 그림
(\ref{algebra-equation-functoriality-NL})에 연관된 사슬 복합체의 표준적(!) 사상
$$
\NL_{S/R} \longrightarrow \NL_{S'/R'},\quad\text{and}\quad
\NL_{S/R} \otimes_S S' \longrightarrow \NL_{S'/R'}
$$
을 정한다. 이 구성은 합성과 양립함에 유의하라. 가환 그림
$$
\xymatrix{
S \ar[r]_{\phi} & S' \ar[r]_{\phi'} & S'' \\
R \ar[r] \ar[u] & R' \ar[u] \ar[r] & R'' \ar[u]
}
$$
이 주어지면,
% P01-E0742: the frozen diagram lacks a comma before the continuing “we see” clause.
사상들의 합성
$$
\NL_{S/R} \longrightarrow \NL_{S'/R'} \longrightarrow \NL_{S''/R''}
$$
은 바깥 정사각형이 주는 사상 $\NL_{S/R} \to \NL_{S''/R''}$임을 알 수 있다.

\medskip\noindent
$\NL(\alpha)$는 $\NL_{S/R}$와 호모토피 동치이고, 위에서 구성한
사상들은 호모토피를 제외하고 잘 정의된다는 것이 밝혀진다
(복합체 사상의 호모토피는
Homology, Section \ref{homology-section-complexes}에서 논하지만,
아래 보조정리의 증명에서도 각 명제의 정확한 뜻을 명시한다).

\begin{lemma}
\label{algebra-lemma-NL-homotopy}
그림 (\ref{algebra-equation-functoriality-NL})이 주어졌다고 하자.
$\alpha : P \to S$와 $\alpha' : P' \to S'$을 표시라 하자.
\begin{enumerate}
\item $\alpha$에서 $\alpha'$으로 가는 표시의 사상이 존재한다.
\item 표시의 임의의 두 사상은 복합체의 호모토픽한 사상
$\NL(\alpha) \to \NL(\alpha')$을 유도한다.
\item 이 구성은 표시의 사상들의 합성과 양립한다
(정확한 명제는 증명을 보라).
\item $R \to R'$과 $S \to S'$이 동형사상이면, $\alpha$에서
$\alpha'$으로 가는 표시의 임의의 사상 $\varphi$에 대하여 유도된 사상
$\NL(\alpha) \to \NL(\alpha')$은 호모토피 동치이고 준동형이다.
\end{enumerate}
특히 $\alpha$를 표준적 표시 (\ref{algebra-equation-canonical-presentation})와
비교하면, 호모토피를 제외하고 잘 정의되며 모든 함자성과
(호모토피를 제외하고) 양립하는 준동형
$\NL(\alpha) \to \NL_{S/R}$가 존재한다고 결론내린다.
\end{lemma}

\begin{proof}
$P$는 $R$ 위의 다항식 대수이므로 어떤 집합 $A$에 대하여
$P = R[x_a, a \in A]$로 쓸 수 있다. $\alpha'$은 전사이므로, 모든
$a \in A$에 대하여 $\alpha'(f_a) = \phi(\alpha(x_a))$를 만족하는
원소 $f_a \in P'$을 택할 수 있다. $\varphi(x_a) = f_a$를 만족하는
유일한 $R$-대수 사상
$\varphi : P = R[x_a, a \in A] \to P'$을 취한다. 이것이 (1)의
사상을 준다.

\medskip\noindent
$\varphi$와 $\varphi'$을 $\alpha$에서 $\alpha'$으로 가는 표시의
사상들이라 하자.
% P01-E0743: the frozen source omits “be” in “Let phi and phi' morphisms”.
$I = \Ker(\alpha)$ 및 $I' = \Ker(\alpha')$로 두자. 그림
$$
\xymatrix{
I/I^2 \ar[r]^-{\text{d}}
\ar@<1ex>[d]^{\varphi'_1} \ar@<-1ex>[d]_{\varphi_1}
&
\Omega_{P/R} \otimes_P S
\ar@<1ex>[d]^{\varphi'_0} \ar@<-1ex>[d]_{\varphi_0}
\ar[ld]_h
\\
I'/(I')^2 \ar[r]^-{\text{d}}
&
\Omega_{P'/R'} \otimes_{P'} S'
}
$$
에서 대각 사상 $h$를 구성해야 한다. 여기서 수직 사상들은
$\varphi$, $\varphi'$이 유도하며, 이 대각 사상은
$$
\varphi_1 - \varphi'_1 = h \circ \text{d}
\quad\text{and}\quad
\varphi_0 - \varphi'_0 = \text{d} \circ h
$$
를 만족해야 한다.
% P01-E0744: the frozen “where ... induced by phi, phi' such that” clause omits a conjunction and the subject h.
사상 $\varphi - \varphi' : P \to P'$을 생각하자. $\varphi$와
$\varphi'$은 모두 $\alpha$ 및 $\alpha'$과 양립하므로
$\varphi - \varphi' : P \to I'$을 얻는다. 이는
$\varphi, \varphi' : P \to P'$이 $I'/(I')^2$ 위에 같은
$P$-가군 구조를 유도함을 뜻한다. 실제로
$\varphi(p)i' - \varphi'(p)i' = (\varphi - \varphi')(p)i' \in (I')^2$이다.
또한 $\varphi - \varphi'$은 $R$-선형이고
$$
(\varphi - \varphi')(fg) =
\varphi(f)(\varphi - \varphi')(g) + (\varphi - \varphi')(f)\varphi'(g)
$$
이다. 따라서 유도된 사상 $D : P \to I'/(I')^2$은 $R$-미분이다.
% P01-E0745: the frozen source says “a R-derivation” instead of “an R-derivation”.
그러므로 $D = h \circ \text{d}$를 만족하는 표준적 사상
$h : \Omega_{P/R} \otimes_P S \to I'/(I')^2$을 얻는다.
계산(생략함)은 $h$가 원하는 호모토피임을 보인다.

\medskip\noindent
가환 그림
$$
\xymatrix{
S \ar[r]_{\phi} & S' \ar[r]_{\phi'} & S'' \\
R \ar[r] \ar[u] & R' \ar[u] \ar[r] & R'' \ar[u]
}
$$
이 주어졌고,
\begin{enumerate}
\item $\alpha : P \to S$,
\item $\alpha' : P' \to S'$, 그리고
\item $\alpha'' : P'' \to S''$
\end{enumerate}
이 표시라고 하자. 또한
\begin{enumerate}
\item $\varphi : P \to P'$이 $\alpha$에서 $\alpha'$으로 가는
표시의 사상이고,
\item $\varphi' : P' \to P''$이 $\alpha'$에서 $\alpha''$으로 가는
표시의 사상이라고 하자.
\end{enumerate}
그러면 $\varphi' \circ \varphi : P \to P''$이 $\alpha$에서
$\alpha''$으로 가는 표시의 사상임은 즉시 알 수 있으며, 이것이 유도하는
소박한 여접 복합체의 사상 $\NL(\alpha) \to \NL(\alpha'')$은
$\varphi$와 $\varphi'$이 유도하는 사상
$\NL(\alpha) \to \NL(\alpha')$과
$\NL(\alpha') \to \NL(\alpha'')$의 합성이다.

\medskip\noindent
두 항을 갖는 복합체의 단순한 경우에 준동형이란 두 항 사이의 사상의
여핵과 핵 모두에 동형을 유도하는 사상일 뿐이다.
% P01-E0746: the frozen compound modifier “2 term” should be hyphenated.
위에서 설명한 것처럼 두 항 복합체의 호모토픽한 사상들은 핵과 여핵에
같은 사상들을 정의함에 유의하라. 따라서 $\varphi$가 $R$ 위의 $S$의
표시 $\alpha$에서 자기 자신으로 가는 사상이면, 유도된 사상
$\NL(\alpha) \to \NL(\alpha)$은 (2)에 의해 항등사상과 호모토픽하므로
준동형이다.
% P01-E0747: the frozen phrase “is a quasi-isomorphism being homotopic” uses an awkward/dangling participial construction.
(4)를 완전히 일반적으로 증명하기 위해, (1)에 의해 존재하는
$\alpha'$에서 $\alpha$로 가는 사상 $\varphi'$을 생각하자. 합성
$\NL(\alpha) \to \NL(\alpha') \to \NL(\alpha)$ 및
$\NL(\alpha') \to \NL(\alpha) \to \NL(\alpha')$은 (3)에 의해 항등사상들과
호모토픽하며, 따라서 이 사상들은 정의상 호모토피 동치이다.
% P01-E0748: compatibility in (3) alone does not give homotopy to the identity; the argument also needs part (2).
형식적으로 두 사상
$\NL(\alpha) \to \NL(\alpha')$과 $\NL(\alpha') \to \NL(\alpha)$이
모두 준동형임이 따른다. 일부 세부사항은 생략한다.
\end{proof}

\begin{lemma}
\label{algebra-lemma-NL-polynomial-algebra}
$A \to B$를 다항식 대수라 하자. 그러면 $\NL_{B/A}$는
$\Omega_{B/A}$가 차수 $0$에 놓인 사슬 복합체
$(0 \to \Omega_{B/A})$와 호모토피 동치이다.
\end{lemma}

\begin{proof}
이는 보조정리 \ref{algebra-lemma-NL-homotopy}와
$\text{id}_B : B \to B$가 핵이 영인 $A$ 위의 $B$의 표시라는 사실에서 따른다.
% P01-E0749: the frozen proof is the sentence fragment “Follows from ...”.
\end{proof}

\noindent
다음 보조정리는 소박한 여접 복합체를 도입하는 동기의 일부이다.
여접 복합체는 이를 진정한 긴 완전 코호몰로지 열로 확장한다.
$B \to C$가 국소 완전 교차이면 이 열의 왼쪽에 영을 붙여 확장할 수 있다.
More on Algebra, Lemma
\ref{more-algebra-lemma-transitive-lci-at-end}를 보라.

\begin{lemma}[Jacobi--Zariski 열]
\label{algebra-lemma-exact-sequence-NL}
$A \to B \to C$를 환 사상이라 하자. 핵이 $I$인 표시
$\alpha : A[x_s, s \in S] \to B$를 택하자. 핵이 $J$인 표시
$\beta : B[y_t, t \in T] \to C$를 택하자.
$\gamma : A[x_s, y_t] \to C$를 핵이 $K$인 유도된 $C$의 표시라 하자.
그러면 행들이 완전인 표준적 가환 그림
$$
\xymatrix{
0 \ar[r] &
\Omega_{A[x_s]/A} \otimes C \ar[r] &
\Omega_{A[x_s, y_t]/A} \otimes C \ar[r] &
\Omega_{B[y_t]/B} \otimes C \ar[r] &
0 \\
&
I/I^2 \otimes C \ar[r] \ar[u] &
K/K^2 \ar[r] \ar[u] &
J/J^2 \ar[r] \ar[u] &
0
}
$$
을 얻는다.
% P01-E0750: the frozen display lacks terminal punctuation before the next sentence.
다음과 같은 호몰로지 군의 완전열
$$
H_1(\NL_{B/A} \otimes_B C) \to
H_1(L_{C/A}) \to
H_1(L_{C/B}) \to
C \otimes_B \Omega_{B/A} \to
\Omega_{C/A} \to
\Omega_{C/B} \to 0
$$
을 얻으며, 이는 보조정리
\ref{algebra-lemma-exact-sequence-differentials}의 열을 확장하는
$C$-가군의 열이다.
% P01-E0751: the frozen display lacks terminal punctuation before the next sentence.
$\text{Tor}_1^B(\Omega_{B/A}, C) = 0$이고
$\text{Tor}_2^B(\Omega_{B/A}, C) = 0$이면
$H_1(\NL_{B/A} \otimes_B C) = H_1(L_{B/A}) \otimes_B C$이다.
\end{lemma}

\begin{proof}
사상들의 정확한 정의는 생략한다. 위 행의 완전성은
$\text{d}x_s$, $\text{d}y_t$가 가운데 가군의 기저를 이루기 때문에
따른다. 사상 $\gamma$는
$$
A[x_s, y_t] \to B[y_t] \to C
$$
로 인수분해되며, 첫 번째 화살표는 전사이고 두 번째 화살표는
$\beta$와 같다. 따라서 $K \to J$가 전사임을 알 수 있다.
더욱이 첫 번째로 표시된 화살표의 핵은 $IA[x_s, y_t]$이다.
따라서 $I/I^2 \otimes C$는 $K/K^2 \to J/J^2$의 핵으로 전사한다.
마지막으로 보조정리 \ref{algebra-lemma-NL-homotopy}를 사용하여 각 항을
소박한 여접 복합체의 호몰로지 군으로 식별할 수 있다.

\medskip\noindent
마지막 주장은 호몰로지 대수의 명제이다. $\NL_{B/A} = (N^{-1} \to N^0)$가
$N^0$가 자유이고 코호몰로지 가군들이
$H^0 = \Omega_{B/A}$와 $H^{-1} = H_1(L_{B/A})$인 두 항
$B$-가군 복합체임을 상기하라. 미분의 상을 $M \subset N^0$으로 쓰자.
$\text{Tor}_1^B(H^0, C) = 0$이면 완전열
$$
0 \to M \otimes_B C \to N^0 \otimes_B C \to H^0 \otimes_B C \to 0
$$
을 얻는다. $N^0$가 자유이므로
$\text{Tor}_2^B(H^0, C) = \text{Tor}_1^B(M, C)$임도 알 수 있다.
따라서 $\text{Tor}_2^B(H^0, C) = 0$이면 완전열
$$
0 \to H^{-1} \otimes_B C \to N^{-1} \otimes_B C \to M \otimes_B C \to 0
$$
도 얻는다. 모두 종합하면
$\text{Tor}_1^B(H^0, C) = 0$이고 $\text{Tor}_2^B(H^0, C) = 0$일 때
$H^{-1} \otimes_B C$는 원하는 대로
$N^{-1} \otimes_B C \to N^0 \otimes_B C$의 핵이다.
\end{proof}

\begin{remark}
\label{algebra-remark-composition-homotopy-equivalent-to-zero}
$A \to B$와 $\phi : B \to C$를 환 사상이라 하자. 그러면 합성
$\NL_{B/A} \to \NL_{C/A} \to \NL_{C/B}$은 영과 호모토피 동치이다.
실제로 이 합성은 정사각형
$$
\xymatrix{
B \ar[r]_\phi & C \\
A \ar[r] \ar[u] & B \ar[u]
}
$$
에 대한 소박한 여접 복합체의 함자성이다.
% P01-E0752: the frozen diagram lacks terminal punctuation before the next sentence.
$J = \Ker(B[C] \to C)$로 쓰자. 명시적 호모토피는 기저원
$\text{d}[b]$를 $J/J^2$에서 $[\phi(b)] - b$의 동치류로 보내는 사상
$\Omega_{A[B]/A} \otimes_{A[B]} B \to J/J^2$로 주어진다.
\end{remark}

\begin{lemma}
\label{algebra-lemma-NL-surjection}
$A \to B$를 핵이 $I$인 전사 환 사상이라 하자. 그러면
$\NL_{B/A}$는 $I/I^2$가 차수 $1$에 놓인 사슬 복합체
$(I/I^2 \to 0)$와 호모토피 동치이다. 특히
$H_1(L_{B/A}) = I/I^2$이다.
\end{lemma}

\begin{proof}
이는 보조정리 \ref{algebra-lemma-NL-homotopy}와 $A \to B$가 $A$ 위의
$B$의 표시라는 사실에서 따른다.
% P01-E0753: the frozen proof is the sentence fragment “Follows from ...”.
\end{proof}

\begin{lemma}
\label{algebra-lemma-application-NL}
$A \to B \to C$를 환 사상이라 하자. $A \to C$가 전사라고 가정하자
(따라서 $B \to C$도 전사이다). $I = \Ker(A \to C)$ 및
$J = \Ker(B \to C)$로 나타내자. 그러면 열
$$
I/I^2 \to J/J^2 \to \Omega_{B/A} \otimes_B B/J \to 0
$$
은 완전하다.
\end{lemma}

\begin{proof}
이는 보조정리 \ref{algebra-lemma-exact-sequence-NL}와 보조정리
\ref{algebra-lemma-NL-surjection}에 있는 소박한 여접 복합체
$\NL_{C/B}$ 및 $\NL_{C/A}$의 기술에서 따른다.
% P01-E0754: the frozen proof is the sentence fragment “Follows from ...”.
\end{proof}

\begin{lemma}[평탄 기저변환]
\label{algebra-lemma-change-base-NL}
$R \to S$를 환 사상이라 하자. $\alpha : P \to S$를 표시라 하자.
$R \to R'$을 평탄 환 사상이라 하자. $\alpha' : P \otimes_R R' \to
S' = S \otimes_R R'$을 유도된 표시라 하자. 그러면
$\NL(\alpha) \otimes_R R' = \NL(\alpha) \otimes_S S' = \NL(\alpha')$이다.
특히 $R \to R'$이 평탄이면 표준적 사상
$$
\NL_{S/R} \otimes_S S' \longrightarrow \NL_{S \otimes_R R'/R'}
$$
은 호모토피 동치이다.
\end{lemma}

\begin{proof}
$R \to R'$이 평탄이므로
$\Ker(\alpha') = R' \otimes_R \Ker(\alpha)$이며, 따라서 이 명제가 성립한다.
\end{proof}

\begin{lemma}
\label{algebra-lemma-colimits-NL}
$R_\lambda \to S_\lambda$를 유향집합 $\Lambda$ 위의 환 사상들의
계라 하자. $R = \colim R_\lambda$ 및 $S = \colim S_\lambda$로 두자.
그러면 $\NL_{S/R} = \colim \NL_{S_\lambda/R_\lambda}$이다.
\end{lemma}

\begin{proof}
$\NL_{S/R}$가 복합체
$I/I^2 \to \bigoplus_{s \in S} S\text{d}[s]$임을 상기하라. 여기서
$I \subset R[S]$는 표준적 표시 $R[S] \to S$의 핵이다. 이제
$R[S] = \colim R_\lambda[S_\lambda]$이고, 마찬가지로
$I_\lambda = \Ker(R_\lambda[S_\lambda] \to S_\lambda)$라 할 때
$I = \colim I_\lambda$임은 분명하다. 따라서 보조정리는 자명하다.
\end{proof}

\begin{lemma}
\label{algebra-lemma-NL-of-localization}
$S \subset A$가 $A$의 곱셈적 부분집합이면,
$\NL_{S^{-1}A/A}$는 영 복합체와 호모토피 동치이다.
\end{lemma}

\begin{proof}
$A \to S^{-1}A$가 평탄이므로 평탄 기저변환
(보조정리 \ref{algebra-lemma-change-base-NL})에 의해
$\NL_{S^{-1}A/A} \otimes_A S^{-1}A \to \NL_{S^{-1}A/S^{-1}A}$가
호모토피 동치임을 알 수 있다. 이 화살표의 정의역은
$\NL_{S^{-1}A/A}$와 동형이고, 공역은 보조정리
\ref{algebra-lemma-NL-surjection}에 의해 영이므로 결론을 얻는다.
\end{proof}

\begin{lemma}
\label{algebra-lemma-NL-localize-bottom}
$S \subset A$를 $A$의 곱셈적 부분집합이라 하자.
% P01-E0755: the frozen statement has “Let S ... is” instead of “be”.
$S^{-1}A \to B$를 환 사상이라 하자. 그러면
$\NL_{B/A} \to \NL_{B/S^{-1}A}$는 호모토피 동치이다.
\end{lemma}

\begin{proof}
$\alpha : P \to B$를 $A$ 위의 $B$의 표시로 택하자. 그러면
$\beta : S^{-1}P \to B$는 $S^{-1}A$ 위의 $B$의 표시이다. 직접 계산하면
$\NL(\alpha) = \NL(\beta)$를 얻는다. 소박한 여접 복합체는 보조정리
\ref{algebra-lemma-NL-homotopy}에 의해 호모토피를 제외하고 잘 정의되므로
이 등식이 보조정리를 증명한다.
% P01-E0756: the frozen concluding clause is weakly joined by an unpunctuated “as”.
\end{proof}

\begin{lemma}
\label{algebra-lemma-principal-localization-NL}
\begin{slogan}
소박한 여접 복합체의 형성은 한 원소에서의 국소화와 가환한다.
\end{slogan}
$A \to B$를 환 사상이라 하자. $g \in B$라 하자. $\alpha : P \to B$가
핵이 $I$인 표시라고 하자. 그러면 $A$ 위의 $B_g$의 한 표시는
$\alpha$를 확장하고 $x$를 $1/g$로 보내는 사상
$$
\beta : P[x] \longrightarrow B_g
$$
이다. $\beta$의 핵 $J$는 $I$와 원소 $f x - 1$로 생성된다. 여기서
$f \in P$는 $\alpha$에 의해 $g \in B$로 보내지는 원소이다. 이 상황에서
\begin{enumerate}
\item $J/J^2 = (I/I^2)_g \oplus B_g (f x - 1)$,
\item $\Omega_{P[x]/A} \otimes_{P[x]} B_g =
\Omega_{P/A} \otimes_P B_g \oplus B_g \text{d}x$,
\item $\NL(\beta) \cong
\NL(\alpha) \otimes_B B_g \oplus (B_g \xrightarrow{g} B_g)$
\end{enumerate}
이다.
% P01-E0757: the frozen enumeration has no terminal punctuation.
따라서 표준적 사상 $\NL_{B/A} \otimes_B B_g \to \NL_{B_g/A}$는
호모토피 동치이다.
\end{lemma}

\begin{proof}
$P[x]/(I, fx - 1) = B[x]/(gx - 1) = B_g$이므로 $I$와 $fx - 1$이
$J$를 생성한다는 명제를 얻는다. 보조정리
\ref{algebra-lemma-exact-sequence-NL}의 완전한 행들을 갖는 가환 그림
$$
\xymatrix{
0 \ar[r] &
\Omega_{P/A} \otimes B_g \ar[r] &
\Omega_{P[x]/A} \otimes B_g \ar[r] &
\Omega_{B[x]/B} \otimes B_g \ar[r] &
0 \\
&
(I/I^2)_g \ar[r] \ar[u] &
J/J^2 \ar[r] \ar[u] &
(gx - 1)/(gx - 1)^2 \ar[r] \ar[u] &
0
}
$$
을 생각하자.
% P01-E0758: several tensor products in the frozen diagram and prose omit their inferable base rings.
$B_g$-가군 $\Omega_{B[x]/B} \otimes B_g$는 $\text{d}x$ 위의 랭크
$1$ 자유가군이다. $B_g$-가군 $\Omega_{P[x]/A} \otimes B_g$의 원소
$\text{d}x$는 위 행의 분할을 준다. 원소
$gx - 1 \in (gx - 1)/(gx - 1)^2$은
$\Omega_{B[x]/B} \otimes B_g$의 $g\text{d}x$로 보내지므로
$(gx - 1)/(gx - 1)^2$은 $B_g$ 위의 랭크 $1$ 자유가군이다.
(이는 $gx - 1$이 가역 상수항을 갖는 다항식이므로 $B[x]$에서
비영인자가 되고, 임의의 비영인자는 보조정리
\ref{algebra-lemma-regular-quasi-regular}에 의해 길이 $1$의 준정칙열을 준다고
논증하여도 알 수 있다.)

\medskip\noindent
$(I/I^2)_g \to J/J^2$가 단사임을 증명하자.
% P01-E0759: the frozen sentence “Let us prove ... injective” omits “is”.
$x$를 $1/f$로 보내는 $P$-대수 사상
$$
\pi : P[x] \to (P/I^2)_f = P_f/I_f^2
$$
을 생각하자. $J$는 $I$와 $fx - 1$로 생성되므로
$\pi(J) \subset (I/I^2)_f = (I/I^2)_g$임을 알 수 있다. 이는 제곱이
영인 이데알이므로 $\pi(J^2) = 0$임을 알 수 있다. $a \in I$의 상이
$J$에서 $J^2$의 원소이면 $\pi(a) = 0$이고, 이는 $a$가
$I_f/I_f^2$에서 영으로 보내짐을 뜻한다.
% P01-E0760: the frozen phrase “maps to an element of J^2 in J” is grammatically malformed.
이로써 원하는 단사성을 증명했다.

\medskip\noindent
따라서 두 항 복합체의 짧은 완전열
% P01-E0761: the frozen compound adjective “two term” is unhyphenated.
$$
0 \to \NL(\alpha) \otimes_B B_g \to \NL(\beta)
\to (B_g \xrightarrow{g} B_g) \to 0
$$
을 얻는다. 그러한 짧은 완전열은 언제나 복합체의 범주에서 분할할 수 있다.
% P01-E0762: the frozen blanket splitting claim is false for arbitrary short exact sequences of complexes; the explicit termwise splittings below justify only this case.
우리의 특수한 경우에는 다음을 분할로 택할 수 있다.
$$
J/J^2 = (I/I^2)_g \oplus B_g (fx - 1)\quad\text{and}\quad
\Omega_{P[x]/A} \otimes B_g = \Omega_{P/A} \otimes B_g \oplus
B_g (g^{-2}\text{d}f + \text{d}x)
$$
% P01-E0763: the frozen display lacks terminal punctuation before “This works”.
이는 $\Omega_{P[x]/A} \otimes B_g$에서
$\text{d}(fx - 1) = x\text{d}f + f \text{d}x =
g(g^{-2}\text{d}f + \text{d}x)$이기 때문에 성립한다.
\end{proof}

\begin{lemma}
\label{algebra-lemma-localize-NL}
$A \to B$를 환 사상이라 하자. $S \subset B$를 곱셈적 부분집합이라 하자.
표준적 사상 $\NL_{B/A} \otimes_B S^{-1}B \to \NL_{S^{-1}B/A}$는
준동형이다.
\end{lemma}

\begin{proof}
$S$를 나눗셈 관계로 순서화한 유향집합으로 생각하면
$S^{-1}B = \colim_{g \in S} B_g$이다. 보조정리
\ref{algebra-lemma-localization-colimit}를 보라. 보조정리
\ref{algebra-lemma-principal-localization-NL}에 의해 각 사상
$\NL_{B/A} \otimes_B B_g \to \NL_{B_g/A}$은 준동형이다.
% P01-E0764: the frozen subject “each” incorrectly takes the plural verb “are”.
보조정리 \ref{algebra-lemma-colimits-NL}로부터 보조정리가 따른다.
\end{proof}

\begin{lemma}
\label{algebra-lemma-sum-two-terms}
$R$을 환이라 하자. $A_1 \to A_0$와 $B_1 \to B_0$를 두 항
복합체라 하자.
% P01-E0765: the frozen sentence has both a spurious comma after the first complex and an unhyphenated “two term”.
복합체의 사상 $\varphi : A_\bullet \to B_\bullet$ 및
$\psi : B_\bullet \to A_\bullet$가 존재하여
$\varphi \circ \psi$와 $\psi \circ \varphi$가 항등사상들과
호모토픽하다고 하자. 그러면 $R$-가군으로서
$A_1 \oplus B_0 \cong B_1 \oplus A_0$이다.
\end{lemma}

\begin{proof}
다음을 만족하는 사상 $h : A_0 \to A_1$을 택하자.
$$
\text{id}_{A_1} - \psi_1 \circ \varphi_1 = h \circ d_A
\text{ and }
\text{id}_{A_0} - \psi_0 \circ \varphi_0 = d_A \circ h.
$$
마찬가지로 다음을 만족하는 사상 $h' : B_0 \to B_1$을 택하자.
$$
\text{id}_{B_1} - \varphi_1 \circ \psi_1 = h' \circ d_B
\text{ and }
\text{id}_{B_0} - \varphi_0 \circ \psi_0 = d_B \circ h'.
$$
간단한 계산으로
$$
\left(
\begin{matrix}
\text{id}_{A_1} & -h' \circ \psi_1 + h \circ \psi_0 \\
0 & \text{id}_{B_0}
\end{matrix}
\right)
=
\left(
\begin{matrix}
\psi_1 & h \\
-d_B & \varphi_0
\end{matrix}
\right)
\left(
\begin{matrix}
\varphi_1 & - h' \\
d_A & \psi_0
\end{matrix}
\right)
$$
임을 알 수 있다.
% P01-E0766: the frozen matrix display lacks terminal punctuation.
이는 오른쪽의 두 행렬이 모두 가역임을 보이며, 보조정리를 증명한다.
% P01-E0767: invertibility of the displayed product alone does not imply both factors are invertible over arbitrary modules; the reverse product or explicit inverses are not supplied.
\end{proof}

\begin{lemma}
\label{algebra-lemma-conormal-module}
$R \to S$를 유한형 환 사상이라 하자. 임의의 표시
$\alpha : R[x_1, \ldots, x_n] \to S$와
$\beta : R[y_1, \ldots, y_m] \to S$에 대하여
% P01-E0768: the frozen statement has a spurious comma before “and” between the two presentations.
$$
I/I^2 \oplus S^{\oplus m} \cong J/J^2 \oplus S^{\oplus n}
$$
이 $S$-가군의 동형으로 성립한다. 여기서 $I = \Ker(\alpha)$이고
$J = \Ker(\beta)$이다.
% P01-E0769: the frozen phrase “as S-modules where” lacks a comma before the defining clause.
\end{lemma}

\begin{proof}
보조정리 \ref{algebra-lemma-NL-homotopy}와
\ref{algebra-lemma-sum-two-terms}를 보라.
\end{proof}

\begin{lemma}
\label{algebra-lemma-conormal-module-localize}
$R \to S$를 유한형 환 사상이라 하자. $g \in S$라 하자. 임의의 표시
$\alpha : R[x_1, \ldots, x_n] \to S$와
$\beta : R[y_1, \ldots, y_m] \to S_g$에 대하여
% P01-E0770: the frozen statement has a spurious comma before “and” between the two presentations.
$$
(I/I^2)_g \oplus S^{\oplus m}_g \cong J/J^2 \oplus S_g^{\oplus n}
$$
이 $S_g$-가군의 동형으로 성립한다. 여기서
$I = \Ker(\alpha)$이고 $J = \Ker(\beta)$이다.
% P01-E0771: the frozen phrase “as S_g-modules where” lacks a comma before the defining clause.
\end{lemma}

\begin{proof}
$\beta' : R[x_1, \ldots, x_n, x] \to S_g$를 $\alpha$에서 시작하여
구성한 보조정리 \ref{algebra-lemma-principal-localization-NL}의 표시라 하자.
그러면 $\NL(\alpha) \otimes_S S_g$가 $\NL(\beta')$와 호모토피
동치임을 안다. 보조정리 \ref{algebra-lemma-NL-homotopy}에 의해
$\NL(\beta)$와 $\NL(\beta')$도 호모토피 동치이다. 따라서
$\NL(\alpha) \otimes_S S_g$는 $\NL(\beta)$와 호모토피 동치라고
결론내린다. 마지막으로 보조정리 \ref{algebra-lemma-conormal-module}를 적용한다.
\end{proof}

\section{국소 완전 교차}
\label{algebra-section-lci}

\noindent
국소 완전 교차라는 성질은 뇌터 국소환의 내재적 성질이다.
이 내용은 나눗셈 거듭제곱 대수, 절 \ref{dpa-section-lci}에서
논의할 것이다. 그러나 지금은 체 위의 유한형 대수에 대해서만
이 성질을 정의한다.

\begin{definition}
\label{algebra-definition-lci-field}
$k$를 체라 하자.
$S$를 유한형 $k$-대수라 하자.
\begin{enumerate}
\item 어떤 표시
$S = k[x_1, \ldots, x_n]/(f_1, \ldots, f_c)$가 존재하여
$\dim(S) = n - c$를 만족하면, $S$를 {\it $k$ 위의 전역 완전 교차}라
한다.
\item 어떤 피복 $\Spec(S) = \bigcup D(g_i)$가 존재하여
각 환 $S_{g_i}$가 $k$ 위의 전역 완전 교차이면,
$S$를 {\it $k$ 위의 국소 완전 교차}라 한다.
\end{enumerate}
또한 영환은 $k$ 위의 전역 완전 교차라는 관례를 사용한다.
\end{definition}

\noindent
% P01-E0772: the frozen paragraph reuses the polynomial-ring maximal ideal symbol for S-localization without distinguishing its image.
이제 $S$가
정의 \ref{algebra-definition-lci-field}와 같이 전역 완전 교차
$S = k[x_1, \ldots, x_n]/(f_1, \ldots, f_c)$라고 하자.
극대 아이디얼 $\mathfrak m \subset k[x_1, \ldots, x_n]$에 대하여
$\dim(k[x_1, \ldots, x_n]_\mathfrak m) = n$이다
(보조정리 \ref{algebra-lemma-dim-affine-space}).
$(f_1, \ldots, f_c) \subset \mathfrak m$이면
보조정리 \ref{algebra-lemma-one-equation}에 의해
$\dim(S_\mathfrak m) \geq n - c$를 얻는다.
정의 \ref{algebra-definition-lci-field}에 의한 $\dim(S) = n - c$로부터,
$S$의 모든 극대 아이디얼에 대하여 $\dim(S_\mathfrak m) = n - c$이고
$\Spec(S)$가 차원 $n - c$의 등차원임을 결론내린다
(위상수학, 정의 \ref{topology-definition-equidimensional}).
이는 보조정리 \ref{algebra-lemma-dimension-at-a-point-finite-type-over-field}를
참조하면 된다. 앞으로 이를 자주 별도 언급 없이 사용한다.

\begin{lemma}
\label{algebra-lemma-localize-lci}
$k$를 체라 하자.
$S$를 유한형 $k$-대수라 하자.
$g \in S$라 하자.
\begin{enumerate}
\item $S$가 전역 완전 교차이면 $S_g$도 전역 완전 교차이다.
\item $S$가 국소 완전 교차이면 $S_g$도 국소 완전 교차이다.
\end{enumerate}
\end{lemma}

\begin{proof}
두 번째 명제는 첫 번째 명제에서 즉시 따른다.
% P01-E0773: the frozen transition “Proof of the first statement.” is a sentence fragment.
첫 번째 명제를 증명하자. $S_g$가 영환이면 명제는 참이다.
$S_g$가 영환이 아니라고 하자.
정의 \ref{algebra-definition-lci-field}에 따라
$n - c = \dim(S)$를 만족하는
$S = k[x_1, \ldots, x_n]/(f_1, \ldots, f_c)$로 쓰자.
정의 다음의 논의에 의해 $\dim(S_g) = n - c$이다.
$g$에 대응하는 잉여류를 갖는 원소
$g' \in k[x_1, \ldots, x_n]$을 택하자. 그러면
$S_g =
k[x_1, \ldots, x_n, x_{n + 1}]/(f_1, \ldots, f_c, x_{n + 1}g' - 1)$
이므로 원하는 바를 얻는다.
\end{proof}

\begin{lemma}
\label{algebra-lemma-lci-CM}
$k$를 체라 하자. $S$를 유한형 $k$-대수라 하자.
$S$가 국소 완전 교차이면 $S$는 Cohen--Macaulay 환이다.
\end{lemma}

\begin{proof}
% P01-E0774: the frozen proof says “maximal prime” where the local Cohen--Macaulay argument requires a maximal ideal.
$S$의 극대 아이디얼 $\mathfrak m$을 하나 택하자.
$S_\mathfrak m$이 Cohen--Macaulay임을 보여야 한다.
가정에 의해 $\dim(S) = n - c$를 만족하는
$S = k[x_1, \ldots, x_n]/(f_1, \ldots, f_c)$라고 해도 좋다.
$\mathfrak m' \subset k[x_1, \ldots, x_n]$에 대응하는
$\mathfrak m$의 극대 아이디얼이라 하자.
명제 \ref{algebra-proposition-finite-gl-dim-polynomial-ring}에 따르면
국소환 $k[x_1, \ldots, x_n]_{\mathfrak m'}$은
차원 $n$의 정칙 국소환이다. 특히 보조정리
\ref{algebra-lemma-regular-ring-CM}에 의해 Cohen--Macaulay이다.
보조정리 \ref{algebra-lemma-one-equation}을 $c$번 적용하면
국소환
$S_{\mathfrak m} =
k[x_1, \ldots, x_n]_{\mathfrak m'}/(f_1, \ldots, f_c)$의
차원은 $\geq n - c$이다. 가정에 의해
$\dim(S_{\mathfrak m}) \leq n - c$이다.
따라서 등식이 성립한다. 이는
$f_1, \ldots, f_c$가
$k[x_1, \ldots, x_n]_{\mathfrak m'}$에서 정칙열이고
$S_{\mathfrak m}$이 Cohen--Macaulay임을 뜻한다.
이는 명제 \ref{algebra-proposition-CM-module}에 따른다.
\end{proof}

\noindent
다음 보조정리는 이 절의 나머지 내용에서 기술적인 핵심이다.
이 보조정리의 중요한 특징은 환 $S$의 어떤 표시든 선택할 수 있지만
(1)의 조건은 그 선택에 의존하지 않는다는 점이다.

\begin{lemma}
\label{algebra-lemma-lci}
$k$를 체라 하자.
$S$를 유한형 $k$-대수라 하자.
$\mathfrak q$를 $S$의 소 아이디얼이라 하자.
임의의 표시 $S = k[x_1, \ldots, x_n]/I$를 택하자.
$\mathfrak q'$를 $\mathfrak q$에 대응하는
$k[x_1, \ldots, x_n]$의 소 아이디얼이라 하자.
$c = \text{height}(\mathfrak q') - \text{height}(\mathfrak q)$
즉 $\dim_{\mathfrak q}(S) = n - c$이다
(보조정리 \ref{algebra-lemma-codimension} 참조).
다음 조건들은 동치이다.
\begin{enumerate}
\item $g \in S$, $g \not \in \mathfrak q$가 존재하여
$S_g$가 $k$ 위의 전역 완전 교차이다.
\item 아이디얼
$I_{\mathfrak q'} \subset k[x_1, \ldots, x_n]_{\mathfrak q'}$는
$c$개의 원소로 생성될 수 있다.
\item conormal 가군 $(I/I^2)_{\mathfrak q}$는
$S_{\mathfrak q}$ 위에서 $c$개의 원소로 생성될 수 있다.
\item conormal 가군 $(I/I^2)_{\mathfrak q}$는
계수 $c$인 자유 $S_{\mathfrak q}$-가군이다.
\item 아이디얼 $I_{\mathfrak q'}$는
정칙 국소환 $k[x_1, \ldots, x_n]_{\mathfrak q'}$에서
정칙열로 생성될 수 있다.
\end{enumerate}
이 경우 $I_{\mathfrak q'}$의
$I_{\mathfrak q'}/\mathfrak q'I_{\mathfrak q'}$를 생성하는
임의의 $c$개 원소는 국소환
$k[x_1, \ldots, x_n]_{\mathfrak q'}$에서 정칙열을 이룬다.
\end{lemma}

\begin{proof}
$R = k[x_1, \ldots, x_n]_{\mathfrak q'}$로 두자. 이는 차원
$\text{height}(\mathfrak q')$인 Cohen--Macaulay 국소환이다
(예를 들어 보조정리 \ref{algebra-lemma-lci-CM} 참조).
더욱이 $\overline{R} = R/IR = R/I_{\mathfrak q'} = S_{\mathfrak q}$는
차원 $\text{height}(\mathfrak q)$인 몫이다.
$(I/I^2)_{\mathfrak q}$를 생성하는
$f_1, \ldots, f_c \in I_{\mathfrak q'}$인 원소를 택하자.
보조정리 \ref{algebra-lemma-NAK}에 의해
$f_1, \ldots, f_c$는 $I_{\mathfrak q'}$를 생성한다.
차원이 일치하므로 명제 \ref{algebra-proposition-CM-module}에 의해
$f_1, \ldots, f_c$는 $R$의 정칙열이다.
보조정리 \ref{algebra-lemma-regular-quasi-regular}에 의해
$(I/I^2)_{\mathfrak q}$는 자유이다.
이 논의로 (2), (3), (4)가 동치이고,
이들이 보조정리의 마지막 주장을 함의하며 따라서 (5)를 함의한다.

\medskip\noindent
(5)가 성립한다고 하자. $I_{\mathfrak q'}$가 길이 $e$인 정칙열로
생성된다고 하자. 그러면 차원론과 절
\ref{algebra-section-dimension}에 의해
$\text{height}(\mathfrak q) = \dim(S_{\mathfrak q}) =
\dim(k[x_1, \ldots, x_n]_{\mathfrak q'}) - e =
\text{height}(\mathfrak q') - e$
이다. 따라서 $e = c$이다. 그러므로 (5)는 (2)를 함의한다.

\medskip\noindent
첫 단락에서 도입한 표기를 계속 사용하자.
각 $f_i$에 대하여 $d_i \in k[x_1, \ldots, x_n]$,
$d_i \not \in \mathfrak q'$를 택하여
$f_i' = d_i f_i \in k[x_1, \ldots, x_n]$가 되게 할 수 있다.
그러면 여전히 $I_{\mathfrak q'} = (f_1', \ldots, f_c')R$이다.
따라서 $g' \in k[x_1, \ldots, x_n]$, $g' \not \in \mathfrak q'$가
존재하여 $I_{g'} = (f_1', \ldots, f_c')$이다.
또한 $g'' \in k[x_1, \ldots, x_n]$, $g'' \not \in \mathfrak q'$를
택하여 $\dim(S_{g''}) = \dim_{\mathfrak q} \Spec(S)$가 되게 하자.
보조정리 \ref{algebra-lemma-codimension}에 의해 이 차원은 $n - c$이다.
마지막으로 $g$를 $S$에서 $g'g''$의 상으로 두자.
그러면
$$
S_g \cong k[x_1, \ldots, x_n, x_{n + 1}]
/
(f_1', \ldots, f_c', x_{n + 1}g'g'' - 1)
$$
이고 $g''$의 선택에 의해 이 환의 차원은 $n - c$이다.
따라서 이는 전역 완전 교차이다.
그러므로 (2), (3), (4) 각각은 (1)을 함의한다.

\medskip\noindent
(1)을 가정하자. $S_g$를 전역 완전 교차로 표시하는
$S_g \cong k[y_1, \ldots, y_m]/(f_1, \ldots, f_t)$를 택하자.
$J = (f_1, \ldots, f_t)$로 쓰자.
$\mathfrak q'' \subset k[y_1, \ldots, y_m]$를 $\mathfrak qS_g$에
대응하는 소 아이디얼이라 하자. 그러면
$t = m - \dim(S_g) =
\text{height}(\mathfrak q'') - \text{height}(\mathfrak q)$이다.
마지막 등식은 보조정리 \ref{algebra-lemma-codimension}을 참조하라.
보조정리 \ref{algebra-lemma-lci-CM}의 증명(및 위의 논의)에서 본 것처럼
$f_1, \ldots, f_t$는 국소환
$k[y_1, \ldots, y_m]_{\mathfrak q''}$에서 정칙열을 이룬다.
보조정리 \ref{algebra-lemma-regular-quasi-regular}에 의해
$(J/J^2)_{\mathfrak q}$는 계수 $t$인 자유 가군이다.
보조정리 \ref{algebra-lemma-conormal-module-localize}에 의해
% P01-E0775: the frozen stable-isomorphism display uses ordinary powers instead of explicit finite direct sums.
$$
J/J^2 \oplus S_g^n \cong (I/I^2)_g \oplus S_g^m
$$
이다. 따라서 $(I/I^2)_{\mathfrak q}$는
계수
$t + n - m = m - \dim(S_g) + n - m = n - \dim(S_g) =
\text{height}(\mathfrak q') - \text{height}(\mathfrak q) = c$
인 자유 가군이다.
따라서 (4)를 얻는다.
\end{proof}

\noindent
보조정리 \ref{algebra-lemma-lci}의 결과는 다음 정의를 제시한다.

\begin{definition}
\label{algebra-definition-lci-local-ring}
$k$를 체라 하자. $S$를 $k$ 위에서 본질적으로 유한형인
국소 $k$-대수라 하자. 국소 $k$-대수 $R$과 원소
$f_1, \ldots, f_c \in \mathfrak m_R$가 존재하여 다음을 만족하면
$S$를 {\it 완전 교차 (over $k$)}라 한다.
\begin{enumerate}
\item $R$은 $k$ 위에서 본질적으로 유한형이다.
\item $R$은 정칙 국소환이다.
\item $f_1, \ldots, f_c$는 $R$에서 정칙열을 이룬다.
\item $k$-대수로서 $S \cong R/(f_1, \ldots, f_c)$이다.
\end{enumerate}
\end{definition}

\noindent
Cohen 구조 정리(정리 \ref{algebra-theorem-cohen-structure-theorem} 참조)에 의해
모든 완비 뇌터 국소환은 어떤 정칙 완비 국소환의 몫으로 쓸 수 있다.
따라서 위 정의를 사용하여 임의의 완비 뇌터 국소환에 대한
완전 교차 환의 개념을 정의할 수 있다.
이에 대해서는 나눗셈 거듭제곱 대수, 절 \ref{dpa-section-lci}에서
논의할 것이다. 그동안 다음 보조정리가 이러한 정의가 타당함을 보여 준다.

\begin{lemma}
\label{algebra-lemma-ci-well-defined}
$A \to B \to C$를 전사하는 국소환 준동형이라 하자.
$A$와 $B$가 정칙 국소환이라고 가정하자. 다음 조건들은 동치이다.
% P01-E0776: the frozen sentence introducing the equivalence list lacks a colon.
\begin{enumerate}
\item $\Ker(A \to C)$는 정칙열로 생성된다.
\item $\Ker(A \to C)$는 $\dim(A) - \dim(C)$개의 원소로 생성된다.
\item $\Ker(B \to C)$는 정칙열로 생성된다.
\item $\Ker(B \to C)$는 $\dim(B) - \dim(C)$개의 원소로 생성된다.
\end{enumerate}
\end{lemma}

\begin{proof}
정칙 국소환은 Cohen--Macaulay이다
(보조정리 \ref{algebra-lemma-regular-ring-CM} 참조).
따라서 (1) $\Leftrightarrow$ (2) 및
(3) $\Leftrightarrow$ (4)가 성립한다
(명제 \ref{algebra-proposition-CM-module} 참조).
보조정리 \ref{algebra-lemma-regular-quotient-regular}에 의해
아이디얼 $\Ker(A \to B)$는
$\dim(A) - \dim(B)$개의 원소로 생성될 수 있다.
따라서 (4)가 (2)를 함의함을 알 수 있다.

\medskip\noindent
(1)이 (4)를 함의함을 보이는 것이 남았다.
$\dim(A) - \dim(B)$에 대한 귀납법으로 이를 보인다.
$\dim(A) - \dim(B) = 0$인 경우는 자명하다.
$\dim(A) > \dim (B)$라고 가정하자.
$I = \Ker(A \to C)$ 및 $J = \Ker(A \to B)$로 쓰자.
$J \subset I$임에 주의하자. 가정에 의해 $I$의 최소 생성원 수는
$\dim(A) - \dim(C)$이다.
$\mathfrak m \subset A$를 극대 아이디얼이라 하고 다음 사상을 생각하자.
$$
J/ \mathfrak m J \to  I / \mathfrak m I \to \mathfrak m /\mathfrak m^2
$$
보조정리 \ref{algebra-lemma-regular-quotient-regular}와 그 증명에 의해
이 합성은 단사이다.
$J /\mathfrak mJ$에서 0이 아닌 임의의 원소 $x \in J$를 택하자.
위 사실과 Nakayama 보조정리에 의해 $x$는 $I$의 최소 생성계에
들어가는 원소이다.
따라서 $A$를 $A/xA$로, $I$를 $I/xA$로 바꾸어도 좋다.
그러면 $\dim(A)$와 $I$의 최소 생성원 수가 모두 $1$씩 감소한다.
따라서 주장이 증명된다.
\end{proof}

\begin{lemma}
\label{algebra-lemma-lci-local}
$k$를 체라 하자. $S$를 $k$ 위에서 본질적으로 유한형인
국소 $k$-대수라 하자. 다음 조건들은 동치이다.
\begin{enumerate}
\item $S$는 $k$ 위의 완전 교차이다.
\item $R \to S$가 전사이고 $R$이 $k$ 위에서 본질적으로 유한 표시인
정칙 국소환이면, 아이디얼 $\Ker(R \to S)$는 정칙열로 생성될 수 있다.
\item $R \to S$가 전사이고 $R$이 $k$ 위에서 본질적으로 유한 표시인
정칙 국소환인 어떤 경우에, 아이디얼 $\Ker(R \to S)$는
$\dim(R) - \dim(S)$개의 원소로 생성될 수 있다.
\item $k$ 위의 전역 완전 교차 $A$와 $A$의 소 아이디얼
$\mathfrak a$가 존재하여 $S \cong A_{\mathfrak a}$이다.
\item $k$ 위의 국소 완전 교차 $A$와 $A$의 소 아이디얼
$\mathfrak a$가 존재하여 $S \cong A_{\mathfrak a}$이다.
\end{enumerate}
\end{lemma}

\begin{proof}
(2)가 (1)을, (1)이 (3)을 함의함은 명백하다.
(4)가 (5)를 함의함도 명백하다. (3)이 (4)를 함의함을 보이자.
따라서 $R$이 $k$ 위에서 본질적으로 유한 표시인 정칙 국소환이고,
아이디얼 $\Ker(R \to S)$가 $\dim(R) - \dim(S)$개의 원소로
생성되는 전사 $R \to S$가 존재한다고 가정하자.
어떤 $J \subset k[x_1, \ldots, x_n]$과
$J \subset \mathfrak q$인 소 아이디얼
$\mathfrak q \subset k[x_1, \ldots, x_n]$에 대하여
$R = (k[x_1, \ldots, x_n]/J)_{\mathfrak q}$로 쓸 수 있다.
$I \subset k[x_1, \ldots, x_n]$를
$k[x_1, \ldots, x_n] \to S$의 핵이라 하자. 그러면
$S \cong (k[x_1, \ldots, x_n]/I)_{\mathfrak q}$이다.
가정에 의해 $(I/J)_{\mathfrak q}$는
$\dim(R) - \dim(S)$개의 원소로 생성된다.
보조정리 \ref{algebra-lemma-ci-well-defined}에 의해
$I_{\mathfrak q}$는
$\dim(k[x_1, \ldots, x_n]_{\mathfrak q}) - \dim(S)$개의 원소로
생성될 수 있다.
보조정리 \ref{algebra-lemma-lci}에 의해 어떤
$g \in k[x_1, \ldots, x_n]$, $g \not \in \mathfrak q$에 대하여
대수 $(k[x_1, \ldots, x_n]/I)_g$는 전역 완전 교차이고
$S$는 그 국소환과 동형이다.

\medskip\noindent
보조정리의 증명을 끝내려면 (5)가 (2)를 함의함을 보여야 한다.
(5)를 가정하고 $\pi : R \to S$를
$R$이 $k$ 위에서 본질적으로 유한형인 정칙 국소 $k$-대수인
전사라 하자.
가정에 의해 $k$ 위의 어떤 국소 완전 교차 $A$에 대하여
$S = A_{\mathfrak a}$이다.
$J \subset \mathfrak q \subset k[y_1, \ldots, y_m]$을 만족하는
표시 $R = (k[y_1, \ldots, y_m]/J)_{\mathfrak q}$를 택하자.
$J$가 $k[y_1, \ldots, y_m] \to R$의 핵이라고 가정해도 좋다.
$I \subset k[y_1, \ldots, y_m]$를
$k[y_1, \ldots, y_m] \to S = A_{\mathfrak a}$의 핵이라 하자.
그러면 $J \subset I$이고 $(I/J)_{\mathfrak q}$는 전사
$\pi : R \to S$의 핵이다. 따라서
$S = (k[y_1, \ldots, y_m]/I)_{\mathfrak q}$이다.

\medskip\noindent
보조정리 \ref{algebra-lemma-isomorphic-local-rings}에 의해
$g \in A$, $g \not \in \mathfrak a$ 및
$g' \in k[y_1, \ldots, y_m]$, $g' \not \in \mathfrak q$가 존재하여
$A_g \cong (k[y_1, \ldots, y_m]/I)_{g'}$이다.
$A$를 $A_g$로, $k[y_1, \ldots, y_m]$을
$k[y_1, \ldots, y_{m + 1}]$로 바꾼 뒤
$A \cong k[y_1, \ldots, y_m]/I$라고 가정할 수 있다.
다음 국소환의 전사 사상을 생각하자.
$$
k[y_1, \ldots, y_m]_{\mathfrak q} \to R \to S.
$$
$R \to S$의 핵이 정칙열로 생성됨을 보여야 한다.
보조정리 \ref{algebra-lemma-lci}에 의해
$k[y_1, \ldots, y_m]_{\mathfrak q} \to A_{\mathfrak a} = S$의
핵은 정칙열로 생성된다
($A$가 $k$ 위의 국소 완전 교차이기 때문이다).
보조정리 \ref{algebra-lemma-ci-well-defined}에 의해 결론을 얻는다.
\end{proof}

\begin{lemma}
\label{algebra-lemma-lci-at-prime}
$k$를 체, $S$를 유한형 $k$-대수, $\mathfrak q$를 $S$의
소 아이디얼이라 하자. 다음 조건들은 동치이다.
\begin{enumerate}
\item 국소환 $S_{\mathfrak q}$는 완전 교차 환이다
(정의 \ref{algebra-definition-lci-local-ring}).
\item $g \in S$, $g \not \in \mathfrak q$가 존재하여
$S_g$가 $k$ 위의 국소 완전 교차이다.
\item $g \in S$, $g \not \in \mathfrak q$가 존재하여
$S_g$가 $k$ 위의 전역 완전 교차이다.
\item $S = k[x_1, \ldots, x_n]/I$의 임의의 표시에 대하여,
$\mathfrak q' \subset k[x_1, \ldots, x_n]$가 $\mathfrak q$에
대응한다고 하자. 그러면 보조정리 \ref{algebra-lemma-lci}의 동치인
조건 (1)--(5) 중 하나가 성립한다.
% P01-E0779: the frozen plural phrase “any of the equivalent conditions ... hold” has number disagreement and is logically clearer as “any one ... holds”.
\end{enumerate}
\end{lemma}

\begin{proof}
이는 보조정리 \ref{algebra-lemma-lci}, 보조정리 \ref{algebra-lemma-lci-local} 및
정의들을 결합한 것이다.
\end{proof}

\begin{lemma}
\label{algebra-lemma-lci-global}
$k$를 체, $S$를 유한형 $k$-대수라 하자.
다음 조건들은 동치이다.
\begin{enumerate}
\item 환 $S$는 $k$ 위의 국소 완전 교차이다.
\item $S$의 모든 국소환은 $k$ 위의 완전 교차 환이다.
\item 극대 아이디얼에서 취한 $S$의 모든 국소화는
$k$ 위의 완전 교차 환이다.
\end{enumerate}
\end{lemma}

\begin{proof}
이는 보조정리 \ref{algebra-lemma-lci-at-prime}, $\Spec(S)$의 준콤팩트성,
그리고 정의들에서 따른다.
\end{proof}

\noindent
다음 보조정리는 완전 교차라는 성질이 기초체를 바꾸어도
(강한 의미에서) 보존됨을 말한다.

\begin{lemma}
\label{algebra-lemma-lci-field-change-local}
$K/k$를 체 확장이라 하자.
$S$를 $k$ 위의 유한형 대수라 하자.
$\mathfrak q_K$를 $S_K = K \otimes_k S$의 소 아이디얼이라 하고
$\mathfrak q$를 이에 대응하는 $S$의 소 아이디얼이라 하자.
그러면 $S_{\mathfrak q}$가 $k$ 위의 완전 교차라는 것
(정의 \ref{algebra-definition-lci-local-ring})과
$(S_K)_{\mathfrak q_K}$가 $K$ 위의 완전 교차라는 것은 동치이다.
\end{lemma}

\begin{proof}
$S = k[x_1, \ldots, x_n]/I$라는 표시를 택하자.
그러면 $I_K = K \otimes_k I$인 표시
$S_K = K[x_1, \ldots, x_n]/I_K$를 얻는다.
$\mathfrak q_K' \subset K[x_1, \ldots, x_n]$와
$\mathfrak q' \subset k[x_1, \ldots, x_n]$를 각각 대응하는
소 아이디얼이라 하자.
보조정리 \ref{algebra-lemma-lci}의 동치 조건들이
쌍 $(S = k[x_1, \ldots, x_n]/I, \mathfrak q)$에 대해 성립하는 것과
쌍 $(S_K = K[x_1, \ldots, x_n]/I_K, \mathfrak q_K)$에 대해
성립하는 것이 동치임을 보이자.
그러면 보조정리가 따른다
(보조정리 \ref{algebra-lemma-lci-at-prime} 참조).

\medskip\noindent
보조정리 \ref{algebra-lemma-dimension-at-a-point-preserved-field-extension}에 의해
$\dim_{\mathfrak q} S = \dim_{\mathfrak q_K} S_K$이다.
따라서 보조정리 \ref{algebra-lemma-lci}에 나타나는 정수 $c$는
두 쌍 $(S = k[x_1, \ldots, x_n]/I, \mathfrak q)$ 및
$(S_K = K[x_1, \ldots, x_n]/I_K, \mathfrak q_K)$에 대해 같다.
한편 다음 등식이 성립한다.
\begin{eqnarray*}
I \otimes_{k[x_1, \ldots, x_n]} \kappa(\mathfrak q')
\otimes_{\kappa(\mathfrak q')} \kappa(\mathfrak q_K')
& = &
I \otimes_{k[x_1, \ldots, x_n]} \kappa(\mathfrak q_K') \\
& = &
I \otimes_{k[x_1, \ldots, x_n]} K[x_1, \ldots, x_n]
\otimes_{K[x_1, \ldots, x_n]} \kappa(\mathfrak q_K') \\
& = &
(K \otimes_k I) \otimes_{K[x_1, \ldots, x_n]} \kappa(\mathfrak q_K') \\
& = &
I_K \otimes_{K[x_1, \ldots, x_n]} \kappa(\mathfrak q'_K).
\end{eqnarray*}
% P01-E0780: the frozen display alternates between q_K' and q'_K for the same extended prime.
따라서
\[
\dim_{\kappa(\mathfrak q')}
I \otimes_{k[x_1, \ldots, x_n]} \kappa(\mathfrak q')
=
\dim_{\kappa(\mathfrak q'_K)}
I_K \otimes_{K[x_1, \ldots, x_n]} \kappa(\mathfrak q_K')
\]
이다. 따라서 Nakayama 보조정리
\ref{algebra-lemma-NAK}에 의해 $I_{\mathfrak q'}$의 최소 생성원 수와
$(I_K)_{\mathfrak q'_K}$의 최소 생성원 수는 같다.
그러므로 보조정리 \ref{algebra-lemma-lci}의 특징짓기 (2)에서
원하는 결론을 얻는다.
\end{proof}

\begin{lemma}
\label{algebra-lemma-lci-field-change}
$k \to K$를 체 확장이라 하자.
$S$를 유한형 $k$-대수라 하자.
그러면 $S$가 $k$ 위의 국소 완전 교차라는 것과
$S \otimes_k K$가 $K$ 위의 국소 완전 교차라는 것은 동치이다.
\end{lemma}

\begin{proof}
이는 보조정리 \ref{algebra-lemma-lci-global}과
보조정리 \ref{algebra-lemma-lci-field-change-local}을 결합하면 따른다.
그러나 여기서는 (같은 원리에 기초한) 다른 증명도 제시한다.

\medskip\noindent
$S' = S \otimes_k K$로 두자.
핵이 $I$인 표시 $\alpha : k[x_1, \ldots, x_n] \to S$를 택하자.
그러면 핵이 $I'$인 유도된 표시
$\alpha' : K[x_1, \ldots, x_n] \to S'$를 얻는다.

\medskip\noindent
$S$가 국소 완전 교차라고 가정하자.
$\mathfrak q \subset S'$의 소 아이디얼을 하나 택하자.
$\mathfrak q'$를 $K[x_1, \ldots, x_n]$의 대응하는 소 아이디얼,
$\mathfrak p$를 $S$의 대응하는 소 아이디얼,
$\mathfrak p'$를 $k[x_1, \ldots, x_n]$의 대응하는 소 아이디얼이라 하자.
다음 네터 국소환의 도식을 생각하자.
% P01-E0781: the frozen designation sentence omits “by/as” after “Denote”.
$$
\xymatrix{
S'_{\mathfrak q} &  K[x_1, \ldots, x_n]_{\mathfrak q'} \ar[l] \\
S_{\mathfrak p}\ar[u] &  k[x_1, \ldots, x_n]_{\mathfrak p'} \ar[u] \ar[l]
}
$$
% P01-E0782: the frozen diagram lacks terminal punctuation before the next sentence.
보조정리 \ref{algebra-lemma-lci}에 의해
$S_{\mathfrak p}$는
$k[x_1, \ldots, x_n]_{\mathfrak p'}$에서 어떤 정칙열
$f_1, \ldots, f_c$로 절단된다.
오른쪽의 수직 사상은 평탄하므로,
$f_1, \ldots, f_c$의 상들은
$K[x_1, \ldots, x_n]_{\mathfrak q'}$에서 정칙열을 이룬다.
$k$ 위에서 $K$와 텐서하는 함자는 완전하므로
$S'_{\mathfrak q} =
K[x_1, \ldots, x_n]_{\mathfrak q'}/(f_1, \ldots, f_c)$이다.
따라서 다시 보조정리 \ref{algebra-lemma-lci}에 의해 $S'$는
$\mathfrak q$의 근방에서 국소 완전 교차이다.
$\mathfrak q$가 임의였으므로 $S'$는 $K$ 위의 국소 완전 교차이다.

\medskip\noindent
$S'$가 국소 완전 교차라고 가정하자.
$S$의 극대 아이디얼 $\mathfrak m$을 하나 택하자.
$\mathfrak m'$을 $k[x_1, \ldots, x_n]$의 대응하는 극대 아이디얼이라 하자.
잔여체를 $\kappa = \kappa(\mathfrak m)$이라 하자.
% P01-E0783: the frozen residue-field designation omits “by/as”.
주의 \ref{algebra-remark-fundamental-diagram}에 의해
$\mathfrak m$ 위에 있는 $S'$의 소 아이디얼들은
$K \otimes_k \kappa$의 소 아이디얼들에 대응한다.
영점 정리, 정리 \ref{algebra-theorem-nullstellensatz}에 의해
$[\kappa : k] < \infty$이다.
따라서 $K \otimes_k \kappa$는 $K$ 위에서 유한이고 영이 아니다.
그러므로 $K \otimes_k \kappa$는 유한 개의, 그리고 $> 0$개인
소 아이디얼을 가지며, 이들은 모두 극대이고 각각의 잔여체는
$K$ 위에서 유한하다 (절 \ref{algebra-section-artinian} 참조).
% P01-E0784: the frozen prose twice says “a finite number > 0” of primes.
따라서 $\mathfrak m$ 위에 있는 소 아이디얼
$\mathfrak n \subset S'$는 유한 개이고, 그 수는 $> 0$이며,
각각 극대이고 잔여체는 $K$ 위에서 유한하다.
그러한 것 중 하나 $\mathfrak n \subset S'$를 택하고,
$\mathfrak n' \subset K[x_1, \ldots, x_n]$을
$K[x_1, \ldots, x_n]$의 대응하는 소 아이디얼이라 하자.
$V(\mathfrak mS')$가 유한이므로 $\mathfrak n$은 그 고립된 폐점이다.
따라서 $\mathfrak mS'_{\mathfrak n}$은
$S'_{\mathfrak n}$의 정의 아이디얼이다.
예를 들어 보조정리
\ref{algebra-lemma-dimension-base-fibre-equals-total}에 의해
$\dim(S_{\mathfrak m}) = \dim(S'_{\mathfrak n})$을 얻는다.
(이는 보조정리
\ref{algebra-lemma-dimension-at-a-point-preserved-field-extension}을 사용해도
알 수 있다.)
대응하는 네터 국소환의 도식을 생각하자.
% P01-E0785: the frozen diagram lacks terminal punctuation before the next sentence.
$$
\xymatrix{
S'_{\mathfrak n} &  K[x_1, \ldots, x_n]_{\mathfrak n'} \ar[l] \\
S_{\mathfrak m}\ar[u] &  k[x_1, \ldots, x_n]_{\mathfrak m'} \ar[u] \ar[l]
}
$$
보조정리 \ref{algebra-lemma-change-base-NL}에 따르면
$\NL(\alpha) \otimes_S S' = \NL(\alpha')$이고, 특히
$I'/(I')^2 = I/I^2 \otimes_S S'$이다.
따라서
$(I/I^2)_{\mathfrak m} \otimes_{S_{\mathfrak m}} \kappa$와
$(I'/(I')^2)_{\mathfrak n}
\otimes_{S'_{\mathfrak n}} \kappa(\mathfrak n)$의 차원은 같다.
$(I'/(I')^2)_{\mathfrak n}$은
$n - \dim S'_{\mathfrak n}$을 계수로 하는 자유 가군이므로,
$(I/I^2)_{\mathfrak m}$는
$n - \dim S'_{\mathfrak n} = n - \dim S_{\mathfrak m}$개의
원소로 생성될 수 있다.
보조정리 \ref{algebra-lemma-lci}에 의해 $S$는
$\mathfrak m$의 근방에서 국소 완전 교차이다.
$\mathfrak m$이 임의의 극대 아이디얼이었으므로
$S$는 국소 완전 교차이다.
\end{proof}

\noindent
마지막으로, 나중에 다음 사실을 증명할 때 사용할 보조정리를 끝으로
제시한다. 환 사상 $T \to A \to B$가 주어지고 $B$가 $T$ 위에서
syntomic이며 $B$가 $A$ 위에서 syntomic이면, $A$는 $T$ 위에서
syntomic이다.

\begin{lemma}
\label{algebra-lemma-lci-permanence-initial}
국소환의 가환도식
$$
\xymatrix{
B & S \ar[l] \\
A \ar[u] & R \ar[l] \ar[u]
}
$$
을 생각하자.
% P01-E0786: the frozen diagram lacks terminal punctuation before “be a commutative square”.
다음을 가정하자.
\begin{enumerate}
\item $R$과 $\overline{S} = S/\mathfrak m_R S$는 정칙 국소환이다.
\item 어떤 아이디얼 $I$, $J$에 대하여 $A = R/I$ 및 $B = S/J$이다.
\item $J \subset S$이고
$\overline{J} = J/\mathfrak m_R \cap J \subset \overline{S}$는
정칙열로 생성된다.
% P01-E0787: the frozen quotient omits parentheses and likely the extension of m_R to S.
\item $A \to B$와 $R \to S$는 평탄하다.
\end{enumerate}
그러면 $I$는 정칙열로 생성된다.
\end{lemma}

\begin{proof}
$\overline{B} = B/\mathfrak m_RB = B/\mathfrak m_AB$로 두자.
그러면 $\overline{B} = \overline{S}/\overline{J}$이다.
$\overline{J}$를 생성하는 정칙열을 이루도록
$\overline{f}_1, \ldots, \overline{f}_{\overline{c}} \in \overline{J}$가 되는
$f_1, \ldots, f_{\overline{c}} \in J$를 택하자.
보조정리 \ref{algebra-lemma-ci-well-defined}에 의해
$\overline{c} = \dim(\overline{S}) - \dim(\overline{B})$임에 주의하자.
보조정리 \ref{algebra-lemma-grothendieck-regular-sequence}에 의해
환 $S/(f_1, \ldots, f_{\overline{c}})$는 $R$ 위에서 평탄하다.
따라서 $S/(f_1, \ldots, f_{\overline{c}}) + IS$는
$A$ 위에서 평탄하다.
% P01-E0788: the frozen quotient by the sum lacks parentheses in both occurrences.
그러므로
$S/(f_1, \ldots, f_{\overline{c}}) + IS \to B$는
$A$ 위에서 평탄한 유한 $S/IS$-가군의 전사이고
$\mathfrak m_A$를 법으로 취하면 동형이므로,
보조정리 \ref{algebra-lemma-mod-injective}에 의해 동형이다.
즉,
$J = (f_1, \ldots, f_{\overline{c}}) + IS$이다.

\medskip\noindent
다시 보조정리 \ref{algebra-lemma-ci-well-defined}에 의해
아이디얼 $J$는
$c = \dim(S) - \dim(B)$개의 원소로 이루어진 정칙열로 생성된다.
따라서 $J/\mathfrak m_SJ$는 차원 $c$인 벡터 공간이다.
위에서 얻은 $J$의 표현에 의해
$g_1, \ldots, g_{c - \overline{c}} \in I$가 존재하여
$J$가
$f_1, \ldots, f_{\overline{c}}, g_1, \ldots, g_{c - \overline{c}}$로
생성된다 (Nakayama 보조정리 \ref{algebra-lemma-NAK} 사용).
환 $A' = R/(g_1, \ldots, g_{c - \overline{c}})$와
전사 $A' \to A$를 생각하자.
위 논의에 의해
$B = S/(f_1, \ldots, f_{\overline{c}}, g_1, \ldots, g_{c - \overline{c}})$는
$A'$ 위에서 평탄하다
($S/(f_1, \ldots, f_{\overline{c}})$가 $R$ 위에서 평탄하기 때문이다).
따라서 $A' \to B$는 단사이다
(이는 충실 평탄이기 때문이다. 보조정리
\ref{algebra-lemma-local-flat-ff} 참조).
이 사상은 $A$를 거쳐 인수분해되므로 $A' = A$를 얻는다.
보조정리 \ref{algebra-lemma-dimension-base-fibre-equals-total}에 의해
$\dim(B) = \dim(A) + \dim(\overline{B})$ 및
$\dim(S) = \dim(R) + \dim(\overline{S})$임에 주의하자.
따라서 초등적인 대수 계산에 의해
$c - \overline{c} = \dim(R) -\dim(A)$이다.
% P01-E0789: the frozen dimension difference omits spacing before the second dim command.
그러므로 보조정리 \ref{algebra-lemma-ci-well-defined}에 의해
$I = (g_1, \ldots, g_{c - \overline{c}})$는 정칙열로 생성된다.
\end{proof}

\section{Syntomic morphisms}
\label{algebra-section-syntomic}

\noindent
Syntomic 환 사상은 평탄하고 유한 표시인 환 사상이며, 모든 섬유가
국소 완전 교차이다.
일반적인 국소 완전 교차 환 사상은
More on Algebra
 Section \ref{more-algebra-section-lci}에서 일반적인 국소 완전 교차 환 사상을 논의한다.

\begin{definition}
\label{algebra-definition-lci}
환 사상 $R \to S$를 {\it syntomic}이라 부른다. $S$가
$R$ 위의 {\it 평탄 국소 완전 교차}라는 말은 이 사상이 평탄하고
유한 표시이며 모든 섬유환 $S \otimes_R \kappa(\mathfrak p)$가
국소 완전 교차라는 뜻이다(Definition \ref{algebra-definition-lci-field} 참조).
\end{definition}

\noindent
체 위의 대수가 그 체 위에서 syntomic이라는 것은 국소 완전 교차라는
것과 동치이다. 이 정의에는 다음의 좋은 성질이 있다.

\begin{lemma}
\label{algebra-lemma-syntomic-descends}
\begin{slogan}
Syntomic이라는 성질은 기저 위에서 fpqc 국소적이다.
\end{slogan}
환 사상 $R \to S$와 충실 평탄한 환 사상 $R \to R'$를 잡고
$S' = R'\otimes_R S$라 두자. 그러면 $R \to S$가 syntomic인 것은
$R' \to S'$가 syntomic인 것과 동치이다.
\end{lemma}

\begin{proof}
보조정리 \ref{algebra-lemma-finite-presentation-descends}와
\ref{algebra-lemma-flatness-descends}에 의해 평탄성과 유한 표시성에 대해서는
이 주장이 성립한다. 사상 $\Spec(R') \to \Spec(R)$는 전사이다
(보조정리 \ref{algebra-lemma-ff-rings}). 따라서
$\mathfrak p' \subset R'$가 $\mathfrak p \subset R$ 위에 놓일 때
$S \otimes_R \kappa(\mathfrak p)$가 국소 완전 교차인 것과
$S' \otimes_{R'} \kappa(\mathfrak p')$가 국소 완전 교차인 것이
동치임을 보이면 충분하다. 다음에 주의하자.
$S' \otimes_{R'} \kappa(\mathfrak p') =
S \otimes_R \kappa(\mathfrak p)
\otimes_{\kappa(\mathfrak p)} \kappa(\mathfrak p')$.
그러므로 보조정리 \ref{algebra-lemma-lci-field-change}를 적용할 수 있다.
\end{proof}

\begin{lemma}
\label{algebra-lemma-base-change-syntomic}
Syntomic 사상의 임의의 기저변환도 syntomic이다.
\end{lemma}

\begin{proof}
평탄성, 유한 표시성, 섬유의 국소 완전 교차성 각각에 대해
보조정리 \ref{algebra-lemma-flat-base-change},
\ref{algebra-lemma-compose-finite-type}, \ref{algebra-lemma-lci-field-change}가 이를 보인다.
\end{proof}

\begin{lemma}
\label{algebra-lemma-local-syntomic}
환 사상 $R \to S$를 생각하자. 단위원을 생성하는
$g_1, \ldots g_m \in S$가 있고 각 $R \to S_{g_i}$가 syntomic이라고
하자. 그러면 $R \to S$도 syntomic이다.
\end{lemma}

\begin{proof}
평탄성 및 유한 표시성에 대해서는 보조정리
\ref{algebra-lemma-flat-localization}과 \ref{algebra-lemma-cover-upstairs}가 이를 보인다.
섬유환이 국소 완전 교차라는 성질은 정의상 $S$ 위에서 국소적이다
(Definition \ref{algebra-definition-lci-field} 참조).
\end{proof}

\begin{definition}
\label{algebra-definition-relative-global-complete-intersection}
$R \to S$를 환 사상이라 하자. 이 환 사상 $R \to S$가
{\it 상대 전역 완전 교차}라는 것은 어떤 표시
$S = R[x_1, \ldots, x_n]/(f_1, \ldots, f_c)$가 존재하고
$\Spec(S) \to \Spec(R)$의 모든 공집합이 아닌 섬유가 차원 $n-c$를
갖는다는 뜻이다. 이 상황을 나타내기 위해
‘$S = R[x_1, \ldots, x_n]/(f_1, \ldots, f_c)$가 상대 전역 완전
교차라고 하자’라고 말할 것이다.
\end{definition}

\noindent
다음 보조정리는 전역 표시를 찾을 때 때때로 유용하다.

\begin{lemma}
\label{algebra-lemma-huber}
$S$가 유한 표시인 $R$-대수이고
$S = R[x_1, \ldots, x_n]/I$라는 표시에서 $I/I^2$가 $S$ 위에서
자유라고 하자. 그러면 $S$는 다음과 같은 표시를 갖는다.
$S = R[y_1, \ldots, y_m]/(f_1, \ldots, f_c)$라는 표시가 존재하며
$(f_1, \ldots, f_c)/(f_1, \ldots, f_c)^2$는
$f_1, \ldots, f_c$의 류들을 기저로 갖는 자유 가군이다.
\end{lemma}

\begin{proof}
보조정리 \ref{algebra-lemma-finite-presentation-independent}에 의해 $I$는
유한 생성 아이디얼이다. $I/I^2$의 기저로 가는 원소
$f_1, \ldots, f_c \in I$를 택하자. 나카야마 보조정리
(보조정리 \ref{algebra-lemma-NAK})에 의해 어떤 $g \in 1 + I$가 존재하여
$$
g \cdot I \subset (f_1, \ldots, f_c)
$$
이고 $I_g \cong (f_1, \ldots, f_c)_g$이다. 따라서
$$
S \cong R[x_1, \ldots, x_n]/(f_1, \ldots, f_c)[1/g]
\cong R[x_1, \ldots, x_n, x_{n + 1}]/(f_1, \ldots, f_c, gx_{n + 1} - 1)
$$
을 얻는다. 이는 원하는 표시이다. 또한
보조정리 \ref{algebra-lemma-principal-localization-NL}을 적용하면
$f_1, \ldots, f_c,gx_{n + 1} - 1$가
$(f_1, \ldots, f_c, gx_{n + 1} - 1)/(f_1, \ldots, f_c, gx_{n + 1} - 1)^2$
의 기저를 이룬다.
\end{proof}
\begin{example}
\label{algebra-example-factor-polynomials}
정수 $n , m \geq 1$을 택하자. 다음 환 사상을 생각하자.
\begin{eqnarray*}
R = \mathbf{Z}[a_1, \ldots, a_{n + m}]
& \longrightarrow &
S = \mathbf{Z}[b_1, \ldots, b_n, c_1, \ldots, c_m] \\
a_1 & \longmapsto & b_1 + c_1 \\
a_2 & \longmapsto & b_2 + b_1 c_1 + c_2 \\
\ldots & \ldots & \ldots \\
a_{n + m} & \longmapsto & b_n c_m
\end{eqnarray*}
즉 이는 다음 다항식 인수분해가 성립하도록 주어진 다항식환 사이의
유일한 환 사상이다.
$$
x^{n + m} + a_1 x^{n + m - 1} + \ldots + a_{n + m}
=
(x^n + b_1 x^{n - 1} + \ldots + b_n)
(x^m + c_1 x^{m - 1} + \ldots + c_m)
$$
$S$는 $R$ 위에서 $n+m$개의 원소($b_i,c_j$)로 생성되고
$n+m$개의 방정식($a_k=a_k(b_i,c_j)$)을 가진다. 따라서 $S$가
$R$ 위의 상대 전역 완전 교차임을 보이려면 모든 섬유의 차원이 $0$임을
보이면 충분하다.

\medskip\noindent
이를 위해 $R \to k$가 체 $k$로 가는 환 사상이라고 하자.
$a_i$의 상을 $\alpha_i \in k$라 쓰자. 섬유환 $S_k = k \otimes_R S$를
생각하자. 체 확장 $k \to K$에 대해 $k$-대수 사상
$S_k \to K$는 다음을 만족하는
$\beta_1, \ldots, \beta_n, \gamma_1, \ldots, \gamma_m \in K$를
찾는 것과 동치이다.
$$
x^{n + m} + \alpha_1 x^{n + m - 1} + \ldots + \alpha_{n + m}
=
(x^n + \beta_1 x^{n - 1} + \ldots + \beta_n)
(x^m + \gamma_1 x^{m - 1} + \ldots + \gamma_m).
$$
그러한 $n+m$-튜플의 $K$에서의 선택은 많아야 유한 개이다. 따라서 모든 섬유는
유한 개의 폐점을 가진다(예를 들어 힐베르트 영점정리를 사용하면
폐점이 $\overline{k}$에서의 해에 대응함을 알 수 있다). 그러므로
$R \to S$는 상대 전역 완전 교차이다.

\medskip\noindent
다른 방법은
$\mathbf{Z}[a_1, \ldots, a_{n + m}] \to
\mathbf{Z}[b_1, \ldots, b_n, c_1, \ldots, c_m]$가 실제로 유한
환 사상임을 보이는 것이다. 보조정리 \ref{algebra-lemma-polynomials-divide}에
의해 각 $b_i,c_j$는 $R$ 위에서 정수적이고, 따라서
보조정리 \ref{algebra-lemma-characterize-integral}에 의해 $R \to S$는 유한이다.
\end{example}

\begin{example}
\label{algebra-example-roots-universal-polynomial}
다음 환 사상을 생각하자.
\begin{eqnarray*}
R = \mathbf{Z}[a_1, \ldots, a_n]
& \longrightarrow &
S = \mathbf{Z}[\alpha_1, \ldots, \alpha_n] \\
a_1 & \longmapsto &
\alpha_1 + \ldots + \alpha_n \\
\ldots & \ldots & \ldots \\
a_n & \longmapsto &
\alpha_1 \ldots \alpha_n
\end{eqnarray*}
이는 다음 등식이 $\mathbf{Z}[\alpha_i, x]$에서 성립하도록 하는
유일한 다항식환 사상이다.
$$
x^n + a_1 x^{n - 1} + \ldots + a_n
=
\prod\nolimits_{i = 1}^n (x + \alpha_i)
$$
다르게 말하면 $a_i$의 상은
$\alpha_1, \ldots, \alpha_n$의 $i$번째 기본 대칭함수이다. 기본 대칭
다항식의 통상적인 이론(세부사항 생략)에 의해 $S$는 $R$ 위에서 유한
자유이고, 기저는
$\alpha_1^{e_1} \alpha_2^{e_2} \ldots \alpha_n^{e_n}$이며
$0 \leq e_i \leq n - i$이다. 특히 $R \to S$의 섬유환은 유한하므로
차원이 $0$이다. 한편 $S$는 $R$ 위에서 $n$개의 원소로 생성되고
$n$개의 방정식을 만족한다. 따라서 $S$는 $R$ 위의 상대 전역 완전
교차이다. $R$ 위에서 $S$의 계수가 양수이므로 $S$는 $R$ 위에서
충실 평탄이기도 하다.
\end{example}

\begin{lemma}
\label{algebra-lemma-base-change-relative-global-complete-intersection}
$S = R[x_1, \ldots, x_n]/(f_1, \ldots, f_c)$를 상대 전역 완전
교차라고 하자(Definition \ref{algebra-definition-relative-global-complete-intersection}).
\begin{enumerate}
\item 임의의 $R \to R'$에 대해 기저변환
$R' \otimes_R S = R'[x_1, \ldots, x_n]/(f_1, \ldots, f_c)$도 상대 전역
완전 교차이다.
\item 어떤 $h \in R[x_1, \ldots, x_n]$의 상인 임의의 $g \in S$에 대해
$S_g = R[x_1, \ldots, x_n, x_{n + 1}]/(f_1, \ldots, f_c, hx_{n + 1} - 1)$도
상대 전역 완전 교차이다.
\item 환 사상 $R \to S$가 어떤 $f \in R$에 대해
$R \to R_f \to S$로 인수분해된다면
$S = \allowbreak R_f[x_1, \ldots, x_n]/\allowbreak(f_1, \ldots, f_c)$는 $R_f$ 위의 상대
전역 완전 교차이다.
\end{enumerate}
\end{lemma}

\begin{proof}
보조정리 \ref{algebra-lemma-dimension-preserved-field-extension}에 의해
기저변환의 섬유는 원래 사상의 섬유와 같은 차원을 가진다. 또한
$R' \otimes_R R[x_1,\ldots,x_n]/(f_1,\ldots,f_c)
= R'[x_1,\ldots,x_n]/(f_1,\ldots,f_c)$이므로 (1)이 성립한다.
(2)는 한 원소에서의 국소화가
$S_g \cong S[x_{n + 1}]/(gx_{n + 1} - 1)$로 표현된다는 사실이다.
(3)은 (1)에서 따른다. 실제로 가정하에서 $R_f\otimes_R S\cong S$이다.
\end{proof}

\begin{lemma}
\label{algebra-lemma-localize-relative-complete-intersection}
환 $R$을 택하고 $S = R[x_1, \ldots, x_n]/(f_1, \ldots, f_c)$라 하자.
다음 각 경우에 대해 $h \in R[x_1, \ldots, x_n]$와 그 상 $g \in S$를
찾아
$$
S_g = R[x_1, \ldots, x_n, x_{n + 1}]/(f_1, \ldots, f_c, hx_{n + 1} - 1)
$$
가 Definition \ref{algebra-definition-relative-global-complete-intersection}의
표시를 갖는 상대 전역 완전 교차가 되게 하겠다.
\begin{enumerate}
\item $I \subset R$가 아이디얼이고 $\Spec(S/IS)\to\Spec(R/I)$의
섬유 차원이 $n-c$이면, $g$의 상이 $1 \in S/IS$가 되도록 $(h,g)$를
택할 수 있다.
\item $\mathfrak p\subset R$가 소 아이디얼이고
$\dim(S\otimes_R\kappa(\mathfrak p))=n-c$이면, 위와 같은 $(h, g)$를
$g$의 상이 $S\otimes_R\kappa(\mathfrak p)$에서 가역원이 되도록
택할 수 있다.
\item $\mathfrak q\subset S$가 $\mathfrak p\subset R$ 위에 놓인 소
아이디얼이고 $\dim_{\mathfrak q}(S/R)=n-c$이면 위와 같은 $(h, g)$를
$g\not \in\mathfrak q$가 되도록 택할 수 있다.
\end{enumerate}
\end{lemma}

\begin{proof}
Ad (1). 보조정리 \ref{algebra-lemma-dimension-fibres-bounded-open-upstairs}에
의해 $V(IS)$를 포함하며 $W\to\Spec(R)$의 모든 섬유 차원이
$\leq n-c$인 열린집합 $W\subset\Spec(S)$가 존재한다.
$W=\Spec(S)\setminus V(J)$라 하면 $V(J)\cap V(IS)=\emptyset$이므로
$1 \in S/IS$로 가는 $g\in J$를 찾을 수 있다.
$h \in R[x_1, \ldots, x_n]$를 $g$의 임의의 원상으로 택한다.

\medskip\noindent
Ad (2). 보조정리 \ref{algebra-lemma-dimension-fibres-bounded-open-upstairs}에 의해
$\Spec(S\otimes_R\kappa(\mathfrak p))$를 포함하고
$W \to \Spec(R)$의 모든 섬유 차원이 $\leq n-c$인 열린집합
$W\subset\Spec(S)$가 존재한다.
$W=\Spec(S)\setminus V(J)$라 하면
$V(J\cdot S\otimes_R\kappa(\mathfrak p))=\emptyset$이다.
따라서 $S\otimes_R\kappa(\mathfrak p)$에서 가역원이 되는 $g\in J$를
찾을 수 있다(세부사항 생략).
$h \in R[x_1, \ldots, x_n]$를 $g$의 임의의 원상으로 택한다.

\medskip\noindent
Ad (3). 보조정리 \ref{algebra-lemma-dimension-fibres-bounded-open-upstairs}에 의해 $g\in S$, $g\not \in\mathfrak q$이고
$R\to S_g$의 모든 공집합이 아닌 섬유 차원이 $\leq n-c$가 되게 할
수 있다. $h \in R[x_1, \ldots, x_n]$를 $g$로 가는 임의의 원소로 택한다.
\end{proof}

\noindent
다음 보조정리는 상대 전역 완전 교차에 대해 절대 네터 근사를 수행할
수 있음을 말한다.

\begin{lemma}
\label{algebra-lemma-relative-global-complete-intersection-Noetherian}
환 $R$과 상대 전역 완전 교차
$S=R[x_1,\ldots,x_n]/(f_1,\ldots,f_c)$를 택하자
(Definition \ref{algebra-definition-relative-global-complete-intersection}).
어떤 유한형 $\mathbf Z$-부분대수 $R_0\subset R$가 존재하여
$f_i\in R_0[x_1,\ldots,x_n]$이고
$$
S_0 = R_0[x_1, \ldots, x_n]/(f_1, \ldots, f_c)
$$
가 상대 전역 완전 교차가 된다.
\end{lemma}

\begin{proof}
$f_1,\ldots,f_c$의 모든 계수로 생성되는 $R$의 $\mathbf Z$-대수를
$R_0\subset R$라 하고
$S_0=R_0[x_1,\ldots,x_n]/(f_1,\ldots,f_c)$라 하자.
분명히 $S=R\otimes_{R_0}S_0$이다. 소 아이디얼 $\mathfrak q \subset S$를
택하고, 그 아래에 놓이는 소 아이디얼들을 각각
$\mathfrak p \subset R$, $\mathfrak q_0 \subset S_0$,
$\mathfrak p_0 \subset R_0$라 하자.
$\dim(S\otimes_R\kappa(\mathfrak p))=n-c$이므로
보조정리 \ref{algebra-lemma-dimension-preserved-field-extension}에 의해
$\dim(S_0\otimes_{R_0}\kappa(\mathfrak p_0))=n-c$이다.
보조정리 \ref{algebra-lemma-dimension-fibres-bounded-open-upstairs}에 의해
$g\in S_0$, $g\notin\mathfrak q_0$이고 $R_0\to(S_0)_g$의 모든
공집합이 아닌 섬유 차원이 $\leq n-c$인 원소를 찾을 수 있다.
$\mathfrak q$가 임의였고 $\Spec(S)$가 준콤팩트이므로
$g_1, \ldots, g_m \in S_0$를 유한 개 택하여 다음이 성립하게 할 수 있다.
$j = 1, \ldots, m$에 대해 $R_0 \to (S_0)_{g_j}$의 공집합이 아닌
모든 섬유의 차원은 $\leq n-c$이고, $\Spec(S) \to \Spec(S_0)$의 상은
$D(g_1) \cup \ldots \cup D(g_m)$에 포함된다. 즉
$S = R \otimes_{R_0} S_0$에서 $g_1, \ldots, g_m$의 상들이 단위원을
생성한다. $R_0$를 키워 $g_1, \ldots, g_m$가 $S_0$에서도 단위원을
생성한다고 할 수 있다. 따라서 $R_0 \to S_0$의 공집합이 아닌 모든
섬유의 차원은 $\leq n-c$이고 결론이 따른다.
\end{proof}

\begin{lemma}
\label{algebra-lemma-relative-global-complete-intersection-conormal}
환 $R$과 상대 전역 완전 교차
$S=R[x_1,\ldots,x_n]/(f_1,\ldots,f_c)$를 택하자
(Definition \ref{algebra-definition-relative-global-complete-intersection}).
$S$의 각 소 아이디얼 $\mathfrak q$에 대해 대응하는
$R[x_1,\ldots,x_n]$의 소 아이디얼을 $\mathfrak q'$라 하자. 그러면
\begin{enumerate}
\item $f_1,\ldots,f_c$는 국소환
$R[x_1,\ldots,x_n]_{\mathfrak q'}$에서 정칙열이다.
\item 각 환
$R[x_1,\ldots,x_n]_{\mathfrak q'}/(f_1,\ldots,f_i)$는 $R$ 위에서 평탄하다.
\item $S$-가군 $(f_1,\ldots,f_c)/(f_1,\ldots,f_c)^2$는
$f_i\bmod(f_1,\ldots,f_c)^2$를 기저로 하는 자유 가군이다.
\end{enumerate}
\end{lemma}

\begin{proof}
$R$이 네터라고 하자. $\mathfrak p = R \cap \mathfrak q'$로 두자.
예를 들어 보조정리 \ref{algebra-lemma-lci}에 의해
$f_1,\ldots,f_c$는
$R[x_1,\ldots,x_n]_{\mathfrak q'}\otimes_R\kappa(\mathfrak p)$에서
정칙열이다. 또한 국소환 $R[x_1,\ldots,x_n]_{\mathfrak q'}$는
$R_{\mathfrak p}$ 위에서 평탄하다. $R$과 따라서
$R[x_1,\ldots,x_n]_{\mathfrak q'}$도 네터이므로 보조정리
\ref{algebra-lemma-grothendieck-regular-sequence}에 의해 (1),(2)가 성립한다.

\medskip\noindent
$R$이 일반적인 경우에는
$R=\colim_{\lambda\in\Lambda}R_\lambda$를 유한형 $\mathbf Z$-부분대수들의
필터링된 귀납적 극한으로 쓴다(Section \ref{algebra-section-colimits-flat} 참조).
모든 $\lambda$에 대해 $f_1,\ldots,f_c\in R_\lambda[x_1,\ldots,x_n]$라고
할 수 있다. 보조정리
\ref{algebra-lemma-relative-global-complete-intersection-Noetherian}의
$R_0 \subset R$를 택하면 모든 $\lambda$에 대해
$R_0 \subset R_\lambda$로 둘 수 있다. 따라서
$S_\lambda = R_\lambda[x_1, \ldots, x_n]/(f_1, \ldots, f_c)$도
$R_0 \to R_\lambda$에 의한 $S_0$의 기저변환이므로 상대 전역 완전
교차이다(보조정리
\ref{algebra-lemma-base-change-relative-global-complete-intersection}).
$\mathfrak p_\lambda$, $\mathfrak q_\lambda$, $\mathfrak q'_\lambda$를
각각 $\mathfrak p$, $\mathfrak q$, $\mathfrak q'$로부터 유도된
$R_\lambda$, $S_\lambda$, $R_\lambda[x_1, \ldots, x_n]$의 소 아이디얼이라 하자. 그러면
각 $\lambda$에서 (1),(2)가 성립한다. 그리고
$$
R[x_1, \ldots, x_n]_{\mathfrak q'}/(f_1, \ldots, f_i)
=
\colim R_\lambda[x_1, \ldots, x_n]_{\mathfrak q_\lambda'}/(f_1, \ldots, f_i)
$$
이므로 보조정리 \ref{algebra-lemma-colimit-rings-flat}에 의해 $R$ 위에서
(2)의 평탄성을 얻는다.
또한
\begin{align*}
R[x_1, \ldots, x_n]_{\mathfrak q'}/(f_1, \ldots, f_i)
\xrightarrow{f_{i + 1}}
R[x_1, \ldots, x_n]_{\mathfrak q'}/(f_1, \ldots, f_i) \\
=
\colim
\left(
R_\lambda[x_1, \ldots, x_n]_{\mathfrak q_\lambda'}/(f_1, \ldots, f_i)
\xrightarrow{f_{i + 1}}
R_\lambda[x_1, \ldots, x_n]_{\mathfrak q_\lambda'}/(f_1, \ldots, f_i)
\right)
\end{align*}
이고 필터링된 귀납적 극한이 완전하므로
(보조정리 \ref{algebra-lemma-directed-colimit-exact}) (1)도 얻는다.

\medskip\noindent
(3)을 증명하자. $N$을
$S$-가군 $(f_1,\ldots,f_c)/(f_1,\ldots,f_c)^2$라 하고
$e_i\in N$을 $f_i$의 상이라 하자. 보조정리
\ref{algebra-lemma-regular-quasi-regular}와 (1)에 의해 모든 $S$의 소 아이디얼
$\mathfrak q$에서 $e_1,\ldots,e_c$는 $N_\mathfrak q$의 기저이다.
보조정리 \ref{algebra-lemma-characterize-zero-local}에 의해 (3)이 성립한다.
\end{proof}
\begin{lemma}
\label{algebra-lemma-relative-global-complete-intersection}
상대 전역 완전교차는 syntomic이다. 즉, 평탄하다.
\end{lemma}

\begin{proof}
$R \to S$가 상대 전역 완전교차라고 하자.
그 섬유들은 전역 완전교차이고 $S$는 $R$ 위에서 유한 표시이다.
따라서 보일 것은 $R \to S$가 평탄하다는 것뿐이다.
이는 보조정리 \ref{algebra-lemma-relative-global-complete-intersection-conormal}의
(2)에 의해 성립한다.
\end{proof}

\begin{lemma}
\label{algebra-lemma-adjoin-roots}
$A$가 환이고
$P(x) = x^n + b_1 x^{n-1} + \ldots + b_n \in A[x]$가
$A$ 위의 최고차항 계수가 1인 다항식이라고 하자. 그러면 어떤
syntomic 유한 자유 충실 평탄 환 확장
$A \subset A'$가 존재하여
$P(x) = \prod_{i = 1, \ldots, n} (x - \beta_i)$가 된다.
여기서 $\beta_i \in A'$이다.
\end{lemma}

\begin{proof}
Example \ref{algebra-example-roots-universal-polynomial}의 $R$과 $S$를 사용하여
$A' = A \otimes_R S$로 놓자. 이때 $R \to A$는 $a_i$를 $b_i$로
보내며, $\beta_i = -1 \otimes \alpha_i$로 둔다.
보조정리 \ref{algebra-lemma-relative-global-complete-intersection}에 의해
$R \to S$는 충실 평탄이고 유한 자유이며 syntomic임을 관찰하자.
이 성질들은 기저변환 $A \to A'$에도 물려온다(세부 사항은 생략한다).
\end{proof}

\begin{lemma}
\label{algebra-lemma-syntomic}
$R \to S$를 환 사상이라 하자.
$\mathfrak q \subset S$를 $R$의 소 아이디얼 $\mathfrak p$ 위에 놓인
소 아이디얼이라 하자.
다음 조건들은 서로 동치이다.
\begin{enumerate}
\item $g \in S$, $g \not\in \mathfrak q$인 원소가 존재하여
$R \to S_g$가 syntomic이다.
\item $g \in S$, $g \not\in \mathfrak q$인 원소가 존재하여
$S_g$가 $R$ 위의 상대 전역 완전교차이다.
\item $g \in S$, $g \not\in \mathfrak q$인 원소가 존재하여
$R \to S_g$가 유한 표시이고, 국소 환 사상
$R_{\mathfrak p} \to S_{\mathfrak q}$가 평탄하며,
국소 환 $S_{\mathfrak q}/\mathfrak pS_{\mathfrak q}$가
$\kappa(\mathfrak p)$ 위의 완전교차 환이다
(정의 \ref{algebra-definition-lci-local-ring} 참조).
\end{enumerate}
\end{lemma}

\begin{proof}
함의 (1) $\Rightarrow$ (3)은 보조정리 \ref{algebra-lemma-lci-at-prime}이다.
함의 (2) $\Rightarrow$ (1)은
보조정리 \ref{algebra-lemma-relative-global-complete-intersection}이다.
이제 (3)이 (2)를 함의함을 보이면 된다.

\medskip\noindent
(3)을 가정하자. 어떤 $g \in S$, $g\not\in\mathfrak q$에 대해
$S$를 $S_g$로 바꾸면 $S$가 $R$ 위에서 유한 표시라고 가정할 수 있다.
표시 $S = R[x_1, \ldots, x_n]/I$를 택하자. $\mathfrak q$에 대응하는
소 아이디얼을 $\mathfrak q' \subset R[x_1, \ldots, x_n]$라 하자.
$\kappa(\mathfrak p)=k$로 쓰자.
그러면 $S \otimes_R k = k[x_1, \ldots, x_n]/\overline{I}$이고,
$\overline{I} \subset k[x_1, \ldots, x_n]$는 $I$의 상이 생성하는
아이디얼이다. $\mathfrak q'$의 상이 생성하는 소 아이디얼을
$\overline{\mathfrak q}' \subset k[x_1, \ldots, x_n]$라 하자.
보조정리 \ref{algebra-lemma-lci-at-prime}와 \ref{algebra-lemma-lci}에 의해
$\overline{I}$와 $\overline{\mathfrak q}'$에 대해 동치 조건들이 성립한다.
$\overline{I}_{\overline{\mathfrak q}'}/
\overline{\mathfrak q}'\overline{I}_{\overline{\mathfrak q}'}$의
$\kappa(\overline{\mathfrak q}')$ 위 차원을 $c$라고 하자.
이 벡터 공간의 기저로 가는 $f_1, \ldots, f_c \in I$를 택한다.
보조정리 \ref{algebra-lemma-lci}에 의해 상 $\overline{f}_j \in \overline{I}$는
$\overline{I}_{\overline{\mathfrak q}'}$를 생성한다.
$S' = R[x_1, \ldots, x_n]/(f_1, \ldots, f_c)$로 놓고,
전사 사상 $S' \to S$의 핵을 $J$라 하자. $S$가 유한 표시이므로
$J$는 유한 생성 아이디얼이다
(보조정리 \ref{algebra-lemma-compose-finite-type}). 다음 짧은 완전열을 보자.
$$
0 \to J \to S' \to S \to 0
$$
 $S_\mathfrak q$가 $R$ 위에서 평탄하므로
$J_{\mathfrak q'} \otimes_R k \to S'_{\mathfrak q'} \otimes_R k$는
단사이다(보조정리 \ref{algebra-lemma-flat-tor-zero}).
한편 구성에 의해 $S'_{\mathfrak q'} \otimes_R k$는
$S_\mathfrak q \otimes_R k$로 동형사상된다. 따라서
$J_{\mathfrak q'} \otimes_R k =
J_{\mathfrak q'}/\mathfrak pJ_{\mathfrak q'} = 0$이다.
Nakayama 보조정리(보조정리 \ref{algebra-lemma-NAK})에 의해
$g \in R[x_1, \ldots, x_n]$, $g \not \in \mathfrak q'$인 원소가 존재하여
$J_g=0$이다. 즉 $S'_g \cong S_g$이다. 더 국소화하면
보조정리 \ref{algebra-lemma-localize-relative-complete-intersection}에 의해
$S'$ (따라서 $S$)가 상대 전역 완전교차가 되므로 원하는 결론을 얻는다.
\end{proof}

\begin{lemma}
\label{algebra-lemma-syntomic-presentation-ideal-mod-squares}
$R$를 환이라 하자. $S = R[x_1, \ldots, x_n]/I$이고 $I$가
유한 생성 아이디얼이라고 하자. $g \in S$에 대해 $S_g$가 $R$ 위에서
syntomic이면 $(I/I^2)_g$는 유한 사영 $S_g$-가군이다.
\end{lemma}

\begin{proof}
보조정리 \ref{algebra-lemma-syntomic}에 의해 $S_g$에서 단위원을 생성하는
유한한 원소 $g_1, \ldots, g_m \in S$를 택할 수 있으며, 각
$S_{gg_j}$는 $R$ 위의 상대 전역 완전교차이다. $(I/I^2)_{gg_j}$가
유한 사영임을 보이는 것으로 충분하므로(보조정리
\ref{algebra-lemma-finite-projective} 참조), $S_g$ 자체가 상대 전역 완전교차라고
가정해도 된다. 이 경우 결론은 보조정리
\ref{algebra-lemma-conormal-module-localize}와
\ref{algebra-lemma-relative-global-complete-intersection-conormal}에서 따른다.
\end{proof}

\begin{lemma}
\label{algebra-lemma-composition-syntomic}
$R \to S$, $S \to S'$를 환 사상이라 하자.
\begin{enumerate}
\item $R \to S$와 $S \to S'$가 syntomic이면 $R \to S'$도 syntomic이다.
\item $R \to S$와 $S \to S'$가 상대 전역 완전교차이면
$R \to S'$도 상대 전역 완전교차이다.
\end{enumerate}
\end{lemma}

\begin{proof}
먼저 (2)를 증명하자. $R \to S$와 $S \to S'$가 상대 전역 완전교차이고
다음 표시를 갖는다고 하자.
$S =  R[x_1, \ldots, x_n]/(f_1, \ldots, f_c)$이고
$S' = S[y_1, \ldots, y_m]/(h_1, \ldots, h_d)$는 정의
\ref{algebra-definition-relative-global-complete-intersection}에 따른 것이다. 그러면
$$
S' \cong
R[x_1, \ldots, x_n, y_1, \ldots, y_m]/(f_1, \ldots, f_c, h'_1, \ldots, h'_d)
$$
여기서 $h_j$의 올림 $h_j' \in R[x_1, \ldots, x_n, y_1, \ldots, y_m]$를
사용했다. 따라서 섬유 환의 차원을 위에서 제한하면 충분하다.
그러므로 $R=k$가 체라고 가정해도 된다.
이때 $S$는 $k$ 위 유한형이고 차원 $n-c$의 등차원 환이며,
$S \to S'$는 모든 비어 있지 않은 섬유 환이 차원 $m-d$의 등차원인
유한형 환 사상이다. 예를 들어 $S'$의 극대 아이디얼에서 국소화하고
보조정리 \ref{algebra-lemma-dimension-base-fibre-total}를 적용하면
$\dim(S') \leq n-c+m-d$를 얻는다.

\medskip\noindent
이제 (1)을 (2)로 환원하자. $R \to S$와 $S \to S'$가 syntomic이라고
하자. $\mathfrak q \subset S$ 위의 소 아이디얼 $\mathfrak q' \subset S$를
택한다. 보조정리 \ref{algebra-lemma-syntomic}에 의해 $g' \in S'$,
$g' \notin\mathfrak q'$인 원소가 존재하여 $S \to S'_{g'}$가 상대 전역
완전교차이다. 마찬가지로 $g \in S$, $g\notin\mathfrak q$를 택하여
$R \to S_g$가 상대 전역 완전교차가 되게 한다.
보조정리 \ref{algebra-lemma-base-change-relative-global-complete-intersection}에
의해 $S_g \to S_{gg'}$도 상대 전역 완전교차이다.
(2)에 의해 $R \to S_{gg'}$가 상대 전역 완전교차이고
$gg'\notin\mathfrak q'$임을 알 수 있다. $\mathfrak q'$가 임의였으므로
보조정리 \ref{algebra-lemma-syntomic}와 \ref{algebra-lemma-local-syntomic}을 결합하면
$R \to S'$가 syntomic이다(여기서 $S'$의 스펙트럼이 준콤팩트라는 사실과
보조정리 \ref{algebra-lemma-quasi-compact}도 사용한다).
\end{proof}

\noindent
다음 보조정리는 뒤에서 개선된다. Smoothing Ring Maps의 명제
\ref{smoothing-proposition-lift-smooth}를 보라.

\begin{lemma}
\label{algebra-lemma-lift-syntomic}
$R$를 환, $I \subset R$를 아이디얼이라 하자.
$R/I \to \overline{S}$가 syntomic 사상이라고 하자.
그러면 $\overline{S}$의 단위원을 생성하는 원소
$\overline{g}_i \in \overline{S}$들이 존재하여, 각
$\overline{S}_{g_i} \cong S_i/IS_i$가 $R$ 위의 어떤 상대 전역
완전교차 $S_i$에 대해 성립한다.
\end{lemma}

\begin{proof}
보조정리 \ref{algebra-lemma-syntomic}에 의해 $\overline{S}$의 단위원을
생성하는 원소들의 모음 $\overline{g}_i \in \overline{S}$를 택할 수 있으며,
각 $\overline{S}_{g_i}$는 $R/I$ 위의 상대 전역 완전교차이다.
따라서 $\overline{S}$ 자체가 상대 전역 완전교차라고 가정할 수 있다.
다음과 같이 쓰자.
$\overline{S} =
(R/I)[x_1, \ldots, x_n]/(\overline{f}_1, \ldots, \overline{f}_c)$
정의 \ref{algebra-definition-relative-global-complete-intersection}에 따른 것이다.
$\overline{f}_1, \ldots, \overline{f}_c$를 올리는
$f_1, \ldots, f_c \in R[x_1, \ldots, x_n]$를 택하자.
$S = R[x_1, \ldots, x_n]/(f_1, \ldots, f_c)$로 놓으면
$S/IS \cong \overline{S}$이다.
보조정리 \ref{algebra-lemma-localize-relative-complete-intersection}에 의해
$\overline{S}$에서 $1$로 사상되는 $g \in S$를 택하여
$S_g$가 $R$ 위의 상대 전역 완전교차가 되게 할 수 있다.
$\overline{S} \cong S_g/IS_g$이므로 증명이 끝난다.
\end{proof}

\section{매끄러운 환 사상}
\label{algebra-section-smooth}

\noindent
매끄러운 환 사상의 정의를 예로 동기 부여하자.
$R$가 환이고 $S = R[x, y]/(f)$이며 $f \in R[x, y]$가 영이 아니라고 하자.
이 경우 다음 완전열이 있다.
$$
S \to
S\text{d}x \oplus S\text{d}y \to
\Omega_{S/R} \to 0
$$
첫 번째 사상은 $1$을
$\frac{\partial f}{\partial x} \text{d}x +
\frac{\partial f}{\partial y} \text{d}y$로 보내며
절 \ref{algebra-section-netherlander}를 참조하라.
$f$의 편미분들이 $S$에서 단위원을 생성하면
$\Omega_{S/R}$는 계수 $1$의 국소 자유 가군이다.
이때 $S$는 $R$ 위 상대 차원 $1$의 매끄러운 대수이다.
그러나 $f$의 두 편미분이 모두 영이면 $\Omega_{S/R}$가 계수 $2$의 국소 자유
가군일 수도 있다. 예를 들어 소수 $p$에 대해
$R$에서 $p=0$이고 $f=x^p+y^p$이면 그렇다.
이 경우 $R \to S$는 상대 차원 $1$의 상대 전역 완전교차이지만 매끄럽지 않다.
따라서 미분 가군이 자유인지 확인하는 것만으로는 매끄러움을 판정할 수 없다.
올바른 조건은 다음과 같다.

\begin{definition}
\label{algebra-definition-smooth}
$R \to S$가 {\it 매끄럽다}는 것은 유한 표시이고,
순진 코탄젠트 복합체 $\NL_{S/R}$가 $0$차에 놓인 유한 사영 $S$-가군과
준동형이라는 뜻이다. 즉 $H_1(\NL_{S/R}) = 0$이고
$\Omega_{S/R}$가 유한 사영 $S$-가군이다.
\end{definition}

\noindent
$R \to S$를 환 사상이라 하자. 보조정리 \ref{algebra-lemma-NL-homotopy}에 의해
임의의 표시 $\alpha : P \to S$에 대해
$H_1(\NL(\alpha)) = H_1(\NL_{S/R})$이고
$H_0(\NL(\alpha)) = \Omega_{S/R}$이다.
따라서 $R \to S$가 매끄러우면, 핵이 $I$인 임의의 전사
$\alpha : R[x_1, \ldots, x_n] \to S$에 대해 다음 사상은 단사이고
그 여핵은 유한 사영 $S$-가군이다.
$$
I/I^2
\longrightarrow
\Omega_{R[x_1, \ldots, x_n]/R} \otimes_{R[x_1, \ldots, x_n]} S
$$
따라서 이 화살표는 분할 단사이다. 즉 $S$-가군으로서
$\bigoplus_{i = 1}^n S \text{d}x_i \cong I/I^2 \oplus \Omega_{S/R}$이고
$I/I^2$도 유한 사영 $S$-가군이다.
역으로 $R \to S$가 유한 표시이고 핵이 $I$인 전사
$\alpha : R[x_1, \ldots, x_n] \to S$에 대해 표시된 화살표가 분할
단사이면 $R \to S$는 매끄럽다.

\begin{lemma}
\label{algebra-lemma-localize-smooth}
$R \to S$가 매끄러운 환 사상이라고 하자.
임의의 국소화 $S_g$는 $R$ 위에서 매끄럽다.
$f \in R$가 $S$에서 가역 원소로 사상되면 $R_f \to S$는 매끄럽다.
\end{lemma}

\begin{proof}
보조정리 \ref{algebra-lemma-localize-NL}에 의해 $R$ 위의 $S_g$의 순진 코탄젠트
복합체는 $R$ 위의 $S$의 순진 코탄젠트 복합체의 기저변환이다.
$S/R$의 복합체가 $\Omega_{S/R}$이고 이것이 유한 사영 $S$-가군이므로
그 기저변환도 유한 사영이다. 따라서 $S_g$는 $R$ 위에서 매끄럽다.

\medskip\noindent
두 번째 주장도 보조정리 \ref{algebra-lemma-NL-localize-bottom}을 같은 방식으로
적용하면 따른다.
\end{proof}

\begin{lemma}
\label{algebra-lemma-base-change-smooth}
\begin{slogan}
매끄러움은 기저변환에 대해 보존된다
\end{slogan}
$R \to S$가 매끄러운 환 사상이고 $R \to R'$가 임의의 환 사상이라고 하자.
그러면 기저변환 $R' \to S' = R' \otimes_R S$는 매끄럽다.
\end{lemma}

\begin{proof}
핵이 $I$인 표시 $\alpha : R[x_1, \ldots, x_n] \to S$를 택하자.
유도된 표시를 $\alpha' : R'[x_1, \ldots, x_n] \to R' \otimes_R S$라 하고
$I' = \Ker(\alpha')$라 하자.
$0 \to I \to R[x_1, \ldots, x_n] \to S \to 0$가 완전이므로
$R' \otimes_R I \to R'[x_1, \ldots, x_n] \to R' \otimes_R S \to 0$도 완전이다.
따라서 $R' \otimes_R I \to I'$는 전사이다.
정의 \ref{algebra-definition-smooth}에 의해 다음 짧은 완전열이 있다.
$$
0 \to I/I^2 \to
\Omega_{R[x_1, \ldots, x_n]/R} \otimes_{R[x_1, \ldots, x_n]} S \to
\Omega_{S/R} \to
0
$$
그리고 $S$-가군 $\Omega_{S/R}$는 유한 사영이다.
따라서 $I/I^2$는
$\Omega_{R[x_1, \ldots, x_n]/R} \otimes_{R[x_1, \ldots, x_n]} S$의
직접 가군합이다. 다음 가환 도식을 보자.
$$
\xymatrix{
R' \otimes_R (I/I^2) \ar[r] \ar[d] &
R' \otimes_R (\Omega_{R[x_1, \ldots, x_n]/R} \otimes_{R[x_1, \ldots, x_n]} S)
\ar[d] \\
I'/(I')^2 \ar[r] &
\Omega_{R'[x_1, \ldots, x_n]/R'}
\otimes_{R'[x_1, \ldots, x_n]} (R' \otimes_R S)
}
$$
오른쪽 세로 사상이 동형이므로 왼쪽 세로 사상은 단사이자 전사이다.
따라서 $\NL(\alpha')$는 $0$차에 놓인
$\Omega_{S'/R'} \cong S' \otimes_S \Omega_{S/R}$와 준동형이다.
이 가군은 유한 사영 가군의 기저변환이므로 유한 사영이다.
\end{proof}

\begin{lemma}
\label{algebra-lemma-smooth-over-field}
$k$를 체라 하자. $S$가 매끄러운 $k$-대수이면
$S$는 국소 완전교차이다.
\end{lemma}

\begin{proof}
보조정리 \ref{algebra-lemma-base-change-smooth}와
\ref{algebra-lemma-lci-field-change}에 의해 $k$가 대수적으로 닫힌 경우만 증명하면 된다.
핵이 $I$인 표시 $\alpha : k[x_1, \ldots, x_n] \to S$를 택하자.
$\mathfrak m$을 $S$의 극대 아이디얼, $\mathfrak m' \supset I$를
이에 대응하는 $k[x_1, \ldots, x_n]$의 극대 아이디얼이라 하자.
$\mathfrak q$ 대신 $\mathfrak m$을 사용하여 보조정리
\ref{algebra-lemma-lci}의 조건 (5)를 보이겠다.
$k$가 대수적으로 닫혔으므로 정리 \ref{algebra-theorem-nullstellensatz}에 의해
어떤 $a_i \in k$에 대해
$\mathfrak m' = (x_1 - a_1, \ldots, x_n - a_n)$로 쓸 수 있다.
$k \to S$가 매끄럽다는 가정에 의해 $S$-가군 사상
$\text{d} : I/I^2 \to \bigoplus_{i = 1}^n S \text{d}x_i$는 분할 단사이다.
따라서 대응하는
$I/\mathfrak m' I \to \bigoplus \kappa(\mathfrak m') \text{d}x_i$도 단사이다.
$\dim_{\kappa(\mathfrak m')}(I/\mathfrak m' I) = c$라 하고
$I/\mathfrak m' I$의 $\kappa(\mathfrak m')$-기저로 사상되는
$f_1, \ldots, f_c \in I$를 택하자.
Nakayama 보조정리 \ref{algebra-lemma-NAK}에 의해
$f_1, \ldots, f_c$는 $k[x_1, \ldots, x_n]_{\mathfrak m'}$ 위에서
$I_{\mathfrak m'}$를 생성한다. 다음 가환 도식을 보자.
$$
\xymatrix{
I \ar[r] \ar[d] & I/I^2 \ar[rr] \ar[d] & &
I/\mathfrak m'I \ar[d] \\
\Omega_{k[x_1, \ldots, x_n]/k} \ar[r] &
\bigoplus S\text{d}x_i \ar[rr]^{\text{d}x_i \mapsto x_i - a_i} & &
\mathfrak m'/(\mathfrak m')^2
}
$$
(가환성의 증명은 생략한다.) 가운데 세로 사상은 $\alpha$의 순진 코탄젠트
복합체를 정의한다. 오른쪽 아래 가로 사상은
$\bigoplus \kappa(\mathfrak m') \text{d}x_i \to \mathfrak m'/(\mathfrak m')^2$의
동형을 유도한다. 따라서 $I_{\mathfrak m'}$의 생성자
$f_1, \ldots, f_c$는 $k[x_1, \ldots, x_n]_{\mathfrak m'}$의 원소들로서
$\mathfrak m'/(\mathfrak m')^2$에서 $\kappa(\mathfrak m')$ 위 선형 독립인
동치류들로 사상된다. 그러므로 보조정리
\ref{algebra-lemma-regular-ring-CM}에 의해 $k[x_1, \ldots, x_n]_{\mathfrak m'}$에서
정칙열을 이룬다.
이는 보조정리 \ref{algebra-lemma-lci}의 조건 (5)를 검증한다.
따라서 $g \in S$, $g \not \in \mathfrak m$인 어떤 원소에 대해
$S_g$는 $k$ 위의 전역 완전교차이다. 이것이 $S$의 모든 극대 아이디얼에 대해
성립하므로 $S$는 $k$ 위의 국소 완전교차이다.
\end{proof}

\begin{definition}
\label{algebra-definition-standard-smooth}
$R$를 환이라 하자. $n \geq c \geq 0$인 정수와
$f_1, \ldots, f_c \in R[x_1, \ldots, x_n]$가 주어졌을 때
다음을 생각하자.
$$
R[x_1, \ldots, x_n]/(f_1, \ldots, f_c)
$$
다항식
$$
g =
\det
\left(
\begin{matrix}
\partial f_1/\partial x_1 &
\partial f_2/\partial x_1 &
\ldots &
\partial f_c/\partial x_1 \\
\partial f_1/\partial x_2 &
\partial f_2/\partial x_2 &
\ldots &
\partial f_c/\partial x_2 \\
\ldots & \ldots & \ldots & \ldots \\
\partial f_1/\partial x_c &
\partial f_2/\partial x_c &
\ldots &
\partial f_c/\partial x_c
\end{matrix}
\right)
$$
의 상이 $R[x_1, \ldots, x_n]/(f_1, \ldots, f_c)$에서 가역 원소이면
이를 $R$ 위의 {\it 표준 매끄러운 대수}라고 한다.
$R$-대수 $S$가 {\it 표준 매끄럽다} 또는 환 사상 $R \to S$가
{\it 표준 매끄럽다}는 것은 어떤 $n \geq c \geq 0$과
$f_1, \ldots, f_c \in R[x_1, \ldots, x_n]$가 존재하여
$R[x_1, \ldots, x_n]/(f_1, \ldots, f_c)$가 $R$ 위의 표준 매끄러운
대수이고 $S$가 $R[x_1, \ldots, x_n]/(f_1, \ldots, f_c)$와
$R$-대수로서 동형이라는 뜻이다.
\end{definition}

\begin{lemma}
\label{algebra-lemma-standard-smooth}
$$
S = R[x_1, \ldots, x_n]/(f_1, \ldots, f_c) = R[x_1, \ldots, x_n]/I
$$
가 표준 매끄러운 대수라고 하자. 그러면 다음이 성립한다.
\begin{enumerate}
\item 환 사상 $R \to S$는 매끄럽다.
\item $S$-가군 $\Omega_{S/R}$는
$\text{d}x_{c + 1}, \ldots, \text{d}x_n$ 위에서 자유이다.
\item $S$-가군 $I/I^2$는 $f_1, \ldots, f_c$의 동치류들 위에서 자유이다.
\item 임의의 $g \in S$에 대해 $R \to S_g$는 표준 매끄럽다.
\item 임의의 환 사상 $R \to R'$에 대해 기저변환
$R' \to R'\otimes_R S$는 표준 매끄럽다.
\item $f \in R$가 $S$에서 가역 원소로 사상되면
$R_f \to S$는 표준 매끄럽다.
\item 환 $S$는 $R$ 위의 상대 전역 완전교차이다.
\end{enumerate}
\end{lemma}

\begin{proof}
주어진 표시의 순진 코탄젠트 복합체를 생각하자.
$$
(f_1, \ldots, f_c)/(f_1, \ldots, f_c)^2
\longrightarrow
\bigoplus\nolimits_{i = 1}^n S \text{d}x_i
$$
이 사상을 직합의 처음 $c$개 직합인자로의 사영과 합성하자.
표준 매끄러운 대수의 정의에 따라 $f_i \bmod (f_1, \ldots, f_c)^2$의
동치류들은 $\bigoplus_{i = 1}^c S\text{d}x_i$의 기저로 사상된다.
따라서 $(f_1, \ldots, f_c)/(f_1, \ldots, f_c)^2$는 계수 $c$의 자유
가군이고, 기저는 $f_i \bmod (f_1, \ldots, f_c)^2$이다.
또한 순진 코탄젠트 복합체의 $0$차 호몰로지, 즉 $\Omega_{S/R}$는
$j = 1, \ldots, n - c$에 대한 $\text{d}x_{c + j}$의 상을 기저로 갖는
자유 $S$-가군이다. 특히 $R \to S$가 매끄럽다는 것이 증명된다.

\medskip\noindent
(4)와 (6)의 증명은 생략한다. 아래의 예와
보조정리 \ref{algebra-lemma-base-change-relative-global-complete-intersection}의
증명을 보라.

\medskip\noindent
$\varphi : R \to R'$를 임의의 환 사상이라 하자.
$S' = R'[x_1, \ldots, x_n]/(f_1^\varphi, \ldots, f_c^\varphi)$로 두자.
여기서 $f^\varphi$는 $f \in R[x_1, \ldots, x_n]$의 모든 계수에
$\varphi$를 적용하여 얻는 다항식이다. 그러면
$S' \cong R' \otimes_R S$이다.
또한 $S'/R'$에 대한 정의 \ref{algebra-definition-standard-smooth}의 행렬식은
$g^\varphi$와 같다. 따라서 그 $S'$에서의 상은
$R[x_1, \ldots, x_n] \to S \to S'$를 통한 $g$의 상이고,
가역이다. 이것으로 (5)가 증명된다.

\medskip\noindent
(7)을 보이려면 모든 소 아이디얼 $\mathfrak p \subset R$에 대해
$S \otimes_R \kappa(\mathfrak p)$의 차원이 $n-c$임을 보이면 충분하다.
(5)에 의해 체 $k$ 위의 표준 매끄러운 대수
$k[x_1, \ldots, x_n]/(f_1, \ldots, f_c)$가 차원 $n-c$임을
보이면 충분하다. 보조정리 \ref{algebra-lemma-smooth-over-field}에 의해
$k[x_1, \ldots, x_n]/(f_1, \ldots, f_c)$는 국소 완전교차임을 이미 안다.
그러므로 $I/I^2$가 계수 $c$의 자유 가군이라는 사실과
보조정리 \ref{algebra-lemma-lci}를 사용하면
$k[x_1, \ldots, x_n]/(f_1, \ldots, f_c)$의 차원은 $n-c$이다.
\end{proof}

\begin{example}
\label{algebra-example-make-standard-smooth}
환 $R$을 잡자.
$f_1, \ldots, f_c \in R[x_1, \ldots, x_n]$라 하자.
다음을 두자.
$$
h =
\det
\left(
\begin{matrix}
\partial f_1/\partial x_1 &
\partial f_2/\partial x_1 &
\ldots &
\partial f_c/\partial x_1 \\
\partial f_1/\partial x_2 &
\partial f_2/\partial x_2 &
\ldots &
\partial f_c/\partial x_2 \\
\ldots & \ldots & \ldots & \ldots \\
\partial f_1/\partial x_c &
\partial f_2/\partial x_c &
\ldots &
\partial f_c/\partial x_c
\end{matrix}
\right).
$$
$S = R[x_1, \ldots, x_{n + 1}]/(f_1, \ldots, f_c, x_{n + 1}h - 1)$로 두자.
이는 표준 매끄러운 대수의 한 예이지만, 표시가 잘못되었고
변수들은 다음 순서여야 한다:
$x_1, \ldots, x_c, x_{n + 1}, x_{c + 1}, \ldots, x_n$.
\end{example}

\begin{lemma}
\label{algebra-lemma-compose-standard-smooth}
표준 매끄러운 환 사상들의 합성은 표준 매끄럽다.
\end{lemma}

\begin{proof}
$R \to S$와 $S \to S'$가 표준 매끄럽다고 하자. 다음 표시들을 택한다.
$S = R[x_1, \ldots, x_n]/(f_1, \ldots, f_c)$,
$S' = S[y_1, \ldots, y_m]/(g_1, \ldots, g_d)$.
$g_j$로 사상되는 원소
$g_j' \in R[x_1, \ldots, x_n, y_1, \ldots, y_m]$를 택하자.
그러면
$S' = R[x_1, \ldots, x_n, y_1, \ldots, y_m]/
(f_1, \ldots, f_c, g'_1, \ldots, g'_d)$임을 알 수 있다.
$S'$가 표준 매끄럽다는 것을 보이려면 다음 행렬식이
가역임을 확인하면 충분하다.
$$
\det
\left(
\begin{matrix}
\partial f_1/\partial x_1 &
\ldots &
\partial f_c/\partial x_1 &
\partial g_1/\partial x_1 &
\ldots &
\partial g_d/\partial x_1 \\
\ldots &
\ldots &
\ldots &
\ldots &
\ldots &
\ldots  \\
\partial f_1/\partial x_c &
\ldots &
\partial f_c/\partial x_c &
\partial g_1/\partial x_c &
\ldots &
\partial g_d/\partial x_c \\
0 &
\ldots &
0 &
\partial g_1/\partial y_1 &
\ldots &
\partial g_d/\partial y_1 \\
\ldots &
\ldots &
\ldots &
\ldots &
\ldots &
\ldots  \\
0 &
\ldots &
0 &
\partial g_1/\partial y_d &
\ldots &
\partial g_d/\partial y_d
\end{matrix}
\right)
$$
$S'$에서 가역이어야 한다. 이는 가정에 따라 가역인 두 행렬식의
곱이므로 자명하다.
\end{proof}

\begin{lemma}
\label{algebra-lemma-smooth-syntomic}
$R \to S$를 매끄러운 환 사상이라 하자.
그러면 $\Spec(S)$에는 각 $S_g$가 $R$ 위에서 표준 매끄럽게 되는
표준 열린집합 $D(g)$들에 의한 열린 덮개가 존재한다.
특히 $R \to S$는 신토믹이다.
\end{lemma}

\begin{proof}
핵이 $I = (f_1, \ldots, f_m)$인 표시
$\alpha : R[x_1, \ldots, x_n] \to S$를 택하자. 각 부분집합
$E \subset \{1, \ldots, m\}$에 대해, $f_e, e\in E$의 류들이 유한
사영 $S$-가군 $I/I^2$를 자유롭게 생성하는 열린집합 $U_E$를 생각하자.
이는 보조정리 \ref{algebra-lemma-cokernel-flat}을 보라.
$\Spec(S)$를, 각 표준 열린집합 $D(g)$가 이 열린집합들 중 하나인 $U_E$에
완전히 포함되도록 덮을 수 있다. 그러한 $g$에 대해 $x_{n+1}$을 $1/g$로
보내는 표시
$$
\beta : R[x_1, \ldots, x_n, x_{n + 1}] \longrightarrow S_g
$$
를 보자. $J = \Ker(\beta)$로 두고 보조정리
\ref{algebra-lemma-principal-localization-NL}을 사용하면
$J/J^2 \cong (I/I^2)_g \oplus S_g$가 자유임을 알 수 있다.
이제 $S$를 $S_g$로 바꾸어도 된다. 보조정리 \ref{algebra-lemma-huber}를
사용하면, 핵 $I = (f_1, \ldots, f_c)$를 갖고
$I/I^2$가 $f_1, \ldots, f_c$의 류들 위에서 자유인 표시
$\alpha : R[x_1, \ldots, x_n] \to S$를 가정할 수 있다.

\medskip\noindent
앞 단락 끝에서 얻은 표시 $\alpha$를 사용하여, 이번에는
$\Omega_{R[x_1, \ldots, x_n]/R}$의 기저 원소
$\text{d}x_1, \ldots, \text{d}x_n$에 대해 같은 논증을 거의 반복한다.
즉, 크기가 $c$인 임의의 부분집합 $E \subset \{1, \ldots, n\}$에 대해
$\NL(\alpha)$의 미분을 다음 사영과 합성한 사상이 동형이 되는
$\Spec(S)$의 열린집합 $U_E$를 생각할 수 있다.
$$
S^{\oplus c} \cong I/I^2
\longrightarrow
\Omega_{R[x_1, \ldots, x_n]/R} \otimes_{R[x_1, \ldots, x_n]} S
\longrightarrow
\bigoplus\nolimits_{i \in E} S\text{d}x_i
$$
다시 $\Spec(S)$를 (유한 개의) 표준 열린집합 $D(g)$로 덮되,
각 $D(g)$가 $U_E$ 중 하나에 완전히 포함되도록 할 수 있다.
번호를 다시 매기면 $E = \{1, \ldots, c\}$라고 가정할 수 있다.
$D(g) \subset U_E$인 $g$에 대해 $x_{n+1}$을 $1/g$로 보내는 표시
$$
\beta : R[x_1, \ldots, x_n, x_{n + 1}] \to S_g
$$
를 보자. $J = \Ker(\beta)$로 두면 보조정리
\ref{algebra-lemma-principal-localization-NL}에 의해
$J = (f_1, \ldots, f_c, fx_{n + 1} - 1)$이고 $\alpha(f) = g$임을 얻는다.
또한 합성
$$
J/J^2 \longrightarrow
\Omega_{R[x_1, \ldots, x_{n + 1}]/R} \otimes_{R[x_1, \ldots, x_{n + 1}]} S_g
\longrightarrow
\bigoplus\nolimits_{i = 1}^c S_g\text{d}x_i \oplus S_g \text{d}x_{n + 1}
$$
은 동형이다. 좌표의 순서를
$x_1, \ldots, x_c, x_{n + 1}, x_{c + 1}, \ldots, x_n$으로 바꾸면
$S_g$가 원하는 대로 $R$ 위에서 표준 매끄럽다는 결론을 얻는다.

\medskip\noindent
표준 매끄러운 대수는 신토믹이므로
(보조정리 \ref{algebra-lemma-standard-smooth} 및
\ref{algebra-lemma-relative-global-complete-intersection}), 그리고 $R$ 위에서
신토믹이라는 성질은 $S$ 위에서 국소적이므로
(보조정리 \ref{algebra-lemma-local-syntomic}), 이것으로 증명이 끝난다.
\end{proof}

\begin{definition}
\label{algebra-definition-smooth-at-prime}
$R \to S$를 환 사상이라 하자.
$S$의 소 아이디얼 $\mathfrak q$를 잡자.
$g \in S$, $g \notin \mathfrak q$인 원소가 존재하여
$R \to S_g$가 매끄러우면 $R \to S$가
$\mathfrak q$에서 \textit{매끄럽다}고 한다.
\end{definition}

\noindent
유한 표시인 환 사상에 대해서는 이를 다음과 같이 특징지을 수 있다.

\begin{lemma}
\label{algebra-lemma-smooth-at-point}
$R \to S$가 유한 표시라고 하자. $S$의 소 아이디얼 $\mathfrak q$에 대해
다음 조건들은 서로 동치이다.
\begin{enumerate}
\item $R \to S$는 $\mathfrak q$에서 매끄럽다,
\item $H_1(L_{S/R})_\mathfrak q = 0$이고
$\Omega_{S/R, \mathfrak q}$는 유한 자유 $S_\mathfrak q$-가군이다,
\item $H_1(L_{S/R})_\mathfrak q = 0$이고
$\Omega_{S/R, \mathfrak q}$는 사영 $S_\mathfrak q$-가군이며,
\item $H_1(L_{S/R})_\mathfrak q = 0$이고
$\Omega_{S/R, \mathfrak q}$는 평탄한 $S_\mathfrak q$-가군이다.
\end{enumerate}
\end{lemma}

\begin{proof}
순진 코탄젠트 복합체의 구성은 국소화와 교환한다는 사실을 더 언급하지
않고 사용하겠다. 이는 절 \ref{algebra-section-netherlander}, 특히 보조정리
\ref{algebra-lemma-localize-NL}에 있다.
$\Omega_{S/R}$는 유한 표시 $S$-가군임을 주의하자. 보조정리
\ref{algebra-lemma-differentials-finitely-presented}를 보라. 따라서 (2), (3),
(4)는 보조정리 \ref{algebra-lemma-finite-projective}에 의해 동치이다.
(1)이 동치인 (2), (3), (4)를 함의한다는 것은 명백하다.
(2)를 가정하자. $S_\mathfrak q$를 주 아이디얼 국소화들의 귀납적 극한으로
쓰고 보조정리 \ref{algebra-lemma-colimit-category-fp-modules}를 사용하면,
$g \in S$, $g \notin \mathfrak q$인 원소를 택하여
$(\Omega_{S/R})_g$가 유한 자유가 되게 할 수 있다.
핵이 $I$인 표시 $\alpha : R[x_1, \ldots, x_n] \to S$를 택하자.
보조정리 \ref{algebra-lemma-NL-homotopy}에 따라 $\NL_{S/R}$ 대신
$\NL(\alpha)$를 사용해도 된다. 전사
$$
\Omega_{R[x_1, \ldots, x_n]/R} \otimes_{R[x_1, \ldots, x_n]} S
\to \Omega_{S/R} \to 0
$$
는 $(\Omega_{S/R})_g$가 사영이므로 $g$를 가역화한 뒤 우역을 갖는다.
따라서
$\text{d} : (I/I^2)_g \to
\Omega_{R[x_1, \ldots, x_n]/R} \otimes_{R[x_1, \ldots, x_n]} S_g$
의 상은 직합인자이고 이 사상도 우역을 갖는다.
그러므로 $H_1(L_{S/R})_g$는 $(I/I^2)_g$의 몫임을 얻는다.
특히 $H_1(L_{S/R})_g$는 유한 $S_g$-가군이다.
따라서 $H_1(L_{S/R})_{\mathfrak q}$의 소멸은 어떤
$g' \in S$, $g' \notin \mathfrak q$에 대해
$H_1(L_{S/R})_{gg'}$의 소멸을 함의한다.
정의에 의해 $R \to S_{gg'}$가 매끄럽다.
\end{proof}

\begin{lemma}
\label{algebra-lemma-locally-smooth}
\begin{slogan}
환 사상은 표적의 모든 소 아이디얼에서 매끄러운 경우 그리고 그 경우에만 매끄럽다
\end{slogan}
$R \to S$를 환 사상이라 하자.
$R \to S$가 매끄러운 것은 $S$의 모든 소 아이디얼 $\mathfrak q$에서
$R \to S$가 매끄러운 것과 동치이다.
\end{lemma}

\begin{proof}
직접 함의는 자명하다. $R \to S$가 $S$의 모든 소 아이디얼
$\mathfrak q$에서 매끄럽다고 하자. 보조정리 \ref{algebra-lemma-quasi-compact}에
따르면 $\Spec(S)$가 준콤팩트이므로, 각 $S_{g_i}$가 매끄러운 유한 덮개
$\Spec(S) = \bigcup D(g_i)$가 존재한다.
보조정리 \ref{algebra-lemma-cover-upstairs}에 의해 이는 $S$가 $R$ 위에서
유한 표시임을 뜻한다. 보조정리 \ref{algebra-lemma-localize-NL}에 따라
$\NL_{S/R} \otimes_S S_{g_i}$는 $0$차에 놓인 유한 사영
$S_{g_i}$-가군과 준동형 동치이다.
보조정리 \ref{algebra-lemma-finite-projective}에 의해 이는
$\NL_{S/R}$가 $0$차에 놓인 유한 사영 $S$-가군과 준동형 동치임을 뜻한다.
\end{proof}

\begin{lemma}
\label{algebra-lemma-compose-smooth}
매끄러운 환 사상들의 합성은 매끄럽다.
\end{lemma}

\begin{proof}
이를 증명하는 방법은 여러 가지가 있다. 한 가지 방법은 뱀 보조정리
(보조정리 \ref{algebra-lemma-snake}), 야코비-자리스키 열
(보조정리 \ref{algebra-lemma-exact-sequence-NL}), 그리고 사영 가군이 자유
가군의 직합인자라는 특징(보조정리 \ref{algebra-lemma-characterize-projective})을
함께 사용하는 것이다.
또 다른 증명은 보조정리 \ref{algebra-lemma-smooth-syntomic},
\ref{algebra-lemma-compose-standard-smooth}, \ref{algebra-lemma-locally-smooth}를
결합하여 얻을 수 있다.
\end{proof}

\begin{lemma}
\label{algebra-lemma-product-smooth}
환 $R$을 잡자. $S = S' \times S''$가 $R$-대수들의 곱이라고 하자.
그러면 $S$가 $R$ 위에서 매끄러운 것은 $S'$와 $S''$가 각각 $R$ 위에서
매끄러운 것과 동치이다.
\end{lemma}

\begin{proof}
생략한다. 힌트: 보조정리 \ref{algebra-lemma-locally-smooth}에 의해 매끄러움은
다음과 같이 확인할 수 있다.
한 소 아이디얼씩 확인하면 된다. 보조정리 \ref{algebra-lemma-spec-product}에 의해
$\Spec(S)$는 $\Spec(S')$와 $\Spec(S'')$의 서로소 합집합이므로,
$\mathfrak q$에서 $R \to S$가 매끄럽다는 것은 대응하는 소 아이디얼에서
$R \to S'$가 매끄럽거나 $R \to S''$가 매끄럽다는 것에 대응한다.
\end{proof}

\begin{lemma}
\label{algebra-lemma-relative-global-complete-intersection-smooth}
환 $R$을 잡자. $S = R[x_1, \ldots, x_n]/(f_1, \ldots, f_c)$가
상대 전역 완전교차라고 하자.
$\mathfrak q \subset S$를 소 아이디얼이라 하자. 그러면 $R \to S$가
$\mathfrak q$에서 매끄러운 것은 다음 조건과 동치이다:
크기가 $c$인 부분집합 $I \subset \{1, \ldots, n\}$가 존재하여 다항식
$$
g_I = \det (\partial f_j/\partial x_i)_{j = 1, \ldots, c, \ i \in I}.
$$
이 $\mathfrak q$의 원소로 사상되지 않는다.
\end{lemma}

\begin{proof}
보조정리 \ref{algebra-lemma-relative-global-complete-intersection-conormal}에
의해 $S$의 주어진 표시에 대응하는 순진 코탄젠트 복합체는
$$
\bigoplus\nolimits_{j = 1}^c S \cdot f_j
\longrightarrow
\bigoplus\nolimits_{i = 1}^n S \cdot \text{d}x_i, \quad
f_j \longmapsto \sum \frac{\partial f_j}{\partial x_i} \text{d}x_i.
$$
이다. 이 사상을 나타내는 행렬의 최대 소행렬식들이 정확히
다항식 $g_I$들이다.

\medskip\noindent
$g_I$가 $g \in S$로 사상되고 $g \notin \mathfrak q$라고 하자.
그러면 대수 $S_g$는 $R$ 위에서 매끄럽다. 실제로 그 순진 코탄젠트
복합체는 위 복합체를 $g$에서 국소화한 것과 준동형 동치이며,
이는 보조정리 \ref{algebra-lemma-localize-NL}을 따른다. 그리고 구성에 의해
이는 $0$차에 놓인 계수 $n-c$의 자유 가군과 준동형 동치이다.

\medskip\noindent
반대로 모든 $g_I$가 결국 $\mathfrak q$에 들어간다고 하자.
이 경우 위 복합체를 $\kappa(\mathfrak q)$와 텐서한 것은 최대 계수를
갖지 않으며, 따라서 $g \in S$, $g \notin \mathfrak q$인 원소로
국소화하여 이 사상이 분할 단사가 되게 할 수 없다.
보조정리 \ref{algebra-lemma-localize-NL}에 의해 다시, $R$ 위에서 매끄러운
그러한 국소화도 존재하지 않는다.
\end{proof}

\begin{lemma}
\label{algebra-lemma-flat-fibre-smooth}
$R \to S$를 환 사상이라 하자.
$\mathfrak q \subset S$가 $R$의 소 아이디얼 $\mathfrak p$ 위에 놓인
소 아이디얼이라고 하자. 다음을 가정하자.
\begin{enumerate}
\item $g \in S$, $g \notin \mathfrak q$인 원소가 존재하여
$R \to S_g$가 유한 표시이다,
\item 국소 환 준동형 $R_{\mathfrak p} \to S_{\mathfrak q}$가 평탄하다,
\item 섬유 $S \otimes_R \kappa(\mathfrak p)$가
$\mathfrak q$에 대응하는 소 아이디얼에서
$\kappa(\mathfrak p)$ 위에서 매끄럽다.
\end{enumerate}
그러면 $R \to S$는 $\mathfrak q$에서 매끄럽다.
\end{lemma}

\begin{proof}
보조정리 \ref{algebra-lemma-syntomic}과 \ref{algebra-lemma-smooth-over-field}에 의해
어떤 $g \in S$가 존재하여 $S_g$가 상대 전역 완전교차가 됨을 안다.
$S$를 $S_g$로 바꾸면
$S = R[x_1, \ldots, x_n]/(f_1, \ldots, f_c)$가 상대 전역 완전교차라고
가정할 수 있다.
크기가 $c$인 각 부분집합 $I \subset \{1, \ldots, n\}$에 대해
보조정리 \ref{algebra-lemma-relative-global-complete-intersection-smooth}의
다항식
$g_I = \det (\partial f_j/\partial x_i)_{j = 1, \ldots, c, i \in I}$를
생각하자.
$g_I$의 상 $\overline{g}_I$를 다항식환
$\kappa(\mathfrak p)[x_1, \ldots, x_n]$에서 생각하면,
이는 $\kappa(\mathfrak p)[x_1, \ldots, x_n]$에서 $f_j$의 상
$\overline{f}_j$들의 편미분으로 이루어진 행렬의 행렬식이다.
따라서 보조정리 \ref{algebra-lemma-relative-global-complete-intersection-smooth}를
$R \to S$와
$\kappa(\mathfrak p) \to S \otimes_R \kappa(\mathfrak p)$에 각각
적용하면 결론을 얻는다.
\end{proof}

\noindent
다음 보조정리에서 $U, V$는 정의에 의해 열린집합임에 주의하자.

\begin{lemma}
\label{algebra-lemma-flat-base-change-locus-smooth}
$R \to S$를 유한 표시인 환 사상이라 하자.
$R \to R'$를 평탄한 환 사상이라 하자.
기저변환을 $S' = R' \otimes_R S$로 쓰자.
$U \subset \Spec(S)$를 $R \to S$가 매끄러운 소 아이디얼들의 집합이라 하자.
$V \subset \Spec(S')$를 $R' \to S'$가 매끄러운 소 아이디얼들의 집합이라 하자.
그러면 $V$는 사상 $f : \Spec(S') \to \Spec(S)$에 의한 $U$의
역상이다.
\end{lemma}

\begin{proof}
보조정리 \ref{algebra-lemma-change-base-NL}에 의해
$\NL_{S/R} \otimes_S S'$는 $\NL_{S'/R'}$와 호모토피 동치임을 안다.
이것으로 이미 $f^{-1}(U) \subset V$가 따른다.

\medskip\noindent
$\mathfrak q' \subset S'$를 $\mathfrak q \subset S$ 위에 놓인 소 아이디얼이라 하자.
$\mathfrak q' \in V$라고 가정하자. $\mathfrak q \in U$임을 보여야 한다.
$S \to S'$가 평탄하므로, 보조정리 \ref{algebra-lemma-local-flat-ff}에 의해
$S_{\mathfrak q} \to S'_{\mathfrak q'}$는 충실히 평탄하다.
따라서 $H_1(L_{S'/R'})_{\mathfrak q'}$의 소멸은
$H_1(L_{S/R})_{\mathfrak q}$의 소멸을 함의한다.
보조정리 \ref{algebra-lemma-finite-projective-descends}를
$S_{\mathfrak q}$-가군 $(\Omega_{S/R})_{\mathfrak q}$와 사상
$S_{\mathfrak q} \to S'_{\mathfrak q'}$에 적용하면
$(\Omega_{S/R})_{\mathfrak q}$가 사영임을 얻는다.
따라서 보조정리 \ref{algebra-lemma-smooth-at-point}에 의해
$R \to S$는 $\mathfrak q$에서 매끄럽다.
\end{proof}

\begin{lemma}
\label{algebra-lemma-smooth-field-change-local}
$K/k$를 체의 확대라 하자. $S$를 $k$ 위의 유한 생성 대수라 하자.
$\mathfrak q_K$를 $S_K = K \otimes_k S$의 소 아이디얼이라 하고,
$\mathfrak q$를 $S$의 대응하는 소 아이디얼이라 하자.
그러면 $S$가 $\mathfrak q$에서 $k$ 위에 매끄러운 것은
$S_K$가 $\mathfrak q_K$에서 $K$ 위에 매끄러운 것과 동치이다.
\end{lemma}

\begin{proof}
이는 보조정리 \ref{algebra-lemma-flat-base-change-locus-smooth}의 특수한 경우이다.
\end{proof}

\begin{lemma}
\label{algebra-lemma-lift-smooth}
$R$을 환이라 하고 $I \subset R$를 아이디얼이라 하자.
$R/I \to \overline{S}$가 매끄러운 환 사상이라고 하자.
$\overline{S}$의 단위원을 생성하는 원소들
$\overline{g}_i \in \overline{S}$가 존재하여, 각
$\overline{S}_{g_i} \cong S_i/IS_i$가 $R$ 위의 어떤 (표준) 매끄러운 환
$S_i$에 대해 성립하게 할 수 있다.
\end{lemma}

\begin{proof}
보조정리 \ref{algebra-lemma-smooth-syntomic}에 의해
$\overline{S}$의 단위원을 생성하는 원소들
$\overline{g}_i \in \overline{S}$를 택하여 각 $\overline{S}_{g_i}$가
$R/I$ 위에서 표준 매끄럽게 되도록 할 수 있다.
따라서 $\overline{S}$ 자체가 $R/I$ 위에서 표준 매끄럽다고 가정해도 된다.
정의 \ref{algebra-definition-standard-smooth}에 따라
$$
\overline{S} =
(R/I)[x_1, \ldots, x_n]/(\overline{f}_1, \ldots, \overline{f}_c)
$$
로 쓰자. $\overline{f}_1, \ldots, \overline{f}_c$의 올림인
$f_1, \ldots, f_c \in R[x_1, \ldots, x_n]$을 택하자.
$$
S = R[x_1, \ldots, x_n, x_{n + 1}]/(f_1, \ldots, f_c, x_{n + 1}\Delta - 1)
$$
로 두자. 여기서
$\Delta = \det(\frac{\partial f_j}{\partial x_i})_{i, j = 1, \ldots, c}$는
예 \ref{algebra-example-make-standard-smooth}에서와 같다.
이것으로 보조정리가 증명된다.
\end{proof}

\section{형식적으로 매끄러운 사상}
\label{algebra-section-formally-smooth}

\noindent
이 절에서는 형식적으로 매끄러운 환 사상을 정의한다. 유한 표시인 환 사상이
형식적으로 매끄러운 것은 매끄러운 것과 동치임이 밝혀진다.
보조정리 \ref{algebra-proposition-smooth-formally-smooth}를 보라.

\begin{definition}
\label{algebra-definition-formally-smooth}
$R \to S$를 환 사상이라 하자.
모든 가환 실선 도식
$$
\xymatrix{
S \ar[r] \ar@{-->}[rd] & A/I \\
R \ar[r] \ar[u] & A \ar[u]
}
$$
에서 $I \subset A$가 제곱 영 아이디얼이면 도식을 가환으로 만드는
점선 사상이 존재할 때, $S$가 $R$ 위에서 \textit{형식적으로 매끄럽다}고 한다.
\end{definition}

\begin{lemma}
\label{algebra-lemma-base-change-fs}
$R \to S$를 형식적으로 매끄러운 환 사상이라 하자.
$R \to R'$를 임의의 환 사상이라 하자.
그러면 기저변환 $S' = R' \otimes_R S$는 $R'$ 위에서 형식적으로 매끄럽다.
\end{lemma}

\begin{proof}
정의 \ref{algebra-definition-formally-smooth}와 같은 실선 도식
$$
\xymatrix{
S \ar[r] \ar@{-->}[rrd] & R' \otimes_R S \ar[r] \ar@{-->}[rd] & A/I \\
R  \ar[u] \ar[r] & R' \ar[r] \ar[u] & A \ar[u]
}
$$
을 주자. 가정에 의해 더 긴 점선 사상이 존재한다.
텐서곱의 보편 성질에 의해 더 짧은 점선 사상을 얻는다.
\end{proof}

\begin{lemma}
\label{algebra-lemma-compose-formally-smooth}
형식적으로 매끄러운 환 사상들의 합성은 형식적으로 매끄럽다.
\end{lemma}

\begin{proof}
생략한다. (힌트: 적절한 도식을 생각하면 완전히 형식적으로 따른다.)
\end{proof}

\begin{lemma}
\label{algebra-lemma-polynomial-ring-formally-smooth}
$R$ 위의 다항식환은 $R$ 위에서 형식적으로 매끄럽다.
\end{lemma}

\begin{proof}
$S = R[x_j; j \in J]$이고 정의 \ref{algebra-definition-formally-smooth}와 같은
도식이 주어졌다고 하자. 위쪽 수평 사상 아래에서 $x_j$가 보내지는
$A/I$의 원소들의 올림 $a_j \in A$를 택하면 점선 사상이 존재한다.
\end{proof}

\begin{lemma}
\label{algebra-lemma-characterize-formally-smooth}
$R \to S$를 환 사상이라 하자.
$P \to S$를 다항식환 $P$에서 $S$로 가는 전사 $R$-대수 사상이라 하자.
그 핵을 $J \subset P$로 표시하자. 그러면 $R \to S$가 형식적으로 매끄러운
것은 $R$-대수 사상 $\sigma : S \to P/J^2$가 존재하여
전사 사상 $P/J^2 \to S$의 우역이 되는 것과 동치이다.
\end{lemma}

\begin{proof}
$R \to S$가 형식적으로 매끄럽다고 하자.
다음 가환 도식을 생각하자.
$$
\xymatrix{
S \ar[r] \ar@{-->}[rd] & P/J \\
R \ar[r] \ar[u] &  P/J^2\ar[u]
}
$$
가정에 의해 점선 사상이 존재한다. 이것이 $\sigma$의 존재를 증명한다.

\medskip\noindent
반대로 보조정리의 $\sigma$가 존재한다고 하자.
정의 \ref{algebra-definition-formally-smooth}와 같은 실선 도식
$$
\xymatrix{
S \ar[r] \ar@{-->}[rd] & A/I \\
R \ar[r] \ar[u] & A \ar[u]
}
$$
을 주자. 보조정리 \ref{algebra-lemma-polynomial-ring-formally-smooth}에 의해
$P$는 형식적으로 매끄러우므로, $P \to S \to A/I$를 올리는
$R$-대수 준동형 $\psi : P \to A$가 존재한다.
분명히 $\psi(J) \subset I$이고 $I^2 = 0$이므로 $\psi(J^2) = 0$이다.
따라서 $\psi$는 $\overline{\psi} : P/J^2 \to A$로 인수분해된다.
원하는 점선 사상은 합성
$\overline{\psi} \circ \sigma : S \to A$이다.
\end{proof}

\begin{remark}
\label{algebra-remark-lemma-characterize-formally-smooth}
보조정리 \ref{algebra-lemma-characterize-formally-smooth}는 $P$가 $R$ 위에서
형식적으로 매끄러운 경우에도 더 일반적으로 성립한다.
\end{remark}

\begin{lemma}
\label{algebra-lemma-characterize-formally-smooth-again}
$R \to S$를 환 사상이라 하자.
$P \to S$를 다항식환 $P$에서 $S$로 가는 전사 $R$-대수 사상이라 하자.
그 핵을 $J \subset P$로 표시하자. 그러면 $R \to S$가 형식적으로 매끄러운
것은 보조정리 \ref{algebra-lemma-differential-seq}의 다음 열이 분할 완전열인 것과 동치이다.
$$
0 \to J/J^2 \to \Omega_{P/R} \otimes_P S \to \Omega_{S/R} \to 0
$$
\end{lemma}

\begin{proof}
$S$가 $R$ 위에서 형식적으로 매끄럽다고 하자. 보조정리
\ref{algebra-lemma-characterize-formally-smooth}에 의해 이는 $S \to P/J^2$인
$R$-대수 사상이 $P/J^2 \to S$의 우역이라는 뜻이다.
보조정리 \ref{algebra-lemma-differential-mod-power-ideal}에 의해
$\Omega_{P/R} \otimes_P S = \Omega_{(P/J^2)/R} \otimes_{P/J^2} S$이다.
보조정리 \ref{algebra-lemma-differential-seq-split}에 의해 열은 분할된다.

\medskip\noindent
이 보조정리의 완전열이 분할 완전이라고 가정하자.
분할 $\sigma : \Omega_{S/R} \to \Omega_{P/R} \otimes_P S$를 택하자.
각 $\lambda \in S$에 대해 $\lambda$로 사상되는 $x_\lambda \in P$를 택하자.
다음으로 각 $\lambda \in S$에 대해
$$
\text{d}f_\lambda = \text{d}x_\lambda - \sigma(\text{d}\lambda)
$$
가 완전열의 중간 항에서 성립하도록 $f_\lambda \in J$를 택하자.
$s : \lambda \mapsto x_\lambda - f_\lambda \mod J^2$가
$R$-대수 준동형 $s : S \to P/J^2$임을 보이자.
이를 위해 $h \in J$이고 $\text{d}h = 0$가
$\Omega_{P/R} \otimes_R S$에서 성립하면 $h \in J^2$라는 사실을 반복해서 사용한다.
$\lambda, \mu \in S$라 하자. 그러면
$\sigma(\text{d}\lambda + \text{d}\mu - \text{d}(\lambda + \mu)) = 0$이다.
이는
$$
\text{d}(x_\lambda + x_\mu - x_{\lambda + \mu}
- f_\lambda - f_\mu + f_{\lambda + \mu}) = 0
$$
을 함의한다. 따라서
$x_\lambda + x_\mu - x_{\lambda + \mu}
- f_\lambda - f_\mu + f_{\lambda + \mu} \in J^2$이고,
결국 $s(\lambda) + s(\mu) = s(\lambda + \mu)$이다.
비슷하게
$\sigma(\lambda \text{d}\mu + \mu \text{d}\lambda - \text{d}\lambda \mu) = 0$이고, 이는
$$
\mu(\text{d}x_\lambda - \text{d}f_\lambda) +
\lambda(\text{d}x_\mu - \text{d}f_\mu) -
\text{d}x_{\lambda\mu} + \text{d}f_{\lambda\mu} = 0
$$
을 완전열의 중간 항에서 함의한다.
또한
$$
\text{d}(x_\lambda x_\mu) =
x_\lambda \text{d}x_\mu + x_\mu \text{d}x_\lambda =
\lambda \text{d}x_\mu + \mu \text{d} x_\lambda
$$
가 다시 중간 항에서 성립한다.
이 식들을 합치면
$x_\lambda x_\mu - x_{\lambda\mu}
- \mu f_\lambda - \lambda f_\mu + f_{\lambda\mu} \in J^2$이고,
$f_\lambda f_\mu \in J^2$이므로
$(x_\lambda - f_\lambda)(x_\mu - f_\mu) -
(x_{\lambda\mu} - f_{\lambda\mu}) \in J^2$이다.
따라서 $s(\lambda)s(\mu) = s(\lambda\mu)$이다.
$\lambda \in R$이면 $\text{d}\lambda = 0$이므로
$\text{d}f_\lambda = \text{d}x_\lambda$이고,
$\lambda - x_\lambda + f_\lambda \in J^2$이다. 따라서 원하는 대로
$s(\lambda) = \lambda$이다. 이제 보조정리
\ref{algebra-lemma-characterize-formally-smooth}를 적용하여 $S/R$가
형식적으로 매끄럽다는 결론을 얻는다.
\end{proof}

\begin{proposition}
\label{algebra-proposition-characterize-formally-smooth}
$R \to S$를 환 사상이라 하자. 형식적으로 매끄러운 $R$-대수 $P$와
핵이 $J$인 전사 $P \to S$를 생각하자. 다음 조건들은 서로 동치이다.
\begin{enumerate}
\item $S$는 $R$ 위에서 형식적으로 매끄럽다,
\item 위와 같은 어떤 $P \to S$에 대해 $P/J^2 \to S$의 단면이 존재한다,
\item 위와 같은 모든 $P \to S$에 대해 $P/J^2 \to S$의 단면이 존재한다,
\item 위와 같은 어떤 $P \to S$에 대해
$0 \to J/J^2 \to \Omega_{P/R} \otimes S \to \Omega_{S/R} \to 0$가 분할 완전이다,
\item 위와 같은 모든 $P \to S$에 대해
$0 \to J/J^2 \to \Omega_{P/R} \otimes S \to \Omega_{S/R} \to 0$가 분할 완전이다,
\item 순진 코탄젠트 복합체 $\NL_{S/R}$가 $0$차에 놓인 사영
$S$-가군과 준동형 동치이다. 이는 $H_1(\NL_{S/R}) = 0$이고
$\Omega_{S/R}$가 사영 $S$-가군이라는 뜻이다.
\end{enumerate}
\end{proposition}

\begin{proof}
(1)이 (3)을 함의하고 (3)이 (2)를 함의한다는 것은 명백하다.
보조정리 \ref{algebra-lemma-characterize-formally-smooth}의 앞부분을 보라.
또한 (3)이 (5)를 함의하고 (5)가 (4)를 함의하며 (2)가 (4)를 함의한다.
이는 보조정리 \ref{algebra-lemma-characterize-formally-smooth-again}의
앞부분을 보라.
마지막으로 보조정리 \ref{algebra-lemma-characterize-formally-smooth-again}을
표준 전사 $R[S] \to S$ (\ref{algebra-equation-canonical-presentation})에 적용하면
(1)이 (6)을 함의함을 알 수 있다.

\medskip\noindent
(4)가 (6)을 함의함을 보이자. 환 사상 $R \to P \to S$에 대응하는
보조정리 \ref{algebra-lemma-exact-sequence-NL}의 열을 생각하자.
위에서 증명한 (1) $\Rightarrow$ (6)에 의해
$\NL_{P/R} \otimes_R S$는 $0$차에 놓인
$\Omega_{P/R} \otimes_P S$와 준동형 동치이다.
따라서 $H_1(\NL_{P/R} \otimes_P S) = 0$이다.
$P \to S$가 전사이므로 보조정리 \ref{algebra-lemma-NL-surjection}에 의해
$\NL_{S/P}$는 $1$차에 놓인 $J/J^2$와 호모토피 동치이다.
따라서 정확한 열
$0 \to H_1(L_{S/R}) \to J/J^2 \to \Omega_{P/R} \otimes_P S \to
\Omega_{S/R} \to 0$를 얻는다.
가정에 의해 $H_1(L_{S/R}) = 0$이고 $\Omega_{S/R}$는 사영
$S$-가군이다. 따라서 (6)이 따른다.

\medskip\noindent
마지막으로 (6)이 (1)을 함의함을 보이자. 이 가정은
$P = R[S]$이고 $P \to S$가 표준 전사
(\ref{algebra-equation-canonical-presentation})일 때의
$J/J^2 \to \Omega_{P/R} \otimes S$ 복합체가 $0$차에 놓인
사영 $S$-가군과 준동형 동치라는 뜻이다.
따라서 보조정리 \ref{algebra-lemma-characterize-formally-smooth-again}에 의해
$S$는 $R$ 위에서 형식적으로 매끄럽다.
\end{proof}

\begin{lemma}
\label{algebra-lemma-ses-formally-smooth}
$A \to B \to C$를 환 사상들이라 하자. $B \to C$가 형식적으로 매끄럽다고 하자.
그러면 보조정리 \ref{algebra-lemma-exact-sequence-differentials}의 열
$$
0 \to \Omega_{B/A} \otimes_B C \to \Omega_{C/A} \to \Omega_{C/B} \to 0
$$
은 분할 짧은 완전열이다.
\end{lemma}

\begin{proof}
명제 \ref{algebra-proposition-characterize-formally-smooth}와
보조정리 \ref{algebra-lemma-exact-sequence-NL}에서 따른다.
\end{proof}

\begin{lemma}
\label{algebra-lemma-differential-seq-formally-smooth}
$A \to B \to C$를 환 사상들이라 하자. $A \to C$가 형식적으로 매끄럽고
$B \to C$가 핵 $J \subset B$를 갖는 전사라고 하자. 그러면 보조정리
\ref{algebra-lemma-differential-seq}의 완전열
$$
0 \to J/J^2 \to \Omega_{B/A} \otimes_B C \to \Omega_{C/A} \to 0
$$
은 분할 완전이다.
\end{lemma}

\begin{proof}
명제 \ref{algebra-proposition-characterize-formally-smooth},
보조정리 \ref{algebra-lemma-exact-sequence-NL}, 보조정리
\ref{algebra-lemma-differential-seq}에서 따른다.
\end{proof}

\begin{lemma}
\label{algebra-lemma-application-NL-formally-smooth}
$A \to B \to C$를 환 사상들이라 하자. $A \to C$가 전사(따라서
$B \to C$도 전사)이고 $A \to B$가 형식적으로 매끄럽다고 하자.
$I = \Ker(A \to C)$, $J = \Ker(B \to C)$로 두자. 그러면 보조정리
\ref{algebra-lemma-application-NL}의 열
$$
0 \to I/I^2 \to J/J^2 \to \Omega_{B/A} \otimes_B B/J \to 0
$$
은 분할 완전이다.
\end{lemma}

\begin{proof}
$A \to B$가 형식적으로 매끄러우므로
환 사상 $\sigma : B \to A/I^2$가 존재하며 $A \to B$와의 합성은
몫 사상 $A \to A/I^2$와 같다. 그러면 $\sigma$는
$J/J^2 \to I/I^2$를 유도하고, 이는 $I/I^2 \to J/J^2$의 역이다.
\end{proof}

\begin{lemma}
\label{algebra-lemma-lift-formal-smoothness}
$R \to S$를 환 사상이라 하자. $I \subset R$를 아이디얼이라 하자.
다음을 가정하자.
\begin{enumerate}
\item $I^2 = 0$,
\item $R \to S$가 평탄하다,
\item $R/I \to S/IS$가 형식적으로 매끄럽다.
\end{enumerate}
그러면 $R \to S$는 형식적으로 매끄럽다.
\end{lemma}

\begin{proof}
(1), (2), (3)을 가정하자.
$P = R[\{x_t\}_{t \in T}] \to S$를 핵 $J$를 갖는 전사 $R$-대수 사상이라 하자.
따라서 $0 \to J \to P \to S \to 0$는 평탄한 $R$-가군들의 짧은 완전열이다.
이는 $I \otimes_R S = IS$, $I \otimes_R P = IP$, $I \otimes_R J = IJ$ 및
$J \cap IP = IJ$를 함의한다.
증명 전체에서
$$
\Omega_{(S/IS)/(R/I)} = \Omega_{S/R} \otimes_S (S/IS)
= \Omega_{S/R} \otimes_R R/I = \Omega_{S/R} / I\Omega_{S/R}
$$
및 $P$에 대한 유사한 식을 사용하겠다. 보조정리 \ref{algebra-lemma-differentials-base-change}를 보라.
보조정리 \ref{algebra-lemma-characterize-formally-smooth-again}에 의해 열
\begin{equation}
\label{algebra-equation-split}
0 \to J/(IJ + J^2) \to
\Omega_{P/R} \otimes_P S/IS \to
\Omega_{S/R} \otimes_S S/IS \to 0
\end{equation}
은 분할 완전이다. 중간 항은
$\bigoplus_{t \in T} S/IS \text{d}x_t$이다.
분할 $\sigma : \Omega_{P/R} \otimes_P S/IS \to J/(IJ + J^2)$를 택하자.
각 $t \in T$에 대해 $J/(IJ + J^2)$에서
$\sigma(\text{d}x_t)$로 사상되는 $f_t \in J$를 택하자. 그러면
$$
\tilde \sigma : \Omega_{P/R} \otimes_R S
= \bigoplus S\text{d}x_t \longrightarrow J/J^2
$$
라는 유일한 $S$-가군 사상이 정해지며
$\tilde\sigma(\text{d}x_t) = f_t$이다.
$\sigma$가 $\text{d}$의 단면이므로 차이
$$
\Delta = \text{id}_{J/J^2} - \tilde \sigma \circ \text{d}
$$
는 $J/J^2 \to J/J^2$ 자기 사상이고 그 상은
$(IJ + J^2)/J^2$에 포함된다. 특히 $I^2 = 0$이므로
$\Delta((IJ + J^2)/J^2) = 0$이다. 따라서 $\Delta$는
$$
J/J^2 \to J/(IJ + J^2) \xrightarrow{\overline{\Delta}}
(IJ + J^2)/J^2 \to J/J^2
$$
로 인수분해되며, 여기서 $\overline{\Delta}$는 $S/IS$-가군 사상이다.
열 (\ref{algebra-equation-split})이 분할이라는 사실을 다시 사용하면
$\overline{\delta} \circ d$가 $\overline{\Delta}$와 같아지는
$S/IS$-가군 사상
$\overline{\delta} : \Omega_{P/R} \otimes_P S/IS \to (IJ + J^2)/J^2$를
찾을 수 있다. 같은 방식으로 $\overline{\delta}$는
$S$-가군 사상
$\delta : \Omega_{P/R} \otimes_P S \to J/J^2$를 정한다.
$\tilde \sigma$를 $\tilde \sigma + \delta$로 바꾸면 간단한 계산으로
$\Delta = 0$을 얻는다. 즉 $\tilde \sigma$는
$J/J^2 \to \Omega_{P/R} \otimes_P S$의 단면이다.
보조정리 \ref{algebra-lemma-characterize-formally-smooth-again}에 의해
$R \to S$가 형식적으로 매끄럽다는 결론을 얻는다.
\end{proof}

\begin{proposition}
\label{algebra-proposition-smooth-formally-smooth}
$R \to S$를 환 사상이라 하자. 다음 조건들은 서로 동치이다.
\begin{enumerate}
\item $R \to S$는 유한 표시이고 형식적으로 매끄럽다,
\item $R \to S$는 매끄럽다.
\end{enumerate}
\end{proposition}

\begin{proof}
명제 \ref{algebra-proposition-characterize-formally-smooth}와 정의
\ref{algebra-definition-smooth}에서 따른다.
($R \to S$가 유한 표시이면 $\Omega_{S/R}$가 유한 표시 $S$-가군이라는
점에 주의하자. 보조정리 \ref{algebra-lemma-differentials-finitely-presented}를 보라.)
\end{proof}

\begin{lemma}
\label{algebra-lemma-finite-presentation-fs-Noetherian}
$R \to S$를 매끄러운 환 사상이라 하자. 그러면 $\mathbf{Z}$ 위에서
유한형인 부분환 $R_0 \subset R$가 존재하고,
$R_0 \to S_0$가 매끄러운 환 사상이며 $S \cong R \otimes_{R_0} S_0$가 된다.
\end{lemma}

\begin{proof}
매끄러움은 유한 표시이고 형식적으로 매끄러운 것과 동치라는 사실을
명제 \ref{algebra-proposition-smooth-formally-smooth}에 의해 사용한다.
$S = R[x_1, \ldots, x_n]/(f_1, \ldots, f_m)$로 쓰고
$I = (f_1, \ldots, f_m)$로 두자.
보조정리 \ref{algebra-lemma-characterize-formally-smooth}에 따라
$S$로의 사영에 대한 우역
$\sigma : S \to R[x_1, \ldots, x_n]/I^2$를 택하자.
$\sigma(x_i \bmod I) = h_i \bmod I^2$가 되도록
$h_i \in R[x_1, \ldots, x_n]$을 택하자. $x_i - h_i \in I$이므로
$$
x_i - h_i = \sum\nolimits_j b_{ij} f_j
$$
가 되는 $b_{ij} \in R[x_1, \ldots, x_n]$가 존재한다.
$\sigma$가 $R$-대수 준동형
$R[x_1, \ldots, x_n]/I \to R[x_1, \ldots, x_n]/I^2$이라는 사실은
다음 조건과 동치이다.
$$
f_j(h_1, \ldots, h_n) = \sum\nolimits_{j_1 j_2} a_{j_1 j_2} f_{j_1} f_{j_2}
$$
일부 $a_{kl} \in R[x_1, \ldots, x_n]$에 대해.
$R_0 \subset R$를 모든 다항식 $f_j,h_i,a_{kl},b_{ij}$의 계수로
$\mathbf{Z}$ 위에서 생성되는 부분환으로 두자.
$S_0 = R_0[x_1, \ldots, x_n]/(f_1, \ldots, f_m)$,
$I_0 = (f_1, \ldots, f_m)$로 둔다. 두 번째 표시식이
$R_0[x_1, \ldots, x_n]$에서 성립하므로
$x_i \mapsto h_i \bmod I_0^2$로 정의되는 $R_0$-대수 사상
$\sigma_0 : S_0 \to R_0[x_1, \ldots, x_n]/I_0^2$를 둘 수 있다.
첫 번째 표시식이 $R_0[x_1, \ldots, x_n]$에서 성립하므로
$\sigma_0$는 사영
$R_0[x_1, \ldots, x_n]/I_0^2 \to R_0[x_1, \ldots, x_n]/I_0 = S_0$의
우역이다. 따라서 보조정리 \ref{algebra-lemma-characterize-formally-smooth}에 의해
$S_0$는 $R_0$ 위에서 형식적으로 매끄럽다.
\end{proof}

\begin{lemma}
\label{algebra-lemma-smooth-descends-through-colimit}
$A = \colim A_i$를 환들의 여과 여극한이라 하자.
$A \to B$를 매끄러운 환 사상이라 하자. 어떤 $i$와 매끄러운 환 사상
$A_i \to B_i$가 존재하여 $B = B_i \otimes_{A_i} A$가 된다.
\end{lemma}

\begin{proof}
보조정리 \ref{algebra-lemma-finite-presentation-fs-Noetherian}에서 따른다.
보조정리 \ref{algebra-lemma-characterize-finite-presentation}에 의해
$R_0 \to A$는 어떤 $i$에 대해 $A_i$를 통해 인수분해되기 때문이다.
\end{proof}

\begin{lemma}
\label{algebra-lemma-descent-formally-smooth}
$R \to S$를 환 사상이라 하자. $R \to R'$를 충실히 평탄한 환 사상이라 하자.
$S' = S \otimes_R R'$로 두자. 그러면 $R \to S$가 형식적으로 매끄러운 것은
$R' \to S'$가 형식적으로 매끄러운 것과 동치이다.
\end{lemma}

\begin{proof}
$R \to S$가 형식적으로 매끄러우면 보조정리
\ref{algebra-lemma-base-change-fs}에 의해 $R' \to S'$도 형식적으로 매끄럽다.
역을 보이자. $R' \to S'$가 형식적으로 매끄럽다고 가정하자.
임의의 $S$-가군 $N$에 대해 $N \otimes_R R' = N \otimes_S S'$임에 주의하자.
특히 $S \to S'$도 충실히 평탄하다.
다항식환 $P = R[\{x_i\}_{i \in I}]$와 핵이 $J$인 전사
$R$-대수 사상 $P \to S$를 택하자.
$P' = P \otimes_R R'$는 $R'$ 위의 다항식 대수이다.
$R \to R'$가 평탄하므로 전사 $P' \to S'$의 핵 $J'$는
$J \otimes_R R'$이다. 따라서 보조정리
\ref{algebra-lemma-characterize-formally-smooth-again}에 따른 분할 완전열
$$
0 \to J'/(J')^2 \to \Omega_{P'/R'} \otimes_{P'} S' \to \Omega_{S'/R'} \to 0
$$
은 대응하는 열
$$
J/J^2 \to \Omega_{P/R} \otimes_P S \to \Omega_{S/R} \to 0
$$
의 $S \to S'$에 의한 기저변환이다. 이는 보조정리
\ref{algebra-lemma-differential-seq}를 따른다.
$S \to S'$가 충실히 평탄하므로 두 가지를 결론지을 수 있다:
(1) 이 열(${ }'$가 없는 것)도 완전이고, (2)
$\Omega_{S/R}$는 사영 $S$-가군이다. 실제로
$\Omega_{S'/R'}$는 자유 가군
$\Omega_{P'/R'} \otimes_{P'} S'$의 직합인자로서 사영이고,
위에서 본 바에 의해 $\Omega_{S/R} \otimes_S {S'} = \Omega_{S'/R'}$이다.
따라서 보조정리 \ref{algebra-theorem-ffdescent-projectivity}에 따른
충실히 평탄한 환 사상을 통한 사영성 하강으로 (2)가 따른다.
그러므로
$$
0 \to J/J^2 \to \Omega_{P/R} \otimes_P S \to \Omega_{S/R} \to 0
$$
도 완전이고, 보조정리 \ref{algebra-lemma-characterize-formally-smooth-again}를
다시 적용하면 결론을 얻는다.
\end{proof}

\noindent
매끄러운 환 사상은 다음과 같은 강한 올림 성질을 만족한다는 것이 밝혀진다.

\begin{lemma}
\label{algebra-lemma-smooth-strong-lift}
$R \to S$를 매끄러운 환 사상이라 하자. $I \subset A$가 국소 멱영
아이디얼인 가환 실선 도식
$$
\xymatrix{
S \ar[r] \ar@{-->}[rd] & A/I \\
R \ar[r] \ar[u] & A \ar[u]
}
$$
이 주어지면 도식을 가환으로 만드는 점선 사상이 존재한다.
\end{lemma}

\begin{proof}
보조정리 \ref{algebra-lemma-finite-presentation-fs-Noetherian}에 의해 이 도식을
다음 가환 도식으로 확장할 수 있다.
$$
\xymatrix{
S_0 \ar[r] & S \ar[r] \ar@{-->}[rd] & A/I \\
R_0 \ar[r] \ar[u] & R \ar[r] \ar[u] & A \ar[u]
}
$$
여기서 $R_0 \to S_0$는 매끄럽고, $R_0$는 $\mathbf{Z}$ 위에서 유한형이며,
$S = S_0 \otimes_{R_0} R$이다. $x_1, \ldots, x_n \in S_0$를
$S_0$의 $R_0$ 위 생성원이라 하자. $x_1, \ldots, x_n$과 $A/I$에서 같은 원소로
사상되는 $A$의 원소들을 $a_1, \ldots, a_n$이라 하자.
$A_0 \subset A$를 $R_0$의 상과 $a_1, \ldots, a_n$으로 생성되는 부분환으로
두자. $I_0 = A_0 \cap I$로 두자. 그러면 $A_0/I_0 \subset A/I$이고
$S_0 \to A/I$의 상은 $A_0/I_0$ 안에 있다. 따라서 다음 도식에서 점선 사상을
찾으면 충분하다.
$$
\xymatrix{
S_0 \ar[r] \ar@{-->}[rd] & A_0/I_0 \\
R_0 \ar[r] \ar[u] & A_0 \ar[u]
}
$$
구성에 의해 $A_0$는 $\mathbf{Z}$ 위에서 유한형이다. 따라서 $A_0$는
뇌터 환이고, 보조정리 \ref{algebra-lemma-Noetherian-power}에 의해 $I_0$는 멱영이다.
$I_0^n = 0$이라고 하자. 명제 \ref{algebra-proposition-smooth-formally-smooth}에 의해
$R_0$-대수 사상 $S_0 \to A_0/I_0$를 차례로
$S_0 \to A_0/I_0^2$, $S_0 \to A_0/I_0^3$, $\ldots$, 그리고
$S_0 \to A_0/I_0^n = A_0$로 올릴 수 있다.
\end{proof}

\section{매끄러움과 미분}
\label{algebra-section-smooth-differential}

\noindent
미분과 매끄러운 환 사상에 관한 몇 가지 결과를 제시한다.

\begin{lemma}
\label{algebra-lemma-triangle-differentials-smooth}
$A \to B \to C$가 환 사상들이고 $B \to C$가 매끄럽다고 하자. 그러면
보조정리 \ref{algebra-lemma-exact-sequence-differentials}의 열
$$
0 \to C \otimes_B \Omega_{B/A} \to \Omega_{C/A} \to \Omega_{C/B} \to 0
$$
은 완전하다.
\end{lemma}

\begin{proof}
이는 보조정리 \ref{algebra-lemma-ses-formally-smooth}의 더 일반적인 결과에서 따른다.
매끄러운 환 사상은 형식적으로 매끄럽기 때문이다. 명제
\ref{algebra-proposition-smooth-formally-smooth}를 보라.
또한 $H_1(L_{C/B}) = 0$가 $B \to C$의 매끄러움 정의에 포함되므로,
보조정리 \ref{algebra-lemma-exact-sequence-NL}에서도 직접 따른다.
\end{proof}

\begin{lemma}
\label{algebra-lemma-differential-seq-smooth}
$A \to B \to C$를 환 사상들이라 하자. $A \to C$가 매끄럽고
$B \to C$가 핵 $J \subset B$를 갖는 전사라고 하자.
그러면 보조정리 \ref{algebra-lemma-differential-seq}의 완전열
$$
0 \to J/J^2 \to \Omega_{B/A} \otimes_B C \to \Omega_{C/A} \to 0
$$
은 분할 완전이다.
\end{lemma}

\begin{proof}
이는 보조정리 \ref{algebra-lemma-differential-seq-formally-smooth}의 더 일반적인
결과에서 따른다. 매끄러운 환 사상은 형식적으로 매끄럽기 때문이다.
명제 \ref{algebra-proposition-smooth-formally-smooth}를 보라.
\end{proof}

\begin{lemma}
\label{algebra-lemma-application-NL-smooth}
$A \to B \to C$를 환 사상들이라 하자. $A \to C$가 전사(따라서
$B \to C$도 전사)이고 $A \to B$가 매끄럽다고 하자.
$I = \Ker(A \to C)$, $J = \Ker(B \to C)$로 두자.
그러면 보조정리 \ref{algebra-lemma-application-NL}의 열
$$
0 \to I/I^2 \to J/J^2 \to \Omega_{B/A} \otimes_B B/J \to 0
$$
은 완전하다.
\end{lemma}

\begin{proof}
이는 보조정리 \ref{algebra-lemma-application-NL-formally-smooth}의 더 일반적인
결과에서 따른다. 매끄러운 환 사상은 형식적으로 매끄럽기 때문이다.
명제 \ref{algebra-proposition-smooth-formally-smooth}를 보라.
\end{proof}

\begin{lemma}
\label{algebra-lemma-section-smooth}
\begin{slogan}
$R$가 $S$의 직합인자이고 $S$가 $R$ 위에서 매끄러우면,
$S$의 $I$-진 완비화는 종종 $R$ 위의 멱급수환이다.
여기서 $I$는 $S$에서 $R$로 가는 사상의 핵이다.
\end{slogan}
$\varphi : R \to S$를 매끄러운 환 사상이라 하자.
$\sigma : S \to R$를 $ \varphi$의 좌역이라 하자.
$I = \Ker(\sigma)$로 두자. 그러면
\begin{enumerate}
\item $I/I^2$는 유한 국소 자유 $R$-가군이고,
\item $I/I^2$가 자유이면 $S^\wedge \cong R[[t_1, \ldots, t_d]]$가
$R$-대수로서 성립한다. 여기서 $S^\wedge$는 $S$의 $I$-진 완비화이다.
\end{enumerate}
\end{lemma}

\begin{proof}
$R \to S \to R$에 보조정리 \ref{algebra-lemma-differential-seq-split}을 적용하면
$I/I^2 = \Omega_{S/R} \otimes_{S, \sigma} R$임을 얻는다.
매끄러운 사상의 정의에 따라 $\Omega_{S/R}$가 $S$ 위의 유한 국소 자유
가군이므로 (1)이 따른다.
$I/I^2$가 자유라면 $I/I^2$에서의 상들이 $R$-기저를 이루는
$f_1, \ldots, f_d \in I$를 택하자. 다음과 같이 정의된 $R$-대수 사상을 생각하자.
$$
\Psi : R[[x_1, \ldots, x_d]] \longrightarrow S^\wedge, \quad
x_i \longmapsto f_i.
$$
$P = R[[x_1, \ldots, x_d]]$, $J = (x_1, \ldots, x_d) \subset P$로 두자.
유도되는 몫환 사상을 $\Psi_n : P/J^n \to S/I^n$로 쓰자.
$S/I^2 = \varphi(R) \oplus I/I^2$임에 주의하자. 따라서 $\Psi_2$는 동형이다.
$\Psi_2$의 역을 $\sigma_2 : S/I^2 \to P/J^2$로 쓰자.
$n$에 대한 귀납법으로 모든 $n > 2$에 대해 $\Psi_n$의 역인
$\sigma_n : S/I^n \to P/J^n$가 존재함을 보이겠다.
실제로 $S$는 $R$ 위에서 형식적으로 매끄럽다
(명제 \ref{algebra-proposition-smooth-formally-smooth}에 의함).
따라서 $R$-대수의 실선 도식
$$
\xymatrix{
S \ar@{..>}[r] \ar[rd]_{\sigma_{n - 1}} & P/J^n \ar[d] \\
 & P/J^{n - 1}
}
$$
에서 점선 사상을 어떤 $R$-대수 사상
$\tau : S \to P/J^n$으로 채워 도식을 가환으로 만들 수 있다.
이는 $R$-대수 사상 $\overline{\tau} : S/I^n \to P/J^n$를 유도하며,
이 사상은 $J^{n-1}$을 법으로 하여 $\sigma_{n-1}$과 같다.
구성에 의해 $\Psi_n$은 전사이고,
$\overline{\tau} \circ \Psi_n$은 $P/J^n$의 $R$-대수 자기 사상이며
$x_i$를 $x_i + \delta_{i,n}$으로 보낸다. 여기서
$\delta_{i,n} \in J^{n-1}/J^n$이다.
따라서 $\Psi_n$은 동형이고 역 $\sigma_n$을 갖는다.
이것으로 보조정리가 증명된다.
\end{proof}

\section{체 위의 매끄러운 대수}
\label{algebra-section-smooth-over-field}

\noindent
경고: 다음 두 보조정리는 일반적으로 비완전체 위에서는 성립하지 않는다.

\begin{lemma}
\label{algebra-lemma-rank-omega}
$k$를 대수적으로 닫힌 체라 하자.
$S$를 유한 생성 $k$-대수라 하자.
$\mathfrak m \subset S$를 극대 아이디얼이라 하자. 그러면
$$
\dim_{\kappa(\mathfrak m)} \Omega_{S/k} \otimes_S \kappa(\mathfrak m)
=
\dim_{\kappa(\mathfrak m)} \mathfrak m/\mathfrak m^2.
$$
\end{lemma}

\begin{proof}
보조정리 \ref{algebra-lemma-differential-seq}의 완전열
$$
\mathfrak m/\mathfrak m^2 \to
\Omega_{S/k} \otimes_S \kappa(\mathfrak m) \to
\Omega_{\kappa(\mathfrak m)/k} \to 0
$$
을 생각하자. 첫 번째 사상이 동형임을 보이면 된다. $k$가 대수적으로
닫혀 있으므로 정리 \ref{algebra-theorem-nullstellensatz}에 의해 합성
$k \to \kappa(\mathfrak m)$은 동형이다.
따라서 전사 $S \to \kappa(\mathfrak m)$은 $k$-대수 사상으로서 분할되고,
보조정리 \ref{algebra-lemma-differential-seq-split}에 의해 위 열은 왼쪽에서도
완전하다. $\Omega_{\kappa(\mathfrak m)/k} = 0$이므로 결론을 얻는다.
\end{proof}

\begin{lemma}
\label{algebra-lemma-characterize-smooth-kbar}
$k$를 대수적으로 닫힌 체라 하자.
$S$를 유한 생성 $k$-대수라 하자.
$\mathfrak m \subset S$를 극대 아이디얼이라 하자.
다음 조건들은 서로 동치이다.
\begin{enumerate}
\item 환 $S_{\mathfrak m}$은 정칙 국소환이다.
\item 다음이 성립한다:
$\dim_{\kappa(\mathfrak m)} \Omega_{S/k} \otimes_S \kappa(\mathfrak m)
\leq \dim(S_{\mathfrak m})$.
\item 다음이 성립한다:
$\dim_{\kappa(\mathfrak m)} \Omega_{S/k} \otimes_S \kappa(\mathfrak m)
= \dim(S_{\mathfrak m})$.
\item $g \in S$, $g \not \in \mathfrak m$인 원소가 존재하여
$S_g$가 $k$ 위에서 매끄럽다. 다시 말해 $S/k$는
$\mathfrak m$에서 매끄럽다.
\end{enumerate}
\end{lemma}

\begin{proof}
(1), (2), (3)은 보조정리 \ref{algebra-lemma-rank-omega}와 정의
\ref{algebra-definition-regular}에 의해 동치임에 주의하자.

\medskip\noindent
$S$가 $\mathfrak m$에서 매끄럽다고 하자.
보조정리 \ref{algebra-lemma-smooth-syntomic}에 의해 적절한
$g \in S$, $g \not \in \mathfrak m$에 대해
$S_g$가 $k$ 위의 표준 매끄러운 대수임을 안다.
따라서 보조정리 \ref{algebra-lemma-standard-smooth}에 의해
$\Omega_{S_g/k}$는 계수 $\dim(S_g)$의 자유 가군이다.
그러므로 보조정리 \ref{algebra-lemma-rank-omega}에 의해
$\dim(S_{\mathfrak m}) = \dim (\mathfrak m/\mathfrak m^2)$이고,
다시 말해 $S_\mathfrak m$은 정칙이다.

\medskip\noindent
반대로 $S_{\mathfrak m}$이 정칙이라고 하자.
$d = \dim(S_{\mathfrak m}) = \dim \mathfrak m/\mathfrak m^2$로 두자.
모든 $i$에 대해 $x_i$가 $\mathfrak m$의 원소로 사상되도록
표시 $S = k[x_1, \ldots, x_n]/I$를 택하자.
다시 말해 $\mathfrak m'' = (x_1, \ldots, x_n)$는
$k[x_1, \ldots, x_n]$의 대응하는 극대 아이디얼이다.
짧은 완전열
$$
I/\mathfrak m''I \to \mathfrak m''/(\mathfrak m'')^2
\to \mathfrak m/(\mathfrak m)^2 \to 0
$$
이 있다.
$\mathfrak m''/(\mathfrak m'')^2$에서의 상들이
$\mathfrak m/\mathfrak m^2$로 가는 사상의 핵을 생성하도록
$c = n - d$개의 원소 $f_1, \ldots, f_c \in I$를 택하자.
이는 분명 가능하다. $J = (f_1, \ldots, f_c)$로 두자. 그러면 $J \subset I$이다.
$S' = k[x_1, \ldots, x_n]/J$로 두면 전사 $S' \to S$가 있다.
$\mathfrak m' = \mathfrak m''S'$를 $S'$의 대응하는 극대 아이디얼이라 하자.
따라서 다음 도식이 있다.
$$
\xymatrix{
k[x_1, \ldots, x_n] \ar[r] & S' \ar[r] & S \\
\mathfrak m'' \ar[u] \ar[r] & \mathfrak m' \ar[r] \ar[u] &
\mathfrak m \ar[u]
}
$$
$J$의 선택에 의해 완전열
$$
J/\mathfrak m''J \to \mathfrak m''/(\mathfrak m'')^2
\to \mathfrak m'/(\mathfrak m')^2 \to 0
$$
은 $\dim( \mathfrak m'/(\mathfrak m')^2 ) = d$임을 보인다.
$S'_{\mathfrak m'}$가 $S_{\mathfrak m}$ 위로 전사하므로
$\dim(S_{\mathfrak m'}) \geq d$이다. 따라서 정의
\ref{algebra-definition-regular-local} 앞의 논의에 의해
$S'_{\mathfrak m'}$도 차원 $d$인 정칙환이다.
$S'$가 $c = n - d$개의 방정식으로 잘렸으므로
보조정리 \ref{algebra-lemma-lci}에 의해
$S'_{g'}$가 $k$ 위의 전역 완전교차가 되는
$g' \in S'$, $g' \not \in \mathfrak m'$가 존재한다.
또한 $S'_{\mathfrak m'} \to S_{\mathfrak m}$는 같은 차원의
뇌터 국소 정역들 사이의 전사이므로 동형이다.
따라서 $S' \to S$는 유한 생성 핵을 갖는 전사이고
$\mathfrak m'$에서 국소화하면 동형이 된다.
그러므로 $S'_{g'} \to S_{g'}$가 동형이 되는
$g' \in S'$, $g \not \in \mathfrak m'$를 찾을 수 있다.
결국 $S$를 주 아이디얼 국소화로 바꾼 뒤 $S$가
전역 완전교차라고 가정할 수 있다.

\medskip\noindent
이제 $S = k[x_1, \ldots, x_n]/(f_1, \ldots, f_c)$이고
$\dim S = n - c$로 쓸 수 있다. 이 대수의 순진 코탄젠트 복합체는
$$
\bigoplus S \cdot f_j
\to
\bigoplus S \cdot \text{d}x_i
$$
로 주어짐을 상기하자. 보조정리
\ref{algebra-lemma-relative-global-complete-intersection-conormal}를 보라.
보조정리 \ref{algebra-lemma-relative-global-complete-intersection-smooth}에 의해
$S$가 $\mathfrak m$에서 매끄러움을 보이려면 위 사상을 나타내는
행렬 “$A$”의 $c \times c$ 소행렬식 $g_I$ 가운데 하나가
$\mathfrak m$에서 소멸하지 않음을 보이면 된다.
보조정리 \ref{algebra-lemma-rank-omega}에 의해
행렬 $A \bmod \mathfrak m$의 계수는 $c$이다. 이것으로 증명된다.
\end{proof}

\begin{lemma}
\label{algebra-lemma-characterize-smooth-over-field}
$k$를 임의의 체라 하자.
$S$를 유한 생성 $k$-대수라 하자.
$X = \Spec(S)$로 두자.
$\mathfrak q \subset S$를 $x \in X$에 대응하는 소 아이디얼이라 하자.
다음 조건들은 서로 동치이다.
\begin{enumerate}
\item $k$-대수 $S$는 $k$ 위에서 $\mathfrak q$에서 매끄럽다.
\item 다음이 성립한다:
$\dim_{\kappa(\mathfrak q)} \Omega_{S/k} \otimes_S \kappa(\mathfrak q)
\leq \dim_x X$.
\item 다음이 성립한다:
$\dim_{\kappa(\mathfrak q)} \Omega_{S/k} \otimes_S \kappa(\mathfrak q)
= \dim_x X$.
\end{enumerate}
더욱이 이 경우 국소환 $S_{\mathfrak q}$는 정칙이다.
\end{lemma}

\begin{proof}
$S$가 $k$ 위에서 $\mathfrak q$에서 매끄러우면
어떤 $g \in S$, $g \not \in \mathfrak q$가 존재하여
$S_g$가 $k$ 위에서 표준 매끄럽다.
보조정리 \ref{algebra-lemma-smooth-syntomic}을 보라.
$k$ 위의 표준 매끄러운 대수의 미분 가군은 그 차원과 같은 계수의
자유 가군이다. 보조정리 \ref{algebra-lemma-standard-smooth}를 보라
(체 위의 상대 전역 완전교차의 차원은 변수 수에서 방정식 수를 뺀 것임을 사용한다).
따라서 (1)은 (3)을 함의한다. 이 보조정리의 증명을 끝내려면 (2)가 (1)을
함의하고 $S_{\mathfrak q}$가 정칙임을 함의한다는 것을 보이면 충분하다.

\medskip\noindent
(2)를 가정하자. 나카야마 보조정리 \ref{algebra-lemma-NAK}에 의해
$\Omega_{S/k, \mathfrak q}$는 $\leq \dim_x X$개의 원소로 생성될 수 있다.
어떤 $g \in S$, $g \not \in \mathfrak q$에 대해 $S$를 $S_g$로 바꾸어
$\Omega_{S/k}$가 많아야 $\dim_x X$개의 원소로 생성된다고 가정할 수 있다.
$K/k$를 대수적으로 닫힌 체의 확대이고
$k$-대수 사상 $\psi : \kappa(\mathfrak q) \to K$가 존재하도록 택하자.
$S_K = K \otimes_k S$를 생각하자.
다음 전사에 대응하는 극대 아이디얼을 $\mathfrak m \subset S_K$라 하자.
$$
\xymatrix{
S_K = K \otimes_k S \ar[r] &
K \otimes_k \kappa(\mathfrak q)
\ar[r]^-{\text{id}_K \otimes \psi} &
K.
}
$$
$\mathfrak m \cap S = \mathfrak q$임에 주의하자. 다시 말해
$\mathfrak m$은 $\mathfrak q$ 위에 놓인다.
보조정리 \ref{algebra-lemma-dimension-at-a-point-preserved-field-extension}에 의해
$\mathfrak m$에 대응하는 점에서 $X_K = \Spec(S_K)$의 차원은 $\dim_x X$이다.
보조정리 \ref{algebra-lemma-dimension-closed-point-finite-type-field}에 의해
이는 $\dim((S_K)_{\mathfrak m})$과 같다.
보조정리 \ref{algebra-lemma-differentials-base-change}에 의해
$K$ 위의 $S_K$의 미분 가군은 $\Omega_{S/k}$의 기저변환이고,
따라서 많아야 $\dim_x X = \dim((S_K)_{\mathfrak m})$개의 원소로 생성된다.
보조정리 \ref{algebra-lemma-characterize-smooth-kbar}에 의해
$S_K$는 $K$ 위에서 $\mathfrak m$에서 매끄럽다.
보조정리 \ref{algebra-lemma-flat-base-change-locus-smooth}에 의해
$S$는 $k$ 위에서 $\mathfrak q$에서 매끄럽다. 이것으로 (1)이 증명된다.
더욱이 보조정리 \ref{algebra-lemma-characterize-smooth-kbar}에 의해
국소환 $(S_K)_{\mathfrak m}$이 정칙임을 안다.
$S_{\mathfrak q} \to (S_K)_{\mathfrak m}$가 평탄하므로
보조정리 \ref{algebra-lemma-flat-under-regular}에 의해
$S_{\mathfrak q}$가 정칙이라는 결론을 얻는다.
\end{proof}

\noindent
다음 보조정리는 여러 가지 방식으로 상당히 일반화할 수 있다.

\begin{lemma}
\label{algebra-lemma-computation-differential}
$k$를 체라 하자.
$R$을 $k$를 포함하는 뇌터 국소환이라 하자.
잉여체 $\kappa = R/\mathfrak m$이 $k$의 유한 생성 분리 확대라고 가정하자.
그러면 사상
$$
\text{d} :
\mathfrak m/\mathfrak m^2
\longrightarrow
\Omega_{R/k} \otimes_R \kappa(\mathfrak m)
$$
은 단사이다.
\end{lemma}

\begin{proof}
$R$을 $R/\mathfrak m^2$으로 바꾸어도 된다. 따라서
$\mathfrak m^2 = 0$이라고 가정할 수 있다. 가정에 의해
$\kappa = k(\overline{x}_1, \ldots, \overline{x}_r, \overline{y})$
로 쓸 수 있다. 여기서 $\overline{x}_1, \ldots, \overline{x}_r$는
$k$ 위의 $\kappa$의 초월 기저이고, $\overline{y}$는
$k(\overline{x}_1, \ldots, \overline{x}_r)$ 위에서 분리 대수적이다.
$P(T)$를 $k(\overline{x}_1, \ldots, \overline{x}_r)[T]$에서
$\overline{y}$의 $k(\overline{x}_1, \ldots, \overline{x}_r)$ 위 최소
다항식이라 하자.
그러면 $P(\overline{y}) = 0$이지만 $P'(\overline{y}) \not = 0$이다.
$\overline{x}_i \in \kappa$의 임의의 올림 $x_i \in R$을 택하자.
그러면 $k$-대수의 가환 도식
$$
\xymatrix{
R \ar[r] & \kappa \\
& k(\overline{x}_1, \ldots, \overline{x}_r) \ar[lu]^\varphi \ar[u]
}
$$
을 얻는다. 왼쪽 위쪽 사상 $\varphi$를 $\kappa$에서 $R$로 가는
$k$-대수 사상으로 확장하고 싶다.
이를 위해 $\overline{y}$의 임의의 올림 $y \in R$을 택하자.
이것이
$\kappa \cong k(\overline{x}_1, \ldots, \overline{x}_r)[T]/(P)$
위에 정의된 $k$-대수 사상을 주려면
$P^\varphi(y) = 0$이 되도록 $y$를 택할 수 있음을 보이면 충분하다.
그렇지 않으면 $\delta \in \mathfrak m$에 대해
$$
P^\varphi(y + \delta) = P^\varphi(y) + (P')^\varphi(y)\delta
$$
를 계산한다. 이는 $\mathfrak m^2 = 0$이기 때문이다.
$(P')^\varphi(y)$가 $R$의 가역 원소이므로 원하는 대로 선택을 조정할 수 있다.
따라서 $R \cong \kappa \oplus \mathfrak m$가 $k$-대수로서 성립한다!
이제 $\Omega_{\kappa \oplus \mathfrak m/k}$를 직접 계산하거나
보조정리 \ref{algebra-lemma-differential-seq-split}을 적용하면 증명이 끝난다.
\end{proof}

\begin{lemma}
\label{algebra-lemma-separable-smooth}
$k$를 체라 하자.
$S$를 유한 생성 $k$-대수라 하자.
$\mathfrak q \subset S$를 소 아이디얼이라 하자.
$\kappa(\mathfrak q)$가 $k$ 위에서 분리라고 가정하자.
다음 조건들은 서로 동치이다.
\begin{enumerate}
\item 대수 $S$는 $k$ 위에서 $\mathfrak q$에서 매끄럽다.
\item 환 $S_{\mathfrak q}$는 정칙이다.
\end{enumerate}
\end{lemma}

\begin{proof}
$R = S_{\mathfrak q}$로 두고, 그 극대 아이디얼을 $\mathfrak m$,
그 잉여체를 $\kappa$로 쓰자.
보조정리 \ref{algebra-lemma-computation-differential}과
\ref{algebra-lemma-differential-seq}에 의해 짧은 완전열
$$
0 \to \mathfrak m/\mathfrak m^2 \to
\Omega_{R/k} \otimes_R \kappa \to
\Omega_{\kappa/k} \to 0
$$
을 얻는다. 보조정리 \ref{algebra-lemma-differentials-localize}에 의해
$\Omega_{R/k} = \Omega_{S/k, \mathfrak q}$임에 주의하자.
더욱이 $\kappa$가 $k$ 위에서 분리이므로
$\dim_{\kappa} \Omega_{\kappa/k} = \text{trdeg}_k(\kappa)$이다.
따라서
$$
\dim_{\kappa} \Omega_{R/k} \otimes_R \kappa
=
\dim_\kappa \mathfrak m/\mathfrak m^2 + \text{trdeg}_k (\kappa)
\geq
\dim R + \text{trdeg}_k (\kappa)
=
\dim_{\mathfrak q} S
$$
를 얻는다. 마지막 등식은 보조정리
\ref{algebra-lemma-dimension-at-a-point-finite-type-field}를 보라.
등호가 성립하는 것은 $R$가 정칙인 것과 동치이다.
따라서 보조정리 \ref{algebra-lemma-characterize-smooth-over-field}를 적용하면 끝난다.
\end{proof}

\begin{lemma}
\label{algebra-lemma-characteristic-zero}
$R \to S$를 $\mathbf{Q}$-대수 사상이라 하자.
어떤 $S$-부분가군 $C$에 대해
$\Omega_{S/R} = S \text{d}f \oplus C$가 되는 $f \in S$를 잡자. 그러면
\begin{enumerate}
\item $f$는 멱영이 아니고,
\item $S$가 뇌터 국소환이면 $f$는 $S$의 영인자가 아니다.
\end{enumerate}
\end{lemma}

\begin{proof}
$a \in S$에 대해 어떤 $\theta(a) \in S$와 $c(a) \in C$가 존재하여
$\text{d}(a) = \theta(a)\text{d}f + c(a)$로 쓰자.
$R$-미분 $S \to S$, $a \mapsto \theta(a)$를 생각하자.
$\theta(f) = 1$임에 주의하자.

\medskip\noindent
$f^n = 0$이고 $n > 1$이 최소이면
$0 = \theta(f^n) = n f^{n - 1}$이므로 $n$의 최소성과 모순이다.
따라서 $f$는 멱영이 아니다.

\medskip\noindent
$fa = 0$이라고 하자. $f$가 가역이면 $a = 0$이고 끝난다.
$f$가 가역이 아니라고 하자. 라이프니츠 법칙에 의해
$0 = \theta(fa) = f\theta(a) + a$이므로 $a \in (f)$이다.
귀납적으로 $fa = 0 \Rightarrow a \in (f^n)$을 보였다고 하자.
$a = f^nb$로 쓰면
$0 = \theta(f^{n + 1}b) = (n + 1)f^nb + f^{n + 1}\theta(b)$이다.
따라서 $a = f^n b = -f^{n + 1}\theta(b)/(n + 1) \in (f^{n + 1})$이다.
뇌터 국소환 $S$에서 $\bigcap (f^n) = 0$이므로 결론을 얻는다.
보조정리 \ref{algebra-lemma-intersect-powers-ideal-module-zero}를 보라.
\end{proof}

\noindent
다음 보조정리는 응용에서는 아마 별로 쓸모가 없다.

\begin{lemma}
\label{algebra-lemma-characteristic-zero-local-smooth}
$k$를 표수 $0$인 체라 하자.
$S$를 유한 생성 $k$-대수라 하자.
$\mathfrak q \subset S$를 소 아이디얼이라 하자.
다음 조건들은 서로 동치이다.
\begin{enumerate}
\item 대수 $S$는 $k$ 위에서 $\mathfrak q$에서 매끄럽다.
\item $S_{\mathfrak q}$-가군 $\Omega_{S/k, \mathfrak q}$는 (유한) 자유이다.
\item 환 $S_{\mathfrak q}$는 정칙이다.
\end{enumerate}
\end{lemma}

\begin{proof}
표수 0에서는 모든 체의 확대가 분리이므로 (1)과 (3)의 동치는
보조정리 \ref{algebra-lemma-separable-smooth}에서 따른다.
또한 매끄러운 대수의 정의에 의해 (1)은 (2)를 함의한다.
$\Omega_{S/k, \mathfrak q}$가 $S_{\mathfrak q}$ 위에서 자유라고 하자.
보조정리 \ref{algebra-lemma-separable-smooth}의 증명에 나타난 표기와 관찰을
사용하겠다. 즉 $R = S_{\mathfrak q}$이고 극대 아이디얼은 $\mathfrak m$,
잉여체는 $\kappa$이다. 목표는 $R$가 정칙임을 증명하는 것이다.

\medskip\noindent
$\mathfrak m/\mathfrak m^2 = 0$이면 $\mathfrak m = 0$이고
$R \cong \kappa$이다. 따라서 $R$는 정칙이다.

\medskip\noindent
$\mathfrak m/ \mathfrak m^2 \not = 0$이면
$\mathfrak m/\mathfrak m^2$에서 상이 영이 아닌 임의의
$f \in \mathfrak m$를 택하자.
보조정리 \ref{algebra-lemma-computation-differential}에 의해 $\text{d}f$의
$\Omega_{R/k}/\mathfrak m\Omega_{R/k}$에서의 상은 영이 아니다.
가정에 의해 $\Omega_{R/k} = \Omega_{S/k, \mathfrak q}$는 유한 자유이므로
나카야마 보조정리 \ref{algebra-lemma-NAK}에 의해 $\text{d}f$는 직합인자를 생성한다.
보조정리 \ref{algebra-lemma-characteristic-zero}를 적용하면 $f$가 $R$의
영인자가 아님을 얻는다.
더욱이 보조정리 \ref{algebra-lemma-differential-seq}에 의해 완전열
$$
(f)/(f^2) \to \Omega_{R/k} \otimes_R R/fR \to \Omega_{(R/fR)/k} \to 0
$$
을 얻는다. 이는 $\Omega_{(R/fR)/k}$도 유한 자유임을 뜻한다.
따라서 귀납법으로 $R/fR$가 정칙 국소환임을 얻는다.
$f \in \mathfrak m$가 영인자가 아니므로 보조정리
\ref{algebra-lemma-regular-mod-x}에 의해 $R$가 정칙이라는 결론을 얻는다.
\end{proof}

\begin{example}
\label{algebra-example-characteristic-p}
보조정리 \ref{algebra-lemma-characteristic-zero-local-smooth}는 표수 $p > 0$에서는
성립하지 않는다. 표준적인 예는 환 사상
$$
\mathbf{F}_p \longrightarrow \mathbf{F}_p[x]/(x^p)
$$
이다. 그 미분 가군은 자유이지만 이 사상은 분명 매끄럽지 않다.
또 다른 예는 환 사상 ($p > 2$)
$$
\mathbf{F}_p(t) \to \mathbf{F}_p(t)[x, y]/(x^p + y^2 + \alpha)
$$
이다. 이는 소 아이디얼 $\mathfrak q = (y, x^p + \alpha)$에서 매끄럽지
않지만 정칙이다.
\end{example}

\noindent
위의 내용을 사용하면 일반점에서의 매끄러움을 체의 확대를 통해
특징지을 수 있다.

\begin{lemma}
\label{algebra-lemma-smooth-at-generic-point}
$R \to S$를 단사 유한 생성 환 사상이라 하고 $R$, $S$가 정역이라고 하자.
그러면 $R \to S$가 $\mathfrak q = (0)$에서 매끄러운 것은 유도된 분수체의
확대 $L/K$가 분리인 것과 동치이다.
\end{lemma}

\begin{proof}
$R \to S$가 $(0)$에서 매끄럽다고 하자.
어떤 영이 아닌 $g \in S$에 대해 $S$를 $S_g$로 바꾸어
$R \to S$가 매끄럽다고 가정할 수 있다.
그러면 $K \to S \otimes_R K$는 매끄럽다
(보조정리 \ref{algebra-lemma-base-change-smooth}).
더욱이 임의의 체의 확대 $K'/K$에 대해
$K' \to S \otimes_R K'$도 매끄럽다. 따라서 보조정리
\ref{algebra-lemma-characterize-smooth-over-field}에 의해
$S \otimes_R K'$는 정칙환이고, 특히 축소환이다.
따라서 $S \otimes_R K$는 $K$ 위에서 기하적으로 축소이다.
그러므로 보조정리 \ref{algebra-lemma-geometrically-reduced-permanence}에 의해
$L$은 $K$ 위에서 기하적으로 축소이다.
보조정리 \ref{algebra-lemma-characterize-separable-field-extensions}에 의해
$L/K$는 분리이다.

\medskip\noindent
반대로 $L/K$가 분리라고 하자.
보조정리 \ref{algebra-lemma-generic-finite-presentation}에 의해
$R \to S$가 유한 표시라고 가정할 수 있다.
보조정리 \ref{algebra-lemma-flat-base-change-locus-smooth}에 의해
$K \to S \otimes_R K$가 $(0)$에서 매끄러움을 보이면 충분하다.
이는 보조정리 \ref{algebra-lemma-separable-smooth}, 체가 정칙환이라는 사실,
그리고 $L/K$가 분리라는 가정에서 따른다.
\end{proof}

\section{뇌터 경우의 매끄러운 환 사상}
\label{algebra-section-smooth-Noetherian}

\begin{definition}
\label{algebra-definition-small-extension}
$\varphi : B' \to B$를 환 사상이라 하자.
$B'$와 $B$가 아르틴 국소환이고, $\varphi$가 전사이며,
$I = \Ker(\varphi)$가 $B'$-가군으로서 길이 $1$을 가지면
$\varphi$를 {\it 작은 확대}라고 한다.
\end{definition}

\noindent
이는 분명 $I^2 = 0$이고, $\mathfrak m' \subset B'$가 극대 아이디얼일 때
$\mathfrak m' x = 0$을 만족하는
어떤 $x \in B'$에 대해 $I = (x)$라는 뜻이다.

\begin{lemma}
\label{algebra-lemma-smooth-test-artinian}
% P01-E0865: 동결 원문의 두 번째 병렬 술어에는 동사 ``is''가 빠져 있다.
$R \to S$를 환 사상이라 하자. $\mathfrak q$를 $S$의 소 아이디얼이라 하고,
$\mathfrak p \subset R$ 위에 놓인다고 하자. $R$가 뇌터이고
$R \to S$가 유한 생성이라고 가정하자.
다음은 서로 동치이다.
\begin{enumerate}
\item $R \to S$는 $\mathfrak q$에서 매끄럽다.
\item 핵 $\Ker(B' \to B)$의 제곱이 영인 모든 국소 $R$-대수의 전사
$(B', \mathfrak m') \to (B, \mathfrak m)$와 모든 실선 가환 그림
$$
\xymatrix{
S \ar[r] \ar@{-->}[rd] & B \\
R \ar[r] \ar[u] & B' \ar[u]
}
$$
% P01-E0866: 동결 원문의 교집합 표기는 일반적인 비단사 사상에 맞지 않으며 역상을 써야 한다.
에 대해, $\mathfrak q = S \cap \mathfrak m$이면 그림을 가환이 되게 하는
점선 화살표가 존재한다.
\item (2)와 같되 $B' \to B$는 작은 확대만을 돈다.
\item (2)와 같되 $B' \to B$는 작은 확대만을 돌고, 추가로 $S \to B$가
동형사상 $\kappa(\mathfrak q) \cong \kappa(\mathfrak m)$을 유도한다.
\end{enumerate}
\end{lemma}

\begin{proof}
(1)을 가정하자. 이는 어떤 $g \in S$, $g \not \in \mathfrak q$에 대해
$R \to S_g$가 매끄럽다는 뜻이다. 명제
\ref{algebra-proposition-smooth-formally-smooth}에 의해 $R \to S_g$는 형식적으로
매끄럽다. (2)와 같은 임의의 그림에서 사상 $S \to B$는 자동으로
$S_{\mathfrak q}$를 통해, 따라서 더욱이 $S_g$를 통해 인수분해된다.
$R$ 위에서 $S_g$의 형식적 매끄러움은 왼쪽 위 꼭짓점에 $S_g$가 놓인
비슷한 그림에 맞는 사상 $S_g \to B'$를 준다. 이를 $S \to S_g$와 합성하면
원하는 화살표를 얻는다. 즉 (1)이 (2)를 함의함을 보였다.

\medskip\noindent
(2)가 (3)을 함의하고 (3)이 (4)를 함의함은 분명하다.

\medskip\noindent
(4)를 가정하자. (1)이 성립함을 보여 보조정리의 증명을 마치겠다.
표시
$S = R[x_1, \ldots, x_n]/(f_1, \ldots, f_m)$
를 택하자. $S$는 $R$ 위에서 유한 생성이고, 따라서 유한 표시이므로
(보조정리 \ref{algebra-lemma-Noetherian-finite-type-is-finite-presentation} 참조)
이는 가능하다. $I = (f_1, \ldots, f_m)$로 두자.
이 표시의 소박한 여접 복합체
$$
\text{d} : I/I^2
\longrightarrow
\bigoplus\nolimits_{j = 1}^n S\text{d}x_j
$$
를 생각하자(절 \ref{algebra-section-netherlander} 참조).
이 복합체를 $\mathfrak q$에서 국소화했을 때 사상이 분할 단사가 됨을
보이면 충분하다. 보조정리 \ref{algebra-lemma-smooth-at-point}를 보라.
% P01-E0867: 동결 원문의 ``Denote'' 구문에는 by/as가 빠져 있다.
$S' = R[x_1, \ldots, x_n]/I^2$로 두자.
보조정리 \ref{algebra-lemma-differential-mod-power-ideal}에 의해
% P01-E0868: 아래 동결 수식의 상한 m은 변수 개수 n이어야 하지만 그대로 보존한다.
$$
S \otimes_{S'} \Omega_{S'/R} =
S \otimes_{R[x_1, \ldots, x_n]} \Omega_{R[x_1, \ldots, x_n]/R} =
\bigoplus\nolimits_{j = 1}^m S\text{d}x_j.
$$
따라서 사상
$$
\text{d} :
I/I^2
\longrightarrow
S \otimes_{S'} \Omega_{S'/R}
$$
은 위 소박한 여접 복합체의 사상과 같다. 특히 증명하려는 명제의 진위는
세 환 $R \to S' \to S$에만 의존한다.
$\mathfrak q$에 대응하는 소 아이디얼을
$\mathfrak q' \subset R[x_1, \ldots, x_n]$라 하자.
국소화는 미분 가군을 취하는 것과 가환하므로
(보조정리 \ref{algebra-lemma-differentials-localize}),
다음 사상이
\begin{equation}
\label{algebra-equation-target-map}
\text{d} :
I_{\mathfrak q'}/I_{\mathfrak q'}^2
\longrightarrow
S_{\mathfrak q} \otimes_{S'_{\mathfrak q'}} \Omega_{S'_{\mathfrak q'}/R}
\end{equation}
은 $R \to S'_{\mathfrak q'} \to S_{\mathfrak q}$에서 오며,
분할 단사임을 보이면 충분하다.

\medskip\noindent
$N \in \mathbf{N}$을 정수라 하자. 환
$$
B'_N = S'_{\mathfrak q'} / (\mathfrak q')^N S'_{\mathfrak q'}
= (S'/(\mathfrak q')^N S')_{\mathfrak q'}
$$
와 그 몫 $B_N = B'_N/IB'_N$을 생각하자.
$B_N \cong S_{\mathfrak q}/\mathfrak q^NS_{\mathfrak q}$임에 유의하라.
$B'_N$은 국소 뇌터 환을 그 극대 아이디얼의 거듭제곱으로 나눈 몫이므로
아르틴 국소환이다. $B'_N$-부분가군들로 된 $B'_N \to B_N$의 핵
$I_N$의 여과
$$
0 \subset J_{N, 1} \subset J_{N, 2} \subset \ldots \subset J_{N, n(N)} = I_N
$$
를 택하되, 각 연속 몫 $J_{N, i}/J_{N, i - 1}$의 길이가 $1$이 되게 하자.
($B'_N$이 아르틴이므로 그러한 여과가 존재한다.) 이는 작은 확대들의 열
$$
B'_N  \to B'_N/J_{N, 1} \to B'_N/J_{N, 2} \to \ldots \to
B'_N/J_{N, n(N)} = B'_N/I_N
= B_N = S_{\mathfrak q}/\mathfrak q^NS_{\mathfrak q}
$$
을 준다. 사상 $S \to B_N$에서 시작하여 이 작은 확대들에 조건 (4)를
차례로 적용하면 가환 그림
% P01-E0869: 동결 원문의 절에는 접속어 ``that''이 빠져 있다.
$$
\xymatrix{
S \ar[r] \ar[rd] & B_N  \\
R \ar[r] \ar[u] & B'_N \ar[u]
}
$$
이 존재함을 알 수 있다. 환 사상 $S \to B'_N$은 분명
$S \to S_{\mathfrak q} \to B'_N$로 인수분해되며, 여기서
$S_{\mathfrak q} \to B'_N$은 국소환들의 국소 준동형사상이다.
더욱이 $B'_N$의 극대 아이디얼의 $N$제곱이 영이므로
$S_{\mathfrak q} \to B'_N$은
$S_{\mathfrak q}/(\mathfrak q)^NS_{\mathfrak q} = B_N$을 통해
인수분해된다. 즉 모든 $N \in \mathbf{N}$에 대해 $R$-대수의 전사
$B'_N \to B_N$에는 단면이 있음을 보였다.

\medskip\noindent
핵이 $I_N$인 전사 $B'_N \to B_N$에서 오는 표시
$$
I_N \to B_N \otimes_{B'_N} \Omega_{B'_N/R} \to \Omega_{B_N/R} \to 0
$$
를 생각하자(보조정리 \ref{algebra-lemma-differential-seq} 참조).
위에서 $R$-대수 사상 $B'_N \to B_N$에는 우역이 있음을 보였다.
따라서 보조정리 \ref{algebra-lemma-differential-seq-split}에 의해 위 열은
분할 완전열이다. 그러므로 모든 $N$에 대해 사상
$$
I_N \longrightarrow B_N \otimes_{B'_N} \Omega_{B'_N/R}
$$
은 분할 단사이다. 증명의 나머지는 이것이 정확히 무엇을 뜻하는지
풀어 쓰면 얻어진다. 다음에 유의하라.
$$
I_N = I_{\mathfrak q'}/
(I_{\mathfrak q'}^2 + (\mathfrak q')^N \cap I_{\mathfrak q'})
$$
아르틴--리스 보조정리(보조정리 \ref{algebra-lemma-Artin-Rees})에 의해
모든 $N \geq c$에 대해
$$
S_{\mathfrak q}/\mathfrak q^{N - c}S_{\mathfrak q}
\otimes_{S_{\mathfrak q}} I_N =
S_{\mathfrak q}/\mathfrak q^{N - c}S_{\mathfrak q}
\otimes_{S_{\mathfrak q}}
I_{\mathfrak q'}/I_{\mathfrak q'}^2
$$
가 되게 하는 $c \geq 0$이 존재한다
% P01-E0870: 동결 원문의 ``these tensor product are''에는 수 일치 오류가 있다.
(이 텐서곱들은 단지 $\mathfrak q^{N - c}$로 나누는 것을 화려하게 쓴 것이다).
물론 $c \geq 1$이라고 가정해도 된다.
보조정리 \ref{algebra-lemma-differential-mod-power-ideal}에 의해
$$
S'_{\mathfrak q'}/(\mathfrak q')^{N - c}S'_{\mathfrak q'}
\otimes_{S'_{\mathfrak q'}} \Omega_{B'_N/R} =
S'_{\mathfrak q'}/(\mathfrak q')^{N - c}S'_{\mathfrak q'}
\otimes_{S'_{\mathfrak q'}} \Omega_{S'_{\mathfrak q'}/R}
$$
% P01-E0871: 동결 원문은 위 수식 뒤 문장 경계가 깨져 있다.
임을 안다. 여기에 $B_N = S_{\mathfrak q}/\mathfrak q^N$을 다시 텐서하면
$$
S_{\mathfrak q}/\mathfrak q^{N - c}S_{\mathfrak q}
\otimes_{S'_{\mathfrak q'}} \Omega_{B'_N/R} =
S_{\mathfrak q}/\mathfrak q^{N - c}S_{\mathfrak q}
\otimes_{S'_{\mathfrak q'}} \Omega_{S'_{\mathfrak q'}/R}.
$$
분할 단사는 무엇을 텐서해도 여전히 분할 단사이므로, 모든 $N \geq c$에 대해
$$
S_{\mathfrak q}/\mathfrak q^{N - c}S_{\mathfrak q}
\otimes_{S_{\mathfrak q}}
(\ref{algebra-equation-target-map}) =
S_{\mathfrak q}/\mathfrak q^{N - c}S_{\mathfrak q}
\otimes_{S_{\mathfrak q}/\mathfrak q^N S_{\mathfrak q}}
(I_N \longrightarrow B_N \otimes_{B'_N} \Omega_{B'_N/R})
$$
은 분할 단사이다. 보조정리
\ref{algebra-lemma-split-injection-after-completion}에 의해
(\ref{algebra-equation-target-map})은 분할 단사이다. 이로써 증명이 끝난다.
\end{proof}








\section{매끄러운 환 사상에 관한 결과의 개관}
\label{algebra-section-smooth-overview}

\noindent
앞 절들에서 증명한 매끄러운 환 사상에 관한 결과들을 열거한다.
더 정확한 명제와 정의는 제시된 참조를 보라.
\begin{enumerate}
\item 환 사상 $R \to S$가 매끄럽다는 것은 유한 표시이고 $S/R$의
소박한 여접 복합체가 차수 $0$의 유한 사영 $S$-가군과 준동형이라는
뜻이다. 정의 \ref{algebra-definition-smooth}를 보라.
\item $S$가 $R$ 위에서 매끄러우면 $\Omega_{S/R}$은 유한 사영
$S$-가군이다. 정의 \ref{algebra-definition-smooth} 뒤의 논의를 보라.
\item 매끄럽다는 성질은 $S$ 위에서 국소적이다. 보조정리
\ref{algebra-lemma-locally-smooth}를 보라.
\item 매끄럽다는 성질은 밑변환 아래에서 안정적이다. 보조정리
\ref{algebra-lemma-base-change-smooth}를 보라.
\item 매끄럽다는 성질은 합성 아래에서 안정적이다. 보조정리
\ref{algebra-lemma-compose-smooth}를 보라.
\item 매끄러운 환 사상은 신토믹이고, 특히 평탄하다. 보조정리
\ref{algebra-lemma-smooth-syntomic}을 보라.
\item 매끄러운 섬유환을 갖는 유한 표시 평탄 환 사상은 매끄럽다.
보조정리 \ref{algebra-lemma-flat-fibre-smooth}를 보라.
\item 유한 표시 환 사상 $R \to S$가 매끄러울 필요충분조건은
형식적으로 매끄러운 것이다. 명제
\ref{algebra-proposition-smooth-formally-smooth}를 보라.
\item $R \to S$가 유한 생성 환 사상이고 $R$가 뇌터이면,
$R \to S$의 매끄러움을 확인하기 위해서는 아르틴 국소환들의 작은 확대를
따르는 형식적 매끄러움의 올림 성질만 확인하면 충분하다. 보조정리
\ref{algebra-lemma-smooth-test-artinian}을 보라.
\item 매끄러운 환 사상 $R \to S$는 $\mathbf{Z}$ 위에서 유한 생성인
$R_0$를 갖는 매끄러운 환 사상 $R_0 \to S_0$의 밑변환이다.
보조정리 \ref{algebra-lemma-finite-presentation-fs-Noetherian}을 보라.
\item 환 사상이 매끄러운 점들의 집합을 취하는 것은 평탄 밑변환과
가환한다. 보조정리 \ref{algebra-lemma-flat-base-change-locus-smooth}를 보라.
% P01-E0872: 동결 원문은 술어와 목록 앞 콜론을 빠뜨린다.
\item $S$가 대수적으로 닫힌 체 $k$ 위에서 유한 생성이고,
$\mathfrak m \subset S$가 극대 아이디얼이면 다음은 서로 동치이다.
\begin{enumerate}
\item $S$는 $\mathfrak m$의 근방에서 $k$ 위에 매끄럽다.
\item $S_{\mathfrak m}$은 정칙 국소환이다.
% P01-E0873: 아래 동결 수식은 정의된 mathfrak m 대신 평문 m을 첨자로 쓴다.
\item $\dim(S_{\mathfrak m}) =
\dim_{\kappa(m)} \Omega_{S/k} \otimes_S \kappa(\mathfrak m)$.
\end{enumerate}
보조정리 \ref{algebra-lemma-characterize-smooth-kbar}를 보라.
% P01-E0874: 동결 원문의 세 개 항목에는 소 아이디얼 술어의 동사가 빠져 있다.
\item $S$가 체 $k$ 위에서 유한 생성이고 $\mathfrak q \subset S$가
소 아이디얼이면 다음은 서로 동치이다.
\begin{enumerate}
\item $S$는 $\mathfrak q$의 근방에서 $k$ 위에 매끄럽다.
\item $\dim_{\mathfrak q}(S/k) =
\dim_{\kappa(\mathfrak q)} \Omega_{S/k} \otimes_S \kappa(\mathfrak q)$.
\end{enumerate}
보조정리 \ref{algebra-lemma-characterize-smooth-over-field}를 보라.
\item $S$가 체 위에서 매끄러우면 모든 국소환은 정칙이다.
보조정리 \ref{algebra-lemma-characterize-smooth-over-field}를 보라.
\item $S$가 체 $k$ 위에서 유한 생성이고, $\mathfrak q \subset S$가
소 아이디얼이며, 체의 확대 $\kappa(\mathfrak q)/k$가 분리이고
$S_{\mathfrak q}$가 정칙이면, $S$는 $\mathfrak q$에서 $k$ 위에 매끄럽다.
보조정리 \ref{algebra-lemma-separable-smooth}를 보라.
\item $S$가 체 $k$ 위에서 유한 생성이고, $k$의 표수가 $0$이며,
$\mathfrak q \subset S$가 소 아이디얼이고,
$\Omega_{S/k, \mathfrak q}$가 자유이면, $S$는 $\mathfrak q$에서
$k$ 위에 매끄럽다. 보조정리
\ref{algebra-lemma-characteristic-zero-local-smooth}를 보라.
\end{enumerate}
이 결과들 중 일부는 표준 매끄러운 환 사상의 개념을 사용해 증명했다.
정의 \ref{algebra-definition-standard-smooth}를 보라. 이는 신토믹 사상의 경우에
상대 전역 완전교차 사상이 맡는 역할과 유사하다. 또한 예를 만드는 가장
쉬운 방법이다.













\section{에탈 환 사상}
\label{algebra-section-etale}

\noindent
에탈 환 사상은 상대 차원이 영인 매끄러운 환 사상이다. 이는 다음의
조금 더 직접적인 정의와 같다.

\begin{definition}
\label{algebra-definition-etale}
$R \to S$를 환 사상이라 하자. $R \to S$가 유한 표시이고 소박한 여접
복합체 $\NL_{S/R}$가 영 복합체와 준동형이면, 즉
$H_1(\NL_{S/R}) = 0$이고 $\Omega_{S/R} = 0$이면 이를
{\it 에탈}이라고 한다. $S$의 소 아이디얼 $\mathfrak q$가 주어졌을 때,
어떤 $g \in S$, $g \not \in \mathfrak q$에 대해 $R \to S_g$가
에탈이면 $R \to S$가 {\it $\mathfrak q$에서 에탈}이라고 한다.
\end{definition}

\noindent
특히 $S$가 $R$ 위에서 에탈이면 $\Omega_{S/R} = 0$임을 알 수 있다.
$R \to S$가 매끄러우면, $R \to S$가 에탈일 필요충분조건은
$\Omega_{S/R} = 0$인 것이다. 매끄러운 환 사상에 관한 결과들로부터
에탈 사상에 관한 많은 결과를 자동으로 얻는다. 이를 아래 보조정리
\ref{algebra-lemma-etale}에 요약한다. 그 전에 {\it 모든} 에탈 환 사상이
표준 매끄러움임을 증명한다.

\begin{lemma}
\label{algebra-lemma-etale-standard-smooth}
모든 에탈 환 사상은 표준 매끄럽다. 더 정확히, $R \to S$가 에탈이면
상의 $\det(\partial f_j/\partial x_i)$가 $S$에서 가역원이 되게 하는 표시
$S = R[x_1, \ldots, x_n]/(f_1, \ldots, f_n)$가 존재한다.
\end{lemma}

\begin{proof}
$R \to S$를 에탈이라 하자. 표시 $S = R[x_1, \ldots, x_n]/I$를 택하자.
$R \to S$가 에탈이므로
$$
\text{d} :
I/I^2
\longrightarrow
\bigoplus\nolimits_{i = 1, \ldots, n} S\text{d}x_i
$$
는 동형사상이고, 특히 $I/I^2$는 자유 $S$-가군이다. 따라서 보조정리
\ref{algebra-lemma-huber}에 의해 (필요하면 표시를 바꾼 뒤)
$I = (f_1, \ldots, f_c)$이고 류 $f_i \bmod I^2$들이 $I/I^2$의 기저를
이룬다고 가정할 수 있다. 위 표시 사상이 동형사상이라는 사실로부터
즉시 $c = n$이고 $\det(\partial f_j/\partial x_i)$가 $S$에서
가역원임을 얻는다.
\end{proof}

\begin{lemma}
\label{algebra-lemma-etale}
에탈 환 사상에 관해 다음이 성립한다.
\begin{enumerate}
\item 임의의 환 $R$와 임의의 $f \in R$에 대해 환 사상
$R \to R_f$는 에탈이다.
\item 에탈 환 사상들의 합성은 에탈이다.
\item 에탈 환 사상의 밑변환은 에탈이다.
\item 에탈이라는 성질은 국소적이다. 환 사상 $R \to S$와 단위 아이디얼을
생성하는 원소 $g_1, \ldots, g_m \in S$가 주어지고,
$j = 1, \ldots, m$마다 $R \to S_{g_j}$가 에탈이면 $R \to S$도 에탈이다.
\item 유한 표시 사상 $R \to S$와 평탄 환 사상 $R \to R'$가 주어졌다고
하고 $S' = R' \otimes_R S$로 두자. $R' \to S'$가 에탈인 소 아이디얼들의
집합은 $\Spec(S') \to \Spec(S)$에 의한, $R \to S$가 에탈인 소
아이디얼들의 집합의 역상이다.
\item 에탈 환 사상은 신토믹이고, 특히 평탄하다.
\item $S$가 체 $k$ 위에서 유한 생성이면, $S$가 $k$ 위에서 에탈일
필요충분조건은 $\Omega_{S/k} = 0$인 것이다.
\item 임의의 에탈 환 사상 $R \to S$는 $\mathbf{Z}$ 위에서 유한 생성인
$R_0$를 갖는 에탈 환 사상 $R_0 \to S_0$의 밑변환이다.
\item $A = \colim A_i$를 환들의 여과 쌍극한이라 하자.
$A \to B$를 에탈 환 사상이라 하자. 그러면 어떤 $i$에 대해
$B \cong A \otimes_{A_i} B_i$가 되게 하는 에탈 환 사상
$A_i \to B_i$가 존재한다.
\item $A$를 환, $S$를 $A$의 곱셈적 부분집합이라 하자.
$S^{-1}A \to B'$가 에탈이라고 하자. 그러면
$B' \cong S^{-1}B$가 되게 하는 에탈 환 사상 $A \to B$가 존재한다.
\item $A$를 환이라 하고 $B = B' \times B''$를 $A$-대수들의 곱이라 하자.
그러면 $B$가 $A$ 위에서 에탈일 필요충분조건은 $B'$와 $B''$가 모두
$A$ 위에서 에탈인 것이다.
\end{enumerate}
\end{lemma}

\begin{proof}
각 경우에 매끄러운 환 사상에 관한 대응하는 결과를 사용하고,
$\Omega_{S/R}$가 영임을 보이는 짧은 논증을 덧붙인다.

\medskip\noindent
(1)의 증명. 환 사상 $R \to R_f$는 매끄럽고 $\Omega_{R_f/R} = 0$이다.

\medskip\noindent
(2)의 증명. 매끄러운 사상 $A \to B$와 $B \to C$의 합성
$A \to C$는 매끄럽다. 보조정리 \ref{algebra-lemma-compose-smooth}를 보라.
보조정리 \ref{algebra-lemma-exact-sequence-differentials}에 의해
$\Omega_{C/B}$와 $\Omega_{B/A}$가 모두 영이므로 $\Omega_{C/A}$도
영이다.

\medskip\noindent
(3)의 증명. $R \to S$가 에탈이고 $R \to R'$가 임의의 사상이라고 하자.
보조정리 \ref{algebra-lemma-base-change-smooth}에 의해
$R' \to S' = R' \otimes_R S$는 매끄럽다. 보조정리
\ref{algebra-lemma-differentials-base-change}에 의해
$\Omega_{S'/R'} = S' \otimes_S \Omega_{S/R}$이므로
$\Omega_{S'/R'} = 0$이다. 따라서 $R' \to S'$는 에탈이다.

\medskip\noindent
(4)의 증명. (4)의 가정을 만족한다고 하자. 보조정리
\ref{algebra-lemma-locally-smooth}에 의해 $R \to S$는 매끄럽다.
또 모든 $i$에 대해
$\Omega_{S_{g_i}/R} = (\Omega_{S/R})_{g_i} = 0$이다.
그러면 보조정리 \ref{algebra-lemma-cover}에 의해 $\Omega_{S/R} = 0$이다.

\medskip\noindent
(5)의 증명. 매끄러운 사상에 관한 결과는 보조정리
\ref{algebra-lemma-flat-base-change-locus-smooth}이다. 그 보조정리의 증명에서는
$\NL_{S/R} \otimes_S S'$가 $\NL_{S'/R'}$와 호모토피 동치임을 사용했다.
따라서 다음을 보이면 된다. $M$이 유한 표시 $S$-가군일 때,
$(M \otimes_S S')_{\mathfrak q'} = 0$인 $S'$의 소 아이디얼
$\mathfrak q'$들의 집합은 $M_{\mathfrak q} = 0$인 $S$의 소 아이디얼
$\mathfrak q$들의 집합의 역상이다. 이는 보조정리
\ref{algebra-lemma-support-base-change}에서 따른다.

\medskip\noindent
(6)의 증명. 매끄러운 환 사상에 관한 대응하는 결과(보조정리
\ref{algebra-lemma-smooth-syntomic})에서 바로 따른다.

\medskip\noindent
(7)의 증명. 보조정리 \ref{algebra-lemma-characterize-smooth-over-field}와
정의들에서 따른다.

\medskip\noindent
(8)의 증명. 보조정리 \ref{algebra-lemma-finite-presentation-fs-Noetherian}은
매끄러운 환 사상에 관한 결과를 준다. 그 결과 얻은 매끄러운 환 사상
$R_0 \to S_0$는 보조정리 \ref{algebra-lemma-relative-dimension-CM}의 가정을
만족하므로, $S_0$를 $R_0$ 위 상대 차원 $0$인 인자로 바꿀 수 있다.

\medskip\noindent
(9)의 증명. 보조정리 \ref{algebra-lemma-characterize-finite-presentation}에 의해
$R_0 \to A$는 어떤 $i$에 대해 $A_i$를 통해 인수분해되므로 (8)에서 따른다.

\medskip\noindent
(10)의 증명. $S^{-1}A$가 $A$의 주 국소화들의 여과 쌍극한이므로
(9), (1), (2)에서 따른다.

\medskip\noindent
(11)의 증명. 보조정리 \ref{algebra-lemma-product-smooth}를 사용해 매끄러움에 관한
결과를 얻은 다음, $\Omega_{B/A}$가 영일 필요충분조건이
$\Omega_{B'/A}$와 $\Omega_{B''/A}$가 모두 영인 것임을 사용한다.
\end{proof}

\noindent
이제 체 위에서 에탈이라는 것이 무엇을 뜻하는지 더 자세히 살펴본다.

\begin{lemma}
\label{algebra-lemma-etale-over-field}
$k$를 체라 하자. 환 사상 $k \to S$가 에탈일 필요충분조건은
$S$가 $k$-대수로서 $k$의 유한 분리 확대들의 유한 곱과 동형인 것이다.
\end{lemma}

\begin{proof}
다음 사실을 더 언급하지 않고 사용하겠다. $S = S_1 \times \ldots \times S_n$가
$k$-대수들의 유한 곱이면, $S$가 $k$ 위에서 에탈일 필요충분조건은 각
$S_i$가 $k$ 위에서 에탈인 것이다. 보조정리 \ref{algebra-lemma-etale}의 (11)을 보라.

\medskip\noindent
$k'/k$가 유한 분리 체의 확대이면
$k' = k(\alpha) \cong k[x]/(f)$로 쓸 수 있다. 여기서 $f$는 원소
$\alpha$의 최소다항식이다. $k'$가 $k$ 위에서 분리이므로
$\gcd(f, f') = 1$이다. 이는
$\text{d} : k'\cdot f \to k' \cdot \text{d}x$가 동형사상임을 함의한다.
따라서 $k \to k'$는 에탈이다. 그러므로 $S$가 $k$의 유한 분리 확대들의
유한 곱이면 $S$는 $k$ 위에서 에탈이다.

\medskip\noindent
반대로 $k \to S$가 에탈이라고 하자. 그러면 $S$는 $k$ 위에서 매끄럽고
$\Omega_{S/k} = 0$이다. 보조정리
\ref{algebra-lemma-characterize-smooth-over-field}에 의해 $S$의 모든 극대 아이디얼
$\mathfrak m$에 대해 $\dim_\mathfrak m \Spec(S) = 0$이다. 따라서
$\dim(S) = 0$이다. 명제 \ref{algebra-proposition-dimension-zero-ring}에 의해
$S$는 아르틴 국소환들의 유한 곱이다. 다시 보조정리
\ref{algebra-lemma-characterize-smooth-over-field}에 의해 이 국소환들은 체이다.
따라서 $S = k'$가 체라고 가정해도 된다. 힐베르트 영점정리(정리
\ref{algebra-theorem-nullstellensatz})에 의해 확대 $k'/k$는 유한하다.
$k \to k'$의 매끄러움과 보조정리
\ref{algebra-lemma-smooth-at-generic-point}에 의해 $k'/k$는 분리 확대이다.
이로써 증명이 끝난다.
\end{proof}

\begin{lemma}
\label{algebra-lemma-etale-at-prime}
$R \to S$를 환 사상이라 하자. $\mathfrak q \subset S$를
$R$의 $\mathfrak p$ 위에 놓이는 소 아이디얼이라 하자.
$S/R$이 $\mathfrak q$에서 에탈이면
\begin{enumerate}
\item $\mathfrak p S_{\mathfrak q} = \mathfrak qS_{\mathfrak q}$는 국소환
$S_{\mathfrak q}$의 극대 아이디얼이고,
\item 체의 확대 $\kappa(\mathfrak q)/\kappa(\mathfrak p)$는 유한 분리이다.
\end{enumerate}
\end{lemma}

\begin{proof}
먼저 어떤 $g \in S$, $g \not \in \mathfrak q$에 대해 $S$를 $S_g$로
바꾸고 $R \to S$가 에탈이라고 가정할 수 있다. 그러면 보조정리
\ref{algebra-lemma-etale-over-field}와 $S \otimes_R \kappa(\mathfrak p)$가
$\kappa(\mathfrak p)$ 위에서 에탈이라는 사실을 풀어 쓰면 결론이 따른다.
\end{proof}

\begin{lemma}
\label{algebra-lemma-etale-quasi-finite}
에탈 환 사상은 준유한이다.
\end{lemma}

\begin{proof}
$R \to S$를 에탈 환 사상이라 하자. 정의에 의해 $R \to S$는 유한 생성이다.
임의의 소 아이디얼 $\mathfrak p \subset R$에 대해 섬유환
$S \otimes_R \kappa(\mathfrak p)$는 $\kappa(\mathfrak p)$ 위에서
에탈이고, 따라서 $\kappa(\mathfrak p)$의 유한 분리 확대들인 체들의
% P01-E0875: 동결 원문에는 단복수 불일치와 연결어 누락이 있다.
유한 곱이다. 특히 $\kappa(\mathfrak p)$ 위에서 유한하다.
그러므로 보조정리 \ref{algebra-lemma-quasi-finite}에 의해 $R \to S$는 준유한이다.
\end{proof}

\begin{lemma}
\label{algebra-lemma-characterize-etale}
$R \to S$를 환 사상이라 하자. $\mathfrak q$를 $S$의 소 아이디얼,
$\mathfrak p$를 $R$의 소 아이디얼이라 하고, 두 소 아이디얼이 서로
대응한다고 하자. 다음을 가정하자.
\begin{enumerate}
\item $R \to S$는 유한 표시이다.
\item $R_{\mathfrak p} \to S_{\mathfrak q}$는 평탄하다.
\item $\mathfrak p S_{\mathfrak q}$는 국소환 $S_{\mathfrak q}$의
극대 아이디얼이다.
\item 체의 확대 $\kappa(\mathfrak q)/\kappa(\mathfrak p)$는 유한 분리이다.
\end{enumerate}
그러면 $R \to S$는 $\mathfrak q$에서 에탈이다.
\end{lemma}

\begin{proof}
보조정리 \ref{algebra-lemma-isolated-point-fibre}를 적용하여 $\mathfrak q$가
$\mathfrak p$ 위에 놓이는 $S_g$의 유일한 소 아이디얼이 되게 하는
$g \in S$, $g \not \in \mathfrak q$를 택하자. $S$를 $S_g$로
바꾸기로 한다. 그러면 $S \otimes_R \kappa(\mathfrak p)$는 유일한
소 아이디얼을 가지므로 국소환이고, 따라서
$S_{\mathfrak q}/\mathfrak pS_{\mathfrak q}
\cong \kappa(\mathfrak q)$와 같다.
보조정리 \ref{algebra-lemma-flat-fibre-smooth}에 의해 $R \to S_g$가 매끄럽게
되는 $g \in S$, $g \not \in \mathfrak q$가 존재한다.
% P01-E0876: 동결 원문의 치환과 결론 사이에는 접속어나 문장부호가 빠져 있다.
$S$를 다시 $S_g$로 바꾸면 $R \to S$가 매끄럽다고 가정할 수 있다.
보조정리 \ref{algebra-lemma-smooth-syntomic}에 의해 나아가 $R \to S$가
표준 매끄럽고
$S = R[x_1, \ldots, x_n]/(f_1, \ldots, f_c)$라고 가정해도 된다.
$S \otimes_R \kappa(\mathfrak p) = \kappa(\mathfrak q)$의 차원이 $0$이므로
$n = c$이다. 즉 $R \to S$는 에탈이다.
\end{proof}

\noindent
여기에는 전혀 새로운 현상이 있다.

\begin{lemma}
\label{algebra-lemma-map-between-etale}
$R \to S$와 $R \to S'$가 에탈이라고 하자. 그러면 임의의
$R$-대수 사상 $S' \to S$는 에탈이다.
\end{lemma}

\begin{proof}
우선 보조정리 \ref{algebra-lemma-compose-finite-type}에 의해 $S' \to S$는
유한 표시이다. $\mathfrak q \subset S$를
$\mathfrak q' \subset S'$와 $\mathfrak p \subset R$ 위에 놓이는
소 아이디얼이라 하자. 보조정리 \ref{algebra-lemma-etale-at-prime}에 의해 환 사상
$S'_{\mathfrak q'}/\mathfrak p S'_{\mathfrak q'} \to
S_{\mathfrak q}/\mathfrak p S_{\mathfrak q}$
은 $\kappa(\mathfrak p)$의 유한 분리 확대들 사이의 사상이다.
특히 이는 평탄하다. 따라서 보조정리
\ref{algebra-lemma-criterion-flatness-fibre}에 의해
$S'_{\mathfrak q'} \to S_{\mathfrak q}$는 평탄하다. 그러므로
$S' \to S$는 평탄하다. 더욱이 위의 사실로부터
$\mathfrak q'S_{\mathfrak q}$가 $S_{\mathfrak q}$의 극대 아이디얼이고,
$S'_{\mathfrak q'} \to S_{\mathfrak q}$의 잉여체 확대가 유한 분리임도
알 수 있다. 따라서 보조정리 \ref{algebra-lemma-characterize-etale}에 의해
$S' \to S$는 $\mathfrak q$에서 에탈이다. 에탈이라는 성질은 국소적이므로
(보조정리 \ref{algebra-lemma-etale} 참조), 결론을 얻는다.
\end{proof}

\begin{lemma}
\label{algebra-lemma-surjective-flat-finitely-presented}
% P01-E0877: 동결 원문의 존재문은 단수 명사구와 동사의 수가 일치하지 않는다.
$\varphi : R \to S$를 환 사상이라 하자. $R \to S$가 전사이고 평탄하며
유한 표시이면, $S = R_e$가 되게 하는 멱등원 $e \in R$가 존재한다.
\end{lemma}

\begin{proof}[첫째 증명]
$I$를 $\varphi$의 핵이라 하자. $\varphi$가 유한 표시이므로 보조정리
\ref{algebra-lemma-finite-presentation-independent}에 의해 $I$는 유한 생성이다.
또한 $S$가 $R$ 위에서 평탄하므로 완전열 $0 \to I \to R \to S \to 0$에
$R$ 위에서 $S$를 텐서하면 $I/I^2 = 0$을 얻는다. 이제 보조정리
\ref{algebra-lemma-ideal-is-squared-union-connected}에 의해 결론을 얻는다.
\end{proof}

\begin{proof}[둘째 증명]
$\Spec(S) \to \Spec(R)$은 닫힌 부분집합 위의 위상동형사상이고
(보조정리 \ref{algebra-lemma-spec-closed} 참조), 열린 사상이므로
(명제 \ref{algebra-proposition-fppf-open} 참조), 상은 어떤 멱등원 $e \in R$에
대한 $D(e)$이다(보조정리 \ref{algebra-lemma-disjoint-decomposition} 참조).
따라서 $R_e \to S$는 스펙트럼들 사이의 전단사를 유도한다.
예를 들어 보조정리 \ref{algebra-lemma-finite-flat-local}과 \ref{algebra-lemma-NAK}에 의해
이 사상은 모든 국소환에서 동형사상을 유도한다. 그러면 예를 들어
보조정리 \ref{algebra-lemma-characterize-zero-local}에 의해 $R_e \to S$는
단사이기도 하다.
\end{proof}

\begin{lemma}
\label{algebra-lemma-lift-etale}
\begin{slogan}
에탈 환 사상은 환의 전사를 따라 올릴 수 있다
\end{slogan}
$R$을 환, $I \subset R$을 아이디얼이라 하자.
$R/I \to \overline{S}$를 에탈 환 사상이라 하자. 그러면
$R/I$-대수로서 $\overline{S} \cong S/IS$가 되게 하는 에탈 환 사상
$R \to S$가 존재한다.
\end{lemma}

\begin{proof}
보조정리 \ref{algebra-lemma-etale-standard-smooth}에 의해 정의
\ref{algebra-definition-standard-smooth}에서와 같이
$\overline{S} =
(R/I)[x_1, \ldots, x_n]/(\overline{f}_1, \ldots, \overline{f}_n)$
로 쓸 수 있고,
$\overline{\Delta} =
\det(\frac{\partial \overline{f}_i}{\partial x_j})_{i, j = 1, \ldots, n}$
는 $\overline{S}$에서 가역원이다. $f_i$들을 임의로 올리고, 예
\ref{algebra-example-make-standard-smooth}에서처럼
$\Delta = \det(\frac{\partial f_i}{\partial x_j})_{i, j = 1, \ldots, n}$로
두어
$S = R[x_1, \ldots, x_n, x_{n+1}]/(f_1, \ldots, f_n, x_{n + 1}\Delta - 1)$
로 놓으면 된다. 이로써 보조정리가 증명된다.
\end{proof}

\begin{lemma}
\label{algebra-lemma-lift-etale-infinitesimal}
완전한 행들을 갖는 가환 그림
$$
\xymatrix{
0 \ar[r] &
J \ar[r] &
B' \ar[r] &
B \ar[r] & 0 \\
0 \ar[r] &
I \ar[r] \ar[u] &
A' \ar[r] \ar[u] &
A \ar[r] \ar[u] & 0
}
$$
을 생각하자. 여기서 $B' \to B$와 $A' \to A$는 제곱이 영인 아이디얼을
핵으로 갖는 전사 환 사상이다. $A \to B$가 에탈이고
$J = I \otimes_A B$이면 $A' \to B'$도 에탈이다.
\end{lemma}

\begin{proof}
보조정리 \ref{algebra-lemma-lift-etale}에 의해 $C/IC = B$가 되게 하는 에탈 환 사상
$A' \to C$가 존재한다. 그러면 명제
\ref{algebra-proposition-smooth-formally-smooth}에 의해 $A' \to C$는 형식적으로
매끄러우므로 $A'$-대수 사상 $\varphi : C \to B'$를 얻는다.
$A' \to C$가 평탄하므로
$I \otimes_A B = I \otimes_A C/IC = IC$이다. 따라서
$J = I \otimes_A B$라는 가정에 의해 $\varphi$는 동형사상
$IC \to J$와 동형사상 $C/IC \to B'/IB'$를 유도하며, 이에 따라
$\varphi$는 동형사상이다.
\end{proof}

\begin{example}
\label{algebra-example-factor-polynomials-etale}
% P01-E0878: 아래 동결 수학식의 n 뒤 쉼표 앞에는 이례적인 공백이 있다.
$n , m \geq 1$을 정수라 하자. 예 \ref{algebra-example-factor-polynomials}의 환 사상
\begin{eqnarray*}
R = \mathbf{Z}[a_1, \ldots, a_{n + m}]
& \longrightarrow &
S = \mathbf{Z}[b_1, \ldots, b_n, c_1, \ldots, c_m] \\
a_1 & \longmapsto & b_1 + c_1 \\
a_2 & \longmapsto & b_2 + b_1 c_1 + c_2 \\
\ldots & \ldots & \ldots \\
a_{n + m} & \longmapsto & b_n c_m
\end{eqnarray*}
을 생각하자. 기호적으로
$$
S = R[b_1, \ldots, c_m]/(\{a_k(b_i, c_j) - a_k\}_{k = 1, \ldots, n + m})
$$
로 쓰자. 예를 들어 $a_1(b_i, c_j) = b_1 + c_1$이다.
편미분들의 행렬은
$$
\left(
\begin{matrix}
1 & c_1 & \ldots & c_m & 0 & \ldots & \ldots & 0 \\
0 & 1 & c_1 & \ldots & c_m & 0 & \ldots & 0 \\
\ldots & \ldots & \ldots & \ldots & \ldots & \ldots & \ldots & \ldots \\
0 & \ldots & 0 & 1 & c_1 & c_2 & \ldots & c_m \\
1 & b_1 & \ldots & b_{n - 1} & b_n & 0 & \ldots & 0 \\
0 & 1 & b_1 & \ldots & b_{n - 1} & b_n & \ldots & 0 \\
\ldots & \ldots & \ldots & \ldots & \ldots & \ldots & \ldots & \ldots \\
0 & \ldots & \ldots & 0 & 1 & b_1 & \ldots & b_n
\end{matrix}
\right)
$$
이다. 이 행렬의 행렬식 $\Delta$는 다항식
$g = x^n + b_1 x^{n - 1} + \ldots + b_n$과
$h = x^m + c_1 x^{m - 1} + \ldots + c_m$의 {\it 종결식}으로 더 잘
알려져 있고, 위 행렬은 $g, h$에 딸린 {\it 실베스터 행렬}이라고 한다.
식으로는 $\Delta = \text{Res}_x(g, h)$이다. 실베스터 행렬은 선형사상
\begin{eqnarray*}
S[x]_{< m} \oplus S[x]_{< n} & \longrightarrow & S[x]_{< n + m} \\
a \oplus b & \longmapsto & ag + bh
\end{eqnarray*}
의 행렬을 전치한 것이다. $\mathfrak q \subset S$를 임의의 소 아이디얼이라
하자. 위의 내용에 의해 다음은 서로 동치이다.
\begin{enumerate}
\item $R \to S$는 $\mathfrak q$에서 에탈이다.
\item $\Delta = \text{Res}_x(g, h) \not \in \mathfrak q$이다.
\item 다항식 $g, h$의 상
$\overline{g}, \overline{h} \in \kappa(\mathfrak q)[x]$는
$\kappa(\mathfrak q)[x]$에서 서로소이다.
\end{enumerate}
(2)와 (3)의 동치는 $\text{Mat}(n + m, \kappa(\mathfrak q))$에서
실베스터 행렬의 상에 핵이 있을 필요충분조건이 다항식
$\overline{g}, \overline{h}$가 공통인 인자를 갖는 것이라는 사실에서
따른다. 따라서 환 사상
$$
R \longrightarrow S[\frac{1}{\Delta}] = S[\frac{1}{\text{Res}_x(g, h)}]
$$
은 에탈이다.
\end{example}


\begin{lemma}
\label{algebra-lemma-factor-mod-lift-etale}
$R$을 환이라 하자. $f \in R[x]$를 최고차항계수가 1인 다항식이라 하자.
$\mathfrak p$를 $R$의 소 아이디얼이라 하자.
$f \bmod \mathfrak p = \overline{g} \overline{h}$를
$\kappa(\mathfrak p)[x]$에서 $f$의 상의 인수분해라 하자.
$\gcd(\overline{g}, \overline{h}) = 1$이면 다음이 존재한다.
\begin{enumerate}
\item 에탈 환 사상 $R \to R'$,
\item $\mathfrak p$ 위에 놓이는 소 아이디얼 $\mathfrak p' \subset R'$,
\item $R'[x]$에서의 인수분해 $f = g h$.
\end{enumerate}
그리고 다음이 성립한다.
\begin{enumerate}
\item $\kappa(\mathfrak p) = \kappa(\mathfrak p')$,
\item $\overline{g} = g \bmod \mathfrak p'$,
$\overline{h} = h \bmod \mathfrak p'$,
\item 다항식 $g, h$는 $R'[x]$에서 단위 아이디얼을 생성한다.
\end{enumerate}
\end{lemma}

\begin{proof}
$\overline{g} = \overline{b}_0 x^n + \overline{b}_1 x^{n - 1} + \ldots
+ \overline{b}_n$, 그리고
$\overline{h} = \overline{c}_0 x^m + \overline{c}_1 x^{m - 1} + \ldots
+ \overline{c}_m$이고
$\overline{b}_0, \overline{c}_0 \in \kappa(\mathfrak p)$가 영이 아니라고 하자.
$\mathfrak p$에 속하지 않는 $R$의 어떤 원소에서 $R$를 국소화한 뒤,
$\overline{b}_0$가 가역원 $b_0 \in R$의 상이라고 가정할 수 있다.
$\overline{g}$를 $\overline{g}/b_0$로, $\overline{h}$를
$b_0\overline{h}$로 바꾸면 $\overline{g}$와 $\overline{h}$의
최고차항계수가 1인 경우로 환원된다(확인은 생략한다).
$\overline{g} = x^n + \overline{b}_1 x^{n - 1} + \ldots + \overline{b}_n$,
$\overline{h} = x^m + \overline{c}_1 x^{m - 1} + \ldots + \overline{c}_m$라고 하자.
$f = x^{n + m} + a_1 x^{n + m - 1} + \ldots + a_{n + m}$로 쓰자.
섬유곱
$$
R' = R \otimes_{\mathbf{Z}[a_1, \ldots, a_{n + m}]}
\mathbf{Z}[b_1, \ldots, b_n, c_1, \ldots, c_m]
$$
을 생각하자. 여기서 사상 $\mathbf{Z}[a_k] \to \mathbf{Z}[b_i, c_j]$는
예 \ref{algebra-example-factor-polynomials}와
\ref{algebra-example-factor-polynomials-etale}에서와 같다. 구성에 의해
$b_i$를 $\overline{b}_i$로, $c_j$를 $\overline{c}_j$로 보내는
$R$-대수 사상
$$
R' = R \otimes_{\mathbf{Z}[a_1, \ldots, a_{n + m}]}
\mathbf{Z}[b_1, \ldots, b_n, c_1, \ldots, c_m]
\longrightarrow
\kappa(\mathfrak p)
$$
이 존재한다.
% P01-E0879: 동결 원문의 ``Denote''에는 by가 빠져 있다.
이 사상의 핵을 $\mathfrak p' \subset R'$라 하자.
가정에 의해 다항식 $\overline{g}, \overline{h}$는 서로소이므로,
표시된 사상 아래에서 원소
$\Delta = \text{Res}_x(g, h) \in \mathbf{Z}[b_i, c_j]$
(예 \ref{algebra-example-factor-polynomials-etale} 참조)의 상은
$\kappa(\mathfrak p)$에서 영이 아니다. 따라서 $R \to R'$는
$\mathfrak p'$에서 에탈이다. 실제로 보조정리에 제시된 문제의 한 해는
환 사상 $R \to R'[1/\Delta]$와 소 아이디얼
$\mathfrak p' R'[1/\Delta]$이다.
% P01-E0880: 아래 동결 수식의 Res_x(f,g)는 논증상 Res_x(g,h)여야 하지만 그대로 보존한다.
이 환에서 $\text{Res}_x(f, g)$가 가역원이므로 실베스터 행렬은
$R'[1/\Delta]$ 위에서 가역이고, 따라서 어떤
$a, b \in R'[1/\Delta][x]$에 대해 $1 = a g +  b h$이다.
예 \ref{algebra-example-factor-polynomials-etale}를 보라.
\end{proof}

\section{에탈 환 사상의 국소 구조}
\label{algebra-section-etale-local-structure}

\noindent
보조정리 \ref{algebra-lemma-etale-standard-smooth}은 상대 차원이 $0$인 표준
매끄러운 사상으로 표준 에탈 사상을 정의하는 것이 실제로는 타당하지
않음을 말해 준다. 에탈 사상의 모형으로 체의 유한 분리 확대
$k'/k$가 주는 예를 택하자. 즉, 어떤 $\alpha \in k'$를 잡아
$k' = k(\alpha)$가 되게 할 수 있고, $\alpha$의 최소 다항식
$f(x) \in k[x]$의 도함수 $f'$가 $f$와 서로소가 되게 할 수 있다.

\begin{definition}
\label{algebra-definition-standard-etale}
% P01-E0881: 동결 정의식에는 g , f 사이의 이례적인 공백과 중복 공백이 있다.
환 $R$을 잡자. $g , f  \in R[x]$라 하자.
 $f$가 최고차항계수가 1인 다항식이고 도함수 $f'$가
국소화 $R[x]_g/(f)$에서 가역원이라고 하자.
이 경우 환 사상 $R \to R[x]_g/(f)$를
{\it 표준 에탈}이라고 한다.
\end{definition}

\noindent
명제 \ref{algebra-proposition-etale-locally-standard}에서 모든 에탈 환 사상이
국소적으로 표준 에탈임을 보일 것이다.

\begin{lemma}
\label{algebra-lemma-standard-etale}
$R \to R[x]_g/(f)$가 표준 에탈이라고 하자.
\begin{enumerate}
\item 환 사상 $R \to R[x]_g/(f)$는 에탈이다.
\item 임의의 환 사상 $R \to R'$에 대해, 표준 에탈 환 사상
$R \to R[x]_g/(f)$의 기저 변환
$R' \to R'[x]_g/(f)$도 표준 에탈이다.
\item $R[x]_g/(f)$의 임의의 주아이디얼 국소화는 $R$ 위에서
표준 에탈이다.
\item 표준 에탈 사상들의 합성은 일반적으로 {\bf 표준 에탈이 아니다}.
\end{enumerate}
\end{lemma}

\begin{proof}
증명은 생략한다. (4)의 예를 들면 다음과 같다.
환 사상 $\mathbf{F}_2 \to \mathbf{F}_{2^2}$는 표준 에탈이다.
환 사상
$\mathbf{F}_{2^2} \to \mathbf{F}_{2^2} \times \mathbf{F}_{2^2}
\times \mathbf{F}_{2^2} \times \mathbf{F}_{2^2}$도 표준 에탈이다.
그러나 환 사상
$\mathbf{F}_2 \to \mathbf{F}_{2^2} \times \mathbf{F}_{2^2}
\times \mathbf{F}_{2^2} \times \mathbf{F}_{2^2}$는 표준 에탈이 아니다.
\end{proof}

\noindent
표준 에탈 사상은 에탈 사상을 만들어 내는 편리한 방법이다.
예를 하나 보자.

\begin{lemma}
\label{algebra-lemma-make-etale-map-prescribed-residue-field}
환 $R$을 잡자. $\mathfrak p$를 $R$의 소 아이디얼이라 하자.
$L/\kappa(\mathfrak p)$가 유한 분리 체 확대라고 하자.
$\mathfrak p$ 위에 놓이는 소 아이디얼 $\mathfrak p'$를 갖는 에탈
환 사상 $R \to R'$가 존재하여
체 확대 $\kappa(\mathfrak p')/\kappa(\mathfrak p)$가
$\kappa(\mathfrak p) \subset L$과 동형이 된다.
\end{lemma}

\begin{proof}
원시 원소 정리에 의해 $L = \kappa(\mathfrak p)[\alpha]$로 쓸 수 있다.
 $\overline{f} \in \kappa(\mathfrak p)[x]$를 $\alpha$의 최소 다항식이라
하자(특히 이는 최고차항계수가 1인 다항식이다).
$\alpha$를 어떤 $c \in R$, $c\not \in \mathfrak p$에 대해
$c\alpha$로 바꾼 뒤에는 $\overline{f}$의 모든 계수가
$R \to \kappa(\mathfrak p)$의 상에 들어간다고 가정할 수 있다
(검증은 생략한다). 따라서 $\kappa(\mathfrak p)[x]$에서
$\overline{f}$의 상이 되는 최고차항계수가 1인 다항식
$f \in R[x]$를 찾을 수 있다.
$\kappa(\mathfrak p) \subset L$가 분리이므로
$\gcd(\overline{f}, \overline{f}') = 1$이다.
따라서 $\overline{f}'(\alpha) \gamma = 1$이 되는 원소
$\gamma \in L$가 존재한다. 그러므로 다음의 $R$-대수 사상을 얻는다.
\begin{eqnarray*}
R[x, 1/f']/(f) & \longrightarrow & L \\
x & \longmapsto & \alpha \\
1/f' & \longmapsto & \gamma
\end{eqnarray*}
왼쪽 환은 $R$ 위의 표준 에탈 대수 $R'$이고, 환 사상의 핵이
원하는 소 아이디얼을 준다.
\end{proof}

\begin{proposition}
\label{algebra-proposition-etale-locally-standard}
$R \to S$를 환 사상이라 하자. $\mathfrak q \subset S$를 소 아이디얼이라 하자.
$R \to S$가 $\mathfrak q$에서 에탈이면, $g \in S$, $g \not \in \mathfrak q$인
원소가 존재하여 $R \to S_g$가 표준 에탈이 된다.
\end{proposition}

\begin{proof}
다음 증명은 다소 우회적이며 더 짧게 하는 방법이 있을 수 있다.

\medskip\noindent
단계 1. 정의 \ref{algebra-definition-etale}에 의해 $g \in S$, $g \not \in \mathfrak q$인
원소가 존재하여 $R \to S_g$가 에탈이 된다. 따라서 $S$가 $R$ 위에서
에탈이라고 가정해도 된다.

\medskip\noindent
% P01-E0882: 동결 원문은 기저변환 항등식의 왼쪽 대수도 R로 적고 있다.
% P01-E0883: 동결 원문의 Denote에는 by가 빠져 있다.
단계 2. 보조정리 \ref{algebra-lemma-etale}에 의해, $R_0$가 $\mathbf{Z}$ 위에서
유한 생성이고 에탈 환 사상 $R_0 \to S_0$인 것과 환 사상
$R_0 \to R$가 존재하여
$R = R \otimes_{R_0} S_0$가 된다. $\mathfrak q_0$를 $\mathfrak q$에
대응하는 $S_0$의 소 아이디얼이라 하자.
$(R_0 \to S_0, \mathfrak q_0)$에 대해 결과를 보이면
기저 변환에 의해 $(R \to S, \mathfrak q)$에 대한 결과가 따른다.
그러므로 $R$이 뇌터 환이라고 가정해도 된다.

\medskip\noindent
단계 3.
보조정리 \ref{algebra-lemma-etale-quasi-finite}에 의해 $R \to S$는 준유한이다.
% P01-E0884: 동결 원문의 such that 절이 반복되고 g'의 소속/비소속
% 조건이 혼동되기 쉽다.
보조정리 \ref{algebra-lemma-quasi-finite-open-integral-closure}에 의해 유한 환 사상
$R \to S'$, $R$-대수 사상 $S' \to S$, 그리고 $g' \in S'$가 존재하여
$g' \not \in \mathfrak q$이고 $S' \to S$가
$S'_{g'} \cong S_{g'}$를 유도한다.
(물론 일반적으로 $S'$가 $R$ 위에서 에탈일 필요는 없다.)
따라서 (a) $R$이 뇌터이고, (b) $R \to S$가 유한이며,
(c) $R \to S$가 $\mathfrak q$에서 에탈이라고 가정해도 된다
(다만 이제 모든 소 아이디얼에서 에탈이라고 할 수는 없다).

\medskip\noindent
단계 4. $\mathfrak p \subset R$를 $\mathfrak q$에 대응하는 소 아이디얼이라 하자.
섬유 환 $S \otimes_R \kappa(\mathfrak p)$를 생각하자.
이는 $\kappa(\mathfrak p)$ 위의 유한 대수이다. 따라서 아르틴 환이고
(보조정리 \ref{algebra-lemma-finite-dimensional-algebra}), 유한 개의 국소 환의
곱이다.
% P01-E0885: 동결 원문에서 곱의 분해식 뒤의 문장 연결에 문장부호가 없다.
$$
S \otimes_R \kappa(\mathfrak p) = \prod\nolimits_{i = 1}^n A_i
$$
명제 \ref{algebra-proposition-dimension-zero-ring}에 따른다. 인수 중 하나,
가령 $A_1$은 $S_{\mathfrak q}/\mathfrak pS_{\mathfrak q}$인 국소 환이고,
보조정리 \ref{algebra-lemma-etale-at-prime}에 의해 $\kappa(\mathfrak q)$와 동형이다.
나머지 인수들은 $\mathfrak p$ 위에 놓이는 $S$의 다른 소 아이디얼들,
가령 $\mathfrak q_2, \ldots, \mathfrak q_n$에 대응한다.

\medskip\noindent
단계 5. 유한 분리 체 확대
$\kappa(\mathfrak q)/\kappa(\mathfrak p)$를 생성하는 0이 아닌 원소
$\alpha \in \kappa(\mathfrak q)$를 택할 수 있다
(체 확대가 자명하더라도 $\alpha = 0$은 허용하지 않는다).
임의의 $\lambda \in \kappa(\mathfrak p)^*$에 대해 원소
$\lambda \alpha$도 $\kappa(\mathfrak p)$ 위에서 $\kappa(\mathfrak q)$를
생성한다는 점을 주목하자. 다음 원소를 생각하자.
$$
\overline{t} =
(\alpha, 0, \ldots, 0) \in
\prod\nolimits_{i = 1}^n A_i =
S \otimes_R \kappa(\mathfrak p).
$$
필요하면 위와 같이 $\alpha$를 $\lambda \alpha$로 바꾸어
$\overline{t}$가 $t \in S$의 상이라고 가정할 수 있다.
$x$를 $t$로 보내는 $R$-대수 사상 $R[x] \to S$의 핵을
$I \subset R[x]$라 하자. $S' = R[x]/I$로 두면 $S' \subset S$이다.
다음은 그림이다.
$$
\xymatrix{
R[x] \ar[r] & S' \ar[r] & S \\
R \ar[u] \ar[ru] \ar[rru] & &
}
$$
구성에 의해 $j \geq 2$인 $S$의 소 아이디얼 $\mathfrak q_j$들은
$R[x]$의 소 아이디얼 $(\mathfrak p, x)$ 위에 놓인다. 반면
$\mathfrak q$는 $\alpha \not = 0$이므로 $R[x]$의 다른 소 아이디얼
위에 놓인다.

\medskip\noindent
% P01-E0886: 동결 원문의 Denote에는 by가 빠져 있다.
% P01-E0887: 동결 원문의 going-up 절은 문법적으로 불완전하다.
단계 6. $\mathfrak q' \subset S'$를 $S'$의 소 아이디얼로서
$\mathfrak q$에 대응하는 것으로 쓰자.
위에서 본 바에 의해 $\mathfrak q'$ 위에 놓이는 $S$의 소 아이디얼은
$\mathfrak q$ 하나뿐이다. 따라서
$S_{\mathfrak q} = S_{\mathfrak q'}$임을 알 수 있다.
보조정리 \ref{algebra-lemma-unique-prime-over-localize-below}를 보라
($S' \to S$가 유한이고 $R \to S$가 유한이므로,
보조정리 \ref{algebra-lemma-integral-going-up}에 의해 $S' \to S$에 대해
상승 성질이 성립한다).
그러므로 $S'_{\mathfrak q'} \to S_{\mathfrak q}$는 유한 단사 사상이다.
이는 유한 단사 환 사상 $S' \to S$의 국소화이기 때문이다.
다음 국소 환 사상들을 생각하자.
$$
R_{\mathfrak p} \to S'_{\mathfrak q'} \to S_{\mathfrak q}
$$
두 번째 사상은 유한 단사이다. 보조정리 \ref{algebra-lemma-etale-at-prime}에 의해
$S_{\mathfrak q}/\mathfrak pS_{\mathfrak q} = \kappa(\mathfrak q)$이다.
따라서 특히 $S_{\mathfrak q}/\mathfrak q'S_{\mathfrak q} =
\kappa(\mathfrak q)$이다. 또
$$
\kappa(\mathfrak p) \subset \kappa(\mathfrak q') \subset \kappa(\mathfrak q)
$$
이고 $\alpha$는 $\kappa(\mathfrak q')$의 $\kappa(\mathfrak q)$로의 상에
들어 있으므로 $\kappa(\mathfrak q') = \kappa(\mathfrak q)$라고 결론 내린다.
그러므로 보조정리 \ref{algebra-lemma-NAK}를 $S'_{\mathfrak q'}$-가군 사상
$S'_{\mathfrak q'} \to S_{\mathfrak q}$에 적용하면 이 사상
$S'_{\mathfrak q'} \to S_{\mathfrak q}$는 전사이다. 즉,
$S'_{\mathfrak q'} \cong S_{\mathfrak q}$이다.

\medskip\noindent
단계 7. 보조정리 \ref{algebra-lemma-isomorphic-local-rings}에 의해
$g \in S$, $g \not \in \mathfrak q$ 및 $g' \in S'$,
$g' \not \in \mathfrak q'$가 존재하여 $S'_{g'} \cong S_g$가 된다.
$R$이 뇌터이므로 $S'$는 유한 $R$-가군 $S$의 $R$-부분가군이어서
$R$ 위에서 유한이다. 따라서 $S$를 $S'$로 바꾼 뒤
(a) $R$이 뇌터이고, (b) $S$가 $R$ 위에서 유한이며,
(c) $S$가 $\mathfrak q$에서 $R$ 위에 에탈이고, (d) $S = R[x]/I$라고
가정할 수 있다.

\medskip\noindent
단계 8. 다음 환을 생각하자.
$S \otimes_R \kappa(\mathfrak p) = \kappa(\mathfrak p)[x]/\overline{I}$
이며, $\overline{I} = I \cdot \kappa(\mathfrak p)[x]$는
$\kappa(\mathfrak p)[x]$에서 $I$가 생성하는 아이디얼이다.
$\kappa(\mathfrak p)[x]$가 주아이디얼 정역이므로
$\overline{I} = (\overline{h})$인 어떤 최고차항계수가 1인
$\overline{h} \in \kappa(\mathfrak p)[x]$가 존재한다.
$\overline{h}$를 어떤 $\lambda \in \kappa(\mathfrak p)$에 대해
$\lambda \cdot \overline{h}$로 바꾸면 $\overline{h}$가
$h \in I \subset R[x]$의 상이라고 가정할 수 있다.
(문제는 $h$를 최고차항계수가 1인 것으로 택할 수 있는지 모른다는 점이다.)
또한 단계 4에서와 같이
$S \otimes_R \kappa(\mathfrak p) = A_1 \times \ldots \times A_n$이고
$A_1 = \kappa(\mathfrak q)$는 $\kappa(\mathfrak p)$의 유한 분리 확대이며
$A_2, \ldots, A_n$은 국소 환이다. 따라서
$$
\overline{h} = \overline{h}_1 \overline{h}_2^{e_2} \ldots \overline{h}_n^{e_n}
$$
가 된다. 여기서 $\overline{h}_i \in \kappa(\mathfrak p)[x]$는 서로
소인 최고차항계수가 1인 기약 다항식들이고 $e_2, \ldots, e_n \geq 1$이다.
번호는 $A_i = \kappa(\mathfrak p)[x]/(\overline{h}_i^{e_i})$가
$\kappa(\mathfrak p)[x]$-대수가 되도록 정한 것이다. $\overline{h}_1$은
$\alpha \in \kappa(\mathfrak q)$의 최소 다항식이므로 분리 다항식이다
(그 도함수는 자기 자신과 서로소이다).

\medskip\noindent
% P01-E0890: 동결 원문의 Denote에는 by가 빠져 있다.
단계 9. $m \in I$인 최고차항계수가 1인 원소를 택하자. 환 확대
$R \to R[x]/I$가 유한이고 따라서 정수적이므로 이러한 원소가 존재한다.
$\overline{m}$을 $\kappa(\mathfrak p)[x]$에서의 상이라 하자.
다음과 같이 인수분해할 수 있다.
$$
\overline{m} = \overline{k}
\overline{h}_1^{d_1} \overline{h}_2^{d_2} \ldots \overline{h}_n^{d_n}
$$
여기서 $d_1 \geq 1$, $d_j \geq e_j$, $j = 2, \ldots, n$이고
$\overline{k} \in \kappa(\mathfrak p)[x]$는 모든 $\overline{h}_i$와 서로소이다.
$l \deg(m) > \deg(h)$이고 $l \geq 2$인 조건을 만족하는 것을 택하고
$f = m^l + h$로 두자. 그러면 $f$는 $R$ 위의 다항식으로서
최고차항계수가 1이다. $\kappa(\mathfrak p)[x]$에서 $f$의 상
$\overline{f}$는 다음과 같이 인수분해된다.
$$
\overline{f} =
\overline{h}_1 \overline{h}_2^{e_2} \ldots \overline{h}_n^{e_n}
+
\overline{k}^l \overline{h}_1^{ld_1} \overline{h}_2^{ld_2}
\ldots \overline{h}_n^{ld_n}
=
\overline{h}_1(\overline{h}_2^{e_2} \ldots \overline{h}_n^{e_n}
+
\overline{k}^l
\overline{h}_1^{ld_1 - 1} \overline{h}_2^{ld_2} \ldots \overline{h}_n^{ld_n})
= \overline{h}_1 \overline{w}
$$
여기서 $\overline{w}$는 $\overline{h}_1$과 서로소인 다항식이다.
$g = f'$ (즉 $x$에 대한 도함수)로 두자.

\medskip\noindent
단계 10. 환 사상 $R[x] \to S = R[x]/I$는 다음 성질을 갖는다.
(1) $f$를 0으로 보낸다.
(2) $g$를 $S \setminus \mathfrak q$의 원소로 보낸다.
첫째 주장은 $f$가 $I$의 원소이므로 자명하다.
둘째 주장을 위해서는 $g$가
$\kappa(\mathfrak q) = \kappa(\mathfrak p)[x]/(\overline{h}_1)$에서
0으로 가지 않음을 보이면 된다. $\kappa(\mathfrak p)[x]$에서 $g$의 상은
$\overline{f}$의 도함수이다. 따라서
$$
\overline{g} =
\frac{\text{d}\overline{f}}{\text{d}x} =
\overline{w}\frac{\text{d}\overline{h}_1}{\text{d}x} +
\overline{h}_1\frac{\text{d}\overline{w}}{\text{d}x},
$$
이고 $\overline{w}$는 $\overline{h}_1$과 서로소이며
$\overline{h}_1$은 분리이므로 (2)가 분명하다.

\medskip\noindent
단계 11.
$\varphi : R[x]/(f) \to S$가 전사 환 사상이고,
$R[x]_g/(f)$가 $R$ 위에서 에탈임을 결론 내린다
(표준 에탈이므로, 보조정리 \ref{algebra-lemma-standard-etale} 참조).
또한 $\varphi(g) \not \in \mathfrak q$이다.
$g' \in R[x]/(f)$를 택하여 $\varphi(g') \not \in \mathfrak q$이고
$S_{\varphi(g')}$가 $R$ 위에서 에탈이 되게 하자
($S$가 $\mathfrak q$에서 $R$ 위에 에탈이므로 그런 원소가 존재한다).
그러면
$R[x]_{gg'}/(f) \to S_{\varphi(gg')}$는 $R$ 위의 에탈 대수 사이의 전사 사상이다.
따라서 보조정리 \ref{algebra-lemma-map-between-etale}에 의해 에탈이다.
보조정리 \ref{algebra-lemma-surjective-flat-finitely-presented}에 의해
이는 국소화이다. 그러므로 $\mathfrak q$에 들어 있지 않은 원소에서
$S$를 국소화한 것은 $R$ 위의 표준 에탈 대수의 국소화와 동형이며,
이것이 보이려던 결과이다.
\end{proof}
\noindent
다음 두 보조정리는 에탈 위상이 자리스키 피복과 유한 평탄 사상들이
생성하는 위상보다 더 성기다는 것을 말한다. 첫 번째 읽기에서는
건너뛰어도 된다.

\begin{lemma}
\label{algebra-lemma-standard-etale-finite-flat-Zariski}
$R \to S$를 표준 에탈 사상이라 하자.
다음 성질을 갖는 환 사상 $R \to S'$가 존재한다.
\begin{enumerate}
\item $R \to S'$는 유한이고, 유한 표시이며, 평탄하다
(다시 말해 $S'$는 $R$-가군으로서 유한 사영이다).
\item $\Spec(S') \to \Spec(R)$는 전사이다.
\item $\mathfrak p \subset R$ 위에 놓이는 $S$의 모든 소 아이디얼
$\mathfrak q \subset S$와, $\mathfrak p$ 위에 놓이는 모든 소 아이디얼
$\mathfrak q' \subset S'$에 대해,
$g' \in S'$, $g' \not \in \mathfrak q'$가 존재하여
환 사상 $R \to S'_{g'}$가 어떤 사상
$\varphi : S \to S'_{g'}$를 통해 인수분해되고
$\varphi^{-1}(\mathfrak q'S'_{g'}) = \mathfrak q$가 된다.
\end{enumerate}
\end{lemma}

\begin{proof}
정의 \ref{algebra-definition-standard-etale}와 같이 $S = R[x]_g/(f)$로
이를 표시하자. $f = x^n + a_1 x^{n - 1} + \ldots + a_n$,
$a_i \in R$로 쓰자.
보조정리 \ref{algebra-lemma-adjoin-roots}에 의해 유한 자유이고 충실히 평탄한
환 사상 $R \to S'$가 존재하여 어떤 $\alpha_i \in S'$에 대해
$f = \prod (x - \alpha_i)$가 된다. 따라서 $R \to S'$는 (1), (2)를 만족한다.
$\mathfrak q \subset R[x]/(f)$를 $g \not \in \mathfrak q$인 소 아이디얼이라 하자
(즉 $S$의 소 아이디얼에 대응한다).
$\mathfrak p = R \cap \mathfrak q$로 두고 $\mathfrak p$ 위에 놓이는
소 아이디얼 $\mathfrak q' \subset S'$를 택하자.
다음과 같은 $R$-대수 사상이 $n$개 있다.
\begin{eqnarray*}
\varphi_i : R[x]/(f) & \longrightarrow & S' \\
x & \longmapsto & \alpha_i
\end{eqnarray*}
증명을 끝내려면 어떤 $i$에 대해 다음이 성립함을 보이면 된다.
(a) $\varphi_i(g)$의 $\kappa(\mathfrak q')$에서의 상이 0이 아니다.
(b) $\varphi_i^{-1}(\mathfrak q') = \mathfrak q$이다.
그러면 그 $i$에 대해 $g' = \varphi_i(g)$, $\varphi = \varphi_i$로
두면 된다.

\medskip\noindent
% P01-E0891: 동결 원문의 목록 도입부에는 열거 앞 문장부호가 없다.
$\overline{f}$를 $\kappa(\mathfrak p)[x]$에서 $f$의 상이라 하자.
$\Spec(\kappa(\mathfrak p)[x]/(\overline{f}))$의 점으로서
$\mathfrak q$는 $\overline{f}$의 기약 인수 $f_1$에 대응한다.
또한 $g \not \in \mathfrak q$라는 것은 $f_1$이
$\kappa(\mathfrak p)[x]$에서 $g$의 상 $\overline{g}$를 나누지 않는다는 뜻이다.
% P01-E0892: 동결 원문의 Denote에는 by가 빠져 있다.
$\alpha_1, \ldots, \alpha_n$의 상을 $\kappa(\mathfrak q')$에서
$\overline{\alpha}_1, \ldots, \overline{\alpha}_n$이라 하자.
그러면 다항식 $\overline{f}$는 $\kappa(\mathfrak q')[x]$에서 완전히
분해되며, 즉
$$
\overline{f} = \prod\nolimits_i (x - \overline{\alpha}_i)
$$
이다. 더구나 $\varphi_i(g)$는 $\overline{g}(\overline{\alpha}_i)$로
내려간다. 따라서 $f_1(\overline{\alpha}_i) = 0$이고
$\overline{g}(\overline{\alpha}_i) \not = 0$인 $i$를 택할 수 있다.
이 $i$에 대해 (a), (b)가 성립한다. 세부 사항은 생략한다.
\end{proof}
\begin{lemma}
\label{algebra-lemma-etale-finite-flat-zariski}
$R \to S$를 환 사상이라 하자. 다음을 가정하자.
\begin{enumerate}
\item $R \to S$는 에탈이다.
\item $\Spec(S) \to \Spec(R)$는 전사이다.
\end{enumerate}
그러면 다음 성질을 갖는 환 사상 $R \to S'$가 존재한다.
\begin{enumerate}
\item $R \to S'$는 유한이고, 유한 표시이며, 평탄하다
(다시 말해 $R$-가군으로서 유한 사영이다).
\item $\Spec(S') \to \Spec(R)$는 전사이다.
\item 모든 소 아이디얼 $\mathfrak q' \subset S'$에 대해
$g' \in S'$, $g' \not \in \mathfrak q'$가 존재하여
환 사상 $R \to S'_{g'}$가 $R \to S \to S'_{g'}$를 통해 인수분해된다.
\end{enumerate}
\end{lemma}

\begin{proof}
명제 \ref{algebra-proposition-etale-locally-standard}와
$\Spec(S)$의 준콤팩트성(보조정리 \ref{algebra-lemma-quasi-compact} 참조)에 의해
$S$의 단위원을 생성하는 $g_1, \ldots, g_n \in S$를 찾아
각 $R \to S_{g_i}$가 표준 에탈이 되게 할 수 있다.
환 사상 $R \to \prod_{i = 1, \ldots, n} S_{g_i}$에 대해 보조정리를
증명하면 환 사상 $R \to S$에 대한 결론이 따른다.
따라서 $S = \prod_{i = 1, \ldots, n} S_i$가 표준 에탈 사상들의
유한 곱이라고 가정할 수 있다.

\medskip\noindent
각 $i$에 대해 보조정리
\ref{algebra-lemma-standard-etale-finite-flat-Zariski}에 따라
표준 에탈 사상 $R \to S_i$에 맞는 환 사상 $R \to S_i'$를 택하자.
이하에서는 $R$-대수 사상 $S_i' \to S'$
$S' = S_1' \otimes_R \ldots \otimes_R S_n'$를
더 언급하지 않고 사용한다.
이 선택이 작동한다고 주장하자. 성질 (1), (2)는 즉시 따른다.
% P01-E0893: 동결 원문의 Denote에는 by가 빠져 있다.
성질 (3)을 위해 $\mathfrak q' \subset S'$를 소 아이디얼이라 하자.
% P01-E0894: 동결 원문의 Set ... the image에는 to be가 빠져 있다.
$\mathfrak p$를 $\Spec(R)$에서의 그 상으로 두자.
$\mathfrak p$가 $\Spec(S_i) \to \Spec(R)$의 상에 들어가는
$i \in \{1, \ldots, n\}$를 택하자. 이는 가정에 의해 가능하다.
$\mathfrak q'$의 $S_i'$의 스펙트럼에서의 상을
$\mathfrak q_i' \subset S_i'$라 하자.
% P01-E0895: 앞의 보조정리에서 얻는 인수분해 조건은 번역에서 명시적으로 보존한다.
$S'_i$의 구성에 의해 $g'_i \in S_i'$가 존재하여
$R \to (S_i')_{g_i'}$가
$R \to S_i \to (S_i')_{g_i'}$를 통해 인수분해된다.
따라서 $R \to S'_{g_i'}$도 다음을 통해 인수분해된다.
$$
R \to S_i \to (S_i')_{g_i'} \to S'_{g_i'}
$$
원하던 바이다.
\end{proof}

\section{준유한 환 사상의 에탈 국소 구조}
\label{algebra-section-etale-local-quasi-finite}

\noindent
다음 보조정리들은 대략 에탈 확대를 거치면 준유한 환 사상이 유한 사상이
된다는 것을 말한다. 결과를 해석하는 데 도움이 되도록, 유한형 환 사상이
준유한인 자리는 열린집합이고(보조정리 \ref{algebra-lemma-quasi-finite-open} 참조),
이 자리의 구성은 임의의 기초변환과 가환한다는 사실
(보조정리 \ref{algebra-lemma-quasi-finite-base-change} 참조)을 상기하자.

\begin{lemma}
\label{algebra-lemma-produce-finite}
환 사상 $R \to S' \to S$가 주어졌다고 하자.
$\mathfrak p \subset R$를 소 아이디얼이라 하고,
$g \in S'$를 원소라 하자. 다음을 가정하자.
\begin{enumerate}
\item $R \to S'$는 정수적이다.
\item $R \to S$는 유한형이다.
\item $S'_g \cong S_g$이다.
% P01-E0896: 동결 원문의 네 번째 조건에는 is가 빠져 있다.
\item $g$는 $S' \otimes_R \kappa(\mathfrak p)$에서 가역이다.
\end{enumerate}
% P01-E0897: 동결 원문은 수학 기호 f 앞에 부정관사 a를 쓴다.
그러면 $f \in R$, $f \not \in \mathfrak p$인 원소가 존재하여
$R_f \to S_f$가 유한이 된다.
\end{lemma}

\begin{proof}
가정에 의해 $V(g) \subset \Spec(S')$의
$\Spec(S') \to \Spec(R)$ 아래에서의 상 $T$는
$\mathfrak p$를 포함하지 않는다.
% P01-E0898: 동결 원문의 especially가 문법적으로 잘못 결합되어 있다.
\ref{algebra-section-going-up}절, 특히 보조정리
\ref{algebra-lemma-going-up-closed}에 의해 $T$는 닫힌집합이다.
$f \in R$, $f \not \in \mathfrak p$이고
$T \cap D(f) = \emptyset$인 것을 택하자. 그러면 $g$는 $S'_f$에서
가역이 된다. 따라서 $S'_f \cong S_f$이다. 그러므로 $S_f$는
$R_f$ 위에서 유한형이면서 정수적이므로 유한이다.
\end{proof}

\begin{lemma}
\label{algebra-lemma-etale-makes-quasi-finite-finite-one-prime}
$R \to S$를 환 사상이라 하자.
$\mathfrak q \subset S$를 $\mathfrak p \subset R$ 위에 놓이는
소 아이디얼이라 하자.
% P01-E0899: 동결 원문의 유한형 가정에는 is가 빠져 있다.
$R \to S$가 유한형이고 $\mathfrak q$에서 준유한이라고 가정하자.
그러면 다음이 존재한다.
\begin{enumerate}
\item 에탈 환 사상 $R \to R'$,
\item $\mathfrak p$ 위에 놓이는 소 아이디얼
$\mathfrak p' \subset R'$,
\item 곱 분해
$$
R' \otimes_R S = A \times B
$$
\end{enumerate}
% P01-E0900: 동결 원문의 성질 목록 도입부에는 문장부호가 없다.
이들은 다음 성질을 갖는다.
\begin{enumerate}
\item $\kappa(\mathfrak p) = \kappa(\mathfrak p')$이다.
\item $R' \to A$는 유한이다.
\item $A$에는 $\mathfrak p'$ 위에 놓이는 소 아이디얼
$\mathfrak r$가 정확히 하나 있다.
\item $\mathfrak r$는 $\mathfrak q$ 위에 놓인다.
\item $B$에는 $\mathfrak q$와 $\mathfrak p'$ 위에 놓이는
소 아이디얼이 없다.
\end{enumerate}
\end{lemma}

\begin{proof}
$S' \subset S$를 $S$ 안에서 $R$의 정수적 폐포라 하자.
$\mathfrak q' = S' \cap \mathfrak q$로 두자.
자리스키 주정리 \ref{algebra-theorem-main-theorem}에 의해
$g \in S'$, $g \not \in \mathfrak q'$가 존재하여
$S'_g \cong S_g$가 된다. 섬유 환
$F = S \otimes_R \kappa(\mathfrak p)$와
$F' = S' \otimes_R \kappa(\mathfrak p)$를 생각하자.
% P01-E0901: 동결 원문의 Denote에는 by가 빠져 있다.
$\overline{\mathfrak q}'$를 $\mathfrak q'$에 대응하는 $F'$의
소 아이디얼이라 하자. $F'$는 $\kappa(\mathfrak p)$ 위에서
정수적이므로 $\overline{\mathfrak q}'$는 $\Spec(F')$의 닫힌점이다
(보조정리 \ref{algebra-lemma-integral-over-field} 참조).
$\mathfrak q$는 $\Spec(F)$의 고립 닫힌점
$\overline{\mathfrak q}$를 정의한다
(정의 \ref{algebra-definition-quasi-finite} 참조).
$S'_g \cong S_g$이므로 $F'_g \cong F_g$이고,
$\overline{\mathfrak q}$와 $\overline{\mathfrak q}'$는
$\Spec(F)$와 $\Spec(F')$에서 서로 동형인 열린 근방을 갖는다.
따라서 집합
$\{\overline{\mathfrak q}'\} \subset \Spec(F')$는 열려 있다.
% P01-E0902: 동결 원문은 닫힌 섬유점을 q-bar'가 아니라 q'로 적는다.
위에서 보인 $\mathfrak q'$의 닫힘성과 합치면,
$\overline{\mathfrak q}'$도 $\Spec(F')$의 고립 닫힌점을 정의한다.

\medskip\noindent
추가로, 사상 $\Spec(F) \to \Spec(F')$ 아래에서
$\overline{\mathfrak q}'$로 가는 점은 $\overline{\mathfrak q}$ 하나뿐이다.
이는 위의 논의에서 따른다.

\medskip\noindent
보조정리 \ref{algebra-lemma-disjoint-implies-product}에 의해
$F' = F'_1 \times F'_2$로 쓸 수 있고,
$\Spec(F'_1) = \{\overline{\mathfrak q}'\}$이다.
$F' = S' \otimes_R \kappa(\mathfrak p)$이므로,
어떤 $r \in R$, $r \not \in \mathfrak p$에 대해
$(r, 0) \in F'_1 \times F'_2 = F'$로 가는 원소
$s' \in S'$가 존재한다. 실제로 우리가 사용할 이 원소의 성질은
$s'$가 $S'$의 원소이고 $\mathfrak q'$에는 들어 있지 않지만
$\mathfrak p$ 위에 놓이는 다른 모든 소 아이디얼에는 들어간다는 것이다.

\medskip\noindent
$f(s') = 0$인 최고차항계수가 1인 다항식 $f(x) \in R[x]$를 택하자.
% P01-E0903: 동결 원문의 Denote에는 by가 빠져 있다.
$\overline{f} \in \kappa(\mathfrak p)[x]$를 그 상이라 하자.
이를 $\overline{f} = x^e \overline{h}$로 인수분해할 수 있고,
$\overline{h}(0) \not = 0$이다. 필요하면 $f$를 $x f$로 바꾸어
$e \geq 1$이라고 가정할 수 있다.
보조정리 \ref{algebra-lemma-factor-mod-lift-etale}에 의해
에탈 환 확대 $R \to R'$, $\mathfrak p$ 위에 놓이는 소 아이디얼
$\mathfrak p'$, 그리고 다음을 만족하는 $R'[x]$에서의 인수분해
$f = h i$를 찾을 수 있다.
$\kappa(\mathfrak p) = \kappa(\mathfrak p')$,
$\overline{h} = h \bmod \mathfrak p'$,
$x^e = i \bmod \mathfrak p'$이고, 적당한 $a, b$에 대해
$R'[x]$에서 $a h + b i = 1$로 쓸 수 있다.

\medskip\noindent
원소 $h(s'), i(s') \in R' \otimes_R S'$를 생각하자.
구성에 의해 $h(s')i(s') = f(s') = 0$이다.
한편 $a(s')h(s') + b(s')i(s') = 1$이므로 이들은 단위 아이디얼을
생성한다. 따라서 $R' \otimes_R S'$는 이 원소들에서의 국소화들의 곱이다.
$$
R' \otimes_R S'
=
(R' \otimes_R S')_{i(s')}
\times
(R' \otimes_R S')_{h(s')}
=
S'_1 \times S'_2
$$
더구나 이 곱 분해는 위에서 찾은 섬유 환 $F'$의 곱 분해와 양립한다.
이는 $s', i, h$의 선택으로부터 따르는데, 그 선택은
$\overline{\mathfrak q}'$가 $F'$에서 $i(s')$의 상을 포함하지 않는
유일한 소 아이디얼임을 보장한다. 여기서는
$\kappa(\mathfrak p) = \kappa(\mathfrak p')$이므로
$\mathfrak p'$에서 $R'$ 위의 $R'\otimes_R S'$의 섬유 환이
$F'$와 같다는 사실을 사용한다. 여기서 $F'$는 위에서 정한 섬유 환이다.
% P01-E0904: 동결 원문의 S'_1 뒤에는 중복 공백이 있다.
따라서 $S'_1$에는 $\mathfrak p'$ 위에 놓이는 소 아이디얼
$\mathfrak r'$가 정확히 하나 있고, 이 소 아이디얼은
$\mathfrak q'$ 위에 놓인다.
그러므로 $g \in S'$의 $S'_1$에서의 상은 $\mathfrak r'$에 들어 있지 않다.

\medskip\noindent
기초변환 $R'\otimes_R S$도 비슷한 곱 분해를 물려받는다.
$$
R' \otimes_R S
=
(R' \otimes_R S)_{i(s')}
\times
(R' \otimes_R S)_{h(s')}
=
S_1 \times S_2
$$
위의 논의로부터 $S_1$에는 $\mathfrak p'$ 위에 놓이는 소 아이디얼
$\mathfrak r$가 정확히 하나 있고(위와 같이 섬유 환을 생각하라),
이 소 아이디얼은 $\mathfrak q$ 위에 놓인다.

\medskip\noindent
이제 환 사상 $R' \to S'_1 \to S_1$, 소 아이디얼 $\mathfrak p'$,
원소 $g$에 보조정리 \ref{algebra-lemma-produce-finite}를 적용할 수 있다.
따라서 $R'$를 주국소화로 바꾼 뒤에는, 원하는 대로 $S_1$이
$R'$ 위에서 유한이라고 가정할 수 있다.
\end{proof}

\begin{lemma}
\label{algebra-lemma-etale-makes-quasi-finite-finite}
$R \to S$를 환 사상이라 하고, $\mathfrak p \subset R$를
소 아이디얼이라 하자. $R \to S$가 유한형이라고 가정하자.
% P01-E0905: 동결 원문의 존재 목록 도입부에는 문장부호가 없다.
그러면 다음이 존재한다.
\begin{enumerate}
\item 에탈 환 사상 $R \to R'$,
\item $\mathfrak p$ 위에 놓이는 소 아이디얼
$\mathfrak p' \subset R'$,
\item 곱 분해
$$
R' \otimes_R S = A_1 \times \ldots \times A_n \times B
$$
\end{enumerate}
% P01-E0906: 동결 원문의 성질 목록 도입부에는 문장부호가 없다.
이들은 다음 성질을 갖는다.
\begin{enumerate}
\item $\kappa(\mathfrak p) = \kappa(\mathfrak p')$이다.
\item 각 $A_i$는 $R'$ 위에서 유한이다.
\item 각 $A_i$에는 $\mathfrak p'$ 위에 놓이는 소 아이디얼
$\mathfrak r_i$가 정확히 하나 있다.
% P01-E0907: 동결 원문의 마지막 조건에는 is가 빠져 있다.
\item $R' \to B$는 $\mathfrak p'$ 위에 놓이는 어느 소 아이디얼에서도
준유한이 아니다.
\end{enumerate}
\end{lemma}

\begin{proof}
% P01-E0908: 동결 원문의 Denote에는 by가 빠져 있다.
$F = S \otimes_R \kappa(\mathfrak p)$를 $\mathfrak p$에서
$S/R$의 섬유 환이라 하자. $F$는 $\kappa(\mathfrak p)$ 위에서
유한형이므로 뇌터이고, 따라서 $\Spec(F)$에는 고립 닫힌점이
유한 개뿐이다. 고립 닫힌점이 하나도 없다면, 즉 $\mathfrak p$ 위에
놓이고 $S/R$가 $\mathfrak q$에서 준유한인 $S$의 소 아이디얼
$\mathfrak q$가 없다면 보조정리는 성립한다.
그런 소 아이디얼 $\mathfrak q$가 하나라도 있으면
보조정리 \ref{algebra-lemma-etale-makes-quasi-finite-finite-one-prime}를
적용할 수 있다. 이로부터 다음 도식을 얻는다.
$$
\xymatrix{
S \ar[r] & R'\otimes_R S \ar@{=}[r] & A_1 \times B' \\
R \ar[r] \ar[u] & R' \ar[u] \ar[ru]
}
$$
이는 앞의 보조정리에 주어진 것과 같다.
$\mathfrak p$와 $\mathfrak p'$에서의 잉여체가 같으므로
$S/R$와 $(A_1 \times B')/R'$의 섬유 환은 같다.
따라서 섬유의 고립 닫힌점의 개수에 대한 귀납법으로
$R' \to B'$와 $\mathfrak p'$에 대해 보조정리가 성립한다고
가정할 수 있다. 그러면 에탈 환 사상 $R' \to R''$,
소 아이디얼 $\mathfrak p'' \subset R''$, 그리고 분해
$$
R'' \otimes_{R'} B' = A_2 \times \ldots \times A_n \times B
$$
를 얻는다.
% P01-E0909: 동결 원문의 위 분해식 뒤에는 문장 종결 부호가 없다.
환 사상 $R \to R''$, 소 아이디얼 $\mathfrak p''$, 그리고 그로부터
생기는 분해
$$
R'' \otimes_R S = (R'' \otimes_{R'} A_1) \times
A_2 \times \ldots \times A_n \times B
$$
가 보조정리의 문제에 대한 해임을 확인하는 과정은 생략한다.
\end{proof}

\begin{lemma}
\label{algebra-lemma-etale-makes-quasi-finite-finite-variant}
$R \to S$를 환 사상이라 하고, $\mathfrak p \subset R$를
소 아이디얼이라 하자.
% P01-E0910: 동결 원문의 유한형 가정에는 is가 빠져 있다.
$R \to S$가 유한형이라고 가정하자.
% P01-E0911: 동결 원문의 존재 목록 도입부에는 문장부호가 없다.
그러면 다음이 존재한다.
\begin{enumerate}
\item 에탈 환 사상 $R \to R'$,
\item $\mathfrak p$ 위에 놓이는 소 아이디얼
$\mathfrak p' \subset R'$,
\item 곱 분해
$$
R' \otimes_R S = A_1 \times \ldots \times A_n \times B
$$
\end{enumerate}
% P01-E0912: 동결 원문의 성질 목록 도입부에는 문장부호가 없다.
이들은 다음 성질을 갖는다.
\begin{enumerate}
\item 각 $A_i$는 $R'$ 위에서 유한이다.
\item 각 $A_i$에는 $\mathfrak p'$ 위에 놓이는 소 아이디얼
$\mathfrak r_i$가 정확히 하나 있다.
\item 유한 체 확대
$\kappa(\mathfrak r_i)/\kappa(\mathfrak p')$는 순수 비분리이다.
% P01-E0913: 동결 원문의 마지막 조건에는 is가 빠져 있다.
\item $R' \to B$는 $\mathfrak p'$ 위에 놓이는 어느 소 아이디얼에서도
준유한이 아니다.
\end{enumerate}
\end{lemma}

\begin{proof}
증명의 전략은 에탈 환 확대를 두 번 만드는 것이다. 먼저 잉여체들을
조절한 뒤 보조정리 \ref{algebra-lemma-etale-makes-quasi-finite-finite}를 적용한다.

\medskip\noindent
% P01-E0914: 동결 원문의 Denote에는 by가 빠져 있다.
$F = S \otimes_R \kappa(\mathfrak p)$를 $\mathfrak p$에서
$S/R$의 섬유 환이라 하자.
보조정리 \ref{algebra-lemma-etale-makes-quasi-finite-finite}의 증명에서와 같이
% P01-E0915: 동결 원문의 finitely may는 finitely many의 오자이다.
% P01-E0916: 동결 원문은 이 섬유의 소 아이디얼들이 p가 아니라 R 위에 놓인다고 적는다.
$R$ 위에 놓이고 환 사상 $R \to S$가 준유한이 되는 $S$의 소 아이디얼은
유한 개뿐이며, 이들을 $\mathfrak q_1, \ldots, \mathfrak q_n$이라 하자.
$\kappa(\mathfrak p) \subset L_i \subset \kappa(\mathfrak q_i)$를
$\kappa(\mathfrak p) \subset L_i$가 분리이고
$\kappa(\mathfrak q_i)/L_i$가 순수 비분리가 되게 하는 부분체라 하자.
$L/\kappa(\mathfrak p)$를 모든 $i = 1, \ldots, n$에 대해
$L_i$가 매장되는 유한 갈루아 확대라 하자.
보조정리 \ref{algebra-lemma-make-etale-map-prescribed-residue-field}에 의해
에탈 환 확대 $R \to R'$와 $\mathfrak p$ 위에 놓이는
소 아이디얼 $\mathfrak p'$를 찾아, 체 확대
$\kappa(\mathfrak p')/\kappa(\mathfrak p)$가
$\kappa(\mathfrak p) \subset L$과 동형이 되게 할 수 있다.
따라서 $\mathfrak p'$에서 $R' \otimes_R S$의 섬유 환은
$F \otimes_{\kappa(\mathfrak p)} L$과 동형이다.
$\mathfrak q_i$ 위에 놓이는 소 아이디얼들은
$\kappa(\mathfrak q_i) \otimes_{\kappa(\mathfrak p)} L$의
소 아이디얼들에 대응한다. $L$의 선택과 초등적인 체론에 의해
이 환은 $L$ 위의 순수 비분리 체들의 곱이다.
보조정리 \ref{algebra-lemma-quasi-finite-base-change}에 의해 이들은
$\mathfrak p'$ 위에서 $R' \to R' \otimes_R S$가 준유한인
소 아이디얼 전부이기도 하다.
그러므로 $R$을 $R'$로, $\mathfrak p$를 $\mathfrak p'$로,
$S$를 $R' \otimes_R S$로 바꾼 뒤에는, $\mathfrak p$ 위에 놓이고
$S/R$가 준유한인 모든 소 아이디얼 $\mathfrak q$에 대해
체 확대 $\kappa(\mathfrak q)/\kappa(\mathfrak p)$가
순수 비분리라고 가정할 수 있다.

\medskip\noindent
이제 보조정리 \ref{algebra-lemma-etale-makes-quasi-finite-finite}를 적용한다.
이 에탈 환 확대 아래에서 체 확대들은 변하지 않으므로 결과는
바로 원하는 것이다.
\end{proof}






\section{국소 준동형사상}
\label{algebra-section-local-homomorphisms}

\noindent
자연스럽게 들어갈 절이 없는 몇 가지 보조정리를 모아 둔다.
첫 번째 보조정리는 대략적으로 말해, 국소환 사이의 에탈 사상은
비단위원이 생성하는 주 아이디얼의 모든 거듭제곱을 법으로 하면
동형사상이 된다는 것이다.

\begin{lemma}
\label{algebra-lemma-lindel}
\begin{reference}
\cite[Lemma on page 321]{Lindel}, \cite[Lemma 4.1.5]{KC}
\end{reference}
$(R, \mathfrak m_R) \to (S, \mathfrak m_S)$를 국소환들의
국소 준동형사상이라 하자. $S$가 $R$의 어떤 에탈 환 확대의
국소화이고, $\kappa(\mathfrak m_R) \to \kappa(\mathfrak m_S)$가
동형사상이라고 가정하자. 그러면 어떤 $t \in \mathfrak m_R$가 존재하여
모든 $n \geq 1$에 대해 $R/t^nR \to S/t^nS$가 동형사상이다.
% P01-E0917: the frozen source uses the wrong indefinite article before t.
% P01-E0918: the frozen source misspells “isomorphism”.
\end{lemma}

\begin{proof}
어떤 에탈 $R$-대수 $T$와 $\mathfrak m_R$ 위에 놓이는 소 아이디얼
$\mathfrak q \subset T$에 대해 $S = T_{\mathfrak q}$로 쓰자.
명제 \ref{algebra-proposition-etale-locally-standard}에 의해
$R \to T$가 표준 에탈이라고 가정해도 된다.
정의 \ref{algebra-definition-standard-etale}에서처럼
$T = R[x]_g/(f)$로 쓰자. 잉여체에 관한 가정에 의해
$x$와 $a$가
$\kappa(\mathfrak q) = \kappa(\mathfrak m_S) = \kappa(\mathfrak m_R)$에서
같은 상을 갖도록 $a \in R$를 고를 수 있다.
따라서 $x$를 $x - a$로 바꾼 뒤에는
$\mathfrak q$가 $T$에서 $x$와 $\mathfrak m_R$로 생성된다고
가정할 수 있다. 특히 $t = f(0) \in \mathfrak m_R$이다.
$t = f(0)$가 조건을 만족함을 보이겠다.

\medskip\noindent
$f = x^d + \sum_{i = 1, \ldots, d - 1} a_i x^i + t$로 쓰자.
$R \to T$가 표준 에탈이므로 $a_1$은 $R$의 단위원이다.
실제로 $f$의 도함수는 $T$에서 가역이므로, 특히
$\mathfrak q$에 들어가지 않는다.
% P01-E0919: the frozen source runs the derivative clauses together.
$h = a_1 + a_2 x + \ldots + a_{d - 1} x^{d - 2} + x^{d - 1} \in R[x]$라
놓으면 $R[x]$에서 $f = t + xh$이다.
$h \not \in \mathfrak q$이므로 $T$를
$R[x]_{hg}/(f)$로 바꾸어도 된다. 이와 같이 바꾸면
$$
T/tT = (R/tR)[x]_{hg}/(f) = (R/tR)[x]_{hg}/(xh) = (R/tR)[x]_{hg}/(x)
$$
이고, 이는 $R/tR$의 몫이다.
보조정리 \ref{algebra-lemma-surjective-mod-locally-nilpotent}에 의해
모든 $n \geq 1$에 대해 $R/t^nR \to T/t^nT$가 전사임을 얻는다.
한편 평탄한 국소환 사상 $R/t^nR \to S/t^nS$는 모든 $n$에 대해
$R/t^nR \to T/t^nT$를 거쳐 인수분해됨을 알고 있으므로,
이 사상들은 단사이기도 하다(국소환들의 평탄한 국소 준동형사상은
충실히 평탄하고 따라서 단사이다. 보조정리
\ref{algebra-lemma-local-flat-ff}와
\ref{algebra-lemma-faithfully-flat-universally-injective}를 보라).
$S$는 $T$의 국소화이므로 $S/t^nS$는
$T/t^nT = R/t^nR$를 극대 아이디얼 위에 놓이는 한 소 아이디얼에서
국소화한 것이다. 그러나 이 환은 이미 국소환이므로 증명이 끝난다.
\end{proof}

\begin{lemma}
\label{algebra-lemma-etale-under-finite-flat}
$(R, \mathfrak m_R) \to (S, \mathfrak m_S)$를 국소환들의
국소 준동형사상이라 하자. $S$가 $R$의 어떤 에탈 환 확대의
국소화라고 가정하자. 그러면 유한이고 유한 표시이며 충실히 평탄한
환 사상 $R \to S'$가 존재하여, $S'$의 모든 극대 아이디얼
$\mathfrak m'$에 대해 환 사상 $R \to S'_{\mathfrak m'}$의
인수분해
$$
R \to S \to S'_{\mathfrak m'}.
$$
가 존재한다.
\end{lemma}

\begin{proof}
어떤 에탈 $R$-대수 $T$에 대해 $S = T_{\mathfrak q}$로 쓰자.
명제 \ref{algebra-proposition-etale-locally-standard}에 의해
$T$가 표준 에탈이라고 가정해도 된다.
환 사상 $R \to T$에 보조정리
\ref{algebra-lemma-standard-etale-finite-flat-Zariski}를 적용하여
$R \to S'$를 얻는다. 그러면 특히 $S'$의 모든 극대 아이디얼
$\mathfrak m'$에 대해, 어떤 $g' \not \in \mathfrak m'$와
$\mathfrak q = \varphi^{-1}(\mathfrak m'S'_{g'})$를 만족하는
인수분해 $\varphi : T \to S'_{g'}$를 얻는다.
따라서 $\varphi$는 원하는 국소환 사상
$S \to S'_{\mathfrak m'}$를 유도한다.
\end{proof}





\section{정수적 폐포와 매끄러운 기초변환}
\label{algebra-section-integral-closure-smooth-base-change}

\begin{lemma}
\label{algebra-lemma-trick}
$R$를 환이라 하자.
$f \in R[x]$를 최고차항계수가 1인 다항식이라 하자.
$R \to B$를 환 사상이라 하자.
$h \in B[x]/(f)$가 $R$ 위에서 정수적이면, 원소
$f' h$를 $f'h = \sum_i b_i x^i$로 쓸 수 있으며
각 $b_i \in B$는 $R$ 위에서 정수적이다.
\end{lemma}

\begin{proof}
환 $B[x]/(f)$에서 어떤 $r_i \in R$에 대해
$h^e + r_1 h^{e - 1} + \ldots + r_e = 0$이라고 하자.
어떤 $\alpha_i \in B'$에 대해
$f = (x - \alpha_1) \ldots (x - \alpha_d)$가 되게 하는
유한 자유 환 확대 $B \subset B'$가 존재한다.
보조정리 \ref{algebra-lemma-adjoin-roots}를 보라.
각 $\alpha_i$는 $R$ 위에서 정수적임에 유의하자.
어떤 $h_i \in B$를 사용하여
$h = h_0 + h_1 x + \ldots + h_{d - 1} x^{d - 1}$로 나타낼 수 있다.
그러면 $B'[x]/(f)$의 원소로서 다음 보편적 항등식이 성립한다.
$$
f' h =
\sum\nolimits_{i = 1, \ldots, d}
h(\alpha_i)
(x - \alpha_1) \ldots \widehat{(x - \alpha_i)} \ldots (x - \alpha_d)
$$
이는 환 $B_{univ} = \mathbf{Z}[\alpha_i, h_j]$ 위의 점
$\alpha_i$들에서 양변을 계산하여 증명할 수 있다
(몇몇 세부사항은 생략한다).
$h$가 환 $B[x]/(f)$에서
$h^e + r_1 h^{e - 1} + \ldots + r_e = 0$을 만족한다는 가정으로부터
$$
h(\alpha_i)^e + r_1 h(\alpha_i)^{e - 1} + \ldots + r_e = 0
$$
이 $B'$에서 성립함을 알 수 있다.
따라서 $h(\alpha_i)$는 $R$ 위에서 정수적이다.
위 공식을 사용하면 $B'[x]/(f)$에서
$f'h \equiv \sum_{j = 0, \ldots, d - 1} b'_j x^j$이고,
각 $b'_j \in B'$는 $R$ 위에서 정수적임을 알 수 있다.
그러나 $f' h \in B[x]/(f)$이고
$1, x, \ldots, x^{d - 1}$은 $B'[x]/(f)$의 $B'$-기저이므로,
원하는 대로 $b'_j \in B$이다.
\end{proof}

\begin{lemma}
\label{algebra-lemma-integral-closure-commutes-etale}
$R \to S$를 에탈 환 사상이라 하자.
$R \to B$를 임의의 환 사상이라 하자.
$A \subset B$를 $B$ 안에서 $R$의 정수적 폐포라 하자.
$A' \subset S \otimes_R B$를 $S \otimes_R B$ 안에서
$S$의 정수적 폐포라 하자. 그러면 표준 사상
$S \otimes_R A \to A'$는 동형사상이다.
\end{lemma}

\begin{proof}
$A \subset B$이고 $R \to S$가 평탄하므로
사상 $S \otimes_R A \to A'$는 단사이다.
정수적 폐포를 취하는 연산은 국소화와 교환한다는 사실을
반복해서 사용할 것이다. 보조정리
\ref{algebra-lemma-integral-closure-localize}를 보라.
따라서 보조정리 \ref{algebra-lemma-cover}($S$-가군 사상이 동형사상인지
판정하는 기준)에 의해 $S$ 위에서 국소화해도 된다.
그러므로 $S = R[x]_g/(f) = (R[x]/(f))_g$가
$R$ 위에서 표준 에탈이라고 가정해도 된다.
명제 \ref{algebra-proposition-etale-locally-standard}를 보라.
한 번 더 국소화하면 $A'$는 $(A'')_g$인데,
여기서 $A''$은 $B[x]/(f)$ 안에서 $R[x]/(f)$의 정수적 폐포이다.
$a \in A''$라 하자. $a$가 $S \otimes_R A$에 속함을 보이면 충분하다.
보조정리 \ref{algebra-lemma-trick}에 의해
$f' a = \sum a_i x^i$이며 $a_i \in A$이다.
$f'$는 $S$에서 가역이므로(표준 에탈 환 사상의 정의에 의해),
원하는 대로 $a \in S \otimes_R A$이다.
\end{proof}

\begin{example}
\label{algebra-example-fourier}
$p$를 소수라 하자. 환 확대
$$
R = \mathbf{Z}[1/p] \subset
R' = \mathbf{Z}[1/p][x]/(x^{p - 1} + \ldots + x + 1)
$$
는 다음 성질을 갖는다. $d < p$이면
$$
\prod\nolimits_{0 \leq i < j < d} (\alpha_i - \alpha_j)
$$
가 $R'$의 단위원이 되게 하는 원소
$\alpha_0, \ldots, \alpha_{d - 1} \in R'$가 존재한다.
구체적으로 $i = 0, \ldots, p - 1$에 대해
$\alpha_i$를 $R'$에서 $x^i$의 동치류로 택한다. 그러면
$R'[T]$에서
$$
T^p - 1 = \prod\nolimits_{i = 0, \ldots, p - 1} (T - \alpha_i)
$$
이다. 실제로 원분다항식
$x^{p - 1} + \ldots + x + 1$이 $\mathbf{Q}$ 위에서 기약이므로
환 $\mathbf{Q}[x]/(x^{p - 1} + \ldots + x + 1)$은 체이고,
$\alpha_i$들은 $T^p - 1$의 서로 다른 근들이므로 이 등식이 성립한다.
양변을 미분하고 $T = \alpha_i$를 대입하면
% P01-E0920: the frozen differentiated product repeats alpha_1 at both endpoints.
$$
p \alpha_i^{p - 1}
=
(\alpha_i - \alpha_1) \ldots
\widehat{(\alpha_i - \alpha_i)} \ldots
(\alpha_i - \alpha_1)
$$
을 얻고, 이것이 $R'$에서 가역임을 알 수 있다.
\end{example}

\begin{lemma}
\label{algebra-lemma-integral-closure-commutes-smooth}
$R \to S$를 매끄러운 환 사상이라 하자.
$R \to B$를 임의의 환 사상이라 하자.
$A \subset B$를 $B$ 안에서 $R$의 정수적 폐포라 하자.
$A' \subset S \otimes_R B$를 $S \otimes_R B$ 안에서
$S$의 정수적 폐포라 하자. 그러면 표준 사상
$S \otimes_R A \to A'$는 동형사상이다.
\end{lemma}

\begin{proof}
보조정리 \ref{algebra-lemma-integral-closure-commutes-etale}의 증명에서처럼
논증하면 $S$ 위에서 국소화해도 된다.
따라서 $R \to S$가 표준 매끄러운 환 사상이라고 가정할 수 있다.
보조정리 \ref{algebra-lemma-smooth-syntomic}을 보라.
표준 매끄러운 환 사상의 정의에 의해 $S$는 다항식환
$R[x_1, \ldots, x_n]$ 위에서 에탈이다.
에탈 환 확대의 경우에는 이미 결과를 보였으므로
(보조정리 \ref{algebra-lemma-integral-closure-commutes-etale}),
$S = R[x]$인 경우로 환원된다. 따라서 다음을 보이면 된다.
$$
f = \sum b_i x^i
\text{ integral over }R[x]
\Leftrightarrow
\text{each }b_i\text{ integral over }R.
$$
오른쪽에서 왼쪽으로의 함의는 $B[x]$에서 $R[x]$ 위에 정수적인
원소들의 집합이 환이고(보조정리
\ref{algebra-lemma-integral-closure-is-ring}) $x$를 포함하기 때문에 성립한다.

\medskip\noindent
$f \in B[x]$가 $R[x]$ 위에서 정수적이고
$f = \sum_{i < d} b_i x^i$의 차수가 $< d$라고 가정하자.
정수적 폐포와 국소화는 교환하므로, 서로 다른 소수 $p, q$가
존재하여 각 $b_i$가 $R[1/p]$와 $R[1/q]$ 양쪽 위에서
정수적임을 보이면 충분하다. 따라서 유한 자유 환 확대
$R \subset R'$를 찾을 수 있는데,
% P01-E0921: the frozen transition does not explicitly pass to R[1/p] or R[1/q].
$R'$는 $\alpha_1, \ldots, \alpha_d$를 포함하고
$\prod_{i < j} (\alpha_i - \alpha_j)$는 $R'$의 단위원이다.
예 \ref{algebra-example-fourier}를 보라.
이 경우 다음 보편적 등식이 성립한다.
$$
f
=
\sum_i
f(\alpha_i)
\frac{(x - \alpha_1) \ldots \widehat{(x - \alpha_i)} \ldots (x - \alpha_d)}
{(\alpha_i - \alpha_1) \ldots \widehat{(\alpha_i - \alpha_i)} \ldots
(\alpha_i - \alpha_d)}.
$$
% P01-E0922: the frozen source begins the next sentence with the editorial artifact “OK,”.
그리고 원소 $f(\alpha_i)$들은 $R'$ 위에서 정수적이다.
실제로 $(R' \otimes_R B)[x] \to R' \otimes_R B$,
$h \mapsto h(\alpha_i)$는 환 사상이다.
따라서 $(R' \otimes_R B)[x]$에서 $f$의 계수들이
$R'$ 위에서 정수적임을 알 수 있다.
$R'$는 $R$ 위에서 유한이므로(따라서 $R$ 위에서 정수적이므로)
그 계수들은 $R$ 위에서도 정수적이다. 이것이 원하는 결론이다.
\end{proof}

\begin{lemma}
\label{algebra-lemma-integral-closure-commutes-colim-smooth}
$R \to S$와 $R \to B$를 환 사상이라 하자.
$A \subset B$를 $B$ 안에서 $R$의 정수적 폐포라 하자.
$A' \subset S \otimes_R B$를 $S \otimes_R B$ 안에서
$S$의 정수적 폐포라 하자. $S$가 매끄러운 $R$-대수들의
여과 여극한이면, 표준 사상 $S \otimes_R A \to A'$는
동형사상이다.
\end{lemma}

\begin{proof}
텐서곱을 취하는 연산과 정수적 폐포를 취하는 연산이
여과 여극한과 교환한다는 직접적인 사실과
보조정리 \ref{algebra-lemma-integral-closure-commutes-smooth}로부터 따른다.
% P01-E0923: the frozen English compound subject has singular “commutes”.
\end{proof}






\section{형식적으로 비분기인 사상}
\label{algebra-section-formally-unramified}

\noindent
형식적 에탈 사상의 개념을 도입하기 전에 형식적으로 비분기인
사상의 개념을 정의하는 편이 논리적으로 더 효율적임이 드러난다.

\begin{definition}
\label{algebra-definition-formally-unramified}
$R \to S$를 환 사상이라 하자.
모든 가환 실선 도식
$$
\xymatrix{
S \ar[r] \ar@{-->}[rd] & A/I \\
R \ar[r] \ar[u] & A \ar[u]
}
$$
에서 $I \subset A$가 제곱 영 아이디얼일 때, 도식을 가환으로 만드는
점선 사상이 기껏해야 하나 존재하면 $S$가 $R$ 위에서
\textit{형식적으로 비분기이다}라고 한다.
\end{definition}

\begin{lemma}
\label{algebra-lemma-base-change-formally-unramified}
$R \to S$를 형식적으로 비분기인 사상이라 하자.
$R \to R'$를 임의의 환 사상이라 하자.
그러면 기저변환 $S' = R' \otimes_R S$는 $R'$ 위에서
형식적으로 비분기이다.
\end{lemma}

\begin{proof}
정의 \ref{algebra-definition-formally-unramified}에서와 같은 실선 도식
$$
\xymatrix{
S \ar[r] \ar@{-->}[rrd] & R' \otimes_R S \ar[r] \ar@{-->}[rd] & A/I \\
R  \ar[u] \ar[r] & R' \ar[r] \ar[u] & A \ar[u]
}
$$
이 주어졌다고 하자. 가정에 의해 더 긴 점선 사상은 기껏해야 하나
존재한다. 텐서곱의 보편 성질에 의해 더 짧은 점선 사상도
기껏해야 하나 존재한다.
\end{proof}

\begin{lemma}
\label{algebra-lemma-characterize-formally-unramified}
$R \to S$를 환 사상이라 하자.
다음 조건들은 서로 동치이다.
\begin{enumerate}
\item $R \to S$는 형식적으로 비분기이다.
\item 미분 가군 $\Omega_{S/R}$는 영이다.
\end{enumerate}
\end{lemma}

\begin{proof}
$J = \Ker(S \otimes_R S \to S)$를 곱셈 사상의 핵이라 하자.
$A_{univ} = S \otimes_R S/J^2$라 하자.
$I_{univ} = J/J^2$가 $\Omega_{S/R}$와 동형임을 상기하자.
보조정리 \ref{algebra-lemma-differentials-diagonal}을 보라.
또한 두 $R$-대수 사상
$\sigma_1, \sigma_2 : S \to A_{univ}$,
$\sigma_1(s) = s \otimes 1 \bmod J^2$와
$\sigma_2(s) = 1 \otimes s \bmod J^2$의 차이는
보편 미분 $\text{d} : S \to \Omega_{S/R} = I_{univ}$이다.

\medskip\noindent
$R \to S$가 형식적으로 비분기라고 가정하자.
그러면 $\sigma_1 = \sigma_2$임을 알 수 있다.
따라서 모든 $s \in S$에 대해 $\text{d}(s) = 0$이다.
그러므로 $\Omega_{S/R} = 0$이다.

\medskip\noindent
$\Omega_{S/R} = 0$이라고 가정하자.
$A, I, R \to A, S \to A/I$가 정의
\ref{algebra-definition-formally-unramified}에서와 같은 실선 도식이라고 하자.
$\tau_1, \tau_2 : S \to A$를 이 도식을 가환으로 만드는 두 점선
사상이라 하자. 규칙
$s_1 \otimes s_2 \mapsto \tau_1(s_1)\tau_2(s_2)$로 정의되는
$R$-대수 사상 $A_{univ} \to A$를 생각하자.
이 사상이 잘 정의된다는 확인은 생략한다.
$I_{univ} = \Omega_{S/R} = 0$이므로 $A_{univ} \cong S$이고,
따라서 $\tau_1 = \tau_2$이다.
\end{proof}

\begin{lemma}
\label{algebra-lemma-formally-unramified-local}
$R \to S$를 환 사상이라 하자.
다음 조건들은 서로 동치이다.
\begin{enumerate}
\item $R \to S$는 형식적으로 비분기이다.
\item $S$의 모든 소 아이디얼 $\mathfrak q$에 대해
$R \to S_{\mathfrak q}$가 형식적으로 비분기이다.
\item $\mathfrak p = R \cap \mathfrak q$인 $S$의 모든 소 아이디얼
$\mathfrak q$에 대해 $R_{\mathfrak p} \to S_{\mathfrak q}$가
형식적으로 비분기이다.
\end{enumerate}
\end{lemma}

\begin{proof}
보조정리 \ref{algebra-lemma-characterize-formally-unramified}에서
(1)이 $\Omega_{S/R} = 0$과 동치임을 보았다.
마찬가지로 보조정리 \ref{algebra-lemma-differentials-localize}에 의해
(2)와 (3)은 모든 $\mathfrak q$에 대해
$(\Omega_{S/R})_{\mathfrak q} = 0$이라는 조건과 동치이다.
따라서 동치성은 보조정리
\ref{algebra-lemma-characterize-zero-local}로부터 따른다.
\end{proof}

\begin{lemma}
\label{algebra-lemma-formally-unramified-localize}
$A \to B$를 형식적으로 비분기인 환 사상이라 하자.
\begin{enumerate}
\item 곱셈적 부분집합 $S \subset A$에 대해
$S^{-1}A \to S^{-1}B$는 형식적으로 비분기이다.
\item 곱셈적 부분집합 $S \subset B$에 대해
$A \to S^{-1}B$는 형식적으로 비분기이다.
\end{enumerate}
\end{lemma}

\begin{proof}
두 명제 모두 보조정리
\ref{algebra-lemma-formally-unramified-local}로부터 따른다.
% P01-E0924: the frozen proof begins with a sentence fragment lacking a subject.
(보조정리 \ref{algebra-lemma-characterize-formally-unramified}와
보조정리 \ref{algebra-lemma-differentials-localize}를 결합하여
도출할 수도 있다.)
\end{proof}

\begin{lemma}
\label{algebra-lemma-colimit-formally-unramified}
$R$를 환이라 하고 $I$를 유향 집합이라 하자.
$(S_i, \varphi_{ii'})$를 $I$ 위의 $R$-대수들의 계라 하자.
각 $R \to S_i$가 형식적으로 비분기이면,
$S = \colim_{i \in I} S_i$는 $R$ 위에서 형식적으로 비분기이다.
% P01-E0925: the frozen lemma statement lacks terminal punctuation.
\end{lemma}

\begin{proof}
정의 \ref{algebra-definition-formally-unramified}에서와 같은 도식을 생각하자.
가정에 의해 합성 $S_i \to S \to A/I$를 올리는
$R$-대수 사상 $S_i \to A$는 기껏해야 하나 존재한다.
$S$의 모든 원소는 사상들 $S_i \to S$ 중 하나의 상에 들어가므로,
도식에 들어맞는 사상 $S \to A$도 기껏해야 하나 존재한다.
\end{proof}



\section{코노멀 가군과 보편 두꺼워짐}
\label{algebra-section-conormal}

\noindent
스킴의 닫힌 몰입에 대해서뿐 아니라 임의의 형식적 비분기 사상에
대해서도 첫 무한소 근방을 정의할 수 있음이 드러난다.
이는 다음 대수적 사실에 바탕을 둔다.

\begin{lemma}
\label{algebra-lemma-universal-thickening}
$R \to S$를 형식적으로 비분기인 환 사상이라 하자.
핵이 제곱 영 아이디얼인 $R$-대수들의 전사
$S' \to S$ 중 다음 보편 성질을 갖는 것이 존재한다.
임의의 가환 도식
$$
\xymatrix{
S \ar[r]_a & A/I \\
R \ar[r]^b \ar[u] & A \ar[u]
}
$$
이 주어지고 $I \subset A$가 제곱 영 아이디얼이면,
$S' \to A \to A/I$가 $S' \to S \to A/I$와 같아지게 하는
유일한 $R$-대수 사상 $a' : S' \to A$가 존재한다.
\end{lemma}

\begin{proof}
$R$-대수로서 $S$의 생성원 집합 $z_i \in S$, $i \in I$를 택하자.
$P = R[\{x_i\}_{i \in I}]$를 생성원 $x_i$, $i \in I$ 위의
다항식환이라 하자. $x_i$를 $z_i$로 보내는
$R$-대수 사상 $P \to S$를 생각하자.
$J = \Ker(P \to S)$라 하자. 사상
$$
\text{d} : J/J^2 \longrightarrow \Omega_{P/R} \otimes_P S
$$
을 생각하자. 보조정리 \ref{algebra-lemma-differential-seq}를 보라.
가정과 보조정리
\ref{algebra-lemma-characterize-formally-unramified}에 의해
$\Omega_{S/R} = 0$이므로 이 사상은 전사이다.
$\Omega_{P/R}$는 $\text{d}x_i$들 위에서 자유이고,
따라서 가군 $\Omega_{P/R} \otimes_P S$는 $S$ 위에서 자유임에 유의하자.
그러므로 위 전사의 분할을 택하여
$$
J/J^2 = K \oplus \Omega_{P/R} \otimes_P S
$$
라고 쓸 수 있다.
% P01-E0926: the frozen proof reuses I for the generator index set and the square-zero ideal.
% P01-E0927: the frozen display lacks terminal punctuation.
$J^2 \subset J' \subset J$를 위 분해에서 $J'/J^2$가
두 번째 직합인자가 되는 $P$의 아이디얼이라 하자.
$S' = P/J'$로 두자. 짧은 완전열
$$
0 \to J/J' \to S' \to S \to 0
$$
을 얻으며, $J/J' \cong K$가 $S'$의 제곱 영 아이디얼임을 알 수 있다.
따라서
$$
\xymatrix{
S \ar[r]_1 & S \\
R \ar[r] \ar[u] & S' \ar[u]
}
$$
는 위와 같은 도식이다. 사실 이것이 도식들의 범주에서 시초 대상임을
주장한다. 구체적으로 임의의 도식 $(I \subset A, a, b)$를 택하자.
다음 도식이 가환이 되게 하는 $R$-대수 사상
$\beta : P \to A$를 택할 수 있다.
$$
\xymatrix{
S \ar[r]_1 & S \ar[r]_a & A/I \\
R \ar[r] \ar@/_/[rr]_b \ar[u] & P \ar[u] \ar[r]^\beta & A \ar[u]
}
$$
이제 $\beta(J') = 0$이라고 할 수는 없으며, 다시 말해
$\beta$가 $S' = P/J'$를 거쳐 인수분해한다고 할 수는 없다.
그러나 $\beta(J') \subset I$임은 분명하고,
$\beta(J) \subset I$와 $I^2 = 0$으로부터
$\beta(J^2) = 0$을 얻는다. 따라서
$(J/J' \subset S', 1, R \to S')$에서
$(I \subset A, a, b)$로 가는 사상을 찾는 데 대한 “장애물”은
대응하는 $S$-선형 사상
$\overline{\beta} : J'/J^2 \to I$이다.
$\beta$를 택할 때의 선택은 $\beta(x_i)$의 선택에 놓여 있다.
$\beta$의 다른 선택 $\beta'$는
$\delta_i \in I$에 대해
$\beta'(x_i) = \beta(x_i) + \delta_i$로 택하여 얻어진다.
이 경우 $g \in J'$에 대해
$$
\beta'(g) =
\beta(g) + \sum\nolimits_i \delta_i \beta(\frac{\partial g}{\partial x_i})
$$
을 얻는다.
% P01-E0928: the frozen Taylor display lacks terminal punctuation.
구성상 사상
$\text{d}|_{J'/J^2} : J'/J^2 \to \Omega_{P/R} \otimes_P S$는
$g \mapsto \frac{\partial g}{\partial x_i}\text{d}x_i \otimes 1$로
주어지는 동형사상이므로, 모든 $g \in J'$에 대해
$\beta'(g) = 0$이 되게 하는 $\delta_i \in I$의 선택이 유일하게
존재함을 알 수 있다.
% P01-E0929: the frozen differential formula has a free i and omits the sum.
(구체적으로 $\delta_i$는 $-\overline{\beta}(g)$인데,
여기서 $g \in J'/J^2$는
$\text{d}|_{J'/J^2}$에 의해 $\text{d}x_i \otimes 1$로
보내지는 유일한 원소이다.)
% P01-E0930: the frozen source leaves the i-dependence of g implicit.
해의 유일성으로부터 보조정리에서 요구한 유일성이 따른다.
\end{proof}

\noindent
보조정리 \ref{algebra-lemma-universal-thickening}의 상황에서
$R$-대수 사상 $S' \to S$는 유일한 동형사상을 제외하고 유일하다.

\begin{definition}
\label{algebra-definition-universal-thickening}
$R \to S$를 형식적으로 비분기인 환 사상이라 하자.
\begin{enumerate}
\item $R$ 위의 $S$의 \textit{보편 1차 두꺼워짐}은
보조정리 \ref{algebra-lemma-universal-thickening}의
$R$-대수 전사 $S' \to S$이다.
\item $R \to S$의 \textit{코노멀 가군}은 보편 1차 두꺼워짐
$S' \to S$의 핵 $I$를 $S$-가군으로 본 것이다.
\end{enumerate}
이 상황에서 코노멀 가군을 흔히 \textit{$C_{S/R}$}로 나타낸다.
\end{definition}

\begin{lemma}
\label{algebra-lemma-universal-thickening-quotient}
$I \subset R$를 환의 아이디얼이라 하자.
$R$ 위의 $R/I$의 보편 1차 두꺼워짐은 전사
$R/I^2 \to R/I$이다. $R$ 위의 $R/I$의 코노멀 가군은
$C_{(R/I)/R} = I/I^2$이다.
\end{lemma}

\begin{proof}
생략한다.
\end{proof}
\begin{lemma}
\label{algebra-lemma-universal-thickening-localize}
$A \to B$를 형식적으로 비분기인 환 사상이라 하자.
$\varphi : B' \to B$를 $A$ 위의 $B$의 보편 1차 두꺼워짐이라 하자.
\begin{enumerate}
\item $S \subset A$를 곱셈적 부분집합이라 하자.
그러면 $S^{-1}B' \to S^{-1}B$는 $S^{-1}A$ 위의
$S^{-1}B$의 보편 1차 두꺼워짐이다.
특히 $S^{-1}C_{B/A} = C_{S^{-1}B/S^{-1}A}$이다.
\item $S \subset B$를 곱셈적 부분집합이라 하자.
그러면 $S' = \varphi^{-1}(S)$는 $B'$의 곱셈적 부분집합이고,
$(S')^{-1}B' \to S^{-1}B$는 $A$ 위의 $S^{-1}B$의
보편 1차 두꺼워짐이다.
특히 $S^{-1}C_{B/A} = C_{S^{-1}B/A}$이다.
\end{enumerate}
보조정리 \ref{algebra-lemma-formally-unramified-localize}에 의해
이 보조정리의 진술은 의미가 있다.
\end{lemma}

\begin{proof}
(1)의 표기와 가정을 사용하자.
% P01-E0931: both frozen proof paragraphs begin with sentence fragments.
$(S^{-1}B)' \to S^{-1}B$를 $S^{-1}A$ 위의 $S^{-1}B$의
보편 1차 두꺼워짐이라 하자.
$S^{-1}B' \to S^{-1}B$는 핵의 제곱이 영인
$S^{-1}A$-대수들의 전사임에 유의하자.
따라서 정의에 의해 $S^{-1}B$로 가는 사상들과 양립하는 사상
$(S^{-1}B)' \to S^{-1}B'$를 얻는다.
임의의 가환 도식
$$
\xymatrix{
B \ar[r] & S^{-1}B \ar[r] & D/I \\
A \ar[r] \ar[u] & S^{-1}A \ar[r] \ar[u] & D \ar[u]
}
$$
을 생각하자. 여기서 $I \subset D$는 제곱 영 아이디얼이다.
$B'$는 $A$ 위의 $B$의 보편 1차 두꺼워짐이므로
$A$-대수 사상 $B' \to D$를 얻는다.
그런데 $D$에서 $S$의 상은 분명히 $D$의 가역원들로 보내지므로
양립하는 사상 $S^{-1}B' \to D$를 얻는다.
이를 $D = (S^{-1}B)'$에 적용하면 사상
$S^{-1}B' \to (S^{-1}B)'$를 얻는다.
이 사상이 위에서 설명한 사상의 역이라는 확인은 생략한다.

\medskip\noindent
(2)의 표기와 가정을 사용하자.
$(S^{-1}B)' \to S^{-1}B$를 $A$ 위의 $S^{-1}B$의
보편 1차 두꺼워짐이라 하자.
$(S')^{-1}B' \to S^{-1}B$는 핵의 제곱이 영인
$A$-대수들의 전사임에 유의하자.
따라서 정의에 의해 $S^{-1}B$로 가는 사상들과 양립하는 사상
$(S^{-1}B)' \to (S')^{-1}B'$를 얻는다.
임의의 가환 도식
$$
\xymatrix{
B \ar[r] & S^{-1}B \ar[r] & D/I \\
A \ar[r] \ar[u] & A \ar[r] \ar[u] & D \ar[u]
}
$$
을 생각하자. 여기서 $I \subset D$는 제곱 영 아이디얼이다.
$B'$는 $A$ 위의 $B$의 보편 1차 두꺼워짐이므로
$A$-대수 사상 $B' \to D$를 얻는다.
그런데 $D$에서 $S'$의 상은 분명히 $D$의 가역원들로 보내지므로
양립하는 사상 $(S')^{-1}B' \to D$를 얻는다.
이를 $D = (S^{-1}B)'$에 적용하면 사상
$(S')^{-1}B' \to (S^{-1}B)'$를 얻는다.
이 사상이 위에서 설명한 사상의 역이라는 확인은 생략한다.
\end{proof}

\begin{lemma}
\label{algebra-lemma-differentials-universal-thickening}
$R \to A  \to B$를 환 사상들이라 하자.
$A \to B$가 형식적으로 비분기라고 가정하자.
$B' \to B$를 $A$ 위의 $B$의 보편 1차 두꺼워짐이라 하자.
그러면 $B'$는 $A$ 위에서 형식적으로 비분기이고, 표준 사상
$\Omega_{A/R} \otimes_A B \to \Omega_{B'/R} \otimes_{B'} B$는
동형사상이다.
\end{lemma}

\begin{proof}
원칙적으로는 이 결과들을 정의로부터 형식적으로 이끌어낼 수 있어야 하지만,
여기서는 보조정리 \ref{algebra-lemma-universal-thickening}의 증명에서
$B'$를 구성한 방식을 사용하겠다.
구체적으로 $P = A[x_i]$가 $A$ 위의 다항식환이 되게 하는 표시
$B = P/J$를 택한다. 다음으로
$\Omega_{P/A} \otimes_P B$에서
$\text{d}f_i = \text{d}x_i \otimes 1$을 만족하는 원소
$f_i \in J$를 택한다. 이렇게 택하면
$J' = (f_i) + J^2$에 대해 $B' = P/J'$이다.
보조정리 \ref{algebra-lemma-universal-thickening}의 증명을 보라.

\medskip\noindent
표준 완전열
$$
J'/(J')^2 \to \Omega_{P/A} \otimes_P B' \to \Omega_{B'/A} \to 0
$$
을 생각하자. 보조정리 \ref{algebra-lemma-differential-seq}를 보라.
구성에 의해 $f_i \in J'$의 동치류들은
$\Omega_{P/A} \otimes_P B'$의 원소들로 보내지며,
그 원소들은 구성상 이 가군을 $J'/J^2$를 법으로 생성한다.
$J'/J^2$가 멱영 아이디얼이므로 이 원소들은 가군 전체를 생성한다
(나카야마 보조정리 \ref{algebra-lemma-NAK}).
% P01-E0932: the frozen Nakayama step uses J'/J^2 where the governing kernel is J/J'.
이는 $\Omega_{B'/A} = 0$임을 증명하고,
따라서 $B'$가 $A$ 위에서 형식적으로 비분기임을 증명한다.
보조정리 \ref{algebra-lemma-characterize-formally-unramified}를 보라.

\medskip\noindent
$P$는 $A$ 위의 다항식환이므로
$\Omega_{P/R} = \Omega_{A/R} \otimes_A P \oplus \bigoplus P\text{d}x_i$이다.
이 분해를 사용하겠다. 다음 완전열을 생각하자.
$$
J'/(J')^2 \to
\Omega_{P/R} \otimes_P B' \to
\Omega_{B'/R} \to 0
$$
보조정리 \ref{algebra-lemma-differential-seq}를 보라.
이를 $B$와 텐서하여 완전열
$$
J'/(J')^2 \otimes_{B'} B \to
\Omega_{P/R} \otimes_P B \to
\Omega_{B'/R} \otimes_{B'} B \to 0
$$
을 얻을 수 있다.
% P01-E0933: the frozen tensored exact sequence lacks terminal punctuation.
$J' = (f_i) + J^2$임을 기억하면 첫 번째 화살표가 부분가군
$J^2/(J')^2$를 영멸함을 알 수 있다.
주어진 직합 분해
$\Omega_{P/R} \otimes_P B =
\Omega_{A/R} \otimes_A B \oplus \bigoplus B\text{d}x_i $의 관점에서 보면,
% P01-E0934: the frozen prose omits “above” and its separating comma.
부분가군 $(f_i)/(J')^2 \otimes_{B'} B$는 직합인자
$\bigoplus B\text{d}x_i$ 위로 동형사상으로 보내진다.
따라서 이 완전열에 남는 것은 원하는 동형사상
$\Omega_{A/R} \otimes_A B \to \Omega_{B'/R} \otimes_{B'} B$이다.
\end{proof}







\section{형식적 에탈 사상}
\label{algebra-section-formally-etale}

\begin{definition}
\label{algebra-definition-formally-etale}
$R \to S$를 환 사상이라 하자.
모든 가환 실선 도식
$$
\xymatrix{
S \ar[r] \ar@{-->}[rd] & A/I \\
R \ar[r] \ar[u] & A \ar[u]
}
$$
에서 $I \subset A$가 제곱 영 아이디얼일 때 도식을 가환으로 만드는
유일한 점선 사상이 존재하면 $S$가 $R$ 위에서
\textit{형식적으로 에탈이다}라고 한다.
\end{definition}

\noindent
환 사상이 형식적으로 에탈인 것은 형식적으로 매끄럽고
형식적으로 비분기인 것과 동치임이 분명하다.

\begin{lemma}
\label{algebra-lemma-base-change-formally-etale}
$R \to S$를 형식적 에탈 사상이라 하자.
$R \to R'$를 임의의 환 사상이라 하자.
그러면 기저변환 $S' = R' \otimes_R S$는
$R'$ 위에서 형식적으로 에탈이다.
\end{lemma}

\begin{proof}
보조정리 \ref{algebra-lemma-base-change-fs}와
\ref{algebra-lemma-base-change-formally-unramified}를 결합하면 된다.
\end{proof}

\begin{lemma}
\label{algebra-lemma-formally-etale-etale}
$R \to S$를 유한 표시인 환 사상이라 하자.
다음 조건들은 서로 동치이다.
\begin{enumerate}
\item $R \to S$는 형식적으로 에탈이다.
\item $R \to S$는 에탈이다.
\end{enumerate}
\end{lemma}

\begin{proof}
$R \to S$가 형식적으로 에탈이라고 가정하자.
그러면 명제 \ref{algebra-proposition-smooth-formally-smooth}에 의해
$R \to S$는 매끄럽다.
보조정리 \ref{algebra-lemma-characterize-formally-unramified}에 의해
$\Omega_{S/R} = 0$이다.
따라서 정의에 의해 $R \to S$는 에탈이다.

\medskip\noindent
$R \to S$가 에탈이라고 가정하자.
그러면 명제 \ref{algebra-proposition-smooth-formally-smooth}에 의해
$R \to S$는 형식적으로 매끄럽다.
보조정리 \ref{algebra-lemma-characterize-formally-unramified}에 의해
형식적으로 비분기이다. 따라서 $R \to S$는 형식적으로 에탈이다.
\end{proof}

\begin{lemma}
\label{algebra-lemma-colimit-formally-etale}
$R$를 환이라 하고 $I$를 유향 집합이라 하자.
$(S_i, \varphi_{ii'})$를 $I$ 위의 $R$-대수들의 계라 하자.
각 $R \to S_i$가 형식적으로 에탈이면,
$S = \colim_{i \in I} S_i$는 $R$ 위에서 형식적으로 에탈이다.
% P01-E0935: the frozen lemma statement lacks terminal punctuation.
\end{lemma}

\begin{proof}
정의 \ref{algebra-definition-formally-etale}에서와 같은 도식을 생각하자.
가정에 의해 각 지표에 대해 합성 $S_i \to S \to A/I$를 올리는
유일한 $R$-대수 사상 $S_i \to A$를 얻는다.
% P01-E0936: the frozen proof obscures the per-index uniqueness of the lifts.
따라서 이 사상들은 전이 사상 $\varphi_{ii'}$들과 양립하고
올림 $S \to A$를 정의한다. 이것으로 존재성이 증명된다.
각 $S_i$로 제한하면 유일성은 분명하다.
\end{proof}

\begin{lemma}
\label{algebra-lemma-localization-formally-etale}
$R$를 환이라 하고 $S \subset R$를 임의의 곱셈적 부분집합이라 하자.
그러면 환 사상 $R \to S^{-1}R$는 형식적으로 에탈이다.
\end{lemma}

\begin{proof}
$I \subset A$를 제곱 영 아이디얼이라 하자.
여기서 주장하는 것은 모든 $f \in S$에 대해
$\varphi(f) \mod I$가 가역이 되게 하는 환 사상
$\varphi : R \to A$가 주어지면, 모든 $f \in S$에 대해
$\varphi(f)$도 $A$에서 가역이라는 것이다.
이는 표준 사상 $A \to A/I$ 아래에서 $A^*$가
$(A/I)^*$의 역상이기 때문에 성립한다.
\end{proof}

\begin{lemma}
\label{algebra-lemma-formally-etale-lift-infinitesimal}
$R \to S$를 환 사상이라 하자. $J \subset S$를
$R \to S/J$가 전사가 되게 하는 아이디얼이라 하고,
$I \subset R$를 그 핵이라 하자.
$R \to S$가 형식적으로 에탈이면, 모든 $n$에 대해
$R/I^n \to S/J^n$는 동형사상이고,
$\bigoplus I^n/I^{n + 1} \to \bigoplus J^n/J^{n + 1}$는
등급환들의 동형사상이다.
\end{lemma}

\begin{proof}
올림 성질을 귀납적으로 사용하면 다음 점선 사상들을 얻는다.
$$
\xymatrix{
S \ar[r] \ar@{-->}[rd] & S/J = R/I \\
R \ar[r] \ar[u] & R/I^2 \ar[u]
}
\quad
\xymatrix{
S \ar[r] \ar@{-->}[rd] & R/I^2 \\
R \ar[r] \ar[u] & R/I^3 \ar[u]
}
\quad
\xymatrix{
S \ar[r] \ar@{-->}[rd] & R/I^3 \\
R \ar[r] \ar[u] & R/I^4 \ar[u]
}
$$
% P01-E0937: the frozen three-diagram display lacks terminal punctuation.
대응하는 사상 $S/J^n \to R/I^n$는 동형사상이다.
실제로 합성 $S/J^n \to R/I^n \to S/J^n$는
형식적 에탈 환 사상의 올림 성질에서 유일성에 의해
(귀납적으로) 항등사상이다.
\end{proof}

\begin{lemma}
\label{algebra-lemma-formally-etale-omega}
$R \to S \to S'$를 환 사상들이라 하자.
$J$와 $J'$를 각각 곱셈 사상
$S \otimes_R S \to S$와
$S' \otimes_{R'} S' \to S'$의 핵이라 하자.
% P01-E0938: the frozen formula uses the undefined base ring R' instead of R.
$S \to S'$가 형식적으로 에탈이면, 모든 $k \geq 0$에 대해 사상
$$
S' \otimes_S \left((S \otimes_R S)/J^{k + 1}\right)
\longrightarrow
(S' \otimes_R S')/(J')^{k + 1}
$$
은 동형사상이다. 특히 사상
$S' \otimes_S \Omega_{S/R} \to \Omega_{S'/R}$는 동형사상이다.
\end{lemma}

\begin{proof}
$S' \otimes_S (S \otimes_R S) = S' \otimes_R S$임을 관찰하자.
또한 $J$가 $S' \otimes_R S$에서 생성하는 아이디얼은
전사 $S' \otimes_R S \to S'$의 핵 $I$이다.
따라서 좌변은 $S' \otimes_R S/I^{k + 1}$와 같다.
환 사상 $S' \otimes_R S \to S' \otimes_R S'$는
보조정리 \ref{algebra-lemma-base-change-formally-etale}에 의해
형식적으로 에탈이다. 그러므로 보조정리
\ref{algebra-lemma-formally-etale-lift-infinitesimal}에 의해
보조정리의 진술에 표시된 화살표가 동형사상임을 얻는다.
마지막 주장은 이 결과와
$\Omega_{S/R} = J/J^2$ 및
$\Omega_{S'/R} = J'/(J')^2$라는 사실로부터 따른다.
보조정리 \ref{algebra-lemma-differentials-diagonal}을 보라.
\end{proof}

\begin{lemma}
\label{algebra-lemma-formally-etale-principal-parts}
$R \to S \to S'$를 환 사상들이라 하고,
$S \to S'$가 형식적으로 에탈이라고 하자(예를 들어 에탈이라고 하자).
$M$을 $S$-가군이라 하고 $M' = S' \otimes_R M$으로 두자.
% P01-E0939: the frozen statement tensors M over R rather than over S.
그러면
$$
S' \otimes_S P^k_{S/R}(M) = P^k_{S'/R}(M')
$$
이다.
% P01-E0940: the frozen display lacks terminal punctuation.
따라서 임의의 $S$-가군 $N$과 임의의 유한 차수 미분 연산자
$D : M \to N$에 대해, $D$를 같은 차수 또는 더 작은 차수의
미분 연산자로 확대하는 유일한 사상
$D' : M' \to S' \otimes_S N$가 존재한다.
\end{lemma}

\begin{proof}
$J$와 $J'$를 보조정리
\ref{algebra-lemma-formally-etale-omega}의 진술에서와 같이 두자. 그러면
\begin{align*}
S' \otimes_S P^k_{S/R}(M)
& =
S' \otimes_S \left((S \otimes_R M)/J^{k + 1}(S \otimes_R M)\right) \\
& =
S' \otimes_S \left((S \otimes_R S)/J^{k + 1}\right) \otimes_S M \\
& =
\left(S' \otimes_R S'/(J')^{k + 1}\right) \otimes_S M \\
& =
\left(S' \otimes_R S'/(J')^{k + 1}\right) \otimes_{S'} M' \\
& =
S' \otimes_R M'/(J')^{k + 1}(S' \otimes_R M') \\
& =
P^k_{S'/R'}(M')
\end{align*}
% P01-E0941: the frozen last term uses the undefined base ring R'.
첫 번째와 마지막 등식은 보조정리
\ref{algebra-lemma-principal-parts-diagonal}에서 따른다.
세 번째 등식은 보조정리 \ref{algebra-lemma-formally-etale-omega}이다.
마지막 주장은 다음에서 따른다. $D$가 선형 사상
$\gamma : P^k_{S/R}(M) \to N$에 대응한다면,
$D' : M' \to S' \otimes_R N$가 선형 사상
$1 \otimes \gamma : S' \otimes_R P^k_{S/R}(M) \to S' \otimes_R N$에
대응하게 할 수 있다.
% P01-E0942: the frozen concluding tensor products use R where the statement requires S.
\end{proof}

\begin{remark}
\label{algebra-remark-formally-etale-differential-operators}
$R \to S \to S'$를 환 사상들이라 하고,
$S \to S'$가 형식적으로 에탈이라고 하자(예를 들어 에탈이라고 하자).
$M_i$, $i = 1, 2, 3$을 $S$-가군들이라 하고,
$D_i : M_i \to M_{i + 1}$, $i = 1, 2$를 유한 차수 미분 연산자들이라 하자.
$D'_i : M'_i \to M'_{i + 1}$, $i = 1, 2$가 보조정리
\ref{algebra-lemma-formally-etale-principal-parts}에서와 같이
$D_i$를 $M'_i = S' \otimes_S M_i$로 확대한 것이라면,
$D'_2 \circ D'_1$은 $D_2 \circ D_1$의 확대이다.
특히 $M$이 $S$-가군이면 $M' = S' \otimes_S M$은
$S$-대수 $\text{Diff}_{S/R}(M, M)$ 위의 가군이다.
\end{remark}





\section{비분기 환 사상}
\label{algebra-section-unramified}

\noindent
비분기 환 사상에 대한 우리의 정의는 \cite{Henselian}의 정의이다.
우리가 G-비분기 환 사상이라 부르는 것은 EGA에서 비분기 사상이라
부르는 것이다.
% P01-E0943: the frozen source has the incorrect article “an G-unramified ring map”.

\begin{definition}
\label{algebra-definition-unramified}
$R \to S$를 환 사상이라 하자.
\begin{enumerate}
\item $R \to S$가 유한형이고 $\Omega_{S/R} = 0$이면
$R \to S$를 {\it 비분기}라고 한다.
\item $R \to S$가 유한 표시이고 $\Omega_{S/R} = 0$이면
$R \to S$를 {\it G-비분기}라고 한다.
\item $S$의 소 아이디얼 $\mathfrak q$가 주어졌다고 하자.
$g \in S$, $g \not \in \mathfrak q$이고 $R \to S_g$가 비분기인
그런 원소가 존재하면 $S$가 {\it $\mathfrak q$에서 비분기}라고 한다.
\item $S$의 소 아이디얼 $\mathfrak q$가 주어졌다고 하자.
$g \in S$, $g \not \in \mathfrak q$이고 $R \to S_g$가 G-비분기인
그런 원소가 존재하면 $S$가 {\it $\mathfrak q$에서 G-비분기}라고 한다.
\end{enumerate}
\end{definition}

\noindent
물론 G-비분기 사상은 비분기이다.

\begin{lemma}
\label{algebra-lemma-formally-unramified-unramified}
$R \to S$를 환 사상이라 하자. 다음은 서로 동치이다.
\begin{enumerate}
\item $R \to S$는 형식적으로 비분기이고 유한형이다.
\item $R \to S$는 비분기이다.
\end{enumerate}
또한 다음도 서로 동치이다.
\begin{enumerate}
\item $R \to S$는 형식적으로 비분기이고 유한 표시이다.
\item $R \to S$는 G-비분기이다.
\end{enumerate}
\end{lemma}

\begin{proof}
보조정리 \ref{algebra-lemma-characterize-formally-unramified}와 정의들에서 따른다.
\end{proof}

\begin{lemma}
\label{algebra-lemma-unramified}
비분기 및 G-비분기 환 사상은 다음 성질들을 갖는다.
\begin{enumerate}
\item 비분기 환 사상의 기저변환은 비분기이다.
G-비분기 환 사상의 기저변환은 G-비분기이다.
\item 비분기 환 사상들의 합성은 비분기이다.
G-비분기 환 사상들의 합성은 G-비분기이다.
\item 모든 주 국소화 $R \to R_f$는 G-비분기이고 비분기이다.
\item $I \subset R$가 아이디얼이면 $R \to R/I$는 비분기이다.
$I \subset R$가 유한 생성 아이디얼이면 $R \to R/I$는 G-비분기이다.
\item 에탈 환 사상은 G-비분기이고 비분기이다.
\item $R \to S$가 유한형(각각 유한 표시)이고,
$\mathfrak q \subset S$가 소 아이디얼이며
$(\Omega_{S/R})_{\mathfrak q} = 0$이면,
$R \to S$는 $\mathfrak q$에서 비분기(각각 G-비분기)이다.
\item $R \to S$가 유한형(각각 유한 표시)이고,
$\mathfrak q \subset S$가 소 아이디얼이며
$\Omega_{S/R} \otimes_S \kappa(\mathfrak q) = 0$이면,
$R \to S$는 $\mathfrak q$에서 비분기(각각 G-비분기)이다.
\item $R \to S$가 유한형(각각 유한 표시)이고,
$\mathfrak p \subset R$ 위에 놓인 $\mathfrak q \subset S$가 소 아이디얼이며
$(\Omega_{S \otimes_R \kappa(\mathfrak p)/\kappa(\mathfrak p)})_{\mathfrak q}
= 0$이면, $R \to S$는 $\mathfrak q$에서 비분기
(각각 G-비분기)이다.
\item $R \to S$가 유한형(각각 표시)이고,
$\mathfrak p \subset R$ 위에 놓인 $\mathfrak q \subset S$가 소 아이디얼이며
$$
(\Omega_{S \otimes_R \kappa(\mathfrak p)/\kappa(\mathfrak p)})
\otimes_{S \otimes_R \kappa(\mathfrak p)} \kappa(\mathfrak q) = 0
$$
이면,
$R \to S$는 $\mathfrak q$에서 비분기(각각 G-비분기)이다.
% P01-E0944: the frozen source says “resp. presentation”, omitting “finite”.
\item $R \to S$가 환 사상이고, $g_1, \ldots, g_m \in S$가
단위 아이디얼을 생성하며, $j = 1, \ldots, m$에 대해
$R \to S_{g_j}$가 비분기(각각 G-비분기)이면,
$R \to S$는 비분기(각각 G-비분기)이다.
\item 환 사상 $R \to S$가 $S$의 모든 소 아이디얼에서 비분기
(각각 G-비분기)이면, $R \to S$는 비분기(각각 G-비분기)이다.
\item $R \to S$가 G-비분기이면, 유한형 $\mathbf{Z}$-대수 $R_0$,
G-비분기 환 사상 $R_0 \to S_0$, 그리고 $S = R \otimes_{R_0} S_0$가
되는 환 사상 $R_0 \to R$가 존재한다.
\item $R \to S$가 비분기이면, 유한형 $\mathbf{Z}$-대수 $R_0$,
비분기 환 사상 $R_0 \to S_0$, 그리고 $S$가
$R \otimes_{R_0} S_0$의 몫이 되는 환 사상 $R_0 \to R$가 존재한다.
\end{enumerate}
\end{lemma}

\begin{proof}
각 항을 순서대로 증명한다.

\medskip\noindent
(1)의 증명. 보조정리 \ref{algebra-lemma-differentials-base-change}와
\ref{algebra-lemma-base-change-finiteness}에서 따른다.

\medskip\noindent
(2)의 증명. 보조정리 \ref{algebra-lemma-exact-sequence-differentials}와
\ref{algebra-lemma-base-change-finiteness}에서 따른다.
% P01-E0945: the frozen proof of composition cites the base-change finiteness lemma; finite-type/finite-presentation composition is lemma-compose-finite-type.

\medskip\noindent
(3)의 증명. $\Omega_{R_f/R}$를 직접 계산하면 되며, 계산은 생략한다.

\medskip\noindent
(4)의 증명. 보조정리 \ref{algebra-lemma-trivial-differential-surjective}에 의해
$\Omega_{(R/I)/R} = 0$이고, 환 사상 $R \to R/I$는 유한형이다.
$I$가 유한 생성 아이디얼이면 $R \to R/I$는 유한 표시이다.

\medskip\noindent
(5)의 증명. 정의 \ref{algebra-definition-etale} 뒤의 논의를 보라.

\medskip\noindent
(6)의 증명. 이 경우 $\Omega_{S/R}$는 유한 $S$-가군이다
(보조정리 \ref{algebra-lemma-differentials-finitely-generated}을 보라).
따라서 $g \in S$, $g \not \in \mathfrak q$이고
$(\Omega_{S/R})_g = 0$인 원소가 존재한다.
보조정리 \ref{algebra-lemma-differentials-localize}에 의해 이는
$\Omega_{S_g/R} = 0$임을 뜻하므로, 원하는 대로 $R \to S_g$는 비분기이다.

\medskip\noindent
(7)의 증명. 나카야마 보조정리(보조정리 \ref{algebra-lemma-NAK})를 사용하면
이 조건이 (6)의 조건과 동치임을 알 수 있다.

\medskip\noindent
(8)과 (9)의 증명. 이 둘은 (6)과 (7)이 동치인 것과 같은 방식으로
서로 동치이다. 또한 보조정리 \ref{algebra-lemma-differentials-base-change}에 의해
$\Omega_{S \otimes_R \kappa(\mathfrak p)/\kappa(\mathfrak p)} =
\Omega_{S/R} \otimes_S (S \otimes_R \kappa(\mathfrak p))$이다.
따라서 양쪽의 $\kappa(\mathfrak q)$-벡터공간들이 표준적으로
동형이므로 (9)는 (7)과 동치이다.

\medskip\noindent
(10)의 증명. 보조정리 \ref{algebra-lemma-cover}와
\ref{algebra-lemma-differentials-localize}에서 따른다.

\medskip\noindent
(11)의 증명. (6), (7), 그리고 $S$의 스펙트럼이 준콤팩트라는
사실에서 따른다.
% P01-E0946: the frozen proof cites (6) and (7), which assume global finite type/presentation; the local-at-every-prime hypothesis plus quasi-compactness instead yields a finite principal cover and invokes (10).

\medskip\noindent
(12)의 증명. $S = R[x_1, \ldots, x_n]/(g_1, \ldots, g_m)$로 쓰자.
$\Omega_{S/R} = 0$이므로 어떤
$h_{ij}, a_{ijk} \in R[x_1, \ldots, x_n]$에 대해
$\Omega_{R[x_1, \ldots, x_n]/R}$에서
$$
\text{d}x_i = \sum h_{ij}\text{d}g_j + \sum a_{ijk}g_j\text{d}x_k
$$
로 쓸 수 있다. 다항식 $g_i, h_{ij}, a_{ijk}$의 모든 계수를 포함하는
유한 생성 $\mathbf{Z}$-부분대수 $R_0 \subset R$를 택하자.
$S_0 = R_0[x_1, \ldots, x_n]/(g_1, \ldots, g_m)$로 두면 된다.

\medskip\noindent
(13)의 증명. $S = R[x_1, \ldots, x_n]/I$로 쓰자.
$\Omega_{S/R} = 0$이므로 어떤
$h_{ij} \in R[x_1, \ldots, x_n]$와 $g_{ij}, g'_{ik} \in I$에 대해
$\Omega_{R[x_1, \ldots, x_n]/R}$에서
$$
\text{d}x_i = \sum h_{ij}\text{d}g_{ij} + \sum g'_{ik}\text{d}x_k
$$
로 쓸 수 있다. 다항식 $g_{ij}, h_{ij}, g'_{ik}$의 모든 계수를 포함하는
유한 생성 $\mathbf{Z}$-부분대수 $R_0 \subset R$를 택하자.
$S_0 = R_0[x_1, \ldots, x_n]/(g_{ij}, g'_{ik})$로 두면 된다.
\end{proof}

\begin{lemma}
\label{algebra-lemma-diagonal-unramified}
$R \to S$를 환 사상이라 하자. $R \to S$가 비분기이면,
$S \otimes_R S \to S$가 $S \otimes_R S \to (S \otimes_R S)_e$와
동형이 되게 하는 멱등원 $e \in S \otimes_R S$가 존재한다.
\end{lemma}

\begin{proof}
$J = \Ker(S \otimes_R S \to S)$로 두자. 가정에 의해
$J/J^2 = 0$이다. 보조정리 \ref{algebra-lemma-differentials-diagonal}을 보라.
$S$가 $R$ 위에서 유한형이므로 $J$는 유한 생성이다.
구체적으로, $x_i$들이 $R$ 위에서 $S$를 생성할 때
$x_i \otimes 1 - 1 \otimes x_i$들이 그 핵을 생성한다.
보조정리 \ref{algebra-lemma-ideal-is-squared-union-connected}에서 결론이 따른다.
\end{proof}

\begin{lemma}
\label{algebra-lemma-unramified-at-prime}
$R \to S$를 환 사상이라 하자. $\mathfrak q \subset S$를
$R$의 $\mathfrak p$ 위에 놓인 소 아이디얼이라 하자.
$S/R$가 $\mathfrak q$에서 비분기이면 다음이 성립한다.
\begin{enumerate}
\item $\mathfrak p S_{\mathfrak q} = \mathfrak qS_{\mathfrak q}$이고,
이는 국소환 $S_{\mathfrak q}$의 극대 아이디얼이다.
% P01-E0947: the frozen source reads “we have pS_q = qS_q is the maximal ideal”, joining two predicates ungrammatically.
\item 체확대 $\kappa(\mathfrak q)/\kappa(\mathfrak p)$는 유한 분리이다.
\end{enumerate}
\end{lemma}

\begin{proof}
먼저 어떤 $g \in S$, $g \not \in \mathfrak q$에 대해 $S$를 $S_g$로
바꾸어 $R \to S$가 비분기라고 가정할 수 있다.
보조정리 \ref{algebra-lemma-unramified}에 의해 기저변환
$S \otimes_R \kappa(\mathfrak p)$는 $\kappa(\mathfrak p)$ 위에서 비분기이다.
보조정리 \ref{algebra-lemma-characterize-smooth-over-field}에 의해 이는 매끄럽고,
따라서 $\kappa(\mathfrak p)$ 위에서 에탈이다.
그러므로 보조정리 \ref{algebra-lemma-etale-over-field}에 의해
$F = S \otimes_R \kappa(\mathfrak p)$는 $\kappa(\mathfrak p)$의
유한 분리 체확대들의 유한 곱이다. 주석
\ref{algebra-remark-local-ring-fibre}의 표기와 결과를 사용하면
$F_{\overline{\mathfrak q}} = S_\mathfrak q/\mathfrak pS_\mathfrak q$가
$\kappa(\mathfrak q)$와 같음을 얻는다. 이로부터 보조정리가 따른다.
\end{proof}

\begin{lemma}
\label{algebra-lemma-unramified-quasi-finite}
$R \to S$를 유한형 환 사상이라 하고, $\mathfrak q$를 $S$의 소 아이디얼이라 하자.
$R \to S$가 $\mathfrak q$에서 비분기이면,
$R \to S$는 $\mathfrak q$에서 준유한이다.
특히 비분기 환 사상은 준유한이다.
\end{lemma}

\begin{proof}
비분기 환 사상은 유한형이다. 따라서 두 번째 주장이 첫 번째에서
따르는 것은 명백하다. 첫 번째 주장을 보이려면 보조정리
\ref{algebra-lemma-isolated-point-fibre}의 특징짓기 (2)를 보조정리
\ref{algebra-lemma-unramified-at-prime}을 사용하여 적용한다.
\end{proof}

\begin{lemma}
\label{algebra-lemma-characterize-unramified}
$R \to S$를 환 사상이라 하자. $\mathfrak q$를 $R$의 소 아이디얼
$\mathfrak p$ 위에 놓인 $S$의 소 아이디얼이라 하자. 다음을 가정하자.
\begin{enumerate}
\item $R \to S$는 유한형이다.
\item $\mathfrak p S_{\mathfrak q}$는 국소환 $S_{\mathfrak q}$의
극대 아이디얼이다.
\item 체확대 $\kappa(\mathfrak q)/\kappa(\mathfrak p)$는 유한 분리이다.
\end{enumerate}
그러면 $R \to S$는 $\mathfrak q$에서 비분기이다.
\end{lemma}

\begin{proof}
보조정리 \ref{algebra-lemma-unramified} (8)에 의해
$\Omega_{S \otimes_R \kappa(\mathfrak p) / \kappa(\mathfrak p)}$를
$\mathfrak q$에서 국소화한 것이 영임을 보이면 충분하다.
따라서 $S$를 $S \otimes_R \kappa(\mathfrak p)$로,
$R$을 $\kappa(\mathfrak p)$로 바꿀 수 있다. 다시 말해 $R = k$는
체이고 $S$는 유한형 $k$-대수라고 가정할 수 있다.
이 경우 가정들로부터 $S_{\mathfrak q} \cong \kappa(\mathfrak q)$이다.
따라서 원하는 대로
$(\Omega_{S/k})_{\mathfrak q} = \Omega_{S_\mathfrak q/k} =
\Omega_{\kappa(\mathfrak q)/k}$는 영이다
(첫 번째 등식은 보조정리 \ref{algebra-lemma-differentials-localize}이다).
\end{proof}

\begin{lemma}
\label{algebra-lemma-etale-flat-unramified-finite-presentation}
$R \to S$를 환 사상이라 하자. 다음은 서로 동치이다.
\begin{enumerate}
\item $R \to S$는 에탈이다.
\item $R \to S$는 평탄하고 G-비분기이다.
\item $R \to S$는 평탄하고 비분기이며 유한 표시이다.
\end{enumerate}
\end{lemma}

\begin{proof}
(2)와 (3)은 정의에 의해 동치이다.
(1) $\Rightarrow$ (3)은 에탈 환 사상이 유한 표시라는 사실,
보조정리 \ref{algebra-lemma-etale}(에탈 사상의 평탄성), 그리고 보조정리
\ref{algebra-lemma-unramified}(에탈 사상은 비분기임)에서 따른다.
역으로, 보조정리 \ref{algebra-lemma-characterize-etale}의 에탈 환 사상의
특징짓기와 보조정리 \ref{algebra-lemma-unramified-at-prime}의 비분기 환 사상의
구조로부터 (3)이 (1)을 함의함을 알 수 있다.
% P01-E0948: the frozen compound subject “the characterization ... and the structure ... shows” has singular agreement.
(여기서는 $R \to S$가 모든 소 아이디얼 $\mathfrak q \subset S$에서
에탈이면 $R \to S$가 에탈이라는 사실을 사용한다.
보조정리 \ref{algebra-lemma-etale}을 보라.)
\end{proof}

\begin{lemma}
\label{algebra-lemma-characterize-etale-over-polynomial-ring}
$k$를 체라 하고
$$
\varphi : k[x_1, \ldots, x_n] \to A, \quad x_i \longmapsto a_i
$$
를 유한형 환 사상이라 하자. 그러면 $\varphi$가 에탈일 필요충분조건은
다음 두 조건이 성립하는 것이다. (a) 극대 아이디얼에서 $A$의 국소환들은
차원 $n$이고, (b) 원소 $\text{d}(a_1), \ldots, \text{d}(a_n)$은
$A$-가군 $\Omega_{A/k}$를 생성한다.
\end{lemma}

\begin{proof}
(a)와 (b)를 가정하자. 조건 (b)로부터
$\Omega_{A/k[x_1, \ldots, x_n]} = 0$이므로 $\varphi$는 비분기이다.
따라서 보조정리 \ref{algebra-lemma-etale-flat-unramified-finite-presentation}에 의해
$\varphi$가 평탄임을 보이면 충분하다.
$\mathfrak m \subset A$를 극대 아이디얼이라 하자.
$X = \Spec(A)$로 두고 $\mathfrak m$에 대응하는 닫힌점을 $x \in X$로 나타내자.
보조정리 \ref{algebra-lemma-dimension-closed-point-finite-type-field}에 의해
$\dim(A_\mathfrak m)$은 $\dim_x X$이다. 따라서 (a)와 (b)가 성립하면
보조정리 \ref{algebra-lemma-characterize-smooth-over-field}에 의해 모든 극대
아이디얼 $\mathfrak m$에 대해 $A_\mathfrak m$은 정칙 국소환이다.
그러면 보조정리 \ref{algebra-lemma-CM-over-regular-flat}에 의해
$k[x_1, \ldots, x_n]_{\varphi^{-1}(\mathfrak m)} \to A_\mathfrak m$은
평탄이다(정칙 국소환이 CM이라는 사실은 보조정리
\ref{algebra-lemma-regular-ring-CM}을 보라). 그러므로 보조정리
\ref{algebra-lemma-flat-localization}에 의해 $\varphi$는 평탄이다.

\medskip\noindent
$\varphi$가 에탈이라고 가정하자. 그러면
$\Omega_{A/k[x_1, \ldots, x_n]} = 0$이므로 (b)가 성립한다.
한편 에탈 환 사상은 평탄하고(보조정리 \ref{algebra-lemma-etale})
준유한이다(보조정리 \ref{algebra-lemma-etale-quasi-finite}).
따라서 $A$의 모든 극대 아이디얼 $\mathfrak m$에 대해
$k[x_1, \ldots, x_n]_{\varphi^{-1}(\mathfrak m)} \to A_\mathfrak m$에
보조정리 \ref{algebra-lemma-dimension-base-fibre-equals-total}을 적용할 수 있어서
$\dim(A_\mathfrak m) = n$을 얻고, 따라서 (a)가 성립한다.
% P01-E0949: the frozen source has the typo “we my apply” instead of “we may apply”.
\end{proof}





\section{비분기 환 사상의 국소 구조}
\label{algebra-section-local-structure-unramified}

\noindent
비분기 사상은 국소적으로(적절한 의미에서) 닫힌 몰입과 에탈 사상의
합성이다. 이 사실의 대수적 토대를 이 절에서 논의한다.

\begin{proposition}
\label{algebra-proposition-unramified-locally-standard}
$R \to S$를 환 사상이라 하고 $\mathfrak q \subset S$를 소 아이디얼이라 하자.
$R \to S$가 $\mathfrak q$에서 비분기이면 다음이 존재한다.
\begin{enumerate}
\item $g \in S$, $g \not \in \mathfrak q$,
\item 표준 에탈 환 사상 $R \to S'$,
\item 전사 $R$-대수 사상 $S' \to S_g$.
\end{enumerate}
\end{proposition}

\begin{proof}
이 증명은 명제 \ref{algebra-proposition-etale-locally-standard}의 증명과
``같다''. 증명이 다소 우회적이며 더 짧게 할 방법이 있을 수도 있다.

\medskip\noindent
1단계. 정의 \ref{algebra-definition-unramified}에 의해
$g \in S$, $g \not \in \mathfrak q$이고 $R \to S_g$가 비분기인
원소가 존재한다. 따라서 $S$가 $R$ 위에서 비분기라고 가정할 수 있다.

\medskip\noindent
2단계. 보조정리 \ref{algebra-lemma-unramified}에 의해,
$\mathbf{Z}$ 위에서 유한형인 $R_0$와 비분기 환 사상
$R_0 \to S_0$, 그리고 $S$가 $R \otimes_{R_0} S_0$의 몫이 되게 하는
환 사상 $R_0 \to R$가 존재한다. $\mathfrak q$에 대응하는 $S_0$의
소 아이디얼을 $\mathfrak q_0$라 하자.
$(R_0 \to S_0, \mathfrak q_0)$에 대해 결과를 보이면 기저변환에 의해
$(R \to S, \mathfrak q)$에 대한 결과가 따른다.
따라서 $R$이 뇌터 환이라고 가정할 수 있다.

\medskip\noindent
3단계. 보조정리 \ref{algebra-lemma-unramified-quasi-finite}에 의해
$R \to S$는 준유한임을 주목하자.
보조정리 \ref{algebra-lemma-quasi-finite-open-integral-closure}에 의해
유한 환 사상 $R \to S'$, $R$-대수 사상 $S' \to S$, 그리고
$S' \to S$가 동형사상 $S'_{g'} \cong S_{g'}$를 유도하게 하는
$g' \in S'$, $g' \not \in \mathfrak q$가 존재한다.
% P01-E0950: the frozen source compares g' in S' directly with q in S and repeats “such that”; it should refer to the image of g' (or the contracted prime) once.
($S'$가 $R$ 위에서 비분기일 필요는 없음을 주의하라.)
따라서 (a) $R$은 뇌터 환이고, (b) $R \to S$는 유한이며,
(c) $R \to S$는 $\mathfrak q$에서 비분기라고 가정할 수 있다
(그러나 이제 모든 소 아이디얼에서 비분기일 필요는 없다).

\medskip\noindent
4단계. $\mathfrak p \subset R$를 $\mathfrak q$에 대응하는 소 아이디얼이라
하자. 섬유환 $S \otimes_R \kappa(\mathfrak p)$를 생각하자.
이는 $\kappa(\mathfrak p)$ 위의 유한 대수이다. 따라서 이는 아르틴 환이고
(보조정리 \ref{algebra-lemma-finite-dimensional-algebra}을 보라), 국소환들의
유한 곱
$$
S \otimes_R \kappa(\mathfrak p) = \prod\nolimits_{i = 1}^n A_i
$$
이다. 명제 \ref{algebra-proposition-dimension-zero-ring}을 보라.
그 인자 중 하나, 이를테면 $A_1$은 국소환
$S_{\mathfrak q}/\mathfrak pS_{\mathfrak q}$이고,
보조정리 \ref{algebra-lemma-unramified-at-prime}에 의해 이는
$\kappa(\mathfrak q)$와 동형이다. 나머지 인자들은 $\mathfrak p$ 위에 놓인
$S$의 다른 소 아이디얼들, 이를테면
$\mathfrak q_2, \ldots, \mathfrak q_n$에 대응한다.

\medskip\noindent
5단계. 유한 분리 체확대 $\kappa(\mathfrak q)/\kappa(\mathfrak p)$를
생성하는 영이 아닌 원소 $\alpha \in \kappa(\mathfrak q)$를 택할 수 있다
(따라서 체확대가 자명하더라도 $\alpha = 0$은 허용하지 않는다).
임의의 $\lambda \in \kappa(\mathfrak p)^*$에 대해 원소
$\lambda \alpha$도 $\kappa(\mathfrak p)$ 위에서
$\kappa(\mathfrak q)$를 생성함을 주의하자. 원소
$$
\overline{t} =
(\alpha, 0, \ldots, 0) \in
\prod\nolimits_{i = 1}^n A_i =
S \otimes_R \kappa(\mathfrak p).
$$
를 생각하자. 필요하면 위와 같이 $\alpha$를 $\lambda \alpha$로 바꾸어,
$\overline{t}$가 어떤 $t \in S$의 상이라고 가정할 수 있다.
$x$를 $t$로 보내는 $R$-대수 사상 $R[x] \to S$의 핵을
$I \subset R[x]$라 하자. $S' = R[x]/I$로 두면 $S' \subset S$이다.
다음 도식이 있다.
$$
\xymatrix{
R[x] \ar[r] & S' \ar[r] & S \\
R \ar[u] \ar[ru] \ar[rru] & &
}
$$
구성에 의해 $S$의 소 아이디얼 $\mathfrak q_j$, $j \geq 2$는 모두
$R[x]$의 소 아이디얼 $(\mathfrak p, x)$ 위에 놓이는 반면,
$\alpha \not = 0$이므로 $\mathfrak q$는 $R[x]$의 다른 소 아이디얼 위에 놓인다.

\medskip\noindent
6단계. $\mathfrak q$에 대응하는 $S'$의 소 아이디얼을
$\mathfrak q' \subset S'$라 하자. 위 논의에 의해 $\mathfrak q$는
$\mathfrak q'$ 위에 놓인 $S$의 유일한 소 아이디얼이다.
따라서 보조정리 \ref{algebra-lemma-unique-prime-over-localize-below}에 의해
$S_{\mathfrak q} = S_{\mathfrak q'}$이다
($S' \to S$에 대해 보조정리 \ref{algebra-lemma-integral-going-up}에 의해
상승 성질이 성립한다. 이는 $R \to S$가 유한이므로
$S' \to S$도 유한이기 때문이다).
그러므로 $S'_{\mathfrak q'} \to S_{\mathfrak q}$는 유한 단사 환 사상
$S' \to S$의 국소화이므로 유한이고 단사이다. 국소환 사상들
$$
R_{\mathfrak p} \to S'_{\mathfrak q'} \to S_{\mathfrak q}
$$
을 생각하자. 두 번째 사상은 유한이고 단사이다.
보조정리 \ref{algebra-lemma-unramified-at-prime}에 의해
$S_{\mathfrak q}/\mathfrak pS_{\mathfrak q} = \kappa(\mathfrak q)$이다.
따라서 더욱이
$S_{\mathfrak q}/\mathfrak q'S_{\mathfrak q} = \kappa(\mathfrak q)$이다.
그리고
$$
\kappa(\mathfrak p) \subset \kappa(\mathfrak q') \subset \kappa(\mathfrak q)
$$
이며 $\alpha$가 $\kappa(\mathfrak q')$에서 $\kappa(\mathfrak q)$로 가는
사상의 상에 속하므로 $\kappa(\mathfrak q') = \kappa(\mathfrak q)$이다.
이제 $S'_{\mathfrak q'}$-가군 사상
$S'_{\mathfrak q'} \to S_{\mathfrak q}$에 나카야마 보조정리
\ref{algebra-lemma-NAK}를 적용하면
$S'_{\mathfrak q'} \to S_{\mathfrak q}$는 전사이다.
다시 말해 $S'_{\mathfrak q'} \cong S_{\mathfrak q}$이다.

\medskip\noindent
7단계. 보조정리 \ref{algebra-lemma-isomorphic-local-rings}에 의해
$g \in S$, $g \not \in \mathfrak q$와
$g' \in S'$, $g' \not \in \mathfrak q'$ 중
$S'_{g'} \cong S_g$가 되게 하는 것들이 존재한다.
$R$이 뇌터 환이고 유한 $R$-가군 $S$의 $R$-부분가군인
$S'$는 $R$ 위에서 유한이다. 따라서 $S$를 $S'$로 바꾸어
(a) $R$은 뇌터 환이고, (b) $S$ 유한 $R$-대수이며,
(c) $S$는 $R$ 위에서 $\mathfrak q$에서 비분기이고,
(d) $S = R[x]/I$라고 가정할 수 있다.
% P01-E0951: the frozen source omits “is” in “(b) S finite over R”.

\medskip\noindent
8단계. 환
$S \otimes_R \kappa(\mathfrak p) = \kappa(\mathfrak p)[x]/\overline{I}$를
생각하자. 여기서 $\overline{I} = I \cdot \kappa(\mathfrak p)[x]$는
$\kappa(\mathfrak p)[x]$에서 $I$가 생성하는 아이디얼이다.
$\kappa(\mathfrak p)[x]$가 PID이므로 어떤 최고차항계수가 1인
$\overline{h} \in \kappa(\mathfrak p)$에 대해
$\overline{I} = (\overline{h})$임을 안다.
% P01-E0952: the frozen source places the monic polynomial h-bar in kappa(p) instead of kappa(p)[x].
어떤 $\lambda \in \kappa(\mathfrak p)$에 대해 $\overline{h}$를
$\lambda \cdot \overline{h}$로 바꾼 뒤, $\overline{h}$가 어떤
$h \in R[x]$의 상이라고 가정할 수 있다.
% P01-E0953: the scaling factor must be nonzero to preserve the principal generator, but the frozen source omits the star on kappa(p).
(문제는 $h$를 최고차항계수가 1이 되게 택할 수 있는지 알 수 없다는 점이다.)
또한 4단계와 마찬가지로
$S \otimes_R \kappa(\mathfrak p) = A_1 \times \ldots \times A_n$이고,
$A_1 = \kappa(\mathfrak q)$는 $\kappa(\mathfrak p)$의 유한 분리 확대이며
$A_2, \ldots, A_n$은 국소환임을 안다. 이는
$$
\overline{h} = \overline{h}_1 \overline{h}_2^{e_2} \ldots \overline{h}_n^{e_n}
$$
임을 뜻한다. 여기서 $\overline{h}_i \in \kappa(\mathfrak p)[x]$는
쌍마다 서로소인 기약 최고차항계수 1인 다항식들이고,
$e_2, \ldots, e_n \geq 1$이다.
% P01-E0954: after the frozen source replaces the monic generator by an arbitrary nonzero scalar multiple, monicity is no longer assured; this factorization into monic factors needs the resulting leading unit as a coefficient.
번호는 $\kappa(\mathfrak p)[x]$-대수로서
$A_i = \kappa(\mathfrak p)[x]/(\overline{h}_i^{e_i})$가 되게 택했다.
% P01-E0955: the frozen formula uses e_i also for i = 1, although only e_2,...,e_n are defined; it should set e_1 = 1 or state the i = 1 case separately.
$\overline{h}_1$은 $\alpha \in \kappa(\mathfrak q)$의 최소다항식이므로
분리다항식임을 주의하자(그 도함수는 자기 자신과 서로소이다).

\medskip\noindent
9단계. $m \in I$를 최고차항계수가 1인 원소라 하자.
그러한 원소는 환확대 $R \to R[x]/I$가 유한이고 따라서 정수적이므로 존재한다.
$\kappa(\mathfrak p)[x]$에서의 상을 $\overline{m}$이라 하자.
어떤 $d_1 \geq 1$, $d_j \geq e_j$, $j = 2, \ldots, n$과
모든 $\overline{h}_i$와 서로소인
$\overline{k} \in \kappa(\mathfrak p)[x]$에 대해
$$
\overline{m} = \overline{k}
\overline{h}_1^{d_1} \overline{h}_2^{d_2} \ldots \overline{h}_n^{d_n}
$$
로 인수분해할 수 있다. $l \deg(m) > \deg(h)$이고 $l \geq 2$인
수를 택하여 $f = m^l + h$로 두자. 그러면 $f$는 $R$ 위의
다항식으로서 최고차항계수가 1이다. 또한 $f$의
$\kappa(\mathfrak p)[x]$에서의 상 $\overline{f}$는 다음과 같이 인수분해된다.
$$
\overline{f} =
\overline{h}_1 \overline{h}_2^{e_2} \ldots \overline{h}_n^{e_n}
+
\overline{k}^l \overline{h}_1^{ld_1} \overline{h}_2^{ld_2}
\ldots \overline{h}_n^{ld_n}
=
\overline{h}_1(\overline{h}_2^{e_2} \ldots \overline{h}_n^{e_n}
+
\overline{k}^l
\overline{h}_1^{ld_1 - 1} \overline{h}_2^{ld_2} \ldots \overline{h}_n^{ld_n})
= \overline{h}_1 \overline{w}
$$
여기서 $\overline{w}$는 $\overline{h}_1$과 서로소인 다항식이다.
$g = f'$로 두자($x$에 관한 도함수이다).

\medskip\noindent
10단계. 환 사상 $R[x] \to S = R[x]/I$는 다음 성질을 갖는다.
(1) $f$를 영으로 보내고, (2) $g$를 $S \setminus \mathfrak q$의
원소로 보낸다. 첫 번째 주장은 $f$가 $I$의 원소이므로 명백하다.
% P01-E0956: the frozen source only says h lies in R[x], but Step 10 needs h in I to conclude f = m^l + h lies in I.
두 번째 주장을 위해서는 $g$가
$\kappa(\mathfrak q) = \kappa(\mathfrak p)[x]/(\overline{h}_1)$에서
영으로 가지 않음을 보이면 된다. $\kappa(\mathfrak p)[x]$에서
$g$의 상은 $\overline{f}$의 도함수이다. 따라서
$$
\overline{g} =
\frac{\text{d}\overline{f}}{\text{d}x} =
\overline{w}\frac{\text{d}\overline{h}_1}{\text{d}x} +
\overline{h}_1\frac{\text{d}\overline{w}}{\text{d}x},
$$
$\overline{w}$가 $\overline{h}_1$과 서로소이고
$\overline{h}_1$이 분리다항식이므로 (2)는 명백하다.

\medskip\noindent
11단계. $\varphi : R[x]/(f) \to S$는 전사 환 사상이고,
$R[x]_g/(f)$는 $R$ 위에서 에탈이며(표준 에탈이기 때문이다.
보조정리 \ref{algebra-lemma-standard-etale}을 보라),
$\varphi(g) \not \in \mathfrak q$라고 결론내린다.
따라서 사상 $(R[x]/(f))_g \to S_{\varphi(g)}$가 원하는 전사이다.
\end{proof}



\begin{lemma}
\label{algebra-lemma-etale-makes-unramified-closed-at-prime}
$R \to S$를 환 사상이라 하자. $\mathfrak q$를
$\mathfrak p \subset R$ 위에 놓인 $S$의 소 아이디얼이라 하자.
$R \to S$가 유한형이고 $\mathfrak q$에서 비분기라고 가정하자.
그러면 다음이 존재한다.
\begin{enumerate}
\item 에탈 환 사상 $R \to R'$,
\item $\mathfrak p$ 위에 놓인 소 아이디얼 $\mathfrak p' \subset R'$,
\item 곱 분해
$$
R' \otimes_R S = A \times B
$$
\end{enumerate}
이들은 다음 성질을 갖는다.
\begin{enumerate}
\item $R' \to A$는 전사이다.
\item $\mathfrak p'A$는 $\mathfrak p'$ 위와 $\mathfrak q$ 위에 놓인
$A$의 소 아이디얼이다.
\end{enumerate}
\end{lemma}

\begin{proof}
$(R \to S, \mathfrak p, \mathfrak q)$를 다음과 같은 임의의
기저변환으로 바꿀 수 있다. $R \to R'$는 에탈 환 사상이고,
$\mathfrak p'$는 $\mathfrak p$ 위에 놓인 소 아이디얼이며,
$\mathfrak q'$는 $\mathfrak q$와 $\mathfrak p'$ 양쪽 위에 놓이는
소 아이디얼로 택한다. 즉
$(R' \to R'\otimes_R S, \mathfrak p', \mathfrak q')$로 바꿀 수 있다.
또한 $R \to R'$와 $\mathfrak p'$가 주어지면 적합한
$\mathfrak q'$를 항상 찾을 수 있음을 주의하자.

\medskip\noindent
$R \to S$가 유한형이라는 가정으로 보조정리
\ref{algebra-lemma-etale-makes-quasi-finite-finite-variant}을 적용할 수 있다.
따라서 $S = A_1 \times \ldots \times A_n \times B$이고,
각 $R \to A_i$는 유한이며 $\mathfrak p$ 위에 놓이는 소 아이디얼
$\mathfrak r_i$를 정확히 하나 갖고,
$\kappa(\mathfrak p) \subset \kappa(\mathfrak r_i)$는 순수 비분리이며,
$R \to B$는 $\mathfrak p$ 위에 놓인 어떤 소 아이디얼에서도
준유한이 아니라고 가정할 수 있다. 비분기 사상은 준유한이므로
(보조정리 \ref{algebra-lemma-unramified-quasi-finite}을 보라), 명백히
$\mathfrak q = \mathfrak r_i$인 $i$가 있다. $\mathfrak q = \mathfrak r_1$이라 하자.
보조정리 \ref{algebra-lemma-unramified-at-prime}에 의해
$\kappa(\mathfrak r_1)/\kappa(\mathfrak p)$는 분리이므로 자명한
체확대이고, $\mathfrak p(A_1)_{\mathfrak r_1}$은 극대 아이디얼이다.
또한 보조정리 \ref{algebra-lemma-unique-prime-over-localize-below}에 의해
$(A_1)_{\mathfrak r_1} = (A_1)_{\mathfrak p}$이다
($R \to A_1$은 유한 환 사상이므로 보조정리
\ref{algebra-lemma-integral-going-up}에 의해 상승 성질을 만족하여 이 보조정리가 적용된다).
나카야마 보조정리 \ref{algebra-lemma-NAK}에 의해 국소환 사상
$R_{\mathfrak p} \to (A_1)_{\mathfrak p} = (A_1)_{\mathfrak r_1}$은
전사이다. $A_1$은 $R$ 위에서 유한이므로,
$f \in R$, $f \not \in \mathfrak p$이고
$R_f \to (A_1)_f$가 전사가 되는 원소가 존재한다.
$R$을 $R_f$로 바꾸면 결론을 얻는다.
\end{proof}

\begin{lemma}
\label{algebra-lemma-etale-makes-unramified-closed}
\begin{slogan}
비분기 환 사상에서는 에탈 근방으로 옮겨 한 섬유의 점들을 분리할 수 있다.
\end{slogan}
$R \to S$를 환 사상이라 하고 $\mathfrak p$를 $R$의 소 아이디얼이라 하자.
$R \to S$가 비분기이면 다음이 존재한다.
\begin{enumerate}
\item 에탈 환 사상 $R \to R'$,
\item $\mathfrak p$ 위에 놓인 소 아이디얼 $\mathfrak p' \subset R'$,
\item 곱 분해
$$
R' \otimes_R S = A_1 \times \ldots \times A_n \times B
$$
\end{enumerate}
이들은 다음 성질을 갖는다.
\begin{enumerate}
\item $R' \to A_i$는 전사이다.
\item $\mathfrak p'A_i$는 $\mathfrak p'$ 위에 놓인 $A_i$의 소 아이디얼이다.
\item $\mathfrak p'$ 위에 놓인 $B$의 소 아이디얼은 없다.
\end{enumerate}
\end{lemma}

\begin{proof}
보조정리 \ref{algebra-lemma-etale-makes-quasi-finite-finite-variant}을 적용할 수 있다.
따라서 에탈 기저변환 뒤에
$S = A_1 \times \ldots \times A_n \times B$이고,
각 $R \to A_i$는 유한이며 $\mathfrak p$ 위에 놓이는 소 아이디얼
$\mathfrak r_i$를 정확히 하나 갖고,
$\kappa(\mathfrak p) \subset \kappa(\mathfrak r_i)$는 순수 비분리이며,
$R \to B$는 $\mathfrak p$ 위에 놓인 어떤 소 아이디얼에서도
준유한이 아니라고 가정할 수 있다.
$R \to S$는 준유한이므로(보조정리
\ref{algebra-lemma-unramified-quasi-finite}을 보라),
$\mathfrak p$ 위에 놓인 $B$의 소 아이디얼은 없다.
보조정리 \ref{algebra-lemma-unramified-at-prime}에 의해
$\kappa(\mathfrak r_i)/\kappa(\mathfrak p)$는 분리이므로 자명한
체확대이고, $\mathfrak p(A_i)_{\mathfrak r_i}$는 극대 아이디얼이다.
또한 보조정리 \ref{algebra-lemma-unique-prime-over-localize-below}에 의해
$(A_i)_{\mathfrak r_i} = (A_i)_{\mathfrak p}$이다
($R \to A_i$는 유한 환 사상이므로 보조정리
\ref{algebra-lemma-integral-going-up}에 의해 상승 성질을 만족하여 이 보조정리가 적용된다).
나카야마 보조정리 \ref{algebra-lemma-NAK}에 의해 국소환 사상
$R_{\mathfrak p} \to (A_i)_{\mathfrak p} = (A_i)_{\mathfrak r_i}$은
전사이다. $A_i$는 $R$ 위에서 유한이므로,
$f \in R$, $f \not \in \mathfrak p$이고
$R_f \to (A_i)_f$가 전사가 되는 원소가 존재한다.
$R$을 $R_f$로 바꾸면 결론을 얻는다.
\end{proof}





\section{헨젤 국소환}
\label{algebra-section-henselian}

\noindent
이 절에서는 헨젤 국소환의 개념을 조금 논의한다.
$(R, \mathfrak m, \kappa)$를 국소환이라 하자.
$a \in R$에 대해 $\kappa$에서 $a$의 상을 $\overline{a}$로 나타낸다.
다항식 $f \in R[T]$에 대해서는 흔히 $\kappa[T]$에서 $f$의 상을
$\overline{f}$로 나타낸다. 다항식 $f \in R[T]$가 주어지면
$T$에 관한 $f$의 도함수를 $f'$로 나타낸다.
$\overline{f}' = \overline{f'}$임을 주의하자.

\begin{definition}
\label{algebra-definition-henselian}
$(R, \mathfrak m, \kappa)$를 국소환이라 하자.
\begin{enumerate}
\item 모든 최고차항계수가 1인 $f \in R[T]$와
$\overline{f}$의 모든 근 $a_0 \in \kappa$ 중
$\overline{f'}(a_0) \not = 0$을 만족하는 것에 대해,
$f(a) = 0$이고 $a_0 = \overline{a}$인 $a \in R$가 존재하면
$R$을 {\it 헨젤}이라고 한다.
\item $R$이 헨젤이고 그 잉여체가 분리 대수적으로 닫혀 있으면
$R$을 {\it 엄밀 헨젤}이라고 한다.
\end{enumerate}
\end{definition}

\noindent
조건 $\overline{f'}(a_0) \not = 0$은 $a_0$가 다항식
$\overline{f}$의 단순근이라는 조건과 동치임을 주의하자.
실제로 이 조건은 올림 $a \in R$가 존재한다면 유일함을 뜻한다.

\begin{lemma}
\label{algebra-lemma-uniqueness}
$(R, \mathfrak m, \kappa)$를 국소환이라 하자.
$f \in R[T]$이고 $a, b \in R$가
$f(a) = f(b) = 0$, $a = b \bmod \mathfrak m$,
$f'(a) \not \in \mathfrak m$을 만족한다고 하자. 그러면 $a = b$이다.
\end{lemma}

\begin{proof}
$R[x, y]$에서
$f(x + y) - f(x) = f'(x)y + g(x, y) y^2$로 쓰자
($f(x + y)$를 전개하면 가능하며, 세부사항은 생략한다).
그러면 어떤 $c \in R$에 대해
$0 = f(b) - f(a) = f(a + (b - a)) - f(a) =
f'(a)(b - a) + c (b - a)^2$이다. 가정에 의해
$f'(a)$는 $R$의 단위원이다. 따라서
$(b - a)(1 + f'(a)^{-1}c(b - a)) = 0$이다.
가정에 의해 $b - a \in \mathfrak m$이므로
$1 + f'(a)^{-1}c(b - a)$는 $R$의 단위원이다.
따라서 $R$에서 $b - a = 0$이다.
\end{proof}

\noindent
다음은 헨젤 국소환의 특징짓기이다.

\begin{lemma}
\label{algebra-lemma-characterize-henselian}
\begin{slogan}
헨젤 국소환의 특징짓기
\end{slogan}
$(R, \mathfrak m, \kappa)$를 국소환이라 하자. 다음은 서로 동치이다.
\begin{enumerate}
\item $R$은 헨젤이다.
\item 모든 $f \in R[T]$와 $\overline{f}$의 모든 근
$a_0 \in \kappa$ 중 $\overline{f'}(a_0) \not = 0$을 만족하는 것에 대해,
$f(a) = 0$이고 $a_0 = \overline{a}$인 $a \in R$가 존재한다.
\item 최고차항계수가 1인 임의의 $f \in R[T]$와
$\gcd(g_0, h_0) = 1$인 임의의 인수분해
$\overline{f} = g_0 h_0$에 대해, $g_0 = \overline{g}$이고
$h_0 = \overline{h}$인 인수분해 $f = gh$가 $R[T]$에 존재한다.
\item 최고차항계수가 1인 임의의 $f \in R[T]$와
$\gcd(g_0, h_0) = 1$인 임의의 인수분해
$\overline{f} = g_0 h_0$에 대해, $g_0 = \overline{g}$,
$h_0 = \overline{h}$이고 또한 $\deg_T(g) = \deg_T(g_0)$인
인수분해 $f = gh$가 $R[T]$에 존재한다.
\item 임의의 $f \in R[T]$와 $\gcd(g_0, h_0) = 1$인 임의의 인수분해
$\overline{f} = g_0 h_0$에 대해, $g_0 = \overline{g}$이고
$h_0 = \overline{h}$인 인수분해 $f = gh$가 $R[T]$에 존재한다.
\item 임의의 $f \in R[T]$와 $\gcd(g_0, h_0) = 1$인 임의의 인수분해
$\overline{f} = g_0 h_0$에 대해, $g_0 = \overline{g}$,
$h_0 = \overline{h}$이고 또한 $\deg_T(g) = \deg_T(g_0)$인
인수분해 $f = gh$가 $R[T]$에 존재한다.
\item 임의의 에탈 환 사상 $R \to S$와 $\mathfrak m$ 위에 놓이며
$\kappa = \kappa(\mathfrak q)$인 $S$의 소 아이디얼 $\mathfrak q$에 대해,
$R \to S$의 수축 $\tau : S \to R$가 존재한다.
\item 임의의 에탈 환 사상 $R \to S$와 $\mathfrak m$ 위에 놓이며
$\kappa = \kappa(\mathfrak q)$인 $S$의 소 아이디얼 $\mathfrak q$에 대해,
$\mathfrak q = \tau^{-1}(\mathfrak m)$을 만족하는 $R \to S$의 유일한
수축 $\tau : S \to R$가 존재한다.
\item 모든 유한 $R$-대수는 국소환들의 곱이다.
\item 모든 유한 $R$-대수는 국소환들의 유한 곱이다.
\item 모든 유한형 $R$-대수 $S$는 $R \to A$가 유한이고
$R \to B$가 $\mathfrak m$ 위에 놓인 어떤 소 아이디얼에서도
준유한이 아니도록 $A \times B$로 쓸 수 있다.
\item 모든 유한형 $R$-대수 $S$는 $R \to A$가 유한이고
$\Spec(B \otimes_R \kappa)$의 각 기약 성분의 차원이 $\geq 1$이 되도록
$A \times B$로 쓸 수 있다.
\item 모든 준유한 $R$-대수 $S$는 $R \to A$가 유한이고
$B \otimes_R \kappa = 0$이 되도록 $S = A \times B$로 쓸 수 있다.
\end{enumerate}
\end{lemma}

\begin{proof}
먼저 쉬운 함의들을 나열한다.
\begin{enumerate}
\item 2$\Rightarrow$1: (2)에서는 모든 다항식을 다루고 (1)에서는
최고차항계수가 1인 것만 다루기 때문이다.
\item 5$\Rightarrow$3: (5)에서는 모든 다항식을 다루고 (3)에서는
최고차항계수가 1인 것만 다루기 때문이다.
\item 6$\Rightarrow$4: (6)에서는 모든 다항식을 다루고 (4)에서는
최고차항계수가 1인 것만 다루기 때문이다.
\item 4$\Rightarrow$3은 명백하다.
\item 6$\Rightarrow$5는 명백하다.
\item 8$\Rightarrow$7은 명백하다.
\item 10$\Rightarrow$9는 명백하다.
\item 11$\Leftrightarrow$12는 한 소 아이디얼에서 준유한이라는 정의에 의한다.
\item 11$\Rightarrow$13은 준유한의 정의에 의한다.
\end{enumerate}

\noindent
1$\Rightarrow$8의 증명. (1)을 가정하자.
$R \to S$가 에탈이고 $\mathfrak q \subset S$가
$\kappa(\mathfrak q) \cong \kappa$인 소 아이디얼이라 하자.
명제 \ref{algebra-proposition-etale-locally-standard}에 의해
$g \in S$, $g \not \in \mathfrak q$이고 $R \to S_g$가 표준 에탈인
원소를 찾을 수 있다. $S$를 $S_g$로 바꾸어
$S = R[t]_g/(f)$가 표준 에탈이라고 가정할 수 있다(세부사항은 생략한다).
소 아이디얼 $\mathfrak q$의 잉여체가 $\kappa$이므로 이는
$\overline{g}$의 근이 아닌 $\overline{f}$의 근 $a_0$에 대응한다.
표준 에탈 대수의 정의에 의해 이는 또한
$\overline{f'}(a_0) \not = 0$임을 뜻한다.
또한 표준 에탈 대수의 정의에 의해 $f$는 최고차항계수가 1이므로,
$R$이 헨젤이라는 사실을 사용하여 $a_0 = \overline{a}$이고
$f(a) = 0$인 $a \in R$가 존재한다고 결론낼 수 있다.
이는 $g(a)$가 $R$의 단위원임을 뜻하고, 규칙 $t \mapsto a$로 원하는 사상
$\tau : S = R[t]_g/(f) \to R$를 얻는다.
구성에 의해 $\tau^{-1}(\mathfrak m) = \mathfrak q$이다.
보조정리 \ref{algebra-lemma-uniqueness}에 의해 사상 $\tau$는 유일하다.
이로써 (8)이 성립함을 보였다.

\medskip\noindent
7$\Rightarrow$8의 증명. (이는 전혀 중요하지 않으므로 건너뛰어도 좋다.)
(7)이 성립하고 $R \to S$가 에탈이라고 가정하자.
$\mathfrak q_1, \ldots, \mathfrak q_r$를 $\mathfrak m$ 위에 놓인
$S$의 다른 소 아이디얼들이라 하자. 그러면
$g \in S$, $g \not \in \mathfrak q$이고 $i = 1, \ldots, r$에 대해
$g \in \mathfrak q_i$인 원소를 찾을 수 있다. 실제로 그렇지 않으면
$$
\bigcap_{i=1}^{r} \mathfrak{q}_{i} \not\subset \mathfrak{q}
$$
임을 다음과 같이 논증할 수 있다. 그렇지 않다면 어떤 $i$에 대해
$$
\mathfrak{q}_{i} \subset \mathfrak{q}
$$
일 텐데, 에탈 사상의 섬유는 이산적이므로 이는 불가능하다
(예를 들어 보조정리 \ref{algebra-lemma-etale-over-field}을 사용하라).
에탈 환 사상 $R \to S_g$와 소 아이디얼 $\mathfrak qS_g$에 (7)을 적용하자.
그러면 합성 $\tau : S \to S_g \to R$가
$\tau^{-1}(\mathfrak m) = \mathfrak q$를 만족하도록 하는 수축
$\tau_g : S_g \to R$를 얻는다. 세부사항은 생략한다.

\medskip\noindent
8$\Rightarrow$11의 증명. (8)을 가정하고 $R \to S$를 유한형 환 사상이라 하자.
보조정리 \ref{algebra-lemma-etale-makes-quasi-finite-finite}을 적용한다.
그러면 에탈 환 사상 $R \to R'$와 $\mathfrak m$ 위에 놓이며
$\kappa = \kappa(\mathfrak m')$인 소 아이디얼
$\mathfrak m' \subset R'$를 얻는다. 이때
$R' \otimes_R S = A' \times B'$이고, $A'$는 $R'$ 위에서 유한이며,
$B'$는 $\mathfrak m'$ 위에 놓인 어떤 소 아이디얼에서도
$R'$ 위에서 준유한이 아니다. (8)을 적용하여
$\mathfrak m = \tau^{-1}(\mathfrak m')$을 만족하는 수축
$\tau : R' \to R$를 얻는다.
% P01-E0957: the frozen inverse-image equality is ill-typed and reversed; for tau : R' -> R it should be m' = tau^{-1}(m).
이제
$$
S = (S \otimes_R R') \otimes_{R', \tau} R
= (A' \times B') \otimes_{R', \tau} R
= (A' \otimes_{R', \tau} R)  \times  (B' \otimes_{R', \tau} R)
$$
임을 사용하면 (11)과 같은 분해를 얻는다.

\medskip\noindent
8$\Rightarrow$10의 증명. (8)을 가정하고 $R \to S$를 유한 환 사상이라 하자.
보조정리 \ref{algebra-lemma-etale-makes-quasi-finite-finite}을 적용한다.
그러면 에탈 환 사상 $R \to R'$와 $\mathfrak m$ 위에 놓이며
$\kappa = \kappa(\mathfrak m')$인 소 아이디얼
$\mathfrak m' \subset R'$를 얻는다. 이때
$R' \otimes_R S = A'_1 \times \ldots \times A'_n \times B'$이고,
$A'_i$는 $R'$ 위에서 유한하고 $\mathfrak m'$ 위의 소 아이디얼을
정확히 하나 가지며, $B'$는 $\mathfrak m'$ 위에 놓인 어떤
소 아이디얼에서도 $R'$ 위에서 준유한이 아니다.
(8)을 적용하여 $\mathfrak m' = \tau^{-1}(\mathfrak m)$을 만족하는
수축 $\tau : R' \to R$를 얻는다. 그러면
\begin{align*}
S & = (S \otimes_R R') \otimes_{R', \tau} R \\
& = (A'_1 \times \ldots \times A'_n \times B') \otimes_{R', \tau} R \\
& = (A'_1 \otimes_{R', \tau} R)  \times
\ldots \times (A'_1 \otimes_{R', \tau} R) \times
(B' \otimes_{R', \tau} R) \\
& = A_1 \times \ldots \times A_n \times B
\end{align*}
% P01-E0958: the frozen third line repeats A'_1 in the final A-factor; it should end with A'_n.
인자 $B$는 $R$ 위에서 유한이지만 $R \to B$는 $\mathfrak m$ 위에
놓인 어떤 소 아이디얼에서도 준유한이 아니다. 따라서 $B = 0$이다.
인자 $A_i$는 $\mathfrak m$ 위에 놓인 소 아이디얼을 정확히 하나 갖는
유한 $R$-대수이므로 국소환이다. 이로써 $S$가 국소환들의 유한 곱임을 보였다.

\medskip\noindent
9$\Rightarrow$10의 증명. $S$가 국소환 $R$ 위에서 유한이면
극대 아이디얼을 유한 개 이하만 가지므로 성립한다.
실제로 $R \to S$의 상승 성질에 의해 $S$의 모든 극대 아이디얼은
$\mathfrak m$ 위에 놓이고, $S/\mathfrak mS$는 아르틴 환이므로
소 아이디얼을 유한 개만 갖는다.

\medskip\noindent
10$\Rightarrow$1의 증명. (10)을 가정하자.
$f \in R[T]$를 최고차항계수가 1인 다항식이라 하고,
$a_0 \in \kappa$를 $\overline{f}$의 단순근이라 하자.
그러면 $S = R[T]/(f)$는 유한 $R$-대수이다. (10)을 적용하면
$S = A_1 \times \ldots \times A_r$는 국소 $R$-대수들의 유한 곱이다.
특히 $S/\mathfrak mS = \prod A_i/\mathfrak mA_i$는
$\kappa[T]/(\overline{f})$의 국소환들의 곱으로의 분해이다.
이는 인자 중 하나, 이를테면 $A_1/\mathfrak mA_1$이 몫사상
$\kappa[T]/(\overline{f}) \to \kappa[T]/(T - a_0)$임을 뜻한다.
$A_1$은 유한 자유 $R$-가군 $S$의 직합항이므로 그 자체도 유한 자유
$R$-가군이다. $A_1/\mathfrak mA_1$은 1차원 $\kappa$-벡터공간이므로
$R$-가군으로서 $A_1 \cong R$이다. 이는 명백히 $R \to A_1$이
동형사상임을 뜻한다. 사상 $R[T] \to S \to A_1 \to R$ 아래에서
$T$의 상을 $a \in R$라 하자. 그러면 원하는 대로
$f(a) = 0$이고 $\overline{a} = a_0$이다.

\medskip\noindent
13$\Rightarrow$1의 증명. (13)을 가정하자.
$f \in R[T]$를 최고차항계수가 1인 다항식이라 하고,
$a_0 \in \kappa$를 $\overline{f}$의 단순근이라 하자.
그러면 $S_1 = R[T]/(f)$는 유한 $R$-대수이다.
$\overline{g} = \overline{f}/(T - a_0)$인 임의의 원소
$g \in R[T]$를 택하자. 그러면 $S = (S_1)_g$는
$$
S \otimes_R \kappa \cong \kappa[T]_{\overline{g}}/(\overline{f})
\cong \kappa[T]/(T - a_0) \cong \kappa
$$
를 만족하는 준유한 $R$-대수이다. (13)을 $S$에 적용하면
$A$가 $R$ 위에서 유한이고 $B \otimes_R \kappa = 0$인 분해
$S = A \times B$를 얻는다. 특히
$\kappa \cong S/\mathfrak mS = A/\mathfrak mA$이다.
$A$는 평탄 $R$-대수 $S$의 직합항이므로 유한 평탄이고,
따라서 $R$ 위에서 자유이다. $A/\mathfrak mA$는 1차원
$\kappa$-벡터공간이므로 $R$-가군으로서 $A \cong R$이다.
이는 명백히 $R \to A$가 동형사상임을 뜻한다.
사상 $R[T] \to S \to A \to R$ 아래에서 $T$의 상을
$a \in R$라 하자. 그러면 원하는 대로 $f(a) = 0$이고
$\overline{a} = a_0$이다.

\medskip\noindent
8$\Rightarrow$2의 증명. (8)을 가정하자. $f \in R[T]$를 임의의
다항식이라 하고 $a_0 \in \kappa$를 단순근이라 하자.
% P01-E0959: the frozen proof calls a_0 “a simple root” without specifying that it is a root of f-bar.
그러면 대수 $S = R[T]_{f'}/(f)$는 $R$ 위에서 에탈이다.
$\mathfrak q \subset S$를 $\mathfrak m$과 $T - b$가 생성하는
소 아이디얼이라 하자. 여기서 $b \in R$는
$\overline{b} = a_0$인 임의의 원소이다. $S$와 $\mathfrak q$에
(8)을 적용하여 $\tau : S \to R$를 얻는다.
그러면 상 $\tau(T) = a \in R$가 (2)에서 원하는 원소이다.

\medskip\noindent
이제 (1), (2), (7), (8), (9), (10), (11), (12), (13)이 모두
서로 동치임을 알았다. (3), (4), (5), (6) 중 가장 약한 명제는 (3)이고
가장 강한 명제는 (6)이다. 따라서 (3)이 (1)을 함의하고 (1)이 (6)을
함의함을 보이는 일만 남았다.

\medskip\noindent
3$\Rightarrow$1의 증명. (3)을 가정하자.
$f \in R[T]$를 최고차항계수가 1인 다항식이라 하고,
$a_0 \in \kappa$를 $\overline{f}$의 단순근이라 하자.
그러면 $h_0(a_0) \not = 0$인 인수분해
$\overline{f} = (T - a_0)h_0$를 얻으므로
$\gcd(T - a_0, h_0) = 1$이다. (3)을 적용하여
$\overline{g} = T - a_0$, $\overline{h} = h_0$인 인수분해
$f = gh$를 얻는다. $S = R[T]/(f)$로 두면 이는 유한 자유 $R$-대수이다.
$g$, $h$는 각각 $g$, $h$의 $S$에서의 상도 나타내기로 한다. 그러면 이 등식이
$\mathfrak m$을 법으로 성립하므로 나카야마 보조정리
\ref{algebra-lemma-NAK}에 의해 $gS + hS = S$이다.
$S$에서 $gh = f = 0$이므로 이는 또한 $gS \cap hS = 0$을 뜻한다.
따라서 중국인의 나머지 정리에 의해 $S = S/(g) \times S/(h)$를 얻는다.
이는 $A = S/(g)$가 유한 자유 $R$-가군의 직합항이고, 따라서 유한 자유임을
뜻한다. 또한 $A/\mathfrak mA = \kappa[T]/(T - a_0)$이므로 $A$의 계수는
$1$이다. 따라서 사상 $R \to A$는 동형사상이다.
사상 $R[T] \to S \to A \to R$ 아래에서 $T$의 상을 $a \in R$라 하면
$R$의 원소를 얻으며, 이는 $f(a) = 0$이고 $\overline{a} = a_0$이다.

\medskip\noindent
1$\Rightarrow$6의 증명. (1), 또는 이와 동치인
(1), (2), (7), (8), (9), (10), (11), (12), (13) 전부를 가정하자.
$f \in R[T]$를 다항식이라 하자. $\gcd(g_0, h_0) = 1$인 인수분해
$\overline{f} = g_0h_0$가 주어졌다고 하자.
$g_0$는 최고차항계수가 1이라고 가정해도 좋고 그렇게 하자.
$S = R[T]/(f)$를 생각하자. 인수분해가 있으므로 $f$의 계수들이
$R$의 단위 아이디얼을 생성함을 알 수 있다.
% P01-E0960: if g_0 = 1 and h_0 = 0, the coprime factorization is valid but f-bar is zero, so the coefficients of f need not generate the unit ideal; the trivial factorization must be handled separately.
이는 $S$가 $R$ 위에서 유한 섬유를 가지며 따라서 $R$ 위에서
준유한임을 뜻한다. 또한 보조정리 \ref{algebra-lemma-grothendieck-general}에 의해
$S$가 $R$ 위에서 평탄임을 뜻한다. (13)과 (10)을 결합하여
$S = A_1 \times \ldots \times A_n \times B$로 쓸 수 있다.
여기서 각 $A_i$는 $R$ 위에서 국소적이고 유한이며,
$B \otimes_R \kappa = 0$이다. 인자 $A_1, \ldots, A_n$의 순서를
바꾸어 다음을 가정할 수 있다.
$$
\kappa[T]/(g_0) =
A_1/\mathfrak m A_1 \times \ldots \times A_r/\mathfrak mA_r,
\ \kappa[T]/(h_0) =
A_{r + 1}/\mathfrak mA_{r + 1} \times \ldots \times A_n/\mathfrak mA_n
$$
이는 $\kappa[T]$의 몫들로서의 등식이다. 유한 평탄 $R$-대수
$A = A_1 \times \ldots \times A_r$는 $R$-가군으로서 자유이다.
보조정리 \ref{algebra-lemma-finite-flat-local}을 보라. 그 계수는
$\deg_T(g_0)$이다. $g \in R[T]$를 $R$-선형 연산자
$T : A \to A$의 특성다항식이라 하자. 그러면 $g$는 차수가
$\deg_T(g) = \deg_T(g_0)$인 최고차항계수가 1인 다항식이고,
또한 $\overline{g} = g_0$이다. 케일리--해밀턴 정리
(보조정리 \ref{algebra-lemma-charpoly})에 의해 $T_A$가 $A$에서 $T$의 상을
나타낼 때 $g(T_A) = 0$이다. 따라서 잘 정의된 전사 사상
$R[T]/(g) \to A$를 얻으며, 나카야마 보조정리
\ref{algebra-lemma-NAK}에 의해 이는 동형사상이다. 구성에 의해 사상
$R[T] \to A$는 $R[T]/(f)$를 통해 인수분해되므로 어떤 $h$에 대해
$f = gh$로 쓸 수 있다. 이로써 증명이 끝난다.
\end{proof}

\begin{lemma}
\label{algebra-lemma-finite-over-henselian}
$(R, \mathfrak m, \kappa)$를 헨젤 국소환이라 하자.
\begin{enumerate}
\item $R \to S$가 유한 환 사상이면, $S$는 각각 $R$ 위에서 유한인
헨젤 국소환들의 유한 곱이다.
\item $R \to S$가 유한 환 사상이고 $S$가 국소환이면,
$S$는 헨젤 국소환이고 $R \to S$는 (유한) 국소환 사상이다.
\item $R \to S$가 유한형 환 사상이고 $\mathfrak q$가
$\mathfrak m$ 위에 놓이는 $S$의 소 아이디얼이며 그곳에서
$R \to S$가 준유한이면, $S_{\mathfrak q}$는 헨젤이고 $R$ 위에서 유한이다.
\item $R \to S$가 준유한이면, $\mathfrak m$ 위에 놓인 모든
소 아이디얼 $\mathfrak q$에 대해 $S_{\mathfrak q}$는 헨젤이고
$R$ 위에서 유한이다.
\end{enumerate}
\end{lemma}

\begin{proof}
(2)는 (1)을 함의한다. 실제로 (1)의 $S$는 보조정리
\ref{algebra-lemma-characterize-henselian} (10)에 의해 $\mathfrak m$ 위에 놓인
소 아이디얼들에서의 국소화들의 유한 곱이다.
또한 모든 유한 $S$-대수는 유한 $R$-대수이므로 보조정리
\ref{algebra-lemma-characterize-henselian} (10)에서 (2)가 따른다
(물론 국소환 사이의 모든 유한 환 사상은 국소적이다).

\medskip\noindent
$R \to S$와 $\mathfrak q$가 (3)에서와 같다고 하자.
$A$는 $R$ 위에서 유한이고 $B$는 $\mathfrak m$ 위에 놓인 어떤
소 아이디얼에서도 $R$ 위에서 준유한이 아니도록
$S = A \times B$로 쓰자. 보조정리
\ref{algebra-lemma-characterize-henselian} (11)을 보라.
따라서 $S_\mathfrak q$는 $A$를 한 극대 아이디얼에서 국소화한 것이고,
(1)에서 (3)이 따른다. (4)는 (3)에서 따른다.
\end{proof}

\begin{lemma}
\label{algebra-lemma-mop-up}
$(R, \mathfrak m, \kappa)$를 헨젤 국소환이라 하자.
모든 유한형 $R$-대수 $S$는
$S = A_1 \times \ldots \times A_n \times B$로 쓸 수 있다.
여기서 $A_i$는 국소환이고 $R$ 위에서 유한이며,
$R \to B$는 $\mathfrak m$ 위에 놓이는 $B$의 어떤 소 아이디얼에서도
준유한이 아니다.
\end{lemma}

\begin{proof}
보조정리 \ref{algebra-lemma-characterize-henselian}의 (11)과 (10)을 결합하면 된다.
\end{proof}

\begin{lemma}
\label{algebra-lemma-mop-up-strictly-henselian}
$(R, \mathfrak m, \kappa)$를 엄밀 헨젤 국소환이라 하자.
모든 유한형 $R$-대수 $S$는
$S = A_1 \times \ldots \times A_n \times B$로 쓸 수 있다.
여기서 $A_i$는 국소환이고 $R$ 위에서 유한이며,
$\kappa \subset \kappa(\mathfrak m_{A_i})$는 유한 순수 비분리이고,
$R \to B$는 $\mathfrak m$ 위에 놓이는 $B$의 어떤 소 아이디얼에서도
준유한이 아니다.
\end{lemma}

\begin{proof}
먼저 보조정리 \ref{algebra-lemma-mop-up}에서와 같이
$S = A_1 \times \ldots \times A_n \times B$로 쓰자.
체확대 $\kappa(\mathfrak m_{A_i})/\kappa$는 유한이고
$\kappa$는 분리 대수적으로 닫혀 있으므로, 이 확대는 유한 순수 비분리이다.
\end{proof}

\begin{lemma}
\label{algebra-lemma-henselian-cat-finite-etale}
$(R, \mathfrak m, \kappa)$를 헨젤 국소환이라 하자.
유한 에탈 환확대 $R \to S$들의 범주는 함자
$S \mapsto S/\mathfrak mS$를 통해 유한 에탈 대수
$\kappa \to \overline{S}$들의 범주와 동치이다.
\end{lemma}

\begin{proof}
명제의 범주들 사이의 함자를 $\mathcal{C} \to \mathcal{D}$로 나타내자.
$R \to S$가 유한 에탈이라고 하자. 그러면
$$
S = A_1 \times \ldots \times A_n
$$
으로 쓸 수 있다. 여기서 $A_i$는 국소환이고 $S$ 위에서 유한 에탈이다.
보조정리 \ref{algebra-lemma-mop-up} 또는 보조정리
\ref{algebra-lemma-characterize-henselian} (10)을 사용하라.
특히 보조정리 \ref{algebra-lemma-etale-at-prime}에 의해
$A_i/\mathfrak mA_i$는 $\kappa$의 유한 분리 체확대이다.
따라서 $\mathcal{C}$와 $\mathcal{D}$의 모든 대상이 주어진 함자 아래에서
서로 대응하는 기약 조각들로 표준적으로 분해됨을 알 수 있다.
다음으로 $S_1$, $S_2$가 $R$ 위에서 유한 에탈이고
$\kappa_1 = S_1/\mathfrak mS_1$ 및
$\kappa_2 = S_2/\mathfrak mS_2$가 체라고 하자
(이들은 $\kappa$의 유한 분리 확대이다).
그러면 $S_1 \otimes_R S_2$는 $R$ 위에서 유한 에탈이고, 앞에서처럼
$$
S_1 \otimes_R S_2 = A_1 \times \ldots \times A_n
$$
으로 쓸 수 있다. 그러면 $\Hom_R(S_1, S_2)$는
$S_2 \to A_i$가 동형사상인 지표 $i \in \{1, \ldots, n\}$들의
집합과 동일시된다. 이를 보이려면 임의의 $R$-대수 사상
$\varphi : S_1 \to S_2$에 대해 사상
$\varphi \times 1 : S_1 \otimes_R S_2 \to S_2$가 전사이고,
따라서 인자 $A_i$ 중 하나로의 사영과 같다는 사실을 사용한다.
정확히 같은 방식으로 $\Hom_\kappa(\kappa_1, \kappa_2)$는
$\kappa_2 \to A_i/\mathfrak mA_i$가 동형사상인 지표
$i \in \{1, \ldots, n\}$들의 집합과 동일시된다.
위 논의에 의해 이 지표 집합들은 일치하므로 함자는 충실충만하다.
마지막으로 $\kappa'/\kappa$를 유한 분리 체확대라 하자.
보조정리 \ref{algebra-lemma-make-etale-map-prescribed-residue-field}에 의해
에탈 환 사상 $R \to S$와 $\mathfrak m$ 위에 놓이는 $S$의 소 아이디얼
$\mathfrak q$ 중에서 $\kappa \subset \kappa(\mathfrak q)$가 주어진
확대와 동형인 것이 존재한다. (1)에 의해
$S = A_1 \times \ldots \times A_n \times B$로 쓸 수 있다.
% P01-E0961: the frozen proof cites an undefined “part (1)” here; the required global finite-plus-non-quasi-finite decomposition is lemma-mop-up.
$R \to S$는 준유한이므로 $\mathfrak m$ 위에 놓이는 $B$의 소 아이디얼은
없다. 따라서 어떤 $i$에 대해 $S_{\mathfrak q}$는 $A_i$와 같다.
그러므로 $R \to A_i$는 유한 에탈이고 주어진 잉여체 확대를 낳는다.
따라서 함자는 본질적으로 전사이며, 결론을 얻는다.
\end{proof}

\begin{lemma}
\label{algebra-lemma-unramified-over-strictly-henselian}
$(R, \mathfrak m, \kappa)$를 엄밀 헨젤 국소환이라 하자.
$R \to S$를 비분기 환 사상이라 하자. 그러면
$$
S = A_1 \times \ldots \times A_n \times B
$$
이고, 각 $R \to A_i$는 전사이며 $\mathfrak m$ 위에 놓이는
$B$의 소 아이디얼은 없다.
\end{lemma}

\begin{proof}
먼저 보조정리 \ref{algebra-lemma-mop-up}에서와 같이
$S = A_1 \times \ldots \times A_n \times B$로 쓰자.
이제 $R \to A_i$는 유한 비분기이고 $A_i$는 국소환이다.
따라서 $A_i$의 극대 아이디얼은 $\mathfrak mA_i$이고, 그 잉여체
$A_i / \mathfrak m A_i$는 $\kappa$의 유한 분리 확대이다.
보조정리 \ref{algebra-lemma-unramified-at-prime}을 보라.
그러나 $R$이 엄밀 헨젤이라는 조건은 $\kappa$가 분리 대수적으로
닫혀 있음을 뜻하므로 $\kappa = A_i / \mathfrak m A_i$이다.
나카야마 보조정리 \ref{algebra-lemma-NAK}에 의해 원하는 대로
$R \to A_i$가 전사라고 결론낸다.
\end{proof}

\begin{lemma}
\label{algebra-lemma-complete-henselian}
\begin{slogan}
완비 국소환은 뉴턴 방법에 의해 헨젤이다
\end{slogan}
$(R, \mathfrak m, \kappa)$를 완비 국소환이라 하자.
정의 \ref{algebra-definition-complete-local-ring}을 보라. 그러면 $R$은 헨젤이다.
\end{lemma}

\begin{proof}
$f \in R[T]$를 최고차항계수가 1인 다항식이라 하자.
$f_n \in R/\mathfrak m^{n + 1}[T]$를 그 상이라 하고,
$T$에 관한 $f_n$의 도함수를 $f'_n$이라 하자.
$a_0 \in \kappa$를 $f_0$의 단순근이라 하자.
다음과 같이 귀납적으로 이를 $R$ 위에서 $f$의 해로 올린다.
$a_n \bmod \mathfrak m = a_0$이고 $f_n(a_n) = 0$인
$a_n \in R/\mathfrak m^{n + 1}$이 주어졌다고 하자.
$a_n = b \bmod \mathfrak m^{n + 1}$인 임의의 원소
$b \in R/\mathfrak m^{n + 2}$를 택하자. 그러면
$f_{n + 1}(b) \in \mathfrak m^{n + 1}/\mathfrak m^{n + 2}$이다. 다음과 두자.
$$
a_{n + 1} = b - f_{n + 1}(b)/f'_{n + 1}(b)
$$
(뉴턴 방법이다.) $a_0$에 대한 조건에 의해
$f'_{n + 1}(b) \in R/\mathfrak m^{n + 2}$가 가역이므로 이는 의미가 있다.
그러면 $R/\mathfrak m^{n + 2}$에서
$f_{n + 1}(a_{n + 1}) = f_{n + 1}(b) - f_{n + 1}(b) = 0$임을 계산한다.
이렇게 구성한 원소들의 계 $a_n \in R/\mathfrak m^{n + 1}$은
양립하므로 원소 $a \in \lim R/\mathfrak m^n = R$를 얻는다
(여기서 $R$이 완비라는 사실을 사용한다).
또한 $f(a) = 0$이다. 이는 각 $R/\mathfrak m^n$에서 영으로 가기 때문이다.
마지막으로 $\overline{a} = a_0$이고, 결론을 얻는다.
\end{proof}

\begin{lemma}
\label{algebra-lemma-local-dimension-zero-henselian}
\begin{slogan}
차원 영인 국소환은 헨젤이다.
\end{slogan}
$(R, \mathfrak m)$을 차원 $0$인 국소환이라 하자. 그러면 $R$은 헨젤이다.
\end{lemma}

\begin{proof}
$R \to S$를 유한 환 사상이라 하자. 보조정리
\ref{algebra-lemma-characterize-henselian}에 의해 $S$가 국소환들의 곱임을
보이면 충분하다. 보조정리 \ref{algebra-lemma-finite-finite-fibres}에 의해
$S$는 유한 개의 소 아이디얼 $\mathfrak m_1, \ldots, \mathfrak m_r$를
가지며, 이들은 모두 $\mathfrak m$ 위에 놓인다.
보조정리 \ref{algebra-lemma-integral-no-inclusion}에 의해 이 소 아이디얼들
사이에는 포함관계가 없으므로, 모두 극대 아이디얼이다.
보조정리 \ref{algebra-lemma-Zariski-topology}에 의해
$\mathfrak m_1 \cap \ldots \cap \mathfrak m_r$의 모든 원소는 멱영이다.
따라서 보조정리 \ref{algebra-lemma-product-local}에 의해 $S$는
소 아이디얼 $\mathfrak m_i$에서 $S$를 국소화한 것들의 곱이다.
\end{proof}

\noindent
다음 보조정리는 헨젤화와 엄밀 헨젤화의 유일성 및 함자적 성질의
핵심이 될 것이다.

\begin{lemma}
\label{algebra-lemma-map-into-henselian}
$R \to S$를 $S$가 헨젤 국소환인 환 사상이라 하자. 다음이 주어졌다고 하자.
\begin{enumerate}
\item 에탈 환 사상 $R \to A$,
\item $\mathfrak p = R \cap \mathfrak m_S$ 위에 놓이는
$A$의 소 아이디얼 $\mathfrak q$,
% P01-E0962: for a possibly noninjective map R -> S, the frozen expression R intersect m_S should be the inverse image of m_S in R.
\item $\kappa(\mathfrak p)$-대수 사상
$\tau : \kappa(\mathfrak q) \to S/\mathfrak m_S$.
\end{enumerate}
그러면 $\mathfrak q = f^{-1}(\mathfrak m_S)$이고 $f$가 잉여체 위에서
사상 $\tau$를 유도하는 유일한 $R$-대수 준동형 $f : A \to S$가 존재한다.
\end{lemma}

\begin{proof}
$A \otimes_R S$를 생각하자. 이는 $S$ 위의 에탈 대수이다.
보조정리 \ref{algebra-lemma-etale}을 보라. 또한 핵
$$
\mathfrak q' = \Ker(A \otimes_R S \to
\kappa(\mathfrak q) \otimes_{\kappa(\mathfrak p)} \kappa(\mathfrak m_S)
\xrightarrow{\tau \otimes 1}
\kappa(\mathfrak m_S))
$$
은 $\mathfrak m_S$ 위에 놓이고 $S$의 잉여체와 같은 잉여체를 갖는
소 아이디얼이다. 따라서 보조정리 \ref{algebra-lemma-characterize-henselian}에 의해
$\sigma^{-1}(\mathfrak m_S) = \mathfrak q'$인 유일한 수축
$\sigma : A \otimes_R S \to S$가 존재한다.
$f$를 합성 $A \to A \otimes_R S \to S$와 같게 두자.
$f$의 성질들의 확인은 생략한다. $f$의 유일성은 $\sigma$의 유일성에서
따른다(세부사항은 생략한다).
\end{proof}

\begin{lemma}
\label{algebra-lemma-strictly-henselian-solutions}
$\varphi : R \to S$를 엄밀 헨젤 국소환들의 국소 준동형사상이라 하자.
$P_1, \ldots, P_n \in R[x_1, \ldots, x_n]$을
$R[x_1, \ldots, x_n]/(P_1, \ldots, P_n)$이 $R$ 위에서 에탈이 되게 하는
다항식들이라 하자. 그러면 사상
$$
R^n \longrightarrow S^n, \quad
(h_1, \ldots, h_n) \longmapsto (\varphi(h_1), \ldots, \varphi(h_n))
$$
은 다음 두 집합 사이의 전단사를 유도한다.
$$
\{
(r_1, \ldots, r_n) \in R^n
\mid
P_i(r_1, \ldots, r_n) = 0, \ i = 1, \ldots, n
\}
$$
및
$$
\{
(s_1, \ldots, s_n) \in S^n
\mid
P^\varphi_i(s_1, \ldots, s_n) = 0, \ i = 1, \ldots, n
\}
$$
여기서 $P^\varphi_i \in S[x_1, \ldots, x_n]$는 $\varphi$ 아래에서
$P_i$의 상이다.
\end{lemma}

\begin{proof}
첫 번째 해집합은 집합
$$
\Hom_R(R[x_1, \ldots, x_n]/(P_1, \ldots, P_n), R).
$$
과 표준적으로 동형이다. $R$이 헨젤이므로 사상
$R \to R/\mathfrak m_R$는 이 집합과 잉여체
$R/\mathfrak m_R$에서의 해집합 사이의 전단사를 유도한다.
보조정리 \ref{algebra-lemma-characterize-henselian}을 보라. $S$에 대해서도 마찬가지이다.
이제 $R[x_1, \ldots, x_n]/(P_1, \ldots, P_n)$은 $R$ 위에서 에탈이고
$R/\mathfrak m_R$는 분리 대수적으로 닫혀 있으므로,
$R/\mathfrak m_R[x_1, \ldots, x_n]/(\overline{P}_1, \ldots, \overline{P}_n)$은
$R/\mathfrak m_R$의 복사본들의 유한 곱이다. 여기서 $\overline{P}_i$는
$R/\mathfrak m_R[x_1, \ldots, x_n]$에서 $P_i$의 상이다.
% P01-E0963: the frozen residue-field polynomial-ring expressions omit parentheses around R/m_R (and later S/m_S), making the displayed quotients formally ill-typed or ambiguous.
따라서 텐서곱
$$
R/\mathfrak m_R[x_1, \ldots, x_n]/(\overline{P}_1, \ldots, \overline{P}_n)
\otimes_{R/\mathfrak m_R} S/\mathfrak m_S
=
S/\mathfrak m_S[x_1, \ldots, x_n]/
(\overline{P}^\varphi_1, \ldots, \overline{P}^\varphi_n)
$$
도 같은 지표 집합을 갖는 $S/\mathfrak m_S$의 복사본들의 유한 곱이다.
이로써 보조정리가 증명된다.
\end{proof}

\begin{lemma}
\label{algebra-lemma-split-ML-henselian}
$R$을 헨젤 국소환이라 하자. $R$ 위의 모든 가산 생성
미타그--레플러 가군은 유한 표시 $R$-가군들의 직합이다.
\end{lemma}

\begin{proof}
$M$을 가산 생성 미타그--레플러 $R$-가군이라 하자.
임의의 원소 $x \in M$에 대해 $x \in N$이고 $N$은 유한 표시이며
$K$는 미타그--레플러인 직합 분해 $M = N \oplus K$가 존재한다고 주장한다.

\medskip\noindent
이 주장이 참이라고 하자. $M$의 생성원들
$x_1, x_2, x_3, \ldots$를 택하자. 주장에 의해 귀납적으로
$$
M = N_1 \oplus N_2 \oplus \ldots \oplus N_n \oplus K_n
$$
인 직합 분해들을 찾을 수 있다. 여기서 $N_i$는 유한 표시이고,
$x_1, \ldots, x_n \in N_1 \oplus \ldots \oplus N_n$이며,
$K_n$은 미타그--레플러이다. 이를 무한히 반복하면
$M = \bigoplus N_i$임을 알 수 있다.

\medskip\noindent
주장을 증명하는 일만 남았다. $x \in M$이라 하자.
보조정리 \ref{algebra-lemma-ML-countable}에 의해 $\alpha$가 유한 표시 가군을 통해
인수분해되고 $\alpha (x) = x$를 만족하는 자기준동형
$\alpha : M \to M$이 존재한다. $\alpha$가 다음과 같이 인수분해된다고 하자.
$$
\xymatrix{
M \ar[r]^\pi & P \ar[r]^i & M
}
$$
$a = \pi \circ \alpha \circ i : P \to P$로 두면
$i \circ a \circ \pi = \alpha^3$이다. 보조정리
\ref{algebra-lemma-charpoly-module}에 의해 $P(a) = 0$인 최고차항계수가 1인
다항식 $P \in R[T]$가 존재한다.
% P01-E0964: the frozen proof reuses P for both the finitely presented module in the diagram and the annihilating monic polynomial, making later references such as “N is a quotient of P” ambiguous.
이는 형식적으로 $\alpha^2 P(\alpha) = 0$임을 뜻함을 주의하자.
따라서 $M$을 $R[T]/(T^2P)$ 위의 가군으로 생각할 수 있다.
$x \not = 0$이라고 가정하자. 그러면 $\alpha(x) = x$로부터
$0 = \alpha^2P(\alpha)x = P(1)x$이고, 따라서
$I = \{r \in R \mid rx = 0\}$가 $x$의 소멸자일 때
$R/I$에서 $P(1) = 0$이다. $x \not = 0$이므로
$I \subset \mathfrak m_R$이고, 따라서 $1$은
$\overline{P} = P \bmod \mathfrak m_R \in R/\mathfrak m_R[T]$의 근이다.
$R$이 헨젤이므로 다음 인수분해를 찾을 수 있다.
$$
T^2P = (T^2 Q_1) Q_2
$$
여기서 $Q_1, Q_2 \in R[T]$이고,
$Q_2 = (T - 1)^e \bmod \mathfrak m_R R[T]$이며
$Q_1(1) \not = 0 \bmod \mathfrak m_R$이다.
보조정리 \ref{algebra-lemma-characterize-henselian}을 보라.
$N = \Im(\alpha^2Q_1(\alpha) : M \to M)$ 및
$K = \Im(Q_2(\alpha) : M \to M)$로 두자.
$T^2Q_1$과 $Q_2$가 $R[T]$의 단위 아이디얼을 생성하므로
직합 분해 $M = N \oplus K$를 얻는다. 또한 $Q_2$는 $N$ 위에서 영으로
작용하고 $T^2Q_1$은 $K$ 위에서 영으로 작용한다.
$N$은 $P$의 몫이므로 유한 생성임을 주의하자. 또한
$\alpha^2Q_1(\alpha)x = Q_1(1)x$이고 $Q_1(1)$이 $R$의 단위원이므로
$x \in N$이다. 보조정리 \ref{algebra-lemma-direct-sum-ML}에 의해
$N$과 $K$는 미타그--레플러이다. 마지막으로 유한 생성
미타그--레플러 가군은 유한 표시이므로, 유한 생성 가군 $N$은 유한 표시이다.
예제 \ref{algebra-example-ML} (1)을 보라.
\end{proof}




\section{에탈 환 사상들의 여과 여극한}
\label{algebra-section-ind-etale}

\noindent
이 절은 ind-에탈 환 사상을 다루는 절의 준비이다
(프로에탈 코호몰로지, 절 \ref{proetale-section-ind-etale}).
이 내용은 국소환의 헨젤화와 엄밀 헨젤화의 유일성 성질을 증명할 때도
유용할 것이다.

\begin{lemma}
\label{algebra-lemma-base-change-colimit-etale}
$R \to A$와 $R \to R'$를 환 사상이라 하자.
$A$가 에탈 $R$-대수들의 여과 여극한이면,
$R' \otimes_R A$도 에탈 $R'$-대수들의 여과 여극한이다.
% P01-E0965: the frozen English statement duplicates its predicate (“then so is ... is a filtered colimit”).
\end{lemma}

\begin{proof}
여극한은 텐서곱과 가환하고, 에탈 환 사상은 기저변환 아래 보존되므로
참이다(보조정리 \ref{algebra-lemma-etale}).
\end{proof}

\begin{lemma}
\label{algebra-lemma-composition-colimit-etale}
$A \to B \to C$를 환 사상이라 하자. $B$가 에탈 $A$-대수들의
여과 여극한이고 $C$가 에탈 $B$-대수들의 여과 여극한이면,
$C$는 에탈 $A$-대수들의 여과 여극한이다.
\end{lemma}

\begin{proof}
보조정리 \ref{algebra-lemma-when-colimit}의 판정법을 사용한다.
$P$가 $A$ 위에서 유한 표시인 $A \to C$의 인수분해
$A \to P \to C$를 택하자. $I$가 유향 집합이고 $B_i$가 에탈
$A$-대수가 되도록 $B = \colim_{i \in I} B_i$로 쓰자.
$J$가 유향 집합이고 $C_j$가 에탈 $B$-대수가 되도록
$C = \colim_{j \in J} C_j$로 쓰자.
보조정리 \ref{algebra-lemma-characterize-finite-presentation}에 의해 어떤 $j$에 대해
$P \to C$를 $P \to C_j \to C$로 인수분해할 수 있다.
보조정리 \ref{algebra-lemma-etale}에 의해 어떤 $i \in I$와 에탈 환 사상
$B_i \to C'_j$ 중 $C_j = B \otimes_{B_i} C'_j$가 되는 것을 찾을 수 있다.
그러면 $C_j = \colim_{i' \geq i} B_{i'} \otimes_{B_i} C'_j$이고,
다시 충분히 큰 지표에 대해 $P \to C_j$가
$P \to B_{i'} \otimes_{B_i} C'_j \to C$로 인수분해됨을 알 수 있다.
에탈 환 사상들의 합성과 텐서곱은 에탈이므로
$A \to C' = B_{i'} \otimes_{B_i} C'_j$는 에탈이다.
% P01-E0966: the frozen English source begins this sentence with “As” and then starts a second causal “as”, leaving a sentence fragment.
따라서 $C'$가 $A$ 위에서 에탈이도록 $P \to C$를
$P \to C' \to C$로 인수분해했고, 보조정리
\ref{algebra-lemma-when-colimit}의 판정법이 적용된다.
\end{proof}

\begin{lemma}
\label{algebra-lemma-colimit-colimit-etale}
$R$을 환이라 하자. $A = \colim A_i$를 $R$-대수들의 여과 여극한이라 하고,
각 $A_i$가 에탈 $R$-대수들의 여과 여극한이라고 하자.
그러면 $A$는 에탈 $R$-대수들의 여과 여극한이다.
\end{lemma}

\begin{proof}
$J_i$가 유향 집합이고 $A_j$가 에탈 $R$-대수가 되도록
$A_i = \colim_{j \in J_i} A_j$로 쓰자. 각 $i \leq i'$와
$j \in J_i$에 대해, 도식
$$
\xymatrix{
A_i \ar[r] & A_{i'} \\
A_j \ar[u] \ar[r]^{\varphi_{jj'}} & A_{j'} \ar[u]
}
$$
을 가환하게 하는 $j' \in J_{i'}$와 $R$-대수 사상
$\varphi_{jj'} : A_j \to A_{j'}$가 존재한다.
$R \to A_j$가 유한 표시이므로 보조정리
\ref{algebra-lemma-characterize-finite-presentation}이 적용되어 이는 참이다.
$\mathcal{J}$를 대상들이 $\coprod_{i \in I} J_i$이고 사상들이 위와 같은
삼중항 $(j, j', \varphi_{jj'})$인 범주로 두자
(합성법은 자명하다). 그러면 $\mathcal{J}$는 여과 범주이고
$A = \colim_\mathcal{J} A_j$이다. 세부사항은 생략한다.
\end{proof}

\begin{lemma}
\label{algebra-lemma-colimit-colimit-etale-better}
$I$를 유향 집합이라 하고, $i \mapsto (R_i \to A_i)$를
$I$ 위의 환 사상들의 계라 하자. $R = \colim R_i$ 및
$A = \colim A_i$로 두자. 각 $A_i$가 에탈 $R_i$-대수들의 여과
여극한이면, $A$는 에탈 $R$-대수들의 여과 여극한이다.
\end{lemma}

\begin{proof}
$A = A \otimes_R R = \colim A_i \otimes_{R_i} R$이므로 참이다.
실제로 보조정리 \ref{algebra-lemma-colimit-colimit-etale}을 적용할 수 있다.
이는 보조정리 \ref{algebra-lemma-base-change-colimit-etale}에 의해
$R \to A_i \otimes_{R_i} R$가 에탈 환 사상들의 여과 여극한이기 때문이다.
\end{proof}

\begin{lemma}
\label{algebra-lemma-colimits-of-etale}
$R$을 환이라 하고 $A \to B$를 $R$-대수 준동형사상이라 하자.
$A$와 $B$가 에탈 $R$-대수들의 여과 여극한이면,
$B$는 에탈 $A$-대수들의 여과 여극한이다.
\end{lemma}

\begin{proof}
$A_i$와 $B_j$가 $R$ 위에서 에탈이도록 $A = \colim A_i$와
$B = \colim B_j$를 여과 여극한으로 쓰자. 각 $i$에 대해
$A_i \to B$가 $B_j$를 통해 인수분해되는 $j$를 찾을 수 있다.
보조정리 \ref{algebra-lemma-characterize-finite-presentation}을 보라.
보조정리 \ref{algebra-lemma-map-between-etale}에 의해 인수분해
$A_i \to B_j$는 에탈이다.
$A \to A \otimes_{A_i} B_j$는 에탈이므로(보조정리 \ref{algebra-lemma-etale}),
$B = \colim A \otimes_{A_i} B_j$임을 보이면 충분하다.
여기서 여극한은 $A_i \to B$의 인수분해
$A_i \to B_j \to B$와 쌍 $(i, j)$ 위에서 취한다
(이는 유향 계이며, 세부사항은 생략한다).
여극한은 텐서곱과 가환하므로 이는 명백하다. 따라서
$\colim A \otimes_{A_i} B_j = A \otimes_A B = B$이다.
\end{proof}

\begin{lemma}
\label{algebra-lemma-map-into-henselian-colimit}
$R \to S$를 $S$가 헨젤 국소환인 환 사상이라 하자. 다음이 주어졌다고 하자.
\begin{enumerate}
\item 에탈 $R$-대수들의 여과 여극한인 $R$-대수 $A$,
\item $\mathfrak p = R \cap \mathfrak m_S$ 위에 놓이는
$A$의 소 아이디얼 $\mathfrak q$,
% P01-E0967: as in P01-E0962, the frozen source uses an intersection for the contraction along a possibly noninjective ring map.
\item $\kappa(\mathfrak p)$-대수 사상
$\tau : \kappa(\mathfrak q) \to S/\mathfrak m_S$.
\end{enumerate}
그러면 $\mathfrak q = f^{-1}(\mathfrak m_S)$이고
$f \bmod \mathfrak q = \tau$인 유일한 $R$-대수 준동형사상
$f : A \to S$가 존재한다.
\end{lemma}

\begin{proof}
$A = \colim A_i$를 에탈 $R$-대수들의 여과 여극한으로 쓰자.
$\mathfrak q_i = A_i \cap \mathfrak q$로 두자.
% P01-E0968: the frozen source again writes an intersection although the transition map A_i -> A need not be injective; the contraction is the inverse image.
보조정리 \ref{algebra-lemma-map-into-henselian}을 적용하여
$f_i : A_i \to S$를 얻는다. $f = \colim f_i$로 두자.
\end{proof}

\begin{lemma}
\label{algebra-lemma-uniqueness-henselian}
$R$을 환이라 하자. 환 사상들의 가환 도식
$$
\xymatrix{
S \ar[r] & K \\
R \ar[u] \ar[r] & S' \ar[u]
}
$$
이 주어졌다고 하자. 여기서 $S$, $S'$는 헨젤 국소환이고,
$S$, $S'$는 에탈 $R$-대수들의 여과 여극한이며, $K$는 체이고,
화살표 $S \to K$와 $S' \to K$는 $K$를 $S$와 $S'$ 양쪽의
잉여체와 동일시한다. 그러면 $K$로 가는 사상들과 양립하는 유일한
$R$-대수 동형사상 $S \to S'$가 존재한다.
\end{lemma}

\begin{proof}
보조정리 \ref{algebra-lemma-map-into-henselian-colimit}에서 바로 따른다.
\end{proof}

\noindent
다음 보조정리는 엄밀히 말해 에탈 환 사상들의 여극한에 관한 것이 아니다.

\begin{lemma}
\label{algebra-lemma-colimit-henselian}
국소 준동형사상들에 따른 (엄밀) 헨젤 국소환들의 여과 여극한은
(엄밀) 헨젤이다.
\end{lemma}

\begin{proof}
범주론, 보조정리 \ref{categories-lemma-directed-category-system}에 의해
이는 실제로 유향 집합 위의 (엄밀) 헨젤 국소환들의 여극한에 관한 문제이다.
각 $\varphi_{ii'}$가 국소적인 그러한 계
$(R_i, \varphi_{ii'})$를 택하자. 그러면 $R = \colim_i R_i$는 국소환이고,
그 잉여체 $\kappa$는 $\colim \kappa_i$이다(논증은 생략한다).
각 $\kappa_i$가 분리 대수적으로 닫혀 있으면
$\colim \kappa_i$도 그러함을 쉽게 알 수 있다. 따라서 각 $R_i$가
헨젤일 때 $R$이 헨젤임을 보이면 충분하다.
$f \in R[T]$가 최고차항계수가 1이고 $a_0 \in \kappa$가
$\overline{f}$의 단순근이라고 하자. 그러면 충분히 큰 어떤 $i$에 대해
$f$로 가는 $f_i \in R_i[T]$와 $a_0$로 가는
$a_{0, i} \in \kappa_i$가 존재한다.
$\overline{f_i}(a_{0, i}) \in \kappa_i$와 각각
$\overline{f_i'}(a_{0, i}) \in \kappa_i$는 각각
$0 = \overline{f}(a_0) \in \kappa$와
$0 \not = \overline{f'}(a_0) \in \kappa$로 가므로,
$a_{0, i}$는 $\overline{f_i}$의 단순근이라고 결론낸다.
$R_i$가 헨젤이므로 $f_i(a_i) = 0$이고
$a_{0, i} = \overline{a_i}$인 $a_i \in R_i$를 찾을 수 있다.
그러면 $a_i$의 상 $a \in R$가 원하는 해이다.
따라서 $R$은 헨젤이다.
\end{proof}




\section{헨젤화와 엄밀 헨젤화}
\label{algebra-section-henselization}

\noindent
이 절에서는 헨젤화를 구성한다. 이 내용을 읽는 동안 이미 증명한
유일성인 보조정리 \ref{algebra-lemma-uniqueness-henselian}과
보조정리 \ref{algebra-lemma-map-into-henselian-colimit}에서 지적한
함자적 거동을 염두에 두기를 권한다.

\begin{lemma}
\label{algebra-lemma-henselization}
$(R, \mathfrak m, \kappa)$를 국소환이라 하자. 다음 성질을 갖는
국소환 사상 $R \to R^h$가 존재한다.
\begin{enumerate}
\item $R^h$는 헨젤이다.
\item $R^h$는 에탈 $R$-대수들의 여과 여극한이다.
\item $\mathfrak m R^h$는 $R^h$의 극대 아이디얼이다.
\item $\kappa = R^h/\mathfrak m R^h$이다.
\end{enumerate}
\end{lemma}

\begin{proof}
$R \to S$가 에탈 환 사상이고 $\mathfrak q$가 $\mathfrak m$ 위에 놓이며
$\kappa = \kappa(\mathfrak q)$인 $S$의 소 아이디얼인 쌍
$(S, \mathfrak q)$들의 범주를 생각하자. 쌍들의 사상
$(S, \mathfrak q) \to (S', \mathfrak q')$은
$\varphi^{-1}(\mathfrak q') = \mathfrak q$를 만족하는 $R$-대수 사상
$\varphi : S \to S'$로 주어진다. 다음과 두자.
$$
R^h = \colim_{(S, \mathfrak q)} S.
$$
쌍들의 범주가 여과됨을 보이자. 범주론, 정의
\ref{categories-definition-directed}를 보라. 이 범주는 쌍
$(R, \mathfrak m)$을 포함하므로 공집합이 아니며, 이는 범주론, 정의
\ref{categories-definition-directed}의 (1)을 증명한다.
임의의 쌍 $(S, \mathfrak q)$에 대해 합성
$\kappa \to S/\mathfrak q \to \kappa(\mathfrak q)$가 동형사상이므로,
소 아이디얼 $\mathfrak q$는 잉여체가 $\kappa$인 극대 아이디얼이다.
$(S, \mathfrak q)$와 $(S', \mathfrak q')$가 두 대상이라고 하자. 다음과 두자.
$S'' = S \otimes_R S'$ 및 $\mathfrak q'' = \mathfrak qS'' + \mathfrak q'S''$.
위에서 말한 바에 의해
$S''/\mathfrak q'' = S/\mathfrak q \otimes_R S'/\mathfrak q' = \kappa$이다.
또한 보조정리 \ref{algebra-lemma-etale}에 의해 $R \to S''$는 에탈이다.
이는 범주론, 정의 \ref{categories-definition-directed}의 (2)를 증명한다.
다음으로
$\varphi, \psi : (S, \mathfrak q) \to (S', \mathfrak q')$가
쌍들의 두 사상이라고 하자. 그러면 보조정리
\ref{algebra-lemma-map-between-etale}에 의해 $\varphi$, $\psi$, 그리고
$S' \otimes_R S' \to S'$는 에탈 환 사상이다. 소 아이디얼
$$
S'' = (S' \otimes_{\varphi, S, \psi} S')
\otimes_{S' \otimes_R S'} S'
$$
와 함께
$$
\mathfrak q'' =
(\mathfrak q' \otimes S' + S' \otimes \mathfrak q') \otimes S'
+
(S' \otimes_{\varphi, S, \psi} S') \otimes \mathfrak q'
$$
를 생각하자. 위와 같이 논증하면(에탈 사상의 기저변환은 에탈이고,
에탈 사상의 합성은 에탈이다) $S''$가 $R$ 위에서 에탈임을 알 수 있다.
또한 표준 사상 $S' \to S''$(예를 들어 맨 오른쪽 인자를 사용한다)은
$\varphi$와 $\psi$를 균등화한다. 이는 범주론, 정의
\ref{categories-definition-directed}의 (3)을 증명한다.
따라서 $R^h$는 $f \in S$인 삼중항 $(S, \mathfrak q, f)$들로 이루어지며,
그러한 두 삼중항 $(S, \mathfrak q, f)$, $(S', \mathfrak q', f')$가
$R^h$의 같은 원소를 정의할 필요충분조건은 쌍 $(S'', \mathfrak q'')$와
쌍들의 사상 $\varphi : (S, \mathfrak q) \to (S'', \mathfrak q'')$ 및
$\varphi' : (S', \mathfrak q') \to (S'', \mathfrak q'')$ 중
$\varphi(f) = \varphi'(f')$를 만족하는 것들이 존재하는 것이다.

\medskip\noindent
$x \in R^h$라고 하자. $x$를 삼중항 $(S, \mathfrak q, f)$로 나타내자.
$\mathfrak q_1, \ldots, \mathfrak q_r$를 $\mathfrak m$ 위에 놓이는
$S$의 다른 소 아이디얼들이라 하자. 위에서 $\mathfrak q$가 극대임을
보았으므로 $\mathfrak q \not \subset \mathfrak q_i$이다.
따라서 $\mathfrak q$가 소 아이디얼이므로, $g \in S$,
$g \not \in \mathfrak q$이고 $i = 1, \ldots, r$에 대해
$g \in \mathfrak q_i$인 원소를 찾을 수 있다. 쌍들의 사상
$(S, \mathfrak q) \to (S_g, \mathfrak qS_g)$를 생각하자.
이 방법으로 $x$가 항상 삼중항 $(S, \mathfrak q, f)$로 주어진다고
가정할 수 있음을 알 수 있다. 여기서 $\mathfrak q$는 $\mathfrak m$ 위에
놓이는 $S$의 유일한 소 아이디얼, 즉
$\sqrt{\mathfrak mS} = \mathfrak q$이다. 그런데 $R \to S$가 에탈이므로
보조정리 \ref{algebra-lemma-etale-at-prime}에 의해
$\mathfrak mS_{\mathfrak q} = \mathfrak qS_{\mathfrak q}$이다.
따라서 실제로 $\mathfrak mS = \mathfrak q$를 얻는다.

\medskip\noindent
$x \not \in \mathfrak mR^h$라고 하자. $x$를
$\mathfrak mS = \mathfrak q$인 삼중항 $(S, \mathfrak q, f)$로 나타내자.
그러면 $f \not \in \mathfrak mS$, 즉 $f \not \in \mathfrak q$이다.
따라서 $(S, \mathfrak q) \to (S_f, \mathfrak qS_f)$는 $f$의 상이
가역이 되는 쌍들의 사상이다. 그러므로 $x$는 가역이고, 그 역원은
삼중항 $(S_f, \mathfrak qS_f, 1/f)$로 나타난다. 따라서 $R^h$는
극대 아이디얼이 $\mathfrak mR^h$인 국소환이다. 잉여체는 $\kappa$이다.
실제로 삼중항 $(S, \mathfrak q, f)$를 $\mathfrak q$를 법으로 한
$f$의 잉여류로 보내어 $R^h/\mathfrak mR^h \to \kappa$를 정의할 수 있다.

\medskip\noindent
$R^h$가 헨젤임을 보이는 일만 남았다. 즉, $P \in R^h[T]$가
최고차항계수가 1인 다항식이고 $a_0 \in \kappa$가 축소
$\overline{P} \in \kappa[T]$의 단순근이라고 하자.
그러면 $P$가 최고차항계수가 1인 다항식 $Q \in S[T]$의 상이 되는 쌍
$(S, \mathfrak q)$를 찾을 수 있다. $S' = S[T]/(Q)$로 두고,
$a' \in S$가 $a_0$의 임의의 올림일 때
$\mathfrak q' = \mathfrak qS' + (T - a')S'$인 극대 아이디얼
$\mathfrak q' \subset S'$를 택하자. 구성에 의해 $S \to S'$는
$\mathfrak q'$에서 에탈이고 $\kappa = \kappa(\mathfrak q')$이다.
$g \in S'$, $g \not \in \mathfrak q'$이고 $S'' = S'_g$가 $S$ 위에서
에탈이 되는 것을 택하자. 그러면
$(S, \mathfrak q) \to (S'', \mathfrak q'S'')$는 쌍들의 사상이다.
이제 삼중항 $(S'', \mathfrak q'S'', \text{class of }T)$는
$P(a) = 0$이고 $\overline{a} = a_0$인 원소 $a \in R^h$를 정한다.
% P01-E0969: “Now that triple” in the frozen English source should read “Now the triple”.
이는 원하는 바이다.
\end{proof}

\begin{lemma}
\label{algebra-lemma-strict-henselization}
$(R, \mathfrak m, \kappa)$를 국소환이라 하고,
$\kappa \subset \kappa^{sep}$를 분리 대수적 폐포라 하자. 다음 가환 도식이 존재한다.
$$
\xymatrix{
\kappa \ar[r] & \kappa \ar[r] & \kappa^{sep} \\
R \ar[r] \ar[u] & R^h \ar[r] \ar[u] & R^{sh} \ar[u]
}
$$
이 도식은 다음 성질을 갖는다.
\begin{enumerate}
\item 사상 $R^h \to R^{sh}$는 국소적이다.
\item $R^{sh}$는 엄밀 헨젤이다.
\item $R^{sh}$는 에탈 $R$-대수들의 여과 여극한이다.
\item $\mathfrak m R^{sh}$는 $R^{sh}$의 극대 아이디얼이다.
\item $\kappa^{sep} = R^{sh}/\mathfrak m R^{sh}$이다.
\end{enumerate}
\end{lemma}

\begin{proof}
보조정리 \ref{algebra-lemma-henselization}에 사용한 것과 정확히 같은 증명으로
보인다. 유일한 차이는 쌍 대신 삼중항 $(S, \mathfrak q, \alpha)$를
사용한다는 것이다. 여기서 $R \to S$는 에탈이고, $\mathfrak q$는
$\mathfrak m$ 위에 놓이는 $S$의 소 아이디얼이며,
$\alpha : \kappa(\mathfrak q) \to \kappa^{sep}$는 $\kappa$의 확대들의 매장이다.
\end{proof}

\begin{definition}
\label{algebra-definition-henselization}
$(R, \mathfrak m, \kappa)$를 국소환이라 하자.
\begin{enumerate}
\item 보조정리 \ref{algebra-lemma-henselization}에서 구성한 국소환 사상
$R \to R^h$를 $R$의 {\it 헨젤화}라고 한다.
\item 분리 대수적 폐포 $\kappa \subset \kappa^{sep}$가 주어졌을 때,
보조정리 \ref{algebra-lemma-strict-henselization}에서 구성한 국소환 사상
$R \to R^{sh}$를
{\it $\kappa \subset \kappa^{sep}$에 관한 $R$의 엄밀 헨젤화}라고 한다.
\item 국소환 사상 $R \to R^{sh}$가 보조정리
\ref{algebra-lemma-strict-henselization}에서 구성한 국소환 사상 중 하나와
동형이면 이를 $R$의 {\it 엄밀 헨젤화}라고 한다
% P01-E0970: the frozen third definition item lacks terminal punctuation.
\end{enumerate}
\end{definition}

\noindent
사상 $R \to R^h \to R^{sh}$는 평탄 국소환 준동형사상이다.
보조정리 \ref{algebra-lemma-uniqueness-henselian}에 의해 $R$-대수 $R^h$와
$R^{sh}$는 잉여체가 각각 $\kappa$와 $\kappa^{sep}$인 에탈 $R$-대수들의
여과 여극한인 헨젤 국소환이라는 조건에 의해 유일한 동형사상까지 잘 정의된다.
이 절의 나머지에서는 주로 (엄밀) 헨젤화의 함자성을 논의한다.
$R$과 그 헨젤화 사이의 관계에 관한 더 섬세한 결과는
더 많은 대수, 절 \ref{more-algebra-section-permanence-henselization}에서 논의한다.

\begin{remark}
\label{algebra-remark-construct-sh-from-h}
$R^{sh}$를 $R^h$로부터 구성할 수도 있다. 실제로 임의의 유한 분리 중간 확대
$\kappa^{sep}/\kappa'/\kappa$에 대해, 주어진 잉여체 확대를 재현하는
유일한(유일한 동형사상까지) 유한 에탈 국소환 확대
$R^h \subset R^h(\kappa')$가 존재한다. 보조정리
\ref{algebra-lemma-henselian-cat-finite-etale}을 보라. 따라서 다음과 둘 수 있다.
$$
R^{sh} =
\bigcup\nolimits_{\kappa \subset \kappa' \subset \kappa^{sep}}
R^h(\kappa')
$$
잉여체 수준의 화살표들과 양립하는 이 계의 화살표들은 보조정리
\ref{algebra-lemma-henselian-cat-finite-etale}에 의해 존재한다.
보조정리 \ref{algebra-lemma-colimit-henselian}에 의해 이는 헨젤 국소환을 만든다.
각 환 $R^h(\kappa')$가 보조정리 \ref{algebra-lemma-finite-over-henselian}에 의해
헨젤이기 때문이다.
구성에 의해 $R^h \to R^{sh}$가 유도하는 잉여체 확대는 체확대
$\kappa^{sep}/\kappa$이다. 따라서 이렇게 구성한 $R^{sh}$는 엄밀 헨젤이다.
보조정리 \ref{algebra-lemma-composition-colimit-etale}에 의해 $R$-대수 $R^{sh}$는
에탈 $R$-대수들의 여극한이다. 그러므로 보조정리
\ref{algebra-lemma-uniqueness-henselian}의 유일성에 의해 $R^{sh}$는 엄밀 헨젤화이다.
\end{remark}

\begin{lemma}
\label{algebra-lemma-henselian-functorial-prepare}
$R \to S$를 국소환들의 국소 사상이라 하자. $S \to S^h$를 헨젤화라 하자.
$R \to A$를 에탈 환 사상이라 하고, $\mathfrak q$를 $\mathfrak m_R$ 위에 놓이며
$R/\mathfrak m_R \cong \kappa(\mathfrak q)$를 만족하는 $A$의 소 아이디얼이라 하자.
그러면 다음 가환 도식에 들어맞는 유일한 환 사상 $f : A \to S^h$가 존재한다.
$$
\xymatrix{
A \ar[r]_f & S^h \\
R \ar[u] \ar[r] & S \ar[u]
}
$$
또한 $f^{-1}(\mathfrak m_{S^h}) = \mathfrak q$이다.
\end{lemma}

\begin{proof}
이는 보조정리 \ref{algebra-lemma-map-into-henselian}의 특수한 경우이다.
\end{proof}

\begin{lemma}
\label{algebra-lemma-henselian-functorial}
$R \to S$를 국소환들의 국소 사상이라 하자.
$R \to R^h$와 $S \to S^h$를 헨젤화들이라 하자.
그러면 다음 가환 도식에 들어맞는 유일한 국소환 사상 $R^h \to S^h$가 존재한다.
$$
\xymatrix{
R^h \ar[r]_f & S^h \\
R \ar[u] \ar[r] & S \ar[u]
}
$$
\end{lemma}

\begin{proof}
보조정리 \ref{algebra-lemma-map-into-henselian-colimit}에서 바로 따른다.
\end{proof}

\noindent
다음은 헨젤화의 조금 다른 구성이다.

\begin{lemma}
\label{algebra-lemma-henselization-different}
$R$을 환이라 하고 $\mathfrak p \subset R$를 소 아이디얼이라 하자.
$R \to S$가 에탈이고 $\mathfrak q$가 $\mathfrak p$ 위에 놓이며
$\kappa(\mathfrak p) = \kappa(\mathfrak q)$를 만족하는 소 아이디얼인
쌍 $(S, \mathfrak q)$들의 범주를 생각하자. 이 범주는 여과되며 표준적으로
$$
(R_{\mathfrak p})^h = \colim_{(S, \mathfrak q)} S
= \colim_{(S, \mathfrak q)} S_{\mathfrak q}
$$
이다.
\end{lemma}

\begin{proof}
쌍들의 사상 $(S, \mathfrak q) \to (S', \mathfrak q')$은
$\varphi^{-1}(\mathfrak q') = \mathfrak q$를 만족하는 $R$-대수 사상
$\varphi : S \to S'$로 주어진다. 쌍들의 범주가 여과됨을 보이자.
범주론, 정의 \ref{categories-definition-directed}를 보라.
이 범주는 쌍 $(R, \mathfrak p)$를 포함하므로 공집합이 아니며,
이는 범주론, 정의 \ref{categories-definition-directed}의 (1)을 증명한다.
$(S, \mathfrak q)$와 $(S', \mathfrak q')$가 두 쌍이라고 하자.
$\mathfrak q$, 각각 $\mathfrak q'$는 잉여체가 $\kappa(\mathfrak p)$인
섬유환 $S \otimes \kappa(\mathfrak p)$, 각각
$S' \otimes \kappa(\mathfrak p)$의 소 아이디얼에 대응한다.
따라서 이들은 $S \otimes \kappa(\mathfrak p)$, 각각
$S' \otimes \kappa(\mathfrak p)$의 극대 아이디얼에 대응한다.
$S'' = S \otimes_R S'$로 두자. 위 사실에 의해 $\mathfrak q$와
$\mathfrak q'$ 위에 놓이고 잉여체가 $\kappa(\mathfrak p)$인 유일한
소 아이디얼 $\mathfrak q'' \subset S''$가 존재한다.
보조정리 \ref{algebra-lemma-etale}에 의해 환 사상 $R \to S''$는 에탈이다.
이는 범주론, 정의 \ref{categories-definition-directed}의 (2)를 증명한다.
다음으로
$\varphi, \psi : (S, \mathfrak q) \to (S', \mathfrak q')$가
쌍들의 두 사상이라고 하자. 보조정리 \ref{algebra-lemma-map-between-etale}에 의해
$\varphi$, $\psi$, 그리고 $S' \otimes_R S' \to S'$는 에탈 환 사상이다.
다음을 생각하자.
$$
S'' = (S' \otimes_{\varphi, S, \psi} S')
\otimes_{S' \otimes_R S'} S'
$$
위와 같이 논증하면(에탈 사상의 기저변환은 에탈이고 에탈 사상의 합성은
에탈이다) $S''$가 $R$ 위에서 에탈임을 알 수 있다. $\mathfrak p$ 위의
$S''$의 섬유환은
$$
F'' = (F' \otimes_{\varphi, F, \psi} F')
\otimes_{F' \otimes_{\kappa(\mathfrak p)} F'} F'
$$
이다. 여기서 $F', F$는 $S'$와 $S$의 섬유환들이다. $\varphi$와 $\psi$가
쌍들의 사상이므로, $\mathfrak p'$에 대응하는 사상
$F' \to \kappa(\mathfrak p)$는 $F'' \to \kappa(\mathfrak p)$인 사상으로
확장되고, 이는 다시 잉여체가 $\kappa(\mathfrak p)$인 소 아이디얼
$\mathfrak q'' \subset S''$에 대응한다.
% P01-E0971: the frozen proof uses the undefined prime p' here; the pair supplies q'.
표준 사상 $S' \to S''$(예를 들어 맨 오른쪽 인자를 사용한다)은
$\varphi$와 $\psi$를 균등화하는 쌍들의 사상
$(S', \mathfrak q') \to (S'', \mathfrak q'')$이다.
이는 범주론, 정의 \ref{categories-definition-directed}의 (3)을 증명한다.
따라서 이 범주는 여과된다.

\medskip\noindent
보조정리 \ref{algebra-lemma-henselization}의 증명에서 우리는 대응하는 여극한으로
$(R_{\mathfrak p})^h$를 구성했지만, $R_{\mathfrak p}$와 그 극대 아이디얼
$\mathfrak pR_{\mathfrak p}$에서 출발했음을 상기하자. 이제
$(R, \mathfrak p)$의 임의의 쌍 $(S, \mathfrak q)$로부터
$(R_{\mathfrak p}, \mathfrak pR_{\mathfrak p})$의 쌍
$(S_{\mathfrak p}, \mathfrak qS_{\mathfrak p})$를 얻는다. 또한 이 상황에서
$$
S_{\mathfrak p} = \colim_{f \in R, f \not \in \mathfrak p} S_f.
$$
따라서 보조정리의 등식들을 보이려면
$(R_{\mathfrak p}, \mathfrak pR_{\mathfrak p})$의 임의의 쌍
$(S_{loc}, \mathfrak q_{loc})$가 어떤 $(R, \mathfrak p)$ 위의 쌍
$(S, \mathfrak q)$에 대해 $(S_{\mathfrak p}, \mathfrak qS_{\mathfrak p})$의
꼴임을 보이면 충분하다(일부 세부사항은 생략한다).
이는 보조정리 \ref{algebra-lemma-etale}에서 따른다.
\end{proof}

\begin{lemma}
\label{algebra-lemma-henselian-functorial-improve}
$R \to S$를 환 사상이라 하고, $\mathfrak q \subset S$를
$\mathfrak p \subset R$ 위에 놓이는 소 아이디얼이라 하자.
$R \to R^h$와 $S \to S^h$를 $R_\mathfrak p$와 $S_\mathfrak q$의
헨젤화들이라 하자. 보조정리 \ref{algebra-lemma-henselian-functorial}의 국소환 사상
$R^h \to S^h$는 $S^h$를 $\mathfrak m^h$와 $\mathfrak q$ 위에 놓이는
유일한 소 아이디얼에서의 $R^h \otimes_R S$의 헨젤화와 동일시한다.
\end{lemma}

\begin{proof}
보조정리 \ref{algebra-lemma-henselization-different}에 의해 $R^h$, 각각 $S^h$는
에탈 $R$-대수들, 각각 에탈 $S$-대수들의 여과 여극한이다.
따라서 $R^h \otimes_R S$는 에탈 $S$-대수 $A_i$들의 여과 여극한이다
(보조정리 \ref{algebra-lemma-etale}). 보조정리 \ref{algebra-lemma-colimits-of-etale}에 의해
$S^h$는 에탈 $R^h \otimes_R S$-대수들의 여과 여극한이다.
또한 $S^h$는 잉여체가 $\kappa(\mathfrak q)$인 헨젤 국소환이므로,
명제는 보조정리 \ref{algebra-lemma-uniqueness-henselian}의 유일성 결과에서 따른다.
\end{proof}

\begin{lemma}
\label{algebra-lemma-strictly-henselian-functorial-prepare}
$\varphi : R \to S$를 국소환들의 국소 사상이라 하자.
$S/\mathfrak m_S \subset \kappa^{sep}$를 분리 대수적 폐포라 하자.
$S \to S^{sh}$를 $S$의 $S/\mathfrak m_S \subset \kappa^{sep}$에 관한
엄밀 헨젤화라 하자. $R \to A$를 에탈 환 사상이라 하고,
$\mathfrak q$를 $\mathfrak m_R$ 위에 놓이는 $A$의 소 아이디얼이라 하자.
임의의 가환 도식
$$
\xymatrix{
\kappa(\mathfrak q) \ar[r]_{\phi} & \kappa^{sep} \\
R/\mathfrak m_R \ar[r]^{\varphi} \ar[u] & S/\mathfrak m_S \ar[u]
}
$$
이 주어졌다고 하자. 그러면 다음 가환 도식에 들어맞는 유일한 환 사상
$f : A \to S^{sh}$가 존재한다.
$$
\xymatrix{
A \ar[r]_f & S^{sh} \\
R \ar[u] \ar[r]^{\varphi} & S \ar[u]
}
$$
또한 $f^{-1}(\mathfrak m_{S^h}) = \mathfrak q$이고, 유도된 사상
$\kappa(\mathfrak q) \to \kappa^{sep}$는 주어진 사상이다.
% P01-E0972: the frozen inverse-image condition uses m_{S^h}, although f has codomain S^{sh}.
\end{lemma}

\begin{proof}
이는 보조정리 \ref{algebra-lemma-map-into-henselian}의 특수한 경우이다.
\end{proof}

\begin{lemma}
\label{algebra-lemma-strictly-henselian-functorial}
$R \to S$를 국소환들의 국소 사상이라 하자. 분리 대수적 폐포
$R/\mathfrak m_R \subset \kappa_1^{sep}$와
$S/\mathfrak m_S \subset \kappa_2^{sep}$를 택하자.
$R \to R^{sh}$와 $S \to S^{sh}$를 대응하는 엄밀 헨젤화들이라 하자.
임의의 가환 도식
$$
\xymatrix{
\kappa_1^{sep} \ar[r]_{\phi} & \kappa_2^{sep} \\
R/\mathfrak m_R \ar[r]^{\varphi} \ar[u] & S/\mathfrak m_S \ar[u]
}
$$
이 주어지면, 다음 가환 도식에 들어맞고 $R^{sh}$와 $S^{sh}$의 잉여체 위에서
$\phi$를 유도하는 유일한 국소환 사상 $R^{sh} \to S^{sh}$가 존재한다.
% P01-E0973: the frozen English source splits the “Given” clause from its existence predicate.
$$
\xymatrix{
R^{sh} \ar[r]_f & S^{sh} \\
R \ar[u] \ar[r] & S \ar[u]
}
$$
\end{lemma}

\begin{proof}
보조정리 \ref{algebra-lemma-map-into-henselian-colimit}에서 바로 따른다.
\end{proof}

\begin{lemma}
\label{algebra-lemma-strict-henselization-different}
$R$을 환이라 하고 $\mathfrak p \subset R$를 소 아이디얼이라 하자.
$\kappa(\mathfrak p) \subset \kappa^{sep}$를 분리 대수적 폐포라 하자.
$R \to S$가 에탈이고, $\mathfrak q$가 $\mathfrak p$ 위에 놓이는
소 아이디얼이며, $\phi : \kappa(\mathfrak q) \to \kappa^{sep}$가
$\kappa(\mathfrak p)$-대수 사상인 삼중항 $(S, \mathfrak q, \phi)$들의
범주를 생각하자. 이 범주는 여과되며 표준적으로
$$
(R_{\mathfrak p})^{sh} =
\colim_{(S, \mathfrak q, \phi)} S =
\colim_{(S, \mathfrak q, \phi)} S_{\mathfrak q}
$$
이다.
\end{lemma}

\begin{proof}
삼중항들의 사상
$(S, \mathfrak q, \phi) \to (S', \mathfrak q', \phi')$은
$\varphi^{-1}(\mathfrak q') = \mathfrak q$이고
$\phi' \circ \varphi = \phi$인 $R$-대수 사상 $\varphi : S \to S'$로
주어진다. 쌍들의 범주가 여과됨을 보이자. 범주론, 정의
\ref{categories-definition-directed}를 보라.
% P01-E0974: the frozen proof says “category of pairs”, but its objects are triples.
이 범주는 삼중항 $(R, \mathfrak p, \kappa(\mathfrak p) \subset \kappa^{sep})$을
포함하므로 공집합이 아니며, 이는 범주론, 정의
\ref{categories-definition-directed}의 (1)을 증명한다.
$(S, \mathfrak q, \phi)$와 $(S', \mathfrak q', \phi')$가 두 삼중항이라고 하자.
$\mathfrak q$, 각각 $\mathfrak q'$는 $\kappa(\mathfrak p)$ 위에서
유한 분리적인 잉여체를 갖는 섬유환 $S \otimes \kappa(\mathfrak p)$,
각각 $S' \otimes \kappa(\mathfrak p)$의 소 아이디얼에 대응하고,
$\phi$, 각각 $\phi'$는 $\kappa^{sep}$로 가는 사상에 대응한다.
따라서 이 자료는 $\kappa(\mathfrak p)$-대수 사상들
$$
\phi : S \otimes_R \kappa(\mathfrak p) \longrightarrow \kappa^{sep},
\quad
\phi' : S' \otimes_R \kappa(\mathfrak p) \longrightarrow \kappa^{sep}.
$$
에 대응한다. $S'' = S \otimes_R S'$로 두자. 위의 사상들을 결합하여
유일한 $\kappa(\mathfrak p)$-대수 사상
$$
\phi'' = \phi \otimes \phi' :
S'' \otimes_R \kappa(\mathfrak p)
\longrightarrow
\kappa^{sep}
$$
을 얻는다.
% P01-E0975: the frozen phrase “the maps the above” has its article and modifier transposed.
그 핵은 $\mathfrak q$와 $\mathfrak q'$ 위에 놓이는 소 아이디얼
$\mathfrak q'' \subset S''$에 대응하고, 그 잉여체는 $\kappa^{sep}$ 안의
$\phi(\kappa(\mathfrak q))$와 $\phi'(\kappa(\mathfrak q'))$의 합성체로
$\phi''$를 통해 사상된다. 보조정리 \ref{algebra-lemma-etale}에 의해
환 사상 $R \to S''$는 에탈이다. 따라서 $(S'', \mathfrak q'', \phi'')$는
$(S, \mathfrak q, \phi)$와 $(S', \mathfrak q', \phi')$ 둘 다를 지배하는
삼중항이다. 이는 범주론, 정의 \ref{categories-definition-directed}의 (2)를 증명한다.
다음으로
$\varphi, \psi : (S, \mathfrak q, \phi) \to (S', \mathfrak q', \phi')$가
쌍들의 두 사상이라고 하자. 그러면 보조정리
\ref{algebra-lemma-map-between-etale}에 의해 $\varphi$, $\psi$, 그리고
$S' \otimes_R S' \to S'$는 에탈 환 사상이다.
% P01-E0976: the frozen proof again calls morphisms between triples “morphisms of pairs”.
다음을 생각하자.
$$
S'' = (S' \otimes_{\varphi, S, \psi} S')
\otimes_{S' \otimes_R S'} S'
$$
위와 같이 논증하면(에탈 사상의 기저변환은 에탈이고 에탈 사상의 합성은
에탈이다) $S''$가 $R$ 위에서 에탈임을 알 수 있다. $\mathfrak p$ 위의
$S''$의 섬유환은
$$
F'' = (F' \otimes_{\varphi, F, \psi} F')
\otimes_{F' \otimes_{\kappa(\mathfrak p)} F'} F'
$$
이다. 여기서 $F', F$는 $S'$와 $S$의 섬유환들이다. $\varphi$와 $\psi$가
삼중항들의 사상이므로, 사상 $\phi' : F' \to \kappa^{sep}$는
$\phi'' : F'' \to \kappa^{sep}$인 사상으로 확장되고, 이는 다시
소 아이디얼 $\mathfrak q'' \subset S''$에 대응한다.
표준 사상 $S' \to S''$(예를 들어 맨 오른쪽 인자를 사용한다)은
$\varphi$와 $\psi$를 균등화하는 삼중항들의 사상
$(S', \mathfrak q', \phi') \to (S'', \mathfrak q'', \phi'')$이다.
이는 범주론, 정의 \ref{categories-definition-directed}의 (3)을 증명한다.
따라서 이 범주는 여과된다.

\medskip\noindent
이 계의 여극한 $R_{colim}$이 $\kappa^{sep}$에 관한
$R_{\mathfrak p}$의 엄밀 헨젤화와 같음을 보이는 일만 남았다.
이를 위해 삼중항 $(S, \mathfrak q, \phi)$들의 계가 보조정리
\ref{algebra-lemma-henselization-different}의 쌍 $(S, \mathfrak q)$들을 부분계로
포함함을 주의하자. 따라서 그 보조정리의 결과에 의해 $R_{colim}$은
$R_{\mathfrak p}^h$를 포함한다. 또한 $R_{\mathfrak p}^h \subset R_{colim}$이
에탈 환 확대들의 유향 여극한임은 명백하다. 보조정리
\ref{algebra-lemma-finite-over-henselian}와 \ref{algebra-lemma-colimit-henselian}에 의해
$R_{colim}$은 헨젤이다. 마지막으로 보조정리
\ref{algebra-lemma-make-etale-map-prescribed-residue-field}에 의해 $R_{colim}$의
잉여체가 $\kappa^{sep}$와 같음을 알 수 있다. 따라서 $R_{colim}$은
엄밀 헨젤이고, 원하는 대로 $R_{\mathfrak p}$의 엄밀 헨젤화와 같다.
일부 세부사항은 생략한다.
\end{proof}

\begin{lemma}
\label{algebra-lemma-strictly-henselian-functorial-improve}
$R \to S$를 환 사상이라 하고, $\mathfrak q \subset S$를
$\mathfrak p \subset R$ 위에 놓이는 소 아이디얼이라 하자.
분리 대수적 폐포 $\kappa(\mathfrak p) \subset \kappa_1^{sep}$와
$\kappa(\mathfrak q) \subset \kappa_2^{sep}$를 택하자.
$R^{sh}$와 $S^{sh}$를 각각 $R_\mathfrak p$와 $S_\mathfrak q$의
대응하는 엄밀 헨젤화들이라 하자. 임의의 가환 도식
$$
\xymatrix{
\kappa_1^{sep} \ar[r]_{\phi} & \kappa_2^{sep} \\
\kappa(\mathfrak p) \ar[r]^{\varphi} \ar[u] & \kappa(\mathfrak q) \ar[u]
}
$$
이 주어지면, 보조정리 \ref{algebra-lemma-strictly-henselian-functorial}의 국소환 사상
$R^{sh} \to S^{sh}$는 $S^{sh}$를 $\mathfrak q$와 극대 아이디얼
$\mathfrak m^{sh} \subset R^{sh}$ 위에 놓이는 한 소 아이디얼에서의
$R^{sh} \otimes_R S$의 엄밀 헨젤화와 동일시한다.
% P01-E1465: the frozen English “Given” clause before this predicate lacks the joining comma and continues as a capitalized new sentence.
\end{lemma}

\begin{proof}
증명은 보조정리 \ref{algebra-lemma-henselian-functorial-improve}의 증명과 동일하되,
보조정리 \ref{algebra-lemma-strict-henselization-different}를
보조정리 \ref{algebra-lemma-henselization-different} 대신 사용한다.
\end{proof}

\begin{lemma}
\label{algebra-lemma-sh-from-h-map}
$R \to S$를 환 사상이라 하자. $\mathfrak q \subset S$를
$\mathfrak p \subset R$ 위에 놓이며
$\kappa(\mathfrak p) \to \kappa(\mathfrak q)$가 동형사상인
소 아이디얼이라 하자. $\kappa(\mathfrak p) = \kappa(\mathfrak q)$의
분리 대수적 폐포 $\kappa^{sep}$를 택하자. 그러면
$$
(S_\mathfrak q)^{sh} =
(S_\mathfrak q)^h \otimes_{(R_\mathfrak p)^h} (R_\mathfrak p)^{sh}
$$
이다.
\end{lemma}

\begin{proof}
이는 주의 \ref{algebra-remark-construct-sh-from-h}에 있는 국소환의 엄밀 헨젤화에 대한
대안적 구성과 잉여체들이 같다는 사실에서 따른다. 일부 세부사항은 생략한다.
\end{proof}




\section{헨젤화와 준유한 환 사상}
\label{algebra-section-henselization-quasi-finite}

\noindent
이 절에서는 준유한 환 사상에 대한 (엄밀) 헨젤화 사이의 함자적 사상에
관한 몇 가지 결과를 증명한다.

\begin{lemma}
\label{algebra-lemma-quasi-finite-henselization}
$R \to S$를 환 사상이라 하자. $\mathfrak q$를 $R$의 $\mathfrak p$ 위에
놓이는 $S$의 소 아이디얼이라 하자. $R \to S$가 $\mathfrak q$에서
준유한이라고 가정하자. 보조정리 \ref{algebra-lemma-henselian-functorial}의 가환 도식
$$
\xymatrix{
R_{\mathfrak p}^h \ar[r] & S_{\mathfrak q}^h \\
R_{\mathfrak p} \ar[u] \ar[r] & S_{\mathfrak q} \ar[u]
}
$$
은 $S_{\mathfrak q}^h$를 $\mathfrak q$가 생성하는 소 아이디얼에서의
$R_{\mathfrak p}^h \otimes_{R_{\mathfrak p}} S_{\mathfrak q}$의
국소화와 동일시한다. 또한 환 사상
$R_{\mathfrak p}^h \to S_{\mathfrak q}^h$는 유한이다.
\end{lemma}

\begin{proof}
$R_{\mathfrak p}^h \otimes_R S$는 $\mathfrak q$에 대응하는 소 아이디얼에서
$R_{\mathfrak p}^h$ 위에 준유한임을 주의하자. 보조정리
\ref{algebra-lemma-four-rings}를 보라. 따라서
$R_{\mathfrak p}^h \otimes_{R_{\mathfrak p}} S_{\mathfrak q}$의 국소화
$S'$는 $R_{\mathfrak p}^h$ 위에서 헨젤이고 유한이다. 보조정리
\ref{algebra-lemma-finite-over-henselian}을 보라. 국소화로서 $S'$는 에탈
$R_{\mathfrak p}^h \otimes_{R_{\mathfrak p}} S_{\mathfrak q}$-대수들의
여과 여극한이다. 보조정리 \ref{algebra-lemma-henselian-functorial-improve}에 의해
$S_\mathfrak q^h$는
$R_{\mathfrak p}^h \otimes_{R_{\mathfrak p}} S_{\mathfrak q}$의 헨젤화이다.
따라서 보조정리 \ref{algebra-lemma-uniqueness-henselian}의 유일성 결과에 의해
$S' = S_\mathfrak q^h$이다.
\end{proof}

\begin{lemma}
\label{algebra-lemma-quotient-henselization}
\begin{slogan}
헨젤화는 몫과 양립한다.
\end{slogan}
$R$을 헨젤화가 $R^h$인 국소환이라 하고 $I \subset \mathfrak m_R$라 하자.
그러면 $R^h/IR^h$는 $R/I$의 헨젤화이다.
\end{lemma}

\begin{proof}
이는 보조정리 \ref{algebra-lemma-quasi-finite-henselization}의 특수한 경우이다.
\end{proof}

\begin{lemma}
\label{algebra-lemma-quasi-finite-strict-henselization}
$R \to S$를 환 사상이라 하자. $\mathfrak q$를 $R$의 $\mathfrak p$ 위에
놓이는 $S$의 소 아이디얼이라 하자. $R \to S$가 $\mathfrak q$에서
준유한이라고 가정하자. $\kappa_2^{sep}/\kappa(\mathfrak q)$를
분리 대수적 폐포라 하고, $\kappa_1^{sep} \subset \kappa_2^{sep}$를
$\kappa(\mathfrak p)$ 위에서 분리 대수적인 원소들의 부분체라 하자
(체론, 보조정리 \ref{fields-lemma-separable-first}). 보조정리
\ref{algebra-lemma-strictly-henselian-functorial}의 가환 도식
$$
\xymatrix{
R_{\mathfrak p}^{sh} \ar[r] & S_{\mathfrak q}^{sh} \\
R_{\mathfrak p} \ar[u] \ar[r] & S_{\mathfrak q} \ar[u]
}
$$
은 $S_{\mathfrak q}^{sh}$를 다음 사상의 핵인 소 아이디얼에서의
$R_{\mathfrak p}^{sh} \otimes_{R_{\mathfrak p}} S_{\mathfrak q}$의
국소화와 동일시한다.
$$
R_{\mathfrak p}^{sh} \otimes_{R_{\mathfrak p}} S_{\mathfrak q}
\longrightarrow
\kappa_1^{sep} \otimes_{\kappa(\mathfrak p)} \kappa(\mathfrak q)
\longrightarrow
\kappa_2^{sep}
$$
또한 환 사상 $R_{\mathfrak p}^{sh} \to S_{\mathfrak q}^{sh}$는
잉여체 확대가 유한이면서 순수 비분리인 확대
$\kappa_2^{sep}/\kappa_1^{sep}$인 국소환들의 유한 국소 준동형사상이다.
\end{lemma}

\begin{proof}
$R \to S$가 $\mathfrak q$에서 준유한이므로 확대
$\kappa(\mathfrak q)/\kappa(\mathfrak p)$는 유한이다. 정의
\ref{algebra-definition-quasi-finite}와 보조정리
\ref{algebra-lemma-isolated-point-fibre}를 보라. 따라서 $\kappa_1^{sep}$는
$\kappa(\mathfrak p)$의 분리 대수적 폐포이다(작은 세부사항은 생략한다).
특히 보조정리 \ref{algebra-lemma-strictly-henselian-functorial}은 실제로 적용된다.
다음으로 $\kappa_2^{sep}$ 안의 $\kappa(\mathfrak q)$와
$\kappa_1^{sep}$의 합성체는 분리 대수적으로 닫혀 있으므로
$\kappa_2^{sep}$와 같다. 따라서 $\kappa_2^{sep}/\kappa_1^{sep}$는
유한이다. 구성에 의해 확대 $\kappa_2^{sep}/\kappa_1^{sep}$는
순수 비분리이다. 환 사상 $R_{\mathfrak p}^{sh} \to S_{\mathfrak q}^{sh}$는
실제로 국소적이고 실제로 유한 순수 비분리인 잉여체 확대
$\kappa_2^{sep}/\kappa_1^{sep}$를 유도한다.

\medskip\noindent
$R_{\mathfrak p}^{sh} \otimes_R S$는 보조정리의 명제에 주어진
소 아이디얼 $\mathfrak q'$에서 $R_{\mathfrak p}^{sh}$ 위에 준유한임을
주의하자. 보조정리 \ref{algebra-lemma-four-rings}를 보라. 따라서
$\mathfrak q'$에서의
$R_{\mathfrak p}^{sh} \otimes_{R_{\mathfrak p}} S_{\mathfrak q}$의 국소화
$S'$는 $R_{\mathfrak p}^{sh}$ 위에서 헨젤이고 유한이다. 보조정리
\ref{algebra-lemma-finite-over-henselian}을 보라. 앞 문단의 논의에 의해 사상
$\kappa_1^{sep} \otimes_{\kappa(\mathfrak p)} \kappa(\mathfrak q)
\to \kappa_2^{sep}$가 전사이므로 $S'$의 잉여체는 $\kappa_2^{sep}$이다.
또한 국소화로서 $S'$는 에탈
$R_{\mathfrak p}^{sh} \otimes_{R_{\mathfrak p}} S_{\mathfrak q}$-대수들의
여과 여극한이다. 보조정리
\ref{algebra-lemma-strictly-henselian-functorial-improve}에 의해
$S_{\mathfrak q}^{sh}$는 $\mathfrak q'$에서의
$R_{\mathfrak p}^{sh} \otimes_{R_{\mathfrak p}} S_{\mathfrak q}$의
엄밀 헨젤화이다. 따라서 보조정리 \ref{algebra-lemma-uniqueness-henselian}의
유일성 결과에 의해 $S' = S_\mathfrak q^{sh}$이다.
\end{proof}

\begin{lemma}
\label{algebra-lemma-quotient-strict-henselization}
$R$을 엄밀 헨젤화가 $R^{sh}$인 국소환이라 하고,
$I \subset \mathfrak m_R$라 하자. 그러면 $R^{sh}/IR^{sh}$는
$R/I$의 엄밀 헨젤화이다.
\end{lemma}

\begin{proof}
이는 보조정리 \ref{algebra-lemma-quasi-finite-strict-henselization}의 특수한 경우이다.
\end{proof}

\begin{lemma}
\label{algebra-lemma-local-tensor-with-integral}
$A \to B$와 $A \to C$를 국소환들의 국소 준동형사상이라 하자.
$A \to C$가 정수적이고
$\kappa(\mathfrak m_C)/\kappa(\mathfrak m_A)$ 또는
$\kappa(\mathfrak m_B)/\kappa(\mathfrak m_A)$ 중 하나가 순수 비분리이면,
$D = B \otimes_A C$는 국소환이고 $B \to D$와 $C \to D$는 국소적이다.
\end{lemma}

\begin{proof}
$B \to D$가 정수적인 환 사상이므로 상승 성질에 의해 $D$의 모든
극대 아이디얼은 $B$의 극대 아이디얼 위에 놓인다(보조정리
\ref{algebra-lemma-integral-going-up}). 이제
$D/\mathfrak m_B D = \kappa(\mathfrak m_B) \otimes_A C =
\kappa(\mathfrak m_B) \otimes_{\kappa(\mathfrak m_A)} C/\mathfrak m_A C$이다.
$C/\mathfrak m_A C$의 스펙트럼은 단 하나의 점, 즉 $\mathfrak m_C$로
이루어진다. 따라서 $D/\mathfrak m_B D$의 스펙트럼은
$\kappa(\mathfrak m_B) \otimes_{\kappa(\mathfrak m_A)} \kappa(\mathfrak m_C)$의
스펙트럼과 같다. 그런데 $\kappa(\mathfrak m_C)/\kappa(\mathfrak m_A)$ 또는
$\kappa(\mathfrak m_B)/\kappa(\mathfrak m_A)$ 중 하나가 순수 비분리라는
가정에 의해 후자는 한 점이다. 이는 $D$가 국소환이고 환 사상
$B \to D$와 $C \to D$가 국소적임을 증명한다.
\end{proof}

\begin{lemma}
\label{algebra-lemma-base-change-strict-henselization-quasi-finite}
$A \to B$와 $A \to C$를 환 사상이라 하자. $\kappa$를 분리 대수적으로
닫힌 체라 하고 $B \otimes_A C \to \kappa$를 환 준동형사상이라 하자.
대응하는 엄밀 헨젤화 사상들을 다음과 나타내자(증명을 보라).
$$
\xymatrix{
B^{sh} \ar[r] & (B \otimes_A C)^{sh} \\
A^{sh} \ar[u] \ar[r] & C^{sh} \ar[u]
}
$$
다음 중 하나가 성립한다고 하자.
\begin{enumerate}
\item $A \to B$가 소 아이디얼 $\mathfrak p_B = \Ker(B \to \kappa)$에서
준유한이다.
\item $B$는 준유한 $A$-대수들의 여과 여극한이다.
\item $B_{\mathfrak p_B}$는 $A_{\mathfrak p_A}$ 위의 준유한 대수들의
여과 여극한이다.
\item $B$는 $A$ 위에서 정수적이다.
\end{enumerate}
그러면 $B^{sh} \otimes_{A^{sh}} C^{sh} \to (B \otimes_A C)^{sh}$는
동형사상이다.
\end{lemma}

\begin{proof}
$D = B \otimes_A C$로 쓰자. $\mathfrak p_A = \Ker(A \to \kappa)$로 두고,
$\mathfrak p_B$, $\mathfrak p_C$, $\mathfrak p_D$도 마찬가지로 나타내자.
$\kappa_A \subset \kappa$를 $\kappa$ 안에서의
$\kappa(\mathfrak p_A)$의 분리 대수적 폐포라 하고,
$\kappa_B$, $\kappa_C$, $\kappa_D$도 마찬가지로 나타내자.
$A^{sh}$를 분리 대수적 폐포 $\kappa_A/\kappa(\mathfrak p_A)$를 사용하여
구성한 $A_{\mathfrak p_A}$의 엄밀 헨젤화라 하자. $B^{sh}$, $C^{sh}$,
$D^{sh}$도 마찬가지이다. 보조정리
\ref{algebra-lemma-strictly-henselian-functorial}의 함자성에서 보조정리의 가환
도식을 얻는다.

\medskip\noindent
그 가환 도식에서 얻는 사상
$$
c : B^{sh} \otimes_{A^{sh}} C^{sh} \to D^{sh} = (B \otimes_A C)^{sh}
$$
을 생각하자. $A \to B$가 $\mathfrak p_B = \Ker(B \to \kappa)$에서
준유한이면, 보조정리 \ref{algebra-lemma-four-rings}에 의해 환 사상 $C \to D$는
$\mathfrak p_D$에서 준유한이다. 따라서 보조정리
\ref{algebra-lemma-quasi-finite-strict-henselization}(그리고 보조정리
\ref{algebra-lemma-base-change-integral})에 의해 환 사상 $c$는 유한
$C^{sh}$-대수들의 준동형사상이고, 어떤 소 아이디얼
$\mathfrak q$와 $\mathfrak r$에 대해
$$
B^{sh} = (B \otimes_A A^{sh})_{\mathfrak q}
\quad\text{and}\quad
D^{sh} = (D \otimes_C C^{sh})_{\mathfrak r} =
(B \otimes_A C^{sh})_{\mathfrak r}
$$
이다. 또한
$$
B^{sh} \otimes_{A^{sh}} C^{sh} =
(B \otimes_A A^{sh})_{\mathfrak q} \otimes_{A^{sh}} C^{sh} =
\text{a localization of }
B \otimes_A C^{sh}
$$
이므로 $c$의 정의역과 공역은 모두 $B \otimes_A C^{sh}$의 국소화이다
(사상과 양립한다). 따라서 $B^{sh} \otimes_{A^{sh}} C^{sh}$가
국소적임을 보이면 충분하다(작은 세부사항은 생략한다).
이는 보조정리 \ref{algebra-lemma-local-tensor-with-integral}과, 이미 사용한 보조정리
\ref{algebra-lemma-quasi-finite-strict-henselization}에 의해
$A^{sh} \to B^{sh}$가 순수 비분리 잉여체 확대를 갖는 유한 사상이라는
사실에서 따른다. 이는 보조정리의 (1)을 증명한다.

\medskip\noindent
(2)의 경우 $B = \colim B_i$를 준유한 $A$-대수들의 여과 여극한으로 쓰자.
이에 따라 $D_i = B_i \otimes_A C$인 $D = \colim D_i$를 얻는다.
$B^{sh} = \colim B_i^{sh}$임을 관찰하자. 실제로 보조정리
\ref{algebra-lemma-colimit-henselian}에 의해 환 $\colim B_i^{sh}$는 엄밀 헨젤
국소환이다. 또한 보조정리 \ref{algebra-lemma-colimit-colimit-etale-better}에 의해
$\colim B_i^{sh}$는 에탈 $B$-대수들의 여과 여극한이다. 마지막으로
$\colim B_i^{sh}$의 잉여체는 $\kappa(\mathfrak p_B)$의 분리 대수적
폐포이다(세부사항은 생략한다). 그러므로 정의
\ref{algebra-definition-henselization} 뒤의 논의를 보아
$B^{sh} = \colim B_i^{sh}$라고 결론낸다. 마찬가지로
$D^{sh} = \colim D_i^{sh}$이다. 이제 (1)에 의해
$$
D^{sh} = \colim D_i^{sh} = \colim B_i^{sh} \otimes_{A^{sh}} C^{sh} =
B^{sh} \otimes_{A^{sh}} C^{sh}
$$
이므로 결론을 얻는다. 이는 여과 여극한들이 텐서곱과 가환하기 때문이다.
% P01-E0977: the frozen English source has singular “filtered colimit commute”.

\medskip\noindent
(3)의 경우 $A$, $B$, $C$를 각각 $\mathfrak p_A$, $\mathfrak p_B$,
$\mathfrak p_C$에서의 국소화로 바꾸어도 좋다. 따라서 (3)은 (2)에서 따른다.

\medskip\noindent
정수적 환 사상은 유한 환 사상들의 여과 여극한이므로 (4)도 (2)에서 따른다.
\end{proof}





















\section{정규성에 대한 세르의 판정법}
\label{algebra-section-serre-criterion}

\noindent
뇌터 환에 관한 다음 성질들을 도입한다.

\begin{definition}
\label{algebra-definition-conditions}
$R$을 뇌터 환이라 하자.
$k \geq 0$를 정수라 하자.
\begin{enumerate}
\item 높이가 $\leq k$인 모든 소 아이디얼 $\mathfrak p$에 대해 국소환
$R_{\mathfrak p}$가 정칙이면 $R$이 성질 {\it $(R_k)$}를 갖는다고 한다.
또한 $R$이 {\it 여차원 $\leq k$에서 정칙}이라고 한다.
\item 모든 소 아이디얼 $\mathfrak p$에 대해 국소환 $R_{\mathfrak p}$의
깊이가 적어도 $\min\{k, \dim(R_{\mathfrak p})\}$이면 $R$이 성질
{\it $(S_k)$}를 갖는다고 한다.
\item $M$을 유한 $R$-가군이라 하자. 모든 소 아이디얼 $\mathfrak p$에 대해
가군 $M_{\mathfrak p}$의 깊이가 적어도
$\min\{k, \dim(\text{Supp}(M_{\mathfrak p}))\}$이면 $M$이 성질
$(S_k)$를 갖는다고 한다.
\end{enumerate}
\end{definition}

\noindent
모든 뇌터 환은 성질 $(S_0)$를 가지며, 그 위의 모든 유한 가군도 그러하다.
영가군의 깊이를 $\infty$로 둔다는 우리의 규약(절 \ref{algebra-section-depth} 참조)과
공집합의 차원을 $-\infty$로 둔다는 규약(위상수학, 절
\ref{topology-section-krull-dimension} 참조)에 의해 영가군은 모든 $k$에 대해
성질 $(S_k)$를 갖는다.

\begin{lemma}
\label{algebra-lemma-criterion-no-embedded-primes}
$R$을 뇌터 환이라 하자.
$M$을 유한 $R$-가군이라 하자.
다음은 동치이다.
\begin{enumerate}
\item $M$에는 매립 연관 소 아이디얼이 없다.
\item $M$은 성질 $(S_1)$을 갖는다.
\end{enumerate}
\end{lemma}

\begin{proof}
$\mathfrak p$를 $M$의 매립 연관 소 아이디얼이라 하자. 그러면
$\mathfrak p \supset \mathfrak q$인 $M$의 다른 연관 소 아이디얼
$\mathfrak q$가 존재한다. 특히 $\mathfrak q$도 지지집합에 속하므로
$\dim(\text{Supp}(M_{\mathfrak p})) \geq 1$이다. 한편
$\mathfrak pR_{\mathfrak p}$는 $M_{\mathfrak p}$에 연관되어 있고
(보조정리 \ref{algebra-lemma-associated-primes-localize}), 따라서
$\text{depth}(M_{\mathfrak p}) = 0$이다
(보조정리 \ref{algebra-lemma-ideal-nonzerodivisor} 참조).
즉 $(S_1)$은 성립하지 않는다. 역으로 $(S_1)$이 성립하지 않으면
$\dim(\text{Supp}(M_{\mathfrak p})) \geq 1$이고
$\text{depth}(M_{\mathfrak p}) = 0$인 소 아이디얼 $\mathfrak p$가 존재한다.
$\text{depth}(M_{\mathfrak p}) = 0$이므로 두 보조정리
\ref{algebra-lemma-associated-primes-localize}와 \ref{algebra-lemma-ideal-nonzerodivisor}에 의해
$\mathfrak p \in \text{Ass}(M)$이다.
$\dim(\text{Supp}(M_{\mathfrak p})) \geq 1$이므로
$\mathfrak q \subset \mathfrak p$, $\mathfrak q \not = \mathfrak p$인
소 아이디얼 $\mathfrak q \in \text{Supp}(M)$가 존재한다.
그러한 $\mathfrak q$를 $\text{Supp}(M)$에서 극소가 되도록 택할 수 있다.
그러면 명제 \ref{algebra-proposition-minimal-primes-associated-primes}에 의해
$\mathfrak q \in \text{Ass}(M)$이고, 따라서 $\mathfrak p$는 매립 연관
소 아이디얼이다.
\end{proof}

\begin{lemma}
\label{algebra-lemma-criterion-reduced}
\begin{slogan}
축소는 R0와 S1의 합이다.
\end{slogan}
$R$을 뇌터 환이라 하자.
다음은 동치이다.
\begin{enumerate}
\item $R$은 축소환이다.
\item $R$은 성질 $(R_0)$와 $(S_1)$을 갖는다.
\end{enumerate}
\end{lemma}

\begin{proof}
$R$이 축소환이라고 가정하자. 보조정리
\ref{algebra-lemma-minimal-prime-reduced-ring}에 의해 $R$의 모든 극소 소 아이디얼
$\mathfrak p$에 대해 $R_{\mathfrak p}$는 체이다. 따라서 $(R_0)$를 얻는다.
$\mathfrak p$를 높이가 $\geq 1$인 소 아이디얼이라 하자. 그러면
$A = R_{\mathfrak p}$는 차원이 $\geq 1$인 축소 국소환이다. 극대 아이디얼을
$\mathfrak m$이라 하자. 만일 이것이 연관 소 아이디얼이면 소멸자가
$\mathfrak m$인 $x \in \mathfrak m$이 존재하여 $x^2 = 0$이 되므로,
그 극대 아이디얼은 연관 소 아이디얼이 아니다. 따라서 보조정리
\ref{algebra-lemma-ass-zero-divisors}에 의해 $A = R_{\mathfrak p}$의 깊이는 적어도
하나이다. 이는 $(S_1)$이 성립함을 보인다.

\medskip\noindent
역으로 $R$이 $(R_0)$와 $(S_1)$을 만족한다고 가정하자.
$\mathfrak p$가 $R$의 극소 소 아이디얼이면 $(R_0)$에 의해
$R_{\mathfrak p}$는 체이고, 따라서 축소환이다. $\mathfrak p$가 극소가
아니면 $(S_1)$에 의해 $R_{\mathfrak p}$의 깊이는 $\geq 1$이다. 따라서
$t \in \mathfrak pR_{\mathfrak p}$이고
$R_{\mathfrak p} \to R_{\mathfrak p}[1/t]$가 단사인 원소가 존재한다.
이제 $R_\mathfrak p[1/t]$는 그 소 아이디얼들에서의 국소화들의 곱에 포함된다.
보조정리 \ref{algebra-lemma-characterize-zero-local}을 보라. 따라서
$R_{\mathfrak p}$는 $\mathfrak p \supset \mathfrak q$이고
$t \not \in \mathfrak q$인 소 아이디얼에서의 $R$의 국소화들의 곱의
부분환이다. 이 소 아이디얼들은 높이가 더 작으므로 높이에
관한 귀납법에 의해 $R$은 축소환이다.
\end{proof}

\begin{lemma}[정규성에 대한 세르의 판정법]
\label{algebra-lemma-criterion-normal}
\begin{reference}
\cite[IV, Theorem 5.8.6]{EGA}
\end{reference}
\begin{slogan}
정규는 R1과 S2의 합이다.
\end{slogan}
$R$을 뇌터 환이라 하자.
다음은 동치이다.
\begin{enumerate}
\item $R$은 정규환이다.
\item $R$은 성질 $(R_1)$과 $(S_2)$를 갖는다.
\end{enumerate}
\end{lemma}

\begin{proof}
(1) $\Rightarrow$ (2)의 증명. $R$이 정규라고, 즉 모든 소 아이디얼에서의
국소화 $R_{\mathfrak p}$가 정규 정역이라고 가정하자.
특히 보조정리 \ref{algebra-lemma-criterion-reduced}에 의해 $R$은 $(R_0)$와
$(S_1)$을 갖는다. 따라서 차원이 $d$인 국소 뇌터 정규 정역 $R$의 깊이가
$\geq \min(2, d)$이고, $d = 1$이면 정칙임을 보이면 충분하다.
$d = 1$인 주장은 보조정리 \ref{algebra-lemma-characterize-dvr}에서 따른다.

\medskip\noindent
$R$을 극대 아이디얼 $\mathfrak m$과 차원 $d \geq 2$를 갖는 국소 뇌터
정규 정역이라 하자. $R$에 보조정리 \ref{algebra-lemma-hart-serre-loc-thm}을
적용한다. $R$은 이 보조정리의 경우 (1)이나 (2)에 해당하지 않음이 분명하다.
보조정리의 (4)에서와 같은 $R \to R'$을 택하자. $R$이 정역이므로
$R \subset R'$이다. $\mathfrak m$은 $R'$의 연관 소 아이디얼이 아니므로
$R'$ 위의 영인자가 아닌 $x \in \mathfrak m$이 존재한다. 그러면
$R_x = R'_x$이므로 $R$과 $R'$은 같은 분수체를 갖는 정역이다. 그런데
$R \subset R'$의 유한성은 $R'$의 모든 원소가 $R$ 위에서 정수적임을
함의한다(보조정리 \ref{algebra-lemma-finite-is-integral}). $R$이 정규이므로
$R = R'$라고 결론낸다. 즉 (4)는 일어나지 않는다. 따라서 남은 가능성 (3),
곧 원하는 $\text{depth}(R) \geq 2$를 얻는다.

\medskip\noindent
(2) $\Rightarrow$ (1)의 증명. $R$이 $(R_1)$과 $(S_2)$를 만족한다고 하자.
보조정리 \ref{algebra-lemma-criterion-reduced}에 의해 $R$은 축소환이다. 따라서
차원이 $d$이고 $(S_2)$와 $(R_1)$을 만족하는 축소 국소 뇌터 환 $R$을
생각하자. $R$이 정규 정역임을 보이면 충분하다. $d = 0$이면 결과는 분명하다.
$d = 1$이면 결과는 보조정리 \ref{algebra-lemma-characterize-dvr}에서 따른다.

\medskip\noindent
$R$을 극대 아이디얼 $\mathfrak m$과 차원 $d \geq 2$를 가지며 $(R_1)$과
$(S_2)$를 만족하는 축소 국소 뇌터 환이라 하자. 보조정리
\ref{algebra-lemma-characterize-reduced-ring-normal}에 의해 $R$이 그 전분수환
$Q(R)$에서 정수적으로 닫혀 있음을 보이면 충분하다. $R$ 위에서 정수적인
$x \in Q(R)$를 택하자. 그러면 $R' = R[x]$는 $R$의 유한 환 확대이다
(보조정리 \ref{algebra-lemma-characterize-finite-in-terms-of-integral}). 모든 비극대
소 아이디얼 $\mathfrak p \subset R$에 대해
$\dim(R_\mathfrak p) < d$이므로 귀납법에 의해
$R_\mathfrak p = R'_\mathfrak p$이다. 따라서 $R'/R$의 지지집합은
$\{\mathfrak m\}$이다. 그러므로 $R'/R$은 $\mathfrak m$의 어떤 거듭제곱에
의해 소멸된다(보조정리 \ref{algebra-lemma-Noetherian-power-ideal-kills-module}).
보조정리 \ref{algebra-lemma-hart-serre-loc-thm}에 의해 이는 $R$의 깊이가
$\geq 2 = \min(2, d)$라는 가정과 모순이고, 증명이 끝난다.
\end{proof}

\begin{lemma}
\label{algebra-lemma-regular-normal}
정칙환은 정규이다.
\end{lemma}

\begin{proof}
$R$을 정칙환이라 하자. 보조정리 \ref{algebra-lemma-criterion-normal}에 의해
$R$이 $(R_1)$과 $(S_2)$를 만족함을 보이면 충분하다. 보조정리
\ref{algebra-lemma-regular-ring-CM}에서 보듯 정칙 국소환은 Cohen--Macaulay이므로
$R$이 $(S_2)$임은 분명하다. 성질 $(R_1)$은 즉시 따른다.
\end{proof}

\begin{lemma}
\label{algebra-lemma-normal-domain-intersection-localizations-height-1}
$R$을 분수체가 $K$인 뇌터 정규 정역이라 하자. 그러면
\begin{enumerate}
\item 영이 아닌 임의의 $a \in R$에 대해 몫 $R/aR$에는 매립 소
아이디얼이 없고, 그 연관 소 아이디얼들은 모두 높이가 $1$이다
% P01-E0978
\item
$$
R = \bigcap\nolimits_{\text{height}(\mathfrak p) = 1} R_{\mathfrak p}
$$
\item 영이 아닌 임의의 $x \in K$에 대해 몫 $R/(R \cap xR)$에는 매립 소
아이디얼이 없고, 그 연관 소 아이디얼들은 모두 높이가 $1$이다.
% P01-E0979
\end{enumerate}
\end{lemma}

\begin{proof}
보조정리 \ref{algebra-lemma-criterion-normal}에 의해 $R$은 $(S_2)$를 갖는다.
따라서 영이 아닌 임의의 원소 $a \in R$에 대해 $R/aR$은 $(S_1)$을 갖는다
(예를 들어 보조정리 \ref{algebra-lemma-depth-in-ses}을 사용하라).
% P01-E0980
그러므로 $R/aR$에는 매립 소 아이디얼이 없다
(보조정리 \ref{algebra-lemma-criterion-no-embedded-primes}). $R/aR$의 연관 소
아이디얼들은 정확히 $(a)$ 위의 극소 소 아이디얼 $\mathfrak p$들이며,
$a$가 영이 아니므로 이들은 높이가 $1$이다(보조정리
\ref{algebra-lemma-minimal-over-1}). 이는 (1)을 증명한다.

\medskip\noindent
따라서 $b \in R$가 주어졌을 때, $b \in aR$인 것과 $(a)$ 위의 모든
극소 소 아이디얼 $\mathfrak p$에 대해 $b \in aR_{\mathfrak p}$인 것은
동치이다(보조정리 \ref{algebra-lemma-zero-at-ass-zero} 참조). 위에서 보았듯
이 소 아이디얼들은 모두 높이가 $1$이므로, $b/a \in R$인 것과 높이 1인
모든 소 아이디얼에 대해 $b/a \in R_{\mathfrak p}$인 것은 동치이다.
따라서 (2)가 성립한다.

\medskip\noindent
(3)에 대해 $x = a/b$라고 쓰자. $\mathfrak p_1, \ldots, \mathfrak p_r$을
$(ab)$ 위의 극소 소 아이디얼들이라 하자. 위에 의해 이들은 모두 높이 1이다.
그러면 보조정리의 (2)에 의해
$R \cap xR = \bigcap_{i = 1, \ldots, r} (R \cap xR_{\mathfrak p_i})$이다.
따라서 $R/(R \cap xR)$은
$\bigoplus R/(R \cap xR_{\mathfrak p_i})$의 부분가군이다.
$R$은 뇌터 정규 정역에 대해 성질 $(R_1)$을 만족하므로 각
$R_{\mathfrak p_i}$는 이산 값매김환이다(보조정리
\ref{algebra-lemma-criterion-normal} 참조). 따라서 어떤 $e_i \in \mathbf{Z}$에 대해
$xR_{\mathfrak p_i} = \mathfrak p_i^{e_i}R_{\mathfrak p_i}$이다. 그러므로
그 직합은 다음과 같다.
$\bigoplus_{e_i > 0} R/\mathfrak p_i^{(e_i)}$; 정의
\ref{algebra-definition-symbolic-power}를 보라. 보조정리
\ref{algebra-lemma-symbolic-power-associated}에 의해 가군
$R/\mathfrak p^{(n)}$의 유일한 연관 소 아이디얼은 $\mathfrak p$이다.
따라서 $R/(R \cap xR)$의 연관 소 아이디얼 집합은
$\{\mathfrak p_i\}$의 부분집합이고, 이들 사이에는 포함관계가 없다.
% P01-E0981
이는 (3)을 증명한다.
\end{proof}















\section{체의 형식적 매끄러움}
\label{algebra-section-p-bases}

\noindent
이 절에서는 체 확대가 형식적으로 매끄러운 것과 분리인 것이 동치임을
보인다. 그러나 먼저 유한 생성 체 확대가 분리 대수적인 것과 형식적으로
비분기인 것이 동치임을 증명한다.

\begin{lemma}
\label{algebra-lemma-characterize-separable-algebraic-field-extensions}
$K/k$를 유한 생성 체 확대라 하자.
다음은 동치이다.
\begin{enumerate}
\item $K$는 $k$의 유한 분리 체 확대이다.
\item $\Omega_{K/k} = 0$이다.
\item $K$는 $k$ 위에서 형식적으로 비분기이다.
\item $K$는 $k$ 위에서 비분기이다.
\item $K$는 $k$ 위에서 형식적 에탈이다.
\item $K$는 $k$ 위에서 에탈이다.
\end{enumerate}
\end{lemma}

\begin{proof}
(2)와 (3)의 동치는 보조정리
\ref{algebra-lemma-characterize-formally-unramified}이다. 보조정리
\ref{algebra-lemma-etale-over-field}에 의해 (1)과 (6)은 동치이다.
성질 (6)은 (5)와 (4)를 함의하고, 이들은 각각 (3)을 함의한다
(보조정리 \ref{algebra-lemma-formally-etale-etale}, \ref{algebra-lemma-unramified},
\ref{algebra-lemma-formally-unramified-unramified}). 따라서 (2)가 (1)을 함의함을
보이면 충분하다. $K$가 정역 $A$의 분수체가 되도록 유한 생성
$k$-부분대수 $A \subset K$를 택하자. $S = A \setminus \{0\}$로 두자.
$0 = \Omega_{K/k} = S^{-1}\Omega_{A/k}$이고
(보조정리 \ref{algebra-lemma-differentials-localize}), $\Omega_{A/k}$가 유한 생성이므로
(보조정리 \ref{algebra-lemma-differentials-finitely-generated}), $A$를 국소화
$A_f$로 바꾸어 $\Omega_{A/k} = 0$인 경우로 환원할 수 있다
(세부사항은 생략한다). 그러면 $A$는 $k$ 위에서 비분기이고, 따라서 예컨대
$\mathfrak q = (0)$으로 보조정리 \ref{algebra-lemma-unramified-at-prime}을 적용하면
$K/k$는 유한 분리이다.
\end{proof}

\begin{lemma}
\label{algebra-lemma-derivative-zero-pth-power}
$k$를 특성이 $p > 0$인 완전체라 하자. $K/k$를 확대라 하자.
$a \in K$라 하자. 그러면 $\Omega_{K/k}$에서 $\text{d}a = 0$인 것과
$a$가 $p$제곱인 것은 동치이다.
\end{lemma}

\begin{proof}
보조정리 \ref{algebra-lemma-colimit-differentials}에 의해 $L/k$가 유한 생성 체 확대이고
$\Omega_{L/k}$에서 $\text{d}a$가 영이 되도록 하는 부분체
$k \subset L \subset K$가 존재한다. 따라서 $K$가 $k$의 유한 생성 체
확대인 경우를 가정해도 좋다.

\medskip\noindent
$K$가 $k(x_1, \ldots, x_r)$ 위에서 유한 분리가 되도록 초월기저
$x_1, \ldots, x_r \in K$를 택하자. 이는 정의에 의해 가능하다. 정의
\ref{algebra-definition-perfect}와 \ref{algebra-definition-separable-field-extension}을 보라.
결과는 순수 초월 부분체 $k(x_1, \ldots, x_r) \subset K$에 대해 성립한다.
실제로
$$
\Omega_{k(x_1, \ldots, x_r)/k} =
\bigoplus\nolimits_{i = 1}^r k(x_1, \ldots, x_r) \text{d}x_i
$$
이고, 모든 편도함수가 영인 임의의 유리함수는 $p$제곱이다. 또한
$k(x_1, \ldots, x_r) \subset K$가 유한 분리이므로(계산은 생략한다)
$$
\Omega_{K/k} =
\bigoplus\nolimits_{i = 1}^r K\text{d}x_i
$$
이다. $a \in K$가 미분 가군에서 $\text{d}a = 0$을 만족하는 원소라고 하자.
$x_i$의 선택에 의해 $a$의 최소다항식
$P(T) \in k(x_1, \ldots, x_r)[T]$는 분리이다. 다음과 같이 쓰자.
$$
P(T) = T^d + \sum\nolimits_{i = 1}^d a_i T^{d - i}
$$
그러면 $\Omega_{K/k}$에서
$$
0 = \text{d}P(a) = \sum\nolimits_{i = 1}^d a^{d - i}\text{d}a_i
$$
이다. 위의 $\Omega_{K/k}$의 기술과 $P$가 $a$의 최소다항식이라는 사실에
의해 이는 $\text{d}a_i = 0$을 함의한다. 따라서 각 $i$에 대해
$a_i = b_i^p$이다. 그러므로 체론, 보조정리
\ref{fields-lemma-pth-root}에 의해 $a$는 $p$제곱이다.
\end{proof}

\begin{lemma}
\label{algebra-lemma-size-extension-pth-roots}
$k$를 특성이 $p > 0$인 체라 하자. $a_1, \ldots, a_n \in k$를
$\Omega_{k/\mathbf{F}_p}$에서 $\text{d}a_1, \ldots, \text{d}a_n$이
일차독립이 되도록 하는 원소들이라 하자. 그러면 체 확대
$k(a_1^{1/p}, \ldots, a_n^{1/p})$의 $k$ 위 차수는 $p^n$이다.
\end{lemma}

\begin{proof}
$n$에 대한 귀납법을 쓴다. $n = 1$이면 결과는 보조정리
\ref{algebra-lemma-derivative-zero-pth-power}이다. 귀납 단계에서는
$k(a_1^{1/p}, \ldots, a_{n - 1}^{1/p})$의 $k$ 위 차수가
$p^{n - 1}$이라고 가정하자. $a_n$이
$k(a_1^{1/p}, \ldots, a_{n - 1}^{1/p})$에서 $p$제곱으로 가지 않음을
보여야 한다. 만일 그렇게 간다면 다음과 같이 쓸 수 있다.
\begin{align*}
a_n & =
\left(\sum\nolimits_{I = (i_1, \ldots, i_{n - 1}),\ 0 \leq i_j \leq p - 1}
\lambda_I a_1^{i_1/p} \ldots a_{n - 1}^{i_{n - 1}/p}\right)^p \\
& = \sum\nolimits_{I = (i_1, \ldots, i_{n - 1}),\ 0 \leq i_j \leq p - 1}
\lambda_I^p a_1^{i_1} \ldots a_{n - 1}^{i_{n - 1}}
\end{align*}
$\text{d}$를 적용하면 $\text{d}a_n$은 $\text{d}a_i$, $i < n$에
일차종속임을 알 수 있다. 이는 모순이다.
\end{proof}

\begin{lemma}
\label{algebra-lemma-separable-differentials}
$k$를 특성이 $p > 0$인 체라 하자.
% P01-E1466
다음은 동치이다.
\begin{enumerate}
\item 체 확대 $K/k$는 분리이다
(정의 \ref{algebra-definition-separable-field-extension} 참조).
\item 사상
$K \otimes_k \Omega_{k/\mathbf{F}_p} \to \Omega_{K/\mathbf{F}_p}$는
단사이다.
\end{enumerate}
\end{lemma}

\begin{proof}
$K$를 유한 생성 체 확대 $K_i/k$들의 유향 여극한
$K = \colim_i K_i$로 쓰자. 정의에 의해 $K$가 분리인 것과 각 $K_i$가
$k$ 위에서 분리인 것은 동치이고, 보조정리
\ref{algebra-lemma-colimit-differentials}에 의해
$K \otimes_k \Omega_{k/\mathbf{F}_p} \to \Omega_{K/\mathbf{F}_p}$가
단사인 것과 각각의
$K_i \otimes_k \Omega_{k/\mathbf{F}_p} \to \Omega_{K_i/\mathbf{F}_p}$가
단사인 것은 동치이다. 따라서 $K/k$가 유한 생성 체 확대라고 가정해도 좋다.

\medskip\noindent
$K/k$가 분리인 유한 생성 체 확대라고 가정하자. 보조정리
\ref{algebra-lemma-generating-finitely-generated-separable-field-extensions}에서와 같은
$x_1, \ldots, x_{r + 1} \in K$를 택하자. 이 경우
$G(x_1, \ldots, x_{r + 1}) = 0$이고 $\partial G/\partial X_{r + 1}$이
항등적으로 영이 아닌 기약다항식
$G(X_1, \ldots, X_{r + 1}) \in k[X_1, \ldots, X_{r + 1}]$이 존재한다.
또한 $K$는 다음 정역의 분수체이다.
% P01-E0982
$S = K[X_1, \ldots, X_{r + 1}]/(G)$.
다음과 같이 쓰자.
$$
G = \sum a_I X^I, \quad X^I = X_1^{i_1}\ldots X_{r + 1}^{i_{r + 1}}.
$$
위의 $S$의 표시에 의해 다음을 얻는다.
% P01-E0983
$$
\Omega_{S/\mathbf{F}_p}
=
\frac{
S \otimes_k \Omega_k \oplus
\bigoplus\nolimits_{i = 1, \ldots, r + 1} S\text{d}X_i
}{
\langle
\sum X^I \text{d}a_I + \sum \partial G/\partial X_i \text{d}X_i
\rangle
}
$$
$\Omega_{K/\mathbf{F}_p}$는 $S$-가군 $\Omega_{S/\mathbf{F}_p}$의
국소화이므로(보조정리 \ref{algebra-lemma-differentials-localize} 참조), 다음을 얻는다.
$$
\Omega_{K/\mathbf{F}_p}
=
\frac{
K \otimes_k \Omega_k \oplus
\bigoplus\nolimits_{i = 1, \ldots, r + 1} K\text{d}X_i
}{
\langle
\sum X^I \text{d}a_I + \sum \partial G/\partial X_i \text{d}X_i
\rangle
}
$$
이제 다항식 $\partial G/\partial X_{r + 1}$이 항등적으로 영이 아니므로
사상 $K \otimes_k \Omega_{k/\mathbf{F}_p} \to \Omega_{S/\mathbf{F}_p}$가
원하는 대로 단사라고 결론낸다.
% P01-E0984

\medskip\noindent
$K/k$가 유한 생성 체 확대이고
$K \otimes_k \Omega_{k/\mathbf{F}_p} \to \Omega_{K/\mathbf{F}_p}$가
단사라고 가정하자. (이 증명 부분은 보조정리
\ref{algebra-lemma-characterize-separable-field-extensions}을 증명하는 논증과 같다.)
$x_1, \ldots, x_r$을 $k$ 위의 $K$의 초월기저 가운데 유한 확대
$k(x_1, \ldots, x_r) \subset K$의 비분리 차수가 최소가 되도록 택하자.
$K$가 $k(x_1, \ldots, x_r)$ 위에서 분리이면 끝난다. 그렇지 않다고
가정하여 모순을 얻자. 그러면 $k(x_1, \ldots, x_r)$ 위에서 분리가 아닌
원소 $\alpha \in K$가 존재한다. 그 최소다항식을
$P(T) \in k(x_1, \ldots, x_r)[T]$라 하자. $\alpha$가 분리가 아니므로
$P$는 실제로 $T^p$에 관한 다항식이다. 분모를 제거하면
$$
G(X_1, \ldots, X_r, T) = \sum a_{I, i} X^I T^i \in k[X_1, \ldots, X_r, T]
$$
인 기약다항식을 얻으며, 이는 $K$에서
$G(x_1, \ldots, x_r, \alpha) = 0$을 만족한다. 이는
$k[X_1, \ldots, X_r, T]/(G) \subset K$임을 뜻한다. 어떤 쌍
$(I_0, i_0)$에 대해 계수 $a_{I_0, i_0} = 1$이라고 가정해도 좋다.
적어도 하나의 $i$에 대해 $\text{d}G/\text{d}X_i$가 항등적으로 영이
아님을 주장한다. 그렇지 않다면 $G$는 실제로
$X_1^p, \ldots, X_r^p, T^p$에 관한 다항식이다. 그러면
$$
\sum\nolimits_{(I, i) \not = (I_0, i_0)} x^I\alpha^i \text{d}a_{I, i}
$$
가 $\Omega_{K/\mathbf{F}_p}$에서 영이라는 뜻이다. $K$의 원소들의 집합
$$
\{x^I\alpha^i \mid a_{I, i} \not = 0 \text{ and } (I, i) \not = (I_0, i_0)\}
$$
에는 $k$-일차 관계가 없음을 주의하자. 따라서
$K \otimes_k \Omega_{k/\mathbf{F}_p} \to \Omega_{K/\mathbf{F}_p}$가
단사라는 가정은 모든 $(I, i)$에 대해 $\Omega_{k/\mathbf{F}_p}$에서
$\text{d}a_{I, i} = 0$임을 함의한다. 보조정리
\ref{algebra-lemma-derivative-zero-pth-power}에 의해 각 $a_{I, i}$는 $p$제곱이고,
이는 $G$가 $p$제곱임을 함의하여 $G$의 기약성과 모순된다. 따라서 번호를
다시 매겨 $\text{d}G/\text{d}X_1$이 영이 아니라고 가정해도 좋다. 그러면
$x_1$은 $k(x_2, \ldots, x_r, \alpha)$ 위에서 분리 대수적이고,
$x_2, \ldots, x_r, \alpha$는 $k$ 위의 $K$의 초월기저이다. 이는 유한 확대
$k(x_2, \ldots, x_r, \alpha) \subset K$의 비분리 차수가 유한 확대
$k(x_1, \ldots, x_r) \subset K$의 비분리 차수보다 작다는 뜻이며,
모순이다.
\end{proof}

\begin{lemma}
\label{algebra-lemma-formally-smooth-implies-separable}
$K/k$를 체 확대라 하자. $K$가 $k$ 위에서 형식적으로 매끄러우면
$K$는 $k$의 분리 확대이다.
\end{lemma}

\begin{proof}
$K$가 $k$ 위에서 형식적으로 매끄럽다고 가정하자. 보조정리
\ref{algebra-lemma-ses-formally-smooth}에 의해
$K \otimes_k \Omega_{k/\mathbf{Z}} \to \Omega_{K/\mathbf{Z}}$는 단사이다.
따라서 보조정리 \ref{algebra-lemma-separable-differentials}에 의해 $K$는 $k$ 위에서
분리이다.
% P01-E0985
\end{proof}

\begin{lemma}
\label{algebra-lemma-characterize-formally-smooth-field-extension}
$K/k$를 체 확대라 하자. 그러면 $K$가 $k$ 위에서 형식적으로 매끄러운 것과
$H_1(L_{K/k}) = 0$인 것은 동치이다.
\end{lemma}

\begin{proof}
이는 명제 \ref{algebra-proposition-characterize-formally-smooth}와 벡터공간이
자유이고 따라서 사영이라는 사실에서 따른다.
% P01-E0986
\end{proof}

\begin{lemma}
\label{algebra-lemma-formally-smooth-extensions-easy}
$K/k$를 체 확대라 하자.
\begin{enumerate}
\item $K$가 $k$ 위에서 순수 초월이면 $K$는 $k$ 위에서 형식적으로
매끄럽다.
\item $K$가 $k$ 위에서 분리 대수적이면 $K$는 $k$ 위에서 형식적으로
매끄럽다.
\item $K$가 $k$ 위에서 분리이면 $K$는 $k$ 위에서 형식적으로 매끄럽다.
\end{enumerate}
\end{lemma}

\begin{proof}
(1)에 대해 $K = k(x_j; j \in J)$라고 쓰자. $A$를 $k$-대수라 하고,
$I \subset A$를 제곱이 영인 아이디얼이라 하자. $\varphi : K \to A/I$를
$k$-대수 사상이라 하자. $a_j \mod I = \varphi(x_j)$인 원소
$a_j \in A$를 택하자. 그러면 $x_j$를 $a_j$로 보내고 $I$로 나눈 뒤
$\varphi$로 환원되는 유일한 $k$-대수 사상 $K \to A$가 존재함을 쉽게 알 수
있다. 따라서 $k \subset K$는 형식적으로 매끄럽다.

\medskip\noindent
(2)의 경우 $k \subset K$는 에탈 환 확대들의 여극한이다. 에탈 환 사상은
형식적 에탈이다(보조정리 \ref{algebra-lemma-formally-etale-etale}). 따라서 이 경우는
보조정리 \ref{algebra-lemma-colimit-formally-etale}와 형식적 에탈 환 사상은
형식적으로 매끄럽다는 자명한 관찰에서 따른다.

\medskip\noindent
(3)의 경우 그 유한 생성 $k$-부분확대들의 여과 여극한을
$K = \colim K_i$로 쓰자.
% P01-E0987
정의 \ref{algebra-definition-separable-field-extension}에 의해 각 $K_i$는 $k$의
순수 초월 확대 위에서 분리 대수적이다. 따라서 경우 (1), (2)와 보조정리
\ref{algebra-lemma-compose-formally-smooth}에 의해 $K_i/k$는 형식적으로 매끄럽다.
그러므로 보조정리 \ref{algebra-lemma-characterize-formally-smooth-field-extension}에
의해 $H_1(L_{K_i/k}) = 0$이다. 보조정리
\ref{algebra-lemma-colimits-NL}에 의해 $H_1(L_{K/k}) = 0$이다. 다시 보조정리
\ref{algebra-lemma-characterize-formally-smooth-field-extension}에 의해 $K/k$는
형식적으로 매끄럽다.
\end{proof}

\begin{lemma}
\label{algebra-lemma-fields-are-formally-smooth}
\begin{slogan}
체 확대에서는 형식적으로 매끄러운 것과 분리인 것이 같다.
\end{slogan}
$k$를 체라 하자.
\begin{enumerate}
\item $k$의 특성이 영이면 $k$의 모든 확대체는 $k$ 위에서 형식적으로
매끄럽다.
\item $k$의 특성이 $p > 0$이면 $K/k$가 형식적으로 매끄러운 것과 분리 체
확대인 것은 동치이다.
\end{enumerate}
\end{lemma}

\begin{proof}
보조정리 \ref{algebra-lemma-formally-smooth-implies-separable}와
\ref{algebra-lemma-formally-smooth-extensions-easy}를 결합하라.
\end{proof}

\noindent
여기서는 분리 체 확대에 관한 여러 특징지음을 한데 모은다.

\begin{proposition}
\label{algebra-proposition-characterize-separable-field-extensions}
$K/k$를 체 확대라 하자. $k$의 특성이 영이면
% P01-E0988
\begin{enumerate}
\item $K$는 $k$ 위에서 분리이다.
\item $K$는 $k$ 위에서 기하적으로 축소이다.
\item $K$는 $k$ 위에서 형식적으로 매끄럽다.
\item $H_1(L_{K/k}) = 0$이다.
\item 사상 $K \otimes_k \Omega_{k/\mathbf{Z}} \to \Omega_{K/\mathbf{Z}}$는
단사이다.
\end{enumerate}
$k$의 특성이 $p > 0$이면 다음은 동치이다.
\begin{enumerate}
\item $K$는 $k$ 위에서 분리이다.
\item 환 $K \otimes_k k^{1/p}$는 축소환이다.
\item $K$는 $k$ 위에서 기하적으로 축소이다.
\item 사상 $K \otimes_k \Omega_{k/\mathbf{F}_p} \to \Omega_{K/\mathbf{F}_p}$는
단사이다.
\item $H_1(L_{K/k}) = 0$이다.
\item $K$는 $k$ 위에서 형식적으로 매끄럽다.
\end{enumerate}
\end{proposition}

\begin{proof}
이는 보조정리 \ref{algebra-lemma-characterize-separable-field-extensions},
\ref{algebra-lemma-fields-are-formally-smooth}
% P01-E0989
\ref{algebra-lemma-formally-smooth-implies-separable}, 그리고
\ref{algebra-lemma-separable-differentials}을 결합한 것이다.
\end{proof}

\noindent
유한 생성 분리 체 확대에 관한 또 하나의 특징지음은 다음과 같다.

\begin{lemma}
\label{algebra-lemma-localization-smooth-separable}
$K/k$를 유한 생성 체 확대라 하자. 그러면 $K$가 $k$ 위에서 분리인 것과
$K$가 매끄러운 $k$-대수의 국소화인 것은 동치이다.
\end{lemma}

\begin{proof}
분수체가 $K$인 정역인 유한형 $k$-대수 $R$을 택하자. 보조정리
\ref{algebra-lemma-smooth-at-generic-point}에 따르면 $k \to R$가 $(0)$에서
매끄러운 것과 $K/k$가 분리인 것은 동치이다. 이는 보조정리를 증명한다.
\end{proof}

\begin{lemma}
\label{algebra-lemma-colimit-syntomic}
$K/k$를 체 확대라 하자. 그러면 $K$는 $k$ 위의 전역 완전 교차 대수들의
여과 여극한이다. $K/k$가 분리이면 $K$는 $k$ 위의 매끄러운 대수들의
여과 여극한이다.
\end{lemma}

\begin{proof}
$E \subset K$를 유한 부분집합이라 하자. $E$를 포함하고 $k$ 위의 전역 완전
교차인(각각 매끄러운) $k$-부분대수 $A \subset K$가 존재함을 보이면 충분하다.
분리/매끄러운 경우는 보조정리
\ref{algebra-lemma-localization-smooth-separable}에서 따른다. 일반적으로 $L \subset K$를
$E$가 생성하는 부분체라 하자. $x_1, \ldots, x_d \in L$을 $k$ 위의
초월기저라 택하자. 확대 $L/k(x_1, \ldots, x_d)$는 유한이다. 다음과 같이 쓰자.
$L = k(x_1, \ldots, x_d)[y_1, \ldots, y_r]$.
다항식 $P_i \in k(x_1, \ldots, x_d)[Y_1, \ldots, Y_r]$를 귀납적으로 택하여
$P_i = P_i(Y_1, \ldots, Y_i)$가
$k(x_1, \ldots, x_d)[Y_1, \ldots, Y_{i - 1}]$ 위에서 $Y_i$에 관한
일계수다항식이고, $y_i$의 최소다항식
$k(x_1, \ldots, x_d)[y_1, \ldots, y_{i - 1}][Y_i]$로 가게 하자.
% P01-E0990
그러면 $P_1, \ldots, P_r$은
$k(x_1, \ldots, x_r)[Y_1, \ldots, Y_r]$에서 정칙열이고
$L = k(x_1, \ldots, x_r)[Y_1, \ldots, Y_r]/(P_1, \ldots, P_r)$임이
분명하다.
% P01-E0991
$h \in k[x_1, \ldots, x_d]$가
$P_i \in k[x_1, \ldots, x_d, 1/h, Y_1, \ldots, Y_r]$을 만족하는
다항식이면, $P_1, \ldots, P_r$은
$k[x_1, \ldots, x_d, 1/h, Y_1, \ldots, Y_r]$에서 정칙열이고
$A = k[x_1, \ldots, x_d, 1/h, Y_1, \ldots, Y_r]/(P_1, \ldots, P_r)$는
전역 완전 교차이다. $h$의 선택을 조정하여 $E \subset A$라고 가정할 수 있고,
이로써 끝난다.
\end{proof}






\section{평탄 환 사상의 구성}
\label{algebra-section-constructing-flat}

\noindent
다음 보조정리는 때때로 유용하다.

\begin{lemma}
\label{algebra-lemma-flat-local-given-residue-field}
$(R, \mathfrak m, k)$를 국소환이라 하자. $K/k$를 체 확대라 하자.
국소환 $(R', \mathfrak m', k')$와 평탄 국소 환 사상 $R \to R'$가 존재하여
$\mathfrak m' = \mathfrak mR'$이고, $k'$는 $k$의 확대체로서 $K$와
동형이다.
\end{lemma}

\begin{proof}
$k' = k(\alpha)$가 $k$의 단순 생성 확대라고 가정하자. 그러면 $k'$는
보조정리에서와 같은 평탄 국소 확대 $R \subset R'$의 잉여체이다. 실제로
$\alpha$가 $k$ 위에서 초월적이면 $R'$을 소 아이디얼 $\mathfrak mR[x]$에서의
$R[x]$의 국소화로 둔다. $\alpha$가 최소다항식
$T^d + \sum \overline{\lambda}_iT^{d - i}$를 갖는 대수적 원소이면
$R' = R[T]/(T^d + \sum \lambda_i T^{d - i})$로 둔다.

\medskip\noindent
삼중항 $(k', R \to R', \phi)$들의 모임을 생각하자. 여기서
$k \subset k' \subset K$는 부분체이고, $R \to R'$는 보조정리에서와 같은
국소 환 사상이며, $\phi : R' \to k'$는 $k$-확대들의 동형사상
$R'/\mathfrak mR' \cong k'$를 유도한다. 이들은 ``큰'' 범주
$\mathcal{C}$를 이룬다. 그 사상
$(k_1, R_1, \phi_1) \to (k_2, R_2, \phi_2)$은 다음 도식이 가환하도록 하는
환 사상 $\psi : R_1 \to R_2$로 주어진다.
$$
\xymatrix{
R_1 \ar[d]_\psi \ar[r]_{\phi_1} & k_1 \ar[r] & K \ar@{=}[d] \\
R_2 \ar[r]^{\phi_2} & k_2 \ar[r] & K
}
$$
이는 $k_1 \subset k_2$를 함의한다.

\medskip\noindent
$I$를 유향집합이라 하고 $((R_i, k_i, \phi_i), \psi_{ii'})$를 $I$ 위의
계라 하자. 범주론, 절 \ref{categories-section-posets-limits}를 보라.
% P01-E0992
이 경우 다음을 생각할 수 있다.
$$
R' = \colim_{i \in I} R_i
$$
이는 극대 아이디얼이 $\mathfrak mR'$이고 잉여체가
$k' = \bigcup_{i \in I} k_i$인 국소환이다. 또한 환 사상 $R \to R'$는
평탄 사상들의 여극한이므로 평탄하다(텐서곱은 유향 여극한과 가환한다).
따라서 $(R', k', \phi')$가 이 계의 ``상계''임을 알 수 있다.

\medskip\noindent
$\mathcal{C}$가 집합이었다면 초른 보조정리의 거의 자명한 적용으로 증명이
끝나겠지만, 그렇지 않다.
% P01-E0993
(실제로 등장하는 국소환들의 기수에 합리적인 상계를 찾아 이 논증을
작동시킬 수 있다.) 이 문제를 피하기 위해 $K$ 위에 정순서를 택한다.
$x \in K$에 대해 $K(x)$를 $\leq x$인 $K$의 모든 원소가 생성하는 $K$의
부분체라 하자. $x \in K$에 대한 초한귀납법으로, 잉여체 확대가 $K(x)/k$인
보조정리에서와 같은 환 사상 $R \subset R(x)$를 만들 것이다. 또한 구성에
의해 모든 $y \leq x$에 대해 $R(x)$가 $R(y)$를 포함하게 된다. 실제로
$x$에 바로 앞 원소 $x'$가 있으면 $K(x) = K(x')[x]$이고, 따라서 증명의
% P01-E0994
첫 문단에서 구성한 국소환 확대를 $R(x') \subset R(x)$로 택할 수 있다.
$x$에 바로 앞 원소가 없으면 먼저 증명의 셋째 문단에서와 같이
$R'(x) = \colim_{x' < x} R(x')$로 둔다. $R'(x)$의 잉여체는
$K'(x) = \bigcup_{x' < x} K(x')$이다. $K(x) = K'(x)[x]$이므로 증명의
첫 문단의 구성을 사용하여 $R'(x) \subset R(x)$를 만들 수 있다.
이로써 보조정리의 증명이 끝난다.
\end{proof}

\begin{lemma}
\label{algebra-lemma-colimit-finite-etale-given-residue-field}
$(R, \mathfrak m, k)$를 국소환이라 하자. $k \subset K$가 분리 대수적
확대이면 유향집합 $I$와 국소환들의 유한 에탈 확대들의 계
$R \subset R_i$, $i \in I$가 존재하여 $R' = \colim R_i$의 잉여체가
$K$이다($k$의 확대체로서).
\end{lemma}

\begin{proof}
$R \subset R'$를 보조정리 \ref{algebra-lemma-flat-local-given-residue-field}의 증명에서
구성한 확대라 하자. 구성에 의해
$R' = \colim_{\alpha \in A} R_\alpha$이고, 여기서 $A$는 정순서집합이며
전이사상 $R_\alpha \to R_{\alpha + 1}$은 유한 에탈이고, $\alpha$가
후속 원소가 아니면 $R_\alpha = \colim_{\beta < \alpha} R_\beta$이다.
초한귀납법으로 결과를 증명하겠다.

\medskip\noindent
$R_\alpha$에 대해 결과가 성립한다고, 즉 $R_i$가 $R$ 위에서 유한 에탈인
$R_\alpha = \colim R_i$가 주어졌다고 하자. $R_\alpha \to R_{\alpha + 1}$이
유한 에탈이므로, 어떤 $i$와 유한 에탈 확대 $R_i \to R_{i, 1}$가 존재하여
$R_{\alpha + 1} = R_\alpha \otimes_{R_i} R_{i, 1}$이다. 따라서
$R_{\alpha + 1} = \colim_{i' \geq i} R_{i'} \otimes_{R_i} R_{i, 1}$이고,
결과는 $\alpha + 1$에 대해 성립한다. 이제 $\alpha$가 후속 원소가 아니고
모든 $\beta < \alpha$에 대해 결과가 $R_\beta$에 성립한다고 하자.
모든 유한 부분집합 $E \subset R_\alpha$는 어떤
$\beta < \alpha$에 대한 $R_\beta$에 포함되므로, 가정에 의해 $E$는 유한
에탈 부분확대에 포함된다. 따라서 결과는 $R_\alpha$에 대해서도 성립한다.
% P01-E0995
\end{proof}

\begin{lemma}
\label{algebra-lemma-finite-free-given-residue-field-extension}
$R$을 환이라 하자. $\mathfrak p \subset R$를 소 아이디얼이라 하고,
$L/\kappa(\mathfrak p)$를 유한 체 확대라 하자. 그러면 유한 자유 환 사상
$R \to S$가 존재하여 $\mathfrak q = \mathfrak pS$가 소 아이디얼이고,
$\kappa(\mathfrak q)/\kappa(\mathfrak p)$는 주어진 확대
$L/\kappa(\mathfrak p)$와 동형이다.
\end{lemma}

\begin{proof}
$\kappa(\mathfrak p) \subset L$의 차수에 대한 귀납법을 쓴다.
% P01-E0996
차수가 $1$이면 $R = S$로 택한다. 일반적으로 부분확대
$\kappa(\mathfrak p) \subset L' \subset L$가 존재하면 차수에 대한
귀납법으로 결론을 얻는다. 먼저 $L'/\kappa(\mathfrak p)$에 대응하는
$R \subset S'$를 구성한 다음 $L/L'$에 대응하는 $S' \subset S$를 구성한다.
% P01-E0997
따라서 $L \supset \kappa(\mathfrak p)$가 하나의 원소 $\alpha \in L$로
생성된다고 가정해도 좋다. $\alpha$의 $\kappa(\mathfrak p)$ 위 최소다항식을
$X^d + \sum_{i < d} a_iX^i$라 하자. 따라서
$a_i \in \kappa(\mathfrak p)$이다. 어떤 $f_i, g \in R$와
$g \not \in \mathfrak p$에 대해 $a_i$를 $f_i/g$의 상으로 쓸 수 있다.
$\alpha$를 $g\alpha$로 바꾸고(이에 따라 $a_i$를
$g^{d - i}a_i$로 바꾸어) $a_i$가 어떤 $f_i \in R$의 상이라고 가정할 수 있다.
이제 단순히 $S = R[x]/(x^d + \sum f_ix^i)$로 택한다.
\end{proof}

\begin{lemma}
\label{algebra-lemma-cofinal-system-flat}
$A$를 환이라 하자. $\kappa = \max(|A|, \aleph_0)$로 두자. 그러면 모든
평탄 $A$-대수 $B$는 기수가 $|B'| \leq \kappa$인 평탄 $A$-부분대수
$B' \subset B$들의 여과 여극한이다. ($B$가 $A$ 위에서 충실히 평탄이면
$B'$도 $A$ 위에서 충실히 평탄임에 주의하라.)
\end{lemma}

\begin{proof}
$B$의 기수가 $\leq \kappa$이면 주장은 참이다. $E \subset B$를
$|E| \leq \kappa$인 $A$-부분대수라 하자. $E$가 기수
$|B'| \leq \kappa$인 평탄 $A$-부분대수 $B'$에 포함됨을 보이겠다.
그러면 보조정리가 따른다. 실제로 (a) $B$의 모든 유한 부분집합은 기수가
많아야 $\kappa$인 $A$-부분대수에 포함되고, (b) 기수가 많아야 $\kappa$인
$B$의 임의의 두 $A$-부분대수는 기수가 많아야 $\kappa$인 하나의
$A$-부분대수에 함께 포함된다. 세부사항은 생략한다.

\medskip\noindent
각각 기수가 $\leq \kappa$인 $A$-부분대수들의 열
$$
E = E_0 \subset E_1 \subset E_2 \subset \ldots
$$
을 귀납적으로 구성하고 $B' = \bigcup E_k$가 $A$ 위에서 평탄임을 보여
증명을 끝내겠다.

\medskip\noindent
구성은 다음과 같다. $E_0 = E$로 둔다. $k \geq 0$에 대해 $E_k$가
주어졌다고 하자. $A$의 계수를 갖는 $E_k$의 원소들 사이의 관계들의 집합을
$S_k$라 하자. 즉 $s \in S_k$는 정수 $n \geq 1$과
$a_1, \ldots, a_n \in A$, $e_1, \ldots, e_n \in E_k$로 주어지며,
$E_k$에서 $\sum a_i e_i = 0$을 만족한다. $A \to B$의 평탄성 및 보조정리
\ref{algebra-lemma-flat-eq}에 의해 모든
$s = (n, a_1, \ldots, a_n, e_1, \ldots, e_n) \in S_k$에 대해 다음을
택할 수 있다.
$$
(m_s, b_{s, 1}, \ldots, b_{s, m_s}, a_{s, 11}, \ldots, a_{s, nm_s})
$$
여기서 $m_s \geq 0$는 정수이고, $b_{s, j} \in B$,
$a_{s, ij} \in A$이며,
$$
e_i = \sum\nolimits_j a_{s, ij} b_{s, j}, \forall i,
\quad\text{and}\quad
0 = \sum\nolimits_i a_i a_{s, ij}, \forall j.
$$
이들을 택한 뒤 $E_{k + 1} \subset B$를 다음 원소들이 생성하는
$A$-부분대수로 둔다.
% P01-E0998
\begin{enumerate}
\item $E_k$.
\item 모든 $s \in S_k$에 대한 원소 $b_{s, 1}, \ldots, b_{s, m_s}$.
\end{enumerate}
몇 가지 집합론(생략함)에 의해 $E_{k + 1}$의 기수는 많아야
$\kappa$이다. 이는 귀납적으로 $|E_k| \leq \kappa$임을 알고, 따라서
$S_k$의 기수도 많아야 $\kappa$라는 사실을 사용한다.
% P01-E0999

\medskip\noindent
$B' = \bigcup E_k$가 $A$ 위에서 평탄임을 보이기 위해, $A$에 계수를 갖는
$B'$에서의 관계 $\sum_{i = 1, \ldots, n} a_i b'_i = 0$을 생각하자.
모든 $i = 1, \ldots, n$에 대해 $b'_i \in E_k$가 되도록 충분히 큰 $k$를
택하자. 그러면 $(n, a_1, \ldots, a_n, b'_1, \ldots, b'_n) \in S_k$이고,
따라서 이 관계는 $E_{k + 1}$에서 자명하며, 더구나 $B'$에서도 자명하다.
그러므로 보조정리 \ref{algebra-lemma-flat-eq}에 의해 $A \to B'$는 평탄하다.
\end{proof}
\section{Cohen 구조 정리}
\label{algebra-section-cohen-structure-theorem}

\noindent
다음은 가환대수학의 기본 개념이다.

\begin{definition}
\label{algebra-definition-complete-local-ring}
$(R, \mathfrak m)$을 국소환이라 하자.
$\mathfrak m$에 관한 $R$의 완비화로 가는 표준 사상
$$
R \longrightarrow \lim_n R/\mathfrak m^n
$$
이 동형사상이면 $R$을 {\it 완비 국소환}이라 한다\footnote{여기에는
$\bigcap \mathfrak m^n = (0)$이라는 조건도 포함된다. 일부 문헌에서는
$R$이 완비이고 분리적이라고 말하여 이를 나타내기도 한다. 주의:
국소환의 완비화 $\lim_n R/\mathfrak m^n$가 완비가 아닐 수도 있다.
Examples, 보조정리 \ref{examples-lemma-noncomplete-completion}를 보라.
$\mathfrak m$이 유한 생성이면 이런 일은 일어나지 않는다. 보조정리
\ref{algebra-lemma-hathat-finitely-generated}를 보라. 이 경우 완비화는
뇌터환이다. 보조정리 \ref{algebra-lemma-completion-Noetherian}을 보라.}.
\end{definition}

\noindent
아르틴 국소환 $R$은 완비 국소환임에 주의하라. 실제로 어떤
$n > 0$에 대해 $\mathfrak m_R^n = 0$이다. 본 절에서는 주로
뇌터 완비 국소환에 집중한다.

\begin{lemma}
\label{algebra-lemma-quotient-complete-local}
$R$을 뇌터 완비 국소환이라 하자. $R$의 모든 상환도 뇌터 완비
국소환이다. 유한 환 사상 $R \to S$가 주어지면, $S$는 뇌터 완비
국소환들의 곱이다.
% P01-E1000
\end{lemma}

\begin{proof}
보조정리 \ref{algebra-lemma-Noetherian-permanence}에 의해 환 $S$는 뇌터이다.
보조정리 \ref{algebra-lemma-completion-tensor}에 의해 $R$-가군으로서 $S$는
완비이다. 따라서 보조정리 \ref{algebra-lemma-completion-finite-extension}에 의해
$S$는 그 극대 아이디얼들에서의 완비화들의 곱이다.
\end{proof}

\begin{lemma}
\label{algebra-lemma-complete-local-ring-Noetherian}
$(R, \mathfrak m)$을 완비 국소환이라 하자. $\mathfrak m$이 유한 생성
아이디얼이면 $R$은 뇌터환이다.
\end{lemma}

\begin{proof}
보조정리 \ref{algebra-lemma-completion-Noetherian}을 보라.
\end{proof}

\begin{definition}
\label{algebra-definition-coefficient-ring}
$(R, \mathfrak m)$을 완비 국소환이라 하자. 부분환
$\Lambda \subset R$가 다음 조건들을 만족하면 이를 {\it 계수환}이라 한다.
\begin{enumerate}
\item $\Lambda$는 극대 아이디얼이 $\Lambda \cap \mathfrak m$인
완비 국소환이다.
\item $\Lambda$의 잉여체는 $R$의 잉여체로 동형적으로 사상된다.
\item $\Lambda \cap \mathfrak m = p\Lambda$이다. 여기서 $p$는
$R$의 잉여체의 표수이다.
\end{enumerate}
\end{definition}

\noindent
이 정의에 관해 몇 가지 언급을 하자. 논의를 다음 경우들로 나눈다.
\begin{enumerate}
\item 국소환 $R$이 체를 포함하는 경우. 이는
$\mathbf{Q} \subset R$이거나, $p$가 $R/\mathfrak m$의 표수일 때
$pR = 0$이면 일어난다. 이 경우 계수환 $\Lambda$는 $R$에 포함되고
$R/\mathfrak m$으로 동형적으로 사상되는 체이다.
\item $R/\mathfrak m$의 표수가 $p > 0$이지만 $R$에서 $p$의 어느
거듭제곱도 영이 아닌 경우. 이 경우 $\Lambda$는 균일화원이 $p$이고
잉여체가 $R/\mathfrak m$인 완비 이산 값매김환이다.
\item $R/\mathfrak m$의 표수가 $p > 0$이고, 어떤 $n > 1$에 대해
$R$에서 $p^{n - 1} \not = 0$, $p^n = 0$인 경우. 이 경우
$\Lambda$는 극대 아이디얼이 $p$로 생성되고 잉여체가
$R/\mathfrak m$인 아르틴 국소환이다.
\end{enumerate}
위에 나타난 균일화원이 $p$인 완비 이산 값매김환들은 특별한 역할을
하므로 다음과 같이 이름 붙인다.

\begin{definition}
\label{algebra-definition-cohen-ring}
{\it Cohen 환}은 소수 $p$를 균일화원으로 갖는 완비 이산 값매김환이다.
\end{definition}

\begin{lemma}
\label{algebra-lemma-cohen-rings-exist}
$p$를 소수라 하자. $k$를 표수가 $p$인 체라 하자.
$\Lambda/p\Lambda \cong k$인 Cohen 환 $\Lambda$가 존재한다.
\end{lemma}

\begin{proof}
먼저 $p$-진 정수 $\mathbf{Z}_p$는 $\mathbf{F}_p$에 대한 Cohen 환을
이룸에 주의하라. $k$를 표수가 $p$인 임의의 체라 하자.
$\mathfrak m_R = pR$이고 $R/pR = k$인 평탄 국소 환 사상
$\mathbf{Z}_p \to R$을 택하자. 보조정리
\ref{algebra-lemma-flat-local-given-residue-field}를 보라.
보조정리 \ref{algebra-lemma-completion-Noetherian}에 의해 완비화
$\Lambda = R^\wedge$는 뇌터이다. $\Lambda/p\Lambda = R/pR$가
체이므로, $\Lambda$는 극대 아이디얼이 $(p)$인 완비 뇌터 국소환이다
(등식과 완비성은 보조정리 \ref{algebra-lemma-hathat-finitely-generated}에서 따른다).
모든 $n$에 대해 사상
$\mathbf{Z}/p^n\mathbf{Z} \to \Lambda/p^n\Lambda = R/p^nR$은
평탄하다(등식은 다시 보조정리
\ref{algebra-lemma-hathat-finitely-generated}에서 따른다).
따라서 예를 들어 보조정리 \ref{algebra-lemma-flat-module-powers}에 의해
$\mathbf{Z}_p \to \Lambda$는 평탄하다. 그러므로 $p$는 $\Lambda$의
비영인자이다. 따라서 $\Lambda$의 차원은 $\geq 1$이고
(보조정리 \ref{algebra-lemma-one-equation}), $\Lambda$는 차원 $1$인
정칙환, 즉 보조정리 \ref{algebra-lemma-characterize-dvr}에 의해
이산 값매김환이다. 따라서 $\Lambda$는 $k$에 대한 Cohen 환이다.
\end{proof}

\begin{lemma}
\label{algebra-lemma-cohen-ring-formally-smooth}
$p > 0$을 소수라 하자. $\Lambda$를 잉여체의 표수가 $p$인
Cohen 환이라 하자. 모든 $n \geq 1$에 대해 환 사상
$$
\mathbf{Z}/p^n\mathbf{Z} \to \Lambda/p^n\Lambda
$$
은 형식적으로 매끄럽다.
\end{lemma}

\begin{proof}
$n = 1$이면 이는 명제
\ref{algebra-proposition-characterize-separable-field-extensions}에서 따른다.
일반적인 $n$에 대해서는 $n$에 관한 귀납법을 쓴다. 즉
$\mathbf{Z}/p^n\mathbf{Z} \to \Lambda/p^n\Lambda$가 형식적으로
매끄러우면, 환 사상
$\mathbf{Z}/p^{n + 1}\mathbf{Z} \to \Lambda/p^{n + 1}\Lambda$과
아이디얼 $I = (p^n) \subset \mathbf{Z}/p^{n + 1}\mathbf{Z}$에
보조정리 \ref{algebra-lemma-lift-formal-smoothness}를 적용할 수 있다.
\end{proof}

\begin{theorem}[Cohen 구조 정리]
\label{algebra-theorem-cohen-structure-theorem}
$(R, \mathfrak m)$을 완비 국소환이라 하자.
\begin{enumerate}
\item $R$은 계수환을 갖는다(정의
\ref{algebra-definition-coefficient-ring}를 보라).
\item $\mathfrak m$이 유한 생성 아이디얼이면 $R$은 다음 상환과
동형이다.
$$
\Lambda[[x_1, \ldots, x_n]]/I
$$
여기서 $\Lambda$는 체이거나 Cohen 환이다.
\end{enumerate}
\end{theorem}

\begin{proof}
먼저 계수환이 존재함을 증명하자. 우선 잉여체 $\kappa$의 표수가
영인 경우를 증명한다. 즉 이 경우 $n > 0$에 관한 귀납법으로,
표준 사상 $R/\mathfrak m^n \to \kappa = R/\mathfrak m$의 절단
$$
\varphi_n : \kappa \longrightarrow R/\mathfrak m^n
$$
이 존재함을 보이겠다. $n = 1$이면 자명하다. $n > 1$이고
$\varphi_{n - 1}$이 주어졌다고 하자. 체 확대
$\kappa/\mathbf{Q}$는 명제
\ref{algebra-proposition-characterize-separable-field-extensions}에 의해
형식적으로 매끄럽다. 따라서 다음 도식에서 점선 화살표를 찾을 수 있다.
$$
\xymatrix{
R/\mathfrak m^{n - 1} &
R/\mathfrak m^n \ar[l] \\
\kappa \ar[u]^{\varphi_{n - 1}} \ar@{..>}[ru] & \mathbf{Q} \ar[l] \ar[u]
}
$$
이는 귀납 단계를 증명한다. 이 사상들을 합치면
$$
\lim_n\ \varphi_n : \kappa \longrightarrow
R = \lim_n\ R/\mathfrak m^n
$$
을 얻고, 그 상이 원하는 계수환이다.

\medskip\noindent
다음으로 잉여체 $\kappa$의 표수가 $p > 0$인 경우 계수환의 존재를
증명한다. 즉 $\kappa = \Lambda/p\Lambda$인 Cohen 환 $\Lambda$를
택하자. 보조정리 \ref{algebra-lemma-cohen-rings-exist}를 보라.
이 경우 $n > 0$에 관한 귀납법으로 다음 사상이 존재함을 보이겠다.
$$
\varphi_n :
\Lambda/p^n\Lambda
\longrightarrow
R/\mathfrak m^n
$$
이 사상과 축약 사상 $R/\mathfrak m^n \to \kappa$의 합성은 주어진
동형사상 $\Lambda/p\Lambda = \kappa$를 준다. $n = 1$이면 자명하다.
$n > 1$이고 $\varphi_{n - 1}$이 주어졌다고 하자. 환 사상
$\mathbf{Z}/p^n\mathbf{Z} \to \Lambda/p^n\Lambda$는 보조정리
\ref{algebra-lemma-cohen-ring-formally-smooth}에 의해 형식적으로 매끄럽다.
따라서 다음 도식에서 점선 화살표를 찾을 수 있다.
$$
\xymatrix{
R/\mathfrak m^{n - 1} &
R/\mathfrak m^n \ar[l] \\
\Lambda/p^n\Lambda \ar[u]^{\varphi_{n - 1}} \ar@{..>}[ru] &
\mathbf{Z}/p^n\mathbf{Z} \ar[l] \ar[u]
}
$$
이는 귀납 단계를 증명한다. 이 사상들을 합치면
$$
\lim_n\ \varphi_n :
\Lambda = \lim_n\ \Lambda/p^n\Lambda
\longrightarrow
R = \lim_n\ R/\mathfrak m^n
$$
을 얻고, 그 상이 원하는 계수환이다.

\medskip\noindent
정리의 마지막 명제도 바로 따른다. 즉
$y_1, \ldots, y_n$이 아이디얼 $\mathfrak m$의 생성원들이면,
방금 구성한 사상 $\Lambda \to R$을 사용하여 다음 사상을 얻는다.
$$
\Lambda[[x_1, \ldots, x_n]] \longrightarrow R,
\quad x_i \longmapsto y_i.
$$
양변은 모두 $(x_1, \ldots, x_n)$-진 완비이고, 구성에 의해 이 사상은
$(x_1, \ldots, x_n)$을 법으로 하여 전사이므로, 보조정리
\ref{algebra-lemma-completion-generalities}에 의해 전사이다.
\end{proof}

\begin{remark}
\label{algebra-remark-Noetherian-complete-local-ring-universally-catenary}
$k$가 체이면 멱급수환 $k[[X_1, \ldots, X_d]]$는 차원이 $d$인
뇌터 완비 정칙 국소환이다. $\Lambda$가 Cohen 환이면
$\Lambda[[X_1, \ldots, X_d]]$는 차원이 $d + 1$인 완비 국소
뇌터 정칙환이다. 따라서 Cohen 구조 정리는 모든 뇌터 완비 국소환이
정칙 국소환의 상환임을 함의한다. 특히 뇌터 완비 국소환은
보편적으로 카테너리이다. 보조정리 \ref{algebra-lemma-CM-ring-catenary}와
보조정리 \ref{algebra-lemma-regular-ring-CM}을 보라.
\end{remark}

\begin{lemma}
\label{algebra-lemma-regular-complete-containing-coefficient-field}
$(R, \mathfrak m)$을 뇌터 완비 국소환이라 하자. $R$이 정칙이라고
가정하자.
\begin{enumerate}
\item $R$이 $\mathbf{F}_p$ 또는 $\mathbf{Q}$를 포함하면, $R$은
그 잉여체 위의 멱급수환과 동형이다.
\item $k$가 체이고 $k \to R$이 동형사상
$k \to R/\mathfrak m$을 유도하는 환 사상이면, $R$은
$k$-대수로서 $k$ 위의 멱급수환과 동형이다.
\end{enumerate}
\end{lemma}

\begin{proof}
(1)의 경우 Cohen 구조 정리(정리
\ref{algebra-theorem-cohen-structure-theorem})에 의해 계수환이 존재하고,
이는 잉여체로 동형적으로 사상되는 체여야 한다. 따라서 (2)를
증명하면 충분하다. (2)의 경우
$f_1, \ldots, f_d \in \mathfrak m$을 택하여
$\mathfrak m/\mathfrak m^2$의 기저로 사상되게 하고, $x_i$를 $f_i$로
보내는 연속 $k$-대수 사상 $k[[x_1, \ldots, x_d]] \to R$을 생각하자.
정의역과 공역이 모두 $(x_1, \ldots, x_d)$-진 완비이므로, 보조정리
\ref{algebra-lemma-completion-generalities}에 의해 이 사상은 전사이다.
한편 이 사상은 단사여야 한다. 그렇지 않으면 보조정리
\ref{algebra-lemma-one-equation}에 의해 $R$의 차원이 $< d$이기 때문이다.
\end{proof}

\begin{lemma}
\label{algebra-lemma-complete-local-Noetherian-domain-finite-over-regular}
$(R, \mathfrak m)$을 뇌터 완비 국소 정역이라 하자. 그러면 다음
성질들을 갖는 $R_0 \subset R$가 존재한다.
% P01-E1001
\begin{enumerate}
\item $R_0$는 정칙 완비 국소환이다.
\item $R_0 \subset R$는 유한이고 잉여체 위의 동형사상을 유도한다.
\item $R_0$는 $k$가 체일 때 $k[[X_1, \ldots, X_d]]$와 동형이거나,
$\Lambda$가 Cohen 환일 때 $\Lambda[[X_1, \ldots, X_d]]$와 동형이다.
\end{enumerate}
\end{lemma}

\begin{proof}
$\Lambda$를 $R$의 계수환이라 하자. $R$이 정역이므로
$\Lambda$는 체이거나, 아니면 $\Lambda$는 Cohen 환이다.

\medskip\noindent
경우 I: $\Lambda = k$가 체이다. $d = \dim(R)$로 두자.
정의 아이디얼 $I \subset R$을 생성하는
$x_1, \ldots, x_d \in \mathfrak m$을 택하자.
(절 \ref{algebra-section-dimension}을 보라.)
보조정리 \ref{algebra-lemma-change-ideal-completion}에 의해 $R$은
$I$-진으로도 완비이다. $X_i$를 $x_i$로 보내는 사상
$R_0 = k[[X_1, \ldots, X_d]] \to R$을 생각하자.
$R_0$는 아이디얼 $I_0 = (X_1, \ldots, X_d)$에 관해 완비이고,
$R/I_0R \cong R/IR$는 $k = R_0/I_0$ 위에서 유한이다
($\dim(R/I) = 0$이기 때문이다. 절 \ref{algebra-section-dimension}을 보라).
따라서 보조정리 \ref{algebra-lemma-finite-over-complete-ring}에 의해
$R_0 \to R$은 유한이다. $\dim(R) = \dim(R_0)$이므로 이는
$R_0 \to R$이 단사임을 함의하며(보조정리
\ref{algebra-lemma-integral-dim-up}을 보라), 보조정리가 증명된다.

\medskip\noindent
경우 II: $\Lambda$가 Cohen 환이다. $d + 1 = \dim(R)$로 두자.
$p > 0$을 잉여체 $k$의 표수라 하자. $R$이 정역이므로 $p$는
$R$의 비영인자이다. 따라서 $\dim(R/pR) = d$이다. 보조정리
\ref{algebra-lemma-one-equation}을 보라.
$R/pR$에서 정의 아이디얼을 생성하는
$x_1, \ldots, x_d \in R$을 택하자. 그러면
$I = (p, x_1, \ldots, x_d)$는 $R$의 정의 아이디얼이다.
보조정리 \ref{algebra-lemma-change-ideal-completion}에 의해 $R$은
$I$-진으로도 완비이다. $X_i$를 $x_i$로 보내는 사상
$R_0 = \Lambda[[X_1, \ldots, X_d]] \to R$을 생각하자.
$R_0$는 아이디얼 $I_0 = (p, X_1, \ldots, X_d)$에 관해 완비이고,
$R/I_0R \cong R/IR$는 $k = R_0/I_0$ 위에서 유한이다
($\dim(R/I) = 0$이기 때문이다. 절 \ref{algebra-section-dimension}을 보라).
따라서 보조정리 \ref{algebra-lemma-finite-over-complete-ring}에 의해
$R_0 \to R$은 유한이다. $\dim(R) = \dim(R_0)$이므로 이는
$R_0 \to R$이 단사임을 함의하며(보조정리
\ref{algebra-lemma-integral-dim-up}을 보라), 보조정리가 증명된다.
\end{proof}


\section{Japanese 환}
\label{algebra-section-japanese}

\noindent
이 절에서는 정수적 폐포의 유한성을 논의하기 시작한다.

\begin{definition}
\label{algebra-definition-N}
\begin{reference}
\cite[Chapter 0, Definition 23.1.1]{EGA}
\end{reference}
$R$을 분수체가 $K$인 정역이라 하자.
\begin{enumerate}
\item $K$에서 $R$의 정수적 폐포가 유한 $R$-가군이면
$R$을 {\it N-1}이라 한다.
\item 모든 체의 유한확대 $L/K$에 대해 $L$에서 $R$의 정수적 폐포가
$R$ 위에서 유한이면 $R$을 {\it N-2} 또는 {\it Japanese}라고 한다.
\end{enumerate}
\end{definition}

\noindent
이 개념들의 주된 관심사는 뇌터환에 있지만, 다음은 뇌터가 아닌 예이다.

\begin{example}
\label{algebra-example-Japanese-not-Noetherian}
$k$를 체라 하자. 정역 $R = k[x_1, x_2, x_3, \ldots]$는 N-2이지만
뇌터가 아니다. 이유는 다음과 같다. $R \subset L$이고 체 $L$이
$R$의 분수체의 유한확대라고 하자. 그러면 어떤 정수 $n$이 존재하여
$L$은 유한확대 $L_0/k(x_1, \ldots, x_n)$에 (초월적인) 원소들
$x_{n + 1}, x_{n + 2}$ 등을 첨가해서 얻어진다.
$S_0$를 $L_0$에서 $k[x_1, \ldots, x_n]$의 정수적 폐포라 하자.
아래 명제 \ref{algebra-proposition-ubiquity-nagata}에 의해
$S_0$는 $k[x_1, \ldots, x_n]$ 위에서 유한이다.
또한 $L$에서 $R$의 정수적 폐포는
$S = S_0[x_{n + 1}, x_{n + 2}, \ldots]$이고(보조정리
\ref{algebra-lemma-polynomial-domain-normal}을 사용하라), 따라서 $R$ 위에서
유한이다. 같은 논증은 $R = \mathbf{Z}[x_1, x_2, x_3, \ldots]$에도
적용된다.
\end{example}

\begin{lemma}
\label{algebra-lemma-localize-N}
$R$을 정역이라 하자. $R$이 N-1이면 $R$의 모든 국소화도 N-1이다.
N-2에 대해서도 같다.
\end{lemma}

\begin{proof}
정수적 폐포를 취하는 것은 국소화와 가환하므로 이 명제들이 성립한다.
보조정리 \ref{algebra-lemma-integral-closure-localize}를 보라.
\end{proof}

\begin{lemma}
\label{algebra-lemma-Japanese-local}
$R$을 정역이라 하자. $f_1, \ldots, f_n \in R$이 단위 아이디얼을
생성한다고 하자. 각 정역 $R_{f_i}$가 N-1이면 $R$도 N-1이다.
N-2에 대해서도 같다.
\end{lemma}

\begin{proof}
$R_{f_i}$가 N-2라고 가정하자(N-1의 경우도 같다).
$L$을 $R$의 분수체의 유한확대라 하자(N-1의 경우에는 분수체 자체).
$S$를 $L$에서 $R$의 정수적 폐포라 하자. 보조정리
\ref{algebra-lemma-integral-closure-localize}에 의해 $S_{f_i}$는 $L$에서
$R_{f_i}$의 정수적 폐포이다. 따라서 가정에 의해 $S_{f_i}$는
$R_{f_i}$ 위에서 유한이다. 그러므로 보조정리 \ref{algebra-lemma-cover}에 의해
$S$는 $R$ 위에서 유한이다.
\end{proof}

\begin{lemma}
\label{algebra-lemma-quasi-finite-over-Noetherian-japanese}
$R$을 정역이라 하자. $R \subset S$를 정역들의 준유한 확대
(예를 들어 유한확대)라 하자. $R$이 N-2이자 뇌터라고 가정하자.
그러면 $S$는 N-2이다.
\end{lemma}

\begin{proof}
$L/K$를 유도된 분수체들의 확대라 하자. 이는 체의 유한확대임에
주의하라(예를 들어 섬유 $S \otimes_R K$에 적용한 보조정리
\ref{algebra-lemma-isolated-point-fibre} (2)와 준유한 환 사상의 정의에 의한다).
$S'$을 $S$에서 $R$의 정수적 폐포라 하자. 그러면 $S'$은
$L$에서 $R$의 정수적 폐포에 포함되며, 이는 가정에 의해 $R$ 위에서
유한이다. $R$이 뇌터이므로 $S'$은 $R$ 위에서 유한이다.
보조정리 \ref{algebra-lemma-quasi-finite-open-integral-closure}에 의해
$S'_{g_i} \cong S_{g_i}$이고 $g_1, \ldots, g_n$이 $S$에서 단위
아이디얼을 생성하도록 하는 원소들 $g_1, \ldots, g_n \in S'$이 존재한다.
따라서 보조정리 \ref{algebra-lemma-localize-N}과
\ref{algebra-lemma-Japanese-local}에 의해 $S'$이 N-2임을 보이면 충분하다.
그러므로 $S$가 $R$ 위에서 유한인 경우로 귀착되었다.

\medskip\noindent
$R \subset S$가 보조정리의 가정들을 만족하고, 더 나아가 $S$가
$R$ 위에서 유한이라고 하자. $M$을 $S$의 분수체의 체의 유한확대라
하자. 그러면 $M$은 $K$의 체의 유한확대이기도 하므로, $M$에서
$R$의 정수적 폐포 $T$는 $R$ 위에서 유한이다. 보조정리
\ref{algebra-lemma-integral-closure-transitive}에 의해 $T$는 $M$에서
$S$의 정수적 폐포이기도 하며, 보조정리
\ref{algebra-lemma-integral-permanence}에 의해 결론을 얻는다.
\end{proof}

\begin{lemma}
\label{algebra-lemma-Laurent-ring-N-1}
$R$을 뇌터 정역이라 하자. $R[z, z^{-1}]$이 N-1이면 $R$도 N-1이다.
\end{lemma}

\begin{proof}
$R'$을 분수체 $K$에서 $R$의 정수적 폐포라 하자.
$S'$을 $R[z, z^{-1}]$의 분수체에서의 정수적 폐포라
하자. 명백히 $R' \subset S'$이다. $K[z, z^{-1}]$이 정규 정역이므로
$S' \subset K[z, z^{-1}]$이다. $f_1, \ldots, f_n \in S'$이
$R[z, z^{-1}]$-가군으로서 $S'$을 생성한다고 하자.
$f_i = \sum a_{ij}z^j$라 쓰자(유한합). 여기서 $a_{ij} \in K$이다.
모든 $x \in R'$에 대해
$$
x = \sum h_i f_i
$$
로 쓸 수 있으며, 여기서 $h_i \in R[z, z^{-1}]$이다. 따라서
$R'$은 유한 $R$-부분가군 $\sum Ra_{ij} \subset K$에 포함된다.
$R$이 뇌터이므로 $R'$은 유한 $R$-가군이다.
\end{proof}

\begin{lemma}
\label{algebra-lemma-finite-extension-N-2}
$R$을 뇌터 정역이라 하고 $R \subset S$를 정역들의 유한확대라 하자.
$S$가 N-1이면 $R$도 N-1이다. $S$가 N-2이면 $R$도 N-2이다.
\end{lemma}

\begin{proof}
생략한다. (힌트: 확대체들에서 $R$의 정수적 폐포들은 확대체들에서
$S$의 정수적 폐포들에 포함된다.)
\end{proof}

\begin{lemma}
\label{algebra-lemma-Noetherian-normal-domain-finite-separable-extension}
$R$을 분수체가 $K$인 뇌터 정규 정역이라 하자.
$L/K$를 체의 유한 분리 확대라 하자. 그러면 $L$에서 $R$의
정수적 폐포는 $R$ 위에서 유한이다.
\end{lemma}

\begin{proof}
대각합 짝지음
(Fields, 정의 \ref{fields-definition-trace-pairing})
$$
L \times L \longrightarrow K,
\quad (x, y) \longmapsto \langle x, y\rangle := \text{Trace}_{L/K}(xy).
$$
을 생각하자. $L/K$가 분리이므로 이는 비퇴화이다
(Fields, 보조정리 \ref{fields-lemma-separable-trace-pairing}).
또한 $x \in L$이 $R$ 위에서 정수적이면
$\text{Trace}_{L/K}(x)$는 $R$에 속한다. 실제로 $K$ 위에서 $x$의
최소다항식의 계수들은 $R$에 속하고(보조정리
\ref{algebra-lemma-minimal-polynomial-normal-domain}),
$\text{Trace}_{L/K}(x)$는 이 계수들 가운데 하나의 정수배이기 때문이다
(Fields, 보조정리
\ref{fields-lemma-trace-and-norm-from-minimal-polynomial}).
$R$ 위에서 정수적이고 $L$의 $K$-기저를 이루는
$x_1, \ldots, x_n \in L$을 택하자. 그러면 정수적 폐포
$S \subset L$은 $R$-가군
$$
M = \{y \in L \mid \langle x_i, y\rangle \in R, \ i = 1, \ldots, n\}
$$
에 포함된다. 선형대수학에 의해 $R$-가군으로서
$M \cong R^{\oplus n}$이다. 따라서 $R$이 뇌터이므로
$S \subset R^{\oplus n}$은 유한 생성 $R$-가군이다.
\end{proof}

\begin{example}
\label{algebra-example-bad-invariants}
환이 뇌터가 아니면 보조정리
\ref{algebra-lemma-Noetherian-normal-domain-finite-separable-extension}은
성립하지 않는다. 예를 들어 $-1$이 $x_i$를 $-x_i$로 보내는
$A = \mathbf{C}[x_1, x_2, x_3, \ldots]$ 위의
$G = \{+1, -1\}$의 작용을 생각하자. 불변환 $R = A^G$는 모든
$x_ix_j$로 생성되는 $\mathbf{C}$-대수이다. 따라서
$R \subset A$는 유한이 아니다. 그러나 $R$은 정규 정역이고,
그 분수체는 $A$의 분수체 $L$ 안의 $G$-불변원들로 이루어진
$K = L^G$이다. 또한 명백히 $A$는 $L$에서 $R$의 정수적 폐포이다.
% P01-E1002
\end{example}

\noindent
순수 비분리 확대의 경우에는 다음 보조정리를 보조정리
\ref{algebra-lemma-Noetherian-normal-domain-finite-separable-extension} 대신
쓸 수 있는 때가 있다.

\begin{lemma}
\label{algebra-lemma-Noetherian-normal-domain-insep-extension}
$R$을 분수체 $K$의 표수가 $p > 0$인 뇌터 정규 정역이라 하자.
$D(a) \not = 0$인 미분 $D : R \to R$가 존재하는 원소
$a \in K$를 택하자. 그러면 $L = K[x]/(x^p - a)$에서
$R$의 정수적 폐포는 $R$ 위에서 유한이다.
\end{lemma}

\begin{proof}
어떤 $f \in R$에 대해 $x$를 $fx$로, $a$를 $f^pa$로 바꾸어
$a \in R$이라고 가정할 수 있다. 따라서 $D(a) \in R$이기도 하다.
$i \leq p - 1$에 관한 귀납법으로 다음을 보이겠다.
$$
y = a_0 + a_1x + \ldots + a_i x^i,\quad a_j \in K
$$
가 $R$ 위에서 정수적이면 $D(a)^i a_j \in R$이다. 따라서 정수적
폐포는 기저 $D(a)^{-p + 1}x^j$, $j = 0, \ldots, p - 1$을 갖는
유한 $R$-가군에 포함된다. $R$이 뇌터이므로 이것으로 보조정리가
증명된다.

\medskip\noindent
$i = 0$이면 $y = a_0$가 $R$ 위에서 정수적일 필요충분조건은
$a_0 \in R$이고, 명제는 참이다. 명제가 어떤 $i < p - 1$에 대해
성립한다고 하고
$$
y = a_0 + a_1x + \ldots + a_{i + 1} x^{i + 1},\quad a_j \in K
$$
가 $R$ 위에서 정수적이라고 하자. 그러면
$$
y^p = a_0^p + a_1^p a + \ldots + a_{i + 1}^pa^{i + 1}
$$
은 $R$의 원소이다($K$에 속하면서 $R$ 위에서 정수적이기 때문이다).
$D$를 적용하면
$$
(a_1^p + 2a_2^p a + \ldots + (i + 1)a_{i + 1}^p a^i)D(a)
$$
가 $R$에 속한다. 따라서
$$
D(a)a_1 + 2D(a) a_2 x + \ldots + (i + 1)D(a) a_{i + 1} x^i
$$
은 $R$ 위에서 정수적이다. 귀납법에 의해 $j = 1, \ldots, i + 1$에
대해 $D(a)^{i + 1}a_j \in R$을 얻는다. (여기서
$1, \ldots, i + 1$이 가역임을 사용한다.) 또한
$D(a)^{i + 1}a_0$도 $R$에 속한다. 실제로 이는
$D(a)^{i + 1}y$와 $\sum_{j > 0} D(a)^{i + 1}a_jx^j$의 차이고,
두 원소 모두 $R$ 위에서 정수적이다($a \in R$이므로 $x$가
$R$ 위에서 정수적이기 때문이다).
\end{proof}

\begin{lemma}
\label{algebra-lemma-domain-char-zero-N-1-2}
분수체의 표수가 영인 뇌터 정역이 N-1일 필요충분조건은 N-2
(즉 Japanese)인 것이다.
\end{lemma}

\begin{proof}
표수 영인 모든 체확대가 분리이므로 보조정리
\ref{algebra-lemma-Noetherian-normal-domain-finite-separable-extension}에서
명백히 따른다.
\end{proof}

\begin{lemma}
\label{algebra-lemma-domain-char-p-N-1-2}
$R$을 분수체 $K$의 표수가 $p > 0$인 뇌터 정역이라 하자.
그러면 $R$이 N-2일 필요충분조건은 모든 유한 순수 비분리 확대
$L/K$에 대해 $L$에서 $R$의 정수적 폐포가 $R$ 위에서 유한인 것이다.
\end{lemma}

\begin{proof}
$K$의 모든 유한 순수 비분리 체확대에서 $R$의 정수적 폐포가
유한이라고 가정하자. $L/K$를 임의의 유한확대라 하자.
$L$에서 $R$의 정수적 폐포가 $R$ 위에서 유한임을 보여야 한다.
$L$을 포함하는 체의 유한 정규 확대 $M/K$를 택하자.
$R$이 뇌터이므로 $M$에서 $R$의 정수적 폐포가 $R$ 위에서 유한임을
보이면 충분하다. Fields, 보조정리 \ref{fields-lemma-normal-case}에 의해
$M_{insep}/K$는 순수 비분리이고 $M/M_{insep}$은 분리인 부분확대
$M/M_{insep}/K$가 존재한다. 가정에 의해 $M_{insep}$에서 $R$의
정수적 폐포 $R'$은 $R$ 위에서 유한이다. 보조정리
\ref{algebra-lemma-Noetherian-normal-domain-finite-separable-extension}에 의해
$M$에서 $R'$의 정수적 폐포 $R''$은 $R'$ 위에서 유한이다.
그러면 보조정리 \ref{algebra-lemma-finite-transitive}에 의해 $R''$은
$R$ 위에서 유한이다. $R''$은 $M$에서 $R$의 정수적 폐포이기도
하므로(보조정리 \ref{algebra-lemma-integral-closure-transitive}를 보라)
결론을 얻는다.
\end{proof}

\begin{lemma}
\label{algebra-lemma-polynomial-ring-N-2}
$R$을 뇌터 정역이라 하자. $R$이 N-1이면 $R[x]$는 N-1이다.
$R$이 N-2이면 $R[x]$는 N-2이다.
\end{lemma}

\begin{proof}
$R$이 N-1이라고 하자. $R'$을 $R$의 정수적 폐포라 하자.
이는 $R$ 위에서 유한이다. 따라서 $R'[x]$도 $R[x]$ 위에서
유한이다. 환 $R'[x]$는 정규이고(보조정리
\ref{algebra-lemma-polynomial-domain-normal}을 보라), 따라서 N-1이다.
이로써 첫 번째 명제가 증명된다.

\medskip\noindent
두 번째 명제에 대해서는 보조정리 \ref{algebra-lemma-finite-extension-N-2}에
의해 $R'[x]$가 N-2임을 보이면 충분하다. 다시 말해 $R$이 정규
N-2 정역이라고 가정해도 좋고 그렇게 하자. 표수가 영이면 보조정리
\ref{algebra-lemma-domain-char-zero-N-1-2}에 의해 끝난다. 표수가
$p > 0$이면 $R$의 분수체를 $K$라 할 때 $L/K(x)$의 임의의 유한
순수 비분리 확대에서 $R[x]$의 정수적 폐포가 유한임을 보여야 한다.
% P01-E1003
어떤 체의 유한 순수 비분리 확대 $L'/K$와 $q = p^e$가 존재하여
$L \subset L'(x^{1/q})$이다. 일부 세부사항은 생략한다.
$R[x]$가 뇌터이므로 $L'(x^{1/q})$에서 $R[x]$의 정수적 폐포가
$R[x]$ 위에서 유한임을 보이면 충분하다. 이 정수적 폐포는
$R \subset R' \subset L'$을 $L'$에서 $R$의 정수적 폐포로 두었을 때
$R'[x^{1/q}]$과 같다.
% P01-E1004
$R$이 N-2이므로 $R'$은 $R$ 위에서 유한이고, 따라서
$R'[x^{1/q}]$은 $R[x]$ 위에서 유한이다.
\end{proof}

\begin{lemma}
\label{algebra-lemma-openness-normal-locus}
$R$을 뇌터 정역이라 하자. $R_f$가 정규인 비영원
$f \in R$가 존재하면
$$
U = \{\mathfrak p \in \Spec(R) \mid R_{\mathfrak p} \text{ is normal}\}
$$
은 $\Spec(R)$에서 열린집합이다.
\end{lemma}

\begin{proof}
표준 열린집합 $D(f)$가 $U$에 포함됨은 명백하다.
세르 판정법인 보조정리 \ref{algebra-lemma-criterion-normal}에 의해
$\mathfrak p \not \in U$이면 어떤
$\mathfrak q \subset \mathfrak p$에 대해 다음 둘 중 하나가 성립한다.
\begin{enumerate}
\item 경우 I: $\text{depth}(R_{\mathfrak q}) < 2$이고
$\dim(R_{\mathfrak q}) \geq 2$이다.
\item 경우 II: $R_{\mathfrak q}$는 정칙이 아니고
$\dim(R_{\mathfrak q}) = 1$이다.
\end{enumerate}
이는 특히 $R_{\mathfrak q}$가 정규가 아님을 뜻하므로
$f \in \mathfrak q$이다. 경우 I에서는
$\text{depth}(R_{\mathfrak q}) =
\text{depth}(R_{\mathfrak q}/fR_{\mathfrak q}) + 1$이다.
따라서 그러한 소 아이디얼 $\mathfrak q$는 $R/fR$의 매립 연관
소 아이디얼과 같은 자료이다. 경우 II에서 $\mathfrak q$는 높이가
1인 $R/fR$의 연관 소 아이디얼이다. 그러므로 그러한 소 아이디얼들
$\mathfrak q$의 유한집합 $E$가 존재하고(보조정리
\ref{algebra-lemma-finite-ass}를 보라),
$$
\Spec(R) \setminus U
=
\bigcup\nolimits_{\mathfrak q \in E} V(\mathfrak q)
$$
가 되어 원하는 결론을 얻는다.
\end{proof}

\begin{lemma}
\label{algebra-lemma-characterize-N-1}
$R$을 뇌터 정역이라 하자. 그러면 $R$이 N-1일 필요충분조건은
다음 두 조건이 성립하는 것이다.
\begin{enumerate}
\item $R_f$가 정규인 비영원 $f \in R$가 존재한다.
\item 모든 극대 아이디얼 $\mathfrak m \subset R$에 대해 국소환
$R_{\mathfrak m}$은 N-1이다.
\end{enumerate}
\end{lemma}

\begin{proof}
먼저 $R$이 N-1이라고 하자. $R'$을 분수체 $K$에서 $R$의 정수적
폐포라 하자. 가정에 의해 $R$-가군으로서 $R'$을 생성하는
$x_1, \ldots, x_n$을 $R'$에서 찾을 수 있다. $R' \subset K$이므로
$f_i x_i \in R$인 비영원 $f_i \in R$를 찾을 수 있다.
$f = f_1 \ldots f_n$이면 $R_f \cong R'_f$이다. 따라서 $R_f$는
정규이고 (1)을 얻는다. (2)는 보조정리 \ref{algebra-lemma-localize-N}에서
따른다.

\medskip\noindent
(1)과 (2)를 가정하자. $K$를 $R$의 분수체라 하자.
$R$의 유한확대이고 $K$에 포함되는
$R \subset R' \subset K$가 주어졌다고 하자. $R_f$가 이미
정규이므로 $R_f = R'_f$임에 주의하라. 따라서 보조정리
\ref{algebra-lemma-openness-normal-locus}에 의해 $R'_{\mathfrak p'}$가
정규가 아닌 소 아이디얼들 $\mathfrak p' \in \Spec(R')$의 집합은
$\Spec(R')$에서 닫혀 있다. $\Spec(R') \to \Spec(R)$가 닫힌 사상이므로
이 집합의 상은 $\Spec(R)$에서 닫혀 있다. 그러한 환 $R'$에 대해
이 상을 $Z_{R'} \subset \Spec(R)$로 나타내자.
% P01-E1005

\medskip\noindent
극대 아이디얼 $\mathfrak m \subset R$을 택하자.
$R_{\mathfrak m} \subset R_{\mathfrak m}'$을 $K$에서 이 국소환의
정수적 폐포라 하자. 가정에 의해 이는 유한 환 확대이다.
보조정리 \ref{algebra-lemma-integral-closure-localize}에 의해 $R$ 위에서
정수적인 유한 개의 원소 $x_1, \ldots, x_n \in K$를 찾아
$R_{\mathfrak m}'$이 $R_{\mathfrak m}$ 위에서
$x_1, \ldots, x_n$로 생성되게 할 수 있다.
$R' = R[x_1, \ldots, x_n] \subset K$로 두자. 이 선택에 대해
$\mathfrak m \not \in Z_{R'}$임은 명백하다.

\medskip\noindent
$\Spec(R)$이 준콤팩트이므로 위 사실에 의해
$\bigcap Z_{R'_i} = \emptyset$인 유한 개의
$R \subset R'_i \subset K$를 찾을 수 있다. $R'$을 이들 모두가
$K$ 안에서 생성하는 부분환이라 하자. 이는 $R$ 위에서 유한이다.
또한 $Z_{R'} = \emptyset$이다. 실제로 모든 소 아이디얼
$\mathfrak p'$는 $(R'_i)_{\mathfrak p'_i}$가 정규인 소 아이디얼
$\mathfrak p'_i$ 위에 놓인다. 이는
$R'_{\mathfrak p'} = (R'_i)_{\mathfrak p'_i}$도 정규임을 함의한다.
따라서 $R'$은 정규이다. 다시 말해 $R'$은 $K$에서 $R$의
정수적 폐포이다.
\end{proof}

\begin{lemma}[Tate]
\label{algebra-lemma-tate-japanese}
\begin{reference}
\cite[Theorem 23.1.3]{EGA}
\end{reference}
$R$을 환이라 하자. $x \in R$라 하자. 다음을 가정하자.
\begin{enumerate}
\item $R$은 정규 뇌터 정역이다.
\item $R/xR$은 정역이고 N-2이다.
\item $R \cong \lim_n R/x^nR$이고 $x$에 관해 완비이다.
\end{enumerate}
그러면 $R$은 N-2이다.
\end{lemma}

\begin{proof}
자명한 경우를 제외하면 $x \not = 0$이라고 가정해도 좋다.
$K$를 $R$의 분수체라 하자. $K$의 표수가 영이면 보조정리는
(1)에서 따른다. 보조정리 \ref{algebra-lemma-domain-char-zero-N-1-2}를 보라.
따라서 $K$의 표수가 $p > 0$이라고 가정해도 좋고, 보조정리
\ref{algebra-lemma-domain-char-p-N-1-2}를 적용할 수 있다. 그러므로
$L/K$가 체의 유한 순수 비분리 확대일 때 $L$에서 $R$의 정수적
폐포 $S$가 $R$ 위에서 유한임을 보여야 한다.

\medskip\noindent
$L^q \subset K$인 $p$의 거듭제곱 $q$를 택하자.
필요하면 $L$을 확대하여 $y^q = x$인 원소 $y \in L$이 존재한다고
가정해도 좋다. $R \to S$가 스펙트럼의 위상동형사상을 유도하므로
(보조정리 \ref{algebra-lemma-p-ring-map}을 보라), 소 아이디얼
$\mathfrak p = xR$ 위에 놓이는 소 아이디얼
$\mathfrak q \subset S$가 유일하게 존재한다. 명백히
$$
\mathfrak q = \{f \in S \mid f^q \in \mathfrak p\} = yS
$$
이다. 왜냐하면 $y^q = x$이기 때문이다. 보조정리
\ref{algebra-lemma-characterize-dvr}에 의해 $R_{\mathfrak p}$가 이산
값매김환임에 주의하라. 그러면 Krull--Akizuki(보조정리
\ref{algebra-lemma-krull-akizuki})에 의해 $S_{\mathfrak q}$는 뇌터이다.
이에 다시 보조정리 \ref{algebra-lemma-characterize-dvr}를 적용하면
$S_{\mathfrak q}$가 이산 값매김환임을 얻는다.
보조정리 \ref{algebra-lemma-finite-extension-residue-fields-dimension-1}에
의해 $\kappa(\mathfrak q)/\kappa(\mathfrak p)$는 체의 유한확대이다.
따라서 $R/xR$의 정수적 폐포 $S' \subset \kappa(\mathfrak q)$는
가정 (2)에 의해 $R/xR$ 위에서 유한이다.
$S/yS \subset S'$이므로 $S/yS$는 $R$ 위에서 유한하다.
$S/y^nS$는 각 몫이
$y^iS/y^{i + 1}S \cong S/yS$인 유한 여과를 갖는다.
따라서 각 $S/y^nS$가 $R$ 위에서 유한임을 알 수 있다.
특히 $S/xS$는 $R$ 위에서 유한이다. 또한 교집합에 속하는 원소의
$q$제곱이 $\bigcap x^nR = (0)$에 포함되므로
$\bigcap x^nS = (0)$임은 명백하다(보조정리
\ref{algebra-lemma-intersect-powers-ideal-module-zero}).
따라서 보조정리 \ref{algebra-lemma-finite-over-complete-ring}을 적용하여
$S$가 $R$ 위에서 유한임을 결론내릴 수 있고, 증명이 끝난다.
\end{proof}

\begin{lemma}
\label{algebra-lemma-power-series-over-N-2}
$R$을 환이라 하자. $R$이 뇌터이고 정역이며 N-2이면
$R[[x]]$도 그러하다.
\end{lemma}

\begin{proof}
보조정리 \ref{algebra-lemma-Noetherian-power-series}에 의해
$R[[x]]$가 뇌터임에 주의하라. $R' \supset R$을 분수체에서 $R$의
정수적 폐포라 하자. $R$이 N-2이므로 이는 $R$ 위에서 유한이다.
따라서 $R'[[x]]$은 $R[[x]]$ 위에서 유한이다. 보조정리
\ref{algebra-lemma-power-series-over-Noetherian-normal-domain}에 의해
$R'[[x]]$은 정규 정역이다. 원소 $x \in R'[[x]]$에 보조정리
\ref{algebra-lemma-tate-japanese}를 적용하면 $R'[[x]]$이 N-2임을 얻는다.
그러면 보조정리 \ref{algebra-lemma-finite-extension-N-2}에 의해
$R[[x]]$은 N-2이다.
\end{proof}



\section{Nagata 환}
\label{algebra-section-nagata}

\noindent
정의는 다음과 같다.

\begin{definition}
\label{algebra-definition-nagata}
$R$을 환이라 하자.
\begin{enumerate}
\item $S$가 정역인 모든 유한형 환 사상 $R \to S$에 대해 $S$가
N-2(즉 Japanese)이면 $R$을 {\it 보편적으로 Japanese}라고 한다.
\item $R$이 뇌터이고 모든 소 아이디얼 $\mathfrak p$에 대해
$R/\mathfrak p$가 N-2이면 $R$을 {\it Nagata 환}이라 한다.
\end{enumerate}
\end{definition}

\noindent
뇌터이면서 보편적으로 Japanese인 환이 Nagata 환임은 명백하다.
우리의 목표는 Nagata 환이 보편적으로 Japanese임을 보이는 것이다.
이는 전혀 자명하지 않으며 약간의 작업이 필요하다. 먼저 유용한
보조정리를 하나 제시한다.

\begin{lemma}
\label{algebra-lemma-nagata-in-reduced-finite-type-finite-integral-closure}
$R$을 Nagata 환이라 하자. $R \to S$가 본질적으로 유한형이고
$S$가 축소환이라 하자. 그러면 $S$에서 $R$의 정수적 폐포 $A$는
$R$ 위에서 유한이다.
\end{lemma}

\begin{proof}
$S$가 $R$ 위에서 본질적으로 유한형이므로 뇌터이고, 유한 개의
극소 소 아이디얼 $\mathfrak q_1, \ldots, \mathfrak q_m$을 갖는다.
보조정리 \ref{algebra-lemma-Noetherian-irreducible-components}를 보라.
$S$가 축소이므로 $S \subset \prod S_{\mathfrak q_i}$이고,
각 $S_{\mathfrak q_i} = K_i$는 체이다. 보조정리
\ref{algebra-lemma-total-ring-fractions-no-embedded-points}와
\ref{algebra-lemma-minimal-prime-reduced-ring}을 보라.
각 $K_i$에서 $R$의 정수적 폐포 $A_i'$이 $R$ 위에서
유한임을 보이면 충분하다. $R$이 뇌터이고
$A \subset \prod A_i'$이므로 실제로 충분하다.
$\mathfrak q_i$에 대응하는 $R$의 소 아이디얼을
$\mathfrak p_i \subset R$라 하자. $S$가 $R$ 위에서 본질적으로
유한형이므로 $K_i = S_{\mathfrak q_i} = \kappa(\mathfrak q_i)$는
$\kappa(\mathfrak p_i)$의 유한 생성 체확대이다. 따라서 $K_i$ 안에서
$\kappa(\mathfrak p_i)$의 대수적 폐포 $L_i$는
$\kappa(\mathfrak p_i)$ 위에서 유한이다. Fields, 보조정리
\ref{fields-lemma-algebraic-closure-in-finitely-generated}를 보라.
$A_i'$이 $L_i$에서 $R/\mathfrak p_i$의 정수적 폐포임은 명백하므로,
Nagata 환의 정의에서 결론을 얻는다.
\end{proof}

\begin{lemma}
\label{algebra-lemma-check-universally-japanese}
$R$을 환이라 하자. $R$이 보편적으로 Japanese임을 확인하려면 다음을
보이면 충분하다. $R \to S$가 유한형이고 $S$가 정역이면 $S$는
N-1이다.
\end{lemma}

\begin{proof}
보조정리의 조건을 가정하자. $R \to S$를 $S$가 정역인 유한형
환 사상이라 하자. $L$을 $S$의 분수체의 유한확대라 하자.
$S \subset S' \subset L$인 유한 환 확대가 존재하여 $L$은
$S'$의 분수체이다. 가정에 의해 $S'$은 N-1이고, 따라서 $L$에서
$S'$의 정수적 폐포 $S''$은 $S'$ 위에서 유한이다.
그러므로 $S''$은 $S$ 위에서 유한이고(보조정리
\ref{algebra-lemma-finite-transitive}), $S''$은 $L$에서 $S$의 정수적
폐포이다(보조정리 \ref{algebra-lemma-integral-closure-transitive}).
따라서 $R$은 보편적으로 Japanese이다.
\end{proof}

\begin{lemma}
\label{algebra-lemma-universally-japanese}
$R$이 보편적으로 Japanese이면 $R$ 위에서 본질적으로 유한형인
모든 대수도 보편적으로 Japanese이다.
\end{lemma}

\begin{proof}
$R$ 위의 유한형 대수인 경우는 정의에서 바로 따른다.
일반적인 경우에는 보조정리 \ref{algebra-lemma-localize-N}을 적용하면 된다.
\end{proof}

\begin{lemma}
\label{algebra-lemma-quasi-finite-over-nagata}
$R$을 Nagata 환이라 하자. $R \to S$가 준유한 환 사상
(예를 들어 유한 사상)이면 $S$도 Nagata 환이다.
\end{lemma}

\begin{proof}
먼저 $R$이 뇌터이고 준유한 환 사상은 유한형이므로 $S$가 뇌터임에
주의하라. $\mathfrak q \subset S$를 소 아이디얼이라 하고
$\mathfrak p = R \cap \mathfrak q$로 두자. 그러면
$R/\mathfrak p \subset S/\mathfrak q$는 준유한이고, 보조정리
\ref{algebra-lemma-quasi-finite-over-Noetherian-japanese}에 의해
$S/\mathfrak q$가 N-2임을 얻는다.
\end{proof}

\begin{lemma}
\label{algebra-lemma-nagata-localize}
Nagata 환의 모든 국소화는 Nagata 환이다.
\end{lemma}

\begin{proof}
보조정리 \ref{algebra-lemma-localize-N}에서 명백히 따른다.
\end{proof}

\begin{lemma}
\label{algebra-lemma-nagata-local}
$R$을 환이라 하자. $f_1, \ldots, f_n \in R$이 단위 아이디얼을
생성한다고 하자.
\begin{enumerate}
\item 각 $R_{f_i}$가 보편적으로 Japanese이면 $R$도 그렇다.
\item 각 $R_{f_i}$가 Nagata이면 $R$도 그렇다.
\end{enumerate}
\end{lemma}

\begin{proof}
$S$가 정역인 유한형 환 사상 $\varphi : R \to S$를 택하자.
$\varphi(f_1), \ldots, \varphi(f_n)$은 $S$에서 단위 아이디얼을
생성한다. 따라서 각 $S_{f_i} = S_{\varphi(f_i)}$가 N-1이면
보조정리 \ref{algebra-lemma-Japanese-local}에 의해 $S$도 N-1이다.
이는 (1)을 증명한다.

\medskip\noindent
각 $R_{f_i}$가 Nagata이면 각 $R_{f_i}$는 뇌터이고, 보조정리
\ref{algebra-lemma-cover}에 의해 $R$도 뇌터이다. 또한
$\mathfrak p \subset R$이 소 아이디얼이면 각
$R_{f_i}/\mathfrak pR_{f_i} = (R/\mathfrak p)_{f_i}$는 N-2이고,
보조정리 \ref{algebra-lemma-Japanese-local}에 의해 $R/\mathfrak p$도
N-2이다. 이는 (2)를 증명한다.
\end{proof}

\begin{lemma}
\label{algebra-lemma-Noetherian-complete-local-Nagata}
뇌터 완비 국소환은 Nagata 환이다.
\end{lemma}

\begin{proof}
$R$을 완비 국소 뇌터환이라 하자. $\mathfrak p \subset R$을
소 아이디얼이라 하자. 그러면 보조정리
\ref{algebra-lemma-quotient-complete-local}에 의해 $R/\mathfrak p$도
완비 국소 뇌터환이다. 따라서 뇌터 완비 국소 정역 $R$이
N-2임을 보이면 충분하다. 보조정리
\ref{algebra-lemma-quasi-finite-over-Noetherian-japanese}와
\ref{algebra-lemma-complete-local-Noetherian-domain-finite-over-regular}에 의해
$k$가 체인 $R = k[[X_1, \ldots, X_d]]$의 경우 또는
$\Lambda$가 Cohen 환인
$R = \Lambda[[X_1, \ldots, X_d]]$의 경우로 귀착된다.

\medskip\noindent
$k[[X_1, \ldots, X_d]]$의 경우 보조정리
\ref{algebra-lemma-power-series-over-N-2}에 의해 체가 N-2라는 명제로
귀착되며, 이는 명백하다. $\Lambda[[X_1, \ldots, X_d]]$의 경우
Cohen 환 $\Lambda$가 N-2라는 명제로 귀착된다. 다시
$x = p \in \Lambda$에 보조정리 \ref{algebra-lemma-tate-japanese}를 적용하면
또다시 체의 경우로 귀착된다. 따라서 결론을 얻는다.
\end{proof}

\begin{definition}
\label{algebra-definition-analytically-unramified}
$(R, \mathfrak m)$을 뇌터 국소환이라 하자. 완비화
$R^\wedge = \lim_n R/\mathfrak m^n$가 축소환이면 $R$을
{\it 해석적으로 비분기}라고 한다. 소 아이디얼
$\mathfrak p \subset R$에 대해 $R/\mathfrak p$가 해석적으로
비분기이면 이 소 아이디얼을 {\it 해석적으로 비분기}라고 한다.
\end{definition}

\noindent
이제 모든 뇌터 국소환 $R$에 대해 다음 사실들을 안다.
$R \to R^\wedge$는 충실히 평탄인 국소 환 준동형이다
(보조정리 \ref{algebra-lemma-completion-faithfully-flat}).
완비화 $R^\wedge$는 뇌터이고(보조정리
\ref{algebra-lemma-completion-Noetherian}) 완비이다(보조정리
\ref{algebra-lemma-completion-complete}). 따라서 완비화 $R^\wedge$는
Nagata 환이다(보조정리 \ref{algebra-lemma-Noetherian-complete-local-Nagata}).
더욱이 절 \ref{algebra-section-cohen-structure-theorem}에서 $R^\wedge$가
정칙 국소환의 상환임을 보았고(정리
\ref{algebra-theorem-cohen-structure-theorem}), 따라서 보편적으로
카테너리이다(주의
\ref{algebra-remark-Noetherian-complete-local-ring-universally-catenary}).

\begin{lemma}
\label{algebra-lemma-analytically-unramified-easy}
$(R, \mathfrak m)$을 뇌터 국소환이라 하자.
\begin{enumerate}
\item $R$이 해석적으로 비분기이면 $R$은 축소환이다.
\item $R$이 해석적으로 비분기이면 $R$의 각 극소 소 아이디얼은
해석적으로 비분기이다.
\item $R$이 극소 소 아이디얼
$\mathfrak q_1, \ldots, \mathfrak q_t$을 갖는 축소환이고 각
$\mathfrak q_i$가 해석적으로 비분기이면 $R$은 해석적으로 비분기이다.
\item $R$이 해석적으로 비분기이면 전체 분수환 $Q(R)$에서 $R$의
정수적 폐포는 $R$ 위에서 유한이다.
\item $R$이 정역이고 해석적으로 비분기이면 $R$은 N-1이다.
\end{enumerate}
\end{lemma}

\begin{proof}
이 증명에서는 정의 \ref{algebra-definition-analytically-unramified} 직후의
언급들을 사용한다. $R \to R^\wedge$가 충실히 평탄인 국소 환
준동형이므로 단사이고, (1)이 따른다.

\medskip\noindent
$\mathfrak q$를 $R$의 극소 소 아이디얼이라 하고 $R$이 해석적으로
비분기라고 가정하자. 그러면 $\mathfrak q$는 $R$의 연관 소
아이디얼이다(명제
\ref{algebra-proposition-minimal-primes-associated-primes}를 보라).
따라서 $\{x \in R \mid fx = 0\} = \mathfrak q$인 $f \in R$가
존재한다. 완비화가 완전하므로(보조정리
\ref{algebra-lemma-completion-flat}) 다음이 성립한다.
$(R/\mathfrak q)^\wedge = R^\wedge/\mathfrak q^\wedge$이고
$\{x \in R^\wedge \mid fx = 0\} = \mathfrak q^\wedge$이다.
$x \in R^\wedge$이고 $x^2 \in \mathfrak q^\wedge$이면
$fx^2 = 0$, 따라서 $(fx)^2 = 0$, 따라서 $fx = 0$, 따라서
$x \in \mathfrak q^\wedge$이다. 그러므로 $\mathfrak q$는
해석적으로 비분기이고 (2)가 성립한다.

\medskip\noindent
$R$이 극소 소 아이디얼 $\mathfrak q_1, \ldots, \mathfrak q_t$을
갖는 축소환이고 각 $\mathfrak q_i$가 해석적으로 비분기라고 하자.
그러면 $R \to R/\mathfrak q_1 \times \ldots \times
R/\mathfrak q_t$은 단사이다. 완비화가 완전하므로(보조정리
\ref{algebra-lemma-completion-flat}를 보라)
$$
R^\wedge \subset (R/\mathfrak q_1)^\wedge \times \ldots \times
(R/\mathfrak q_t)^\wedge
$$
이다. 따라서 (3)은 명백하다.

\medskip\noindent
$R$이 해석적으로 비분기라고 하자. $\mathfrak p_1, \ldots,
\mathfrak p_s$를 $R^\wedge$의 극소 소 아이디얼이라 하자. 그러면
$$
Q(R^\wedge) =
R^\wedge_{\mathfrak p_1} \times \ldots \times R^\wedge_{\mathfrak p_s}
$$
이고, $R^\wedge$가 축소환이므로 각
$R^\wedge_{\mathfrak p_i}$는 체이다(보조정리
\ref{algebra-lemma-total-ring-fractions-no-embedded-points}를 보라).
$Q(R^\wedge)$에서 $R^\wedge$의 정수적 폐포 $S$는
$S = S_1 \times \ldots \times S_s$와 같고, 여기서 $S_i$는
$R^\wedge/\mathfrak p_i$의 분수체에서의 정수적 폐포이다.
보조정리 \ref{algebra-lemma-Noetherian-complete-local-Nagata}에 의해
$S_i$는 $R^\wedge/\mathfrak p_i$ 위에서 유한이다. 따라서
$S$는 $R^\wedge$ 위에서 유한이다.
$R'$을 $Q(R)$에서 $R$의 정수적 폐포라 하자.
% P01-E1006
$R \to R^\wedge$가 평탄이므로
$R' \otimes_R R^\wedge \subset Q(R) \otimes_R R^\wedge \subset
Q(R^\wedge)$.
또한 $R' \otimes_R R^\wedge$는 $R^\wedge$ 위에서 정수적이다
(보조정리 \ref{algebra-lemma-base-change-integral}).
따라서 $R' \otimes_R R^\wedge \subset S$는 $R^\wedge$-부분가군이다.
% P01-E1007
$R^\wedge$가 뇌터이므로 이는 유한 $R^\wedge$-가군이다.
따라서 $R' \otimes_R R^\wedge$가 $R^\wedge$-가군으로서
$f_i \otimes 1$들로 생성되도록 하는
$f_1, \ldots, f_n \in R'$을 찾을 수 있다. 충실한 평탄성에 의해
$R'$은 $R$-가군으로서 $f_1, \ldots, f_n$로 생성된다.
이는 (4)를 증명한다.

\medskip\noindent
(5)는 (4)의 특수한 경우이다.
\end{proof}

\begin{lemma}
\label{algebra-lemma-codimension-1-analytically-unramified}
$R$을 뇌터 국소환이라 하자. $\mathfrak p \subset R$을 소
아이디얼이라 하자. 다음을 가정하자.
\begin{enumerate}
\item $R_{\mathfrak p}$는 이산 값매김환이다.
\item $\mathfrak p$는 해석적으로 비분기이다.
\end{enumerate}
그러면 $R^\wedge/\mathfrak pR^\wedge$의 모든 연관 소 아이디얼
$\mathfrak q$에 대해 국소환 $(R^\wedge)_{\mathfrak q}$는
이산 값매김환이다.
\end{lemma}

\begin{proof}
가정 (2)는 $R^\wedge/\mathfrak pR^\wedge$가 축소환임을 말한다.
따라서 $R^\wedge/\mathfrak pR^\wedge$의 연관 소 아이디얼
$\mathfrak q \subset R^\wedge$는 $\mathfrak pR^\wedge$ 위의
극소 소 아이디얼과 같은 자료이다. 특히
$(R^\wedge)_{\mathfrak q}$의 극대 아이디얼은
$\mathfrak p(R^\wedge)_{\mathfrak q}$이다.
$xR_{\mathfrak p} = \mathfrak pR_{\mathfrak p}$인
$x \in R$을 택하자. 위 사실에 의해
$x \in (R^\wedge)_{\mathfrak q}$는 극대 아이디얼을 생성한다.
$R \to R^\wedge$가 충실히 평탄이므로 $x$는
$(R^\wedge)_{\mathfrak q}$의 비영인자이다. 따라서 결론을 얻는다.
\end{proof}

\begin{lemma}
\label{algebra-lemma-criterion-analytically-unramified}
$(R, \mathfrak m)$을 뇌터 국소 정역이라 하자.
$x \in \mathfrak m$라 하자. 다음을 가정하자.
\begin{enumerate}
\item $x \not = 0$.
\item $R/xR$은 매립 소 아이디얼을 갖지 않는다.
\item $R/xR$의 각 연관 소 아이디얼 $\mathfrak p \subset R$에 대해
\begin{enumerate}
\item 국소환 $R_{\mathfrak p}$는 정칙이다.
\item $\mathfrak p$는 해석적으로 비분기이다.
\end{enumerate}
\end{enumerate}
그러면 $R$은 해석적으로 비분기이다.
\end{lemma}

\begin{proof}
$\mathfrak p_1, \ldots, \mathfrak p_t$를 $R$-가군 $R/xR$의
연관 소 아이디얼들이라 하자. $R/xR$이 매립 소 아이디얼을 갖지
않으므로 각 $\mathfrak p_i$는 높이 $1$이고 $(x)$ 위의 극소
소 아이디얼이다. 각 $i$에 대해
$\mathfrak q_{i1}, \ldots, \mathfrak q_{is_i}$를
$R^\wedge$-가군 $R^\wedge/\mathfrak p_iR^\wedge$의 연관
소 아이디얼들이라 하자. 보조정리
\ref{algebra-lemma-codimension-1-analytically-unramified}에 의해
$(R^\wedge)_{\mathfrak q_{ij}}$는 정칙이다.
보조정리 \ref{algebra-lemma-bourbaki}에 의해
$$
\text{Ass}_{R^\wedge}(R^\wedge/xR^\wedge)
=
\bigcup\nolimits_{\mathfrak p \in \text{Ass}_R(R/xR)}
\text{Ass}_{R^\wedge}(R^\wedge/\mathfrak pR^\wedge)
=
\{\mathfrak q_{ij}\}.
$$
이다. $y \in R^\wedge$이고 $y^2 = 0$이라 하자.
$(R^\wedge)_{\mathfrak q_{ij}}$가 정칙이고 따라서 정역이므로
(보조정리 \ref{algebra-lemma-regular-domain}),
$y$는 $(R^\wedge)_{\mathfrak q_{ij}}$에서 영으로 사상된다.
따라서 보조정리 \ref{algebra-lemma-zero-at-ass-zero}에 의해
$y$는 $R^\wedge/xR^\wedge$에서 영으로 사상된다.
그러므로 $y = xy'$이다. $R \to R^\wedge$가 평탄하므로 $x$는
비영인자이고, 따라서 $(y')^2 = 0$이다. 이로부터
$y \in \bigcap x^nR^\wedge = (0)$을 얻는다
(보조정리 \ref{algebra-lemma-intersect-powers-ideal-module-zero}).
\end{proof}

\begin{lemma}
\label{algebra-lemma-local-nagata-domain-analytically-unramified}
$(R, \mathfrak m)$을 국소환이라 하자. $R$이 뇌터이고 정역이며
Nagata이면 $R$은 해석적으로 비분기이다.
\end{lemma}

\begin{proof}
$\dim(R)$에 관해 귀납한다. $\dim(R) = 0$인 경우는 자명하다.
따라서 $\dim(R) = d$라고 하고 차원이 $< d$인 모든 뇌터 Nagata
정역에 대해 보조정리가 성립한다고 가정한다.

\medskip\noindent
$R \subset S$를 $R$의 분수체에서 $R$의 정수적 폐포라 하자.
가정에 의해 $S$는 유한 $R$-가군이다. 보조정리
\ref{algebra-lemma-quasi-finite-over-nagata}에 의해 $S$는 Nagata이다.
보조정리 \ref{algebra-lemma-integral-sub-dim-equal}에 의해
$\dim(R) = \dim(S)$이다. $\mathfrak m_1, \ldots, \mathfrak m_t$를
$S$의 극대 아이디얼들이라 하자. 이들은 모두 $R$의 극대
아이디얼 $\mathfrak m$ 위에 놓인다. 더욱이 $S/\mathfrak mS$가
아르틴이므로 충분히 큰 $n$에 대해
$$
(\mathfrak m_1 \cap \ldots \cap \mathfrak m_t)^n \subset \mathfrak mS
$$
이다. 보조정리 \ref{algebra-lemma-completion-flat}에 의해
$R^\wedge \to S^\wedge$는 단사이고, 중국인의 나머지 정리인
보조정리 \ref{algebra-lemma-chinese-remainder}와 보조정리
\ref{algebra-lemma-change-ideal-completion}을 함께 쓰면
$S^\wedge = \prod S^\wedge_i$이다. 여기서 $S^\wedge_i$는
극대 아이디얼 $\mathfrak m_i$에 관한 $S$의 완비화이다.
따라서 $S_{\mathfrak m_i}$가 해석적으로 비분기임을 보이면 충분하다.
다시 말해 $R$이 뇌터 정규 Nagata 정역인 경우로 귀착되었다.

\medskip\noindent
$R$을 뇌터 정규 국소 Nagata 정역이라 하자. 극대 아이디얼에서
비영원 $x \in \mathfrak m$을 택하자. 보조정리
\ref{algebra-lemma-criterion-analytically-unramified}를 적용하겠다.
성질 (1), (2), (3)(a), (3)(b)를 확인해야 한다.
성질 (1)은 명백하다. 보조정리
\ref{algebra-lemma-normal-domain-intersection-localizations-height-1}에 의해
$R/xR$은 매립 소 아이디얼을 갖지 않으므로 성질 (2)가 성립한다.
같은 보조정리에 의해 $R/xR$의 각 연관 소 아이디얼
$\mathfrak p$는 높이 $1$이다. 따라서 $R_{\mathfrak p}$는
$1$차원 정규 정역이고, 그러므로 정칙이다(보조정리
\ref{algebra-lemma-characterize-dvr}). 따라서 (3)(a)가 성립한다.
마지막으로 귀납가정에 의해 (3)(b)가 성립한다. 실제로
$R/\mathfrak p$는 Nagata이다(보조정리
\ref{algebra-lemma-quasi-finite-over-nagata} 또는 정의에서 직접 따른다).
따라서 $R$은 해석적으로 비분기이다.
\end{proof}

\begin{lemma}
\label{algebra-lemma-local-nagata-and-analytically-unramified}
$(R, \mathfrak m)$을 뇌터 국소환이라 하자. 다음은 동치이다.
% P01-E1008
\begin{enumerate}
\item $R$은 Nagata이다.
\item $R \to S$가 $S$가 정역인 유한 사상이고
$\mathfrak m' \subset S$가 극대이면 국소환
$S_{\mathfrak m'}$은 해석적으로 비분기이다.
\item $(R, \mathfrak m) \to (S, \mathfrak m')$이 $S$가 정역인
유한 국소 준동형사상이면 $S$는 해석적으로 비분기이다.
% P01-E1009
\end{enumerate}
\end{lemma}

\begin{proof}
$R$이 Nagata이고 $R \to S$와 $\mathfrak m' \subset S$가
(2)에서와 같다고 하자. 보조정리
\ref{algebra-lemma-quasi-finite-over-nagata}에 의해 $S$는 Nagata이다.
따라서 국소환 $S_{\mathfrak m'}$은 Nagata이고(보조정리
\ref{algebra-lemma-nagata-localize}), 보조정리
\ref{algebra-lemma-local-nagata-domain-analytically-unramified}에 의해
해석적으로 비분기이다. (2)가 (3)을 함의함은 명백하다.

\medskip\noindent
(3)이 성립한다고 하자. $\mathfrak p \subset R$을 소 아이디얼이라
하고 $L/\kappa(\mathfrak p)$를 체의 유한확대라 하자.
(1)을 증명하려면 $R/\mathfrak p$의 정수적 폐포가
$R/\mathfrak p$ 위에서 유한임을 보여야 한다.
% P01-E1010
$L$을 $\kappa(\mathfrak p)$ 위에서 생성하는
$x_1, \ldots, x_n \in L$을 택하자. 각 $i$에 대해
$P_i(T) = T^{d_i} + a_{i, 1} T^{d_i - 1} + \ldots + a_{i, d_i}$를
$\kappa(\mathfrak p)$ 위에서 $x_i$의 최소다항식이라 하자.
적절한 $f_i \in R$, $f_i \not \in \mathfrak p$에 대해 $x_i$를
$f_i x_i$로 바꾸어 $a_{i, j} \in R/\mathfrak p$라고 가정해도 좋다.
실제로 $\mathfrak m$의 원소들을 더 곱하여 모든 $i, j$에 대해
$a_{i, j} \in \mathfrak m/\mathfrak p \subset R/\mathfrak p$라고
가정해도 좋다. 이렇게 한 뒤
$S = R/\mathfrak p[x_1, \ldots, x_n] \subset L$로 두자.
% P01-E1011
그러면 $S$는 $R$ 위에서 유한이고 정역이며,
$S/\mathfrak m S$는
$R/\mathfrak m[T_1, \ldots, T_n]/(T_1^{d_1}, \ldots, T_n^{d_n})$의
상환이다.
% P01-E1012
따라서 $S$는 국소환이다. (3)에 의해 $S$는 해석적으로 비분기이고,
보조정리 \ref{algebra-lemma-analytically-unramified-easy}에 의해
$L$에서 그 정수적 폐포 $S'$은 $S$ 위에서 유한이다.
$S'$은 $L$에서 $R/\mathfrak p$의 정수적 폐포이기도 하므로
결론을 얻는다.
\end{proof}

\noindent
다음 명제는 특히 Nagata 환 위의 유한형 대수가 Nagata 환임을 말한다.

\begin{proposition}[Nagata]
\label{algebra-proposition-nagata-universally-japanese}
$R$을 환이라 하자. 다음은 동치이다.
\begin{enumerate}
\item $R$은 Nagata 환이다.
\item 모든 유한형 $R$-대수는 Nagata이다.
\item $R$은 보편적으로 Japanese이고 뇌터이다.
\end{enumerate}
\end{proposition}

\begin{proof}
뇌터이고 보편적으로 Japanese인 환이 보편적으로 Nagata임
(즉 조건 (2)가 성립함)은 명백하다. $R$을 Nagata 환이라 하자.
모든 유한 생성 $R$-대수 $S$가 Nagata임을 보이겠다.
그러면 명제가 증명된다.

\medskip\noindent
단계 1. 각 $R_i \to R_{i + 1}$이 한 원소로 생성되는 환 사상들의 열
$R = R_0 \to R_1 \to R_2 \to \ldots \to R_n = S$가 존재한다.
따라서 귀납법에 의해 $S \cong R[x]/I$일 때 $S$가 Nagata임을
증명하면 충분하다.

\medskip\noindent
단계 2. $\mathfrak q \subset S$를 $S$의 소 아이디얼이라 하고
$\mathfrak p \subset R$을 이에 대응하는 $R$의 소 아이디얼이라 하자.
$S/\mathfrak q$가 N-2임을 보여야 한다. 따라서 다음 명제로
귀착되었다.
(*) Nagata 정역 $R$과 정역들의 단일 생성 확대 $R \subset S$가
주어지면 $S$는 N-2이다.

\medskip\noindent
단계 3: $R$을 Nagata 정역이라 하고 $R \subset S$를 정역들의
단일 생성 확대라 하자. $R$과 $S$의 분수체들의 유도된 확대가
순수 초월이라고 하자. 이 경우 $S = R[x]$이다. 보조정리
\ref{algebra-lemma-polynomial-ring-N-2}에 의해 $S$는 N-2이다.
따라서 다음 명제로 귀착되었다.
(**) Nagata 정역 $R$과 분수체들의 유한확대를 유도하는 정역들의
단일 생성 확대 $R \subset S$가 주어지면 $S$는 N-2이다.

\medskip\noindent
단계 4. $R$을 정규 Nagata 정역이라 하고 $R \subset S$를
분수체들의 유한확대 $L/K$를 유도하는 정역들의 단일 생성 확대라 하자.
$R$-대수로서 $S$를 생성하는 원소 $x \in S$를 택하자.
$M/L$을 체의 유한확대라 하자. $R'$을 $M$에서 $R$의 정수적
폐포라 하자. 그러면 $M$에서 $S$의 정수적 폐포 $S'$은
$M$에서 $R'[x]$의 정수적 폐포와 같다. 또한 $R'$의 분수체는
$M$이고 $R \subset R'$은 유한이다($R$의 Nagata 성질에 의한다).
이는 $R'$이 Nagata 환임을 함의한다(보조정리
\ref{algebra-lemma-quasi-finite-over-nagata}). $S'$이 $S$ 위에서
유한임을 보이는 것은 $S'$이 $R'[x]$ 위에서 유한임을 보이는 것과 같다.
$R$을 $R'$으로, $S$를 $R'[x]$으로 바꾸어 다음 명제로 귀착한다.
(***) 분수체가 $K$인 정규 Nagata 정역 $R$과 $x \in K$가
주어질 때, $R$과 $x$가 생성하는 환 $S \subset K$는 N-1이다.

\medskip\noindent
단계 5. $R$을 분수체가 $K$인 정규 Nagata 정역이라 하자.
$x = b/a \in K$라 하자. $R$과 $x$가 생성하는 환
$S \subset K$가 N-1임을 보여야 한다. $S_a \cong R_a$는
정규임에 주의하라. 따라서 보조정리 \ref{algebra-lemma-characterize-N-1}에
의해 $S$의 모든 극대 아이디얼 $\mathfrak m$에 대해
$S_{\mathfrak m}$이 N-1임을 보이면 충분하다.

\medskip\noindent
앞 문단의 가정 아래 그러한 극대 아이디얼을 택하고
$\mathfrak n = R \cap \mathfrak m$으로 두자. 잉여체 확대
$\kappa(\mathfrak m)/\kappa(\mathfrak n)$는 유한하고(정리
\ref{algebra-theorem-nullstellensatz}) $x$의 상으로 생성된다.
어떤 $a \in R$, $a \not \in \mathfrak n$에 대해 $R$을 $R_a$로,
$S$를 $S_a$로 바꾸어 다음을 가정할 수 있다.
% P01-E1013
$R[x]$ 안에 일계수다항식
$f(X) = X^d + \sum_{i = 1, \ldots, d} a_iX^{d - i}$이 존재하여
$f(x) \in \mathfrak m$이다(세부사항은 생략한다).
% P01-E1014 P01-E1015
$K''/K$를 $f(X)$가 $K''[X]$에서 완전히 분해되는 체의 유한확대라
하자. $R'$을 $K''$에서 $R$의 정수적 폐포라 하자.
$S' \subset K''$을 $R'$과 $x$가 생성하는 부분환이라 하자.
$R$이 Nagata이므로 $R'$은 $R$ 위에서 유한이고 Nagata이다
(보조정리 \ref{algebra-lemma-quasi-finite-over-nagata}).
또한 $S'$은 $S$ 위에서 유한이다. $S'$의 모든 극대 아이디얼
$\mathfrak m'$에 대해 국소환 $S'_{\mathfrak m'}$이 N-1이면,
보조정리 \ref{algebra-lemma-characterize-N-1}에 의해
$S'_{\mathfrak m}$은 N-1이고, 다시 보조정리
\ref{algebra-lemma-finite-extension-N-2}에 의해 $S_{\mathfrak m}$은 N-1이다.
$R$을 $R'$으로, $S$를 $S'$으로, $\mathfrak m$을
$\mathfrak m$ 위에 놓이는 임의의 극대 아이디얼 $\mathfrak m'$으로
바꾸면, 위의 다항식 $f$가 완전히 분해되는 상황에 이른다.
$f(X) = \prod_{i = 1, \ldots, d} (X - a_i)$이고 $a_i \in R$이다.
% P01-E1313
$f(x) \in \mathfrak m$이므로 어떤 $i$에 대해
$x - a_i \in \mathfrak m$이다. 마지막으로 $x$를 $x - a_i$로
바꾸어 $x \in \mathfrak m$이라고 가정해도 좋다.

\medskip\noindent
요약하면 다음과 같다. $R$은 분수체가 $K$인 정규 Nagata 정역이고,
$x \in K$이며, $S$는 $K$ 안에서 $x$와 $R$이 생성하는 부분환이다.
끝으로 $\mathfrak m \subset S$는 $x \in \mathfrak m$인 극대
아이디얼이다. $S_{\mathfrak m}$이 N-1임을 보여야 한다.

\medskip\noindent
국소환 $S_{\mathfrak m}$과 원소 $x$에 보조정리
\ref{algebra-lemma-criterion-analytically-unramified}가 적용됨을 보이겠다.
그러면 $S_{\mathfrak m}$은 해석적으로 비분기이고, 이어서 보조정리
\ref{algebra-lemma-analytically-unramified-easy}에 의해 N-1임을 알 수 있다.

\medskip\noindent
성질 (1), (2), (3)(a), (3)(b)를 확인해야 한다.
성질 (1)은 자명하다. $X \mapsto x$인
$I = \Ker(R[X] \to S)$로 두자. $K$에서 $ax = b$인 모든
일차식 $aX - b$가 $I$를 생성한다고 주장한다. 명백히 이 일차식들은
모두 $I$에 속한다. $g = a_d X^d + \ldots a_1 X + a_0 \in I$이면
% P01-E1314
$a_dx$는 $R$ 위에서 정수적이고(보조정리
\ref{algebra-lemma-make-integral-trivial}), $R$이 정규이므로
$b := a_dx \in R$이다. 그러면
$g - (a_dX - b)X^{d - 1} \in I$이고, 차수에 관한 귀납법으로
결론을 얻는다. 그 결과
$$
S/xS = R[X]/(X, I) = R/J
$$
이고, 여기서
$$
J = \{b \in R \mid ax = b \text{ for some }a \in R\} = xR \cap R
$$
이다. 보조정리
\ref{algebra-lemma-normal-domain-intersection-localizations-height-1}에 의해
$S/xS = R/J$는 $R$-가군으로서 매립 소 아이디얼을 갖지 않고,
따라서 $R/J$-가군으로서, 따라서 $S/xS$-가군으로서, 따라서
$S$-가군으로서도 매립 소 아이디얼을 갖지 않는다.
이는 성질 (2)를 증명한다.
$\mathfrak q \subset \mathfrak m$인 그러한 연관 소 아이디얼
$\mathfrak q \subset S$를 택하자(그러면 이는
$S_{\mathfrak m}/xS_{\mathfrak m}$의 연관 소 아이디얼이지만,
논증에는 중요하지 않다). 그러면 $\mathfrak q$는 $xS$ 위에서
극소이고, 따라서 높이가 $1$이다. 위 등식들에 의해
$\mathfrak p = R \cap \mathfrak q$는 $R/J$의 연관 소
아이디얼이고, 따라서 높이가 $1$이다(보조정리
\ref{algebra-lemma-normal-domain-intersection-localizations-height-1}를 보라).
그러므로 $R_{\mathfrak p}$는 이산 값매김환이고,
$R_{\mathfrak p} \subset S_{\mathfrak q}$는 등식이다.
이는 $S_{\mathfrak q}$가 정칙임을 보인다. 따라서 성질 (3)(a)가
성립한다. 마지막으로 $(S/\mathfrak q)_{\mathfrak m}$은
$S/\mathfrak q$의 국소화이고, 이는 $S/xS = R/J$의 상환이다.
따라서 $(S/\mathfrak q)_{\mathfrak m}$은 Nagata 환 $R$의 상환의
국소화이고, 그러므로 Nagata이며(보조정리
\ref{algebra-lemma-quasi-finite-over-nagata}와
\ref{algebra-lemma-nagata-localize}), 따라서 해석적으로 비분기이다
(보조정리 \ref{algebra-lemma-local-nagata-domain-analytically-unramified}).
이는 (3)(b)를 증명하고, 증명이 끝난다.
\end{proof}

\begin{proposition}
\label{algebra-proposition-ubiquity-nagata}
다음 종류의 환들은 Nagata이고, 특히 보편적으로 Japanese이다.
\begin{enumerate}
\item 체.
\item 뇌터 완비 국소환.
\item $\mathbf{Z}$.
\item 분수체의 표수가 영인 데데킨트 정역.
\item 위 환들 가운데 하나의 유한형 환 확대.
\end{enumerate}
\end{proposition}

\begin{proof}
뇌터 완비 국소환의 경우는 보조정리
\ref{algebra-lemma-Noetherian-complete-local-Nagata}이다.
다른 경우에는 환의 모든 소 아이디얼 $\mathfrak p$에 대해
$R/\mathfrak p$가 N-2인지 확인하면 된다. $R/\mathfrak p$가 체인
경우, 즉 $\mathfrak p$가 극대인 경우에는 명백하다.
따라서 데데킨트 환의 경우에는 $\mathfrak p = (0)$일 때만 확인하면
된다. 그러나 분수체의 표수가 영이라고 가정했으므로 보조정리
\ref{algebra-lemma-domain-char-zero-N-1-2}를 적용할 수 있다.
\end{proof}

\begin{example}
\label{algebra-example-nonjapanese-dvr}
이산 값매김환이 Nagata일 필요충분조건은 N-2인 것이다
(극대 아이디얼에 의한 상환이 체이고 따라서 N-2이기 때문이다).
예 \ref{algebra-example-bad-dvr-char-p}의 이산 값매김환 $A$는 Nagata가
아니다. 즉 N-2가 아니다. 실제로 유한확대
$A \subset R = A[f]$는 N-1이 아니다. 이를 보이기 위해
$f = \sum a_i x^i$라 쓰자. 모든 $n \geq 1$에 대해
$g_n = \sum_{i < n} a_i x^i \in A$로 두자.
그러면 $h_n = (f - g_n)/x^n$은 $R$의 분수체의 원소이고
$h_n^p \in k^p[[x]] \subset A$이다. 따라서 $R$의 정수적 폐포
$R'$은 $h_1, h_2, h_3, \ldots$를 포함한다. 이제 $R'$이
$R$ 위에서, 따라서 $A$ 위에서 유한이라면, 모든 $n$에 대해
$f  = x^n h_n + g_n$은 부분가군 $A + x^nR'$에 포함될 것이다.
Artin--Rees에 의해 이는 $f \in A$를 함의하고(보조정리
\ref{algebra-lemma-intersect-powers-ideal-module-zero}), 모순이다.
\end{example}

\begin{lemma}
\label{algebra-lemma-nagata-pth-roots}
$(A, \mathfrak m)$을 Nagata이고 분수체의 표수가 $p$인 뇌터 국소
정역이라 하자. $a \in A$가 $A^\wedge$에서 $p$제곱근을 가지면,
$a$는 $A$에서도 $p$제곱근을 갖는다.
\end{lemma}

\begin{proof}
$\alpha \in A^\wedge$를 $a$의 $p$제곱근이라 하자. 모순을 얻기 위해
$a$가 $A$에서 $p$제곱근을 갖지 않는다고 가정하자. 그러면 $a$는
$A$의 분수체 $K$에서도 $p$제곱근을 갖지 않는다. 실제로
$b, c \in A$, $c \not = 0$인 $\beta = b/c$가
$\beta^p = a$를 만족하면 $A^\wedge$에서 $\alpha = b/c$이다.
이는 $A \to A^\wedge$의 충실한 평탄성에 의해 $A$에서 $c$가
$b$를 나눔을 함의하므로 $\beta \in A$이고, 모순이다.
따라서 $K[x]/(x^p - a)$가 체이므로
$B = A[x]/(x^p - a)$는 정역이다. 그러나 $\alpha$가 존재한다는
가정에 의해 $B$의 완비화는 축소환이 아니다. 한편 $B$는 Nagata
환이고(명제 \ref{algebra-proposition-nagata-universally-japanese}),
따라서 보조정리
\ref{algebra-lemma-local-nagata-domain-analytically-unramified}에 의해
해석적으로 비분기이므로, 이는 앞의 결과들과 모순이다.
\end{proof}

\section{성질의 상승}
\label{algebra-section-ascending-properties}

\noindent
이 절에서는 환의 성질이 ``상승''하는 것에 관한 몇 가지 대수적
사실을 증명하기 시작한다. 환의 깊이에 이를 적용하려면 가군의
깊이 상승에 관한 다음 결과를 사용한다. \cite[IV, Proposition
6.3.1]{EGA}을 보라.

\begin{lemma}
\label{algebra-lemma-apply-grothendieck-module}
\begin{reference}
\cite[IV, Proposition 6.3.1]{EGA}
\end{reference}
다음 등식이 성립한다.
$$
\text{depth}(M \otimes_R N)
=
\text{depth}(M) + \text{depth}(N/\mathfrak m_RN)
$$
여기서 $R \to S$는 뇌터 국소환들의 국소 준동형사상이고,
$M$은 유한 $R$-가군이며, $N$은 $R$ 위에서 평탄인 유한
$S$-가군이다.
\end{lemma}

\begin{proof}
명제와 아래 증명에서 $M$의 깊이는 $R$-가군으로서 취하고,
$M \otimes_R N$의 깊이는 $S$-가군으로서 취하며,
$N/\mathfrak m_RN$의 깊이는 $S/\mathfrak m_RS$-가군으로서 취한다.
% P01-E1016
우변을 $n$이라 하자. 먼저 $n$이 영이라고 가정한다. 그러면
$\text{depth}(M) = 0$이고 $\text{depth}(N/\mathfrak m_RN) = 0$이다.
이는 소멸자가 $\mathfrak m_R$인 $z \in M$과 소멸자가
$\mathfrak m_S/\mathfrak m_RS$인
$\overline{y} \in N/\mathfrak m_RN$이 존재한다는 뜻이다.
$y \in N$을 $\overline{y}$의 올림이라 하자. $N$이 $R$ 위에서
평탄이므로 사상 $z : R/\mathfrak m_R \to M$은 단사사상
$N/\mathfrak m_RN \to M \otimes_R N$을 유도한다. 따라서
$z \otimes y$의 소멸자는 $\mathfrak m_S$이다. 그러므로
$\text{depth}(M \otimes_R N) = 0$이기도 하다.

\medskip\noindent
$n > 0$이라고 가정한다. $\text{depth}(N/\mathfrak m_RN) > 0$이면,
$N/\mathfrak m_RN$ 위의 영인자가 아닌 원소
$\overline{f} \in S/\mathfrak m_RS$로 가는
$f \in \mathfrak m_S$를 택할 수 있다. 그러면 보조정리
\ref{algebra-lemma-depth-drops-by-one}에 의해
$\text{depth}(N/\mathfrak m_RN) =
\text{depth}(N/(f, \mathfrak m_R)N) + 1$이다.
보조정리 \ref{algebra-lemma-mod-injective}에 따라 원소 $f \in S$는
$N$ 위의 영인자가 아니고 $N/fN$은 $R$ 위에서 평탄이다.
따라서 $n$에 관한 귀납법으로
$$
\text{depth}(M \otimes_R N/fN) =
\text{depth}(M) + \text{depth}(N/(f, \mathfrak m_R)N).
$$
$N/fN$이 $R$ 위에서 평탄이므로 다음 열은 완전하다.
$$
0 \to M \otimes_R N \to M \otimes_R N \to M \otimes_R N/fN \to 0
$$
여기서 첫 사상은 $f$를 곱하는 사상이다(보조정리
\ref{algebra-lemma-flat-tor-zero}). 따라서 보조정리
\ref{algebra-lemma-depth-drops-by-one}에 의해
$\text{depth}(M \otimes_R N) = \text{depth}(M \otimes_R N/fN) + 1$이고,
보조정리의 공식에서 등식이 성립한다.

\medskip\noindent
$n > 0$이지만 $\text{depth}(N/\mathfrak m_RN) = 0$이면,
$M$ 위의 영인자가 아닌 원소 $f \in \mathfrak m_R$을 택할 수 있다.
$N$이 $R$ 위에서 평탄이므로 $f$는 $M \otimes_R N$ 위에서도
영인자가 아니다. 다시 $n$에 관한 귀납법으로
$$
\text{depth}(M/fM \otimes_R N) =
\text{depth}(M/fM) + \text{depth}(N/\mathfrak m_RN).
$$
이 경우 보조정리 \ref{algebra-lemma-depth-drops-by-one}에 의해
$$
\text{depth}(M \otimes_R N) = \text{depth}(M/fM \otimes_R N) + 1
$$
이고
$$
\text{depth}(M) = \text{depth}(M/fM) + 1
$$
이므로, 보조정리의
공식에서 등식이 성립한다.
\end{proof}


\begin{lemma}
\label{algebra-lemma-apply-grothendieck}
$R \to S$를 뇌터 국소환들의 평탄한 국소 준동형사상이라 하자.
그러면
$$
\text{depth}(S) = \text{depth}(R) + \text{depth}(S/\mathfrak m_RS).
$$
\end{lemma}

\begin{proof}
이는 보조정리 \ref{algebra-lemma-apply-grothendieck-module}의 특수한 경우이다.
\end{proof}

\begin{lemma}
\label{algebra-lemma-CM-goes-up}
$R \to S$를 뇌터 국소환들의 평탄한 국소 준동형사상이라 하자.
% P01-E1017
그러면 다음 두 조건은 동치이다.
\begin{enumerate}
\item $S$는 Cohen--Macaulay이다.
\item $R$과 $S/\mathfrak m_RS$는 Cohen--Macaulay이다.
\end{enumerate}
\end{lemma}

\begin{proof}
정의와 보조정리 \ref{algebra-lemma-apply-grothendieck} 및
\ref{algebra-lemma-dimension-base-fibre-equals-total}에서 따른다.
\end{proof}

\begin{lemma}
\label{algebra-lemma-Sk-goes-up}
$\varphi : R \to S$를 환 준동형사상이라 하자. 다음을 가정한다.
\begin{enumerate}
\item $R$은 뇌터이다.
\item $S$는 뇌터이다.
\item $\varphi$는 평탄이다.
\item 올환 $S \otimes_R \kappa(\mathfrak p)$들은 $(S_k)$이다.
\item $R$은 성질 $(S_k)$를 갖는다.
\end{enumerate}
그러면 $S$는 성질 $(S_k)$를 갖는다.
\end{lemma}

\begin{proof}
$\mathfrak q$를 $R$의 소 아이디얼 $\mathfrak p$ 위에 놓이는
$S$의 소 아이디얼이라 하자. 보조정리
\ref{algebra-lemma-apply-grothendieck}에 의해
$$
\text{depth}(S_{\mathfrak q}) =
\text{depth}(S_{\mathfrak q}/\mathfrak pS_{\mathfrak q}) +
\text{depth}(R_{\mathfrak p}).
$$
한편 보조정리 \ref{algebra-lemma-dimension-base-fibre-total}에 의해
$$
\dim(R_{\mathfrak p})
+
\dim(S_{\mathfrak q}/\mathfrak pS_{\mathfrak q})
\geq
\dim(S_{\mathfrak q})
$$
이다. (실제로 보조정리
\ref{algebra-lemma-dimension-base-fibre-equals-total}에 의해 등식이 성립하지만,
엄밀히 말해 여기에는 필요하지 않다.) 끝으로 사상의 올환들이
$(S_k)$라고 가정했으므로
$\text{depth}(S_{\mathfrak q}/\mathfrak pS_{\mathfrak q})
\geq \min(k, \dim(S_{\mathfrak q}/\mathfrak pS_{\mathfrak q}))$이다.
따라서 다음 부등식들로 보조정리가 증명된다.
\begin{eqnarray*}
\text{depth}(S_{\mathfrak q}) & = &
\text{depth}(S_{\mathfrak q}/\mathfrak pS_{\mathfrak q}) +
\text{depth}(R_{\mathfrak p}) \\
& \geq &
\min(k, \dim(S_{\mathfrak q}/\mathfrak pS_{\mathfrak q})) +
\min(k, \dim(R_{\mathfrak p})) \\
& = &
\min(2k, \dim(S_{\mathfrak q}/\mathfrak pS_{\mathfrak q}) + k,
k + \dim(R_\mathfrak p),
\dim(S_{\mathfrak q}/\mathfrak pS_{\mathfrak q}) +
\dim(R_{\mathfrak p})) \\
& \geq &
\min(k, \dim(S_{\mathfrak q}))
\end{eqnarray*}
이고, 이것이 원하는 결론이다.
\end{proof}

\begin{lemma}
\label{algebra-lemma-Rk-goes-up}
$\varphi : R \to S$를 환 준동형사상이라 하자. 다음을 가정한다.
\begin{enumerate}
\item $R$은 뇌터이다.
% P01-E1018
\item $S$는 뇌터이다.
\item $\varphi$는 평탄이다.
\item 올환 $S \otimes_R \kappa(\mathfrak p)$들은 성질 $(R_k)$를 갖는다.
\item $R$은 성질 $(R_k)$를 갖는다.
\end{enumerate}
그러면 $S$는 성질 $(R_k)$를 갖는다.
\end{lemma}

\begin{proof}
$\mathfrak q$를 $R$의 소 아이디얼 $\mathfrak p$ 위에 놓이는
$S$의 소 아이디얼이라 하자. $\dim(S_{\mathfrak q}) \leq k$라고
가정한다. 보조정리 \ref{algebra-lemma-dimension-base-fibre-equals-total}에 의해
$\dim(S_{\mathfrak q}) = \dim(R_{\mathfrak p})
+ \dim(S_{\mathfrak q}/\mathfrak pS_{\mathfrak q})$이므로
$\dim(R_{\mathfrak p}) \leq k$이고
$\dim(S_{\mathfrak q}/\mathfrak pS_{\mathfrak q}) \leq k$이다.
따라서 가정에 의해 $R_{\mathfrak p}$와
$S_{\mathfrak q}/\mathfrak pS_{\mathfrak q}$는 정칙이다.
그러므로 보조정리 \ref{algebra-lemma-flat-over-regular-with-regular-fibre}에 의해
$S_{\mathfrak q}$는 정칙이다.
\end{proof}

\begin{lemma}
\label{algebra-lemma-reduced-goes-up-noetherian}
$\varphi : R \to S$를 환 준동형사상이라 하자. 다음을 가정한다.
\begin{enumerate}
\item $R$은 뇌터이다.
% P01-E1019
\item $S$는 뇌터이다.
\item $\varphi$는 평탄이다.
\item 올환 $S \otimes_R \kappa(\mathfrak p)$들은 축소환이다.
\item $R$은 축소환이다.
\end{enumerate}
그러면 $S$는 축소환이다.
\end{lemma}

\begin{proof}
뇌터환에서 축소인 것은 성질 $(S_1)$과 $(R_0)$를 갖는 것과 같다.
보조정리 \ref{algebra-lemma-criterion-reduced}를 보라. 따라서 $R$과 올환들이
이 성질들을 갖는다. 보조정리 \ref{algebra-lemma-Sk-goes-up}과
\ref{algebra-lemma-Rk-goes-up}을 적용하면 $S$가 $(S_1)$과 $(R_0)$임을 알 수
있고, 다시 보조정리 \ref{algebra-lemma-criterion-reduced}에 의해 곧
축소환이다.
\end{proof}

\begin{lemma}
\label{algebra-lemma-reduced-goes-up}
$\varphi : R \to S$를 환 준동형사상이라 하자. 다음을 가정한다.
\begin{enumerate}
\item $\varphi$는 매끄럽다.
\item $R$은 축소환이다.
\end{enumerate}
그러면 $S$는 축소환이다.
\end{lemma}

\begin{proof}
$R \to S$는 정칙인 올환들을 갖는 평탄한 사상이다(절
\ref{algebra-section-smooth-overview}의 매끄러운 환 사상에 관한 결과 목록을
보라). 특히 올환들은 축소환이다. 따라서 $R$이 뇌터이면 $S$도
뇌터이고, 보조정리 \ref{algebra-lemma-reduced-goes-up-noetherian}에서 결론을
얻는다.

\medskip\noindent
일반적인 경우, 유한 생성 $\mathbf{Z}$-부분대수 $R_0 \subset R$과
매끄러운 환 사상 $R_0 \to S_0$ 중에서
$S \cong R \otimes_{R_0} S_0$인 것을 찾을 수 있다. 절
\ref{algebra-section-smooth-overview}의 주의 (10)을 보라. 이제 $x \in S$이고
$x^2 = 0$이면, $R_0$를 확대하여 $x$가 어떤 $x_0 \in S_0$에서
온다고 가정할 수 있다. $R_0$를 한 번 더 확대하여 $S_0$에서
$x_0^2 = 0$이라고 가정할 수 있다.
% P01-E1020
그러나 $R_0 \subset R$은 축소환이므로 $S_0$도 축소환이고,
따라서 원하는 대로 $x_0 = 0$이다.
\end{proof}

\begin{lemma}
\label{algebra-lemma-normal-goes-up-noetherian}
$\varphi : R \to S$를 환 준동형사상이라 하자. 다음을 가정한다.
\begin{enumerate}
\item $R$은 뇌터이다.
\item $S$는 뇌터이다.
\item $\varphi$는 평탄이다.
\item 올환 $S \otimes_R \kappa(\mathfrak p)$들은 정규환이다.
\item $R$은 정규환이다.
\end{enumerate}
그러면 $S$는 정규환이다.
\end{lemma}

\begin{proof}
뇌터환에서 정규인 것은 성질 $(S_2)$와 $(R_1)$을 갖는 것과 같다.
보조정리 \ref{algebra-lemma-criterion-normal}을 보라. 따라서 $R$과 올환들은
이 성질들을 갖는다. 보조정리 \ref{algebra-lemma-Sk-goes-up}과
\ref{algebra-lemma-Rk-goes-up}을 적용하면 $S$가 $(S_2)$와 $(R_1)$임을 알 수
있고, 다시 보조정리 \ref{algebra-lemma-criterion-normal}에 의해 곧
정규환이다.
\end{proof}

\begin{lemma}
\label{algebra-lemma-normal-goes-up}
$\varphi : R \to S$를 환 준동형사상이라 하자. 다음을 가정한다.
\begin{enumerate}
\item $\varphi$는 매끄럽다.
\item $R$은 정규환이다.
\end{enumerate}
그러면 $S$는 정규환이다.
\end{lemma}

\begin{proof}
$R \to S$는 정칙인 올환들을 갖는 평탄한 사상이다(절
\ref{algebra-section-smooth-overview}의 매끄러운 환 사상에 관한 결과 목록을
보라). 특히 올환들은 정규환이다. 따라서 $R$이 뇌터이면 $S$도
뇌터이고, 보조정리 \ref{algebra-lemma-normal-goes-up-noetherian}에서 결론을
얻는다.

\medskip\noindent
일반적인 경우를 보자. 먼저 $R$은 축소환이고, 따라서 보조정리
\ref{algebra-lemma-reduced-goes-up}에 의해 $S$도 축소환이다. $\mathfrak q$를
$S$의 소 아이디얼이라 하고 $\mathfrak p$를 이에 대응하는 $R$의
소 아이디얼이라 하자. $R_{\mathfrak p}$는 정규 정역이다.
$S_{\mathfrak q}$가 정규 정역임을 보여야 한다. 이를 위해 $R$을
$R_{\mathfrak p}$로, $S$를 $S_{\mathfrak p}$로 바꾸어도 된다.
따라서 $R$이 정규 정역이라고 가정할 수 있다.

\medskip\noindent
$R \to S$가 매끄럽고 $R$이 정규 정역이라고 가정한다. 유한 생성
$\mathbf{Z}$-부분대수 $R_0 \subset R$과 매끄러운 환 사상
$R_0 \to S_0$ 중에서 $S \cong R \otimes_{R_0} S_0$인 것을 찾을 수
있다. 절 \ref{algebra-section-smooth-overview}의 주의 (10)을 보라.
$R_0$는 Nagata 정역이므로(명제 \ref{algebra-proposition-ubiquity-nagata}를
보라) 그 정수적 폐포 $R_0'$은 $R_0$ 위에서 유한이다. 더욱이 $R$이
정규 정역이므로 $R_0' \subset R$임은 명백하다. 따라서 $R_0$를
$R_0'$으로, $S_0$를 $R_0' \otimes_{R_0} S_0$로 바꾸어 $R_0$가
뇌터 정규 정역이라고 가정할 수 있다. 증명의 첫 문단에서 $S_0$가
정규환임을 얻는다(물론 정역일 필요는 없다). 이와 같이
$R = \bigcup R_\lambda$는 뇌터 정규 정역들의 합집합이고, 이에
대응하여 $S = \colim R_\lambda \otimes_{R_0} S_0$는 정규환들의
여극한이다. 이는 $S$가 정규환임을 함의한다. 일부 세부사항은
생략한다.
\end{proof}

\begin{lemma}
\label{algebra-lemma-regular-goes-up}
\begin{slogan}
정칙성은 매끄러운 환 사상을 따라 상승한다.
\end{slogan}
$\varphi : R \to S$를 환 준동형사상이라 하자. 다음을 가정한다.
\begin{enumerate}
\item $\varphi$는 매끄럽다.
\item $R$은 정칙환이다.
\end{enumerate}
그러면 $S$는 정칙환이다.
\end{lemma}

\begin{proof}
% P01-E1021
모든 $(R_k)$에 보조정리 \ref{algebra-lemma-Rk-goes-up}을 적용하면 결론을
얻는다. 가정이 성립함은 보조정리
\ref{algebra-lemma-characterize-smooth-over-field}에서 따른다.
\end{proof}

\section{성질의 하강}
\label{algebra-section-descending-properties}

\noindent
이 절에서는 환의 성질이 ``하강''하는 것에 관한 몇 가지 대수적
사실을 증명하기 시작한다. 성질을 상승시키는 것보다 하강시키는 것이
흔히 ``더 쉽다''. 다시 말해, 환 사상 $R \to S$에 관한 가정은 흔히
앞 절의 대응하는 보조정리에 있는 가정들보다 약하다.
% P01-E1022
그러나 대응하는 상승도 증명할 수 없으면 하강 결과들은 흔히
쓸모가 없다는 점을 독자에게 경고한다! 다음은 이 현상을 보여 주는
전형적인 결과이다.

\begin{lemma}
\label{algebra-lemma-descent-Noetherian}
$R \to S$를 환 준동형사상이라 하자. 다음을 가정한다.
\begin{enumerate}
\item $R \to S$는 충실히 평탄이다.
\item $S$는 뇌터이다.
\end{enumerate}
그러면 $R$은 뇌터이다.
\end{lemma}

\begin{proof}
$I_0 \subset I_1 \subset I_2 \subset \ldots$를 $R$의 아이디얼들의
증가열이라 하자. 가정에 의해 어떤 $n$에 대해
$I_nS = I_{n+1}S = I_{n+2}S = \ldots$이다. 충실한 평탄성에 의해
확대하고 축약하면 같은 아이디얼이 된다. 즉 $R$의 각 아이디얼
$I$에 대해 $I = R \cap IS$이다(보조정리
\ref{algebra-lemma-faithfully-flat-universally-injective}). 따라서 원하는 대로
$I_n = I_{n+1} = I_{n+2} = \ldots$이다.
\end{proof}

\begin{lemma}
\label{algebra-lemma-descent-reduced}
$R \to S$를 환 준동형사상이라 하자. 다음을 가정한다.
\begin{enumerate}
\item $R \to S$는 충실히 평탄이다.
\item $S$는 축소환이다.
\end{enumerate}
그러면 $R$은 축소환이다.
\end{lemma}

\begin{proof}
보조정리 \ref{algebra-lemma-faithfully-flat-universally-injective}에 의해
$R \to S$가 단사이므로 명백하다.
\end{proof}

\begin{lemma}
\label{algebra-lemma-descent-normal}
$R \to S$를 환 준동형사상이라 하자. 다음을 가정한다.
\begin{enumerate}
\item $R \to S$는 충실히 평탄이다.
\item $S$는 정규환이다.
\end{enumerate}
그러면 $R$은 정규환이다.
\end{lemma}

\begin{proof}
$S$가 축소환이므로 $R$도 축소환이다. $\mathfrak p$를 $R$의
소 아이디얼이라 하자. $R_{\mathfrak p}$가 정규 정역임을 보여야 한다.
% P01-E1023
$S_{\mathfrak p}$도 $R_{\mathfrak p}$ 위에서 충실히 평탄이므로,
$R$이 극대 아이디얼 $\mathfrak m$을 갖는 국소환이라고 가정할 수 있다.
$\mathfrak q$를 $\mathfrak m$ 위에 놓이는 $S$의 소 아이디얼이라 하자.
그러면 $R \to S_{\mathfrak q}$가 충실히 평탄이다(보조정리
\ref{algebra-lemma-local-flat-ff}). 따라서 $S$도 국소환이라고 가정할 수 있다.
특히 $S$는 정규 정역이다. $R \to S$가 충실히 평탄이고 $S$가
정규 정역이므로 $R$은 정역이다. 다음으로 $a, b \in R$이고
$a/b$가 $R$ 위에서 정수적이라고 하자. $S$가 정규이므로
$a/b \in S$이다. 따라서 $a \in bS$이다. 이는 사상
$a : R \to R/bR$가 $S$로 기초변환한 뒤 영사상이 된다는 뜻이다.
충실한 평탄성에 의해 $a \in bR$이고, 따라서 $a/b \in R$이다.
그러므로 $R$은 정규이다.
\end{proof}

\begin{lemma}
\label{algebra-lemma-descent-regular}
$R \to S$를 환 준동형사상이라 하자. 다음을 가정한다.
\begin{enumerate}
\item $R \to S$는 충실히 평탄이다.
\item $S$는 정칙환이다.
\end{enumerate}
그러면 $R$은 정칙환이다.
\end{lemma}

\begin{proof}
보조정리 \ref{algebra-lemma-descent-Noetherian}에 의해 $R$은 뇌터이다.
$\mathfrak p \subset R$를 소 아이디얼이라 하자. $\mathfrak p$ 위에
놓이는 소 아이디얼 $\mathfrak q \subset S$를 택한다. 그러면
보조정리 \ref{algebra-lemma-flat-under-regular}를
$R_\mathfrak p \to S_\mathfrak q$에 적용할 수 있으므로
$R_\mathfrak p$는 정칙이다. $\mathfrak p$가 임의였으므로 $R$은
정칙이다.
\end{proof}

\begin{lemma}
\label{algebra-lemma-descent-Sk}
$R \to S$를 환 준동형사상이라 하자. 다음을 가정한다.
\begin{enumerate}
\item $R \to S$는 충실히 평탄이다.
\item $S$는 뇌터이고 성질 $(S_k)$를 갖는다.
\end{enumerate}
그러면 $R$은 뇌터이고 성질 $(S_k)$를 갖는다.
\end{lemma}

\begin{proof}
(1)과 (2)가 $R$이 뇌터임을 함의한다는 것은 이미 보조정리
\ref{algebra-lemma-descent-Noetherian}에서 보았다. $\mathfrak p \subset R$를
소 아이디얼이라 하자. 올환 $S \otimes_R \kappa(\mathfrak p)$의
극소 소 아이디얼에 대응하는, $\mathfrak p$ 위에 놓이는 소 아이디얼
$\mathfrak q \subset S$를 택한다. 그러면
$A = R_{\mathfrak p} \to S_{\mathfrak q} = B$는 뇌터 국소환들의
평탄한 국소 준동형사상이고 $\mathfrak m_AB$는 $B$의 정의
아이디얼이다. 따라서 $\dim(A) = \dim(B)$이고(보조정리
\ref{algebra-lemma-dimension-base-fibre-equals-total}),
$\text{depth}(A) = \text{depth}(B)$이다(보조정리
\ref{algebra-lemma-apply-grothendieck}). $B$가 $(S_k)$이므로 $A$도
$(S_k)$이다.
\end{proof}

\begin{lemma}
\label{algebra-lemma-descent-Rk}
$R \to S$를 환 준동형사상이라 하자. 다음을 가정한다.
\begin{enumerate}
\item $R \to S$는 충실히 평탄이다.
\item $S$는 뇌터이고 성질 $(R_k)$를 갖는다.
\end{enumerate}
그러면 $R$은 뇌터이고 성질 $(R_k)$를 갖는다.
\end{lemma}

\begin{proof}
(1)과 (2)가 $R$이 뇌터임을 함의한다는 것은 이미 보조정리
\ref{algebra-lemma-descent-Noetherian}에서 보았다. $\mathfrak p \subset R$를
소 아이디얼이라 하고 $\dim(R_{\mathfrak p}) \leq k$라고 가정하자.
올환 $S \otimes_R \kappa(\mathfrak p)$의 극소 소 아이디얼에
대응하는, $\mathfrak p$ 위에 놓이는 소 아이디얼
$\mathfrak q \subset S$를 택한다. 그러면
$A = R_{\mathfrak p} \to S_{\mathfrak q} = B$는 뇌터 국소환들의
평탄한 국소 준동형사상이고 $\mathfrak m_AB$는 $B$의 정의
아이디얼이다. 따라서 $\dim(A) = \dim(B)$이다(보조정리
\ref{algebra-lemma-dimension-base-fibre-equals-total}). $S$가 $(R_k)$이므로
$B$는 정칙 국소환이다. 보조정리 \ref{algebra-lemma-flat-under-regular}에 의해
$A$도 정칙이다.
\end{proof}

\begin{lemma}
\label{algebra-lemma-descent-nagata}
$R \to S$를 환 준동형사상이라 하자. 다음을 가정한다.
\begin{enumerate}
\item $R \to S$는 매끄럽고 스펙트럼 위에서 전사이다.
\item $S$는 Nagata 환이다.
\end{enumerate}
그러면 $R$은 Nagata 환이다.
\end{lemma}

\begin{proof}
Nagata 환은 뇌터이고 보편적으로 Japanese인 환과 같은 것이다
(명제 \ref{algebra-proposition-nagata-universally-japanese}). 보조정리
\ref{algebra-lemma-descent-Noetherian}에서 이미 $R$이 뇌터임을 보았다.
$R \to A$를 정역으로 가는 유한형 환 준동형사상이라 하자.
보조정리 \ref{algebra-lemma-check-universally-japanese}에 의해 $A$가 N-1임을
확인하면 충분하다. $B = A \otimes_R S$가 유한형 $S$-대수이고
따라서 Nagata임은 명백하다(명제
\ref{algebra-proposition-nagata-universally-japanese}). $A \to B$가
매끄러우므로(보조정리 \ref{algebra-lemma-base-change-smooth}), $B$는
축소환이다(보조정리 \ref{algebra-lemma-reduced-goes-up}). $B$가 뇌터이므로
유한 개의 극소 소 아이디얼 $\mathfrak q_1, \ldots, \mathfrak q_t$만
갖는다(보조정리 \ref{algebra-lemma-Noetherian-irreducible-components}를 보라).
$A \to B$가 평탄이므로 이들은 각각 $(0) \subset A$ 위에 놓인다
(하강 정리에 의한다. 보조정리 \ref{algebra-lemma-flat-going-down}를 보라).
% P01-E1024
그 전체 분수환 $Q(B)$는 $L_i = \kappa(\mathfrak q_i)$들의 곱이다
(보조정리 \ref{algebra-lemma-total-ring-fractions-no-embedded-points}와
\ref{algebra-lemma-minimal-prime-reduced-ring}). 더욱이 $Q(B)$에서 $B$의
정수적 폐포 $B'$은 각 인자 $L_i$에서 $B/\mathfrak q_i$의 정수적
폐포 $B_i'$들의 곱이다(보조정리
\ref{algebra-lemma-characterize-reduced-ring-normal}과 비교하라). $B$가
보편적으로 Japanese이므로 환 확대
$B/\mathfrak q_i \subset B_i'$들은 유한이고, 따라서
$B' = \prod B_i'$은 $B$ 위에서 유한이다. $A \to B$가 평탄이므로
$A$의 영인자가 아닌 원소는 $B$의 영인자가 아닌 원소로 간다.
이에 대응하는 사상
$$
Q(A) \otimes_A B = (A \setminus \{0\})^{-1}A \otimes_A B
= (A \setminus \{0\})^{-1}B \to Q(B)
$$
은 단사이다(보조정리 \ref{algebra-lemma-tensor-localization}을 사용했다).
% P01-E1025
이 사상을 통해 $A'$은 $B'$으로 사상된다. 이는 단사사상
$$
A' \otimes_A B \longrightarrow B'
$$
을 유도한다(위의 사실과 $A \to B$의 평탄성에 의한다).
$B'$은 유한 $B$-가군이고 $B$는 뇌터이므로
$A' \otimes_A B$는 유한 $B$-가군이다. 따라서 유한 개의 원소
$x_i \in A'$가 존재하여 $x_i \otimes 1$들이
$A' \otimes_A B$를 $B$-가군으로 생성한다. 끝으로 $A \to B$의
충실한 평탄성에 의해
% P01-E1026
$x_i$들이 $A'$을 $A$-가군으로도 생성한다고 결론내리고, 증명이
끝난다.
\end{proof}

\begin{remark}
\label{algebra-remark-universally-catenary-does-not-descend}
``보편적으로 카테너리''인 성질은 하강하지 않는다. 에탈 환 사상을
따라서도 하강하지 않는다. 예들의 절
\ref{examples-section-non-catenary-Noetherian-local}에는 유한 환 사상
$A \to B$의 구성이 있다. 여기서 $A$는 국소 뇌터환이지만
보편적으로 카테너리가 아니고, $B$는 두 극대 아이디얼
$\mathfrak m$, $\mathfrak n$을 갖는 반국소환이며,
$B_{\mathfrak m}$과 $B_{\mathfrak n}$은 각각 차원이 $2$와 $1$인
정칙환이고 $A$의 잉여체와 같은 잉여체를 갖는다.
% P01-E1027
더욱이 $\mathfrak m_A$는 $B_{\mathfrak m}$과
$B_{\mathfrak n}$ 모두에서 극대 아이디얼을 생성한다(따라서
$A \to B$는 유한일 뿐 아니라 비분기이다). 보조정리
\ref{algebra-lemma-etale-makes-unramified-closed}에 의해 국소 에탈 환 사상
$A \to A'$이 존재하여 $B \otimes_A A' = B_1 \times B_2$로 분해되고,
$A' \to B_i$는 전사이다. 이는 $A'$이 두 극소 소 아이디얼
$\mathfrak q_i$를 가지며 $A'/\mathfrak q_i \cong B_i$임을 보인다.
$B_i$는 정칙 국소환이므로($B_{\mathfrak m}$ 또는
$B_{\mathfrak n}$ 위에서 에탈이기 때문이다), $A'$은 보편적으로
카테너리라고 결론내린다.
\end{remark}

\section{기하적으로 정규인 대수}
\label{algebra-section-geometrically-normal}

\noindent
이 절에서는 환의 성질의 상승과 하강을 응용한 몇 가지 결과를
제시한다.

\begin{lemma}
\label{algebra-lemma-geometrically-normal}
$k$를 체라 하고 $A$를 $k$-대수라 하자. 다음과 같은 $A$의
성질들은 동치이다.
\begin{enumerate}
\item 모든 체 확대 $k'/k$에 대해 $k' \otimes_k A$는 정규환이다.
\item 모든 유한 생성 체 확대 $k'/k$에 대해 $k' \otimes_k A$는
정규환이다.
\item 모든 유한 순수 비분리 확대 $k'/k$에 대해
$k' \otimes_k A$는 정규환이다.
\item $k^{perf} \otimes_k A$는 정규환이다.
\end{enumerate}
여기서 정규환은 정의 \ref{algebra-definition-ring-normal}에 정의되어 있다.
\end{lemma}

\begin{proof}
(1) $\Rightarrow$ (2) $\Rightarrow$ (3)이고 (1) $\Rightarrow$ (4)임은
명백하다.

\medskip\noindent
$k'/k$가 유한 순수 비분리 확대이면 $k$-확대의 매장
$k' \to k^{perf}$가 존재한다. 환 준동형사상
$k' \otimes_k A \to k^{perf} \otimes_k A$는 충실히 평탄하므로,
$k^{perf} \otimes_k A$가 정규이면 보조정리
\ref{algebra-lemma-descent-normal}에 의해 $k' \otimes_k A$도 정규이다.
이로써 (4) $\Rightarrow$ (3)을 얻는다.

\medskip\noindent
(2)를 가정하고 $k'/k$를 임의의 체 확대라 하자. 그러면 이를
유한 생성 체 확대들의 유향 여극한 $k' = \colim_i k_i$로 쓸 수 있다.
따라서 $k' \otimes_k A = \colim_i k_i \otimes_k A$는 정규환들의
유향 여극한이다. 그러므로 보조정리
\ref{algebra-lemma-colimit-normal-ring}에 의해 $k' \otimes_k A$는
정규환이다. 따라서 (1)이 성립한다.

\medskip\noindent
(3)을 가정하고 $K/k$를 유한 생성 체 확대라 하자. 보조정리
\ref{algebra-lemma-make-separable}에 의해 다음 그림을 찾을 수 있다.
$$
\xymatrix{
K \ar[r] & K' \\
k \ar[u] \ar[r] & k' \ar[u]
}
$$
여기서 $k'/k$, $K'/K$는 유한 순수 비분리 체 확대이고 $K'/k'$은
분리 확대이다. 보조정리 \ref{algebra-lemma-localization-smooth-separable}에
의해 매끄러운 $k'$-대수 $B$가 존재하고, $B$의 분수체는 $K'$이다. 이제
다음과 같이 논증할 수 있다.
단계 1: (3)을 가정했으므로 $k' \otimes_k A$는 정규환이다.
단계 2: $k' \otimes_k A \to B \otimes_{k'} k' \otimes_k A$가
매끄러우므로(보조정리 \ref{algebra-lemma-base-change-smooth})
$B \otimes_{k'} k' \otimes_k A$는 매끄러운 사상에 따른 정규성의
상승에 의해 정규환이다(보조정리 \ref{algebra-lemma-normal-goes-up}).
단계 3. $K' \otimes_{k'} k' \otimes_k A = K' \otimes_k A$는 정규환의
국소화이므로 정규환이다(보조정리
\ref{algebra-lemma-localization-normal-ring}).
단계 4. 끝으로 충실히 평탄인 환 준동형사상
$K \otimes_k A \to K' \otimes_k A$을 따른 정규성의 하강에 의해
$K \otimes_k A$는 정규환이다(보조정리
\ref{algebra-lemma-descent-normal}). 이것으로 보조정리가 증명되었다.
\end{proof}

\begin{definition}
\label{algebra-definition-geometrically-normal}
$k$를 체라 하자. $k$-대수 $R$이 보조정리
\ref{algebra-lemma-geometrically-normal}의 동치조건들을 만족하면 이를
$k$ 위에서 {\it 기하적으로 정규}라고 한다.
\end{definition}

\begin{lemma}
\label{algebra-lemma-localization-geometrically-normal-algebra}
\begin{slogan}
국소화는 기하적 정규성을 보존한다.
\end{slogan}
$k$를 체라 하자. 기하적으로 정규인 $k$-대수의 국소화는
기하적으로 정규이다.
\end{lemma}

\begin{proof}
정규환인지는 소 아이디얼에서의 국소화들에서 확인할 수 있으므로
명백하다.
\end{proof}

\begin{lemma}
\label{algebra-lemma-separable-field-extension-geometrically-normal}
$k$를 체라 하고 $K/k$를 분리 체 확대라 하자. 그러면 $K$는
$k$ 위에서 기하적으로 정규이다.
\end{lemma}

\begin{proof}
$k^{perf} \otimes_k K$가 체이기 때문에 성립한다. 실제로 보조정리
\ref{algebra-lemma-separable-extension-preserves-reducedness}에 의해 이는
축소환이다. 보조정리 \ref{algebra-lemma-perfection}에 의해(또는 정의
\ref{algebra-definition-perfection}에 의해) 체 확대 $k^{perf}/k$는 순수
비분리 확대이다. 따라서 보조정리
\ref{algebra-lemma-radicial-integral-bijective}에 의해 환
$k^{perf} \otimes_k K$는 소 아이디얼을 하나만 갖는다. 소
아이디얼을 하나만 갖는 축소환은 체이다.
\end{proof}

\begin{lemma}
\label{algebra-lemma-geometrically-normal-tensor-normal}
$k$를 체라 하고 $A, B$를 $k$-대수라 하자. $A$는 $k$ 위에서
기하적으로 정규이고 $B$는 정규환이라고 가정하자. 그러면
$A \otimes_k B$는 정규환이다.
\end{lemma}

\begin{proof}
$\mathfrak r$를 $A \otimes_k B$의 소 아이디얼이라 하자.
% P01-E1028
$\mathfrak p$, 그리고 각각 $\mathfrak q$를 이에 대응하는 $A$와
$B$의 소 아이디얼이라 하자. 그러면
$(A \otimes_k B)_{\mathfrak r}$은
$A_{\mathfrak p} \otimes_k B_{\mathfrak q}$의 국소화이다. 따라서
보조정리 \ref{algebra-lemma-localization-normal-ring}과 보조정리
\ref{algebra-lemma-localization-geometrically-normal-algebra}에 의해 환
$A_{\mathfrak p} \otimes_k B_{\mathfrak q}$에 대해 결과를 증명하면
충분하다. 그러므로 $A$와 $B$가 정역이라고 가정할 수 있다.

\medskip\noindent
% P01-E1029
$A$와 $B$가 각각 분수체 $K$와 $L$을 갖는 정역이라고 가정하자.
% P01-E1030
$B$는 유한형 정규 $k$-부분대수들의 여과 여극한이다($k$가 Nagata
환이고, 명제 \ref{algebra-proposition-ubiquity-nagata}를 보라. 따라서
유한형 $k$-부분대수의 정수적 폐포는 명제
\ref{algebra-proposition-nagata-universally-japanese}에 의해 다시 유한형
$k$-부분대수이다). 보조정리 \ref{algebra-lemma-colimit-normal-ring}에 의해
$B$가 $k$ 위에서 유한형인 경우로 귀착된다.

\medskip\noindent
$A$와 $B$가 분수체 $K$와 $L$을 갖는 정역이라고 가정하고,
% P01-E1031
$B$가 $k$ 위에서 유한형이라고 하자. 이 경우 환
$K \otimes_k B$는 $K$ 위에서 유한형이고, 따라서 뇌터이다
(보조정리 \ref{algebra-lemma-Noetherian-permanence}). 특히
$K \otimes_k B$는 유한 개의 극소 소 아이디얼만 갖는다(보조정리
\ref{algebra-lemma-Noetherian-irreducible-components}). $A \to A \otimes_k B$가
평탄이므로, 이는 $A \otimes_k B$도 유한 개의 극소 소 아이디얼만
갖는다는 것을 함의한다(평탄한 환 사상에 대한 하강 정리인 보조정리
\ref{algebra-lemma-flat-going-down}에 의해 이 소 아이디얼들은 모두
$(0) \subset A$ 위에 놓인다). 따라서 $A \otimes_k B$가 그 전체
분수환 안에서 정수적으로 닫혀 있음을 증명하면 충분하다(보조정리
\ref{algebra-lemma-characterize-reduced-ring-normal}).

\medskip\noindent
$K \otimes_k B$와 $A \otimes_k L$이 모두 정규환이라고 주장한다.
이것이 참이면 $A \otimes_k B$ 위에서 정수적인
$Q(A \otimes_k B)$의 임의의 원소 $x$는 보조정리
\ref{algebra-lemma-normal-ring-integrally-closed}에 의해
$K \otimes_k B \cap A \otimes_k L = A \otimes_k B$에 속하므로
증명이 끝난다. 가정에 의해
% P01-E1032
$A \otimes_K L$이 정규환이므로, $K \otimes_k B$가 정규임을
증명하면 충분하다.

\medskip\noindent
$A$가 $k$ 위에서 기하적으로 정규이므로 보조정리
\ref{algebra-lemma-localization-geometrically-normal-algebra}에 의해 $K$는
$k$ 위에서 기하적으로 정규이고, 따라서 $K$는 $k$ 위에서
기하적으로 축소이다.
그러므로 $K = \bigcup K_i$는 $k$의 기하적으로 축소인 유한 생성
체 확대들의 합집합이다(보조정리 \ref{algebra-lemma-subalgebra-separable}).
각 $K_i$는 매끄러운 $k$-대수의 국소화이다(보조정리
\ref{algebra-lemma-localization-smooth-separable}). 따라서
$K_i \otimes_k B$는 매끄러운 $B$-대수의 국소화이고, 그러므로
정규이다(보조정리 \ref{algebra-lemma-normal-goes-up}). 따라서
$K \otimes_k B$는 정규환이고(보조정리
\ref{algebra-lemma-colimit-normal-ring}), 증명이 끝난다.
\end{proof}

\begin{lemma}
\label{algebra-lemma-geometrically-normal-over-separable-algebraic}
$k'/k$를 분리 대수적 체 확대라 하자. $A$를 $k'$ 위의 대수라 하자.
그러면 $A$가 $k$ 위에서 기하적으로 정규일 필요충분조건은 $k'$ 위에서
기하적으로 정규인 것이다.
\end{lemma}

\begin{proof}
$L/k$를 유한 순수 비분리 체 확대라 하자. 그러면
$L' = k' \otimes_k L$은 체이고(체들, 절
\ref{fields-section-algebraic}의 내용을 보라),
$A \otimes_k L = A \otimes_{k'} L'$이다. 따라서 $A$가 $k'$ 위에서
기하적으로 정규이면 $A$는 $k$ 위에서도 기하적으로 정규이다.

\medskip\noindent
$A$가 $k$ 위에서 기하적으로 정규라고 가정하자. $K/k'$을 체
확대라 하자. 그러면
$$
K \otimes_{k'} A = (K \otimes_k A) \otimes_{(k' \otimes_k k')} k'
$$
이다. 보조정리 \ref{algebra-lemma-separable-algebraic-diagonal}에 의해
$k' \otimes_k k' \to k'$이 국소화이므로 $K \otimes_{k'} A$는
정규환의 국소화이고, 따라서 정규이다.
\end{proof}

\section{기하적으로 정칙인 대수}
\label{algebra-section-geometrically-regular}

\noindent
$k$를 체라 하고 $A$를 뇌터 $k$-대수라 하자. $K/k$를 유한 생성
체 확대라 하자. 그러면 환 $K \otimes_k A$도 뇌터이다. 보조정리
\ref{algebra-lemma-Noetherian-field-extension}를 보라. 따라서 다음 보조정리의
진술은 의미가 있다.

\begin{lemma}
\label{algebra-lemma-geometrically-regular}
$k$를 체라 하고 $A$를 $k$-대수라 하자. $A$가 뇌터라고 가정하자.
다음과 같은 $A$의 성질들은 동치이다.
\begin{enumerate}
\item 모든 유한 생성 체 확대 $k'/k$에 대해 $k' \otimes_k A$는
정칙이고,
\item 모든 유한 순수 비분리 확대 $k'/k$에 대해 $k' \otimes_k A$는
정칙이다.
\end{enumerate}
여기서 정칙환은 정의 \ref{algebra-definition-regular}에서와 같은 뜻이다.
\end{lemma}

\begin{proof}
보조정리 앞의 설명에 의해 진술은 의미가 있다.
(1) $\Rightarrow$ (2)는 명백하다.

\medskip\noindent
(2)를 가정하고 $K/k$를 유한 생성 체 확대라 하자. 보조정리
\ref{algebra-lemma-make-separable}에 의해 다음 그림을 찾을 수 있다.
$$
\xymatrix{
K \ar[r] & K' \\
k \ar[u] \ar[r] & k' \ar[u]
}
$$
여기서 $k'/k$, $K'/K$는 유한 순수 비분리 체 확대이고 $K'/k'$은
분리 확대이다. 보조정리 \ref{algebra-lemma-localization-smooth-separable}에
의해 매끄러운 $k'$-대수 $B$가 존재하며 $K'$은 $B$의 분수체이다.
이제 다음과 같이 논증할 수 있다.
단계 1: (2)를 가정했으므로 $k' \otimes_k A$는 정칙환이다.
단계 2: $k' \otimes_k A \to B \otimes_{k'} k' \otimes_k A$가
매끄럽고(보조정리 \ref{algebra-lemma-base-change-smooth}), 매끄러운 사상을
따라 정칙성이 상승하므로(보조정리 \ref{algebra-lemma-regular-goes-up})
$B \otimes_{k'} k' \otimes_k A$는 정칙환이다.
단계 3. $K' \otimes_{k'} k' \otimes_k A = K' \otimes_k A$는 정칙환의
국소화이므로 정칙환이다(정의에서 바로 따른다).
단계 4. 끝으로 충실히 평탄인 환 준동형사상
$K \otimes_k A \to K' \otimes_k A$을 따라 정칙성이 하강하므로
$K \otimes_k A$는 정칙환이다(보조정리
\ref{algebra-lemma-descent-regular}). 이것으로 보조정리가 증명되었다.
\end{proof}

\begin{definition}
\label{algebra-definition-geometrically-regular}
$k$를 체라 하고 $R$을 뇌터 $k$-대수라 하자. $k$-대수 $R$이
보조정리 \ref{algebra-lemma-geometrically-regular}의 동치조건들을 만족하면
이를 $k$ 위에서 {\it 기하적으로 정칙}이라고 한다.
\end{definition}

\noindent
정의로부터, $k$의 임의의 유한 생성 체 확대 $K$에 대해
$K \otimes_k R$이 $K$ 위에서 기하적으로 정칙인 대수임은 명백하다.
뒤에서(더 많은 대수, 명제
\ref{more-algebra-proposition-characterization-geometrically-regular})
$k \subset k' \subset k^{1/p}$가 유한일 때마다
$R \otimes_k k'$이 정칙인지만 확인하면 충분함을 보일 것이다.

\begin{lemma}
\label{algebra-lemma-geometrically-regular-descent}
\begin{slogan}
기하적 정칙성은 충실히 평탄인 대수 사상을 따라 하강한다.
\end{slogan}
$k$를 체라 하고 $A \to B$를 충실히 평탄인 $k$-대수 사상이라 하자.
$B$가 $k$ 위에서 기하적으로 정칙이면 $A$도 그러하다.
\end{lemma}

\begin{proof}
$B$가 $k$ 위에서 기하적으로 정칙이라고 가정하자. $k'/k$를 유한
순수 비분리 확대라 하자. 그러면 $A \to B$의 밑변환인
$A \otimes_k k' \to B \otimes_k k'$은 충실히 평탄이다(보조정리
\ref{algebra-lemma-surjective-spec-radical-ideal}와
\ref{algebra-lemma-flat-base-change}에 의한다). 또한 $B$에 대한 가정에 의해
$B \otimes_k k'$은 $k$ 위에서 정칙이다. 따라서 보조정리
\ref{algebra-lemma-descent-regular}에 의해 $A \otimes_k k'$은 정칙이다.
\end{proof}

\begin{lemma}
\label{algebra-lemma-geometrically-regular-goes-up}
$k$를 체라 하고 $A \to B$를 $k$-대수들의 매끄러운 환 사상이라
하자. $A$가 $k$ 위에서 기하적으로 정칙이면 $B$도 $k$ 위에서
기하적으로 정칙이다.
\end{lemma}

\begin{proof}
$k'/k$를 유한 생성 체 확대라 하자. 그러면
$A \otimes_k k' \to B \otimes_k k'$은 매끄러운 환 사상이고(보조정리
\ref{algebra-lemma-base-change-smooth}), $A \otimes_k k'$은 정칙이다.
따라서 보조정리 \ref{algebra-lemma-regular-goes-up}에 의해
$B \otimes_k k'$은 정칙이다.
\end{proof}

\begin{lemma}
\label{algebra-lemma-geometrically-regular-over-subfields}
$k$를 체라 하고 $A$를 $k$ 위의 대수라 하자. $k = \colim k_i$를
부분체들의 유향 여극한이라 하자. $A$가 각 $k_i$ 위에서
기하적으로 정칙이면 $A$는 $k$ 위에서 기하적으로 정칙이다.
\end{lemma}

\begin{proof}
$k'/k$를 유한 순수 비분리 체 확대라 하자. 유한 개의 변수를
$k$에 첨가하고 유한 개의 다항식 관계를 부과하여 $k'$을 얻을 수
있다. 따라서 어떤 $i$와 유한 순수 비분리 체 확대 $k_i'/k_i$가
존재하여
% P01-E1033
$k_i = k \otimes_{k_i} k_i'$임을 알 수 있다. 그러므로
$A \otimes_k k' = A \otimes_{k_i} k_i'$이고 보조정리는 명백하다.
\end{proof}

\begin{lemma}
\label{algebra-lemma-geometrically-regular-over-separable-algebraic}
$k'/k$를 분리 대수적 체 확대라 하고 $A$를 $k'$ 위의 대수라 하자.
그러면 $A$가 $k$ 위에서 기하적으로 정칙일 필요충분조건은 $k'$ 위에서
기하적으로 정칙인 것이다.
\end{lemma}

\begin{proof}
$L/k$를 유한 순수 비분리 체 확대라 하자. 그러면
$L' = k' \otimes_k L$은 체이고(체들, 절
\ref{fields-section-algebraic}의 내용을 보라),
$A \otimes_k L = A \otimes_{k'} L'$이다. 따라서 $A$가 $k'$ 위에서
기하적으로 정칙이면 $A$는 $k$ 위에서도 기하적으로 정칙이다.

\medskip\noindent
$A$가 $k$ 위에서 기하적으로 정칙이라고 가정하자. $k'$은 $k$의
유한 확대들의 여과 여극한이므로 보조정리
\ref{algebra-lemma-geometrically-regular-over-subfields}에 의해 $k'/k$가
유한 분리 확대인 경우를 가정할 수 있다. 다음 환 사상들을 생각하자.
$$
k' \to A \otimes_k k' \to A.
$$
$A \otimes_k k'$은 $k'$ 위에서 기하적으로 정칙임에 유의하자.
이는 $A$를 $k'$으로 밑변환한 것이다. 또한
$A \otimes_k k' \to A$는 사상 $k' \otimes_k k' \to k'$을
$k' \to A$로 밑변환한 것이다. $k'/k$는 환들의 에탈 확대이므로
$k' \otimes_k k' \to k'$은 에탈이다(보조정리
\ref{algebra-lemma-etale}). 따라서 보조정리
\ref{algebra-lemma-geometrically-regular-goes-up}에 의해 $A$는 $k'$ 위에서
기하적으로 정칙이다.
\end{proof}

\section{기하적으로 Cohen--Macaulay인 대수}
\label{algebra-section-geometrically-CM}

\noindent
이 절의 제목은 다소 부정확하다. Cohen--Macaulay 대수는 자동으로
기하적으로 Cohen--Macaulay이기 때문이다. 아래의 보조정리
\ref{algebra-lemma-extend-field-CM-locus}와 보조정리
\ref{algebra-lemma-CM-geometrically-CM}을 보라.

\begin{lemma}
\label{algebra-lemma-tensor-fields-CM}
$k$를 체라 하고 $K/k$와 $L/k$를 두 체 확대라 하자. 그중 하나가
유한형 체 확대이면 $K \otimes_k L$은 뇌터 Cohen--Macaulay 환이다.
\end{lemma}

\begin{proof}
보조정리 \ref{algebra-lemma-Noetherian-field-extension}에 의해 환
$K \otimes_k L$은 뇌터이다. $K$가 순수 초월 확대
$k(t_1, \ldots, t_r)$의 유한 확대라고 하자. 그러면
$k(t_1, \ldots, t_r) \otimes_k L \to K \otimes_k L$은 유한 자유
환 사상이다. 보조정리 \ref{algebra-lemma-finite-flat-over-regular-CM}에 의해
$k(t_1, \ldots, t_r) \otimes_k L$이 Cohen--Macaulay임을 보이면
충분하다. 이는 다항식환 $L[t_1, \ldots, t_r]$의 국소화이므로
명백하다. 다항식환이 Cohen--Macaulay이라는 사실은 예를 들어
보조정리 \ref{algebra-lemma-CM-polynomial-algebra}를 보라.
\end{proof}

\begin{lemma}
\label{algebra-lemma-CM-geometrically-CM}
$k$를 체라 하고 $S$를 뇌터 $k$-대수라 하자. $K/k$를 유한 생성
체 확대라 하고 $S_K = K \otimes_k S$로 놓자.
$\mathfrak q \subset S$를 $S$의 소 아이디얼이라 하고,
$\mathfrak q$ 위에 놓이는 $S_K$의 소 아이디얼을
$\mathfrak q_K \subset S_K$라 하자. 그러면 $S_{\mathfrak q}$가
Cohen--Macaulay일 필요충분조건은 $(S_K)_{\mathfrak q_K}$가
Cohen--Macaulay인 것이다.
\end{lemma}

\begin{proof}
보조정리 \ref{algebra-lemma-Noetherian-field-extension}에 의해 환 $S_K$는
뇌터이다. 따라서
$S_{\mathfrak q} \to (S_K)_{\mathfrak q_K}$는 뇌터 국소환들의 평탄
국소 준동형사상이다. 그 올
$$
(S_K)_{\mathfrak q_K} / \mathfrak q (S_K)_{\mathfrak q_K}
\cong (\kappa(\mathfrak q) \otimes_k K)_{\mathfrak q'}
$$
은 Cohen--Macaulay 환(보조정리 \ref{algebra-lemma-tensor-fields-CM})
$\kappa(\mathfrak q) \otimes_k K$을 적당한 소 아이디얼
$\mathfrak q'$에서 국소화한 것임에 유의하자. 따라서 보조정리
\ref{algebra-lemma-CM-goes-up}에서 결론이 따른다.
\end{proof}

\section{여극한과 유한 표시 사상, II}
\label{algebra-section-colimits-finite-presentation}

\noindent
이 절은 절 \ref{algebra-section-colimits-flat}의 연속이다.

\medskip\noindent
평탄성의 열린 성질을 응용하는 것으로 시작한다. 이는 평탄 가군을
평탄 가군으로 근사할 수 있다는 유용한 결과이다.

\begin{lemma}
\label{algebra-lemma-flat-finite-presentation-limit-flat}
$R \to S$를 환 사상이라 하고 $M$을 $S$-가군이라 하자. 다음을
가정하자.
\begin{enumerate}
\item $R \to S$는 유한 표시이고,
\item $M$은 유한 표시 $S$-가군이며,
\item $M$은 $R$ 위에서 평탄이다.
\end{enumerate}
이 경우 다음이 성립한다.
\begin{enumerate}
\item 유한형 $\mathbf{Z}$-대수 $R_0$, 유한형 환 사상
$R_0 \to S_0$, 유한 $S_0$-가군 $M_0$가 존재하여 $M_0$는
$R_0$ 위에서 평탄이고, 또한
% P01-E1034
환 사상들 $R_0 \to R$과 $S_0 \to S$ 및 $S_0$-가군 사상
$M_0 \to M$이 존재하여
$S \cong R \otimes_{R_0} S_0$이고 $M = S \otimes_{S_0} M_0$이다.
\item $R = \colim_{\lambda \in \Lambda} R_\lambda$가 유향
여극한으로 주어졌다면, 어떤 $\lambda$와 유한 표시 환 사상
$R_\lambda \to S_\lambda$, 그리고 유한 표시 $S_\lambda$-가군
$M_\lambda$가 존재하여 $M_\lambda$는 $R_\lambda$ 위에서 평탄이고
$S = R \otimes_{R_\lambda} S_\lambda$ 및
$M = S \otimes_{S_{\lambda}} M_\lambda$가 성립한다.
\item
$$
(R \to S, M) =
\colim_{\lambda \in \Lambda}
(R_\lambda \to S_\lambda, M_\lambda)
$$
가 다음을 만족하는 유향 여극한으로 주어졌다고 하자.
\begin{enumerate}
\item $\mu \geq \lambda$일 때
$R_\mu \otimes_{R_\lambda} S_\lambda \to S_\mu$와
$S_\mu \otimes_{S_\lambda} M_\lambda \to M_\mu$는 동형이고,
\item $R_\lambda \to S_\lambda$는 유한 표시이며,
\item $M_\lambda$는 유한 표시 $S_\lambda$-가군이다.
\end{enumerate}
그러면 충분히 큰 모든 $\lambda$에 대해 가군 $M_\lambda$는
$R_\lambda$ 위에서 평탄이다.
\end{enumerate}
\end{lemma}

\begin{proof}
먼저 $(R \to S, M)$을 보조정리
\ref{algebra-lemma-limit-module-finite-presentation}에서와 같은 계
$(R_\lambda \to S_\lambda, M_\lambda)$의 유향 여극한으로 쓰자.
$\mathfrak q \subset S$를 소 아이디얼이라 하자. 이에 대응하는
소 아이디얼들을 각각 $\mathfrak p \subset R$,
$\mathfrak q_\lambda \subset S_\lambda$,
$\mathfrak p_\lambda \subset R_\lambda$라 하자. 정리
\ref{algebra-theorem-openness-flatness}의 증명에서 본 것처럼
$$
((R_\lambda)_{\mathfrak p_\lambda},
(S_\lambda)_{\mathfrak q_\lambda},
(M_\lambda)_{\mathfrak q_{\lambda}})
$$
은 보조정리 \ref{algebra-lemma-limit-module-essentially-finite-presentation}에서와
같은 계이다. 따라서 보조정리 \ref{algebra-lemma-colimit-eventually-flat}에
의해 어떤 $\lambda_{\mathfrak q} \in \Lambda$가 존재하여 모든
$\lambda \geq \lambda_{\mathfrak q}$에 대해 $M_\lambda$는
소 아이디얼 $\mathfrak q_{\lambda}$에서 $R_\lambda$ 위에 평탄이다.

\medskip\noindent
정리 \ref{algebra-theorem-openness-flatness}에 의해 열린 부분집합
$U_\lambda \subset \Spec(S_\lambda)$가 존재하여 $M_\lambda$는
$U_\lambda$의 모든 소 아이디얼에서 $R_\lambda$ 위에 평탄이다.
$\Spec(S) \to \Spec(S_\lambda)$ 아래에서 $U_\lambda$의 역상을
$V_\lambda \subset \Spec(S)$라 하자. 위 논증은 모든
$\mathfrak q \in \Spec(S)$에 대해 어떤
$\lambda_{\mathfrak q}$가 존재하여 모든
$\lambda \geq \lambda_{\mathfrak q}$에 대해
$\mathfrak q \in V_\lambda$임을 보인다. $\Spec(S)$가 준콤팩트이므로
이는 하나의 $\lambda_0 \in \Lambda$가 존재하여
$V_{\lambda_0} = \Spec(S)$임을 함의한다.

\medskip\noindent
여집합 $\Spec(S_{\lambda_0}) \setminus U_{\lambda_0}$는 어떤
아이디얼 $I \subset S_{\lambda_0}$에 대한 $V(I)$이다.
$V_{\lambda_0} = \Spec(S)$이므로 $IS = S$이다.
% P01-E1035
$f_1, \ldots, f_r \in I$와 $s_1, \ldots, s_n \in S$를
$\sum f_i s_i = 1$이 되도록 택하자. $\colim S_\lambda = S$이므로
$\lambda_0$를 증가시킨 뒤
$s_{i, \lambda_0} \in S_{\lambda_0}$가 존재하여
$\sum f_i s_{i, \lambda_0} = 1$이라고 가정할 수 있다. 따라서 이
$\lambda_0$에 대해 $U_{\lambda_0} = \Spec(S_{\lambda_0})$이다.
이것으로 (1)이 증명되었다.

\medskip\noindent
(2)의 증명. $(R_0 \to S_0, M_0)$를 (1)에서와 같이 택하고
$R = \colim R_\lambda$라고 하자. $R_0$는 유한형
$\mathbf{Z}$-대수이므로 어떤 $\lambda$와 사상
$R_0 \to R_\lambda$가 존재하여 $R_0 \to R_\lambda \to R$은 주어진
사상 $R_0 \to R$이다(보조정리
\ref{algebra-lemma-characterize-finite-presentation}를 보라). 이제
$S_\lambda = R_\lambda \otimes_{R_0} S_0$ 및
$M_\lambda = S_\lambda \otimes_{S_0} M_0$로 택하면 (2)가 따른다.

\medskip\noindent
끝으로 (3)을 증명하자. $(R_\lambda \to S_\lambda, M_\lambda)$를
(3)에서와 같이 택하고, $(R_0 \to S_0, M_0)$와 $R_0 \to R$을
(1)에서와 같이 택하자. (2)의 증명에서처럼 어떤 $\lambda_0$와
환 사상 $R_0 \to R_{\lambda_0}$가 존재하여
$R_0 \to R_{\lambda_0} \to R$은 주어진 사상 $R_0 \to R$이다.
$S_0$는 $R_0$ 위에서 유한 표시이고 $S = \colim S_\lambda$이므로,
어떤 $\lambda_1 \geq \lambda_0$에 대해 $R_0$-대수 사상
$S_0 \to S_{\lambda_1}$을 얻어서 합성
$S_0 \to S_{\lambda_1} \to S$가 주어진 사상 $S_0 \to S$가 되게
할 수 있다(보조정리 \ref{algebra-lemma-characterize-finite-presentation}를
보라). 모든 $\lambda \geq \lambda_1$에 대해 이는 사상
$$
\Psi_{\lambda} :
R_\lambda \otimes_{R_0} S_0
\longrightarrow
R_\lambda \otimes_{R_{\lambda_1}} S_{\lambda_1}
\cong
S_\lambda
$$
을 준다. 마지막 동형은 가정에서 따른다. 구성에 의해
$\colim_\lambda \Psi_\lambda$는 동형이다. 따라서 보조정리
\ref{algebra-lemma-colimit-category-fp-algebras}에 의해 충분히 큰 모든
$\lambda$에 대해 $\Psi_\lambda$는 동형이다. 같은 방식으로 어떤
$\lambda_2 \geq \lambda_1$와 $S_0$-가군 사상
$M_0 \to M_{\lambda_2}$가 존재하여
$M_0 \to M_{\lambda_2} \to M$이 주어진 사상 $M_0 \to M$이 되게
할 수 있다(보조정리 \ref{algebra-lemma-module-map-property-in-colimit}를 보라).
$\lambda \geq \lambda_2$에 대해 유도된 사상
$$
S_\lambda \otimes_{S_0} M_0
\longrightarrow
S_\lambda \otimes_{S_{\lambda_2}} M_{\lambda_2}
\cong
M_\lambda
$$
이 있고, 보조정리 \ref{algebra-lemma-colimit-category-fp-modules}에 의해
충분히 큰 $\lambda$에 대해 이 사상은 동형이다. $M_0$가 $R_0$
위에서 평탄이므로 이는 (3)을 함의한다.
\end{proof}

\begin{lemma}
\label{algebra-lemma-descend-faithfully-flat-finite-presentation}
$R \to A \to B$를 환 사상들이라 하자. $A \to B$가 유한 표시인
충실히 평탄 사상이라고 가정하자. 그러면 가환 그림
$$
\xymatrix{
R \ar[r] \ar@{=}[d] &
A_0 \ar[d] \ar[r] &
B_0 \ar[d] \\
R \ar[r] & A \ar[r] & B
}
$$
이 존재하며, $R \to A_0$는 유한 표시이고
$A_0 \to B_0$는 유한 표시인 충실히 평탄 사상이며
$B = A \otimes_{A_0} B_0$이다.
\end{lemma}

\begin{proof}
% P01-E1036
먼저 $R$을 $\mathbf{Z}$로 바꾼 경우의 보조정리를 증명하자.
보조정리 \ref{algebra-lemma-flat-finite-presentation-limit-flat}에 의해 그림
$$
\xymatrix{
A_0 \ar[r] \ar[d] & A \ar[d] \\
B_0 \ar[r] & B
}
$$
이 존재한다. 여기서 $A_0$는 $\mathbf{Z}$ 위에서 유한형이고,
$B_0$는 $A_0$ 위에서 평탄하고 유한 표시이며
$B = A \otimes_{A_0} B_0$이다. $A_0 \to B_0$는 평탄하고 유한
표시이므로 $\Spec(B_0) \to \Spec(A_0)$의 상은 열린집합이다. 명제
\ref{algebra-proposition-fppf-open}을 보라. 따라서 그 상의 여집합은 어떤
아이디얼 $I_0 \subset A_0$에 대한 $V(I_0)$이다. $A \to B$가 충실히
평탄이므로 사상 $\Spec(B) \to \Spec(A)$은 전사이다. 보조정리
\ref{algebra-lemma-ff-rings}를 보라. 이제 밑변환의 상은 상의 밑변환이라는
사실을 쓴다. 따라서 $I_0A = A$이다. $r_i \in A$,
$f_i \in I_0$인 관계식 $\sum f_i r_i = 1$을 택하자. 원소
$r_i$들을 포함하도록 $A_0$를 확대하고(이에 맞추어 $B_0$도 확대한)
뒤에는 $A_0 \to B_0$가 스펙트럼 위에서도 전사, 즉 충실히 평탄임을
알 수 있다.

\medskip\noindent
따라서 $R = \mathbf{Z}$인 경우 보조정리가 성립한다. 일반적인 경우,
방금 얻은 해 $A_0' \to B_0'$을 택하고
$A_0 = A_0' \otimes_{\mathbf{Z}} R$,
$B_0 = B_0' \otimes_{\mathbf{Z}} R$로 놓는다.
\end{proof}

\begin{lemma}
\label{algebra-lemma-colimit-finite}
% P01-E1037
$A = \colim_{i \in I} A_i$를 환들의 유향 여극한이라 하자.
$0 \in I$라 하고 $\varphi_0 : B_0 \to C_0$를 $A_0$-대수들의
사상이라 하자. 다음을 가정하자.
\begin{enumerate}
\item $A \otimes_{A_0} B_0 \to A \otimes_{A_0} C_0$는 유한 사상이고,
\item $C_0$는 $B_0$ 위에서 유한형이다.
\end{enumerate}
그러면 어떤 $i \geq 0$가 존재하여 사상
$A_i \otimes_{A_0} B_0 \to A_i \otimes_{A_0} C_0$는 유한이다.
\end{lemma}

\begin{proof}
$x_1, \ldots, x_m$을 $B_0$ 위의 $C_0$의 생성원들이라 하자.
$A \otimes_{A_0} C_0$에서 $P_j(1 \otimes x_j) = 0$이 되도록
일계수 다항식 $P_j \in A \otimes_{A_0} B_0[T]$를 택하자.
어떤 $i \geq 0$에 대해 $P_j$로 가는
$P_{j, i} \in A_i \otimes_{A_0} B_0[T]$를 찾을 수 있다.
$\otimes$는 여극한과 가환하므로, 필요하면 $i$를 증가시킨 뒤
$P_{j, i}(1 \otimes x_j)$가 $A_i \otimes_{A_0} C_0$에서 영임을
알 수 있다. 이 $i$가 원하는 것이다.
\end{proof}

\begin{lemma}
\label{algebra-lemma-colimit-surjective}
% P01-E1037
$A = \colim_{i \in I} A_i$를 환들의 유향 여극한이라 하자.
$0 \in I$라 하고 $\varphi_0 : B_0 \to C_0$를 $A_0$-대수들의
사상이라 하자. 다음을 가정하자.
\begin{enumerate}
\item $A \otimes_{A_0} B_0 \to A \otimes_{A_0} C_0$는 전사이고,
\item $C_0$는 $B_0$ 위에서 유한형이다.
\end{enumerate}
그러면 어떤 $i \geq 0$에 대해 사상
$A_i \otimes_{A_0} B_0 \to A_i \otimes_{A_0} C_0$는 전사이다.
\end{lemma}

\begin{proof}
$x_1, \ldots, x_m$을 $B_0$ 위의 $C_0$의 생성원들이라 하자.
$A \otimes_{A_0} C_0$에서 $1 \otimes x_j$로 가는
$b_j \in A \otimes_{A_0} B_0$를 택하자. 어떤 $i \geq 0$에 대해
$b_j$로 가는 $b_{j, i} \in A_i \otimes_{A_0} B_0$를 찾을 수 있다.
$i$를 증가시킨 뒤 모든 $j = 1, \ldots, m$에 대해 $b_{j, i}$가
$A_i \otimes_{A_0} C_0$에서 $1 \otimes x_j$로 간다고 가정할 수
있다. 이 $i$가 원하는 것이다.
\end{proof}

\begin{lemma}
\label{algebra-lemma-colimit-unramified}
% P01-E1037
$A = \colim_{i \in I} A_i$를 환들의 유향 여극한이라 하자.
$0 \in I$라 하고 $\varphi_0 : B_0 \to C_0$를 $A_0$-대수들의
사상이라 하자. 다음을 가정하자.
\begin{enumerate}
\item $A \otimes_{A_0} B_0 \to A \otimes_{A_0} C_0$는 비분기이고,
\item $C_0$는 $B_0$ 위에서 유한형이다.
\end{enumerate}
그러면 어떤 $i \geq 0$에 대해 사상
$A_i \otimes_{A_0} B_0 \to A_i \otimes_{A_0} C_0$는 비분기이다.
\end{lemma}

\begin{proof}
$B_i = A_i \otimes_{A_0} B_0$, $C_i = A_i \otimes_{A_0} C_0$,
$B = A \otimes_{A_0} B_0$, $C = A \otimes_{A_0} C_0$로 놓자.
$x_1, \ldots, x_m$을 $B_0$ 위의 $C_0$의 생성원들이라 하자.
그러면 $\text{d}x_1, \ldots, \text{d}x_m$은 $C_0$ 위의
$\Omega_{C_0/B_0}$를 생성하고, 그 상들은 $C_i$ 위의
$\Omega_{C_i/B_i}$를 생성한다(보조정리
\ref{algebra-lemma-differentials-polynomial-ring}과
\ref{algebra-lemma-differential-seq}). 또한
$0 = \Omega_{C/B} = \colim \Omega_{C_i/B_i}$이다(보조정리
\ref{algebra-lemma-colimit-differentials}). 따라서 어떤 $i$가 존재하여
$\text{d}x_1, \ldots, \text{d}x_m$은 영으로 가고, 따라서 원한 대로
$\Omega_{C_i/B_i} = 0$이다.
\end{proof}

\begin{lemma}
\label{algebra-lemma-colimit-isomorphism}
% P01-E1037
$A = \colim_{i \in I} A_i$를 환들의 유향 여극한이라 하자.
$0 \in I$라 하고 $\varphi_0 : B_0 \to C_0$를 $A_0$-대수들의
사상이라 하자. 다음을 가정하자.
\begin{enumerate}
\item $A \otimes_{A_0} B_0 \to A \otimes_{A_0} C_0$는 동형이고,
\item $B_0 \to C_0$는 유한 표시이다.
\end{enumerate}
그러면 어떤 $i \geq 0$에 대해 사상
$A_i \otimes_{A_0} B_0 \to A_i \otimes_{A_0} C_0$는 동형이다.
\end{lemma}

\begin{proof}
보조정리 \ref{algebra-lemma-colimit-surjective}에 의해 어떤 $i$가 존재하여
$A_i \otimes_{A_0} B_0 \to A_i \otimes_{A_0} C_0$는 전사이다.
이 사상은 유한 표시이므로 그 핵은 유한 생성 아이디얼이다.
$g_1, \ldots, g_r \in A_i \otimes_{A_0} B_0$가 그 핵을 생성한다고
하자. 그러면 $g_j$들이 $A_{i'} \otimes_{A_0} B_0$에서 영으로 가도록
어떤 $i' \geq i$를 택할 수 있다. 이때
$A_{i'} \otimes_{A_0} B_0 \to A_{i'} \otimes_{A_0} C_0$는 동형이다.
\end{proof}

\begin{lemma}
\label{algebra-lemma-colimit-etale}
% P01-E1037
$A = \colim_{i \in I} A_i$를 환들의 유향 여극한이라 하자.
$0 \in I$라 하고 $\varphi_0 : B_0 \to C_0$를 $A_0$-대수들의
사상이라 하자. 다음을 가정하자.
\begin{enumerate}
\item $A \otimes_{A_0} B_0 \to A \otimes_{A_0} C_0$는 에탈이고,
\item $B_0 \to C_0$는 유한 표시이다.
\end{enumerate}
그러면 어떤 $i \geq 0$에 대해 사상
$A_i \otimes_{A_0} B_0 \to A_i \otimes_{A_0} C_0$는 에탈이다.
\end{lemma}

\begin{proof}
$C_0 = B_0[x_1, \ldots, x_n]/(f_{1, 0}, \ldots, f_{m, 0})$로 쓰자.
$B_i = A_i \otimes_{A_0} B_0$ 및 $C_i = A_i \otimes_{A_0} C_0$로
쓰자. 여기서 $f_{j, i}$가 $B_i$ 위 다항식환에서 $f_{j, 0}$의
상일 때
$C_i = B_i[x_1, \ldots, x_n]/(f_{1, i}, \ldots, f_{m, i})$임에
유의하자. $B = A \otimes_{A_0} B_0$ 및
$C = A \otimes_{A_0} C_0$로 쓰자. 여기서 $f_j$가 $B$ 위 다항식환에서
$f_{j, 0}$의 상일 때
$C = B[x_1, \ldots, x_n]/(f_1, \ldots, f_m)$임에 유의하자.
가정은 사상
$$
\text{d} :
(f_1, \ldots, f_m)/(f_1, \ldots, f_m)^2
\longrightarrow
\bigoplus C \text{d}x_k
$$
이 동형이라는 것이다. 따라서 충분히 큰 $i$에 대해 원소
$$
\xi_{k, i} \in (f_{1, i}, \ldots, f_{m, i})/(f_{1, i}, \ldots, f_{m, i})^2
$$
를 찾아서 $\bigoplus C_i \text{d}x_k$에서
$\text{d}\xi_{k, i} = \text{d}x_k$가 되게 할 수 있다. 또한 필요하면
$i$를 증가시켜
$\sum (\partial f_{j, i}/\partial x_k) \xi_{k, i} =
f_{j, i} \bmod (f_{1, i}, \ldots, f_{m, i})^2$라고 할 수 있다.
이는 여극한에서 참이기 때문이다. 이 $i$가 원하는 것이다.
\end{proof}

\begin{lemma}
\label{algebra-lemma-colimit-smooth}
% P01-E1037
$A = \colim_{i \in I} A_i$를 환들의 유향 여극한이라 하자.
$0 \in I$라 하고 $\varphi_0 : B_0 \to C_0$를 $A_0$-대수들의
사상이라 하자. 다음을 가정하자.
\begin{enumerate}
\item $A \otimes_{A_0} B_0 \to A \otimes_{A_0} C_0$는 매끄럽고,
\item $B_0 \to C_0$는 유한 표시이다.
\end{enumerate}
그러면 어떤 $i \geq 0$에 대해 사상
$A_i \otimes_{A_0} B_0 \to A_i \otimes_{A_0} C_0$는 매끄럽다.
\end{lemma}

\begin{proof}
$C_0 = B_0[x_1, \ldots, x_n]/(f_{1, 0}, \ldots, f_{m, 0})$로 쓰자.
$B_i = A_i \otimes_{A_0} B_0$ 및 $C_i = A_i \otimes_{A_0} C_0$로
쓰자. 여기서 $f_{j, i}$가 $B_i$ 위 다항식환에서 $f_{j, 0}$의
상일 때
$C_i = B_i[x_1, \ldots, x_n]/(f_{1, i}, \ldots, f_{m, i})$임에
유의하자. $B = A \otimes_{A_0} B_0$ 및
$C = A \otimes_{A_0} C_0$로 쓰자. 여기서 $f_j$가 $B$ 위 다항식환에서
$f_{j, 0}$의 상일 때
$C = B[x_1, \ldots, x_n]/(f_1, \ldots, f_m)$임에 유의하자.
가정은 사상
$$
\text{d} :
(f_1, \ldots, f_m)/(f_1, \ldots, f_m)^2
\longrightarrow
\bigoplus C \text{d}x_k
$$
이 분할 단사라는 것이다. 원소
$\xi_k \in (f_1, \ldots, f_m)/(f_1, \ldots, f_m)^2$를
$\sum (\partial f_j/\partial x_k) \xi_k =
f_j \bmod (f_1, \ldots, f_m)^2$가 되도록 택하자. 그러면 충분히 큰
$i$에 대해 원소
$$
\xi_{k, i} \in (f_{1, i}, \ldots, f_{m, i})/(f_{1, i}, \ldots, f_{m, i})^2
$$
를 찾아서
$\sum (\partial f_{j, i}/\partial x_k) \xi_{k, i} =
f_{j, i} \bmod (f_{1, i}, \ldots, f_{m, i})^2$가 되게 할 수 있다.
이는 여극한에서 참이기 때문이다. 이 $i$가 원하는 것이다.
\end{proof}

\begin{lemma}
\label{algebra-lemma-colimit-lci}
% P01-E1037
$A = \colim_{i \in I} A_i$를 환들의 유향 여극한이라 하자.
$0 \in I$라 하고 $\varphi_0 : B_0 \to C_0$를 $A_0$-대수들의
사상이라 하자. 다음을 가정하자.
\begin{enumerate}
\item $A \otimes_{A_0} B_0 \to A \otimes_{A_0} C_0$는 신토믹
(각각, 상대 전역 완전 교차)이고,
\item $C_0$는 $B_0$ 위에서 유한 표시이다.
\end{enumerate}
그러면 어떤 $i \geq 0$가 존재하여 사상
$A_i \otimes_{A_0} B_0 \to A_i \otimes_{A_0} C_0$는 신토믹
(각각, 상대 전역 완전 교차)이다.
\end{lemma}

\begin{proof}
$A \otimes_{A_0} B_0 \to A \otimes_{A_0} C_0$가 상대 전역 완전
교차라고 가정하자. 보조정리
\ref{algebra-lemma-relative-global-complete-intersection-Noetherian}에 의해
유한형 $\mathbf{Z}$-대수 $R$, 환 사상
$R \to A \otimes_{A_0} B_0$, 상대 전역 완전 교차 $R \to S$, 그리고
동형
$$
(A \otimes_{A_0} B_0) \otimes_R S
\longrightarrow
A \otimes_{A_0} C_0
$$
이 존재한다. $R$은 $\mathbf{Z}$ 위에서 유한형(따라서 유한 표시)이므로,
어떤 $i$와 사상 $R \to A_i \otimes_{A_0} B_0$가 존재하여 사상
$R \to A \otimes_{A_0} B_0$를 올린다. 보조정리
\ref{algebra-lemma-characterize-finite-presentation}를 보라. 같은 보조정리를
사용하면 어떤 $i' \geq i$가 존재하여
$(A_i \otimes_{A_0} B_0) \otimes_R S \to A \otimes_{A_0} C_0$가 사상
$(A_i \otimes_{A_0} B_0) \otimes_R S \to A_{i'} \otimes_{A_0} C_0$에서
온다. 따라서 $i$를 $i'$으로 바꾼 뒤, 표시된 사상이
$A_i \otimes_{A_0} B_0$-대수 사상
$$
(A_i \otimes_{A_0} B_0) \otimes_R S
\longrightarrow
A_i \otimes_{A_0} C_0
$$
에서 온다고 가정할 수 있다. 보조정리
\ref{algebra-lemma-colimit-isomorphism}에 의해 $i$를 증가시키면 이 사상은
동형이다. 상대 전역 완전 교차의 밑변환은 보조정리
\ref{algebra-lemma-base-change-relative-global-complete-intersection}에 의해
다시 상대 전역 완전 교차이므로, 이 경우의 증명이 끝난다.

\medskip\noindent
$A \otimes_{A_0} B_0 \to A \otimes_{A_0} C_0$가 신토믹이라고
가정하자. 단위원 아이디얼을 생성하는 원소들
$g_1, \ldots, g_m$이 $A \otimes_{A_0} C_0$에 존재하여
$A \otimes_{A_0} B_0 \to (A \otimes_{A_0} C_0)_{g_j}$는 상대 전역
완전 교차이다. 보조정리 \ref{algebra-lemma-syntomic}을 보라. 어떤 $i$와
$g_j$로 가는 원소들 $g_{i, j} \in A_i \otimes_{A_0} C_0$를 찾을 수
있다. $i$를 증가시킨 뒤 $g_{i, 1}, \ldots, g_{i, m}$이
$A_i \otimes_{A_0} C_0$의 단위원 아이디얼을 생성한다고 가정할 수
있다. 앞 문단의 결과에 의해 $i$를 증가시킨 뒤 사상들
$A_i \otimes_{A_0} B_0 \to (A_i \otimes_{A_0} C_0)_{g_{i, j}}$가
상대 전역 완전 교차라고 가정할 수 있다. 그러면 보조정리
\ref{algebra-lemma-local-syntomic}에 의해(그리고 이미 사용한 보조정리
\ref{algebra-lemma-syntomic}에 의해)
$A_i \otimes_{A_0} B_0 \to A_i \otimes_{A_0} C_0$는 신토믹이다.
\end{proof}

\noindent
다음 보조정리는 위 결과들의 응용인데, 다른 어디에도 잘 들어맞지
않는 듯하다.

\begin{lemma}
\label{algebra-lemma-fppf-fpqf}
$R \to S$를 유한 표시인 충실히 평탄 환 사상이라 하자. 그러면 가환
그림
$$
\xymatrix{
S \ar[rr] & & S' \\
& R \ar[lu] \ar[ru]
}
$$
이 존재하며, 여기서 $R \to S'$은 준유한이고 충실히 평탄이며 유한
표시이다.
\end{lemma}

\begin{proof}
첫 단계로 이 보조정리를 $R$이 $\mathbf{Z}$ 위에서 유한형인 경우로
귀착시킨다. 보조정리
\ref{algebra-lemma-descend-faithfully-flat-finite-presentation}에 의해 그림
$$
\xymatrix{
S_0 \ar[r] & S \\
R_0 \ar[u] \ar[r] & R \ar[u]
}
$$
이 존재한다. 여기서 $R_0$는 $\mathbf{Z}$ 위에서 유한형이고,
$S_0$는 $R_0$ 위에서 유한 표시인 충실히 평탄 대수이며
$S = R \otimes_{R_0} S_0$이다. 환 사상 $R_0 \to S_0$에 대해
보조정리를 증명하면, 준유한 환 사상의 밑변환이 준유한이라는
보조정리 \ref{algebra-lemma-quasi-finite-base-change}에 의해 밑변환을 취하여
$R \to S$에 대한 보조정리가 따른다. (물론 평탄 사상의 밑변환은
평탄이고 유한 표시 사상의 밑변환은 유한 표시라는 사실도 사용한다.)

\medskip\noindent
$R \to S$가 유한 표시인 충실히 평탄 환 사상이고 $R$이 뇌터라고
가정하자(앞 문단에 의해 이렇게 가정해도 된다). 보조정리
\ref{algebra-lemma-finite-presentation-flat-CM-locus-open}의 열린집합을
$W \subset \Spec(S)$라 하자. $R \to S$가 충실히 평탄이므로 사상
$\Spec(S) \to \Spec(R)$은 전사이다. 보조정리 \ref{algebra-lemma-ff-rings}를
보라. 보조정리 \ref{algebra-lemma-generic-CM-flat-finite-presentation}에 의해
사상 $W \to \Spec(R)$도 전사이다. 따라서 $S$를 곱
$S_{g_1} \times \ldots \times S_{g_m}$으로 바꾸어
$W = \Spec(S)$라고 가정할 수 있다. 여기서는 $\Spec(R)$이
준콤팩트라는 사실(보조정리 \ref{algebra-lemma-quasi-compact})과 사상
$\Spec(S) \to \Spec(R)$이 열린 사상이라는 사실(명제
\ref{algebra-proposition-fppf-open})을 사용한다.
$\mathfrak p \subset R$을 소 아이디얼이라 하자. $\mathfrak p$ 위에
놓이고 올환 $S \otimes_R \kappa(\mathfrak p)$의 한 극대 아이디얼에
대응하는 소 아이디얼 $\mathfrak q \subset S$를 택하자. 뇌터 국소환
$\overline{S}_{\mathfrak q} = S_{\mathfrak q}/\mathfrak pS_{\mathfrak q}$는
Cohen--Macaulay이고, 그 차원을 $d$라 하자. $S_{\mathfrak q}$의
극대 아이디얼에서 $f_1, \ldots, f_d$를 택하여 그 상들이
$\overline{S}_{\mathfrak q}$에서 정칙열을 이루게 할 수 있다.
$f_1, \ldots, f_d$의 공통 분모 $g \in S$, $g \not \in \mathfrak q$를
택하고 $R$-대수
$$
S' = S_g/(f_1, \ldots, f_d).
$$
를 생각하자. 구성에 의해 $\mathfrak p$ 위에 놓이고
$\mathfrak q$에 대응하는 소 아이디얼 $\mathfrak q' \subset S'$이
존재한다(사상 $S_g \to S'_g$을 통해서이다). 또한 구성에 의해
% P01-E1038
환 사상 $R \to S'$은 $\mathfrak q$에서 준유한이다. 실제로 국소환
% P01-E1039
$$
S'_{\mathfrak q'}/\mathfrak pS'_{\mathfrak q'} =
S_{\mathfrak q}/(f_1, \ldots, f_d) + \mathfrak pS_{\mathfrak q} =
\overline{S}_{\mathfrak q}/(\overline{f}_1, \ldots, \overline{f}_d)
$$
은 차원이 영이다. 보조정리 \ref{algebra-lemma-isolated-point-fibre}를 보라.
또한 구성에 의해 $R \to S'$은 유한 표시이다. 끝으로 보조정리
\ref{algebra-lemma-grothendieck-regular-sequence}에 의해 국소환 사상
$R_{\mathfrak p} \to S'_{\mathfrak q'}$은 평탄이다(여기서 $R$이
뇌터라는 사실을 사용한다). 따라서 평탄성의 열린 성질(정리
\ref{algebra-theorem-openness-flatness})과 준유한성의 열린 성질(보조정리
\ref{algebra-lemma-quasi-finite-open})에 의해, 어떤 $g' \in S$,
$g' \not \in \mathfrak q$에 대해 $g$를 $gg'$으로 바꾸어
$R \to S'$이 평탄이고 준유한이라고 가정할 수 있다. 사상
$\Spec(S') \to \Spec(R)$의 상은 열린집합이며 $\mathfrak p$를 포함한다.
달리 말하면, 보조정리의 진술에서와 같은(충실성만은 아직 모를 수
있는) 환 $S'$이 존재하고 그 상이 임의로 주어진 소 아이디얼을
포함함을 보였다. $\Spec(R)$의 준콤팩트성을 한 번 더 사용하면,
그러한 환들의 유한 곱이 원하는 조건을 만족함을 알 수 있다.
\end{proof}
